\documentclass[11pt,a4paper,openany]{book}

\usepackage{fontspec}
\setmainfont{TeX Gyre Termes}
\usepackage[english]{babel}
\usepackage[margin=2.4cm]{geometry}
\usepackage{amsmath,amssymb,amsthm,mathtools}
\usepackage{booktabs,tabularx,array,longtable}
\usepackage{enumitem}
\usepackage{xcolor}
\usepackage{microtype}
\usepackage{needspace}
\usepackage{titlesec}
\usepackage{hyperref}

\hypersetup{
  colorlinks=true,
  linkcolor=blue!55!black,
  citecolor=green!40!black,
  urlcolor=blue!65!black,
  pdftitle={A Collection of Open Problems in Mathematics},
  pdfauthor={Compiled from the supplied source files}
}

\setlist[itemize]{leftmargin=2em,itemsep=0.25em,topsep=0.4em}
\setlist[enumerate]{leftmargin=2.2em,itemsep=0.25em,topsep=0.4em}
\renewcommand{\arraystretch}{1.22}
\setlength{\emergencystretch}{2em}
\newcolumntype{Y}{>{\raggedright\arraybackslash}X}
\setcounter{tocdepth}{1}
\setcounter{secnumdepth}{2}
\titlespacing*{\subsection}{0pt}{2.2ex plus 0.6ex minus 0.3ex}{0.7ex}
\newcounter{problemitem}[section]
\renewcommand{\theproblemitem}{\thesection.\arabic{problemitem}}
\newcommand{\problemsection}[1]{%
  \Needspace{5\baselineskip}%
  \refstepcounter{problemitem}%
  \subsection*{\theproblemitem\quad #1}%
}
\newcommand{\plainproblemsection}{%
  \Needspace{5\baselineskip}%
  \refstepcounter{problemitem}%
  \subsection*{Problem \arabic{problemitem}}%
}

\newtheorem{conjecture}{Conjecture}[section]
\newtheorem{problem}[conjecture]{Open Problem}
\newtheorem{theorem}[conjecture]{Theorem}
\theoremstyle{definition}
\newtheorem{definition}[conjecture]{Definition}
\newtheorem{remark}[conjecture]{Remark}

\newenvironment{openproblem}{%
  \par\smallskip\noindent
  \textbf{Problem statement.}%
  \par\smallskip
}{\par\medskip}

\newcommand{\R}{\mathbb{R}}
\newcommand{\T}{\mathbb{T}}
\newcommand{\Sph}{\mathbb{S}}
\newcommand{\eps}{\varepsilon}
\newcommand{\dd}{\,\mathrm{d}}
\newcommand{\supp}{\operatorname{supp}}
\newcommand{\dist}{\operatorname{dist}}
\newcommand{\dimH}{\dim_{\mathrm H}}
\newcommand{\dimM}{\dim_{\mathrm M}}
\newcommand{\one}{\mathbf{1}}
\newcommand{\lesssimE}{\lesssim_{\eps}}
\newcommand{\Ric}{\operatorname{Ric}}
\newcommand{\curl}{\operatorname{curl}}
\newcommand{\diver}{\operatorname{div}}
\newcommand{\Haus}{\mathcal{H}}
\newcommand{\StatusOpen}{\textcolor{red!70!black}{\textbf{Open}}}
\newcommand{\StatusPartial}{\textcolor{orange!80!black}{\textbf{Partially resolved}}}
\newcommand{\StatusSolved}{\textcolor{green!45!black}{\textbf{Resolved}}}
\newcommand{\StatusPreprint}{\textcolor{blue!65!black}{\textbf{Claimed in a recent preprint}}}
\providecommand{\C}{\mathbb{C}}
\providecommand{\cA}{\mathcal{A}}
\providecommand{\cB}{\mathcal{B}}
\providecommand{\cC}{\mathcal{C}}
\providecommand{\cF}{\mathcal{F}}
\providecommand{\cH}{\mathcal{H}}
\providecommand{\cT}{\mathcal{T}}
\providecommand{\conn}{_{\mathrm{conn}}}
\providecommand{\DD}{\Delta\!\Delta}
\providecommand{\e}{\varepsilon}
\providecommand{\FF}{\mathbb{F}}
\providecommand{\hLMO}{\widehat Z^{\mathrm{LMO}}}
\providecommand{\hZ}{Z}
\providecommand{\K}{\mathbb{K}}
\providecommand{\LL}{\mathbb{L}}
\providecommand{\LMO}{Z^{\mathrm{LMO}}}
\providecommand{\lmo}{z^{\mathrm{LMO}}}
\providecommand{\mboxsm}[1]{\mbox{\small #1}}
\providecommand{\mboxss}[1]{\mbox{\scriptsize #1}}
\providecommand{\M}{\mathbb{M}}
\providecommand{\MK}{\mathbb{M}\mathbb{K}}
\providecommand{\N}{\mathbb{N}}
\providecommand{\ora}[1]{\ensuremath{\overrightarrow{\mathrm{#1}}}}
\providecommand{\Q}{\mathbb{Q}}
\providecommand{\R}{\mathbb{R}}
\providecommand{\subkz}{_{\mathrm{KZ}}}
\providecommand{\simsub}[1]{\mathrel{\mathop{\sim}\limits_{#1}}}
\providecommand{\simsubC}[1]{\mathrel{\mathop{\sim}\limits_{C_{#1}}}}
\providecommand{\simsubHL}[1]{\mathrel{\mathop{\sim}\limits_{HL_{#1}}}}
\providecommand{\simsubY}[1]{\mathrel{\mathop{\sim}\limits_{Y_{#1}}}}
\providecommand{\vcA}{\overrightarrow{\mathcal A}}
\providecommand{\Z}{\mathbb{Z}}
\providecommand{\ZHS}{\mathbb{Z}\mathrm{HS}}
\providecommand{\ZHSs}{\mathbb{Z}\mathrm{HS's}}
\providecommand{\ZTV}{Z^{\mathrm{TV}}}
\providecommand{\ext}[1]{\bigwedge^{#1}}
\providecommand{\ThKeCob}{\mathcal{C}\mathrm{ob}}
\providecommand{\ThKeV}{\mathcal V}
\providecommand{\ThKev}{\mathcal V}
\providecommand{\ThKeQc}{\mathcal Q}
\title{\bfseries Open Problems in Mathematics}
\author{Moonshine Tech, Shanghai}
\date{}

\begin{document}
\frontmatter

\begin{titlepage}
  \thispagestyle{empty}
  \centering

  \vspace*{0.13\textheight}

  {\color{blue!55!black}\rule{\textwidth}{1.2pt}\par}

  \vspace{2.2cm}

  {\fontsize{34}{41}\selectfont\bfseries
    OPEN PROBLEMS\\[0.22em]
    IN MATHEMATICS\par}

  \vspace{1.3cm}

  {\color{blue!55!black}\rule{0.24\textwidth}{0.8pt}\par}

  \vspace{1.3cm}


  \vfill

  {\large Moonshine Tech, Shanghai\par}

  \vspace{1.5cm}

  {\color{blue!55!black}\rule{\textwidth}{1.2pt}\par}
\end{titlepage}

\chapter*{Abstract}
\addcontentsline{toc}{chapter}{Abstract}

Artificial intelligence is advancing rapidly and is playing an
increasingly important role in mathematical research. Motivated
by these developments, we have collected and systematically organized a
selection of open problems from several areas of mathematics. This volume
is intended to serve as a structured reference for researchers, students,
and others interested in current mathematical questions. In particular,
the clear statements and subject-based organization are intended to make
the collection useful as a structured problem resource for AI-assisted
mathematical research, including problem retrieval, systematic
exploration, and collaboration between human researchers and AI systems.

Further information and updated materials are available from the GitHub
source
\href{https://github.com/DeepMathLLM/Moonshine/blob/main/Open\%20Problems\%20in\%20Mathematics/Open\%20Problems\%20in\%20Mathematics.tex}
{\texttt{Open Problems in Mathematics.tex}}(https://github.com/DeepMathLLM/Moonshine)
or from the official website
\href{https://www.deepmath.cn/books.php}
{https://www.deepmath.cn/books.php}.

\tableofcontents

\mainmatter
\chapter{Analysis}
\clearpage
\section{Harmonic Analysis}

\paragraph{Notation.}
The Fourier-transform convention used in the supplied source is
\[
\widehat f(\xi)=\int_{\R^n}f(x)e^{-2\pi i x\cdot\xi}\dd x,
\qquad
f(x)=\int_{\R^n}\widehat f(\xi)e^{2\pi i x\cdot\xi}\dd\xi.
\]
For $1\leq p\leq\infty$, $p'$ denotes the conjugate exponent.  We write
\[
A\lesssim B \Longleftrightarrow A\leq CB,
\qquad
A\lesssimE B \Longleftrightarrow A\leq C_\eps B,
\]
and, for $K\subset\R^n$,
\[
K_\delta=\{x\in\R^n:\dist(x,K)<\delta\}.
\]

\plainproblemsection

The turning-needle problem asks for the least area needed to rotate a unit
line segment through all directions in the plane.  Besicovitch constructions
show that one may place a unit segment in every direction inside a set of
Lebesgue measure zero\cite{Besicovitch1928}; the relevant invariant is
therefore fractal dimension rather than volume.

\begin{definition}[Kakeya or Besicovitch set]
	A compact set $K\subset\R^n$ is a Kakeya set if, for every direction
	$e\in\Sph^{n-1}$, there is an $a_e\in\R^n$ such that
	\[
	\{a_e+te:0\leq t\leq1\}\subset K.
	\]
\end{definition}

\begin{conjecture}[Kakeya set conjecture]
	Every Kakeya set $K\subset\R^n$ has full Hausdorff and Minkowski dimension:
	\[
	\dimH K=\dimM K=n.
	\]
	Equivalently in the Minkowski formulation, for every $\eps>0$ there is
	$C_\eps>0$ such that, for all sufficiently small $\delta>0$,
	\[
	|K_\delta|\geq C_\eps\delta^\eps.
	\]
\end{conjecture}

Davies proved the two-dimensional result\cite{Davies1971}; Wolff's hairbrush
method gave the lower bound $5/2$ in dimension three\cite{Wolff1995}.  The
source records that Wang--Zahl proved full Hausdorff and Minkowski dimension
in $\R^3$ in 2025, followed by a streamlined proof in
2026\cite{WangZahl2025,GuthWangZahl2026}.  The set conjecture remains open
for $n\geq4$.

\plainproblemsection

Let $T_e^\delta(a)$ be a tube of length $1$, cross-sectional radius $\delta$,
and long-axis direction $e$.  Define
\[
\mathcal M_\delta f(e)
=\sup_{a\in\R^n}
\frac{1}{|T_e^\delta(a)|}
\int_{T_e^\delta(a)}|f(x)|\dd x.
\]

\begin{conjecture}[Kakeya maximal-operator conjecture]
	For every $\eps>0$,
	\[
	\|\mathcal M_\delta f\|_{L^n(\Sph^{n-1})}
	\lesssimE
	\delta^{-\eps}\|f\|_{L^n(\R^n)}.
	\]
\end{conjecture}

The maximal-operator conjecture remains open for $n\geq3$.

\plainproblemsection

For a compact $C^2$ hypersurface $S\subset\R^n$ whose principal curvatures are
nonzero and have the same sign, define
\[
E_Sg(x)=\int_S e^{2\pi i x\cdot\xi}g(\xi)\dd\sigma(\xi).
\]
By duality, $E_S:L^r(S)\to L^q(\R^n)$ is equivalent to the restriction map
\[
f\longmapsto\widehat f|_S:
L^{q'}(\R^n)\longrightarrow L^{r'}(S).
\]

\begin{conjecture}[Stein's Fourier restriction conjecture]
	If
	\[
	q>\frac{2n}{n-1},
	\qquad
	q\geq\frac{n+1}{n-1}\,r',
	\]
	then
	\[
	\|E_Sg\|_{L^q(\R^n)}
	\lesssim_{r,q,S}
	\|g\|_{L^r(S)}.
	\]
	Apart from endpoint qualifications, both conditions are necessary.
\end{conjecture}

The first condition reflects stationary-phase decay, while the second comes
from the Knapp example.  In the diagonal case $r=q=p$, the statement is
\[
\|E_Sg\|_{L^p(\R^n)}\lesssim_p\|g\|_{L^p(S)},
\qquad p>\frac{2n}{n-1}.
\]
The Stein--Tomas theorem includes
\[
E_S:L^2(S)\longrightarrow
L^{\frac{2(n+1)}{n-1}}(\R^n)
\]
and the region obtained by interpolation\cite{Tomas1975}.  The full
two-dimensional conjecture is known, while $n\geq3$ remains open.  The source
records the best diagonal result for the positively curved paraboloid in
$\R^3$ as $p>22/7$, compared with the conjectured $p>3$\cite{WangWu2024}.

\plainproblemsection

The sharp frequency-ball cutoff is
\[
S_R^0f(x)=\int_{|\xi|\leq R}\widehat f(\xi)e^{2\pi i x\cdot\xi}\dd\xi.
\]
Fefferman proved that for $n\geq2$ and $p\neq2$ these ball multipliers are not
uniformly bounded on $L^p$; the counterexample uses Kakeya tube
geometry\cite{Fefferman1971}.  The Bochner--Riesz means introduce smoothing:
\[
S_R^\delta f(x)
=\int_{\R^n}
\left(1-\frac{|\xi|^2}{R^2}\right)_+^\delta
\widehat f(\xi)e^{2\pi i x\cdot\xi}\dd\xi,
\qquad (t)_+=\max\{t,0\}.
\]
Set
\[
\delta_c(p)
=\max\left\{
n\left|\frac1p-\frac12\right|-\frac12,\,0
\right\}.
\]

\begin{conjecture}[Bochner--Riesz conjecture]
	Let $n\geq2$ and $1<p<\infty$.  If $\delta>\delta_c(p)$, then
	\[
	\sup_{R>0}\|S_R^\delta f\|_{L^p(\R^n)}
	\lesssim_{p,\delta,n}
	\|f\|_{L^p(\R^n)}.
	\]
	The assertion fails for $\delta<\delta_c(p)$; the line
	$\delta=\delta_c(p)$ is a separate endpoint problem.
\end{conjecture}

Carleson--Sj\"olin resolved the conjecture in two dimensions by more general
oscillatory-integral estimates\cite{CarlesonSjolin1972}; the full conjecture
remains open for $n\geq3$.

\plainproblemsection

For the half-wave propagator
\[
u(x,t)=e^{it\sqrt{-\Delta}}f(x),
\]
fixed-time Fourier-integral estimates for $p\geq2$ require
\[
\|u(\cdot,t)\|_{L^p(\R^n)}
\lesssim
\|f\|_{W^{s_p,p}(\R^n)},
\qquad
s_p=(n-1)\left(\frac12-\frac1p\right).
\]
Sogge proposed that integration over a time interval of length one should
yield almost $1/p$ of an additional derivative\cite{Sogge1991}.

\begin{conjecture}[Euclidean local smoothing for the wave equation]
	Let $n\geq2$ and
	\[
	p\geq\frac{2n}{n-1}.
	\]
	For every $\eps>0$,
	\[
	\|e^{it\sqrt{-\Delta}}f\|_{L^p(\R^n\times[1,2])}
	\lesssim_{p,\eps}
	\|f\|_{W^{s_p-1/p+\eps,p}(\R^n)}.
	\]
	Equivalently, one may use $W^{s_p-\sigma,p}$ for any $\sigma<1/p$.
\end{conjecture}

For the full wave equation, the corresponding right-hand side is
\[
\|f_0\|_{W^{s_p-\sigma,p}}
+\|f_1\|_{W^{s_p-1-\sigma,p}}.
\]
The lower bound on $p$ is forced by focusing wave-packet examples.
Bourgain--Demeter decoupling gives the important higher-dimensional
range
\[
p\geq\frac{2(n+1)}{n-1}
\]
with the customary non-endpoint or $\eps$ qualification\cite{BourgainDemeter2015}.
Guth--Wang--Zhang resolved the conjecture in two spatial dimensions,
$p\geq4$\cite{GuthWangZhang2020}.  The source records that in spatial
dimension three a 2026 preprint reaches $p\geq10/3$, while the conjectured
threshold is $p\geq3$\cite{GanHeLiWu2026}.

\plainproblemsection

\begin{openproblem}
	Suppose that $T\left( x\right)  = \mathop{\sum }\limits_{{n = 0}}^{N}{c}_{n}\cos {n}^{2}x$ , that $T\left( 0\right)  = 1$ , and that $T\left( x\right)  \geq  0$ for all $x$ . Among all such trigonometric polynomials, let $\delta \left( N\right)$ denote the minimum possible value of ${c}_{0}$ . It is known, via indirect reasoning, that $\delta \left( N\right)  \rightarrow  0$ as $N \rightarrow  \infty$ . How fast, as a function of $N$ , does $\delta \left( N\right)$ tend to 0 ? Like ${N}^{-a}$ ? Like ${\left( \log N\right) }^{-a}$ ? ...
\end{openproblem}

\plainproblemsection

\begin{openproblem}
	Suppose that $0 \leq  {t}_{1} < {t}_{2} < \ldots  < {t}_{R} \leq  T$ with ${t}_{r + 1} - {t}_{r} \geq  1$ for all $r$ , and $\left| {a}_{n}\right|  \leq  1$ for all $n$ .
	
	Show that
	
	\[
	\mathop{\sum }\limits_{{r = 1}}^{R}{\left| \mathop{\sum }\limits_{{n = 1}}^{N}{a}_{n}{n}^{-i{t}_{r}}\right| }^{2} \ll  \left( {N + R}\right) {N}^{1 + \varepsilon }.
	\]
\end{openproblem}

\plainproblemsection

\begin{openproblem}
	Let ${\lambda }_{1},\ldots ,{\lambda }_{N}$ be distinct real numbers, and put ${\delta }_{n} = \mathop{\min }\limits_{\substack{m \\  {m \neq  n} }}\left| {{\lambda }_{m} - {\lambda }_{n}}\right|$ . Determine the best constant
	
	\[
	\mathop{\sum }\limits_{\substack{{m, n} \\  {m \neq  n} }}\frac{{x}_{m}\overline{{x}_{n}}}{{\lambda }_{m} - {\lambda }_{n}} \leq  c\mathop{\sum }\limits_{{n = 1}}^{N}\frac{{\left| {x}_{n}\right| }^{2}}{{\delta }_{n}}.
	\]
	
	Certainly the best constant is no smaller than $\pi$ . Montgomery and Vaughan [11] proved the above inequality with $c = {3\pi }/2$ .
\end{openproblem}

\plainproblemsection

\begin{openproblem}
	(Littlewood) Let $\parallel \theta \parallel$ denote the distance from $\theta$ to the nearest integer. Is it true that
	
	\[
	\mathop{\liminf }\limits_{{n \rightarrow  \infty }}n\parallel {n\theta }\parallel \parallel {n\phi }\parallel  = 0
	\]
	
	for any real numbers $\theta$ and $\phi$ ?
\end{openproblem}

\plainproblemsection

\begin{openproblem}
	(Lehmer) The Mahler measure of a polynomial $P\left( x\right)  = {a}_{n}{x}^{n} + \cdots  + {a}_{0} = {a}_{n}\left( {x - {\alpha }_{1}}\right) \cdots \left( {x - {\alpha }_{n}}\right)$ is the quantity
	
	\[
	M\left( P\right)  = \left| {a}_{n}\right| \mathop{\prod }\limits_{{i = 1}}^{n}\max \left( {1,\left| {\alpha }_{i}\right| }\right) .
	\]
	
	Show that there is a constant $c > 1$ such that if $P \in  \mathbb{Z}\left\lbrack  x\right\rbrack$ and $M\left( P\right)  < c$ then $P$ is a product of cyclotomic polynomials (i.e. $M\left( P\right)  = 1$ ).
	
	The least known value of $M\left( P\right)  > 1$ was found by Lehmer [9] long ago, in an exhaustive search. Lehmer’s polynomial $P\left( x\right)  = {x}^{10} + {x}^{9} - {x}^{7} - {x}^{6} - {x}^{5} - {x}^{4} - {x}^{3} + x + 1$ has arisen many times since. It achieves $M\left( P\right)  = {1.1762808}$ .
	
	Let $\theta  > 1$ satisfy ${\theta }^{3} - \theta  - 1 = 0$ . Siegel proved that $\theta$ is the least PV number. Smyth [15] proved that if $P$ is irreducible and $M\left( P\right)  < \theta$ then $P$ is a reciprocal polynomial, that is, $P\left( x\right)  = {x}^{n}P\left( {1/x}\right)$ . More recent numerical studies have been conducted by Boyd [2].
\end{openproblem}

\plainproblemsection

\begin{openproblem}
	(Salem) An algebraic integer $\alpha  > 1$ is called a Salem number if all conjugates of $\alpha$ lie in the closed unit disk $\left| z\right|  \leq  1$ , with at least one conjugate on the unit circle. Let $\mathcal{S}$ denote the set of all PV numbers, and $\mathcal{T}$ the set of all Salem numbers. Salem proved that every member of $\mathcal{S}$ is a limit of members of $\mathcal{T}$ . Does $\mathcal{T}$ have any other limit points? Prove that 1 is not a limit point of $\mathcal{T}$ . Prove that $\mathcal{S} \cup  \mathcal{T}$ is closed.
	
	For basic properties of PV and Salem numbers, see Salem [13] or Bertin, Decomps-Guilloux, Grandet-Hugot, Pathiaux-Delefosse and Schreiber [1]. See also numerous papers of Boyd, and papers cited therein.
\end{openproblem}

\plainproblemsection

\begin{openproblem}
	Show that there is a constant $c > 0$ such that if an algebraic integer $\alpha$ and all its conjugates lie in the disk $\left| z\right|  \leq  1 + c/n$ then $\alpha$ is a root of unity.
	
	The best result in this direction thus far is due to Dobrowolski [4], who proved that there is a positive constant $c$ such that if
	
	\[
	M\left( \alpha \right)  < 1 + c{\left( \frac{\log \log n}{\log n}\right) }^{3}
	\]
	
	then $\alpha$ is a root of unity. The value of $c$ has been improved, and the proof simplified, by Cantor and Straus [3], Rausch [12] and Louboutin [10].
\end{openproblem}

\plainproblemsection

\begin{openproblem}
	(D.J. Newman) Let ${\varepsilon }_{0},{\varepsilon }_{1},\ldots ,{\varepsilon }_{N}$ be independent random variables with $P\left( {{\varepsilon }_{n} = 0}\right)  = P\left( {{\varepsilon }_{n} = }\right.$ $1) = 1/2$ , and put $K\left( z\right)  = \mathop{\sum }\limits_{{n = 0}}^{N}{\varepsilon }_{n}{z}^{n}$ . Show that there is an absolute constant $c > 0$ such that
	
	\[
	P\left( {\mathop{\min }\limits_{{\left| z\right|  = 1}}\left| {K\left( z\right) }\right|  < 1}\right)  > c.
	\]
	
	Does there exist a choice of the ${\varepsilon }_{n} \in  \{ 0,1\}$ such that $\mathop{\min }\limits_{{\left| z\right|  = 1}}\left| {K\left( z\right) }\right|  > c$ ? What if $c$ is replaced by $c\sqrt{N}$ ?
\end{openproblem}

\plainproblemsection

\begin{openproblem}
	(Hardy-Littlewood) Suppose that $f \in  {L}^{2}\left( \mathbb{T}\right)$ , and that $\widehat{f}\left( k\right)  \neq  0$ only when $k$ is a perfect square. Does it follow that $f \in  {L}^{p}\left( \mathbb{T}\right)$ for all $p < 4$ ?
	
	A. Córdoba has shown that the answer is in the affirmative if one assumes that the numbers $\widehat{f}\left( {k}^{2}\right)$ are positive and monotonically decreasing.
\end{openproblem}

\plainproblemsection

\begin{openproblem}
	Consider a simple closed curve in the plane. A point $P$ is called equichordal if every line through $P$ meets the curve at two points that are a constant distance apart. Can a curve have two distinct equichordal points?
\end{openproblem}

\plainproblemsection

\begin{openproblem}
	(Erdős) Let $\mathcal{A}$ be an infinite set of positive integers, and let $r\left( n\right)$ be the number of ways of writing $n$ as a sum of two members of $\mathcal{A}$ . If $r\left( n\right)  > 0$ for all $n$ , does it follow that $\limsup r\left( n\right)  =  + \infty$ ?
\end{openproblem}

\plainproblemsection

\begin{openproblem}
	(Fürstenberg) Let $p$ and $q$ be positive integers, not powers of the same number. Suppose that $x$ is a real number, and let ${\mathrm{S}}_{p}$ be the set of fractional parts of the numbers ${p}^{n}x$ . Similarly, let ${S}_{q}$ be the set of fractional parts of the numbers ${q}^{n}x$ . Show that if $x$ is irrational then the sum of the Hausdorff dimensions of the closures of these sets is $\geq  1$ .
	
	From this it would follow that ${2}^{n}$ has a 7 in its base 10 expansion, for all large $n$ .
\end{openproblem}

\plainproblemsection

\begin{openproblem}
	(Veech) Suppose that $\alpha$ is given, $0 < \alpha  < 1/2$ . Does there exist a function $f$ , not identically 0, such that $f$ has period 1, is even, $f\left( 0\right)  = 0,\mathop{\sum }\limits_{{a = 1}}^{q}f\left( {a/q}\right)  = 0$ for every positive integer $q$ , and $f \in  \operatorname{Lip}\left( \alpha \right)$ ?
\end{openproblem}

\plainproblemsection

\begin{openproblem}
	For $P \in  \mathbb{C}\left\lbrack  z\right\rbrack$ let $\parallel P\parallel$ denote the ${\ell }_{1}$ norm of the coefficients of $P$ . Suppose that $A$ and $B$ are polynomials of degrees $a$ and $b$ , respectively. It is known that the ratio $\parallel A\parallel \parallel B\parallel /\parallel {AB}\parallel$ always lies between 1 and ${2}^{a + b}$ provided that ${AB} \neq  0$ . What is the ’usual’ or ’average’ size of this ratio, as one averages over polynomials in a suitable way?
\end{openproblem}

\plainproblemsection

\begin{openproblem}
	(H. Maier) Let ${\Phi }_{n}\left( z\right)$ denote the ${n}^{\text{th }}$ cyclotomic polynomial,
	
	\[
	{\Phi }_{n}\left( z\right)  = \mathop{\prod }\limits_{{a = 1}}^{n}\left( {z - {e}^{{2\pi ia}/n}}\right)  = \mathop{\sum }\limits_{{m = 0}}^{{\phi \left( n\right) }}c\left( {m, n}\right) {z}^{m},
	\]
	
	Put $C\left( n\right)  = \mathop{\max }\limits_{m}\left| {c\left( {m, n}\right) }\right|$ . Show that $\left( {\log C\left( n\right) }\right) /\log n$ has a limiting distribution as $n \rightarrow  \infty$ .
\end{openproblem}

\plainproblemsection

\begin{openproblem}
	(Schinzel) Let ${\Phi }_{n}\left( z\right)$ denote the ${n}^{\text{th }}$ cyclotomic polynomial, as above. Show that if $n$ is squarefree then ${\Phi }_{n}$ has $\gg  {n}^{1/2}$ non-zero coefficients.
\end{openproblem}

\plainproblemsection

\begin{openproblem}
	(Balazard) Does there exist a Dirichlet series $D\left( s\right)  = \mathop{\sum }\limits_{{n = 1}}^{\infty }{a}_{n}{n}^{-s}$ that converges in a half-plane $\Re s > \alpha$ with the property that the equation $D\left( s\right)  = 0$ has exactly one solution in this half-plane?
	
	If the Riemann Hypothesis is true then $1/\zeta \left( s\right)$ is such Dirichlet series, since then one can take $\alpha  = 1/2$ and the only zero is at $s = 1$ . The object is to give an unconditional construction.
\end{openproblem}

\plainproblemsection

\begin{openproblem}
	(Schinzel) Suppose that $f\left( x\right)  = \mathop{\sum }\limits_{{i = 0}}^{n}{a}_{i}{x}^{i}$ with ${a}_{i} \geq  0$ for all $i$ , and that $f{\left( x\right) }^{2} = \mathop{\sum }\limits_{{i = 0}}^{{2n}}{b}_{i}{x}^{i}$ with ${b}_{i} \leq  1$ for all $i$ . Let $\mathrm{S}\left( n\right)$ denote the set of all such polynomials. It is known that $\mathop{\max }\limits_{{f \in  \mathcal{S}\left( n\right) }}f{\left( 1\right) }^{2} \sim  {Cn}$ . Show that $C = 4/\pi$ .
	
	Schmidt proved that $4/\pi  \leq  C \leq  2$ and Schinzel showed that $C \leq  7/4$ .
\end{openproblem}

\plainproblemsection

\begin{openproblem}
	(Ruzsa) Let $\mathcal{A}$ be a set of positive integers with positive asymptotic density. It follows that gaps between members of the set $\mathcal{A} - \mathcal{A}$ are bounded. It does not follow that gaps between members of the set $\mathcal{A} + \mathcal{A}$ are bounded. Does it follow that gaps between members of the set $\mathcal{A} + \mathcal{A} + \mathcal{A}$ are bounded?
\end{openproblem}

\plainproblemsection

\begin{openproblem}
	(Melfi-Erdös) Let $\mathcal{A} = \left\{  {\sum {\epsilon }_{i}{3}^{i} : {\epsilon }_{i} = 0,1}\right\}$ , and let $\mathcal{B} = \left\{  {\sum {\epsilon }_{i}{4}^{i} : {\epsilon }_{i} = 0,1}\right\}$ . Does the set $\mathcal{A} + \mathcal{B}$ have positive lower density?
\end{openproblem}

\plainproblemsection

\begin{openproblem}
	(Brüdern) Let $\mathcal{P}$ denote the set of prime numbers. Suppose that a set $\mathcal{A}$ is chosen so that $\mathcal{A} + \mathcal{P}$ has density 1. Let $A\left( N\right)$ denote the number of members of $\mathcal{A}$ not exceeding $N$ . How slowly can $A\left( N\right)$ grow? Certainly one can have $A\left( N\right)  < {\left( \log N\right) }^{C}$ . What is the least possible value of $C$ ?
\end{openproblem}

\plainproblemsection

\begin{openproblem}
	(Skinner) Let $f$ be a non-constant polynomial with rational coefficients such that $f\left( i\right)$ is 0 or 1 for $i = 0,1,\ldots , n$ . How small can the degree of $f$ be?
	
	Clearly $\left\lbrack  {n/2}\right\rbrack$ is a lower bound, and $n$ is an upper bound.
\end{openproblem}

\plainproblemsection

\begin{openproblem}
	Let $\parallel \theta \parallel$ denote the distance from $\theta$ to the nearest integer. From a classical theorem of Dirichlet we know that for any real $\theta$ and any positive integer $N$ there is an $n,1 \leq  n \leq  N$ , such that $\parallel {n\theta }\parallel  \leq  1/\left( {N + 1}\right)$ . We want something similar, but using only squares: For every real $\theta$ and every positive integer $N$ there is an $n,1 \leq  n \leq  N$ , such that $\begin{Vmatrix}{{n}^{2}\theta }\end{Vmatrix} \leq  f\left( N\right)$ , with $f\left( N\right)$ as small as possible.
	
	Heilbronn showed that one may take $f\left( N\right)  = {N}^{-1/2 + \varepsilon }$ , and more recently A. Zaharescu [16] showed that this can be improved to $f\left( N\right)  = {N}^{-4/7 + \varepsilon }$ . Can this be improved to $f\left( N\right)  =$ ${N}^{-1 + \varepsilon }$ ?
\end{openproblem}

\plainproblemsection

\begin{openproblem}
	Show that if $P\left( x\right)  = \mathop{\sum }\limits_{{j = 1}}^{k}{\alpha }_{j}{x}^{j},\left| {{\alpha }_{k} - a/q}\right|  \leq  1/{q}^{2},\left( {a, q}\right)  = 1$ , then
	
	\[
	\mathop{\sum }\limits_{{n = 1}}^{N}e\left( {P\left( n\right) }\right) { \ll  }_{k}{N}^{1 + \varepsilon }{\left( \frac{1}{q} + \frac{q}{{N}^{k}}\right) }^{1/k}.
	\]
	
	Alternatively, construct examples that violate this.
	
	Even small improvements of existing bounds would be interesting. For example when $k = 3$ and $q \approx  {N}^{3/2}$ , derive an upper bound that is $o\left( {N}^{3/4}\right)$ .
\end{openproblem}

\plainproblemsection

\begin{openproblem}
	Show that for any positive number $B$ there exist complex numbers ${z}_{1},\ldots ,{z}_{N}$ such that $\left| {z}_{n}\right|  = 1$ for all $n$ , and $\mathop{\sum }\limits_{{n = 1}}^{N}{z}_{n}^{\nu }{ \ll  }_{B}{N}^{1/2}$ uniformly for $1 \leq  \nu  \leq  {N}^{B}$ .
	
	This is known when $B = 2$ .
\end{openproblem}

\plainproblemsection

\begin{openproblem}
	Suppose that $\left| {z}_{n}\right|  \geq  1$ and that ${b}_{n} \geq  0$ for all $n$ . Prove that
	
	\[
	\mathop{\max }\limits_{{1 \leq  \nu  \leq  {2N}}}\left| {\mathop{\sum }\limits_{{n = 1}}^{N}{b}_{n}{z}_{n}^{\nu }}\right|  \gg  {\left( \mathop{\sum }\limits_{{n = 1}}^{N}{\left| {b}_{n}\right| }^{2}\right) }^{1/2}
	\]
	
	or give a counter-example.
	
	This was proved by Gonek [5] under the stronger hypothesis that $\left| {z}_{n}\right|  = 1$ for all $n$ .
\end{openproblem}

\plainproblemsection

\begin{openproblem}
	Suppose that $N$ points ${\mathbf{u}}_{n}$ are given in $\mathbb{T}$ . Let $\mathcal{B}$ be a box with one corner at0and the diagonally opposite corner at $\mathbf{\alpha }$ , and let $D\left( \mathbf{\alpha }\right)  = \mathop{\sum }\limits_{{n = 1}}^{N}{\chi }_{\mathfrak{B}}\left( {\mathbf{u}}_{n}\right)  - N\operatorname{area}\left( \mathfrak{B}\right)$ denote the discrepancy of the ${\mathbf{u}}_{n}$ with respect to this box. Finally, let $D$ denote the supremum of $\left| {D\left( \mathbf{\alpha }\right) }\right|$ over all boxes. It is known that points ${\mathbf{u}}_{n}$ can be chosen so that $\left| {D\left( \mathbf{\alpha }\right) }\right|$ is uniformly bounded by ${c}_{k}{\left( \log N\right) }^{k - 1}$ . Is it always the case that $\parallel D{\parallel }_{\infty }$ is at least this order of magnitude?
	
	This is known to be the case when $k = 2$ , first proved by W. M. Schmidt [14] in 1971, and then by Halász [6] later by a modified form of Roth’s method. But it is not known for larger $k$ .
\end{openproblem}

\plainproblemsection

\begin{openproblem}
	Let $N$ points be given in ${\mathbb{T}}^{k}$ , and let $D\left( \mathbf{\alpha }\right)$ denote the associated discrepancy function, as defined above. Show that $\parallel D{\parallel }_{1} > {c}_{k}{\left( \log N\right) }^{\left( {k - 1}\right) /2}$ .
	
	This was proved by Halász [6] when $k = 2$ , but the case $k > 2$ is open.
\end{openproblem}

\plainproblemsection

\begin{openproblem}
	Show that if $\left| {a}_{n}\right|  \leq  1$ for all $n$ , then
	
	\[
	{\int }_{0}^{T}{\left| \mathop{\sum }\limits_{{n = 1}}^{N}{a}_{n}{n}^{-{it}}\right| }^{2p}{dt} \ll  \left( {T + {N}^{p}}\right) {N}^{p + \varepsilon }
	\]
	
	uniformly for $1 \leq  p \leq  2$ .
	
	This is easy to prove when $p = 1$ or $p = 2$ . The Density Hypothesis concerning zeros of the Riemann zeta function follows from the above.
\end{openproblem}

\begingroup
\renewcommand{\bibname}{References}
\let\chapter\subsection
\begin{thebibliography}{99}
	
	\bibitem{Besicovitch1928}
	A.~S. Besicovitch,
	\newblock On Kakeya's problem and a similar one,
	\newblock \emph{Mathematische Zeitschrift} \textbf{27} (1928), 312--320.
	
	\bibitem{Davies1971}
	R.~O. Davies,
	\newblock Some remarks on the Kakeya problem,
	\newblock \emph{Proc. Cambridge Philos. Soc.} \textbf{69} (1971), 417--421.
	\newblock \url{https://doi.org/10.1017/S0305004100046867}
	
	\bibitem{Fefferman1971}
	C. Fefferman,
	\newblock The multiplier problem for the ball,
	\newblock \emph{Ann. of Math.} (2) \textbf{94} (1971), 330--336.
	\newblock \url{https://doi.org/10.2307/1970864}
	
	\bibitem{CarlesonSjolin1972}
	L. Carleson and P. Sj\"olin,
	\newblock Oscillatory integrals and a multiplier problem for the disc,
	\newblock \emph{Studia Math.} \textbf{44} (1972), 287--299.
	\newblock \url{https://eudml.org/doc/217689}
	
	\bibitem{Tomas1975}
	P.~A. Tomas,
	\newblock A restriction theorem for the Fourier transform,
	\newblock \emph{Bull. Amer. Math. Soc.} \textbf{81} (1975), 477--478.
	\newblock \url{https://doi.org/10.1090/S0002-9904-1975-13790-6}
	
	\bibitem{Wolff1995}
	T. Wolff,
	\newblock An improved bound for Kakeya type maximal functions,
	\newblock \emph{Rev. Mat. Iberoamericana} \textbf{11} (1995), 651--674.
	
	\bibitem{TaoVargasVega1998}
	T. Tao, A. Vargas and L. Vega,
	\newblock A bilinear approach to the restriction and Kakeya conjectures,
	\newblock \emph{J. Amer. Math. Soc.} \textbf{11} (1998), 967--1000.
	
	\bibitem{Tao1999}
	T. Tao,
	\newblock The Bochner--Riesz conjecture implies the restriction conjecture,
	\newblock \emph{Duke Math. J.} \textbf{96} (1999), 363--375.
	\newblock \url{https://doi.org/10.1215/S0012-7094-99-09610-2}
	
	\bibitem{Sogge1991}
	C.~D. Sogge,
	\newblock Propagation of singularities and maximal functions in the plane,
	\newblock \emph{Invent. Math.} \textbf{104} (1991), 349--376.
	\newblock \url{https://doi.org/10.1007/BF01245080}
	
	\bibitem{BourgainDemeter2015}
	J. Bourgain and C. Demeter,
	\newblock The proof of the $\ell^2$ decoupling conjecture,
	\newblock \emph{Ann. of Math.} (2) \textbf{182} (2015), 351--389.
	\newblock \url{https://doi.org/10.4007/annals.2015.182.1.9}
	
	\bibitem{Guth2016}
	L. Guth,
	\newblock A restriction estimate using polynomial partitioning,
	\newblock \emph{J. Amer. Math. Soc.} \textbf{29} (2016), 371--413.
	
	\bibitem{GuthWangZhang2020}
	L. Guth, H. Wang and R. Zhang,
	\newblock A sharp square function estimate for the cone in $\R^3$,
	\newblock \emph{Ann. of Math.} (2) \textbf{192} (2020), 551--581.
	\newblock \url{https://doi.org/10.4007/annals.2020.192.2.6}
	
	\bibitem{WangWu2024}
	H. Wang and S. Wu,
	\newblock Restriction estimates using decoupling theorems and two-ends Furstenberg inequalities,
	\newblock arXiv:2411.08871, version 3, 2024.
	\newblock \url{https://arxiv.org/abs/2411.08871}
	
	\bibitem{WangZahl2025}
	H. Wang and J. Zahl,
	\newblock Volume estimates for unions of convex sets, and the Kakeya set conjecture in three dimensions,
	\newblock arXiv:2502.17655, 2025.
	\newblock \url{https://arxiv.org/abs/2502.17655}
	
	\bibitem{GuthWangZahl2026}
	L. Guth, H. Wang and J. Zahl,
	\newblock A streamlined proof of the Kakeya set conjecture in $\R^3$,
	\newblock arXiv:2601.14411, 2026.
	\newblock \url{https://arxiv.org/abs/2601.14411}
	
	\bibitem{GanHeLiWu2026}
	S. Gan, D. He, X. Li and S. Wu,
	\newblock On local smoothing estimates for wave equations,
	\newblock arXiv:2502.05973, version 2, 2026.
	\newblock \url{https://arxiv.org/abs/2502.05973}
	
\end{thebibliography}
\endgroup

\clearpage
\section{Partial Differential Equations and Inverse Problems}

\plainproblemsection

On $\R^3$ or $\T^3$, the incompressible Navier--Stokes equations are
\begin{equation}\label{eq:collection-NS}
	\begin{cases}
		\partial_t u+(u\cdot\nabla)u+\nabla p=\nu\Delta u,\\
		\nabla\cdot u=0,\\
		u|_{t=0}=u_0,
	\end{cases}
	\qquad \nu>0.
\end{equation}
They are invariant under
\[
u_\lambda(t,x)=\lambda u(\lambda^2t,\lambda x),
\qquad
p_\lambda(t,x)=\lambda^2p(\lambda^2t,\lambda x).
\]
Thus $L_x^3$ and $\dot H_x^{1/2}$ are scale-critical, whereas the
three-dimensional energy norm $L_x^2$ is supercritical.

\begin{conjecture}[Global existence and smoothness in three dimensions]
	For every smooth, divergence-free, rapidly decreasing initial datum $u_0$
	(or every smooth periodic datum), equation \eqref{eq:collection-NS} has a
	unique global smooth solution whose required higher-order norms remain finite
	on every finite time interval.
\end{conjecture}

The Clay problem also permits a resolution in the opposite direction by
constructing admissible data for which a smooth solution fails to exist for
all positive time\cite{Fefferman2000}.  Two-dimensional global regularity,
three-dimensional global finite-energy weak solutions, small critical data,
and local smooth solutions are known.  General smoothness and uniqueness for
large three-dimensional data remain unknown.  The source records the status
as \StatusOpen\ and cites the Clay page as still marking the problem
unsolved\cite{ClayNS2026}.

\plainproblemsection

Setting $\nu=0$ gives
\begin{equation}\label{eq:collection-Euler}
	\begin{cases}
		\partial_t u+(u\cdot\nabla)u+\nabla p=0,\\
		\nabla\cdot u=0.
	\end{cases}
\end{equation}
Here the standard problem means the Cauchy problem on $\R^3$ or $\T^3$ with
smooth finite-energy data and no solid boundary.

\begin{problem}[Global regularity or finite-time blowup for boundaryless Euler]
	Does every such smooth datum generate a global smooth solution?  If not, can
	one construct a smooth datum for which regularity fails at a finite time
	$T_*$?
\end{problem}

With $\omega=\curl u$,
\[
\partial_t\omega+(u\cdot\nabla)\omega=(\omega\cdot\nabla)u.
\]
The Beale--Kato--Majda criterion says that breakdown at $T_*<\infty$ requires
\[
\int_0^{T_*}\|\omega(t)\|_{L^\infty}\dd t=\infty;
\]
finiteness of this integral permits continuation\cite{BealeKatoMajda1984}.
The source notes the 2025 Chen--Hou construction with a solid boundary, but
the boundary stagnation point and reflection geometry are part of that
mechanism, so it does not settle the standard problem on $\R^3$ or
$\T^3$\cite{ChenHou2025}.  Recorded status: \StatusOpen.

\plainproblemsection

For smooth Euler solutions on $\T^3$, the kinetic energy
\[
E(t)=\frac12\int_{\T^3}|u(t,x)|^2\dd x
\]
is conserved.  For rough weak solutions, high-frequency energy flux may
produce anomalous dissipation.

\begin{conjecture}[Onsager, standard non-endpoint form]
	Let $u$ be a weak solution of the three-dimensional incompressible Euler
	equations.
	\begin{enumerate}
		\item If $u\in L^\infty_tC^\alpha_x$ and $\alpha>1/3$, then total kinetic
		energy is conserved.
		\item For every $\alpha<1/3$, there is a $C^\alpha$ weak solution that does
		not conserve kinetic energy.
	\end{enumerate}
\end{conjecture}

Constantin--E--Titi proved the rigidity direction for
$\alpha>1/3$\cite{ConstantinETiti1994}; Isett proved the flexibility direction
for every $\alpha<1/3$\cite{Isett2018}.  The standard non-endpoint conjecture
is therefore \StatusSolved.  The exact $C^{1/3}$ endpoint remains a separate
open problem; the cited endpoint work approaches but does not supply an exact
$C^{1/3}$ dissipative example\cite{Isett2024}.

\plainproblemsection

The dissipative surface quasi-geostrophic equation is
\begin{equation}\label{eq:collection-SQG}
	\partial_t\theta+u\cdot\nabla\theta+\kappa\Lambda^\gamma\theta=0,
	\qquad
	u=R^\perp\theta,
	\qquad
	\Lambda=(-\Delta)^{1/2},
\end{equation}
where $\kappa>0$.  With this notation, $\gamma>1$ is subcritical,
$\gamma=1$ critical, and $0<\gamma<1$ supercritical.  The scaling
\[
\theta_\lambda(t,x)
=\lambda^{\gamma-1}\theta(\lambda^\gamma t,\lambda x)
\]
leaves the equation invariant.

\begin{conjecture}[Supercritical global regularity]
	For every $0<\gamma<1$ and every smooth initial datum, the classical solution
	of \eqref{eq:collection-SQG} exists globally and remains smooth.
\end{conjecture}

Global regularity for arbitrary data in the critical case $\gamma=1$ is
known\cite{KiselevNazarovVolberg2007,CaffarelliVasseur2010}.  For
$0<\gamma<1$, the source lists small-data theory, slightly supercritical
dissipation, eventual and conditional regularity, and special large-data
classes, including small perturbations of possibly large radial
data\cite{BulutDong2024}.  Arbitrary large smooth data remain
\StatusOpen.

\plainproblemsection

Soliton resolution is a program for long-time dynamics across nonlinear
dispersive and wave equations.  For the focusing energy-critical wave
equation
\[
\partial_t^2u-\Delta u=|u|^{\frac{4}{d-2}}u,
\qquad d\geq3,
\]
the expected form is as follows.

\begin{conjecture}[Soliton resolution, schematic form]
	For a generic global finite-energy solution in an appropriate sense, as
	$t\to+\infty$,
	\[
	\vec u(t)
	=\vec u_{\mathrm{lin}}(t)
	+\sum_{j=1}^{J}\vec Q_j(t)
	+o_{\mathcal H}(1),
	\]
	where $\vec u_{\mathrm{lin}}$ is free radiation, the $\vec Q_j$ are
	asymptotically separated solitons or traveling waves, and the remainder tends
	to zero in the energy space $\mathcal H$.  Finite-time blowup solutions should
	have a corresponding multi-bubble decomposition.
\end{conjecture}

The formulation depends on the equation: time-dependent scales, breathers,
threshold resonances, or infinite soliton trains may have to be allowed.
The full nonradial, nonintegrable focusing problem remains
\StatusPartial\cite{Tao2008,Chatterjee2013}.  The source records complete
decompositions for several equivariant or radial critical models and
integrable equations\cite{DuyckaertsEtAl2021,JendrejLawrie2022}, as well as a
2026 preprint claim for sufficiently regular decaying Benjamin--Ono
data\cite{GassotGerardMiller2026}.

\plainproblemsection

Consider
\begin{equation}\label{eq:collection-AllenCahn}
	\Delta u+u-u^3=0
	\qquad\text{in }\R^n,
\end{equation}
with
\[
|u|\leq1,
\qquad
\partial_{x_n}u>0.
\]

\begin{conjecture}[De Giorgi's monotone one-dimensionality conjecture]
	If $n\leq8$, every bounded entire solution of
	\eqref{eq:collection-AllenCahn} that is strictly monotone in one direction
	depends on only one Euclidean direction; equivalently, all of its level sets
	are hyperplanes.
\end{conjecture}

One often adds the end-state condition
\[
\lim_{x_n\to\pm\infty}u(x',x_n)=\pm1,
\]
and the versions with and without this condition must be distinguished.
The full monotone statement is known for $n=2,3$\cite{AmbrosioCabre2000}.
For $4\leq n\leq8$, it is known with the end-state condition
\cite{Savin2009}, while the unrestricted monotone statement remains open.
For every $n\geq9$, nonplanar monotone solutions give
counterexamples\cite{DelPinoKowalczykWei2011}.  Recorded status:
\StatusPartial.

\plainproblemsection

Consider positive classical solutions of
\begin{equation}\label{eq:collection-LaneEmden}
	\begin{cases}
		-\Delta u=v^p,\\
		-\Delta v=u^q,\\
		u,v>0
	\end{cases}
	\qquad\text{in }\R^N,
	\qquad p,q>0.
\end{equation}
The critical curve is
\[
\frac{1}{p+1}+\frac{1}{q+1}=\frac{N-2}{N}.
\]

\begin{conjecture}[Lane--Emden system conjecture]
	If
	\[
	\frac{1}{p+1}+\frac{1}{q+1}>\frac{N-2}{N},
	\]
	then \eqref{eq:collection-LaneEmden} has no positive solution in
	$C^2(\R^N)\times C^2(\R^N)$.
\end{conjecture}

When $p=q$, the condition becomes
\[
p<\frac{N+2}{N-2},
\]
the scalar Sobolev critical exponent.  Souplet's result in $N=4$, together
with earlier work, completes $N\leq4$\cite{Souplet2009}.  The source records
the full conjecture as open for $N\geq5$, with results only for restricted
parameter ranges or classes of solutions\cite{HuangZou2025}.  Recorded
status: \StatusPartial.

\plainproblemsection

Let $(M^n,g)$ be a compact connected smooth Riemannian manifold without
boundary and let
\[
-\Delta_g\varphi_\lambda=\lambda\varphi_\lambda,
\qquad
Z_{\varphi_\lambda}=\{x\in M:\varphi_\lambda(x)=0\}.
\]

\begin{conjecture}[Yau's nodal-set measure conjecture]
	There are constants $c,C>0$, depending only on $(M,g)$, such that every
	nonzero eigenfunction satisfies
	\[
	c\sqrt{\lambda}
	\leq
	\Haus^{n-1}(Z_{\varphi_\lambda})
	\leq
	C\sqrt{\lambda}.
	\]
\end{conjecture}

For real-analytic metrics, Donnelly--Fefferman proved both
bounds\cite{DonnellyFefferman1988}.  For general smooth metrics, Logunov
proved the sharp lower bound and a polynomial upper
bound\cite{LogunovLower2018,LogunovUpper2018}; the sharp
$C\sqrt{\lambda}$ upper bound remains open.  Recorded status:
\StatusPartial.

\plainproblemsection

The vacuum Einstein equation is
\[
\Ric(g)=0.
\]
Given constraint-satisfying initial data $(\Sigma,h,k)$, local well-posedness
produces a maximal globally hyperbolic development (MGHD).  Any precise
cosmic-censorship statement must specify the matter model, asymptotics,
topology used for genericity, and the regularity class of extensions.

\begin{conjecture}[Weak cosmic censorship, schematic form]
	For generic reasonable asymptotically flat initial data, singularities formed
	by the evolution are not visible from future null infinity; roughly, event
	horizons hide them and future null infinity remains complete.
\end{conjecture}

\plainproblemsection

\begin{conjecture}[Strong cosmic censorship, $C^k$ form]
	For generic initial data, the MGHD cannot be extended as part of a larger
	spacetime with a $C^k$ Lorentzian metric.
\end{conjecture}

The source records several model-dependent results: instability of naked
singularities for the spherically symmetric Einstein--scalar-field
model\cite{Christodoulou1999}; $C^2$ strong cosmic censorship but failure of
the $C^0$ version in a spherically symmetric
Einstein--Maxwell--real-scalar-field model\cite{LukOh2019}; and a 2025
nonsymmetric advance showing that a class of Christodoulou-type naked
singularities is hidden by a horizon under anisotropic
perturbations\cite{An2025}.  General nonsymmetric vacuum and broad matter
models remain \StatusPartial.

\plainproblemsection

Kerr spacetimes are parametrized by mass $M>0$ and angular-momentum parameter
$a$; the subextremal range is
\[
|a|<M.
\]

\begin{conjecture}[Nonlinear stability of Kerr]
	If vacuum Einstein initial data are sufficiently close, in a suitable
	weighted topology, to a subextremal Kerr initial slice, then the future
	development exists globally and, in the black-hole exterior, approaches
	another Kerr spacetime with slightly different parameters after radiation is
	emitted.
\end{conjecture}

The statement must be understood modulo coordinate gauge and with simultaneous
control of the event horizon, null infinity, radiation tails, and final mass
and angular momentum.  The source records nonlinear stability for slowly
rotating Kerr\cite{KlainermanSzeftel2021,GiorgiKlainermanSzeftel2022}.
It also records a June 2026 preprint claiming the full subextremal range
$|a|<M$ for its specified class of asymptotic data\cite{Hintz2026}.  Because
that preprint was only about one month old at the survey's verification date,
the source labels the full-range claim \StatusPreprint, rather than treating
it as an independently verified theorem.

\plainproblemsection

\begin{openproblem}
	\textbf{Parabolic Inverse Problem.}
	
	\textbf{Background:}
	
	Consider, for a bounded, connected open set $\Omega  \subset  {\mathbb{R}}^{n}$ , the following initial-boundary value problem for a parabolic equation
	
	(7.1)
	
	\[
	\left\{  \begin{array}{ll} {u}_{t} - \operatorname{div}\left( {a\nabla u}\right)  = f, & \text{ in }\;\Omega  \times  \left( {0, T}\right) , \\  u = 0, & \text{ on }\;\partial \Omega  \times  \left( {0, T}\right)  \cup  \Omega  \times  \{ 0\} . \end{array}\right.
	\]
	
	Here $a = a\left( x\right)$ is assumed to be a time independent scalar coefficient, satisfying uniform ellipticity
	
	\[
	{K}^{-1} \leq  a\left( x\right)  \leq  K.
	\]
	
	This background provides the framework for the following problem.
	The problem to be solved is: (Parabolic inverse problem). Let $n = 2$ . Given $f$ on all of $\Omega  \times  \left( {0, T}\right)$ , and given $u\left( {x, t}\right)$ for all $x \in  \Omega$ and for one or more fixed values of $t > 0$ , is a uniquely determined?
	
	The limitation of the space dimension to $n = 2$ was motivated by applications, more specifically: identification of transmissivity $a$ in groundwater flow by piezometric head measurements $u$ . The origin of the problem should be searched in the hydrogeology literature, and it might be dated much earlier, see for instance [35].
	
	The problem however maintains its interest in any dimension, and is still open in general. A solution under some special assumptions on boundary data and when $n = 1$ was given by Victor Isakov in [23].
	
	Here I propose a variant to Problem 7.1.
\end{openproblem}

\plainproblemsection

\begin{openproblem}
	\textbf{Parabolic Inverse Problem with Final Data.}
	
	\textbf{Background:}
	
	Consider, for a bounded, connected open set $\Omega  \subset  {\mathbb{R}}^{n}$ , the following initial-boundary value problem for a parabolic equation
	
	(7.1)
	
	\[
	\left\{  \begin{array}{ll} {u}_{t} - \operatorname{div}\left( {a\nabla u}\right)  = f, & \text{ in }\;\Omega  \times  \left( {0, T}\right) , \\  u = 0, & \text{ on }\;\partial \Omega  \times  \left( {0, T}\right)  \cup  \Omega  \times  \{ 0\} . \end{array}\right.
	\]
	
	Here $a = a\left( x\right)$ is assumed to be a time independent scalar coefficient, satisfying uniform ellipticity
	
	\[
	{K}^{-1} \leq  a\left( x\right)  \leq  K.
	\]
	
	This background provides the framework for the following problem.
	The problem to be solved is: (Parabolic inverse problem with final data). Consider the following initial-boundary value parabolic problem
	
	(7.2)
	
	\[
	\left\{  \begin{array}{lll} {u}_{t} - \operatorname{div}\left( {a\nabla u}\right)  = 0, & \text{ in } & \Omega  \times  \left( {0, T}\right) , \\  u = g\left( {x, t}\right) , & \text{ on } & \partial \Omega  \times  \left( {0, T}\right) , \\  u = h\left( x\right) , & \text{ on } & \Omega  \times  \{ 0\} . \end{array}\right.
	\]
	
	Given appropriately chosen functions $g$ and $h$ , and assuming a known on $\partial \Omega$ , does the knowledge of $u\left( {\cdot , T}\right)$ uniquely determines a?
	
	This problem can be seen as an extension to a non-stationary setting of an inverse elliptic problem with interior data, for which uniqueness can be proven. In such an elliptic context, it can be seen that some boundary information on the coefficient $a$ is needed, that is why I have added that kind of information also in the parabolic question.
\end{openproblem}

\plainproblemsection

\begin{openproblem}
	\textbf{Nodal Line of a Second Eigenfunction.}
	
	\textbf{Background:}
	
	In this Section and in the following one I shall discuss problems arising from Courant's Nodal Domain Theorem [17]. A brief introduction may be useful. Consider, in a bounded and connected open set $\Omega  \subset  {\mathbb{R}}^{n}, n \geq  2$ , the eigenvalue problem
	
	(5.1)
	
	\[
	\left\{  \begin{array}{lll} {\Delta u} + {\lambda u} = 0, & \text{ in } & \Omega , \\  u = 0, & \text{ on } & \partial \Omega . \end{array}\right.
	\]
	
	It is well known that there exists a complete set of eigenfunctions $\left\{  {u}_{m}\right\}$ , orthonormal in ${L}^{2}\left( \Omega \right)$ , with corresponding eigenvalues $\left\{  {\lambda }_{m}\right\}$ . The eigenvalues are all positive and can be arranged in a nondecreasing, diverging sequence $0 < {\lambda }_{1} < {\lambda }_{2} \leq  {\lambda }_{3} \leq  \ldots$ .
	
	Courant’s Nodal Domain Theorem states that every eigenfunction ${u}_{m}$ corresponding to the $m$ -th eigenvalue ${\lambda }_{m}$ has at most $m$ nodal domains, that is the set $\left\{  {x \in  \Omega  \mid  {u}_{m}\left( x\right)  \neq  0}\right\}$ has at most $m$ connected components.
	
	In particular, ${u}_{1}$ has constant sign, and any second eigenfunction ${u}_{2}$ (which, being orthogonal to ${u}_{1}$ , must change its sign) has exactly two nodal domains.
	
	L. E. Payne [32, Conjecture 5] posed the following conjecture.
	
	Conjecture 5.1 (Payne). Let $n = 2$ , for any second eigenfunction ${u}_{2}$ the nodal line
	
	\[
	\overline{\left\{  x \in  \Omega  \mid  {u}_{2}\left( x\right)  = 0\right\}  }
	\]
	
	is not a closed curve.
	
	Melas [29] proved the conjecture if $\Omega$ is convex and has smooth boundary, the smoothness assumption was removed in [5]. Hoffmann-Ostenhof, Hoffmann-Ostenhof and Nadirashvili [21] showed by an example that the conjecture may fail if $\Omega$ is not simply connected.
	
	The following remains an open problem.
	The problem to be solved is: Assume $\Omega$ be simply connected and with smooth boundary. Prove that the nodal line of any second eigenfunction is a simple open curve whose endpoints are two distinct points of $\partial \Omega$ .
\end{openproblem}

\plainproblemsection

\begin{openproblem}
	\textbf{The Insulating Crack Problem.}
	
	\textbf{Background:}
	
	Consider, in a bounded and connected open set $\Omega  \subset  {\mathbb{R}}^{3}$ a two-dimensional ori-entable smooth surface $\sum  \Subset  \Omega$ having a simple closed curve as boundary. We view $\sum$ as a fracture in a homogeneous electrically conducting body $\Omega$ . If we prescribe a stationary current density $\psi$ (having zero average) on $\partial \Omega$ then the electrostatic potential $u$ is governed by the following boundary value problem (suitably interpreted in a weak sense)
	
	(4.1)
	
	\[
	\left\{  \begin{array}{lll} {\Delta u} = 0, & \text{ in } & \Omega  \smallsetminus  \sum , \\  \nabla u \cdot  \nu  = 0, & \text{ on either side of } & \sum , \\  \nabla u \cdot  \nu  = \psi , & \text{ on } & \partial \Omega . \end{array}\right.
	\]
	
	The inverse problem is:
	The problem to be solved is: (The insulating crack problem). To find appropriately chosen current profiles ${\psi }_{1},\ldots {\psi }_{N}$ such that, letting ${u}_{1},\ldots {u}_{N}$ be the corresponding potentials, the boundary measurements ${\left. {u}_{1}\right| }_{\partial \Omega },\ldots {\left. {u}_{N}\right| }_{\partial \Omega }$ uniquely determine $\sum$ .
	
	When $\sum$ is a portion of a plane, Kubo [25] found a triple of suitable current profiles for which uniqueness holds, in [7] suitable pairs were found.
	
	When $\sum$ is allowed to be curved and the full Neumann-to-Dirichlet map
	
	\[
	{\mathcal{N}}_{\sum } : \psi  \rightarrow  {\left. u\right| }_{\partial \Omega }
	\]
	
	is known, uniqueness was proven by Eller [19].
\end{openproblem}

\plainproblemsection

\begin{openproblem}
	\textbf{Unique Continuation along Level Surfaces.}
	
	\textbf{Background:}
	
	This is probably the only problem of this collection for which I can claim some form of paternity. It originates from my work with Emmanuele Di Benedetto on the inverse problems of cracks [7], and it could be used to extend some of those results to the case of cracks in inhomogeneous media. It was first formulated in a paper with Alberto Favaron [10].
	
	Consider an elliptic equation
	
	(2.1)
	
	\[
	\operatorname{div}\left( {A\nabla u}\right)  = 0
	\]
	
	in a connected open set $\Omega  \subset  {\mathbb{R}}^{n}, n \geq  2$ . Here $A = A\left( x\right)$ is a symmetric matrix valued function, satisfying uniform ellipticity
	
	(2.2)
	
	\[
	{K}^{-1}{\left| \xi \right| }^{2} \leq  A\left( x\right) \xi  \cdot  \xi  \leq  K{\left| \xi \right| }^{2},\forall x,\xi  \in  {\mathbb{R}}^{n},
	\]
	
	and Lipschitz continuity
	
	(2.3)
	
	\[
	\left| {A\left( x\right)  - A\left( y\right) }\right|  \leq  E\left| {x - y}\right| ,\forall x, y \in  {\mathbb{R}}^{n},
	\]
	
	for some positive constants $K, E$ . These assumptions are stated in order to guarantee that the standard unique continuation property applies, see again [20]. I propose the following less standard form of unique continuation.
	
	This background provides the framework for the following problem.
	The problem to be solved is: (Unique continuation along level surfaces). Let $u, v$ be two non-costant solutions to (2.1), let $S \subset  \Omega$ be a connected(n - 1)-dimensional smooth hypersurface, and let $\sum$ be an open proper subset of $S$ . Suppose $u = 0$ on $S$ and $v = 0$ on $\sum$ . Is it true that $v = 0$ on all of $S$ ?
	
	If $A \equiv  I$ , then $u, v$ are harmonic and in this case, $S$ is an analytic hypersurface, see for instance [7, Appendix A], hence the answer is affirmative, by analytic continuation. Some smoothness on $S$ must be indeed assumed. This can be easily seen in the harmonic two-dimensional case:
	
	Let $\Omega  = {\mathbb{R}}^{2}, u = {xy}, v = y$ . In this case, we may pick $S = \{ x = 0, y \geq$ $0\}  \cup  \{ x \geq  0, y = 0\}$ , which is a simple curve with a corner point at the origin and $\sum  = \{ x > 0, y = 0\}  \subset  S$ is the positive horizontal semiaxis. Clearly $v$ vanishes on $\sum$ , but it does not vanish on all of $S \subset  \{ u = 0\}$ .
	
	From another point of view, it may be noted that a parallel could be drawn with the issue of unique continuation at fixed time for parabolic equations, which was treated jointly with Sergio Vessella [12].
\end{openproblem}

\plainproblemsection

\begin{openproblem}
	\textbf{The Inclusion Problem with Finite Data.}
	
	\textbf{Background:}
	
	\paragraph{Problem Setup.}
	
	\begin{itemize}
		\item \textbf{Domain:} $\Omega \subset \mathbb{R}^n$ ($n \geq 2$), smooth boundary.
		\item \textbf{Inclusion:} $D \Subset \Omega$, conductivity $k \neq 1$ (background conductivity $=1$).
	\end{itemize}
	
	
	\paragraph{PDE Model.}
	
	\[
	\begin{cases}
		\operatorname{div}\!\left(\big(1+(k-1)\chi_D\big)\nabla u\right)=0, & x \in \Omega,\\
		u=\varphi, & x \in \partial\Omega.
	\end{cases}
	\]
	
	
	\paragraph{Dirichlet-to-Neumann Map.}
	
	\[
	\Lambda_D : \varphi \mapsto \nabla u \cdot \nu \quad
	\text{on } \partial \Omega.
	\]
	
	Well-defined for $\varphi \in W^{1/2,2}(\partial \Omega)$.
	
	
	\paragraph{Inverse Problem.}
	
	\begin{itemize}
		\item \textbf{Goal:} Recover $D$ from $\Lambda_D$.
		\item \textbf{Known:} Uniqueness holds (Isakov).
		\item \textbf{Question:} Can limited data determine $D$?
	\end{itemize}
	
	
	\paragraph{Special Case ($n=2$).}
	
	\begin{itemize}
		\item \textbf{Seo’s Result:} If $D$ is a simply connected polygon, two suitable boundary data $\varphi_1, \varphi_2$ with their Neumann traces uniquely determine $D$.
	\end{itemize}
	
	
	\paragraph{Clever Choice of Data (Radó’s Condition).}
	
	\begin{itemize}
		\item Pair $(\varphi_1, \varphi_2)$ satisfies Radó’s condition if
	\end{itemize}
	
	\[
	\Phi = (\varphi_1, \varphi_2): \partial \Omega \to \Gamma
	\]
	
	is a homeomorphism onto the boundary of a convex set $G \subset \mathbb{R}^2$.
	
	\begin{itemize}
		\item Then $U=(u_1,u_2)$ is a homeomorphism $\Omega \to G$.
	\end{itemize}
	
	This background provides the framework for the following problem.
	The problem to be solved is: (The inclusion problem with finite data). Let $n = 2$ , let us assume a priori $D$ be simply connected and with smooth boundary. Let the pair of Dirichlet data ${\varphi }_{1},{\varphi }_{2}$ satisfy Radó’s condition, is D uniquely determined by the Neumann data $\nabla {u}_{1} \cdot  \nu ,\nabla {u}_{2} \cdot  \nu$ ?
	
	One clue that Radó's condition might be the appropriate one comes from the fact that it has shown to be effective in the germane inverse problem of cracks in dimension 2. The main results, and more references, can be found in the paper by Alvaro Diaz Valenzuela and myself [6], and in the one by Jin Keun Seo with Kim 24.
	
	Of course it would be meaningful to formulate Problem 3.2 also in higher dimensions. However it is not at all clear what might be a clever choice of Dirichlet data in such a case.
\end{openproblem}

\plainproblemsection

\begin{openproblem}
	\textbf{Weak Unique Continuation for the $p$-Laplace Equation.}
	
	\textbf{Background:}
	
	Given $1 < p < \infty$ we consider the so-called $p$ -Laplace equation
	
	(1.1)
	
	\[
	\operatorname{div}\left( {{\left| \nabla u\right| }^{p - 2}\nabla u}\right)  = 0
	\]
	
	in a connected open set $\Omega  \subset  {\mathbb{R}}^{n}, n \geq  2$ . The natural space in which we may search for weak (variational) solutions is ${W}^{1, p}\left( \Omega \right)$ . The problem that I want to address here is whether solutions to (1.1) satisfy the unique continuation property. However, this formulation is somewhat vague, because the unique continuation property can be expressed in many ways. More precise formulations are the following ones.
	
	This background provides the framework for the following problem.
	The problem to be solved is: (Weak unique continuation). Suppose $u$ solves (1.1) and vanishes on a nonempty open set $U \subsetneq  \Omega$ , does $u \equiv  0$ on all of $\Omega$ ?
\end{openproblem}

\plainproblemsection

\begin{openproblem}
	\textbf{Strong Unique Continuation for the $p$-Laplace Equation.}
	
	\textbf{Background:}
	
	Given $1 < p < \infty$ we consider the so-called $p$ -Laplace equation
	
	(1.1)
	
	\[
	\operatorname{div}\left( {{\left| \nabla u\right| }^{p - 2}\nabla u}\right)  = 0
	\]
	
	in a connected open set $\Omega  \subset  {\mathbb{R}}^{n}, n \geq  2$ . The natural space in which we may search for weak (variational) solutions is ${W}^{1, p}\left( \Omega \right)$ . The problem that I want to address here is whether solutions to (1.1) satisfy the unique continuation property. However, this formulation is somewhat vague, because the unique continuation property can be expressed in many ways. More precise formulations are the following ones.
	
	This background provides the framework for the following problem.
	The problem to be solved is: (Strong unique continuation). Suppose $u$ solves (1.1) and vanishes of infinite order at an interior point ${x}_{0}$ of $\Omega$ , does $u \equiv  0$ on all of $\Omega$ ?
	
	If $p = 2$ then (1.1) reduces to the classical Laplace’s equation, the solutions are harmonic functions which are well-known to be real analytic, and hence the strong unique continuation property holds. If $n = 2$ the strong unique continuation property holds true for any $p > 1$ . This was proved by Bojarski and Iwaniec [14] and by myself [3] by two different methods, a gap in the proof of Bojarski and Iwaniec was later filled by Manfredi [28].
	
	Thus the problem, in both formulations 1.1 and 1.2 remains open in the case $n \geq  3, p \neq  2$ .
	
	I should also mention that this problem was first proposed to me by Gene Fabes, in 1981, in an even stronger form:
\end{openproblem}

\plainproblemsection

\begin{openproblem}
	\textbf{The Muckenhoupt-Weight Question.}
	
	\textbf{Background:}
	
	Given $1 < p < \infty$ we consider the so-called $p$ -Laplace equation
	
	(1.1)
	
	\[
	\operatorname{div}\left( {{\left| \nabla u\right| }^{p - 2}\nabla u}\right)  = 0
	\]
	
	in a connected open set $\Omega  \subset  {\mathbb{R}}^{n}, n \geq  2$ . The natural space in which we may search for weak (variational) solutions is ${W}^{1, p}\left( \Omega \right)$ . The problem that I want to address here is whether solutions to (1.1) satisfy the unique continuation property. However, this formulation is somewhat vague, because the unique continuation property can be expressed in many ways. More precise formulations are the following ones.
	
	This background provides the framework for the following problem.
	The problem to be solved is: (Mukenhoupt). Suppose $u$ solves (1.1) and is not identically constant, is $\left| {\nabla u}\right|$ a Mukenhoupt weight?
	
	For a definition of Mukenhoupt weights, and the basics of the theory, I refer to Coifman and Feffermann [16].
	
	Again, the answer is affirmative if $p = 2$ , this can be viewed within the general theory of unique continuation for linear elliptic equations, a benchmark in this respect is due to Nico Garofalo and F.H. Lin [20]. When $n = 2$ , the positive answer can be found in a joint paper with Daniela Lupo and Edi Rosset [4].
\end{openproblem}

\plainproblemsection

\begin{openproblem}
	\textbf{Equivalence of Properties PI and PII.}
	
	\textbf{Background.}
	Let $D \subset \mathbb{R}^n$ be a region with regular boundary. Consider two properties:
	
	\begin{itemize}
		\item \textbf{(PI)} There exists a bounded, nontrivial temperature $u(x,t)$ solving
	\end{itemize}
	\[
	u_t = \Delta u, \quad (x,t) \in D \times [0,1],
	\]
	with
	\[
	u(x,1) \equiv 0, \quad x \in D.
	\]
	
	\begin{itemize}
		\item \textbf{(PII)} There exists a nontrivial harmonic function $v(x)$ in $D$ such that
	\end{itemize}
	\[
	v(x) = O\!\left(\exp(-|x|^2)\right), \quad |x|\to\infty.
	\]
	
	According to Gurarii and Matsaev [5], these two properties are \textbf{equivalent}, but no published proof exists.
	
	
	\textbf{Known results.}
	
	\begin{itemize}
		\item \textbf{For $n=1$:} PI and PII hold if and only if $D$ is a bounded interval (PI: Tychonov [7], PII: trivial).
		\item \textbf{For $n=2$:}
		\item If $D$ is a sector $\{ z: |\arg z| < \alpha \}$, then PII holds iff $\alpha < 45^\circ$.
		\item Escauriaza proved that for such sectors PI also holds, and more generally PII $\implies$ PI for cones (via the Appell transform).
		\item \textbf{Obstructions:} Li and Šverák [6] showed that PI fails in rotational cones with opening angle larger than
	\end{itemize}
	\[
	2\arccos(1/\sqrt{3}) \approx 109.52^\circ.
	\]
	\begin{itemize}
		\item \textbf{For PII:}
		\item Holds in wedges extended by $\mathbb{R}^{n-1}$.
		\item For rotational cones in $\mathbb{R}^n$, PII holds iff the opening angle is less than $90^\circ$.
	\end{itemize}
	
	
	\textbf{Open Question.}
	Are properties \textbf{PI} and \textbf{PII} equivalent for general regions $D \subset \mathbb{R}^n$?
\end{openproblem}

\begingroup
\renewcommand{\bibname}{References}
\let\chapter\subsection
\begin{thebibliography}{99}
	
	\bibitem{ClayNS2026}
	Clay Mathematics Institute,
	\emph{Navier--Stokes Equation},
	\url{https://www.claymath.org/millennium/navier-stokes-equation/}.
	
	\bibitem{Fefferman2000}
	C.~L. Fefferman,
	\emph{Existence and Smoothness of the Navier--Stokes Equation},
	Clay Millennium Problem official description:
	\url{https://www.claymath.org/wp-content/uploads/2022/06/navierstokes.pdf}.
	
	\bibitem{BealeKatoMajda1984}
	J.~T. Beale, T. Kato and A. Majda,
	\emph{Remarks on the breakdown of smooth solutions for the 3-D Euler equations},
	Commun. Math. Phys. \textbf{94} (1984), 61--66.
	\url{https://doi.org/10.1007/BF01212349}.
	
	\bibitem{ChenHou2025}
	J. Chen and T.~Y. Hou,
	\emph{Singularity formation in 3D Euler equations with smooth initial data and boundary},
	Proc. Natl. Acad. Sci. USA \textbf{122} (2025), e2500940122.
	\url{https://doi.org/10.1073/pnas.2500940122}.
	
	\bibitem{ConstantinETiti1994}
	P. Constantin, W. E and E.~S. Titi,
	\emph{Onsager's conjecture on the energy conservation for solutions of Euler's equation},
	Commun. Math. Phys. \textbf{165} (1994), 207--209.
	\url{https://doi.org/10.1007/BF02099744}.
	
	\bibitem{Isett2018}
	P. Isett,
	\emph{A proof of Onsager's conjecture},
	Ann. of Math. \textbf{188} (2018), 871--963.
	\url{https://annals.math.princeton.edu/2018/188-3/p04}.
	
	\bibitem{Isett2024}
	P. Isett,
	\emph{On the endpoint regularity in Onsager's conjecture},
	Anal. PDE \textbf{17} (2024), 2123--2159.
	\url{https://doi.org/10.2140/apde.2024.17.2123}.
	
	\bibitem{KiselevNazarovVolberg2007}
	A. Kiselev, F. Nazarov and A. Volberg,
	\emph{Global well-posedness for the critical 2D dissipative quasi-geostrophic equation},
	Invent. Math. \textbf{167} (2007), 445--453.
	\url{https://doi.org/10.1007/s00222-006-0010-3}.
	
	\bibitem{CaffarelliVasseur2010}
	L.~A. Caffarelli and A. Vasseur,
	\emph{Drift diffusion equations with fractional diffusion and the quasi-geostrophic equation},
	Ann. of Math. \textbf{171} (2010), 1903--1930.
	\url{https://annals.math.princeton.edu/wp-content/uploads/annals-v171-n3-p10-p.pdf}.
	
	\bibitem{BulutDong2024}
	A. Bulut and H. Dong,
	\emph{Global well-posedness for supercritical SQG with perturbations of radially symmetric data},
	arXiv:2402.19439 (2024).
	\url{https://arxiv.org/abs/2402.19439}.
	
	\bibitem{Tao2008}
	T. Tao,
	\emph{Why are solitons stable?},
	Bull. Amer. Math. Soc. \textbf{46} (2009), 1--33;
	arXiv:0802.2408.
	\url{https://arxiv.org/abs/0802.2408}.
	
	\bibitem{Chatterjee2013}
	S. Chatterjee,
	\emph{Invariant measures and the soliton resolution conjecture},
	Commun. Pure Appl. Math. \textbf{67} (2014), 1737--1842;
	arXiv:1203.4027.
	\url{https://arxiv.org/abs/1203.4027}.
	
	\bibitem{DuyckaertsEtAl2021}
	T. Duyckaerts, C. Kenig, Y. Martel and F. Merle,
	\emph{Soliton resolution for critical co-rotational wave maps and radial cubic wave equation},
	arXiv:2103.01293 (2021).
	\url{https://arxiv.org/abs/2103.01293}.
	
	\bibitem{JendrejLawrie2022}
	J. Jendrej and A. Lawrie,
	\emph{Soliton resolution for the energy-critical nonlinear wave equation in the radial case},
	arXiv:2203.09614 (2022).
	\url{https://arxiv.org/abs/2203.09614}.
	
	\bibitem{GassotGerardMiller2026}
	L. Gassot, P. G\'erard and P.~D. Miller,
	\emph{A proof of the soliton resolution conjecture for the Benjamin--Ono equation},
	arXiv:2601.10488 (2026).
	\url{https://arxiv.org/abs/2601.10488}.
	
	\bibitem{AmbrosioCabre2000}
	L. Ambrosio and X. Cabr\'e,
	\emph{Entire solutions of semilinear elliptic equations in $\R^3$ and a conjecture of De Giorgi},
	J. Amer. Math. Soc. \textbf{13} (2000), 725--739.
	\url{https://doi.org/10.1090/S0894-0347-00-00345-3}.
	
	\bibitem{Savin2009}
	O. Savin,
	\emph{Regularity of flat level sets in phase transitions},
	Ann. of Math. \textbf{169} (2009), 41--78.
	\url{https://annals.math.princeton.edu/2009/169-1/p02}.
	
	\bibitem{DelPinoKowalczykWei2011}
	M. del Pino, M. Kowalczyk and J. Wei,
	\emph{On De Giorgi's conjecture in dimension $N\geq9$},
	Ann. of Math. \textbf{174} (2011), 1485--1569.
	\url{https://annals.math.princeton.edu/2011/174-3/p03}.
	
	\bibitem{ChanEtAl2025}
	H. Chan, X. Fern\'andez-Real, A. Figalli and J. Serra,
	\emph{Global stable solutions to the free boundary Allen--Cahn and Bernoulli problems in 3D are one-dimensional},
	arXiv:2503.21245 (2025).
	\url{https://arxiv.org/abs/2503.21245}.
	
	\bibitem{Souplet2009}
	P. Souplet,
	\emph{The proof of the Lane--Emden conjecture in four space dimensions},
	Adv. Math. \textbf{221} (2009), 1409--1427.
	\url{https://doi.org/10.1016/j.aim.2009.02.014}.
	
	\bibitem{HuangZou2025}
	L.-H. Huang and W. Zou,
	\emph{Liouville-type theorems for stable solutions of the H\'enon--Lane--Emden system},
	arXiv:2512.16566 (2025); accepted version, J. Lond. Math. Soc.
	\url{https://arxiv.org/abs/2512.16566}.
	
	\bibitem{DonnellyFefferman1988}
	H. Donnelly and C. Fefferman,
	\emph{Nodal sets of eigenfunctions on Riemannian manifolds},
	Invent. Math. \textbf{93} (1988), 161--183.
	\url{https://doi.org/10.1007/BF01393691}.
	
	\bibitem{LogunovUpper2018}
	A. Logunov,
	\emph{Nodal sets of Laplace eigenfunctions: polynomial upper estimates of the Hausdorff measure},
	Ann. of Math. \textbf{187} (2018), 221--239.
	\url{https://annals.math.princeton.edu/2018/187-1/p04}.
	
	\bibitem{LogunovLower2018}
	A. Logunov,
	\emph{Nodal sets of Laplace eigenfunctions: proof of Nadirashvili's conjecture and of the lower bound in Yau's conjecture},
	Ann. of Math. \textbf{187} (2018), 241--262.
	\url{https://annals.math.princeton.edu/2018/187-1/p05}.
	
	\bibitem{Christodoulou1999}
	D. Christodoulou,
	\emph{The instability of naked singularities in the gravitational collapse of a scalar field},
	Ann. of Math. \textbf{149} (1999), 183--217.
	\url{https://arxiv.org/abs/math/9901147}.
	
	\bibitem{LukOh2019}
	J. Luk and S.-J. Oh,
	\emph{Strong cosmic censorship in spherical symmetry for two-ended asymptotically flat initial data I},
	Ann. of Math. \textbf{190} (2019), 1--111.
	\url{https://annals.math.princeton.edu/2019/190-1/p01}.
	
	\bibitem{An2025}
	X. An,
	\emph{Naked singularity censoring with anisotropic apparent horizon},
	Ann. of Math. \textbf{201} (2025), 775--908.
	\url{https://annals.math.princeton.edu/2025/201-3/p03}.
	
	\bibitem{KlainermanSzeftel2021}
	S. Klainerman and J. Szeftel,
	\emph{Kerr stability for small angular momentum},
	Pure Appl. Math. Q. \textbf{19} (2023), 791--1678;
	arXiv:2104.11857.
	\url{https://arxiv.org/abs/2104.11857}.
	
	\bibitem{GiorgiKlainermanSzeftel2022}
	E. Giorgi, S. Klainerman and J. Szeftel,
	\emph{Wave equations estimates and the nonlinear stability of slowly rotating Kerr black holes},
	arXiv:2205.14808 (2022).
	\url{https://arxiv.org/abs/2205.14808}.
	
	\bibitem{Hintz2026}
	P. Hintz,
	\emph{Nonlinear stability of subextremal Kerr black holes},
	arXiv:2606.28253 (2026), preprint.
	\url{https://arxiv.org/abs/2606.28253}.
	
\end{thebibliography}
\endgroup

\clearpage
\section{Complex Analysis}

\plainproblemsection

\begin{openproblem}
	\textbf{Zeros of Meromorphic Functions of a Special Form.}
	
	Let $\{z_n\} \to \infty$ be a sequence in the complex plane, and consider the meromorphic function
	
	\[
	f(z) = \sum_{k=1}^{\infty} \frac{a_k}{z - z_k},
	\qquad \sum_{k=1}^{\infty} |a_k| < \infty,
	\qquad \sum_{k=1}^{\infty} a_k \neq 0.
	\]
	
	\textbf{Open Question:}
	Is it true that every such function $f$ must have zeros?
	
	
	\textbf{Known results:}
	\begin{itemize}
		\item (a) The answer is yes if $f$ is of finite order [1].
		\item (b) The answer is yes if $a_k > 0$ for all $k$ [2].
	\end{itemize}
\end{openproblem}

\plainproblemsection

\begin{openproblem}
	\textbf{Wronskians of More Than Two Real Polynomials.}
	
	Known theorem: Let $f, g$ be real polynomials, and suppose that their Wronskian
	
	$W\left( {f, g}\right)  = {f}^{\prime }g - f{g}^{\prime }$
	
	has all real zeros. If $I$ is a real interval containing no zeros of $W(f,g)$, then any linear combination $a f + b g$ has at most one root on $I$.
	
	My question is: can this result be generalized to more than two polynomials?
	
	For example, consider three real polynomials $f, g, h$: if their Wronskian
	
	$W\left( {f, g, h}\right)$
	
	has all real zeros, and $I$ is a real interval containing no zeros of the Wronskian, does any linear combination $a f + b g + c h$ have at most two roots on $I$? This case seems to be an open conjecture.
\end{openproblem}

\plainproblemsection

\begin{openproblem}
	\textbf{Polynomial-Exponential Solutions of a Second-Order Differential Equation.}
	
	\textbf{Background:}
	
	Let $h$ and $p$ be two polynomials. When $y = p\left( x\right) {e}^{h\left( x\right) }$ satisfies a second order differential equation
	
	\begin{equation*}
		{y}^{\prime \prime } + {Py} = 0, \tag{1}
	\end{equation*}
	
	where $P$ is a polynomial? Substitution gives
	
	\begin{equation*}
		\frac{{p}^{\prime \prime }}{p} + 2\frac{{p}^{\prime }}{p}{h}^{\prime } + {h}^{\prime \prime } + {h}^{\prime 2} + P = 0. \tag{2}
	\end{equation*}
	
	Such $P$ exists if and only if ${p}^{\prime \prime } + 2{p}^{\prime }{h}^{\prime }$ is divisible by $p$ .
	
	Another criterion is obtained if we consider the second solution ${y}_{1}$ of (1) which is linearly independent of $y$ . This second solution can be found from the condition
	
	\begin{equation*}
		y{y}_{1}^{\prime } - {y}^{\prime }{y}_{1} = 1 \tag{3}
	\end{equation*}
	
	Solving (3) with respect to ${y}_{1}$ we obtain
	
	\begin{equation*}
		{y}_{1} = p{e}^{h}\int {p}^{-2}{e}^{-{2h}}. \tag{4}
	\end{equation*}
	
	As all solutions of (1) must be entire functions, we conclude that all residues of ${p}^{-2}{e}^{-{2h}}$ must vanish. This condition is necessary and sufficient for $y = p{e}^{h}$ to satisfy equation (1) with some $P$ . Indeed, if ${y}_{1}$ defined by (4) is entire, then $y$ and ${y}_{1}$ is a pair of entire functions whose Wronski deterinant is 1, so this pair must satisfy a differential equation (1) with entire $P$ and asymptotics at infinity show that $P$ must be a polynomial.
	
	Now suppose that $h$ is an odd polynomial of degree 3, which we write in the form
	
	\begin{equation*}
		h\left( x\right)  = {x}^{3}/3 + {bx} \tag{5}
	\end{equation*}
	
	Suppose that all residues of ${p}^{-2}{e}^{-{2h}}$ vanish. Then the integral $\int {p}^{-2}{e}^{-{2h}}$ is a meromorphic function in the plane. Surprisingly, the integral of some linear
	
	combination
	
	\[
	\int \left( {{p}^{2}\left( {-x}\right) {e}^{-{2h}\left( x\right) } - c{p}^{-2}\left( x\right) {e}^{-{2h}\left( x\right) }}\right)
	\]
	
	is not only meromorphic but is an elementary function! Here $c$ is a constant depending on $p$ .
	
	\textbf{Conjecture}.
	Let $h$ be given by (5). Let $p$ be a polynomial. All residues of ${p}^{-2}{e}^{-{2h}}$ vanish if and only if there exist a constant $c$ and a polynomial $q$ such
	
	that
	
	\[
	\left( {{p}^{2}\left( {-x}\right)  - \frac{c}{{p}^{2}\left( x\right) }}\right) {e}^{-{2h}\left( x\right) } = \frac{d}{dx}\left( {\frac{q\left( x\right) }{p\left( x\right) }{e}^{-{2h}\left( x\right) }}\right) .
	\]
	
	In other words:
	
	\[
	{p}^{2}\left( x\right) {p}^{2}\left( {-x}\right)  - c = {q}^{\prime }\left( x\right) p\left( x\right)  - q\left( x\right) {p}^{\prime }\left( x\right)  - {2q}\left( x\right) p\left( x\right) {h}^{\prime }\left( x\right) .
	\]
	
	It is known [1] that for given $h$ of the form (5) there exist polynomials $p$ of any given degree such that all residues of ${p}^{-2}{e}^{-{2h}}$ vanish. These polynomials $p$ have simple roots. We verified the Conjecture for $1 \leq  n \leq  4$ where $n = \deg p$ by symbolic computation using Maple. We don't know whether there is any analog of this conjecture for other polynomials $h$ .
	
	Substituting $p\left( x\right)  = {x}^{n} + a{x}^{n - 1} + \ldots$ into (2) and using (5), we conclude that
	
	\begin{equation*}
		P\left( x\right)  =  - {h}^{\prime 2} - {h}^{\prime \prime } - {2nx} + {2a} =  - {x}^{4} - 2{x}^{2}b - 2\left( {n + 1}\right) x - {b}^{2} + {2a}. \tag{6}
	\end{equation*}
	
	Substituting $y = {p}_{n}{e}^{h}$ with $h$ as in (5) to (1) with this $P$ we obtain a polynomial relation between $b$ and $\lambda  \mathrel{\text{:=}} {b}^{2} - {2a}$ ,
	
	\begin{equation*}
		{Q}_{n}\left( {b,\lambda }\right)  = 0. \tag{7}
	\end{equation*}
	
	We have ${\deg }_{\lambda }{Q}_{n} = n + 1$ ,[1]. For every $b$ and every $\lambda$ satisfying this equation, the differential equation (1), with $P$ as in (6), has a solution $y = {p}_{n}{e}^{h}$ where $\deg {p}_{n} = n$ .
	
	Functions ${y}_{n} = {p}_{n}{e}^{h}$ are eigenfunctions of the operator
	
	\begin{equation*}
		{y}^{\prime \prime } - \left( {{x}^{4} + {2b}{x}^{2} + 2\left( {n + 1}\right) x}\right) y \tag{8}
	\end{equation*}
	
	with eigenvalue $\lambda$ . For each non-negative integer $n$ , and generic $b$ , the operator (8) has $n + 1$ eigenfunctions of the form ${p}_{n}{e}^{h}$ with eigenvalues $\lambda$ which are solutions of (7).
	
	We assume that ${Q}_{n}$ is monic as a polynomial of $\lambda$ , and ${p}_{n}$ is a monic polynomial of $x$ . Constant $c$ in the Conjecture, turns out to be
	
	\[
	c\left( {b,\lambda }\right)  = {\left( -1\right) }^{n}{2}^{-{2n}}\frac{\partial }{\partial \lambda }{Q}_{n}.
	\]
	
	This is also confirmed only by symbolic computation for small $n$ .
\end{openproblem}

\subsection*{References}

\begin{enumerate}[label={[\arabic*]},leftmargin=*,itemsep=0.4em]
	\item The three supplied complex-analysis problem files. Full
	bibliographic data for their source-local numerical citations were not
	included in the supplied files and are therefore not reconstructed here.
\end{enumerate}

\clearpage
\section{Matrix Theory and Matrix Inequalities}

\plainproblemsection

\textbf{The permanental dominance conjecture.}

(The permanental dominance conjecture) Suppose \( G \) is a subgroup of \( {S}_{n} \) and \( \chi \) is an irreducible character of \( G \) . Then for any positive semidefinite matrix \( A \) of order \( n \) ,

\[
\text{ per }A \geq  {\bar{d}}_{\chi }\left( A\right) \text{ . }
\]

\textbf{Background.} Let \( {S}_{n} \) denote the symmetric group on \( \{ 1,2,\ldots , n\} \) and \( {M}_{n} \) denote the set of complex matrices of order \( n \) . Suppose \( G \) is a subgroup of \( {S}_{n} \) and \( \chi \) is a character of \( G \) . The generalized matrix function \( {d}_{\chi } : {M}_{n} \rightarrow  \mathbf{C} \) is defined by

\[
{d}_{\chi }\left( A\right)  = \mathop{\sum }\limits_{{\sigma  \in  G}}\chi \left( \sigma \right) \mathop{\prod }\limits_{{i = 1}}^{n}{a}_{{i\sigma }\left( i\right) },
\]

where \( A = \left( {a}_{ij}\right) \) . Incidental to his work on group representation theory, Schur introduced this notion. For \( G = {S}_{n} \) , if \( \chi \) is the alternating character then \( {d}_{\chi } \) is the determinant while if \( \chi \) is the principal character then \( {d}_{\chi } \) is the permanent

\[
\operatorname{per}A = \mathop{\sum }\limits_{{\sigma  \in  {S}_{n}}}\mathop{\prod }\limits_{{i = 1}}^{n}{a}_{{i\sigma }\left( i\right) }.
\]

When \( \chi \) is the principal character of \( G = \{ e\} \) where \( e \) is the identity permutation in \( {S}_{n},{d}_{\chi } \) is Hadamard’s function \( h\left( A\right) \) .

In 1907 Fischer proved that if the matrix

\[
A = \left( \begin{array}{ll} {A}_{1} & B \\  {B}^{ * } & {A}_{2} \end{array}\right)
\]

is positive semidefinite with \( {A}_{1} \) and \( {A}_{2} \) square, then

\[
\det A \leq  \left( {\det {A}_{1}}\right) \left( {\det {A}_{2}}\right) .
\]

Hadamard's inequality follows from this inequality immediately. In 1918 Schur obtained the following generalization of Fischer's inequality:

\[
\chi \left( e\right) \det A \leq  {d}_{\chi }\left( A\right)
\]

for positive semidefinite \( A \) . Let \( G \) be a subgroup of \( {S}_{n} \) and let \( \chi \) be an irreducible character of \( G \) . The normalized generalized matrix function is defined as

\[
{\bar{d}}_{\chi }\left( A\right)  = {d}_{\chi }\left( A\right) /\chi \left( e\right) .
\]

Since any character of \( G \) is a sum of irreducible characters, Schur’s inequality is equivalent to

\[
\det A \leq  {\bar{d}}_{\chi }\left( A\right)
\]

for positive semidefinite \( A \) . In 1963, M. Marcus proved the permanental analog of Hadamard's inequality

\[
\text{ per }A \geq  h\left( A\right)
\]

and E.H. Lieb proved the permanental analog of Fischer's inequality

\[
\text{ per }A \geq  \left( {\operatorname{per}{A}_{1}}\right) \left( {\operatorname{per}{A}_{2}}\right)
\]

three years later, where \( A \) is positive semidefinite. These results naturally led to the following conjecture which was first published by Lieb [45] in 1966:

\textbf{Conjecture (The permanental dominance conjecture).} Suppose \( G \) is a subgroup of \( {S}_{n} \) and \( \chi \) is an irreducible character of \( G \) . Then for any positive semidefinite matrix \( A \) of order \( n \) ,

\[
\text{ per }A \geq  {\bar{d}}_{\chi }\left( A\right) \text{ . }
\]

A lot of work has been done on this conjecture. It has been confirmed for every irreducible character of \( {S}_{n} \) with \( n \leq  {13} \) . The reader is referred to [22 section 3] and the references therein for more details and recent progress.




\plainproblemsection

\textbf{The Marcus-de Oliveira conjecture.}

Let \( {S}_{n} \) denote the symmetric group on \( \{ 1,2,\ldots , n\} \) and co \( \Omega \) denote the convex hull of a set \( \Omega \) in the complex plane. In 1973 Marcus [50] and in 1982 de Oliveira [56] independently made the following

\textbf{Conjecture.} Let \( A, B \) be normal complex matrices of order \( n \) with eigenvalues \( {x}_{1},\ldots ,{x}_{n} \) and \( {y}_{1},\ldots ,{y}_{n} \) respectively. Then

\[
\det \left( {A + B}\right)  \in  \operatorname{co}\left\{  {\mathop{\prod }\limits_{{i = 1}}^{n}\left( {{x}_{i} + {y}_{\sigma \left( i\right) }}\right)  : \sigma  \in  {S}_{n}}\right\}  .
\]

It is known that Conjecture 8 is true in many special cases, e.g., (1) \( A, B \) are Hermitian [30]; (2) all the eigenvalues have the same modulus, \( \left| {x}_{1}\right|  = \cdots  = \left| {x}_{n}\right|  = \left| {y}_{1}\right|  = \cdots  = \left| {y}_{n}\right| \) [8]; (3) \( A + B \) is singular [26]. See \( \left\lbrack  {6,7}\right\rbrack \) for more verified cases.

\plainproblemsection

\textbf{Permanents of Hadamard matrices.}

In 1974 Wang [65] posed the following

\textbf{Conjecture.} Can the permanent of a Hadamard matrix of order \( n \) vanish for \( n > 2 \) ?
Here A Hadamard matrix is a square matrix with entries equal to \( \pm  1 \) whose rows and hence columns are mutually orthogonal. In other words, a Hadamard matrix of order \( n \) is a \( \{ 1, - 1\} \) -matrix \( A \) satisfying

\[
A{A}^{T} = {nI}
\]

where \( I \) is the identity matrix.


\plainproblemsection

\textbf{The Bessis-Moussa-Villani trace conjecture.}

In 1975, while studying partition functions of quantum mechanical systems, Bessis, Moussa and Villani [10] formulated the conjecture that if \( A, B \) are Hermitian matrices of the same order with \( B \) positive semidefinite then the function

\[
f\left( t\right)  = \operatorname{Tr}\exp \left( {A - {tB}}\right)
\]

is the Laplace transform of a positive measure on \( \lbrack 0,\infty ) \) , where \( t \) is a real variable and Tr means trace. Recently, Lieb and Seiringer [46] has proved that this conjecture is equivalent to the following

\textbf{Conjecture (Bessis-Moussa-Villani).} Let \( A, B \) be positive semidefinite matrices of order \( n \) and let \( k \) be a positive integer. Then the polynomial \( p\left( t\right)  = \operatorname{Tr}{\left( A + tB\right) }^{k} \) has all nonnegative coefficients.

The following cases of this conjecture are proved: (1) \( k \leq  5 \) and all \( n;\left( 2\right) n = 2 \) and any \( k \) . See [38] and the references therein for many partial results and recent advances.

\plainproblemsection

\textbf{The S-matrix conjecture.}

An S-matrix of order \( n \) is a 0-1 matrix formed by taking a Hadamard matrix of order \( n + 1 \) in which the entries in the first row and column are 1, changing 1's to 0's and -1's to 1's, and deleting the first row and column. Let \( \parallel  \cdot  {\parallel }_{F} \) denote the Frobenius norm. In 1976 Sloane and Harwit [60] made the following conjecture. See also [37, p.59].

\textbf{Conjecture.} If \( A \) is a nonsingular matrix of order \( n \) all of whose entries are in the interval \( \left\lbrack  {0,1}\right\rbrack \) , then

\[
{\begin{Vmatrix}{A}^{-1}\end{Vmatrix}}_{F} \geq  \frac{2n}{n + 1}.
\]

Equality holds if and only if \( A \) is an \( S \) -matrix.

I do not know of any results on Conjecture 11. This problem arose from weighing designs in optics and statistics.

\plainproblemsection

\textbf{Ryser's conjecture on minimum values of permanents.}

Let \( {\Lambda }_{n}^{k} \) be the set of those \( n \times  n \) 0-1 matrices with each row and column having exactly \( k \) 1’s. In 1978 Ryser (in [53]) posed the following

\textbf{Conjecture.} If \( {\Lambda }_{v}^{k} \) contains incidence matrices of \( \left( {v, k,\lambda }\right) \) -designs, then the permanent takes its minimum in \( {\Lambda }_{v}^{k} \) at one of these incidence matrices.

Wanless [67] showed by computer enumeration that this conjecture is true for \( v \leq  {12} \) .

\plainproblemsection

\textbf{Foregger's conjecture on minimum values of permanents.}

A fully indecomposable square matrix \( A \) is called nearly decomposable if whenever a nonzero entry of \( A \) is replaced with a 0, the resulting matrix is partly decomposable. In 1980 Foregger [31] made the following

\textbf{Conjecture.} If \( A \) is a nearly decomposable doubly stochastic matrix of order \( n \) , then

\[
\text{ per }A \geq  {2}^{1 - n}\text{ . }
\]

Note that this lower bound can be attained at \( A = \left( {I + P}\right) /2 \) where \( P \) is the permutation matrix corresponding to the permutation cycle \( \left( {{1234}\cdots n}\right) \) . Foregger [31] proved the cases \( 2 \leq  n \leq  9 \) .

\plainproblemsection

\textbf{Dittert's conjecture on permanents.}

Let \( {K}_{n} \) be the set of \( n \times  n \) nonnegative matrices with the sum of their entries equal to \( n \) . Define the function \( \phi \) on \( n \times  n \) matrices by

\[
\phi \left( A\right)  = \mathop{\prod }\limits_{{i = 1}}^{n}{r}_{i} + \mathop{\prod }\limits_{{j = 1}}^{n}{c}_{j} - \operatorname{per}A
\]

where \( {r}_{1},\ldots ,{r}_{n} \) and \( {c}_{1},\ldots ,{c}_{n} \) are the row and column sums of \( A \) respectively. In 1983 Dittert (in [54]) posed the following

\textbf{Conjecture.}

\[
\max \left\{  {\phi \left( A\right)  : A \in  {K}_{n}}\right\}   = 2 - \frac{n!}{{n}^{n}},
\]

and the maximum is attained only for the matrix with each entry equal to \( 1/n \) .

Sinkhorn [59] proved the case \( n = 2 \) and Hwang [40] proved the case \( n = 3 \) .

\plainproblemsection

\textbf{The Brualdi-Li conjecture on tournament matrices.}

A tournament matrix is a square 0-1 matrix A satisfying \( A + {A}^{T} = \; J - I \) where \( J \) is the all ones matrix. Such matrices arise from the results of round robin competitions. An \( n \times  n \) tournament matrix \( A \) is called regular if each of the row sums of \( A \) is \( \left( {n - 1}\right) /2 \) . For even \( n \) , an \( n \times  n \) tournament matrix is called almost regular if half of its row sums are \( \left( {n - 2}\right) /2 \) and the other half are \( n/2 \) . It is known [19] that for odd \( n \) the regular tournament matrices maximize the Perron root over the class of \( n \times  n \) tournament matrices. For even \( n \) , it is not known which tournament matrices maximize the Perron root. Let \( {U}_{k} \) be the strictly upper triangular matrix of order \( k \) with ones above the main diagonal. In 1983, Brualdi and Li [20] made the following

\textbf{Conjecture.} For even \( n \) , the matrix

\[
\left\lbrack  \begin{matrix} {U}_{n/2} & {U}_{n/2}^{T} \\  {U}_{n/2}^{T} + I & {U}_{n/2} \end{matrix}\right\rbrack
\]

maximizes the Perron root over the class of tournament matrices of order \( n \) .

See [28] and the references therein for some partial results. Kirkland [43] has proved that for all sufficiently large even orders \( n \) the maximizers are almost regular.



\plainproblemsection

\textbf{The Grone-Merris conjecture on Laplacian spectra.}

Let \( G \) be a graph of order \( n \) and let \( d\left( G\right)  = \left( {{d}_{1},\ldots ,{d}_{n}}\right) \) be the degree sequence of \( G \) , where \( {d}_{1},\ldots ,{d}_{n} \) are the degrees of the vertices of \( G \) . Let \( A\left( G\right) \) be the adjacency matrix of \( G \) and denote \( D\left( G\right)  = \; \operatorname{diag}\left( {{d}_{1},\ldots ,{d}_{n}}\right) \) . Then the matrix \( L\left( G\right)  = D\left( G\right)  - A\left( G\right) \) is called the Laplacian matrix of \( G \) , and \( s\left( G\right)  = \left( {{\lambda }_{1},\ldots ,{\lambda }_{n}}\right) \) is called the Laplacian spectrum, where \( {\lambda }_{1},\ldots ,{\lambda }_{n} \) are the eigenvalues of \( L\left( G\right) \) .

For a sequence \( x = \left( {{x}_{1},\ldots ,{x}_{n}}\right) \) with nonnegative integer components, the conjugate sequence of \( x \) is \( {x}^{ \star  } = \left( {{x}_{1}^{ \star  },\ldots ,{x}_{n}^{ \star  }}\right) \) , where

\[
{x}_{j}^{ \star  } = \left| \left\{  {i : {x}_{i} \geq  j}\right\}  \right| .
\]

Denote by \( {d}^{ \star  }\left( G\right) \) the conjugate sequence of the degree sequence \( d\left( G\right) \) . Use \( x \prec  y \) to mean that \( x \) is majorized by \( y \) . In 1994, Grone and Merris [35] made the following

\textbf{Conjecture.} Let \( G \) be a connected graph. Then \( s\left( G\right)  \prec  {d}^{ \star  }\left( G\right) \) .

If we arrange the components of \( s\left( G\right) \) and \( {d}^{ \star  }\left( G\right) \) in decreasing order, then it is known that

\[
{\lambda }_{1} \leq  {d}_{1}^{ \star  },\;{\lambda }_{1} + {\lambda }_{2} \leq  {d}_{1}^{ \star  } + {d}_{2}^{ \star  }.
\]

\plainproblemsection

\textbf{The CP-rank conjecture.}

An \( n \times  n \) real matrix \( A \) is called completely positive (CP) if, for some \( m \) , there exists an \( n \times  m \) nonnegative matrix \( B \) such that \( A = \; B{B}^{T} \) . The smallest such \( m \) is called the \( {CP} \) -rank of \( A \) . CP matrices have applications in block designs. In 1994, Drew, Johnson and Loewy [25] posed the following

\textbf{Conjecture.} If \( A \) is a CP matrix of order \( n \geq  4 \) , then

\[
{CP}\text{ -rank }\left( A\right)  \leq  \left\lfloor  {{n}^{2}/4}\right\rfloor  \text{ . }
\]

This conjecture is true in the following cases: (1) \( n = 4 \) [51]; (2) \( n = 5 \) and \( A \) has at least one zero entry [49]; (3) the graph of \( A \) does not contain an odd cycle of length greater than 4 [24, 9]. It is known [25, p.309] that for each \( n \geq  4 \) the conjectured upper bound \( \left\lfloor  {{n}^{2}/4}\right\rfloor \) can be attained.

\plainproblemsection

\textbf{Bhatia-Kittaneh's question on singular values.}

Denote by \( {s}_{1}\left( X\right)  \geq  {s}_{2}\left( X\right)  \geq  \cdots \) the ordered singular values of a complex matrix \( X \) . In 2000, Bhatia and Kittaneh [17] asked the following

\textbf{Question.} Let \( A, B \) be positive semidefinite matrices of order n. Is it true that

\[
{s}_{j}^{1/2}\left( {AB}\right)  \leq  \frac{1}{2}{s}_{j}\left( {A + B}\right) ,\;j = 1,2,\ldots , n?
\]

The case \( n = 2 \) is known to be true [17]. Since the square function \( f\left( t\right)  = {t}^{2} \) is operator convex on \( \mathbf{R} \) , this inequality is stronger than the known inequality

\[
2{s}_{j}\left( {X{Y}^{ * }}\right)  \leq  {s}_{j}\left( {{X}^{ * }X + {Y}^{ * }Y}\right) ,\;j = 1,2,\ldots , n
\]

for any complex matrices \( X, Y \) of order \( n \) due to the same authors [18].

\plainproblemsection

\textbf{Convergence of the iterated Aluthge transforms.}

Every square complex matrix \( A \) has the polar decomposition \( A = \; {UP} \) where \( U \) is unitary and \( P \) is positive semidefinite. The Aluthge transform of \( A \) is

\[
\Delta \left( A\right)  = {P}^{1/2}U{P}^{1/2}.
\]

Though the unitary factor in the polar decomposition is not unique when \( A \) is singular, the Aluthge transform is well defined, that is, it does not depend on the choice made for the unitary factor. If \( A \) is normal, then \( \Delta \left( A\right)  = A \) . For \( 0 < \lambda  < 1 \) , the \( \lambda \) -Aluthge transform

\[
{\Delta }_{\lambda }\left( A\right)  = {P}^{\lambda }U{P}^{1 - \lambda }
\]

is also well defined. Note that \( \Delta  = {\Delta }_{1/2} \) . Let \( B\left( H\right) \) be the algebra of bounded linear operators on a Hilbert space \( H \) . The Aluthge transform can also be defined for operators in \( B\left( H\right) \) . In 2000, Jung, Ko and Pearcy [41] conjectured that for any \( T \in  B\left( H\right) \) , the sequence \( {\left\{  {\Delta }^{m}\left( T\right) \right\}  }_{m = 1}^{\infty } \) is norm convergent to an operator. Here \( {\Delta }^{1}\left( T\right)  = \Delta \left( T\right) \) and \( {\Delta }^{m}\left( T\right)  = \; \Delta \left( {{\Delta }^{m - 1}\left( T\right) }\right) , m = 2,3,\ldots \) However Cho, Jung and Lee [23] showed that this conjecture is false for infinite dimensional Hilbert spaces. So there remains the possibility that it holds in finite dimensions:

\textbf{Conjecture.} Let \( A \) be a square complex matrix. Then the sequence \( {\left\{  {\Delta }^{m}\left( A\right) \right\}  }_{m = 1}^{\infty } \) converges.

Note that \( \Delta \left( A\right) \) and \( A \) have the same eigenvalues. It is known ([41, Prop. 1.10], [42, Prop. 3.1], [2, Thm 1]) that if the Aluthge sequence of a matrix converges, then the limit matrix is normal. These two properties make the Aluthge transform more interesting.

Ando and Yamazaki [4] verified Conjecture 20 when \( A \) is of order 2. See [42] for some special cases. Huang and Tam [39] proved that if the nonzero eigenvalues of \( A \) have distinct moduli, then the \( \lambda \) -Aluthge sequence \( {\left\{  {\Delta }_{\lambda }^{m}\left( A\right) \right\}  }_{m = 1}^{\infty } \) converges. Huang and Tam [39] also posed the following

\textbf{Conjecture.} For any square complex matrix \( A \) and \( 0 < \lambda  < 1 \) ,

\[
{\begin{Vmatrix}{A}^{ * }A - A{A}^{ * }\end{Vmatrix}}_{F} \geq  {\begin{Vmatrix}{\Delta }_{\lambda }{\left( A\right) }^{ * }{\Delta }_{\lambda }\left( A\right)  - {\Delta }_{\lambda }\left( A\right) {\Delta }_{\lambda }{\left( A\right) }^{ * }\end{Vmatrix}}_{F}.
\]



\plainproblemsection

\textbf{Monotonicity of a geometric mean of positive definite matrices.}

The following geometric mean of three or more positive definite matrices has recently been defined in [55] and [15] independently using a geometric approach. See also [12, 16]. Here we give only the basic idea. Denote by \( {\mathbf{P}}_{n} \) the set of positive definite matrices of order \( n \) . We consider \( {\mathbf{P}}_{n} \) as a differentiable manifold. The distance \( \delta \left( {A, B}\right) \) between \( A, B \in  {\mathbf{P}}_{n} \) is the infimum of lengths of curves in \( {\mathbf{P}}_{n} \) that connect \( A \) to \( B \) . It can be proved that \( \delta \left( {A, B}\right)  = {\begin{Vmatrix}\log \left( {A}^{-\frac{1}{2}}B{A}^{-\frac{1}{2}}\right) \end{Vmatrix}}_{F} \) . Given \( {A}_{i} \in  {\mathbf{P}}_{n} \) , \( i = 1,2,\ldots , k \) , there is a unique matrix in \( {\mathbf{P}}_{n} \) , denoted \( G\left( {{A}_{1},\ldots ,{A}_{k}}\right) \) , that minimizes the function

\[
f\left( X\right)  = \mathop{\sum }\limits_{{i = 1}}^{k}{\delta }^{2}\left( {{A}_{i}, X}\right) .
\]

Then \( G\left( {{A}_{1},\ldots ,{A}_{k}}\right) \) is called the geometric mean of \( {A}_{1},\ldots ,{A}_{k} \) . This geometric mean is symmetric, invariant under congruence, and continuous. Let " \( \leq \) " be the Loewner partial order. In 2006, Bhatia and Holbrook [15, p.616] made the following

\textbf{Conjecture.} The mean \( G \) is monotone with respect to its arguments, i.e., if \( {A}_{i},{B}_{i} \in  {\mathbf{P}}_{n} \) satisfy \( {A}_{i} \leq  {B}_{i}, i = 1,\ldots , k \) , then

\[
G\left( {{A}_{1},\ldots ,{A}_{k}}\right)  \leq  G\left( {{B}_{1},\ldots ,{B}_{k}}\right) .
\]

Another geometric mean of three or more positive definite matrices has been defined in [3].

\subsection*{References}

\begin{enumerate}[label={[\arabic*]},leftmargin=*,itemsep=0.4em]
	\item The supplied matrix-inequality problem file. It contains
	source-local numerical citations but does not include the corresponding
	bibliography; no bibliographic entries are reconstructed here.
\end{enumerate}

\chapter{Algebra}
\clearpage
\section{Quantum Algebra}

% Problem statements extracted from the three PDFs in 04_Quantum_Groups.
% Source-local problem numbers, page locators, and citation callouts are omitted.

\plainproblemsection

\begin{openproblem}
	If the antipode of \(H\) is bijective, is \(H\) faithfully flat over any Hopf subalgebra with bijective antipode?
\end{openproblem}

\plainproblemsection

\begin{openproblem}
	If \(H\) is Noetherian, is it residually finite-dimensional?
\end{openproblem}

\plainproblemsection

\begin{openproblem}
	Does \(\operatorname{GKdim} H\) belong to \(\mathbb{Z}\cup\{\infty\}\)?
\end{openproblem}

\plainproblemsection

\begin{openproblem}
	If \(U(\mathfrak{g})\) is Noetherian, is the Lie algebra \(\mathfrak{g}\) finite-dimensional?
\end{openproblem}

\plainproblemsection

\begin{openproblem}
	If \(\Bbbk\Gamma\) is Noetherian, is \(\Gamma\) polycyclic-by-finite?
\end{openproblem}

\plainproblemsection

\begin{openproblem}
	If \(H\) is Noetherian, is it necessarily affine?
\end{openproblem}

\plainproblemsection

\begin{openproblem}
	If \(H\) is affine and \(\operatorname{GKdim}H<\infty\), is \(H\) necessarily Noetherian?
\end{openproblem}

\plainproblemsection

\begin{openproblem}
	If \(H\) is Noetherian, is any Hopf subalgebra also Noetherian?
\end{openproblem}

\plainproblemsection

\begin{openproblem}
	If \(H\) is affine and PI, is it necessarily Noetherian?
\end{openproblem}

\plainproblemsection

\begin{openproblem}
	Classify the affine co-Frobenius, or cosemisimple, Hopf algebras that either have finite Gelfand--Kirillov dimension or are Noetherian.
\end{openproblem}

\plainproblemsection

\begin{openproblem}
	The antipode of a Noetherian Hopf algebra is conjectured to be bijective.
\end{openproblem}

\plainproblemsection

\begin{openproblem}
	If \(H\) is affine, Noetherian, and PI, does \(S\) have finite order? Is some even power of \(S\) an inner automorphism of \(H\)?
\end{openproblem}

\plainproblemsection

\begin{openproblem}
	\textbf{Kaplansky's conjecture.}
	If \(\Gamma\) is torsion-free, then \(\Bbbk\Gamma\) is a domain.
	
	This question appeared earlier in the literature, where it was proved for locally indicable groups.
\end{openproblem}

\plainproblemsection

\begin{openproblem}
	Classify affine Hopf algebras with finite Gelfand--Kirillov dimension that are domains, prime, or semiprime.
\end{openproblem}

\plainproblemsection

\begin{openproblem}
	If \(H\) is a Noetherian domain, is it necessarily semiprimitive, that is, does \(\operatorname{Jac}H=0\)?
\end{openproblem}

\plainproblemsection

\begin{openproblem}
	If \(H\) is a Noetherian PI domain, is it necessarily module-finite over its centre? What if it is also affine?
	
	The original version with the hypothesis ``prime'' instead of ``domain'' was disproved.
\end{openproblem}

\plainproblemsection

\begin{openproblem}
	If \(H\) is affine, Noetherian, and a domain, does \(\operatorname{GKdim}H\leq 2\) imply that it is generated by grouplikes and skew primitives?
\end{openproblem}

\plainproblemsection

\begin{openproblem}
	Let \((V,c)\) be a braided vector space with \(\operatorname{GKdim}\mathcal{B}(V)<\infty\). Is \((V,c)\) a locally finite braided vector space?
\end{openproblem}

\plainproblemsection

\begin{openproblem}
	If \(H\) is affine, \(H_0\) is a Hopf subalgebra, and \(\operatorname{GKdim}H<\infty\), is \(\operatorname{gr}H\) affine?
\end{openproblem}

\plainproblemsection

\begin{openproblem}
	Determine \(\operatorname{Post}_{\mathrm{fGK},H}(V)\) for \(V\in{}^{H}_{H}\mathcal{YD}\) with \(\operatorname{GKdim}\mathcal{B}(V)<\infty\).
\end{openproblem}

\plainproblemsection

\begin{openproblem}
	Determine \(\operatorname{Pre}_{\mathrm{fGK}}(V)\) when \(V\) is of diagonal type with braiding matrix of Cartan type \(A_\theta\) or \(D_\theta\), with \(q=-1\).
\end{openproblem}

\plainproblemsection

\begin{openproblem}
	Let \(\Gamma\) be a nilpotent-by-finite group. Determine all post-Nichols algebras \(\mathcal{R}\) of \(V\in{}^{\Bbbk\Gamma}_{\Bbbk\Gamma}\mathcal{YD}\) such that \(\operatorname{GKdim}\mathcal{R}<\infty\).
\end{openproblem}

\plainproblemsection

\begin{openproblem}
	Determine the pairs \((X,\mathbf{q})\), where \(X\) is a finite rack and \(\mathbf{q}:X\times X\to\operatorname{GL}(n,\Bbbk)\) is a \(2\)-cocycle, such that \(\operatorname{GKdim}\mathcal{B}(X,\mathbf{q})<\infty\).
\end{openproblem}

\plainproblemsection

\begin{openproblem}
	Let \(\mathcal{O}\) be a finite rack of type \(C\). It is conjectured that
	\[
	\operatorname{GKdim}\mathcal{B}(\mathcal{O},\mathbf{q})=\infty
	\]
	for every faithful \(2\)-cocycle \(\mathbf{q}\).
\end{openproblem}

\plainproblemsection

\begin{openproblem}
	Let \(X\) be an indecomposable rack with \(1<|X|<\infty\), and let
	\[
	\mathbf{1}:X\times X\longrightarrow\Bbbk^\times
	\]
	be the constant cocycle \(1\). Is \(\operatorname{GKdim}\mathcal{B}(X,\mathbf{1})<\infty\)?
\end{openproblem}

\plainproblemsection

\begin{openproblem}
	Let \(\mathcal{B}\) be a pre-Nichols algebra of a finite-dimensional braided vector space \((V,c)\). Assume that \(\operatorname{GKdim}\mathcal{B}<\infty\). Is
	\[
	\operatorname{GKdim}\mathcal{B}
	=
	\lim_{n\to\infty}\log_n\dim\mathcal{B}^{n}?
	\]
	Is the Hilbert series of \(\mathcal{B}\) a rational function?
\end{openproblem}

\plainproblemsection

\begin{openproblem}
	Is
	\[
	\operatorname{GKdim}\mathcal{B}(V\otimes W)
	=
	\operatorname{GKdim}\mathcal{B}(V)
	+
	\operatorname{GKdim}\mathcal{B}(W)
	\]
	provided that \(V,W\in{}^{H}_{H}\mathcal{YD}\) satisfy
	\[
	c_{V,W}c_{W,V}=\operatorname{id}_{W\otimes V}?
	\]
\end{openproblem}

\plainproblemsection

\begin{openproblem}
	Classify the affine compact Hopf algebras that are either Noetherian or of finite Gelfand--Kirillov dimension.
\end{openproblem}

\plainproblemsection

\begin{openproblem}
	Classify the CQG algebras, or more generally the cosemisimple Hopf algebras, of Lie type.
\end{openproblem}

\plainproblemsection

\begin{openproblem}
	Is a simple cosemisimple Hopf algebra \(H\), where ``simple'' means that \(H\) has no normal Hopf subalgebras, with
	\[
	0<\operatorname{GKdim}H<\infty
	\]
	necessarily of Lie type?
\end{openproblem}

\plainproblemsection

\begin{openproblem}
	If \(\operatorname{Jac}H\) is a Hopf ideal, when does it satisfy
	\[
	\bigcap_{n\in\mathbb{N}}(\operatorname{Jac}H)^n=0?
	\]
\end{openproblem}

\plainproblemsection

\begin{openproblem}
	For \(H\) semiprimitive with \(\operatorname{GKdim}H<\infty\), classify all pre-Nichols algebras \(\mathcal{B}\) in \({}^{H}_{H}\mathcal{YD}\) with \(\operatorname{GKdim}\mathcal{B}<\infty\), and all their liftings, that is, all Hopf algebras \(L\) such that
	\[
	\operatorname{Gr}L\simeq\mathcal{B}\mathbin{\#}H.
	\]
\end{openproblem}

\plainproblemsection

\begin{openproblem}
	When does the ideal
	\[
	\operatorname{res}H
	=
	\bigcap_{V\in\operatorname{rep}H}\ker\rho_V
	\]
	satisfy
	\[
	\bigcap_{n\in\mathbb{N}}(\operatorname{res}H)^n=0?
	\]
\end{openproblem}

\plainproblemsection

\begin{openproblem}
	Assume that \(V\) is a multiplicative unitary with manageable antipodes. The \(\sigma\)-weak closures of its legs are von Neumann algebras \(M\) and \(\widehat{M}\), and they carry coproducts. On \(M\),
	\[
	\Delta(x)=V(x\otimes 1)V^*.
	\]
	Formulate conditions under which it is possible to prove the existence of the Haar weights. It is sufficient to have either the left or the right Haar weight, since the unitary antipode \(R\) is available from the polar decomposition of \(S\) on \(M\).
\end{openproblem}

\plainproblemsection

\begin{openproblem}
	Find examples of locally compact quantum groups with an involutive antipode that are not Kac algebras.
\end{openproblem}

\plainproblemsection

\begin{openproblem}
	Given the right invariant Haar weight \(\psi\) on the pair \((M,\Delta)\), standard techniques define the canonical map \(V\) on the Hilbert space \(H_\psi\otimes H_\psi\) and show that it is an isometry. However, it is not trivial to show that it is actually a unitary. This is easy in the algebraic setting, using the surjectivity of the canonical map
	\[
	a\otimes a'\longmapsto\Delta(a)(1\otimes a'),
	\]
	but not in the topological setting.
	
	In the original work of Kustermans and Vaes, ``Kustermans's trick'' does the job, and the same method is used elsewhere. Find a more direct and natural proof of this property. Such a proof is also important for the passage from \(C^*\)-algebras to von Neumann algebras. An elegant solution was presented by P.\ Kasprzak at a workshop in Marseille, but no record of the proof is available.
\end{openproblem}

\plainproblemsection

\begin{openproblem}
	\textbf{Finite Generation of Cohomology for Bosonizations.}
	
	\textit{Participants.}
	Nicol\'as Andruskiewitsch, David Jaklitsch, Van Nguyen, Amrei Oswald, Julia Plavnik, Anne Shepler, Harshit Yadav, Xingting Wang.
	
	\textit{Project description.}
	Nicol\'as Andruskiewitsch posed the following question concerning the finite generation of cohomology conjecture (FGC):
	
	Consider an exact sequence of finite-dimensional Hopf algebras over a field,
	\[
	\longrightarrow A\longrightarrow C\longrightarrow B\longrightarrow.
	\]
	If \(A\) and \(B\) satisfy the FGC, does \(C\) as well? In other words, which extensions of finite-dimensional Hopf algebras preserve the FGC property?
	
	The group started by considering the case of abelian extensions, that is, extensions with \(A\) commutative and \(B\) cocommutative, following ideas of Andruskiewitsch and Natale. It then considered general conditions on extensions that would allow the use of ideas of Negron on equivariant deformations and of Negron and Pevtsova. The discussion focused especially on arguments applicable when \(C\) is a smash product algebra and also considered the more general case of crossed products. Certain desired conditions on extensions were identified, and an argument that may establish the FGC for these extensions was outlined.
\end{openproblem}

\plainproblemsection

\begin{openproblem}
	\textbf{Tame Hopf Algebras and the Balmer Spectrum.}
	
	\textit{Participants.}
	Karin Erdmann, Julia Pevtsova, Vera Serganova, Sarah Witherspoon.
	
	\textit{Project description.}
	Discussions began with the concrete goal of finding the Balmer spectrum of a specific example, namely modules for the small quantum supergroup \(u_q(\mathfrak{sl}(1|1))\), where \(q\) is a primitive \(n\)th root of unity. The answer was found on the first day to be the projective line modulo the action of the cyclic group \(C_n\). The team then considered other Hopf algebras of tame representation type, except in characteristic \(2\).
	
	After considering known examples, the team conjectured that the Balmer spectrum of each tame finite-dimensional Hopf algebra is homeomorphic to a projective line. In view of known classification results for such Hopf algebras, this appears possible to prove.
\end{openproblem}

\plainproblemsection

\begin{openproblem}
	\textbf{Rank Varieties for Quantum Complete Intersections.}
	
	\textit{Participants.}
	Karin Erdmann, Cris Negron, Julia Pevtsova, Vera Serganova, Sarah Witherspoon.
	
	\textit{Project description.}
	The team from the preceding project worked on a second question, motivated by the frequent utility of rank varieties in proving the tensor product property for support varieties. Let \(\Bbbk\) be a field and let
	\[
	\Lambda_{\mathbf{q}}
	=
	\left\langle x_1,\ldots,x_m\right\rangle
	\big/
	\left(x_ix_j-q_{ij}x_jx_i,\ x_i^{n_i}\right)
	\]
	be a quantum complete intersection associated with a matrix \(\mathbf{q}=(q_{ij})\) and positive integers \(n_i\geq 2\). Assume that these parameters arise from a braiding on a vector space of dimension \(m\) via an action of a finite group \(G\), and that \(\Lambda_{\mathbf{q}}\) is a Nichols algebra. Let
	\[
	A_{\mathbf{q}}=\Lambda_{\mathbf{q}}\mathbin{\#}G
	\]
	be its bosonization. Define a rank variety that is homeomorphic to the cohomological support variety.
	
	The homogeneous case, in which all \(n_i\) are equal and all \(q_{ij}\) are equal for \(i<j\), is known. The nonhomogeneous case is unknown and not immediately obvious, while at the same time arising naturally in some quantum-group and pointed-Hopf-algebra settings. The team first discussed a rank-\(2\) example arising from \(u_q(\mathfrak{sl}(1|2))\), in which \(A_{\mathbf{q}}\) can alternatively be viewed as a semidirect product \(R\mathbin{\#}G\) for the commutative algebra
	\[
	R=\Bbbk[s,t]/(s^N,t^2).
	\]
	Several possible definitions of rank variety were explored:
	\begin{enumerate}
		\item \(\pi\)-points;
		\item noncommutative hypersurfaces;
		\item commutative hypersurfaces;
		\item matrix factorizations.
	\end{enumerate}
	The first two ideas do not seem to work. The team spent the most time discussing the last two in combination, converting the problem to a commutative one and using known commutative hypersurface support theory. The team is optimistic about both projects and plans to continue the discussion remotely.
\end{openproblem}

\plainproblemsection

\begin{openproblem}
	\textbf{Cohomology of \(\mathrm{FK}(4)\) and Smaller Nichols Algebras over Nonabelian Groups of Group-Likes.}
	
	\textit{Participants.}
	Istv\'an Heckenberger, Stanislas Herscovich, H\'ector Pe\~na Pollestri, and Milen Yakimov on the first day.
	
	\textit{Project description.}
	The group discussed techniques for computing projective resolutions of quadratic algebras, with special focus on the Fomin--Kirillov algebras \(\mathrm{FK}(3)\) and \(\mathrm{FK}(4)\). Although the Yoneda algebra of \(\mathrm{FK}(3)\) is known---it is the polynomial extension
	\[
	\mathrm{FK}(3)^{!}[\omega]
	\]
	of the quadratic dual \(\mathrm{FK}(3)^{!}\)---the origin of the element \(\omega\) is not understood.
	
	Because \(\mathrm{FK}(4)\) was too complicated, the group moved to other finite-dimensional Nichols algebras. More particularly, after considering possible geometric tools that could simplify computations, it discussed the extra algebraic structure of the quadratic dual \(A^{!}\) of quadratic algebras \(A\) obtained as the quadratic part of known examples of finite-dimensional Nichols algebras in characteristic zero. The examples are listed at
	\[
	\text{\url{http://mate.dm.uba.ar/~lvendram/zoo/}}.
	\]
	The group noted, in particular, the role in the algebraic structure of \(A^{!}\) played by the groups and racks naturally used in the definition of the Nichols algebras: \(A^{!}\) can essentially be regarded as a twisted monoid bialgebra.
	
	The group decided to perform more detailed computations in order to understand better which structures appear and what role they play in determining the corresponding Yoneda algebra. In particular, GAP computations for the cohomology of the quadratic part \(A\) of the second example in the preceding list suggest that
	\[
	\operatorname{Ext}^{\bullet}_{A}(\Bbbk,\Bbbk)
	=
	A^{!}[\omega],
	\]
	where \(\omega\) has cohomological degree \(4\) and internal degree \(-6\), exactly as for \(\mathrm{FK}(3)\).
\end{openproblem}

\plainproblemsection

\begin{openproblem}
	\textbf{Cohomology of Group Schemes in the Verlinde Category.}
	
	\textit{Participants.}
	Iv\'an Angiono, Agustina Czenky, Alexei Davydov, Eric Friedlander, Pablo Ocal, Guillermo Sanmarco, Antoine Touz\'e.
	
	\textit{Project description.}
	The group considered the following problem, suggested by Vera Serganova. Let \(p\) be a prime, let \(\operatorname{Ver}_p\) be the Verlinde category, and let \(X\in\operatorname{Ver}_p\).
	\begin{enumerate}
		\item Compute the cohomology of the group scheme \(\operatorname{GL}(X)\).
		\item Compute the cohomology of some finite group schemes in \(\operatorname{Ver}_p\).
	\end{enumerate}
	
	The group considered finite generation of the cohomology of group schemes in the Verlinde category \(\operatorname{Ver}_p\). The goal was to prove this property for \(\operatorname{GL}(X)\), where \(X\in\operatorname{Ver}_p\), and to use this as a springboard for proving the finite generation conjecture for finite group schemes in \(\operatorname{Ver}_p\).
	
	The group began by studying the cohomology of \(\operatorname{GL}(X)\). Its complexity led the group to study instead the restricted Lie algebra \(u(\mathfrak{gl}(X))\), a strategy that had succeeded in the classical case. The group observed that in some cases
	\[
	u(\mathfrak{g})\simeq U(\mathfrak{g})
	\]
	because of certain vanishing in positive characteristic. Since the latter appeared simpler, it switched to understanding \(U(\mathfrak{g})\).
	
	For this purpose, a categorical Chevalley--Eilenberg complex \(C_{\bullet}^{\mathrm{CE}}(\mathfrak{g})\) was constructed. This complex is well defined in any symmetric monoidal category, and its terms are projective \(U(\mathfrak{g})\)-modules. When specialized to
	\[
	-\operatorname{Vec}\simeq\operatorname{Ver}_2
	\qquad\text{and}\qquad
	-\operatorname{sVec}\simeq\operatorname{Ver}_3,
	\]
	the complex recovers the usual Chevalley--Eilenberg resolution and the May resolution, respectively. However, \(C_{\bullet}^{\mathrm{CE}}(\mathfrak{g})\) is not a resolution in general.
	
	The group proved that \(C_{\bullet}^{\mathrm{CE}}(\mathfrak{g})\) is a projective resolution of the monoidal identity in \(\operatorname{Ver}_p\) if and only if the underlying object of \(\mathfrak{g}\) in \(\operatorname{Ver}_p\) has only \(L_1\) and \(L_{p-1}\) factors. It also computed \(H^\bullet(C_{\bullet}^{\mathrm{CE}}(\mathfrak{g}))\) explicitly. The main tools were a spectral sequence with vanishing differentials, the existence of a PBW basis of \(\mathfrak{g}\), and the K\"unneth formula.
	
	The group plans to continue working on this project. It has already made some progress on modifying \(C_{\bullet}^{\mathrm{CE}}(\mathfrak{g})\) to obtain an actual resolution, as well as developed ideas for understanding the structure of the symmetric algebra of the simple objects \(S(L_m)\) for \(1\leq m\leq p-1\).
\end{openproblem}

\plainproblemsection

\begin{openproblem}
	\textbf{Noetherianity of the Balmer Spectrum of Stable and Derived Categories.}
	
	\textit{Participants.}
	Paul Balmer, Karthik Ganapathy, Juan Omar G\'omez Rodr\'iguez, Cris Negron, Kent Vashaw, Milen Yakimov.
	
	\textit{Project description.}
	Let \(\mathcal{C}\) be a finite symmetric tensor category. The goal was to consider the relationship among:
	\begin{enumerate}
		\item cohomological support satisfying the tensor product property;
		\item the Etingof--Ostrik finite generation conjecture holding for \(\mathcal{C}\);
		\item the Noetherianity of the Balmer spectrum of the stable category \(\underline{\mathcal{C}}\) or of the derived category \(D^b(\mathcal{C})\).
	\end{enumerate}
	A very ambitious form of this question asks whether the comparison map
	\[
	\operatorname{Spc}D^b(\mathcal{C})
	\longrightarrow
	\operatorname{Spec}^{h}H^\bullet(\mathcal{C})
	\]
	is a homeomorphism.
	
	Tensor-triangular localization techniques make it possible to reduce to the local case, where one seeks to prove that the comparison map
	\[
	\operatorname{Spc}\mathcal{K}
	\longrightarrow
	\operatorname{Spec}^{h}H^\bullet(\mathcal{K})
	\]
	for a tensor-triangulated category \(\mathcal{K}\) with a Noetherian, local, periodic cohomology ring has the property that the fiber of the irrelevant ideal of \(H^\bullet(\mathcal{K})\) is a single point.
	
	The group proved this fact for \(D^b(\mathcal{C})\) by a direct analysis of the perfect complexes. Although the comparison map cannot be a homeomorphism in the broadest generality, since the stable homotopy category gives a counterexample, the group hopes that there is a describable class of tensor-triangulated categories, closed under localization, for which it is a homeomorphism.
	
	As a special case, the group asked: if \(f:\mathbf{1}\to X\) and \(g:\mathbf{1}\to Y\) are two morphisms such that \(f\otimes g=0\), is either \(f\) or \(g\) nilpotent on some nonzero object? This is true if \(X\) or \(Y\) is invertible, since this implies that \(f\), respectively \(g\), is nilpotent on its own cone. The group attempted to generalize this proof to arbitrary morphisms in \(D^b(\mathcal{C})\).
\end{openproblem}

\subsection*{References}
\begin{enumerate}[label={[\arabic*]},leftmargin=*,itemsep=0.4em]
	\item Nicol\'{a}s Andruskiewitsch, \emph{On Infinite-Dimensional Hopf Algebras}, arXiv:2308.13120v2 [math.QA], 2023. \url{https://doi.org/10.48550/arXiv.2308.13120}
	\item Alfons Van Daele, \emph{From Hopf Algebras to Topological Quantum Groups: A Short History, Various Aspects and Some Problems}, arXiv:1901.04328 [math.QA], 2019. \url{https://doi.org/10.48550/arXiv.1901.04328}
	\item Nicol\'{a}s Andruskiewitsch, Iv\'{a}n Angiono, Julia Pevtsova, and Julia Plavnik (organizers), \emph{Finite Tensor Categories: Their Cohomology and Geometry}, Workshop Summary, American Institute of Mathematics, 2024. \url{https://aimath.org/pastworkshops/finitetensorrep.pdf}
\end{enumerate}

\clearpage

\section{Group Theory}

\plainproblemsection

\begin{openproblem}
	(Sidki) Can the odometer acting on the rooted binary tree be embedded in a free pro 2-group?
	
	It can be embedded in a free group.
\end{openproblem}

\plainproblemsection

\begin{openproblem}
	(many people) Does the Baumslag-Pride theorem hold for pro $p$ -groups? That is, is it true that if $G$ is a pro $p$ -group with $r + 2$ generators and $r$ relators, then $G$ virtually surjects onto a non-Abelian free pro p-group?
	
	It is true in the discrete case.
\end{openproblem}

\plainproblemsection

\begin{openproblem}
	Let $G$ be a closed transitive subgroup of the automorphism group of a rooted tree. Is it true that every level stabilizer of $G$ contains a fixedpoint-free element?
	
	Transitive and fixedpoint-free are understood with respect to the boundary of the tree. This is equivalent to the following old question of Jehne in field arithmetic: Do there exist fields $K \leq  L \leq  M$ with $K$ is global, $L/K$ finite separable and $M/K$ infinite separable, such that all intermediate fields $L \leq$ ${M}_{0} \leq  M$ of finite degree over $K$ are Kronecker-conjugate to $L$ ? For pro-p groups it is true, though that is not so interesting from the number theoretic point of view.
\end{openproblem}

\plainproblemsection

\begin{openproblem}
	Let $\Gamma$ be a countable subgroup of $S{L}_{2}\left( {\mathbb{Q}}_{p}\right)$ that does not contain parabolic elements and let $\gamma$ be a random element of $S{L}_{2}\left( {\mathbb{Q}}_{p}\right)$ . Show that the group generated by $\Gamma$ and $\gamma$ does not contain any parabolic elements a.s.
	
	The analogous result is true for rooted tree actions.
\end{openproblem}

\plainproblemsection

\begin{openproblem}
	(A, Virag) What is the length of the shortest law in the $n$ times iterated wreath product of ${C}_{2}$ (the automorphism group of the rooted binary tree of length $n$ )?
	
	Probably the shortest law is ${x}^{{2}^{n}}$ .
\end{openproblem}

\plainproblemsection

\begin{openproblem}
	(Benjamini) Can the set of balls in an infinite Cayley graph form an expander family?
	
	Most likely not. I do not know why, but I really like this problem.
\end{openproblem}

\plainproblemsection

\begin{openproblem}
	Suppose $G$ and $H$ are Cayley (or vertex transitive) expanders on the same number of vertices and you can almost match them in edge distance. Is it true that they are isomorphic?
	
	An analogous rigidity result holds for graphings of property (T) groups.
\end{openproblem}

\plainproblemsection

\begin{openproblem}
	Suppose $G$ is a finite vertex transitive expander. Is it true that every almost automorphism of $G$ is almost an automorphism?
	
	An almost automorphism is a bijection of the vertex set that almost sends edges to edges. The statement is true for Cayley diagrams. Expansion is necessary.
\end{openproblem}

\plainproblemsection

\begin{openproblem}
	Let $\Gamma$ be a finitely generated group and let $\left\{  {{H}_{n} \mid  n \geq  1}\right\}$ be a property (τ) family of normal subgroups of finite index in $\Gamma$ . Does the chain ${\Gamma }_{n} = { \cap  }_{k = 1}^{n}{H}_{k}$ have property $\left( \tau \right)$ ?
	
	On graph theory language, this asks if the diagonal product of expander Cayley graphs keeps being expander on every connected component. It is not true if one omits normality.
\end{openproblem}

\plainproblemsection

\begin{openproblem}
	Let $\Gamma$ be a finitely presented group with a chain of normal subgroups of finite index and with trivial intersection. Assume that the chain has property (τ). Does $\Gamma$ have a free subgroup?
	
	I hope to have a positive answer for this one - if that helps, assume that every index is a prime power. One would need an 'asymptotic ping-pong lemma' to start with.
\end{openproblem}

\plainproblemsection

\begin{openproblem}
	Let $G$ be an infinite d-regular Ramanujan graph (meaning, the Markov operator has the same spectral radius as of the d-regular tree). Is it true that for every $C$ , the probability that a random walk on $G$ ends touching a C-cycle tends to zero? That is, is it true that the random walk neighbourhood sampling on $G$ converges to the $d$ -regular tree?
	
	It is known that a unimodular random network that is $d$ -regular, infinite and Ramanujan is the $d$ -regular tree a.s.
\end{openproblem}

\plainproblemsection

\begin{openproblem}
	Let $\left( {G}_{n}\right)$ be a locally convergent sequence of bounded degree, integer labeled graphs. Is the normalized mod $p$ rank of the adjacency matrix convergent?
	
	The same over $\mathbb{Q}$ is true and is equivalent to the Lück Approximation Theorem.
\end{openproblem}

\plainproblemsection

\begin{openproblem}
	(Elek) For a graph sequence $\left( {G}_{n}\right)$ let the edge measure be
	
	$$
	e\left( \left( {G}_{n}\right) \right)  = \liminf \frac{\left| E\left( {G}_{n}\right) \right| }{\left| V\left( {G}_{n}\right) \right| }
	$$
	
	A graph sequence $\left( {H}_{n}\right)$ defined on the same vertex set is equivalent to $\left( {G}_{n}\right)$ , if the bi-Lipschitz constants of the identity maps are bounded. Let the combinatorial cost of $\left( {G}_{n}\right)$ be the infimum of edge measures of graph sequences equivalent to $\left( {G}_{n}\right)$ . Let $\left( {G}_{n}\right)$ and $\left( {H}_{n}\right)$ be graph sequences that locally converge to the same limit. Do they have the same combinatorial cost?
	
	This is 'morally' equivalent to the Fixed Price problem, in the sense that once one is solved, the other is expected to be solved soon as well.
\end{openproblem}

\plainproblemsection

\begin{openproblem}
	(Lovasz) By compactness, for any $\varepsilon  > 0$ there exists $K > 0$ such that any finite graph can be approximated within the error $\varepsilon$ by a finite graph of size $K$ . Can we give an estimate for $K$ in terms of $\varepsilon$ ?
\end{openproblem}

\plainproblemsection

\begin{openproblem}
	Let $\left( {G}_{n}\right)$ be a locally convergent graph sequence and let ${\mu }_{n}$ be the probability distribution of the roots of the chromatic polynomial of ${G}_{n}$ . Is it true that for all $k$ , the limit
	
	$$
	\mathop{\lim }\limits_{{n \rightarrow  \infty }}{\int }_{\mathbb{C}}{z}^{k}d{\mu }_{n}
	$$
	
	exists?
\end{openproblem}

\plainproblemsection

\begin{openproblem}
	What can be the shape of the eigenvalue distribution of a finite d-regular graph? Find natural restrictions.
	
	One such non-trivial restriction is: if most of the measure is inside the Alon-Boppana bound, then we roughly know its shape.
\end{openproblem}

\plainproblemsection

\begin{openproblem}
	Let $G$ be a d-regular Cayley graph with spectral measure $\mu$ . Is $\mu$ a weak limit of spectral measures of finite d-regular graphs?
	
	A negative answer would solve the sofic problem negatively.
\end{openproblem}

\plainproblemsection

\begin{openproblem}
	(Bollobas) Let ${G}_{n}$ be a random d-regular graph on $n$ points. Is there a limit of the independence ratio of ${G}_{n}$ ?
	
	It is known that for every $n$ it is very concentrated. So, maybe the value its concentrated around changes with $n$ ?
\end{openproblem}

\plainproblemsection

\begin{openproblem}
	(Szegedy) Let ${G}_{n}$ be a random d-regular graph on $n$ points. Does $\left( {G}_{n}\right)$ converges to the (weak closure of) i.i.d.-s in local-global convergence?
\end{openproblem}

\plainproblemsection

\begin{openproblem}
	Is the i.i.d. action of the free group ${F}_{2}$ a local-global limit of finite actions of ${F}_{2}$ ?
\end{openproblem}

\plainproblemsection

\begin{openproblem}
	(Bowen) Let $\Gamma$ be a Property (T) group and let ${G}_{n}$ be a sofic approximation to $G$ , that is, a sequence of finite graphs that converges to a Cayley graph of $G$ . Surely, ${G}_{n}$ does not have to be an expander family. But can we modify the sequence by an asymptotically vanishing amount (in the edit distance) to a sequence ${G}_{n}^{\prime }$ such that any subsequence of connected components of ${G}_{n}^{\prime }$ is an expander family? To put it another way, can one ’pullback’ the ergodic decomposition of the invariant measure on the ultraproduct space?
\end{openproblem}

\plainproblemsection

\begin{openproblem}
	Is it true that any ergodic random unimodular network that is an infinite tree a.s. can be obtained as the limit of an expander family?
\end{openproblem}

\plainproblemsection

\begin{openproblem}
	Let $G$ be an infinite vertex transitive graph, let $A$ be a finite set of vertices of $G$ and $b$ a vertex. Show that
	
	$$
	\mathop{\sum }\limits_{{x \in  \partial A}}d\left( {b, x}\right)  \geq  \left| A\right|
	$$
	
	where $\partial A$ is the set of points of distance 1 from $A$ .
	
	Known for unimodular vertex transitive graphs.
\end{openproblem}

\plainproblemsection

\begin{openproblem}
	One can define the first ${L}^{2}$ Betti number of an arbitrary vertex transitive graph $G$ , from the expected degree of a free spanning forest of $G$ . Does $G$ and ${G}^{2}$ have the same first ${L}^{2}$ Betti number?
	
	When $G$ is a Cayley graph for a group $\Gamma$ , this number only depends on $\Gamma$ , and not on the generating set, hence it is true there.
\end{openproblem}

\plainproblemsection

\begin{openproblem}
	Can one show using free spanning forests any basic properties of the first ${L}^{2}$ Betti number? For instance, that it is multiplicative for taking a finite index subgroup.
\end{openproblem}

\plainproblemsection

\begin{openproblem}
	(Lyons) Let $G$ be an infinite Cayley graph of a group that is not virtually cyclic. Show that there exists $p < 1$ such that the $p$ -edge percolation of $G$ has some infinite clusters.
	
	Not virtually cyclic is the same as not having linear word growth.
\end{openproblem}

\plainproblemsection

\begin{openproblem}
	Let $G$ be an infinite, connected Cayley graph. Then $G$ has a perfect matching. Does $G$ have an invariant random subgraph that is a perfect matching a.s.?
	
	Known if $G$ is bipartite.
\end{openproblem}

\plainproblemsection

\begin{openproblem}
	(Babai) Let $G$ be an infinite Cayley (vertex transitive) graph. Does $G$ have a spanning tree without leaves?
\end{openproblem}

\plainproblemsection

\begin{openproblem}
	Let $G$ be an infinite Cayley (vertex transitive) graph. Does the set of dead ends in $G$ have density 0 ?
\end{openproblem}

\plainproblemsection

\begin{openproblem}
	Are factors of i.i.d. on the 3-regular tree closed in the weak topology?
	
	Most likely not. In particular, look at the weak limit of majority functions on $n$ -balls. Is that a factor of i.i.d.?
\end{openproblem}

\plainproblemsection

\begin{openproblem}
	Let $G$ be an infinite Cayley graph. Does there exist a $G$ -invariant proper coloring with $\chi \left( G\right)$ colors?
\end{openproblem}

\plainproblemsection

\begin{openproblem}
	Let $X$ be the space of $k$ -regular Cayley graphs with the local convergence topology. It is easy to see that the set $T$ of transient graphs is closed. Is the Green function (the expected number of returns to the identity) continuous on $T$ ?
\end{openproblem}

\plainproblemsection

\begin{openproblem}
	Let $k > 1$ . Does there exist $C\left( k\right)  < 1$ such that if $G$ is a $k$ -regular transient Cayley graph, then the probability of return is at most $C\left( k\right)$ ?
\end{openproblem}

\plainproblemsection

\begin{openproblem}
	Let $\Gamma$ be a nonamenable group. Does $\{ 0,1{\} }^{\Gamma }$ factor onto $\{ 0,1,2{\} }^{\Gamma }$ ?
\end{openproblem}

\plainproblemsection

\begin{openproblem}
	Can every d-regular graphing without multiple edges be properly edge colored by $d + 1$ colors (measurable Vizing theorem)?
\end{openproblem}

\plainproblemsection

\begin{openproblem}
	Let ${G}_{n}$ and ${H}_{n}$ be graph sequences that converge to the same limit. Let ${T}_{n}$ be a spanning tree in ${G}_{n}$ that is convergent. Does there exist a spanning tree ${P}_{n}$ in ${H}_{n}$ that converges to the same limit as ${T}_{n}$ ?
\end{openproblem}

\plainproblemsection

\begin{openproblem}
	Let $G$ be an expander (strongly ergodic) graphing of bounded degree that can be properly colored by $C$ colors with arbitrarily small error. Can it be properly $C$ -colored?
\end{openproblem}

\plainproblemsection

\begin{openproblem}
	Let $G$ be an expander (strongly ergodic) graphing of bounded degree that weakly contains a finite graph $H$ . Does $G$ factor on $H$ ?
	
	True when $H$ has two vertices.
\end{openproblem}

\plainproblemsection

\begin{openproblem}
	Let $\Gamma$ be a higher rank semisimple real lattice. Is the rank gradient of $\Gamma$ zero?
	
	Known for non-uniform lattices.
\end{openproblem}

\plainproblemsection

\begin{openproblem}
	(Gaboriau) Let $A$ and $B$ be countably infinite groups. Does $A \times  B$ have fixed price 1 ?
	
	Known if $A$ or $B$ contains an infinite amenable (or arbitrarily large finite) subgroup.
\end{openproblem}

\plainproblemsection

\begin{openproblem}
	(Gaboriau) Does every countable group have fixed price?
\end{openproblem}

\plainproblemsection

\begin{openproblem}
	Let $\Gamma$ be a residually finite group with property (T). Does $\Gamma$ have zero rank gradient?
	
	This is a baby version of the question, whether these groups have fixed price 1.
\end{openproblem}

\plainproblemsection

\begin{openproblem}
	Let $\Gamma$ act on $X$ by p.m.p. maps and let $H$ be subgroup of finite index in $\Gamma$ . Is it true that
	
	$$
	\operatorname{cost}\left( {H, X}\right)  - 1 = \left( {\operatorname{cost}\left( {\Gamma , X}\right)  - 1}\right) \left| {\Gamma  : H}\right| ?
	$$
	
	Already for i.i.d., this would show that every profinite free action of $\Gamma$ has the same cost, so the rank gradient is independent of the chain.
\end{openproblem}

\plainproblemsection

\begin{openproblem}
	Is it true that $\operatorname{cost}\left( {\Gamma , X}\right)  = \operatorname{cost}\left( {\Gamma ,{X}^{2}}\right)$ for free actions?
	
	By ${X}^{2}$ we mean the diagonal action of $\Gamma$ on $X \times  X$ .
\end{openproblem}

\plainproblemsection

\begin{openproblem}
	Let $\Gamma$ be amenable, acting ergodically and essentially faithfully on $X$ (every nontrivial element moves a set of positive measure). Is the groupoid cost of the action 1 ? If $\Gamma$ is finitely presented, is it true for any infinite ergodic action?
	
	The second theorem is true for profinite actions.
\end{openproblem}

\plainproblemsection

\begin{openproblem}
	(Gaboriau) Given a graphing of an equivalence relation, is there a subgraphing whose cost is almost the cost of the relation?
	
	What does this suggest for profinite actions?
\end{openproblem}

\plainproblemsection

\begin{openproblem}
	Can a non-Abelian free group $F$ have a nontrivial pseudocharac-ter that is invariant under $\operatorname{Aut}\left( F\right)$ ?
	
	Most likely not.
\end{openproblem}

\plainproblemsection

\begin{openproblem}
	Are two independent random subsets of $\mathbb{Z}$ as metric spaces quasi-isometric a.s.? How about other Cayley graphs, like for $S{L}_{3}\left( \mathbb{Z}\right)$ ?
	
	The answer is expected to be positive for $\mathbb{Z}$ and negative for $S{L}_{3}\left( \mathbb{Z}\right)$ .
\end{openproblem}

\plainproblemsection

\begin{openproblem}
	(Amit) Is it true that if $P$ is a finite $p$ -group of order $n$ and $w$ is any word then the probability that $w$ is satisfied is at least $1/n$ ?
	
	I tend to believe that its true. There are partial results, but nothing deep yet.
\end{openproblem}

\plainproblemsection

\begin{openproblem}
	Let $M$ be the set of measurable real functions. For $f \in  M$ let ${H}_{f},{V}_{f} : {\mathbb{R}}^{2} \rightarrow  {\mathbb{R}}^{2}$ be defined by
	
	$$
	{H}_{f}\left( {x, y}\right)  = \left( {x, y + f\left( x\right) }\right) \text{ and }{V}_{f}\left( {x, y}\right)  = \left( {x + f\left( y\right) , y}\right) .
	$$
	
	Let $H = \left\{  {{H}_{f} \mid  f \in  M}\right\}$ and $V = \left\{  {{V}_{f} \mid  f \in  M}\right\}$ . What is the group generated by $H$ and $V$ ? Is there $N$ such that every element of this group can be obtained as the product of $N$ elements of $H$ and $V$ ?
	
	For the class of all functions over any Abelian group, it generates the full symmetric group in boundedly many steps. For continuous functions over $\mathbb{R}$ , it is not bounded. For polynomials it is bounded.
\end{openproblem}

\plainproblemsection

\begin{openproblem}
	(A, Virag) What is the length of the shortest nontrivial word which is satisfied in all groups of size ${2}^{n}$ ?
	
	Probably it is ${2}^{n}$ . However, opposed to the similar question above, there are many candidates for such a law, using commutators.
\end{openproblem}

\plainproblemsection

\begin{openproblem}
	1.1. Do there exist non-trivial finitely-generated divisible groups? Equivalently, do there exist non-trivial finitely-generated divisible simple groups? Yu. A. Bogan Yes, there exist (V. S. Guba, Math. USSR-Izv., 29 (1987), 233-277).
\end{openproblem}

\plainproblemsection

\begin{openproblem}
	1.2. Let $G$ be a group, $F$ a free group with free generators ${x}_{1},\ldots ,{x}_{n}$ , and $R$ the free product of $G$ and $F$ . An equation (in the unknowns ${x}_{1},\ldots ,{x}_{n}$ ) over $G$ is an expression of the form $v\left( {{x}_{1},\ldots {x}_{n}}\right)  = 1$ , where on the left is an element of $R$ not conjugate in $R$ to any element of $G$ . We call $G$ algebraically closed if every equation over $G$ has a solution in $G$ . Do there exist algebraically closed groups? L. A. Bokut’ Yes, there exist (S. D. Brodskii, Dep. no. 2214-80, VINITI, Moscow, 1980 (Russian)).
\end{openproblem}

\plainproblemsection

\begin{openproblem}
	1.3. (Well-known problem). Can the group ring $\mathbb{Z}\left\lbrack  G\right\rbrack$ of a torsion-free group $G$ contain zero divisors? L. A. Bokut’
\end{openproblem}

\plainproblemsection

\begin{openproblem}
	1.4. (A.I. Mal'cev). Does there exist a ring without zero divisors which is not embeddable in a skew field, while the multiplicative semigroup of its non-zero elements is embeddable in a group? L. A. Bokut’
	
	Yes, there exists (L. A. Bokut', Soviet Math. Dokl., 8 (1967), 901-904; A. Bowtell, J. Algebra, 7 (1967), 126-139; A. Klein, J. Algebra, 7 (1967), 100-125).
\end{openproblem}

\plainproblemsection

\begin{openproblem}
	1.5. (Well-known problem). Does there exist a group whose group ring does not contain zero divisors and is not embeddable into a skew field? L. A. Bokut’
\end{openproblem}

\plainproblemsection

\begin{openproblem}
	1.6. (A. I. Mal'cev). Is the group ring of a right-ordered group embeddable into a skew field? L. A. Bokut’
\end{openproblem}

\plainproblemsection

\begin{openproblem}
	1.9. Can the factor-group of a locally normal group by the second term of its upper central series be embedded (isomorphically) in a direct product of finite groups? Yu. M. Gorchakov
	
	Yes, it can (Yu. M. Gorchakov, Algebra and Logic, 15 (1976), 386-390).
\end{openproblem}

\plainproblemsection

\begin{openproblem}
	1.10. An automorphism $\varphi$ of a group $G$ is called splitting if $g{g}^{\varphi }\cdots {g}^{{\varphi }^{n - 1}} = 1$ for any element $g \in  G$ , where $n$ is the order of $\varphi$ . Is a soluble group admitting a regular splitting automorphism of prime order necessarily nilpotent? Yu. M. Gorchakov Yes, it is (E.I. Khukhro, Algebra and Logic, 17 (1978), 402-406) and even without the regularity condition on the automorphism (E. I. Khukhro, Algebra and Logic, 19 (1980), 77-84).
\end{openproblem}

\plainproblemsection

\begin{openproblem}
	1.11. (E. Artin). The conjugacy problem for the braid groups ${\mathcal{B}}_{n}, n > 4$ . M. D. Greendlinger
	
	The solution is affirmative (G. S. Makanin, Soviet Math. Dokl., 9 (1968), 1156-1157; F. A. Garside, Quart. J. Math., 20 (1969), 235-254).
\end{openproblem}

\plainproblemsection

\begin{openproblem}
	1.12. (W. Magnus). The problem of the isomorphism to the trivial group for all groups with $n$ generators and $n$ defining relations, where $n > 2.\;M.D.$ Greendlinger
\end{openproblem}

\plainproblemsection

\begin{openproblem}
	1.13. (J. Stallings). If a finitely presented group is trivial, is it always possible to replace one of the defining words by a primitive element without changing the group? M. D. Greendlinger
	
	No, not always (S. V. Ivanov, Invent. Math., 165, no. 3 (2006), 525-549).
\end{openproblem}

\plainproblemsection

\begin{openproblem}
	1.14. (B. H. Neumann). Does there exist an infinite simple finitely presented group? M. D. Greendlinger
	
	Yes, there exists (G. Higman, Finitely presented infinite simple groups (Notes on Pure Mathematics, no. 8) Dept. of Math., Inst. Adv. Studies Austral. National Univ., Canberra, 1974), where there is also a reference to R. Thompson.
\end{openproblem}

\plainproblemsection

\begin{openproblem}
	1.17. Write an explicit set of generators and defining relations for a universal finitely presented group. M. D. Greendlinger
	
	This was done (M. K. Valiev, Soviet Math. Dokl., 14 (1973), 987-991).
\end{openproblem}

\plainproblemsection

\begin{openproblem}
	1.18. (A. Tarski). a) Does there exist an algorithm for determining the solubility of equations in a free group?
	
	b) Describe the structure of all solutions of an equation when it has at least one solution. Yu. L. Ershov
	
	a) Yes, there exists (G. S. Makanin, Math. USSR-Izv., 21 (1982), 546-582; 25 (1985), 75-88).
	
	b) This was described (A. A. Razborov, Math. USSR-Izv., 25 (1984), 115-162).
\end{openproblem}

\plainproblemsection

\begin{openproblem}
	1.19. (A.I. Mal'cev). Which subgroups (subsets) are first order definable in a free group? Which subgroups are relatively elementarily definable in a free group? In particular, is the derived subgroup first order definable (relatively elementarily definable) in a free group? Yu. L. Ershov
	
	All definable (and relatively elementarily definable) sets are described; in particular, a subgroup is relatively elementarily definable if and only if it is cyclic (O. Kharlampovich, A. Myasnikov, Int. J. Algebra Comput., 23, no. 1 (2013), 91- 110).
\end{openproblem}

\plainproblemsection

\begin{openproblem}
	1.20. For which groups (classes of groups) is the lattice of normal subgroups first order definable in the lattice of all subgroups? Yu. L. Ershov
\end{openproblem}

\plainproblemsection

\begin{openproblem}
	1.21. Are there only finitely many finite simple groups of a given exponent $n$ ? M. I. Kargapolov
	
	Yes, modulo CFSG, see, for example, (G. A. Jones, J. Austral. Math. Soc., 17 (1974), 162-173).
\end{openproblem}

\plainproblemsection

\begin{openproblem}
	1.23. Does there exist an infinite simple locally finite group of finite rank? No (V. P. Shunkov, Algebra and Logic, 10 (1971), 127-142). M. I. Kargapolov
\end{openproblem}

\plainproblemsection

\begin{openproblem}
	1.24. Does every infinite group possess an infinite abelian subgroup? M. I. Kargapolov
	
	No (P. S. Novikov, S. I. Adian, Math. USSR-Izv., 2 (1968), 1131-1144).
\end{openproblem}

\plainproblemsection

\begin{openproblem}
	1.25. a) Is the universal theory of the class of finite groups decidable? b) Is the universal theory of the class of finite nilpotent groups decidable?
	
	a) No (A. M. Slobodskoi, Algebra and Logic, 20 (1981), 139-156). M. I. K. M. I. Kargapolov
	
	b) No (O. G. Kharlampovich, Math. Notes, 33 (1983), 254-263).
\end{openproblem}

\plainproblemsection

\begin{openproblem}
	1.26. Does elementary equivalence of two finitely generated nilpotent groups imply that they are isomorphic? M. I. Kargapolov
	
	No (B. I. Zil’ber, Algebra and Logic, 10 (1971), 192-197).
\end{openproblem}

\plainproblemsection

\begin{openproblem}
	1.27. Describe the universal theory of free groups. M. I. Kargapolov
\end{openproblem}

\plainproblemsection

\begin{openproblem}
	1.28. Describe the universal theory of a free nilpotent group. M. I. Kargapolov
\end{openproblem}

\plainproblemsection

\begin{openproblem}
	1.29. (A. Tarski). Is the elementary theory of a free group decidable? M. I. Kargapolov
	
	Yes, it is (O. Kharlampovich, A. Myasnikov, J. Algebra, 302, no. 2 (2006), 451-552).
\end{openproblem}

\plainproblemsection

\begin{openproblem}
	1.30. Is the universal theory of the class of soluble groups decidable?
	
	No (O. G. Kharlampovich, Math. USSR-Izv., 19 (1982), 151-169). M. I. Kargapolov
\end{openproblem}

\plainproblemsection

\begin{openproblem}
	1.31. Is a residually finite group with the maximum condition for subgroups almost polycyclic? M. I. Kargapolov
\end{openproblem}

\plainproblemsection

\begin{openproblem}
	1.32. Is the Frattini subgroup of a finitely generated matrix group over a field nilpotent? M. I. Kargapolov
	
	Yes, it is (V. P. Platonov, Soviet Math. Dokl., 7 (1966), 1557-1560).
\end{openproblem}

\plainproblemsection

\begin{openproblem}
	1.33. (A. I. Mal'cev). Describe the automorphism group of a free solvable group. M. I. Kargapolov
\end{openproblem}

\plainproblemsection

\begin{openproblem}
	1.34. Has every orderable polycyclic group a faithful representation by matrices over the integers? M. I. Kargapolov
	
	Yes, it has (L. Auslander, Ann. Math. (2), 86 (1967), 112-117; R. G. Swan, Proc. Amer. Math. Soc., 18 (1967), 573-574).
\end{openproblem}

\plainproblemsection

\begin{openproblem}
	1.35. A group is called pro-orderable if every partial ordering of the group extends to a linear ordering.
	
	a) Is the wreath product of two arbitrary pro-orderable groups again pro-orderable?
	
	b) (A. I. Mal'cev). Is every subgroup of a pro-orderable group again pro-orderable? M. I. Kargapolov
	
	Not always, in both cases (V. M. Kopytov, Algebra i Logika, 5, no. 6 (1966), 27-31 (Russian)).
\end{openproblem}

\plainproblemsection

\begin{openproblem}
	1.35. c) (A.I. Mal'cev, L. Fuchs). Do there exist simple pro-orderable groups? A group is said to be pro-orderable if each of its partial orderings can be extended to a linear ordering. M. I. Kargapolov
\end{openproblem}

\plainproblemsection

\begin{openproblem}
	1.36. If a group $G$ is factorizable by $p$ -subgroups, that is, $G = {AB}$ , where $A$ and $B$ are $p$ -subgroups, does it follows that $G$ is itself a $p$ -group? Sh. S. Kemkhadze Not always (Ya. P. Sysak, Products of infinite groups, Math. Inst. Akad. Nauk Ukrain. SSR, Kiev, 1982 (Russian)).
\end{openproblem}

\plainproblemsection

\begin{openproblem}
	1.37. Is it true that every subgroup of a locally nilpotent group is quasi-invariant? Sh. S. Kemkhadze
	
	No, not always; for example, in the inductive limit of Sylow $p$ -subgroups of symmetric groups of degrees ${p}^{i}, i = 1,2,\ldots$ (V.I. Sushchanskii, Abstracts of 15th All-USSR Algebraic Conf., Krasnoyarsk, 1979, 154 (Russian)).
\end{openproblem}

\plainproblemsection

\begin{openproblem}
	1.38. An ${N}^{0}$ -group is a group in which every cyclic subgroup is a term of some normal system of the group. Is every ${N}^{0}$ -group an $\widetilde{N}$ -group? Sh. S. Kemkhadze No. Every group $G$ with a central system $\left\{  {H}_{i}\right\}$ is an ${N}^{0}$ -group, since for any $g \in  G$ the system $\left\{  \left\langle  {g,{H}_{i}}\right\rangle  \right\}$ refined by the trivial subgroup and the intersections of all the subsystems is a normal system of $G$ containing $\langle g\rangle$ . Free groups have central systems but are not $\widetilde{N}$ -groups. (Yu. I. Merzlyakov,1973.)
\end{openproblem}

\plainproblemsection

\begin{openproblem}
	1.39. Is a group binary nilpotent if it is the product of two normal binary nilpotent subgroups? Sh. S. Kemkhadze
	
	No, not always (A.I. Sozutov, Algebra and Logic, 30, no. 1 (1991), 70-72).
\end{openproblem}

\plainproblemsection

\begin{openproblem}
	1.40. Is a group a nilgroup if it is the product of two normal nilsubgroups? By definition, a nilgroup is a group consisting of nilelements, in other words, of (not necessarily boundedly) Engel elements. Sh. S. Kemkhadze
\end{openproblem}

\plainproblemsection

\begin{openproblem}
	1.41. A subgroup of a linearly orderable group is called relatively convex if it is convex with respect to some linear ordering of the group. Under what conditions is a subgroup of an orderable group relatively convex?
	
	Necessary and sufficient conditions are given in (A.I. Kokorin, V. M. Kopytov, Fully ordered groups, John Wiley, New York, 1974).
\end{openproblem}

\plainproblemsection

\begin{openproblem}
	1.42. Is the centre of a relatively convex subgroup relatively convex? A. I. Kokorin Not always (A. I. Kokorin, V. M. Kopytov, Fully ordered groups, John Wiley, New York, 1974).
\end{openproblem}

\plainproblemsection

\begin{openproblem}
	1.43. Is the centralizer of a relatively convex subgroup relatively convex? A. I. Kokorin
	
	Not always (A. I. Kokorin, V. M. Kopytov, Siberian Math. J., 9 (1968), 622-626).
\end{openproblem}

\plainproblemsection

\begin{openproblem}
	1.44. Is a maximal abelian normal subgroup relatively convex? A. I. Kokorin Not always (D. M. Smirnov, Algebra i Logika, 5, no. 6 (1966), 41-60 (Russian)).
\end{openproblem}

\plainproblemsection

\begin{openproblem}
	1.45. Is the largest locally nilpotent normal subgroup relatively convex? A. I. Kokorin
	
	Not always (D. M. Smirnov, Algebra i Logika, 5, no. 6 (1966), 41-60 (Russian)).
\end{openproblem}

\plainproblemsection

\begin{openproblem}
	1.46. What conditions ensure the normalizer of a relatively convex subgroup to be relatively convex? A. I. Kokorin
\end{openproblem}

\plainproblemsection

\begin{openproblem}
	1.47. A subgroup $H$ of a group $G$ is said to be strictly isolated if, whenever $x{g}_{1}^{-1}x{g}_{1}\cdots {g}_{n}^{-1}x{g}_{n}$ belongs to $H$ , so do $x$ and each ${g}_{i}^{-1}x{g}_{i}$ . A group in which the identity subgroup is strictly isolated is called an $S$ -group. Do there exist $S$ -groups that are not orderable groups? A. I. Kokorin
	
	Yes, there exist (V.V.Bludov, Algebra and Logic, 13 (1974), 343-360).
\end{openproblem}

\plainproblemsection

\begin{openproblem}
	1.48. Is a free product of two orderable groups with an amalgamated subgroup that is relatively convex in each of the factors again an orderable group? A. I. Kokorin Not always (M. I. Kargapolov, A. I. Kokorin, V. M. Kopytov, Algebra i Logika, 4, no. 6 (1965), 21-27 (Russian)).
\end{openproblem}

\plainproblemsection

\begin{openproblem}
	1.49. An automorphism $\varphi$ of a linearly ordered group is called order preserving if $x < y$ implies ${x}^{\varphi } < {y}^{\varphi }$ . Is it possible to order an abelian strictly isolated normal subgroup of an $S$ -group (see 1.47) in such a way that the order is preserved under the action of the inner automorphisms of the whole group? A. I. Kokorin
	
	Not always (V. V. Bludov, Algebra and Logic, 11 (1972), 341-349).
\end{openproblem}

\plainproblemsection

\begin{openproblem}
	1.50. Do the order-preserving automorphisms of a linearly ordered group form an orderable group? A. I. Kokorin
	
	Not always (D. M. Smirnov, Algebra i Logika, 5, no. 6 (1966), 41-60 (Russian)).
\end{openproblem}

\plainproblemsection

\begin{openproblem}
	1.51. What conditions ensure a matrix group over a field (of complex numbers) to be orderable? A. I. Kokorin
\end{openproblem}

\plainproblemsection

\begin{openproblem}
	1.52. Describe the groups that can be ordered linearly in a unique way (reversed orderings are not regarded as different). A. I. Kokorin
	
	This was done (V. V. Bludov, Algebra and Logic, 13 (1974), 343-360).
\end{openproblem}

\plainproblemsection

\begin{openproblem}
	1.53. Describe all possible linear orderings of a free nilpotent group with finitely many generators. A. I. Kokorin
	
	This was done (V. F. Kleimenov, in: Algebra, Logic and Applications, Irkutsk, 1994, 22-27 (Russian)).
\end{openproblem}

\plainproblemsection

\begin{openproblem}
	1.54. Describe all linear orderings of a free metabelian group with a finite number of generators. A. I. Kokorin
\end{openproblem}

\plainproblemsection

\begin{openproblem}
	1.55. Give an elementary classification of linearly ordered free groups with a fixed number of generators. A. I. Kokorin
\end{openproblem}

\plainproblemsection

\begin{openproblem}
	1.56. Is a torsion-free group pro-orderable (see 1.35) if the factor-group by its centre is pro-orderable? A. I. Kokorin
	
	Not always (M. I. Kargapolov, A. I. Kokorin, V. M. Kopytov, Algebra i Logika, 4, no. 6 (1965), 21-27 (Russian)).
\end{openproblem}

\plainproblemsection

\begin{openproblem}
	1.60. Can an orderable metabelian group be embedded in a radicable orderable group? A. I. Kokorin
	
	Yes, it can (V. V. Bludov, N. Ya. Medvedev, Algebra and Logic, 13 (1974), 207-209).
\end{openproblem}

\plainproblemsection

\begin{openproblem}
	1.61. Can any orderable group be embedded into an orderable group
	
	a) with a radicable maximal locally nilpotent normal subgroup?
	
	b) with a radicable maximal abelian normal subgroup? A. I. Kokorin
	
	Yes, it can, in both cases (S. A. Gurchenkov, Math. Notes, 51, no. 2 (1992), 129-132).
\end{openproblem}

\plainproblemsection

\begin{openproblem}
	1.63. A group $G$ is called dense if it has no proper isolated subgroups other than its trivial subgroup.
	
	a) Do there exist dense torsion-free groups that are not locally cyclic?
	
	b) Suppose that any two non-trivial elements $x$ and $y$ of a torsion-free group $G$ satisfy the relation ${x}^{k} = {y}^{l}$ , where $k$ and $l$ are non-zero integers depending on $x$ and $y$ . Does it follow that $G$ is abelian?
	
	a) Yes, there exist; b) Not necessarily (S. I. Adian, Math. USSR-Izv., 5 (1971), 475- 484).
\end{openproblem}

\plainproblemsection

\begin{openproblem}
	1.64. A torsion-free group is said to be separable if it can be represented as the set-theoretic union of two of its proper subsemigroups. Is every $R$ -group separable? P. G. Kontorovich
	
	No (S. J. Pride, J. Wiegold, Bull. London Math. Soc., 9 (1977), 36-37).
\end{openproblem}

\plainproblemsection

\begin{openproblem}
	1.65. Is the class of groups of abelian extensions of abelian groups closed under taking direct sums $\left( {A, B}\right)  \mapsto  A \oplus  B$ ? L. Ya. Kulikov
\end{openproblem}

\plainproblemsection

\begin{openproblem}
	1.66. Suppose that $T$ is a periodic abelian group, and $\mathfrak{m}$ an uncountable cardinal number. Does there always exist an abelian torsion-free group $U\left( {T,\mathfrak{m}}\right)$ of cardinality $\mathfrak{m}$ with the following property: for any abelian torsion-free group $A$ of cardinality $\leq  \mathfrak{m}$ , the equality $\operatorname{Ext}\left( {A, T}\right)  = 0$ holds if and only if $A$ is embeddable in $U\left( {T,\mathfrak{m}}\right)$ ?
	
	L. Ya. Kulikov
	
	No, not always. There is a model of ZFC in which for a certain class of cardinals $\mathfrak{m}$ the answer is negative (S. Shelah, L. Strüngmann, J. London Math. Soc., 67, no. 3 (2003), 626-642). On the other hand, under the assumption of Gödel's constructivity hypothesis $\left( {V = L}\right)$ the answer is affirmative for any cardinal if $T$ has only finitely many non-trivial bounded $p$ -components (L. Strüngmann, Ill. J. Math.,46, no. 2 (2002), 477-490).
\end{openproblem}

\plainproblemsection

\begin{openproblem}
	1.67. Suppose that $G$ is a finitely presented group, $F$ a free group whose rank is equal to the minimal number of generators of $G$ , with a fixed homomorphism of $F$ onto $G$ with kernel $N$ . Find a complete system of invariants of the factor-group of $N$ by the commutator subgroup $\left\lbrack  {F, N}\right\rbrack$ . L. Ya. Kulikov
\end{openproblem}

\plainproblemsection

\begin{openproblem}
	1.68. (A. Tarski). Let $\mathfrak{K}$ be a class of groups and $Q\mathfrak{K}$ the class of all homomorphic images of groups from $\mathfrak{K}$ . If $\mathfrak{K}$ is axiomatizable, does it follow that $Q\mathfrak{K}$ is? Not always (F. Clare, Algebra Universalis, 5 (1975), 120-124). Yu. I. Merzlyakov
\end{openproblem}

\plainproblemsection

\begin{openproblem}
	1.69. (B. I. Plotkin). Do there exist locally nilpotent torsion-free groups without the property $R{N}^{ * }$ ? Yu. I. Merzlyakov
	
	Yes, there exist (E. M. Levich, A. I. Tokarenko, Siberian Math. J., 11 (1970), 1033- 1034; A. I. Tokarenko, Trudy Rizhsk. Algebr. Sem., Riga Univ., Riga, 1969, 280-281 (Russian)).
\end{openproblem}

\plainproblemsection

\begin{openproblem}
	1.70. Let $p$ be a prime number and let $G$ be the group of all matrices of the form $\left( \begin{matrix} 1 + {p\alpha } & {p\beta } \\  {p\gamma } & 1 + {p\delta } \end{matrix}\right)$ , where $\alpha ,\beta ,\gamma ,\delta$ are rational numbers with denominators coprime to $p$ . Does $G$ have the property $\overline{RN}$ ?
	
	Yes, it does (G. A. Noskov, Siberian Math. J., 14 (1973), 475-477).
\end{openproblem}

\plainproblemsection

\begin{openproblem}
	1.71. Let $G$ be a connected algebraic group over an algebraically closed field. Is the number of conjugacy classes of maximal soluble subgroups of $G$ finite?
	
	Yes, it is (V. P. Platonov, Siberian Math. J., 10 (1969), 800-804). V. P. Platonov
\end{openproblem}

\plainproblemsection

\begin{openproblem}
	1.72. D. Hertzig has shown that a connected algebraic group over an algebraically closed field is soluble if it has a rational regular automorphism. Is this result true for an arbitrary field? V. P. Platonov
	
	No (V. P. Platonov, Soviet Math. Dokl., 7 (1966), 825-829).
\end{openproblem}

\plainproblemsection

\begin{openproblem}
	1.73. Are there only finitely many conjugacy classes of maximal periodic subgroups in a finitely generated linear group over the integers?
	
	Not always. The extension $H$ of the free group on free generators $a, b$ by the automorphism $\varphi  : a \rightarrow  {a}^{-1}, b \rightarrow  b$ is a linear group over the integers. For every $n \in  \mathbb{Z}$ the element ${c}_{n} = \varphi {b}^{-n}a{b}^{n}$ has order 2 and ${C}_{H}\left( {c}_{n}\right)  = \left\langle  {c}_{n}\right\rangle$ . Let the dash denote an isomorphism of $H$ onto its copy ${H}^{\prime }$ . As a counterexample one can take the subgroup $G \leq  H \times  {H}^{\prime }$ generated by the elements $a,\varphi ,{a}^{\prime }{\varphi }^{\prime }, b{b}^{\prime }$ . Indeed, $G$ contains the subgroups ${T}_{n} = \left\langle  {{c}_{0},{c}_{n}^{\prime }}\right\rangle  , n \in  \mathbb{Z}$ , which are maximal periodic (even in $H \times  {H}^{\prime }$ ). If ${T}_{n}$ and ${T}_{m}$ are conjugate by an element $x{y}^{\prime } \in  G$ (where $x \in  H$ and ${y}^{\prime } \in  {H}^{\prime }$ ), then ${c}_{0}^{x} = {c}_{0}$ and ${c}_{n}^{y} = {c}_{m}$ , whence $x \in  \left\langle  {c}_{0}\right\rangle  , y \in  {b}^{m - n}\left\langle  {c}_{m}\right\rangle$ . Since $x{y}^{\prime } \in  G$ , the sums of the exponents at the occurrences of $b$ in $x$ and $y$ (in any expression) must coincide; hence $m = n$ . (Yu. I. Merzlyakov, 1973.)
\end{openproblem}

\plainproblemsection

\begin{openproblem}
	1.74. Describe all minimal topological groups, that is, non-discrete groups all of whose closed subgroups are discrete. The minimal locally compact groups can be described without much effort. At the same time, the problem is probably complicated in the general case. V. P. Platonov
\end{openproblem}

\plainproblemsection

\begin{openproblem}
	1.75. Classify the infinite simple periodic linear groups over a field of characteristic $p > 0$ . V. P. Platonov
	
	They are classified for every $p{\;\operatorname{mod}\;\operatorname{CFSG}}(\mathrm{V}.\mathrm{V}$ . Belyaev, in: Investigations in Group Theory, Sverdlovsk, UNC AN SSSR, 1984, 39-50 (Russian); A. V. Borovik, Siberian Math. J., 24, no. 6 (1983), 843-851; B. Hartley, G. Shute, Quart. J. Math. Oxford (2), 35 (1984), 49-71; S. Thomas, Arch. Math., 41 (1983), 103-116). For $p > 2$ there is a proof without using the CFSG (A. V. Borovik, Siberian Math. J., 25, no. 2 (1984), 221-235). For all $p$ , without using CFSG, the classification also follows from the paper (M. J. Larsen, R. Pink, J. Amer. Math. Soc., 24 (2011), 1105-1158) also known as a preprint of 1998.
\end{openproblem}

\plainproblemsection

\begin{openproblem}
	1.76. Does there exist a simple, locally nilpotent, locally compact, topological group? No (I. V. Protasov, Soviet Math. Dokl., 19 (1978), 487-489). V. P. Platonov
\end{openproblem}

\plainproblemsection

\begin{openproblem}
	1.78. Let a group $G$ be the product of two divisible abelian $p$ -groups of finite rank. Is then $G$ itself a divisible abelian $p$ -group of finite rank? N. F. Sesekin
	
	Yet, it is (N. F. Sesekin, Siberian Math. J., 9 (1968), 1070-1072).
\end{openproblem}

\plainproblemsection

\begin{openproblem}
	1.80. Does there exist a finite simple group whose Sylow 2-subgroup is a direct product of quaternion groups? A. I. Starostin
	
	No (G. Glauberman, J. Algebra, 4 (1966), 403-420).
\end{openproblem}

\plainproblemsection

\begin{openproblem}
	1.81. The width of a group $G$ is, by definition, the smallest cardinal $m = m\left( G\right)$ with the property that any subgroup of $G$ generated by a finite set $S \subseteq  G$ is generated by a subset of $S$ of cardinality at most $m$ .
	
	a) Does a group of finite width satisfy the minimum condition for subgroups?
	
	b) Does a group with the minimum condition for subgroups have finite width?
	
	c) The same questions under the additional condition of local finiteness. In particular, is a locally finite group of finite width a Chernikov group? $L.N.$ Shevrin a) Not always (S. V. Ivanov, Geometric methods in the study of groups with given subgroup properties, Cand. Diss., Moscow Univ., Moscow, 1988 (Russian)).
	
	b) Not always (G. S. Deryabina, Math. USSR-Sb., 52 (1985), 481-490).
	
	c) Yes, it is (V. P. Shunkov, Algebra and Logic, 9 (1970), 137-151; 10 (1971), 127-142).
\end{openproblem}

\plainproblemsection

\begin{openproblem}
	1.82. Two sets of identities are said to be equivalent if they determine the same variety of groups. Construct an infinite set of identities which is not equivalent to any finite one. A. L. Shmel'kin
	
	This was done (S. I. Adian, Math. USSR-Izv., 4 (1970), 721-739).
\end{openproblem}

\plainproblemsection

\begin{openproblem}
	1.83. Does there exist a simple group, the orders of whose elements are unbounded, in which a non-trivial identity relation holds? A. L. Shmel'kin
	
	Yes, there exists (V. S. Atabekyan, Dep. no. 5381-V86, VINITI, Moscow, 1986 (Russian)).
\end{openproblem}

\plainproblemsection

\begin{openproblem}
	1.84. Is it true that a polycyclic group $G$ is residually a finite $p$ -group if and only if $G$ has a nilpotent normal torsion-free subgroup of $p$ -power index? A. L. Shmel’kin No, it is not (K. Seksenbaev, Algebra i Logika, 4, no. 3 (1965), 79-83 (Russian)).
\end{openproblem}

\plainproblemsection

\begin{openproblem}
	1.85. Is it true that the identity relations of a metabelian group have a finite basis? Yes, it is (D. E. Cohen, J. Algebra, 5 (1967), 267-273). A. L. Shmel'kin
\end{openproblem}

\plainproblemsection

\begin{openproblem}
	1.86. Is it true that the identical relations of a polycyclic group have a finite basis? A. L. Shmel'kin
\end{openproblem}

\plainproblemsection

\begin{openproblem}
	1.87. The same question for matrix groups (at least over a field of characteristic 0 ). A. L. Shmel'kin
\end{openproblem}

\plainproblemsection

\begin{openproblem}
	1.88. Is it true that if a matrix group over a field of characteristic 0 does not satisfy any non-trivial identity relation, then it contains a non-abelian free subgroup?
	
	Yes, it is (J. Tits, J. Algebra, 20 (1972), 250-270). A. L. Shmel'kin
\end{openproblem}

\plainproblemsection

\begin{openproblem}
	1.89. Is the following assertion true? Let $G$ be a free soluble group, and $a$ and $b$ elements of $G$ whose normal closures coincide. Then there is an element $x \in  G$ such that ${b}^{\pm 1} = {x}^{-1}{ax}$ . A. L. Shmel'kin
	
	No (A. L. Shmel'kin, Algebra i Logika, 6, no. 2 (1967), 95-109 (Russian)).
\end{openproblem}

\plainproblemsection

\begin{openproblem}
	1.90. A subgroup $H$ of a group $G$ is called 2-infinitely isolated in $G$ if, whenever the centralizer ${C}_{G}\left( h\right)$ in $G$ of some element $h \neq  1$ of $H$ contains at least one involution and intersects $H$ in an infinite subgroup, it follows that ${C}_{G}\left( h\right)  \leq  H$ . Let $G$ be an infinite simple locally finite group whose Sylow 2-subgroups are Chernikov groups, and suppose $G$ has a proper 2-infinitely isolated subgroup $H$ containing some Sylow 2-subgroup of $G$ . Does it follow that $G$ is isomorphic to a group of the type ${PS}{L}_{2}\left( k\right)$ , where $k$ is a field of odd characteristic? V. P. Shunkov
	
	Yes, it does (V.P. Shunkov, Algebra and Logic, 11 (1972), 260-272).
\end{openproblem}

\plainproblemsection

\begin{openproblem}
	2.1. Classify the finite groups having a Sylow $p$ -subgroup as a maximal subgroup. V. A. Belonogov, A. I. Starostin
	
	A description can be extracted from (B. Baumann, J. Algebra, 38 (1976), 119-135) and (L. A. Shemetkov, Math. USSR-Izv., 2 (1968), 487-513).
\end{openproblem}

\plainproblemsection

\begin{openproblem}
	2.2. A quasigroup is a groupoid $Q\left( \cdot \right)$ in which the equations ${ax} = b$ and ${ya} = b$ have a unique solution for any $a, b \in  Q$ . Two quasigroups $Q\left( \cdot \right)$ and $Q\left( \circ \right)$ are isotopic if there are bijections $\alpha ,\beta ,\gamma$ of the set $Q$ onto itself such that $x \circ  y = \gamma \left( {{\alpha x} \cdot  {\beta y}}\right)$ for all $x, y \in  Q$ . It is well-known that all quasigroups that are isotopic to groups form a variety $\mathfrak{G}$ . Let $\mathfrak{V}$ be a variety of quasigroups. Characterize the class of groups isotopic to quasigroups in $\mathfrak{G} \cap  \mathfrak{V}$ . For which identities characterizing $\mathfrak{V}$ is every group isotopic to a quasigroup in $\mathfrak{G} \cap  \mathfrak{V}$ ? Under what conditions on $\mathfrak{V}$ does any group isotopic to a quasigroup in $\mathfrak{G} \cap  \mathfrak{V}$ consist of a single element? V. D. Belousov, A. A. Gvaramiya Every part is answered (A. A. Gvaramiya, Dep. no. 6704-V84, VINITI, Moscow, 1984 (Russian)).
\end{openproblem}

\plainproblemsection

\begin{openproblem}
	2.3. A finite group is called quasi-nilpotent ( $\Gamma$ -quasi-nilpotent) if any two of its (maximal) subgroups $A$ and $B$ satisfy one of the conditions 1) $A \leq  B,\;2)B \leq  A$ , 3) ${N}_{A}\left( {A \cap  B}\right)  \neq  A \cap  B \neq  {N}_{B}\left( {A \cap  B}\right)$ . Do the classes of quasi-nilpotent and $\Gamma$ -quasi-nilpotent groups coincide? Ya. G. Berkovich, M. I. Kravchuk
	
	No. The group $G = \left\langle  {x, y, z, t \mid  {x}^{4} = {y}^{4} = {z}^{2} = {t}^{3} = 1,\left\lbrack  {x, y}\right\rbrack   = z,\left\lbrack  {x, z}\right\rbrack   = \left\lbrack  {y, z}\right\rbrack   = 1}\right.$ , ${x}^{t} = y,{y}^{t} = {x}^{-1}{y}^{-1}\rangle$ is $\Gamma$ -quasi-nilpotent, but not quasi-nilpotent. Since $G/\Phi \left( G\right)  \cong$ ${\mathbb{A}}_{4}$ , the intersection of any two maximal subgroups $A$ and $B$ of $G$ equals $\Phi \left( G\right)$ , whence ${N}_{A}\left( {A \cap  B}\right)  \neq  A \cap  B \neq  {N}_{B}\left( {A \cap  B}\right)$ ; thus $G$ is $\Gamma$ -quasi-nilpotent. On the other hand, if ${A}_{1} = \left\langle  {z{x}^{2}, z{y}^{2}, t}\right\rangle$ and ${B}_{1} = \langle z, t\rangle$ , then ${N}_{{A}_{1}}\left( {{A}_{1} \cap  {B}_{1}}\right)  = {A}_{1} \cap  {B}_{1} = \langle t\rangle$ ; hence $G$ is not quasi-nilpotent. (V. D. Mazurov, 1973.)
\end{openproblem}

\plainproblemsection

\begin{openproblem}
	2.4. (S. Chase). Suppose that an abelian group $A$ can be written as the union of pure subgroups ${A}_{\alpha },\alpha  \in  \Omega$ , where $\Omega$ is the first non-denumerable ordinal, ${A}_{\alpha }$ is a free abelian group of denumerable rank, and, if $\beta  < \alpha$ , then ${A}_{\beta }$ is a direct summand of ${A}_{\alpha }$ . Does it follow that $A$ is a free abelian group? Yu. A. Bogan Not always (P. A. Griffith, Pacific J. Math., 29 (1969), 279-284).
\end{openproblem}

\plainproblemsection

\begin{openproblem}
	2.5. According to Plotkin, a group is called an ${NR}$ -group if the set of its nil-elements coincides with the locally nilpotent radical, or, which is equivalent, if every inner automorphism of it is locally stable. Can an ${NR}$ -group have a nil-automorphism that is not locally stable? V. G. Vilyatser
\end{openproblem}

\plainproblemsection

\begin{openproblem}
	2.6. In an ${NR}$ -group, the set of generalized central elements coincides with the nil-kernel. Is the converse true, that is, must a group be an ${NR}$ -group if the set of generalized central elements coincides with the nil-kernel? V. G. Vilyatser
\end{openproblem}

\plainproblemsection

\begin{openproblem}
	2.7. Find the cardinality of the set of all polyverbal operations acting on the class of all groups. O. N. Golovin
	
	It has the cardinality of the continuum (A. Yu. Ol'shanskiĭ, Math. USSR-Izv., 4 (1970), 381-389).
\end{openproblem}

\plainproblemsection

\begin{openproblem}
	2.8. (A. I. Mal'cev). Do there exist regular associative operations having the heredity property with respect to transition from the factors to their subgroups? $O.N$ . Golovin Yes, they exist (S. V. Ivanov, Trans. Moscow Math. Soc., 1993, 217-249).
\end{openproblem}

\plainproblemsection

\begin{openproblem}
	2.9. Do there exist regular associative operations on the class of groups satisfying the weakened Mal'cev condition (that is, monomorphisms of the factors of an arbitrary product can be glued together, generally speaking, into a homomorphism of the whole product), but not satisfying the analogous condition for epimorphisms of the factors?
\end{openproblem}

\plainproblemsection

\begin{openproblem}
	2.10. Prove an analogue of the Remak-Shmidt theorem for decompositions of a group into nilpotent products. O. N. Golovin
	
	Such an analogue is proved (V. V. Limanskiǐ, Trudy Moskov. Mat. Obshch., 39 (1979), 135-155 (Russian)).
\end{openproblem}

\plainproblemsection

\begin{openproblem}
	2.13. (Well-known problem). Let $G$ be a periodic group in which every $\pi$ -element commutes with every ${\pi }^{\prime }$ -element. Does $G$ decompose into the direct product of a maximal $\pi$ -subgroup and a maximal ${\pi }^{\prime }$ -subgroup? S. G. Ivanov
	
	Not always. S.I. Adian (Math. USSR-Izv., 5 (1971), 475-484) has constructed a group $A = A\left( {m, n}\right)$ which is torsion-free and has a central element $d$ such that $A\left( {m, n}\right) /\langle d\rangle  \cong  B\left( {m, n}\right)$ , the free $m$ -generator Burnside group of odd exponent $n \geq  {4381}$ . Given a prime $p$ coprime to $n$ , a counterexample with $\pi  = \{ p\}$ can be found in the form $G = A/\left\langle  {d}^{{p}^{k}}\right\rangle$ for some positive integer $k$ . Indeed, $\langle d\rangle /\left\langle  {d}^{{p}^{k}}\right\rangle$ is a maximal $p$ -subgroup of $G$ . Suppose that $A/\left\langle  {d}^{{p}^{k}}\right\rangle   = \langle d\rangle /\left\langle  {d}^{{p}^{k}}\right\rangle   \times  {H}_{k}/\left\langle  {d}^{{p}^{k}}\right\rangle$ for every $k$ . Then $\langle d\rangle  \cap  {H}_{k} = \left\langle  {d}^{{p}^{k}}\right\rangle$ for all $k$ and therefore $\langle d\rangle  \cap  H = 1$ , where $H = \mathop{\bigcap }\limits_{k}{H}_{k}$ . Since $H$ is torsion-free and is isomorphic to a subgroup of $A/\langle d\rangle  \cong  B\left( {m, n}\right)$ , we obtain $H = 1$ . This implies that $A$ is isomorphic to a subgroup of the Cartesian product of the abelian groups $A/{H}_{k}$ , a contradiction. (Yu. I. Merzlyakov,1973.)
\end{openproblem}

\plainproblemsection

\begin{openproblem}
	2.15. Does there exist a torsion-free group such that the factor group by some term of its upper central series is nontrivial periodic, with a bound on the orders of the elements? G. A. Karasëv
	
	Yes, there does (S. I. Adian, Proc. Steklov Inst. Math., 112 (1971), 61-69).
\end{openproblem}

\plainproblemsection

\begin{openproblem}
	2.16. A group $G$ is called conjugacy separable if any two of its elements are conjugate in $G$ if and only if their images are conjugate in every finite homomorphic image of $G$ . Is $G$ conjugacy separable in the following cases:
	
	a) $G$ is a polycyclic group,
	
	b) $G$ is a free soluble group,
	
	c) $G$ is a group of (all) integral matrices,
	
	d) $G$ is a finitely generated group of matrices,
	
	e) $G$ is a finitely generated metabelian group? M. I. Kargapolov
	
	a) Yes (V. N. Remeslennikov, Algebra and Logic, 8 (1969), 404-411); E. Formanek, J. Algebra, 42 (1976), 1-10).
	
	b) Yes (V. N. Remeslennikov, V. G. Sokolov, Algebra and Logic, 9 (1970), 342-349). c), d) Not always (V. P. Platonov, G. V. Matveev, Dokl. Akad. Nauk BSSR, 14 (1970), 777-779 (Russian); V. N. Remeslennikov, V. G. Sokolov, Algebra and Logic, 9 (1970), 342-349; V. N. Remeslennikov, Siberian Math. J., 12 (1971), 783-792).
	
	e) Not always. Let $p$ be a prime and let ${A}_{1},{A}_{2}$ be two copies of the additive group $\left\{  {m/{p}^{k} \mid  m, k \in  \mathbb{Z}}\right\}$ . Let ${b}_{1},{b}_{2}$ be the automorphisms of the direct sum $A = {A}_{1} \oplus  {A}_{2}$ defined by ${a}^{{b}_{2}} = {pa}$ for any $a \in  A$ and ${a}_{2}^{{b}_{1}} = {a}_{1} + {a}_{2}$ and ${a}_{1}^{{b}_{1}} = {a}_{1}$ for some fixed elements ${a}_{1} \in  {A}_{1},{a}_{2} \in  {A}_{2}$ . Let $G$ be the semidirect product of $A$ and the direct product $\left\langle  {b}_{1}\right\rangle   \times  \left\langle  {b}_{2}\right\rangle$ (of two infinite cyclics). It can be shown that the elements ${a}_{2}$ and ${a}_{2} + {a}_{1}/p$ are not conjugate in $G$ , but their images are conjugate in any finite quotient of $G$ . (M. I. Kargapolov, E. I. Timoshenko, Abstracts of the 4th All-Union Sympos. Group Theory, Akademgorodok, 1973, Novosibirsk, 1973, 86-88 (Russian).)
\end{openproblem}

\plainproblemsection

\begin{openproblem}
	2.17. Is it true that the wreath product $A\wr B$ of two groups that are conjugacy separable is itself conjugacy separable if and only if either $A$ is abelian or $B$ is finite? No (V. N. Remeslennikov, Siberian Math. J., 12 (1971), 783-792). M. I. Kargapolov
\end{openproblem}

\plainproblemsection

\begin{openproblem}
	2.18. Compute the ranks of the factors of the lower central series of a free soluble group. M. I. Kargapolov
	
	They were computed (V. G. Sokolov, Algebra and Logic, 8 (1969), 212-215; Yu. M. Gorchakov, G. P. Egorychev, Soviet Math. Dokl., 13 (1972), 565-568).
\end{openproblem}

\plainproblemsection

\begin{openproblem}
	2.19. Are finitely generated subgroups of a free soluble group finitely separable? M. I. Kargapolov
	
	Not always (S. A. Agalakov, Algebra and Logic, 22 (1983), 261-268).
\end{openproblem}

\plainproblemsection

\begin{openproblem}
	2.20. Is it true that the wreath products $A\wr B$ and ${A}_{1}\wr {B}_{1}$ are elementarily equivalent if and only if $A, B$ are elementarily equivalent to ${A}_{1},{B}_{1}$ , respectively?
	
	M. I. Kargapolov
	
	No, but if the word "elementarily" is replaced by the word "universally," then it is true (E. I. Timoshenko, Algebra and Logic, 7, no. 4 (1968), 273-276).
\end{openproblem}

\plainproblemsection

\begin{openproblem}
	2.21. Do the classes of Baer and Fitting groups coincide? Sh. S. Kemkhadze
	
	No (R. S. Dark, Math. Z., 105 (1968), 294-298).
\end{openproblem}

\plainproblemsection

\begin{openproblem}
	2.22. a) An abstract group-theoretic property $\sum$ is said to be radical (in our sense) if, in any group $G$ , the subgroup $\sum \left( G\right)$ generated by all normal $\sum$ -subgroups is a $\sum$ - subgroup itself (called the $\sum$ -radical of $G$ ). A radical property $\sum$ is said to be strongly radical if, for any group $G$ , the factor-group $G/\sum \left( G\right)$ contains no non-trivial normal $\sum$ -subgroups. Is the property $\overline{RN}$ radical? strongly radical? Sh. S. Kemkhadze
\end{openproblem}

\plainproblemsection

\begin{openproblem}
	2.22. b) An abstract group-theoretical property $\sum$ is called radical (in our sense) if, in any group $G$ , the join $\sum \left( G\right)$ of all normal $\sum$ -subgroups is a $\sum$ -subgroup. Is the property ${N}^{0}$ (see Archive,1.38) radical? Sh. S. Kemkhadze
	
	No. The group $S{L}_{n}\left( \mathbb{Z}\right)$ for sufficiently large $n \geq  3$ contains a non-abelian finite simple group and therefore is not an ${N}^{0}$ -group. On the other hand, as shown in (M. I. Kargapolov, Yu. I. Merzlyakov, in: Itogi Nauki. Algebra. Topologiya. Ge-ometriya, 1966, VINITI, Moscow, 1968, 57-90 (Russian)), it is the product of its congruence-subgroups mod 2 and mod 3 , which have central systems and hence are ${N}^{0}$ -groups (see Archive,1.38). (Yu. I. Merzlyakov,1973.)
\end{openproblem}

\plainproblemsection

\begin{openproblem}
	2.23. a) A subgroup $H$ of a group $G$ is called quasisubinvariant if there is a normal system of $G$ passing through $H$ . Let $\mathfrak{K}$ be a class of groups that is closed with respect to taking homomorphic images. A group $G$ is called an ${R}^{0}\left( \mathfrak{K}\right)$ -group if each of its non-trivial homomorphic images has a non-trivial quasisubinvariant $\mathfrak{K}$ -subgroup. Do ${R}^{0}\left( \mathfrak{K}\right)$ and $\overline{RN}$ coincide when $\mathfrak{K}$ is the class of all abelian groups?
	
	Sh. S. Kemkhadze
	
	No (J. S. Wilson, Arch. Math., 25 (1974), 574-577).
\end{openproblem}

\plainproblemsection

\begin{openproblem}
	2.24. Are Engel torsion-free groups orderable? A. I. Kokorin
\end{openproblem}

\plainproblemsection

\begin{openproblem}
	2.25. a) (L. Fuchs). Describe the groups which are linearly orderable in only finitely many ways. A. I. Kokorin
\end{openproblem}

\plainproblemsection

\begin{openproblem}
	2.25. b) Do there exist groups that are linearly orderable in countably many ways? A. I. Kokorin
	
	Yes, there do (R. N. Buttsworth, Bull. Austral. Math. Soc., 4 (1971), 97-104).
\end{openproblem}

\plainproblemsection

\begin{openproblem}
	2.26. (L. Fuchs). Characterize as abstract groups the multiplicative groups of orderable skew fields. A. I. Kokorin
\end{openproblem}

\plainproblemsection

\begin{openproblem}
	2.28. Can every orderable group be embedded in a pro-orderable group? (See 1.35 for the definition of pro-orderable groups.) A. I. Kokorin
\end{openproblem}

\plainproblemsection

\begin{openproblem}
	2.29. Does the class of finite groups in which every proper abelian subgroup is contained in a proper normal subgroup coincide with the class of finite groups in which every proper abelian subgroup is contained in a proper normal subgroup of prime index? P. G. Kontorovich, V. T. Nagrebetskii
	
	No. Let $B$ be a finite group such that $B = \left\lbrack  {B, B}\right\rbrack   \neq  1$ and let $r$ be the rank of $B$ . Put $A = \bigoplus_{p \mid |B|} (\mathbb{Z}/p\mathbb{Z})^{r+1}$ and $G = A\wr B$ . We define a homomorphism $\varphi  : G \rightarrow  A$ by setting ${\left( bf\right) }^{\varphi } = \mathop{\sum }\limits_{{x \in  B}}f\left( x\right)$ , where $b \in  B$ and $f \in  F = \operatorname{Fun}\left( {B, A}\right)$ . Suppose that ${f}^{b} = f$ for $b \in  B$ and $f \in  F$ . It is clear that ${f}^{\varphi } \in  {nA}$ , where $n = \left| b\right|$ . In particular, ${C}_{F}{\left( b\right) }^{\varphi } \leq  {pA} \leq  {O}_{{p}^{\prime }}\left( A\right)$ if $p \mid  n$ . Every proper abelian subgroup $H$ of $G$ is contained in a proper normal subgroup. Indeed, we may assume that $H \nleq  F$ . We fix an element ${bf} \in  H$ , where $f \in  F, b \in  B, b \neq  1$ . Let $p$ be a prime dividing $\left| b\right|$ . For any $h \in  H \cap  F$ we have ${bfh} = {hbf} = b{h}^{b}f = {bf}{h}^{b}$ , whence $h = {h}^{b}$ . Hence ${\left( H \cap  F\right) }^{\varphi } \leq  {C}_{F}{\left( b\right) }^{\varphi } \leq  {O}_{{p}^{\prime }}\left( A\right)$ . If $T$ is the full preimage of ${O}_{{p}^{\prime }}\left( A\right)$ in $G$ , then $G/T \cong  {\left( \mathbb{Z}/p\mathbb{Z}\right) }^{r + 1}$ and therefore $H/H \cap  T$ is an elementary abelian $p$ -group. The rank of it is $\leq  r$ , since $H/H \cap  T$ embeds into $B$ and $H \cap  F \leq  H \cap  T$ . Thus, ${HT}$ is a proper normal subgroup containing $H$ . On the other hand, $F$ is a proper abelian subgroup that is not contained in any proper normal subgroup of prime index, since $G/F \cong  B = \left\lbrack  {B, B}\right\rbrack$ . (G. M. Bergman, I. M. Isaacs, Letter of June,17,1974.)
\end{openproblem}

\plainproblemsection

\begin{openproblem}
	2.30. Does there exist a finite group in which a Sylow $p$ -subgroup is covered by other Sylow $p$ -subgroups? P. G. Kontorovich, A. L. Starostin
	
	Yes, there exists, for any prime $p$ (V.D.Mazurov, Ural Gos. Univ. Mat. Zap.,7, no. 3 (1969/70), 129-132 (Russian)).
\end{openproblem}

\plainproblemsection

\begin{openproblem}
	2.31. Can every group admitting an ordering with only finitely many convex subgroups be represented by matrices over a field? No. For example, the group $G = \left\langle  {{a}_{n},{b}_{n}\left( {n \in  \mathbb{Z}}\right) , c, d \mid  \left\lbrack  {{a}_{n},{b}_{n}}\right\rbrack   = c,{a}_{n}^{d} = {a}_{n + 1}}\right.$ , ${b}_{n}^{d} = {b}_{n + 1},\left\lbrack  {{a}_{n},{a}_{m}}\right\rbrack   = \left\lbrack  {{b}_{n},{b}_{m}}\right\rbrack   = \left\lbrack  {{a}_{n}, c}\right\rbrack   = \left\lbrack  {{b}_{n}, c}\right\rbrack   = 1\;\left( {n, m \in  \mathbb{Z}}\right) \rangle$ is soluble and orderable with finitely many convex subgroups; but it is not residually finite and hence is not linear over a field. (V. A. Churkin, 1969.)
\end{openproblem}

\plainproblemsection

\begin{openproblem}
	2.33. Is a direct summand of a direct sum of finitely generated modules over a Noetherian ring again a direct sum of finitely generated modules? V. I. Kuz'minov Not always (P. A. Linnell, Bull. London Math. Soc., 14 (1982), 124-126).
\end{openproblem}

\plainproblemsection

\begin{openproblem}
	2.35. An inverse spectrum $\xi$ of abelian groups is said to be acyclic if $\mathop{\lim }\limits_{ \leftarrow  }^{\left( p\right) }\xi  = 0$ for $p > 0$ . Here $\mathop{\lim }\limits_{ \leftarrow  }\left( p\right)$ denotes the right derived functor of the projective limit functor. Let $\xi$ be an acyclic spectrum of finitely generated groups. Is the spectrum $\bigoplus {\xi }_{\alpha }$ also acyclic, where each spectrum ${\xi }_{\alpha }$ coincides with $\xi$ ? The answer depends on the axioms of Set Theory (A. A. Khusainov, Siberian Math. J., 37, no. 2 (1996), 405-413).
\end{openproblem}

\plainproblemsection

\begin{openproblem}
	2.36. (de Groot). Is the group of all continuous integer-valued functions on a compact space free abelian? V. I. Kuz'minov
	
	Yes, it is. By (G. Nöbeling, Invent. Math., 6 (1968) 41-55) the additive group of all bounded integer-valued functions on an arbitrary set is free. Hence the group of all continuous integer-valued functions on the Chech compactification of an arbitrary discrete space is also free. For any compact space $X$ there exists a continuous mapping of the Chech compactification $Y$ of a discrete space onto $X$ . This induces an embedding of the group of continuous integer-valued functions on $X$ into the free group of continuous integer-valued functions on $Y$ . (V. I. Kuz’minov,1969.)
\end{openproblem}

\plainproblemsection

\begin{openproblem}
	2.37. Describe the finite simple groups whose Sylow $p$ -subgroups are cyclic for all odd $p$ . V. D. Mazurov
	
	This was done (M. Aschbacher, J. Algebra, 54 (1978), 50-152).
\end{openproblem}

\plainproblemsection

\begin{openproblem}
	2.38. (Old problem). The class of rings embeddable in associative division rings is universally axiomatizable. Is it finitely axiomatizable? A. I. Mal’cev
	
	No (P. M. Cohn, Bull. London Math. Soc., 6 (1974), 147-148).
\end{openproblem}

\plainproblemsection

\begin{openproblem}
	2.39. Does there exist a non-finitely-axiomatizable variety of
	
	a) (H. Neumann) groups?
	
	b) of associative rings (the Specht problem)?
	
	c) Lie rings? A. I. Mal’cev
	
	a) Yes, there exists (S.I. Adian, Math. USSR-Izv., 4 (1970), 721-739; A. Yu. Ol'- shanskii, Math. USSR-Izv., 4 (1970), 381-389).
	
	b) Yes, there exists (A. Ya. Belov, Fundam. Prikl. Mat., 5, no. 1 (1999), 47-66 (Russian), Sb. Math., 191, no. 3-4 (2000), 329-340; A. V. Grishin, Fundam. Prikl. Mat., 5, no. 1 (1999),101-118 (Russian); V.V. Shchigolev, Fundam. Prikl. Mat.,5, no. 1 (1999), 307-312 (Russian)). But every variety of associative algebras over a field of characteristic 0 has a finite basis of identities (A. R. Kemer, Algebra and Logic, 26, no. 5 (1987), 362-397).
	
	c) Yes, there exists (M. R. Vaughan-Lee, Quart. J. Math., 21 (1970), 297-308).
\end{openproblem}

\plainproblemsection

\begin{openproblem}
	2.40. The $I$ -theory (Q-theory) of a class $\mathfrak{K}$ of universal algebras is the totality of all identities (quasi-identities) that are true on all the algebras in $\mathfrak{K}$ . Does there exist a finitely axiomatizable variety of
	
	a) groups,
	
	b) semigroups,
	
	c) (1) rings
	
	(2) of associative rings
	
	(i) whose $I$ -theory is non-decidable?
	
	(ii) whose $Q$ -theory is non-decidable?
	
	(3) of Lie rings
	
	(ii) whose $Q$ -theory is non-decidable?
	
	whose $I$ -theory ( $Q$ -theory) is non-decidable? A. I. Mal’cev
	
	a) Yes (Yu. G. Kleiman, Trans. Moscow Math. Soc., 1983, no. 2, 63-110).
	
	b) Yes (V. L. Murskiĭ, Math. Notes, 3 (1968), 423-427).
	
	c) (1) Yes (V. Yu. Popov, Math. Notes, 67 (2000), 495-504). (2i) No, it does not (A. Ya. Belov, Izv. Math. 74, no. 1 (2010), 1-126).
	
	(2ii) and (3ii): Yes, it exists (A. I. Budkin, Izv. Altai Univ., 65, no. 1 (2010), 15-17 (Russian)).
\end{openproblem}

\plainproblemsection

\begin{openproblem}
	2.40. c) The $I$ -theory ( $Q$ -theory) of a class $\mathfrak{K}$ of universal algebras is the totality of all identities (quasi-identities) that are valid on all algebras in $\mathfrak{K}$ . Does there exist a finitely axiomatizable variety
	
	(3) of Lie rings
	
	(i) whose $I$ -theory is non-decidable?
	
	Remark: it is not difficult to find a recursively axiomatizable variety of semigroups with identity whose $I$ -theory is non-recursive (see also A. I. Mal’cev, Mat. Sbornik, 69, no. 1 (1966), 3-12 (Russian)). A. I. Mal’cev
\end{openproblem}

\plainproblemsection

\begin{openproblem}
	2.41. Is the variety generated by
	
	a) a finite associative ring;
	
	b) a finite Lie ring;
	
	c) a finite quasigroup
	
	finitely axiomatizable?
	
	d) What is the cardinality $n$ of the smallest semigroup generating a non-finitely-axiomatizable variety? A. I. Mal’cev
	
	a) Yes (I. V. L'vov, Algebra and Logic, 12 (1973), 156-167; R. L. Kruse, J. Algebra, 26 (1973), 298-318).
	
	b) Yes (Yu. A. Bakhturin, A. Yu. Ol'shanskii, Math. USSR-Sb., 25 (1975), 507-523).
	
	c) Not always (M. R. Vaughan-Lee, Algebra Universalis, 9 (1979), 269-280).
	
	d) It is proved that $n = 6$ (P. Perkins, J. Algebra,11 (1969),298-314; A. N. Trakht-man, Semigroup Forum, 27 (1983), 387-389).
\end{openproblem}

\plainproblemsection

\begin{openproblem}
	2.42. What is the structure of the groupoid of quasivarieties
	
	a) of all semigroups?
	
	b) of all rings?
	
	c) of all associative rings?
	
	Compare with A. I. Mal’cev, Siberian Math. J., 8, no. 2 (1967), 254-267).
	
	A. I. Mal’cev
\end{openproblem}

\plainproblemsection

\begin{openproblem}
	2.43. A group $G$ is called an ${FN}$ -group if the groups ${\gamma }_{i}G/{\gamma }_{i + 1}G$ are free abelian and $\mathop{\bigcap }\limits_{{i = 1}}^{\infty }{\gamma }_{i}G = 1$ , where ${\gamma }_{i + 1}G = \left\lbrack  {{\gamma }_{i}G, G}\right\rbrack$ . A variety of groups $\mathfrak{M}$ is called a $\sum$ -variety (where $\sum$ is an abstract property) if the $\mathfrak{M}$ -free groups have the property $\sum$ .
	
	a) Which properties $\sum$ are preserved under multiplication and intersection of varieties? Is the ${FN}$ property preserved under these operations?
	
	b) Are all varieties obtained by multiplication and intersection from the nilpotent varieties ${\mathfrak{N}}_{1},{\mathfrak{N}}_{2},\ldots$ (where ${\mathfrak{N}}_{1}$ is the variety of abelian groups) ${FN}$ -varieties?
	
	A. I. Mal’cev
	
	a) The property ${FN}$ is preserved by multiplication of varieties (A. L. Shmel’kin, Trans. Moscow Math. Soc., $\mathbf{{29}}\left( {1973}\right) ,{239} - {252}$ ). The property ${FN}$ is not preserved by the intersection of varieties. The following example is due to L. G. Kovács. Let $\mathfrak{U}$ and $\mathfrak{V}$ be the varieties of all nilpotent groups of class at most 4 satisfying the identities $\left\lbrack  {x, y, y, x}\right\rbrack   \equiv  1$ and $\left\lbrack  {\left\lbrack  {x, y}\right\rbrack  ,\left\lbrack  {z, t}\right\rbrack  }\right\rbrack   \equiv  1$ , respectively. Then $\mathfrak{U}$ and $\mathfrak{V}$ are ${FN}$ -varieties because 1) both of them are nilpotent of class 4 and contain all nilpotent groups of class $\leq  3$ , and 2) relatively free groups in $\mathfrak{U}$ and $\mathfrak{V}$ are torsion-free (well-known for $\mathfrak{V}$ and follows for $\mathfrak{U}$ from (P. Fitzpatrick, L. G. Kovács, J. Austral. Math. Soc. Ser. A, 35, no. $1\left( {1983}\right) ,{59} - {73}))$ . On the other hand, $\mathfrak{U} \cap  \mathfrak{V}$ is not an ${FN}$ -variety because the relatively free group in $\mathfrak{U} \cap  \mathfrak{V}$ of rank 3 is not torsion-free. Indeed, every torsion-free group in $\mathfrak{U} \cap  \mathfrak{V}$ is of class $\leq  3$ (follows from the paper by Fitzpatrick and Kovács cited above), and there is a 3-generated (exponent 2)-by-(exponent 2) group of class precisely 4 in $\mathfrak{U} \cap  \mathfrak{V}$ . (A. N. Krasil'nikov, Letter of July,17,1998.)
	
	b) Yes (Yu. M. Gorchakov, Algebra i Logika, 6, no. 3 (1967), 25-30 (Russian)).
\end{openproblem}

\plainproblemsection

\begin{openproblem}
	2.44. Let $\mathfrak{A}$ and $\mathfrak{B}$ be subvarieties of a variety of groups $\mathfrak{M}$ ; then $\left( {\mathfrak{A}\mathfrak{B}}\right)  \cap  \mathfrak{M}$ is called the $\mathfrak{M}$ -product of $\mathfrak{A}$ by $\mathfrak{B}$ , where $\mathfrak{{AB}}$ is the usual product. Does there exist a non-abelian variety $\mathfrak{M}$ with an infinite lattice of subvarieties and commutative $\mathfrak{M}$ - multiplication? A. I. Mal’cev
	
	Yes, there exists; for example, $\mathfrak{M} = {\mathfrak{A}}_{p}{\mathfrak{A}}_{p}$ , where ${\mathfrak{A}}_{p}$ is the variety of all abelian groups of prime exponent $p$ (Yu. M. Gorchakov, Talk of June 21,1967, Krasnoyarsk).
\end{openproblem}

\plainproblemsection

\begin{openproblem}
	2.45. (P. Hall). Prove or refute the following conjectures:
	
	a) If a word $v$ takes only finitely many values in a group $G$ , then the verbal subgroup ${vG}$ is finite.
	
	c) If $G$ satisfies the maximum condition and ${vG}$ is finite, then ${v}^{ * }G$ has finite index in $G$ . Yu. I. Merzlyakov
	
	These conjectures are refuted in:
	
	a) (S. V. Ivanov, Soviet Math. (Izv. VUZ), 33, no. 6 (1989), 59-70;
	
	c) I. S. Ashmanov, A. Yu. Ol'shanskiĭ, Soviet Math. (Izv. VUZ), 29, no. 11 (1985),
	
	65-82).
\end{openproblem}

\plainproblemsection

\begin{openproblem}
	2.45. b) (P. Hall). Prove or refute the following conjecture: If the marginal subgroup ${v}^{ * }G$ has finite index $m$ in $G$ , then the order of ${vG}$ is finite and divides a power of $m$ .
	
	Editors’ comment: There are examples when $\left| {vG}\right|$ does not divide a power of $m$ (Yu. G. Kleiman, Trans. Moscow Math. Soc., 1983, no. 2, 63-110), but the question of finiteness remains open. Yu. I. Merzlyakov
\end{openproblem}

\plainproblemsection

\begin{openproblem}
	2.46. Find conditions under which a finitely-generated matrix group is almost residually a finite $p$ -group for some prime $p$ .
	
	A finitely generated matrix group over a field of characteristic zero has a subgroup of finite index which is residually a finite $p$ -group for almost all primes $p$ (M. I. Kargapolov, Algebra i Logika, 6, no. 5 (1967), 17-20 (Russian); Yu. I. Merzlyakov, Soviet Math. Dokl., 8 (1967), 1538-1541; V. P. Platonov, Dokl. Akad. Nauk Belarus. SSR, 12, no. 6 (1968), 492-494 (Russian)). A finitely generated matrix group over a field of characteristic $p > 0$ has a subgroup of finite index which is residually a finite $p$ -group (V. P. Platonov, Dokl. Akad. Nauk Belarus. SSR, 12, no. 6 (1968), 492-494 (Russian); A.I. Tokarenko, Proc. Riga Algebr. Sem., Latv. Univ., Riga, 1969, 280-281 (Russian)). See also some generalizations in (V. N. Remeslennikov, Algebra i Logika, 7, no. 4 (1968), 106-113 (Russian); B. A. F. Wehrfritz, Proc. London Math. Soc., 20, no. 1 (1970), 101-122).
\end{openproblem}

\plainproblemsection

\begin{openproblem}
	2.47. In which abelian groups is the lattice of all fully invariant subgroups a chain? A. P. Mishina
	
	Such groups were characterized (M.-P. Brameret, Séminaire P. Dubreil, M.-L. Dubreil-Jacotin, L. Lesieur, et C. Pisot, 16 (1962/63), no. 13, Paris, 1967).
\end{openproblem}

\plainproblemsection

\begin{openproblem}
	*2.48. (N. Aronszajn). Let $G$ be a connected topological group locally satisfying some identical relation ${\left. f\right| }_{U} = 1$ , where $U$ is a neighborhood of the identity element of $G$ . Is it then true that ${\left. f\right| }_{G} = 1$ ? V. P. Platonov
	
	*No, not necessarily: for any odd integer $n \geq  {10}^{10}$ there is a connected topological group such that the identity ${x}^{n} = 1$ holds in some neighborhood of unity, but not in the entire group (E. Reznichenko, I. Zyabrev, Preprint, 2024, https://arxiv.org/ abs/2406.05203).
\end{openproblem}

\plainproblemsection

\begin{openproblem}
	2.49. (A. Selberg). Let $G$ be a connected semisimple linear Lie group whose corresponding symmetric space has rank greater than 1, and let $\Gamma$ be an irreducible discrete subgroup of $G$ such that $G/\Gamma$ has finite volume. Does it follow that $\Gamma$ is an arithmetic subgroup? V. P. Platonov
	
	Yes, it does (G. A. Margulis, Functional Anal. Appl., 8 (1974), 258-259).
\end{openproblem}

\plainproblemsection

\begin{openproblem}
	2.50. (A. Selberg). Let $\Gamma$ be an irreducible discrete subgroup of a connected Lie group $G$ such that the factor space $G/\Gamma$ is non-compact but has finite volume in the Haar measure. Prove that $\Gamma$ contains a non-trivial unipotent element. V.P. Platonov This is proved (D. A. Kazhdan, G. A. Margulis, Math. USSR-Sb., 4 (1968), 147-152).
\end{openproblem}

\plainproblemsection

\begin{openproblem}
	2.51. (A. Borel, R. Steinberg). Let $G$ be a semisimple algebraic group and ${R}_{G}$ the set of classes of conjugate unipotent elements of $G$ . Is ${R}_{G}$ finite? $\;V.P$ . Platonov Yes, it is (G. Lusztig, Invent. Math., 34 (1976), 201-213).
\end{openproblem}

\plainproblemsection

\begin{openproblem}
	2.52. (F. Bruhat, N. Iwahori, M. Matsumoto). Let $G$ be a semisimple algebraic group over a locally compact, totally disconnected field. Do the maximal compact subgroups of $G$ fall into finitely many conjugacy classes? If so, estimate this number. V. P. Platonov
	
	Yes, they do, and a formula for this number is found (F. Bruhat, J. Tits, Publ. Math. IHES, 41 (1972), 5-251).
\end{openproblem}

\plainproblemsection

\begin{openproblem}
	2.54. Can ${SL}\left( {n, k}\right)$ have maximal subgroups that are not closed in the Zariski topology? V. P. Platonov
	
	Yes, it can if $k$ is either $\mathbb{Q}$ or an algebraically closed field of characteristic zero (N. S. Romanovskii', Algebra i Logika, 6, no. 4 (1967), 75-82; 7, no. 3 (1968), 123 (Russian)).
\end{openproblem}

\plainproblemsection

\begin{openproblem}
	2.55. b) Does $S{L}_{n}\left( \mathbb{Z}\right) , n \geq  2$ , have maximal subgroups of infinite index? V. P. Platonov
	
	Yes, it does (G. A. Margulis, G. A. Soifer, Soviet Math. Dokl., 18 (1977), 847-851).
\end{openproblem}

\plainproblemsection

\begin{openproblem}
	2.56. Classify up to isomorphism the abelian connected algebraic unipotent linear groups over a field of positive characteristic. This is not difficult in the case of a field of characteristic zero. On the other hand, C. Chevalley has solved the classification problem for such groups up to isogeny. V. P. Platonov
\end{openproblem}

\plainproblemsection

\begin{openproblem}
	2.58. Let $G$ be a vector space and $\Gamma$ a group of automorphisms of $G.\Gamma$ is called locally finitely stable if, for any finitely generated subgroup $\Delta$ of $\Gamma , G$ has a finite series stable relative to $\Delta$ . If the characteristic of the field is zero and $\Gamma$ is locally finitely stable, then $\Gamma$ is locally nilpotent and torsion-free. Is it true that every locally nilpotent torsion-free group can be realized in this way? B. I. Plotkin
	
	No (L. A. Simonyan, Siberian Math. J., 12 (1971), 602-606).
\end{openproblem}

\plainproblemsection

\begin{openproblem}
	2.59. Let $\Gamma$ be any nilpotent group of class $n - 1$ . Does $\Gamma$ always admit a faithful representation as a group of automorphisms of an abelian group with a series of length $n$ stable relative to $\Gamma$ ? B. I. Plotkin
	
	Not always (E. Rips, Israel J. Math., 12 (1972), 342-346).
\end{openproblem}

\plainproblemsection

\begin{openproblem}
	2.61. Let $\Gamma$ be a Noetherian group of automorphisms of a vector space $G$ such that every element of $\Gamma$ is unipotent. Is $\Gamma$ necessarily a stable group of automorphisms? B. I. Plotkin
	
	Not if the field has non-zero characteristic (A. Yu. Ol'shanskiĭ, The geometry of defining relations in groups, Kluwer, Dordrecht, 1991). This is true for fields of characteristic zero if the indices of unipotence are uniformly bounded (B. I. Plotkin, S. M. Vovsi, Varieties of group representations, Zinatne, Riga, 1983 (Russian)).
\end{openproblem}

\plainproblemsection

\begin{openproblem}
	2.62. If $G$ is a finite-dimensional vector space over a field and $\Gamma$ is a group of automorphisms of $G$ in which every element is stable, then the whole of $\Gamma$ is stable (E. Kolchin). Is Kolchin’s theorem true for spaces over skew fields? B. I. Plotkin Yes, it is, if the characteristic of the skew field is zero or sufficiently large compared to the dimension of the space (H.Y.Mochizuki, Canad. Math. Bull., 21 (1978), 249-250).
\end{openproblem}

\plainproblemsection

\begin{openproblem}
	2.63. Let $G$ be a group of automorphisms of a vector space over a field of characteristic zero, and suppose that all elements of $G$ are unipotent, with uniformly bounded unipotency indices. Must such a group be locally finitely stable? B. I. Plotkin Yes, it must, which follows from (H. Heineken, Arch. Math. (Basel), 13 (1962), 29- 37). Moreover, such a group is even finitely stable, as follows from (E. I. Zel'manov, Math. USSR-Sb., 66, no. 1 (1990), 159-168).
\end{openproblem}

\plainproblemsection

\begin{openproblem}
	2.64. Does the set of nil-elements of a finite-dimensional linear group coincide with its locally nilpotent radical? B. I. Plotkin
	
	Not always. The non-nilpotent 3-generator nilgroup of E. S. Golod constructed by an algebra over $\mathbb{Q}$ is residually torsion-free nilpotent; hence it is orderable and therefore it is embeddable into $G{L}_{n}$ over some skew field by Mal’cev’s theorem (V. A. Roman’kov, 1978).
\end{openproblem}

\plainproblemsection

\begin{openproblem}
	2.65. Does the adjoint group of a radical ring (in the sense of Jacobson) have a central series? B. I. Plotkin
	
	Not always (O. M. Neroslavskii, Vesci Akad. Navuk BSSR, Ser. Fiz.-Mat. Navuk, 1973, No. 2, 5-10 (Russian)).
\end{openproblem}

\plainproblemsection

\begin{openproblem}
	2.66. Is an $R$ -group determined by its subgroup lattice? Is every lattice isomorphism of $R$ -groups induced by a group isomorphism? L. E. Sadovskiĭ Not always, in both cases (A. Yu. Ol'shanskiĭ, C. R. Acad. Bulgare Sci., 32, no. 9 (1979), 1165-1166).
\end{openproblem}

\plainproblemsection

\begin{openproblem}
	2.67. Find conditions which ensure that the nilpotent product of pure nilpotent groups (from certain classes) is determined by the lattice of its subgroups. This is known to be true if the product is torsion-free. L. E. Sadovskiĭ
\end{openproblem}

\plainproblemsection

\begin{openproblem}
	2.68. What can one say about lattice isomorphisms of a pure soluble group? Is such a group strictly determined by its lattice? It is well-known that the answer is affirmative for free soluble groups. L. E. Sadovskiǐ
\end{openproblem}

\plainproblemsection

\begin{openproblem}
	2.69. Let a group $G$ be the product of two subgroups $A$ and $B$ , each of which is nilpotent and satisfies the minimum condition. Prove or refute the following: a) $G$ is soluble; b) the divisible parts of $A$ and $B$ commute elementwise. N.F. Sesekin Both parts are proved (N. S. Chernikov, Soviet Math. Dokl., 21 (1980), 701-703).
\end{openproblem}

\plainproblemsection

\begin{openproblem}
	2.70. a) Let a group $G$ be the product of two subgroups $A$ and $B$ , each of which is locally cyclic and torsion-free. Prove that either $A$ or $B$ has a non-trivial subgroup that is normal in $G$ .
	
	b) Characterize the groups that can be factorized in this way. N. F. Sesekin a) This was proved (D. I. Zaitsev, Algebra and Logic, 19 (1980), 94-106).
	
	b) They were characterized (Ya. P. Sysak, Algebra and Logic, 25 (1986), 425-433).
\end{openproblem}

\plainproblemsection

\begin{openproblem}
	2.71. Does there exist a finitely generated right-orderable group which coincides with its derived subgroup and, therefore, does not have the property ${RN}$ ? D. M. Smirnov Yes, there exists (G. M. Bergman, Pacific J. Math., 147, no. 2 (1991), 243-248).
\end{openproblem}

\plainproblemsection

\begin{openproblem}
	2.72. (G. Baumslag). Suppose that $F$ is a finitely generated free group, $N$ its normal subgroup and $V$ a fully invariant subgroup of $N$ . Is $F/V$ necessarily Hopfian if $F/N$ is Hopfian? D. M. Smirnov
	
	No, not necessarily (S. V. Ivanov, A. M. Storozhev, Geom. Dedicata, 114 (2005), 209-228).
\end{openproblem}

\plainproblemsection

\begin{openproblem}
	2.73. (Well-known problem). Does there exist an infinite group all of whose proper subgroups have prime order? A. I. Starostin
	
	Yes, there does (A. Yu. Ol'shanskiĭ, Algebra and Logic, 21 (1982), 369-418).
\end{openproblem}

\plainproblemsection

\begin{openproblem}
	2.74. (Well-known problem). Describe the finite groups all of whose involutions have soluble centralizers. A. I. Starostin
\end{openproblem}

\plainproblemsection

\begin{openproblem}
	2.75. Let $G$ be a periodic group containing an infinite family of finite subgroups whose intersection contains non-trivial elements. Does $G$ contain a non-trivial element with infinite centralizer? S. P. Strunkov
	
	Not always (K. I. Lossov, Dep. no. 5528-V88, VINITI, Moscow, 1988 (Russian)).
\end{openproblem}

\plainproblemsection

\begin{openproblem}
	2.76. Let $\Gamma$ be the holomorph of an abelian group $A$ . Find conditions for $A$ to be maximal among the locally nilpotent subgroups of $\Gamma$ . D. A. Suprunenko
	
	These are found (A. V. Yagzhev, Mat. Zametki, 46, no. 6 (1989), 118 (Russian)).
\end{openproblem}

\plainproblemsection

\begin{openproblem}
	2.77. Let $A$ and $B$ be abelian groups. Find conditions under which every extension of $A$ by $B$ is nilpotent. D. A. Suprunenko
	
	These are found (A. V. Yagzhev, Math. Notes, 43 (1988), 244-245).
\end{openproblem}

\plainproblemsection

\begin{openproblem}
	2.78. Any set of all subgroups of the same given order of a finite group $G$ that contains at least one non-normal subgroup is called an $I{E}_{\bar{n}}$ -system of $G$ . A positive integer $k$ is called a soluble (non-soluble; simple; composite; absolutely simple) group-theoretic number if every finite group having exactly $k{I}_{\bar{n}}$ -systems is soluble (respectively, if there is at least one non-soluble finite group having $k{I}_{\bar{n}}$ -systems; if there is at least one simple finite group having $k\;I{E}_{\bar{n}}$ -systems; if there are no simple finite groups having $kI{E}_{\bar{n}}$ -systems; if there is at least one simple finite group having $k{I}_{\bar{n}}$ -systems and there are no non-soluble non-simple finite groups having $\left. {k{I}_{\bar{n}}\text{-systems }}\right)$ .
	
	Are the sets of all soluble and of all absolutely simple group-theoretic numbers finite or infinite? Do there exist composite, but not soluble group-theoretic numbers?
\end{openproblem}

\plainproblemsection

\begin{openproblem}
	2.79. Do there exist divisible (simple) groups with maximal subgroups? M. S. Tsalenko
	
	Yes, there exist, for example, (V. G. Sokolov, Algebra and Logic, 7 (1968), 122-126).
\end{openproblem}

\plainproblemsection

\begin{openproblem}
	2.80. Does every non-trivial group satisfying the normalizer condition contain a non-trivial abelian normal subgroup? S. N. Chernikov
\end{openproblem}

\plainproblemsection

\begin{openproblem}
	2.81. a) Does there exist an axiomatizable class of lattices $\mathfrak{K}$ such that the lattice of all subsemigroups of a semigroup $S$ is isomorphic to some lattice in $\mathfrak{K}$ if and only if $S$ is a free group?
	
	b) The same question for free abelian groups.
	
	Analogous questions have affirmative answers for torsion-free groups, for nonperiodic groups, for abelian torsion-free groups, for abelian non-periodic groups, for orderable groups (the corresponding classes of lattices are even finitely axiomatizable). Thus, in posed questions one may assume from the outset that the semigroup $S$ is a torsion-free group (respectively, a torsion-free abelian group). L. N. Shevrin
\end{openproblem}

\plainproblemsection

\begin{openproblem}
	2.82. Can the class of groups with the $n$ th Engel condition $\left\lbrack  {x,\underset{n}{\underbrace{y,\ldots , y}}}\right\rbrack   = 1$ be
	
	defined by identical relations of the form $u = v$ , where $u$ and $v$ are words without negative powers of variables? This can be done for $n = 1,2,3$ (A. I. Shirshov, Algebra i Logika, 2, no. 5 (1963), 5-18 (Russian)).
	
	Editors’ comment: This has also been done for $n = 4$ (G. Traustason, J. Group Theory, 2, no. 1 (1999), 39-46). As remarked by O. Macedońska, this is also true for the class of locally graded $n$ -Engel groups because by (Y. Kim, A. Rhemtulla, in Groups-Korea'94, de Gruyter, Berlin, 1995, 189-197) such a group is locally nilpotent and then by (R. G. Burns, Yu. Medvedev, J. Austral. Math. Soc., 64 (1998), 92-100) such a group is an extension of a nilpotent group of $n$ -bounded class by a group of $n$ - bounded exponent; then a classical result of Mal'cev implies that such groups satisfy a positive law. A. I. Shirshov
\end{openproblem}

\plainproblemsection

\begin{openproblem}
	2.83. Suppose that a periodic group $G$ is the product of two locally finite subgroups. Is then $G$ locally finite? V. P. Shunkov
	
	Not always (V. I. Sushchanskii, Math. USSR-Sb., 67, no. 2 (1990), 535-553).
\end{openproblem}

\plainproblemsection

\begin{openproblem}
	2.84. Suppose that a locally finite group $G$ is a product of two locally nilpotent subgroups. Is $G$ necessarily locally soluble? V. P. Shunkov
\end{openproblem}

\plainproblemsection

\begin{openproblem}
	2.88. Is every Hall $\pi$ -subgroup of an arbitrary group a maximal $\pi$ -subgroup? M. I. Éidinov
	
	Not always (R. Baer, J. Reine Angew. Math., 239/240 (1969), 109-144).
\end{openproblem}

\plainproblemsection

\begin{openproblem}
	3.2. Classify the faithful irreducible (infinite-dimensional) representations of the nilpotent group defined by generators $a, b, c$ and relations $\left\lbrack  {a, b}\right\rbrack   = c,{ac} = {ca},{bc} = {cb}$ . (A condition for a representation to be monomial is given in (A. E. Zalesskii, Math. Notes, 9 (1971), 117-123).) S. D. Berman, A. E. Zalesskiĭ
	
	Classification of these representations up to equivalence does not seem to be feasible. D. Segal (Math. Proc. Cambridge Phil. Soc., 81 (1977), 201-208) proved that there exist primitive irreducible representations of this group (that is, representations that are not induced from a representation of any proper subgroup). There are results concerning classification of primitive ideals of the groups algebra of this group (over various fields). For instance, Zalesskii (ibid.) showed that every primitive ideal is maximal.
\end{openproblem}

\plainproblemsection

\begin{openproblem}
	3.3. (Well-known problem). Describe the automorphism group of the free associative algebra with $n$ generators, $n \geq  2$ .
	
	It is known that all automorphisms are tame for $n = 2$ (A. J. Czerniakiewicz, Trans. Amer. Math. Soc., 160 (1971), 393-401; 171 (1972), 309-315; L. G. Makar-Limanov, Funct. Anal. Appl.,4 (1971),262-264). For $n = 3$ , the Anick automorphism is wild (U. U. Umirbaev, J. Reine Angew. Math., 605 (2007), 165-178). There is an algorithm for recognizing automorphisms among endomorphisms of free associative algebras of finite rank over constructive fields (A. V. Jagzhev, Siberian. Math. J., 21 (1980), 142-146). L. A. Bokut’
\end{openproblem}

\plainproblemsection

\begin{openproblem}
	3.4. (Well-known problems). a) Is there an algorithm that decides, for any set of group words ${f}_{1},\ldots ,{f}_{m}$ (in a fixed set of variables ${x}_{1},{x}_{2},\ldots$ ) and a separate word $f$ , whether $f = 1$ is a consequence of ${f}_{1} = 1,\ldots ,{f}_{m} = 1$ ?
	
	b) Given words ${f}_{1},\ldots ,{f}_{m}$ , is there an algorithm that decides, for any word $f$ , whether $f = 1$ is a consequence of ${f}_{1} = 1,\ldots ,{f}_{m} = 1$ ?
	
	No, in both cases (Yu. G. Kleiman, Soviet Math. Dokl., 20 (1979), 115-119).
\end{openproblem}

\plainproblemsection

\begin{openproblem}
	3.5. Can a subgroup of a relatively free group be radicable? In particular, can a verbal subgroup of a relatively free group be radicable? N. R. Brumberg
\end{openproblem}

\plainproblemsection

\begin{openproblem}
	3.6. Describe the insoluble finite groups in which every soluble subgroup is either 2-closed, ${2}^{\prime }$ -closed, or isomorphic to ${\mathbb{S}}_{4}$ . V. A. Vedernikov
	
	This follows from (D. Gorenstein, J. H. Walter, Illinois J. Math., 6 (1962) 553-593; V. D. Mazurov, Soviet Math. Dokl., 7 (1966), 681-682; V. D. Mazurov, V. M. Sitnikov, S. A. Syskin, Algebra and Logic, 9 (1970), 187-204).
\end{openproblem}

\plainproblemsection

\begin{openproblem}
	3.7. An automorphism $\sigma$ of a group $G$ is called algebraic if, for any $g \in  G$ , the minimal $\sigma$ -invariant subgroup of $G$ containing $g$ is finitely generated. An automorphism $\sigma$ of $G$ is called an $e$ -automorphism if, for any $\sigma$ -invariant subgroups $A$ and $B$ , where $A$ is a proper subgroup of $B, A \smallsetminus  B$ contains an element $x$ such that $\left\lbrack  {x,\sigma }\right\rbrack   \in  B$ . Is every $e$ -automorphism algebraic? V. G. Vilyatser
	
	Yes, it is (A. V. Yagzhev, Algebra and Logic, 28 (1989), 83-85).
\end{openproblem}

\plainproblemsection

\begin{openproblem}
	3.8. Let $G$ be the free product of free groups $A$ and $B$ , and $V$ the verbal subgroup of $G$ corresponding to the equation ${x}^{4} = 1$ . Is it true that, if $a \in  A \smallsetminus  V$ and $b \in  B \smallsetminus  V$ , then ${\left( ab\right) }^{2} \notin  V$ ? V. G. Vilyatser
	
	No, it is not. If $A$ and $B$ are free groups with bases ${a}_{1},{a}_{2}$ and ${b}_{1},{b}_{2}$ , respectively, then, for example, the commutators $a = \left\lbrack  {\left\lbrack  {{a}_{1},{a}_{2},{a}_{2}}\right\rbrack  ,\left\lbrack  {{a}_{1},{a}_{2}}\right\rbrack  }\right\rbrack$ and $b = \left\lbrack  {\left\lbrack  {{b}_{1},{b}_{2},{b}_{2}}\right\rbrack  ,\left\lbrack  {{b}_{1},{b}_{2}}\right\rbrack  }\right\rbrack$ do not belong to $V$ , but ${\left( ab\right) }^{2} \in  V$ . Indeed, it follows from (C.R.B.Wright, Pacif. J. Math.,11, no. 1 (1961),387-394) that ${a}^{2} \in  V,{b}^{2} \in  V$ , and $\left\lbrack  {a, b}\right\rbrack   \in  V$ . (V. A. Roman'kov, Talk at the seminar Algebra and Logic, March, 17, 1970.)
\end{openproblem}

\plainproblemsection

\begin{openproblem}
	3.9. Does there exist an infinite periodic group with a finite maximal subgroup? Yu. M. Gorchakov
	
	Yes, there does (S. I. Adian, Proc. Steklov Inst. Math., 112 (1971), 61-69).
\end{openproblem}

\plainproblemsection

\begin{openproblem}
	3.10. Need the number of finite non-abelian simple groups contained in a proper variety of groups be finite? Yu. M. Gorchakov
	
	Yes, mod CFSG; see, for example, (G. A. Jones, J. Austral. Math. Soc., 17 (1974), 162-173).
\end{openproblem}

\plainproblemsection

\begin{openproblem}
	3.11. An element $g$ of a group $G$ is said to be generalized periodic if there exist ${x}_{1},\ldots ,{x}_{n} \in  G$ such that ${x}_{1}^{-1}g{x}_{1}\cdots {x}_{n}^{-1}g{x}_{n} = 1$ . Does there exist a finitely generated torsion-free group all of whose elements are generalized periodic? Yu. M. Gorchakov Yes, there does (A. P. Goryushkin, Siberian Math. J., 14 (1973), 146-148). Another example is $G = \left\langle  {a, b \mid  {\left( {b}^{2}\right) }^{a} = {b}^{-2},{\left( {a}^{2}\right) }^{b} = {a}^{-2}}\right\rangle$ . Then $N = \left\langle  {{a}^{2},{b}^{2},{\left( ab\right) }^{2}}\right\rangle$ is an abelian normal subgroup of $G$ and $G/N$ is non-cyclic of order 4 . If ${a}^{2l}{b}^{2m}{\left( ab\right) }^{2n} = 1$ , then after conjugating by $a$ we get ${a}^{2l}{b}^{-{2m}}{\left( ab\right) }^{-{2n}} = 1$ , whence ${a}^{4l} = 1$ . In view of the obvious homomorphism $G \rightarrow  \mathbb{Z}\left( {\operatorname{Aut}\mathbb{Z}}\right)$ that maps $a$ to the number $1 \in  \mathbb{Z}$ we have $l = 0$ . Similarly, $m = n = 0$ . Hence $N$ is free abelian of rank 3 . The squares of elements outside of $N$ are non-trivial; for example, ${\left( a{a}^{2l}{b}^{2m}{\left( ab\right) }^{2n}\right) }^{2} =$ ${a}^{2}{a}^{-1}\left( {{a}^{2l}{b}^{2m}{\left( ab\right) }^{2n}}\right) a{a}^{2l}{b}^{2m}{\left( ab\right) }^{2n} = {a}^{{4l} + 2} \neq  1$ . Hence $G$ is torsion-free. Since the square of any element $x \in  G$ belongs to $N$ , we have ${x}^{2}{\left( {x}^{2}\right) }^{a}{\left( {x}^{2}\right) }^{b}{\left( {x}^{2}\right) }^{ab} = 1$ . (V. A. Churkin, 1973.)
\end{openproblem}

\plainproblemsection

\begin{openproblem}
	3.12. (Well-known problem). Is a locally finite group with a full Sylow basis locally soluble? Yu. M. Gorchakov
\end{openproblem}

\plainproblemsection

\begin{openproblem}
	3.13. Is the elementary theory of subgroup lattices of finite abelian groups decidable? Yu. L. Ershov, A. I. Kokorin
	
	No (G. T. Kozlov, Algebra and Logic, 9 (1970), 104-107).
\end{openproblem}

\plainproblemsection

\begin{openproblem}
	3.14. Let $G$ be a finite group of $n \times  n$ matrices over a skew field $T$ of characteristic zero. Prove that $G$ has a soluble normal subgroup $H$ whose index in $G$ is not greater than some number $f\left( n\right)$ not depending on $G$ or $T$ . This problem is related to the representation theory of finite groups (the Schur index). A. E. Zalesskiǐ This is proved mod CFSG (B. Hartley, M. A. Shahabi Shojaei, Math. Proc. Cambridge Phil. Soc., 92 (1982), 55-64).
\end{openproblem}

\plainproblemsection

\begin{openproblem}
	3.15. A group $G$ is said to be $U$ -embeddable in a class $\mathfrak{K}$ of groups if, for any finite submodel $M \subset  G$ , there is a group $A \in  \mathfrak{K}$ such that $M$ is isomorphic to some submodel of $A$ . Are the following groups $U$ -embeddable in the class of finite groups:
	
	a) every group with one defining relation?
	
	b) every group defined by one relation in the variety of soluble groups of a given derived length? M. I. Kargapolov
	
	No, in both cases (A. I. Budkin, V. A. Gorbunov, Algebra and Logic, 14 (1975), 73- 84). Another example. The group $G = \left\langle  {a, b \mid  {\left( {b}^{2}\right) }^{a} = {b}^{3}}\right\rangle$ is non-Hopfian as proved in (G. Baumslag, D. Solitar, Bull. Amer. Math. Soc., 68, no. 3 (1962), 199-201) and therefore is not metabelian. We choose elements $a, b,{x}_{1},{x}_{2},{x}_{3},{x}_{4} \in  G$ such that $w = \left\lbrack  {\left\lbrack  {{x}_{1},{x}_{2}}\right\rbrack  ,\left\lbrack  {{x}_{3},{x}_{4}}\right\rbrack  }\right\rbrack   \neq  1$ , add to them 1, their inverses, and all the initial segments of the word $w$ in the alphabet $\left\{  {x}_{i}\right\}$ , and all the initial segments of the word ${a}^{-1}{b}^{2}a{b}^{-3}$ and of each of the words ${x}_{i}$ in the alphabet of $\{ a, b\}$ . Let $M$ be the resulting model. If $M$ was embeddable into a finite group ${G}_{0}$ , then ${G}_{0}$ would have to be metacyclic, and also would have to have elements satisfying $\left\lbrack  {\left\lbrack  {{x}_{1},{x}_{2}}\right\rbrack  ,\left\lbrack  {{x}_{3},{x}_{4}}\right\rbrack  }\right\rbrack   \neq  1$ , which is impossible. Hence $G$ is not $U$ -embeddable into finite groups. Since the group $G/{G}^{\left( k\right) }$ is also non-Hopfian for a suitable $k$ (ibid.), it is not $U$ -embeddable into finite groups for similar reasons. (Yu. I. Merzlyakov, 1969.)
\end{openproblem}

\plainproblemsection

\begin{openproblem}
	3.16. The word problem for a group admitting a single defining relation in the variety of soluble groups of derived length $n, n \geq  3$ . M. I. Kargapolov
\end{openproblem}

\plainproblemsection

\begin{openproblem}
	3.17. a) Is a non-abelian group with a unique linear ordering necessarily simple?
	
	b) Is a non-commutative orderable group simple if it has no non-trivial normal relatively convex subgroups? A. I. Kokorin
	
	Not always, in both cases (V. V. Bludov, Algebra and Logic, 13 (1974), 343-360).
\end{openproblem}

\plainproblemsection

\begin{openproblem}
	3.18. (B. H. Neumann). A group $U$ is called universal for a class $\mathfrak{K}$ of groups if $U$ contains an isomorphic image of every member of $\mathfrak{K}$ . Does there exist a countable group that is universal for the class of countable orderable groups? A. I. Kokorin No (M. I. Kargapolov, Algebra and Logic, 9 (1970), 257-261; D. B. Smith, Pacific J. Math., 35 (1970), 499-502).
\end{openproblem}

\plainproblemsection

\begin{openproblem}
	3.19. (A.I.Mal’cev). For an arbitrary linearly ordered group $G$ , does there exist a linearly ordered abelian group with the same order-type as $G$ ? A. I. Kokorin No (W. C. Holland, A. H. Mekler, S. Shelah, Order, 1 (1985), 383-397).
\end{openproblem}

\plainproblemsection

\begin{openproblem}
	3.20. Does the class of orderable groups coincide with the smallest axiomatizable class containing pro-orderable groups? (See 1.35.) A. I. Kokorin
\end{openproblem}

\plainproblemsection

\begin{openproblem}
	3.21. Let $\mathfrak{K}$ be the class of one-based models of signature $\sigma ,\vartheta$ a property that makes sense for models in $\mathfrak{K}$ , and ${\mathfrak{K}}_{0}$ the class of two-based models whose first base is a set $M$ taken from $\mathfrak{K}$ and whose second base consists of all submodels of $M$ with property $\vartheta$ and whose signature consists of the symbols in $\sigma$ together with $\in$ and $\subseteq$ in the usual set-theoretic sense. The elementary theory of ${\mathfrak{K}}_{0}$ is called the element- $\vartheta$ -submodel theory of $\mathfrak{K}$ . Are either of the following theories decidable:
	
	a) the element-pure-subgroup theory of abelian groups?
	
	b) the element-pure-subgroup theory of abelian torsion-free groups?
	
	c) the element- $\vartheta$ -subgroup theory of abelian groups, when the set of $\vartheta$ -subgroups is linearly ordered by inclusion? A. I. Kokorin
	
	No, in all cases (for a), c): G. T. Kozlov, Algebra and Logic, 9 (1970), 104-107, Algebra, no. 1, Irkutsk Univ., 1972, 21-23 (Russian); for b): É.I. Fridman, Algebra, no. 1, Irkutsk Univ., 1972 , 97-100 (Russian)).
\end{openproblem}

\plainproblemsection

\begin{openproblem}
	*3.22. Let $\xi  = \left\{  {{G}_{\alpha },{\pi }_{\beta }^{\alpha } \mid  \alpha ,\beta  \in  I}\right\}$ be a projective system (over a directed set $I$ ) of finitely generated free abelian groups. If all the projections ${\pi }_{\beta }^{\alpha }$ are epimorphisms and all the ${G}_{\alpha }$ are non-zero, does it follow that $\mathop{\lim }\limits_{ \rightarrow  }\xi  \neq  0$ ? Equivalently, suppose every finite set of elements of an abelian group $A$ is contained in a pure finitely-generated free subgroup of $A$ . Then does it follow that $A$ has a direct summand isomorphic to the infinite cyclic group? V. I. Kuz'minov
	
	*Under the assumption of the continuum hypothesis, not always (E. A. Palyutin, Siberian Math. J., 19 (1978), 1415-1417). Another example, not using the continuum hypothesis, is given by the abelian group $A = {\mathbb{Z}}^{{\aleph }_{1}}/N$ where $N$ is the subgroup consisting of elements having countable support. Every countable subgroup of $A$ is free abelian, which means that each finite subset of $A$ is contained in a pure finitely generated free abelian subgroup. But $A$ has no direct summand isomorphic to $\mathbb{Z}$ . Both properties can be verified using the standard ZFC axioms of set theory. (S. Corson, Letter of 26 August 2024).
\end{openproblem}

\plainproblemsection

\begin{openproblem}
	3.26. (F. Gross). Is it true that finite groups of exponent ${p}^{\alpha }{q}^{\beta }$ have nilpotent length $\leq  \alpha  + \beta$ ? V. D. Mazurov
	
	No, not always (E.I. Khukhro, Algebra and Logic, 17 (1978), 473-482).
\end{openproblem}

\plainproblemsection

\begin{openproblem}
	3.27. (J. G. Thompson). Is every finite simple group with a nilpotent maximal subgroup isomorphic to some ${PS}{L}_{2}\left( q\right)$ ? V. D. Mazurov
	
	Yes, it is (B. Baumann, J. Algebra, 38 (1976), 119-135).
\end{openproblem}

\plainproblemsection

\begin{openproblem}
	3.28. If $G$ is a finite 2-group with cyclic centre and every abelian normal subgroup 2-generated, then is every abelian subgroup of $G$ 3-generated? V. D. Mazurov Not always (Ya. G. Berkovich, Algebra and Logic, 9 (1970), 75).
\end{openproblem}

\plainproblemsection

\begin{openproblem}
	3.29. Under what conditions can a wreath product of matrix groups over a field be represented by matrices over a field? Yu. I. Merzlyakov
	
	Conditions have been found (Yu. E. Vapne, Soviet Math. Dokl., 11 (1970), 1396- 1399).
\end{openproblem}

\plainproblemsection

\begin{openproblem}
	3.30. A torsion-free abelian group is called factor-decomposable if, in all its factor groups, the periodic part is a direct summand. Characterize these groups. A. P. Mishina
	
	This is done (L. Bican, Commentat. Math. Univ. Carolinae, 19 (1978), 653-672).
\end{openproblem}

\plainproblemsection

\begin{openproblem}
	3.31. Find necessary and sufficient conditions under which every pure subgroup of a completely decomposable torsion-free abelian group is itself completely decomposable. A. P. Mishina
	
	These have been found (L. Bican, Czech. Math. J., 24 (1974), 176-191; A. A. Krav-chenko, Vestnik Moskov. Univ. Ser. 1 Mat. Mekh., 1980, no. 3, p. 104 (Russian)).
\end{openproblem}

\plainproblemsection

\begin{openproblem}
	3.33. Are two groups necessarily isomorphic if each of them can be defined by a single relation and is a homomorphic image of the other one? D. I. Moldavanskii No, not necessarily (A. V. Borshchev, D. I. Moldavanskii, Math. Notes, 79, no. 1 (2006), 31-40).
\end{openproblem}

\plainproblemsection

\begin{openproblem}
	3.34. (Well-known problem). The conjugacy problem for groups with a single defining relation. D. I. Moldavanskiĭ
\end{openproblem}

\plainproblemsection

\begin{openproblem}
	3.35. (K. Ross). Suppose that a group $G$ admits two topologies $\sigma$ and $\tau$ yielding locally compact topological groups ${G}_{\sigma }$ and ${G}_{\tau }$ . If the sets of closed subgroups in ${G}_{\sigma }$ and ${G}_{\tau }$ are the same, does it follow that ${G}_{\sigma }$ and ${G}_{\tau }$ are topologically isomorphic? Not always (A. I. Moskalenko, Ukrain. Math. J., 30 (1978), 199-201). Yu. N. Mukhin
\end{openproblem}

\plainproblemsection

\begin{openproblem}
	3.37. Suppose that every finitely generated subgroup of a locally compact group $G$ is pronilpotent. Then is it true that every maximal closed subgroup of $G$ contains the derived subgroup ${G}^{\prime }$ ? Yu. N. Mukhin
	
	Yes, it is (I. V. Protasov, Soviet Math. Dokl., 19 (1978), 1208-1210).
\end{openproblem}

\plainproblemsection

\begin{openproblem}
	3.38. Describe the topological groups which have no proper closed subgroups.
	
	Yu. N. Mukhin
\end{openproblem}

\plainproblemsection

\begin{openproblem}
	3.39. Describe the finite groups with a self-centralizing subgroup of prime order. V. T. Nagrebetskii
	
	They are described. Self-centralizing subgroups of prime order are CC-subgroups. Finite groups with a CC-subgroup were fully classified in (Z. Arad, W. Herfort, Com-mun. in Algebra, 32 (2004), 2087-2098).
\end{openproblem}

\plainproblemsection

\begin{openproblem}
	3.40. (I. R. Shafarevich). Let $S{L}_{2}\left( \mathbb{Z}\right)$ and $S{L}_{2}{\left( \mathbb{Z}\right) }^{ - }$ denote the completions of $S{L}_{2}\left( \mathbb{Z}\right)$ determined by all subgroups of finite index and all congruence subgroups, respectively, and let $\psi  : S{L}_{2}{\left( \mathbb{Z}\right) }^{ \land  } \rightarrow  S{L}_{2}{\left( \mathbb{Z}\right) }^{ - }$ be the natural homomorphism. Is $\operatorname{Ker}\psi$ a free profinite group? V. P. Platonov
	
	Yes, it is (O. V. Mel'nikov, Soviet Math. Dokl., 17 (1976), 867-870).
\end{openproblem}

\plainproblemsection

\begin{openproblem}
	3.41. Is every compact periodic group locally finite? V. P. Platonov
	
	Yes, it is (E. I. Zel'manov, Israel J. Math., 77 (1992), 83-95).
\end{openproblem}

\plainproblemsection

\begin{openproblem}
	3.42. (Kneser-Tits conjecture). Let $G$ be a simply connected $k$ -defined simple algebraic group, and ${E}_{k}\left( G\right)$ the subgroup generated by unipotent $k$ -elements. If ${E}_{k}\left( G\right)  \neq  1$ , then ${G}_{k} = {E}_{k}\left( G\right)$ . The proof is known for $k$ -decomposable groups (C. Chevalley) and for local fields (V. P. Platonov). The conjecture was refuted (V. P. Platonov, Math. USSR-Izv., 10, no. 2 (1976), 211- 243).
\end{openproblem}

\plainproblemsection

\begin{openproblem}
	3.43. Let $\mu$ be an infinite cardinal number. A group $G$ is said to be $\mu$ -overnilpotent if every cyclic subgroup of $G$ is a member of some ascending normal series of length less than $\mu$ reaching $G$ . It is not difficult to show that the class of $\mu$ -overnilpotent groups is a radical class. Is it true that if ${\mu }_{1} < {\mu }_{2}$ for two infinite cardinal numbers ${\mu }_{1}$ and ${\mu }_{2}$ , then there exists a group $G$ which is ${\mu }_{2}$ -overnilpotent and ${\mu }_{1}$ -semisimple?
	
	Remark of 2001: In (S. Vovsi, Sov. Math. Dokl., 13 (1972), 408-410) it was proved that for any two infinite cardinals ${\mu }_{1} < {\mu }_{2}$ there exists a group that is ${\mu }_{2}$ -overnilpotent but not ${\mu }_{1}$ -overnilpotent. B. I. Plotkin
\end{openproblem}

\plainproblemsection

\begin{openproblem}
	3.44. Suppose that a group is generated by its subinvariant soluble subgroups. Is it necessarily locally soluble? B. I. Plotkin
\end{openproblem}

\plainproblemsection

\begin{openproblem}
	3.45. Let $\mathfrak{X}$ be a hereditary radical. Is it true that, in a locally nilpotent torsion-free group $G$ , the subgroup $\mathfrak{X}\left( G\right)$ is isolated? B. I. Plotkin
\end{openproblem}

\plainproblemsection

\begin{openproblem}
	3.46. Does there exist a group having more than one, but finitely many maximal locally soluble normal subgroups? B. I. Plotkin
\end{openproblem}

\plainproblemsection

\begin{openproblem}
	3.47. (Well-known problem of A. I. Maltsev). Is it true that every locally nilpotent group is a homomorphic image of some torsion-free locally nilpotent group?
	
	Editors' comment: An affirmative answer is known for periodic groups (E. M. Le-vich, A. I. Tokarenko, Siberian Math. J., 11 (1970), 1033-1034) and for countable groups (N. S. Romanovskii, Preprint, 1969). B. I. Plotkin
\end{openproblem}

\plainproblemsection

\begin{openproblem}
	3.48. It can be shown that hereditary radicals form a semigroup with respect to taking the products of classes. It is an interesting problem to find all indecomposable elements of this semigroup. In particular, we point out the problem of finding all indecomposable radicals contained in the class of locally finite $p$ -groups.
\end{openproblem}

\plainproblemsection

\begin{openproblem}
	3.49. Does the semigroup generated by all indecomposable radicals satisfy any identity? B. I. Plotkin
\end{openproblem}

\plainproblemsection

\begin{openproblem}
	3.50. Let $G$ be a group of order ${p}^{\alpha } \cdot  m$ , where $p$ is a prime, $p$ and $m$ are coprime, and let $k$ be an algebraically closed field of characteristic $p$ . Is it true that if the indecomposable projective module corresponding to the 1-representation of $G$ has $k$ -dimension ${p}^{\alpha }$ , then $G$ has a Hall ${p}^{\prime }$ -subgroup? The converse is trivially true.
	
	A. I. Saksonov
	
	No, not always: see Example 4.5 in (W. Willems, Math. Z., 171 (1980), 163-174).
\end{openproblem}

\plainproblemsection

\begin{openproblem}
	*3.51. Is it true that every finite group with a group of automorphisms $\Phi$ which acts regularly on the set of conjugacy classes of $G$ (that is, leaves only the identity class fixed) is soluble? The answer is known to be affirmative in the case where $\Phi$ is a cyclic group generated by a regular automorphism. A. I. Saksonov
	
	*No, not always (Y. Fine, J. Group Theory, 22, no. 6 (2019), 1077-1087).
\end{openproblem}

\plainproblemsection

\begin{openproblem}
	3.52. Can the quasivariety generated by the free group of rank 2 be defined by a system of quasi-identities in finitely many variables? D. M. Smirnov
	
	No (A. I. Budkin, Algebra and Logic, 15 (1976), 25-33).
\end{openproblem}

\plainproblemsection

\begin{openproblem}
	3.53. Let $L\left( {\mathfrak{N}}_{4}\right)$ denote the lattice of subvarieties of the variety ${\mathfrak{N}}_{4}$ of nilpotent groups of class at most 4 . Is $L\left( {\mathfrak{N}}_{4}\right)$ distributive? D. M. Smirnov No (Yu. A. Belov, Algebra and Logic, 9 (1970), 371-374).
\end{openproblem}

\plainproblemsection

\begin{openproblem}
	3.55. Is every binary soluble group all of whose abelian subgroups have finite rank, locally soluble? S. P. Strunkov
\end{openproblem}

\plainproblemsection

\begin{openproblem}
	3.56. Is a 2-group with the minimum condition for abelian subgroups locally finite? Yes, it is (V. P. Shunkov, Algebra and Logic, 9 (1970), 291-297). S. P. Strunkov
\end{openproblem}

\plainproblemsection

\begin{openproblem}
	3.57. Determine the laws of distribution of non-soluble and simple group-theoretic numbers in the sequence of natural numbers. See 2.78. P. I. Trofimov
\end{openproblem}

\plainproblemsection

\begin{openproblem}
	3.58. Let $G$ be a compact 0 -dimensional topological group all of whose Sylow $p$ -subgroups are direct products of cyclic groups of order $p$ . Then is every normal subgroup of $G$ complementable? V. S. Charin
	
	Yes, it is (M. I. Kabenyuk, Siberian Math. J., 13 (1972), 654-657).
\end{openproblem}

\plainproblemsection

\begin{openproblem}
	3.59. Let $H$ be an insoluble minimal normal subgroup of a finite group $G$ , and suppose that $H$ has cyclic Sylow $p$ -subgroups for every prime $p$ dividing $\left| {G : H}\right|$ . Prove that $H$ has at least one complement in $G$ . A. Shemetkov
	
	This was proved (S. A. Syskin, Siberian Math. J., 12 (1971), 342-344; L. A. Shemet-kov, Soviet Math. Dokl., 11 (1970), 1436-1438).
\end{openproblem}

\plainproblemsection

\begin{openproblem}
	3.60. The notion of the $p$ -length of an arbitrary finite group was introduced in (L. A. Shemetkov, Math. USSR Sbornik, 1 (1968), 83-92). Investigate the relations between the $p$ -length of a finite group and the invariants ${c}_{p},{d}_{p},{e}_{p}$ of its Sylow $p$ -subgroup. L. A. Shemetkov
\end{openproblem}

\plainproblemsection

\begin{openproblem}
	3.61. Let $\sigma$ be an automorphism of prime order $p$ of a finite group $G$ , which has a Hall $\pi$ -subgroup with cyclic Sylow subgroups. Suppose that $p \in  \pi$ . Does the centralizer ${C}_{G}\left( \sigma \right)$ have at least one Hall $\pi$ -subgroup?
	
	Yes, it does, mod CFSG (V. D. Mazurov, Algebra and Logic, 31, no. 6 (1992), 360- 366).
\end{openproblem}

\plainproblemsection

\begin{openproblem}
	3.62. (Well-known problem). A finite group is said to be a ${D}_{\pi }$ -group if any two of its maximal $\pi$ -subgroups are conjugate. Is an extension of a ${D}_{\pi }$ -group by a ${D}_{\pi }$ -group always a ${D}_{\pi }$ -group? L. A. Shemetkov
	
	Yes, it is mod CFSG (E. P. Vdovin, D. O. Revin, Contemp. Math., 402 (2006), 229- 263).
\end{openproblem}

\plainproblemsection

\begin{openproblem}
	3.64. Describe the finite simple groups with a Sylow 2-subgroup of the following type: $\left\langle  {a, t \mid  {a}^{{2}^{n}} = {t}^{2} = 1,\text{ tat } = {a}^{{2}^{n - 1} - 1}}\right\rangle  , n > 2$ . This was done (J.L. Alperin, R. Brauer, D. Gorenstein, Trans. Amer. Math. Soc., 151 (1970), 1-261).
\end{openproblem}

\plainproblemsection

\begin{openproblem}
	4.1. Find an infinite finitely generated group with an identical relation of the form ${x}^{{2}^{n}} = 1$ . S. I. Adian
	
	It has been found (S.V.Ivanov, Int. J. Algebra Comput., 4, no. 1-2 (1994), 1-308; I. G. Lysënok, Izv. Math., 60, no. 3 (1996), 453-654).
\end{openproblem}

\plainproblemsection

\begin{openproblem}
	4.2. a) Find an infinite finitely generated group of exponent $< {100}$ . b) Do there exist such groups of exponent 5 ? S. I. Adian
\end{openproblem}

\plainproblemsection

\begin{openproblem}
	4.3. Construct a finitely presented group with insoluble word problem and satisfying a non-trivial law. S. I. Adian
	
	This has been done (O. G. Kharlampovich, Math. USSR-Izv., 19 (1982), 151-169).
\end{openproblem}

\plainproblemsection

\begin{openproblem}
	4.4. Construct a finitely presented group with undecidable word problem all of whose non-trivial defining relations have the form ${A}^{2} = 1$ . This problem is interesting for topologists. S. I. Adian
	
	It is constructed (O. A. Sarkisyan, The word problem for some classes of groups and semigroups, Candidate disser., Moscow Univ., 1983 (Russian)).
\end{openproblem}

\plainproblemsection

\begin{openproblem}
	4.5. P. F. Pickel has proved that there are only finitely many non-isomorphic finitely-generated nilpotent groups having the same family of finite homomorphic images. Can Pickel's theorem be extended to polycyclic groups? V. N. Remeslennikov
	
	Yes, it can (F. J. Grunewald, P. F. Pickel, D. Segal, Ann. of Math. (2), 111 (1980), 155-195).
\end{openproblem}

\plainproblemsection

\begin{openproblem}
	4.5. a) (J. Milnor). Is it true that an arbitrary finitely generated group has either polynomial or exponential growth?
	
	c) Is it true that every finitely generated group with undecidable word problem has exponential growth? S. I. Adian
	
	a) No, c) No (R. I. Grigorchuk, Soviet Math. Dokl., 28 (1983), 23-26).
\end{openproblem}

\plainproblemsection

\begin{openproblem}
	4.5. b) Is it true that an arbitrary finitely presented group has either polynomial or exponential growth? S. I. Adian
\end{openproblem}

\plainproblemsection

\begin{openproblem}
	4.6. (P. Hall). Are the projective groups in the variety of metabelian groups free? V. A. Artamonov
\end{openproblem}

\plainproblemsection

\begin{openproblem}
	4.7. (Well-known problem). For which ring epimorphisms $R \rightarrow  Q$ is the corresponding group homomorphism $S{L}_{n}\left( R\right)  \rightarrow  S{L}_{n}\left( Q\right)$ an epimorphism (for a fixed $n \geq  2$ )? In particular, for what rings $R$ does the equality $S{L}_{n}\left( R\right)  = {E}_{n}\left( R\right)$ hold?
	
	V. A. Artamonov
\end{openproblem}

\plainproblemsection

\begin{openproblem}
	4.8. Suppose $G$ is a finitely generated free-by-cyclic group. Is $G$ finitely presented? G. Baumslag
	
	Yes, it is; for infinite cyclic extension this is proved in (M. Feighn, M. Handel, Ann. Math., 149, no. 3 (1999), 1061-1077), while another easier argument applies for finite cyclic extension.
\end{openproblem}

\plainproblemsection

\begin{openproblem}
	4.9. Let $G$ be a finitely generated torsion-free nilpotent group. Are there only a finite number of non-isomorphic groups in the sequence ${\alpha G},{\alpha }^{2}G,\ldots$ ? Here ${\alpha G}$ denotes the automorphism group of $G$ and ${\alpha }^{n + 1}G = \alpha \left( {{\alpha }^{n}G}\right)$ for $n = 1,2,\ldots G$ . Ball G. Baumslag
\end{openproblem}

\plainproblemsection

\begin{openproblem}
	4.10. A group $G$ is called locally indicable if every non-trivial finitely generated subgroup of $G$ has an infinite cyclic factor group. Is every torsion-free one-relator group locally indicable? G. Baumslag
	
	Yes, it is (S. D. Brodskiĭ, Dep. no. 2214-80, VINITI, Moscow, 1980 (Russian)).
\end{openproblem}

\plainproblemsection

\begin{openproblem}
	4.11. Let $F$ be the free group of rank 2 in some variety of groups. If $F$ is not finitely presented, is the multiplicator of $F$ necessarily infinitely generated? G. Baumslag
\end{openproblem}

\plainproblemsection

\begin{openproblem}
	4.12. Let $G$ be a finite group and $A$ a group of automorphisms of $G$ stabilizing a series of subgroups beginning with $G$ and ending with its Frattini subgroup. Then is $A$ nilpotent? Ya. G. Berkovich
	
	Yes, it is (P. Schmid, Math. Ann., 202 (1973), 57-69).
\end{openproblem}

\plainproblemsection

\begin{openproblem}
	4.13. Prove that every finite non-abelian $p$ -group admits an automorphism of order $p$ which is not an inner one. Ya. G. Berkovich
\end{openproblem}

\plainproblemsection

\begin{openproblem}
	4.14. Let $p$ be a prime number. What are necessary and sufficient conditions for a finite group $G$ in order that the group algebra of $G$ over a field of characteristic $p$ be indecomposable as a two-sided ideal? There exist some nontrivial examples, for instance, the group algebra of the Mathieu group ${M}_{24}$ is indecomposable when $p$ is 2 . R. Brauer
	
	A necessary and sufficient condition can be extracted from (G. R. Robinson, J. Algebra, 84 (1983), 493-502); another solution, which requires considering fewer subgroups, can be extracted from (B. Külshammer, Arch. Math. (Basel), 56 (1991), 313-319).
\end{openproblem}

\plainproblemsection

\begin{openproblem}
	4.16. Suppose that $\mathfrak{K}$ is a class of groups meeting the following requirements: 1) subgroups and epimorphic images of $\mathfrak{K}$ -groups are $\mathfrak{K}$ -groups; 2 ) if the group $G = {UV}$ is the product of $\mathfrak{K}$ -subgroups $U$ and $V$ (neither of which need be normal), then $G \in  \mathfrak{K}$ . If $\pi$ is a set of primes, then the class of all finite $\pi$ -groups meets these requirements. Are these the only classes $\mathfrak{K}$ with these properties? R. Baer
	
	No (S. A. Syskin, Siberian Math. J., 20 (1979), 475-476).
\end{openproblem}

\plainproblemsection

\begin{openproblem}
	4.17. If $\mathfrak{V}$ is a variety of groups, denote by $\widetilde{\mathfrak{V}}$ the class of all finite $\mathfrak{V}$ -groups. How to characterize the classes of finite groups of the form $\mathfrak{V}$ for $\mathfrak{V}$ a variety? $B$ Bac
\end{openproblem}

\plainproblemsection

\begin{openproblem}
	4.18. Characterize the classes $\mathfrak{K}$ of groups, meeting the following requirements: subgroups, epimorphic images and groups of automorphisms of $\mathfrak{K}$ -groups are $\mathfrak{K}$ -groups, but not every countable group is a $\mathfrak{K}$ -group. Note that the class of all finite groups and the class of all almost cyclic groups meet these requirements. R. Baer
\end{openproblem}

\plainproblemsection

\begin{openproblem}
	4.19. Denote by $\mathfrak{C}$ the class of all groups $G$ with the following property: if $U$ and $V$ are maximal locally soluble subgroups of $G$ , then $U$ and $V$ are conjugate in $G$ (or at least isomorphic). It is almost obvious that a finite group $G$ belongs to $\mathfrak{C}$ if and only if $G$ is soluble. What can be said about the locally finite groups in $\mathfrak{C}$ ? R. Baer
\end{openproblem}

\plainproblemsection

\begin{openproblem}
	4.20. a) Let $F$ be a free group, and $N$ a normal subgroup of it. Is it true that the Cartesian square of $N$ is $m$ -reducible to $N$ (that is, there is an algorithm that from a pair of words ${w}_{1},{w}_{2} \in  F$ constructs a word $w \in  F$ such that ${w}_{1} \in  N$ and ${w}_{2} \in  N$ if and only if $w \in  N$ )?
	
	b) (Well-known problem). Do there exist finitely presented groups in which the word problem has an arbitrary pre-assigned recursively enumerable $m$ -degree of un-solvability? M. K. Valiev
	
	a) No, it is not true (O. V. Belegradek, Siberian Math. J., 19 (1978), 867-870).
	
	b) No, there exist $m$ -degrees which do not contain the word problem of any recursively presented cancellation semigroup (C. G. Jockusch, jr., Z. Math. Logik Grundlag. Math., 26, no. 1 (1980), 93-95).
\end{openproblem}

\plainproblemsection

\begin{openproblem}
	4.21. Let $G$ be a finite group, $p$ an odd prime number, and $P$ a Sylow $p$ -subgroup of $G$ . Let the order of every non-identity normal subgroup of $G$ be divisible by $p$ . Suppose $P$ has an element $x$ that is conjugate to no other from $P$ . Does $x$ belong to the centre of $G$ ? For $p = 2$ , the answer is positive (G. Glauberman, J. Algebra,4 (1966), 403-420). G. Glauberman
	
	Yes, it does, mod CFSG (O. D. Artemovich, Ukrain. Math. J., 40 (1988), 343-345).
\end{openproblem}

\plainproblemsection

\begin{openproblem}
	4.22. (J. G. Thompson). Let $G$ be a finite group, $A$ a group of automorphisms of $G$ such that $\left| A\right|$ and $\left| G\right|$ are coprime. Does there exist an $A$ -invariant soluble subgroup $H$ of $G$ such that ${C}_{A}\left( H\right)  = 1$ ? G. Glauberman
	
	Yes, it does (S. A. Syskin, Siberian Math. J., 32, no. 6 (1991), 1034-1037).
\end{openproblem}

\plainproblemsection

\begin{openproblem}
	4.23. Let $G$ be a finite simple group, $\tau$ some element of prime order, and $\alpha$ an automorphism of $G$ whose order is coprime to $\left| G\right|$ . Suppose $\alpha$ centralizes ${C}_{G}\left( \tau \right)$ . Is $\alpha  = 1$ ? G. Glauberman
	
	Not always; for example, $G = {Sz}\left( 8\right) ,\left| \tau \right|  = 5,\left| \alpha \right|  = 3$ (N. D. Podufalov, Letter of September, 3, 1975).
\end{openproblem}

\plainproblemsection

\begin{openproblem}
	4.24. Suppose $T$ is a non-abelian Sylow 2-subgroup of a finite simple group $G$ .
	
	a) Suppose $T$ has nilpotency class $n$ . The best possible bound for the exponent of the center of $T$ is ${2}^{n - 1}$ . This easily implies the bound ${2}^{n\left( {n - 1}\right) }$ for the exponent of $T$ , however, this is almost certainly too crude. What is the best possible bound?
	
	d) Find a "small number" of subgroups ${T}_{1},\ldots ,{T}_{n}$ of $T$ which depend only on the isomorphism class of $T$ such that $\left\{  {{N}_{G}\left( {T}_{1}\right) ,\ldots ,{N}_{G}\left( {T}_{n}\right) }\right\}$ together control fusion in $T$ with respect to $G$ (in the sense of Alperin). D. M. Goldschmidt
\end{openproblem}

\plainproblemsection

\begin{openproblem}
	4.24. Suppose that $T$ is a non-abelian Sylow 2-subgroup of a finite simple group $G$ . b) Is it possible that $T$ is the direct product of two proper subgroups?
	
	c) Is ${T}^{\prime } = \Phi \left( T\right)$ ? D. Goldschmidt
	
	b) Yes, it is. For example, the Sylow 2-subgroups of the alternating groups ${\mathbb{A}}_{14}$ and ${\mathbb{A}}_{15}$ are isomorphic to the Sylow 2-subgroup of ${\mathbb{S}}_{4} \times  {\mathbb{S}}_{8}$ . The Sylow 2-subgroups of ${D}_{4}\left( q\right)$ for $q$ odd are also decomposable into direct products (A. S. Kondratiev, Letter of October, 13, 1977).
	
	c) Not always; for example, for $G = {\operatorname{PSL}}_{3}\left( q\right)$ with $q \equiv  1\left( {\;\operatorname{mod}\;4}\right)$ (A. S. Kondratiev).
\end{openproblem}

\plainproblemsection

\begin{openproblem}
	4.27. Describe all finite simple groups $G$ which can be represented in the form $G = {ABA}$ , where $A$ and $B$ are abelian subgroups.
	
	They are described mod CFSG (D. L. Zagorin, L. S. Kazarin, Dokl. Math., 53, no. 2 (1996), 237-239; D. L. Zagorin, Some problems of ABA-factorizations of permutation groups and classical groups (Cand. Diss.), Yaroslavl', 1994 (Russian); D. L. Zagorin, L. S. Kazarin, Problems in algebra, Gomel' Univ., Gomel', no. 11 (1997), 27-41 (Russian); E. P. Vdovin, Algebra and Logic, 38, no. 2 (1999), 67-83).
\end{openproblem}

\plainproblemsection

\begin{openproblem}
	4.28. For a given field $k$ of characteristic $p > 0$ , characterize the locally finite groups with semisimple group algebras over $k$ .
	
	They are characterized: simple groups in (D. S. Passman, A. E. Zalesskii, Proc. London Math. Soc. (3), 67 (1993), 243-276); the general case in (D. S. Passman, Proc. London Math. Soc. (3), 73 (1996), 323-357).
\end{openproblem}

\plainproblemsection

\begin{openproblem}
	4.29. Classify the irreducible matrix groups over a finite field that are generated by reflections, that is, by matrices with Jordan form $\operatorname{diag}\left( {-1,1,\ldots ,1}\right)$ . A. E. Zalesskiĭ These are classified (A. E. Zalesskiĭ, V. N. Serězhkin, Math. USSR-Izv., 17 (1981), 477-503; A. Wagner, Geom. Dedicata, 9 (1980), 239-253; 10 (1981), 191-203, 475- 523).
\end{openproblem}

\plainproblemsection

\begin{openproblem}
	4.30. Describe the groups (finite groups, abelian groups) that are the full automorphism groups of topological groups. M. I. Kargapolov
\end{openproblem}

\plainproblemsection

\begin{openproblem}
	4.31. Describe the lattice of quasivarieties of nilpotent groups of nilpotency class 2 . M. I. Kargapolov
\end{openproblem}

\plainproblemsection

\begin{openproblem}
	4.32. The conjugacy problem for metabelian groups. M. I. Kargapolov
	
	This is algorithmically soluble (G. A. Noskov, Math. Notes, 31 (1982), 252-258).
\end{openproblem}

\plainproblemsection

\begin{openproblem}
	4.33. Let ${\mathfrak{K}}_{n}$ be the class of all groups with a single defining relation in the variety of soluble groups of derived length $n$ .
	
	a) Under what conditions does a ${\mathfrak{K}}_{n}$ -group have non-trivial center? Can a ${\mathfrak{K}}_{n}$ - group, $n \geq  2$ , that cannot be generated by two elements have non-trivial centre? Editors’ comment (1998): These questions were answered for $n = 2$ (E. I. Timoshenko, Siberian Math. J., 14, no. 6 (1973), 954-957; Math. Notes, 64, no. 6 (1998), 798-803).
	
	b) Describe the abelian subgroups of ${\mathfrak{K}}_{n}$ -groups.
	
	c) Investigate the periodic subgroups of ${\mathfrak{K}}_{n}$ -groups. M. I. Kargapolov
\end{openproblem}

\plainproblemsection

\begin{openproblem}
	4.34. Let $v$ be a group word, and let ${\mathfrak{K}}_{v}$ be the class of groups $G$ such that there exists a positive integer $n = n\left( G\right)$ such that each element of the verbal subgroup ${vG}$ can be represented as a product of $n$ values of the word $v$ on the group $G$ .
	
	a) For which $v$ do all finitely generated soluble groups belong to the class ${\mathfrak{K}}_{v}$ ?
	
	b) Does the word $v\left( {x, y}\right)  = {x}^{-1}{y}^{-1}{xy}$ satisfy this condition? M.I. Kargapolov
\end{openproblem}

\plainproblemsection

\begin{openproblem}
	4.35. (Well-known problem). Is there an infinite locally finite simple group satisfying the minimum condition for $p$ -subgroups for every prime $p$ ? O. H. Kegel No (V. V. Belyaev, Algebra and Logic, 20 (1981), 393-401).
\end{openproblem}

\plainproblemsection

\begin{openproblem}
	4.36. Is there an infinite locally finite simple group $G$ with an involution $i$ such that the centralizer ${C}_{G}\left( i\right)$ is a Chernikov group? O. H. Kegel
	
	No (A. O. Asar, Proc. London Math. Soc., 45 (1982), 337-364).
\end{openproblem}

\plainproblemsection

\begin{openproblem}
	4.37. Is there an infinite locally finite simple group $G$ that cannot be represented by matrices over a field and is such that for some prime $p$ the $p$ -subgroups of $G$ are either of bounded derived length or of finite exponent? O. H. Kegel
	
	No (mod CFSG) by Theorem 4.8 of (O. H. Kegel, B. A. F. Wehrfritz, Locally finite groups, North Holland, Amsterdam, 1973).
\end{openproblem}

\plainproblemsection

\begin{openproblem}
	4.38. Classify the composition factors of automorphism groups of finite (nilpotent of class 2) groups of prime exponent. V. D. Mazurov
	
	Any finite simple group can be such a composition factor (P. M. Beletskiĭ, Russian Math. Surveys, 33, no. 6 (1980), 85-86).
\end{openproblem}

\plainproblemsection

\begin{openproblem}
	4.39. A countable group $U$ is said to be ${SQ}$ -universal if every countable group is isomorphic to a subgroup of a quotient group of $U$ . Let $G$ be a group that has a presentation with $r \geq  2$ generators and at most $r - 2$ defining relations. Is $G{S}_{Q}$ - universal? A. M. Macbeath, P. M. Neumann
	
	Yes, it is (B. Baumslag, S. J. Pride, J. London Math. Soc., 17 (1978), 425-426).
\end{openproblem}

\plainproblemsection

\begin{openproblem}
	4.40. Let $C$ be a fixed non-trivial group (for instance, $C = \mathbb{Z}/2\mathbb{Z}$ ). As it is shown in (Yu. I. Merzlyakov, Algebra and Logic, 9, no. 5 (1970), 326-337), for any two groups $A$ and $B$ , all split extensions of $B$ by $A$ can be imbedded in a certain unified way into the direct product $A \times  \operatorname{Aut}\left( {B\wr C}\right)$ . How are they situated in it? Yu. I. Merzlyakov
\end{openproblem}

\plainproblemsection

\begin{openproblem}
	4.41. A point $z$ in the complex plane is called free if the matrices $\left( \begin{array}{ll} 1 & 2 \\  0 & 1 \end{array}\right)$ and $\left( \begin{array}{ll} 1 & 0 \\  z & 1 \end{array}\right)$ generate a free group. Are all points outside the rhombus with vertices $\pm  2$ , $\pm  i$ free? Yu. I. Merzlyakov
	
	No (Yu. A. Ignatov, Math. USSR-Sb., 35 (1979), 49-56; A. I. Shkuratskiĭ, Math. Notes, 24 (1978), 720-721).
\end{openproblem}

\plainproblemsection

\begin{openproblem}
	4.42. For what natural numbers $n$ does the following equality hold:
	
	$$
	G{L}_{n}\left( \mathbb{R}\right)  = {D}_{n}\left( \mathbb{R}\right)  \cdot  {O}_{n}\left( \mathbb{R}\right)  \cdot  U{T}_{n}\left( \mathbb{R}\right)  \cdot  G{L}_{n}\left( \mathbb{Z}\right) ?
	$$
	
	For notation, see, for instance, (M. I. Kargapolov, Ju. I. Merzljakov, Fundamentals of the Theory of Groups, Springer, New York,1979). For given $n$ this equality implies the affirmative solution of Minkowski’s problem on the product of $n$ linear forms (A. M. Macbeath, Proc. Glasgow Math. Assoc., 5, no. 2 (1961), 86-89), which remains open for $n \geq  6$ . It is known that the equality holds for $n \leq  3$ (Kh. N. Narzullayev, Math. Notes, 18 (1975), 713-719); on the other hand, it does not hold for all sufficiently large $n$ (N. S. Akhmedov, Zapiski Nauchn. Seminarov LOMI,67 (1977), 86-107 (Russian)). Yu. I. Merzlyakov
\end{openproblem}

\plainproblemsection

\begin{openproblem}
	4.44. (Well-known problem). Describe the groups whose automorphism groups are abelian. V. T. Nagrebetskii
\end{openproblem}

\plainproblemsection

\begin{openproblem}
	4.45. Let $G$ be a free product amalgamating proper subgroups $H$ and $K$ of $A$ and $B$ , respectively.
	
	a) Suppose that $H, K$ are finite and $\left| {A : H}\right|  > 2,\left| {B : K}\right|  \geq  2$ . Is $G\;{SQ}$ -universal?
	
	b) Suppose that $A, B, H, K$ are free groups of finite ranks. Can $G$ be simple? P. M. Neumann a) Yes, it is (K. I. Lossov, Siberian Math. J., 27, no. 6 (1986), 890-899). b) Yes, it can (M. Burger, S. Mozes, Inst. Hautes Études Sci. Publ. Math., 92 (2000), 151-194). But if one of the indices $\left| {A : H}\right| ,\left| {B : K}\right|$ is infinite, then no, it cannot (S. V. Ivanov, P. E. Schupp, in: Algorithmic problems in groups and semigroups, Int. Conf., Lincoln, NE, 1998, Boston MA, Birkhäuser, 2000, 139-142).
\end{openproblem}

\plainproblemsection

\begin{openproblem}
	4.46. a) We call a variety of groups a limit variety if it cannot be defined by finitely many laws, while each of its proper subvarieties has a finite basis of identities. It follows from Zorn's lemma that every variety that has no finite basis of identities contains a limit subvariety. Give explicitly (by means of identities or by a generating group) at least one limit variety. A. Yu. Ol'shanskii
\end{openproblem}

\plainproblemsection

\begin{openproblem}
	4.46. b) We call a variety of groups a limit variety if it cannot be defined by finitely many laws, while each of its proper subvarieties has a finite basis of identities. It follows from Zorn's lemma that every variety that has no finite basis of identities contains a limit subvariety. Is the set of limit varieties countable? A. Yu. Ol'shanski No, there are continuum of such varieties (P. A. Kozhevnikov, On varieties of groups of large odd exponent, Dep. 1612-V00, VINITI, Moscow, 2000 (Russian); S. V. Ivanov, A. M. Storozhev, Contemp. Math., 360 (2004), 55-62).
\end{openproblem}

\plainproblemsection

\begin{openproblem}
	4.47. Does there exist a countable family of groups such that every variety is generated by a subfamily of it? A. Yu. Ol'shanskiĭ
	
	No (Yu. G. Kleiman, Math. USSR-Izv., 22 (1984), 33-65).
\end{openproblem}

\plainproblemsection

\begin{openproblem}
	4.48. A locally finite group is said to be an $A$ -group if all of its Sylow subgroups are abelian. Does every variety of $A$ -groups possess a finite basis of identities?
	
	A. Yu. Ol'shanskiĭ
\end{openproblem}

\plainproblemsection

\begin{openproblem}
	4.49. Let $G$ and $H$ be finitely generated torsion-free nilpotent groups such that Aut $G \cong$ Aut $H$ . Does it follow that $G \cong  H$ ? V. N. Remeslennikov
	
	No (G. A. Noskov, V. A. Roman'kov, Algebra and Logic, 13 (1974), 306-311).
\end{openproblem}

\plainproblemsection

\begin{openproblem}
	4.50. What are the soluble varieties of groups all of whose finitely generated subgroups are residually finite? V. N. Remeslennikov
\end{openproblem}

\plainproblemsection

\begin{openproblem}
	4.51. (Well-known problem). Are knot groups residually finite? V. N. Remeslennikov Yes, they are (W. P. Thurston, Bull. Amer. Math. Soc., 6 (1982), 357-381).
\end{openproblem}

\plainproblemsection

\begin{openproblem}
	4.52. Let $G$ be a finitely generated torsion-free group that is an extension of an abelian group by a nilpotent group. Then is $G$ almost a residually finite $p$ -group for almost all primes $p$ ? V. N. Remeslennikov
	
	Yes, it is (G. A. Noskov, Algebra and Logic, 13 (1974), 388-393; D. Segal, J. Algebra, 32 (1974), 389-399).
\end{openproblem}

\plainproblemsection

\begin{openproblem}
	4.54. Are any two minimal relation-modules of a finite group isomorphic? Not always (P. A. Linnell, Ph. D. Thesis, Cambridge, 1979; K. W. Roggenkamp A. J. Sieradski, M. N. Dyer, J. Pure Appl. Algebra, 15 (1979), 199-217). 4.57. Let a group $G$ be the product of two of its abelian minimax subgroups $A$ and $B$ . Prove or refute the following statements:
	
	a) ${A}_{0}{B}_{0} \neq  1$ , where ${A}_{0} = \mathop{\bigcap }\limits_{{x \in  G}}{A}^{x}$ and similarly for ${B}_{0}$ ;
	
	b) the derived subgroup of $G$ is a minimax subgroup. N. F. Sesekin
	
	Both parts were proved (D. I. Zaitsev, Algebra and Logic, 19 (1980), 94-106).
\end{openproblem}

\plainproblemsection

\begin{openproblem}
	*4.55. Let $G$ be a finite group and ${\mathbb{Z}}_{p}$ the localization at $p$ . Does the Krull-Schmidt theorem hold for projective ${\mathbb{Z}}_{p}G$ -modules? K. W. Roggenkamp
	
	*No, not always (S. M. Woods, Canadian J. Math., 26, no. 1 (1974), 121-129).
\end{openproblem}

\plainproblemsection

\begin{openproblem}
	4.56. Let $R$ be a commutative Noetherian ring with 1, and $\Lambda$ an $R$ -algebra, which is finitely generated as $R$ -module. Put $T = \{ U \in  \operatorname{Mod}\Lambda  \mid  \exists$ an exact $\Lambda$ -sequence $0 \rightarrow$ $P \rightarrow  {\Lambda }^{\left( n\right) } \rightarrow  U \rightarrow  0$ for some $n$ , with ${P}_{m} \cong  {\Lambda }_{m}^{\left( n\right) }$ for every maximal ideal $m$ of $R\}$ . Denote by $\mathbb{G}\left( T\right)$ the Grothendieck group of $T$ relative to short exact sequences.
	
	a) Describe $\mathbb{G}\left( T\right)$ , in particular, what does it mean: $\left\lbrack  U\right\rbrack   = \left\lbrack  V\right\rbrack$ in $\mathbb{G}\left( T\right)$ ?
	
	b) Conjecture: if $\dim \left( {\max \left( R\right) }\right)  = d < \infty$ , and there are two epimorphisms $\varphi$ : ${\Lambda }^{\left( n\right) } \rightarrow  U,\psi  : {\Lambda }^{\left( n\right) } \rightarrow  V, n > d$ , and $\left\lbrack  U\right\rbrack   = \left\lbrack  V\right\rbrack$ in $\mathbb{G}\left( T\right)$ , then $\operatorname{Ker}\varphi  = \operatorname{Ker}\psi$ . K. W. Roggenkamp
\end{openproblem}

\plainproblemsection

\begin{openproblem}
	4.58. Let a finite group $G$ be the product of two subgroups $A$ and $B$ , where $A$ is abelian and $B$ is nilpotent. Find the dependence of the derived length of $G$ on the nilpotency class of $B$ and the order of its derived subgroup. N. F. Sesekin
	
	This was found (D. I. Zaitsev, Math. Notes, 33 (1983), 414-419).
\end{openproblem}

\plainproblemsection

\begin{openproblem}
	4.59. (P. Hall). Find the smallest positive integer $n$ such that every countable group can be embedded in a simple group with $n$ generators. D. M. Smirnov It is proved that $n = 2$ (A. P. Goryushkin, Math. Notes,16 (1974),725-727).
\end{openproblem}

\plainproblemsection

\begin{openproblem}
	4.60. (P. Hall). What is the cardinality of the set of simple groups generated by two elements, one of order 2 and the other of order 3? D. M. Smirnov
	
	The cardinality of the continuum. Every 2-generated group $G$ is embeddable into a simple group $H$ with two generators of orders 2 and 3 (P.E. Schupp, J. London Math. Soc., 13, no. 1 (1976), 90-94). Such a group $H$ has at most countably many 2-generator subgroups, while there are continuum of groups $G$ . (Yu. I. Merzlyakov, 1976.)
\end{openproblem}

\plainproblemsection

\begin{openproblem}
	4.61. Does there exist a linear function $f$ with the following property: if every abelian subgroup of a finite 2-group $G$ is generated by $n$ elements, then $G$ is generated by $f\left( n\right)$ elements? S. A. Syskin
	
	No (A. Yu. Ol'shanskiĭ, Math. Notes, 23 (1978), 183-185).
\end{openproblem}

\plainproblemsection

\begin{openproblem}
	4.62. Does there exist a finitely based variety of groups whose universal theory is undecidable? A. Tarski
	
	Yes, there does; for example, ${\mathfrak{A}}^{5}$ . Indeed, it is shown in (V. N. Remeslennikov, Algebra and Logic, 12, no. 5 (1975), 327-346) that there exists a finitely presented group $G = \left\langle  {{x}_{1},\ldots ,{x}_{n} \mid  {w}_{1},\ldots ,{w}_{m},{\;\operatorname{mod}\;{\mathfrak{A}}^{5}}}\right\rangle$ in ${\mathfrak{A}}^{5}$ with undecidable word problem. We put ${\Phi }_{w} = \left( {\forall {x}_{1},\ldots ,{x}_{n}}\right) \left( {\left( {{w}_{1} = 1}\right)  \land  \ldots  \land  \left( {{w}_{m} = 1}\right)  \rightarrow  \left( {w = 1}\right) }\right)$ , where $w$ runs over all the words in ${x}_{1},\ldots ,{x}_{n}$ . Clearly, there is no algorithm to decide whether a formula ${\Phi }_{w}$ is true in ${\mathfrak{A}}^{5}$ . (V. N. Remeslennikov,1976.)
\end{openproblem}

\plainproblemsection

\begin{openproblem}
	4.63. Does there exist a non-abelian variety of groups (in particular, one that contains the variety of all abelian groups) whose elementary theory is decidable? A. Tarski No (A. P. Zamyatin, Algebra and Logic, 27 (1978), 13-17).
\end{openproblem}

\plainproblemsection

\begin{openproblem}
	4.64. Does there exist a variety of groups that does not admit an independent system of defining identities? A. Tarski
	
	Yes, there does (Yu. G. Kleiman, Math. USSR-Izv., 22 (1984), 33-65).
\end{openproblem}

\plainproblemsection

\begin{openproblem}
	4.65. Conjecture: $\frac{{p}^{q} - 1}{p - 1}$ never divides $\frac{{q}^{p} - 1}{q - 1}$ if $p, q$ are distinct primes. The validity of this conjecture would simplify the proof of solvability of groups of odd order (W. Feit, J. G. Thompson, Pacific J. Math., 13, no. 3 (1963), 775-1029), rendering unnecessary the detailed use of generators and relations. J. G. Thompson
\end{openproblem}

\plainproblemsection

\begin{openproblem}
	4.66. Let $P$ be a presentation of a finite group $G$ on ${m}_{p}$ generators and ${r}_{p}$ relations. The deficiency $\operatorname{def}\left( G\right)$ is the maximum of ${m}_{p} - {r}_{p}$ over all presentations $P$ . Let $G$ be a finite group such that $G = {G}^{\prime } \neq  1$ and the multiplicator $M\left( G\right)  = 1$ . Prove that $\operatorname{def}\left( {G}^{n}\right)  \rightarrow   - \infty$ as $n \rightarrow  \infty$ , where ${G}^{n}$ is the $n$ th direct power of $G$ . J. Wiegold
\end{openproblem}

\plainproblemsection

\begin{openproblem}
	4.67. Let $G$ be a finite $p$ -group. Show that the rank of the multiplicator $M\left( G\right)$ of $G$ is bounded in terms of the rank of $G$ . J. Wiegold
	
	This was done (A. Lubotzky, A. Mann, J. Algebra, 105 (1987), 484-505).
\end{openproblem}

\plainproblemsection

\begin{openproblem}
	4.68. Construct a finitely generated (infinite) characteristically simple group that is not a direct power of a simple group. J. Wiegold
	
	This was done (J. S. Wilson, Math. Proc. Cambridge Phil. Soc., 80 (1976), 19-35).
\end{openproblem}

\plainproblemsection

\begin{openproblem}
	4.69. Let $G$ be a finite $p$ -group, and suppose that $\left| {G}^{\prime }\right|  > {p}^{n\left( {n - 1}\right) /2}$ for some nonnegative integer $n$ . Prove that $G$ is generated by the elements of breadths $\geq  n$ . The breadth of an element $x$ of $G$ is $b\left( x\right)$ where $\left| {G : {C}_{G}\left( x\right) }\right|  = {p}^{b\left( x\right) }$ . J. Wiegold This has been proved (A. Skutin, J. Algebra, 526 (2019), 1-5).
\end{openproblem}

\plainproblemsection

\begin{openproblem}
	4.70. Let $k$ be a field of characteristic different from 2, and ${G}_{k}$ the group of transformations $A = \left( {a,\alpha }\right)  : x \rightarrow  {ax} + \alpha \;\left( {a,\alpha  \in  k, a \neq  0}\right)$ . Extend ${G}_{k}$ to the projective plane by adjoining the symbols $\left( {0,\alpha }\right)$ and a line at infinity. Then the lines are just the centralizers ${C}_{G}\left( A\right)$ of elements $A \in  {G}_{k}$ and their cosets. Do there exist other groups $G$ complementable to the projective plane such that the lines are just the cosets of the centralizers of elements of $G$ ? H. Schwerdtfeger No (E. A. Kuznetsov, Dep. no. 7028-V89, VINITI, Moscow, 1989 (Russian)).
\end{openproblem}

\plainproblemsection

\begin{openproblem}
	4.71. Let $A$ be a group of automorphisms of a finite group $G$ which has a series of $A$ -invariant subgroups $G = {G}_{0} > \cdots  > {G}_{k} = 1$ such that every $\left| {{G}_{i} : {G}_{i + 1}}\right|$ is prime. Prove that $A$ is supersoluble. L. A. Shemetkov
	
	This was proved (L. A. Shemetkov, Math. USSR-Sb., 23 (1974), 593-611).
\end{openproblem}

\plainproblemsection

\begin{openproblem}
	4.72. Is it true that every variety of groups whose free groups are residually nilpotent torsion-free, is either soluble or coincides with the variety of all groups? For an affirmative answer, it is sufficient to show that every variety of Lie algebras over the field of rational numbers which does not contain any finite-dimensional simple algebras is soluble. A. L. Shmel'kin
\end{openproblem}

\plainproblemsection

\begin{openproblem}
	4.73. (Well-known problem). Does there exist a non-abelian variety of groups
	
	a) all of whose finite groups are abelian?
	
	b) all of whose periodic groups are abelian? A. L. Shmel'kin
	
	a) Yes, there exists (A. Yu. Ol'shanskii, Math. USSR-Sb., 54 (1986), 57-80).
	
	b) Yes, there exists (P. A. Kozhevnikov, On varieties of groups of large odd exponent, Dep. 1612-V00, VINITI, Moscow, 2000 (Russian); S. V. Ivanov, A. M. Storozhev, Contemp. Math., 360 (2004), 55-62).
\end{openproblem}

\plainproblemsection

\begin{openproblem}
	4.74. a) Is every 2-group of order greater than 2 non-simple? V.P. Shunkov
	
	No, the existence of simple infinite 2-groups follows from the existence of finitely generated infinite groups of exponent ${2}^{n}$ (S. V. Ivanov, Int. J. Algebra Comput.,4, no. 1-2 (1994), 1-308; I. G. Lysënok, Izv. Math., 60, no. 3 (1996), 453-654) and from the positive solution to the Restricted Burnside Problem for groups of exponent ${2}^{n}$ (E. I. Zel'manov, Math. USSR-Sb., 72 (1992), 543-565).
\end{openproblem}

\plainproblemsection

\begin{openproblem}
	4.74. b) Is every binary-finite 2-group of order greater than 2 non-simple?
	
	V. P. Shunkov
\end{openproblem}

\plainproblemsection

\begin{openproblem}
	4.75. Let $G$ be a periodic group containing an involution $i$ and suppose that the Sylow 2-subgroups of $G$ are either locally cyclic or generalized quaternion. Does the element $i{O}_{{2}^{\prime }}\left( G\right)$ of the factor-group $G/{O}_{{2}^{\prime }}\left( G\right)$ always lie in its centre?
	
	V. P. Shunkov
\end{openproblem}

\plainproblemsection

\begin{openproblem}
	4.76. Let $G$ be a locally finite group containing an element $a$ of prime order such that the centralizer ${C}_{G}\left( a\right)$ is finite. Is $G$ almost soluble? Yes, it is. P. Fong (Osaka J. Math., 13 (1976), 483-489) has shown (mod CFSG) that $G$ is almost locally soluble. Given this, B. Hartley and T. Meixner (Arch. Math., 36 (1981),211-213) have proved that $G$ is almost locally nilpotent. Therefore $G$ is almost soluble by (J.L. Alperin, Proc. Amer. Math. Soc., 13 (1962), 175-180) and (G. Higman, J. London Math. Soc., 32 (1957), 321-334). Moreover, by (E. I. Khukhro, Math. Notes, 38, no. 5-6 (1985), 867-870; E. I. Khukhro, Math. USSR-Sb., 71, no. 1 (1992), ${51} - {63}$ ) then $G$ is almost nilpotent.
\end{openproblem}

\plainproblemsection

\begin{openproblem}
	4.77. In 1972, A. Rudvalis discovered a new simple group $R$ of order ${2}^{14} \cdot  {3}^{3} \cdot  {5}^{3} \cdot  7 \cdot  {13} \cdot  {29}$ . He has shown that $R$ possesses an involution $i$ such that ${C}_{R}\left( i\right)  = V \times  F$ , where $V$ is a 4-group (an elementary abelian group of order 4) and $F \cong  {Sz}\left( 8\right)$ .
	
	a) Show that $R$ is the only finite simple group $G$ that possesses an involution $i$ such that ${C}_{G}\left( i\right)  = V \times  F$ , where $V$ is a 4-group and $F \cong  {Sz}\left( 8\right)$ .
	
	b) Let $G$ be a non-abelian finite simple group that possesses an involution $i$ such that ${C}_{G}\left( i\right)  = V \times  F$ , where $V$ is an elementary abelian 2-group of order ${2}^{n}, n \geq  1$ , and $F \cong  {Sz}\left( {2}^{m}\right) , m \geq  3$ . Show that $n = 2$ and $m = 3$ .
	
	a) This has been shown (V. D. Mazurov, Math. Notes, 31 (1982), 165-173).
	
	b) This follows from the CFSG.
\end{openproblem}

\plainproblemsection

\begin{openproblem}
	5.1. a) Is every locally finite minimal non- ${FC}$ -group non-simple?
	
	V. V. Belyaev, N. F. Sesekin
	
	Yes, it is (M. Kuzucuogīu, R. E. Phillips, Proc. Cambridge Philos. Soc., 105, no. 3 (1989), 417-420).
\end{openproblem}

\plainproblemsection

\begin{openproblem}
	5.1. b) Is every locally finite minimal non- ${FC}$ -group distinct from its derived subgroup? The question has an affirmative answer for minimal non- ${BFC}$ -groups.
	
	V. V. Belyaev, N. F. Sesekin
\end{openproblem}

\plainproblemsection

\begin{openproblem}
	5.2. Is it true that the growth function $f$ of any infinite finitely generated group satisfies the inequality $f\left( n\right)  \leq  \left( {f\left( {n - 1}\right)  + f\left( {n + 1}\right) }\right) /2$ for all sufficiently large $n$ (for a fixed finite system of generators)? No (P. de la Harpe, R. I. Grigorchuk, Algebra and Logic, 37, no. 6 (1998), 353-356).
\end{openproblem}

\plainproblemsection

\begin{openproblem}
	5.3. (Well-known problem). Can every finite lattice $L$ be embedded in the lattice of subgroups of a finite group? G. M. Bergman
	
	Yet, it can (P. Pudlák, J. Tuma, Algebra Universalis, 10 (1980), 74-95).
\end{openproblem}

\plainproblemsection

\begin{openproblem}
	5.4. Let $g$ and $h$ be positive elements of a linearly ordered group $G$ . Can one always embed $G$ in a linearly ordered group $\bar{G}$ in such a way that $g$ and $h$ are conjugate in $\bar{G}$ ? V. V. Bludov
	
	No, not always (V.V. Bludov, Algebra and Logic, 44, no. 6 (2005), 370-380).
\end{openproblem}

\plainproblemsection

\begin{openproblem}
	5.5. If $G$ is a finitely generated abelian-by-polycyclic-by-finite group, does there exist a finitely generated metabelian group $M$ such that $G$ is isomorphic to a subgroup of the automorphism group of $M$ ? If so, many of the tricky properties of $G$ like its residual finiteness would become transparent. B. A. F. Wehrfritz
\end{openproblem}

\plainproblemsection

\begin{openproblem}
	5.6. If $R$ is a finitely generated integral domain of characteristic $p > 0$ , does the profinite (ideal) topology of $R$ induce the profinite topology on the group of units of $R$ ? It does for $p = 0$ . B. A. F. Wehrfritz
	
	Yes, it does (D. Segal, Bull. London Math. Soc., 11 (1979), 186-190).
\end{openproblem}

\plainproblemsection

\begin{openproblem}
	5.7. An algebraic variety $X$ over a field $k$ is called rational if the field of functions $k\left( X\right)$ is purely transcendental over $k$ , and it is called stably rational if $k\left( X\right)$ becomes purely transcendental after adjoining finitely many independent variables. Let $T$ be a stably rational torus over a field $k$ . Is $T$ rational? Other formulations of this question and some related results see in (V. E. Voskresenskiĭ, Russ. Math. Surveys, 28, no. 4 (1973), 79-105). V. E. Voskresenskiĭ
	
	Yes, it is (V. E. Voskresenskii, J. Math. Sci. (N. Y.), 161, no. 1 (2009), 176-180).
\end{openproblem}

\plainproblemsection

\begin{openproblem}
	5.8. Let $G$ be a torsion-free soluble group of finite cohomological dimension $\operatorname{cd}G$ . If ${hG}$ denotes the Hirsch number of $G$ , then it is known that ${hG} \leq  \operatorname{cd}G \leq  {hG} + 1$ . Find a purely group-theoretic criterion for ${hG} = \operatorname{cd}G$ . K. W. Gruenberg
	
	This has been found (P. H. Kropholler, J. Pure Appl. Algebra, 43 (1986), 281-287).
\end{openproblem}

\plainproblemsection

\begin{openproblem}
	5.9. Let $1 \rightarrow  {R}_{i}\xrightarrow[]{{\pi }_{i}}F \rightarrow  G \rightarrow  1, i = 1,2$ , be two exact sequences of groups with $G$ finite and $F$ free of finite rank $d\left( F\right)$ . If we assume that $d\left( F\right)  = d\left( G\right)  + 1$ (where $d\left( G\right)$ is the minimum number of generators of $G$ ), are the corresponding abelianized extensions isomorphic? $K.W.{Gruenberg}$
	
	Yes, they are (P. A. Linnell, J. Pure Appl. Algebra, 22 (1981), 143-166).
\end{openproblem}

\plainproblemsection

\begin{openproblem}
	5.10. Let $E = H * K$ and denote the augmentation ideals of the groups $E, H, K$ by $\mathfrak{e},\mathfrak{h},\mathfrak{k}$ , respectively. If $I$ is a right ideal in $\mathbb{Z}E$ , let ${d}_{E}\left( I\right)$ denote the minimum number of generators of $I$ as a right ideal. Assuming $H$ and $K$ finitely generated, is it true that ${d}_{E}\left( \mathfrak{e}\right)  = {d}_{E}\left( {\mathfrak{h}E}\right)  + {d}_{E}\left( {\mathfrak{k}E}\right)$ ? Here $\mathfrak{h}E$ is, of course, the right ideal generated by $\mathfrak{h}$ ; similarly for $\mathfrak{k}E$ . K. W. Gruenberg
	
	Not always (P. A. Linnell, unpublished). A simple example: $H$ is elementary abelian group of order $4, K$ is elementary abelian of order 9 . Here ${d}_{E}\left( \mathfrak{c}\right)  = 3$ , although ${d}_{E}\left( {\mathfrak{h}E}\right)  = {d}_{E}\left( {\mathfrak{k}E}\right)  = 2$ . (K. W. Gruenberg, Letter of July,27,1995.)
\end{openproblem}

\plainproblemsection

\begin{openproblem}
	5.11. Let $G$ be a finite group and suppose that there exists a non-empty proper subset $\pi$ of the set of all primes dividing $\left| G\right|$ such that the centralizer of every non-trivial $\pi$ -element is a $\pi$ -subgroup. Does it follow that $G$ contains a subgroup $U$ such that ${U}^{g} \cap  U = 1$ or $U$ for every $g \in  G$ , and the centralizer of every non-trivial element of $U$ is contained in $U$ ? $K.W.{Gruenberg}$
	
	Yes, it does (mod CFSG) (J. S. Williams, J. Algebra, 69 (1981), 487-513).
\end{openproblem}

\plainproblemsection

\begin{openproblem}
	5.12. Let $G$ be a finite group with trivial soluble radical in which there are Sylow 2-subgroups having non-trivial intersection. Suppose that, for any two Sylow 2-subgroups $P$ and $Q$ of $G$ with $P \cap  Q \neq  1$ , the index $\left| {P : P \cap  Q}\right|$ does not exceed ${2}^{n}$ . Is it then true that $\left| P\right|  \leq  {2}^{2n}$ ? V. V. Kabanov
	
	Yes, it is, mod CFSG (V.I. Zenkov, Algebra and Logic, 36, no. 2 (1997), 93-98).
\end{openproblem}

\plainproblemsection

\begin{openproblem}
	5.13. Suppose that $K$ and $L$ are distinct conjugacy classes of involutions in a finite group $G$ and $\langle x, y\rangle$ is a 2-group for all $x \in  K$ and $y \in  L$ . Does it follow that $G \neq  \left\lbrack  {K, L}\right\rbrack$ ? V. V. Kabanov
	
	Not always; for example, $G = S{p}_{4}\left( {2}^{n}\right) , n \geq  2, K, L$ being classes of involutions that have non-trivial intersections with the centres of ${N}_{G}{\left( M\right) }^{\prime }$ and ${N}_{G}{\left( N\right) }^{\prime }$ , respectively, where $M, N$ are distinct elementary abelian subgroups of order 8 in a Sylow 2-subgroup of $G$ and the dash means taking the derived subgroup (A. A. Makhnëv, Letter of October, 10, 1981).
\end{openproblem}

\plainproblemsection

\begin{openproblem}
	5.14. An intersection of some Sylow 2-subgroups is called a Sylow intersection and an intersection of a pair of Sylow 2-subgroups is called a paired Sylow intersection.
	
	a) Describe the finite groups all of whose 2-local subgroups have odd indices.
	
	b) Describe the finite groups all of whose normalizers of Sylow intersections have odd indices.
	
	c) Describe the finite groups all of whose normalizers of paired Sylow intersections have odd indices.
	
	d) Describe the finite groups in which for any two Sylow 2-subgroups $P$ and $Q$ , the intersection $P \cap  Q$ is normal in some Sylow 2-subgroup of $\langle P, Q\rangle$ .
	
	V. V. Kabanov, A. A. Makhnëv, A. I. Starostin
\end{openproblem}

\plainproblemsection

\begin{openproblem}
	5.15. Do there exist finitely presented residually finite groups with recursive, but not primitive recursive, solution of the word problem? F. B. Cannonito
\end{openproblem}

\plainproblemsection

\begin{openproblem}
	5.16. Is every countable locally linear group embeddable in a finitely presented group? In (G. Baumslag, F. B. Cannonito, C. F. Miller III, Math. Z., 153 (1977), 117-134) it is proved that every countable group which is locally linear of bounded degree can be embedded in a finitely presented group with solvable word problem.
	
	F. B. Cannonito, C. Miller
\end{openproblem}

\plainproblemsection

\begin{openproblem}
	5.17. If the finite group $G$ has the form $G = {AB}$ where $A$ and $B$ are nilpotent of classes $\alpha$ and $\beta$ , respectively, then $G$ is soluble. Is ${G}^{\left( \alpha  + \beta \right) } = 1$ ? One can show that ${G}^{\left( \alpha  + \beta \right) }$ is nilpotent (E. Pennington,1973). One can ask the same question for infinite groups (or Lie algebras), but there is nothing known beyond Ito's theorem: ${A}^{\prime } = {B}^{\prime } = 1$ implies ${G}^{\left( 2\right) } = 1$ . O. H. Kegel
	
	Not always (J. Cossey, S. Stonehewer, Bull. London Math. Soc., 30, no. 144 (1998), 247-250). See also new problem 14.43.
\end{openproblem}

\plainproblemsection

\begin{openproblem}
	5.18. Let $G$ be an infinite locally finite simple group. Is the centralizer of every element of $G$ infinite? O. H. Kegel
	
	Yes, it is (B. Hartley, M. Kuzucuoglu, Proc. London Math. Soc. (3), 62, no. 2 (1991), 301-324).
\end{openproblem}

\plainproblemsection

\begin{openproblem}
	5.19. a) Let $G$ be an infinite locally finite simple group satisfying the minimum condition for 2-subgroups. Is $G = {PS}{L}_{2}\left( F\right) , F$ some locally finite field of odd characteristic, if the centralizer of every involution of $G$ is almost locally soluble?
	
	b) Can one characterize the simple locally finite groups with the min-2 condition containing a maximal radical non-trivial 2-subgroup of rank $\leq  2$ as linear groups of small rank? O. H. Kegel
	
	a) Yes, it is (N. S. Chernikov, Ukrain. Math. J., 35 (1983), 230-231).
	
	b) Yes, one can (mod CFSG) by Theorem 4.8 of (O. H. Kegel, B. A. F. Wehrfritz,
	
	Locally finite groups, North Holland, Amsterdam, 1973).
\end{openproblem}

\plainproblemsection

\begin{openproblem}
	5.20. Is the elementary theory of lattices of $l$ -ideals of lattice-ordered abelian groups decidable? A. I. Kokorin
	
	No, it is undecidable (N. Ya. Medvedev, Algebra and Logic, 44, no. 5 (2005), 302-312).
\end{openproblem}

\plainproblemsection

\begin{openproblem}
	5.21. Can every torsion-free group with solvable word problem be embedded in a group with solvable conjugacy problem? An example due to A. Macintyre shows that this question has a negative answer when the condition of torsion-freeness is omitted. D. J. Collins
	
	No, not every (A. Darbinyan, Invent. Math., 224 (2021), 987-997).
\end{openproblem}

\plainproblemsection

\begin{openproblem}
	5.22. Does there exist a version of the Higman embedding theorem in which the degree of unsolvability of the conjugacy problem is preserved? D. J. Collins Yes, it exists (A. Yu. Ol'shanskii, M. V. Sapir, The conjugacy problem and Higman embeddings (Memoirs AMS, 804), 2004, 133 p.).
\end{openproblem}

\plainproblemsection

\begin{openproblem}
	5.23. Is it true that a free lattice-ordered group of the variety of lattice-ordered groups defined by the law ${x}^{-1}\left| y\right| x \ll  {\left| y\right| }^{2}$ (or, equivalently, by the law $\left| \left\lbrack  {x, y}\right\rbrack  \right|  \ll  \left| x\right|$ ), is residually linearly ordered nilpotent? V. M. Kopytov
	
	No, it is not (V. V. Bludov, A. M. W. Glass, Trans. Amer. Math. Soc., 358 (2006), 5179-5192).
\end{openproblem}

\plainproblemsection

\begin{openproblem}
	5.24. Is it true that a free lattice-ordered group of the variety of the lattice-ordered groups which are residually linearly ordered, is residually soluble linearly ordered? V. M. Kopytov
	
	No, it is not (N. Ya. Medvedev, Algebra and Logic, 44, no. 3 (2005), 197-204).
\end{openproblem}

\plainproblemsection

\begin{openproblem}
	5.25. Prove that the factor-group of any soluble linearly ordered group by its derived subgroup is non-periodic. V. M. Kopytov
\end{openproblem}

\plainproblemsection

\begin{openproblem}
	5.26. Let $G$ be a finite $p$ -group with the minimal number of generators $d$ , and let ${r}_{1}$ (respectively, ${r}_{2}$ ) be the minimal number of defining relations on $d$ generators in the sense of representing $G$ as a factor-group of a free discrete group (pro- $p$ -group). It is well known that always ${r}_{2} > {d}^{2}/4$ . For each prime number $p$ denote by $c\left( p\right)$ the exact upper bound for the numbers $b\left( p\right)$ with the property that ${r}_{2} \geq  b\left( p\right) {d}^{2}$ for all finite $p$ -groups.
	
	a) It is obvious that ${r}_{1} \geq  {r}_{2}$ . Find a $p$ -group with ${r}_{1} > {r}_{2}$ .
	
	b) Conjecture: $\mathop{\lim }\limits_{{p \rightarrow  \infty }}c\left( p\right)  = 1/4$ . It is proved (J. Wisliceny, Math. Nachr.,102 (1981), ${57} - {78}$ ) that $\mathop{\lim }\limits_{{d \rightarrow  \infty }}{r}_{2}/{d}^{2} = 1/4$ .H. Koch
\end{openproblem}

\plainproblemsection

\begin{openproblem}
	5.27. Prove that if $G$ is a torsion-free pro- $p$ -group with a single defining relation, then $\operatorname{cd}G = 2$ . This has been proved for a large class of groups in (J. P. Labute, Inv. Math., 4, no. 2 (1967), 142-158). H. Koch
\end{openproblem}

\plainproblemsection

\begin{openproblem}
	5.28. Let $G$ be a group and $H$ a torsion-free subgroup of $G$ such that the augmentation ideal ${I}_{G}$ of the integral group-ring $\mathbb{Z}G$ can be decomposed as ${I}_{G} = {I}_{H}\mathbb{Z}G \oplus  M$ for some $\mathbb{Z}G$ -submodule $M$ . Prove that $G$ is a free product of the form $G = H * K$ . D. E. Cohen
	
	This has been proved (W. Dicks, M. J. Dunwoody, Groups acting on graphs, Cambridge Univ. Press, Cambridge-New York, 1989).
\end{openproblem}

\plainproblemsection

\begin{openproblem}
	5.29. Consider a group $G$ given by a presentation with $m$ generators and $n$ defining relations, where $m \geq  n$ . Do some $m - n$ of the given generators generate a free subgroup of $G$ ? R. C. Lyndon
	
	Yes, they do (N. S. Romanovskii, Algebra and Logic, 16 (1977), 62-68).
\end{openproblem}

\plainproblemsection

\begin{openproblem}
	5.30. (Well-known problem). Suppose that $G$ is a finite soluble group, $A \leq  \operatorname{Aut}G$ , ${C}_{G}\left( A\right)  = 1$ , the orders of $G$ and $A$ are coprime, and let $\left| A\right|$ be the product of $n$ not necessarily distinct prime numbers. Is the nilpotent length of $G$ bounded above by $n$ ? This is proved for large classes of groups (E. Shult, F. Gross, T. Berger, A. Turull), and there is a bound in terms of $n$ if $A$ is soluble (J. G. Thompson’s $\leq  {5}^{n}$ , H. Kurzweil’s $\leq  {4n}$ , A. Turull’s $\leq  {2n}$ ), but the problem remains open. V. D. Mazurov
\end{openproblem}

\plainproblemsection

\begin{openproblem}
	5.32. Let $p$ be a prime, $C$ a conjugacy class of $p$ -elements of a finite group $G$ and suppose that for any two elements $x$ and $y$ of $C$ the product $x{y}^{-1}$ is a $p$ -element. Is the subgroup generated by the class $C$ a $p$ -group? V. D. Mazurov
	
	Yes, it is (for $p > 2$ : W. Xiao, Sci. in China,34, no. 9 (1991),1025-1031; for $p = 2$ : V. I. Zenkov, in: Algebra. Proc. IIIrd Int. Conf. on Algebra, Krasnoyarsk, 1993, Berlin, De Gruyter, 1996, 297-302).
\end{openproblem}

\plainproblemsection

\begin{openproblem}
	5.33. (Y. Ihara). Consider the quaternion algebra $Q$ with norm $f = {x}^{2} - \tau {y}^{2} -$ $\rho {z}^{2} + {\rho \tau }{u}^{2},\rho ,\tau  \in  \mathbb{Z}$ . Assume that $f$ is indefinite and of $\mathbb{Q}$ -rank 0, i. e. $f = 0$ for $x, y, z, u \in  \mathbb{Q}$ implies $x = y = z = u = 0$ . Consider $Q$ as the algebra of the matrices
	
	$$
	X = \left( \begin{matrix} x + \sqrt{\tau }y & \rho \left( {z + \sqrt{\tau }u}\right) \\  z - \sqrt{\tau }u & x - \sqrt{\tau }y \end{matrix}\right)
	$$
	
	with $x, y, z, u \in  \mathbb{Q}$ . Let $p$ be a prime, $p \nmid  {\rho \tau }$ . Consider the group $G$ of all $X$ with $x, y, z, u \in  {\mathbb{Z}}^{\left( p\right) }$ , det $X = 1$ , where ${\mathbb{Z}}^{\left( p\right) } = \left\{  {m/{p}^{t} \mid  m, t \in  \mathbb{Z}}\right\}$ .
	
	Conjecture: $G$ has the congruence subgroup property, i. e. every non-central normal subgroup $N$ of $G$ contains a full congruence subgroup $N\left( \mathfrak{A}\right)  = \{ X \in  G \mid  X \equiv$ $E\left( {\;\operatorname{mod}\;\mathfrak{A}}\right) \}$ for some $\mathfrak{A}$ . Notice that the congruence subgroup property is independent of the matrix representation of $Q$ . J. Mennicke
\end{openproblem}

\plainproblemsection

\begin{openproblem}
	5.34. Let $\mathfrak{o}$ be a commutative ring with identity in which 2 is invertible and which is not generated by zero divisors. Do there exist non-standard automorphisms of $G{L}_{n}\left( \mathfrak{o}\right)$ for $n \geq  3$ ? Yu. I. Merzlyakov
	
	No (V. Ya. Bloshchitsyn, Algebra and Logic, 17 (1978), 415-417).
\end{openproblem}

\plainproblemsection

\begin{openproblem}
	5.35. Let $V$ be a vector space of dimension $n$ over a field. A subgroup $G$ of $G{L}_{n}\left( V\right)$ is said to be rich in transvections if $n \geq  2$ and for every hyperplane $H \subseteq  V$ and every line $L \subseteq  H$ there is at least one transvection in $G$ with residual line $L$ and fixed space $H$ . Describe the automorphisms of the subgroups of $G{L}_{2}\left( V\right)$ which are rich in transvections. Yu. I. Merzlyakov
\end{openproblem}

\plainproblemsection

\begin{openproblem}
	5.36. What profinite groups satisfy the maximum condition for closed subgroups?
	
	Yu. N. Mukhin
\end{openproblem}

\plainproblemsection

\begin{openproblem}
	5.38. Is it the case that if $A, B$ are finitely generated soluble Hopfian groups then $A \times  B$ is Hopfian? Peter M. Neumann
\end{openproblem}

\plainproblemsection

\begin{openproblem}
	5.39. Prove that every countable group can act faithfully as a group of automorphisms of a finitely generated soluble group (of derived length at most 4). For background to this problem, in particular its relationship with problem 8.50, see (Peter M. Neumann, in: Groups-Korea, Pusan, 1988 (Lect. Notes Math., 1398), Springer, Berlin, 1989, 124-139). Peter M. Neumann
\end{openproblem}

\plainproblemsection

\begin{openproblem}
	5.40. Let $G$ be a countable group acting on a set $\Omega$ . Suppose that $G$ is $k$ -fold transitive for every finite $k$ , and $G$ contains no non-trivial permutations of finite support. Is it true that $\Omega$ can be identified with the rational line $\mathbb{Q}$ in such a way that $G$ becomes a group of autohomeomorphisms? P. M. Neumann
	
	Not always (A. H. Mekler, J. London Math. Soc., 33 (1986), 49-58).
\end{openproblem}

\plainproblemsection

\begin{openproblem}
	5.41. Does every non-trivial finite group, which is free in some variety, contain a non-trivial abelian normal subgroup?
	
	Yes, it does (mod CFSG). Otherwise, if $G$ is a counterexample, the centralizer of the product $E$ of all minimal normal subgroups of $G$ is trivial and there is an element of $G/E$ whose order is $m$ , where $m$ is the maximum of the orders of 2-elements of $G$ . By Theorem 1 in (M. Aschbacher, P. B. Kleidman, M. W. Liebeck, Math. Z., 208, no. 3 (1991),401-409), which is proved using CFSG, there is an element of order ${2m}$ in $G$ , which contradicts the choice of $m$ . (S. A. Syskin, Letter of August,20,1992.)
\end{openproblem}

\plainproblemsection

\begin{openproblem}
	5.42. Does the free group of rank 2 have an infinite ascending chain of verbal subgroups each being generated as a verbal subgroup by a single element? A. Yu. Ol'shanskiĭ
\end{openproblem}

\plainproblemsection

\begin{openproblem}
	5.43. Does there exist a soluble variety of groups that is not generated by its finite groups? A. Yu. Ol'shanskii
	
	Yes, there exists (Yu. G. Kleiman, Trans. Moscow Math. Soc., 1983, no. 2, 63-110).
\end{openproblem}

\plainproblemsection

\begin{openproblem}
	5.44. A union $\mathfrak{A} = \bigcup {\mathfrak{V}}_{\alpha }$ of varieties of groups (in the lattice of varieties) is called irreducible if $\mathop{\bigcup }\limits_{{\beta  \neq  \alpha }}{\mathfrak{V}}_{\beta } \neq  \mathfrak{A}$ for each index $\alpha$ . Is every variety an irreducible union of (finitely or infinitely many) varieties each of which cannot be decomposed into a union of two proper subvarieties? A. Yu. Ol'shanskii
\end{openproblem}

\plainproblemsection

\begin{openproblem}
	5.46. Is every recursively presented soluble group embeddable in a group finitely presented in the variety of all soluble groups of derived length $n$ , for suitable $n$ ?
	
	Yes, it is (G. P. Kukin, Soviet Math. Dokl., 21 (1980), 378-381). V. N. Remeslennikov
\end{openproblem}

\plainproblemsection

\begin{openproblem}
	*5.47. Is every countable abelian group embeddable in the center of some finitely presented group? V. N. Remeslennikov
	
	*Yes, it is (A. O. Houcine, J. Algebra, 307, no. 1 (2007), 1-23).
\end{openproblem}

\plainproblemsection

\begin{openproblem}
	5.48. Suppose that $G$ and $H$ are finitely generated residually finite groups having the same set of finite homomorphic images. Are $G$ and $H$ isomorphic if one of them is a free (free soluble) group? V. N. Remeslennikov
\end{openproblem}

\plainproblemsection

\begin{openproblem}
	5.49. Let ${A}_{m, n}$ be the group of all automorphisms of the free soluble group of derived length $n$ and rank $m$ . Is it true that
	
	a) ${A}_{m, n}$ is finitely generated for any $m$ and $n$ ?
	
	b) every automorphism in ${A}_{m, n}$ is induced by some automorphism in ${A}_{m, n + 1}$ ? V. N. Remeslennikov
	
	a) No: ${A}_{3,2}$ is not finitely generated, although ${A}_{m,2}$ is whenever $m \neq  3$ (S. Bachmuth, H. Y. Mochizuki, Trans. Amer. Math. Soc., 270 (1982), 693-700, 292 (1985), 81- 101).
	
	b) No (V. Shpilrain, Int. J. Algebra and Comput., 1, no. 2 (1991), 177-184).
\end{openproblem}

\plainproblemsection

\begin{openproblem}
	5.50. Is there a finite group whose set of quasi-laws does not have an independent basis? D. M. Smirnov
	
	Yes, there is (A. N. Fedorov, Siberian Math. J., 21 (1980), 840-850).
\end{openproblem}

\plainproblemsection

\begin{openproblem}
	5.51. Does a non-abelian free group have an independent basis for its quasi-laws? Yes, it does (A. I. Budkin, Math. Notes, 31 (1982), 413-417). D. M. Smirnov
\end{openproblem}

\plainproblemsection

\begin{openproblem}
	5.52. It is not hard to show that a finite perfect group is the normal closure of a single element. Is the same true for infinite finitely generated groups? J. Wiegold
\end{openproblem}

\plainproblemsection

\begin{openproblem}
	5.53. (P. Scott). Let $p, q, r$ be distinct prime numbers. Prove that the free product $G = {C}_{p} * {C}_{q} * {C}_{r}$ of cyclics of orders $p, q, r$ is not the normal closure of a single element. By Lemma 3.1 of (J. Wiegold, J. Austral. Math. Soc., 17, no. 2 (1974), 133-141), every soluble image, and every finite image, of $G$ is the normal closure of a single element. J. Wiegold
	
	This is proved in (J. Howie, J. Pure Appl. Algebra, 173, no. 2 (2002), 167-176).
\end{openproblem}

\plainproblemsection

\begin{openproblem}
	5.54. Let $p, q, r$ be distinct primes and $u\left( {x, y, z}\right)$ a commutator word in three variables. Prove that there exist (infinitely many?) natural numbers $n$ such that the alternating group ${\mathbb{A}}_{n}$ can be generated by three elements $\xi ,\eta ,\zeta$ satisfying ${\xi }^{p} = {\eta }^{q} =$ ${\zeta }^{r} = {\xi \eta \zeta } \cdot  u\left( {\xi ,\eta ,\zeta }\right)  = 1$ J. Wiegold
\end{openproblem}

\plainproblemsection

\begin{openproblem}
	5.55. Find a finite $p$ -group that cannot be embedded in a finite $p$ -group with trivial multiplicator. Notice that every finite group can be embedded in a group with trivial multiplicator. J. Wiegold
\end{openproblem}

\plainproblemsection

\begin{openproblem}
	5.56. a) Let $p$ be a prime greater than 3 . Is it true that every finite group of exponent $p$ can be embedded in the commutator subgroup of a finite group of exponent $p$ ? J. Wiegold
\end{openproblem}

\plainproblemsection

\begin{openproblem}
	5.56. b) Does there exist a locally nilpotent group of prime exponent that coincides with its derived subgroup (and hence has no maximal subgroups)? $J$ . Yes, there exists. Every non-soluble variety $\mathfrak{V}$ contains a non-trivial group coinciding with the derived subgroup, the direct limit of the spectrum $F\overset{\varphi }{ \rightarrow  }F\overset{\varphi }{ \rightarrow  }\ldots$ , where $F$ is the free group in $\mathfrak{V}$ on the free generators ${x}_{1},{x}_{2},\ldots$ and $\varphi$ is the homomorphism given by ${x}_{i} \rightarrow  \left\lbrack  {{x}_{{2i} - 1},{x}_{2i}}\right\rbrack$ . As shown in (Yu. P. Razmyslov, Algebra and Logic,10 (1971), 21-29) the Kostrikin variety of locally nilpotent groups of prime exponent $p \geq  5$ is unsoluble. (E. I. Khukhro, I. V. L'vov, Letter of June, 19, 1976.) The same was also proved in (Yu. A. Kolmakov, Math. Notes, 35, no. 5-6 (1984), 389-391). An example answering the question was also produced in (M. R. Vaughan-Lee, J. Wiegold, Bull. London Math. Soc., 13, no. 1 (1981), 45-46).
\end{openproblem}

\plainproblemsection

\begin{openproblem}
	5.59. Let $G$ be a locally finite group which is the product of a $p$ -subgroup and a $q$ -subgroup, where $p$ and $q$ are distinct primes. Is $G$ a $\{ p, q\}$ -group? B. Hartley
\end{openproblem}

\plainproblemsection

\begin{openproblem}
	5.60. Is an arbitrary soluble group that satisfies the minimum condition for normal subgroups countable? B. Hartley
	
	Not always (B. Hartley, Proc. London Math. Soc., 33 (1977), 55-75).
\end{openproblem}

\plainproblemsection

\begin{openproblem}
	5.61. (Well-known problem). Does an arbitrary uncountable locally finite group have only one end? B. Hartley
	
	Yes, it does (D. Holt, Bull. London Math. Soc., 13 (1981), 557-560).
\end{openproblem}

\plainproblemsection

\begin{openproblem}
	5.62. Is a group locally finite if it contains infinite abelian subgroups and all of them are complementable? S. N. Chernikov
	
	Not always (N. S. Chernikov, Math. Notes, 28 (1980), 788-792).
\end{openproblem}

\plainproblemsection

\begin{openproblem}
	5.63. Prove that a finite group is not simple if it contains two non-identity elements whose centralizers have coprime indices. S. A. Chunikhin
	
	This has been proved mod CFSG (L. S. Kazarin, in: Studies in Group Theory, Sverdlovsk, 1984, 81-99 (Russian)).
\end{openproblem}

\plainproblemsection

\begin{openproblem}
	5.64. Suppose that a finite group $G$ is the product of two subgroups ${A}_{1}$ and ${A}_{2}$ . Prove that if ${A}_{i}$ contains a nilpotent subgroup of index $\leq  2$ for $i = 1,2$ , then $G$ is soluble. L. A. Shemetkov
	
	This has been proved (L. S. Kazarin, Math. USSR-Sb., 38 (1981), 47-59).
\end{openproblem}

\plainproblemsection

\begin{openproblem}
	5.65. Is the class of all finite groups that have Hall $\pi$ -subgroups closed under taking finite subdirect products? L. A. Shemetkov
	
	Yes, it is (D. O. Revin, E. P. Vdovin, J. Group Theory, 14, no. 1 (2011), 93-101).
\end{openproblem}

\plainproblemsection

\begin{openproblem}
	5.67. Is a periodic residually finite group finite if it satisfies the weak minimum condition for subgroups? V. P. Shunkov
\end{openproblem}

\plainproblemsection

\begin{openproblem}
	5.68. Let $G$ be a finitely-presented group, and assume that $G$ has polynomial growth in the sense of Milnor. Show that $G$ has soluble word problem. P. E. Schup This has been shown (M. Gromov, Publ. Math. IHES, 53 (1981), 53-73).
\end{openproblem}

\plainproblemsection

\begin{openproblem}
	5.69. Is every lattice isomorphism between torsion-free groups having no non-trivial cyclic normal subgroups induced by a group isomorphism? $B.V$ . Yakovlev No (A. Yu. Ol'shanskii, C. R. Acad. Bulgare Sci., 32 (1979), 1165-1166).
\end{openproblem}

\plainproblemsection

\begin{openproblem}
	6.1. A subgroup $H$ of an arbitrary group $G$ is said to be $C$ -closed if $H = {C}^{2}\left( H\right)  =$ $C\left( {C\left( H\right) }\right)$ , and weakly $C$ -closed if ${C}^{2}\left( x\right)  \leq  H$ for any element $x$ of $H$ . The structure of the finite groups all of whose proper subgroups are $C$ -closed was studied by Gaschütz. Describe the structure of the locally finite groups all of whose proper subgroups are weakly $C$ -closed. V. A. Antonov, N. F. Sesekin
\end{openproblem}

\plainproblemsection

\begin{openproblem}
	6.2. The totality of $C$ -closed subgroups of an arbitrary group $G$ is a complete lattice with respect to the operations $A \land  B = A \cap  B, A \vee  B = C\left( {C\left( A\right)  \cap  C\left( B\right) }\right)$ . Describe the groups whose lattice of $C$ -closed subgroups is a sublattice of the lattice of all subgroups. V. A. Antonov, N. F. Sesekin
\end{openproblem}

\plainproblemsection

\begin{openproblem}
	6.3. A group $G$ is called of type ${\left( FP\right) }_{\infty }$ if the trivial $G$ -module $\mathbb{Z}$ has a resolution by finitely generated projective $G$ -modules. The class of all groups of type ${\left( FP\right) }_{\infty }$ has a couple of excellent closure properties with respect to extensions and amalgamated products (R. Bieri, Homological dimension of discrete groups, Queen Mary College Math. Notes, London,1976). Is every periodic group of type ${\left( FP\right) }_{\infty }$ finite? This is related to the question whether there is an infinite periodic group with a finite presentation. R. Bieri
\end{openproblem}

\plainproblemsection

\begin{openproblem}
	6.4. A group $G$ is called of type ${\left( FP\right) }_{\infty }$ if the trivial $G$ -module $\mathbb{Z}$ has a resolution by finitely generated projective $G$ -modules. Is it true that every torsion-free group of type ${\left( FP\right) }_{\infty }$ has finite cohomological dimension?
	
	No (K. S. Brown, R. Geoghegan, Invent. Math., 77 (1984), 367-381).
\end{openproblem}

\plainproblemsection

\begin{openproblem}
	6.5. Is it the case that every soluble group $G$ of type ${\left( FP\right) }_{\infty }$ is constructible in the sense of (G. Baumslag, R. Bieri, Math. Z., 151, no. 3 (1976), 249-257)? In other words, can $G$ be built up from the trivial group in terms of finite extensions and ${HNN}$ -extensions? R. Bieri
\end{openproblem}

\plainproblemsection

\begin{openproblem}
	6.6. Let $P, Q$ be permutation representations of a finite group $G$ with the same character. Suppose $P\left( G\right)$ is a primitive permutation group. Is $Q\left( G\right)$ necessarily primitive? The answer is known to be affirmative if $G$ is soluble. H. Wielandt No, not necessarily (mod CFSG) (R. M. Guralnick, J. Saxl, in: Groups, Combinatorics, Geometry, Proc. Durham, 1990, Cambridge Univ. Press, 1992, 364-367).
\end{openproblem}

\plainproblemsection

\begin{openproblem}
	6.7. Suppose that $P$ is a finite 2-group. Does there exist a characteristic subgroup $L\left( P\right)$ of $P$ such that $L\left( P\right)$ is normal in $H$ for every finite group $H$ that satisfies the following conditions: 1) $P$ is a Sylow 2-subgroup of $H,2)H$ is ${\mathbb{S}}_{4}$ -free, and 3) ${C}_{H}\left( {{O}_{2}\left( H\right) }\right)  \leq  {O}_{2}\left( H\right)$ ? G. Glauberman
	
	Yes, there exists (B. Stellmacher, Israel J. Math., 94 (1996), 367-379).
\end{openproblem}

\plainproblemsection

\begin{openproblem}
	6.8. Can a residually finite locally normal group be embedded in a Cartesian product of finite groups in such a way that each element of the group has at most finitely many central projections? Yu. M. Gorchakov
	
	Yes, it can (M. J. Tomkinson, Bull. London Math. Soc., 13 (1981), 133-137).
\end{openproblem}

\plainproblemsection

\begin{openproblem}
	6.9. Is the derived subgroup of any locally normal group decomposable into a product of at most countable subgroups which commute elementwise? Yu. M. Gorchakov
\end{openproblem}

\plainproblemsection

\begin{openproblem}
	6.10. Let $p$ be a prime and $n$ be an integer with $p > {2n} + 1$ . Let $x, y$ be $p$ -elements in $G{L}_{n}\left( \mathbb{C}\right)$ . If the subgroup $\langle x, y\rangle$ is finite, then it is an abelian $p$ -group (W. Feit, J. G. Thompson, Pacif. J. Math., 11, no. 4 (1961), 1257-1262). What can be said about $\langle x, y\rangle$ in the case it is an infinite group? J. D. Dixon
\end{openproblem}

\plainproblemsection

\begin{openproblem}
	6.11. Let $\mathfrak{L}$ be the class of locally compact groups with no small subgroups (see D. Montgomery, L. Zippin, Topological transformation groups, New York, 1955; V. M. Glushkov, Uspekhi Matem. Nauk, 12, no. 2 (1957), 3-41 (Russian)). Study extensions of groups in this class with the objective of giving a direct proof of the following: For each $G \in  \mathfrak{L}$ there exists an $H \in  \mathfrak{L}$ and a (continuous) homomorphism $\vartheta  : G \rightarrow  H$ with a discrete kernel and an image $\operatorname{Im}\vartheta$ satisfying $\operatorname{Im}\vartheta  \cap  Z\left( H\right)  = 1$ $\left( {Z\left( H\right) \text{is the center of}H}\right)$ . This result follows from the Gleason-Montgomery-Zippin solution of Hilbert’s 5th Problem (since $\mathfrak{L}$ is the class of finite-dimensional Lie groups and the latter are locally linear). On the other hand a direct proof of this result would give a substantially shorter proof of the 5th Problem since the adjoint representation of $H$ is faithful on $\operatorname{Im}\vartheta$ . J. D. Dixon
\end{openproblem}

\plainproblemsection

\begin{openproblem}
	6.12. a) Is a metabelian group countable if it satisfies the weak minimum condition for normal subgroups?
	
	b) Is such a group minimax if it is torsion-free?D. I. Zaitsev
	
	Yes, it is, in both cases (D. I. Zaitsev, L. A. Kurdachenko, A. V. Tushev, Algebra and Logic, 24 (1985), 412-436).
\end{openproblem}

\plainproblemsection

\begin{openproblem}
	6.13. Is it true that if a non-abelian Sylow 2-subgroup of a finite group $G$ has a non-trivial abelian direct factor, then $G$ is not simple? A. S. Kondratiev Yes, it is, mod CFSG (V. V. Kabanov, A. S. Kondratiev, Sylow 2-subgroups of finite groups (a survey), Inst. Math. Mech. UNC AN SSSR, Sverdlovsk, 1979 (Russian)).
\end{openproblem}

\plainproblemsection

\begin{openproblem}
	6.14. Are the following lattices locally finite: the lattice of all locally finite varieties of groups? the lattice of all varieties of groups? No, they are not (M. I. Anokhin, Izv. Math., 63, no. 4 (1999), 649-665).
\end{openproblem}

\plainproblemsection

\begin{openproblem}
	6.15. A variety is said to be pro-locally-finite if it is not locally finite while all of its proper subvarieties are locally finite. An example - the variety of abelian groups. How many pro-locally-finite varieties of groups are there? $\;A.V.{Kuznetsov}$ There are continuum of such varieties (P. A. Kozhevnikov, On varieties of groups of large odd exponent, Dep. 1612-V00, VINITI, Moscow, 2000 (Russian)).
\end{openproblem}

\plainproblemsection

\begin{openproblem}
	6.16. A variety is called sparse if it has at most countably many subvarieties. How many sparse varieties of groups are there?
	
	There are continuum of such varieties (P. A. Kozhevnikov, On varieties of groups of large odd exponent, Dep. 1612-V00, VINITI, Moscow, 2000 (Russian); S. V. Ivanov, A. M. Storozhev, Contemp. Math., 360 (2004), 55-62).
\end{openproblem}

\plainproblemsection

\begin{openproblem}
	6.17. Is every variety of groups generated by its finitely generated groups that have soluble word problem? A. V. Kuznetsov
	
	No (Yu. G. Kleiman, Math. USSR-Izv., 22 (1984), 33-65).
\end{openproblem}

\plainproblemsection

\begin{openproblem}
	6.18. (Well-known problem). Suppose that a class $\mathfrak{K}$ of 2-generator groups generates the variety of all groups. Is a non-cyclic free group residually in $\mathfrak{K}$ ? V. M. Levchuk Not always (S. J. Pride, Math. Z., 131 (1973), 245-248).
\end{openproblem}

\plainproblemsection

\begin{openproblem}
	6.19. Let $R$ be a nilpotent associative ring. Are the following two statements for a subgroup $H$ of the adjoint group of $R$ always equivalent: 1) $H$ is a normal subgroup; 2) $H$ is an ideal of the groupoid $R$ with respect to Lie multiplication? $V.M.{Levchuk}$ Not always. Let $R$ be the free nilpotent of index 3 associative algebra over ${\mathbb{F}}_{2}$ on the free generators $x, y$ . Let $S$ be the subalgebra generated by the elements $\left\lbrack  {x,{y}^{2}}\right\rbrack   = x * {y}^{2}$ , $\left\lbrack  {x,{xy}}\right\rbrack   = x * \left( {xy}\right)  = x\left( {x * y}\right) ,\left\lbrack  {x,{yx}}\right\rbrack   = x * \left( {yx}\right)  = \left( {x * y}\right) x$ , where $\left\lbrack  {a, b}\right\rbrack$ denotes the commutator in the adjoint group with multiplication $a \circ  b = a + b + {ab}$ , and $a * b = {ab} - {ba}$ is Lie multiplication. Let $M, N$ be the subalgebras generated by $S$ and the elements $x * y$ and $\left\lbrack  {x, y}\right\rbrack$ , respectively. The minimal subgroup of $\left( {R, \circ  }\right)$ that contains $x$ and is an ideal of the groupoid $\left( {R, * }\right)$ equals $\langle x\rangle  \circ  M = \langle x\rangle  + M$ , where $\langle x\rangle  = \left\{  {0, x,{x}^{2}, x + {x}^{2} + {x}^{3}}\right\}$ , while the minimal normal subgroup containing $x$ equals $\langle x\rangle  \circ  N = \langle x\rangle  + N$ . Neither is contained in the other. (E.I. Khukhro, Letter of July, 23, 1979.)
\end{openproblem}

\plainproblemsection

\begin{openproblem}
	6.20. Does there exist a supersoluble group of odd order, all of whose automorphisms are inner? V. D. Mazurov
	
	Yes, there does (B. Hartley, D. J. S. Robinson, Arch. Math., 35 (1980), 67-74).
\end{openproblem}

\plainproblemsection

\begin{openproblem}
	6.21. G. Higman proved that, for any prime number $p$ , there exists a natural number $\chi \left( p\right)$ such that the nilpotency class of any finite group $G$ having an automorphism of order $p$ without non-trivial fixed points does not exceed $\chi \left( p\right)$ . At the same time he showed that $\chi \left( p\right)  \geq  \left( {{p}^{2} - 1}\right) /4$ for any such Higman’s function $\chi$ . Find the best possible Higman’s function. Is it the function defined by equalities $\chi \left( p\right)  = \left( {{p}^{2} - 1}\right) /4$ for $p > 2$ and $\chi \left( 2\right)  = 1$ ? This is known to be true for $p \leq  7$ . V. D. Mazurov 6.24. The membership problem for the braid group on four strings. It is known (T. A. Makanina, Math. Notes, 29, no. 1 (1981), 16-17) that the membership problem is undecidable for braid groups with more than four strings. G. S. Makanin 6.26. Let $D$ be a normal set of involutions in a finite group $G$ and let $\Gamma \left( D\right)$ be the graph with vertex set $D$ and edge set $\{ \left( {a, b}\right)  \mid  a, b \in  D,{ab} = {ba} \neq  1\}$ . Describe the finite groups $G$ with non-connected graph $\Gamma \left( D\right)$ . A. A. Makhnëv
\end{openproblem}

\plainproblemsection

\begin{openproblem}
	6.22. Construct a braid that belongs to the derived subgroup of the braid group but is not a commutator. G. S. Makanin
	
	Let $x = {\sigma }_{1}, y = {\sigma }_{2}$ be the standard generators of the braid group ${B}_{3}$ ; then the braid ${\left( xyxyx\right) }^{12}{\left( xy\right) }^{-{12}}{\left( xyx\right) }^{-{12}}$ belongs to the derived subgroup of ${B}_{3}$ but is not a commutator (Yu. S. Semënov, Abstracts of the 10th All-Union Sympos. on Group Theory, Minsk, 1986, p. 207 (Russian)).
\end{openproblem}

\plainproblemsection

\begin{openproblem}
	6.23. A braid $K$ of the braid group ${\mathfrak{B}}_{n + 1}$ is said to be smooth if removing any of the threads in $K$ transforms $K$ into a braid that is equal to 1 in ${\mathfrak{B}}_{n}$ . It is known that smooth braids form a free subgroup. Describe generators of this subgroup. G. S. Makanin
	
	They are described in (D. L. Johnson, Math. Proc. Camb. Philos. Soc., 92 (1982), 425-427).
\end{openproblem}

\plainproblemsection

\begin{openproblem}
	6.25. (Well-known problem). Find an algorithm for calculating the rank of coefficient-free equations in a free group. The rank of an equation is the maximal rank of the free subgroup generated by a solution of this equation. G. S. Makanin This was found (A. A. Razborov, Math. USSR-Izv., 25 (1984), 115-162).
\end{openproblem}

\plainproblemsection

\begin{openproblem}
	6.28. Let $A$ be an elementary abelian 2-group which is a ${TI}$ -subgroup of a finite group $G$ . Investigate the structure of $G$ under the hypothesis that the weak closure of $A$ in a Sylow 2-subgroup of $G$ is abelian. A. A. Makhnëv
\end{openproblem}

\plainproblemsection

\begin{openproblem}
	6.29. Suppose that a finite group $A$ is isomorphic to the group of all topological automorphisms of a locally compact group $G$ . Does there always exist a discrete group whose automorphism group is isomorphic to $A$ ? This is true if $A$ is cyclic; the condition of local compactness of $G$ is essential (R. J. Wille, Indag. Math.,25, no. 2 (1963), 218-224). O. V. Mel'nikov
\end{openproblem}

\plainproblemsection

\begin{openproblem}
	6.30. Let $G$ be a residually finite Hopfian group, and let $\widehat{G}$ be its profinite completion. Is $\widehat{G}$ necessarily Hopfian (in topological sense)? O. V. Mel'nikov
\end{openproblem}

\plainproblemsection

\begin{openproblem}
	6.31. a) Suppose that $G$ is a finitely-generated residually-finite group, $d\left( G\right)$ the minimal number of generators of $G$ , and $\delta \left( G\right)$ the minimal number of topological generators of the profinite completion of $G$ . Is $\delta \left( G\right)  = d\left( G\right)$ always true? Not always (G. A. Noskov, Math. Notes, 33 (1983), 249-254). O. V. Mel'nikov
\end{openproblem}

\plainproblemsection

\begin{openproblem}
	6.31. b) Let $G$ be a finitely generated residually finite group, $d\left( G\right)$ the minimal number of its generators, and $\delta \left( G\right)$ the minimal number of topological generators of the profinite completion of $G$ . It is known that there exist groups $G$ for which $d\left( G\right)  > \delta \left( G\right)$ (G. A. Noskov, Math. Notes,33, no. 4 (1983),249-254). Is the function $d$ bounded on the set of groups $G$ with the fixed value of $\delta \left( G\right)  \geq  2$ ? O. V. Mel’nikov
\end{openproblem}

\plainproblemsection

\begin{openproblem}
	6.32. Let ${F}_{n}$ be the free profinite group of finite rank $n > 1$ . Is it true that for each normal subgroup $N$ of the free profinite group of countable rank, there exists a normal subgroup of ${F}_{n}$ isomorphic to $N$ ? O. V. Mel'nikov
\end{openproblem}

\plainproblemsection

\begin{openproblem}
	6.34. Let $\mathfrak{o}$ be an associative ring with identity. A system of its ideals $\mathfrak{A} = \left\{  {{\mathfrak{A}}_{ij} \mid  i, j \in  }\right.$ $\mathbb{Z}\}$ is called a carpet of ideals if ${\mathfrak{A}}_{ik}{\mathfrak{A}}_{kj} \subseteq  {\mathfrak{A}}_{ij}$ for all $i, j, k \in  \mathbb{Z}$ . If $\mathfrak{o}$ is commutative, then the set ${\Gamma }_{n}\left( A\right)  = \left\{  {x \in  S{L}_{n}\left( \mathfrak{o}\right)  \mid  {x}_{ij} \equiv  {\delta }_{ij}\left( {\;\operatorname{mod}\;{\mathfrak{A}}_{ij}}\right) }\right\}$ is a group, the (special) congruenz-subgroup modulo the carpet $\mathfrak{A}$ (the "carpet subgroup"). Under quite general conditions, it was proved in (Yu. I. Merzlyakov, Algebra i Logika, 3, no. 4 (1964), 49-59 (Russian); see also M. I. Kargapolov, Yu. I. Merzlyakov, Fundamentals of the Theory of Groups,3rd Ed., Moscow, Nauka,1982, p. 145 (Russian)) that in the groups $G{L}_{n}$ and $S{L}_{n}$ the mutual commutator subgroup of the congruenz-subgroups modulo a carpet of ideals shifted by $k$ and $l$ steps is again the congruenz-subgroup modulo the same carpet shifted by $k + l$ steps. Prove analogous theorems a) for orthogonal groups; b) for unitary groups. Yu. I. Merzlyakov
	
	These are proved (V. M. Levchuk, Sov. Math. Dokl., 42, no. 1 (1991), 82-86; Ukrain. Math. J., 44, no. 6 (1992), 710-718).
\end{openproblem}

\plainproblemsection

\begin{openproblem}
	6.35. (R. Bieri, R. Strebel). Let $\mathfrak{o}$ be an associative ring with identity distinct from zero. A group $G$ is said to be almost finitely presented over $\mathfrak{o}$ if it has a presentation $G = F/R$ where $F$ is a finitely generated free group and the $\mathfrak{o}G$ -module $R/\left\lbrack  {R, R}\right\rbrack  { \otimes  }_{\mathbb{Z}}\mathfrak{o}$ is finitely generated. It is easy to see that every finitely presented group $G$ is almost finitely presented over $\mathbb{Z}$ and therefore also over an arbitrary ring $\mathfrak{o}$ . Is the converse true? Yu. I. Merzlyakov
	
	No, it is not true (M. Bestvina, N. Brady, Invent. Math., 129, no. 3 (1997), 445-470).
\end{openproblem}

\plainproblemsection

\begin{openproblem}
	6.36. (J. W. Grossman). The nilpotent-completion diagram of a group $G$ is as follows: $G/{\gamma }_{1}G \leftarrow  G/{\gamma }_{2}G \leftarrow  \cdots$ , where ${\gamma }_{i}G$ is the $i$ th term of the lower central series and the arrows are natural homomorphisms. It is easy to see that every nilpotent-completion diagram ${G}_{1} \leftarrow  {G}_{2} \leftarrow  \cdots$ is a $\gamma$ -diagram, that is, every sequence $1 \rightarrow  {\gamma }_{s}{G}_{s + 1} \rightarrow$ ${G}_{s + 1} \rightarrow  {G}_{s} \rightarrow  1, s = 1,2,\ldots$ , is exact. Do $\gamma$ -diagrams exist that are not nilpotent-completion diagrams? Yu. I. Merzlyakov
	
	Yes, they exists (N. S. Romanovskii, Sibirsk. Mat. Zh., 26, no. 4 (1985), 194-195 (Russian)).
\end{openproblem}

\plainproblemsection

\begin{openproblem}
	6.37. (H. Wielandt, O. H. Kegel). Is a finite group $G$ soluble if it has soluble subgroups $A, B, C$ such that $G = {AB} = {AC} = {BC}$ ?
	
	Yes, it is (mod CFSG) (L. S. Kazarin, Ukrain. Math. J., 43 (1991), 883-886).
\end{openproblem}

\plainproblemsection

\begin{openproblem}
	6.38. a) Let $k$ be a (commutative) field. Find all irreducible subgroups $G$ of $G{L}_{n}\left( k\right)$ having the property that $G \cap  C \neq  \varnothing$ for every conjugacy class $C$ of $G{L}_{n}\left( k\right)$ . I conjecture that $G = G{L}_{n}\left( k\right)$ except in case $n = \operatorname{char}k = 2$ , the field $k$ is quadratically closed, and $G$ is conjugate to the group of all matrices of the form $\left( \begin{matrix} \alpha & 0 \\  0 & \beta  \end{matrix}\right) ,\left( \begin{matrix} 0 & \alpha \\  \beta & 0 \end{matrix}\right)$ where $\alpha  \neq  0$ and $\beta  \neq  0$ . Peter M. Neumann a) The conjecture is refuted (S. A. Zyubin, Algebra and Logic, 45, no. 5 (2006), 296- 305).
\end{openproblem}

\plainproblemsection

\begin{openproblem}
	6.38. b) Is it true that an arbitrary subgroup of $G{L}_{n}\left( k\right)$ that intersects every conjugacy class is parabolic? How far is the same statement true for subgroups of other groups of Lie type? Peter M. Neumann
\end{openproblem}

\plainproblemsection

\begin{openproblem}
	6.39. A class of groups $\mathfrak{K}$ is said to be radical if it is closed under taking homomorphic images and normal subgroups and if every group generated by its normal $\mathfrak{K}$ -subgroups also belongs to $\mathfrak{K}$ . The question proposed relates to the topic "Radical classes and formulae of Narrow (first order) Predicate Calculus (NPC)". One can show that the only non-trivial radical class definable by universal formulae of NPC is the class of all groups. Recently (in a letter to me), G. M. Bergman constructed a family of locally finite radical classes definable by formulae of NPC. Do there exist similar classes which are not locally finite and different from the class of all groups? In particular, does there exist a radical class of groups which is closed under taking Cartesian products, contains an infinite cyclic group, and is different from the class of all groups? B. I. Plotkin
\end{openproblem}

\plainproblemsection

\begin{openproblem}
	6.42. Let $H$ be a strongly 3-embedded subgroup of a finite group $G$ . Suppose that $Z\left( {H/{O}_{{3}^{\prime }}\left( H\right) }\right)$ contains an element of order 3 . Does $Z\left( {G/{O}_{{3}^{\prime }}\left( G\right) }\right)$ necessarily contain an element of order 3 ? N. D. Podufalov
	
	Yes, it does (mod CFSG) (W. Xiao, Sci. China (A), 33 (1990), 1172-1181).
\end{openproblem}

\plainproblemsection

\begin{openproblem}
	6.43. Does the set of quasi-identities holding in the class of all finite groups possess a basis in finitely many variables? D. M. Smirnov
	
	No (A. K. Rumyantsev, Algebra and Logic, 19 (1980), 297-311).
\end{openproblem}

\plainproblemsection

\begin{openproblem}
	6.44. Construct a finitely generated infinite simple group requiring more than two generators. J. Wiegold
	
	This has been done (V. S. Guba, Siberian Math. J., 27 (1986), 670-684).
\end{openproblem}

\plainproblemsection

\begin{openproblem}
	*6.45. Construct a characteristic subgroup $N$ of a finitely generated free group $F$ such that $F/N$ is infinite and simple. If no such exists, it would follow that $d\left( {S}^{2}\right)  = d\left( S\right)$ for every infinite finitely generated simple group $S$ . There is reason to believe that this is false. J. Wiegold
	
	*It is proved that for all $n \geq  2$ , the free group ${F}_{n}$ admits continuum many pairwise non-isomorphic infinite simple characteristic quotients (R. Coulon, F. Fournier-Facio, https://arxiv.org/abs/2312.11684).
\end{openproblem}

\plainproblemsection

\begin{openproblem}
	6.46. If $G$ is $d$ -generator group having no non-trivial finite homomorphic images (in particular, if $G$ is an infinite simple $d$ -generator group) for some integer $d \geq  2$ , must $G \times  G$ be a $d$ -generator group? J. S. Wilson
	
	No, it must not (V. N. Obraztsov, Proc. Roy. Soc. Edinburgh, 123 (1993), 839-855).
\end{openproblem}

\plainproblemsection

\begin{openproblem}
	6.47. (C. D. H. Cooper). Let $G$ be a group, $v$ a group word in two variables such that the operation $x \odot  y = v\left( {x, y}\right)$ defines the structure of a new group ${G}_{v} = \langle G, \odot  \rangle$ on the set $G$ . Does ${G}_{v}$ always lie in the variety generated by $G$ ? E. I. Khukhro
\end{openproblem}

\plainproblemsection

\begin{openproblem}
	6.48. Is every infinite binary finite $p$ -group non-simple? Here $p$ is a prime number.
	
	N. S. Chernikov
\end{openproblem}

\plainproblemsection

\begin{openproblem}
	6.49. Is the minimal condition for abelian normal subgroups inherited by subgroups of finite index? This is true for the minimal condition for (all) abelian subgroups (J. S. Wilson, Math. Z., 114 (1970), 19-21).S. A. Chechin
	
	No, not always. Let $G = \left\lbrack  {\left( {A \times  B \times  C \times  D}\right)  \rtimes  \left( {\langle g\rangle \times \langle t\rangle }\right) }\right\rbrack   \rtimes  \left( {Y\times \langle x\rangle }\right)$ , where $A, B, C, D, Y$ are quasicyclic $p$ -groups, ${x}^{2} = 1$ , while $g$ and $t$ are of infinite order. Let $A = \mathop{\bigcup }\limits_{{n = 1}}^{\infty }\left\langle  {a}_{n}\right\rangle$ , $B = \mathop{\bigcup }\limits_{{n = 1}}^{\infty }\left\langle  {b}_{n}\right\rangle  ,\;C = \mathop{\bigcup }\limits_{{n = 1}}^{\infty }\left\langle  {c}_{n}\right\rangle  ,\;D = \mathop{\bigcup }\limits_{{n = 1}}^{\infty }\left\langle  {d}_{n}\right\rangle  ,\;Y = \mathop{\bigcup }\limits_{{n = 1}}^{\infty }\left\langle  {y}_{n}\right\rangle$ with ${a}_{n + 1}^{p} = {a}_{n},$ ${b}_{n + 1}^{p} = {b}_{n},{c}_{n + 1}^{p} = {c}_{n},{d}_{n + 1}^{p} = {d}_{n},{y}_{n + 1}^{p} = {y}_{n}$ . We impose the relations $\left\lbrack  {{ABCD}, Y}\right\rbrack   =$ $\left\lbrack  {{ABC}, g}\right\rbrack   = \left\lbrack  {{ABD}, t}\right\rbrack   = 1;\left\lbrack  {x, g}\right\rbrack   = g{t}^{-1};\left\lbrack  {x,{a}_{n}}\right\rbrack   = {a}_{n}{b}_{n}^{-1};\left\lbrack  {x,{c}_{n}}\right\rbrack   = {c}_{n}{d}_{n}^{-1};\left\lbrack  {g,{d}_{n}}\right\rbrack   = {b}_{n};$ $\left\lbrack  {t,{c}_{n}}\right\rbrack   = {a}_{n};\left\lbrack  {{y}_{n}, g}\right\rbrack   = {c}_{n};\left\lbrack  {{y}_{n}, t}\right\rbrack   = {d}_{n}$ . Then all abelian normal subgroups of $G$ satisfy the minimal condition for subgroups. The subgroup $H = \lbrack \left( {A \times  B \times  C \times  D}\right)  \rtimes  (\langle g\rangle  \times$ $\langle t\rangle )\rbrack  \rtimes  Y$ does not satisfy the minimal condition for abelian normal subgroups, since the subgroups ${E}_{n} = A \times  B \times  C \times  \left\langle  {g}^{{2}^{n}}\right\rangle$ are normal in $H$ and form a strictly decreasing chain. (S. A. Chechin, Abstracts of 15th All-USSR Algebraic Conf., Krasnoyarsk, 1979).
\end{openproblem}

\plainproblemsection

\begin{openproblem}
	6.51. Let $\mathfrak{F}$ be a local subformation of some formation $\mathfrak{X}$ of finite groups and let $\Omega$ be the set of all maximal homogeneous $\mathfrak{X}$ -screens of $\mathfrak{F}$ . Find a way of constructing the elements of $\Omega$ with the help of the maximal inner local screen of $\mathfrak{F}$ . What can be said about the cardinality of $\Omega$ ? For definitions see (L. A. Shemetkov, Formations of finite groups, Moscow, Nauka, 1978 (Russian)). L. A. Shemetkov
\end{openproblem}

\plainproblemsection

\begin{openproblem}
	6.52. Let $f$ be a local screen of a formation which contains all finite nilpotent groups and let $A$ be a group of automorphisms of a finite group $G$ . Suppose that $A$ acts $f$ -stably on the socle of $G/\Phi \left( G\right)$ . Is it true that $A$ acts $f$ -stably on $\Phi \left( G\right)$ ?
	
	No, not always (A. N. Skiba, Siber. Math. J., 34 (1993), 953-958). L. A. Shemetkov
\end{openproblem}

\plainproblemsection

\begin{openproblem}
	6.53. A group $G$ of the form $G = F \smallsetminus  H$ is said to be a Frobenius group with kernel $F$ and complement $H$ if $H \cap  {H}^{g} = 1$ for any $g \in  G \smallsetminus  H$ and $F \smallsetminus  \{ 1\}  = G \smallsetminus  \bigcup {H}^{g}$ .
	
	What can be said about the kernel and the complement of a Frobenius group? In particular, which groups can be kernels? complements?
	
	Every group can be embedded into the kernel of a Frobenius group, and every right-orderable group can be a complement in a Frobenius group (V. V. Bludov, Siberian Math. J., 38, no. 6 (1997), 1054-1056)
\end{openproblem}

\plainproblemsection

\begin{openproblem}
	6.54. Are there infinite finitely generated Frobenius groups? V. P. Shunkov Yes, such groups exist (A. I. Sozutov, Siberian Math. J., 35, no. 4 (1994), 795-801).
\end{openproblem}

\plainproblemsection

\begin{openproblem}
	6.55. A group $G$ of the form $G = F \leftthreetimes  H$ is said to be a Frobenius group with kernel $F$ and complement $H$ if $H \cap  {H}^{g} = 1$ for any $g \in  G \smallsetminus  H$ and $F \smallsetminus  \{ 1\}  = G \smallsetminus  \bigcup {H}^{g}$ .
	
	Do there exist Frobenius $p$ -groups? V. P. Shunkov
\end{openproblem}

\plainproblemsection

\begin{openproblem}
	6.56. Let $G = F \cdot  \langle a\rangle$ be a Frobenius group with the complement $\langle a\rangle$ of prime order.
	
	a) Is $G$ locally finite if it is binary finite?
	
	b) Is the kernel $F$ locally finite if the groups $\left\langle  {a,{a}^{g}}\right\rangle$ are finite for all $g \in  G$ ?
	
	V. P. Shunkov
\end{openproblem}

\plainproblemsection

\begin{openproblem}
	6.57. A group $G$ is said to be (conjugacy, $p$ -conjugacy) biprimitively finite if, for any finite subgroup $H$ , any two elements of prime order (any two conjugate elements of prime order, of prime order $p$ ) in ${N}_{G}\left( H\right) /H$ generate a finite subgroup. Do the elements of finite order in a (conjugacy) biprimitively finite group $G$ form a subgroup (the periodic part of $G$ )? V. P. Shunkov
	
	Not always (A. A. Cherep, Algebra and Logic, 29 (1987), 311-313).
\end{openproblem}

\plainproblemsection

\begin{openproblem}
	6.58. Are
	
	a) the Alëshin $p$ -groups and
	
	b) the 2-generator Golod $p$ -groups
	
	conjugacy biprimitively finite groups? V. P. Shunkov
	
	a) Not always (A. V. Rozhkov, Math. USSR-Sb., 57 (1987), 437-448).
	
	b) Not always (A. V. Timofeenko, Algebra and Logic, 24 (1985), 129-139).
\end{openproblem}

\plainproblemsection

\begin{openproblem}
	6.59. A group $G$ is said to be (conjugacy, $p$ -conjugacy) biprimitively finite if, for any finite subgroup $H$ , any two elements of prime order (any two conjugate elements of prime order, of prime order $p$ ) in ${N}_{G}\left( H\right) /H$ generate a finite subgroup. Prove that an arbitrary periodic (conjugacy) biprimitively finite group (in particular, having no involutions) of finite rank is locally finite. V. P. Shunkov
\end{openproblem}

\plainproblemsection

\begin{openproblem}
	6.60. Do there exist infinite finitely generated simple periodic (conjugacy) biprim-itively finite groups which contain both involutions and non-trivial elements of odd order? V. P. Shunkov
\end{openproblem}

\plainproblemsection

\begin{openproblem}
	6.61. Is every infinite periodic (conjugacy) biprimitively finite group without involutions non-simple? V. P. Shunkov
\end{openproblem}

\plainproblemsection

\begin{openproblem}
	6.62. Is a (conjugacy, $p$ -conjugacy) biprimitively finite group finite if it has a finite maximal subgroup ( $p$ -subgroup)? There are affirmative answers for conjugacy biprim-itively finite $p$ -groups and for 2-conjugacy biprimitively finite groups (V. P. Shunkov, Algebra and Logic, 9, no. 4 (1970), 291-297; 11, no. 4 (1972), 260-272; 12, no. 5 (1973), 347-353), while the statement does not hold for arbitrary periodic groups, see Archive, 3.9. V. P. Shunkov
\end{openproblem}

\plainproblemsection

\begin{openproblem}
	6.63. An infinite group $G$ is called a monster of the first kind if it has elements of order $> 2$ and for any such an element $a$ and for any proper subgroup $H$ of $G$ , there is an element $g$ in $G \smallsetminus  H$ , such that $\left\langle  {a,{a}^{g}}\right\rangle   = G$ . Classify the monsters of the first kind all of whose proper subgroups are finite.
	
	The centre of such a group coincides with the set of elements of order $\leq  2$ (V. P. Shunkov, Algebra and Logic, 7 (1968), no. 1 (1970), 66-69). An infinite group all of whose proper subgroups are finite is a monster of the first kind if its centre coincides with the set of elements of order $\leq  2$ (A. I. Sozutov, Algebra and Logic,36, no. 5 (1997), 336-348). There are continuously many such groups (A. Yu. Ol'shanskii, Geometry of defining relations in groups, Kluwer, Dordrecht, 1991).
\end{openproblem}

\plainproblemsection

\begin{openproblem}
	6.64. A group $G$ is called a monster of the second kind if it has elements of order $> 2$ and if for any such element $a$ and any proper subgroup $H$ of $G$ there exists an infinite subset ${\mathfrak{M}}_{a, H}$ consisting of conjugates of $a$ by elements of $G \smallsetminus  H$ such that $\langle a, c\rangle  = G$ for all $c \in  {\mathfrak{M}}_{a, H}$ . Do mixed monsters (that is, with elements of both finite and infinite orders) of the second kind exist? Do there exist torsion-free monsters of the second kind? V. P. Shunkov
	
	Yes, they exist, in both cases (A. Yu. Ol'shanskiĭ, The geometry of defining relations in groups, Kluwer, Dordrecht, 1991).
\end{openproblem}

\plainproblemsection

\begin{openproblem}
	7.1. The free periodic groups $B\left( {m, p}\right)$ of prime exponent $p > {665}$ are known to possess many properties similar to those of absolutely free groups (see S. I. Adian, The Burnside Problem and Identities in Groups, Springer, Berlin, 1979). Is it true that all normal subgroups of $B\left( {m, p}\right)$ are not free periodic groups? $\;S.I.{Adian}$ Yes, it is true for all sufficiently large $p$ (A. Yu. Ol’shanskii, in: Groups, rings, Lie and Hopf algebras, Int. Workshop, Canada, 2001, Dordrecht, Kluwer, 2003, 179-187); this is also proved for all primes $p \geq  {1003}$ (V. S. Atabekyan, Fund. Prikl. Mat.,15, no. 1 (2009), 3-21 (Russian)).
\end{openproblem}

\plainproblemsection

\begin{openproblem}
	7.2. Prove that the free periodic groups $B\left( {m, n}\right)$ of odd exponent $n \geq  {665}$ with $m \geq  2$ generators are non-amenable and that random walks on these groups do not have the recurrence property. S. I. Adian
	
	Both assertions are proved (S. I. Adian, Math. USSR-Izv., 21 (1982), 425-434).
\end{openproblem}

\plainproblemsection

\begin{openproblem}
	7.3. Prove that the periodic product of odd exponent $n \geq  {665}$ of non-trivial groups ${F}_{1},\ldots ,{F}_{k}$ which do not contain involutions cannot be generated by less than $k$ elements. This would imply, on the basis of (S. I. Adian, Sov. Math. Doklady, 19, (1978), 910-913), the existence of $k$ -generated simple groups which cannot be generated by less than $k$ elements, for any $k > 0$ . S. I. Adian
\end{openproblem}

\plainproblemsection

\begin{openproblem}
	7.4. Is a finitely generated group with quadratic growth almost abelian? V. V. Belyaev
	
	Yes, it is (M. Gromov, Publ. Math. IHES, 53 (1981), 53-73; J. A. Wolf, Diff. Geometry, 2 (1968), 421-446).
\end{openproblem}

\plainproblemsection

\begin{openproblem}
	7.5. We say that a group is indecomposable if any two of its proper subgroups generate a proper subgroup. Describe the indecomposable periodic metabelian groups.
	
	V. V. Belyaev
\end{openproblem}

\plainproblemsection

\begin{openproblem}
	7.6. Describe the infinite simple locally finite groups with a Chernikov Sylow 2-subgroup. In particular, are such groups the Chevalley groups over locally finite fields of odd characteristic? V. V. Belyaev, N. F. Sesekin
	
	No, as every countably infinite locally finite $p$ -group can be embedded as a maximal $p$ -subgroup of the simple Hall’s universal group $U$ , the unique countable existentially closed group in the class of all locally finite groups (K. Hickin, Proc. London Math. Soc. (3), 52, no. 1 (1986), 53-72). But if all Sylow 2-subgroups are Chernikov, then this is true mod CFSG (O. H. Kegel, Math. Z., 95 (1967), 169-195; V. V. Belyaev, in: Investigations in Group Theory, Sverdlovsk, UNC AN SSSR, 1984, 39-50 (Russian); A. V. Borovik, Siberian Math. J., 24, no. 6 (1983), 843-851; B. Hartley, G. Shute, Quart. J. Math. Oxford (2), 35 (1984), 49-71; S. Thomas, Arch. Math., 41 (1983), 103-116). When all Sylow 2-subgroups are Chernikov, the same result, without using CFSG, follows from (M. J. Larsen, R. Pink, J. Amer. Math. Soc., 24 (2011), 1105- 1158 (also known as a preprint of 1998)).
\end{openproblem}

\plainproblemsection

\begin{openproblem}
	7.7. (Well-known problem). Is the group $G = \left\langle  {a, b \mid  {a}^{9} = 1,{ab} = {b}^{2}{a}^{2}}\right\rangle$ finite? This group contains $F\left( {2,9}\right)$ , the only Fibonacci group for which it is not yet known whether it is finite or infinite. R. G. Burns
	
	No, it is infinite, since $F\left( {2,9}\right)$ is infinite (M. F. Newman, Arch. Math.,54, no. 3 (1990), 209-212).
\end{openproblem}

\plainproblemsection

\begin{openproblem}
	7.8. Suppose that $H$ is a normal subgroup of a group $G$ , where $H$ and $G$ are subdirect products of the same $n$ groups ${G}_{1},\ldots ,{G}_{n}$ . Does the nilpotency class of $G/H$ increase with $n$ ? Yu. M. Gorchakov
	
	Yes, it does (E. I. Khukhro, Sibirsk. Mat. Zh., 23, no. 6 (1982), 178-180 (Russian)).
\end{openproblem}

\plainproblemsection

\begin{openproblem}
	7.12. Find all groups with a Hall ${2}^{\prime }$ -subgroup.R. L. Griess
	
	These are found mod CFSG (Z. Arad, M. B. Ward, J. Algebra, 77 (1982), 234-246).
\end{openproblem}

\plainproblemsection

\begin{openproblem}
	7.15. Prove that if ${F}^{ * }\left( G\right)$ is quasisimple and $\alpha  \in  \operatorname{Aut}G,\left| \alpha \right|  = 2$ , then ${C}_{G}\left( \alpha \right)$ contains an involution outside $Z\left( {{F}^{ * }\left( G\right) }\right)$ , except when ${F}^{ * }\left( G\right)$ has quaternion Sylow 2-subgroups. R. Griess
\end{openproblem}

\plainproblemsection

\begin{openproblem}
	7.16. If $H$ is a proper subgroup of the finite group $G$ , is there always an element of prime-power order not conjugate to an element of $H$ ? R. L. Griess
	
	Yes, always (mod CFSG) (B. Fein, W. M. Kantor, M. Schacher, J. Reine Angew. Math., 328 (1981), 39-57).
\end{openproblem}

\plainproblemsection

\begin{openproblem}
	*7.17. Is the number of maximal subgroups of the finite group $G$ at most $\left| G\right|  - 1$ ?
	
	Editors’ remarks: This is proved for $G$ soluble (G. E. Wall, J. Austral. Math. Soc., $\mathbf{2}\left( {{1961} - {62}}\right) ,{35} - {59})$ and for symmetric groups ${\mathbb{S}}_{n}$ for sufficiently large $n$ (M. Liebeck, A. Shalev, J. Combin. Theory, Ser. A, 75 (1996), 341-352); it is also proved (M. W. Liebeck, L. Pyber, A. Shalev, J. Algebra, 317 (2007), 184-197) that any finite group $G$ has at most ${2C}{\left| G\right| }^{3/2}$ maximal subgroups, where $C$ is an absolute constant. R. Griess
	
	*No, not always (R. Guralnick, F. Lübeck, L. Scott, T. Sprowl; see (C. P. Bendel, B. D. Boe, C. M. Drupieski, D. K. Nakano, B. J. Parshall, C. Pillen, C. B. Wright, in: Developments and retrospectives in Lie theory. Algebraic methods. Retrospective selected papers based on the presentations at the seminar "Lie groups, Lie algebras and their representations", 1991-2014, Springer, Cham, 2014, 51-69). Infinitely many counterexamples were found in (F. Lübeck, Trans. Amer. Math. Soc., 373, no. 4 (2020), 2331-2347).
\end{openproblem}

\plainproblemsection

\begin{openproblem}
	7.19. Construct an explicit example of a finitely presented simple group with word problem not solvable by a primitive recursive function. F. B. Cannonito
\end{openproblem}

\plainproblemsection

\begin{openproblem}
	7.21. Is there an algorithm which decides if an arbitrary finitely presented solvable group is metabelian? F. B. Cannonito
\end{openproblem}

\plainproblemsection

\begin{openproblem}
	7.22. Suppose that a finite group $G$ is realized as the automorphism group of some torsion-free abelian group. Is it true that for every infinite cardinal $\mathfrak{m}$ there exist ${2}^{\mathfrak{m}}$ non-isomorphic torsion-free abelian groups of cardinality $\mathfrak{m}$ whose automorphism groups are isomorphic to $G$ ? S. F. Kozhukhov
	
	Yes, this is true in the Zermelo-Frenkel system with axioms of choice and 'weak diamond' (M. Dugas, R. Göbel, Proc. London Math. Soc. (3), 45, no. 2 (1982), 319-336), or if $\mathfrak{m}$ is smaller than the first measurable cardinal (V. A. Nikiforov, Mat. Zametki, 39, no. 5 (1986), 641-646 (Russian)).
\end{openproblem}

\plainproblemsection

\begin{openproblem}
	7.23. Does there exist an algorithm which decides, for a given list of group identities, whether the elementary theory of the variety given by this list is soluble, that is, by (A. P. Zamyatin, Algebra and Logic, 17, no. 1 (1978), 13-17), whether this variety is abelian?
\end{openproblem}

\plainproblemsection

\begin{openproblem}
	7.24. We say that a group is sparse if the variety generated by it has at most countably many subvarieties. Does there exist a finitely generated sparse group that has undecidable word problem? A. V. Kuznetsov
	
	Yes, such groups do exist; for example, a free group of ${\mathfrak{N}}_{3}\mathfrak{A}$ . By (A. N. Krasil’nikov, Math. USSR-Izv., 37 (1991), 539-553) every subvariety of ${\mathfrak{N}}_{3}\mathfrak{A}$ has a finite basis for its laws; hence there are only countably many of them. It has undecidable word problem by (O. G. Kharlampovich, Sov. Math., 32, no. 11 (1988), 136-140).
\end{openproblem}

\plainproblemsection

\begin{openproblem}
	7.25. We associate with words in the alphabet $x,{x}^{-1}, y,{y}^{-1}, z,{z}^{-1},\ldots$ operations which are understood as functions in variables $x, y, z,\ldots$ . We say that a word $A$ is expressible in words ${B}_{1},\ldots ,{B}_{n}$ on the group $G$ if $A$ can be constructed from the words ${B}_{1},\ldots ,{B}_{n}$ and variables by means of finitely many substitutions of words one into another and replacements of a word by another word which is identically equal to it on $G$ . A list of words is said to be functionally complete on $G$ if every word can be expressed on $G$ in words from this list. A word $B$ is called a Schaeffer word on $G$ if every word can be expressed on $G$ in $B$ (compare with A.V. Kuznetsov, Matem. Issledovaniya, Kishinëv, 6, no. 4 (1971), 75-122 (Russian)). Does there exist an algorithm which decides
	
	a) by a word $B$ , whether it is a Schaeffer word on every group? Compare with the problem of describing all such words in (A. G. Kurosh, Theory of Groups, Moscow, Nauka,1967, p. 435 (Russian)); here are examples of such words: $x{y}^{-1},{x}^{-1}{y}^{2}z$ , ${x}^{-1}{y}^{-1}{zx}$ .
	
	b) by a list of words, whether it is functionally complete on every group?
	
	c) by words $A,{B}_{1},\ldots ,{B}_{n}$ , whether $A$ is expressible in ${B}_{1},\ldots ,{B}_{n}$ on every group? (This is a problem from A. V. Kuznetsov, ibid., p. 112.)
	
	d) the same for every finite group? (For finite groups, for example, ${x}^{-1}$ is expressible in ${xy}$ .) For a fixed finite group an algorithm exists (compare with A. V. Kuznetsov, ibid., §8). A. V. Kuznetsov
\end{openproblem}

\plainproblemsection

\begin{openproblem}
	7.26. The ordinal height of a variety of groups is, by definition, the supremum of order types (ordinals) of all well-ordered by inclusion chains of its proper subvarieties. It is clear that its cardinality is either finite or countably infinite, or equal to ${\omega }_{1}$ . Is every countable ordinal number the ordinal height of some variety? A. V. Kuznetsov
\end{openproblem}

\plainproblemsection

\begin{openproblem}
	7.27. Is it true that the group $S{L}_{n}\left( q\right)$ contains, for sufficiently large $q$ , a diagonal matrix which is not contained in any proper irreducible subgroup of $S{L}_{n}\left( q\right)$ with the exception of block-monomial ones? For $n = 2,3$ the answer is known to be affirmative (V. M. Levchuk, in: Some problems of the theory of groups and rings, Inst. of Physics SO AN SSSR, Krasnoyarsk, 1973 (Russian)). Similar problems are interesting for other Chevalley groups. V. M. Levchuk
\end{openproblem}

\plainproblemsection

\begin{openproblem}
	7.28. Let $G\left( K\right)$ be the Chevalley group over a commutative ring $K$ associated with the root system $\Phi$ as defined in (R. Steinberg, Lectures on Chevalley groups, Yale Univ., New Haven, Conn., 1967). This group is generated by the root subgroups ${x}_{r}\left( K\right) , r \in  \Phi$ . We define an elementary carpet of type $\Phi$ over $K$ to be any collection of additive subgroups $\left\{  {{\mathfrak{A}}_{r} \mid  r \in  \Phi }\right\}$ of $K$ satisfying the condition
	
	$$
	{c}_{{ij},{rs}}{\mathfrak{A}}_{r}^{i}{\mathfrak{A}}_{s}^{j} \subseteq  {\mathfrak{A}}_{{ir} + {js}}\;\text{ for }r, s,{ir} + {js} \in  \Phi , i > 0, j > 0,
	$$
	
	where ${c}_{{ij},{rs}}$ are constants defined by the Chevalley commutator formula and ${\mathfrak{A}}_{r}^{i} =$ $\left\{  {{a}^{i} \mid  a \in  {\mathfrak{A}}_{r}}\right\}$ . What are necessary and sufficient conditions (in terms of the ${\mathfrak{A}}_{r}$ ) on the elementary carpet to ensure that the subgroup $\left\langle  {{x}_{r}\left( {\mathfrak{A}}_{r}\right)  \mid  r \in  \Phi }\right\rangle$ of $G\left( K\right)$ intersects with ${x}_{r}\left( K\right)$ in ${x}_{r}\left( {\mathfrak{A}}_{r}\right)$ ? See also 15.46. V. M. Levchuk
\end{openproblem}

\plainproblemsection

\begin{openproblem}
	7.30. Which finite simple groups can be generated by three involutions, two of which commute? V. D. Mazurov
	
	The answer is known mod CFSG. For the alternating groups and groups of Lie type see (Ya. N. Nuzhin, Algebra and Logic, 36, no. 4 (1997), 245-256). For sporadic groups B. L. Abasheev, A. V. Ershov, N. S. Nevmerzhitskaya, S. Norton, Ya. N. Nuzhin, A. V. Timofeenko have shown that the groups ${M}_{11},{M}_{22},{M}_{23}$ , and ${McL}$ cannot be generated as required, while the others can; see more details in (V. D. Mazurov, Siberian Math. J., 44, no. 1 (2003), 160-164).
\end{openproblem}

\plainproblemsection

\begin{openproblem}
	7.31. Let $A$ be a group of automorphisms of a finite non-abelian 2-group $G$ acting transitively on the set of involutions of $G$ . Is $A$ necessarily soluble?
	
	Such groups $G$ were divided into several classes in (F. Gross, J. Algebra,40, no. 2 (1976), 348-353). For one of these classes a positive answer was given by E. G. Bryukhanova (Algebra and Logic, 20, no. 1 (1981), 1-12). V. D. Mazurov
\end{openproblem}

\plainproblemsection

\begin{openproblem}
	7.33. An elementary ${TI}$ -subgroup $V$ of a finite group $G$ is said to be a subgroup of non-root type if $1 \neq  {N}_{V}\left( {V}^{g}\right)  \neq  V$ for some $g \in  G$ . Describe the finite groups $G$ which contain a 2-subgroup $V$ of non-root type such that $\left\lbrack  {V,{V}^{g}}\right\rbrack   = 1$ implies that all involutions in $V{V}^{g}$ are conjugate to elements of $V$ . A. A. Makhnëv
\end{openproblem}

\plainproblemsection

\begin{openproblem}
	7.34. In many of the sporadic finite simple groups the 2-ranks of the centralizers of 3-elements are at most 2. Describe the finite groups satisfying this condition.
	
	A. A. Makhnëv
\end{openproblem}

\plainproblemsection

\begin{openproblem}
	7.35. What varieties of groups $\mathfrak{V}$ have the following property: the group $G/\mathfrak{V}\left( G\right)$ is residually finite for any residually finite group $G$ ? O. V. Mel'nikov
\end{openproblem}

\plainproblemsection

\begin{openproblem}
	7.36. Is it true that every residually finite group in which every subgroup of finite index (including the group itself) is defined by a single defining relation is either free or isomorphic to the fundamental group of a compact surface? O. V. Mel'niko No; for example, let ${H}_{n} = \left\langle  {x, y \mid  {y}^{-1}{xy} = {x}^{n}}\right\rangle  , n = 2,3,\ldots$ ; then every subgroup of finite index in ${H}_{n}$ is isomorphic to a group ${H}_{m}$ for some $m(\mathrm{\;V}$ . A. Churkin, Abstracts of 8th All-USSR Symp. on Group Theory, Kiev, 1982, 139-140 (Russian)).
\end{openproblem}

\plainproblemsection

\begin{openproblem}
	7.37. We say that a profinite group is strictly complete if each of its subgroups of finite index is open. It is known (B. Hartley, Math. Z., 168, no. 1 (1979), 71-76) that finitely generated profinite groups having a finite series with pronilpotent factors are strictly complete. Is a profinite group strictly complete if it is
	
	a) finitely generated?
	
	b) finitely generated and prosoluble? O. V. Mel'nikov a) Yes, it is (N. Nikolov, D. Segal, Ann. Math. (2), 165, no. 1 (2007), 171-273). b) Yes, it is (D. Segal, Proc. London Math. Soc. (3), 81 (2000), 29-54).
\end{openproblem}

\plainproblemsection

\begin{openproblem}
	7.38. A variety of profinite groups is a non-empty class of profinite groups closed under taking subgroups, factor-groups, and Tikhonov products. A subvariety $\mathfrak{V}$ of the variety $\mathfrak{N}$ of all pronilpotent groups is said to be (locally) nilpotent if all (finitely generated) groups in $\mathfrak{V}$ are nilpotent.
	
	a) Is it true that any non-nilpotent subvariety of $\mathfrak{N}$ contains a non-nilpotent locally nilpotent subvariety?
	
	b) The same problem for the variety of all pro- $p$ -groups (for a given prime $p$ ).
	
	O. V. Mel'nikov
\end{openproblem}

\plainproblemsection

\begin{openproblem}
	7.39. Let $G = \left\langle  {a, b \mid  {a}^{p} = {\left( ab\right) }^{3} = {b}^{2} = {\left( {a}^{\sigma }b{a}^{2/\sigma }b\right) }^{2} = 1}\right\rangle$ , where $p$ is a prime, $\sigma$ is an integer not divisible by $p$ . The group ${PS}{L}_{2}\left( p\right)$ is a factor-group of $G$ so that there is a short exact sequence $1 \rightarrow  N \rightarrow  G \rightarrow  {PS}{L}_{2}\left( p\right)  \rightarrow  1$ . For each $p > 2$ there is $\sigma$ such that $N = 1$ , for example, $\sigma  = 4$ . Let ${N}^{ab}$ denote the factor-group of $N$ by its commutator subgroup. It is known that for some $p$ there is $\sigma$ such that ${N}^{ab}$ is infinite (for example, for $p = {41}$ one can take ${\sigma }^{2} \equiv  2\left( {\;\operatorname{mod}\;{41}}\right)$ ), whereas for some other $p$ (for example, for $p = {43}$ ) the group ${N}^{ab}$ is finite for every $\sigma$ .
	
	a) Is the set of primes $p$ for which ${N}^{ab}$ is finite for every $\sigma$ infinite?
	
	b) Is there an arithmetic condition on $\sigma$ which ensures that ${N}^{ab}$ is finite?
	
	J. Mennicke
\end{openproblem}

\plainproblemsection

\begin{openproblem}
	7.40. Describe the (lattice of) subgroups of a given classical matrix group over a ring which contain the subgroup consisting of all matrices in that group with coefficients in some subring (see Ju. I. Merzljakov, J. Soviet Math., 1 (1973), 571-593).
	
	Yu. I. Merzlyakov
	
	If the root system has single edges $\left( {{A}_{l},{D}_{l},{E}_{l}}\right)$ and the larger ring is not quasi-algebraic over the smaller one, then it is shown that the lattice in question contains the lattice of subgroups of a free product containing one of the factors (A. V. Stepanov, J. Algebra, 324, no. 7 (2010), 1549-1557), which means that a reasonable description is unfeasible. In almsot all other cases a good description was obtained in (R. A. Schmidt, Zap. Nauchn. Semin. Leningr. Otd. Mat. Inst. Steklova, 94 (1979), 119-130 (Russian); Ya. N. Nuzhin, Algebra Logic, 22 (1983), 378-389; Ya. N. Nuzhin, A. V. Yakushevich, Algebra Logic, 39 (2000), 199-206; A. Stepanov, J. Algebra, 362 (2012), 12-29; Ya. N. Nuzhin, Siberian Math. J., 54, no. 1 (2013), 119-123; A. Stepanov, A. Bak, St. Petersburg Math. J., 28, no. 4 (2017), 47-61).
\end{openproblem}

\plainproblemsection

\begin{openproblem}
	7.41. (J. S. Wilson). Is every linear $\overline{SI}$ -group an $\overline{SN}$ -group? $\;{Yu}$ . Yu. I. Merzlyakov
\end{openproblem}

\plainproblemsection

\begin{openproblem}
	7.42. A group $U$ is called an ${F}_{q}$ -group (where $q \in  \pi \left( U\right)$ ) if, for each finite subgroup $K$ of $U$ and for any two elements $a, b$ of order $q$ in $T = {N}_{U}\left( K\right) /K$ , there exists $c \in  T$ such that the group $\left\langle  {a,{b}^{c}}\right\rangle$ is finite. A group $U$ is called an ${F}^{ * }$ -group if each subgroup $H$ of $U$ is an ${F}_{q}$ -group for every $q \in  \pi \left( H\right)$ (V. P. Shunkov,1977).
	
	a) Is every primary ${F}^{ * }$ -group satisfying the minimum condition for subgroups almost abelian?
	
	b) Does every ${F}^{ * }$ -group satisfying the minimum condition for (abelian) subgroups possess the radicable part? A. N. Ostylovskiǐ
	
	No. A counterexample to both questions is given by an infinite group all of whose subgroups are conjugate and have prime order (A. Yu. Ol'shanskiĭ, Math. USSR-Izv., 16 (1981), 279-289).
\end{openproblem}

\plainproblemsection

\begin{openproblem}
	7.44. Does a normal subgroup $H$ in a finite group $G$ possess a complement in $G$ if each Sylow subgroups of $H$ is a direct factor in some Sylow subgroup of $G$ ? V. I. Sergiyenko Yes, it does, mod CFSG (W. Xiao, J. Pure Appl. Algebra, 87, no. 1 (1993), 97-98).
\end{openproblem}

\plainproblemsection

\begin{openproblem}
	7.45. Does the $Q$ -theory of the class of all finite groups (in the sense of 2.40) coincide with the $Q$ -theory of a single finitely presented group? D. M. Smirnov
\end{openproblem}

\plainproblemsection

\begin{openproblem}
	7.48. (Well-known problem). Suppose that, in a finite group $G$ , each two elements of the same order are conjugate. Is then $\left| G\right|  \leq  6$ ?
	
	Yes, it is, mod CFSG (P. Fitzpatrick, Proc. Roy. Irish Acad. Sect. A, 85, no. 1 (1985), 53-58); see also (W. Feit, G. Seitz, Illinois J. Math., 33, no. 1 (1988), 103-131) and (R. W. van der Waall, A. Bensaid, Simon Stevin, 65 (1991), 361-374).
\end{openproblem}

\plainproblemsection

\begin{openproblem}
	7.49. Let $G$ be a finitely generated group, and $N$ a minimal normal subgroup of $G$ which is an elementary abelian $p$ -group. Is it true that either $N$ is finite or the growth function for the elements of $N$ with respect to a finite generating set of $G$ is bounded below by an exponential function? V. I. Trofimov
\end{openproblem}

\plainproblemsection

\begin{openproblem}
	7.50. Study the structure of the primitive permutation groups (finite and infinite), in which the stabilizer of any three pairwise distinct points is trivial. This problem is closely connected with the problem of describing those group which have a Frobenius group as one of maximal subgroups. A. N. Fomin
\end{openproblem}

\plainproblemsection

\begin{openproblem}
	7.51. What are the primitive permutation groups (finite and infinite) which have a regular sub-orbit, that is, in which a point stabilizer acts faithfully and regularly on at least one of its orbits? A. N. Fomin
\end{openproblem}

\plainproblemsection

\begin{openproblem}
	7.52. Describe the locally finite primitive permutation groups in which the centre of any Sylow 2-subgroup contains involutions stabilizing precisely one symbol. The case of finite groups was completely determined by D. Holt in 1978. A. N. Fomin
\end{openproblem}

\plainproblemsection

\begin{openproblem}
	7.53. Let $p$ be a prime. The law $x \cdot  {x}^{\varphi }\cdots {x}^{{\varphi }^{p - 1}} = 1$ from the definition of a splitting automorphism (see Archive,1.10) gives rise to a variety of groups with operators $\langle \varphi \rangle$ consisting of all groups that admit a splitting automorphism of order $p$ . Does the analogue of Kostrikin's theorem hold for this variety, that is, do the locally nilpotent groups in this variety form a subvariety? E. I. Khukhro
	
	Yes, it does (E. I. Khukhro, Math. USSR-Sb., 58 (1987), 119-126).
\end{openproblem}

\plainproblemsection

\begin{openproblem}
	7.54. Does there exist a group of infinite special rank which can be represented as a product of two subgroups of finite special rank? N. S. Chernikov
\end{openproblem}

\plainproblemsection

\begin{openproblem}
	7.55. Is it true that a group which is a product of two almost abelian subgroups is almost soluble? N. S. Chernikov
\end{openproblem}

\plainproblemsection

\begin{openproblem}
	7.56. (B. Amberg). Does a group satisfy the minimum (respectively, maximum) condition on subgroups if it is a product of two subgroups satisfying the minimum (respectively, maximum) condition on subgroups? N. S. Chernikov
\end{openproblem}

\plainproblemsection

\begin{openproblem}
	7.57. a) A set of generators of a finitely generated group $G$ consisting of the least possible number $d\left( G\right)$ of elements is called a basis of $G$ . Let ${r}_{M}\left( G\right)$ be the least number of relations necessary to define $G$ in the basis $M$ and let $r\left( G\right)$ be the minimum of the numbers ${r}_{M}\left( G\right)$ over all bases $M$ of $G$ . It is known that ${r}_{M}\left( G\right)  \leq  d\left( G\right)  + r\left( G\right)$ for any basis $M$ . Does there exist a finitely presented group $G$ for which the inequality becomes equality for some basis $M$ ? V. A. Churkin
\end{openproblem}

\plainproblemsection

\begin{openproblem}
	7.57. A set of generators of a finitely presented group $G$ that consists of the least possible number $d\left( G\right)$ of generators is called a basis for $G$ . Let ${r}_{M}\left( G\right)$ be the least number of relations necessary to define $G$ in the basis $M$ , and $r\left( G\right)$ the minimum of ${r}_{M}\left( G\right)$ over all bases $M$ for $G$ . Let ${G}_{1},{G}_{2}$ be any non-trivial groups.
	
	b) Is it true that ${r}_{{M}_{1} \cup  {M}_{2}}\left( {{G}_{1} * {G}_{2}}\right)  = {r}_{{M}_{1}}{G}_{1} + {r}_{{M}_{2}}\left( {G}_{2}\right)$ for any bases ${M}_{1},{M}_{2}$ of ${G}_{1},{G}_{2}$ , respectively?
	
	c) Is it true that $r\left( {{G}_{1} * {G}_{2}}\right)  = r\left( {G}_{1}\right)  + r\left( {G}_{2}\right)$ ?V. A. Churkinb) Not always; c) not always (C. Hog, M. Lusztig, W. Metzler, in: Presentation classes, 3-manifolds and free products (Lecture Notes in Math., 1167), Springer, Berlin, 1985, 154-167).
\end{openproblem}

\plainproblemsection

\begin{openproblem}
	7.58. Suppose that $F$ is an absolutely free group, $R$ a normal subgroup of $F$ , and let $\mathfrak{V}$ be a variety of groups. It is well-known (H. Neumann, Varieties of Groups, Springer, Berlin,1967) that the group $F/\mathfrak{V}\left( R\right)$ is isomorphically embeddable in the $\mathfrak{V}$ -verbal wreath product of a $\mathfrak{V}$ -free group of the same rank as $F$ with $F/R$ . Find a criterion indicating which elements of this wreath product belong to the image of this embedding. A criterion is known in the case where $\mathfrak{V}$ is the variety of all abelian groups (V. N. Remeslennikov, V. G. Sokolov, Algebra and Logic, 9, no. 5 (1970), 342- 349). G. G. Yabanzhi
\end{openproblem}

\plainproblemsection

\begin{openproblem}
	8.1. Characterize all groups (or at least all soluble groups) $G$ and fields $F$ such that every irreducible ${FG}$ -module has finite dimension over the centre of its endomorphism ring. For some partial answers see (B. A. F. Wehrfritz, Glasgow Math. J., 24, no. 1
	
	(1983), 169-176). B. A. F. Wehrfritz
\end{openproblem}

\plainproblemsection

\begin{openproblem}
	8.2. Let $G$ be the amalgamated free product of two polycyclic groups, amalgamating a normal subgroup of each. Is $G$ isomorphic to a linear group? $G$ is residually finite (G. Baumslag, Trans. Amer. Math. Soc., 106, no. 2 (1963), 193-209) and the answer is "yes" if the amalgamated part is torsion-free, abelian (B. A. F. Wehrfritz, Proc. London Math. Soc., 27, no. 3 (1973), 402-424) and nilpotent (M. Shirvani, 1981,
	
	unpublished). B. A. F. Wehrfritz
\end{openproblem}

\plainproblemsection

\begin{openproblem}
	8.3. Let $m$ and $n$ be positive integers and $p$ a prime. Let ${P}_{n}\left( {\mathbb{Z}}_{{p}^{m}}\right)$ be the group of all $n \times  n$ matrices $\left( {a}_{ij}\right)$ over the integers modulo ${p}^{m}$ such that ${a}_{ii} \equiv  1$ for all $i$ and ${a}_{ij} \equiv  0\left( {\;\operatorname{mod}\;p}\right)$ for all $i > j$ . The group ${P}_{n}\left( {\mathbb{Z}}_{{p}^{m}}\right)$ is a finite $p$ -group. For which $m$ and $n$ is it regular?
	
	Editors’ comment (2001): The group ${P}_{n}\left( {\mathbb{Z}}_{{p}^{m}}\right)$ is known to be regular if ${mn} < p$ (Yu. I. Merzlyakov, Algebra i Logika, 3, no. 4 (1964), 49-59 (Russian)). The case $m = 1$ is completely done by A. V. Yagzhev in (Math. Notes,56, no. 5-6 (1995), 1283-1290), although this paper is erroneous for $m > 1$ . The case $m = 2$ is completed in (S. G. Kolesnikov, Issledovaniya in analysis and algebra, no. 3, Tomsk Univ., Tomsk, 2001,117-125 (Russian)). Editors’ comment (2009): The group ${P}_{n}\left( {\mathbb{Z}}_{{p}^{m}}\right)$ was shown to be regular for any $m$ if ${n}^{2} < p$ (S. G. Kolesnikov, Siberian Math. J.,47, no. 6 (2006),
	
	1054-1059). B. A. F. Wehrfritz
\end{openproblem}

\plainproblemsection

\begin{openproblem}
	8.4. Construct a finite nilpotent loop with no finite basis for its laws.
	
	M. R. Vaughan-Lee
\end{openproblem}

\plainproblemsection

\begin{openproblem}
	8.5. Prove that if $X$ is a finite group, $F$ is any field, and $M$ is a non-trivial irreducible ${FX}$ -module then $\frac{1}{\left| X\right| }\mathop{\sum }\limits_{{x \in  X}}\dim \operatorname{fix}\left( x\right)  \leq  \frac{1}{2}\dim M$ .
	
	Proved, even with $\ldots  \leq  \frac{1}{p}\dim M$ , where $p$ is the least prime divisor of $\left| G\right|$ (R. M. Guralnick, A. Maróti, Adv. Math., 226, no. 1 (2011), 298-308).
\end{openproblem}

\plainproblemsection

\begin{openproblem}
	8.6. (M. M. Day). Do the classes of amenable groups and elementary groups coincide? The latter consists of groups that can be obtained from commutative and finite groups by forming subgroups, factor-groups, extensions, and direct limits. R. I. Grigorchuk No (R. I. Grigorchuk, Soviet Math. Dokl., 28 (1983), 23-26).
\end{openproblem}

\plainproblemsection

\begin{openproblem}
	8.7. Does there exist a non-amenable finitely presented group which has no free subgroups of rank 2? R. I. Grigorchuk
	
	Yes, it exists (A. Yu. Ol'shanskii, M. V. Sapir, Publ. Math. Inst. Hautes Étud. Sci., 96 (2002), 43-169).
\end{openproblem}

\plainproblemsection

\begin{openproblem}
	8.8. (D.V.Anosov). a) Does there exist a non-cyclic finitely-generated group $G$ containing an element $a$ such that every element of $G$ is conjugate to a power of $a$ ? Yes, it exists (V. S. Guba, Math. USSR-Izv., 29 (1986), 233-277). R. I. Grigorchuk
\end{openproblem}

\plainproblemsection

\begin{openproblem}
	8.8. b) (D. V. Anosov). Does there exist a non-cyclic finitely presented group $G$ which contains an element $a$ such that each element of $G$ is conjugate to some power of $a$ ?
	
	R. I. Grigorchuk
\end{openproblem}

\plainproblemsection

\begin{openproblem}
	8.9. (C. Chou). We say that a group $G$ has property $P$ if for every finite subset $F$ of $G$ , there is a finite subset $S \supset  F$ and a subset $X \subset  G$ such that ${x}_{1}S \cap  {x}_{2}S$ is empty for any ${x}_{1},{x}_{2} \in  X,{x}_{1} \neq  {x}_{2}$ , and $G = \mathop{\bigcup }\limits_{{x \in  X}}{Sx}$ . Does every group have property $P$ ? R. I. Grigorchuk
\end{openproblem}

\plainproblemsection

\begin{openproblem}
	8.10. Is the group $G = \left\langle  {a, b \mid  {a}^{n} = 1,{ab} = {b}^{3}{a}^{3}}\right\rangle$ finite or infinite for $n = 7, n = 9$ , and $n = {15}$ ? All other cases known, see Archive,7.7. Infinite in all three cases. For $n = {15}$ by (D. J. Seal, Proc. Roy. Soc. Edinburgh (A), 92 (1982), 181-192). For $n = 9$ (and 15) by (M. I. Prishchepov, Commun. Algebra, 23 (1995),5095-5117). For $n = 7$ , because it contains the Fibonacci group $F\left( {3,7}\right)$ as an index 7 subgroup, as follows from Theorem 3.0 of (C.P. Chalk, Commun. Algebra 26, no. 5 (1998), 1511-1546) by standard technique for working with Fibonacci groups (G. Williams, Letter of 6 October 2015).
\end{openproblem}

\plainproblemsection

\begin{openproblem}
	*8.11. Consider the group
	
	$$
	M = \langle x, y, z, t \mid  \left\lbrack  {x, y}\right\rbrack   = \left\lbrack  {y, z}\right\rbrack   = \left\lbrack  {z, x}\right\rbrack   = \left( {x, t}\right)  = \left( {y, t}\right)  = \left( {z, t}\right)  = 1\rangle
	$$
	
	where $\left\lbrack  {x, y}\right\rbrack   = {x}^{-1}{y}^{-1}{xy}$ and $\left( {x, t}\right)  = {x}^{-1}{t}^{-1}{x}^{-1}{txt}$ . The subgroup $H = \langle x, t, y\rangle$ is isomorphic to the braid group ${\mathfrak{B}}_{4}$ , and is normally complemented by $N = \left\langle  {\left( z{x}^{-1}\right) }^{M}\right\rangle$ . Is $N$ a free group (?) of countably infinite rank (?) on which $H$ acts faithfully by conjugation? D. L. Johnson
	
	*As pointed by $\mathrm{Y}$ . Antolín, the group $M$ is the Artin group of type ${D}_{4}$ ; it is proved in (B. Perron, J. P. Vannier, Math. Ann.,306, no. 2 (1996),231-246) that $N$ is a free group of rank 3 and $H$ acts on $N$ faithfully by conjugation.
\end{openproblem}

\plainproblemsection

\begin{openproblem}
	8.12. Let ${D0}$ denote the class of finite groups of deficiency zero, i. e. having a presentation $\langle X \mid  R\rangle$ with $\left| X\right|  = \left| R\right|$ .
	
	a) Does ${D0}$ contain any 3-generator $p$ -group for $p \geq  5$ ?
	
	c) Do soluble ${D0}$ -groups have bounded derived length?
	
	d) Which non-abelian simple groups can occur as composition factors of ${D0}$ - groups? D. L. Johnson, E. F. Robertson
\end{openproblem}

\plainproblemsection

\begin{openproblem}
	8.12. b) Let ${D0}$ denote the class of finite groups of deficiency zero, i. e. having a presentation $\langle X \mid  R\rangle$ with $\left| X\right|  = \left| R\right|$ . Are the central factors of nilpotent ${D0}$ -groups 3-generated? D. L. Johnson, E. F. Robertson
	
	No, not always (G. Havas, E. F. Robertson, Commun. Algebra, 24 (1996), 3483-3487).
\end{openproblem}

\plainproblemsection

\begin{openproblem}
	8.13. Let $G$ be a simple algebraic group over an algebraically closed field of characteristic $p$ and $\mathfrak{G}$ be the Lie algebra of $G$ . Is the number of orbits of nilpotent elements of $\mathfrak{G}$ under the adjoint action of $G$ finite? This is known to be true if $p$ is not too small (i. e. if $p$ is not a "bad prime" for $G$ ). Yes, it is (D. F. Holt, N. Spaltenstein, J. Austral. Math. Soc. (A), 38 (1985), 330- 350).
\end{openproblem}

\plainproblemsection

\begin{openproblem}
	8.14. a) Assume a group $G$ is existentially closed in the class $L{\mathfrak{N}}_{p}$ of all locally finite $p$ -groups. Is it true that $G$ is characteristically simple? This is true for $G$ countable in $L{\mathfrak{N}}_{p}$ (Berthold Maier, Freiburg); in fact, up to isomorphism, there is only one such countable locally finite $p$ -group. O. H. Kegel
	
	Not always (S. R. Thomas, Arch. Math., 44 (1985), 98-109).
\end{openproblem}

\plainproblemsection

\begin{openproblem}
	8.14. b) Assume a group $G$ is existentially closed in one of the classes $L{\mathfrak{N}}^{ + }, L{\mathfrak{S}}_{\pi }$ , $L{\mathfrak{S}}^{ + }, L\mathfrak{S}$ of, respectively, all locally nilpotent torsion-free groups, all locally soluble $\pi$ -groups, all locally soluble torsion-free groups, or all locally soluble groups. Is it true that $G$ is characteristically simple?
	
	The existential closedness of $G$ is a local property, thus it seems difficult to obtain global properties of $G$ from it. See also Archive,8.14 a. O. H. Kegel
\end{openproblem}

\plainproblemsection

\begin{openproblem}
	8.15. (Well-known problem). Is every $\widetilde{N}$ -group a $\bar{Z}$ -group? For definitions, see (M. I. Kargapolov, Yu. I. Merzlyakov, Fundamentals of the Theory of Groups, 3rd Edition, Moscow, Nauka, 1982, p. 205 (Russian)). Sh. S. Kemkhadze
\end{openproblem}

\plainproblemsection

\begin{openproblem}
	8.16. Is the class of all $\bar{Z}$ -groups closed under taking normal subgroups?
	
	Sh. S. Kemkhadze
\end{openproblem}

\plainproblemsection

\begin{openproblem}
	8.17. An $\overset{⏜}{RN}$ -group is one whose every homomorphic image is an ${RN}$ -group. Is the class of $\overset{⏜}{RN}$ -groups closed under taking normal subgroups? Sh. S. Kemkhadze No (J. S. Wilson, Arch. Math., 25 (1974), 574-577).
\end{openproblem}

\plainproblemsection

\begin{openproblem}
	8.18. Is every countably infinite abelian group a verbal subgroup of some finitely generated (soluble) relatively free group?
	
	Yes, it is (A. Storozhev, Commun. Algebra, 22, no. 7 (1994), 2677-2701).
\end{openproblem}

\plainproblemsection

\begin{openproblem}
	8.19. Which of the following properties of (soluble) varieties of groups are equivalent to one another:
	
	1) to satisfy the minimum condition for subvarieties;
	
	2) to have at most countably many subvarieties;
	
	3) to have no infinite independent system of identities? Yu. G. Kleiman
\end{openproblem}

\plainproblemsection

\begin{openproblem}
	8.20. What is the cardinality of the set of all varieties covering an abelian (nilpotent? Cross? hereditarily finitely based?) variety of groups? The question is related to 4.46 and 4.73 . Yu. G. Kleiman
	
	There are continually many varieties covering the variety $A$ of abelian groups, as well as the variety ${A}_{n}$ of abelian groups of sufficiently large odd exponent $n$ (P. A. Kozhevnikov, On varieties of groups of large odd exponent, Dep. 1612-V00, VINITI, Moscow, 2000 (Russian); S. V. Ivanov, A. M. Storozhev, Contemp. Math., 360 (2004), 55-62).
\end{openproblem}

\plainproblemsection

\begin{openproblem}
	8.21. If the commutator subgroup ${G}^{\prime }$ of a relatively free group $G$ is periodic, must ${G}^{\prime }$ have finite exponent? L. G. Kovács
\end{openproblem}

\plainproblemsection

\begin{openproblem}
	8.22. If $G$ is a (non-abelian) finite group contained in a join $\mathfrak{A} \vee  \mathfrak{B}$ of two varieties $\mathfrak{A},\mathfrak{B}$ of groups, must there exist finite groups $A \in  \mathfrak{A}, B \in  \mathfrak{B}$ such that $G$ is a section of the direct product $A \times  B$ ? L. G. Kovács
	
	Not always (A. Storozhev, Bull. Austral. Math. Soc., 51, no. 2 (1995), 287-290).
\end{openproblem}

\plainproblemsection

\begin{openproblem}
	8.23. If the dihedral group $D$ of order 18 is a section of a direct product $A \times  B$ , must at least one of $A$ and $B$ have a section isomorphic to $D$ ? L. G. Kovács
\end{openproblem}

\plainproblemsection

\begin{openproblem}
	8.24. Prove that every linearly ordered group of finite rank is soluble.
	
	V. M. Kopytov
\end{openproblem}

\plainproblemsection

\begin{openproblem}
	8.25. Does there exist an algorithm which recognizes by an identity whether it defines a variety of groups that has at most countably many subvarieties? A. V. Kuznetsov
\end{openproblem}

\plainproblemsection

\begin{openproblem}
	8.26. We call a variety passable if there exists an unrefinable chain of its subvarieties which is well-ordered by inclusion. For example, every variety generated by its finite groups is passable - this is an easy consequence of (H. Neumann, Varieties of groups, Berlin et al., Springer, 1967, Chapter 5). Do there exist non-passable varieties of groups? A. V. Kuznetsov
	
	Yes, they exist (M.I. Anokhin, Moscow Univ. Math. Bull., 51, no. 1 (1996), 48-49).
\end{openproblem}

\plainproblemsection

\begin{openproblem}
	8.27. Does the lattice of all varieties of groups possess non-trivial automorphisms? A. V. Kuznetsov
\end{openproblem}

\plainproblemsection

\begin{openproblem}
	8.28. Is the variety of groups finitely based if it is generated by a finitely based quasivariety of groups? A. V. Kuznetsov
	
	No, not always (M. I. Anokhin, Sb. Math., 189, no. 7-8 (1998), 1115-1124).
\end{openproblem}

\plainproblemsection

\begin{openproblem}
	8.29. Do there exist locally nilpotent groups with trivial centre satisfying the weak maximal condition for normal subgroups? L. A. Kurdachenko
\end{openproblem}

\plainproblemsection

\begin{openproblem}
	8.30. Let $\mathfrak{X},\mathfrak{Y}$ be Fitting classes of soluble groups which satisfy the Lockett condition, i. e. $\mathfrak{X} \cap  {\mathfrak{S}}_{ * } = {\mathfrak{X}}_{ * },\mathfrak{Y} \cap  {\mathfrak{S}}_{ * } = {\mathfrak{Y}}_{ * }$ where $\mathfrak{S}$ denotes the Fitting class of all soluble groups and the lower star the bottom group of the Lockett section determined by the given Fitting class. Does $\mathfrak{X} \cap  \mathfrak{Y}$ satisfy the Lockett condition?
	
	Editors’ comment (2001): The answer is affirmative if $\mathfrak{X}$ and $\mathfrak{Y}$ are local (A. Gry-tczuk, N. T. Vorob'ev, Tsukuba J. Math., 18, no. 1 (1994), 63-67). H. Lausch
\end{openproblem}

\plainproblemsection

\begin{openproblem}
	*8.31. Describe the finite groups in which every proper subgroup has a complement in some larger subgroup. Among these groups are, for example, ${PS}{L}_{2}\left( 7\right)$ and all Sylow subgroups of symmetric groups. V. M. Levchuk
	
	*These groups are described independently in (V.M. Levchuk, A. G. Likharev, Siberian Math. J., 47, no. 4 (2006), 659-668) and in (V. N. Tyutyanov, Proc. Gomel' State Univ., 3 (2006), 178-183 (in Russian)).
\end{openproblem}

\plainproblemsection

\begin{openproblem}
	8.32. Suppose $G$ is a finitely generated group such that, for any set $\pi$ of primes and any subgroup $H$ of $G$ , if $G/\left\langle  {H}^{G}\right\rangle$ is a finite $\pi$ -group then $\left| {G : H}\right|$ is a finite $\pi$ -number. Is $G$ nilpotent? This is true for finitely generated soluble groups. J.C. Lennox Not always (V. N. Obraztsov, J. Austral. Math. Soc. (A), 61, no. 2 (1996), 267-288).
\end{openproblem}

\plainproblemsection

\begin{openproblem}
	8.33. Let $a, b$ be two elements of a group, $a$ having infinite order. Find a necessary and sufficient condition for $\mathop{\bigcap }\limits_{{n = 1}}^{\infty }\left\langle  {{a}^{n}, b}\right\rangle   = \langle b\rangle$ . F. N. Liman
\end{openproblem}

\plainproblemsection

\begin{openproblem}
	8.34. Let $G$ be a finite group. Is it true that indecomposable projective $\mathbb{Z}G$ -modules are finitely generated (and hence locally free)? P. A. Linnell
	
	Yes, it is; see Remark 2.13 and Corollary 4.2 in (P. Príhoda, Rend. Semin. Mat. Univ. Padova, 123 (2010), 141-167).
\end{openproblem}

\plainproblemsection

\begin{openproblem}
	8.35. Determine the conjugacy classes of maximal subgroups in the sporadic simple groups
	
	a) ${F}_{24}^{\prime }$ ;
	
	b) ${F}_{2}$ .
	
	a) They are determined mod CFSG (S. A. Linton, R. A. Wilson, Proc. London Math. Soc., 63, no. 1 (1991), 113-164).
	
	b) They are determined mod CFSG (R. A. Wilson, J. Algebra, 211 (1999), 1-14).
\end{openproblem}

\plainproblemsection

\begin{openproblem}
	8.35. c) Determine the conjugacy classes of maximal subgroups in the sporadic simple group ${F}_{1}$ . See the current status in Atlas of Finite Group Representations (http://brauer.maths.qmul.ac.uk/Atlas/spor/M/).V.D.Mazurov
\end{openproblem}

\plainproblemsection

\begin{openproblem}
	8.37. (R. Griess). a) Is ${M}_{11}$ a section in ${O}^{\prime }N$ ?
	
	b) The same question for ${M}_{24}$ in ${F}_{2};{J}_{1}$ in ${F}_{1}$ and in ${F}_{2};{J}_{2}$ in ${F}_{23},{F}_{24}^{\prime }$ , and ${F}_{2}$ . V. D. Mazurov
	
	a) Yes (A. A. Ivanov, S. V. Shpektorov, Abs. 18 All-Union Algebra Conf., Part 1, Kishinev, 1985, 209 (Russian); S. Yoshiara, J. Fac. Sci. Univ. Tokyo (IA), 32 (1985), 105-141).
	
	b) No (R. A. Wilson, Bull. London Math. Soc., 18 (1986), 349-350).
\end{openproblem}

\plainproblemsection

\begin{openproblem}
	8.39. b) Describe the irreducible subgroups of $S{L}_{6}\left( q\right)$ . V. D. Mazurov
	
	They are described mod CFSG (A. S. Kondratiev, Algebra and Logic, 28 (1989), 122- 138; P. Kleidman, Low-dimensional finite classical groups and their subgroups, Harlow, Essex, 1989).
\end{openproblem}

\plainproblemsection

\begin{openproblem}
	8.40. Describe the finite groups generated by a conjugacy class $D$ of involutions which satisfies the following property: if $a, b \in  D$ and $\left| {ab}\right|  = 4$ , then $\left\lbrack  {a, b}\right\rbrack   \in  D$ . This condition is satisfied, for example, in the case where $D$ is a conjugacy class of involutions of a known finite simple group such that $\left\langle  {{C}_{D}\left( a\right) }\right\rangle$ is a 2-group for every $a \in  D$ . A. A. Makhnëv
\end{openproblem}

\plainproblemsection

\begin{openproblem}
	8.41. What finite groups $G$ contain a normal set of involutions $D$ which contains a non-empty proper subset $T$ satisfying the following properties:
	
	1) ${C}_{D}\left( a\right)  \subseteq  T$ for any $a \in  T$ ;
	
	2) if $a, b \in  T$ and ${ab} = {ba} \neq  1$ , then ${C}_{D}\left( {ab}\right)  \subseteq  T$ ?A. A. Makhnëv
\end{openproblem}

\plainproblemsection

\begin{openproblem}
	8.42. Describe the finite groups all of whose soluble subgroups of odd indices have 2-length 1. For example, the groups ${L}_{n}\left( {2}^{m}\right)$ are known to satisfy this condition.
	
	A. A. Makhnëv
\end{openproblem}

\plainproblemsection

\begin{openproblem}
	8.43. (F. Timmesfeld). Let $T$ be a Sylow 2-subgroup of a finite group $G$ . Suppose that $\langle N\left( B\right)  \mid  B$ is a non-trivial characteristic subgroup of $T\rangle$ is a proper subgroup of $G$ . Describe the group $G$ if ${F}^{ * }\left( M\right)  = {O}_{2}\left( M\right)$ for any 2-local subgroup $M$ containing $T$ . A. A. Makhnëv
\end{openproblem}

\plainproblemsection

\begin{openproblem}
	8.44. Prove or disprove that for all, but finitely many, primes $p$ the group
	
	$$
	{G}_{p} = \left\langle  {a, b \mid  {a}^{2} = {b}^{p} = {\left( ab\right) }^{3} = {\left( {b}^{r}a{b}^{-{2r}}a\right) }^{2} = 1}\right\rangle  ,
	$$
	
	where ${r}^{2} + 1 \equiv  0\left( {\;\operatorname{mod}\;p}\right)$ , is infinite. A solution of this problem would have interesting topological applications. It was proved with the aid of computer that ${G}_{p}$ is finite for $p \leq  {17}$ . J. Mennicke
\end{openproblem}

\plainproblemsection

\begin{openproblem}
	8.45. When it follows that a group is residually a $\left( {\mathfrak{X} \cap  \mathfrak{Y}}\right)$ -group, if it is both residually a ±-group and residually a Y-group? Yu. I. Merzlyakov
\end{openproblem}

\plainproblemsection

\begin{openproblem}
	8.46. Describe the automorphisms of the symplectic group $S{p}_{2n}$ over an arbitrary commutative ring. Conjecture: they are all standard. The conjecture was proved (V. M. Petechuk, Algebra and Logic, 22 (1983), 397-405).
\end{openproblem}

\plainproblemsection

\begin{openproblem}
	8.47. Do there exist finitely presented soluble groups in which the maximum condition for normal subgroups fails but all central sections are finitely generated?
	
	Yes (Yu. V. Sosnovskii, Math. Notes, 36 (1984), 577-580). Yu. I. Merzlyakov
\end{openproblem}

\plainproblemsection

\begin{openproblem}
	8.48. If a finite group $G$ can be written as the product of two soluble subgroups of odd index, then is $G$ soluble? V. S. Monakhov
	
	Yes, it is (mod CFSG) (L. S. Kazarin, Ukrain. Math. J., 43 (1991), 883-886).
\end{openproblem}

\plainproblemsection

\begin{openproblem}
	8.49. Let $G$ be a $p$ -group acting transitively as a permutation group on a set $\Omega$ , let $F$ be a field of characteristic $p$ , and regard ${F\Omega }$ as an ${FG}$ -module. Then do the descending and ascending Loewy series of ${F\Omega }$ coincide? $\;P.M.$ Neumann Yes, they do (J. L. Alperin, Quart. J. Math. Oxford (2), 39, no. 154 (1988), 129-133).
\end{openproblem}

\plainproblemsection

\begin{openproblem}
	8.50. At a conference in Oberwolfach in 1979 I exhibited a finitely generated soluble group $G$ (of derived length 3) and a non-zero cyclic $\mathbb{Z}G$ -module $V$ such that $V \cong  V \oplus  V$ . Can this be done with $G$ metabelian? I conjecture that it cannot. For background of this problem and for details of the construction see (Peter M. Neumann, in: Groups-Korea, Pusan, 1988 (Lect. Notes Math., 1398), Springer, Berlin, 1989, 124-139).
	
	Peter M. Neumann
\end{openproblem}

\plainproblemsection

\begin{openproblem}
	8.51. (J. McKay). If $G$ is a finite group and $p$ a prime let ${m}_{p}\left( G\right)$ denote the number of ordinary irreducible characters of $G$ whose degree is not divisible by $p$ . Let $P$ be a Sylow $p$ -subgroup of $G$ . Is it true that ${m}_{p}\left( G\right)  = {m}_{p}\left( {{N}_{G}\left( P\right) }\right)$ ? J. B. Olsson
\end{openproblem}

\plainproblemsection

\begin{openproblem}
	8.52. (Well-known problem). Do there exist infinite finitely presented periodic groups? Compare with 6.3. A. Yu. Ol'shanskii
\end{openproblem}

\plainproblemsection

\begin{openproblem}
	8.53. a) Let $n$ be a sufficiently large odd number. Describe the automorphisms of the free Burnside group $B\left( {m, n}\right)$ of exponent $n$ on $m$ generators. A. Yu. Ol’shanski
\end{openproblem}

\plainproblemsection

\begin{openproblem}
	8.53. b) Let $n$ be a sufficiently large odd number. Is it true that every non-cyclic subgroup of $B\left( {m, n}\right)$ has a subgroup isomorphic to $B\left( {2, n}\right)$ ? $A.{Yu}.{Ol}$ ’shanskii Yes, it is true for odd $n \geq  {1003}$ (V. S. Atabekyan, Izv. Math.,73, no. 5 (2009), 861-892).
\end{openproblem}

\plainproblemsection

\begin{openproblem}
	8.54. a) (Well-known problem). Classify metabelian varieties of groups (or show that this is, in a certain sense, a "wild" problem).
	
	b) Describe the identities of 2-generated metabelian groups, that is, classify varieties generated by such groups. A. Yu. Ol'shanskiĭ
\end{openproblem}

\plainproblemsection

\begin{openproblem}
	8.55. It is easy to see that the set of all quasivarieties of groups, in each of which quasivarieties a non-trivial identity holds, is a semigroup under multiplication of qua-sivarieties. Is this semigroup free? A. Yu. Ol'shanskiĭ
\end{openproblem}

\plainproblemsection

\begin{openproblem}
	8.56. Let $X$ be a finite set and $f$ a mapping from the set of subsets of $X$ to the positive integers. Under the requirement that in a group generated by $X$ , every subgroup $\langle Y\rangle , Y \subseteq  X$ , be nilpotent of class $\leq  f\left( Y\right)$ , is it true that the free group ${G}_{f}$ relative to this condition is torsion-free?
	
	Not always (V. V. Bludov, V. F. Kleimenov, E. V. Khlamov, Algebra and Logic, 29 (1990), 95-96).
\end{openproblem}

\plainproblemsection

\begin{openproblem}
	8.58. Does a locally compact locally nilpotent nonabelian group without elements of finite order contain a proper closed isolated normal subgroup? V. M. Poletskikh
\end{openproblem}

\plainproblemsection

\begin{openproblem}
	8.59. Suppose that, in a locally compact locally nilpotent group $G$ , all proper closed normal subgroups are compact. Does $G$ contain an open compact normal subgroup? V. M. Poletskikh, I. V. Protasov
\end{openproblem}

\plainproblemsection

\begin{openproblem}
	8.60. Describe the locally compact primary locally soluble groups which are covered by compact subgroups and all of whose closed abelian subgroups have finite rank.
	
	V. M. Poletskikh, V. S. Charin
\end{openproblem}

\plainproblemsection

\begin{openproblem}
	8.61. Suppose that a locally compact group $G$ contains a subgroup that is topologically isomorphic to the additve group of the field of real numbers with natural topology. Is the space of all closed subgroups of $G$ connected in the Chabauty topology? I. V. Protasov
	
	No, not always: let $H$ be the group of matrices of the form $\left( \begin{matrix} 1 & {z}_{1} & r \\  0 & 1 & {z}_{2} \\  0 & 0 & 1 \end{matrix}\right)$ , where ${z}_{1},{z}_{2} \in  \mathbb{Z}$ and $r \in  \mathbb{R}$ . Then $H$ is locally compact in the natural topology and contains a central subgroup topologically isomorphic to $\mathbb{R}$ , but $L\left( H\right)$ is not connected (Yu. V. Tsybenko, Abstracts of 17th All-USSR Algebraic Conf., Part 1, Minsk, 1983, 213 (Russian)).
\end{openproblem}

\plainproblemsection

\begin{openproblem}
	8.62. Describe the locally compact locally pronilpotent groups for which the space of closed normal subgroups is compact in the $E$ -topology. The corresponding problem has been solved for discrete groups. I. V. Protasov
\end{openproblem}

\plainproblemsection

\begin{openproblem}
	8.63. Suppose that the space of all closed subgroups of a locally compact group $G$ is $\sigma$ -compact in the $E$ -topology. Is it true that the set of closed non-compact subgroups of $G$ is at most countable? I. V. Protasov
	
	Yes, it is (A. G. Piskunov, Ukrain. J. Math., 40 (1988), 679-683).
\end{openproblem}

\plainproblemsection

\begin{openproblem}
	8.64. Does the class of all finite groups possess an independent basis of quasiidenti-ties? A. K. Rumyantsev, D. M. Smirnov
\end{openproblem}

\plainproblemsection

\begin{openproblem}
	8.66. Construct examples of residually finite groups which would separate Shunkov’s classes of groups with(a, b)-finiteness condition,(weakly) conjugacy biprimitively finite groups and (weakly) biprimitively finite groups (see, in particular, 6.57). Can one derive such examples from Golod's construction? A. I. Sozutov
	
	Such examples are constructed (L. Hammoudi, Nil-algèbres non-nilpotentes et groupes périodiques infinis (Doctor Thesis), Strasbourg, 1996; A. V. Rozhkov, Finiteness conditions in groups of automorphisms of trees (Doctor of Sci. Thesis), Krasnoyarsk, 1997; Algebra and Logic, 37, no. 3 (1998), 192-203).
\end{openproblem}

\plainproblemsection

\begin{openproblem}
	8.67. Do there exist Golod groups all of whose abelian subgroups have finite ranks? Here a Golod group means a finitely generated non-nilpotent subgroup of the adjoint group of a nil-ring. The answer is negative for the whole adjoint group of a nil-ring (Ya. P. Sysak, Abstracts of the 17th All-Union Algebraic Conf., part 1, Minsk, 1983, p. 180 (Russian)). A. I. Sozutov
\end{openproblem}

\plainproblemsection

\begin{openproblem}
	8.68. Let $G = \langle a, b \mid  r = 1\rangle$ where $r$ is a cyclically reduced word that is not a proper power of any word in $a, b$ . If $G$ is residually finite, is ${G}_{t} = \left\langle  {a, b \mid  {r}^{t} = 1}\right\rangle  , t > 1$ , residually finite?
	
	Editors' comment: An analogous fact was proved for relative one-relator groups (S. J. Pride, Proc. Amer. Math. Soc., 136, no. 2 (2008), 377-386). C. Y. Tang Yes, it is (D. T. Wise, The structure of groups with a quasiconvex hierarchy, Annals of Mathematics Studies, 209, Princeton University Press, Princeton, NJ, 2021).
\end{openproblem}

\plainproblemsection

\begin{openproblem}
	8.69. Is every 1-relator group with non-trivial torsion conjugacy separable? Yes, it is (A. Minasyan, P. Zalesskii, J. Algebra, 382 (2013), 39-45).C. Y. Tang
\end{openproblem}

\plainproblemsection

\begin{openproblem}
	8.70. Let $A, B$ be polycyclic-by-finite groups. Let $G = A{ * }_{H}B$ where $H$ is cyclic. Is $G$ conjugacy separable? C. Y. Tang
	
	Yes, it is (L. Ribes, D. Segal, P. A. Zalesskii, J. London Math. Soc. (2), 57, no. 3 (1998), 609-628).
\end{openproblem}

\plainproblemsection

\begin{openproblem}
	8.71. Is every countable conjugacy-separable group embeddable in a 2-generator conjugacy-separable group? C. Y. Tang
	
	Yes, it is (V. A. Roman'kov, Embedding theorems for residually finite groups, Preprint 84-515, Comp. Centre, Novosibirsk, 1984 (Russian)).
\end{openproblem}

\plainproblemsection

\begin{openproblem}
	8.72. Does there exist a finitely presented group $G$ which is not free or cyclic of prime order, having the property that every proper subgroup of $G$ is free? C. Y. Tang
\end{openproblem}

\plainproblemsection

\begin{openproblem}
	8.73. We say that a finite group $G$ separates cyclic subgroups if, for any cyclic subgroups $A$ and $B$ of $G$ , there is a $g \in  G$ such that $A \cap  {B}^{g} = 1$ . Is it true that $G$ separates cyclic subgroups if and only if $G$ has no non-trivial cyclic normal subgroups? J. G. Thompson
	
	Not always. For ${p}_{1} = 2,{p}_{2} = 5,{p}_{3} = {11},{p}_{4} = {17}$ let ${R}_{i}$ be an elementary abelian group of order ${p}_{i}^{2}$ and ${\varphi }_{i}$ a regular automorphism of order 3 of ${R}_{i}, i = 1,2,3,4$ . In the direct product ${R}_{1}\left\langle  {\varphi }_{1}\right\rangle   \times  {R}_{2}\left\langle  {\varphi }_{2}\right\rangle   \times  {R}_{3}\left\langle  {\varphi }_{3}\right\rangle   \times  {R}_{4}\left\langle  {\varphi }_{4}\right\rangle$ let $G$ be the subgroup generated by all the ${R}_{i}$ and the elements ${\varphi }_{2}{\varphi }_{3}{\varphi }_{4}$ and ${\varphi }_{1}{\varphi }_{3}{\varphi }_{4}^{-1}$ . Then $G$ has no non-trivial cyclic normal subgroups and every (cyclic) subgroup of order $2 \cdot  5 \cdot  {11} \cdot  {17}$ intersects any of its conjugates non-trivially. (N. D. Podufalov, Abstracts of the 9th All-Union Group Theory Symp., Moscow, 1984, 113-114 (Russian).)
\end{openproblem}

\plainproblemsection

\begin{openproblem}
	8.74. A subnormal subgroup $H \vartriangleleft   \vartriangleleft  G$ of a group $G$ is said to be good if and only if $\langle H, J\rangle  \vartriangleleft   \vartriangleleft  G$ whenever $J \vartriangleleft   \vartriangleleft  G$ . Is it true that when $H$ and $K$ are good subnormal subgroups of $G$ , then $H \cap  K$ is good? J. G. Thompson
\end{openproblem}

\plainproblemsection

\begin{openproblem}
	*8.75. (A known problem). Suppose $G$ is a finite primitive permutation group on $\Omega$ , and $\alpha ,\beta$ are distinct points of $\Omega$ . Does there exist an element $g \in  G$ such that ${\alpha g} = \beta$ and $g$ fixes no point of $\Omega$ ? J. G. Thompson
	
	*Not always (P. Müller, Preprint, 2023, https://arxiv.org/pdf/2304.08459.pdf).
\end{openproblem}

\plainproblemsection

\begin{openproblem}
	8.76. Give a realistic upper bound for the torsion-free rank of a finitely generated nilpotent group in terms of the ranks of its abelian subgroups. More precisely, for each integer $n$ let $f\left( n\right)$ be the largest integer $h$ such that there is a finitely generated nilpotent group of torsion-free rank $h$ with the property that all abelian subgroups have torsion-free rank at most $n$ . It is easy to see that $f\left( n\right)$ is bounded above by $n\left( {n + 1}\right) /2$ . Describe the behavior of $f\left( n\right)$ for large $n$ . Is $f\left( n\right)$ bounded below by a quadratic in $n$ ? J. S. Wilson
	
	Yes, it is (M. V. Milenteva, J. Group Theory, 7, no. 3 (2004), 403-408).
\end{openproblem}

\plainproblemsection

\begin{openproblem}
	8.77. Do there exist strongly regular graphs with parameters $\lambda  = 0,\mu  = 2$ of degree $k > {10}$ ? Such graphs are known for $k = 5$ and $k = {10}$ , their automorphism groups are primitive permutation groups of rank 3 . D. G. Fon-Der-Flaass
\end{openproblem}

\plainproblemsection

\begin{openproblem}
	8.78. It is known that there exists a countable locally finite group that contains a copy of every other countable locally finite group. For which other classes of countable groups does a similar "largest" group exist? In particular, what about periodic locally soluble groups? periodic locally nilpotent groups? B. Hartley
\end{openproblem}

\plainproblemsection

\begin{openproblem}
	8.79. Does there exist a countable infinite locally finite group $G$ that is complete, in the sense that $G$ has trivial centre and no outer automorphisms? An uncountable one exists (K. K. Hickin, Trans. Amer. Math. Soc., 239 (1978), 213-227).
	
	B. Hartley
\end{openproblem}

\plainproblemsection

\begin{openproblem}
	8.80. Let $G$ be a locally finite group containing a maximal subgroup which is Cher-nikov. Is $G$ almost soluble? B. Hartley
	
	Yes, it is (B. Hartley, Algebra and Logic, 37, no. 1 (1998), 101-106).
\end{openproblem}

\plainproblemsection

\begin{openproblem}
	8.81. Let $G$ be a finite $p$ -group admitting an automorphism $\alpha$ of prime order $q$ with $\left| {{C}_{G}\left( a\right) }\right|  \leq  n$ .
	
	a) If $p = q$ , then does $G$ have a nilpotent subgroup of class at most 2 and index bounded by a function of $n$ ?
	
	b) If $p \neq  q$ , then does $G$ have a nilpotent subgroup of class bounded by a function of $q$ and index bounded by a function of $n$ (and, possibly, $q$ )?
	
	a) Not necessarily (E. I. Khukhro, Math. Notes, 38 (1985), 867-870).
	
	b) Yes, it does (E. I. Khukhro, Math. USSR-Sb., 71, no. 1 (1992), 51-63).
\end{openproblem}

\plainproblemsection

\begin{openproblem}
	8.82. Let $\mathfrak{H} = \mathbb{C} \times  \mathbb{R} = \{ \left( {z, r}\right)  \mid  z \in  \mathbb{C}, r > 0\}$ be the three-dimensional Poincaré space which admits the following action of the group $S{L}_{2}\left( \mathbb{C}\right)$ :
	
	$$
	\left( {z, r}\right) \left( \begin{array}{ll} a & b \\  c & d \end{array}\right)  = \left( {\frac{\left( {{az} + b}\right) \left( \overline{{cz} + d}\right)  + a\bar{c}{r}^{2}}{{\left| cz + d\right| }^{2} + {\left| c\right| }^{2}{r}^{2}},\frac{r}{{\left| cz + d\right| }^{2} + {\left| c\right| }^{2}{r}^{2}}}\right) .
	$$
	
	Let $\mathfrak{o}$ be the ring of integers of the field $K = \mathbb{Q}\left( \sqrt{D}\right)$ , where $D < 0,\pi$ a prime such that $\pi \bar{\pi }$ is a prime in $\mathbb{Z}$ and let $\Gamma  = \left\{  {\left. {\left( \begin{array}{ll} a & b \\  c & d \end{array}\right)  \in  S{L}_{2}\left( \mathfrak{o}\right) }\right| \;b \equiv  0\left( {\;\operatorname{mod}\;\pi }\right) }\right\}$ . Adding to the space $\Gamma  \smallsetminus  {\mathfrak{H}}^{3}$ two vertices we get a three-dimensional compact space $\overline{\Gamma  \smallsetminus  {\mathfrak{H}}^{3}}$ .
	
	Calculate $r\left( \pi \right)  = {\dim }_{\mathbb{Q}}{H}_{1}\left( {\overline{\Gamma  \smallsetminus  {\mathfrak{H}}^{3}},\mathbb{Q}}\right)  = {\dim }_{\mathbb{Q}}\left( {{\Gamma }^{ab} \otimes  \mathbb{Q}}\right)$ .
	
	For example, if $D =  - 3$ then $r\left( \pi \right)$ is distinct from zero for the first time for $\pi  \mid  {73}$ (then $r\left( \pi \right)  = 1$ ), and if $D =  - 4$ then $r\left( \pi \right)  = 1$ for $\pi  \mid  {137}$ . H. Helling
\end{openproblem}

\plainproblemsection

\begin{openproblem}
	8.83. In the notation of 8.82, for $r\left( \pi \right)  > 0$ , one can define via Hecke algebras a formal Dirichlet series with Euler multiplication (see G. Shimura, Introduction to the arithmetic theory of automorphic functions, Princeton Univ. Press, 1971). Does there exist an algebraic Hasse-Weil variety whose $\zeta$ -function is this Dirichlet series? There are several conjectures. H. Helling
\end{openproblem}

\plainproblemsection

\begin{openproblem}
	8.84. We say that an automorphism $\varphi$ of the group $G$ is a pseudo-identity if, for all $x \in  G$ , there exists a finitely generated subgroup ${K}_{x}$ of $G$ such that $x \in  {K}_{x}$ and ${\left. \varphi \right| }_{{K}_{x}}$ is an automorphism of ${K}_{x}$ . Let $G$ be generated by subgroups $H, K$ and let $G$ be locally nilpotent. Let $\varphi  : G \rightarrow  G$ be an endomorphism such that ${\left. \varphi \right| }_{H}$ is pseudo-identity of $H$ and ${\left. \varphi \right| }_{K}$ is an automorphism of $K$ . Does it follow that $\varphi$ is an automorphism of $G$ ? It is known that $\varphi$ is a pseudo-identity of $G$ if, additionally, ${\left. \varphi \right| }_{K}$ is pseudo-identity of $K$ ; it is also known that $\varphi$ is an automorphism if, additionally, $K$ is normal in $G$ . P. Hilton
	
	No, not necessarily (A. V. Yagzhev, Math. Notes, 56, no. 5 (1994), 1205-1207).
\end{openproblem}

\plainproblemsection

\begin{openproblem}
	8.85. Construct a finite $p$ -group $G$ whose Hughes subgroup ${H}_{p}\left( G\right)  = \langle x \in  G \mid$ $\left| x\right|  \neq  p\rangle$ is non-trivial and has index ${p}^{3}$ . E. I. Khukhro
\end{openproblem}

\plainproblemsection

\begin{openproblem}
	8.86. (Well-known problem). Suppose that all proper closed subgroups of a locally compact locally nilpotent group $G$ are compact. Is $G$ abelian if it is non-compact? V. S. Charin
\end{openproblem}

\plainproblemsection

\begin{openproblem}
	8.87. Find all hereditary local formations $\mathfrak{F}$ of finite groups satisfying the following condition: every finite minimal non- $\mathfrak{F}$ -group is biprimary. A finite group is said to be biprimary if its order is divisible by precisely two distinct primes. L. A. Shemetkov They are found (V. N. Semenchuk, Problems of Algebra: Proc. of the Gomel' State Univ., no. 1 (15) (1999), 92-102 (Russian)).
\end{openproblem}

\plainproblemsection

\begin{openproblem}
	9.1. A group $G$ is said to be potent if, to each $x \in  G$ and each positive integer $n$ dividing the order of $x$ (we suppose $\infty$ is divisible by every positive integer), there exists a finite homomorphic image of $G$ in which the image of $x$ has order precisely $n$ . Is the free product of two potent groups again potent? R. B. J. T. Allenby
\end{openproblem}

\plainproblemsection

\begin{openproblem}
	9.2. Is $G = \left\langle  {a, b \mid  {a}^{l} = {b}^{m} = {\left( ab\right) }^{n} = 1}\right\rangle$ conjugacy separable? $R.B.J.T$ . Allenby Yes, it is (B. Fine, G. Rosenberger, Contemp. Math., 109 (1990), 11-18).
\end{openproblem}

\plainproblemsection

\begin{openproblem}
	9.3. Suppose that a countable locally finite group $G$ contains no proper subgroups isomorphic to $G$ itself and suppose that all Sylow subgroups of $G$ are finite. Does $G$ possess a non-trivial finite normal subgroup?
	
	Yes, it does (S. D. Bell, Locally finite groups with Černikov Sylow subgroups, (Ph. D. Thesis), University of Manchester, 1994).
\end{openproblem}

\plainproblemsection

\begin{openproblem}
	9.4. It is known that if $\mathfrak{M}$ is a variety (quasivariety, pseudovariety) of groups, then the class $I\mathfrak{M}$ of all quasigroups that are isotopic to groups in $\mathfrak{M}$ is also a variety (quasivariety, pseudovariety), and $I\mathfrak{M}$ is finitely based if and only if $\mathfrak{M}$ is finitely based. Is it true that if $\mathfrak{M}$ is generated by a single finite group then $I\mathfrak{M}$ is generated by a single finite quasigroup? A. A. Gvaramiya
\end{openproblem}

\plainproblemsection

\begin{openproblem}
	9.5. A variety of groups is called primitive if each of its subquasivarieties is a variety. Describe all primitive varieties of groups. Is every primitive variety of groups locally finite? V. A. Gorbunov
\end{openproblem}

\plainproblemsection

\begin{openproblem}
	9.6. Is it true that an independent basis of quasiidentities of any finite group is finite? V. A. Gorbunov
\end{openproblem}

\plainproblemsection

\begin{openproblem}
	9.7. (A. M. Stëpin). Does there exist an infinite finitely generated amenable group of bounded exponent? R. I. Grigorchuk
\end{openproblem}

\plainproblemsection

\begin{openproblem}
	9.8. Does there exist a finitely generated simple group of intermediate growth? R. I. Grigorchuk
	
	Yes, it does (V. Nekrashevych, Ann. Math., 87, no. 3 (2018), 667-719).
\end{openproblem}

\plainproblemsection

\begin{openproblem}
	9.9. Does there exist a finitely generated group which is not nilpotent-by-finite and whose growth function has as a majorant a function of the form ${c}^{\sqrt{n}}$ where $c$ is a constant greater than 1 ? R. I. Grigorchuk
\end{openproblem}

\plainproblemsection

\begin{openproblem}
	9.10. Do there exist finitely generated groups different from $\mathbb{Z}/2\mathbb{Z}$ which have precisely two conjugacy classes? V. S. Guba
	
	Yes, they exist (D. Osin, Ann. Math., 172, no. 1 (2010), 1-39).
\end{openproblem}

\plainproblemsection

\begin{openproblem}
	9.11. Is an abelian minimal normal subgroup $A$ of a group $G$ an elementary $p$ -group if the factor-group $G/A$ is a soluble group of finite rank? The answer is affirmative under the additional hypothesis that $G/A$ is locally polycyclic (D. I. Zaitsev, Algebra and Logic, 19, no. 2 (1980), 94-105). D. I. Zaitsev
\end{openproblem}

\plainproblemsection

\begin{openproblem}
	9.12. Is there a soluble group with the following properties: torsion-free of finite rank, not finitely generated, and having a faithful irreducible representation over a finite field? D. I. Zaitsev
	
	Yes, there is (A. V. Tushev, Ukrain. Math. J., 42 (1990), 1233-1238).
\end{openproblem}

\plainproblemsection

\begin{openproblem}
	9.13. Is a soluble torsion-free group minimax if it satisfies the weak minimum condition for normal subgroups? D. I. Zaitsev
\end{openproblem}

\plainproblemsection

\begin{openproblem}
	9.14. Is a group locally finite if it decomposes into a product of periodic abelian subgroups which commute pairwise? The answer is affirmative in the case of two subgroups. D. I. Zaitsev
\end{openproblem}

\plainproblemsection

\begin{openproblem}
	9.15. Describe, without using CFSG, all subgroups $L$ of a finite Chevalley group $G$ such that $G = {PL}$ for some parabolic subgroup $P$ of $G$ . A. S. Kondratiev
\end{openproblem}

\plainproblemsection

\begin{openproblem}
	9.16. (Well-known problem). The prime graph of a finite group $G$ is the graph with vertex set $\pi \left( G\right)$ and an edge joining $p$ and $q$ if and only if $G$ has an element of order ${pq}$ . Describe all finite Chevalley groups over a field of characteristic 2 whose prime graph is not connected and describe the connected components. A. S. Kondratiev
	
	This was done in (A. S. Kondratiev, Math. USSR-Sb., 67 (1990), 235-247).
\end{openproblem}

\plainproblemsection

\begin{openproblem}
	9.17. a) Let $G$ be a locally normal residually finite group. Can $G$ be embedded in a direct product of finite groups when the factor group $G/\left\lbrack  {G, G}\right\rbrack$ is a direct product of cyclic groups? L. A. Kurdachenko
	
	Not always (L. A. Kurdachenko, Math. Notes, 39 (1986), 273-279).
\end{openproblem}

\plainproblemsection

\begin{openproblem}
	9.17. b) Let $G$ be a locally normal residually finite group. Does there exist a normal subgroup $H$ of $G$ which is embeddable in a direct product of finite groups and is such that $G/H$ is a divisible abelian group? L. A. Kurdachenko
\end{openproblem}

\plainproblemsection

\begin{openproblem}
	9.18. Let ${\mathfrak{S}}_{ * }$ be the smallest normal Fitting class. Are there Fitting classes which are maximal in ${\mathfrak{S}}_{ * }$ (with respect to inclusion)?
	
	No, there are no such classes (N. T. Vorob'yëv, Dokl. Akad. Nauk Belorus. SSR, 35, no. 6 (1991), 485-487 (Russian)).
\end{openproblem}

\plainproblemsection

\begin{openproblem}
	9.19. a) Let $n\left( X\right)$ denote the minimum of the indices of proper subgroups of a group $X$ . A subgroup $A$ of a finite group $G$ is called wide if $A$ is a maximal element by inclusion of the set $\{ X \mid  X$ is a proper subgroup of $G$ and $n\left( X\right)  = n\left( G\right) \}$ . Find all wide subgroups in finite projective special linear, symplectic, orthogonal, and unitary groups.
\end{openproblem}

\plainproblemsection

\begin{openproblem}
	9.19. b) Let $n\left( X\right)$ denote the minimum of the indices of proper subgroups of a group $X$ . A subgroup $A$ of a finite group $G$ is called wide if $A$ is a maximal element by inclusion of the set $\{ X \mid  X$ is a proper subgroup of $G$ and $n\left( X\right)  = n\left( G\right) \}$ . Prove, without using CFSG, that $n\left( {F}_{1}\right)  = \left| {{F}_{1} : 2{F}_{2}}\right|$ , where ${F}_{1}$ and ${F}_{2}$ are the Fischer simple groups and $2{F}_{2}$ is an extension of a group of order 2 by ${F}_{2}$ . V. D. Mazurov This was proved (S. V. Zharov, V. D. Mazurov, in: Matematicheskoye programmirova-niye i prilozheniya, Ekaterinburg, 1995, 96-97 (Russian)).
\end{openproblem}

\plainproblemsection

\begin{openproblem}
	9.21. Let $P$ be a maximal parabolic subgroup of the smallest index in a finite group $G$ of Lie type ${E}_{6},{E}_{7},{E}_{8}$ , or ${}^{2}{E}_{6}$ and let $X$ be a subgroup such that ${PX} = G$ . Is it true that $X = G$ ? V. D. Mazurov
	
	Yes, it is true (mod CFSG) (C. Hering, M. W. Liebeck, J. Saxl, J. Algebra, 106, no. 2 (1987), 517-527).
\end{openproblem}

\plainproblemsection

\begin{openproblem}
	9.23. Let $G$ be a finite group, $B$ a block of characters of $G, D\left( B\right)$ its defect group and $k\left( B\right)$ (respectively, ${k}_{0}\left( B\right)$ ) the number of all irreducible complex characters (of height 0) lying in $B$ . Conjectures:
	
	a) (R. Brauer) $k\left( B\right)  \leq  \left| {D\left( B\right) }\right|$ ; this has been proved for $p$ -soluble groups (D. Gluck, K. Magaard, U. Riese, P. Schmid, J. Algebra, 279 (2004), 694-719);
	
	b) (J. B. Olsson) ${k}_{0}\left( B\right)  \leq  \left| {D\left( B\right)  : D{\left( B\right) }^{\prime }}\right|$ , where $D{\left( B\right) }^{\prime }$ is the derived subgroup of $D\left( B\right)$ ;
	
	c) (R. Brauer) $D\left( B\right)$ is abelian if and only if ${k}_{0}\left( B\right)  = k\left( B\right)$ . V. D. Mazurov
\end{openproblem}

\plainproblemsection

\begin{openproblem}
	9.24. (J. G. Thompson). Conjecture: every finite simple non-abelian group $G$ can be represented in the form $G = {CC}$ , where $C$ is some conjugacy class of $G$ .
	
	V. D. Mazurov
\end{openproblem}

\plainproblemsection

\begin{openproblem}
	9.25. Find an algorithm which recognizes, by an equation $w\left( {{x}_{1},\ldots ,{x}_{n}}\right)  = 1$ in a free group $F$ and by a list of finitely generated subgroups ${H}_{1},\ldots ,{H}_{n}$ of $F$ , whether there is a solution of this equation satisfying the condition ${x}_{1} \in  {H}_{1},\ldots ,{x}_{n} \in  {H}_{n}$ .
	
	G. S. Makanin
	
	Such an algorithm is found in (V. Diekert, C. Gutiérrez, C. Hagenah, Inf. Comput., 202, no. 2 (2005), 105-140).
\end{openproblem}

\plainproblemsection

\begin{openproblem}
	9.26. a) Describe the finite groups of 2-local 3-rank 1 which have 3-rank at least 3. Described in (A. A. Makhnëv, Siber. Math. J., 29 (1988), 951-959). A. A. Makhnëv
\end{openproblem}

\plainproblemsection

\begin{openproblem}
	9.26. b) Describe the finite groups of 2-local 3-rank 1 which have non-cyclic Sylow 3-subgroups. A. A. Makhnëv
\end{openproblem}

\plainproblemsection

\begin{openproblem}
	9.27. Let $M$ be a subgroup of a finite group $G, A$ an abelian 2-subgroup of $M$ , and suppose that ${A}^{g}$ is not contained in $M$ for some $g$ from $G$ . Determine the structure of $G$ under the hypothesis that $\left\langle  {A,{A}^{x}}\right\rangle   = G$ whenever the subgroup ${A}^{x}, x \in  G$ , is not contained in $M$ . A. A. Makhnëv
	
	It is described mod CFSG (V. I. Zenkov, Algebra and Logic, 35, no. 3 (1996), 160-163).
\end{openproblem}

\plainproblemsection

\begin{openproblem}
	9.28. Suppose that a finite group $G$ is generated by a conjugacy class $D$ of involutions and let ${D}_{i} = \left\{  {{d}_{1}\cdots {d}_{i} \mid  {d}_{1},\ldots ,{d}_{i}\text{are different pairwise commuting elements of}D}\right\}$ . What is $G$ , if ${D}_{1},\ldots ,{D}_{n}$ are all its different conjugacy classes of involutions? For example, the Fischer groups ${F}_{22}$ and ${F}_{23}$ satisfy this condition with $n = 3$ .
	
	A. A. Makhnëv
\end{openproblem}

\plainproblemsection

\begin{openproblem}
	9.29. (Well-known problem). According to a classical theorem of Magnus, the word problem is soluble in 1-relator groups. Do there exist 2-relator groups with insoluble word problem? Yu. I. Merzlyakov
\end{openproblem}

\plainproblemsection

\begin{openproblem}
	9.30. (Well-known problem). A finite set of reductions ${u}_{i} \rightarrow  {v}_{i}$ of words on a finite alphabet $\sum  = {\sum }^{-1}$ is called a group set of reductions if length $\left( {u}_{i}\right)  >$ length $\left( {v}_{i}\right)$ or length $\left( {u}_{i}\right)  =$ length $\left( {v}_{i}\right)$ and ${u}_{i} > {v}_{i}$ in the lexicographical ordering, and every word in $\sum$ can be reduced to the unique reduced form which does not depend on the sequence of reductions. Do there exist group sets of reductions satisfying the condition length $\left( {v}_{i}\right)  \leq  1$ for all $i$ , which are different from 1 ) sets of trivial reductions ${x}^{-\varepsilon }{x}^{\varepsilon } \rightarrow  1,\varepsilon  =  \pm  1,2$ ) multiplication tables ${xy} \rightarrow  z$ of finite groups, and 3) their finite unions? Yu. I. Merzlyakov
	
	Yes, they exist (J. Avenhaus, K. Madlener, F. Otto, Trans. Amer. Math. Soc., 297 (1986), 427-443).
\end{openproblem}

\plainproblemsection

\begin{openproblem}
	9.31. Let $k$ be a field of characteristic 0 . According to (Yu. I. Merzlyakov, Proc. Stek-lov Inst. Math., $\mathbf{{167}}\left( {1986}\right) ,{263} - {266}$ ), the family ${\operatorname{Rep}}_{k}\left( G\right)$ of all canonical matrix representations of a $k$ -powered group $G$ over $k$ may be regarded as an affine $k$ -variety. Find an explicit form of equations defining this variety. Yu. I. Merzlyakov
\end{openproblem}

\plainproblemsection

\begin{openproblem}
	*9.32. What locally compact groups satisfy the following condition: the product of any two closed subgroups is also a closed subgroup? Abelian groups with this property were described in (Yu. N. Mukhin, Math. Notes, 8 (1970), 755-760).
	
	*Such groups are described (W. Herfort, K. H. Hofmann, F. G. Russo, Adv. Math., 390 (2021), 107894). Yu. N. Mukhin
\end{openproblem}

\plainproblemsection

\begin{openproblem}
	9.33. (F. Kümmich, H. Scheerer). If $H$ is a closed subgroup of a connected locally-compact group $G$ such that $\overline{HX} = \bar{H}\bar{X}$ for every closed subgroup $X$ of $G$ , then is $H$ normal? Yu. N. Mukhin
	
	Yes, it is (C. Scheiderer, Monatsh. Math., 98 (1984), 75-81).
\end{openproblem}

\plainproblemsection

\begin{openproblem}
	9.34. (S. K. Grosser, W. N. Herfort). Does there exist an infinite compact $p$ -group in which the centralizers of all elements are finite?
	
	No, it does not exist, in view of the connection of this problem with the Restricted Burnside Problem (S. K. Grosser, W. N. Herfort, Trans. Amer. Math. Soc., 283, no. 1 (1984), 211-224) and because of the positive solution to the latter (E. I. Zel'manov, Math. USSR-Izv., 36 (1991), 41-60; Math. USSR-Sb., 72 (1992), 543-565).
\end{openproblem}

\plainproblemsection

\begin{openproblem}
	9.35. A topological group is said to be inductively compact if any finite set of its elements is contained in a compact subgroup. Is this property preserved under lattice isomorphisms in the class of locally compact groups? Yu. N. Mukhin
\end{openproblem}

\plainproblemsection

\begin{openproblem}
	9.36. Characterize the lattices of closed subgroups in locally compact groups. In the discrete case this was done in (B. V. Yakovlev, Algebra and Logic, 13, no. 6 (1974), 400-412). Yu. N. Mukhin
\end{openproblem}

\plainproblemsection

\begin{openproblem}
	9.37. Is a compactly generated inductively prosoluble locally compact group pro-soluble? Yu. N. Mukhin
\end{openproblem}

\plainproblemsection

\begin{openproblem}
	9.38. A group is said to be compactly covered if it is the union of its compact subgroups. In a null-dimensional locally compact group, are maximal compactly covered subgroups closed? Yu. N. Mukhin
\end{openproblem}

\plainproblemsection

\begin{openproblem}
	9.39. Let $\Omega$ be a countable set and $\mathfrak{m}$ a cardinal number such that ${\aleph }_{0} \leq  \mathfrak{m} \leq  {2}^{{\aleph }_{0}}$ (we assume Axiom of Choice but not Continuum Hypothesis). Does there exist a permutation group $G$ on $\Omega$ that has exactly $\mathfrak{m}$ orbits on the power set $\mathfrak{P}\left( \Omega \right)$ ?
	
	Comment of 2001: It is proved (S. Shelah, S. Thomas, Bull. London Math. Soc., 20, no. 4 (1988), 313-318) that the answer is positive in set theory with Martin's Axiom. The question is still open in ZFC. Peter M. Neumann
\end{openproblem}

\plainproblemsection

\begin{openproblem}
	9.40. Let $\Omega$ be a countably infinite set. Define a moiety of $\Omega$ to be a subset $\sum$ such that both $\sum$ and $\Omega  \smallsetminus  \sum$ are infinite. Which permutation groups on $\Omega$ are transitive on moieties? Peter M. Neumann
\end{openproblem}

\plainproblemsection

\begin{openproblem}
	9.41. a) Let $\Omega$ be a countably infinite set. For $k \geq  2$ , we define a $k$ -section of $\Omega$ to be a partition of $\Omega$ into a union of $k$ infinite subsets. Then does there exist a transitive permutation group on $\Omega$ that is transitive on $k$ -sections but intransitive on ordered $k$ -sections? P. M. Neumann
	
	Yes, there does. Let $U$ be a non-principal ultrafilter in $\mathfrak{P}\left( \Omega \right)$ and let $G = \{ g \in$ $\operatorname{Sym}\left( \Omega \right)  \mid  \operatorname{Fix}\left( g\right)  \in  U\}$ . It is not hard to prove that $G$ is transitive on $k$ -sections but not on ordered $k$ -sections for any $k$ in the range $2 \leq  k \leq  {\aleph }_{0}$ . (P. M. Neumann, Letter of October, 5, 1989.)
\end{openproblem}

\plainproblemsection

\begin{openproblem}
	9.41. Let $\Omega$ be a countably infinite set. For $k \geq  2$ define a $k$ -section of $\Omega$ to be a partition of $\Omega$ as union of $k$ infinite sets.
	
	b) Does there exist a group that is transitive on $k$ -sections but not on $\left( {k + 1}\right)$ - sections?
	
	*c) Does there exist a transitive permutation group on $\Omega$ that is transitive on ${\aleph }_{0}$ -sections but which is a proper subgroup of $\operatorname{Sym}\left( \Omega \right)$ ?
	
	*c) An affirmative answer is consistent with the ZFC axioms of set theory (S. M. Corson, S. Shelah, Preprint, 2025, https://arxiv.org/abs/2503.12997).
\end{openproblem}

\plainproblemsection

\begin{openproblem}
	*9.42. Let $\Omega$ be a countable set and let $D$ be the set of total order relations on $\Omega$ for which $\Omega$ is order-isomorphic with $\mathbb{Q}$ . Does there exist a transitive proper subgroup $G$ of $\operatorname{Sym}\left( \Omega \right)$ which is transitive on $D$ ? Peter M. Neumann
	
	*An affirmative answer is consistent with the ZFC axioms of set theory (S. M. Corson, S. Shelah, Preprint, 2025, https://arxiv.org/abs/2503.12997).
\end{openproblem}

\plainproblemsection

\begin{openproblem}
	9.43. a) The group $G$ described in the solution of 8.73 (Archive) enables us to construct a projective plane of order 3 in which the lines are the elements of any conjugacy class of subgroups of order $2 \cdot  5 \cdot  7 \cdot  {11}$ together with four lines added in a natural way. In a similar way, we can construct a projective plane of order ${p}^{n}$ for any prime $p$ and any positive integer $n$ . Does the resulting projective plane have the Galois property? N. D. Podufalov
	
	Not always. The answer is affirmative for $n = 1$ . But for $n > 1$ the planes in question may not have the Galois property. For example, if there exists a near-field of ${p}^{n}$ elements, then among the planes of order ${p}^{n}$ indicated there will definitely be some non-Desarguesian planes. (N. D. Podufalov, Letter of November, 24, 1986.)
\end{openproblem}

\plainproblemsection

\begin{openproblem}
	9.43. b) The group $G$ indicated in (N.D. Podufalov, Abstracts of the 9th All-Union Symp. on Group Theory, Moscow, 1984, 113-114 (Russian)) allows us to construct a projective plane of order 3: one can take as lines any class of subgroups of order $2 \cdot  5 \cdot  {11} \cdot  {17}$ conjugate under $S$ and add four more lines in a natural way. In a similar way, one can construct projective planes of order ${p}^{n}$ for any prime $p$ and any natural $n$ . Could this method be adapted for constructing new planes? N. D. Podufalov
\end{openproblem}

\plainproblemsection

\begin{openproblem}
	9.44. A topological group is said to be layer compact if the full inverse images of all of its compacts under mappings $x \rightarrow  {x}^{n}, n = 1,2,\ldots$ , are compacts. Describe the locally compact locally soluble layer compact groups. V. M. Poletskikh
\end{openproblem}

\plainproblemsection

\begin{openproblem}
	9.45. Let $a$ be a vector in ${\mathbb{R}}^{n}$ with rational coordinates and set
	
	$$
	S\left( a\right)  = \left\{  {{ka} + b \mid  k \in  \mathbb{Z}, b \in  {\mathbb{Z}}^{n}}\right\}  .
	$$
	
	It is obvious that $S\left( a\right)$ is a discrete subgroup in ${\mathbb{R}}^{n}$ of rank $n$ . Find necessary and sufficient conditions in terms of the coordinates of $a$ for $S\left( a\right)$ to have an orthogonal basis with respect to the standard scalar product in ${\mathbb{R}}^{n}$ .
	
	Yu. D. Popov, I. V. Protasov
\end{openproblem}

\plainproblemsection

\begin{openproblem}
	9.46. Let $G$ be a locally compact group of countable weight and $L\left( G\right)$ the space of all its closed subgroups equipped with the $E$ -topology. Then is $L\left( G\right)$ a $k$ -space? I. V. Protasov
	
	Yes (I. V. Protasov, Dokl. Akad. Nauk Ukr. SSR Ser. A, 10 (1986), 64-66 (Russian)).
\end{openproblem}

\plainproblemsection

\begin{openproblem}
	9.47. Is it true that every scattered compact can be homeomorphically embedded into the space of all closed non-compact subgroups with $E$ -topology of a suitable locally compact group? I. V. Protasov
\end{openproblem}

\plainproblemsection

\begin{openproblem}
	9.48. In a null-dimensional locally compact group, is the set of all compact elements closed? I. V. Protasov
	
	Yes, it is (G. A. Willis, Math. Ann., 300 (1994), 341-363).
\end{openproblem}

\plainproblemsection

\begin{openproblem}
	*9.49. Let $G$ be a compact group of weight $> {\omega }_{2}$ . Is it true that the space of all closed subgroups of $G$ with respect to $E$ -topology is non-dyadic?
	
	I. V. Protasov, Yu. V. Tsybenko
	
	*Yes, it is true (Yu. V. Tsybenko, Ukrain. Math. J., 38 (1986), 542-545).
\end{openproblem}

\plainproblemsection

\begin{openproblem}
	9.50. Is every 4-Engel group
	
	a) without elements of order 2 and 5 necessarily soluble?
	
	b) satisfying the identity ${x}^{5} = 1$ necessarily locally finite?
	
	c) (R. I. Grigorchuk) satisfying the identity ${x}^{8} = 1$ necessarily locally finite?
	
	Yu. P. Razmyslov
	
	a) Yes; this follows from the local nilpotency of 4-Engel groups (G. Traustason, Int. J. Algebra Comput., 15 (2005), 309-316; G. Havas, M. R. Vaughan-Lee, Int. J. Algebra Comput. 15 (2005), 649-682) and the affirmative answer for locally nilpotent groups in (A. Abdollahi, G. Traustason, Proc. Amer. Math. Soc., 130 (2002), 2827-2836). b) Yes, it is (M. R. Vaughan-Lee, Proc. London Math. Soc. (3), 74 (1997), 306-334). c) Yes; moreover, every 4-Engel 2-group is locally finite (G. Traustason, J. Algebra, 178, no. 2 (1995), 414-429). 9.54. If $G = {HK}$ is a soluble group and $H, K$ are minimax groups, is it true that $G$ is a minimax group? D. J. S. Robinson
	
	Yes, it is true (J. S. Wilson, J. Pure Appl. Algebra, 53, no. 3 (1988), 297-318; Ya. P. Sysak, Radical modules over groups of finite rank, Inst. of Math. Acad. Sci. Ukrain. SSR, Kiev, 1989 (Russian)).
\end{openproblem}

\plainproblemsection

\begin{openproblem}
	9.51. Does there exist a finitely presented soluble group satisfying the maximum condition on normal subgroups which has insoluble word problem? D. J. S. Robinson
\end{openproblem}

\plainproblemsection

\begin{openproblem}
	9.52. Does a finitely presented soluble group of finite rank have soluble conjugacy problem? Note: there is an algorithm to decide conjugacy to a given element of the group. D. J. S. Robinson
\end{openproblem}

\plainproblemsection

\begin{openproblem}
	9.53. Is the isomorphism problem soluble for finitely presented soluble groups of finite rank? D. J. S. Robinson
\end{openproblem}

\plainproblemsection

\begin{openproblem}
	9.55. Does there exist a finite $p$ -group $G$ and a central augmented automorphism $\varphi$ of $\mathbb{Z}G$ such that $\varphi$ , if extended to ${\widehat{\mathbb{Z}}}_{p}G$ , the $p$ -adic group ring, is not conjugation by unit in ${\widehat{\mathbb{Z}}}_{p}G$ followed by a homomorphism induced from a group automorphism?
	
	K. W. Roggenkamp
\end{openproblem}

\plainproblemsection

\begin{openproblem}
	9.56. Find all finite groups with the property that the tensor square of any ordinary irreducible character is multiplicity free. J. Saxl
	
	Such groups are shown to be soluble (L. S. Kazarin, E. I. Chankov, Sb. Math., 201, no. 5 (2010), 655-668), and abundance of examples shows that a classification of such soluble groups is not feasible.
\end{openproblem}

\plainproblemsection

\begin{openproblem}
	9.57. The set of subformations of a given formation is a lattice with respect to operations of intersection and generation. What formations of finite groups have distributive lattices of subformations? A. N. Skiba
\end{openproblem}

\plainproblemsection

\begin{openproblem}
	9.58. Can a product of non-local formations of finite groups be local?
	
	A. N. Skiba, L. A. Shemetkov
	
	Yes, it can (V. A. Vedernikov, Math. Notes, 46 (1989), 910-913; N. T. Vorob'ev, in: Proc. 7th Regional Sci. Session Math., Zielona Goru, 1989, 79-86).
\end{openproblem}

\plainproblemsection

\begin{openproblem}
	9.59. (W. Gaschütz). Prove that the formation generated by a finite group has finite lattice of subformations. A. N. Skiba, L. A. Shemetkov
	
	Counterexamples are constructed (V. P. Burichenko, J. Algebra, 372 (2012), 428-458).
\end{openproblem}

\plainproblemsection

\begin{openproblem}
	9.60. Let $\mathfrak{F}$ and $\mathfrak{H}$ be local formations of finite groups and suppose that $\mathfrak{F}$ is not contained in $\mathfrak{H}$ . Does $\mathfrak{F}$ necessarily have at least one minimal local non- $\mathfrak{H}$ -subformation? A. N. Skiba, L. A. Shemetkov
	
	No, not necessarily (V.P. Burichenko, Trudy Inst. Math. National Akad. Sci. Be-larus', 21, no. 1 (2013), 15-24 (Russian)).
\end{openproblem}

\plainproblemsection

\begin{openproblem}
	9.61. Two varieties are said to be $S$ -equivalent if they have the same Mal’cev theory (D. M. Smirnov, Algebra and Logic, 22, no. 6 (1983), 492-501). What is the cardinality of the set of $S$ -equivalent varieties of groups? D. M. Smirnov
\end{openproblem}

\plainproblemsection

\begin{openproblem}
	9.62. In any group $G$ , the cosets of all of its normal subgroups together with the empty set form the block lattice $C\left( G\right)$ with respect to inclusion, which, for infinite $\left| G\right|$ , is subdirectly irreducible and, for finite $\left| G\right|  \geq  3$ , is even simple (D. M. Smirnov, A. V. Reibol'd, Algebra and Logic, 23, no. 6 (1984), 459-470). How large is the class of such lattices? Is every finite lattice embeddable in the lattice $C\left( G\right)$ for some finite group $G$ ? D. M. Smirnov
	
	For each $g \in  G$ the filter $\{ x \in  C\left( G\right)  \mid  g \in  x\}$ is modular (and even arguesian); the variety generated by the block lattices of groups does not contain all finite lattices (D. M. Smirnov, Siberian Math. J., 33, no. 4 (1992), 663-668).
\end{openproblem}

\plainproblemsection

\begin{openproblem}
	9.63. Is a finite group of the form $G = {ABA}$ soluble if $A$ is an abelian subgroup and $B$ is a cyclic subgroup? Ya. P. Sysak
	
	Yes, it is soluble (mod CFSG) (D. L. Zagorin, L. S. Kazarin, Dokl. Math., 53, no. 2 (1996), 237-239).
\end{openproblem}

\plainproblemsection

\begin{openproblem}
	9.64. Is it true that, in a group of the form $G = {AB}$ , every subgroup $N$ of $A \cap  B$ which is subnormal both in $A$ and in $B$ is subnormal in $G$ ? The answer is affirmative in the case of finite groups (H. Wielandt). Ya. P. Sysak
	
	No, not always (C. Casolo, U. Dardano, J. Group Theory, 7 (2004), 507-520).
\end{openproblem}

\plainproblemsection

\begin{openproblem}
	9.65. Is a locally soluble group periodic if it is a product of two periodic subgroups? Ya. P. Sysak
\end{openproblem}

\plainproblemsection

\begin{openproblem}
	9.66. a) B. Jonsson's Conjecture: elementary equivalence is preserved under taking free products in the class of all groups, that is, if ${Th}\left( {G}_{1}\right)  = {Th}\left( {G}_{2}\right)$ and ${Th}\left( {H}_{1}\right)  =$ ${Th}\left( {H}_{2}\right)$ for groups ${G}_{1},{G}_{2},{H}_{1},{H}_{2}$ , then ${Th}\left( {{G}_{1} * {H}_{1}}\right)  = {Th}\left( {{G}_{2} * {H}_{2}}\right)$ .
	
	b) It may be interesting to consider also the following weakened conjecture: if $\operatorname{Th}\left( {G}_{1}\right)  = \operatorname{Th}\left( {G}_{2}\right)$ and $\operatorname{Th}\left( {H}_{1}\right)  = \operatorname{Th}\left( {H}_{2}\right)$ for countable groups ${G}_{1},{G}_{2},{H}_{1},{H}_{2}$ , then for any numerations of the groups ${G}_{i},{H}_{i}$ there are $m$ -reducibility ${T}_{1} \equiv  {T}_{2}$ and Turing reducibility ${T}_{1}\underset{T}{ \equiv  }{T}_{2}$ , where ${T}_{i}$ is the set of numbers of all theorems in $\operatorname{Th}\left( {{G}_{i} * {H}_{i}}\right)$ . A proof of this weakened conjecture would be a good illustration of application of the reducibility theory to solving concrete mathematical problems. A. D. Taimanov
\end{openproblem}

\plainproblemsection

\begin{openproblem}
	9.67. (A. Tarski). Let ${F}_{n}$ be a free group of rank $n$ ; is it true that $\operatorname{Th}\left( {F}_{2}\right)  = \operatorname{Th}\left( {F}_{3}\right)$ ? A. D. Taimanov
	
	Yes, it is (O. Kharlampovich, A. Myasnikov, J. Algebra, 302, (2006), 451-552; Z. Sela, Geom. Funct. Anal., 16 (2006), 707-730).
\end{openproblem}

\plainproblemsection

\begin{openproblem}
	9.68. Let $\mathfrak{V}$ be a variety of groups which is not the variety of all groups, and let $p$ be a prime. Is there a bound on the $p$ -lengths of the finite $p$ -soluble groups whose Sylow $p$ -subgroups are in $\mathfrak{V}$ ? J. S. Wilson
\end{openproblem}

\plainproblemsection

\begin{openproblem}
	9.69. (P. Cameron). Let $G$ be a finite primitive permutation group and suppose that the stabilizer ${G}_{\alpha }$ of a point $\alpha$ induces on some of its orbits $\Delta  \neq  \{ \alpha \}$ a regular permutation group. Is it true that $\left| {G}_{\alpha }\right|  = \left| \Delta \right|$ ? A. N. Fomin
\end{openproblem}

\plainproblemsection

\begin{openproblem}
	9.70. Prove or disprove the conjecture of P. Cameron (Bull. London Math. Soc., 13, no. 1 (1981),1-22): if $G$ is a finite primitive permutation group of subrank $m$ , then either the rank of $G$ is bounded by a function of $m$ or the order of a point stabilizer is at most $m$ . The subrank of a transitive permutation group is defined to be the maximum rank of the transitive constituents of a point stabilizer. A. N. Fomin
\end{openproblem}

\plainproblemsection

\begin{openproblem}
	9.71. Is it true that every infinite 2-transitive permutation groups with locally soluble point stabilizer has a non-trivial irreducible finite-dimensional representation over some field? A. N. Fomin
\end{openproblem}

\plainproblemsection

\begin{openproblem}
	9.72. Let ${G}_{1}$ and ${G}_{2}$ be Lie groups with the following property: each ${G}_{i}$ contains a nilpotent simply connected normal Lie subgroup ${B}_{i}$ such that ${G}_{i}/{B}_{i} \cong  S{L}_{2}\left( K\right)$ , where $K = \mathbb{R}$ or $\mathbb{C}$ . Assume that ${G}_{1}$ and ${G}_{2}$ are contained as closed subgroups in a topological group $G$ , that ${G}_{1} \cap  {G}_{2} \geq  {B}_{1}{B}_{2}$ , and that no non-identity Lie subgroup of ${B}_{1} \cap  {B}_{2}$ is normal in $G$ . Can it then be shown (perhaps by using the method of "amalgams" from the theory of finite groups) that the nilpotency class and the dimension of ${B}_{1}{B}_{2}$ is bounded? A. L. Chermak
\end{openproblem}

\plainproblemsection

\begin{openproblem}
	9.73. Let $\mathfrak{F}$ be any local formation of finite soluble groups containing all finite nilpotent groups. Prove that ${H}^{\mathfrak{F}}K = K{H}^{\mathfrak{F}}$ for any two subnormal subgroups $H$ and $K$ of an arbitrary finite group $G$ . L. A. Shemetkov
	
	This has been proved, even without the hypothesis of solubility of $\mathfrak{F}$ (S. F. Kamorni-kov, Dokl. Akad. Nauk BSSR, 33, no. 5 (1989), 396-399 (Russian)).
\end{openproblem}

\plainproblemsection

\begin{openproblem}
	9.74. Find all local formations $\mathfrak{F}$ of finite groups such that every finite minimal non- $\mathfrak{F}$ -group is either a Shmidt group (that is, a non-nilpotent finite group all of whose proper subgroups are nilpotent) or a group of prime order. L. A. Shemetkov They are found (S. F. Kamornikov, Siberian Math. J., 35, no. 4 (1994), 713-721).
\end{openproblem}

\plainproblemsection

\begin{openproblem}
	9.75. Find all local formations $\mathfrak{F}$ of finite groups such that, in every finite group, the set of $\mathfrak{F}$ -subnormal subgroups forms a lattice.
	
	Comment of 2015: Subgroup-closed formations with this property are described in (Xiaolan Yi, S. F. Kamornikov, J. Algebra, 444 (2015), 143-151. L. A. Shemetkov
\end{openproblem}

\plainproblemsection

\begin{openproblem}
	9.76. We define a Golod group to be the $r$ -generated, $r \geq  2$ , subgroup $\left\langle  {1 + {x}_{1} + I}\right.$ , $\left. {1 + {x}_{2} + I,\ldots ,1 + {x}_{r} + I}\right\rangle$ of the adjoint group $1 + F/I$ of the factor-algebra $F/I$ , where $F$ is a free algebra of polynomials without constant terms in non-commuting variables ${x}_{1},{x}_{2},\ldots ,{x}_{r}$ over a field of characteristic $p > 0$ , and $I$ is an ideal of $F$ such that $F/I$ is a non-nilpotent nil-algebra (see E. S. Golod, Amer. Math. Soc. Transl. (2), 48 (1965), 103-106). Prove that in Golod groups the centralizer of every element is infinite.
	
	Note that Golod groups with infinite centre were constructed by A. V. Timofeenko (Math. Notes, 39, no. 5 (1986), 353-355); independently by other methods the same result was obtained in the 90s by L. Hammoudi. V. P. Shunkov
\end{openproblem}

\plainproblemsection

\begin{openproblem}
	9.77. Does there exist an infinite finitely generated residually finite binary finite group all of whose Sylow subgroups are finite? A. V. Rozhkov (Dr. of Sci. Disser., 1997) showed that such a group does exist if the condition of finiteness of the Sylow subgroups is weakened to local finiteness. V. P. Shunkov
\end{openproblem}

\plainproblemsection

\begin{openproblem}
	9.78. Does there exist a periodic residually finite ${F}^{ * }$ -group (see Archive,7.42) all of whose Sylow subgroups are finite and which is not binary finite? V. P. Shunkov
\end{openproblem}

\plainproblemsection

\begin{openproblem}
	9.79. (A. G. Kurosh). Is every group with the minimum condition countable? No (V. N. Obraztsov, Math. USSR-Sb., 66 (1989), 541-553). V. P. Shunkov
\end{openproblem}

\plainproblemsection

\begin{openproblem}
	9.80. Are the 2-elements of a group with the minimum condition contained in its locally finite radical? V. P. Shunkov
	
	Not necessarily (A. Yu. Ol'shanskii, The geometry of defining relations in groups, Kluwer, Dordrecht, 1991).
\end{openproblem}

\plainproblemsection

\begin{openproblem}
	9.81. Does there exist a simple group with the minimum condition possessing a non-trivial quasi-cyclic subgroup? V. P. Shunkov
	
	Yes, there exists (V. N. Obraztsov, Math. USSR-Sb., 66 (1989), 541-553).
\end{openproblem}

\plainproblemsection

\begin{openproblem}
	9.82. An infinite group, all of whose proper subgroups are finite, is called quasi-finite. Is it true that an element of a quasi-finite group $G$ is central if and only if it is contained in infinitely many subgroups of $G$ ? V. P. Shunkov
	
	No (K. I. Lossov, Dep. no. 5529-V89, VINITI, Moscow, 1988 (Russian)).
\end{openproblem}

\plainproblemsection

\begin{openproblem}
	9.83. Suppose that $G$ is a (periodic) $p$ -conjugacy biprimitively finite group (see 6.59) which has a finite Sylow $p$ -subgroup. Is it then true that all Sylow $p$ -subgroups of $G$ are conjugate? V. P. Shunkov
\end{openproblem}

\plainproblemsection

\begin{openproblem}
	9.84. a) Is every binary finite 2-group of finite exponent locally finite?
	
	b) The same question for $p$ -groups for $p > 2$ . V. P. Shunkov
\end{openproblem}

\plainproblemsection

\begin{openproblem}
	10.1. Let $p$ be a prime number. Describe the groups of order ${p}^{9}$ of nilpotency class 2 which contain subgroups $X$ and $Y$ such that $\left| X\right|  = \left| Y\right|  = {p}^{3}$ and any non-identity elements $x \in  X, y \in  Y$ do not commute. An answer to this question would yield a description of the semifields of order ${p}^{3}$ . They are described (V. A. Antonov, Preprint, Chelyabinsk, 1999 (Russian)).
\end{openproblem}

\plainproblemsection

\begin{openproblem}
	10.2. A mixed identity of a group $G$ is, by definition, an identity of an algebraic system obtained from $G$ by supplementing its signature by some set of 0 -ary operations. One can develop a theory of mixed varieties of groups on the basis of this notion (see, for example, V. S. Anashin, Math. USSR Sbornik, 57 (1987), 171-182). Construct an example of a class of groups which is not a mixed variety, but which is closed under taking factor-groups, Cartesian products, and those subgroups of Cartesian powers that contain the diagonal subgroup. V. S. Anashin
\end{openproblem}

\plainproblemsection

\begin{openproblem}
	10.3. Characterize (in terms of bases of mixed identities or in terms of generating groups) minimal mixed varieties of groups. V. S. Anashin
\end{openproblem}

\plainproblemsection

\begin{openproblem}
	10.4. Let $p$ be a prime number. Is it true that a mixed variety of groups generated by an arbitrary finite $p$ -group of sufficiently large nilpotency class is a variety of groups?
	
	V. S. Anashin
\end{openproblem}

\plainproblemsection

\begin{openproblem}
	10.5. Construct an example of a finite group whose mixed identities do not have a finite basis. Is the group constructed in (R. Bryant, Bull. London Math. Soc., 14, no. 2 (1982), 119-123) such an example? V. S. Anashin
\end{openproblem}

\plainproblemsection

\begin{openproblem}
	10.6. Is it true that in an abelian group every non-discrete group topology can be strengthened up to a non-discrete group topology such that the group becomes a complete topological group? V. I. Arnautov
	
	This is true for group topologies satisfying the first axiom of countability (E. I. Marin, in: Modules, Algebras, Topologies (Mat. Issledovaniya, 105), Kishinëv, 1988, 105-119 (Russian)), but this may not be true in general, see Archive, 12.2.
\end{openproblem}

\plainproblemsection

\begin{openproblem}
	10.7. Is it true that in a countable group $G$ , any non-discrete group topology satisfying the first axiom of countability can be strengthened up to a non-discrete group topology such that $G$ becomes a complete topological group? V. I. Arnautov Yes, it is (V.I. Arnautov, E. I. Kabanova, Siberian Math. J., 31 (1990), 1-10).
\end{openproblem}

\plainproblemsection

\begin{openproblem}
	10.8. Does there exist a topological group which cannot be embedded in the multiplicative semigroup of a topological ring? V. I. Arnautov, A. V. Mikhalëv
\end{openproblem}

\plainproblemsection

\begin{openproblem}
	10.9. Let $p$ be a prime number and let ${L}_{p}$ denote the set of all quasivarieties each of which is generated by a finite $p$ -group. Is ${L}_{p}$ a sublattice of the lattice of all quasivarieties of groups? A. I. Budkin
	
	No, not always (S. A. Shakhova, Math. Notes, 53, no. 3 (1993), 345-347).
\end{openproblem}

\plainproblemsection

\begin{openproblem}
	10.10. Is it true that a quasivariety generated by a finitely generated torsion-free soluble group and containing a non-abelian free metabelian group, can be defined in the class of torsion-free groups by an independent system of quasiidentities?
	
	A. I. Budkin
\end{openproblem}

\plainproblemsection

\begin{openproblem}
	10.11. Is it true that every finitely presented group contains either a free subsemi-group on two generators or a nilpotent subgroup of finite index? R. I. Grigorchuk
\end{openproblem}

\plainproblemsection

\begin{openproblem}
	10.12. Does there exist a finitely generated semigroup $S$ with cancellation having non-exponential growth and such that its group of left quotients $G = {S}^{-1}S$ (which exists) is a group of exponential growth? An affirmative answer would give a positive solution to Problem 12 in (S. Wagon, The Banach-Tarski Paradox, Cambridge Univ. Press, 1985). R. I. Grigorchuk 10.13. Does there exist a massive set of independent elements in a free group ${F}_{2}$ on free generators $a, b$ , that is, a set $E$ of irreducible words on the alphabet $a, b,{a}^{-1},{b}^{-1}$ such that
	
	1) no element $w \in  E$ belongs to the normal closure of $E \smallsetminus  \{ w\}$ , and
	
	2) the massiveness condition is satisfied: $\mathop{\lim }\limits_{{n \rightarrow  \infty }}\sqrt[n]{\left| {E}_{n}\right| } = 3$ , where ${E}_{n}$ is the set of all words of length $n$ in $E$ ? R. I. Grigorchuk
\end{openproblem}

\plainproblemsection

\begin{openproblem}
	10.14. a) Does every group satisfying the minimum condition on subgroups satisfy the weak maximum condition on subgroups?
	
	b) Does every group satisfying the maximum condition on subgroups satisfy the weak minimum condition on subgroups? D. I. Zaitsev a) No, not always (V. N. Obraztsov, Math. USSR-Sb., 66, no. 2 (1990), 541-553). b) No, not always (V. N. Obraztsov, Siberian Math. J., 32, no. 1 (1991), 79-84).
\end{openproblem}

\plainproblemsection

\begin{openproblem}
	10.15. For every (known) finite quasisimple group and every prime $p$ , find the faithful $p$ -modular absolutely irreducible linear representations of minimal degree.
	
	A. S. Kondratiev
\end{openproblem}

\plainproblemsection

\begin{openproblem}
	10.16. A class of groups is called a direct variety if it is closed under taking subgroups, factor-groups, and direct products (Yu. M. Gorchakov, Groups with Finite Classes of Conjugate Elements, Moscow, Nauka, 1978 (Russian)). It is obvious that the class of ${FC}$ -groups is a direct variety. P. Hall (J. London Math. Soc.,34, no. 3 (1959), 289-304) showed that the class of finite groups and the class of abelian groups taken together do not generate the class of ${FC}$ -groups as a direct variety, and it was shown in (L. A. Kurdachenko, Ukrain. Math. J., 39, no. 3 (1987), 255-259) that the direct variety of ${FC}$ -groups is also not generated by the class of groups with finite derived subgroups. Is the direct variety of ${FC}$ -groups generated by the class of groups with finite derived subgroups together with the class of ${FC}$ -groups having quasicyclic derived subgroups? L. A. Kurdachenko
\end{openproblem}

\plainproblemsection

\begin{openproblem}
	10.17. (M. J. Tomkinson). Let $G$ be an ${FC}$ -group whose derived subgroup is embeddable in a direct product of finite groups. Must $G/Z\left( G\right)$ be embeddable in a direct product of finite groups? L. A. Kurdachenko
\end{openproblem}

\plainproblemsection

\begin{openproblem}
	10.18. (M.J. Tomkinson). Let $G$ be an ${FC}$ -group which is residually in the class of groups with finite derived subgroups. Must $G/Z\left( G\right)$ be embeddable in a direct product of finite groups? L. A. Kurdachenko
\end{openproblem}

\plainproblemsection

\begin{openproblem}
	10.19. Characterize the radical associative rings such that the set of all normal subgroups of the adjoint group coincides with the set of all ideals of the associated Lie ring. V. M. Levchuk
\end{openproblem}

\plainproblemsection

\begin{openproblem}
	10.20. In a Chevalley group of rank $\leq  6$ over a finite field of order $\leq  9$ , describe all subgroups of the form $H = \langle H \cap  U, H \cap  V\rangle$ that are not contained in any proper parabolic subgroup, where $U$ and $V$ are opposite unipotent subgroups.
	
	V. M. Levchuk
\end{openproblem}

\plainproblemsection

\begin{openproblem}
	10.23. Is it true that extraction of roots in braid groups is unique up to conjugation? G. S. Makanin
	
	Yes, it is true (J. González-Meneses, Algebr. Geom. Topology, 3 (2003), 1103-1118).
\end{openproblem}

\plainproblemsection

\begin{openproblem}
	10.24. A braid is said to be coloured if its strings represent the identity permutation. Is it true that links obtained by closing coloured braids are equivalent if and only if the original braids are conjugate in the braid group? G. S. Makanin No. For non-oriented links: the braid words ${\sigma }_{1}^{2}$ and ${\sigma }_{1}^{-2}$ both represent the Hopf link but are not conjugate in the braid group (easy to see using the Burau representation). For oriented links see Example 4 in (J. Birman, Braids, links, and mapping class groups (Ann. Math. Stud. Princeton, 82), Princeton, NJ, 1975, p. 100). (V. Shpil-rain, Letters of May, 28, 1998 and January, 20, 2002).
\end{openproblem}

\plainproblemsection

\begin{openproblem}
	10.25. (Well-known problem). Does there exist an algorithm which decides for a given automorphism of a free group whether this automorphism has a non-trivial fixed point? G. S. Makanin
	
	Yes, it does (O. Bogopolski, O. Maslakova, Int. J. Algebra Comput., 26, no. 1 (2016), 29-67).
\end{openproblem}

\plainproblemsection

\begin{openproblem}
	10.26. a) Does there exist an algorithm which decides, for given elements $a, b$ and an automorphism $\varphi$ of a free group, whether the equation $a{x}^{\varphi } = {xb}$ is soluble in this group? This question seems to be useful for solving the problem of equivalence of two knots. G. S. Makanin
	
	a) Yes, it does (O. Bogopolskii, A. Martino, O. Maslakova, E. Ventura, Bull. London Math. Soc., 38, no. 5 (2006), 787-794); further generalizations and applications are in (O. Bogopolski, A. Martino, E. Ventura, Trans. Amer. Math. Soc., 362 (2010), 2003-2036).
\end{openproblem}

\plainproblemsection

\begin{openproblem}
	10.26. b) Does there exist an algorithm which decides whether the equation of the form $w\left( {{x}_{{i}_{1}}^{{\varphi }_{1}},\ldots ,{x}_{{i}_{n}}^{{\varphi }_{n}}}\right)  = 1$ is soluble in a free group where ${\varphi }_{1},\ldots ,{\varphi }_{n}$ are automor-phisms of this group? G. S. Makanin
\end{openproblem}

\plainproblemsection

\begin{openproblem}
	10.27. a) Let $t$ be an involution of a finite group $G$ and suppose that the set $D =$ ${t}^{G} \cup  \left\{  {{t}^{x}{t}^{y} \mid  x, y \in  G,\left| {{t}^{x}{t}^{y}}\right|  = 2}\right\}$ does not intersect ${O}_{2}\left( G\right)$ . Prove that if $t \in  {O}_{2}\left( {C\left( d\right) }\right)$ for any involution $d$ from ${C}_{D}\left( t\right)$ , then $D = {t}^{G}$ (and in this case the structure of the group $\langle D\rangle$ is known).
	
	b) A significantly more general question. Let $t$ be an involution of a finite group $G$ and suppose that $t \in  {Z}^{ * }\left( {N\left( X\right) }\right)$ for every non-trivial subgroup $X$ of odd order which is normalized, but not centralized, by $t$ . What is the group $\left\langle  {t}^{G}\right\rangle$ ? A. A. I
\end{openproblem}

\plainproblemsection

\begin{openproblem}
	10.28. Is it true that finite strongly regular graphs with $\lambda  = 1$ have rank 3 ? A. A. Makhnëv
	
	No, not always (A. Cossidente, T. Penttila, J. London Math. Soc., 72 (2005), 731- 741).
\end{openproblem}

\plainproblemsection

\begin{openproblem}
	10.29. Describe the finite groups which contain a set of involutions $D$ such that, for any subset ${D}_{0}$ of $D$ generating a 2-subgroup, the normalizer ${N}_{D}\left( {D}_{0}\right)$ also generates a 2-subgroup. A. A. Makhnëv
\end{openproblem}

\plainproblemsection

\begin{openproblem}
	10.30. Does there exist a non-right-orderable group which is residually finite $p$ -group for a finite set of prime numbers $p$ containing at least two different primes? If a group is residually finite $p$ -group for an infinite set of primes $p$ , then it admits a linear order (A. H. Rhemtulla, Proc. Amer. Math. Soc., 41, no. 1 (1973), 31-33).
	
	Yes, it exists (P. A. Linnell, J. Algebra, 248 (2002) 605-607). N. Ya. Medvedev
\end{openproblem}

\plainproblemsection

\begin{openproblem}
	10.31. Suppose that $G$ is an algebraic group, $H$ is a closed normal subgroup of $G$ and $f : G \rightarrow  G/H$ is the canonical homomorphism. What conditions ensure that there exists a rational section $s : G/H \rightarrow  G$ such that ${sf} = 1$ ? For partial results, see (Yu. I. Merzlyakov, Rational Groups, Nauka, Moscow, 1987, § 37 (Russian)).
	
	Yu. I. Merzlyakov
\end{openproblem}

\plainproblemsection

\begin{openproblem}
	10.32. A group word is said to be universal on a group $G$ if its values on $G$ run over the whole of $G$ . For an arbitrary constant $c > 8/5$ , J. L. Brenner, R. J. Evans, and D. M. Silberger (Proc. Amer. Math. Soc., 96, no. 1 (1986), 23-28) proved that there exists a number ${N}_{0} = {N}_{0}\left( c\right)$ such that the word ${x}^{r}{y}^{s},{rs} \neq  0$ , is universal on the alternating group ${\mathbb{A}}_{n}$ for any $n \geq  \max \left\{  {{N}_{0}, c \cdot  \log m\left( {r, s}\right) }\right\}$ , where $m\left( {r, s}\right)$ is the product of all primes dividing ${rs}$ if $r, s \notin  \{  - 1,1\}$ , and $m\left( {r, s}\right)  = 1$ if $r, s \in  \{  - 1,1\}$ . For example, one can take ${N}_{0}\left( {5/2}\right)  = 5$ and ${N}_{0}\left( 2\right)  = {29}$ . Find an analogous bound for the degree of the symmetric group ${\mathbb{S}}_{n}$ under hypothesis that at least one of the $r, s$ is odd: up to now, here only a more crude bound $n \geq  \max \{ 6,{4m}\left( {r, s}\right)  - 4\}$ is known (M. Droste, Proc. Amer. Math. Soc., 96, no. 1 (1986), 18-22). Yu. J. Merzlvai 10.34. Does there exist a non-soluble finite group which coincides with the product of any two of its non-conjugate maximal subgroups?
	
	Editors' comment: the answer is negative for almost simple groups (T. V. Tikho-nenko, V. N. Tyutyanov, Siberian Math. J., 51, no. 1 (2010), 174-177).
	
	V. S. Monakhov
\end{openproblem}

\plainproblemsection

\begin{openproblem}
	10.33. For ${HNN}$ -extensions of the form $G = \left\langle  {t, A \mid  {t}^{-1}{Bt} = C,\varphi }\right\rangle$ , where $A$ is a finitely generated abelian group and $\varphi  : B \rightarrow  C$ is an isomorphism of two of its subgroups, find
	
	a) a criterion to be residually finite;
	
	b) a criterion to be Hopfian. Yu. I. Merzlyakov
	
	a) Such a criterion is found (S. Andreadakis, E. Raptis, D. Varsos, Arch. Math., 50 (1988), 495-501.
	
	b) Such a criterion is found (S. Andreadakis, E. Raptis, D. Varsos, Commun. Algebra, 20 (1992), 1511-1533).
\end{openproblem}

\plainproblemsection

\begin{openproblem}
	10.35. Is it true that every finitely generated torsion-free subgroup of $G{L}_{n}\left( \mathbb{C}\right)$ is residually in the class of torsion-free subgroups of $G{L}_{n}\left( \overline{\mathbb{Q}}\right)$ ? G. A. Noskov
\end{openproblem}

\plainproblemsection

\begin{openproblem}
	10.36. Is it true that $S{L}_{n}\left( {\mathbb{Z}\left\lbrack  {{x}_{1},\ldots ,{x}_{r}}\right\rbrack  }\right)$ , for $n$ sufficiently large, is a group of type ${\left( FP\right) }_{m}$ (which is defined in the same way as ${\left( FP\right) }_{\infty }$ was defined in 6.3, but with that weakening that the condition of being finitely generated is not imposed on the terms of the resolution with numbers $> m$ ). An affirmative answer is known for $m = 0$ , $n \geq  3$ (A. A. Suslin, Math. USSR Izvestiya,11(1977), 221-238) and for $m = 1, n \geq  5$ (M. S. Tulenbayev, Math. USSR Sbornik, 45 (1983), 139-154; U. Rehmann, C. Soulé, in: Algebraic K-theory, Proc. Conf., Northwestern Univ., Evanston, Ill., 1976, (Lect. Notes Math., 551), Springer, Berlin, 1976, 164-169). G. A. Noskov
\end{openproblem}

\plainproblemsection

\begin{openproblem}
	10.37. Suppose that $G$ is a finitely generated metabelian group all of whose integral homology groups are finitely generated. Is it true that $G$ is a group of finite rank? The answer is affirmative if $G$ splits over the derived subgroup (J. R. J. Groves, Quart. J. Math., 33, no. 132 (1982), 405-420). G. A. Noskov
	
	Yes, it is (D. H. Kochloukova, Groups St. Andrews 2001 in Oxford, Vol. II, Cambridge Univ. Press, 2003, 332-343).
\end{openproblem}

\plainproblemsection

\begin{openproblem}
	10.38. (W. van der Kallen). Does ${E}_{n + 3}\left( {\mathbb{Z}\left\lbrack  {{x}_{1},\ldots ,{x}_{n}}\right\rbrack  }\right) , n \geq  1$ , have finite breadth with respect to the set of transvections? G. A. Noskov
\end{openproblem}

\plainproblemsection

\begin{openproblem}
	10.39. a) Is the membership problem soluble for the subgroup ${E}_{n}\left( R\right)$ of $S{L}_{n}\left( R\right)$ where $R$ is a commutative ring?
	
	b) Is the word problem soluble for the groups ${K}_{i}\left( R\right)$ where ${K}_{i}$ are the Quillen $K$ -functors and $R$ is a commutative ring? G. A. Noskov
\end{openproblem}

\plainproblemsection

\begin{openproblem}
	10.40. (H. Bass). Let $G$ be a group, $e$ an idempotent matrix over $\mathbb{Z}G$ and let $\operatorname{tr}e = \mathop{\sum }\limits_{{g \in  G}}{e}_{g}g,\;{e}_{g} \in  \mathbb{Z}$ . Strong conjecture: for any non-trivial $x \in  G$ , the equation $\mathop{\sum }\limits_{{g \sim  x}}{e}_{g} = 0$ holds where $\sim$ denotes conjugacy in $G$ . Weak conjecture: $\mathop{\sum }\limits_{{g \in  G}}{e}_{g} = 0$ . The strong conjecture has been proved for finite, abelian and linear groups and the weak conjecture has been proved for residually finite groups.
	
	H. Bass' Comment of 2005: Significant progress has been made; a good up-to-date account is in the paper (A. J. Berrick, I. Chatterji, G. Mislin, Math. Ann., 329 (2004), 597-621). G. A. Noskov
\end{openproblem}

\plainproblemsection

\begin{openproblem}
	10.41. (Well-known problem). Let $\Gamma$ be an almost polycyclic group with no nontrivial finite normal subgroups, and let $k$ be a field. The complete ring of quotients $Q\left( {k\Gamma }\right)$ is a matrix ring ${M}_{n}\left( D\right)$ over a skew field. Conjecture: $n$ is the least common multiple of the orders of the finite subgroups of $\Gamma$ . An equivalent formulation (M. Lorenz) is as follows: $\rho \left( {{G}_{0}\left( {k\Gamma }\right) }\right)  = \rho \left( {{G}_{0}{\left( k\Gamma \right) }_{\mathcal{F}}}\right)$ , where ${G}_{0}\left( {k\Gamma }\right)$ is the Grothendieck group of the category of finitely-generated ${k\Gamma }$ -modules, ${G}_{0}{\left( k\Gamma \right) }_{\mathcal{F}}$ is the subgroup generated by classes of modules induced from finite subgroups of $\Gamma$ , and $\rho$ is the Goldie rank. There is a stronger conjecture: ${G}_{0}\left( {k\Gamma }\right)  = {G}_{0}{\left( k\Gamma \right) }_{\mathcal{F}}.G.A$ . Noskov The strong conjecture ${G}_{0}\left( {k\Gamma }\right)  = {G}_{0}{\left( k\Gamma \right) }_{\mathcal{F}}$ has been proved (J. A. Moody, Bull. Amer. Math. Soc., 17 (1987), 113-116).
\end{openproblem}

\plainproblemsection

\begin{openproblem}
	10.42. (J.T.Stafford). Let $G$ be a poly- $\mathbb{Z}$ -group and let $k$ be a field. Is it true that every finitely generated projective ${kG}$ -module is either a free module or an ideal?
	
	G. A. Noskov
\end{openproblem}

\plainproblemsection

\begin{openproblem}
	10.43. Let $R$ and $S$ be associative rings with identity such that 2 is invertible in $S$ . Let ${\Lambda }_{I} : G{L}_{n}\left( R\right)  \rightarrow  G{L}_{n}\left( {R/I}\right)$ be the homomorphism corresponding to an ideal $I$ of $R$ and let ${E}_{n}\left( R\right)$ be the subgroup of $G{L}_{n}\left( R\right)$ generated by elementary transvections ${t}_{ij}\left( x\right)$ . Let ${a}_{ij} = {t}_{ij}\left( 1\right) {t}_{ji}\left( {-1}\right) {t}_{ij}\left( 1\right)$ , let the bar denote images in the factor-group of $G{L}_{n}\left( R\right)$ by the centre and let ${PG} = \bar{G}$ for $G \leq  G{L}_{n}\left( R\right)$ . A homomorphism
	
	$$
	\Lambda  : {E}_{n}\left( R\right)  \rightarrow  {GL}\left( W\right)  = G{L}_{m}\left( S\right)
	$$
	
	is called standard if ${S}^{m} = P \oplus  \cdots  \oplus  P \oplus  Q$ (a direct sum of $S$ -modules in which there are $n$ summands $P$ ) and
	
	$$
	{\Lambda x} = {g}^{-1}\tau \left( {{\delta }^{ * }\left( x\right) f + {\left( {}^{t}{\delta }^{ * }{\left( x\right) }^{\nu }\right) }^{-1}\left( {1 - f}\right) }\right) g,\;x \in  {E}_{n}\left( R\right) ,
	$$
	
	where ${\delta }^{ * } : G{L}_{n}\left( R\right)  \rightarrow  G{L}_{n}\left( {\operatorname{End}P}\right)$ is the homomorphism induced by a ring homomorphism $\delta  : R \rightarrow$ End $P$ taking identity to identity, $g$ is an isomorphism of the module $W$ onto ${S}^{m},\tau  : G{L}_{n}\left( {\operatorname{End}P}\right)  \rightarrow  {GL}\left( {gW}\right)$ is an embedding, $f$ is a central idempotent of ${\delta R}, t$ denotes transposition and $\nu$ is an antiisomorphism of ${\delta R}$ . Let $n \geq  3, m \geq  2$ . One can show that the homomorphism
	
	$$
	{\Lambda }_{0} : P{E}_{n}\left( R\right)  \rightarrow  {GL}\left( W\right)  = G{L}_{m}\left( S\right)
	$$
	
	is induced by a standard homomorphism $\Lambda$ if ${\Lambda }_{0}{\bar{a}}_{ij} = {g}^{-1}\tau \left( {a}_{ij}^{ * }\right) g$ for some $g$ and $\tau$ and for any $i \neq  j$ where ${a}_{ij}^{ * }$ denotes the matrix obtained from ${a}_{ij}$ by replacing 0 and 1 from $R$ by 0 and 1 from End $P$ . Find (at least in particular cases) the form of a homomorphism ${\Lambda }_{0}$ that does not satisfy the last condition. V. M. Petechuk
\end{openproblem}

\plainproblemsection

\begin{openproblem}
	10.44. Prove that every standard homomorphism $\Lambda$ of ${E}_{n}\left( R\right)  \subset  G{L}_{n}\left( R\right)  = {GL}\left( V\right)$ into ${E}_{m}\left( S\right)  \subset  G{L}_{m}\left( S\right)  = {GL}\left( W\right)$ originates from a collineation and a correlation, that is, has the form ${\Lambda x} = \left( {{g}^{-1}{xg}}\right) f + \left( {{h}^{-1}{xh}}\right) \left( {1 - f}\right)$ , where $g$ is a semilinear isomorphism ${V}_{R} \rightarrow  {W}_{S}$ (collineation) and $h$ is a semilinear isomorphism ${V}_{R}{ \rightarrow  }_{S}{W}_{{S}^{0}}^{\prime }$ (correlation). V. M. Petechuk
\end{openproblem}

\plainproblemsection

\begin{openproblem}
	10.45. Let $n \geq  3$ and suppose that $N$ is a subgroup of $G{L}_{n}\left( R\right)$ which is normalized by ${E}_{n}\left( R\right)$ . Prove that either $N$ contains ${E}_{n}\left( R\right)$ or ${\Lambda }_{I}\left\lbrack  {N,{E}_{n}\left( R\right) }\right\rbrack   = 1$ for a suitable ideal $I \neq  R$ of $R$ . V. M. Petechuk
\end{openproblem}

\plainproblemsection

\begin{openproblem}
	10.46. Prove that for every element $\sigma  \in  G{L}_{n}\left( R\right) , n \geq  3$ , there exist transvections ${\tau }_{1},{\tau }_{2}$ such that $\left\lbrack  {\left\lbrack  {\sigma ,{\tau }_{1}}\right\rbrack  ,{\tau }_{2}}\right\rbrack$ is a unipotent element. V. M. Petechuk
\end{openproblem}

\plainproblemsection

\begin{openproblem}
	10.47. Describe the automorphisms of $P{E}_{2}\left( R\right)$ in the case where the ring $R$ is commutative and 2 and 3 have inverses in it. V. M. Petechuk
\end{openproblem}

\plainproblemsection

\begin{openproblem}
	10.48. Let $V$ be a vector space of finite dimension over a field of prime order. A subset $R$ of ${GL}\left( V\right)  \cup  \{ 0\}$ is called regular if $\left| R\right|  = \left| V\right| ,0,1 \in  R$ and ${vx} \neq  {vy}$ for any non-trivial vector $v \in  V$ and any distinct elements $x, y \in  R$ . It is obvious that $\tau ,\varepsilon ,{\mu }_{g}$ transform a regular set into a regular one, where ${x}^{\tau } = {x}^{-1}$ for $x \neq  0$ and ${0}^{\tau } = 0$ , ${x}^{\varepsilon } = 1 - x,{x}^{{\mu }_{g}} = x{g}^{-1}$ and $g$ is a non-zero element of the set being transformed. We say that two regular subsets are equivalent if one can be obtained from the other by a sequence of such transformations.
	
	a) Study the equivalence classes of regular subsets.
	
	b) Is every regular subset equivalent to a subgroup of ${GL}\left( V\right)$ together with 0 ? N. D. Podufalov
	
	a) They were studied (N. D. Podufalov, Algebra and Logic, 30, no. 1 (1991), 62-69). b) No. A regular set is closed with respect to multiplication if and only if the corresponding $\left( {\gamma ,\gamma }\right)$ -transitive plane is defined over a near-field. (N. D. Podufalov, Letter of February, 13, 1989.)
\end{openproblem}

\plainproblemsection

\begin{openproblem}
	10.49. Does there exist a group $G$ satisfying the following four conditions:
	
	1) $G$ is simple, moreover, there is an integer $n$ such that $G = {C}^{n}$ for any conjugacy class $C$ ,
	
	2) all maximal abelian subgroups of $G$ are conjugate in $G$ ,
	
	3) every maximal abelian subgroup of $G$ is self-normalizing and it is the centralizer of any of its nontrivial elements,
	
	4) there is an integer $m$ such that if $H$ is a maximal abelian subgroup of $G$ and $a \in  G \smallsetminus  H$ then every element in $G$ is a product of $m$ elements in ${aH}$ ?
	
	Comment of 2005: There is a partial solution in (E. Jaligot, A. Ould Houcine, J. Algebra, 280 (2004), 772-796). B. Poizat
\end{openproblem}

\plainproblemsection

\begin{openproblem}
	10.50. Complex characters of a finite group $G$ induced by linear characters of cyclic subgroups are called induced cyclic characters of $G$ . Computer-aided computations (A. V. Rukolaine, Abstracts of the 10th All-Union Sympos. on Group Theory, Gomel', 1986, p. 199 (Russian)) show that there exist groups (for example, ${\mathbb{S}}_{5}, S{L}_{2}\left( {13}\right) ,{M}_{11}$ ) all of whose irreducible complex characters are integral linear combinations of induced cyclic characters and the principal character of the group. Describe all finite groups with this property. A. V. Rukolaine
\end{openproblem}

\plainproblemsection

\begin{openproblem}
	10.51. Describe the structure of thin abelian $p$ -groups. Such a description is known in the class of separable $p$ -groups (C. Megibben, Mich. Math. J.,13, no. 2 (1966), 153-160). S. V. Rychkov
\end{openproblem}

\plainproblemsection

\begin{openproblem}
	10.52. (R. Mines). It is well known that the topology on the completion of an abelian group under the $p$ -adic topology is the $p$ -adic topology. R. Warfield has shown that this is also true in the category of nilpotent groups. Find a categorical setting for this theorem which includes the case of nilpotent groups. S. V. Rychkov
\end{openproblem}

\plainproblemsection

\begin{openproblem}
	10.53. (M. Dugas). Let $\mathfrak{R}$ be the Reid class, i. e. the smallest containing $\mathbb{Z}$ and closed under direct sums and direct products. Is $\mathfrak{R}$ closed under direct summands? S. V. Rychkov
\end{openproblem}

\plainproblemsection

\begin{openproblem}
	10.54. (R. Göbel). For a cardinal number $\mu$ , let
	
	$$
	{\mathbb{Z}}^{ < \mu } = \left\{  {f \in  {\mathbb{Z}}^{\mu }\left| \right| \operatorname{supp}\left( f\right)  \mid   < \mu }\right\}  \;\text{ and }\;{G}_{\mu } = {\mathbb{Z}}^{\mu }/{\mathbb{Z}}^{ < \mu }.
	$$
	
	a) Find a non-zero direct summand $D$ of ${G}_{{\omega }_{1}}$ such that $D ≆ {G}_{{\omega }_{1}}$ .
	
	b) Investigate the structure of ${G}_{\mu }$ (the structure of ${G}_{{\omega }_{0}}$ is well known).
	
	S. V. Rychkov
\end{openproblem}

\plainproblemsection

\begin{openproblem}
	10.55. (A. Mader). "Standard $B$ ", that is, $B = \mathbb{Z}\left( p\right)  \oplus  \mathbb{Z}\left( {p}^{2}\right)  \oplus  \cdots$ , is slender as a module over its endomorphism ring (A. Mader, in: Abelian Groups and Modules, Proc., Udine,1984, Springer,1984,315-327). Which abelian $p$ -groups are slender as modules over their endomorphism rings? S. V. Rychkov
\end{openproblem}

\plainproblemsection

\begin{openproblem}
	10.56. Is the lattice of formations of finite nilpotent groups of class $\leq  4$ distributive? A. N. Skiba
	
	No. The lattice $L$ of all varieties of nilpotent 2-groups of class $\leq  4$ is non-distributive (Yu. A. Belov, Algebra and Logic, 9, no. 6 (1970), 371-374; R. A. Bryce, Philos. Trans. Roy. Soc. London (A), 266 (1970), 281-355, footnote on p. 335). The mapping $V \rightarrow  V \cap  F$ , where $F$ is the class of all finite groups, is an embedding of $L$ into the lattice of formations. (L. Kovács, Letter of May, 12, 1988.) See also (L. A. Shemetkov, A. N. Skiba, Formations of algebraic systems, Nauka, Moscow, 1989 (Russian)).
\end{openproblem}

\plainproblemsection

\begin{openproblem}
	10.57. What are the minimal non- $A$ -formations? An $A$ -formation is, by definition, the formation of the finite groups all of whose Sylow subgroups are abelian.
	
	A. N. Skiba
\end{openproblem}

\plainproblemsection

\begin{openproblem}
	10.58. Is the subsemigroup generated by the undecomposable formations in the semigroup of formations of finite groups free? A. N. Skiba
\end{openproblem}

\plainproblemsection

\begin{openproblem}
	10.59. Is a ${p}^{\prime }$ -group $G$ locally nilpotent if it admits a splitting automorphism $\varphi$ of prime order $p$ such that all subgroups of the form $\left\langle  {g,{g}^{\varphi },\ldots ,{g}^{{\varphi }^{p - 1}}}\right\rangle$ are nilpotent? An automorphism $\varphi$ of order $p$ is called splitting if $g{g}^{\varphi }{g}^{{\varphi }^{2}}\cdots {g}^{{\varphi }^{p - 1}} = 1$ for all $g \in  G$ .
	
	A. I. Sozutov
\end{openproblem}

\plainproblemsection

\begin{openproblem}
	10.60. Does every periodic group $\left( {p\text{-group}}\right) A$ of regular automorphisms of an abelian group have non-trivial centre?
	
	Editors’ comment (2001): The answer is affirmative if $A$ contains an element of order 2 or 3 (A. Kh. Zhurtov, Siberian Math. J., 41, no. 2 (2000), 268-275).
	
	A. I. Sozutov
\end{openproblem}

\plainproblemsection

\begin{openproblem}
	10.61. Suppose that $H$ is a proper subgroup of a group $G, a \in  H,{a}^{2} \neq  1$ and for every $g \in  G \smallsetminus  H$ the subgroup $\left\langle  {a,{a}^{g}}\right\rangle$ is a Frobenius group whose complement contains $a$ . Does the set-theoretic union of the kernels of all Frobenius subgroups of $G$ with complement $\langle a\rangle$ constitute a subgroup? For definitions see 6.55; see also (A. I. Sozutov, Algebra and Logic, 34, no. 5 (1995), 295-305).
	
	Editors’ comment (2005): The answer is affirmative if the order of $a$ is even (A. M. Popov, A. I. Sozutov, Algebra and Logic, 44, no. 1 (2005), 40-45) or if the order of $a$ is not 3 or 5 and the group $\left\langle  {a,{a}^{g}}\right\rangle$ is finite for any $g \notin  H$ (A. M. Popov, Algebra and Logic, 43, no. 2 (2004), 123-127). A. I. Sozutov
\end{openproblem}

\plainproblemsection

\begin{openproblem}
	10.62. Construct an example of a (periodic) group without subgroups of index 2 which is generated by a conjugacy class of involutions $X$ such that the order of the product of any two involutions from $X$ is odd. A. I. Sozutov
\end{openproblem}

\plainproblemsection

\begin{openproblem}
	10.63. Is there a doubly transitive permutation group in which the stabilizer of a point is infinite cyclic? Ya. P. Sysak
	
	No (V. D. Mazurov, Siberian Math. J., 31 (1990), 615-617; Gy. Károlyi, S. J. Kovács, P. P. Pálfy, Aequationes Math., 39, no. 1 (1990), 161-166).
\end{openproblem}

\plainproblemsection

\begin{openproblem}
	10.64. Does there exist a non-periodic doubly transitive permutation group with a periodic stabilizer of a point? Ya. P. Sysak
\end{openproblem}

\plainproblemsection

\begin{openproblem}
	10.65. Determine the structure of infinite 2-transitive permutation groups $\left( {G,\Omega }\right)$ in which the stabilizer of a point $\alpha  \in  \Omega$ has the form ${G}_{\alpha } = A \cdot  {G}_{\alpha \beta }$ where ${G}_{\alpha \beta }$ is the stabilizer of two points $\alpha ,\beta ,\alpha  \neq  \beta$ , such that ${G}_{\alpha \beta }$ contains an element inverting the subgroup $A$ . Suppose, in particular, that $A \smallsetminus  \{ 1\}$ contains at most two conjugacy classes of ${G}_{\alpha }$ ; does $G$ then possess a normal subgroup isomorphic to ${PS}{L}_{2}$ over a field? A. N. Fomin
\end{openproblem}

\plainproblemsection

\begin{openproblem}
	10.66. Is a group $G$ non-simple if it contains two non-trivial subgroups $A$ and $B$ such that ${AB} \neq  G$ and $A{B}^{g} = {B}^{g}A$ for any $g \in  G$ ? This is true if $G$ is finite (O. H. Kegel, Arch. Math., 12, no. 2 (1961), 90-93). A. N. Fomin, V. P. Shunkov
	
	No; for example, for $G$ being any simple linearly ordered group and $A$ and $B$ arbitrary proper convex subgroups of $G$ (V. V. Bludov, Letter of February,12,1997).
\end{openproblem}

\plainproblemsection

\begin{openproblem}
	10.67. The class ${LN}{\mathfrak{M}}_{p}$ of locally nilpotent groups admitting a splitting automorphism of prime order $p$ (for definition see 10.59) is a variety of groups with operators (E. I. Khukhro, Math. USSR Sbornik, 58 (1987), 119-126). Is it true that
	
	$$
	{LN}{\mathfrak{M}}_{p} = \left( {{\mathfrak{N}}_{c\left( p\right) } \cap  {LN}{\mathfrak{M}}_{p}}\right)  \vee  \left( {{\mathfrak{B}}_{p} \cap  {LN}{\mathfrak{M}}_{p}}\right)
	$$
	
	where ${\mathfrak{N}}_{c\left( p\right) }$ is the variety of nilpotent groups of some $p$ -bounded class $c\left( p\right)$ and ${\mathfrak{B}}_{p}$ is the variety of groups of exponent $p$ ? E. I. Khukhro
\end{openproblem}

\plainproblemsection

\begin{openproblem}
	10.68. Suppose that a finite $p$ -group $G$ admits an automorphism of order $p$ with exactly ${p}^{m}$ fixed points. Is it true that $G$ has a subgroup whose index is bounded by a function of $p$ and $m$ and whose nilpotency class is bounded by a function of $m$ only? The results on $p$ -groups of maximal class give an affirmative answer in the case $m = 1$ . E. I. Khukhro
	
	Yes, it is (Yu. Medvedev, J. London Math. Soc. (2), 59, no. 3 (1999), 787-798).
\end{openproblem}

\plainproblemsection

\begin{openproblem}
	10.69. Suppose that $S$ is a closed oriented surface of genus $g > 1$ and $G = {\pi }_{1}S =$ ${G}_{1}{ * }_{A}{G}_{2}$ , where ${G}_{1} \neq  A \neq  {G}_{2}$ and the subgroup $A$ is finitely generated (and hence free). Is this decomposition geometrical, that is, do there exist connected surfaces ${S}_{1},{S}_{2}, T$ such that $S = {S}_{1} \cup  {S}_{2},{S}_{1} \cap  {S}_{2} = T,{G}_{i} = {\pi }_{1}{S}_{i}, A = {\pi }_{1}T$ and embeddings ${S}_{i} \subset  S, T \subset  S$ induce, in a natural way, embeddings ${G}_{i} \subset  G, A \subset  G$ ? This is true for $A = \mathbb{Z}$ . H. Zieschang
	
	No, not always (O. V. Bogopol'skiĭ, Geom. Dedicata, 94 (2002), 63-89).
\end{openproblem}

\plainproblemsection

\begin{openproblem}
	10.70. Find a geometrical justification for the Whitehead method for free products similar to the substantiation given for free groups by Whitehead himself and more visual than given in (D. J. Collins, H. Zieschang, Math. Z., 185, no. 4 (1984), 487-504; 186, no. 3 (1984), 335-361).
	
	D. J. Collins' comment: there is a partial solution in (D. McCullough, A. Miller, Symmetric automorphisms of free products, Mem. Amer. Math. Soc. 582 (1996)).
	
	H. Zieschang
\end{openproblem}

\plainproblemsection

\begin{openproblem}
	10.71. Is it true that the centralizer of any automorphism (any finite set of automor-phisms) in the automorphism group of a free group of finite rank is a finitely presented group? This is true in the case of rank 2 and in the case of inner automorphisms for any finite rank. V. A. Churkin
\end{openproblem}

\plainproblemsection

\begin{openproblem}
	10.72. Prove that the formation of all finite $p$ -groups does not decompose into a product of two non-trivial subformations. L. A. Shemetkov
	
	This has been proved (A. N. Skiba, L. A. Shemetkov, Dokl. Akad. Nauk BSSR, 33 (1989), 581-582 (Russian)).
\end{openproblem}

\plainproblemsection

\begin{openproblem}
	10.73. Enumerate all formations of finite groups all of whose subformations are ${S}_{n}$ -closed. L. A. Shemetkov
\end{openproblem}

\plainproblemsection

\begin{openproblem}
	10.74. Suppose that a group $G$ contains an element $a$ of prime order such that its centralizer ${C}_{G}\left( a\right)$ is finite and all subgroups $\left\langle  {a,{a}^{g}}\right\rangle  , g \in  G$ , are finite and almost all of them are soluble. Is $G$ locally finite? This problem is closely connected with 6.56. The question was solved in the positive for a number of very important partial cases by the author (Abstracts on Group Theory of the Mal'cev Int. Conf. on Algebra, Novosibirsk, 1989, p. 145 (Russian)). V. P. Shunkov
\end{openproblem}

\plainproblemsection

\begin{openproblem}
	10.75. Suppose that a group $G$ contains an element $a$ of prime order $p$ such that the normalizer of every finite subgroup containing $a$ has finite periodic part and all subgroups $\left\langle  {a,{a}^{g}}\right\rangle  , g \in  G$ , are finite and almost all of them are soluble. Does $G$ possess a periodic part if $p > 2$ ? It was proved in (V.P. Shunkov, Groups with involutions, Preprints no.4,5,12of the Comput. centre of SO AN SSSR, Krasnoyarsk,1986 (Russian)) that if $a$ is a point, then the answer is affirmative; on the other hand, a group with a point $a$ of order 2 satisfying the given hypothesis, which has no periodic part, was exhibited in the same works. For the definition of a point see (V. I. Senashov, V. P. Shunkov, Algebra and Logic, 22, no. 1 (1983), 66-81). V. P. Shunkov
\end{openproblem}

\plainproblemsection

\begin{openproblem}
	10.76. Suppose that $G$ is a periodic group having an infinite Sylow 2-subgroup $S$ which is either elementary abelian or a Suzuki 2-group, and suppose that the normalizer ${N}_{G}\left( S\right)$ is strongly embedded in $G$ and is a Frobenius group (see 6.55) with locally cyclic complement. Must $G$ be locally finite?
	
	Yes, it must (A. I. Sozutov, Algebra and Logic, 39, no. 5 (2000), 345-353; A. I. Sozu-tov, N. M. Suchkov, Math. Notes, 68, no. 2 (2000), 237-247).
\end{openproblem}

\plainproblemsection

\begin{openproblem}
	10.77. Suppose that $G$ is a periodic group containing an elementary abelian subgroup $R$ of order 4 . Must $G$ be locally finite
	
	a) if ${C}_{G}\left( R\right)$ is finite?
	
	b) if the centralizer of every involution of $R$ in $G$ is a Chernikov group?
	
	Remark of 1999: P. V. Shumyatsky (Quart. J. Math. Oxford (2), 49, no. 196 (1998), 491-499) gave a positive answer to the question a) in the case where $G$ is residually finite. V. P. Shunkov
\end{openproblem}

\plainproblemsection

\begin{openproblem}
	10.78. Does there exist a non-Chernikov group which is a product of two Chernikov subgroups? V. P. Shunkov
\end{openproblem}

\plainproblemsection

\begin{openproblem}
	11.1. (Well-known problem). Describe the structure of the centralizers of unipotent elements in almost simple groups of Lie type. R. Zh. Aleev
\end{openproblem}

\plainproblemsection

\begin{openproblem}
	*11.2. Classify the simple groups that are isomorphic to the multiplicative groups of finite rings, in particular, of the group rings of finite groups over finite fields and over $\mathbb{Z}/n\mathbb{Z}, n \in  \mathbb{Z}$ . R. Zh. Aleev
	
	*Such a classification is obtained in (C. Davis, T. Occhipinti, J. Pure Appl. Algebra, 218, no. 4 (2014), 743-744).
\end{openproblem}

\plainproblemsection

\begin{openproblem}
	11.3. (M. Aschbacher). A $p$ -local subgroup $H$ in a group $G$ is said to be a superlocal if $H = {N}_{G}\left( {{O}_{p}\left( H\right) }\right)$ . Describe the superlocals in alternating groups and in groups of Lie type. R. Zh. Aleev
\end{openproblem}

\plainproblemsection

\begin{openproblem}
	11.4. Is it true that the lattice of centralizers in a group is modular if it is a sublattice of the lattice of all subgroups? This is true for finite groups. V. A. Antonov No, not always (V. N. Obraztsov, J. Algebra, 199, no. 1 (1998), 337-343).
\end{openproblem}

\plainproblemsection

\begin{openproblem}
	11.5. Let $k$ be a commutative ring and let $G$ be a torsion-free almost polycyclic group. Suppose that $P$ is a finitely generated projective module over the group ring ${kG}$ and $P$ contains two elements independent over ${kG}$ . Is $P$ a free module?
	
	V. A. Artamonov
\end{openproblem}

\plainproblemsection

\begin{openproblem}
	11.6. Let $p$ be an odd prime. Is it true that every finite $p$ -group possesses a set of generators of equal orders? C. Bagiński
	
	No, it is not true (E. A. O'Brien, C. M. Scoppola, M. R. Vaughan-Lee, Proc. Amer. Math. Soc., 134, no. 12 (2006), 3457-3464).
\end{openproblem}

\plainproblemsection

\begin{openproblem}
	11.7. a) Find conditions for a group $G$ , given by its presentation, under which the property some term of the lower central series of $G$ is a free group(*)implies the residual finiteness of $G$ .
	
	b) Find conditions on the structure of a residually finite group $G$ which ensure property $\left( *\right)$ for $G$ . K. Bencsáth
\end{openproblem}

\plainproblemsection

\begin{openproblem}
	11.8. a) For a finite group $X$ , let ${\chi }_{1}\left( X\right)$ denote the totality of the degrees of all irreducible complex characters of $X$ with allowance for their multiplicities. Suppose that ${\chi }_{1}\left( G\right)  = {\chi }_{1}\left( H\right)$ for groups $G$ and $H$ . Clearly, then $\left| G\right|  = \left| H\right|$ . Is it true that $H$ is simple if $G$ is simple? Ya. G. Berkovich
	
	a) Yes, it is (H. P. Tong-Viet, J. Algebra, 357 (2012), 61-68; Monatsh. Math., 166, no. 3-4 (2012), 559-577; Algebr. Represent. Theory, 15, no. 2 (2012), 379-389).
\end{openproblem}

\plainproblemsection

\begin{openproblem}
	11.8. b) For a finite group $X$ , let ${\chi }_{1}\left( X\right)$ denote the totality of the degrees of all irreducible complex characters of $X$ with allowance for their multiplicities. Suppose that ${\chi }_{1}\left( G\right)  = {\chi }_{1}\left( H\right)$ for groups $G$ and $H$ . Clearly, then $\left| G\right|  = \left| H\right|$ . Is it true that $H$ is soluble if $G$ is soluble?
	
	It is known that $H$ is a Frobenius group if $G$ is a Frobenius group.
	
	Ya. G. Berkovich
\end{openproblem}

\plainproblemsection

\begin{openproblem}
	11.9. (I. I. Pyatetskii-Shapiro). Does there exist a finite non-soluble group $G$ such that the set of characters induced by the trivial characters of representatives of all conjugacy classes of subgroups of $G$ is linearly independent? Ya. G. Berkovich
\end{openproblem}

\plainproblemsection

\begin{openproblem}
	11.10. (R. C. Lyndon). a) Does there exist an algorithm that, given a group word $w\left( {a, x}\right)$ , recognizes whether $a$ is equal to the identity element in the group $\left\langle  {a, x \mid  {a}^{n} = 1, w\left( {a, x}\right)  = 1}\right\rangle  ?$ V. V. Bludov
\end{openproblem}

\plainproblemsection

\begin{openproblem}
	11.10. (R. C. Lyndon). b) Is it true that $a \neq  1$ in $G = \left\langle  {a, x \mid  {a}^{5} = 1,{a}^{{x}^{2}} = \left\lbrack  {a,{a}^{x}}\right\rbrack  }\right\rangle$ ? V. V. Bludov
	
	Yes, it is true, since the mapping $a \rightarrow  \left( {13524}\right) , x \rightarrow  \left( {12435}\right)$ can be extended to a homomorphism of the group $G$ onto the alternating group ${\mathbb{A}}_{5}$ (D.N. Azarov, in: Algebraicheskije Systemy, Ivanovo, 1991, 4-5 (Russian)).
\end{openproblem}

\plainproblemsection

\begin{openproblem}
	11.11. The well-known Baer-Suzuki theorem states that if every two conjugates of an element $a$ of a finite group $G$ generate a finite $p$ -subgroup, then $a$ is contained in a normal $p$ -subgroup.
	
	a) Does such a theorem hold in the class of periodic groups? The case $p = 2$ is of particular interest.
	
	b) Does such a theorem hold in the class of binary finite groups?
	
	A. V. Borovik
	
	a) No, it does not hold for $p = 2$ (V. D. Mazurov, A. Yu. Ol’shanskiĭ, A. I. Sozutov, Algebra and Logic, 54, no. 2 (2015), 161-166).
	
	b) Yes, such a theorem does hold for binary finite groups. By the Baer-Suzuki theorem $\left\langle  {{a}_{1}, b}\right\rangle$ is a finite $p$ -group for any element ${a}_{1} \in  {a}^{G} = \left\{  {{a}^{g} \mid  g \in  G}\right\}$ and for any $p$ - element $b \in  G$ . Now induction on $n$ yields that the product ${a}_{1}\cdots {a}_{n}$ is a $p$ -element for any ${a}_{1},\ldots ,{a}_{n} \in  {a}^{G}$ . (A. I. Sozutov, Siberian Math. J.,41, no. 3 (2000),561-562.)
\end{openproblem}

\plainproblemsection

\begin{openproblem}
	11.12. a) Suppose that $G$ is a simple locally finite group in which the centralizer of some element is a linear group, that is, a group admitting a faithful matrix representation over a field. Is $G$ itself a linear group?
	
	b) The same question with replacement of the word "linear" by "finitary linear". A group $H$ is finitary linear if it admits a faithful representation on an infinite-dimensional vector space $V$ such that the residue subspaces $V\left( {1 - h}\right)$ have finite dimensions for all $h \in  H$ . A. V. Borovik
\end{openproblem}

\plainproblemsection

\begin{openproblem}
	11.13. Suppose that $G$ is a periodic group with an involution $i$ such that ${i}^{g} \cdot  i$ has odd order for any $g \in  G$ . Is it true that the image of $i$ in $G/O\left( G\right)$ belongs to the centre of $G/O\left( G\right)$ ? A. V. Borovik
	
	No, not always (E. B. Durakov, A. I. Sozutov, Algebra Logic, 52, no. 5 (2013), 422- 425).
\end{openproblem}

\plainproblemsection

\begin{openproblem}
	11.14. Is every finite simple group characterized by its Cartan matrix over an algebraically closed field of characteristic 2? In (R. Brandl, Arch. Math., 38 (1982), 322-323) it is shown that for any given finite group there exist only finitely many finite groups with the same Cartan matrix. R. Brandl
\end{openproblem}

\plainproblemsection

\begin{openproblem}
	11.15. It is known that for each prime number $p$ there exists a series ${a}_{1},{a}_{2},\ldots$ of words in two variables such that the finite group $G$ has abelian Sylow $p$ -subgroups if and only if ${a}_{k}\left( G\right)  = 1$ for almost all $k$ . For $p = 2$ such a series is known explicitly (R. Brandl, J. Austral. Math. Soc., 31 (1981), 464-469). What about $p > 2$ ?
	
	R. Brandl
\end{openproblem}

\plainproblemsection

\begin{openproblem}
	11.16. Let ${V}_{r}$ be the class of all finite groups $G$ satisfying a law $\left\lbrack  {x,{}_{r}y}\right\rbrack   = \left\lbrack  {x,{}_{s}y}\right\rbrack$ for some $s = s\left( G\right)  > r$ . Here $\left\lbrack  {x,{}_{1}y}\right\rbrack   = \left\lbrack  {x, y}\right\rbrack$ and $\left\lbrack  {x,{}_{i + 1}y}\right\rbrack   = \left\lbrack  {\left\lbrack  {x,{}_{i}y}\right\rbrack  , y}\right\rbrack$ .
	
	a) Is there a function $f$ such that every soluble group in ${V}_{r}$ has Fitting length $< f\left( r\right)$ ? For $r < 3$ see (R. Brandl, Bull. Austral. Math. Soc.,28 (1983),101-110).
	
	b) Is it true that ${V}_{r}$ contains only finitely many nonabelian simple groups? This is true for $r < 4$ . R. Brandl
\end{openproblem}

\plainproblemsection

\begin{openproblem}
	11.17. Let $G$ be a finite group and let $d = d\left( G\right)$ be the least positive integer such that $G$ satisfies a law $\left\lbrack  {x,{}_{r}y}\right\rbrack   = \left\lbrack  {x,{}_{r + d}y}\right\rbrack$ for some nonnegative integer $r = r\left( G\right)$ .
	
	a) Let $e = 1$ if $d\left( G\right)$ is even and $e = 2$ otherwise. Is it true that the exponent of $G/F\left( G\right)$ divides $e \cdot  d\left( G\right)$ ?
	
	b) If $G$ is a nonabelian simple group, does the exponent of $G$ divide $d\left( G\right)$ ?
	
	Part a) is true for soluble groups (N. D. Gupta, H. Heineken, Math. Z., 95 (1967), 276-287). I have checked part b) for ${\mathbb{A}}_{n},{PSL}\left( {2, q}\right)$ , and a number of sporadic groups.
	
	R. Brandl
\end{openproblem}

\plainproblemsection

\begin{openproblem}
	11.18. Let $G\left( {a, b}\right)  = \left\langle  {x, y \mid  x = \left\lbrack  {x,{}_{a}y}\right\rbrack  , y = \left\lbrack  {y,{}_{b}x}\right\rbrack  }\right\rangle$ . Is $G\left( {a, b}\right)$ finite?
	
	It is easy to show that $G\left( {1, b}\right)  = 1$ and one can show that $G\left( {2,2}\right)  = 1$ . Nothing is known about $G\left( {2,3}\right)$ . If one could show that every minimal simple group is a quotient of some $G\left( {a, b}\right)$ , then this would yield a very nice sequence of words in two variables to characterize soluble groups, see (R. Brandl, J. S. Wilson, J. Algebra, 116 (1988),
	
	334-341.) R. Brandl
\end{openproblem}

\plainproblemsection

\begin{openproblem}
	11.19. (C. Sims). Is the $n$ th term of the lower central series of an absolutely free group the normal closure of the set of basic commutators (in some fixed free generators) of weight exactly $n$ ?
	
	D. Jackson announced a positive answer for $n \leq  5$ . A. Gaglione, D. Spellman
\end{openproblem}

\plainproblemsection

\begin{openproblem}
	11.20. Suppose we have $\left\lbrack  {a, b}\right\rbrack   = \left\lbrack  {c, d}\right\rbrack$ in an absolutely free group, where $a, b,\left\lbrack  {a, b}\right\rbrack$ are basic commutators (in some fixed free generators). If $c$ and $d$ are arbitrary (proper) commutators, does it follow that $a = c$ and $b = d$ ? A. Gaglione, D. Spellman No, not always. For example, if ${x}_{1},{x}_{2},{x}_{3}$ are free generators, ${x}_{1} < {x}_{2} < {x}_{3}$ , $a = \left\lbrack  {\left\lbrack  {{x}_{2},{x}_{1}}\right\rbrack  ,{x}_{3}}\right\rbrack  , b = \left\lbrack  {{x}_{2},{x}_{1}}\right\rbrack  , c = {b}^{-1}, d = {b}^{{x}_{3}b}$ (V. G. Bardakov, Abstracts of the IIIrd Intern. Conf. on Algebra, Krasnoyarsk, 1993, p. 33 (Russian)).
\end{openproblem}

\plainproblemsection

\begin{openproblem}
	11.21. Let ${\mathfrak{N}}_{p}$ denote the formation of all finite $p$ -groups, for a given prime number $p$ . Is it true that, for every subformation $\mathfrak{F}$ of ${\mathfrak{N}}_{p}$ , there exists a variety $\mathfrak{M}$ such that $\mathfrak{F} = {\mathfrak{N}}_{p} \cap  \mathfrak{M}?$ A. F. Vasil’yev
	
	No, it is not true. There is a natural one-to-one correspondence between formations of finite $p$ -groups and varieties of pro- $p$ -groups: for every variety $\mathfrak{V}$ of pro- $p$ -groups the class of all finite groups in $\mathfrak{V}$ is a formation of $p$ -groups and every formation of $p$ -groups arises in this way. There are continuously many varieties of nilpotent pro- $p$ -groups of class at most 6 (A. N. Zubkov, Siberian Math. J., 29, no. 3 (1988), 491-494) and only countably many varieties of nilpotent groups of class at most 6 . (A. N. Krasil'nikov, Letter of July, 17th, 1998.)
\end{openproblem}

\plainproblemsection

\begin{openproblem}
	11.22. Characterize all $p$ -groups, $p$ a prime, that can be faithfully represented as $n \times  n$ triangular matrices over a division ring of characteristic $p$ .
	
	If $p \geq  n$ the solution is given in (B. A. F. Wehrfritz, Bull. London Math. Soc.,19 (1987),320-324). The corresponding question with $p = 0$ was solved by B. Hartley and P. Menal (Bull. London Math. Soc., 15 (1983), 378-383), but see above reference for a second proof. B. A. F. Wehrfritz
\end{openproblem}

\plainproblemsection

\begin{openproblem}
	11.23. An automorphism $\varphi$ of a group $G$ is called a nil-automorphism if, for every $a \in  G$ , there exists $n$ such that $\left\lbrack  {a,{}_{n}\varphi }\right\rbrack   = 1$ . Here $\left\lbrack  {x,{}_{1}y}\right\rbrack   = \left\lbrack  {x, y}\right\rbrack$ and $\left\lbrack  {x,{}_{i + 1}y}\right\rbrack   =$ $\left\lbrack  {\left\lbrack  {x,{}_{i}y}\right\rbrack  , y}\right\rbrack$ . An automorphism $\varphi$ is called an $e$ -automorphism if, for any two $\varphi$ -invariant subgroups $A$ and $B$ such that $A \nsubseteq  B$ , there exists $a \in  A \smallsetminus  B$ such that $\left\lbrack  {a,\varphi }\right\rbrack   \in  B$ . Is every $e$ -automorphism of a group a nil-automorphism? V. G. Vilyatser
\end{openproblem}

\plainproblemsection

\begin{openproblem}
	11.24. A Fitting class $\mathfrak{F}$ is said to be local if there exists a group function $f$ (for definition see (L. A. Shemetkov, Formations of Finite Groups, Moscow, Nauka, 1978 (Russian)) such that $f\left( p\right)$ is a Fitting class for every prime number $p$ and $\mathfrak{F} = {\mathfrak{G}}_{\pi \left( \mathfrak{F}\right) }\bigcap \left( {\mathop{\bigcap }\limits_{{p \in  \pi \left( \mathfrak{F}\right) }}f\left( p\right) {\mathfrak{N}}_{p}{\mathfrak{G}}_{{p}^{\prime }}}\right)$ . Is every hereditary Fitting class of finite groups local? N. T. Vorob'ëv
	
	No, not every (L. A. Shemetkov, A. F. Vasil'yev, Abstracts of the Conf. of Mathematicians of Belarus', Part 1, Grodno, 1992, p. 56 (Russian); S. F. Kamornikov, Math. Notes, 55, no. 6 (1994), 586-588).
\end{openproblem}

\plainproblemsection

\begin{openproblem}
	11.25. a) Does there exist a local product (different from the class of all finite groups and from the class of all finite soluble groups) of Fitting classes each of which is not local and is not a formation? See the definition of the product of Fitting classes in (N. T. Vorob'ev, Math. Notes, 43, No 1-2 (1988), 91-94). N. T. Vorob'ev Yes, there exists (N. T. Vorob'ev, A. N. Skiba, Problems in Algebra, 8, Gomel', 1995, 55-58 (Russian)).
\end{openproblem}

\plainproblemsection

\begin{openproblem}
	11.25. b) Do there exist local Fitting classes which are decomposable into a nontrivial product of Fitting classes and in every such a decomposition all factors are non-local? For the definition of the product of Fitting classes see (N. T. Vorob'ev, Math. Notes, 43, no. 2 (1988), 91-94). N. T. Vorob'ëv
\end{openproblem}

\plainproblemsection

\begin{openproblem}
	11.26. Does there exist a group which is not isomorphic to outer automorphism group of a metabelian group with trivial center?
	
	No, given any group $G$ there is a metabelian group $M$ with trivial center such that Out $M \cong  G$ (R. Göbel, A. Paras, J. Pure Appl. Algebra, 149, no. 3 (2000) 251-266), and if $G$ is finite or countable then $M$ above can be chosen countable (R. Göbel, A. Paras, in: Abelian Groups and Modules, Proc. Int. Conf. Dublin, 1998, Birkhäuser, Basel, 1999, 309-317).
\end{openproblem}

\plainproblemsection

\begin{openproblem}
	11.27. What are the minimum numbers of generators for groups $G$ satisfying $S \leq$ $G \leq$ Aut $S$ where $S$ is a finite simple non-abelian group? $\;K$ . Gruenberg The number $d\left( G\right)$ is found (mod CFSG) for every such a group $G$ . In particular, $d\left( G\right)  = \max \{ 2, d\left( {G/S}\right) \}$ and $d\left( G\right)  \leq  3$ (F. Dalla Volta, A. Lucchini, J. Algebra,178, no. 1 (1995), 194-223).
\end{openproblem}

\plainproblemsection

\begin{openproblem}
	11.28. Suppose the prime graph of the finite group $G$ is disconnected. (This means that the set of prime divisors of the order of $G$ is the disjoint union of non-empty subsets $\pi$ and ${\pi }^{\prime }$ such that $G$ contains no element of order ${pq}$ where $p \in  \pi , q \in  {\pi }^{\prime }$ .) Then P. A. Linnell (Proc. London Math. Soc., 47, no. 1 (1983), 83-127) has proved mod CFSG that there is a decomposition of $\mathbb{Z}G$ -modules $\mathbb{Z} \oplus  \mathbb{Z}G = A \oplus  B$ with $A$ and $B$ non-projective. Find a proof independent of CFSG. $K$ . Gruenberg
\end{openproblem}

\plainproblemsection

\begin{openproblem}
	11.29. Let $F$ be a free group and $\mathfrak{f} = \mathbb{Z}F\left( {F - 1}\right)$ the augmentation ideal of the integral group ring $\mathbb{Z}F$ . For any normal subgroup $R$ of $F$ define the corresponding ideal $\mathfrak{r} = \mathbb{Z}F\left( {R - 1}\right)  = {}_{\mathrm{{id}}}\left( {r - 1 \mid  r \in  R}\right)$ . One may identify, for instance, $F \cap  \left( {1 + \mathfrak{r}\mathfrak{f}}\right)  = {R}^{\prime }$ , where $F$ is naturally imbedded into $\mathbb{Z}F$ and $1 + \mathfrak{r}\mathfrak{f} = \{ 1 + a \mid  a \in  \mathfrak{r}\mathfrak{f}\}$ .
	
	Identify in an analogous way in terms of corresponding subgroups of $F$
	
	a) $F \cap  \left( {1 + {\mathfrak{r}}_{1}{\mathfrak{r}}_{2}\cdots {\mathfrak{r}}_{n}}\right)$ , where ${R}_{i}$ are normal subgroups of $F, i = 1,2,\ldots , n$ ;
	
	b) $F \cap  \left( {1 + {\mathfrak{r}}_{1}{\mathfrak{r}}_{2}{\mathfrak{r}}_{3}}\right)$ ;
	
	c) $F \cap  \left( {1 + \mathfrak{{fs}} + {\mathfrak{f}}^{n}}\right)$ , where $F/S$ is finitely generated nilpotent;
	
	d) $F \cap  \left( {1 + \mathfrak{{fs}}f + {\mathfrak{f}}^{n}}\right)$ ;
	
	e) $F \cap  \left( {1 + \mathfrak{r}\left( k\right)  + {\mathfrak{f}}^{n}}\right) , n > k \geq  2$ , where $\mathfrak{r}\left( k\right)  = \mathfrak{r}{\mathfrak{f}}^{k - 1} + \mathfrak{f}\mathfrak{r}{\mathfrak{f}}^{k - 2} + \cdots  + {\mathfrak{f}}^{k - 1}\mathfrak{r}$ .
	
	N. D. Gupta
\end{openproblem}

\plainproblemsection

\begin{openproblem}
	11.29. f) Let $F$ be a free group and $f = \mathbb{Z}F\left( {F - 1}\right)$ the augmentation ideal of the integral group ring $\mathbb{Z}F$ . For any normal subgroup $R$ of $F$ define the corresponding ideal $\mathfrak{r} = \mathbb{Z}F\left( {R - 1}\right)  = {}_{\text{id }}\left( {r - 1 \mid  r \in  R}\right)$ . One may identify, for instance, $F \cap  \left( {1 + \mathfrak{r}\mathfrak{f}}\right)  =$ ${R}^{\prime }$ , where $F$ is naturally imbedded into $\mathbb{Z}F$ and $1 + \mathfrak{{rf}} = \{ 1 + a \mid  a \in  \mathfrak{{rf}}\}$ . Is the quotient group $\left( {F \cap  \left( {1 + \mathfrak{r} + {\mathfrak{f}}^{n}}\right) }\right) /R \cdot  {\gamma }_{n}\left( F\right)$ always abelian? Yes, it is (N.D. Gupta, Yu. V. Kuz'min, J. Pure Appl. Algebra, 78, no. 1 (1992), 165-172).
\end{openproblem}

\plainproblemsection

\begin{openproblem}
	11.30. Is it true that the rank of a torsion-free soluble group is equal to the rank of any of its subgroups of finite index? The answer is affirmative for groups having a rational series (D.I. Zaitsev, in: Groups with restrictions on subgroups, Naukova dumka, Kiev, 1971, 115-130 (Russian)). We note also that every torsion-free soluble group of finite rank contains a subgroup of finite index which has a rational series.
	
	D. I. Zaitsev
\end{openproblem}

\plainproblemsection

\begin{openproblem}
	11.31. Is a radical group polycyclic if it is a product of two polycyclic subgroups? The answer is affirmative for soluble and for hyperabelian groups (D. I. Zaitsev, Math. Notes, 29, no. 4 (1981), 247-252; J. C. Lennox, J. E. Roseblade, Math. Z., 170 (1980), 153-154).
\end{openproblem}

\plainproblemsection

\begin{openproblem}
	11.32. Describe the primitive finite linear groups that contain a matrix with simple spectrum, that is, a matrix all of whose eigenvalues are of multiplicity 1 . See partial results in (A. Zalesskii, I. D. Suprunenko, Commun. Algebra, 26, no. 3 (1998), 863- 888; 28, no. 4 (2000), 1789-1833; Ch. Rudloff, A. Zalesski, J. Group Theory, 10 (2007), 585-612; L. Di Martino, A. E. Zalesski, J. Algebra Appl., 11, no. 2 (2012), 1250038; L. Di Martino, M. Pellegrini, A. Zalesski, Commun. Algebra, 42 (2014), 880-908, where the problem is considered in a more general context, when all but one eigenvalue have multiplicity 1.). A. E. Zalesskiǐ
\end{openproblem}

\plainproblemsection

\begin{openproblem}
	11.33. a) Let $G\left( q\right)$ be a simple Chevalley group over a field of order $q$ . Prove that there exists $m$ such that the restriction of every non-one-dimensional complex representation of $G\left( {q}^{m}\right)$ to $G\left( q\right)$ contains all irreducible representations of $G\left( q\right)$ as composition factors. A. E. Zalesskii
	
	This is proved in (D. Gluck, J. Algebra, 155, no. 2 (1993), 221-237).
\end{openproblem}

\plainproblemsection

\begin{openproblem}
	11.33. b) Let $G\left( q\right)$ be a simple Chevalley group over a field of order $q$ . Prove that there exists $m$ such that the restriction of every non-one-dimensional representation of $G\left( {q}^{m}\right)$ over a field of prime characteristic not dividing $q$ to $G\left( q\right)$ contains all irreducible representations of $G\left( q\right)$ as composition factors. A. E. Zalesskiǐ
\end{openproblem}

\plainproblemsection

\begin{openproblem}
	11.34. Describe the complex representations of quasisimple finite groups which remain irreducible after reduction modulo any prime number $q$ . An important example: representations of degree $\left( {{p}^{k} - 1}\right) /2$ of the symplectic group ${Sp}\left( {{2k}, p}\right)$ where $k \in  \mathbb{N}$ and $p$ is an odd prime. Partial results see in (P. H. Tiep, A. E. Zalesski, Proc. London Math. Soc., 84 (2002), 439-472; P. H. Tiep, A. E. Zalesski, Proc. Amer. Math. Soc., 130 (2002), 3177-3184). A. E. Zalesskiǐ
\end{openproblem}

\plainproblemsection

\begin{openproblem}
	11.35. Suppose that $H$ is a finite linear group over $\mathbb{C}$ and $h$ is an element of $H$ of prime order $p$ which is not contained in any abelian normal subgroup. Is it true that $h$ has at least $\left( {p - 1}\right) /2$ different eigenvalues? A. E. Zalesskii
	
	Yes, it is (G. R. Robinson, J. Algebra, 178, no. 2 (1995), 635-642).
\end{openproblem}

\plainproblemsection

\begin{openproblem}
	11.36. Let $G = B\left( {m, n}\right)$ be the free Burnside group of rank $m$ and of odd exponent $n \gg  1$ . Are the following statements true?
	
	a) Every 2-generated subgroup of $G$ is isomorphic to the Burnside $n$ -product of two cyclic groups.
	
	b) Every automorphism $\varphi$ of $G$ such that ${\varphi }^{n} = 1$ and ${b}^{\varphi } \cdot  {b}^{{\varphi }^{2}}\cdots {b}^{{\varphi }^{n}} = 1$ for all $b \in  G$ is an inner automorphism (here $m > 1$ ).
	
	Editors’ comment: for large prime $n$ this follows from (E. A. Cherepanov, Int. J. Algebra Comput.,16 (2006), 839-847), for prime $n \geq  {1009}$ from (V. S. Atabekyan, Izv. Math.,75, no. 2 (2011),223-237); this is also true for odd $n \geq  {1003}$ if in addition the order of $\varphi$ is a prime power (V. S. Atabekyan, Math. Notes,95, no. 5 (2014), ${586} - {589})$ .
	
	c) The group $G$ is Hopfian if $m < \infty$ .
	
	d) All retracts of $G$ are free. S. V. Ivanov
\end{openproblem}

\plainproblemsection

\begin{openproblem}
	11.36. e) Let $G = B\left( {m, n}\right)$ be the free Burnside group of rank $m$ and of odd exponent $n \gg  1$ . Is it true that all zero divisors in the group ring $\mathbb{Z}G$ are trivial? which means that if ${ab} = 0$ then $a = {a}_{1}c, b = d{b}_{1}$ where ${a}_{1}, c,{b}_{1}, d \in  \mathbb{Z}G,{cd} = 0$ , and the set $\operatorname{supp}c \cup  \operatorname{supp}d$ is contained in a cyclic subgroup of $G$ . e) No, there are nontrivial zero divisors (S. V. Ivanov, R. Mikhailov, Canad. Math. Bull., 57 (2014), 326-334).
\end{openproblem}

\plainproblemsection

\begin{openproblem}
	11.37. a) Can the free Burnside group $B\left( {m, n}\right)$ , for any $m$ and $n$ , be given by defining relations of the form ${v}^{n} = 1$ such that for any natural divisor $d$ of $n$ distinct from $n$ the element ${v}^{d}$ is not trivial in $B\left( {m, n}\right)$ ? This is true for odd $n \geq  {665}$ , and for all $n \geq  {2}^{48}$ divisible by ${2}^{9}$ . S. V. Ivanov
\end{openproblem}

\plainproblemsection

\begin{openproblem}
	11.37. b) Can the free Burnside group $B\left( {m, n}\right)$ , for any $m$ and $n = {2}^{l} \gg  1$ , be given by defining relations of the form ${v}^{n} = 1$ such that for any natural divisor $d$ of $n$ distinct from $n$ the element ${v}^{d}$ is not trivial in $B\left( {m, n}\right)$ ? $\;S.V$ . Ivanov Yes, it can (S. V. Ivanov, Int. J. Algebra Comput., 4, no. 1-2 (1994), 1-308).
\end{openproblem}

\plainproblemsection

\begin{openproblem}
	11.38. Does there exist a finitely presented Noetherian group which is not almost polycyclic? S. V. Ivanov
\end{openproblem}

\plainproblemsection

\begin{openproblem}
	11.39. (Well-known problem). Does there exist a group which is not almost polycyclic and whose integral group ring is Noetherian? S. V. Ivanov
\end{openproblem}

\plainproblemsection

\begin{openproblem}
	11.40. Prove or disprove that a torsion-free group $G$ with the small cancellation condition ${C}^{\prime }\left( \lambda \right)$ where $\lambda  \ll  1$ necessarily has the $\mathcal{U}P$ -property (and therefore ${KG}$ has no zero divisors). S. V. Ivanov
\end{openproblem}

\plainproblemsection

\begin{openproblem}
	11.42. Does there exist a torsion-free group having exactly three conjugacy classes and containing a subgroup of index 2? A. V. Izosov
	
	Yes, it exists (A. Minasyan, Comment. Math. Helv., 84 (2009), 259-296).
\end{openproblem}

\plainproblemsection

\begin{openproblem}
	11.43. For a finite group $X$ , we denote by $k\left( X\right)$ the number of its conjugacy classes. Is it true that $k\left( {AB}\right)  \leq  k\left( A\right) k\left( B\right)$ ?L. S. Kazarin
	
	No, it is not true in general: let $G = \left\langle  {a, b \mid  {a}^{30} = {b}^{2} = 1,{a}^{b} = {a}^{-1}}\right\rangle   \cong  {D}_{60}$ be the dihedral group of order 60, and let $A = \left\langle  {{a}^{10},{ba}}\right\rangle   \cong  {D}_{6}$ and $B = \left\langle  {{a}^{6}, b}\right\rangle   \cong  {D}_{10}$ . Then $G = {AB}$ , but $G$ has 18 conjugacy classes and $A$ and $B$ only 3 and 4, respectively. From (P. Gallagher, Math. Z.,118 (1970), 175-179) a positive answer follows if $A$ or $B$ is normal. The problem remains open in the case where $A$ and $B$ have coprime orders, see new problem 14.44. (J. Sangroniz, Letter of December, 17, 1998.)
\end{openproblem}

\plainproblemsection

\begin{openproblem}
	11.44. For a finite group $X$ , we denote by $r\left( X\right)$ its sectional rank. Is it true that the sectional rank of a finite $p$ -group, which is a product ${AB}$ of its subgroups $A$ and $B$ , is bounded by some linear function of $r\left( A\right)$ and $r\left( B\right)$ ? L. S. Kazarin
\end{openproblem}

\plainproblemsection

\begin{openproblem}
	11.45. A $t - \left( {v, k,\lambda }\right)$ design $\mathcal{D} = \left( {X,\mathcal{B}}\right)$ contains a set $X$ of $v$ points and a set $\mathcal{B}$ of $k$ -element subsets of $X$ called blocks such that each $t$ -element subset of $X$ is contained in $\lambda$ blocks. Prove that there are no nontrivial block-transitive 6-designs. (We have shown that there are no nontrivial block-transitive 8-designs and there are certainly some block-transitive, even flag-transitive, 5-designs.)
	
	Comment of 2009: In (Finite Geometry and Combinatorics (Deinze 1992), Cambridge Univ. Press,1993,103-119) we showed that a block-transitive group $G$ on a nontrivial 6-design is either an affine group ${AGL}\left( {d,2}\right)$ or is between ${PSL}\left( {2, q}\right)$ and ${P\Gamma L}\left( {2, q}\right)$ ; in (M. Huber, J. Combin. Theory Ser. A, 117, no. 2 (2010), 196-203) it is shown that for the case $\lambda  = 1$ the group $G$ may only be ${P\Gamma L}\left( {2,{p}^{e}}\right)$ , where $p$ is 2 or 3 and $e$ is an odd prime power. Comment of 2013: In the case $\lambda  = 1$ there are no block-transitive 7-designs (M. Huber, Discrete Math. Theor. Comput. Sci., 12, no. 1 (2010),123-132). Comment of 2021: There are no nontrivial 6-designs with $k \leq  {10}^{4}$ in the case where the automorphism group is almost simple (Q. Tan, W. Liu, J. Chen, Algebra Colloq., 21, no. 2 (2014), 231-234). P. J. Cameron, C. E. Praeger
\end{openproblem}

\plainproblemsection

\begin{openproblem}
	11.46. a) Does there exist a finite 3-group $G$ of nilpotency class 3 with the property $\left\lbrack  {a,{a}^{\varphi }}\right\rbrack   = 1$ for all $a \in  G$ and all endomorphisms $\varphi$ of $G$ ? (See A. Caranti, J. Algebra, 97, no. 1 (1985), 1-13.) A. Caranti
\end{openproblem}

\plainproblemsection

\begin{openproblem}
	11.46. b) Does there exist a finite $p$ -group of nilpotency class greater than 2, with Aut $G = {\operatorname{Aut}}_{c}G \cdot  \operatorname{Inn}G$ , where ${\operatorname{Aut}}_{c}G$ is the group of central automorphisms of $G$ ?
	
	c) Does there exist a 2-Engel finite $p$ -group $G$ of nilpotency class greater than 2 such that $\operatorname{Aut}G = {\operatorname{Aut}}_{c}G \cdot  \operatorname{Inn}G$ , where ${\operatorname{Aut}}_{c}G$ is the group of central automorphisms of $G$ ? A. Caranti
	
	b) Yes, it exists (I. Malinowska, Rend. Sem. Mat. Padova, 91 (1994), 265-271).
	
	c) Yes, it does (A. Abdollahi, A. Faghihi, S. A. Linton, E. A. O'Brien, Arch. Math.
	
	(Basel), 95, no. 1 (2010), 1-7).
\end{openproblem}

\plainproblemsection

\begin{openproblem}
	11.47. Let ${\mathcal{L}}_{d}$ be the homogeneous component of degree $d$ in a free Lie algebra $\mathcal{L}$ of rank 2 over the field of order 2 . What is the dimension of the fixed point space in ${\mathcal{L}}_{d}$ for the automorphism of $\mathcal{L}$ which interchanges two elements of a free generating set of $\mathcal{L}$ ? L. G. Kovács
	
	It is found; R. M. Bryant and R. Stöhr (Arch. Math., 67, no. 4 (1996), 281-289) confirmed the conjecture from (M. W. Short, Commun. Algebra, 23, no. 8 (1995), 3051-3057).
\end{openproblem}

\plainproblemsection

\begin{openproblem}
	11.48. Is the commutator $\left\lbrack  {x, y, y, y, y, y, y}\right\rbrack$ a product of fifth powers in the free group $\langle x, y\rangle$ ? If not, then the Burnside group $B\left( {2,5}\right)$ is infinite. A. I. Kostrikin
\end{openproblem}

\plainproblemsection

\begin{openproblem}
	11.49. B. Hartley (Proc. London Math. Soc., 35, no. 1 (1977), 55-75) constructed an example of a non-countable Artinian $\mathbb{Z}G$ -module where $G$ is a metabelian group with the minimum condition for normal subgroups. It follows that there exists a non-countable soluble group (of derived length 3) satisfying Min- $n$ . The following question arises in connection with this result and with the study of some classes of soluble groups with the weak minimum condition for normal subgroups. Is an Artinian $\mathbb{Z}G$ -module countable if $G$ is a soluble group of finite rank (in particular, a minimax group)? L. A. Kurdachenko
\end{openproblem}

\plainproblemsection

\begin{openproblem}
	11.50. Let $A, C$ be abelian groups. If $A\left\lbrack  n\right\rbrack   = 0$ , i. e. for $a \in  A,{na} = 0$ implies $a = 0$ , then the sequence $\frac{\operatorname{Hom}\left( {C, A}\right) }{n\operatorname{Hom}\left( {C, A}\right) } \rightarrowtail  \operatorname{Hom}\left( {\frac{C}{nC},\frac{A}{nA}}\right)  \rightarrow  \operatorname{Ext}\left( {C, A}\right) \left\lbrack  n\right\rbrack$ is exact. Given $f \in  \operatorname{Hom}\left( {C,\frac{A}{nA}}\right)  = \operatorname{Hom}\left( {\frac{C}{nC},\frac{A}{nA}}\right)$ the corresponding extension ${X}_{f}$ is obtained as a pull-back
	
	$$
	\begin{array}{l} A \mapsto  {X}_{f} \rightarrow  C \\   \downarrow  \\  A \mapsto  A \rightarrow  A \rightarrow  \frac{A}{nA} \\  \end{array}
	$$
	
	Use this scheme to classify certain extensions of $A$ by $C$ . The case ${nC} = 0, A$ being torsion-free is interesting. Here $\operatorname{Ext}\left( {C, A}\right) \left\lbrack  n\right\rbrack   = \operatorname{Ext}\left( {C, A}\right)$ . (See E. L. Lady, A. Mader, J. Algebra, 140 (1991), 36-64.) A. Mader
\end{openproblem}

\plainproblemsection

\begin{openproblem}
	11.51. Are there (large, non-trivial) classes $\mathfrak{X}$ of torsion-free abelian groups such that for $A, C \in  \mathfrak{X}$ the group $\operatorname{Ext}\left( {C, A}\right)$ is torsion-free? It is a fact (E.L. Lady, A. Mader, J. Algebra, 140 (1991), 36-64) that two groups of such a class are nearly isomorphic if and only if they have equal $p$ -ranks for all $p$ . A. Mader
\end{openproblem}

\plainproblemsection

\begin{openproblem}
	11.52. (Well-known problem). A permutation group on a set $\Omega$ is called sharply doubly transitive if for any two pairs $\left( {\alpha ,\beta }\right)$ and $\left( {\gamma ,\delta }\right)$ of elements of $\Omega$ such that $\alpha  \neq  \beta$ and $\gamma  \neq  \delta$ , there is exactly one element of the group taking $\alpha$ to $\gamma$ and $\beta$ to $\delta$ . Does every sharply doubly transitive group possess a non-trivial abelian normal subgroup? A positive answer is well known for finite groups. No, not every (E. Rips, Y. Segev, K. Tent, J. Europ. Math. Soc., 19, no. 10 (2017), 2895-2910).
\end{openproblem}

\plainproblemsection

\begin{openproblem}
	11.53. (P. Kleidman). Do the sporadic simple groups of Rudvalis ${Ru}$ , Mathieu ${M}_{22}$ , and Higman-Sims ${HS}$ embed into the simple group ${E}_{7}\left( 5\right)$ ? $\;V.D.{Mazurov}$ Yes, they do (P. B. Kleidman, R. A. Wilson, J. Algebra, 157, no. 2 (1993), 316-330).
\end{openproblem}

\plainproblemsection

\begin{openproblem}
	11.54. Is it true that in the group of coloured braids only the identity braid is a conjugate to its inverse? (For definition see Archive, 10.24.) G. S. Makanin 11.55. Is it true that extraction of roots in the group of coloured braids is uniquely determined? G. S. Makanin
	
	The answers to both 11.54 and 11.55 are affirmative, since the group ${K}_{n}$ of coloured braids is embeddable into the group of those automorphisms of the free group ${F}_{n}$ that act trivially modulo the derived subgroup of ${F}_{n}$ , and this group is residually in the class of torsion-free nilpotent groups, for which the corresponding assertions are true (V. A. Roman'kov, Letter of October, 3, 1990). See also (V. G. Bardakov, Russ. Acad. Sci. Sbornik Math., 76, no. 1 (1993), 123-153).
\end{openproblem}

\plainproblemsection

\begin{openproblem}
	11.56. a) Does every infinite residually finite group contain an infinite abelian subgroup? This is equivalent to the following: does every infinite residually finite group contain a non-identity element with an infinite centralizer?
	
	By a famous theorem of Shunkov a torsion group with an involution having a finite centralizer is a virtually soluble group. Therefore we may assume that in our group all elements have odd order. One should start, perhaps, with the following:
	
	b) Does every infinite residually $p$ -group contain an infinite abelian subgroup?
	
	A. Mann
\end{openproblem}

\plainproblemsection

\begin{openproblem}
	11.57. An upper composition factor of a group $G$ is a composition factor of some finite quotient of $G$ . Is there any restriction on the set of non-abelian upper composition factors of a finitely generated group?
	
	There are no restrictions: any set of non-abelian finite simple groups can be the set of upper composition factors of a 63-generator group; if we allow the group to have also abelian upper composition factors, the number of generators can be reduced to 3 (D. Segal, Proc. London Math. Soc. (3), 82 (2001), 597-613).
\end{openproblem}

\plainproblemsection

\begin{openproblem}
	11.58. Describe the finite groups which contain a tightly embedded subgroup $H$ such that a Sylow 2-subgroup of $H$ is a direct product of a quaternion group of order 8 and a non-trivial elementary group. A. A. Makhnëv
\end{openproblem}

\plainproblemsection

\begin{openproblem}
	11.59. A ${TI}$ -subgroup $A$ of a group $G$ is called a subgroup of root type if $\left\lbrack  {A,{A}^{g}}\right\rbrack   = 1$ whenever ${N}_{A}\left( {A}^{g}\right)  \neq  1$ . Describe the finite groups containing a cyclic subgroup of order 4 as a subgroup of root type. A. A. Makhnëv
\end{openproblem}

\plainproblemsection

\begin{openproblem}
	11.60. Is it true that the hypothetical Moore graph with 3250 vertices of valence 57 has no automorphisms of order 2? M. Aschbacher (J. Algebra, 19, no. 4 (1971), 538-540) proved that this graph is not a graph of rank 3 . A. A. Makhnëv
\end{openproblem}

\plainproblemsection

\begin{openproblem}
	11.61. Let $F$ be a non-abelian free pro- $p$ -group. Is it true that the subset $\{ r \in  F \mid$ $\operatorname{cd}\left( {F/\left( r\right) }\right)  \leq  2\}$ is dense in $F$ ? Here(r)denotes the closed normal subgroup of $F$ generated by $r$ . O. V. Mel'nikov
\end{openproblem}

\plainproblemsection

\begin{openproblem}
	11.62. Describe the groups over which any equation is soluble. In particular, is it true that this class of groups coincides with the class of torsion-free groups? It is easy to see that a group, over which every equation is soluble, is torsion-free. On the other hand, S. D. Brodskiĭ (Siberian Math. J., 25, no. 2 (1984), 235-251) showed that any equation is soluble over a locally indicable group. D. I. Moldavanskiĭ
\end{openproblem}

\plainproblemsection

\begin{openproblem}
	11.63. Suppose that $G$ is a one-relator group containing non-trivial elements of finite order and $N$ is a subgroup of $G$ generated by all elements of finite order. Is it true that any subgroup of $G$ that intersects $N$ trivially is a free group? One can show that the answer is affirmative in the cases where $G/N$ has non-trivial centre or satisfies a non-trivial identity. D. I. Moldavanskii
\end{openproblem}

\plainproblemsection

\begin{openproblem}
	11.64. Let $\pi \left( G\right)$ denote the set of prime divisors of the order of a finite group $G$ . Are there only finitely many finite simple groups $G$ , different from alternating groups, which have a proper subgroup $H$ such that $\pi \left( H\right)  = \pi \left( G\right)$ ? V. S. Monakhov No, there are infinitely many such groups. If $G = {S}_{4k}\left( {2}^{s}\right)$ and $H \cong  {\Omega }_{4k}^{ - }\left( {2}^{s}\right)$ , then $\pi \left( G\right)  = \pi \left( H\right)$ (V. I. Zenkov, Letter of March,10,1994.)
\end{openproblem}

\plainproblemsection

\begin{openproblem}
	11.65. Conjecture: any finitely generated soluble torsion-free pro- $p$ -group with decidable elementary theory is an analytic pro- $p$ -group.
	
	A. G. Myasnikov, V. N. Remeslennikov
\end{openproblem}

\plainproblemsection

\begin{openproblem}
	11.66. (Yu. L. Ershov). Is the elementary theory of a free pro- $p$ -group decidable? A. G. Myasnikov, V. N. Remeslennikov
\end{openproblem}

\plainproblemsection

\begin{openproblem}
	11.67. Does there exist a torsion-free group with exactly 3 classes of conjugate elements such that no non-trivial conjugate class contains a pair of inverse elements?
	
	B. Neumann
\end{openproblem}

\plainproblemsection

\begin{openproblem}
	11.68. Can every fully ordered group be embedded in a fully ordered group (continuing the given order) with only 3 classes of conjugate elements? No, not every (V. V. Bludov, Algebra and Logic, 44, no. 6 (2005), 370-380).
\end{openproblem}

\plainproblemsection

\begin{openproblem}
	11.69. A group $G$ acting on a set $\Omega$ will be said to be1- $\{ 2\}$ -transitive if it acts transitively on the set ${\Omega }^{1,\{ 2\} } = \{ \left( {\alpha ,\{ \beta ,\gamma \} }\right)  \mid  \alpha ,\beta ,\gamma$ distinct $\}$ . Thus $G$ is 1- $\{ 2\}$ -transitive if and only if it is transitive and a stabilizer ${G}_{\alpha }$ is 2-homogeneous on $\Omega  \smallsetminus  \{ \alpha \}$ . The problem is to classify all (infinite) permutation groups that are $1 - \{ 2\}$ -transitive but not 3-transitive. Peter M. Neumann
\end{openproblem}

\plainproblemsection

\begin{openproblem}
	11.70. Let $F$ be an infinite field or a skew-field.
	
	a) Find all transitive subgroups of ${PGL}\left( {2, F}\right)$ acting on the projective line $F \cup  \{ \infty \}$ .
	
	c) What are the flag-transitive subgroups of ${PGL}\left( {d + 1, F}\right)$ ?
	
	d) What subgroups of ${PGL}\left( {d + 1, F}\right)$ are 2-transitive on the points of ${PG}\left( {d, F}\right)$ ?
	
	Peter M. Neumann, C. E. Praeger
\end{openproblem}

\plainproblemsection

\begin{openproblem}
	11.70. b) Let $F$ be an infinite field or a skew-field. What conditions on the subring $R$ of $F$ will ensure that ${PGL}\left( {d + 1, R}\right)$ is flag-transitive on the projective $d$ -space ${PG}\left( {d, F}\right)$ ? Peter M. Neumann, C. E. Praeger
	
	b) Such necessary and sufficient conditions are given in (S. A. Zyubin, Siberian Electron. Math. Rep., 11 (2014), 64-69 (Russian)).
\end{openproblem}

\plainproblemsection

\begin{openproblem}
	11.71. Let $\mathrm{A}$ be a finite group with a normal subgroup $H$ . A subgroup $U$ of $H$ is called an $A$ -covering subgroup of $H$ if $\mathop{\bigcup }\limits_{{a \in  A}}{U}^{a} = H$ . Is there a function $f : \mathbb{N} \rightarrow  \mathbb{N}$ such that whenever $U < H < A$ , where $A$ is a finite group, $H$ is a normal subgroup of $A$ of index $n$ , and $U$ is an $A$ -covering subgroup of $H$ , the index $\left| {H : U}\right|  \leq  f\left( n\right)$ ? (We have shown that the answer is "yes" if $U$ is a maximal subgroup of $H$ .)
	
	Peter M. Neumann, C. E. Praeger
\end{openproblem}

\plainproblemsection

\begin{openproblem}
	11.72. Suppose that a variety of groups $\mathfrak{V}$ is non-regular, that is, the free group ${F}_{n + 1}\left( \mathfrak{V}\right)$ is embeddable in ${F}_{n}\left( \mathfrak{V}\right)$ for some $n$ . Is it true that then every countable group in $\mathfrak{V}$ is embeddable in an $n$ -generated group in $\mathfrak{V}$ ? A. Yu. Ol'shanskii
\end{openproblem}

\plainproblemsection

\begin{openproblem}
	11.73. If a relatively free group is finitely presented, is it almost nilpotent?
	
	A. Yu. Ol'shanskii
\end{openproblem}

\plainproblemsection

\begin{openproblem}
	11.74. Let $G$ be a non-elementary hyperbolic group and let ${G}^{n}$ be the subgroup generated by the $n$ th powers of the elements of $G$ .
	
	a) (M. Gromov). Is it true that $G/{G}^{n}$ is infinite for some $n = n\left( G\right)$ ?
	
	b) Is it true that $\mathop{\bigcap }\limits_{{n = 1}}^{\infty }{G}^{n} = \{ 1\}$ ?
	
	a) Yes, it is; b) Yes, it is (S. V. Ivanov, A. Yu. Ol'shanskii, Trans. Amer. Math. Soc., 348, N 6 (1996), 2091-2138).
\end{openproblem}

\plainproblemsection

\begin{openproblem}
	11.75. Let us consider the class of groups with $n$ generators and $m$ relators. A subclass of this class is called dense if the ratio of the number of presentations of the form $\left\langle  {{a}_{1},\ldots ,{a}_{n} \mid  {R}_{1},\ldots ,{R}_{m}}\right\rangle$ (where $\left| {R}_{i}\right|  = {d}_{i}$ ) for groups from this subclass to the number of all such presentations converges to 1 when ${d}_{1} + \cdots  + {d}_{m}$ tends to infinity. Prove that for every $k < m$ and for any $n$ the subclass of groups all of whose $k$ -generator subgroups are free is dense. A. Yu. Ol'shanskiĭ
	
	This is proved (G. N. Arzhantseva, A. Yu. Ol'shanskiĭ, Math. Notes, 59, no. 4 (1996), 350-355).
\end{openproblem}

\plainproblemsection

\begin{openproblem}
	11.76. (Well-known problem). Is the group of collineations of a finite non-Desar-guasian projective plane defined over a semi-field soluble? (The hypothesis on the plane means that the corresponding regular set, see Archive, 10.48, is closed under addition.) N. D. Podufalov
\end{openproblem}

\plainproblemsection

\begin{openproblem}
	11.77. (Well-known problem). Describe the finite translation planes whose collineation groups act doubly transitively on the set of points of the line at infinity.
	
	N. D. Podufalov
\end{openproblem}

\plainproblemsection

\begin{openproblem}
	11.78. An isomorphism of groups of points of algebraic groups is called semialgebraic if it can be represented as a composition of an isomorphism of translation of the field of definition and a rational morphism.
	
	a) Is it true that the existence of an isomorphism of groups of points of two directly undecomposable algebraic groups with trivial centres over an algebraically closed field implies the existence of a semialgebraic isomorphism of the groups of points?
	
	b) Is it true that every isomorphism of groups of points of directly undecomposable algebraic groups with trivial centres defined over algebraic fields is semialgebraic? A field is called algebraic if all its elements are algebraic over the prime subfield.
\end{openproblem}

\plainproblemsection

\begin{openproblem}
	11.79. Let $G$ be a finite group of automorphisms of an infinite field $F$ of characteristic $p$ . Taking integral powers of the elements of $F$ and the action of $G$ define the action of the group ring $\mathbb{Z}G$ of $G$ on the multiplicative group of $F$ . Is it true that any subfield of $F$ that contains the images of all elements of $F$ under the action of some fixed element of $\mathbb{Z}G \smallsetminus  p\mathbb{Z}G$ contains infinitely many $G$ -invariant elements of $F$ ?
	
	Yes (K. N. Ponomarëv, Siber. Math. J., 33 (1992), 1094-1099). K. N. Ponomarëv
\end{openproblem}

\plainproblemsection

\begin{openproblem}
	11.80. Let $G$ be a primitive permutation group on a finite set $\Omega$ and suppose that, for $\alpha  \in  \Omega ,{G}_{\alpha }$ acts 2-transitively on one of its orbits in $\Omega  \smallsetminus  \{ \alpha \}$ . By (C. E. Praeger, J. Austral. Math. Soc. (A), 45, 1988, 66-77) either
	
	(a) $T \leq  G \leq  \operatorname{Aut}T$ for some nonabelian simple group $T$ , or
	
	(b) $G$ has a unique minimal normal subgroup which is regular on $\Omega$ . For what classes of simple groups in (a) is a classification feasible? Describe the examples in as explicit a manner as possible. Classify all groups in (b).
	
	Remark of 1999: Two papers by X. G. Fang and C. E. Praeger (Commun. Algebra, 27 (1999), 3727-3754 and 3755-3769) show that this is feasible for $T$ a Suzuki and Ree simple group, and a paper of J. Wang and C. E. Praeger (J. Algebra, 180 (1996), 808-833) suggests that this may not be the case for $T$ an alternating group. Remark of 2001: In (X. G. Fang, G. Havas, J. Wang, European J. Combin., 20, no. 6 (1999), 551-557) new examples are constructed with $G = {PSU}\left( {3, q}\right)$ . Remark of 2009: It was shown in (D. Leemans, J. Algebra, 322, no. 3 (2009), 882-892) that for part (a) classification is feasible for self-paired 2-transitive suborbits with $T$ a sporadic simple group; in this case these suborbits correspond to 2-arc-transitive actions on undirected graphs. C. E. Praeger
\end{openproblem}

\plainproblemsection

\begin{openproblem}
	11.81. A topological group is said to be $F$ -balanced if for any subset $X$ and any neighborhood of the identity $U$ there is a neighborhood of the identity $V$ such that ${VX} \subseteq  {XU}$ . Is every $F$ -balanced group balanced, that is, does it have a basis of neighborhoods of the identity consisting of invariant sets? I. V. Protasov
\end{openproblem}

\plainproblemsection

\begin{openproblem}
	11.82. Let $R$ be the normal closure of an element $r$ in a free group $F$ with the natural length function and suppose that $s$ is an element of minimal length in $R$ . Is it true that $s$ is conjugate to one of the following elements: $r,{r}^{-1},\left\lbrack  {r, f}\right\rbrack  ,\left\lbrack  {{r}^{-1}, f}\right\rbrack$ for some $f \in  {F}^{\prime }$ V. N. Remeslennikov
	
	No, not always (J. McCool, Glasgow Math. J., 43, no. 1 (2001), 123-124).
\end{openproblem}

\plainproblemsection

\begin{openproblem}
	11.83. Is the conjugacy problem soluble for finitely generated abelian-by-polycyclic groups? V. N. Remeslennikov
\end{openproblem}

\plainproblemsection

\begin{openproblem}
	11.84. Is the isomorphism problem soluble
	
	a) for finitely generated metabelian groups?
	
	b) for finitely generated soluble groups of finite rank? V. N. Remeslennikov
\end{openproblem}

\plainproblemsection

\begin{openproblem}
	11.85. Let $F$ be a free pro- $p$ -group with a basis $X$ and let $R = {r}^{F}$ be the closed normal subgroup generated by an element $r$ . We say that an element $s$ of $F$ is associated to $r$ if ${s}^{F} = R$ .
	
	a) Suppose that none of the elements associated to $r$ is a $p$ th power of an element in $F$ . Is it true that $F/R$ is torsion-free?
	
	b) Among the elements associated to $r$ there is one that depends on the minimal subset ${X}^{\prime }$ of the basis $X$ . Let $x \in  {X}^{\prime }$ . Is it true that the images of the elements of $X \smallsetminus  \{ x\}$ in $F/R$ freely generate a free pro- $p$ -group? N. S. Romanovskii
\end{openproblem}

\plainproblemsection

\begin{openproblem}
	11.86. Does every group $G = \left\langle  {{x}_{1},\ldots ,{x}_{n} \mid  {r}_{1} = \cdots  = {r}_{m} = 1}\right\rangle$ possess, in a natural way, a homomorphic image $H = \left\langle  {{x}_{1},\ldots ,{x}_{n} \mid  {s}_{1} = \cdots  = {s}_{m} = 1}\right\rangle$
	
	a) such that $H$ is a torsion-free group?
	
	b) such that the integral group ring of $H$ is embeddable in a skew field?
	
	If the stronger assertion "b" is true, then this will give an explicit method of finding elements ${x}_{{i}_{1}},\ldots ,{x}_{{i}_{n - m}}$ which generate a free group in $G$ . Such elements exist by N. S. Romanovskiĭ's theorem (Algebra and Logic, 16, no. 1 (1977), 62-67).
	
	V. A. Roman'kov
\end{openproblem}

\plainproblemsection

\begin{openproblem}
	11.87. (Well-known problem). Is the automorphism group of a free metabelian group of rank $n \geq  4$ finitely presented? V. A. Roman'kov
\end{openproblem}

\plainproblemsection

\begin{openproblem}
	11.88. We define the length $l\left( g\right)$ of an Engel element $g$ of a group $G$ to be the smallest number $l$ such that $\left\lbrack  {h, g;l}\right\rbrack   = 1$ for all $h \in  G$ . Here $\left\lbrack  {h, g;1}\right\rbrack   = \left\lbrack  {h, g}\right\rbrack$ and $\left\lbrack  {h, g;i + 1}\right\rbrack   = \left\lbrack  {\left\lbrack  {h, g;i}\right\rbrack  , g}\right\rbrack$ . Does there exist a polynomial function $\varphi \left( {x, y}\right)$ such that $l\left( {uv}\right)  \leq  \varphi \left( {l\left( u\right) , l\left( v\right) }\right)$ ? Up to now, it is unknown whether a product of Engel elements is again an Engel element. V. A. Roman'kov
	
	No, it does not (L. V. Dolbak, Siberian Math. J., 47, no. 1 (2006), 55-57).
\end{openproblem}

\plainproblemsection

\begin{openproblem}
	11.89. Let $k$ be an infinite cardinal number. Describe the epimorphic images of the Cartesian power $\mathop{\prod }\limits_{k}\mathbb{Z}$ of the group $\mathbb{Z}$ of integers. Such a description is known for
	
	$k = {\omega }_{0}$ (R. Nunke, Acta Sci. Math. (Szeged), 23 (1963), 67-73). S. V. Rychkov
\end{openproblem}

\plainproblemsection

\begin{openproblem}
	11.90. Let $\mathfrak{V}$ be a variety of groups and let $k$ be an infinite cardinal number. A group $G \in  \mathfrak{V}$ of rank $k$ is called almost free in $\mathfrak{V}$ if each of its subgroups of rank less than $k$ is contained in a subgroup of $G$ which is free in $\mathfrak{V}$ . For what $k$ do there exist almost free but not free in $\mathfrak{B}$ groups of rank $k$ ? S. V. Rychkov
\end{openproblem}

\plainproblemsection

\begin{openproblem}
	11.91. Prove that a hereditary formation $\mathfrak{F}$ of finite soluble groups is local if every finite soluble non-simple minimal non- $\xi$ -group is a Shmidt group (that is, a non-nilpotent finite group all of whose proper subgroups are nilpotent). V. N. Semenchuk This is proved (A. N. Skiba, Dokl. Akad. Nauk Belorus. SSR, 34, no. 11 (1990), 982-985 (Russian)).
\end{openproblem}

\plainproblemsection

\begin{openproblem}
	11.92. What are the soluble hereditary non-one-generator formations of finite groups all of whose proper hereditary subformations are one-generated? A. N. Skiba
\end{openproblem}

\plainproblemsection

\begin{openproblem}
	11.93. Is the variety of all lattices generated by the block lattices (see Archive, 9.62) of finite groups? D. M. Smirnov
	
	No, it is not (D. M. Smirnov, Siberian Math. J., 33, no. 4 (1992), 663-668).
\end{openproblem}

\plainproblemsection

\begin{openproblem}
	11.94. Describe all simply reducible groups, that is, groups such that all their characters are real and the tensor product of any two irreducible representations contains no multiple components. This question is interesting for physicists. Is every finite simply reducible group soluble? S. P. Strunkov
	
	Yes, it is (L. S. Kazarin, E. I. Chankov, Sb. Math., 201, no. 5 (2010), 655-668).
\end{openproblem}

\plainproblemsection

\begin{openproblem}
	11.95. Suppose that $G$ is a $p$ -group $G$ containing an element $a$ of order $p$ such that the subgroup $\left\langle  {a,{a}^{g}}\right\rangle$ is finite for any $g$ and the set ${C}_{G}\left( a\right)  \cap  {a}^{G}$ is finite. Is it true that $G$ has non-trivial centre? This is true for 2-groups. S. P. Strunkov
\end{openproblem}

\plainproblemsection

\begin{openproblem}
	*11.96. (a) Is it true that, for a given number $n$ , there exist only finitely many finite simple groups each of which contains an involution which commutes with at most $n$ involutions of the group?
	
	(b) Is it true that there are no infinite simple groups satisfying this condition? S. P. Strunkov
	
	*(a) Yes, it is true (mod CFSG); moreover, if a locally finite group has such an involution, then it is finite of $n$ -bounded order (S. V. Skresanov, https://arxiv.org/ pdf/2402.05437.pdf).
	
	*(b) No, since ${PS}{L}_{2}\left( \mathbb{R}\right)$ is an infinite simple group, which contains an involution but does not contain a subgroup isomorphic to the Klein 4-group (S. V. Skresanov, https://arxiv.org/pdf/2402.05437.pdf).
\end{openproblem}

\plainproblemsection

\begin{openproblem}
	11.97. Are there only finitely many finite simple groups with a given set of all different values of irreducible characters on a single element? S. P. Strunkov No; all complex irreducible characters of the groups ${L}_{2}\left( {2}^{m}\right) , m \geq  2$ , take the values $0, \pm  1$ on involutions (V. D. Mazurov).
\end{openproblem}

\plainproblemsection

\begin{openproblem}
	11.98. a) (R. Brauer). Find the best-possible estimate of the form $\left| G\right|  \leq  f\left( r\right)$ where $r$ is the number of conjugacy classes of elements in a finite (simple) group $G$ .
	
	S. P. Strunkov
\end{openproblem}

\plainproblemsection

\begin{openproblem}
	11.98. b) Let $r$ be the number of conjugacy classes of elements in a finite (simple) group $G$ . Is it true that $\left| G\right|  \leq  \exp \left( r\right)$ ?
	
	No, this is not true: the group ${M}_{22}$ has order 443520 and contains 12 conjugacy classes (T. Plunkett, Letter of May, 9, 2000).
\end{openproblem}

\plainproblemsection

\begin{openproblem}
	11.99. Find (in group-theoretic terms) necessary and sufficient conditions for a finite group to have complex irreducible characters having defect 0 for more than one prime number dividing the order of the group. Express the number of such characters in the same terms. S. P. Strunkov
\end{openproblem}

\plainproblemsection

\begin{openproblem}
	11.100. Is it true that every periodic conjugacy biprimitively finite group (see 6.59) can be obtained from 2-groups and binary finite groups by taking extensions? This question is of independent interest for $p$ -groups, $p$ an odd prime. S. P. Strunkov
\end{openproblem}

\plainproblemsection

\begin{openproblem}
	11.101. Does there exist a Golod group (see 9.76) with finite centre?
	
	A. V. Timofeenko
	
	Yes, it exists (V. A. Sereda, A. I. Sozutov, Algebra and Logic, 45, no. 2 (2006), 134- 138).
\end{openproblem}

\plainproblemsection

\begin{openproblem}
	11.102. Does there exist a residually soluble but insoluble group satisfying the maximum condition on subgroups? J. Wiegold
\end{openproblem}

\plainproblemsection

\begin{openproblem}
	11.103. Is a 2-group satisfying the minimum condition for centralizers necessarily locally finite? J. Wilson
	
	Yes, it is (F. O. Wagner, J. Algebra, 217, no. 2 (1999), 448-460).
\end{openproblem}

\plainproblemsection

\begin{openproblem}
	11.104. Let $G$ be a finite group of order ${p}^{a} \cdot  {q}^{b}\cdots$ , where $p, q,\ldots$ are distinct primes. Introduce distinct variables ${x}_{p},{x}_{q},\ldots$ corresponding to $p, q,\ldots$ . Define functions $f,\varphi$ from the lattice of subgroups of $G$ to the polynomial ring $\mathbb{Z}\left\lbrack  {{x}_{p},{x}_{q},\ldots }\right\rbrack$ as follows:
	
	(1) if $H$ has order ${p}^{\alpha } \cdot  {q}^{\beta }\cdots$ , then $f\left( H\right)  = {x}_{p}^{\alpha } \cdot  {x}_{q}^{\beta }\cdots$ ;
	
	(2) for all $H \leq  G$ , we have $\mathop{\sum }\limits_{{K \leq  H}}\varphi \left( K\right)  = f\left( H\right)$ .
	
	Then $f\left( G\right) ,\varphi \left( G\right)$ may be called the order and Eulerian polynomials of $G$ . Substituting ${p}^{m},{q}^{m},\ldots$ for ${x}_{p},{x}_{q},\ldots$ in these polynomials we get the $m$ th power of the order of $G$ and the number of ordered $m$ -tuples of elements that generate $G$ respectively.
	
	It is known that if $G$ is $p$ -solvable, then $\varphi \left( G\right)$ is a product of a polynomial in ${x}_{p}$ and a polynomial in the remaining variables. Consequently, if $G$ is solvable, $\varphi \left( G\right)$ is the product of a polynomial in ${x}_{p}$ by a polynomial in ${x}_{q}$ by . . . . Are the converses of these statements true a) for solvable groups? b) for $p$ -solvable groups? G. E. Wall Yes, the converses are true: a) for solvable groups (E. Detomi, A. Lucchini, J. London Math. Soc. (2), 70 (2004), 165-181); b) for p-solvable groups (E. Damian, A. Lucchini, Commun. Algebra, 35 (2007), 3451-3472).
\end{openproblem}

\plainproblemsection

\begin{openproblem}
	11.105. a) Let $\mathfrak{V}$ be a variety of groups. Its relatively free group of given rank has a presentation $F/N$ , where $F$ is absolutely free of the same rank and $N$ fully invariant in $F$ . The associated Lie ring $\mathcal{L}\left( {F/N}\right)$ has a presentation $L/J$ , where $L$ is the free Lie ring of the same rank and $J$ an ideal of $L$ . Is $J$ always fully invariant in $L$ ?
	
	G. E. Wall
	
	No, not always; the ideal $J$ is not fully invariant for $F/{\left( {F}^{2}\right) }^{4}$ , that is, for $\mathfrak{V} = {\mathfrak{B}}_{4}{\mathfrak{B}}_{2}$ (D. Groves, J. Algebra, 211, no. 1 (1999), 15-25).
\end{openproblem}

\plainproblemsection

\begin{openproblem}
	11.105. b) Let $\mathfrak{V}$ be a variety of groups. Its relatively free group of given rank has a presentation $F/N$ , where $F$ is absolutely free of the same rank and $N$ fully invariant in $F$ . The associated Lie ring $\mathcal{L}\left( {F/N}\right)$ has a presentation $L/J$ , where $L$ is the free Lie ring of the same rank and $J$ an ideal of $L$ . Is $J$ fully invariant in $L$ if $\mathfrak{V}$ is the Burnside variety of all groups of given exponent $q$ , where $q$ is a prime-power, $q \geq  4$ ?
	
	G. E. Wall
\end{openproblem}

\plainproblemsection

\begin{openproblem}
	11.106. Can every periodic group be embedded in a simple periodic group? Yes (A. Yu. Ol'shanskii, Ukrain. Math. J., 44 (1992), 761-763). R. Phillips
\end{openproblem}

\plainproblemsection

\begin{openproblem}
	11.107. Is there a non-linear locally finite simple group each of whose proper subgroups is residually finite? R. Phillips
\end{openproblem}

\plainproblemsection

\begin{openproblem}
	11.108. Is every locally finite simple group absolutely simple? A group $G$ is said to be absolutely simple if the only composition series of $G$ is $\{ 1, G\}$ . For equivalent formulations see (R. E. Phillips, Rend. Sem. Mat. Univ. Padova, 79 (1988), 213- 220). R. Phillips
	
	No, not every (U. Meierfrankenfeld, in: Proc. Int. Conf. Finite and Locally Finite Groups, Istanbul, 1994, Klüwer, 1995, 189-212).
\end{openproblem}

\plainproblemsection

\begin{openproblem}
	11.109. Is it true that the union of an ascending series of groups of Lie type of equal ranks over some fields should be also a group of Lie type of the same rank over a field? (For locally finite fields the answer is "yes", the theorem in this case is due to V. Belyaev, A. Borovik, B. Hartley \& G. Shute, and S. Thomas.) R. Phillips
\end{openproblem}

\plainproblemsection

\begin{openproblem}
	11.110. Is it possible to embed ${SL}\left( {2,\mathbb{Q}}\right)$ in the multiplicative group of some division ring? B. Hartley
	
	No, it is not (W. Dicks, B. Hartley, Commun. Algebra, 19 (1991), 1919-1943).
\end{openproblem}

\plainproblemsection

\begin{openproblem}
	11.111. Does every infinite locally finite simple group $G$ contain a nonabelian finite simple subgroup? Is there an infinite tower of such subgroups in $G$ ? B. Hartley
\end{openproblem}

\plainproblemsection

\begin{openproblem}
	11.112. Let $L = L\left( {K\left( p\right) }\right)$ be the associated Lie ring of a free countably generated group $K\left( p\right)$ of the Kostrikin variety of locally finite groups of a given prime exponent $p$ . Is it true that
	
	a) $L$ is a relatively free Lie ring?
	
	c) all identities of $L$ follow from a finite number of identities of $L$ ?
	
	E. I. Khukhro
\end{openproblem}

\plainproblemsection

\begin{openproblem}
	11.112. b) Let $L = L\left( {K\left( p\right) }\right)$ be the associated Lie ring of a free countably generated group $K\left( p\right)$ of the Kostrikin variety of locally finite groups of a given prime exponent $p$ . Is it true that all identities of $L$ follow from multilinear identities of $L$ ? E. I. Khukhro No, there are two relations of weight 29 which hold in $L\left( {K\left( 7\right) }\right)$ (in two generators, of multiweights(14,15)and(15,14)) that are not consequences of multilinear relations (E. O'Brien, M. R. Vaughan-Lee, Int. J. Algebra Comput., 12 (2002), 575-592; M. F. Newman, M. R. Vaughan-Lee, Electron. Res. Announc. Amer. Math. Soc., 4, no. 1 (1998), 1-3).
\end{openproblem}

\plainproblemsection

\begin{openproblem}
	11.113. (B. Hartley). Is it true that the derived length of a nilpotent periodic group admitting a regular automorphism of prime-power order ${p}^{n}$ is bounded by a function of $p$ and $n$ ? E. I. Khukhro
\end{openproblem}

\plainproblemsection

\begin{openproblem}
	11.114. Is every locally graded group of finite special rank almost hyperabelian? This is true in the class of periodic locally graded groups (Ukrain. Math. J., 42, no. 7 (1990), 855-861). N. S. Chernikov
\end{openproblem}

\plainproblemsection

\begin{openproblem}
	11.115. Suppose that $X$ is a free non-cyclic group, $N$ is a non-trivial normal subgroup of $X$ and $T$ is a proper subgroup of $X$ containing $N$ . Is it true that $\left\lbrack  {T, N}\right\rbrack   < \left\lbrack  {X, N}\right\rbrack$ if $N$ is not maximal in $X$ ? V. P. Shaptala
\end{openproblem}

\plainproblemsection

\begin{openproblem}
	11.116. The dimension of a partially ordered set $\langle P, \leq  \rangle$ is, by definition, the least cardinal number $\delta$ such that the relation $\leq$ is an intersection of $\delta$ relations of linear order on $P$ . Is it true that, for any Chernikov group which does not contain a direct product of two quasicyclic groups over the same prime number, the subgroup lattice has finite dimension? Expected answer: yes. L. N. Shevrin
\end{openproblem}

\plainproblemsection

\begin{openproblem}
	11.117. Let $\mathfrak{X}$ be a soluble non-empty Fitting class. Is it true that every finite non-soluble group possesses an $\mathfrak{X}$ -injector? L. A. Shemetkov
\end{openproblem}

\plainproblemsection

\begin{openproblem}
	11.118. What are the hereditary soluble local formations $\mathfrak{F}$ of finite groups such that every finite group has an $\mathfrak{F}$ -covering subgroup? L. A. Shemetkov
\end{openproblem}

\plainproblemsection

\begin{openproblem}
	11.119. Is it true that, for any non-empty set of primes $\pi$ , the $\pi$ -length of any finite $\pi$ -soluble group does not exceed the derived length of its Hall $\pi$ -subgroup?
	
	L. A. Shemetkov
\end{openproblem}

\plainproblemsection

\begin{openproblem}
	11.120. Let $\mathfrak{F}$ be a soluble saturated Fitting formation. Is it true that ${l}_{\mathfrak{F}}\left( G\right)  \leq$ $f\left( {{c}_{\mathfrak{F}}\left( G\right) }\right)$ , where ${c}_{\mathfrak{F}}\left( G\right)$ is the length of a composition series of some $\mathfrak{F}$ -covering subgroup of a finite soluble group $G$ ? This is Problem 17 in (L. A. Shemetkov, Formations of Finite Groups, Moscow, Nauka,1978 (Russian)). For the definition of the $\mathfrak{F}$ -length ${l}_{\mathfrak{F}}\left( G\right)$ , see ibid. L. A. Shemetkov
\end{openproblem}

\plainproblemsection

\begin{openproblem}
	11.121. Does there exist a local formation of finite groups which has a non-trivial decomposition into a product of two formations and in any such a decomposition both factors are non-local? Local products of non-local formations do exist.
	
	L. A. Shemetkov
\end{openproblem}

\plainproblemsection

\begin{openproblem}
	11.122. Does every non-zero submodule of a free module over the group ring of a torsion-free group contain a free cyclic submodule? A. L. Shmel'kin
\end{openproblem}

\plainproblemsection

\begin{openproblem}
	11.123. (Well-known problem). For a given group $G$ , define the following sequence of groups: ${A}_{1}\left( G\right)  = G,{A}_{i + 1}\left( G\right)  = \operatorname{Aut}\left( {{A}_{i}\left( G\right) }\right)$ . Does there exist a finite group $G$ for which this sequence contains infinitely many non-isomorphic groups? M. Short
\end{openproblem}

\plainproblemsection

\begin{openproblem}
	11.124. Let $F$ be a non-cyclic free group and $R$ a non-cyclic subgroup of $F$ . Is it true that if $\left\lbrack  {R, R}\right\rbrack$ is a normal subgroup of $F$ then $R$ is also a normal subgroup of $F$ ? V. E. Shpilrain
\end{openproblem}

\plainproblemsection

\begin{openproblem}
	11.125. Let $G$ be a finite group admitting a regular elementary abelian group of automorphisms $V$ of order ${p}^{n}$ . Is it true that the subgroup $H = \mathop{\bigcap }\limits_{{v \in  V\smallsetminus \{ 1\} }}\left\lbrack  {G, v}\right\rbrack$ is
	
	nilpotent? In the case of an affirmative answer, does there exist a function depending only on $p, n$ , and the derived length of $G$ which bounds the nilpotency class of $H$ ? P. V. Shumyatskiĭ
	
	V. P. Shunkov
\end{openproblem}

\plainproblemsection

\begin{openproblem}
	11.126. Do there exist a constant $h$ and a function $f$ with the following property: if a finite soluble group $G$ admits an automorphism $\varphi$ of order 4 such that $\left| {{C}_{G}\left( \varphi \right) }\right|  \leq  m$ , then $G$ has a normal series $1 \leq  M \leq  N \leq  G$ such that the index $\left| {G : N}\right|$ does not exceed $f\left( m\right)$ , the group $N/M$ is nilpotent of class $\leq  2$ , and the group $M$ is nilpotent of class $\leq  h$ ? P. V. Shumyatskiĭ
	
	Yes, they exist (N. Yu. Makarenko, E. I. Khukhro, Algebra and Logic, 45 (2006), 326- 343).
\end{openproblem}

\plainproblemsection

\begin{openproblem}
	11.128. A group $G$ is said to be a $K$ -group if for every subgroup $A \leq  G$ there exists a subgroup $B \leq  G$ such that $A \cap  B = 1$ and $\langle A, B\rangle  = G$ . Is it true that normal subgroups of $K$ -groups are also $K$ -groups? M. Emaldi No, it is not (V. N. Obraztsov, J. Austral. Math. Soc. (Ser. A), 61, no. 2 (1996), 267-288).
\end{openproblem}

\plainproblemsection

\begin{openproblem}
	12.1. a) H. Bass (Topology, 4, no. 4 (1966), 391-400) has constructed explicitly a proper subgroup of finite index in the group of units of the integer group ring of a finite cyclic group. Calculate the index of Bass' subgroup. R. Zh. Aleev
\end{openproblem}

\plainproblemsection

\begin{openproblem}
	12.1. b) (A. A. Bovdi). H. Bass (Topology, 4, no. 4 (1966), 391-400) has constructed explicitly a proper subgroup of finite index in the group of units of the integer group ring of a finite cyclic group. Construct analogous subgroups for finite abelian noncyclic groups. R. Zh. Aleev
	
	They are constructed: for $p$ -groups, $p \neq  2$ , in (K. Hoechsmann, J. Ritter, J. Pure Appl. Algebra, 68, no. 3 (1990), 325-333); for the general case of groups of central units of integral group rings of arbitrary (not necessarily abelian) finite groups in (R. Zh. Aleev, Mat. Tr. Novosibirsk Inst. Math., 3, no. 1 (2000), 3-37 (Russian)).
\end{openproblem}

\plainproblemsection

\begin{openproblem}
	12.2. Is it true that every non-discrete group topology (in an abelian group) can be strengthened up to a maximal complete group topology? $\;V.I.{Arnautov}$ No, not always, since under the assumption of the Continuum Hypothesis there exist non-complete maximal topologies (V. I. Arnautov, E. G. Zelenyuk, Ukrain. Math. J., 43, no. 1 (1991), 15-20).
\end{openproblem}

\plainproblemsection

\begin{openproblem}
	12.3. A $p$ -group is called thin if every set of pairwise incomparable (by inclusion) normal subgroups contains $\leq  p + 1$ elements. Is the number of thin pro- $p$ -groups finite? R. Brandl
\end{openproblem}

\plainproblemsection

\begin{openproblem}
	12.4. Let $G$ be a group and assume that we have ${\left\lbrack  x, y\right\rbrack  }^{3} = 1$ for all $x, y \in  G$ . Is ${G}^{\prime }$ of finite exponent? Is $G$ soluble? By a result of N. D. Gupta and N. S. Mendelssohn, 1967, we know that ${G}^{\prime }$ is a 3-group. R. Brandl
\end{openproblem}

\plainproblemsection

\begin{openproblem}
	12.5. Does there exist a countable non-trivial filter in the lattice of quasivarieties of metabelian torsion-free groups? A. I. Budkin
	
	No, there is no such filter (S. V. Lenyuk, Siberian Math. J., 39, no. 1 (1998), 57-62).
\end{openproblem}

\plainproblemsection

\begin{openproblem}
	12.6. Let $G$ be a finitely generated group and suppose that $H$ is a $p$ -subgroup of $G$ such that $H$ contains no non-trivial normal subgroups of $G$ and ${HX} = {XH}$ for any subgroup $X$ of $G$ . Is then $G/{C}_{G}\left( {H}^{G}\right)$ a $p$ -group where ${H}^{G}$ is the normal closure of $H$ ? G. Busetto
\end{openproblem}

\plainproblemsection

\begin{openproblem}
	12.7. Is it true that every radical hereditary formation of finite groups is a composition one? A. F. Vasil’ev
	
	No, it is not (S. F. Kamornikov, Math. Notes, 55, no. 6 (1994), 586-588).
\end{openproblem}

\plainproblemsection

\begin{openproblem}
	12.8. Let $\mathcal{V}$ be a non-trivial variety of groups and let ${a}_{1},\ldots ,{a}_{r}$ freely generate a free group ${F}_{r}\left( V\right)$ in $\mathcal{V}$ . We say ${F}_{r}\left( \mathcal{V}\right)$ strongly discriminates $\mathcal{V}$ just in case every finite system of inequalities ${w}_{i}\left( {{a}_{1},\ldots ,{a}_{r},{x}_{1},\ldots ,{x}_{k}}\right)  \neq  1$ for $1 \leq  i \leq  n$ having a solution in some ${F}_{s}\left( \mathcal{V}\right)$ containing ${F}_{r}\left( \mathcal{V}\right)$ as a varietally free factor in the sense of $\mathcal{V}$ , already has a solution in ${F}_{r}\left( \mathcal{V}\right)$ . Does there exist $\mathcal{V}$ such that for some integer $r > 0,{F}_{r}\left( \mathcal{V}\right)$ discriminates but does not strongly discriminate $\mathcal{V}$ ? What about the analogous question for general algebras in the context of universal algebra?
	
	A. M. Gaglione, D. Spellman
\end{openproblem}

\plainproblemsection

\begin{openproblem}
	12.9. Following Bass, call a group tree-free if there is an ordered Abelian group $\Lambda$ and a $\Lambda$ -tree $X$ such that $G$ acts freely without inversion on $X$ .
	
	a) Must every finitely generated tree-free group satisfy the maximal condition for Abelian subgroups?
	
	b) The same question for finitely presented tree-free groups.
	
	A. M. Gaglione, D. Spellman
\end{openproblem}

\plainproblemsection

\begin{openproblem}
	12.10. (P. Neumann). Can the free group on two generators be embedded in Sym ( $\mathbb{N}$ ) so that the image of every non-identity element has only a finite number of orbits? A. M. W. Glass
	
	Yes, it can (H. D. Macpherson, in: Ordered groups and infinite permutation groups, Partially based on Conf. Luminy, France, 1993 (Mathematics and its Applications, 354), Kluwer, Dordrecht, 1996, 221-230).
\end{openproblem}

\plainproblemsection

\begin{openproblem}
	12.11. Suppose that $G, H$ are countable or finite groups and $A$ is a proper subgroup of $G$ and $H$ containing no non-trivial subgroup normal in both $G$ and $H$ . Can $G{ * }_{A}H$ be embedded in $\operatorname{Sym}\left( \mathbb{N}\right)$ so that the image is highly transitive?
	
	Comment of 2013: for partial results, see (S. V. Gunhouse, Highly transitive representations of free products on the natural numbers, Ph. D. Thesis Bowling Green State Univ., 1993; K. K. Hickin, J. London Math. Soc. (2), 46, no. 1 (1992), 81-91).
	
	A. M. W. Glass
\end{openproblem}

\plainproblemsection

\begin{openproblem}
	12.12. Is the conjugacy problem for nilpotent finitely generated lattice-ordered groups soluble? A. M. W. Glass
\end{openproblem}

\plainproblemsection

\begin{openproblem}
	12.13. If $\langle \Omega , \leq  \rangle$ is the countable universal poset, then $G = \operatorname{Aut}\left( {\langle \Omega , \leq  \rangle }\right)$ is simple (A. M. W. Glass, S. H. McCleary, M. Rubin, Math. Z., 214, no. 1 (1993), 55-66). If $H$ is a subgroup of $G$ such that $\left| {G : H}\right|  < {2}^{{\aleph }_{0}}$ and $H$ is transitive on $\Omega$ , does $H = G$ ?
	
	A. M. W. Glass
\end{openproblem}

\plainproblemsection

\begin{openproblem}
	12.14. If $T$ is a countable theory, does there exist a model $\mathfrak{A}$ of $T$ such that the theory of Aut(A)is undecidable? M. Giraudet, A. M. W. Glass
	
	Yes, moreover, every first order theory having infinite models has a model whose automorphism group has undecidable existential theory (V. V. Bludov, M. Giraudet, A. M. W. Glass, G. Sabbagh, in Models, Modules and Abelian Groups, de Gruyter, Berlin, 2008, 325-328).
\end{openproblem}

\plainproblemsection

\begin{openproblem}
	12.15. Suppose that, in a finite 2-group $G$ , any two elements are conjugate whenever their normal closures coincide. Is it true that the derived subgroup of $G$ is abelian?
	
	E. A. Golikova, A. I. Starostin
\end{openproblem}

\plainproblemsection

\begin{openproblem}
	12.16. Is the class of groups of recursive automorphisms of arbitrary models closed with respect to taking free products? S. S. Goncharov
\end{openproblem}

\plainproblemsection

\begin{openproblem}
	12.17. Find a description of autostable periodic abelian groups. S. S. Goncharov
\end{openproblem}

\plainproblemsection

\begin{openproblem}
	12.19. Is it true that, for every $n \geq  2$ and every two epimorphisms $\varphi$ and $\psi$ of a free group ${F}_{2n}$ of rank ${2n}$ onto ${F}_{n} \times  {F}_{n}$ , there exists an automorphism $\alpha$ of ${F}_{2n}$ such that ${\alpha \varphi } = \psi$ ? R. I. Grigorchuk
\end{openproblem}

\plainproblemsection

\begin{openproblem}
	12.20. (Well-known problem). Is R. Thompson's group
	
	$$
	F = \left\langle  {{x}_{0},{x}_{1},\ldots  \mid  {x}_{n}^{{x}_{i}} = {x}_{n + 1}, i < n, n = 1,2,\ldots }\right\rangle   =
	$$
	
	$$
	= \left\langle  {{x}_{0},\ldots ,{x}_{4} \mid  {x}_{1}^{{x}_{0}} = {x}_{2},{x}_{2}^{{x}_{0}} = {x}_{3},{x}_{2}^{{x}_{1}} = {x}_{3},{x}_{3}^{{x}_{1}} = {x}_{4},{x}_{3}^{{x}_{2}} = {x}_{4}}\right\rangle
	$$
	
	amenable? R. I. Grigorchuk
\end{openproblem}

\plainproblemsection

\begin{openproblem}
	12.21. Let $R$ be a commutative ring with identity and let $G$ be a finite group. Prove that the ring $a\left( {RG}\right)$ of $R$ -representations of $G$ has no non-trivial idempotents.
	
	P. M. Gudivok, V. P. Rud'ko
\end{openproblem}

\plainproblemsection

\begin{openproblem}
	12.22. a) Let $\Delta \left( G\right)$ be the augmentation ideal of the integer group ring of an arbitrary group $G$ . Then ${D}_{n}\left( G\right)  = G \cap  \left( {1 + {\Delta }^{n}\left( G\right) }\right)$ contains the $n$ th lower central subgroup ${\gamma }_{n}\left( G\right)$ of $G$ . Is it true that ${D}_{n}\left( G\right) /{\gamma }_{n}\left( G\right)$ is central in $G/{\gamma }_{n}\left( G\right)$ ?
	
	N. D. Gupta, Yu. V. Kuz'min
	
	No, not always; moreover, ${D}_{n}\left( G\right) /{\gamma }_{n}\left( G\right)$ need not be contained in any term of the upper central series of $G/{\gamma }_{n}\left( G\right)$ with fixed number (N. D. Gupta, Yu. V. Kuz’min, J. Pure Appl. Algebra, 104, no. 1 (1995), 191-197).
\end{openproblem}

\plainproblemsection

\begin{openproblem}
	*12.22. b) Let $\Delta \left( G\right)$ be the augmentation ideal of the integer group ring of an arbitrary group $G$ . Then ${D}_{n}\left( G\right)  = G \cap  \left( {1 + {\Delta }^{n}\left( G\right) }\right)$ contains the $n$ th lower central subgroup ${\gamma }_{n}\left( G\right)$ of $G$ . Is it true that ${D}_{n}\left( G\right) /{\gamma }_{n}\left( G\right)$ has exponent dividing 2 ?
	
	N. D. Gupta, Yu. V. Kuz’min
	
	*No, not always (L. Bartholdi, R. Mikhailov, J. Topol., 16, no. 2 (2023), 822-853).
\end{openproblem}

\plainproblemsection

\begin{openproblem}
	12.23. The index of permutability of a group $G$ is defined to be the minimal integer $k \geq  2$ such that for each $k$ -tuple ${x}_{1},\ldots ,{x}_{k}$ of elements in $G$ there is a non-identity permutation $\sigma$ on $k$ symbols such that ${x}_{1}\cdots {x}_{k} = {x}_{\sigma \left( 1\right) }\cdots {x}_{\sigma \left( k\right) }$ . Determine the index of permutability of the symmetric group ${\mathbb{S}}_{n}$ . M. Gutsan
\end{openproblem}

\plainproblemsection

\begin{openproblem}
	12.24. Given a ring $R$ with identity, the automorphisms of $R\left\lbrack  \left\lbrack  x\right\rbrack  \right\rbrack$ sending $x$ to $x\left( {1 + \mathop{\sum }\limits_{{i = 1}}^{\infty }{a}_{i}{x}^{i}}\right) ,{a}_{i} \in  R$ , form a group $N\left( R\right)$ . We know that $N\left( \mathbb{Z}\right)$ contains a copy of the free group ${F}_{2}$ of rank 2 and, from work of A. Weiss, that $N\left( {\mathbb{Z}/p\mathbb{Z}}\right)$ contains a copy of every finite $p$ -group (but not of ${\mathbb{Z}}_{{p}^{\infty }}$ ), $p$ a prime. Does $N\left( {\mathbb{Z}/p\mathbb{Z}}\right)$ contain a copy of ${F}_{2}$ ? D. L. Johnson
	
	Yes, it does (R. Camina, J. Algebra, 196, no. 1 (1997), 101-113). I. B. Fesenko noted that this fact might have been known to specialists in the theory of fields of norms back in 1985 (J.-P. Wintenberger, J.-M. Fontaine, F. Laubie); however a proof based on this theory first appeared only in (I. B. Fesenko, Preprint, Nottingham Univ., 1998).
\end{openproblem}

\plainproblemsection

\begin{openproblem}
	12.25. Let $G$ be a finite group acting irreducibly on a vector space $V$ . An orbit ${\alpha }^{G}$ for $\alpha  \in  V$ is said to be $p$ -regular if the stabilizer of $\alpha$ in $G$ is a ${p}^{\prime }$ -subgroup. Does $G$ have a regular orbit on $V$ if it has a $p$ -regular orbit for every prime $p$ ? Jiping Zhang No, not always. For example, let $H = {\mathbb{A}}_{4}\wr {\mathbb{Z}}_{5}, V = {O}_{2}\left( H\right)$ , and $G$ be a complement to $V$ in $H, G$ acting on $V$ by conjugation. (V.I. Zenkov, Letter of February,11,1994.)
\end{openproblem}

\plainproblemsection

\begin{openproblem}
	12.26. (Shi Shengming). Is it true that a finite $p$ -soluble group $G$ has a $p$ -block of defect zero if and only if there exists an element $x \in  {O}_{{p}^{\prime }}\left( G\right)$ such that ${C}_{G}\left( x\right)$ is a ${p}^{\prime }$ -subgroup? Jiping Zhang
	
	No, it is not. The group ${3}^{2}.G{L}_{2}\left( 3\right)$ has a 2-block of defect 0, but the centralizer of every element from $O\left( G\right)$ has even order (V.I. Zenkov, in: Trudy Inst. Matem. i Mekh. UrO RAN, 3 (1995), 36-40 (Russian)).
\end{openproblem}

\plainproblemsection

\begin{openproblem}
	12.27. Let $G$ be a simple locally finite group. We say that $G$ is a group of finite type if there is a non-trivial permutational representation of $G$ such that some finite subgroup of $G$ has no regular orbits. Investigate and, perhaps, classify the simple locally finite groups of finite type. Partial results see in (B. Hartley, A. Zalesski, Isr. J. Math., 82 (1993), 299-327, J. London Math. Soc., 55 (1997), 210-230; F. Leinen, O. Puglisi, Illinois J. Math., 47 (2003), 345-360). A. E. Zalesskiǐ
\end{openproblem}

\plainproblemsection

\begin{openproblem}
	12.28. Let $G$ be a group. A function $f : G \rightarrow  \mathbb{C}$ is called
	
	1) normed if $f\left( 1\right)  = 1$ ;
	
	2) central if $f\left( {gh}\right)  = f\left( {hg}\right)$ for all $g, h \in  G$ ;
	
	3) positive-definite if $\mathop{\sum }\limits_{{l \in  J}}f\left( {{g}_{k}^{-1}{g}_{l}}\right) {\bar{c}}_{k}{c}_{l} \geq  0$ for any ${g}_{1},\ldots ,{g}_{n} \in  G$ and any ${c}_{1},\ldots ,{c}_{n} \in  \mathbb{C}$ .
	
	Classify the infinite simple locally finite groups $G$ which possess functions satisfying 1)-3). The simple Chevalley groups are known to have no such functions, while such functions exist on locally matrix (or stable) classical groups over finite fields. The question is motivated by the theory of ${C}^{ * }$ -algebras, see $§9$ in (A.M. Vershik, S. V. Kerov, J. Sov. Math., 38 (1987), 1701-1733). Partial results see in (F. Leinen, O. Puglisi, J. Pure Appl. Algebra, 208 (2007), 1003-1021, J. London Math. Soc., 70 (2004), 678-690). A. E. Zalesskiǐ 12.29. Classify the locally finite groups for which the augmentation ideal of the complex group algebra is a simple ring. The problem goes back to I. Kaplansky (1965). For partial results, see (K. Bonvallet, B. Hartley, D. S. Passman, M. K. Smith, Proc. Amer. Math. Soc., 56, no. 1 (1976), 79-82; A. E. Zalesskiĭ, Algebra i Analiz, 2, no. 6 (1990), 132-149 (Russian); Ch. Praeger, A. E. Zalesskii, Proc. London Math. Soc., 70, no. 2 (1995), 313-335; B. Hartley, A. Zalesski, J. London Math. Soc., 55 (1997), 210-230). A. E. Zalesskiǐ 12.30. (O. N. Golovin). On the class of all groups, do there exist associative operations which satisfy the postulates of MacLane and Mal'cev (that is, which are free functorial and hereditary) and which are different from taking free and direct products? S. V. Ivanov
\end{openproblem}

\plainproblemsection

\begin{openproblem}
	12.31. For relatively free groups $G$ , prove or disprove the following conjecture of P. Hall: if a word $v$ takes only finitely many values on $G$ then the verbal subgroup ${vG}$ is finite. S. V. Ivanov
	
	The conjecture is disproved (A. Storozhev, Proc. Amer. Math. Soc., 124, no. 10 (1996), 2953-2954).
\end{openproblem}

\plainproblemsection

\begin{openproblem}
	*12.32. Prove an analogue of Higman’s theorem for the Burnside variety ${\mathfrak{B}}_{n}$ of groups of odd exponent $n \gg  1$ , that is, prove that every recursively presented group of exponent $n$ can be embedded in a finitely presented (in ${\mathfrak{B}}_{n}$ ) group of exponent $n$ . S. V. Ivanov
	
	*It is proved (A. Yu. Olshanskii, J. Algebra, 560 (2020), 960-1052).
\end{openproblem}

\plainproblemsection

\begin{openproblem}
	12.33. Suppose that $G$ is a finite group and $x$ is an element of $G$ such that the subgroup $\langle x, y\rangle$ has odd order for any $y$ conjugate to $x$ in $G$ . Prove, without using CFSG, that the normal closure of $x$ in $G$ is a group of odd order. L. S. Kazarin
\end{openproblem}

\plainproblemsection

\begin{openproblem}
	12.34. Describe the finite groups $G$ such that the sum of the cubes of the degrees of all irreducible complex characters is at most $\left| G\right|  \cdot  {\log }_{2}\left| G\right|$ . The question is interesting for applications in the theory of signal processing. L. S. Kazarin
\end{openproblem}

\plainproblemsection

\begin{openproblem}
	12.35. Suppose that $\mathfrak{F}$ is a radical composition formation of finite groups. Prove that $\langle H, K{\rangle }^{\mathfrak{F}} = \left\langle  {{H}^{\mathfrak{F}},{K}^{\mathfrak{F}}}\right\rangle$ for every finite group $G$ and any subnormal subgroups $H$ and $K$ of $G$ . S. F. Kamornikov
\end{openproblem}

\plainproblemsection

\begin{openproblem}
	12.36. Let $p$ be a prime, $V$ an $n$ -dimensional vector space over the field of $p$ elements, and let $G$ be a subgroup of ${GL}\left( V\right)$ . Let $S = S\left\lbrack  {V}^{ * }\right\rbrack$ be the symmetric algebra on ${V}^{ * }$ , the dual of $V$ . Let $T = {S}^{G}$ be the ring of invariants and let ${b}_{m}$ be the dimension of the homogeneous component of degree $m$ . Then the Poincaré series $\mathop{\sum }\limits_{{m \geq  0}}{b}_{m}{t}^{m}$ is a rational function with a Laurent power series expansion $\mathop{\sum }\limits_{{i \geq   - n}}{a}_{i}{\left( 1 - t\right) }^{i}$ about $t = 1$ where ${a}_{-n} = \frac{1}{\left| G\right| }$ .
	
	Conjecture: ${a}_{-n + 1} = \frac{r}{2\left| G\right| }$ where $r = \mathop{\sum }\limits_{W}\left( {\left( {p - 1}\right) {\alpha }_{W} + {s}_{W} - 1}\right)$ , the sum is taken over all maximal subspaces $W$ of $V$ , and ${\alpha }_{W},{s}_{W}$ are defined by $\left| {G}_{W}\right|  = {p}^{{\alpha }_{W}} \cdot  {s}_{W}$ where $p \nmid  {s}_{W}$ and ${G}_{W}$ denotes the pointwise stabilizer of $W$ .
	
	D. Carlisle, P. H. Kropholler
	
	The conjecture is proved (D. J. Benson, W. W. Crawley-Boevey, Bull. London Math. Soc., 27, no. 5 (1995), 435-440); another proof based on the Grothendieck-Riemann-Roch Theorem was later obtained in (A. Neeman, Comment. Math. Helv., 70, no. 3 (1995), 339-349).
\end{openproblem}

\plainproblemsection

\begin{openproblem}
	*12.37. (J. G. Thompson). For a finite group $G$ and natural number $n$ , set $G\left( n\right)  =$ $\left\{  {x \in  G \mid  {x}^{n} = 1}\right\}$ and define the type of $G$ to be the function whose value at $n$ is the order of $G\left( n\right)$ . Is it true that a group is soluble if its type is the same as that of a soluble one? A. S. Kondratiev, W. J. Shi
	
	*No, not always (P. Piwek, https://arxiv.org/abs/2403.02197).
\end{openproblem}

\plainproblemsection

\begin{openproblem}
	*12.38. (J. G. Thompson). For a finite group $G$ , we denote by $N\left( G\right)$ the set of all orders of the conjugacy classes of $G$ . Is it true that if $G$ is a finite non-abelian simple group, $H$ a finite group with trivial centre and $N\left( G\right)  = N\left( H\right)$ , then $G$ and $H$ are isomorphic? A. S. Kondratiev, W. J. Shi
	
	*Yes, it is true. The final step of the proof is in the paper (I. B. Gorshkov, Com-mun. Algebra, 47, no. 12 (2019), 5192-5206), which contains references to the previous steps by M. Ahanjideh, N. Ahanjideh, S. H. Alavi, G. Y. Chen, A. Daneshkhah, M. R. Darafsheh, I. B. Gorshkov, A. Iranmanesh, I. Kaygorodov, Behn. Khosravi, Behr. Khosravi, A. Kukharev, W. Shi, A. Shlepkin, A. V. Vasil’ev, L. Wang, M. Xu.
\end{openproblem}

\plainproblemsection

\begin{openproblem}
	12.39. (W. J. Shi). Must a finite group and a finite simple group be isomorphic if they have equal orders and the same set of orders of elements? A. S. Kondratiev Yes, they must (M. C. Xu, W. J. Shi, Algebra Colloq., 10 (2003), 427-443; A.V. Vasil'ev, M.A. Grechkoseeva, V.D. Mazurov, Algebra Logic, 48 (2009), 385-409).
\end{openproblem}

\plainproblemsection

\begin{openproblem}
	12.40. Let $\varphi$ be an irreducible $p$ -modular character of a finite group $G$ . Find the best-possible estimate of the form $\varphi {\left( 1\right) }_{p} \leq  f\left( {\left| G\right| }_{p}\right)$ . Here ${n}_{p}$ is the $p$ -part of a positive integer $n$ . A. S. Kondratiev
\end{openproblem}

\plainproblemsection

\begin{openproblem}
	*12.41. Let $F$ be a free group on two generators $x, y$ and let $\varphi$ be the automorphism of $F$ defined by $x \rightarrow  y, y \rightarrow  {xy}$ . Let $G$ be a semidirect product of $F/\left( {{F}^{\prime \prime }{\left( {F}^{\prime }\right) }^{2}}\right)$ and $\langle \varphi \rangle$ . Then $G$ is just-non-polycyclic. What is the cohomological dimension of $G$ over $\mathbb{Q}$ ? (It is either 3 or 4.) P. H. Kropholler
	
	*It is 4, because $G$ is not constructible, and then its cohomological dimension is equal to the homological dimension plus 1 by Theorem III. 8 in (M. R. Bridson, P. H. Kropholler, J. Reine Angew. Math., 699 (2015), 217-243), while the homological dimension is equal to the Hirsch length by (U. Stammbach, J. London Math. Soc. (2), 2 (1970), 567-570).
\end{openproblem}

\plainproblemsection

\begin{openproblem}
	12.42. Describe the automorphisms of the Sylow $p$ -subgroup of a Chevalley group of normal type over $\mathbb{Z}/{p}^{m}\mathbb{Z}, m \geq  2$ , where $p$ is a prime. Described in (S. G. Kolesnikov, Algebra and Logic, 43, no. 1 (2004), 17-33); Izv. Gomel' Univ., 36, no. 3 (2006), 137-146 (Russian); J. Math. Sci., New York, 152 (2008), 220-246).
\end{openproblem}

\plainproblemsection

\begin{openproblem}
	12.43. (Well-known problem). Does there exist an infinite finitely-generated residually finite $p$ -group such that each subgroup is either finite or of finite index?
	
	J. C. Lennox
\end{openproblem}

\plainproblemsection

\begin{openproblem}
	12.44. (P. Hall). Is there a non-trivial group which is isomorphic with every proper extension of itself by itself? J. C. Lennox
	
	Yes, such groups exist of cardinality any regular cardinal (R. Göbel, S. Shelah, Math. Proc. Cambridge Philos. Society (2), 134, no. 1 (2003), 23-31).
\end{openproblem}

\plainproblemsection

\begin{openproblem}
	12.45. (P. Hall). Must a non-trivial group, which is isomorphic to each of its nontrivial normal subgroups, be either free of infinite rank, simple, or infinite cyclic? (Lennox, Smith and Wiegold, 1992, have shown that a finitely generated group of this kind which has a proper normal subgroup of finite index is infinite cyclic.) J. C. Lennox No (V. N. Obraztsov, Proc. London Math. Soc., 75 (1997), 79-98).
\end{openproblem}

\plainproblemsection

\begin{openproblem}
	12.46. Let $F$ be the nonabelian free group on two generators $x, y$ . For $a, b \in  \mathbb{C},\left| a\right|  =$ $\left| b\right|  = 1$ , let ${\vartheta }_{a, b}$ be the automorphism of $\mathbb{C}F$ defined by ${\vartheta }_{a, b}\left( x\right)  = {ax},{\vartheta }_{a, b}\left( y\right)  = {by}$ . Given $0 \neq  \alpha  \in  \mathbb{C}F$ , can we always find $a, b \in  \mathbb{C} \smallsetminus  \{ 1\}$ with $\left| a\right|  = \left| b\right|  = 1$ such that $\alpha \mathbb{C}F \cap  {\vartheta }_{a, b}\left( \alpha \right) \mathbb{C}F \neq  0?$ P. A. Linnell
	
	No, not always (A. V. Tushev, Ukrain. Mat. J., 47, no. 4 (1995), 571-572).
\end{openproblem}

\plainproblemsection

\begin{openproblem}
	12.47. Let $k$ be a field, let $p$ be a prime, and let $G$ be the Wreath product ${\mathbb{Z}}_{p}\wr \mathbb{Z}$ (so the base group has exponent $p$ ). Does ${kG}$ have a classical quotient ring? (i. e. do the non-zero-divisors of ${kG}$ form an Ore set?)
	
	No, it does not (P. A. Linnell, W. Lück, T. Schick, in: High-dimensional manifold topology. Proc. of the school, ICTP (Trieste, 2001), World Scientific, River Edge, NJ, 2003, 315-321).
\end{openproblem}

\plainproblemsection

\begin{openproblem}
	12.48. Let $G$ be a sharply doubly transitive permutation group on a set $\Omega$ (see Archive 11.52 for a definition).
	
	(a) Does $G$ possess a regular normal subgroup if a point stabilizer is locally finite?
	
	(b) Does $G$ possess a regular normal subgroup if a point stabilizer has an abelian subgroup of finite index?
	
	Comments of 2022: an affirmative answer to part (b) is obtained in permutation characteristic 0 (F. O. Wagner, Preprint, 2022, https://hal.archives-ouvertes.fr/hal-03590818).V.D.Mazurov
\end{openproblem}

\plainproblemsection

\begin{openproblem}
	12.49. Construct all non-split extensions of elementary abelian 2-groups $V$ by $H =$ ${PS}{L}_{2}\left( q\right)$ for which $H$ acts irreducibly on $V$ . V. D. Mazurov
	
	They are constructed (V.P. Burichenko, Algebra and Logic, 39 (2000), 160-183).
\end{openproblem}

\plainproblemsection

\begin{openproblem}
	12.50. (Well-known problem). Find an algorithm which decides, by a given finite set of matrices in $S{L}_{3}\left( \mathbb{Z}\right)$ , whether the first matrix of this set is contained in the subgroup generated by the remaining matrices. An analogous problem for $S{L}_{4}\left( \mathbb{Z}\right)$ is insoluble since a direct product of two free groups of rank 2 embeds into $S{L}_{4}\left( \mathbb{Z}\right)$ .
\end{openproblem}

\plainproblemsection

\begin{openproblem}
	*12.51. Does the free group ${F}_{\eta }\left( {1 < \eta  < \infty }\right)$ have a finite subset $S$ for which there is a unique total order of ${F}_{\eta }$ making all elements of $S$ positive? Equivalently, does the free representable $l$ -group of rank $\eta$ have a basic element? (The analogs for right orders of ${F}_{\eta }$ and free $l$ -groups have negative answers.) S. McCleary
	
	*No, it does not (S. Dowhyi, K. Muliarchyk, Groups Geom. Dyn., 17 (2023), 613- 632; K. Muliarchyk, Preprint, 2024, https://arxiv.org/abs/2403.14779).
\end{openproblem}

\plainproblemsection

\begin{openproblem}
	12.52. Is the free $l$ -group ${\mathcal{F}}_{\eta }$ of rank $\eta \left( {1 < \eta  < \infty }\right)$ Hopfian? That is, are $l$ -homomorphisms from ${\mathcal{F}}_{\eta }$ onto itself necessarily one-to-one? S. McCleary
\end{openproblem}

\plainproblemsection

\begin{openproblem}
	12.53. Is it decidable whether or not two elements of a free $l$ -group are conjugate? S. McCleary
\end{openproblem}

\plainproblemsection

\begin{openproblem}
	12.54. Is there a normal valued $l$ -group $G$ for which there is no abelian $l$ -group $A$ with $C\left( A\right)  \cong  C\left( G\right)$ ? (Here $C\left( G\right)$ denotes the lattice of convex $l$ -subgroups of $G$ . If $G$ is not required to be normal valued, this question has an affirmative answer.)
	
	S. McCleary
\end{openproblem}

\plainproblemsection

\begin{openproblem}
	12.55. Let $f\left( {n, p}\right)$ be the number of groups of order ${p}^{n}$ . Is $f\left( {n, p}\right)$ an increasing function of $p$ for any fixed $n \geq  5$ ? A. Mann
\end{openproblem}

\plainproblemsection

\begin{openproblem}
	12.56. Let $\langle X \mid  R\rangle$ be a finite presentation (the words in $R$ are assumed cyclically reduced). Define the length of the presentation to be the sum of the number of generators and the lengths of the relators. Let $F\left( n\right)$ be the number of (isomorphism types of) groups that have a presentation of length at most $n$ . What can one say about the function $F\left( n\right)$ ? It can be shown that $F\left( n\right)$ is not recursive (D. Segal) and it is at least exponential (L. Pyber). A. Mann
\end{openproblem}

\plainproblemsection

\begin{openproblem}
	12.57. Suppose that a cyclic group $A$ of order 4 is a ${TI}$ -subgroup of a finite group $G$ . a) Does $A$ centralize a component $L$ of $G$ if $A$ intersects $L \cdot  C\left( L\right)$ trivially?
	
	b) What is the structure of the normal closure of $A$ in the case where $A$ centralizes each component of $G$ ? A. A. Makhnëv
\end{openproblem}

\plainproblemsection

\begin{openproblem}
	12.58. A generalized quadrangle ${GQ}\left( {s, t}\right)$ with parameters $s, t$ is by definition an incidence system consisting of points and lines in which every line consists of $s + 1$ points, every two different lines have at most one common point, each point belongs to $t + 1$ lines, and for any point $a$ not lying on a line $L$ there is a unique line containing $a$ and intersecting $L$ .
	
	a) Does ${GQ}\left( {4,{11}}\right)$ exist?
	
	b) Can the automorphism group of a hypothetical ${GQ}\left( {5,7}\right)$ contain involutions?
	
	A. A. Makhnëv
\end{openproblem}

\plainproblemsection

\begin{openproblem}
	12.59. Does there exist a strongly regular graph with parameters(85,14,3,2)and a non-connected neighborhood of some vertex? A. A. Makhnëv
\end{openproblem}

\plainproblemsection

\begin{openproblem}
	12.60. Describe the strongly regular graphs in which the neighborhoods of vertices are generalized quadrangles (see 12.58). A. A. Makhnëv
\end{openproblem}

\plainproblemsection

\begin{openproblem}
	12.62. A Frobenius group is a transitive permutation group in which the stabilizer of any two points is trivial. Does there exist a Frobenius group of infinite degree which is primitive as permutation group, and in which the stabilizer of a point is cyclic and has only finitely many orbits?
	
	Conjecture: No. Note that such a group cannot be 2-transitive (Gy. Károlyi, S. J. Kovács, P. P. Pálfy, Aequationes Math., 39 (1990), 161-166; V. D. Mazurov, Siberian Math. J., 31 (1990), 615-617; Peter M. Neumann, P. J. Rowley, in: Geometry and cohomology in group theory (London Math. Soc. Lect. Note Ser., 252), Cambridge Univ. Press, 1998, 291-295). Peter M. Neumann
\end{openproblem}

\plainproblemsection

\begin{openproblem}
	12.63. Does there exist a soluble permutation group of infinite degree that has only finitely many orbits on triples? Conjecture: No.
	
	Comment of 2001: It is proved (D. Macpherson, in: Advances in algebra and model theory (Algebra Logic Appl., 9), Gordon \& Breach, Amsterdam, 1997, 87-92) that every infinite soluble permutation group has infinitely many orbits on quadruples.
	
	Peter M. Neumann
\end{openproblem}

\plainproblemsection

\begin{openproblem}
	12.64. Is it true that for a given number $k \geq  2$ and for any (prime) number $n$ , there exists a number $N = N\left( {k, n}\right)$ such that every finite group with generators $A = \left\{  {{a}_{1},\ldots ,{a}_{k}}\right\}$ has exponent $\leq  n$ if ${\left( {x}_{1}\cdots {x}_{N}\right) }^{n} = 1$ for any ${x}_{1},\ldots ,{x}_{N} \in  A \cup  \{ 1\}$ ?
	
	For given $k$ and $n$ , a negative answer implies, for example, the infiniteness of the free Burnside group $B\left( {k, n}\right)$ , and a positive answer, in the case of sufficiently large $n$ , gives, for example, an opportunity to find a hyperbolic group which is not residually finite (and in this group a hyperbolic subgroup of finite index which has no proper subgroups of finite index). A. Yu. Ol'shanskiĭ
\end{openproblem}

\plainproblemsection

\begin{openproblem}
	12.65. Let $\mathcal{P} = \left( {{P}_{0},{P}_{1},{P}_{2}}\right)$ be a parabolic system in a finite group $G$ , belonging to the ${C}_{3}$ Coxeter diagram $0 \rightarrow  0 \rightarrow  0 \rightarrow  0$ and let the Borel subgroup have
	
	index at least 3 in ${P}_{0}$ and ${P}_{1}$ . It is known that, if we furthermore assume that the chamber system of $\mathcal{P}$ is geometric and that the projective planes arising as $\{ 0,1\}$ -residues from $\mathcal{P}$ are desarguesian, then either ${G}^{\left( \infty \right) }$ is a Chevalley group of type ${C}_{3}$ or ${B}_{3}$ or $G = {A}_{7}$ . Can we obtain the same conclusion in general, without assuming the previous two hypothesis? A. Pasini
	
	Yes, we can (S. Yoshiara, J. Algebraic Combin., 5, no. 3 (1996), 251-284).
\end{openproblem}

\plainproblemsection

\begin{openproblem}
	12.66. Describe the finite translation planes whose collineation groups act doubly transitively on the affine points. N. D. Podufalov
\end{openproblem}

\plainproblemsection

\begin{openproblem}
	12.67. Describe the structure of the locally compact soluble groups with the maximum (minimum) condition for closed non-compact subgroups. V. M. Poletskikh
\end{openproblem}

\plainproblemsection

\begin{openproblem}
	12.68. Describe the locally compact abelian groups in which any two closed subgroups of finite rank generate a subgroup of finite rank. V. M. Poletskikh
\end{openproblem}

\plainproblemsection

\begin{openproblem}
	12.69. Let $F$ be a countably infinite field and $G$ a finite group of automorphisms of $F$ . The group ring $\mathbb{Z}G$ acts on $F$ in a natural way. Suppose that $S$ is a subfield of $F$ satisfying the following property: for any $x \in  F$ there is a non-zero element $f \in  \mathbb{Z}G$ such that ${x}^{f} \in  S$ . Is it true that either $F = S$ or the field extension $F/S$ is purely inseparable? One can show that for an uncountable field an analogous question has an affirmative answer. K. N. Ponomarëv
\end{openproblem}

\plainproblemsection

\begin{openproblem}
	12.70. Let $p$ be a prime number, $F$ a free pro- $p$ -group of finite rank, and $\Delta  \neq  1$ an automorphism of $F$ whose order is a power of $p$ . Is the rank of ${\operatorname{Fix}}_{F}\left( \Delta \right)  = \{ x \in  F \mid$ $\Delta \left( x\right)  = x\}$ finite? If the order of $\Delta$ is prime to $p$ , then ${\operatorname{Fix}}_{F}\left( \Delta \right)$ has infinite rank (W. Herfort, L. Ribes, Proc. Amer. Math. Soc., 108 (1990), 287-295). L. Ribes Yes, it is finite (C. Scheiderer, Proc. Amer. Math. Soc., 127, no. 3 (1999), 695-700). Moreover, in (W. N. Herfort, L. Ribes, P. A. Zalesskii, Forum Math., 11, no. 1 (1999), 49-61) it is proved that if $P$ is a finite $p$ -group of automorphisms of a free pro- $p$ group $F$ of arbitrary rank, then the group of fixed points of $P$ is a free factor of $F$ and the latter result has been extended in (P. A. Zalesskii, J. Reine Angew. Math., 572 (2004), 97-110) to the situation of a free pro- $p$ group of arbitrary rank.
\end{openproblem}

\plainproblemsection

\begin{openproblem}
	12.71. Let $d\left( G\right)$ denote the smallest cardinality of a generating set of the group $G$ . Let $A$ and $B$ be finite groups. Is there a finite group $G$ such that $A, B \leq  G, G = \langle A, B\rangle$ , and $d\left( G\right)  = d\left( A\right)  + d\left( B\right)$ ? The corresponding question has negative answer in the class of solvable groups (L. G. Kovács, H.-S. Sim, Indag. Math., 2, no. 2 (1991), 229-232). No, not always (A. Lucchini, J. Group Theory, 4, no. 1 (2001), 53-58). L. Ribes
\end{openproblem}

\plainproblemsection

\begin{openproblem}
	12.72. Let $\mathfrak{F}$ be a soluble local hereditary formation of finite groups. Prove that $\mathfrak{F}$ is radical if every finite soluble minimal non- $\mathfrak{F}$ -group $G$ is a minimal non- ${\mathfrak{N}}^{l\left( G\right)  - 1}$ -group. Here $\mathfrak{N}$ is the formation of all finite nilpotent groups, and $l\left( G\right)$ is the nilpotent length of $G$ . V. N. Semenchuk
\end{openproblem}

\plainproblemsection

\begin{openproblem}
	12.73. A formation $\mathfrak{H}$ of finite groups is said to have length $t$ if there is a chain of formations $\varnothing  = {\mathfrak{H}}_{0} \subset  {\mathfrak{H}}_{1} \subset  \cdots  \subset  {\mathfrak{H}}_{t} = \mathfrak{H}$ in which ${\mathfrak{H}}_{i - 1}$ is a maximal subformation of ${\mathfrak{H}}_{i}$ . Is the lattice of soluble formations of length $\leq  4$ distributive? A. N. Skiba
\end{openproblem}

\plainproblemsection

\begin{openproblem}
	12.74. Let $\mathfrak{F}$ be a non-primary one-generator composition formation of finite groups. Is it true that if $\mathfrak{F} = \mathfrak{{MH}}$ and the formations $\mathfrak{M}$ and $\mathfrak{H}$ are non-trivial, then $\mathfrak{M}$ is a composition formation? A. N. Skiba
	
	This is true if $\mathfrak{F} \neq  \mathfrak{H}$ (A. N. Skiba, Algebra of formations, Belaruskaja Navuka, Minsk, 1997, p. 144 (Russian)). In general the answer is negative (W. Guo, Commun. Algebra, 28, no. 10 (2000), 4767-4782).
\end{openproblem}

\plainproblemsection

\begin{openproblem}
	12.75. (B. Jonsson). Is the class $N$ of the lattices of normal subgroups of groups a variety? It is proved (C. Herrman, W. Poguntke, Algebra Univers., 4, no. 3 (1974), 280-286 that $N$ is an infinitely based quasivariety. D. M. Smirnov
\end{openproblem}

\plainproblemsection

\begin{openproblem}
	12.76. Is every group generated by a set of 3-transpositions locally finite? A set of 3-transpositions is, by definition, a normal set of involutions such that the orders of their pairwise products are at most 3 . A. I. Sozutov
	
	Yes, it is (H. Cuypers, J. I. Hall, in: Groups, combinatorics and geometry. Proc. L. M. S. Durham symp., July 5-15, 1990 (London Math. Soc. Lecture Note Ser., 165), Cambridge Univ. Press, 1992, 121-138).
\end{openproblem}

\plainproblemsection

\begin{openproblem}
	12.77. (Well-known problem). Does the order (if it is greater than ${p}^{2}$ ) of a finite non-cyclic $p$ -group divide the order of its automorphism group? A. I. Starostin No, not always (J. González-Sánchez, A. Jaikin-Zapirain, Forum Math. Sigma, 3, Article ID e7, 11 p., electronic only (2015)).
\end{openproblem}

\plainproblemsection

\begin{openproblem}
	12.78. (M.J.Curran). a) Does there exist a group of order ${p}^{6}$ (here $p$ is a prime number), whose automorphism group has also order ${p}^{6}$ ?
	
	b) Is it true that for $p \equiv  1\left( {\;\operatorname{mod}\;3}\right)$ , the smallest order of a $p$ -group that is the automorphism group of a $p$ -group is ${p}^{7}$ ?
	
	c) The same question for $p = 3$ with ${3}^{9}$ replacing ${p}^{7}$ . A. I. Starostin
	
	a) No, there is no such a group; b) yes, it is true; c) no, it is not true (S. Yu, G. Ban, J. Zhang, Algebra Colloq., 3, no. 2 (1996), 97-106).
\end{openproblem}

\plainproblemsection

\begin{openproblem}
	*12.79. Suppose that $a$ and $b$ are two elements of a finite group $G$ such that the
	
	function
	
	$$
	\varphi \left( g\right)  = {1}^{G}\left( g\right)  - {1}_{\langle a\rangle }^{G}\left( g\right)  - {1}_{\langle b\rangle }^{G}\left( g\right)  - {1}_{\langle {ab}\rangle }^{G}\left( g\right)  + 2
	$$
	
	is a character of $G$ . Is it true that $G = \langle a, b\rangle$ ? The converse statement is true. S. P. Strunkov
	
	*No, not always: a counterexample is given by $G = {A}_{4}, a = b =$ (123). (S. V. Skresanov, Algebra Logic, 58, no. 3 (2019), 249-253).
\end{openproblem}

\plainproblemsection

\begin{openproblem}
	12.80. (K. W. Roggenkamp). a) Is it true that the number of $p$ -blocks of defect 0 of a finite group $G$ is equal to the number of the conjugacy classes of elements $g \in  G$ such that the number of solutions of the equation $\left\lbrack  {x, y}\right\rbrack   = g$ in $G$ is not divisible by $p$ ? S. P. Strunkov
	
	No; for example, if $p = 2$ and $G = {\mathbb{Z}}_{3} \times  {\mathbb{S}}_{3}$ (L. Barker, Letter of May,27,1996.)
\end{openproblem}

\plainproblemsection

\begin{openproblem}
	*12.80. b) (K. W. Roggenkamp). Is it true that the number of $p$ -blocks of defect 0 of a finite simple group $G$ is equal to the number of the conjugacy classes of elements $g \in  G$ such that the number of solutions of the equation $\left\lbrack  {x, y}\right\rbrack   = g$ in $G$ is not divisible by $p$ ? S. P. Strunkov
	
	*No, not always. For example in the alternating group ${A}_{5}$ for $p = 2$ , the number of 2-blocks of defect 0 is 1, but there are 3 conjugacy classes of elements $g$ such that the number of solutions $\left\lbrack  {x, y}\right\rbrack   = g$ is odd, namley, $g = \left( {1,2,3}\right) ,\left( {1,2,3,4,5}\right)$ , and (1,3,5,2,4). (B. Sambale, Letter of 5 May 2021).
\end{openproblem}

\plainproblemsection

\begin{openproblem}
	12.81. What is the cardinality of the set of subvarieties of the group variety ${\mathfrak{A}}_{p}^{3}$ (where ${\mathfrak{A}}_{p}$ is the variety of abelian groups of prime exponent $p$ )? V. I. Sushchanski
\end{openproblem}

\plainproblemsection

\begin{openproblem}
	12.82. Find all pairs(n, r)such that the symmetric group ${\mathbb{S}}_{n}$ contains a maximal subgroup isomorphic to ${\mathbb{S}}_{r}$ . V. I. Sushchanskii
	
	They are found (B. Newton, B. Benesh, J. Algebra, 304, no. 2 (2006), 1108-1113).
\end{openproblem}

\plainproblemsection

\begin{openproblem}
	12.84. (Well-known problem). Is it true that if there exist two non-isomorphic groups with the given set of orders of the elements, then there are infinitely many groups with this set of orders of the elements? S. A. Syskin
	
	No, it is not (V. D. Mazurov, Algebra and Logic, 33, no. 1 (1994), 49-55).
\end{openproblem}

\plainproblemsection

\begin{openproblem}
	12.85. Does every variety which is generated by a (known) finite simple group containing a soluble subgroup of derived length $d$ contain a 2-generator soluble group of derived length $d$ ? S. A. Syskin
\end{openproblem}

\plainproblemsection

\begin{openproblem}
	12.86. For each known finite simple group, find its maximal 2-generator direct power.
	
	S. A. Syskin
\end{openproblem}

\plainproblemsection

\begin{openproblem}
	12.87. Let $\Gamma$ be a connected undirected graph without loops or multiple edges and suppose that the automorphism group $\operatorname{Aut}\left( \Gamma \right)$ acts transitively on the vertex set of $\Gamma$ . Is it true that at least one of the following assertions holds?
	
	1. The stabilizer of a vertex of $\Gamma$ in $\operatorname{Aut}\left( \Gamma \right)$ is finite.
	
	2. The group $\operatorname{Aut}\left( \Gamma \right)$ as a permutation group on the vertex set of $\Gamma$ admits an imprimitivity system $\sigma$ with finite blocks for which the stabilizer of a vertex of the factor-graph $\Gamma /\sigma$ in $\operatorname{Aut}\left( {\Gamma /\sigma }\right)$ is finite.
	
	3. There exists a natural number $n$ such that the graph obtained from $\Gamma$ by adding edges joining distinct vertices the distance between which in $\Gamma$ is at most $n$ contains a regular tree of valency 3 . V. I. Trofimov
\end{openproblem}

\plainproblemsection

\begin{openproblem}
	12.88. An undirected graph is called a locally finite Cayley graph of a group $G$ if its vertex set can be identified with the set of elements of $G$ in such a way that, for some finite generating set $X = {X}^{-1}$ of $G$ not containing 1, two vertices $g$ and $h$ are adjacent if and only if ${g}^{-1}h \in  X$ . Do there exist two finitely generated groups with the same locally finite Cayley graph, one of which is periodic and the other is not periodic?
\end{openproblem}

\plainproblemsection

\begin{openproblem}
	12.89. Describe the infinite connected graphs to which the sequences of finite connected graphs with primitive automorphism groups converge. For definitions, see (V. I. Trofimov, Algebra and Logic, 28, no. 3 (1989), 220-237). V. I. Trofimov
\end{openproblem}

\plainproblemsection

\begin{openproblem}
	12.90. Let $G$ be a finitely generated soluble minimax group and let $H$ be a finitely generated residually finite group which has precisely the same finite images as $G$ . Must $H$ be a minimax group? J. S. Wilson
	
	Yes, it must (A. Mann, D. Segal, Proc. London Math. Soc. (3), 61 (1990), 529-545; see also A. V. Tushev, Math. Notes, 56, no. 5 (1994), 1190-1192).
\end{openproblem}

\plainproblemsection

\begin{openproblem}
	12.91. Every metabelian group belonging to a Fitting class of (finite) supersoluble groups is nilpotent. Does the following generalization also hold: Every group belonging to a Fitting class of supersoluble groups has only central minimal normal subgroups? H. Heineken
	
	No, it does not (H. Heineken, Rend. Sem. Mat. Univ. Padova, 98 (1997), 241-251).
\end{openproblem}

\plainproblemsection

\begin{openproblem}
	12.92. For a field $K$ of characteristic 2, each finite group $G = \left\{  {{g}_{1},\ldots ,{g}_{n}}\right\}$ of odd order $n$ is determined up to isomorphism by its group determinant (Formanek-Sibley, 1990) and even by its reduced norm which is defined as the last coefficient ${s}_{m}\left( x\right)$ of the minimal polynomial $\varphi \left( {\lambda ;x}\right)  = {\lambda }^{m} - {s}_{1}\left( x\right) {\lambda }^{m - 1} + \cdots  + {\left( -1\right) }^{m}{s}_{m}\left( x\right)$ for the generic element $x = {x}_{1}{g}_{1} + \cdots  + {x}_{n}{g}_{n}$ of the group ring ${KG}$ (Hoehnke,1991). Is it possible in this theorem to replace ${s}_{m}\left( x\right)$ by some other coefficients ${s}_{i}\left( x\right) , i < m$ ? For notation see (G. Frobenius, Sitzungsber. Preuss. Akad. Wiss. Berlin, 1896, 1343-1382) and (K. W. Johnson, Math. Proc. Cambridge Phil. Soc., 109 (1991), 299-311).
	
	H.-J. Hoehnke
\end{openproblem}

\plainproblemsection

\begin{openproblem}
	12.93. Let $N \rightarrowtail  G \rightarrow  Q$ be an extension of nilpotent groups, with $Q$ finitely generated, which splits at every prime. Does the extension split? This is known to be true if $N$ is finite or commutative. P. Hilton
	
	No, not always (K. Lorensen, Math. Proc. Cambridge Philos. Soc., 123, no. 2 (1998), 213-215).
\end{openproblem}

\plainproblemsection

\begin{openproblem}
	*12.94. Let $G$ be a finitely generated pro- $p$ -group not involving the wreath product ${C}_{p}\wr {\mathbb{Z}}_{p}$ as a closed section (where ${C}_{p}$ is a cyclic group of order $p$ and ${\mathbb{Z}}_{p}$ is the group of $p$ -adic integers). Does it follow that $G$ is $p$ -adic analytic? A. Shalev
	
	*No, it does not (A. Jaikin-Zapirain, B. Klopsch, J. London Math. Soc. (2), 76, no. 2 (2007), 365-383).
\end{openproblem}

\plainproblemsection

\begin{openproblem}
	12.95. Let $G$ be a finitely generated pro- $p$ -group and let ${g}_{1},\ldots ,{g}_{n} \in  G$ . Let $H$ be an open subgroup of $G$ , and suppose there exists a non-trivial word $w = w\left( {{X}_{1},\ldots ,{X}_{n}}\right)$ such that $w\left( {{a}_{1},\ldots ,{a}_{n}}\right)  = 1$ whenever ${a}_{1} \in  {g}_{1}H,\ldots ,{a}_{n} \in  {g}_{n}H$ (that is, $G$ satisfies a coset identity). Does it follow that $G$ satisfies some non-trivial identity?
	
	A positive answer to this question would imply that an analogue of the Tits Alternative holds for finitely generated pro- $p$ -groups. Note that J. S. Wilson and E. I. Zelmanov (J. Pure Appl. Algebra, 81 (1992), 103-109) showed that the graded Lie algebra ${L}_{p}\left( G\right)$ with respect to the dimension subgroups of $G$ in characteristic $p$ satisfies a polynomial identity.
	
	Comment of 2017: This has been shown to be true for finitely generated linear groups (M. Larsen, A. Shalev, Algebra Number Theory, 10, no. 6 (2016), 1359-1371).
	
	A. Shalev
\end{openproblem}

\plainproblemsection

\begin{openproblem}
	12.96. Find a non-empty Fitting class $\mathfrak{F}$ and a non-soluble finite group $G$ such that $G$ has no $\mathfrak{F}$ -injectors. L. A. Shemetkov
	
	Found by E. Salomon (Mainz Univ., unpublished); his example is presented in §7.1 of (A. Ballester-Bolinches, L. M. Ezquerro, Classes of finite groups, Springer, 2006).
\end{openproblem}

\plainproblemsection

\begin{openproblem}
	12.97. Let $\mathfrak{F}$ be the formation of all finite groups all of whose composition factors are isomorphic to some fixed simple non-abelian group $T$ . Prove that $\mathfrak{F}$ is indecomposable into a product of two non-trivial subformations. L. A. Shemetkov
	
	It is proved (O. V. Mel'nikov, Problems in Algebra, 9, Gomel', 1996, 42-47 (Russian)).
\end{openproblem}

\plainproblemsection

\begin{openproblem}
	12.98. Let $F$ be a free group of finite rank, $R$ its recursively defined normal subgroup. Is it true that
	
	a) the word problem for $F/R$ is soluble if and only if it is soluble for $F/\left\lbrack  {R, R}\right\rbrack$ ?
	
	b) the conjugacy problem for $F/R$ is soluble if and only if it is soluble for $F/\left\lbrack  {R, R}\right\rbrack$ ? c) the conjugacy problem for $F/\left\lbrack  {R, R}\right\rbrack$ is soluble if the word problem is soluble for $F/\left\lbrack  {R, R}\right\rbrack$ ? V. E. Shpil'rain
	
	a) Yes, it is; b) No, it is not; c) No, it is not (M. I. Anokhin, Math. Notes, 61, no. 1-2 (1997), 3-8).
\end{openproblem}

\plainproblemsection

\begin{openproblem}
	12.100. Is every periodic group with a regular automorphism of order 4 locally finite?
	
	P. V. Shumyatsky
\end{openproblem}

\plainproblemsection

\begin{openproblem}
	12.101. We call a group $G$ containing an involution $i$ a ${T}_{0}$ -group if
	
	1) the order of the product of any two involutions conjugate to $i$ is finite;
	
	2) all 2-subgroups of $G$ are either cyclic or generalized quaternion;
	
	3) the centralizer $C$ of the involution $i$ in $G$ is infinite, distinct from $G$ , and has finite periodic part;
	
	4) the normalizer of any non-trivial $i$ -invariant finite subgroup in $G$ either is contained in $C$ or has periodic part which is a Frobenius group (see 6.55) with abelian kernel and finite complement of even order;
	
	5) for every element $c$ not contained in $C$ for which ${ci}$ is an involution there is an element $s$ of $C$ such that $\left\langle  {c,{c}^{s}}\right\rangle$ is an infinite subgroup.
	
	Does there exist a simple ${T}_{0}$ -group? V. P. Shunkov
\end{openproblem}

\plainproblemsection

\begin{openproblem}
	12.102. Is every proper factor-group of a group of Golod (see 9.76) residually finite? No, not every (L. Hammoudi, Algebra Colloq., 5 (1998), 371-376). V. P. Shunkov
\end{openproblem}

\plainproblemsection

\begin{openproblem}
	13.1. Let $U\left( K\right)$ denote the group of units of a ring $K$ . Let $G$ be a finite group, $\mathbb{Z}G$ the integral group ring of $G$ , and ${\mathbb{Z}}_{p}G$ the group ring of $G$ over the residues modulo a prime number $p$ . Describe the homomorphism from $U\left( {\mathbb{Z}G}\right)$ into $U\left( {{\mathbb{Z}}_{p}G}\right)$ induced by reducing the coefficients modulo $p$ . More precisely, find the kernel and the image of this homomorphism and an explicit transversal over the kernel. R. Zh. Aleev
\end{openproblem}

\plainproblemsection

\begin{openproblem}
	13.2. Does there exist a finitely based variety of groups $\mathfrak{V}$ such that the word problem is solvable in ${F}_{n}\left( \mathfrak{V}\right)$ for all natural $n$ , but is unsolvable in ${F}_{\infty }\left( \mathfrak{V}\right)$ (with respect to a system of free generators)? M. I. Anokhin
\end{openproblem}

\plainproblemsection

\begin{openproblem}
	13.3. Let $\mathfrak{M}$ be an arbitrary variety of groups. Is it true that every infinitely generated projective group in $\mathfrak{M}$ is an $\mathfrak{M}$ -free product of countably generated projective groups in M? V. A. Artamonov
\end{openproblem}

\plainproblemsection

\begin{openproblem}
	13.4. Let $G$ be a group with a normal (pro-) 2-subgroup $N$ such that $G/N$ is isomorphic to $G{L}_{n}\left( 2\right)$ , inducing its natural module on $N/\Phi \left( N\right)$ , the Frattini factor-group of $N$ . For $n = 3$ determine $G$ such that $N$ is as large as possible. (For $n > 3$ it can be proved that already the Frattini subgroup $\Phi \left( N\right)$ of $N$ is trivial; for $n = 2$ there exists $G$ such that $N$ is the free pro-2-group generated by two elements). B. Baumann
\end{openproblem}

\plainproblemsection

\begin{openproblem}
	13.5. Groups $A$ and $B$ are said to be locally equivalent if for every finitely generated subgroup $X \leq  A$ there is a subgroup $Y \leq  B$ isomorphic to $X$ , and conversely, for every finitely generated subgroup $Y \leq  B$ there is a subgroup $X \leq  A$ , isomorphic to $Y$ . We call a group $G$ categorical if $G$ is isomorphic to any group that is locally equivalent to $G$ . Is it true that a periodic locally soluble group $G$ is categorical if and only if $G$ is hyperfinite with Chernikov Sylow subgroups? (A group is hyperfinite if it has a well ordered ascending normal series with finite factors.) V. V. Belyaev
\end{openproblem}

\plainproblemsection

\begin{openproblem}
	13.6. (B. Hartley). Is it true that a locally finite group containing an element with Chernikov centralizer is almost soluble? V. V. Belyaev
\end{openproblem}

\plainproblemsection

\begin{openproblem}
	13.7. (B. Hartley). Is it true that a simple locally finite group containing a finite subgroup with finite centralizer is linear? V. V. Belyaev
\end{openproblem}

\plainproblemsection

\begin{openproblem}
	13.8. (B. Hartley). Is it true that a locally soluble periodic group has a finite normal series with locally nilpotent factors if it contains an element
	
	*a) with finite centralizer?
	
	b) with Chernikov centralizer? V. V. Belyaev
	\par\smallskip\hrule\smallskip
	*Yes, it is true (E. Khukhro, Preprint, 2025, http://arxiv.org/abs/2505.20999).
\end{openproblem}

\plainproblemsection

\begin{openproblem}
	13.9. Is it true that a locally soluble periodic group containing a finite nilpotent subgroup with Chernikov centralizer has a finite normal series with locally nilpotent factors? V. V. Belyaev, B. Hartley
	\par\smallskip\hrule\smallskip
\end{openproblem}

\plainproblemsection

\begin{openproblem}
	13.10. Is there a function $f : \mathbb{N} \rightarrow  \mathbb{N}$ such that, for every soluble group $G$ of derived length $k$ generated by a set $A$ , the validity of the identity ${x}^{4} = 1$ on each subgroup generated by at most $f\left( k\right)$ elements of $A$ implies that $G$ is a group of exponent 4 ? V. V. Bludov
	
	No, there is not (G. S. Deryabina, A. N. Krasil'nikov, Siberian Math. J., 44, no. 1 (2003), 58-60).
\end{openproblem}

\plainproblemsection

\begin{openproblem}
	13.11. Is a torsion-free group almost polycyclic if it has a finite set of generators ${a}_{1},\ldots ,{a}_{n}$ such that every element of the group has a unique presentation in the form ${a}_{1}^{{k}_{1}}\cdots {a}_{n}^{{k}_{n}}$ , where ${k}_{1},\ldots ,{k}_{n} \in  \mathbb{Z}$ ? V. V. Bludov
	
	No, not always (A. Muranov, Trans. Amer. Math. Soc., 359 (2007), 3609-3645).
\end{openproblem}

\plainproblemsection

\begin{openproblem}
	13.12. Is the group of all automorphisms of an arbitrary hyperbolic group finitely presented?
	
	Comment of 2001: This is proved for the group of automorphisms of a torsion-free hyperbolic group that cannot be decomposed into a free product (Z. Sela, Geom. Funct. Anal., 7, no. 3 (1997) 561-593). O. V. Bogopol'skiǐ
\end{openproblem}

\plainproblemsection

\begin{openproblem}
	13.13. Let $G$ and $H$ be finite $p$ -groups with isomorphic Burnside rings. Is the nilpo-tency class of $H$ bounded by some function of the class of $G$ ? There is an example where $G$ is of class 2 and $H$ is of class 3 . R. Brandl
\end{openproblem}

\plainproblemsection

\begin{openproblem}
	13.14. Is the lattice of quasivarieties of nilpotent torsion-free groups of nilpotency class $\leq  2$ distributive? A. I. Budkin
\end{openproblem}

\plainproblemsection

\begin{openproblem}
	13.15. Does a free non-abelian nilpotent group of class 3 possess an independent basis of quasiidentities in the class of torsion-free groups? A. I. Budkin
\end{openproblem}

\plainproblemsection

\begin{openproblem}
	13.16. Is every locally nilpotent group with minimum condition on centralizers hyper-central? F. O Wagner
	
	Yes, it is (V. V. Bludov, Algebra and Logic, 37, no. 3 (1998), 151-156).
\end{openproblem}

\plainproblemsection

\begin{openproblem}
	13.17. A representation of a group $G$ on a vector space $V$ is called a nil-representation if for every $v$ in $V$ and $g$ in $G$ there exists $n = n\left( {v, g}\right)$ such that $v{\left( g - 1\right) }^{n} = 0$ . Is it true that in zero characteristic every irreducible nil-representation is trivial? (In prime characteristic it is not true.) S. Vovsi
\end{openproblem}

\plainproblemsection

\begin{openproblem}
	13.18. Let $F$ be a finitely generated non-Abelian free group and let $G$ be the Cartesian (unrestricted) product of countable infinity of copies of $F$ . Must the Abelianiza-tion $G/{G}^{\prime }$ of $G$ be torsion-free? A. M. Gaglione, D. Spellman
	
	No, it may contain elements of order 2 (O. Kharlampovich, A. Myasnikov, in: Knots, braids, and mapping class groups - papers dedicated to Joan S. Birman (New York, 1998), Amer. Math. Soc., Providence, RI, 2001, 77-83).
\end{openproblem}

\plainproblemsection

\begin{openproblem}
	13.19. Suppose that both a group $Q$ and its normal subgroup $H$ are subgroups of the direct product ${G}_{1} \times  \cdots  \times  {G}_{n}$ such that for each $i$ the projections of both $Q$ and $H$ onto ${G}_{i}$ coincide with ${G}_{i}$ . If $Q/H$ is a $p$ -group, is $Q/H$ a regular $p$ -group?
	
	Yu. M. Gorchakov
\end{openproblem}

\plainproblemsection

\begin{openproblem}
	13.20. Is it true that the generating series of the growth function of every finitely generated group with one defining relation represents an algebraic (or even a rational) function? R. I. Grigorchuk
\end{openproblem}

\plainproblemsection

\begin{openproblem}
	13.21. a) Is there an infinite finitely generated residually finite $p$ -group, in which the order $\left| g\right|$ of an arbitrary element $g$ does not exceed $f\left( {\delta \left( g\right) }\right)$ , where $\delta \left( g\right)$ is the length of $g$ with respect to a fixed set of generators and $f\left( n\right)$ is a function growing at $n \rightarrow  \infty$ slower than any power function ${n}^{\lambda },\lambda  > 0$ ? R. I. Grigorchuk
	
	Yes, there is (A. V. Rozhkov, Dokl. Math., 58, no. 2 (1998), 234-237).
\end{openproblem}

\plainproblemsection

\begin{openproblem}
	13.21. b) What is the minimal possible rate of growth of the function $\pi \left( n\right)  = \mathop{\max }\limits_{{\delta \left( g\right)  \leq  n}}\left| g\right|$ for the class of groups indicated in part "a" of this problem (see Archive)? It is known (R. I. Grigorchuk, Math. USSR-Izv., $\mathbf{{25}}\left( {1985}\right) ,{259} - {300}$ ) that there exist $p$ -groups with $\pi \left( n\right)  \leq  {n}^{\lambda }$ for some $\lambda  > 0$ . At the same time, it follows from a result of E. I. Zel’manov that $\pi \left( n\right)$ is not bounded if $G$ is infinite. R. I. Grigorchuk
\end{openproblem}

\plainproblemsection

\begin{openproblem}
	13.22. Let the group $G = {AB}$ be the product of two polycyclic subgroups $A$ and $B$ , and assume that $G$ has an ascending normal series with locally nilpotent factors (i. e. $G$ is radical). Is it true that $G$ is polycyclic? F. de Giovanni
\end{openproblem}

\plainproblemsection

\begin{openproblem}
	13.23. Let the group $G$ have a finite normal series with infinite cyclic factors (containing of course $G$ and $\{ 1\} )$ . Is it true that $G$ has a non-trivial outer automorphism?
	
	F. de Giovanni
\end{openproblem}

\plainproblemsection

\begin{openproblem}
	13.24. Let $G$ be a non-discrete topological group with only finitely many ultrafilters that converge to the identity. Is it true that $G$ contains a countable open subgroup of exponent 2?
	
	Comment of 2005: This is proved for $G$ countable (E. G. Zelenyuk, Mat. Studii, 7 (1997), 139-144).
\end{openproblem}

\plainproblemsection

\begin{openproblem}
	13.25. Let $\left( {G,\tau }\right)$ be a topological group with finite semigroup $\tau \left( G\right)$ of ultrafilters converging to the identity (I. V. Protasov, Siberian Math. J., 34, no. 5 (1993), 938- 951). Is it true that $\tau \left( G\right)$ is a semigroup of idempotents?
	
	Comment of 2005: This is proved for $G$ countable (E. G. Zelenyuk, Mat. Studii, 14 (2000), 121-140). E. G. Zelenyuk
\end{openproblem}

\plainproblemsection

\begin{openproblem}
	13.26. Is it true that a countable topological group of exponent 2 with unique free ultrafilter converging to the identity has a basis of neighborhoods of the identity consisting of subgroups? E. G. Zelenyuk, I. V. Protasov
	
	Assuming Martin's Axiom, the answer is "No" (Ye. Zelenyuk, Adv. Math., 229, no. 4 (2012), 2415-2426).
\end{openproblem}

\plainproblemsection

\begin{openproblem}
	13.27. (B. Amberg). Suppose that $G = {AB} = {AC} = {BC}$ for a group $G$ and its subgroups $A, B, C$ . a) Is $G$ a Chernikov group if $A, B, C$ are Chernikov groups?
	
	b) Is $G$ almost polycyclic if $A, B, C$ are almost polycyclic? L. S. Kazarin
\end{openproblem}

\plainproblemsection

\begin{openproblem}
	13.28. (D. M. Evans). A permutation group on an infinite set is cofinitary if its non-identity elements fix only finitely many points. Is it true that a closed cofinitary permutation group is locally compact (in the topology of pointwise convergence)?
	
	No, not always (G. Hjorth, J. Algebra, 200, no. 2 (1998), 439-448). P. J. Cameron
\end{openproblem}

\plainproblemsection

\begin{openproblem}
	13.29. Given an infinite set $\Omega$ , define an algebra $A$ (the reduced incidence algebra of finite subsets) as follows. Let ${V}_{n}$ be the set of functions from the set of $n$ -element subsets of $\Omega$ to the rationals $\mathbb{Q}$ . Now let $A = \bigoplus {V}_{n}$ , with multiplication as follows: for $f \in  {V}_{n}, g \in  {V}_{m}$ , and $\left| X\right|  = m + n$ , let $\left( {fg}\right) \left( X\right)  = \sum f\left( Y\right) g\left( {X \smallsetminus  Y}\right)$ , where the sum is over the $n$ -element subsets $Y$ of $X$ . If $G$ is a permutation group on $\Omega$ , let ${A}^{G}$ be the algebra of $G$ -invariants in $A$ .
	
	Conjecture: If $G$ has no finite orbits on $\Omega$ , then ${A}^{G}$ is an integral domain. P. J. Cameron
	
	Conjecture is proved (M. Pouzet, Theor. Inf. App., 42, no. 1 (2008), 83-103).
\end{openproblem}

\plainproblemsection

\begin{openproblem}
	13.30. A group $G$ is called a $B$ -group if every primitive permutation group which contains the regular representation of $G$ is doubly transitive. Are there any countable
	
	$B$ -groups? P. J. Cameron
\end{openproblem}

\plainproblemsection

\begin{openproblem}
	13.31. Let $G$ be a permutation group on a set $\Omega$ . A sequence of points of $\Omega$ is a base for $G$ if its pointwise stabilizer in $G$ is the identity. The greedy algorithm for a base chooses each point in the sequence from an orbit of maximum size of the stabilizer of its predecessors. Is it true that there is a universal constant $c$ with the property that, for any finite primitive permutation group, the greedy algorithm produces a base whose size is at most $c$ times the minimal base size? P. J. Cameron
\end{openproblem}

\plainproblemsection

\begin{openproblem}
	13.32. (Well-known problem). On a group, a partial order that is directed upwards and has linearly ordered cone of positive elements is said to be semilinear if this order is invariant under right multiplication by the group elements. Can every semilinear order of a group be extended to a right order of the group? V. M. Kopytov
\end{openproblem}

\plainproblemsection

\begin{openproblem}
	13.33. (F. Gross). Is a normal subgroup of a finite ${D}_{\pi }$ -group (see Archive,3.62) always a ${D}_{\pi }$ -group? V. D. Mazurov
	
	Yes, it is mod CFSG (E.P. Vdovin, D. O. Revin, in: Ischia Group Theory 2004 (Contemp. Math., 402), Amer. Math. Soc., 2006, 229-263).
\end{openproblem}

\plainproblemsection

\begin{openproblem}
	13.34. (I. D. Macdonald). If the identity ${\left\lbrack  x, y\right\rbrack  }^{n} = 1$ holds on a group, is the derived subgroup of the group periodic? V. D. Mazurov
	
	No, not always (G. S. Deryabina, P. A. Kozhevnikov, Commun. Algebra, 27, no. 9 (1999), 4525-4530; S. I. Adyan, Dokl. Math., 62, no. 2 (2000), 174-176).
\end{openproblem}

\plainproblemsection

\begin{openproblem}
	13.35. Does every non-soluble pro- $p$ -group of cohomological dimension 2 contain a free non-abelian pro- $p$ -subgroup? O. V. Mel'nikov
\end{openproblem}

\plainproblemsection

\begin{openproblem}
	13.36. For a finitely generated pro- $p$ -group $G$ set ${a}_{n}\left( G\right)  = {\dim }_{{\mathbb{F}}_{p}}{I}^{n}/{I}^{n + 1}$ , where $I$ is the augmentation ideal of the group ring ${\mathbb{F}}_{p}\left\lbrack  \left\lbrack  G\right\rbrack  \right\rbrack$ . We define the growth of $G$ to be the growth of the sequence ${\left\{  {a}_{n}\left( G\right) \right\}  }_{n \in  \mathbb{N}}$ .
	
	a) If the growth of $G$ is exponential, does it follow that $G$ contains a free pro- $p$ - subgroup of rank 2?
	
	c) Do there exist pro- $p$ -groups of finite cohomological dimension which are not $p$ -adic analytic, and whose growth is slower than an exponential one?
	
	O. V. Mel'nikov
\end{openproblem}

\plainproblemsection

\begin{openproblem}
	13.36. b) (A. Lubotzky, A. Shalev). For a finitely generated pro- $p$ -group $G$ set ${a}_{n}\left( G\right)  = {\dim }_{{\mathbb{F}}_{p}}{I}^{n}/{I}^{n + 1}$ , where $I$ is the augmentation ideal of the group ring ${\mathbb{F}}_{p}\left\lbrack  \left\lbrack  G\right\rbrack  \right\rbrack$ . We define the growth of $G$ to be the growth of the sequence ${\left\{  {a}_{n}\left( G\right) \right\}  }_{n \in  \mathbb{N}}$ . Is the growth of $G$ exponential if $G$ contains a finitely generated closed subgroup of exponential growth? O. V. Mel'nikov
	
	b) Not necessarily. Let $G$ be the Nottingham group over ${\mathbb{F}}_{p}$ . The growth of ${c}_{n}\left( G\right)  \mathrel{\text{:=}}$ ${\log }_{p}\left| {G : {\omega }_{n}\left( G\right) }\right|$ , where $\left\{  {{\omega }_{n}\left( G\right) }\right\}$ is the Zassenhaus filtration, is linear, so by Quillen’s theorem the growth of $G$ is subexponential. But $G$ has a 2-generator free pro- $p$ subgroup (R. Camina, J. Algebra, 196 (1997), 101-113) with exponential growth. (M. Ershov, Letter of 19.10.2009.)
\end{openproblem}

\plainproblemsection

\begin{openproblem}
	13.37. Let $G$ be a torsion-free pro- $p$ -group, $U$ an open subgroup of $G$ . Suppose that $U$ is a pro- $p$ -group with a single defining relation. Is it true that then $G$ is also a pro- $p$ -group with a single defining relation? O. V. Mel'nikov
\end{openproblem}

\plainproblemsection

\begin{openproblem}
	13.38. Let $G = F/R$ be a pro- $p$ -group with one defining relation, where $R$ is the normal subgroup of a free pro- $p$ -group $F$ generated by a single element $r \in  {F}^{p}\left\lbrack  {F, F}\right\rbrack$ .
	
	a) Suppose that $r = {t}^{p}$ for some $t \in  F$ ; can $G$ contain a Demushkin group as a subgroup?
	
	b) Do there exist two pro- $p$ -groups ${G}_{1} \supset  {G}_{2}$ with one defining relation, where ${G}_{1}$ has elements of finite order, while the subgroup ${G}_{2}$ is torsion-free? $O.V$ . Mel’nikov a) Yes, it can. b) Yes, they exist. The pro- $p$ -group with presentation $\langle a, b \mid  {\left\lbrack  a, b\right\rbrack  }^{p} = 1\rangle$ contains the Demushkin group $\left\langle  {{x}_{1},\ldots ,{x}_{2g} \mid  \mathop{\prod }\limits_{{i = 1}}^{g}\left\lbrack  {{x}_{{2i} - 1},{x}_{2i}}\right\rbrack   = 1}\right\rangle$ for some $g > 1$ . (O. V. Mel'nikov, Letter of February, 28, 1999).
\end{openproblem}

\plainproblemsection

\begin{openproblem}
	13.39. Let $A$ be an associative ring with unity and with torsion-free additive group, and let ${F}^{A}$ be the tensor product of a free group $F$ by $A$ (A. G. Myasnikov, V. N. Re-meslennikov, Siberian Math. J.,35, no. 5 (1994),986-996); then ${F}^{A}$ is a free exponential group over $A$ ; in (A. G. Myasnikov, V. N. Remeslennikov, Int. J. Algebra Comput., $6\left( {1996}\right) ,{687} - {711})$ , it is shown how to construct ${F}^{A}$ in terms of free products with amalgamation.
	
	*a) (G. Baumslag). Is ${F}^{A}$ residually nilpotent torsion-free?
	
	b) Is ${F}^{A}$ a linear group?
	
	*c) (G. Baumslag). Is the Magnus homomorphism of ${F}^{\mathbb{Q}}$ into the group of power series over the rational number field $\mathbb{Q}$ faithful or not?
	
	In the case where $\langle 1\rangle$ is a pure subgroup of the additive group of $A$ , there are positive answers to questions "a" and "b" (A. M. Gaglione, A. G. Myasnikov, V. N. Remes-lennikov, D. Spellman, Commun. Algebra, 25 (1997), 631-648). In the same paper it is shown that "a" is equivalent to "c". It is also known that the Magnus homomorphism is one-to-one on any subgroup of ${F}^{\mathbb{Q}}$ of the type $\left\langle  {F, t \mid  u = {t}^{n}}\right\rangle$ (G. Baumslag, Commun. Pure Appl. Math., 21 (1968), 491-506).
	
	d) Is the universal theory of ${F}^{A}$ decidable?
	
	e) (G. Baumslag). Can free $A$ -groups be characterized by a length function?
	
	f) (G. Baumslag). Does a free $\mathbb{Q}$ -group admit a free action on some $\Lambda$ -tree? See definition in (R. Alperin, H. Bass, in: Combinatorial group theory and topology, Alta, Utah, 1984 (Ann. Math. Stud., 111), Princeton Univ. Press, 1987, 265-378).
	
	A. G. Myasnikov, V. N. Remeslennikov
	
	*a) Yes, it is (A. Jaikin-Zapirain, Preprint, 2020, http://matematicas.uam.es/ ~andrei.jaikin/preprints/baumslag.pdf)
	
	*c) Yes, it is faithful (A. Jaikin-Zapirain, Preprint, 2020, http://matematicas.uam.es/~andrei.jaikin/preprints/baumslag.pdf).
\end{openproblem}

\plainproblemsection

\begin{openproblem}
	13.40. A group $G$ is said to be $\omega$ -residually free if, for every finite set of non-trivial elements of $G$ , there is a homomorphism of $G$ into a free group such that the images of all these elements are non-trivial. Is every finitely-generated $\omega$ -residually free group embeddable in a free $\mathbb{Z}\left\lbrack  x\right\rbrack$ -group? A. G. Myasnikov, V. N. Remeslennikov
	
	Yes, it is (O. Kharlampovich, A. Myasnikov, J. Algebra, 200 (1998), 472-570).
\end{openproblem}

\plainproblemsection

\begin{openproblem}
	13.41. Is the elementary theory of the class of all groups acting freely on $\Lambda$ -trees decidable? A. G. Myasnikov, V. N. Remeslennikov
\end{openproblem}

\plainproblemsection

\begin{openproblem}
	13.42. Prove that the tensor $A$ -completion (see A. G. Myasnikov, V. N. Remeslenni-kov, Siberian Math. J., 35, no. 5 (1994), 986-996) of a free nilpotent group can be non-nilpotent. A. G. Myasnikov, V. N. Remeslennikov
\end{openproblem}

\plainproblemsection

\begin{openproblem}
	13.43. (G. R. Robinson). Let $G$ be a finite group and $B$ be a $p$ -block of characters of $G$ . Conjecture: If the defect group $D = D\left( B\right)$ of the block $B$ is non-abelian, and if $\left| {D : Z\left( D\right) }\right|  = {p}^{a}$ , then each character in $B$ has height strictly less than $a$ .
	
	G. R. Robinson’s comment of 2020: the conjecture is proved mod CFSG for $p \neq  2$ in (Z. Feng, C. Li, Y. Liu, G. Malle, J. Zhang, Compos. Math., 155, no. 6 (2019), 1098-1117). J. Olsson
\end{openproblem}

\plainproblemsection

\begin{openproblem}
	13.44. For any partition of an arbitrary group $G$ into finitely many subsets $G =$ ${A}_{1} \cup  \cdots  \cup  {A}_{n}$ , there exists a subset of the partition ${A}_{i}$ and a finite subset $F \subseteq  G$ , such that $G = {A}_{i}^{-1}{A}_{i}F$ (I. V. Protasov, Siberian Math. J.,34, no. 5 (1993),938-952). Can the subset $F$ always be chosen so that $\left| F\right|  \leq  n$ ? This is true for amenable groups.
	
	I. V. Protasov
\end{openproblem}

\plainproblemsection

\begin{openproblem}
	13.45. Every infinite group $G$ of regular cardinality $\mathfrak{m}$ can be partitioned into two subsets $G = {A}_{1} \cup  {A}_{2}$ so that ${A}_{1}F \neq  G$ and ${A}_{2}F \neq  G$ for every subset $F \subset  G$ of cardinality less than $\mathfrak{m}$ . Is this statement true for groups of singular cardinality? I. V. Protasov
	
	No, not always. The answer depends on the algebraic structure of $G$ . In particular, this is true for a free group, but the statement does not hold for every Abelian group $G$ of singular cardinality (I. Protasov, S. Slobodianiuk, Quest. Answers Gen. Topology, 33, no. 2 (2015), 61-70).
\end{openproblem}

\plainproblemsection

\begin{openproblem}
	13.46. Can every uncountable abelian group of finite odd exponent be partitioned into two subsets so that neither of them contains cosets of infinite subgroups? Among countable abelian groups, such partitions exist for groups with finitely many involutions. I. V. Protasov
	
	Yes, it can (E. G. Zelenyuk, Math. Notes, 67 (2000), 599-602).
\end{openproblem}

\plainproblemsection

\begin{openproblem}
	13.47. Can every countable abelian group with finitely many involutions be partitioned into two subsets that are dense in every group topology? I. V. Protasov Yes, it can (E. G. Zelenyuk, Ukrain. Math. J., 51, no. 1 (1999), 44-50).
\end{openproblem}

\plainproblemsection

\begin{openproblem}
	13.48. (W. W. Comfort, J. van Mill). A topological group is said to be irresolvable if every two of its dense subsets intersect non-trivially. Does every non-discrete irresolvable group contain an infinite subgroup of exponent 2? I. V. Protasov
\end{openproblem}

\plainproblemsection

\begin{openproblem}
	13.49. (V.I. Malykhin). Can a topological group be partitioned into two dense subsets, if there are infinitely many free ultrafilters on the group converging to the identity? I. V. Protasov
\end{openproblem}

\plainproblemsection

\begin{openproblem}
	13.50. Let $\mathfrak{F}$ be a local Fitting class. Is it true that there are no maximal elements in the partially ordered by inclusion set of the Fitting classes contained in $\mathfrak{F}$ and distinct from $\mathfrak{F}$ ? A. N. Skiba
	
	No, there may exist maximal elements (N. V. Savel'eva, N. T. Vorob'ev, Siberian Math. J., 49, no. 6 (2008), 1124-1130).
\end{openproblem}

\plainproblemsection

\begin{openproblem}
	13.51. Is every finite modular lattice embeddable in the lattice of formations of finite groups? A. N. Skiba
\end{openproblem}

\plainproblemsection

\begin{openproblem}
	13.52. The dimension of a finitely based variety of algebras $\mathcal{V}$ is defined to be the maximal length of a basis (that is, an independent generating set) of the ${SC}$ -theory ${SC}\left( \mathcal{V}\right)$ , which consists of the strong Mal’cev conditions satisfied on $\mathcal{V}$ . The dimension is defined to be infinite if the lengths of bases in ${SC}\left( \mathcal{V}\right)$ are not bounded. Does every finite abelian group generate a variety of finite dimension? D. M. Smirnov
\end{openproblem}

\plainproblemsection

\begin{openproblem}
	13.53. Let $a, b$ be elements of finite order of the infinite group $G = \langle a, b\rangle$ . Is it true that there are infinitely many elements $g \in  G$ such that the subgroup $\left\langle  {a,{b}^{g}}\right\rangle$ is infinite? A. I. Sozutov
\end{openproblem}

\plainproblemsection

\begin{openproblem}
	13.54. a) Is it true that, for $p$ sufficiently large, every (finite) $p$ -group can be a complement in some Frobenius group (see 6.55)? A. I. Sozutov
\end{openproblem}

\plainproblemsection

\begin{openproblem}
	13.54. b) Is it true that every group is embeddable in the kernel of some Frobenius group (see Archive, 6.53)? A. I. Sozutov
	
	Yes, it is true (V. V. Bludov, Siberian Math. J., 38, no. 6 (1997), 1054-1056).
\end{openproblem}

\plainproblemsection

\begin{openproblem}
	13.55. Does there exist a Golod group (see 9.76), which is isomorphic to an ${AT}$ - group? For a definition of an ${AT}$ -group see (A. V. Rozhkov, Math. Notes,40, no. 5 (1986), 827-836). A. V. Timofeenko
\end{openproblem}

\plainproblemsection

\begin{openproblem}
	13.56. (A. Shalev). Let $G$ be a finite $p$ -group of sectional rank $r$ , and $\varphi$ an automorphism of $G$ having exactly $m$ fixed points. Is the derived length of $G$ bounded by a function depending on $r$ and $m$ only? E. I. Khukhro
	
	Yes, it is (A. Jaikin-Zapirain, Israel J. Math., 129 (2002), 209-220).
\end{openproblem}

\plainproblemsection

\begin{openproblem}
	13.57. Let $\varphi$ be an automorphism of prime order $p$ of a finite group $G$ such that ${C}_{G}\left( \varphi \right)  \leq  Z\left( G\right)$ .
	
	a) Is $G$ soluble if $p = 3$ ? V. D. Mazurov and T. L. Nedorezov proved in (Algebra and Logic,35, no. 6 (1996),392-397) that the group $G$ is soluble for $p = 2$ , and there are examples of unsoluble $G$ for all $p > 3$ .
	
	b) If $G$ is soluble, is the derived length of $G$ bounded in terms of $p$ ?
	
	c) (V. K. Kharchenko). If $G$ is a $p$ -group, is the derived length of $G$ bounded in terms of $p$ ? (V. V. Bludov produced a simple example showing that the nilpotency class cannot be bounded.) E. I. Khukhro
\end{openproblem}

\plainproblemsection

\begin{openproblem}
	13.58. Let $\varphi$ be an automorphism of prime order $p$ of a nilpotent (periodic) group $G$ such that ${C}_{G}\left( \varphi \right)$ is a group of finite sectional rank $r$ . Does $G$ possess a normal subgroup $N$ which is nilpotent of class bounded by a function of $p$ only and is such that $G/N$ is a group of finite sectional rank bounded in terms of $r$ and $p$ ?
	
	This was proved for $p = 2$ in (P. Shumyatsky, Arch. Math., $\mathbf{{71}}\left( {1998}\right) ,{425} - {432}$ ). E. I. Khukhro
	
	Yes, it does (E. I. Khukhro, J. London Math. Soc., 77 (2008), 130-148).
\end{openproblem}

\plainproblemsection

\begin{openproblem}
	13.59. One can show that any extension of shape ${\mathbb{Z}}^{d}.\mathcal{S}{L}_{d}\left( \mathbb{Z}\right)$ is residually finite, unless possibly $d = 3$ or $d = 5$ . Are there in fact any extensions of this shape that fail to be residually finite when $d = 5$ ? As shown in (P. R. Hewitt, Groups/St. Andrews' 93 in Galway, Vol. 2 (London Math. Soc. Lecture Note Ser., 212), Cambridge Univ. Press, ${1995},{305} - {313}$ ), there is an extension of ${\mathbb{Z}}^{3}$ by $\mathcal{S}{L}_{3}\left( \mathbb{Z}\right)$ that is not residually finite. Usually it is true that an extension of an arithmetic subgroup of a Chevalley group over a rational module is residually finite. Is it ever false, apart from the examples of shape ${\mathbb{Z}}^{3}.\mathcal{S}{L}_{3}\left( \mathbb{Z}\right)$ ? P. R. Hewitt
\end{openproblem}

\plainproblemsection

\begin{openproblem}
	13.60. If a locally graded group $G$ is a product of two subgroups of finite special rank, is $G$ of finite special rank itself? The answer is affirmative if both factors are periodic groups. (N. S. Chernikov, Ukrain. Math. J., 42, no. 7 (1990), 855-861).
\end{openproblem}

\plainproblemsection

\begin{openproblem}
	13.61. We call a metric space narrow if it is quasiisometric to a subset of the real line, and wide otherwise. Let $G$ be a group with the finite set of generators $A$ , and let $\Gamma  = \Gamma \left( {G, A}\right)$ be the Cayley graph of $G$ with the natural metric. Suppose that, after deleting any narrow subset $L$ from $\Gamma$ , at most two connected components of the graph $\Gamma  \smallsetminus  L$ can be wide, and there exists at least one such a subset $L$ yielding exactly two wide components in $\Gamma  \smallsetminus  L$ . Is it true that $\Gamma$ is quasiisometric to an Euclidean or a hyperbolic plane? V. A. Churkin
	
	No, it is not true in general (O. V. Bogopol'skii, Preprint, Novosibirsk, 1998 (Russian)). See also new problem 14.98.
\end{openproblem}

\plainproblemsection

\begin{openproblem}
	13.62. Let $U$ and $V$ be non-cyclic subgroups of a free group. Does the inclusion $\left\lbrack  {U, U}\right\rbrack   \leq  \left\lbrack  {V, V}\right\rbrack$ imply that $U \leq  V$ ? V. P. Shaptala
	
	No, it does not (M. J. Dunwoody, Arch. Math., 16 (1965), 153-157).
\end{openproblem}

\plainproblemsection

\begin{openproblem}
	13.63. Let ${\pi }_{e}\left( G\right)$ denote the set of orders of elements of a group $G$ . For $\Gamma  \subseteq  \mathbb{N}$ let $h\left( \Gamma \right)$ denote the number of non-isomorphic finite groups $G$ with ${\pi }_{e}\left( G\right)  = \Gamma$ . Is there a number $k$ such that, for every $\Gamma$ , either $h\left( \Gamma \right)  \leq  k$ , or $h\left( \Gamma \right)  = \infty$ ? $\;W.J.{Sh}$ No, there is no such number: $h\left( {{\pi }_{e}\left( {{L}_{3}\left( {73}^{r}\right) }\right) }\right)  = r + 1$ for any $r \geq  0$ (A. V. Zavarnitsine, J. Group Theory, 7, no. 1 (2004), 81-97).
\end{openproblem}

\plainproblemsection

\begin{openproblem}
	13.64. Let ${\pi }_{e}\left( G\right)$ denote the set of orders of elements of a group $G$ . A group $G$ is said to be an $O{C}_{n}$ -group if ${\pi }_{e}\left( G\right)  = \{ 1,2,\ldots , n\}$ . Is every $O{C}_{n}$ -group locally finite? Do there exist infinite $O{C}_{n}$ -groups for $n \geq  7$ ?W. J. Shi
\end{openproblem}

\plainproblemsection

\begin{openproblem}
	13.65. A finite simple group is called a ${K}_{n}$ -group if its order is divisible by exactly $n$ different primes. The number of ${K}_{3}$ -groups is known to be 8 . The ${K}_{4}$ -groups are classified mod CFSG (W. J. Shi, in: Group Theory in China (Math. Appl., 365), Kluwer, 1996, 163-181) and some significant further results are obtained in (Yann Bugeaud, Zhenfu Cao, M. Mignotte, J. Algebra, 241 (2001), 658-668). But the question remains: is the number of ${K}_{4}$ -groups finite or infinite?W. J. Shi
\end{openproblem}

\plainproblemsection

\begin{openproblem}
	13.66. Let $F$ be a
	
	a) free;
	
	b) free metabelian
	
	group of finite rank. Let $M$ denote the set of all endomorphisms of $F$ with noncyclic images. Can one choose two elements $g, h \in  F$ such that, for every $\varphi ,\psi  \in  M$ , equalities $\varphi \left( g\right)  = \psi \left( g\right)$ and $\varphi \left( h\right)  = \psi \left( h\right)$ imply that $\varphi  = \psi$ , that is, the endomorphisms in $M$ are uniquely determined by their values at $g$ and $h$ ? $\;V.E$ . Shpil’rain
	
	a) Yes, one can (D. Lee, J. Algebra, 247 (2002), no. 2, 509-540).
	
	b) No, not always (E. I. Timoshenko, Math. Notes, 62, no. 5-6 (1998), 767-770).
\end{openproblem}

\plainproblemsection

\begin{openproblem}
	13.67. Let $G$ be a ${T}_{0}$ -group (see 12.101), $i$ an involution in $G$ and $G = \left\langle  {{i}^{g} \mid  g \in  G}\right\rangle$ . Is the centralizer ${C}_{G}\left( i\right)$ residually periodic? V. P. Shunkov
\end{openproblem}

\plainproblemsection

\begin{openproblem}
	*14.1. Suppose that $G$ is a finite group with no non-trivial normal subgroups of odd order, and $\varphi$ is its 2-automorphism centralizing a Sylow 2-subgroup of $G$ . Is it true that ${\varphi }^{2}$ is an inner automorphism of $G$ ? R. Zh. Aleev
	
	*Yes, it is true (G. Glauberman, Math. Z., 107 (1968), 1-20).
\end{openproblem}

\plainproblemsection

\begin{openproblem}
	14.2. (S. D. Berman). Prove that every automorphism of the centre of the integral group ring of a finite group induces a monomial permutation on the set of the class sums. R. Zh. Aleev
\end{openproblem}

\plainproblemsection

\begin{openproblem}
	14.3. Is it true that every central unit of the integral group ring of a finite group is a product of a central element of the group and a symmetric central unit? (A unit is symmetric if it is fixed by the canonical antiinvolution that transposes the coefficients at the mutually inverse elements.) R. Zh. Aleev
\end{openproblem}

\plainproblemsection

\begin{openproblem}
	*14.4. a) Is it true that there exists a nilpotent group $G$ for which the lattice $\mathcal{L}\left( G\right)$ of all group topologies is not modular? (It is known that for abelian groups the lattice $\mathcal{L}\left( G\right)$ is modular and that there are groups for which this lattice is not modular: V. I. Arnautov, A. G. Topale, Izv. Akad. Nauk Moldova Mat., 1997, no. 1, 84-92 (Russian).)
	
	b) Is it true that for every countable nilpotent non-abelian group $G$ the lattice $\mathcal{L}\left( G\right)$ of all group topologies is not modular? V. I. Arnautov
	\par\smallskip\hrule\smallskip
	*a) Yes, it exists (V. Arnautov, A. Topalà, Bul. Acad. Stiint. Repub. Moldova, Mat., 1998, no. 2(27), 130-131).
	
	*b) No, there are such groups with modular $\mathcal{L}\left( G\right)$ (Dekui Peng, Preprint,2023, https://arxiv.org/abs/2310.08269).
	\par\smallskip\hrule\smallskip
\end{openproblem}

\plainproblemsection

\begin{openproblem}
	14.5. Let $G$ be an infinite group admitting non-discrete Hausdorff group topologies, and $\mathcal{L}\left( G\right)$ the lattice of all group topologies on $G$ .
	a) Is it true that for any natural number $k$ there exists a non-refinable chain ${\tau }_{0} < {\tau }_{1} < \cdots  < {\tau }_{k}$ of length $k$ of Hausdorff topologies in $\mathcal{L}\left( G\right)$ ? (For countable nilpotent groups this is true (A. G. Topale, Deposited in VINITI, 25.12.98, no. 3849- V 98 (Russian)).)
	
	b) Let $k, m, n$ be natural numbers and let $G$ be a nilpotent group of class $k$ . Suppose that ${\tau }_{0} < {\tau }_{1} < \cdots  < {\tau }_{m}$ and ${\tau }_{0}^{\prime } < {\tau }_{1}^{\prime } < \cdots  < {\tau }_{n}^{\prime }$ are non-refinable chains of Hausdorff topologies in $\mathcal{L}\left( G\right)$ such that ${\tau }_{0} = {\tau }_{0}^{\prime }$ and ${\tau }_{m} = {\tau }_{n}^{\prime }$ . Is it true that $m \leq  n \cdot  k$ and this inequality is best-possible? (This is true if $k = 1$ , since for $G$ abelian the lattice $\mathcal{L}\left( G\right)$ is modular.)
	
	c) Is it true that there exists a countable $G$ such that in the lattice $\mathcal{L}\left( G\right)$ there are a finite non-refinable chain ${\tau }_{0} < {\tau }_{1} < \cdots  < {\tau }_{k}$ of Hausdorff topologies and an infinite chain $\left\{  {{\tau }_{\gamma }^{\prime } \mid  \gamma  \in  \Gamma }\right\}$ of topologies such that ${\tau }_{0} < {\tau }_{\gamma }^{\prime } < {\tau }_{k}$ for any $\gamma  \in  \Gamma$ ?
	
	*d) Let $G$ be an abelian group, $k$ a natural number. Let ${A}_{k}$ be the set of all those Hausdorff group topologies on $G$ that, for every topology $\tau  \in  {A}_{k}$ , any non-refinable chain of topologies starting from $\tau$ and terminating at the discrete topology has length $k$ . Is it true that ${A}_{k}\bigcap \left\{  {{\tau }_{\gamma }^{\prime } \mid  \gamma  \in  \Gamma }\right\}   \neq  \varnothing$ for any infinite non-refinable chain $\left\{  {{\tau }_{\gamma }^{\prime } \mid  \gamma  \in  \Gamma }\right\}$ of Hausdorff topologies containing the discrete topology? (This is true for $k = 1$ .)
	
	*d) No; moreover, no infinite abelian group satisfies this property with $k = 2$ (Dekui Peng, Preprint, 2023, https://arxiv.org/abs/2310.08269).
\end{openproblem}

\plainproblemsection

\begin{openproblem}
	14.6. A group $\Gamma$ is said to have Property ${P}_{\text{nai }}$ if, for any finite subset $F$ of $\Gamma  \smallsetminus  \{ 1\}$ , there exists an element ${y}_{0} \in  \Gamma$ of infinite order such that, for each $x \in  F$ , the canonical epimorphism from the free product $\langle x\rangle  * \left\langle  {y}_{0}\right\rangle$ onto the subgroup $\left\langle  {x,{y}_{0}}\right\rangle$ of $\Gamma$ generated by $x$ and ${y}_{0}$ is an isomorphism. For $n \in  \{ 2,3,\ldots \}$ , does ${PS}{L}_{n}\left( \mathbb{Z}\right)$ have Property ${P}_{\text{nai }}$ ? More generally, if $\Gamma$ is a lattice in a connected real Lie group $G$ which is simple and with centre reduced to $\{ 1\}$ , does $\Gamma$ have Property ${P}_{\text{nai }}$ ?
	
	Answers are known to be "yes" if $n = 2$ , and more generally if $G$ has real rank 1 (M. Bekka, M. Cowling, P. de la Harpe, Publ. Math. IHES, 80 (1994), 117-134).
	
	Comments of 2018: Property ${P}_{\text{nai }}$ has been established for relatively hyperbolic groups (G. Arzhantseva, A. Minasyan, J. Funct. Anal., 243, no. 1 (2007), 345-351) and for groups acting on $\operatorname{CAT}\left( 0\right)$ cube complexes with appropriate conditions (A. Kar, M. Sageev, Comment. Math. Helv., 91 (2016), 543-561). Some progress has been made in the preprint (T. Poznansky, http://arxiv.org/abs/0812.2486).
	
	P. de la Harpe
\end{openproblem}

\plainproblemsection

\begin{openproblem}
	14.7. Let ${\Gamma }_{g}$ be the fundamental group of a closed surface of genus $g \geq  2$ . For each finite system $S$ of generators of ${\Gamma }_{g}$ , let ${\beta }_{S}\left( n\right)$ denote the number of elements of ${\Gamma }_{g}$ which can be written as products of at most $n$ elements of $S \cup  {S}^{-1}$ , and let $\omega \left( {{\Gamma }_{g}, S}\right)  = \mathop{\limsup }\limits_{{n \rightarrow  \infty }}\sqrt[n]{{\beta }_{S}\left( n\right) }$ be the growth rate of the sequence ${\left( {\beta }_{S}\left( n\right) \right) }_{n \geq  0}$ . Compute the infimum $\omega \left( {\Gamma }_{g}\right)$ of the $\omega \left( {{\Gamma }_{g}, S}\right)$ over all finite sets of generators of ${\Gamma }_{g}$ .
	
	It is easy to see that, for a free group ${F}_{k}$ of rank $k \geq  2$ , the corresponding infimum is $\omega \left( {F}_{k}\right)  = {2k} - 1$ (M. Gromov, Structures métriques pour les variétés riemanniennes, Cedic/F. Nathan, Paris,1981, Ex. 5.13). As any generating set of ${\Gamma }_{g}$ contains a subset of ${2g} - 1$ elements generating a subgroup of infinite index in ${\Gamma }_{g}$ with abelianization ${\mathbb{Z}}^{{2g} - 1}$ , hence a subgroup which is free of rank ${2g} - 1$ , it follows that $\omega \left( {\Gamma }_{g}\right)  \geq  {4g} - 3$ . P. de la Harpe
\end{openproblem}

\plainproblemsection

\begin{openproblem}
	14.8. Let $G$ denote the group of germs at $+ \infty$ of orientation-preserving homeomorphisms of the real line $\mathbb{R}$ . Let $\alpha  \in  G$ be the germ of $x \mapsto  x + 1$ . What are the germs $\beta  \in  G$ for which the subgroup $\langle \alpha ,\beta \rangle$ of $G$ generated by $\alpha$ and $\beta$ is free of rank 2 ?
	
	If $\beta$ is the germ of $x \mapsto  {x}^{k}$ for an odd integer $k \geq  3$ , it is known that $\langle \alpha ,\beta \rangle$ is free of rank 2. The proofs of this rely on Galois theory (for $k$ an odd prime: S. White, J. Algebra,118 (1988),408-422; for any odd $k \geq  3$ : S. A. Adeleke, A. M. W. Glass, L. Morley, J. London Math. Soc., $\mathbf{{43}}\left( {1991}\right) ,{255} - {268}$ , and for any odd $k \neq   \pm  1$ and any even $k > 0$ in several papers by S. D. Cohen and A. M. W. Glass). P. de la Harpe
\end{openproblem}

\plainproblemsection

\begin{openproblem}
	14.9. (Well-known problem). Let ${W}^{ * }\left( {F}_{k}\right)$ denote the von Neumann algebra of the free group of rank $k \in  \left\{  {2,3,\ldots ,{\aleph }_{0}}\right\}$ . Is ${W}^{ * }\left( {F}_{k}\right)$ isomorphic to ${W}^{ * }\left( {F}_{l}\right)$ for $k \neq  l$ ?
	
	For a group $G$ , recall that ${W}^{ * }\left( G\right)$ is an appropriate completion of the group algebra $\mathbb{C}G$ (see e. g. S. Sakai, ${C}^{ * }$ -algebras and ${W}^{ * }$ -algebras, Springer,1971, in particular Problem 4.4.44). It is known that either ${W}^{ * }\left( {F}_{k}\right)  \cong  {W}^{ * }\left( {F}_{l}\right)$ for all $k, l \in  \left\{  {2,3,\ldots ,{\aleph }_{0}}\right\}$ , or that the ${W}^{ * }\left( {F}_{k}\right)$ are pairwise non-isomorphic (F. Radulescu, Invent. Math.,115 (1994), 347-389, Corollary 4.7). P. de la Harpe
\end{openproblem}

\plainproblemsection

\begin{openproblem}
	14.10. b) Find an explicit embedding of $\mathbb{Q}$ in a finitely generated group; such a group exists by Theorem IV in (G. Higman, B. H. Neumann, H. Neumann, J. London Math. Soc., 24 (1949), 247-254). P. de la Harpe
	
	Found (V. H. Mikaelian, Int. J. Math. Math. Sci., 2005 (2005), 2119-2123).
\end{openproblem}

\plainproblemsection

\begin{openproblem}
	14.10. *a) (Well-known problem). It is known that any recursively presented group embeds in a finitely presented group (G. Higman, Proc. Royal Soc. London Ser. A, 262 (1961),455-475). Find an explicit and "natural" finitely presented group $\Gamma$ and an embedding of the additive group of the rationals $\mathbb{Q}$ in $\Gamma$ .
	
	c) Find an explicit and "natural" finitely presented group ${\Gamma }_{n}$ and an embedding of $G{L}_{n}\left( \mathbb{Q}\right)$ in ${\Gamma }_{n}$ .
	
	Another phrasing of the same problems is: find a simplicial complex $X$ which covers a finite complex such that the fundamental group of $X$ is $\mathbb{Q}$ or, respectively, $G{L}_{n}\left( \mathbb{Q}\right)$ . P. de la Harpe
	
	*a) Such an embedding is found in (J. Belk, J. Hyde, F. Matucci, Bull. Amer. Math. Soc., 59, no. 4 (2022), 561-567).
\end{openproblem}

\plainproblemsection

\begin{openproblem}
	14.11. (Yu. I. Merzlyakov). It is a well-known fact that for the ring $R = \mathbb{Q}\left\lbrack  {x, y}\right\rbrack$ the elementary group ${E}_{2}\left( R\right)$ is distinct from $S{L}_{2}\left( R\right)$ . Find a minimal subset $A \subseteq  S{L}_{2}\left( R\right)$ such that $\left\langle  {{E}_{2}\left( R\right) , A}\right\rangle   = S{L}_{2}\left( R\right)$ . V. G. Bardakov
\end{openproblem}

\plainproblemsection

\begin{openproblem}
	14.12. (Yu. I. Merzlyakov, J. S. Birman). Is it true that all braid groups ${B}_{n}, n \geq  3$ , are conjugacy separable? V. G. Bardakov
\end{openproblem}

\plainproblemsection

\begin{openproblem}
	14.13. a) By definition the commutator length of an element $z$ of the derived subgroup of a group $G$ is the least possible number of commutators from $G$ whose product is equal to $z$ . Does there exist a simple group on which the commutator length is not bounded?
	
	b) Does there exist a finitely presented simple group on which the commutator length is not bounded? V. G. Bardakov
	
	a) Simple groups with this property were constructed in (J. Barge, É. Ghys, Math. Ann., 294 (1992), 235-265), and finitely generated simple ones in (A. Muranov, Int. J. Algebra Comput., 17 (2007), 607-659).
	
	b) Yes, it exists (P.-E. Caprace, K. Fujiwara, Geom. Funct. Anal., 19 (2010), 1296- 1319).
\end{openproblem}

\plainproblemsection

\begin{openproblem}
	14.14. (C. C. Edmunds, G. Rosenberger). We call a pair of natural numbers(k, m) admissible if in the derived subgroup ${F}_{2}^{\prime }$ of a free group ${F}_{2}$ there is an element $w$ such that the commutator length of the element ${w}^{m}$ is equal to $k$ . Find all admissible pairs.
	
	It is known that every pair $\left( {k,2}\right) , k \geq  2$ , is admissible (V. G. Bardakov, Algebra and Logic, 39 (2000), 224-251). V. G. Bardakov
\end{openproblem}

\plainproblemsection

\begin{openproblem}
	14.15. For the automorphism group ${A}_{n} = \operatorname{Aut}{F}_{n}$ of a free group ${F}_{n}$ of rank $n \geq  3$ find the supremum ${k}_{n}$ of the commutator lengths of the elements of ${A}_{n}^{\prime }$ .
	
	It is easy to show that ${k}_{2} = \infty$ . On the other hand, the commutator length of any element of the derived subgroup of $\mathop{\lim }\limits_{ \rightarrow  }{A}_{n}$ is at most 2 (R. K. Dennis, L. N. Vaserstein, $K$ -Theory, 2, N 6 (1989), 761-767).
\end{openproblem}

\plainproblemsection

\begin{openproblem}
	14.16. Following Yu. I. Merzlyakov we define the width of a verbal subgroup $V\left( G\right)$ of a group $G$ with respect to the set of words $V$ as the smallest $m \in  \mathbb{N} \cup  \{ \infty \}$ such that every element of $V\left( G\right)$ can be written as a product of $\leq  m$ values of words from $V \cup  {V}^{-1}$ . It is known that the width of any verbal subgroup of a finitely generated group of polynomial growth is finite. Is this statement true for finitely generated groups of intermediate growth? V. G. Bardakov
\end{openproblem}

\plainproblemsection

\begin{openproblem}
	14.17. Let $\operatorname{IMA}\left( G\right)$ denote the subgroup of the automorphism group Aut $G$ consisting of all automorphisms that act trivially on the second derived quotient $G/{G}^{\prime \prime }$ . Find generators and defining relations for $\operatorname{IMA}\left( {F}_{n}\right)$ , where ${F}_{n}$ is a free group of rank $n \geq  3$ .
	
	V. G. Bardakov
\end{openproblem}

\plainproblemsection

\begin{openproblem}
	14.18. We say that a family of groups $\mathcal{D}$ discriminates a group $G$ if for any finite subset $\left\{  {{a}_{1},\ldots ,{a}_{n}}\right\}   \subseteq  G \smallsetminus  \{ 1\}$ there exists a group $D \in  \mathcal{D}$ and a homomorphism $\varphi  : G \rightarrow  D$ such that ${a}_{j}\varphi  \neq  1$ for all $j = 1,\ldots , n$ . Is every finitely generated group acting freely on some $\Lambda$ -tree discriminated by torsion-free hyperbolic groups?
	
	G. Baumslag, A. G. Myasnikov, V. N. Remeslennikov
\end{openproblem}

\plainproblemsection

\begin{openproblem}
	*14.19. We say that a group $G$ has the Noetherian Equation Property if every system of equations over $G$ in finitely many variables is equivalent to some finite part of it. Does an arbitrary hyperbolic group have the Noetherian Equation Property?
	
	Comment of 2013: the conjecture holds for torsion-free hyperbolic groups (Z. Sela, Proc. London Math. Soc., 99, no. 1 (2009), 217-273) and for the larger class of toral relatively hyperbolic groups (D. Groves, J. Geom. Topol., 9 (2005), 2319-2358).
	
	G. Baumslag, A. G. Myasnikov, V. N. Remeslennikov
	
	*Yes, it does (R. Weidmann, C. Reinfeldt, Ann. Math. Blaise Pascal, 26, no. 2 (2019), 125-214).
\end{openproblem}

\plainproblemsection

\begin{openproblem}
	14.20. Does a free product of groups have the Noetherian Equation Property if this property is enjoyed by the factors?
	
	G. Baumslag, A. G. Myasnikov, V. N. Remeslennikov
\end{openproblem}

\plainproblemsection

\begin{openproblem}
	14.21. Does a free pro- $p$ -group have the Noetherian Equation Property?
	
	G. Baumslag, A. G. Myasnikov, V. N. Remeslennikov
\end{openproblem}

\plainproblemsection

\begin{openproblem}
	14.22. Prove that any irreducible system of equations $S\left( {{x}_{1},\ldots ,{x}_{n}}\right)  = 1$ with coefficients in a torsion-free linear group $G$ is equivalent over $G$ to a finite system $T\left( {{x}_{1},\ldots ,{x}_{n}}\right)  = 1$ satisfying an analogue of Hilbert’s Nullstellensatz, i. e. ${\operatorname{Rad}}_{G}\left( T\right)  = \sqrt{T}$ . This is true if $G$ is a free group (O. Kharlampovich, A. Myasnikov, J. Algebra, 200 (1998), 472-570).
	
	Here both $S$ and $T$ are regarded as subsets of $G\left\lbrack  X\right\rbrack   = G * F\left( X\right)$ , a free product of $G$ and a free group on $X = \left\{  {{x}_{1},\ldots ,{x}_{n}}\right\}$ . By definition, ${\operatorname{Rad}}_{G}\left( T\right)  =$ $\left\{  {w\left( {{x}_{1},\ldots ,{x}_{n}}\right)  \in  G\left\lbrack  X\right\rbrack   \mid  w\left( {{g}_{1},\ldots ,{g}_{n}}\right)  = 1}\right.$ for any solution ${g}_{1},\ldots ,{g}_{n} \in  G$ of the system $T\left( X\right)  = 1\}$ , and $\sqrt{T}$ is the minimal normal isolated subgroup of $G\left\lbrack  X\right\rbrack$ containing $T$ . G. Baumslag, A. G. Myasnikov, V. N. Remeslennikov
\end{openproblem}

\plainproblemsection

\begin{openproblem}
	14.23. Let ${F}_{n}$ be a free group with basis $\left\{  {{x}_{1},\ldots ,{x}_{n}}\right\}$ , and let $\left| \cdot \right|$ be the length function with respect to this basis. For $\alpha  \in  \operatorname{Aut}{F}_{n}$ we put $\parallel \alpha \parallel  = \max \left\{  {\left| {\alpha \left( {x}_{1}\right) }\right| ,\ldots ,\left| {\alpha \left( {x}_{n}\right) }\right| }\right\}$ . Is it true that there is a recursive function $f : \mathbb{N} \rightarrow  \mathbb{N}$ with the following property: for any $\alpha  \in  \operatorname{Aut}{F}_{n}$ there is a basis $\left\{  {{y}_{1},\ldots ,{y}_{k}}\right\}$ of $\operatorname{Fix}\left( \alpha \right)  = \{ x \mid  \alpha \left( x\right)  = x\}$ such that $\left| {y}_{i}\right|  \leq  f\left( {\parallel \alpha \parallel }\right)$ for all $i = 1,\ldots , k$ ? O. V. Bogopol'skiǐ
\end{openproblem}

\plainproblemsection

\begin{openproblem}
	14.24. Let Aut ${F}_{n}$ be the automorphism group of a free group of rank $n$ with norm $\parallel  \cdot  \parallel$ as in 14.23. Does there exist a recursive function $f : \mathbb{N} \times  \mathbb{N} \rightarrow  \mathbb{N}$ with the following property: for any two conjugate elements $\alpha ,\beta  \in$ Aut ${F}_{n}$ there is an element $\gamma  \in  \operatorname{Aut}{F}_{n}$ such that ${\gamma }^{-1}{\alpha \gamma } = \beta$ and $\parallel \gamma \parallel  \leq  f\left( {\parallel \alpha \parallel ,\parallel \beta \parallel }\right)$ ? O. V. Bogopol'skiǐ
\end{openproblem}

\plainproblemsection

\begin{openproblem}
	14.25. Let ${qG}$ denote the quasivariety generated by a group $G$ . Is it true that there exists a finitely generated group $G$ such that the set of proper maximal subquasiva-rieties in ${qG}$ is infinite? A. I. Budkin
	
	Yes, it is (V. V. Bludov, Algebra and Logic, 41, no. 1 (2002), 1-8).
\end{openproblem}

\plainproblemsection

\begin{openproblem}
	14.26. A quasivariety $\mathcal{M}$ is closed under direct $\mathbb{Z}$ -wreath products if the direct wreath product $G\wr \mathbb{Z}$ belongs to $\mathcal{M}$ for every $G \in  \mathcal{M}$ (here $\mathbb{Z}$ is an infinite cyclic group). Is the quasivariety generated by the class of all nilpotent torsion-free groups closed under direct $\mathbb{Z}$ -wreath products? A. I. Budkin
\end{openproblem}

\plainproblemsection

\begin{openproblem}
	14.27. Let $\Gamma$ be a group generated by a finite set $S$ . Assume that there exists a nested sequence ${F}_{1} \subset  {F}_{2} \subset  \cdots$ of finite subsets of $\Gamma$ such that (i) ${F}_{k} \neq  {F}_{k + 1}$ for all $k \geq  1,\;$ (ii) $\Gamma  = \mathop{\bigcup }\limits_{{k \geq  1}}{F}_{k},\;$ (iii) $\;\mathop{\lim }\limits_{{k \rightarrow  \infty }}\frac{\left| \partial {F}_{k}\right| }{\left| {F}_{k}\right| } = 0$ , where, by definition,
	
	$\partial {F}_{k} = \left\{  {\gamma  \in  \Gamma  \smallsetminus  {F}_{k} \mid  \text{there exists}s \in  S\text{such that}{\gamma s} \in  {F}_{k}}\right\}$ , and (iv) there exist constants $c \geq  0, d \geq  1$ such that $\left| {F}_{k}\right|  \leq  c{k}^{d}$ for all $k \geq  1$ . Does it follow that $\Gamma$ has polynomial growth? A. G. Vaillant
	
	No, not always. These properties are enjoyed by every group of intermediate growth (V. G. Bardakov, Algebra and Logic, 40, no. 1 (2001), 12-16).
\end{openproblem}

\plainproblemsection

\begin{openproblem}
	14.28. Let $\mathfrak{F}$ be a soluble Fitting formation of finite groups with Kegel’s property, that is, $\mathfrak{F}$ contains every finite group of the form $G = {AB} = {BC} = {CA}$ if $A, B, C$ are in $\mathfrak{F}$ . Is $\mathfrak{F}$ a saturated formation?
	
	Comment of 2013: Some progress was made in (A. Ballester-Bolinches, L. M. Ez-querro, J. Group Theory, 8, no. 5 (2005), 605-611). A. F. Vasiliev
\end{openproblem}

\plainproblemsection

\begin{openproblem}
	14.29. Is there a soluble Fitting class of finite groups $\mathfrak{F}$ such that $\mathfrak{F}$ is not a formation and ${A}_{\mathfrak{F}} \cap  {B}_{\mathfrak{F}} \subseteq  {G}_{\mathfrak{F}}$ for every finite soluble group of the form $G = {AB}$ ?
	
	A. F. Vasiliev
\end{openproblem}

\plainproblemsection

\begin{openproblem}
	14.30. Let lFit $\mathfrak{X}$ be the local Fitting class generated by a set of groups $\mathfrak{X}$ and let $\Psi \left( G\right)$ be the smallest normal subgroup of a finite group $G$ such that $\operatorname{lFit}\left( {\Psi \left( G\right)  \cap  M}\right)  = \operatorname{lFit}M$ for every $M \vartriangleleft   \vartriangleleft  G$ (K. Doerk, P. Hauck, Arch. Math.,35, no. 3 (1980),218-227). We say that a Fitting class $\mathfrak{F}$ is saturated if $\Psi \left( G\right)  \in  \mathfrak{F}$ implies that $G \in  \mathfrak{F}$ . Is it true that every non-empty soluble saturated Fitting class is local? N. T. Vorob'ev
\end{openproblem}

\plainproblemsection

\begin{openproblem}
	14.31. Is the lattice of Fitting subclasses of the Fitting class generated by a finite soluble group finite? N. T. Vorob'ev
\end{openproblem}

\plainproblemsection

\begin{openproblem}
	14.32. Extending the classical definition of formations, let us define a formation of (not necessarily finite) groups as a nonempty class of groups closed under taking homomorphic images and subdirect products with finitely many factors. Must every first order axiomatizable formation of groups be a variety? A. M. Gaglione, D. Spellman No, every variety of groups that contains a finite nonsolvable member contains an axiomatic subformation that is not a variety (K. A. Kearnes, J. Group Theory, 13, No. $2\left( {2010}\right) ,{233} - {241})$ .
\end{openproblem}

\plainproblemsection

\begin{openproblem}
	14.33. Does there exist a finitely presented pro- $p$ -group ( $p$ being a prime) which contains an isomorphic copy of every countably based pro- $p$ -group?
	
	R. I. Grigorchuk
	
	Yes, it exists: the Nottingham group over ${\mathbb{F}}_{p}$ for $p > 2$ , which was shown to be finitely presented in (M. V. Ershov, J. London Math. Soc., 71 (2005), 362-378). (M. Ershov, Letter of 19.10.2009.)
\end{openproblem}

\plainproblemsection

\begin{openproblem}
	14.34. By definition, a locally compact group has the Kazhdan T-property if the trivial representation is an isolated point in the natural topological space of unitary representations of the group. Does there exist a profinite group with two dense discrete subgroups one of which is amenable, and the other has the Kazhdan $T$ -property?
	
	See (A. Lubotzky, Discrete groups, expanding graphs and invariant measures (Progress in Mathematics, Boston, Mass., 125), Birkhäuser, Basel, 1994) for further motivation. R. I. Grigorchuk, A. Lubotzky
	
	Yes, it exists (M. Kassabov, Invent. Math., 170 (2007), 297-326; M. Ershov, A. Jaikin-Zapirain, Invent. Math., 179, 303-347).
\end{openproblem}

\plainproblemsection

\begin{openproblem}
	14.35. Is every finitely presented group of prime exponent finite? N. D. Gupta
\end{openproblem}

\plainproblemsection

\begin{openproblem}
	14.36. A group $G$ is called a $T$ -group if every subnormal subgroup of $G$ is normal, while $G$ is said to be a $\bar{T}$ -group if all of its subgroups are $T$ -groups. Is it true that every non-periodic locally graded $\bar{T}$ -group must be abelian? F. de Giovanni
\end{openproblem}

\plainproblemsection

\begin{openproblem}
	14.37. Let $G\left( n\right)$ be one of the classical groups (special, orthogonal, or symplectic) of $\left( {n \times  n}\right)$ -matrices over an infinite field $K$ of non-zero characteristic, and $M\left( n\right)$ the space of all $\left( {n \times  n}\right)$ -matrices over $K$ . The group $G\left( n\right)$ acts diagonally by conjugation on the space $M{\left( n\right) }^{m} = \underset{m}{\underbrace{M\left( n\right)  \oplus  \cdots  \oplus  M\left( n\right) }}$ . Find generators of the algebra of invariants
	
	In characteristic 0 they were found in (C. Procesi, Adv. Math., 19 (1976), 306- 381). Comment of 2001: in positive characteristic the problem is solved for all cases excepting the orthogonal groups in characteristic 2 and special orthogonal groups of even degree (A. N. Zubkov, Algebra and Logic, 38, no. 5 (1999), 299-318). Comment of 2009: ...and for special orthogonal groups of even degree over (infinite) fields of odd characteristic (A. A. Lopatin, J. Algebra, 321 (2009), 1079-1106). A. N. Zubkov
\end{openproblem}

\plainproblemsection

\begin{openproblem}
	14.38. For every pro- $p$ -group $G$ of $\left( {2 \times  2}\right)$ -matrices for $p \neq  2$ an analogue of the Tits Alternative holds: either $G$ is soluble, or the variety of pro- $p$ -groups generated by $G$ contains the group $\overline{\left\langle  \left( \begin{array}{ll} 1 & t \\  0 & 1 \end{array}\right) ,\left( \begin{array}{ll} 1 & 0 \\  t & 1 \end{array}\right) \right\rangle  } \leq  S{L}_{2}\left( {{\mathbb{F}}_{p}\left\lbrack  \left\lbrack  t\right\rbrack  \right\rbrack  }\right)$ (A. N. Zubkov, Algebra and Logic,29, no. 4 (1990),287-301). Is the same result true for matrices of size $\geq  3$ for $p \neq  2?$ A. N. Zubkov
\end{openproblem}

\plainproblemsection

\begin{openproblem}
	14.39. Let $F\left( V\right)$ be a free group of some variety of pro- $p$ -groups $V$ . Is there a uniform bound for the exponents of periodic elements in $F\left( V\right)$ ? This is true if the variety $V$ is metabelian (A. N. Zubkov, Dr. Sci. Diss., Omsk, 1997 (Russian)). A. N. Zubkov
\end{openproblem}

\plainproblemsection

\begin{openproblem}
	14.40. Is a free pro- $p$ -group representable as an abstract group by matrices over a commutative-associative ring with 1 ? A. N. Zubkov
\end{openproblem}

\plainproblemsection

\begin{openproblem}
	14.41. A group $G$ is said to be para-free if all factors ${\gamma }_{i}\left( G\right) /{\gamma }_{i + 1}\left( G\right)$ of its lower central series are isomorphic to the corresponding lower central factors of some free group, and $\mathop{\bigcap }\limits_{{i = 1}}^{\infty }{\gamma }_{i}\left( G\right)  = 1$ . Is an arbitrary para-free group representable by matrices over a commutative-associative ring with 1 ? A. N. Zubkov
\end{openproblem}

\plainproblemsection

\begin{openproblem}
	14.42. Is a free pro- $p$ -group representable by matrices over an associative-commutative profinite ring with 1 ?
	
	A negative answer is equivalent to the fact that every linear pro- $p$ -group satisfies a non-trivial pro- $p$ -identity; this is known to be true in dimension 2 (for $p \neq  2$ : A. N. Zubkov, Siberian Math. J., 28, no. 5 (1987), 742-747; for $p = 2$ : D. E.-C. Ben-Ezra, E. Zelmanov, Trans. Amer. Math. Soc., 374, no. 6 (2021), 4093-4128).
	
	A. N. Zubkov, V. N. Remeslennikov
\end{openproblem}

\plainproblemsection

\begin{openproblem}
	14.43. Suppose that a finite group $G$ has the form $G = {AB}$ , where $A$ and $B$ are nilpotent subgroups of classes $\alpha$ and $\beta$ respectively; then $G$ is soluble (O. H. Kegel, H. Wielandt,1961). Although the derived length $\operatorname{dl}\left( G\right)$ of $G$ need not be bounded by $\alpha  + \beta$ (see Archive,5.17), can one bound $\operatorname{dl}\left( G\right)$ by a (linear) function of $\alpha$ and $\beta$ ?
\end{openproblem}

\plainproblemsection

\begin{openproblem}
	14.44. Let $k\left( X\right)$ denote the number of conjugacy classes of a finite group $X$ . Suppose that a finite group $G = {AB}$ is a product of two subgroups $A, B$ of coprime orders. Is it true that $k\left( {AB}\right)  \leq  k\left( A\right) k\left( B\right)$ ?
	
	Note that one cannot drop the coprimeness condition and the answer is positive if one of the subgroups is normal, see Archive, 11.43. L. S. Kazarin, J. Sangroniz
\end{openproblem}

\plainproblemsection

\begin{openproblem}
	14.45. Does there exist a (non-abelian simple) linearly right-orderable group all of whose proper subgroups are cyclic? U. E. Kaljulaid
\end{openproblem}

\plainproblemsection

\begin{openproblem}
	14.46. A finite group $G$ is said to be almost simple if $T \leq  G \leq  \operatorname{Aut}\left( T\right)$ for some nonabelian simple group $T$ . By definition, a finite linear space consists of a set $V$ of points, together with a collection of $k$ -element subsets of $V$ , called lines $\left( {k \geq  3}\right)$ , such that every pair of points is contained in exactly one line. Classify the finite linear spaces which admit a line-transitive almost simple subgroup $G$ of automorphisms which acts transitively on points.
	
	A. Camina, P. Neumann and C. E. Praeger have solved this problem in the case where $T$ is an alternating group. In (A. Camina, C. E. Praeger, Aequat. Math.,61 (2001), 221-232) it is shown that a line-transitive group of automorphisms of a finite linear space which is point-quasiprimitive (i. e. all of whose non-trivial normal subgroups are point-transitive) is almost simple or affine.
	
	Remarks of 2005: The cases of the following groups, but not almost simple groups with this socle, were dealt with: $\operatorname{PSU}\left( {3, q}\right)$ (W. Liu, Linear Algebra Appl., $\mathbf{{374}}\left( {2003}\right)$ , 291-305); PSL(2, q) (W. Liu, J. Combin. Theory (A), 103 (2003), 209-222); Sz(q) (W. Liu, Discrete Math., 269 (2003), 181-190); Ree(q) (W. Liu, Europ. J. Combin., 25 (2004), 311-325); sporadic simple groups were done in (A. R. Camina, F. Spiezia, J. Combin. Des., 8 (2000), 353-362). Remark of 2009: It was shown (N. Gill, Trans. Amer. Math. Soc., 368 (2016), 3017-3057) that for non-desarguesian projective planes a point-quasiprimitive group would be affine. Remark of 2013: The problem is solved for large-dimensional classical groups (A. Camina, N. Gill, A. E. Zalesski, Bull. Belg. Math. Soc. Simon Stevin, 15, no. 4 (2008), 705-731). A. Camina, C. E. Praeger
\end{openproblem}

\plainproblemsection

\begin{openproblem}
	14.47. Is the lattice of all soluble Fitting classes of finite groups modular?
	
	S. F. Kamornikov, A. N. Skiba
\end{openproblem}

\plainproblemsection

\begin{openproblem}
	14.48. If an equation over a free group $F$ has no solution in $F$ , is there a finite quotient of $F$ in which the equation has no solution? No, not always (T. Coulbois, A. Khelif, Proc. Amer. Math. Soc., 127, no. 4 (1999), 963-965).
\end{openproblem}

\plainproblemsection

\begin{openproblem}
	14.49. Is $S{L}_{3}\left( \mathbb{Z}\right)$ a factor group of the modular group ${PS}{L}_{2}\left( \mathbb{Z}\right)$ ? Since the latter is isomorphic to the free product of two cyclic groups of orders 2 and 3 , the question asks if $S{L}_{3}\left( \mathbb{Z}\right)$ can be generated by two elements of orders 2 and 3. M. Cond No, it is not (M. C. Tamburini, P. Zucca, Atti Accad. Naz. Lincei Cl. Sci. Fis. Mat. Natur. Rend. Lincei (9) Mat. Appl., 11, no. 1 (2000), 5-7; Ya. N. Nuzhin, Math. Notes,70, no. 1-2 (2001),71-78). M. Conder has shown, however, that $S{L}_{3}\left( \mathbb{Z}\right)$ has a subgroup of index 57 that is a factor group of the modular group ${\operatorname{PSL}}_{2}\left( \mathbb{Z}\right)$ . Also it has been shown in (M. C. Tamburini, J. S. Wilson, N. Gavioli, J. Algebra, 168 (1994), 353-370) that $S{L}_{d}\left( \mathbb{Z}\right)$ is a factor group of the modular group ${PS}{L}_{2}\left( \mathbb{Z}\right)$ for all $d \geq  {28}$ .
\end{openproblem}

\plainproblemsection

\begin{openproblem}
	14.50. (Z.I. Borevich). A subgroup $A$ of a group $G$ is said to be paranormal (respectively, polynormal) if ${A}^{x} \leq  \left\langle  {{A}^{u} \mid  u \in  \left\langle  {A,{A}^{x}}\right\rangle  }\right\rangle$ (respectively, ${A}^{x} \leq  \left\langle  {{A}^{u} \mid  u \in  {A}^{\langle x\rangle }}\right\rangle$ ) for any $x \in  G$ . Is every polynormal subgroup of a finite group paranormal?
	
	Not always (V. I. Mysovskikh, Dokl. Math., 60 (1999), 71-72). A. S. Kondratiev
\end{openproblem}

\plainproblemsection

\begin{openproblem}
	14.51. (Well-known problems). Is there a finite basis for the identities of any
	
	a) abelian-by-nilpotent group?
	
	b) abelian-by-finite group?
	
	c) abelian-by-(finite nilpotent) group? A. N. Krasil'nikov
\end{openproblem}

\plainproblemsection

\begin{openproblem}
	14.52. It is known that if a finitely generated group is residually torsion-free nilpotent, then the group is residually finite $p$ -group, for every prime $p$ . Is the converse true? Yu. V. Kuz’min
	
	Not always (B. Hartley, in: Symposia Mathematica, Bologna, Vol. XVII, Convegno sui Gruppi Infiniti, INDAM, Roma, 1973, Academic Press, London, 1976, 225-234).
\end{openproblem}

\plainproblemsection

\begin{openproblem}
	14.53. Conjecture: Let $G$ be a profinite group such that the set of solutions of the equation ${x}^{n} = 1$ has positive Haar measure. Then $G$ has an open subgroup $H$ and an element $t$ such that all elements of the coset ${tH}$ have order dividing $n$ .
	
	This is true in the case $n = 2$ . It would be interesting to see whether similar results hold for profinite groups in which the set of solutions of some equation has positive measure.
	
	Comment of 2021: This is also proved for $n = 3$ (A. Abdollahi, M. S. Malekan, Adv. Group Theory Appl., 13 (2022), 71-81). L. Levai, L. Pyber
\end{openproblem}

\plainproblemsection

\begin{openproblem}
	14.54. Let $k\left( G\right)$ denote the number of conjugacy classes of a finite group $G$ . Is it true that $k\left( G\right)  \leq  \left| N\right|$ for some nilpotent subgroup $N$ of $G$ ? This is true if $G$ is simple; besides, always $k\left( G\right)  \leq  \left| S\right|$ for some soluble subgroup $S \leq  G$ .
	
	M. W. Liebeck, L. Pyber
\end{openproblem}

\plainproblemsection

\begin{openproblem}
	14.55. a) Prove that the Nottingham group $J = N\left( {\mathbb{Z}/p\mathbb{Z}}\right)$ (as defined in Archive, 12.24) is finitely presented for $p > 2$ . C. R. Leedham-Green
	
	Proved for $p > 2$ (M.V. Ershov, J. London Math. Soc.,71 (2005),362-378).
\end{openproblem}

\plainproblemsection

\begin{openproblem}
	14.55. b) Prove that the Nottingham group $J = N\left( {\mathbb{Z}/p\mathbb{Z}}\right)$ (as defined in Archive, 12.24) is finitely presented for $p = 2$ . C. R. Leedham-Green
\end{openproblem}

\plainproblemsection

\begin{openproblem}
	14.56. Prove that if $G$ is an infinite pro- $p$ -group with $G/{\gamma }_{{2p} + 1}\left( G\right)$ isomorphic to $J/{\gamma }_{{2p} + 1}\left( J\right)$ then $G$ is isomorphic to $J$ , where $J$ is the Nottingham group.
	
	C. R. Leedham-Green
\end{openproblem}

\plainproblemsection

\begin{openproblem}
	14.57. Describe the hereditarily just infinite pro- $p$ -groups of finite width.
	
	Definitions: A pro- $p$ -group $G$ has finite width ${p}^{d}$ if $\left| {{\gamma }_{i}\left( G\right) /{\gamma }_{i + 1}\left( G\right) }\right|  \leq  {p}^{d}$ for all $i$ , and $G$ is hereditarily just infinite if $G$ is infinite and each of its open subgroups has no closed normal subgroups of infinite index. C. R. Leedham-Green
\end{openproblem}

\plainproblemsection

\begin{openproblem}
	14.58. a) Suppose that $A$ is a periodic group of regular automorphisms of an abelian group. Is $A$ cyclic if $A$ has prime exponent? V. D. Mazurov
\end{openproblem}

\plainproblemsection

\begin{openproblem}
	14.58. b) Suppose that $A$ is a periodic group of regular automorphisms of an abelian group. Is $A$ finite if $A$ is generated by elements of order3? $\;V.D.$ Mazurov Yes, it is (A. Kh. Zhurtov, Siberian Math. J., 41, no. 2 (2000), 268-275).
\end{openproblem}

\plainproblemsection

\begin{openproblem}
	14.59. Suppose that $G$ is a triply transitive group in which a stabilizer of two points contains no involutions, and a stabilizer of three points is trivial. Is it true that $G$ is similar to ${PG}{L}_{2}\left( P\right)$ in its natural action on the projective line $P \cup  \{ \infty \}$ , for some field $P$ of characteristic 2? This is true under the condition that the stabilizer of two points is periodic. V. D. Mazurov
\end{openproblem}

\plainproblemsection

\begin{openproblem}
	*14.60. Suppose that $H$ is a non-trivial normal subgroup of a finite group $G$ such that the factor-group $G/H$ is isomorphic to one of the simple groups ${L}_{n}\left( q\right) , n \geq  3$ . Is it true that $G$ has an element whose order is distinct from the order of any element in $G/H$ ? V. D. Mazurov
	
	*Yes, it is true: for $n \neq  4$ proved in (A.V. Zavarnitsine, Siberian Math. J.,49 (2008), 246-256), and for $n = 4$ in (M. A. Grechkoseeva, S. V. Skresanov, Siberian Electron. Math. Rep., 17 (2020), 585-589). Editors' comment: previous claim that it is not true for $n = 4$ was erroneous.
\end{openproblem}

\plainproblemsection

\begin{openproblem}
	14.61. Determine all pairs(S, G), where $\mathcal{S}$ is a semipartial geometry and $G$ is an almost simple flag-transitive group of automorphisms of $\mathcal{S}$ . A system of points and lines(P, B)is a semipartial geometry with parameters $\left( {\alpha , s, t,\mu }\right)$ if every point belongs to exactly $t + 1$ lines (two different points belong to at most one line); every line contains exactly $s + 1$ points; for any anti-flag $\left( {a, l}\right)  \in  \left( {P, B}\right)$ the number of lines containing $a$ and intersecting $l$ is either 0 or $\alpha$ ; and for any non-collinear points $a, b$ there are exactly $\mu$ points collinear with $a$ and with $b$ . A. A. Makhnëv
\end{openproblem}

\plainproblemsection

\begin{openproblem}
	14.62. Suppose that $H$ is a non-soluble normal subgroup of a finite group $G$ . Does there always exist a maximal soluble subgroup $S$ of $H$ such that $G = H \cdot  {N}_{G}\left( S\right)$ ? V. S. Monakhov
	
	Yes, it always exists mod CFSG (V. I. Zenkov, V. S. Monakhov, D. O. Revin, Algebra and Logic, 43, no. 2 (2004), 102-108).
\end{openproblem}

\plainproblemsection

\begin{openproblem}
	14.63. What are the composition factors of non-soluble finite groups all of whose normalizers of Sylow subgroups are 2-nilpotent, in particular, supersoluble? V. S. Monakhov
	
	Described (L. S. Kazarin, A. A. Volochkov, Mathematics in Yaroslavl' Univ.: Coll. of Surveys to 30-th Anniv. of Math. Fac., Yaroslavl', 2006, 243-255 (Russian)).
\end{openproblem}

\plainproblemsection

\begin{openproblem}
	14.64. (M. F. Newman). Classify the finite 5-groups of maximal class; computer calculations suggest some conjectures (M. F. Newman, in: Groups-Canberra, 1989 (Lecture Notes Math., 1456), Springer, Berlin, 1990, 49-62). A. Moretó
\end{openproblem}

\plainproblemsection

\begin{openproblem}
	14.65. (Well-known problem). For a finite group $G$ let $\rho \left( G\right)$ denote the set of prime numbers dividing the order of some conjugacy class, and $\sigma \left( G\right)$ the maximum number of primes dividing the order of some conjugacy class. Is it true that $\left| {\rho \left( G\right) }\right|  \leq  {3\sigma }\left( G\right)$ ?
	
	A possible linear bound for $\left| {\rho \left( G\right) }\right|$ in terms of $\sigma \left( G\right)$ cannot be better than ${3\sigma }\left( G\right)$ , since there is a family of groups $\left\{  {G}_{n}\right\}$ such that $\mathop{\lim }\limits_{{n \rightarrow  \infty }}\left| {\rho \left( {G}_{n}\right) }\right| /\sigma \left( {G}_{n}\right)  = 3$ (C. Casolo, S. Dolfi, Rend. Sem. Mat. Univ. Padova, 96 (1996), 121-130).
	
	Comment of 2009: A quadratic bound was obtained in (A. Moretó, Int. Math. Res. Not., 2005, no. 54, 3375-3383), and a linear bound in (C. Casolo, S. Dolfi, J. Group Theory, 10, no. 5 (2007), 571-583). A. Moretó
\end{openproblem}

\plainproblemsection

\begin{openproblem}
	14.66. (Well-known problem). Let $G$ be a finite soluble group, $\pi \left( G\right)$ the set of primes dividing the order of $G$ , and $\nu \left( G\right)$ the maximum number of primes dividing the order of some element. Does there exist a linear bound for $\left| {\pi \left( G\right) }\right|$ in terms of $\nu \left( G\right)$ ?
	
	Yes, it does (T. M. Keller, J. Algebra, 178, no. 2 (1995), 643-652). A. Moretó
\end{openproblem}

\plainproblemsection

\begin{openproblem}
	14.67. Suppose that $a$ is a non-trivial element of a finite group $G$ such that $\left| {{C}_{G}\left( a\right) }\right|  \geq$ $\left| {{C}_{G}\left( x\right) }\right|$ for every non-trivial element $x \in  G$ , and let $H$ be a nilpotent subgroup of $G$ which is normalized by ${C}_{G}\left( a\right)$ . Is it true that $H \leq  {C}_{G}\left( a\right)$ ? This is true if $H$ is abelian, or if $H$ is a ${p}^{\prime }$ -group for some prime number $p \in  \pi \left( {Z\left( {{C}_{G}\left( a\right) }\right) }\right)$ .
	
	I. T. Mukhamet'yanov, A. N. Fomin
\end{openproblem}

\plainproblemsection

\begin{openproblem}
	14.68. (Well-known problem). Suppose that $F$ is an automorphism of order 2 of the polynomial ring ${R}_{n} = \mathbb{C}\left\lbrack  {{x}_{1},\ldots ,{x}_{n}}\right\rbrack  , n > 2$ . Does there exist an automorphism $G$ of ${R}_{n}$ such that ${G}^{-1}{FG}$ is a linear automorphism?
	
	This is true for $n = 2$ . M. V. Neshchadim
\end{openproblem}

\plainproblemsection

\begin{openproblem}
	14.69. For every finite simple group find the minimum of the number of generating involutions satisfying an additional condition, in each of the following cases.
	
	a) The product of the generating involutions equals 1 .
	
	b) (Malle-Saxl-Weigel). All generating involutions are conjugate.
	
	c) (Malle-Saxl-Weigel). The conditions a) and b) are simultaneously satisfied. Editors' comment: partial progress was obtained in (J. Ward, PhD Thesis, QMW, 2009, http://www.maths.qmul.ac.uk/~raw/JWardPhD.pdf).
	
	d) All generating involutions are conjugate and two of them commute.
	
	Ya. N. Nuzhin
\end{openproblem}

\plainproblemsection

\begin{openproblem}
	14.70. A group $G$ is called $n$ -Engel if it satisfies the identity $\left\lbrack  {x, y,\ldots , y}\right\rbrack   = 1$ , where $y$ is taken $n$ times. Are there non-nilpotent finitely generated $n$ -Engel groups? Comment of 2017: For $n = 2,3,4$ , finitely generated $n$ -Engel groups are nilpotent. A course of lectures devoted to constructing non-nilpotent finitely generated $n$ -Engel groups for sufficiently large $n$ is given by E. Rips and is available on the Internet.
	
	B. I. Plotkin
\end{openproblem}

\plainproblemsection

\begin{openproblem}
	14.71. Consider a free group $F$ of finite rank and an arbitrary group $G$ . Define the $G$ -closure ${\operatorname{cl}}_{G}\left( T\right)$ of any subset $T \subseteq  F$ as the intersection of the kernels of all those homomorphisms $\mu  : F \rightarrow  G$ of $F$ into $G$ that vanish on $T : {\operatorname{cl}}_{G}\left( T\right)  = \bigcap \{ \operatorname{Ker}\mu  \mid$ $\mu  : F \rightarrow  G;T \subseteq  \operatorname{Ker}\mu \}$ . Groups $G$ and $H$ are called geometrically equivalent if for every free group $F$ and every subset $T \subset  F$ the $G$ - and $H$ -closures of $T$ coincide: ${\operatorname{cl}}_{G}\left( T\right)  = {\operatorname{cl}}_{H}\left( T\right)$ . It is easy to see that if $G$ and $H$ are geometrically equivalent then they have the same quasiidentities. Is it true that if two groups have the same quasiidentities then they are geometrically equivalent? This is true for nilpotent groups. B. I. Plotkin
	
	No, not always (V. V. Bludov, Abstracts of the 7th Int. Conf. Groups and Group Rings, Suprasl, Poland, 1999, p. 6; R. Göbel, S. Shelah, Proc. Amer. Math. Soc., 130 (2002), 673-674 (electronic); A. G. Myasnikov, V. N. Remeslennikov, J. Algebra, 234, no. 1 (2000), 225-276). The latter paper contains also necessary and sufficient conditions for geometrical equivalence.
\end{openproblem}

\plainproblemsection

\begin{openproblem}
	14.72. Let $X$ be a regular algebraic variety over a field of arbitrary characteristic, and $G$ a finite cyclic group of automorphisms of $X$ . Suppose that the fixed point variety ${X}^{G}$ of $G$ is a regular hypersurface of $X$ (of codimension 1). Is the quotient variety $X/G$ regular?
\end{openproblem}

\plainproblemsection

\begin{openproblem}
	14.73. Conjecture: There is a function $f$ on the natural numbers such that, if $\Gamma$ is a finite, vertex-transitive, locally-quasiprimitive graph of valency $v$ , then the number of automorphisms fixing a given vertex is at most $f\left( v\right)$ . (By definition, a vertex-transitive graph $\Gamma$ is locally-quasiprimitive if the stabilizer in $\operatorname{Aut}\left( \Gamma \right)$ of a vertex $\alpha$ is quasiprimitive (see 14.46) in its action on the set of vertices adjacent to $\alpha$ .)
	
	To prove the conjecture above one need only consider the case where $\operatorname{Aut}\left( \Gamma \right)$ has the property that every non-trivial normal subgroup has at most two orbits on vertices (C. E. Praeger, Ars Combin., 19 A (1985), 149-163). The analogous conjecture for finite, vertex-transitive, locally-primitive graphs was made by R. Weiss in 1978 and is still open. For non-bipartite graphs, there is a "reduction" of Weiss' conjecture to the case where the automorphism group is almost simple (see 14.46 for definition) (M. Conder, C. H. Li, C. E. Praeger, Proc. Edinburgh Math. Soc. (2), 43, no. 1 (2000), ${129} - {138})$ .
	
	Comment of 2013: These conjectures have been proved in the case where there is an upper bound on the degree of any alternating group occurring as a quotient of a subgroup (C. E. Praeger, L. Pyber, P. Spiga, E. Szabó, Proc. Amer. Math. Soc., 140, no. 7 (2012), 2307-2318). Comment of 2021: This has been proved in the case where the group induced on the neighbourhood of a vertex has an abelian regular normal subgroup (P. Spiga, Bull. London Math. Soc., 48, no. 1 (2016), 12-18.) C. E. Praeger
\end{openproblem}

\plainproblemsection

\begin{openproblem}
	14.74. Let $k\left( G\right)$ denote the number of conjugacy classes of a finite group $G$ . Is it true that $k\left( G\right)  \leq  k\left( {P}_{1}\right) \cdots k\left( {P}_{s}\right)$ , where ${P}_{1},\ldots ,{P}_{s}$ are Sylow subgroups of $G$ such that $\left| G\right|  = \left| {P}_{1}\right| \cdots \left| {P}_{s}\right|$ ? L. Pyber
\end{openproblem}

\plainproblemsection

\begin{openproblem}
	14.75. Suppose that $\mathfrak{G} = \left\{  {{G}_{1},{G}_{2},\ldots }\right\}$ is a family of finite 2-generated groups which generates the variety of all groups. Is it true that a free group of rank 2 is residually in (8? L. Pyber
\end{openproblem}

\plainproblemsection

\begin{openproblem}
	14.76. Does there exist an absolute constant $c$ such that any finite $p$ -group $P$ has an abelian section $A$ satisfying ${\left| A\right| }^{c} > \left| P\right|$ ?
	
	By a result of A. Yu. Ol'shanskiĭ (Math. Notes, 23 (1978), 183-185) we cannot require $A$ to be a subgroup. By a result of J. G. Thompson (J. Algebra,13 (1969), 149-151) the existence of such a section $A$ would imply the existence of a class 2 subgroup $H$ of $P$ with ${\left| H\right| }^{c} > \left| P\right|$ . L. Pyber
\end{openproblem}

\plainproblemsection

\begin{openproblem}
	14.77. Let $p$ be a prime number and $X$ a finite set of powers of $p$ containing 1 . Is it true that $X$ is the set of all lengths of the conjugacy classes of some finite $p$ -group? J. Sangroniz
	
	Yes, it is (J. Cossey, T. Hawkes, Proc. Amer. Math. Soc., 128, no. 1 (2000), 49-51).
\end{openproblem}

\plainproblemsection

\begin{openproblem}
	14.78. Suppose that $\mathfrak{H} \subseteq  {\mathfrak{F}}_{1} \subseteq  \mathfrak{M}$ and $H \subseteq  {\mathfrak{F}}_{2} \subseteq  \mathfrak{M}$ where $\mathfrak{H}$ and $\mathfrak{M}$ are local formations of finite groups and ${\mathfrak{F}}_{1}$ is a complement for ${\mathfrak{F}}_{2}$ in the lattice of all formations between $\mathfrak{H}$ and $\mathfrak{M}$ . Is it true that ${\mathfrak{F}}_{1}$ and ${\mathfrak{F}}_{2}$ are local formations? This is true in the case $\mathfrak{H} = \left( 1\right)$ . A. N. Skiba
\end{openproblem}

\plainproblemsection

\begin{openproblem}
	14.79. Suppose that $\mathfrak{F} = \mathfrak{M}H = \operatorname{lFit}\left( G\right)$ is a soluble one-generated local Fitting class of finite groups where $\mathfrak{H}$ and $\mathfrak{M}$ are Fitting classes and $\mathfrak{M} \neq  \mathfrak{F}$ . Is $\mathfrak{H}$ a local Fitting class? A. N. Skiba
\end{openproblem}

\plainproblemsection

\begin{openproblem}
	14.80. Is the lattice of all totally local formations of finite groups modular? The definition of a totally local formation see in (L. A. Shemetkov, A. N. Skiba, Formatsii algebraicheskikh system, Moscow, Nauka, 1989 (Russian)). A. N. Skiba, L. A. Shemetkov
	
	Yes, even distributive (V. G. Safonov, Commun. Algebra, 35 (2007), 3495-3502).
\end{openproblem}

\plainproblemsection

\begin{openproblem}
	14.81. Prove that the formation generated by a finite group has only finitely many ${S}_{n}$ -closed subformations. A. N. Skiba, L. A. Shemetkov
\end{openproblem}

\plainproblemsection

\begin{openproblem}
	14.82. (Well-known problem). Describe the finite simple groups in which every element is a product of two involutions. A. I. Sozutov
	
	Described in (E. P. Vdovin, A. A. Gal’t, Siberian Math. J., 51 (2010), 610-615).
\end{openproblem}

\plainproblemsection

\begin{openproblem}
	14.83. We say that an infinite simple group $G$ is a monster of the third kind if for every non-trivial elements $a, b$ , of which at least one is not an involution, there are infinitely many elements $g \in  G$ such that $\left\langle  {a,{b}^{g}}\right\rangle   = G$ . (Compare with V. P. Shunkov’s definitions in Archive, 6.63, 6.64.) Is it true that every simple quasi-Chernikov group is a monster of the third kind? This is true for quasi-finite groups (A. I. Sozutov, Algebra and Logic, 36, no. 5 (1997), 336-348). (We say that a non- $\sigma$ group is quasi- $\sigma$ if all of its proper subgroups have the property $\sigma$ .) A. I. Sozutov
\end{openproblem}

\plainproblemsection

\begin{openproblem}
	14.84. An element $g$ of a (relatively) free group ${F}_{r}\left( \mathfrak{M}\right)$ of rank $r$ of a variety $\mathfrak{M}$ is said to be primitive if it can be included in a basis of ${F}_{r}\left( \mathfrak{M}\right)$ .
	
	a) Do there exist a group ${F}_{r}\left( \mathfrak{M}\right)$ and a non-primitive element $h \in  {F}_{r}\left( \mathfrak{M}\right)$ such that for some monomorphism $\alpha$ of ${F}_{r}\left( \mathfrak{M}\right)$ the element $\alpha \left( h\right)$ is primitive?
	
	b) Do there exist a group ${F}_{r}\left( \mathfrak{M}\right)$ and a non-primitive element $h \in  {F}_{r}\left( \mathfrak{M}\right)$ such that for some $n > r$ the element $h$ is primitive in ${F}_{n}\left( \mathfrak{M}\right)$ ? E. I. Timoshenko
\end{openproblem}

\plainproblemsection

\begin{openproblem}
	14.85. Suppose that an endomorphism $\varphi$ of a free metabelian group of rank $r$ takes every primitive element to a primitive one. Is $\varphi$ necessarily an automorphism? This is true for $r \leq  2$ . E. I. Timoshenko, V. Shpilrain
\end{openproblem}

\plainproblemsection

\begin{openproblem}
	14.86. Does there exist an infinite locally nilpotent $p$ -group that is equal to its commutator subgroup and in which every proper subgroup is nilpotent? J. Wiegold No, it does not exist (A. O. Asar, J. Lond. Math. Soc. (2), 61, no. 2 (2000), 412-422).
\end{openproblem}

\plainproblemsection

\begin{openproblem}
	14.88. We say that an element $u$ of a group $G$ is a test element if for any endomorphism $\varphi$ of $G$ the equality $\varphi \left( u\right)  = u$ implies that $\varphi$ is an automorphism of $G$ . Does a free soluble group of rank 2 and derived length $d > 2$ have any test elements? B. Fine, V. Shpilrain
	
	Yes, it does (E. I. Timoshenko, Algebra and Logic, 45, no. 4 (2006), 254-260.)
\end{openproblem}

\plainproblemsection

\begin{openproblem}
	14.89. (E. A. O’Brien, A. Shalev). Let $P$ be a finite $p$ -group of order ${p}^{m}$ and let $m = {2n} + e$ with $e = 0$ or 1 . By a theorem of P. Hall the number of conjugacy classes of $P$ has the form $n\left( {{p}^{2} - 1}\right)  + {p}^{e} + a\left( {{p}^{2} - 1}\right) \left( {p - 1}\right)$ for some integer $a \geq  0$ , which is called the abundance of $P$ .
	
	a) Is there a bound for the coclass of $P$ which depends only on $a$ ? Note that $a = 0$ implies coclass 1 and that all known examples with $a = 1$ have coclass $\leq  3$ . (The group $P$ has coclass $r$ if $\left| P\right|  = {p}^{c + r}$ where $c$ is the nilpotency class of $P$ .)
	
	b) Is there an element $s \in  P$ such that $\left| {{C}_{P}\left( s\right) }\right|  \leq  {p}^{f\left( a\right) }$ for some $f\left( a\right)$ depending only on $a$ ? We already know that we can take $f\left( 0\right)  = 2$ and it seems that $f\left( 1\right)  = 3$ .
	
	Note that A. Jaikin-Zapirain (J. Group Theory, 3, no. 3 (2000), 225-231) has proved that $\left| P\right|  \leq  {p}^{f\left( {p, a}\right) }$ for some function $f$ of $p$ and $a$ only.
	
	G. Fernández-Alcober
\end{openproblem}

\plainproblemsection

\begin{openproblem}
	14.90. Let $P$ be a finite $p$ -group of abundance $a$ and nilpotency class $c$ . Does there exist an integer $t = t\left( a\right)$ such that ${\gamma }_{i}\left( G\right)  = {\zeta }_{c - i + 1}\left( G\right)$ for $i \geq  t$ ? This holds for $a = 0$ with $t = 1$ , since $P$ has maximal class; it can be proved that for $a = 1$ one can take $t = 3$ . G. Fernández-Alcober
\end{openproblem}

\plainproblemsection

\begin{openproblem}
	14.91. Let $p$ be a fixed prime. Do there exist finite $p$ -groups of abundance $a$ for any $a \geq  0$ ? G. Fernández-Alcober
\end{openproblem}

\plainproblemsection

\begin{openproblem}
	14.92. (I. D. Macdonald). Every finite $p$ -group has at least $p - 1$ conjugacy classes of maximum size. I. D. Macdonald (Proc. Edinburgh Math. Soc., 26 (1983), 233- 239) constructed groups of order ${2}^{n}$ for any $n \geq  7$ with just one conjugacy class of maximum size. Are there any examples with exactly $p - 1$ conjugacy classes of maximum size for odd $p$ ? G. Fernández-Alcober
	
	Yes, there are, for $p = 3$ (A. Jaikin-Zapirain, M. F. Newman, E. A. O’Brien, Israel J. Math., 172, no. 1 (2009), 119-123).
\end{openproblem}

\plainproblemsection

\begin{openproblem}
	14.93. Let $N\left( {\mathbb{Z}/p\mathbb{Z}}\right)$ be the group defined in Archive,12.24 (the so-called "Nottingham group", or the "Wild group"). Find relations of $N\left( {\mathbb{Z}/p\mathbb{Z}}\right)$ as a pro- $p$ -group (it has two generators, e.g., $x + {x}^{2}$ and $x/\left( {1 - x}\right)$ ). Comment of 2009: for $p > 2$ see progress in (M. V. Ershov, J. London Math. Soc., 71 (2005), 362-378). I. B. Fesenko
\end{openproblem}

\plainproblemsection

\begin{openproblem}
	14.94. For each positive integer $r$ find the $p$ -cohomological dimension ${\operatorname{cd}}_{p}\left( {H}_{r}\right)$ where ${H}_{r}$ is the closed subgroup of $N\left( {\mathbb{Z}/p\mathbb{Z}}\right)$ consisting of the series $x\left( {1 + \mathop{\sum }\limits_{{i = 1}}^{\infty }{a}_{i}{x}^{{p}^{r}i}}\right)$ , ${a}_{i} \in  \mathbb{Z}/p\mathbb{Z}$ . I. B. Fesenko
\end{openproblem}

\plainproblemsection

\begin{openproblem}
	14.95. (C. R. Leedham-Green, Peter M. Neumann, J. Wiegold). For a finite $p$ - group $P$ , denote by $c = c\left( P\right)$ its nilpotency class and by $b = b\left( P\right)$ its breadth, that is, ${p}^{b}$ is the maximum size of a conjugacy class in $P$ . Class-Breadth Problem: Is it true that $c \leq  b + 1$ if $p \neq  2$ ?
	
	So far, the best known bound is $c < \frac{p}{p - 1}b + 1$ (C. R. Leedham-Green, P. M. Neumann, J. Wiegold, J. London Math. Soc. (2), 1 (1969), 409-420). For $p = 2$ for every $n \in  \mathbb{N}$ there exists a 2-group ${T}_{n}$ such that $c\left( {T}_{n}\right)  \geq  b\left( {T}_{n}\right)  + n$ (W. Felsch, J. Neubüser, W. Plesken, J. London Math. Soc. (2), 24 (1981), 113-122). A. Jaikin-Zapira 14.97. Is it true that for any two different prime numbers $p$ and $q$ there exists a non-primary periodic locally soluble $\{ p, q\}$ -group that can be represented as the product of two of its $p$ -subgroups? N. S. Chernikov
\end{openproblem}

\plainproblemsection

\begin{openproblem}
	14.96. Suppose that a finite $p$ -group $P$ admits an automorphism of order ${p}^{n}$ having exactly ${p}^{m}$ fixed points. By (E.I. Khukhro, Russ. Acad. Sci. Sbornik Math.,80 (1995), 435-444) then $P$ has a subgroup of index bounded in terms of $p, n$ and $m$ which is soluble of derived length bounded in terms of ${p}^{n}$ . Is it true that $P$ has also a subgroup of index bounded in terms of $p, n$ and $m$ which is soluble of derived length bounded in terms of $m$ ? There are positive answers in the cases of $m = 1$ (S. McKay, Quart. J. Math. Oxford, Ser. (2), 38 (1987), 489-502; I. Kiming, Math. Scand., 62 (1988), 153-172) and $n = 1$ (Yu. A. Medvedev, see Archive, 10.68). E. I. Khukhro Yes, it is true (A. Jaikin-Zapirain, Adv. Math., 153, no. 2 (2000), 391-402).
\end{openproblem}

\plainproblemsection

\begin{openproblem}
	14.98. We say that a metric space is a 2-end one (a narrow one), if it is quasiisometric to the real line $\mathbb{R}$ (respectively, to a subset of $\mathbb{R}$ ). All other spaces are said to be wide. Suppose that the Caley graph $\Gamma  = \Gamma \left( {G, A}\right)$ of a group $G$ with a finite set of generators $A$ in the natural metric contains a 2-end subset, and suppose that there is $\varepsilon  > 0$ such that the complement in $\Gamma$ to the $\varepsilon$ -neighbourhood of any connected 2-end subset contains exactly two wide connected components. Is it true that the group $G$ in the word metric is quasiisometric to the Euclidean or hyperbolic plane?
	
	V. A. Churkin
\end{openproblem}

\plainproblemsection

\begin{openproblem}
	14.99. a) A formation $\mathfrak{F}$ of finite groups is called superradical if it is ${S}_{n}$ -closed and contains every finite group of the form $G = {AB}$ where $A$ and $B$ are $\mathfrak{F}$ -subnormal $F$ -subgroups. Find all superradical local formations. L. A. Shemetkov
\end{openproblem}

\plainproblemsection

\begin{openproblem}
	14.99. b) A formation $\mathfrak{F}$ of finite groups is called superradical if it is ${S}_{n}$ -closed and contains every finite group of the form $G = {AB}$ where $A$ and $B$ are $\mathfrak{F}$ -subnormal $F$ - subgroups. Prove that every $S$ -closed superradical formation is a solubly saturated formation. L. A. Shemetkov
	
	b) A counterexample is constructed (S. Yi, S. F. Kamornikov, Siberian Math. J., 57, no. 2 (2016), 260-264).
\end{openproblem}

\plainproblemsection

\begin{openproblem}
	14.100. Is it true that in a Shunkov group (i. e. conjugately biprimitively finite group, see 6.59) having infinitely many elements of finite order every element of prime order is contained in some infinite locally finite subgroup? This is true under the additional condition that any two conjugates of this element generate a soluble subgroup (V. P. Shunkov, ${M}_{p}$ -groups, Moscow, Nauka,1990 (Russian)). A. K. Shlëpkin
\end{openproblem}

\plainproblemsection

\begin{openproblem}
	14.101. A group $G$ is saturated with groups from a class $X$ if every finite subgroup $K \leq  G$ is contained in a subgroup $L \leq  G$ isomorphic to some group from $X$ . Is it true that a periodic group saturated with finite simple groups of Lie type of uniformly bounded ranks is itself a simple group of Lie type of finite rank? A. K. Shlëpkin
\end{openproblem}

\plainproblemsection

\begin{openproblem}
	14.102. (V.Lin). Let ${B}_{n}$ be the braid group on $n$ strings, and let $n > 4$ .
	
	a) Does ${B}_{n}$ have any non-trivial non-injective endomorphisms?
	
	*b) Is it true that every non-trivial endomorphism of the derived subgroup $\left\lbrack  {{B}_{n},{B}_{n}}\right\rbrack$ is an automorphism? V. Shpilrain
	\par\smallskip\hrule\smallskip
	*b) Yes, it is true for $n \geq  7$ (K. Kordek, D. Margalit, Bull. London Math. Soc.,54, no. 1 (2022), 95-111).
\end{openproblem}

\plainproblemsection

\begin{openproblem}
	14.102. c) (V. Lin). Let ${B}_{n}$ be the braid group on $n$ strings, and let $n > 4$ . Does ${B}_{n}$ have proper non-abelian torsion-free factor-groups?
	
	Comment of 2001: S.P. Humphries, (Int. J. Algebra Comput., 11, no. 3 (2001), 363-373) has constructed a representation of ${B}_{n}$ which is shown to provide torsion-free non-abelian factor groups of ${B}_{n}$ as well as of $\left\lbrack  {{B}_{n},{B}_{n}}\right\rbrack$ for $n < 7$ . It is likely that the same representation should work for other values of $n$ as well. $\;V$ Yes, it does; moreover, every braid group ${B}_{n}$ is residually torsion-free nilpotent-by-finite (P. Linnell, T. Schick, J. Amer. Math. Soc., 20, no. 4 (2007), 1003-1051).
\end{openproblem}

\plainproblemsection

\begin{openproblem}
	14.103. Let $H$ be a proper subgroup of a group $G$ and let elements $a, b \in  H$ have distinct prime orders $p, q$ . Suppose that, for every $g \in  G \smallsetminus  H$ , the subgroup $\left\langle  {a,{b}^{g}}\right\rangle$ is a finite Frobenius group with complement of order ${pq}$ . Does the subgroup generated by the union of the kernels of all Frobenius subgroups of $G$ with complement $\langle a\rangle$ intersect $\langle a\rangle$ trivially? The case where all groups $\left\langle  {a,{b}^{g}}\right\rangle  , g \in  G$ , are finite is of special interest. V. P. Shunkov
	
	No, in general case not necessarily. As a counter-example one can take $G =$ $\left\langle  {x, y, z, a, b \mid  {x}^{7} = {y}^{7} = \left\lbrack  {x, y, y}\right\rbrack   = \left\lbrack  {x, y, x}\right\rbrack   = {a}^{3} = {b}^{2} = \left\lbrack  {a, b}\right\rbrack   = 1,\;{x}^{a} = {x}^{2},\;{y}^{a} = }\right\rangle$ ${y}^{2},{x}^{b} = {x}^{-1},{y}^{b} = {y}^{-1}\rangle$ with $H = \langle a, b\rangle$ . Analogous examples exist also for $p = 2$ and every odd prime $q$ (A. I. Sozutov, Letter of 2002).
\end{openproblem}

\plainproblemsection

\begin{openproblem}
	14.104. An infinite group $G$ is called a monster of the first kind if it has elements of order $> 2$ and for any such an element $a$ and for any proper subgroup $H$ of $G$ , there is an element $g$ in $G \smallsetminus  H$ , such that $\left\langle  {a,{a}^{g}}\right\rangle   = G$ . A. Yu. Ol’shanskiĭ showed that there are continuously many monsters of the first kind (see Archive, 6.63). Does there exist, for any such a monster, a torsion-free group which is a central extension of a cyclic group by the given monster? V. P. Shunkov
	
	No, not always, since for such an extension to exist it is necessary that every finite subgroup of the monster is cyclic, but this is not always true (A. I. Sozutov, Letter of
	
	November, 20, 2001).
\end{openproblem}

\plainproblemsection

\begin{openproblem}
	15.1. (P. Longobardi, M. Maj, A. H. Rhemtulla). Let $w = w\left( {{x}_{1},\ldots ,{x}_{n}}\right)$ be a group word in $n$ variables ${x}_{1},\ldots ,{x}_{n}$ , and $V\left( w\right)$ the variety of groups defined by the law $w = 1$ . Let $V\left( {w}^{ * }\right)$ (respectively, $V\left( {w}^{\# }\right)$ ) be the class of all groups $G$ in which for every $n$ infinite subsets ${S}_{1},\ldots ,{S}_{n}$ there exist ${s}_{i} \in  {S}_{i}$ such that $w\left( {{s}_{1},\ldots ,{s}_{n}}\right)  = 1$ (respectively, $\left. {\left\langle  {{s}_{1},\ldots ,{s}_{n}}\right\rangle   \in  V\left( w\right) }\right)$ .
	
	a) Is there some word $w$ and an infinite group $G$ such that $G \in  V\left( {w}^{\# }\right)$ but $G \notin  V\left( w\right)$ ?
	
	b) Is there some word $w$ and an infinite group $G$ such that $G \in  V\left( {w}^{ * }\right)$ but $G \notin  V\left( {w}^{\# }\right)$ ?
	
	The answer to both of these questions is likely to be "yes". It is known that in a) $w$ cannot be any of several words such as ${x}_{1}^{n},\left\lbrack  {{x}_{1},\ldots ,{x}_{n}}\right\rbrack  ,{\left\lbrack  {x}_{1},{x}_{2}\right\rbrack  }^{2},{\left( {x}_{1}{x}_{2}\right) }^{3}{x}_{2}^{-3}{x}_{1}^{-3}$ , and ${x}_{1}^{{a}_{1}}\cdots {x}_{n}^{{a}_{n}}$ for any non-zero integers ${a}_{1},\ldots ,{a}_{n}$ . A. Abdollahi
\end{openproblem}

\plainproblemsection

\begin{openproblem}
	15.2. By a theorem of W. Burnside, if $\chi  \in  \operatorname{Irr}\left( G\right)$ and $\chi \left( 1\right)  > 1$ , then there exists $x \in  G$ such that $\chi \left( x\right)  = 0$ , that is, only the linear characters are "nonvanishing". It is interesting to consider the dual notion of nonvanishing elements of a finite group $G$ , that is, the elements $x \in  G$ such that $\chi \left( x\right)  \neq  0$ for all $\chi  \in  \operatorname{Irr}\left( G\right)$ .
	
	a) It is proved (I. M. Isaacs, G. Navarro, T. R. Wolf, J. Algebra, 222, no. 2 (1999), 413-423) that if $G$ is solvable and $x \in  G$ is a nonvanishing element of odd order, then $x$ belongs to the Fitting subgroup $\mathbf{F}\left( G\right)$ . Is this true for elements of even order too? (M. Miyamoto showed in 2008 that every nontrivial abelian normal subgroup of a finite group contains a nonvanishing element.)
	
	b) Which nonabelian simple groups have nonidentity nonvanishing elements? (For example, ${\mathbb{A}}_{7}$ has.) I. M. Isaacs
\end{openproblem}

\plainproblemsection

\begin{openproblem}
	15.3. Let $\alpha$ and $\beta$ be faithful non-linear irreducible characters of a finite group $G$ . There are non-solvable groups $G$ giving examples when the product ${\alpha \beta }$ is again an irreducible character (for some of such $\alpha ,\beta$ ). One example is $G = S{L}_{2}\left( 5\right)$ with two irreducible characters of degree 2. In (I. Zisser, Israel J. Math., 84, no. 1-2 (1993), 147-151) it is proved that such an example exists in an alternating group ${\mathbb{A}}_{n}$ if and only if $n$ is a square exceeding 4 . But do solvable examples exist? Evidence (but no proof) that they do not is given in (I. M. Isaacs, J. Algebra, 223, no. 2 (2000), 630-646). I. M. Isaacs
\end{openproblem}

\plainproblemsection

\begin{openproblem}
	15.4. Is it true that large growth implies non-amenability? More precisely, consider a number $\epsilon  > 0$ , an integer $k \geq  2$ , a group $\Gamma$ generated by a set $S$ of $k$ elements, and the corresponding exponential growth rate $\omega \left( {\Gamma , S}\right)$ defined as in 14.7. For $\epsilon$ small enough, does the inequality $\omega \left( {\Gamma , S}\right)  \geq  {2k} - 1 - \epsilon$ imply that $\Gamma$ is non-amenable?
	
	If $\omega \left( {\Gamma , S}\right)  = {2k} - 1$ , it is easy to show that $\Gamma$ is free on $S$ , and in particular non-amenable; see Section 2 in (R. I. Grigorchuk, P. de la Harpe, J. Dynam. Control Syst., 3, no. 1 (1997), 51-89). P. de la Harpe
	
	No, it does not. Counterexamples can be found even in the classes of abelian-by-nilpotent and metabelian-by-finite groups. (G. N. Arzhantseva, V. S. Guba, L. Guyot, J. Group Theory, 8 (2005), 389-394).
\end{openproblem}

\plainproblemsection

\begin{openproblem}
	15.5. (Well-known problem). Does there exist an infinite finitely generated group which is simple and amenable? P. de la Harpe
	
	Yes, it does (K. Juschenko, N. Monod, Ann. Math., 178, no. 2 (2013), 775-787). Another example was later constructed in (V. Nekrashevych, Ann. Math., 187, no. 3 (2018), 667-719).
\end{openproblem}

\plainproblemsection

\begin{openproblem}
	15.6. (Well-known problem). Is it true that Golod $p$ -groups are non-amenable?
	
	These infinite finitely generated torsion groups are defined in (E. S. Golod, Amer. Math. Soc. Transl. (2), 48 (1965), 103-106).
	
	Yes, moreover, such groups have infinite quotients with Kazhdan's property (T) and are uniformly non-amenable (M. Ershov, A. Jaikin-Zapirain, Proc. London Math. Soc. (3), 102, no. 4 (2011), 599-636).
\end{openproblem}

\plainproblemsection

\begin{openproblem}
	15.7. (Well-known problem). Is it true that the reduced ${C}^{ * }$ -algebra of a countable group without amenable normal subgroups distinct from $\{ 1\}$ is always a simple ${C}^{ * }$ - algebra with unique trace? See (M. B. Bekka, P. de la Harpe, Expos. Math., 18, no. 3 (2000), 215-230).
	
	Comment of 2009: This was proved for linear groups (T. Poznansky, https:// arxiv.org/pdf/0812.2486.pdf). P. de la Harpe
	
	Yes, it is true (E. Breuillard, M. Kalantar, M. Kennedy, N. Ozawa, Publ. Math. IHES, 126, no. 1 (2017),35-71). If $\Gamma$ is assumed to be linear, then ${C}_{\text{red }}^{ * }\left( \Gamma \right)$ is simple if and only if it has a unique trace (Theorem 1.6 ibid., and also Poznansky's 2009 preprint). However, there are infinite countable groups $\Gamma$ having the following two properties: the only amenable normal subgroup is $\{ 1\}$ , and ${C}_{\text{red }}^{ * }\left( \Gamma \right)$ is not simple, as shown in (A. Le Boudec, Invent. Math., 209, no. 1 (2017), 159-174).
\end{openproblem}

\plainproblemsection

\begin{openproblem}
	15.8. a) (S.M.Ulam). Consider the usual compact group ${SO}\left( 3\right)$ of all rotations of a 3-dimensional Euclidean space, and let $G$ denote this group viewed as a discrete group. Can $G$ act non-trivially on a countable set?
	
	Yes, it can (S. Thomas, J. Group Theory, 2 (1999), 401-434); another proof is in (Yu. L. Ershov, V. A. Churkin, Dokl. Math., 70, no. 3 (2004), 896-898).
\end{openproblem}

\plainproblemsection

\begin{openproblem}
	15.8. b) Let $G$ be any Lie group (indeed, any separable continuous group) made discrete: $\operatorname{can}G$ act faithfully on a countable set?
	
	Editors' comment of 2021: an affirmative answer was obtained for the nilpotent case (N. Monod, J. Group Theory, 25, no. 5 (2022), 851-865). P. de la Harpe
\end{openproblem}

\plainproblemsection

\begin{openproblem}
	15.9. An automorphism $\varphi$ of the free group ${F}_{n}$ on the free generators ${x}_{1},{x}_{2},\ldots ,{x}_{n}$ is called conjugating if ${x}_{i}^{\varphi } = {t}_{i}^{-1}{x}_{\pi \left( i\right) }{t}_{i}, i = 1,2,\ldots , n$ , for some permutation $\pi  \in  {\mathbb{S}}_{n}$ and some elements ${t}_{i} \in  {F}_{n}$ . The set of conjugating automorphisms fixing the product ${x}_{1}{x}_{2}\cdots {x}_{n}$ forms the braid group ${B}_{n}$ . The group ${B}_{n}$ is linear for any $n \geq  2$ , while the group of all automorphisms Aut ${F}_{n}$ is not linear for $n \geq  3$ . Is the group of all conjugating automorphisms linear for $n \geq  3$ ?
\end{openproblem}

\plainproblemsection

\begin{openproblem}
	15.10. (Yu. I. Merzlyakov). Is the group of all automorphisms of the free group ${F}_{n}$ that act trivially on ${F}_{n}/\left\lbrack  {{F}_{n},{F}_{n}}\right\rbrack$ linear for $n \geq  3$ ? $\;V.G$ . Bardakov No, it is not: for $n \geq  5$ (A. Pettet, Cohomology of some subgroups of the automorphism group of a free group, Ph.D. Thesis,2006), and for $n \geq  3$ (V.G.Bardakov, R. Mikhailov, Commun. Algebra, 36, no. 4 (2008), 1489-1499).
\end{openproblem}

\plainproblemsection

\begin{openproblem}
	15.11. (M. Morigi). An automorphism of a group is called a power automorphism if it leaves every subgroup invariant. Is every finite abelian $p$ -group the group of all power automorphisms of some group? V. G. Bardakov
\end{openproblem}

\plainproblemsection

\begin{openproblem}
	15.12. Let $G$ be a group acting faithfully and level-transitively by automorphisms on a rooted tree $\mathcal{T}$ . For a vertex $v$ of $\mathcal{T}$ , the rigid vertex stabilizer at $v$ consists of those elements of $G$ whose support in $\mathcal{T}$ lies entirely in the subtree ${\mathcal{T}}_{v}$ rooted at $v$ . For a non-negative integer $n$ , the $n$ -th rigid level stabilizer is the subgroup of $G$ generated by all rigid vertex stabilizers corresponding to the vertices at the level $n$ of the tree $\mathcal{T}$ . The group $G$ is a branch group if all rigid level stabilizers have finite index in $G$ . For motivation, examples and known results see (R. I. Grigorchuk, in: New horizons in pro-p groups, Birkhäuser, Boston, 2000, 121-179).
	
	Do there exist branch groups with Kazhdan’s $T$ -property? (See A14.34.)
	
	L. Bartholdi, R. I. Grigorchuk, Z. Sunik
\end{openproblem}

\plainproblemsection

\begin{openproblem}
	15.13. Do there exist finitely presented branch groups?
	
	L. Bartholdi, R. I. Grigorchuk, Z. Sunik
\end{openproblem}

\plainproblemsection

\begin{openproblem}
	15.14. Do there exist finitely generated branch groups (see 15.12)
	
	a) that are non-amenable?
	
	c) that contain ${F}_{2}$ ?
	
	d) that have exponential growth? L. Bartholdi, R. I. Grigorchuk, Z. Šunik a), c), d) Such groups do exist (S. Sidki, J. S. Wilson, Arch. Math., 80, no. 5 (2003), 458-463).
\end{openproblem}

\plainproblemsection

\begin{openproblem}
	15.14. b) Do there exist finitely generated branch groups that are non-amenable and do not contain the free group ${F}_{2}$ on two generators?
	
	L. Bartholdi, R. I. Grigorchuk, Z. Sunik
\end{openproblem}

\plainproblemsection

\begin{openproblem}
	*15.15. Is every maximal subgroup of a finitely generated branch group necessarily of finite index? L. Bartholdi, R. I. Grigorchuk, Z. Šunik
	
	*No, not necessarily (I. V. Bondarenko, Arch. Math. (Basel), 95, no. 4 (2010), 301-308).
\end{openproblem}

\plainproblemsection

\begin{openproblem}
	15.16. Do there exist groups whose rate of growth is ${e}^{\sqrt{n}}$ a) in the class of finitely generated branch groups? b) in the whole class of finitely generated groups? This question is related to 9.9 . L. Bartholdi, R. I. Grigorchuk, Z. Sunik
\end{openproblem}

\plainproblemsection

\begin{openproblem}
	15.17. An infinite group is just infinite if all of its proper quotients are finite. Is every finitely generated just infinite group of intermediate growth necessarily a branch group? L. Bartholdi, R. I. Grigorchuk, Z. Šunik
	
	No, not every; there exist simple finitely generated groups of intermediate growth (V. Nekrashevych, Ann. Math., 87, no. 3 (2018), 667-719).
\end{openproblem}

\plainproblemsection

\begin{openproblem}
	*15.18. A group is hereditarily just infinite if it is residually finite and all of its non-trivial normal subgroups are just infinite.
	
	a) Do there exist finitely generated hereditarily just infinite torsion groups? (It is believed there are none.)
	
	b) Is every finitely generated hereditarily just infinite group necessarily linear?
	
	A positive answer to the question b) would imply a negative answer to a). L. Bartholdi, R. I. Grigorchuk, Z. Sunik
	\par\smallskip\hrule\smallskip
	*a) Yes, such groups exist (M. Ershov, A. Jaikin-Zapirain, J. Reine Angew. Math., 677 (2013), 71-134).
	
	*b) No, not every (M. Ershov, A. Jaikin-Zapirain, J. Reine Angew. Math., 677 (2013), 71-134).
	\par\smallskip\hrule\smallskip
\end{openproblem}

\plainproblemsection

\begin{openproblem}
	15.19. Let $p$ be a prime, and ${\mathcal{F}}_{p}$ the class of finitely generated groups acting faithfully on a $p$ -regular rooted tree by finite automata. Any group in ${\mathcal{F}}_{p}$ is residually- $p$ (residually in the class of finite $p$ -groups) and has word problem that is solvable in (at worst) exponential time. There exist therefore groups that are residually- $p$ , have a solvable word problem, and do not belong to ${\mathcal{F}}_{p}$ ; though no concrete example is known. For instance:
	
	a) Is it true that some (or even all) the groups given in (R. I. Grigorchuk, Math. USSR-Sb., $\mathbf{{54}}\left( {1986}\right) ,{185} - {205}$ ) do not belong to ${\mathcal{F}}_{p}$ when the sequence $\omega$ is computable, but not periodic?
	
	*b) Does $\mathbb{Z}\wr \left( {\mathbb{Z}\wr \mathbb{Z}}\right)$ belong to ${\mathcal{F}}_{2}$ ? (Here the wreath products are restricted.)
	
	See (A. M. Brunner, S. Sidki, J. Algebra, 257 (2002), 51-64) and (R. I. Grigorchuk, V. V. Nekrashevich, V. I. Sushchanskii; in: Dynamical systems, automata, and infinite groups. Proc. Steklov Inst. Math., 231 (2000), 128-203). L. Bartholdi, S. Sidki
	
	*b) Yes, it does (A. C. Dantas, J. R. Oliveira, T. M. G. Santos, Preprint, 2024, https://arxiv.org/pdf/2405.16678).
\end{openproblem}

\plainproblemsection

\begin{openproblem}
	15.20. (B. Hartley). An infinite transitive permutation group is said to be barely transitive if each of its proper subgroups has only finite orbits. Can a locally finite barely transitive group coincide with its derived subgroup?
	
	Note that there are no simple locally finite barely transitive groups (B. Hartley, M. Kuzucuoğlu, Proc. Edinburgh Math. Soc., 40 (1997), 483-490), any locally finite barely transitive group is a $p$ -group for some prime $p$ , and if the stabilizer of a point in a locally finite barely transitive group $G$ is soluble of derived length $d$ , then $G$ is soluble of derived length bounded by a function of $d$ (V.V. Belyaev, M. Kuzucuoğlu, Algebra and Logic, 42 (2003), 147-152). V. V. Belyaev, M. Kuzucuoğlu
\end{openproblem}

\plainproblemsection

\begin{openproblem}
	15.21. (B. Hartley). Do there exist torsion-free barely transitive groups?
	
	V. V. Belyaev
\end{openproblem}

\plainproblemsection

\begin{openproblem}
	15.22. A permutation group is said to be finitary if each of its elements moves only finitely many points. Do there exist finitary barely transitive groups? $V, V$ Belvagu
\end{openproblem}

\plainproblemsection

\begin{openproblem}
	15.23. A transitive permutation group is said to be totally imprimitive if every finite set of points is contained in some finite block of the group. Do there exist totally imprimitive barely transitive groups that are not locally finite? V. V. Belyaev
\end{openproblem}

\plainproblemsection

\begin{openproblem}
	15.24. Suppose that a finite $p$ -group $G$ has a subgroup of exponent $p$ and order ${p}^{n}$ . Is it true that if $p$ is sufficiently large relative to $n$ , then $G$ contains a normal subgroup of exponent $p$ and order ${p}^{n}$ ?
	
	J. L. Alperin and G. Glauberman (J. Algebra, 203, no. 2 (1998), 533-566) proved that if a finite $p$ -group contains an elementary abelian subgroup of order ${p}^{n}$ , then it contains a normal elementary abelian subgroup of the same order provided $p >$ ${4n} - 7$ , and the analogue for arbitrary abelian subgroups is proved in (G. Glauberman, J. Algebra, 272 (2004), 128-153). Ya. G. Berkovich
	
	Yes, it is true if $p > n$ . By induction, $G$ contains a subgroup $H$ of exponent $p$ and order ${p}^{n}$ which is normal in a maximal subgroup $M$ of $G$ . Then $H \leq  {\zeta }_{p - 1}\left( M\right)$ . The elements of order $p$ of ${\zeta }_{p - 1}\left( M\right)$ constitute a normal subgroup of $G$ , which contains $H$ . (A. Mann, Letter of 1 October 2002.)
\end{openproblem}

\plainproblemsection

\begin{openproblem}
	15.25. A finite group $G$ is said to be rational if every irreducible character of $G$ takes only rational values. Are the Sylow 2-subgroups of the symmetric groups ${\mathbb{S}}_{{2}^{n}}$ rational? Ya. G. Berkovich
	
	Yes, they are. A Sylow 2-subgroup ${T}_{n}$ of ${\mathbb{S}}_{{2}^{n}}$ is the wreath product of the one for ${\mathbb{S}}_{{2}^{n - 1}}$ with the group $C$ of order 2 . If ${T}_{n}$ is rational, then the wreath product of ${T}_{n}$ with $C$ is also rational by Corollary 70 in (D. Kletzing Structure and representations of $Q$ -groups (Lecture Notes in Math.,1084), Springer, Berlin,1984). The result follows by induction. (A. Mann, Letter of 1 October 2002).
\end{openproblem}

\plainproblemsection

\begin{openproblem}
	15.26. A partition of a group is a representation of it as a set-theoretic union of some of its proper subgroups (components) that intersect pairwise trivially. Is it true that every nontrivial partition of a finite $p$ -group has an abelian component?
	
	Ya. G. Berkovich
\end{openproblem}

\plainproblemsection

\begin{openproblem}
	15.27. Is it possible that Aut $G \cong$ Aut $H$ for a finite $p$ -group $G$ of order $> 2$ and a proper subgroup $H < G$ ? Ya. G. Berkovich
	
	Yes, it is possible (T. Li, Arch. Math., $\mathbf{{92}}\left( {2009}\right) ,{287} - {290}$ ) for $\left| G\right|  = 2\left| H\right|  = {32}$ .
\end{openproblem}

\plainproblemsection

\begin{openproblem}
	15.28. Suppose that a finite $p$ -group $G$ is the product of two subgroups: $G = {AB}$ .
	
	a) Is the exponent of $G$ bounded in terms of the exponents of $A$ and $B$ ?
	
	b) Is the exponent of $G$ bounded if $A$ and $B$ are groups of exponent $p$ ?
	
	Ya. G. Berkovich
\end{openproblem}

\plainproblemsection

\begin{openproblem}
	15.29. (A. Mann). The dihedral group of order 8 is isomorphic to its own automorphism group. Are there other non-trivial finite $p$ -groups with this property?
	
	Ya. G. Berkovich
\end{openproblem}

\plainproblemsection

\begin{openproblem}
	15.30. Is it true that every finite abelian $p$ -group is isomorphic to the Schur multiplier of some nonabelian finite $p$ -group? Ya. G. Berkovich
\end{openproblem}

\plainproblemsection

\begin{openproblem}
	15.31. M. R. Vaughan-Lee and J. Wiegold (Proc. R. Soc. Edinburgh Sect. A, 95 (1983),215-221) proved that if a finite $p$ -group $G$ is generated by elements of breadth $\leq  n$ (that is, having at most ${p}^{n}$ conjugates), then $G$ is nilpotent of class $\leq  {n}^{2} + 1$ ; the bound for the class was later improved by A. Mann (J. Group Theory, 4, no. 3 (2001),241-246) to $\leq  {n}^{2} - n + 1$ . Is there a linear bound for the class of $G$ in terms
	
	of $n$ ? Ya. G. Berkovich 15.32. Does there exist a function $f\left( k\right)$ (possibly depending also on $p$ ) such that if a finite $p$ -group $G$ of order ${p}^{m}$ with $m \geq  f\left( k\right)$ has an automorphism of order ${p}^{m - k}$ , then $G$ possesses a cyclic subgroup of index ${p}^{k}$ ? Ya. G. Berkovich
\end{openproblem}

\plainproblemsection

\begin{openproblem}
	15.33. Suppose that all 2-generator subgroups of a finite 2-group $G$ are metacyclic. Is the derived length of $G$ bounded? This is true for finite $p$ -groups if $p \neq  2$ , see 1.1.8 in (M. Suzuki, Structure of a group and the structure of its lattice of subgroups, Springer, Berlin, 1956). Ya. G. Berkovich
	
	Yes, it is at most 2 (E. Crestani, F. Menegazzo, J. Group Theory, 15, no. 3 (2012), 359-383).
\end{openproblem}

\plainproblemsection

\begin{openproblem}
	15.34. Is any free product of linearly ordered groups with an amalgamated subgroup right-orderable? V. V. Bludov
	
	Yes, it is (V. V. Bludov, A. M. W. Glass, Math. Proc. Cambridge Philos. Soc., 146 (2009), 591-601).
\end{openproblem}

\plainproblemsection

\begin{openproblem}
	15.35. Let $F$ be the free group of finite rank $r$ with basis $\left\{  {{x}_{1},\ldots ,{x}_{r}}\right\}$ . Is it true that there exists a number $C = C\left( r\right)$ such that any reduced word of length $n > 1$ in the ${x}_{i}$ lies outside some subgroup of $F$ of index at most $C\log n$ ? O. V. Bogopol’skii No, it is not true (K. Bou-Rabee, D. B. McReynolds, Bull. London Math. Soc., 43, no. 6 (2011), 1059-1068).
\end{openproblem}

\plainproblemsection

\begin{openproblem}
	15.36. For a class $M$ of groups let $L\left( M\right)$ be the class of groups $G$ such that the normal subgroup $\left\langle  {a}^{G}\right\rangle$ generated by an element $a$ belongs to $M$ for any $a \in  G$ . Is it true that the class $L\left( M\right)$ is finitely axiomatizable if $M$ is the quasivariety generated
	
	by a finite group? A. I. Budkin
\end{openproblem}

\plainproblemsection

\begin{openproblem}
	15.37. Let $G$ be a group satisfying the minimum condition for centralizers. Suppose that $X$ is a normal subset of $G$ such that $\left\lbrack  {x, y,\ldots , y}\right\rbrack   = 1$ for any $x, y \in  X$ with $y$ repeated $f\left( {x, y}\right)$ times. Does $X$ generate a locally nilpotent (whence hypercentral) subgroup?
	
	If the numbers $f\left( {x, y}\right)$ can be bounded, the answer is affirmative (F. O. Wagner, J. Algebra, 217, no. 2 (1999), 448-460). F. O. Wagner
\end{openproblem}

\plainproblemsection

\begin{openproblem}
	15.38. Does there exist a non-local hereditary composition formation $\mathfrak{F}$ of finite groups such that the set of all $\mathfrak{F}$ -subnormal subgroups is a sublattice of the subgroup lattice in any finite group? A. F. Vasil’ev, S. F. Kamornikov
	
	No, it does not (S. F. Kamornikov, Dokl. Math., 81, no. 1 (2010), 97-100).
\end{openproblem}

\plainproblemsection

\begin{openproblem}
	15.39. Axiomatizing the basic properties of subnormal subgroups, we say that a functor $\tau$ associating with every finite group $G$ some non-empty set $\tau \left( G\right)$ of its subgroups is an ${ETP}$ -functor if
	
	1) $\tau {\left( A\right) }^{\varphi } \subseteq  \tau \left( B\right)$ and $\tau {\left( B\right) }^{{\varphi }^{-1}} \subseteq  \tau \left( A\right)$ for any epimorphism $\varphi  : A \mapsto  B$ , as well as $\{ H \cap  R \mid  R \in  \tau \left( G\right) \}  \subseteq  \tau \left( H\right)$ for any subgroup $H \leq  G$ ;
	
	2) $\tau \left( H\right)  \subseteq  \tau \left( G\right)$ for any subgroup $H \in  \tau \left( G\right)$ ;
	
	3) $\tau \left( G\right)$ is a sublattice of the lattice of all subgroups of $G$ .
	
	Let $\tau$ be an ${ETP}$ -functor. Does there exist a hereditary formation $\mathfrak{F}$ such that $\tau \left( G\right)$ coincides with the set of all $\mathfrak{F}$ -subnormal subgroups in any finite group $G$ ? This is true for finite soluble groups (A. F. Vasil’ev, S. F. Kamornikov, Siberian Math. J., 42, no. 1 (2001), 25-33). A. F. Vasil’ev, S. F. Kamornikov
	
	Not always (S. F. Kamornikov, Math. Notes, 89, no. 3-4 (2011), 340-348).
\end{openproblem}

\plainproblemsection

\begin{openproblem}
	15.40. Let $N$ be a nilpotent subgroup of a finite simple group $G$ . Is it true that there exists a subgroup ${N}_{1}$ conjugate to $N$ such that $N \cap  {N}_{1} = 1$ ?
	
	The answer is known to be affirmative if $N$ is a $p$ -group. Comment of 2013: an affirmative answer for alternating groups is obtained in (R. K. Kurmazov, Siberian Math. J., 54, no. 1 (2013), 73-77). E. P. Vdovin
\end{openproblem}

\plainproblemsection

\begin{openproblem}
	15.41. Let $R\left( {m, p}\right)$ denote the largest finite $m$ -generator group of prime exponent $p$ . a) Can the nilpotency class of $R\left( {m, p}\right)$ be bounded by a polynomial in $m$ ? (This is true for $p = 2,3,5,7$ .)
	
	b) Can the nilpotency class of $R\left( {m, p}\right)$ be bounded by a linear function in $m$ ? (This is true for $p = 2,3,5$ .)
	
	c) In particular, can the nilpotency class of $R\left( {m,7}\right)$ be bounded by a linear function in $m$ ?
	
	My guess is "no" to the first two questions for general $p$ , but "yes" to the third. By contrast, a beautiful and simple argument of Mike Newman shows that if $m \geq  2$ and $k \geq  2$ ( $k \geq  3$ for $p = 2$ ), then the order of $R\left( {m,{p}^{k}}\right)$ is at least ${p}^{p}$ , with $p$ appearing $k$ times in the tower; see (M. Vaughan-Lee, E. I. Zelmanov, J. Austral. Math. Soc. (A), 67, no. 2 (1999), 261-271). M. R. Vaughan-Lee
\end{openproblem}

\plainproblemsection

\begin{openproblem}
	15.42. Is it true that the group algebra $k\left\lbrack  F\right\rbrack$ of R. Thompson’s group $F$ (see 12.20) over a field $k$ satisfies the Ore condition, that is, for any $a, b \in  k\left\lbrack  F\right\rbrack$ there exist $u, v \in  k\left\lbrack  F\right\rbrack$ such that ${au} = {bv}$ and either $u$ or $v$ is nonzero? If the answer is negative, then $F$ is not amenable. V. S. Guba
\end{openproblem}

\plainproblemsection

\begin{openproblem}
	15.43. Let $G$ be a finite group of order $n$ .
	
	a) Is it true that $\left| {\operatorname{Aut}G}\right|  \geq  \varphi \left( n\right)$ where $\varphi$ is Euler’s function?
	
	b) Is it true that $G$ is cyclic if $\left| {\operatorname{Aut}G}\right|  = \varphi \left( n\right)$ ? M. Deaconescu
	
	Both questions have negative answers; moreover, $\left| {\operatorname{Aut}G}\right| /\varphi \left( \left| G\right| \right)$ can be made arbitrarily small (J. N. Bray, R. A. Wilson, Bull. London Math. Soc., 37, no. 3 (2005), 381-385).
\end{openproblem}

\plainproblemsection

\begin{openproblem}
	15.44. a) Let $G$ be a reductive group over an algebraically closed field $K$ of arbitrary characteristic. Let $X$ be an affine $G$ -variety such that, for a fixed Borel subgroup $B \leq  G$ , the coordinate algebra $K\left\lbrack  X\right\rbrack$ as a $G$ -module is the union of an ascending chain of submodules each of whose factors is an induced module ${\operatorname{Ind}}_{B}^{G}V$ of some one-dimensional $B$ -module $V$ . (See S. Donkin, Rational representations of algebraic groups. Tensor products and flltration (Lect. Notes Math., 1140), Springer, Berlin, 1985). Suppose in addition that $K\left\lbrack  X\right\rbrack$ is a Cohen-Macaulay ring, that is, a free module over the subalgebra generated by any homogeneous system of parameters. Is then the ring of invariants $K{\left\lbrack  X\right\rbrack  }^{G}$ Cohen-Macaulay? Comment of 2005: This is proved in the case where $X$ is a rational $G$ -module (M. Hashimoto, Math. Z.,236 (2001), 605-623). A. N. Zubkov
\end{openproblem}

\plainproblemsection

\begin{openproblem}
	15.44. b) Is the ring of invariants $K{\left\lbrack  M{\left( n\right) }^{m}\right\rbrack  }^{{GL}\left( n\right) }$ Cohen-Macaulay in all characteristics? (Here $M{\left( n\right) }^{m}$ is the direct sum of $m$ copies of the space of $n \times  n$ matrices.) Yes, it is (M. Hashimoto, Math. Z., 236 (2001), 605-623). A. N. Zubkov
\end{openproblem}

\plainproblemsection

\begin{openproblem}
	15.45. We define the class of hierarchically decomposable groups in the following way. First, if $\mathfrak{X}$ is any class of groups, then let ${\mathbf{H}}_{1}\mathfrak{X}$ denote the class of groups which admit an admissible action on a finite-dimensional contractible complex in such a way that every cell stabilizer belongs to $\mathfrak{X}$ . Then the "big" class $\mathbf{H}\mathfrak{X}$ is defined to be the smallest ${\mathbf{H}}_{1}$ -closed class containing $\mathfrak{X}$ .
	
	*a) Let $\mathfrak{F}$ be the class of all finite groups. Find an example of an $\mathbf{H}\mathfrak{F}$ group which is not in ${\mathbf{H}}_{3}\mathfrak{F}\left( { = {\mathbf{H}}_{1}{\mathbf{H}}_{1}{\mathbf{H}}_{1}F}\right)$ .
	
	Such an example is found; moreover, for any countable ordinal $\alpha$ there are groups in ${\mathbf{H}}_{\alpha  + 1}\mathfrak{F}$ that are not in ${\mathbf{H}}_{\alpha }\mathfrak{F}$ (T. Januszkiewicz, P. H. Kropholler, I. J. Leary, Bull. London Math. Soc., 42, No. 5 (2010), 896-904).
	
	b) Prove or disprove that there is an ordinal $\alpha$ such that ${\mathbf{H}}_{\alpha }\mathfrak{F} = \mathbf{H}\mathfrak{F}$ , where ${\mathbf{H}}_{\alpha }$ is the operator on classes of groups defined by transfinite induction in the obvious way starting from ${\mathbf{H}}_{1}$ . Editor’s comment: It is proved that such an ordinal cannot be countable (T. Januszkiewicz, P. H. Kropholler, I. J. Leary, Bull. London Math. Soc., 42, No. 5 (2010), 896-904). P. Kropholler
\end{openproblem}

\plainproblemsection

\begin{openproblem}
	15.46. Can the question 7.28 on conditions for admissibility of an elementary carpet $\mathfrak{A} = \left\{  {{\mathfrak{A}}_{r} \mid  r \in  \Phi }\right\}$ be reduced to Lie rank 1 if $K$ is a field? A carpet $\mathfrak{A}$ of type $\Phi$ of additive subgroups of $K$ is called addmissible if in the Chevalley group over $K$ associated with the root system $\Phi$ the subgroup $\left\langle  {{x}_{r}\left( {\mathfrak{A}}_{r}\right)  \mid  r \in  \Phi }\right\rangle$ intersects ${x}_{r}\left( K\right)$ in ${x}_{r}\left( {\mathfrak{A}}_{r}\right)$ . More precisely, is it true that the carpet $\mathfrak{A}$ is admissible if and only if the subcarpets $\left\{  {{\mathfrak{A}}_{r},{\mathfrak{A}}_{-r}}\right\}  , r \in  \Phi$ , of rank 1 are admissible? The answer is known to be affirmative if the field $K$ is locally finite (V. M. Levchuk, Algebra and Logic,22, no. 5 (1983), 362-371). V. M. Levchuk
\end{openproblem}

\plainproblemsection

\begin{openproblem}
	15.47. Let $M < G \leq  \operatorname{Sym}\left( \Omega \right)$ , where $\Omega$ is finite, be such that $M$ is transitive on $\Omega$ and there is a $G$ -invariant partition $\mathcal{P}$ of $\Omega  \times  \Omega  \smallsetminus  \{ \left( {\alpha ,\alpha }\right)  \mid  \alpha  \in  \Omega \}$ such that $G$ is transitive on the set of parts of $\mathcal{P}$ and $M$ fixes each part of $\mathcal{P}$ setwise. (Here $\mathcal{P}$ can be identified with a decomposition of the complete directed graph with vertex set $\Omega$ into edge-disjoint isomorphic directed graphs.) If $G$ induces a cyclic permutation group on $\mathcal{P}$ , then we showed (Trans. Amer. Math. Soc.,355, no. 2 (2003), 637-653) that the numbers $n = \left| \Omega \right|$ and $k = \left| \mathcal{P}\right|$ are such that the $r$ -part ${n}_{r}$ of $n$ satisfies ${n}_{r} \equiv  1\left( {\;\operatorname{mod}\;k}\right)$ for each prime $r$ . Are there examples with $G$ inducing a non-cyclic permutation group on $\mathcal{P}$ for any $n, k$ not satisfying this congruence condition?
\end{openproblem}

\plainproblemsection

\begin{openproblem}
	15.48. Let $G$ be any non-trivial finite group, and let $X$ be any generating set for $G$ . Is it true that every element of $G$ can be obtained from $X$ using fewer than $2{\log }_{2}\left| G\right|$ multiplications? (When counting the number of mutiplications on a path from the generators to a given element, at each step one can use the elements obtained at previous steps.) C. R. Leedham-Green
\end{openproblem}

\plainproblemsection

\begin{openproblem}
	*15.49. A group $G$ is a unique product group if, for any finite nonempty subsets $X, Y$ of $G$ , there is an element of $G$ which can be written in exactly one way in the form ${xy}$ with $x \in  X$ and $y \in  Y$ . Does there exist a unique product group which is not left-orderable? P. Linnell
	
	*Yes, it exists (N. Dunfield, Appendix B in S. Kionke, J. Raimbault, Doc. Math., 21 (2016), 873-915).
\end{openproblem}

\plainproblemsection

\begin{openproblem}
	15.50. Let $G$ be a group of automorphisms of an abelian group of prime exponent. Suppose that there exists $x \in  G$ such that $x$ is regular of order 3 and the order of $\left\lbrack  {x, g}\right\rbrack$ is finite for every $g \in  G$ . Is it true that $\left\langle  {x}^{G}\right\rangle$ is locally finite?
\end{openproblem}

\plainproblemsection

\begin{openproblem}
	15.51. Suppose that $G$ is a periodic group satisfying the identity ${\left\lbrack  x, y\right\rbrack  }^{5} = 1$ . Is then the derived subgroup $\left\lbrack  {G, G}\right\rbrack$ a 5-group? V. D. Mazurov
\end{openproblem}

\plainproblemsection

\begin{openproblem}
	15.52. (Well-known problem). By a famous theorem of Wielandt the sequence ${G}_{0} =$ $G,{G}_{1},\ldots$ , where ${G}_{i + 1} = \operatorname{Aut}{G}_{i}$ , stabilizes for any finite group $G$ with trivial centre. Does there exist a function $f$ of natural argument such that $\left| {G}_{i}\right|  \leq  f\left( \left| G\right| \right)$ for all $i =$ $0,1,\ldots$ for an arbitrary finite group $G$ and the same kind of sequence? $V.D.{Mazurov}$
\end{openproblem}

\plainproblemsection

\begin{openproblem}
	*15.53. Let $S$ be the set of all prime numbers $p$ for which there exists a finite simple group $G$ and an absolutely irreducible $G$ -module $V$ over a field of characteristic $p$ such that the order of any element in the natural semidirect product ${VG}$ coincides with the order of some element in $G$ . Is $S$ finite or infinite? V. D. Mazurov
	\par\smallskip\hrule\smallskip
	*The set $S$ is finite and consists of 2 and 3 . This follows from the complete list of finite simple groups $G$ with a module $V$ in characteristic $p$ satisfying those properties, which was determined in a number of papers culminating in (M. A. Grechkoseeva, S. V. Skresanov, Sibirsk. Elektron. Matem. Izv., 17 (2020), 585-589). Then the group $G$ must be either $3{D}_{4}\left( 2\right)$ with $p = 2$ (I. B. Gorshkov, Commun. Algebra,47, no. 12 (2019), 5192-5206), or ${U}_{5}\left( 2\right)$ with $p = 3$ (A. A. Buturlakin, Algebra Logic, 47, no. 2 (2008), 91-99, and A. V. Zavarnitsine, Siberian Math. J., 49, no. 2 (2008), 246- 256), or ${U}_{3}\left( q\right)$ , where $q$ is a special Mersenne prime, with $p = 2$ (A.V.Zavarnitsine, Algebra Logic, 45, no. 2 (2006), 106-116).
	\par\smallskip\hrule\smallskip
\end{openproblem}

\plainproblemsection

\begin{openproblem}
	15.54. Suppose that $G$ is a periodic group containing an involution $i$ such that the centralizer ${C}_{G}\left( i\right)$ is a locally cyclic 2-group. Does the set of all elements of odd order in $G$ that are inverted by $i$ form a subgroup? V. D. Mazurov
\end{openproblem}

\plainproblemsection

\begin{openproblem}
	15.55. a) The Monster, $M$ , is a 6-transposition group. Pairs of Fischer transpositions generate ${9M}$ -classes of dihedral groups. The order of the product of a pair is the coefficient of the highest root of affine type ${E}_{8}$ . Similar properties hold for Baby $B$ , and ${F}_{24}^{\prime }$ with respect to ${E}_{7}$ and ${E}_{6}$ when the product is read modulo centres $\left( {2.B,3.{F}_{24}^{\prime }}\right)$ . Explain this. Editors’ Comment of 2005: Some progress was made in (C. H. Lam, H. Yamada, H. Yamauchi, Trans. Amer. Math. Soc., 359, no. 9 (2007), 4107-4123).
	
	b) Note that the Schur multiplier of the sporadic simple groups $M, B,{F}_{24}^{\prime }$ is the fundamental group of type ${E}_{8},{E}_{7},{E}_{6}$ , respectively. Why?
	
	See (J. McKay, in: Finite groups, Santa Cruz Conf. 1979 (Proc. Symp. Pure Math., 37), Amer. Math. Soc., Providence, RI, 1980, 183-186) and (R. E. Borcherds, Doc. Math., J. DMV Extra Vol. ICM Berlin 1998, Vol. I (1998), 607-616). J. McKay
\end{openproblem}

\plainproblemsection

\begin{openproblem}
	15.56. Is there a spin manifold such that the Monster acts on its loop space? Perhaps 24-dimensional? hyper-Kähler, non-compact? See (F. Hirzebruch, T. Berger, R. Jung, Manifolds and modular forms, Vieweg, Braunschweig, 1992). J. McKay
\end{openproblem}

\plainproblemsection

\begin{openproblem}
	15.57. Suppose that $H$ is a subgroup of $S{L}_{2}\left( \mathbb{Q}\right)$ that is dense in the Zariski topology and has no nontrivial finite quotients. Is then $H = S{L}_{2}\left( \mathbb{Q}\right)$ ? J. McKay, J.-P. Serre
\end{openproblem}

\plainproblemsection

\begin{openproblem}
	15.58. Suppose that a free profinite product $G * H$ is a free profinite group of finite rank. Must $G$ and $H$ be free profinite groups?
	
	By (J. Neukirch, Arch. Math., 22, no. 4 (1971), 337-357) this may not be true if the rank of $G * H$ is infinite. O. V. Mel'nikov
\end{openproblem}

\plainproblemsection

\begin{openproblem}
	15.59. Does there exist a profinite group $G$ that is not free but can be represented as a projective limit $G = \mathop{\lim }\limits_{ \leftarrow  }\left( {G/{N}_{\alpha }}\right)$ , where all the $G/{N}_{\alpha }$ are free profinite groups of finite ranks?
	
	The finiteness condition on the ranks of the $G/{N}_{\alpha }$ is essential. Such a group $G$ cannot satisfy the first axiom of countability (O. V. Mel'nikov, Dokl. AN BSSR, 24, no. 11 (1980), 968-970 (Russian)). O. V. Mel'nikov
\end{openproblem}

\plainproblemsection

\begin{openproblem}
	15.60. Is it true that any finitely generated ${p}^{\prime }$ -isolated subgroup of a free group is separable in the class of finite $p$ -groups? It is easy to see that this is true for cyclic subgroups. D. I. Moldavanskiǐ
	
	No, it is not (V. G. Bardakov, Siberian Math. J., 45, no. 3 (2004), 416-419).
\end{openproblem}

\plainproblemsection

\begin{openproblem}
	15.61. Is it true that ${l}_{\pi }^{n}\left( G\right)  \leq  n\left( {G}_{\pi }\right)  - 1 + \mathop{\max }\limits_{{p \in  \pi }}{l}_{p}\left( G\right)$ for any $\pi$ -soluble group $G$ ? Here $n\left( {G}_{\pi }\right)$ is the nilpotent length of a Hall $\pi$ -subgroup ${G}_{\pi }$ of the group $G$ and ${l}_{\pi }^{n}\left( G\right)$ is the nilpotent $\pi$ -length of $G$ , that is, the minimum number of $\pi$ -factors in those normal series of $G$ whose factors are either ${\pi }^{\prime }$ -groups, or nilpotent $\pi$ -groups. The answer is known to be affirmative in the case when all proper subgroups of ${G}_{\pi }$ are supersoluble. V. S. Monakhov
\end{openproblem}

\plainproblemsection

\begin{openproblem}
	15.62. Given an ordinary irreducible character $\chi$ of a finite group $G$ write ${p}^{{e}_{p}\left( \chi \right) }$ to denote the $p$ -part of $\chi \left( 1\right)$ and put ${e}_{p}\left( G\right)  = \max \left\{  {{e}_{p}\left( \chi \right)  \mid  \chi  \in  \operatorname{Irr}\left( G\right) }\right\}$ . Suppose that $P$ is a Sylow $p$ -subgroup of a group $G$ . Is it true that ${e}_{p}\left( P\right)$ is bounded above by a function of ${e}_{p}\left( G\right)$ ?
	
	Comment of 2005: An affirmative answer in the case of solvable groups has been given in (A. Moretó, T. R. Wolf, Adv. Math., 184 (2004), 18-36). A. Moreto Yes, it is true (Yong Yang, Guohua Qian, Adv. Math., 328 (2018), 356-366).
\end{openproblem}

\plainproblemsection

\begin{openproblem}
	15.63. a) Let ${F}_{n}$ be the free group of finite rank $n$ on the free generators ${x}_{1},\ldots ,{x}_{n}$ . An element $u \in  {F}_{n}$ is called positive if $u$ belongs to the semigroup generated by the ${x}_{i}$ . An element $u \in  {F}_{n}$ is called potentially positive if $\alpha \left( u\right)$ is positive for some automorphism $\alpha$ of ${F}_{n}$ . Is the property of an element to be potentially positive algorithmically recognizable?
	
	Comment of 2013: In (R. Goldstein, Contemp. Math., Amer. Math. Soc., 421 (2006),157-168) the problem was solved in the affirmative in the special case $n = 2$ .
\end{openproblem}

\plainproblemsection

\begin{openproblem}
	15.63. b) Let ${F}_{n}$ be the free group of finite rank $n$ on the free generators ${x}_{1},\ldots ,{x}_{n}$ . An element $u \in  {F}_{n}$ is called positive if $u$ belongs to the semigroup generated by the ${x}_{i}$ . An element $u \in  {F}_{n}$ is called potentially positive if $\alpha \left( u\right)$ is positive for some automorphism $\alpha$ of ${F}_{n}$ . Finally, $u \in  {F}_{n}$ is called stably potentially positive if it is potentially positive as an element of ${F}_{m}$ for some $m \geq  n$ . Are there stably potentially positive elements that are not potentially positive? A. G. Myasnikov, V. E. Shpilrain No, there are none (A. Clark, R. Goldstein, Commun. Algebra, 33, no. 11 (2005), 4097-4104).
\end{openproblem}

\plainproblemsection

\begin{openproblem}
	15.64. For finite groups $G, X$ define $r\left( {G;X}\right)$ to be the number of inequivalent actions of $G$ on $X$ , that is, the number of equivalence classes of homomorphisms $G \rightarrow$ Aut $X$ , where equivalence is defined by conjugation by an element of Aut $X$ . Now define ${r}_{G}\left( n\right)  \mathrel{\text{:=}} \max \{ r\left( {G;X}\right)  : \left| X\right|  = n\} .$
	
	Is it true that ${r}_{G}\left( n\right)$ may be bounded as a function of $\lambda \left( n\right)$ , the total number (counting multiplicities) of prime factors of $n$ ? Peter M. Neumann
\end{openproblem}

\plainproblemsection

\begin{openproblem}
	15.65. A square matrix is said to be separable if its minimal polynomial has no repeated roots, and cyclic if its minimal and characteristic polynomials are equal. For a matrix group $G$ over a finite field define $s\left( G\right)$ and $c\left( G\right)$ to be the proportion of separable and of cyclic elements respectively in $G$ . For a classical group $\mathrm{X}\left( {d, q}\right)$ of dimension $d$ defined over the field with $q$ elements let $S\left( {\mathrm{X};q}\right)  \mathrel{\text{:=}} \mathop{\lim }\limits_{{d \rightarrow  \infty }}s\left( {\mathrm{\;X}\left( {d, q}\right) }\right)$ and $C\left( {\mathrm{X};q}\right)  \mathrel{\text{:=}} \mathop{\lim }\limits_{{d \rightarrow  \infty }}c\left( {\mathrm{X}\left( {d, q}\right) }\right)$ . Independently G. E. Wall (Bull. Austral. Math. Soc., 60, no. 2 (1999), 253-284) and J. Fulman (J. Group Theory, 2, no. 3 (1999), 251-289) have evaluated $S\left( {\mathrm{{GL}};q}\right)$ and $C\left( {\mathrm{{GL}};q}\right)$ , and have found them to be rational functions of $q$ . Are $S\left( {\mathrm{X};q}\right)$ and $C\left( {\mathrm{X};q}\right)$ rational functions of $q$ also for the unitary, symplectic, and orthogonal groups? Peter M. Neumann
\end{openproblem}

\plainproblemsection

\begin{openproblem}
	15.66. For a class $\mathfrak{X}$ of groups let ${g}_{\mathfrak{X}}\left( n\right)$ be the number of groups of order $n$ in the class $\mathfrak{X}$ (up to isomorphism). Many years ago I formulated the following problem: find good upper bounds for the quotient ${g}_{\mathfrak{V}}\left( n\right) /{g}_{\mathfrak{U}}\left( n\right)$ , where $\mathfrak{V}$ is a variety that is defined by its finite groups and $\mathfrak{U}$ is a subvariety of $\mathfrak{V}$ . (This quotient is not defined for all $n$ but only for those for which there are groups of order $n$ in $\mathfrak{U}$ .) Some progress has been made by G. Venkataraman (Quart. J. Math. Oxford (2), 48, no. 189 (1997), 107-125) when $\mathfrak{V}$ is a variety generated by finite groups all of whose Sylow subgroups are abelian. Conjecture: if $\mathfrak{V}$ is a locally finite variety of $p$ -groups and $\mathfrak{U}$ is a non-abelian subvariety of $\mathfrak{V}$ , then ${g}_{\mathfrak{V}}\left( {p}^{m}\right) /{g}_{\mathfrak{U}}\left( {p}^{m}\right)  < {p}^{O\left( {m}^{2}\right) }$ .
	
	Moreover, this seems a possible way to attack the Sims Conjecture that when we write the number of groups of order ${p}^{m}$ as ${p}^{\frac{2}{27}{m}^{3} + \varepsilon \left( m\right) }$ the error term $\varepsilon \left( m\right)$ is $O\left( {m}^{2}\right)$ .
	
	Peter M. Neumann
\end{openproblem}

\plainproblemsection

\begin{openproblem}
	15.67. Which adjoint Chevalley groups (of normal type) over the integers are generated by three involutions two of which commute?
	
	The groups $S{L}_{n}\left( \mathbb{Z}\right) , n > {13}$ , satisfy this condition (M. C. Tamburini, P. Zucca, J. Algebra,195, no. 2 (1997),650-661). The groups ${\operatorname{PSL}}_{n}\left( \mathbb{Z}\right)$ satisfy it if and only if $n > 4$ (N. Ya. Nuzhin, Vladikavkaz. Mat. Zh.,10, no. 1 (2008), 68-74 (Russian)). The group ${PS}{p}_{4}\left( \mathbb{Z}\right)$ does not satisfy it, which follows from the corresponding fact for ${PS}{p}_{4}\left( 3\right)$ , see Archive,7.30. Ya. N. Nuzhin
\end{openproblem}

\plainproblemsection

\begin{openproblem}
	15.68. Does there exist an infinite finitely generated 2-group (of finite exponent) all of whose proper subgroups are locally finite? A. Yu. Ol'shanskiĭ
\end{openproblem}

\plainproblemsection

\begin{openproblem}
	15.69. Is it true that every hyperbolic group has a free normal subgroup with the factor-group of finite exponent? A. Yu. Ol'shanskiĭ
\end{openproblem}

\plainproblemsection

\begin{openproblem}
	*15.70. Do there exist groups of arbitrarily large cardinality that satisfy the minimum condition for subgroups? A. Yu. Ol'shanskiĭ
	
	*A positive answer is consistent with the ZFC axioms of set theory, and the consistency of a negative answer follows from the consistency of certain large cardinal assumptions (S. M. Corson, S. Shelah, Preprint, 2024, https://arxiv.org/pdf/2408.03201).
\end{openproblem}

\plainproblemsection

\begin{openproblem}
	15.71. (B. Huppert). Let $G$ be a finite solvable group, and let $\rho \left( G\right)$ denote the set of prime divisors of the degrees of irreducible characters of $G$ . Is it true that there always exists an irreducible character of $G$ whose degree is divisible by at least $\left| {\rho \left( G\right) }\right| /2$ different primes? P. P. Pálfy
\end{openproblem}

\plainproblemsection

\begin{openproblem}
	15.72. For a fixed prime $p$ does there exist a sequence of groups ${P}_{n}$ of order ${p}^{n}$ such that the number of conjugacy classes $k\left( {P}_{n}\right)$ satisfies $\mathop{\lim }\limits_{{n \rightarrow  \infty }}\log k\left( {P}_{n}\right) /\sqrt{n} = 0$ ?
	
	Note that J. M. Riedl (J. Algebra, 218 (1999), 190-215) constructed p-groups for which the above limit is $2\log p$ . P. P. Pálfy
\end{openproblem}

\plainproblemsection

\begin{openproblem}
	15.73. Is it true that for every finite lattice $L$ there exists a finite group $G$ and a subgroup $H \leq  G$ such that the interval $\operatorname{Int}\left( {H;G}\right)$ in the subgroup lattice of $G$ is isomorphic to $L$ ? (Probably not.) P. P. Pálfy
\end{openproblem}

\plainproblemsection

\begin{openproblem}
	15.74. For every prime $p$ find a finite $p$ -group of nilpotence class $p$ such that its lattice of normal subgroups cannot be embedded into the subgroup lattice of any abelian group. (Solved for $p = 2$ ,3.) P. P. Pálfy
\end{openproblem}

\plainproblemsection

\begin{openproblem}
	15.75. a) Does there exist a sequence of identities in two variables ${u}_{1} = 1,{u}_{2} = 1,\ldots$ with the following properties: 1) each of these identities implies the next one, and 2) an arbitrary finite group is soluble if and only if it satisfies one of the identities ${u}_{n} = 1?$ B. I. Plotkin
	
	Yes, such sequences exist (T. Bandman, F. Grunewald, G.-M. Greuel, B. Kunyavskii, G. Pfister, E. Plotkin, Compositio Math., 142 (2006), 734-764; J. N. Bray, J. S. Wilson, R. A. Wilson, Bull. Lond. Math. Soc., 37 (2005) 179-186; E. Ribnere, Monatsh. Math., 157 (2009), 387-401).
\end{openproblem}

\plainproblemsection

\begin{openproblem}
	15.75. b) Consider the sequence ${u}_{1} = \left\lbrack  {x, y}\right\rbrack  ,\ldots ,{u}_{n + 1} = \left\lbrack  {\left\lbrack  {{u}_{n}, x}\right\rbrack  ,\left\lbrack  {{u}_{n}, y}\right\rbrack  }\right\rbrack$ . Is it true that an arbitrary finite group is soluble if and only if it satisfies one of these identities ${u}_{n} = 1?$ B. I. Plotkin
\end{openproblem}

\plainproblemsection

\begin{openproblem}
	15.76. a) If $\Theta$ is a variety of groups, then let ${\Theta }^{0}$ denote the category of all free groups of finite rank in $\Theta$ . It is proved (G. Mashevitzky, B. Plotkin, E. Plotkin J. Algebra, 282 (2004),490-512) that if $\Theta$ is the variety of all groups, then every automorphism of the category ${\Theta }^{0}$ is an inner one. The same is true if $\Theta$ is the variety of all abelian groups. Is this true for the variety of nilpotent groups of class 2 ?
	
	An automorphism $\varphi$ of a category is called inner if it is isomorphic to the identity automorphism. Let $s : 1 \rightarrow  \varphi$ be a function defining this isomorphism. Then for every object $A$ we have an isomorphism ${s}_{A} : A \rightarrow  \varphi \left( A\right)$ and for any morphism of objects $\mu  : A \rightarrow  B$ we have $\varphi \left( \mu \right)  = {s}_{B}\mu {s}_{A}^{-1}$ .
	
	Yes, it is, even for the variety of nilpotent groups of any class $n$ (A. Tsurkov, Int. J. Algebra Comput., 17 (2007), 1273-1281).
\end{openproblem}

\plainproblemsection

\begin{openproblem}
	15.76. b) If $\Theta$ is a variety of groups, then let ${\Theta }^{0}$ denote the category of all free groups of finite rank in $\Theta$ . It is proved (G. Mashevitzky, B. Plotkin, E. Plotkin J. Algebra, 282 (2004),490-512) that if $\Theta$ is the variety of all groups, then every automorphism of the category ${\Theta }^{0}$ is an inner one. The same is true if $\Theta$ is the variety of all nilpotent groups of class $\leq  d$ , for $d \geq  2$ (A. Tsurkov, Int. J. Algebra Comput., $\mathbf{{17}}\left( {2007}\right)$ , 1273-1281). Is this true for the varieties of solvable groups? of metabelian groups?
	
	An automorphism $\varphi$ of a category is called inner if it is isomorphic to the identity automorphism. Let $s : 1 \rightarrow  \varphi$ be a function defining this isomorphism. Then for every object $A$ we have an isomorphism ${s}_{A} : A \rightarrow  \varphi \left( A\right)$ and for any morphism of objects $\mu  : A \rightarrow  B$ we have $\varphi \left( \mu \right)  = {s}_{B}\mu {s}_{A}^{-1}$ . B. I. Plotkin
\end{openproblem}

\plainproblemsection

\begin{openproblem}
	15.77. Elements $a, b$ of a group $G$ are said to be symmetric with respect to an element $g \in  G$ if $a = g{b}^{-1}g$ . Let $G$ be an infinite abelian group, $\alpha$ a cardinal, $\alpha  < \left| G\right|$ . Is it true that for any $n$ -colouring $\chi  : G \rightarrow  \{ 0,1,\ldots , n - 1\}$ there exists a monochrome subset of cardinality $\alpha$ that is symmetric with respect to some element of $G$ ? This is known to be true for $n \leq  3$ . I. V. Protasov
\end{openproblem}

\plainproblemsection

\begin{openproblem}
	15.78. (R.I. Grigorchuk). Is it true that for any $n$ -colouring of a free group of any rank there exists a monochrome subset that is symmetric with respect to some element of the group? This is true for $n = 2$ . I. V. Protasov
\end{openproblem}

\plainproblemsection

\begin{openproblem}
	*15.79. Does there exist a Hausdorff group topology on $\mathbb{Z}$ such that the sequence $\left\{  {{2}^{n} + {3}^{n}}\right\}$ converges to zero? I. V. Protasov
	
	*Yes, it exists, as follows from Theorem 3 in (I. Z. Ruzsa, Proc. Conf. Number theory (Budapest, 1987). Vol.I: Elementary and analytic, Colloq. Math. Soc. János Bolyai, 51, North-Holland, Amsterdam, 1990, 473-504). Another solution is in (S. V. Skresanov, Siberian Math. J., 61, no. 3 (2020), 542-544).
\end{openproblem}

\plainproblemsection

\begin{openproblem}
	*15.80. A sequence $\left\{  {F}_{n}\right\}$ of pairwise disjoint finite subsets of a topological group is called expansive if for every open subset $U$ there is a number $m$ such that ${F}_{n} \cap  U \neq  \varnothing$ for all $n > m$ . Suppose that a group $G$ can be partitioned into countably many dense subsets. Is it true that in $G$ there exists an expansive sequence? I. V. Protasov
	
	*Not necessarily. One example is $G = \mathop{\prod }\limits_{\mathbb{N}}\mathbb{R}$ under the box topology. It is wellknown that $G$ does not have a countable dense subset (hence has no expansive sequence). However, the topology on $G$ has a basis having ${2}^{{\aleph }_{0}}$ elements, each element of the basis having ${2}^{{\aleph }_{0}}$ points, and one can partition $G$ into countably many dense subsets (via a transfinite induction of length ${2}^{{\aleph }_{0}}$ ). (Letter of $S$ . Corson of 19 May 2025). Another example: any infinite abelian group $G$ with finite set of elements of order 2 can be partitioned into countably many subsets dense in every non-discrete group topology on $G$ according to Corollary 12.21 in (Y. Zelenyuk, Ultrafilters and topologies on groups, De Gruyter, 2011). (Letter of I. V. Protasov of 19 May 2025).
\end{openproblem}

\plainproblemsection

\begin{openproblem}
	15.81. Let $G$ be a finite non-supersoluble group. Is it true that $G$ has a non-cyclic Sylow subgroup $P$ such that some maximal subgroup of $P$ has no proper complement in $G$ ? A. N. Skiba
	
	Yes, it is (Wang, Yanming, Wei, Huaquan, Sci. China Ser. A, 47 (2004), no. 1, 96-103).
\end{openproblem}

\plainproblemsection

\begin{openproblem}
	15.82. Suppose that a periodic group $G$ contains a strongly isolated 2-subgroup $U$ . Is it true that either $G$ is locally finite, or $U$ is a normal subgroup of $G$ ?
	
	A. I. Sozutov, N. M. Suchkov
\end{openproblem}

\plainproblemsection

\begin{openproblem}
	15.83. (Yu. I. Merzlyakov). Does there exist a rational number $\alpha$ such that $\left| \alpha \right|  < 2$ and the matrices $\left( \begin{array}{ll} 1 & \alpha \\  0 & 1 \end{array}\right)$ and $\left( \begin{array}{ll} 1 & 0 \\  \alpha & 1 \end{array}\right)$ generate a free group? ${Yu}.V$ . Sosnovskiĭ
\end{openproblem}

\plainproblemsection

\begin{openproblem}
	15.84. Yu. I. Merzlyakov (Sov. Math. Dokl., 19 (1978), 64-68) proved that if the complex numbers $\alpha ,\beta ,\gamma$ are each at least 3 in absolute value, then the matrices $\left( \begin{array}{ll} 1 & \alpha \\  0 & 1 \end{array}\right) ,\left( \begin{array}{ll} 1 & 0 \\  \beta & 1 \end{array}\right)$ , and $\left( \begin{matrix} 1 - \gamma &  - \gamma \\  \gamma & 1 + \gamma  \end{matrix}\right)$ generate a free group of rank three. Are there rational numbers $\alpha ,\beta ,\gamma$ , each less than 3 in absolute value, with the same property? Yu. V. Sosnovskii
\end{openproblem}

\plainproblemsection

\begin{openproblem}
	15.85. A torsion-free group all of whose subgroups are subnormal is nilpotent (H. Smith, Arch. Math., 76, no. 1 (2001), 1-6). Is a torsion-free group with the normalizer condition
	
	a) hyperabelian?
	
	b) hypercentral? Yu. V. Sosnovskii
\end{openproblem}

\plainproblemsection

\begin{openproblem}
	15.86. A group $G$ is called discriminating if for any finite set of nontrivial elements of the direct square $G \times  G$ there is a homomorphism $G \times  G \rightarrow  G$ which does not annihilate any of them (G. Baumslag, A. G. Myasnikov, V. N. Remeslennikov). A group $G$ is called squarelike if $G$ is universally equivalent (in the sense of first order logic) to a discriminating group (B. Fine, A. M. Gaglione, A. G. Myasnikov, D. Spellman). Must every squarelike group be elementarily equivalent to a discriminating group?
	
	Yes, it must (O. Belegradek, J. Group Theory, 7, no. 4 (2004), 521-532; B. Fine, A. M. Gaglione, D. Spellman, Archiv Math., 83, no. 2 (2004), 106-112).
\end{openproblem}

\plainproblemsection

\begin{openproblem}
	15.87. Suppose that a 2-group $G$ admits a regular periodic locally cyclic group of automorphisms that transitively permutes the set of involutions of $G$ . Is $G$ locally finite? N. M. Suchkov
\end{openproblem}

\plainproblemsection

\begin{openproblem}
	15.88. Let $\mathfrak{A}$ and ${\mathfrak{N}}_{c}$ denote the varieties of abelian groups and nilpotent groups of class $\leq  c$ , respectively. Let ${F}_{r} = {F}_{r}\left( {\mathfrak{{AN}}}_{c}\right)$ be a free group of rank $r$ in ${\mathfrak{{AN}}}_{c}$ . An element of ${F}_{r}$ is called primitive if it can be included in a basis of ${F}_{r}$ . Does there exist an algorithm recognizing primitive elements in ${F}_{r}$ ? E. I. Timoshenko
\end{openproblem}

\plainproblemsection

\begin{openproblem}
	15.89. Let $\Gamma$ be an infinite undirected connected vertex-symmetric graph of finite valency without loops or multiple edges. Is it true that every complex number is an eigenvalue of the adjacency matrix of $\Gamma$ under its natural action as a linear operator on the complex vector space of all complex-valued functions on the vertex set of $\Gamma$ ?
	
	V. I. Trofimov
\end{openproblem}

\plainproblemsection

\begin{openproblem}
	15.90. Let $\Gamma$ be an infinite directed graph, and $\bar{\Gamma }$ the underlying undirected graph. Suppose that the graph $\Gamma$ admits a vertex-transitive group of automorphisms, and the graph $\bar{\Gamma }$ is connected and of finite valency. Does there exist a positive integer $k$ (possibly depending on $\Gamma$ ) such that for any positive integer $n$ there is a directed path of length at most $k \cdot  n$ in the graph $\Gamma$ whose initial and terminal vertices are at distance at least $n$ in the graph $\bar{\Gamma }$ ?
	
	Comment of 2005: It is proved that there exists a positive integer ${k}^{\prime }$ (depending only on the valency of $\bar{\Gamma }$ ) such that for any positive integer $n$ there is a directed path of length at most ${k}^{\prime } \cdot  {n}^{2}$ in the graph $\Gamma$ whose initial and terminal vertices are at distance at least $n$ in the graph $\bar{\Gamma }$ (V. I. Trofimov, Europ. J. Combinatorics,27, no. 5 (2006), 690-700). V. I. Trofimov
\end{openproblem}

\plainproblemsection

\begin{openproblem}
	15.91. Is it true that any irreducible faithful representation of a linear group $G$ of finite rank over a finitely generated field of characteristic zero is induced from an irreducible representation of a finitely generated dense subgroup of the group $G$ ?
	
	A. V. Tushev
\end{openproblem}

\plainproblemsection

\begin{openproblem}
	15.92. A group $\left\langle  {x, y \mid  {x}^{l} = {y}^{m} = {\left( xy\right) }^{n}}\right\rangle$ is called the triangle group with parameters(l, m, n). It is proved in (A. M. Brunner, R. G. Burns, J. Wiegold, Math. Scientist, 4 (1979), 93-98) that the triangle group (2,3,30) has uncountably many non-isomorphic homomorphic images that are residually finite alternating groups. Is the same true for the triangle group(2,3, n)for all $n > 6$ ? J. Wiegold
\end{openproblem}

\plainproblemsection

\begin{openproblem}
	15.93. Let $G$ be a pro- $p$ group, $p > 2$ , and $\varphi$ an automorphism of $G$ of order 2 . Suppose that the centralizer ${C}_{G}\left( \varphi \right)$ is abelian. Is it true that $G$ satisfies a pro- $p$ identity?
	
	An affirmative answer would generalize a result of A. N. Zubkov (Siberian Math. J., 28, no. 5 (1987), 742-747) saying that non-abelian free pro- $p$ groups cannot be represented by $2 \times  2$ matrices. A. Jaikin-Zapirain
\end{openproblem}

\plainproblemsection

\begin{openproblem}
	15.94. Define the weight of a group $G$ to be the minimum number of generators of $G$ as a normal subgroup of itself. Let $G = {G}_{1} * \cdots  * {G}_{n}$ be a free product of $n$ nontrivial groups.
	
	a) Is it true that the weight of $G$ is at least $n/2$ ?
	
	b) Is it true if the ${G}_{i}$ are cyclic?
	
	c) Is it true for $n = 3$ (without assuming the ${G}_{i}$ cyclic)?
	
	Problem 5.53 in Archive (now answered in the affirmative) is the special case with $n = 3$ and all the ${G}_{i}$ cyclic. J. Howie
\end{openproblem}

\plainproblemsection

\begin{openproblem}
	15.95. (A. Mann, Ch. Praeger). Suppose that all fixed-point-free elements of a transitive permutation group $G$ have prime order $p$ . If $G$ is a finite $p$ -group, must $G$ have exponent $p$ ?
	
	This is true for $p = 2,3$ ; it is also known that the exponent of $G$ is bounded in terms of $p$ (A. Hassani, M. Khayaty, E. I. Khukhro, C. E. Praeger, J. Algebra,214, no. 1 (1999), 317-337). E. I. Khukhro
\end{openproblem}

\plainproblemsection

\begin{openproblem}
	15.96. An automorphism $\varphi$ of a group $G$ is called a splitting automorphism of order $n$ if ${\varphi }^{n} = 1$ and $x{x}^{\varphi }{x}^{{\varphi }^{2}}\cdots {x}^{{\varphi }^{n - 1}} = 1$ for any $x \in  G$ .
	
	a) Is it true that the derived length of a $d$ -generated nilpotent $p$ -group admitting a splitting automorphism of order ${p}^{n}$ is bounded by a function of $d, p$ , and $n$ ? This is true for $n = 1$ , see Archive,7.53.
	
	b) The same question for ${p}^{n} = 4$ . E. I. Khukhro
\end{openproblem}

\plainproblemsection

\begin{openproblem}
	15.97. Let $p$ be a prime. A group $G$ satisfies the $p$ -minimal condition if there are no infinite descending chains ${G}_{1} > {G}_{2} > \ldots$ of subgroups of $G$ such that each difference ${G}_{i} \smallsetminus  {G}_{i + 1}$ contains a $p$ -element (S. N. Chernikov). Suppose that a locally finite group $G$ satisfies the $p$ -minimal condition and has a subnormal series each of whose factors is finite or a ${p}^{\prime }$ -group. Is it true that all $p$ -elements of $G$ generate a Chernikov subgroup?
	
	N. S. Chernikov
\end{openproblem}

\plainproblemsection

\begin{openproblem}
	15.98. Let $\mathfrak{F}$ be a saturated formation, and $G$ a finite soluble minimal non- $\mathfrak{F}$ -group such that ${G}^{\mathfrak{F}}$ is a Sylow $p$ -subgroup of $G$ . Is it true that ${G}^{\mathfrak{F}}$ is isomorphic to a factor group of the Sylow $p$ -subgroup of ${PSU}\left( {3,{p}^{2n}}\right)$ ? This is true for the formation of finite nilpotent groups (V. D. Mazurov, S. A. Syskin, Math. Notes, 14, no. 2 (1973), 683-686; A. Kh. Zhurtov, S. A. Syskin, Siberian Math. J., 26, no. 2 (1987), 235-239).
	
	L. A. Shemetkov
\end{openproblem}

\plainproblemsection

\begin{openproblem}
	15.99. Let $f\left( n\right)$ be the number of isomorphism classes of finite groups of order $n$ . Is it true that the equation $f\left( n\right)  = k$ has a solution for any positive integer $k$ ? The answer is affirmative for all $k \leq  {1000}$ (G. M. Wei, Southeast Asian Bull. Math.,22, no. 1 (1998), 93-102). W. J. Shi
\end{openproblem}

\plainproblemsection

\begin{openproblem}
	15.100. Is a periodic group locally finite if it has a non-cyclic subgroup of order 4 that coincides with its centralizer? A. K. Shlëpkin
\end{openproblem}

\plainproblemsection

\begin{openproblem}
	15.101. Is a periodic group locally finite if it has an involution whose centralizer is a locally finite group with Sylow 2-subgroup of order 2? A. K. Shlëpkin
\end{openproblem}

\plainproblemsection

\begin{openproblem}
	15.102. (W. Magnus). An element $r$ of a free group ${F}_{n}$ is called a normal root of an element $u \in  {F}_{n}$ if $u$ belongs to the normal closure of $r$ in ${F}_{n}$ . Can an element $u$ that lies outside the commutator subgroup $\left\lbrack  {{F}_{n},{F}_{n}}\right\rbrack$ have infinitely many non-conjugate normal roots? V. E. Shpilrain
\end{openproblem}

\plainproblemsection

\begin{openproblem}
	15.103. (Well-known problem). Is the group $\operatorname{Out}\left( {F}_{3}\right)$ of outer automorphisms of a free group of rank 3 linear? V. E. Shpilrain
\end{openproblem}

\plainproblemsection

\begin{openproblem}
	15.104. Let $n$ be a positive integer and let $w$ be a group word in the variables ${x}_{1},{x}_{2},\ldots$ Suppose that a residually finite group $G$ satisfies the identity ${w}^{n} = 1$ . Does it follow that the verbal subgroup $w\left( G\right)$ is locally finite?
	
	This is the Restricted Burnside Problem if $w = {x}_{1}$ . A positive answer was obtained also in a number of other particular cases (P. V. Shumyatsky, Quart. J. Math., 51, no. 4 (2000), 523-528). P. V. Shumyatsky
\end{openproblem}

\plainproblemsection

\begin{openproblem}
	16.1. a) Let $G$ be a finite non-abelian group, and $Z\left( G\right)$ its centre. One can associate a graph ${\Gamma }_{G}$ with $G$ as follows: take $G \smallsetminus  Z\left( G\right)$ as vertices of ${\Gamma }_{G}$ and join two vertices $x$ and $y$ if ${xy} \neq  {yx}$ . Let $H$ be a finite non-abelian group such that ${\Gamma }_{G} \cong  {\Gamma }_{H}$ . If $H$ is simple, is it true that $G \cong  H$ ? Comment of 2009: This is true for groups with disconnected prime graphs (L. Wang, W. Shi, Commun. Algebra, 36 (2008), 523-528). A. Abdollahi, S. Akbari, H. R. Maimani
	
	a) Yes, it is mod CFSG (Ch. Khan', G. Chèn', S. Go, Siberian Math. J., 49, no. 6 (2008), 1138-1146, for sporadic simple groups; A. Abdollahi, H. Shahverdi, J. Algebra, 357 (2012), 203-207, for alternating groups; R. M. Solomon, A. J. Woldar, J. Group Theory, 16, no. 6 (2013), 793-824, for simple groups of Lie type).
\end{openproblem}

\plainproblemsection

\begin{openproblem}
	16.1. Let $G$ be a finite non-abelian group, and $Z\left( G\right)$ its centre. One can associate a graph ${\Gamma }_{G}$ with $G$ as follows: take $G \smallsetminus  Z\left( G\right)$ as vertices of ${\Gamma }_{G}$ and join two vertices $x$ and $y$ if ${xy} \neq  {yx}$ . Let $H$ be a finite non-abelian group such that ${\Gamma }_{G} \cong  {\Gamma }_{H}$ .
	
	b) If $H$ is nilpotent, is it true that $G$ is nilpotent?
	
	Comment of 2009: This is true if $\left| H\right|  = \left| G\right|$ (A. Abdollahi, S. Akbari, H. R. Maimani, J. Algebra, 298 (2006), 468-492).
	
	Comments of 2023: This is true if $\left| {Z\left( H\right) }\right|  \geq  \left| {Z\left( G\right) }\right|$ (H. Shahverdi, J. Algebra, 642 (2024), 60-64), using the earlier result under the additional assumption that the centralizer of every non-central element in $H$ is abelian (V. Grazian, C. Monetta, J. Algebra, 633 (2023), 389-402).
	
	c) If $H$ is solvable, is it true that $G$ is solvable?
	
	A. Abdollahi, S. Akbari, H. R. Maimani
\end{openproblem}

\plainproblemsection

\begin{openproblem}
	16.2. A group $G$ is subgroup-separable if for any subgroup $H \leq  G$ and element $x \in  G \smallsetminus  H$ there is a homomorphism to a finite group $f : G \rightarrow  F$ such that $f\left( x\right)  \notin  f\left( H\right)$ . Is it true that a finitely generated solvable group is locally subgroup-separable if and only if it does not contain a solvable Baumslag-Solitar group? Solvable Baumslag-Solitar groups are ${BS}\left( {1, n}\right)  = \left\langle  {a, b \mid  {ba}{b}^{-1} = {a}^{n}}\right\rangle$ for $n > 1$ .
	
	Background: It is known that a finitely generated solvable group is subgroup-separable if and only if it is polycyclic (R. C. Alperin, in: Groups-Korea'98 (Pusan), de Gruyter, Berlin, 2000, 1-5). R. C. Alperin
	
	No, it is not (J. O. Button, Ricerche di Matematica, 61, no. 1 (2012), 139-145).
\end{openproblem}

\plainproblemsection

\begin{openproblem}
	16.3. Is it true that if $G$ is a finite group with all conjugacy classes of distinct sizes, then $G \cong  {\mathbb{S}}_{3}$ ?
	
	This is true if $G$ is solvable (J. Zhang, J. Algebra, $\mathbf{{170}}\left( {1994}\right) ,{608} - {624}$ ); it is also known that $F\left( G\right)$ is nontrivial (Z. Arad, M. Muzychuk, A. Oliver, J. Algebra,280 (2004), 537-576). Z. Arad
\end{openproblem}

\plainproblemsection

\begin{openproblem}
	16.4. Let $G$ be a finite group with $C, D$ two nontrivial conjugacy classes such that ${CD}$ is also a conjugacy class. Can $G$ be a non-abelian simple group? Z. Arad
\end{openproblem}

\plainproblemsection

\begin{openproblem}
	16.5. A group is said to be perfect if it coincides with its derived subgroup. Does there exist a perfect locally finite $p$ -group
	
	a) all of whose proper subgroups are hypercentral?
	
	b) all of whose proper subgroups are solvable?
	
	c) all of whose proper subgroups are hypercentral and solvable? A. O. Asar
\end{openproblem}

\plainproblemsection

\begin{openproblem}
	16.6. Can a perfect locally finite $p$ -group be generated by a subset of bounded exponent
	
	a) if all of its proper subgroups are hypercentral?
	
	b) if all of its proper subgroups are solvable? A. O. Asar
\end{openproblem}

\plainproblemsection

\begin{openproblem}
	16.7. Is it true that the membership problem is undecidable for any semidirect product ${F}_{n} \leftthreetimes  {F}_{n}$ of non-abelian free groups ${F}_{n}$ ?
	
	An affirmative answer would imply an answer to 6.24. Recall that the membership problem is decidable for ${F}_{n}$ (M. Hall,1949) and undecidable for ${F}_{n} \times  {F}_{n}$ (K. A. Mikhailova, 1958). V. G. Bardakov
\end{openproblem}

\plainproblemsection

\begin{openproblem}
	16.8. The width $w\left( {G}^{\prime }\right)$ of the derived subgroup ${G}^{\prime }$ of a finite non-abelian group $G$ is the smallest positive integer $m$ such that every element of ${G}^{\prime }$ is a product of $\leq  m$ commutators. Is it true that the maximum value of the ratio $w\left( {G}^{\prime }\right) /\left| G\right|$ is $1/6$ (attained at the symmetric group ${\mathbb{S}}_{3}$ )? V. G. Bardakov
	
	Yes, it is (T. Bonner, J. Algebra, 320 (2008), 3165-3171).
\end{openproblem}

\plainproblemsection

\begin{openproblem}
	16.9. An element $g$ of a free group ${F}_{n}$ on the free generators ${x}_{1},\ldots ,{x}_{n}$ is called a palindrome with respect to these generators if the reduced word representing $g$ is the same when read from left to right or from right to left. The palindromic length of an element $w \in  {F}_{n}$ is the smallest number of palindromes in ${F}_{n}$ whose product is $w$ . Is there an algorithm for finding the palindromic length of a given element of ${F}_{n}$ ?
	
	V. G. Bardakov, V. A. Tolstykh, V. E. Shpilrain
\end{openproblem}

\plainproblemsection

\begin{openproblem}
	16.10. Is there an algorithm for finding the primitive length of a given element of ${F}_{n}$ ? The definition of a primitive element is given in 14.84; the primitive length is defined similarly to the palindromic length in 16.9.
	
	V. G. Bardakov, V. A. Tolstykh, V. E. Shpilrain
\end{openproblem}

\plainproblemsection

\begin{openproblem}
	16.11. Let $G$ be a finite $p$ -group. Does there always exist a finite $p$ -group $H$ such that $\Phi \left( H\right)  \cong  \left\lbrack  {G, G}\right\rbrack$ ? Ya. G. Berkovich
\end{openproblem}

\plainproblemsection

\begin{openproblem}
	16.12. Given a finite $p$ -group $G$ , we define a $\Phi$ -extension of $G$ as any finite $p$ -group $H$ containing a normal subgroup $N$ of order $p$ such that $H/N \cong  G$ and $N \leq  \Phi \left( H\right)$ . Is it true that for every finite $p$ -group $G$ there exists an infinite sequence $G = {G}_{1},{G}_{2},\ldots$ such that ${G}_{i + 1}$ is a $\Phi$ -extension of ${G}_{i}$ for all $i = 1,2,\ldots$ ? Ya. G. Berkovich Yes, it is (S. F. Kamornikov, Izv. Gomel' Univ., 5 (2008), 200-201 (Russian)).
\end{openproblem}

\plainproblemsection

\begin{openproblem}
	16.13. Does there exist a finite $p$ -group $G$ all of whose maximal subgroups $H$ are special, that is, satisfy $Z\left( H\right)  = \left\lbrack  {H, H}\right\rbrack   = \Phi \left( H\right)$ ? Ya. G. Berkovich
	
	Yes, for all primes there are groups of arbitrarily large size with this property (J. Cossey, Bull. Austral. Math. Soc., 89 (2014), 415-419); an example for $p = 2$ given by a Sylow 2-subgroup of ${L}_{3}\left( 4\right)$ was independently presented by V. I. Zenkov at Mal'cev Meeting-2014, 10-14 November, 2014, Novosibirsk.
\end{openproblem}

\plainproblemsection

\begin{openproblem}
	16.14. Let $G$ be a finite 2-group such that ${\Omega }_{1}\left( G\right)  \leq  Z\left( G\right)$ . Is it true that the rank of $G/{G}^{2}$ is at most double the rank of $Z\left( G\right)$ ? Ya. G. Berkovich
\end{openproblem}

\plainproblemsection

\begin{openproblem}
	16.15. An element $g$ of a group $G$ is an Engel element if for every $h \in  G$ there exists $k$ such that $\left\lbrack  {h, g,\ldots , g}\right\rbrack   = 1$ , where $g$ occurs $k$ times; if there is such $k$ independent of $h$ , then $g$ is said to be boundedly Engel.
	
	*a) (B.I.Plotkin). Does the set of boundedly Engel elements of a group form a subgroup?
	
	b) The same question for torsion-free groups.
	
	c) The same question for right-ordered groups.
	
	d) The same question for linearly ordered groups. V. V. Bludov
\end{openproblem}

\plainproblemsection

\begin{openproblem}
	16.16. a) Does the set of (not necessarily boundedly) Engel elements of a group without elements of order 2 form a subgroup?
	
	b) The same question for torsion-free groups.
	
	c) The same question for right-ordered groups.
	
	d) The same question for linearly ordered groups.
	
	There are examples of 2-groups where a product of two (unboundedly) Engel elements is not an Engel element. V. V. Bludov
\end{openproblem}

\plainproblemsection

\begin{openproblem}
	16.17. Is it true that a non-abelian simple group cannot contain Engel elements other than the identity element? V. V. Bludov
	
	No, it can: since an involution is an Engel element in any 2-group, counterexamples are provided by non-abelian simple 2-groups; see Archive, 4.74a). (V. V. Bludov,
	
	Letter of 30 March 2006.)
\end{openproblem}

\plainproblemsection

\begin{openproblem}
	16.18. Does there exist a linearly orderable soluble group of derived length exactly $n$ that has a single proper normal relatively convex subgroup
	
	a) for $n = 3$ ?
	
	b) for $n > 4$ ?
	
	Such groups do exist for $n = 2$ and for $n = 4$ .V. V. Bludov
\end{openproblem}

\plainproblemsection

\begin{openproblem}
	16.19. Is the variety of lattice-ordered groups generated by nilpotent groups finitely based? (Here the variety is considered in the signature of group and lattice operations.) V. V. Bludov
\end{openproblem}

\plainproblemsection

\begin{openproblem}
	16.20. Let $M$ be a quasivariety of groups. The dominion ${\operatorname{dom}}_{A}^{M}\left( H\right)$ of a subgroup $H$ of a group $A$ (in $M$ ) is the set of all elements $a \in  A$ such that for any two homomorphisms $f, g : A \rightarrow  B \in  M$ , if $f, g$ coincide on $H$ , then $f\left( a\right)  = g\left( a\right)$ . Suppose that the set $\left\{  {{\operatorname{dom}}_{A}^{N}\left( H\right)  \mid  N\text{is a quasivariety,}N \subseteq  M}\right\}$ forms a lattice with respect to set-theoretic inclusion. Can this lattice be modular and non-distributive? A. I. Budkin
\end{openproblem}

\plainproblemsection

\begin{openproblem}
	16.21. Given a non-central matrix $\alpha  \in  S{L}_{n}\left( F\right)$ over a field $F$ for $n > 2$ , is it true that every non-central matrix in $S{L}_{n}\left( F\right)$ is a product of $n$ matrices, each similar to $\alpha$ ? L. Vaserstein
	\par\smallskip\hrule\smallskip
	*No, not always (A.I. Sozutov, Siberian Math. J., 60, no. 6 (2019), 1099-1100).
	\par\smallskip\hrule\smallskip
\end{openproblem}

\plainproblemsection

\begin{openproblem}
	16.22. (Well-known problem.) Let ${E}_{n}\left( A\right)$ be the subgroup of $G{L}_{n}\left( A\right)$ generated by elementary matrices. Is $S{L}_{2}\left( A\right)  = {E}_{2}\left( A\right)$ when $A = \mathbb{Z}\left\lbrack  {x,1/x}\right\rbrack$ ? L. Vaserstein
\end{openproblem}

\plainproblemsection

\begin{openproblem}
	16.23. Is there, for some $n > 2$ and a ring $A$ with 1, a matrix in ${E}_{n}\left( A\right)$ that is nonscalar modulo any proper ideal and is not a commutator? L. Vaserstein
\end{openproblem}

\plainproblemsection

\begin{openproblem}
	16.24. The spectrum of a finite group is the set of orders of its elements. Does there exist a finite group $G$ whose spectrum coincides with the spectrum of a finite simple exceptional group $L$ of Lie type, but $G$ is not isomorphic to $L$ ? $A.V$ . Vasil’ev Yes, for example, for $L = {}^{3}{D}_{4}\left( 2\right)$ (V. D. Mazurov, Algebra Logic,52, no. 5 (2013),400- 403). Further examples are given in M. A. Grechkoseeva, M. A. Zvezdina, J. Algebra Appl., 15, no. 9 (2016), Article ID 1650168, 13 pp.).
\end{openproblem}

\plainproblemsection

\begin{openproblem}
	16.25. Do there exist three pairwise non-isomorphic finite non-abelian simple groups with the same spectrum? A. V. Vasil’ev
	
	No (A. A. Buturlakin, Sibirsk. Electron. Math. Reports, 7 (2010), 111-114).
\end{openproblem}

\plainproblemsection

\begin{openproblem}
	16.26. We say that the prime graphs of finite groups $G$ and $H$ coincide if the sets of primes dividing their orders are the same, $\pi \left( G\right)  = \pi \left( H\right)$ , and for any distinct $p, q \in  \pi \left( G\right)$ there is an element of order ${pq}$ in $G$ if and only if there is such an element in $H$ . Does there exist a positive integer $k$ such that there are no $k$ pairwise non-isomorphic finite non-abelian simple groups with the same graphs of primes? Conjecture: $k = 5$ . A. V. Vasil’ev
\end{openproblem}

\plainproblemsection

\begin{openproblem}
	16.27. Suppose that a finite group $G$ has the same spectrum as an alternating group. Is it true that $G$ has at most one non-abelian composition factor?
	
	For any finite simple group other than alternating the answer to the corresponding question is affirmative (mod CFSG). A. V. Vasil'ev and V. D. Mazurov Yes, it is (I. B. Gorshkov, Algebra and Logic, 52, no. 1 (2013), 41-45).
\end{openproblem}

\plainproblemsection

\begin{openproblem}
	16.28. Let $G$ be a connected linear reductive algebraic group over a field of positive characteristic, $X$ a closed subset of $G$ , and let ${X}^{k} = \left\{  {{x}_{1}\cdots {x}_{k} \mid  {x}_{i} \in  X}\right\}$ .
	
	a) Is it true that there always exists a positive integer $c = c\left( X\right)  > 1$ such that ${X}^{c}$ is closed?
	
	b) If $X$ is a conjugacy class of $G$ such that ${X}^{2}$ contains an open subset of $G$ , then is ${X}^{2} = G$ ? E. P. Vdovin
\end{openproblem}

\plainproblemsection

\begin{openproblem}
	16.29. Which finite simple groups of Lie type $G$ have the following property: for every semisimple abelian subgroup $A$ and proper subgroup $H$ of $G$ there exists $x \in  G$ such that ${A}^{x} \cap  H = 1$ ?
	
	This property holds when $A$ is a cyclic subgroup (J. Siemons, A. Zalesskii, J. Algebra, $\mathbf{{256}}\left( {2002}\right) ,{611} - {625}$ ), as well as when $A$ is contained in some maximal torus and $G = {PS}{L}_{n}\left( q\right)$ (J. Siemons, A. Zalesskii, J. Algebra,226 (2000),451-478). Note that if $G = {L}_{2}\left( 5\right) , A = 2 \times  2$ , and $H = 5 : 2$ (in Atlas notation), then ${A}^{x} \cap  H > 1$ for every $x \in  G$ (this example was communicated to the author by V. I. Zenkov). E. P. Vdovin
\end{openproblem}

\plainproblemsection

\begin{openproblem}
	16.30. Suppose that $A$ and $B$ are subgroups of a group $G$ and $G = {AB}$ . Will $G$ have composition (principal) series if $A$ and $B$ have composition (respectively, principal) series? V. A. Vedernikov
\end{openproblem}

\plainproblemsection

\begin{openproblem}
	16.31. Suppose that a group $G$ has a composition series and let $\mathfrak{F}\left( G\right)$ be the formation generated by $G$ . Is the set of all subformations of $\mathfrak{F}\left( G\right)$ finite? $V.A$ . Vedernikov No, not always, even for $G$ finite (V. P. Burichenko, J. Algebra, $\mathbf{{372}}\left( {2012}\right) ,{428} - {458}$ ).
\end{openproblem}

\plainproblemsection

\begin{openproblem}
	16.32. Suppose that a group $G$ has a composition series and let $\operatorname{Fit}\left( G\right)$ be the Fitting class generated by $G$ . Is the set of all Fitting subclasses of $\operatorname{Fit}\left( G\right)$ finite? Cf. 14.31.
	
	V. A. Vedernikov
\end{openproblem}

\plainproblemsection

\begin{openproblem}
	16.33. Suppose that a finite $p$ -group $G$ has an abelian subgroup $A$ of order ${p}^{n}$ . Does $G$ contain an abelian subgroup $B$ of order ${p}^{n}$ that is normal in $\left\langle  {B}^{G}\right\rangle$
	
	a) if $p = 3$ ?
	
	b) if $p = 2$ ?
	
	c) If $p = 3$ and $A$ is elementary abelian, does $G$ contain an elementary abelian subgroup $B$ of order ${3}^{n}$ that is normal in $\left\langle  {B}^{G}\right\rangle$ ?
	
	The corresponding results have been proved for greater primes $p$ , based on extensions of Thompson's Replacement Theorem (and the dihedral group of order 32 shows that the third question has negative answer for $p = 2$ ). See (G. Glauberman, J. Algebra, 196 (1997), 301-338; J. Alperin, G. Glauberman, J. Algebra, 203 (1998), 533-566; G. Glauberman, J. Algebra, 272 (2004), 128-153). G. Glauberman
\end{openproblem}

\plainproblemsection

\begin{openproblem}
	16.34. Suppose that $G$ is a finitely generated group acting faithfully on a regular rooted tree by finite-state automorphisms. Is the conjugacy problem decidable for $G$ ?
	
	See the definitions in (R. I. Grigorchuk, V. V. Nekrashevich, V. I. Sushchanskii, Proc. Steklov Inst. Math., 2000, no. 4 (231), 128-203).
	
	R. I. Grigorchuk, V. V. Nekrashevich, V. I. Sushchanskii
\end{openproblem}

\plainproblemsection

\begin{openproblem}
	16.35. Is every finitely presented soluble group nilpotent-by-nilpotent-by-finite? J. R. J. Groves
	
	No: for a prime $p$ the group $G = \left\langle  {a, b, c, d \mid  {b}^{a} = {b}^{p},{c}^{a} = {c}^{p},{d}^{a} = d,{c}^{b} = c}\right.$ , ${d}^{c} = d{d}^{b},\left\lbrack  {d,{d}^{b}}\right\rbrack   = 1,{d}^{p} = 1\rangle$ is soluble of derived length 3, but is not metanilpotent-by-finite. (V. V. Bludov, Letter of 27 April 2007.)
\end{openproblem}

\plainproblemsection

\begin{openproblem}
	16.36. (Well-known problem). We call a finite group rational if all of its ordinary characters are rational-valued. Is every Sylow 2-subgroup of a rational group also a rational group? M. R. Darafsheh
	
	No, not always (I. M. Isaacs, G. Navarro, Math. Z., 272 (2012), 937-945).
\end{openproblem}

\plainproblemsection

\begin{openproblem}
	16.37. Let $G$ be a solvable rational finite group with an extra-special Sylow 2-subgroup. Is it true that either $G$ is a 2-nilpotent group, or there is a normal subgroup $E$ of $G$ such that $G/E$ is an extension of a normal 3-subgroup by an elementary abelian 2-group? M. R. Darafsheh
	
	Yes, it is (M. R. Darafsheh, H. Sharifi, Extracta Math., 22, no. 1 (2007), 83-91).
\end{openproblem}

\plainproblemsection

\begin{openproblem}
	16.38. Let $G$ be a soluble group, and let $A$ and $B$ be periodic subgroups of $G$ . Is it true that any subgroup of $G$ contained in the set ${AB} = \{ {ab} \mid  a \in  A, b \in  B\}$ is periodic? This is known to be true if ${AB}$ is a subgroup of $G$ . F. de Giovanni
\end{openproblem}

\plainproblemsection

\begin{openproblem}
	16.39. (J. E. Humphreys, D. N. Verma). Let $G$ be a semisimple algebraic group over an algebraically closed field $k$ of characteristic $p > 0$ . Let $\mathfrak{g}$ be the Lie algebra of $G$ and let $u = u\left( \mathfrak{g}\right)$ be the restricted enveloping algebra of $\mathfrak{g}$ . By a theorem of Curtis every irreducible restricted $u$ -module (i. e. every irreducible restricted $\mathfrak{g}$ -module) is the restriction to $\mathfrak{g}$ of a (rational) $G$ -module. Is it also true that every projective indecomposable $u$ -module is the restriction of a rational $G$ -module? This is true if $p \geq  {2h} - 2$ (where $h$ is the Coxeter number of $G$ ) by results of Jantzen.
\end{openproblem}

\plainproblemsection

\begin{openproblem}
	16.40. Let $\Delta$ be a subgroup of the automorphism group of a free pro- $p$ group of finite rank $F$ such that $\Delta$ is isomorphic (as a profinite group) to the group ${Z}_{p}$ of $p$ -adic integers. Is the subgroup of fixed points of $\Delta$ in $F$ finitely generated (as a profinite group)? P. A. Zalesskiĭ
\end{openproblem}

\plainproblemsection

\begin{openproblem}
	16.41. Let $F$ be a free pro- $p$ group of finite rank $n > 1$ . Does Aut $F$ possess an open subgroup of finite cohomological dimension? P. A. Zalesskiĭ
\end{openproblem}

\plainproblemsection

\begin{openproblem}
	*16.42. Is a topological Abelian group $\left( {G,\tau }\right)$ compact if every group topology ${\tau }^{\prime } \subseteq  \tau$ on $G$ is complete? (The answer is yes if every continuous homomorphic image of $\left( {G,\tau }\right)$ is complete.) E. G. Zelenyuk
	
	*Yes, it is (T. Banakh, Topology Appl., 271 (2020), Article ID 106983, 17 p.).
\end{openproblem}

\plainproblemsection

\begin{openproblem}
	16.43. Is there a partition of the group ${\bigoplus }_{{\omega }_{1}}\left( {\mathbb{Z}/3\mathbb{Z}}\right)$ into three subsets whose complements do not contain cosets modulo infinite subgroups? (There is a partition into two such subsets.) E. G. Zelenyuk
	
	Yes, moreover, every infinite abelian group with finitely many involutions can be partitioned into infinitely many subsets such that every coset modulo an infinite subgroup meets each subset of the partition (Y. Zelenyuk, J. Combin. Theory (A), 115 (2008), 331-339).
\end{openproblem}

\plainproblemsection

\begin{openproblem}
	16.44. Is there in ZFC a countable non-discrete topological group not containing discrete subsets with a single accumulation point? (Such a group is known to exist under Martin's Axiom.) E. G. Zelenyuk
\end{openproblem}

\plainproblemsection

\begin{openproblem}
	16.45. Let $G$ be a permutation group on a set $\Omega$ . A sequence of points of $\Omega$ is a base for $G$ if its pointwise stabilizer in $G$ is the identity; it is minimal if no point may be removed. Let $b\left( G\right)$ be the maximum, over all permutation representations of the finite group $G$ , of the maximum size of a minimal base for $G$ . Let ${\mu }^{\prime }\left( G\right)$ be the maximum size of an independent set in $G$ , a set of elements with the property that no element belongs to the subgroup generated by the others. Is it true that $b\left( G\right)  = {\mu }^{\prime }\left( G\right)$ ? (It is known that $b\left( G\right)  \leq  {\mu }^{\prime }\left( G\right)$ , and that equality holds for the symmetric groups.)
	
	Remark. An equivalent question is the following. Suppose that the Boolean lattice $B\left( n\right)$ of subsets of an $n$ -element set is embeddable as a meet-semilattice of the subgroup lattice of $G$ , and suppose that $n$ is maximal with this property. Is it true that then there is such an embedding of $B\left( n\right)$ with the property that the least element of $B\left( n\right)$ is a normal subgroup of $G$ ? P. J. Cameron
\end{openproblem}

\plainproblemsection

\begin{openproblem}
	16.46. Among the finitely-presented groups that act arc-transitively on the (infinite) 3-valent tree with finite vertex-stabilizer are the two groups ${G}_{3} = \left\langle  {h, a, P, Q \mid  {h}^{3}}\right.$ , ${a}^{2},{P}^{2},{Q}^{2},\left\lbrack  {P, Q}\right\rbrack  ,\left\lbrack  {h, P}\right\rbrack  ,{\left( hQ\right) }^{2},{a}^{-1}{PaQ}\rangle$ and ${G}_{4} = \left\langle  {h, a, p, q, r \mid  {h}^{3},{a}^{2},{p}^{2}}\right.$ , ${q}^{2},{r}^{2},\left\lbrack  {p, q}\right\rbrack  ,\left\lbrack  {p, r}\right\rbrack  , p{\left( qr\right) }^{2},{h}^{-1}{phq},{h}^{-1}{qhpq},{\left( hr\right) }^{2},\left\lbrack  {a, p}\right\rbrack  ,{a}^{-1}{qar}\rangle$ , each of which contains the modular group ${G}_{1} = \left\langle  {h, a \mid  {h}^{3},{a}^{2}}\right\rangle   \cong  {\operatorname{PSL}}_{2}\left( \mathbb{Z}\right)$ as a subgroup of finite index. The free product of ${G}_{3}$ and ${G}_{4}$ with subgroup ${G}_{1} = \langle h, a\rangle$ amalgamated has a normal subgroup $K$ of index 8 generated by $A = h, B = {aha}, C = p$ , $D = {PpP}, E = {QpQ}$ , and $F = {PQpPQ}$ , with dihedral complement $\langle a, P, Q\rangle$ . The group $K$ has presentation $\langle A, B, C, D, E, F \mid  {A}^{3},{B}^{3},{C}^{2},{D}^{2},{E}^{2},{F}^{2},{\left( AC\right) }^{3}$ , ${\left( AD\right) }^{3},{\left( AE\right) }^{3},{\left( AF\right) }^{3},{\left( BC\right) }^{3},{\left( BD\right) }^{3},{\left( BE\right) }^{3},{\left( BF\right) }^{3},{\left( AB{A}^{-1}C\right) }^{2},{\left( AB{A}^{-1}D\right) }^{2},$ ${\left( {A}^{-1}BAE\right) }^{2},{\left( {A}^{-1}BAF\right) }^{2},{\left( BA{B}^{-1}C\right) }^{2},{\left( {B}^{-1}ABD\right) }^{2},{\left( BA{B}^{-1}E\right) }^{2},{\left( {B}^{-1}ABF\right) }^{2}\rangle .$ Does this group have a non-trivial finite quotient? M. Conder
\end{openproblem}

\plainproblemsection

\begin{openproblem}
	16.47. (P. Conrad). Is it true that every torsion-free abelian group admits an Archimedean lattice ordering? V. M. Kopytov, N. Ya. Medvedev
\end{openproblem}

\plainproblemsection

\begin{openproblem}
	16.48. A group $H$ is said to have generalized torsion if there exists an element $h \neq  1$ such that ${h}^{{x}_{1}}{h}^{{x}_{2}}\cdots {h}^{{x}_{n}} = 1$ for some $n$ and some ${x}_{i} \in  H$ . Is every group without generalized torsion right-orderable? V. M. Kopytov, N. Ya. Medvedev
\end{openproblem}

\plainproblemsection

\begin{openproblem}
	*16.49. Is it true that a free product of groups without generalized torsion is a group without generalized torsion? V. M. Kopytov, N. Ya. Medvedev
	
	*Yes, it is true; moreover, the generalized torsion in a free product of torsion-free groups is conjugate to a generalized torsion of one of its factor groups (T. Ito, K. Motegi, M. Teragaito, Proc. Amer. Math. Soc., 147, no. 11 (2019), 4999-5008).
\end{openproblem}

\plainproblemsection

\begin{openproblem}
	*16.50. Do there exist simple finitely generated right-orderable groups?
	
	There exist finitely generated right-orderable groups coinciding with the derived subgroup (G. M. Bergman). V. M. Kopytov, N. Ya. Medvedev
	
	*Yes, such groups do exist (J. Hyde, Y. Lodha, Invent. Math., 218 (2019), 83-112).
\end{openproblem}

\plainproblemsection

\begin{openproblem}
	*16.51. Do there exist groups that can be right-ordered in infinitely countably many ways?
	
	There exist groups that can be fully ordered in infinitely countably many ways (R. Buttsworth). V. M. Kopytov, N. Ya. Medvedev
	\par\smallskip\hrule\smallskip
	*No, there are no such groups (A. Clay, Preprint, 2009, http://arxiv.org/abs/ 0909.0273; A. Navas, C. Rivas, Algebr. Geom. Topol., 9 (2009), 2079-2100, with an appendix by A. Clay; P. A. Linnell, Bull. London Math. Soc., 43 (2011), 200- 202, based on (A. V. Zenkov, Sibirsk. Mat. Zh., 38 (1997), 90-92 (Russian), English transl., Siberian Math. J., 38, no. 1 (1997), 76-77), where this was proved for locally indicable groups.
\end{openproblem}

\plainproblemsection

\begin{openproblem}
	16.52. Is every finitely presented elementary amenable group solvable-by-finite? P. Linnell, T. Schick
	
	No, not every. I. Belegradek and Y. Cornulier have pointed out that the groups in (C. H. Houghton, Arch. Math. (Basel), 31, no. 3 (1978/79), 254-258) are finitely presented as shown in (K. S. Brown, J. Pure Appl. Algebra, 44, no. 1-3 (1987), 45- 75) and elementary amenable, but not virtually solvable; http://mathoverflow.net/questions/107996
\end{openproblem}

\plainproblemsection

\begin{openproblem}
	16.53. Let $d\left( G\right)$ denote the smallest cardinality of a generating set of the group $G$ . Suppose that $G = \langle A, B\rangle$ , where $A$ and $B$ are two $d$ -generated finite groups of coprime orders. Is it true that $d\left( G\right)  \leq  d + 1$ ?
	
	See Archive 12.71 and (A. Lucchini, J. Algebra, 245 (2001), 552-561).
	
	Comment of 2021: the answer is yes if $A$ and $B$ have prime power order (E. Detomi, A. Lucchini, J. London Math. Soc. (2), 87, no. 3 (2013), 689-706). A. Lucchini
	\par\smallskip\hrule\smallskip
\end{openproblem}

\plainproblemsection

\begin{openproblem}
	16.54. We say that a group $G$ acts freely on a group $V$ if ${vg} \neq  v$ for any nontrivial elements $g \in  G, v \in  V$ . Is it true that a group $G$ that can act freely on a non-trivial abelian group is embeddable in the multiplicative group of some skew-field?
	
	No, it is not. For example, the group $2.{A}_{5}{.2}$ with a quaternion Sylow 2-subgroup can act freely on an elementary abelian group of order ${7}^{4}$ but is not embeddable in the multiplicative group of any skew-field by Theorem 7 in (S. A. Amitsur, Trans. Amer. Math. Soc., 80, no. 2 (1955), 361-386). (D. Nedrenko, Letter of 20 January 2014).
\end{openproblem}

\plainproblemsection

\begin{openproblem}
	16.55. (Well-known problem). Let $V$ be a faithful absolutely irreducible module for a finite group $G$ . Is it true that $\dim {H}^{1}\left( {G, V}\right)  \leq  2$ ? V. D. Mazurov No, it is not (L.L. Scott, J. Algebra, 260 (2003), 416-425; see also J. N. Bray, R. A. Wilson, J. Group Theory, 11 (2008), 669-673).
\end{openproblem}

\plainproblemsection

\begin{openproblem}
	16.56. The spectrum $\omega \left( G\right)$ of a group $G$ is the set of orders of elements of $G$ . Suppose that $\omega \left( G\right)  = \{ 1,2,3,4,5,6\}$ . Is $G$ locally finite? V. D. Mazurov
\end{openproblem}

\plainproblemsection

\begin{openproblem}
	16.57. Suppose that $\omega \left( G\right)  = \omega \left( {{L}_{2}\left( 7\right) }\right)  = \{ 1,2,3,4,7\}$ . Is $G \cong  {L}_{2}\left( 7\right)$ ? This is true for finite $G$ . V. D. Mazurov
	
	Yes, it is (D. V. Lytkina, A. A. Kuznetsov, Siberian Electr. Math. Rep., 4 (2007), 136-140; http://semr.math.nsc.ru).
\end{openproblem}

\plainproblemsection

\begin{openproblem}
	16.58. Is $S{U}_{2}\left( \mathbb{C}\right)$ the only group that has just one irreducible complex representation of dimension $n$ for each $n = 1,2,\ldots$ ?
	
	(If $R\left\lbrack  n\right\rbrack$ is the $n$ -dimensional irreducible complex representation of $S{U}_{2}\left( \mathbb{C}\right)$ , then $R\left\lbrack  2\right\rbrack$ is the natural two-dimensional representation, and $R\left\lbrack  2\right\rbrack   \otimes  R\left\lbrack  n\right\rbrack   = R\left\lbrack  {n - 1}\right\rbrack   + R\left\lbrack  {n + 1}\right\rbrack$ for $n > 1$ .) J. McKay
	
	No, it is not. It is known that $\mathbb{C}$ has infinitely many (discrete) automorphisms. For $\varphi  \in  \operatorname{Aut}\left( \mathbb{C}\right)$ and a matrix $x$ , let ${x}^{\varphi }$ denote the matrix obtained from $x$ by applying $\varphi$ to each element. Then ${T}_{\varphi } : x \mapsto  {x}^{\varphi }$ is an irreducible representation of $S{U}_{2}\left( \mathbb{C}\right)$ . It is easy to show that among these representations there are infinitely many pairwise non-equivalent ones. Therefore the group $S{U}_{2}\left( \mathbb{C}\right)$ itself does not satisfy the condition of the problem. (This observation belongs to von Neumann.) One can show that any group satisfying the condition of the problem is isomorphic to a subgroup of $S{L}_{2}\left( \mathbb{Q}\right)$ satisfying the condition of Problem 15.57. On the other hand, E. Cartan proved that a unitary representation of a simple compact Lie group is always continuous. Hence, if we restrict ourselves to the unitary representations, then $S{U}_{2}\left( \mathbb{C}\right)$ does satisfy the condition of the problem. (V. P. Burichenko, Letter of 16 July 2013.)
\end{openproblem}

\plainproblemsection

\begin{openproblem}
	*16.59. Given a finite group $K$ , does there exist a finite group $G$ such that $K \cong$ Out $G = \operatorname{Aut}G/\operatorname{Inn}G$ ? (It is known that an infinite group $G$ exists with this property.) D. MacHale
	
	*Yes, it does (Y. Cornulier, https://mathoverflow.net/questions/372480/is-every-finite-group-the-outer-automorphism-group-of-a-finite-group/372563).
\end{openproblem}

\plainproblemsection

\begin{openproblem}
	16.60. If $G$ is a finite group, let $T\left( G\right)$ be the sum of the degrees of the irreducible complex representations of $G, T\left( G\right)  = \mathop{\sum }\limits_{{i = 1}}^{{k\left( G\right) }}{d}_{i}$ , where $G$ has $k\left( G\right)$ conjugacy classes. If $\alpha  \in  \operatorname{Aut}G$ , let ${S}_{\alpha } = \left\{  {g \in  G \mid  \alpha \left( g\right)  = {g}^{-1}}\right\}$ . Is it true that $T\left( G\right)  \geq  \left| {S}_{\alpha }\right|$ for all $\alpha  \in  \operatorname{Aut}G$ ? D. MacHale
\end{openproblem}

\plainproblemsection

\begin{openproblem}
	16.61. A subgroup $H$ of a group $G$ is fully invariant if $\vartheta \left( H\right)  \leq  H$ for every endomorphism $\vartheta$ of $G$ . Let $G$ be a finite group such that $G$ has a fully invariant subgroup of order $d$ for every $d$ dividing $\left| G\right|$ . Must $G$ be cyclic? D. MacHale
	
	No: take $G = {C}_{p} \times  {C}_{{p}^{2}}$ , where ${C}_{p},{C}_{{p}^{2}}$ are cyclic $p$ -groups of orders $p,{p}^{2}$ (A. Abdol-lahi, Letter of 4 March 2009).
\end{openproblem}

\plainproblemsection

\begin{openproblem}
	16.62. Let $G$ be a group such that every $\alpha  \in  \operatorname{Aut}G$ fixes every conjugacy class of $G$ (setwise). Must Aut $G = \operatorname{Inn}G$ ? D. MacHale
	
	Not necessarily: the paper (H. Heineken, Arch. Math. (Basel), 33 (1979/80), no. 6, 497-503), in particular, produces finite $p$ -groups all of whose automorphisms preserve all the conjugacy classes, while by Gaschütz’ theorem every finite $p$ -group has outer automorphisms (A. Mann, Letter of 20 April 2006).
\end{openproblem}

\plainproblemsection

\begin{openproblem}
	16.63. Is there a non-trivial finite $p$ -group $G$ of odd order such that $\left| {\operatorname{Aut}G}\right|  = \left| G\right|$ ? See also Archive 12.77. D. MacHale
\end{openproblem}

\plainproblemsection

\begin{openproblem}
	16.64. A non-abelian variety in which all finite groups are abelian is called pseudo-abelian. A group variety is called a $t$ -variety if for all groups in this variety the relation of being a normal subgroup is transitive. By (O. Macedońska, A. Storozhev, Commun. Algebra,25, no. 5 (1997),1589-1593) each non-abelian $t$ -variety is pseudo-abelian, and the pseudo-abelian varieties constructed in (A. Yu. Ol'shanskiĭ, Math. USSR-Sb., 54 (1986), 57-80) are $t$ -varieties. Is every pseudo-abelian variety a $t$ -variety?
	
	O. Macedońska
\end{openproblem}

\plainproblemsection

\begin{openproblem}
	16.65. Does there exist a finitely presented residually torsion-free nilpotent group with a free presentation $G = F/R$ such that the group $F/\left\lbrack  {F, R}\right\rbrack$ is not residually nilpotent? R. Mikhailov, I. B. S. Passi
	
	Yes, it exists (V. V. Bludov, J. Group Theory, 12, no. 4 (2009), 579-590).
\end{openproblem}

\plainproblemsection

\begin{openproblem}
	16.66. For a group $G$ , let ${D}_{n}\left( G\right)$ denote the $n$ -th dimension subgroup of $G$ , and ${\zeta }_{n}\left( G\right)$ the $n$ -th term of its upper central series. For a given integer $n \geq  1$ , let $f\left( n\right)  = \max \{ m \mid  \exists$ a nilpotent group $G$ of class $n$ with ${D}_{m}\left( G\right)  \neq  1\}$ and $g\left( n\right)  = \max \left\{  {m\mid \exists \text{ a nilpotent group }G}\right.$ of class $\left. {n\text{ such that }{D}_{n}\left( G\right)  \nsubseteq  {\zeta }_{m}\left( G\right) }\right\}$ .
	
	a) What is $f\left( 3\right)$ ?
	
	b) (B.I.Plotkin). Is it true that $f\left( n\right)$ is finite for all $n$ ?
	
	c) Is the growth of $f\left( n\right)$ and $g\left( n\right)$ polynomial, exponential, or intermediary?
	
	It is known that both $f\left( n\right)  - n$ and $g\left( n\right)$ tend to infinity as $n \rightarrow  \infty$ (N. Gupta, Yu. V. Kuz'min, J. Pure Appl. Algebra, 104 (1995), 191-197).
	
	R. Mikhailov, I. B. S. Passi
\end{openproblem}

\plainproblemsection

\begin{openproblem}
	16.67. Conjecture: Given any integer $k$ , there exists an integer ${n}_{0} = {n}_{0}\left( k\right)$ such that if $n \geq  {n}_{0}$ then the symmetric group of degree $n$ has at least $k$ different ordinary irreducible characters of equal degrees. A. Moretó
	
	This is proved (D. Craven, Proc. London Math. Soc., 96 (2008), 26-50).
\end{openproblem}

\plainproblemsection

\begin{openproblem}
	16.68. Let $W\left( {x, y}\right)$ be a non-trivial reduced group word, and $G$ one of the groups ${PSL}\left( {2,\mathbb{R}}\right) ,{PSL}\left( {2,\mathbb{C}}\right)$ , or ${SO}\left( {3,\mathbb{R}}\right)$ . Are all the maps $W : G \times  G \rightarrow  G$ surjective?
	
	For ${SL}\left( {2,\mathbb{R}}\right) ,{SL}\left( {2,\mathbb{C}}\right) ,{GL}\left( {2,\mathbb{C}}\right)$ , and for the group ${S}^{3}$ of quaternions of norm 1 there exist non-surjective words; see (J. Mycielski, Amer. Math. Monthly, 84 (1977), 723-726; 85 (1978), 263-264). J. Mycielski
\end{openproblem}

\plainproblemsection

\begin{openproblem}
	16.69. Let $W$ be as above.
	
	a) Must the range of the function $\operatorname{Tr}\left( {W\left( {x, y}\right) }\right)$ for $x, y \in  {GL}\left( {2,\mathbb{R}}\right)$ include the interval $\lbrack  - 2, + \infty )$ ?
	
	b) For $x, y$ being non-zero quaternions, must the range of the function $\operatorname{Re}\left( {W\left( {x, y}\right) }\right)$ include the interval $\left\lbrack  {-5/{27},1}\right\rbrack$ ?
\end{openproblem}

\plainproblemsection

\begin{openproblem}
	16.70. Suppose that a finitely generated group $G$ acts freely on an $A$ -tree, where $A$ is an ordered abelian group. Is it true that $G$ acts freely on a ${\mathbb{Z}}^{n}$ -tree for some $n$ ?
	
	Comment of 2013: It was proved that every finitely presented group acting freely on an $A$ -tree acts freely on some ${\mathbb{R}}^{n}$ -tree for a suitable $n$ , where $\mathbb{R}$ has the lexicographical order (O. Kharlampovich, A. Myasnikov, D. Serbin, Int. J. Algebra Comput., 23 (2013), 325-345).
\end{openproblem}

\plainproblemsection

\begin{openproblem}
	16.71. Is the elementary theory of a torsion-free hyperbolic group decidable? A. G. Myasnikov, O. G. Kharlampovich
	
	Yes, it is (O. Kharlampovich, A. Myasnikov, https://arxiv.org/pdf/1303.0760.pdf).
\end{openproblem}

\plainproblemsection

\begin{openproblem}
	16.72. Does there exist an exponential-time algorithm for obtaining a JSJ-decomposition of a finitely generated fully residually free group?
	
	Some algorithm was found in (O. Kharlampovich, A. Myasnikov, in: Contemp. Math. AMS (Algorithms, Languages, Logic), 378 (2005), 87-212).
	
	A. G. Myasnikov, O. G. Kharlampovich
\end{openproblem}

\plainproblemsection

\begin{openproblem}
	16.73. a) Let $G$ be a group generated by a finite set $S$ , and let $l\left( g\right)$ denote the word length function of $g \in  G$ with respect to $S$ . The group $G$ is said to be contracting if there exist a faithful action of $G$ on the set ${X}^{ * }$ of finite words over a finite alphabet $X$ and constants $0 < \lambda  < 1$ and $C > 0$ such that for every $g \in  G$ and $x \in  X$ there exist $h \in  G$ and $y \in  X$ such that $l\left( h\right)  < {\lambda l}\left( g\right)  + C$ and $g\left( {xw}\right)  = {yh}\left( w\right)$ for all $w \in  {X}^{ * }$ .
	
	Can a contracting group have a non-abelian free subgroup? V. V. Nekrashevich a) No, it cannot (V. V. Nekrashevich, Groups Geom. Dynam., 4, no. 4 (2010), 847- 862).
\end{openproblem}

\plainproblemsection

\begin{openproblem}
	16.73. b) Let $G$ be a group generated by a finite set $S$ , and let $l\left( g\right)$ denote the word length function of $g \in  G$ with respect to $S$ . The group $G$ is said to be contracting if there exist a faithful action of $G$ on the set ${X}^{ * }$ of finite words over a finite alphabet $X$ and constants $0 < \lambda  < 1$ and $C > 0$ such that for every $g \in  G$ and $x \in  X$ there exist $h \in  G$ and $y \in  X$ such that $l\left( h\right)  < {\lambda l}\left( g\right)  + C$ and $g\left( {xw}\right)  = {yh}\left( w\right)$ for all $w \in  {X}^{ * }$ .
	
	Do there exist non-amenable contracting groups? V. V. Nekrashevich
\end{openproblem}

\plainproblemsection

\begin{openproblem}
	16.74. a) Let $G = \langle \alpha ,\beta \rangle$ be the group generated by the following two permutations of $\mathbb{Z} : \alpha \left( n\right)  = n + 1;\beta \left( 0\right)  = 0,\beta \left( {{2}^{k}m}\right)  = {2}^{k}\left( {m + 2}\right)$ , where $m$ is odd and $n$ is a positive integer. Is $G$ amenable? V. V. Nekrashevich
	
	a) Yes, it is (G. Amir, O. Angel, B. Virág, J. Eur. Math. Soc., 15, no. 3 (2013), 705-730).
\end{openproblem}

\plainproblemsection

\begin{openproblem}
	16.74. b) Is it true that all groups generated by automata of polynomial growth in the sense of S. Sidki (Geom. Dedicata, 108 (2004), 193-204) are amenable?
	
	V. V. Nekrashevich
\end{openproblem}

\plainproblemsection

\begin{openproblem}
	16.75. Can a non-abelian one-relator group be the group of all automorphisms of some group? M. V. Neshchadim
	
	Yes, it can: in (D. J. Collins, Proc. London Math. Soc. (3), 36 (1978), no. 3, 480-493) it was proved that for integers $r, s$ such that $\left( {r, s}\right)  = 1,\left| r\right|  \neq  \left| s\right|$ , and $r - s$ is even, the group $G = \left\langle  {a, b \mid  {a}^{-1}{b}^{r}a = {b}^{s}}\right\rangle$ is isomorphic to $\operatorname{Aut}\left( {\operatorname{Aut}\left( G\right) }\right)$ (V. A. Churkin, Letter of 5 April 2006).
\end{openproblem}

\plainproblemsection

\begin{openproblem}
	16.76. We call a group $G$ strictly real if each of its non-trivial elements is conjugate to its inverse by some involution in $G$ . In which groups of Lie type over a field of characteristic 2 the maximal unipotent subgroups are strictly real? Ya. N. Nuzhin
\end{openproblem}

\plainproblemsection

\begin{openproblem}
	16.77. It is known that in every Noetherian group the nilpotent radical coincides with the collection of all Engel elements (R. Baer, Math. Ann., 133 (1957), 256-270; B. I. Plotkin, Izv. Vysš. Učebn. Zaved. Mat., 1958, no. 1(2), 130-135 (Russian)). It would be nice to find a similar characterization of the solvable radical of a finite group. More precisely, let $u = u\left( {x, y}\right)$ be a sequence of words satisfying 15.75 . We say that an element $g \in  G$ is $u$ -Engel if there exists $n = n\left( g\right)$ such that ${u}_{n}\left( {x, g}\right)  = 1$ for every element $x \in  G$ . Does there exist a sequence $u = u\left( {x, y}\right)$ such that the solvable radical of a finite group coincides with the set of all $u$ -Engel elements? B. I. Plotkin
\end{openproblem}

\plainproblemsection

\begin{openproblem}
	16.78. Do there exist linear non-abelian simple groups without involutions?
	
	B. Poizat
\end{openproblem}

\plainproblemsection

\begin{openproblem}
	*16.79. Is it true that in any finitely generated ${AT}$ -group over a sequence of cyclic groups of uniformly bounded orders all Sylow subgroups are locally finite? For the definition of an ${AT}$ -group see (A. V. Rozhkov, Math. Notes,40 (1986),827-836)
	\par\smallskip\hrule\smallskip
	*No, it is not true (A. V. Rozhkov, in Group theory and its applications, Proc. XXII School-Conf. on Group Theory, Kuban' Univ., Krasnodar, 2018, 126-131 (Russian)).
	\par\smallskip\hrule\smallskip
\end{openproblem}

\plainproblemsection

\begin{openproblem}
	16.80. Suppose that a group $G$ is obtained from the free product of torsion-free groups ${A}_{1},\ldots ,{A}_{n}$ by imposing $m$ additional relations, where $m < n$ . Is it true that the free product of some $n - m$ of the ${A}_{i}$ embeds into $G$ ? N. S. Romanovskii
\end{openproblem}

\plainproblemsection

\begin{openproblem}
	16.82. Let $\mathcal{X}$ be a non-empty class of finite groups closed under taking homomorphic images, subgroups, and direct products. With every group $G \in  \mathcal{X}$ we associate some set $\tau \left( G\right)$ of subgroups of $G$ . We say that $\tau$ is a subgroup functor on $\mathcal{X}$ if:
	
	1) $G \in  \tau \left( G\right)$ for all $G \in  \mathcal{X}$ , and
	
	2) for each epimorphism $\varphi  : A \mapsto  B$ , where $A, B \in  \mathcal{X}$ , and for any $H \in  \tau \left( A\right)$ and
	
	$T \in  \tau \left( B\right)$ we have ${H}^{\varphi } \in  \tau \left( B\right)$ and ${T}^{{\varphi }^{-1}} \in  \tau \left( A\right)$ .
	
	A subgroup functor $\tau$ is closed if for each group $G \in  \mathcal{X}$ and for every subgroup $H \in  \mathcal{X} \cap  \tau \left( G\right)$ we have $\tau \left( H\right)  \subseteq  \tau \left( G\right)$ . The set $F\left( \mathcal{X}\right)$ consisting of all closed subgroup functors on $\mathcal{X}$ is a lattice (in which ${\tau }_{1} \leq  {\tau }_{2}$ if and only if ${\tau }_{1}\left( G\right)  \subseteq  {\tau }_{2}\left( G\right)$ for every group $G \in  \mathcal{X}$ ). It is known that $F\left( \mathcal{X}\right)$ is a chain if and only if $\mathcal{X}$ a class of $p$ -groups for some prime $p$ (Theorem 1.5.17 in S. F. Kamornikov and M. V. Sel’kin, Subgroups functors and classes of finite groups, Belaruskaya Navuka, Minsk, 2001 (Russian)).
	
	Is there a non-nilpotent class $\mathcal{X}$ such that the width of the lattice $F\left( \mathcal{X}\right)$ is at most $\left| {\pi \left( \mathcal{X}\right) }\right|$ where $\pi \left( \mathcal{X}\right)$ is the set of all prime divisors of the orders of the groups in X? A. N. Skiba
	
	No such classes exist (S. F. Kamornikov, Siberian Math. J., 51, no. 5 (2010), 824-829).
\end{openproblem}

\plainproblemsection

\begin{openproblem}
	16.83. a) Let ${E}_{n}$ be a free locally nilpotent $n$ -Engel group on countably many generators, and let $\pi \left( {E}_{n}\right)$ be the set of prime divisors of the orders of elements of the periodic part of ${E}_{n}$ . It is known that $2,3,5 \in  \pi \left( {E}_{4}\right)$ .
	
	Does there exist $n$ for which $7 \in  \pi \left( {E}_{n}\right)$ ? Yu. V. Sosnovskii
	
	a) Yes, $n = 6$ , since the 2-generator free nilpotent 6-Engel group has an element of order 7 (W. Nickel, J. Austral. Math. Soc., Ser. A, 67, no. 2 (1999), 214-222; http:// www.mathematik.tu-darmstadt.de/~nickel/Engel/Engel.html). (A. Abdollahi, Letter of 19 August 2011.)
\end{openproblem}

\plainproblemsection

\begin{openproblem}
	16.83. b) Let ${E}_{n}$ be a free locally nilpotent $n$ -Engel group on countably many generators, and let $\pi \left( {E}_{n}\right)$ be the set of prime divisors of the orders of elements of the periodic part of ${E}_{n}$ . It is known that $2,3,5 \in  \pi \left( {E}_{4}\right)$ .
	
	Is it true that $\pi \left( {E}_{n}\right)  = \pi \left( {E}_{n + 1}\right)$ for all sufficiently large $n$ ? Yu. V. Sosnovskii
\end{openproblem}

\plainproblemsection

\begin{openproblem}
	16.84. Can the braid group ${\mathfrak{B}}_{n}, n \geq  4$ , act faithfully on a regular rooted tree by finite-state automorphisms? Such action is known for ${\mathfrak{B}}_{3}$ .
	
	See the definitions in (R. I. Grigorchuk, V. V. Nekrashevich, V. I. Sushchanskii, Proc. Steklov Inst. Math., 2000, no. 4 (231), 128-203). V. I. Sushchanskiǐ
\end{openproblem}

\plainproblemsection

\begin{openproblem}
	*16.85. Suppose that groups $G, H$ act faithfully on a regular rooted tree by finite-state automorphisms. Can their free product $G * H$ act faithfully on a regular rooted tree by finite state automorphisms? V. I. Sushchanskiĭ
	
	*Yes, it can (M. Fedorova, A. Oliynyk, Algebra Discrete Math., 23, no. 2 (2017), 230-236).
\end{openproblem}

\plainproblemsection

\begin{openproblem}
	*16.86. Does the group of all finite-state automorphisms of a regular rooted tree possess an irreducible system of generators? V. I. Sushchanskii
	
	*Yes, it does (Ya. Lavrenyuk, Geometriae Dedicata, 183, no. 1 (2016), 59-67).
\end{openproblem}

\plainproblemsection

\begin{openproblem}
	16.87. Let $\mathfrak{M}$ be a variety of groups and let ${G}_{r}$ be a free $r$ -generator group in $\mathfrak{M}$ . A subset $S \subseteq  {G}_{r}$ is called a test set if every endomorphism of ${G}_{r}$ identical on $S$ is an automorphism. The minimum of the cardinalities of test sets is called the test rank of ${G}_{r}$ . Suppose that the test rank of ${G}_{r}$ is $r$ for every $r \geq  1$ .
	
	a) Is it true that $\mathfrak{M}$ is an abelian variety?
	
	b) Suppose that $\mathfrak{M}$ is not a periodic variety. Is it true that $\mathfrak{M}$ is the variety of all abelian groups? E. I. Timoshenko
\end{openproblem}

\plainproblemsection

\begin{openproblem}
	16.88. (G.M.Bergman). The width of a group $G$ with respect to a generating set $X$ means the supremum over all $g \in  G$ of the least length of a group word in elements of $X$ expressing $g$ . A group $G$ has finite width if the width of $G$ with respect to any generating set is finite. Does there exist a countably infinite group of finite width? All known infinite groups of finite width (infinite permutation groups, infinite-dimensional general linear groups, and some other groups) are uncountable.
\end{openproblem}

\plainproblemsection

\begin{openproblem}
	16.89. (G. M. Bergman). Is it true that the automorphism group of an infinitely generated free group $F$ has finite width? The answer is affirmative if $F$ is countably generated. V. Tolstykh
\end{openproblem}

\plainproblemsection

\begin{openproblem}
	16.90. Is it true that the automorphism group Aut $F$ of an infinitely generated free group $F$ is
	
	a) the normal closure of a single element?
	
	b) the normal closure of some involution in Aut $F$ ?
	
	For a free group ${F}_{n}$ of finite rank $n$ M. Bridson and K. Vogtmann have recently shown that Aut ${F}_{n}$ is the normal closure of some involution, which permutes some basis of ${F}_{n}$ . It is also known that the automorphism groups of infinitely generated free nilpotent groups have such involutions. V. Tolstykh
\end{openproblem}

\plainproblemsection

\begin{openproblem}
	16.91. Let $F$ be an infinitely generated free group. Is there an ${IA}$ -automorphism of $F$ whose normal closure in Aut $F$ is the group of all ${IA}$ -automorphisms of $F$ ?
	
	V. Tolstykh
\end{openproblem}

\plainproblemsection

\begin{openproblem}
	16.92. Let $F$ be an infinitely generated free group. Is Aut $F$ equal to its derived subgroup? This is true if $F$ is countably generated (R. Bryant, V. A. Roman’kov, J. Algebra, 209 (1998), 713-723). V. Tolstykh
\end{openproblem}

\plainproblemsection

\begin{openproblem}
	16.93. Let ${F}_{n}$ be a free group of finite rank $n \geq  2$ . Is the group $\operatorname{Inn}{F}_{n}$ of inner automorphisms of ${F}_{n}$ a first-order definable subgroup of Aut ${F}_{n}$ ? It is known that the set of inner automorphisms induced by powers of primitive elements is definable in Aut ${F}_{n}$ . V. Tolstykh
\end{openproblem}

\plainproblemsection

\begin{openproblem}
	16.94. If $G = \left\lbrack  {G, G}\right\rbrack$ , then the commutator width of the group $G$ is its width relative to the set of commutators. Let $V$ be an infinite-dimensional vector space over a division ring. It is known that the commutator width of ${GL}\left( V\right)$ is finite. Is it true that the commutator width of ${GL}\left( V\right)$ is one? V. Tolstykh
\end{openproblem}

\plainproblemsection

\begin{openproblem}
	16.95. Conjecture: If $F$ is a field and $A$ is in ${GL}\left( {n, F}\right)$ , then there is a permutation matrix $P$ such that ${AP}$ is cyclic, that is, the minimal polynomial of ${AP}$ is also its characteristic polynomial. J. G. Thompson
\end{openproblem}

\plainproblemsection

\begin{openproblem}
	16.96. Let $G$ be a locally finite $n$ -Engel $p$ -group where $p$ a prime greater than $n$ . Is $G$ then a Fitting group? (Examples of N. Gupta and F. Levin show that the condition $p > n$ is necessary in general.) G. Traustason
\end{openproblem}

\plainproblemsection

\begin{openproblem}
	16.97. Let $G$ be a torsion-free group with all subgroups subnormal of defect at most $n$ . Must $G$ then be nilpotent of class at most $n$ ? (This is known to be true for $n < 5)$ . $G$ . Traustason
\end{openproblem}

\plainproblemsection

\begin{openproblem}
	16.98. Suppose that $G$ is a solvable finite group and $A$ is a group of automorphisms of $G$ of relatively prime order. Is there a bound for the Fitting height $h\left( G\right)$ of $G$ in terms of $A$ and $h\left( {{C}_{G}\left( A\right) }\right)$ , or even in terms of the length $l\left( A\right)$ of the longest chain of nested subgroups of $A$ and $h\left( {{C}_{G}\left( A\right) }\right)$ ?
	
	When $A$ is solvable, it is proved in (A. Turull, J. Algebra,86 (1984),555-566) that $h\left( G\right)  \leq  h\left( {{C}_{G}\left( A\right) }\right)  + {2l}\left( A\right)$ and this bound is best possible for $h\left( {{C}_{G}\left( A\right) }\right)  > 0$ . For $A$ non-solvable some results are in (H. Kurzweil, Manuscripta Math.,41(1983), 233-305). A. Turull
\end{openproblem}

\plainproblemsection

\begin{openproblem}
	16.99. Suppose that $G$ is a finite solvable group, $A \leq  \operatorname{Aut}G,{C}_{G}\left( A\right)  = 1$ , the orders of $G$ and $A$ are coprime, and let $l\left( A\right)$ be the length of the longest chain of nested subgroups of $A$ . Is the Fitting height of $G$ bounded above by $l\left( A\right)$ ?
	
	For $A$ solvable the question coincides with 5.30. It is proved that for any finite group $A$ , first, there exist $G$ with $h\left( G\right)  = l\left( A\right)$ and, second, there is a finite set of primes $\pi$ (depending on $A$ ) such that if $\left| G\right|$ is coprime to each prime in $\pi$ , then $h\left( G\right)  \leq  l\left( A\right)$ . See (A. Turull, Math. Z.,187 (1984),491-503). A. Turull
\end{openproblem}

\plainproblemsection

\begin{openproblem}
	16.100. Is there an (infinite) 2-generator simple group $G$ such that Aut ${F}_{2}$ is transitive on the set of normal subgroups $N$ of the free group ${F}_{2}$ such that ${F}_{2}/N \cong  G$ ? Cf. 6.45.
	
	J. Wiegold
\end{openproblem}

\plainproblemsection

\begin{openproblem}
	16.101. Do there exist uncountably many infinite 2-groups that are quotients of the group $\left\langle  {x, y \mid  {x}^{2} = {y}^{4} = {\left( xy\right) }^{8} = 1}\right\rangle$ ? There certainly exists one, namely the subgroup of finite index in Grigorchuk’s first group generated by $b$ and ${ad}$ ; see (R. I. Grigorchuk, Functional Anal. Appl., 14 (1980), 41-43). J. Wiegold
	
	Yes, there exist (R. I. Grigorchuk, Algebra Discrete Math., 2009, no. 4 (2009), 78- 96 (subm. 25 November 2009); A. Minasyan, A. Yu. Olshanskii, D. Sonkin, Groups Geom. Dynam., 3, no. 3 (2009) 423-452 (subm. 21 April 2008)).
\end{openproblem}

\plainproblemsection

\begin{openproblem}
	16.102. We say that a group $G$ is rational if any two elements $x, y \in  G$ satisfying $\langle x\rangle  = \langle y\rangle$ are conjugate. Is it true that for any $d$ there exist only finitely many finite rational $d$ -generated 2-groups? A. Jaikin-Zapirain
\end{openproblem}

\plainproblemsection

\begin{openproblem}
	16.103. Is there a rank analogue of the Leedham-Green-McKay-Shepherd theorem on $p$ -groups of maximal class? More precisely, suppose that $P$ is a 2-generator finite $p$ -group whose lower central quotients ${\gamma }_{i}\left( P\right) /{\gamma }_{i + 1}\left( P\right)$ are cyclic for all $i \geq  2$ . Is it true that $P$ contains a normal subgroup $N$ of nilpotency class $\leq  2$ such that the rank of $P/N$ is bounded in terms of $p$ only? E. I. Khukhro
	
	No, moreover, there are no functions $d\left( p\right)$ and $r\left( p\right)$ such that a group with these properties would necessarily have a normal subgroup of derived length $\leq  d\left( p\right)$ with quotient of rank $\leq  r\left( p\right)$ (E. I. Khukhro, Siber. Math. J.,54, no. 1 (2013),174-184).
\end{openproblem}

\plainproblemsection

\begin{openproblem}
	*16.104. If $G$ is a finite group, then every element $a$ of the rational group algebra $\mathbb{Q}\left\lbrack  G\right\rbrack$ has a unique Jordan decomposition $a = {a}_{s} + {a}_{n}$ , where ${a}_{n} \in  \mathbb{Q}\left\lbrack  G\right\rbrack$ is nilpotent, ${a}_{s} \in  \mathbb{Q}\left\lbrack  G\right\rbrack$ is semisimple over $\mathbb{Q}$ , and ${a}_{s}{a}_{n} = {a}_{n}{a}_{s}$ . The integral group ring $\mathbb{Z}\left\lbrack  G\right\rbrack$ is said to have the additive Jordan decomposition property (AJD) if ${a}_{s},{a}_{n} \in  \mathbb{Z}\left\lbrack  G\right\rbrack$ for every $a \in  \mathbb{Z}\left\lbrack  G\right\rbrack$ . If $a \in  \mathbb{Q}\left\lbrack  G\right\rbrack$ is invertible, then ${a}_{s}$ is also invertible and so $a = {a}_{s}{a}_{u}$ with ${a}_{u} = 1 + {a}_{s}^{-1}{a}_{n}$ unipotent and ${a}_{s}{a}_{u} = {a}_{u}{a}_{s}$ . Such a decomposition is again unique. We say that $\mathbb{Z}\left\lbrack  G\right\rbrack$ has multiplicative Jordan decomposition property (MJD) if ${a}_{s},{a}_{u} \in  \mathbb{Z}\left\lbrack  G\right\rbrack$ for every invertible $a \in  \mathbb{Z}\left\lbrack  G\right\rbrack$ . See the survey (A. W. Hales, I. B. S. Passi, in: Algebra, Some Recent Advances, Birkhäuser, Basel, 1999, 75-87).
	
	Is it true that there are only finitely many isomorphism classes of finite 2-groups $G$ such that $\mathbb{Z}\left\lbrack  G\right\rbrack$ has MJD but not AJD? A. W. Hales, I. B. S. Passi
	
	*Yes, it is true (A. W. Hales, I. B. S. Passi, L. E. Wilson, J. Algebra, 316, no. 1 (2007), 109-132; 371 (2012), 665-666).
\end{openproblem}

\plainproblemsection

\begin{openproblem}
	16.105. Is it true that a locally graded group which is a product of two almost polycyclic subgroups (equivalently, of two almost soluble subgroups with the maximal condition) is almost polycyclic?
	
	By (N. S. Chernikov, Ukrain. Math. J., 32 (1980) 476-479) a locally graded group which is a product of two Chernikov subgroups (equivalently, of two almost soluble subgroups with the minimal condition) is Chernikov. N. S. Chernikov
\end{openproblem}

\plainproblemsection

\begin{openproblem}
	16.106. Let ${\pi }_{e}\left( G\right)$ denote the set of orders of elements of a group $G$ , and $h\left( \Gamma \right)$ the number of non-isomorphic finite groups $G$ with ${\pi }_{e}\left( G\right)  = \Gamma$ . Do there exist two finite groups ${G}_{1},{G}_{2}$ such that ${\pi }_{e}\left( {G}_{1}\right)  = {\pi }_{e}\left( {G}_{2}\right) , h\left( {{\pi }_{e}\left( {G}_{1}\right) }\right)  < \infty$ , and neither of the two groups ${G}_{1},{G}_{2}$ is isomorphic to a subgroup or a quotient of a normal subgroup of the other? W. J. Shi
	
	Yes, there exist: for example, ${G}_{1} = {L}_{15}\left( {2}^{60}\right) {.3}$ and ${G}_{2} = {L}_{15}\left( {2}^{60}\right) {.5}$ (M. A. Grechko-seeva, Algebra and Logic, 47, no. 4 (2008), 229-241).
\end{openproblem}

\plainproblemsection

\begin{openproblem}
	16.107. Is it true that almost every alternating group ${A}_{n}$ is uniquely determined in the class of finite groups by its set of element orders, i.e. that $h\left( {{\pi }_{e}\left( {A}_{n}\right) }\right)  = 1$ for all large enough $n$ ? W. J. Shi
	
	Yes, it is (I. B. Gorshkov, Algebra and Logic, 52, no. 1 (2013), 41-45).
\end{openproblem}

\plainproblemsection

\begin{openproblem}
	16.108. Do braid groups ${B}_{n}, n > 4$ , have non-elementary hyperbolic factor groups?
	
	V. E. Shpilrain
\end{openproblem}

\plainproblemsection

\begin{openproblem}
	16.109. Is there a polynomial time algorithm for solving the word problem in the group Aut ${F}_{n}$ (with respect to some particular finite presentation), where ${F}_{n}$ is the free group of rank $n \geq  2$ ? V. E. Shpilrain
	
	Yes, there is (S. Schleimer, Comment. Math. Helv., 83 (2008), 741-765).
\end{openproblem}

\plainproblemsection

\begin{openproblem}
	16.110. (I. Kapovich, P. Schupp). Is there an algorithm which, when given two elements $u, v$ of a free group ${F}_{n}$ , decides whether or not the cyclic length of $\phi \left( u\right)$ equals the cyclic length of $\phi \left( v\right)$ for every automorphism $\phi$ of ${F}_{n}$ ?
	
	Comment of 2009: The answer was shown to be positive for $n = 2$ in (D. Lee, J. Group Theory, 10 (2007), 561-569). V. E. Shpilrain
\end{openproblem}

\plainproblemsection

\begin{openproblem}
	16.111. Must an infinite simple periodic group with a dihedral Sylow 2-subgroup be isomorphic to ${L}_{2}\left( P\right)$ for a locally finite field $P$ of odd characteristic? $V$ , V. P. Shunkov
\end{openproblem}

\plainproblemsection

\begin{openproblem}
	17.1. (I. M. Isaacs). Does there exist a finite group partitioned by subgroups of equal order not all of which are abelian? (Cf. 15.26.) It is known that such a group must be of prime exponent (I. M. Isaacs, Pacific J. Math., 49 (1973), 109-116). A. Abdollahi
\end{openproblem}

\plainproblemsection

\begin{openproblem}
	17.2. (P. Schmid). Does there exist a finite non-abelian $p$ -group $G$ such that ${H}^{1}\left( {G/\Phi \left( G\right) , Z\left( {\Phi \left( G\right) }\right) }\right)  = 0?$ A. Abdollahi
	
	Yes, it does (A. Abdollahi, J. Algebra, 342 (2011), 154-160).
\end{openproblem}

\plainproblemsection

\begin{openproblem}
	*17.3. Let $G$ be a group in which every 4-element subset contains two elements generating a nilpotent subgroup. Is it true that every 2-generated subgroup of $G$ is nilpotent? A. Abdollahi
	
	*No, not always (A.I. Sozutov, Siberian Math. J., 60, no. 6 (2019), 1099-1100).
\end{openproblem}

\plainproblemsection

\begin{openproblem}
	17.4. Let $x$ be a right 4-Engel element of a group $G$ .
	
	a) Is it true that the normal closure $\langle x{\rangle }^{G}$ of $x$ in $G$ is nilpotent if $G$ is locally nilpotent?
	
	b) If the answer to a) is affirmative, is there a bound on the nilpotency class of $\langle x{\rangle }^{G}$ ?
	
	c) Is it true that $\langle x{\rangle }^{G}$ is always nilpotent? A. Abdollahi
\end{openproblem}

\plainproblemsection

\begin{openproblem}
	17.5. Is the nilpotency class of a nilpotent group generated by $d$ left 3-Engel elements bounded in terms of $d$ ? A group generated by two left 3-Engel elements is nilpotent of class at most 4 (J. Pure Appl. Algebra, 188 (2004), 1-6.) A. Abdollahi
\end{openproblem}

\plainproblemsection

\begin{openproblem}
	17.6. a) Is there a function $f : \mathbb{N} \times  \mathbb{N} \rightarrow  \mathbb{N}$ such that every nilpotent group generated by $d$ left $n$ -Engel elements is nilpotent of class at most $f\left( {n, d}\right)$ ?
	
	b) Is there a function $g : \mathbb{N} \times  \mathbb{N} \rightarrow  \mathbb{N}$ such that every nilpotent group generated by $d$ right $n$ -Engel elements is nilpotent of class at most $g\left( {n, d}\right)$ ? A. Abdollahi
\end{openproblem}

\plainproblemsection

\begin{openproblem}
	17.7. a) Is there a group in which the set of right Engel elements does not form a subgroup?
	
	b) Is there a group in which the set of bounded right Engel elements does not form a subgroup? A. Abdollahi
\end{openproblem}

\plainproblemsection

\begin{openproblem}
	17.8. a) Is there a group containing a right Engel element which is not a left Engel element?
	
	b) Is there a group containing a bounded right Engel element which is not a left Engel element? A. Abdollahi
\end{openproblem}

\plainproblemsection

\begin{openproblem}
	17.9. Is there a group containing a left Engel element whose inverse is not a left Engel element? A. Abdollahi
\end{openproblem}

\plainproblemsection

\begin{openproblem}
	17.10. Is there a group containing a right 4-Engel element which does not belong to the Hirsch-Plotkin radical? A. Abdollahi
\end{openproblem}

\plainproblemsection

\begin{openproblem}
	17.11. Is there a group containing a left 3-Engel element which does not belong to the Hirsch-Plotkin radical? A. Abdollahi
\end{openproblem}

\plainproblemsection

\begin{openproblem}
	17.12. Are there functions $e, c : \mathbb{N} \rightarrow  \mathbb{N}$ such that if in a nilpotent group $G$ a normal subgroup $H$ consists of right $n$ -Engel elements of $G$ , then ${H}^{e\left( n\right) } \leq  {\zeta }_{c\left( n\right) }\left( G\right)$ ? A. Abdollahi
	
	Yes, there are (P. G. Crosby, G. Traustason, J. Algebra, 324 (2010), 875-883).
\end{openproblem}

\plainproblemsection

\begin{openproblem}
	17.13. Let $G$ be a totally imprimitive $p$ -group of finitary permutations on an infinite set. Suppose that the support of any cycle in the cyclic decomposition of every element of $G$ is a block for $G$ . Does $G$ necessarily contain a minimal non- ${FC}$ -subgroup?
	
	A. O. Asar
\end{openproblem}

\plainproblemsection

\begin{openproblem}
	17.14. Can the braid group ${B}_{n}$ for $n \geq  4$ be embedded into the automorphism group $\operatorname{Aut}\left( {F}_{n - 2}\right)$ of a free group ${F}_{n - 2}$ of rank $n - 2$ ?
	
	Artin’s theorem implies that ${B}_{n}$ can be embedded into $\operatorname{Aut}\left( {F}_{n}\right)$ , and an embedding of ${B}_{n}, n \geq  3$ , into $\operatorname{Aut}\left( {F}_{n - 1}\right)$ was constructed in (B. Perron, J. P. Vannier, Math. Ann., 306 (1996), 231-245). V. G. Bardakov
	
	No, it cannot for $n = 4$ (V. G. Bardakov, P. Bellingeri, in Knot theory and its applications. ICTS program knot theory and its applications (KTH-2013), IISER Mohali, India, December 10-20, 2013, Amer. Math. Soc., Contemporary Mathematics, 670 (2016), 285-298).
\end{openproblem}

\plainproblemsection

\begin{openproblem}
	17.15. Construct an algorithm which, for a given polynomial $f \in  \mathbb{Z}\left\lbrack  {{x}_{1},{x}_{2},\ldots ,{x}_{n}}\right\rbrack$ , finds (explicitly, in terms of generators) the maximal subgroup ${G}_{f}$ of the group $\operatorname{Aut}\left( {\mathbb{Z}\left\lbrack  {{x}_{1},{x}_{2},\ldots ,{x}_{n}}\right\rbrack  }\right)$ that leaves $f$ fixed. This question is related to description of the solution set of the Diophantine equation $f = 0$ . V. G. Bardakov
\end{openproblem}

\plainproblemsection

\begin{openproblem}
	17.16. Let $A$ be an Artin group of finite type, that is, the corresponding Coxeter group is finite. It is known that the group $A$ is linear (S. Bigelow, $J$ . Amer. Math. Soc., 14 (2001), 471-486; D. Krammer, Ann. Math., 155 (2002), 131-156; F. Digne, J. Algebra, 268 (2003), 39-57; A. M. Cohen, D. B. Wales, Israel J. Math., 131 (2002), 101-123). Is it true that the automorphism group $\operatorname{Aut}\left( A\right)$ is also linear?
	
	An affirmative answer for $A$ being a braid group ${B}_{n}, n \geq  2$ , was obtained in (V. G. Bardakov, O. V. Bryukhanov, Vestnik Novosibirsk. Univ. Ser. Mat. Mekh. Inform., 7, no. 3 (2007), 45-58 (Russian)). V. G. Bardakov
\end{openproblem}

\plainproblemsection

\begin{openproblem}
	17.17. If a finitely generated group $G$ has $n < \infty$ maximal subgroups, must $G$ be finite? In particular, what if $n = 3$ ? G. M. Bergman
	
	No, it need not. An example of an infinite group with 3 maximal subgroups is given by a 2-generated 2-LERF group in §7 of (M. Ershov, A. Jaikin-Zapirain, J. Reine Angew. Math. 677 (2013), 71-134) as its maximal subgroups have index 2.
\end{openproblem}

\plainproblemsection

\begin{openproblem}
	17.18. Let $\mathbf{A}$ be the class of compact groups $A$ with the property that whenever two compact groups $B$ and $C$ contain $A$ , they can be embedded in a common compact group $D$ by embeddings agreeing on $A$ . I showed (Manuscr. Math.,58 (1987) 253- 281) that all members of $\mathbf{A}$ are (not necessarily connected) finite-dimensional compact Lie groups satisfying a strong "local simplicity" property, and that all finite groups do belong to $\mathbf{A}$ .
	
	a) Is it true that $\mathbb{R}/\mathbb{Z} \in  \mathbf{A}$ ?
	
	b) Do any nonabelian connected compact Lie groups belong to $\mathbf{A}$ ?
	
	c) If $A$ belongs to $\mathbf{A}$ , must the connected component of the identity in $A$ belong to $\mathbf{A}$ ? G. M. Bergman
\end{openproblem}

\plainproblemsection

\begin{openproblem}
	*17.19. If $F$ is a free group of finite rank, $R$ a retract of $F$ , and $H$ a subgroup of $F$ of finite rank, must $H \cap  R$ be a retract of $H$ ? G. M. Bergman
	
	*No, it must not (I. Snopce, S. Tanushevski, P. Zalesskii, Int. Math. Res. Notices, 2022, no. 11 (2022), 8280-8294).
\end{openproblem}

\plainproblemsection

\begin{openproblem}
	*17.20. If $M$ is a real manifold with nonempty boundary, and $G$ the group of self-homeomorphisms of $M$ which fix the boundary pointwise, is $G$ right-orderable?
	
	G. M. Bergman
	
	*No, not always (J. Hyde, Ann. of Math. (2), 190, no. 2 (2019), 657-661).
\end{openproblem}

\plainproblemsection

\begin{openproblem}
	17.21. a) If $A, B, C$ are torsion-free abelian groups with $A \cong  A \oplus  B \oplus  C$ , must $A \cong  A \oplus  B?$
	
	b) What if, furthermore, $B \cong  C$ ? G. M. Bergman
	
	No, not necessarily, even in case b). See (A. L. S. Corner, in Proc. Colloq. Abelian Groups (Tihany,1963), Budapest,1964,43-48); I also have an example with $B$ and $C$ of rank 1 as torsion-free abelian groups. (G. M. Bergman, Letter of 20 February 2011.)
\end{openproblem}

\plainproblemsection

\begin{openproblem}
	17.22. Suppose $A$ is a group which belongs to a variety $\mathbf{V}$ of groups, and which is embeddable in the full symmetric group $S$ on an infinite set. Must the coproduct in $\mathbf{V}$ of two copies of $A$ also be embeddable in $S$ ? (N. G. de Bruijn proved that this is true if $\mathbf{V}$ is the variety of all groups.) G. M. Bergman
\end{openproblem}

\plainproblemsection

\begin{openproblem}
	17.23. Suppose the full symmetric group $S$ on a countably infinite set is generated by the union of two subgroups $G$ and $H$ . Must $S$ be finitely generated over one of these subgroups? G. M. Bergman
\end{openproblem}

\plainproblemsection

\begin{openproblem}
	17.24. (A. Blass, J. Irwin, G. Schlitt). Does ${\mathbb{Z}}^{\omega }$ have a subgroup whose dual is free of uncountable rank? G. M. Bergman
	
	Yes, it does (G. M. Bergman, Portugaliae Math., 69 (2012) 69-84).
\end{openproblem}

\plainproblemsection

\begin{openproblem}
	17.25. (S. P. Farbman). Let $X$ be the set of complex numbers $\alpha$ such that the group generated by the two $2 \times  2$ matrices $I + \alpha {e}_{12}$ and $I + {e}_{21}$ is not free on those generators.
	
	a) Does $X$ contain all rational numbers in the interval(-4,4)?
	
	b) Does $X$ contain any rational number in the interval $\lbrack {27}/7,4)$ ?
	
	Cf. 4.41 (Archive), 15.83, and (S. P. Farbman, Publ. Mat., 39 (1995) 379-391).
	
	G. M. Bergman
\end{openproblem}

\plainproblemsection

\begin{openproblem}
	17.26. Are the classes of right-orderable and right-ordered groups closed under taking solutions of equations $w\left( {{a}_{1},\ldots ,{a}_{k},{x}_{1},\ldots ,{x}_{n}}\right)  = 1$ ? (Here the closures are under group embeddings and order-preserving embeddings, respectively.) This is true for equations with a single constant, when $k = 1$ (J. Group Theory,11 (2008),623-633).
	
	V. V. Bludov, A. M. W. Glass
\end{openproblem}

\plainproblemsection

\begin{openproblem}
	17.27. Can the free product of two ordered groups with order-isomorphic amalgamated subgroups be lattice orderable but not orderable?
	
	V. V. Bludov, A. M. W. Glass
\end{openproblem}

\plainproblemsection

\begin{openproblem}
	*17.28. Is there a soluble right-orderable group with insoluble word problem? V. V. Bludov, A. M. W. Glass
	
	*Yes, there is (A. Darbinyan, J. Symbolic Logic, 85, no. 4 (2020), 1588-1598).
\end{openproblem}

\plainproblemsection

\begin{openproblem}
	17.29. Construct a finitely presented orderable group with insoluble word problem. V. V. Bludov, A. M. W. Glass
	
	Such a group is constructed (V. V. Bludov, A. M. W. Glass, Bull. London Math. Soc., 44 (2012), 85-98).
\end{openproblem}

\plainproblemsection

\begin{openproblem}
	17.30. Is there a constant $l$ such that every finitely presented soluble group has a subgroup of finite index of nilpotent length at most $l$ ? Cf. 16.35.
	
	V. V. Bludov, J. R. J. Groves
\end{openproblem}

\plainproblemsection

\begin{openproblem}
	17.31. Can a soluble right-orderable group have finite quotient by the derived subgroup? V. V. Bludov, A. H. Rhemtulla
\end{openproblem}

\plainproblemsection

\begin{openproblem}
	17.32. Is the following analogue of the Cayley-Hamilton theorem true for the free group ${F}_{n}$ of rank $n$ : If $w \in  {F}_{n}$ and $\varphi  \in  \operatorname{Aut}{F}_{n}$ are such that $\left\langle  {w,\varphi \left( w\right) ,\ldots ,{\varphi }^{n}\left( w\right) }\right\rangle   =$ ${F}_{n}$ , then $\left\langle  {w,\varphi \left( w\right) ,\ldots ,{\varphi }^{n - 1}\left( w\right) }\right\rangle   = {F}_{n}$ ? O. V. Bogopol'skiǐ
\end{openproblem}

\plainproblemsection

\begin{openproblem}
	17.33. Can the quasivariety generated by the group $\left\langle  {a, b \mid  {a}^{-1}{ba} = {b}^{-1}}\right\rangle$ be defined by a set of quasi-identities in a finite set of variables? The answer is known for all other groups with one defining relation. A. I. Budkin
\end{openproblem}

\plainproblemsection

\begin{openproblem}
	17.34. Let ${N}_{c}$ be the quasivariety of nilpotent torsion-free groups of class at most $c$ . Is it true that the dominion in ${N}_{c}$ (see the definition in 16.20) of a divisible subgroup in every group in ${N}_{c}$ is equal to this subgroup? A. I. Budkin
\end{openproblem}

\plainproblemsection

\begin{openproblem}
	17.35. Suppose we have a finite two-colourable triangulation of the sphere, with triangles each coloured either black or white so that no pair of triangles with the same colour share an edge. On each vertex we write a generator, and we assume the generators commute. We use these generators to generate an abelian group ${G}_{W}$ with relations stating that the generators around each white triangle add to zero. Doing the same thing with the black triangles, we generate a group ${G}_{B}$ .
	
	Conjecture: ${G}_{W}$ is isomorphic to ${G}_{B}$ . Conjecture is proved (S. R. Blackburn, T. A. McCourt, Combinatorica, 34, no. 5 (2014), 527-546).
\end{openproblem}

\plainproblemsection

\begin{openproblem}
	17.36. Two groups are called isospectral if they have the same set of element orders. Find all finite non-abelian simple groups $G$ for which there is a finite group $H$ isospectral to $G$ and containing a proper normal subgroup isomorphic to $G$ . For every simple group $G$ determine all groups $H$ satisfying this condition.
	
	It is easy to show that a group $H$ must satisfy the condition $G < H \leq  \operatorname{Aut}G$ . A. V. Vasil’ev
	
	All such groups are determined: for exceptional groups of Lie type in (M. A. Zvezdina, Algebra Logic, 55, no. 5 (2016), 354-366); for the other simple groups in (M. A. Grechkoseeva, Siberian Math. J., 59, no. 4 (2018), 623-640).
\end{openproblem}

\plainproblemsection

\begin{openproblem}
	*17.37. Is there an integer $n$ such that for all $m > n$ the alternating group ${A}_{m}$ has no non-trivial ${A}_{m}$ -permutable subgroups? (See the definition in 17.112.) A. F. Vasil’ev, A. N. Skiba
	\par\smallskip\hrule\smallskip
	*Yes, there is (N. Yang, A. Galt, Bull. Malays. Math. Sci. Soc. (2), 46, no. 5 (2023), Paper no. 177, 16 p.).
	\par\smallskip\hrule\smallskip
\end{openproblem}

\plainproblemsection

\begin{openproblem}
	17.38. A formation $\mathfrak{F}$ is called radical in a class $\mathfrak{H}$ if $\mathfrak{F} \subseteq  \mathfrak{H}$ and in every $\mathfrak{H}$ -group the product of any two normal $\mathfrak{F}$ -subgroups belongs to $\mathfrak{F}$ . Let $\mathfrak{M}$ be the class of all saturated hereditary formations of finite groups such that the formation of all finite supersoluble groups is radical in every element of $\mathfrak{M}$ . Is it true that $\mathfrak{M}$ has the largest (by inclusion) element? A. F. Vasil’ev, L. A. Shemetkov
\end{openproblem}

\plainproblemsection

\begin{openproblem}
	17.39. Is there a positive integer $n$ such that the hypercenter of any finite soluble group coincides with the intersection of $n$ system normalizers of that group? What is the least number with this property?
	
	Comment of 2017: This is proved for groups with trivial Frattini subgroup (A. Ballester-Bolinches, J. Cossey, S. F. Kamornikov, H. Meng, J. Algebra, 499 (2018), 438-449). A. F. Vasil’ev, L. A. Shemetkov
\end{openproblem}

\plainproblemsection

\begin{openproblem}
	*17.40. Let $N$ be a nilpotent subgroup of a finite group $G$ . Do there always exist elements $x, y \in  G$ such that $N \cap  {N}^{x} \cap  {N}^{y} \leq  F\left( G\right)$ ?
	
	*Yes, such elements always exist, mod CFSG (V. I. Zenkov, Siberian Math. J., 62, no. 4 (2021), 621-637).
\end{openproblem}

\plainproblemsection

\begin{openproblem}
	17.41. Let $S$ be a solvable subgroup of a finite group $G$ that has no nontrivial solvable normal subgroups.
	
	a) (L. Babai, A. J. Goodman, L. Pyber). Do there always exist seven conjugates of $S$ whose intersection is trivial?
	
	b) Do there always exist five conjugates of $S$ whose intersection is trivial?
	
	Comment of 2013: reduction to almost simple group $G$ is obtained in (E. P. Vdo-vin, J. Algebra Appl., 11, no. 1 (2012), 1250015 (14 pages)). E. P. Vdovin
\end{openproblem}

\plainproblemsection

\begin{openproblem}
	*17.42. Let $\bar{G}$ be a simple algebraic group of adjoint type over the algebraic closure ${\bar{F}}_{p}$ of a finite field ${F}_{p}$ of prime order $p$ , and $\sigma$ a Frobenius map (that is, a surjective homomorphism such that ${\bar{G}}_{\sigma } = {C}_{\bar{G}}\left( \sigma \right)$ is finite). Then $G = {O}^{{p}^{\prime }}\left( {\bar{G}}_{\sigma }\right)$ is a finite group of Lie type. For a maximal $\sigma$ -stable torus $\bar{T}$ of $\bar{G}$ , let $N = {N}_{\bar{G}}\left( \bar{T}\right)  \cap  G$ . Assume also that $G$ is simple and $G ≆ {\mathrm{{SL}}}_{3}\left( 2\right)$ . Does there always exist $x \in  G$ such that $N \cap  {N}^{x}$ is a $p$ -group? E. P. Vdovin
	\par\smallskip\hrule\smallskip
	*A definitive answer for a stronger question on the size of a base has been obtained in (T. C. Burness, A. R. Thomas, J. Algebra, 619 (2023), 459-504), which implies the following answer (in the notation therein): there is $x \in  G$ such that $N \cap  {N}^{x}$ is a $p$ -group if and only if(G, N)is not one of the following: $\left( {{\mathrm{L}}_{3}\left( 2\right) ,7 : 3}\right) ,\left( {{\mathrm{U}}_{4}\left( 2\right) ,{3}^{3} : {S}_{4}}\right)$ , $\left( {{\mathrm{U}}_{5}\left( 2\right) ,{3}^{4} : {S}_{5}}\right)$ .
\end{openproblem}

\plainproblemsection

\begin{openproblem}
	17.43. Let $\pi$ be a set of primes. Find all finite simple ${D}_{\pi }$ -groups (see Archive,3.62) in which
	
	a) every subgroup is a ${D}_{\pi }$ -group (H. Wielandt);
	
	b) every subgroup possessing a Hall $\pi$ -subgroup is a ${D}_{\pi }$ -group.
	
	All finite simple ${D}_{\pi }$ -groups are known (D. O. Revin, Algebra and Logic,47, no. 3 (2008), 210-227). Comment of 2013: alternating and sporadic groups for both parts are found in (N. Ch. Manzaeva, Siberian Electr. Math. Rep., 9 (2012), 294-305 (Russian)).
	\par\smallskip\hrule\smallskip
\end{openproblem}

\plainproblemsection

\begin{openproblem}
	17.44. Let $\pi$ be a set of primes. A finite group is called a ${C}_{\pi }$ -group if it possesses exactly one class of conjugate Hall $\pi$ -subgroups. A finite group is called a ${D}_{\pi }$ -group if any two of its maximal $\pi$ -subgroups are conjugate.
	
	a) In a ${C}_{\pi }$ -group, is an overgroup of a Hall $\pi$ -subgroup always a ${C}_{\pi }$ -group?
	
	An affirmative answer in the case $2 \notin  \pi$ follows mod CFSG from (F. Gross, Bull. London Math. Soc., 19, no. 4 (1987), 311-319).
	
	b) In a ${D}_{\pi }$ -group, is an overgroup of a Hall $\pi$ -subgroup always a ${D}_{\pi }$ -group? E. P. Vdovin, D. O. Revin
	
	a) Yes, it is (E. P. Vdovin, D. O. Revin, Siberian Math. J., 54, no. 1 (2013), 22-28). b) Yes, it is (N. Ch. Manzaeva, https://arxiv.org/pdf/1504.03137.pdf) and (E. P. Vdovin, N. Ch. Manzaeva, D. O. Revin, Sb. Math., 211, no. 3 (2020), 309-335).
\end{openproblem}

\plainproblemsection

\begin{openproblem}
	17.45. A subgroup $H$ of a group $G$ is called pronormal if $H$ and ${H}^{g}$ are conjugate in $\left\langle  {H,{H}^{g}}\right\rangle$ for every $g \in  G$ . We say that $H$ is strongly pronormal if ${L}^{g}$ is conjugate to a subgroup of $H$ in $\left\langle  {H,{L}^{g}}\right\rangle$ for every $L \leq  H$ and $g \in  G$ .
	
	a) In a finite simple group, are Hall subgroups always pronormal?
	
	An affirmative answer is known modulo CFSG for Hall subgroups of odd order
	
	(F. Gross, Bull. London Math. Soc., 19, no. 4 (1987), 311-319).
	
	b) In a finite simple group, are Hall subgroups always strongly pronormal?
	
	Notice that there exist finite (non-simple) groups with a non-pronormal Hall sub-
	
	group. Hall subgroups with a Sylow tower are known to be pronormal.
	
	E. P. Vdovin, D. O. Revin
	
	a) Yes, they are (E. P. Vdovin, D. O. Revin, Siberian Math. J., 53, no. 3 (2012), 419- 430).
	
	b) No, not always (M. N. Nesterov, Siberian Math. J., 58, no. 1 (2017), 128-133).
\end{openproblem}

\plainproblemsection

\begin{openproblem}
	17.45. c) A subgroup $H$ of a group $G$ is called pronormal if $H$ and ${H}^{g}$ are conjugate in $\left\langle  {H,{H}^{g}}\right\rangle$ for every $g \in  G$ . We say that $H$ is strongly pronormal if ${L}^{g}$ is conjugate to a subgroup of $H$ in $\left\langle  {H,{L}^{g}}\right\rangle$ for every $L \leq  H$ and $g \in  G$ .
	
	In a finite group, is a Hall subgroup with a Sylow tower always strongly pronormal?
	
	Notice that there exist finite (non-simple) groups with a non-pronormal Hall subgroup. Hall subgroups with a Sylow tower are known to be pronormal.
	
	E. P. Vdovin, D. O. Revin
\end{openproblem}

\plainproblemsection

\begin{openproblem}
	17.46. Let $G$ be a finite $p$ -group in which every 2-generator subgroup has cyclic derived subgroup. Is the derived length of $G$ bounded?
	
	If $p \neq  2$ , then ${G}^{\left( 2\right) } = 1$ (J. Alperin), but for $p = 2$ there are examples with ${G}^{\left( 2\right) }$ cyclic and elementary abelian of arbitrary order.
	
	Yes, it is at most 4 (B. Wilkens, J. Group Theory, 17, no. 1 (2014), 151-174).
\end{openproblem}

\plainproblemsection

\begin{openproblem}
	17.47. Let $G$ be a nilpotent group in which every coset $x\left\lbrack  {G, G}\right\rbrack$ for $x \notin  \left\lbrack  {G, G}\right\rbrack$ coincides with the conjugacy class ${x}^{G}$ . Is there a bound for the nilpotency class of $G$ ?
	
	If $G$ is finite, then its class is at most 3 (R. Dark, C. Scoppola, J. Algebra,181, no. 3 (1996), 787-802). S. H. Ghate, A. S. Muktibodh
\end{openproblem}

\plainproblemsection

\begin{openproblem}
	17.48. Does the free product of two groups with stable first-order theory also have stable first-order theory? E. Jaligot
\end{openproblem}

\plainproblemsection

\begin{openproblem}
	17.49. Is it consistent with ZFC that every non-discrete topological group contains a nonempty nowhere dense subset without isolated points? Under Martin's Axiom, there are countable non-discrete topological groups in which each nonempty nowhere dense subset has an isolated point. E. G. Zelenyuk
\end{openproblem}

\plainproblemsection

\begin{openproblem}
	17.50. Is it true that for every finite group $G$ , there is a finite group $F$ and a surjective homomorphism $f : F \rightarrow  G$ such that for each nontrivial subgroup $H$ of $F$ , the restriction ${\left. f\right| }_{H}$ is not injective?
	
	It is known that for every finite group $G$ there is a finite group $F$ and a surjective homomorphism $f : F \rightarrow  G$ such that for each subgroup $H$ of $F$ the restriction ${\left. f\right| }_{H}$ is not bijective. E. G. Zelenyuk
	
	Yes, it is (V. P. Burichenko, Math. Notes, 92, no. 3 (2012), 327-332).
\end{openproblem}

\plainproblemsection

\begin{openproblem}
	17.51. Is it true that every non-discrete topological group containing no countable open subgroup of exponent 2 can be partitioned into two (infinitely many) dense subsets? If the group is Abelian or countable, the answer is yes.
	
	E. G. Zelenyuk, I. V. Protasov
\end{openproblem}

\plainproblemsection

\begin{openproblem}
	17.52. (N. Eggert). Let $R$ be a commutative associative finite-dimensional nilpotent algebra over a finite field $F$ of characteristic $p$ . Let ${R}^{\left( p\right) }$ be the subalgebra of all elements of the form ${r}^{p}\left( {r \in  R}\right)$ . Is it true that $\dim R \geq  p\dim {R}^{\left( p\right) }$ ?
	
	An affirmative answer gives a solution to 11.44 for the case where the factors $A$ and $B$ are abelian. L. S. Kazarin
\end{openproblem}

\plainproblemsection

\begin{openproblem}
	17.53. A finite group $G$ is called $\operatorname{simply}$ reducible (SR-group) if every element of $G$ is conjugate to its inverse, and the tensor product of any two irreducible representations of $G$ decomposes into a sum of irreducible representations of $G$ with coefficients 0 or 1 .
	
	a) Is it true that the nilpotent length of a soluble SR-group is at most 5 ?
	
	b) Is it true that the derived length of a soluble SR-group is bounded by some constant $c$ ? L. S. Kazarin
\end{openproblem}

\plainproblemsection

\begin{openproblem}
	17.54. Does there exist a non-hereditary local formation $\mathfrak{F}$ of finite groups such that in every finite group the set of all $\mathfrak{F}$ -subnormal subgroups is a sublattice of the subgroup lattice?
	
	A negative answer would solve 9.75 , since all the hereditary local formations with the same property are already known. S. F. Kamornikov
\end{openproblem}

\plainproblemsection

\begin{openproblem}
	17.55. Does there exist an absolute constant $k$ such that for any prefrattini subgroup $H$ in any finite soluble group $G$ there exist $k$ conjugates of $H$ whose intersection is $\Phi \left( G\right)$ ? S. F. Kamornikov
	
	Yes, it exists and is equal to 3 (S. F. Kamornikov, Int. J. Group Theory, 6, no. 2 (2017), 1-5).
\end{openproblem}

\plainproblemsection

\begin{openproblem}
	*17.56. Suppose that a subgroup $H$ of a finite group $G$ is such that ${HM} = {MH}$ for every minimal non-nilpotent subgroup $M$ of $G$ . Must $H/{H}_{G}$ be nilpotent, where ${H}_{G}$ is the largest normal subgroup of $G$ contained in $H$ ? V. N. Knyagina, V. S. Monakhov
	
	*No, not always (A. Ballester-Bolinches, R. Esteban-Romero, S. F. Kamornikov, V. Pérez-Calabuig, J. Algebra, 608 (2022), 382-387).
\end{openproblem}

\plainproblemsection

\begin{openproblem}
	17.57. Let $r\left( m\right)  = \{ r + {km} \mid  k \in  \mathbb{Z}\}$ for integers $0 \leq  r < m$ . For ${r}_{1}\left( {m}_{1}\right)  \cap  {r}_{2}\left( {m}_{2}\right)  = \varnothing$ define the class transposition ${\tau }_{{r}_{1}\left( {m}_{1}\right) ,{r}_{2}\left( {m}_{2}\right) }$ as the involution which interchanges ${r}_{1} + k{m}_{1}$ and ${r}_{2} + k{m}_{2}$ for each integer $k$ and fixes everything else. The group $\mathrm{{CT}}\left( \mathbb{Z}\right)$ generated by all class transpositions is simple (Math. Z.,264, no. 4 (2010),
	
	927-938). Is $\operatorname{Out}\left( {\mathrm{{CT}}\left( \mathbb{Z}\right) }\right)  = \left\langle  {\sigma  \mapsto  {\sigma }^{n \mapsto   - n - 1}}\right\rangle   \cong  {\mathrm{C}}_{2}$ ?S. Kohl
\end{openproblem}

\plainproblemsection

\begin{openproblem}
	17.58. Does $\mathrm{{CT}}\left( \mathbb{Z}\right)$ have subgroups of intermediate (word-) growth?S. Kohl
\end{openproblem}

\plainproblemsection

\begin{openproblem}
	17.59. A permutation of $\mathbb{Z}$ is called residue-class-wise affine if there is a positive integer $m$ such that its restrictions to the residue classes $\left( {\;\operatorname{mod}\;m}\right)$ are all affine. Is $\mathrm{{CT}}\left( \mathbb{Z}\right)$ the group of all residue-class-wise affine permutations of $\mathbb{Z}$ which fix the nonnegative integers setwise?S. Kohl
\end{openproblem}

\plainproblemsection

\begin{openproblem}
	17.60. Given a set $\mathcal{P}$ of odd primes, let ${\mathrm{{CT}}}_{\mathcal{P}}\left( \mathbb{Z}\right)$ denote the subgroup of $\mathrm{{CT}}\left( \mathbb{Z}\right)$ which is generated by all class transpositions which interchange residue classes whose moduli have only prime factors in $\mathcal{P} \cup  \{ 2\}$ . The groups ${\mathrm{{CT}}}_{\mathcal{P}}\left( \mathbb{Z}\right)$ are simple (Math. Z., 264, no. 4 (2010), 927-938). Are they pairwise nonisomorphic?S. Kohl
\end{openproblem}

\plainproblemsection

\begin{openproblem}
	17.61. The group ${\mathrm{{CT}}}_{\mathcal{P}}\left( \mathbb{Z}\right)$ is finitely generated if and only if $\mathcal{P}$ is finite. If $\mathcal{P} = \varnothing$ , then it is isomorphic to the finitely presented (first) Higman-Thompson group (J. P. McDermott, see Remark 1.4 in (S. Kohl, J. Group Theory, 20, no. 5 (2017), 1025-1030). Is it always finitely presented if $\mathcal{P}$ is finite?
\end{openproblem}

\plainproblemsection

\begin{openproblem}
	*17.62. Given a free group ${F}_{n}$ and a proper characteristic subgroup $C$ , is it ever possible to generate the quotient ${F}_{n}/C$ by fewer than $n$ elements? J. Conrad
	
	*Yes, it is possible; moreover, for all $n \geq  2$ , the free group ${F}_{n}$ admits continuum many pairwise non-isomorphic infinite simple characteristic 2-generated quotients (R. Coulon, F. Fournier-Facio, https://arxiv.org/abs/2312.11684).
\end{openproblem}

\plainproblemsection

\begin{openproblem}
	17.63. Prove that if $p$ is an odd prime and $s$ is a positive integer, then there are only finitely many $p$ -adic space groups of finite coclass with point group of coclass $s$ .
	
	C. R. Leedham-Green
\end{openproblem}

\plainproblemsection

\begin{openproblem}
	17.64. Say that a group $G$ is an $n$ -approximation to the Nottingham group $J =$ $N\left( p\right)$ (as defined in Archive,12.24) if $G$ is an infinite pro- $p$ group, and $G/{\gamma }_{n}\left( G\right)$ is isomorphic to $J/{\gamma }_{n}\left( J\right)$ . Does there exist a function $f\left( p\right)$ such that, if $G$ is an $f\left( p\right)$ -approximation to the Nottingham group, then ${\gamma }_{i}\left( G\right) /{\gamma }_{i + 1}\left( G\right)$ is isomorphic to ${\gamma }_{i}\left( J\right) /{\gamma }_{i + 1}\left( J\right)$ for all $i$ ?
	
	Cf. 14.56. Note that an affirmative solution to this problem trivially implies the now known fact that $J$ is finitely presented as a pro- $p$ group (if $p > 2$ ), see 14.55.
\end{openproblem}

\plainproblemsection

\begin{openproblem}
	17.65. If Problem 17.64 has an affirmative answer, does it follow that, for some function $g\left( p\right)$ , the number of isomorphism classes of $g\left( p\right)$ -approximations to $J$ is
	
	a) countable?
	
	b) one?
	
	Note that if, for some $g\left( p\right)$ , there are only finitely many isomorphism classes, then for some $h\left( p\right)$ there is only one. C. R. Leedham-Green
\end{openproblem}

\plainproblemsection

\begin{openproblem}
	17.66. (R. Guralnick). Does there exist a positive integer $d$ such that $\dim {H}^{1}\left( {G, V}\right)  \leq$ $d$ for any faithful absolutely irreducible module $V$ for any finite group $G$ ? Cf. 16.55 in Archive. V. D. Mazurov
\end{openproblem}

\plainproblemsection

\begin{openproblem}
	17.67. (H. Zassenhaus). Conjecture: Every invertible element of finite order of the integral group ring $\mathbb{Z}G$ of a finite group $G$ is conjugate in the rational group ring $\mathbb{Q}G$ to an element of $\pm  G$ . V. D. Mazurov
	
	A counterexample was constructed (F. Eisele, L. Margolis, Adv. Math. 339 (2018), 599-641).
\end{openproblem}

\plainproblemsection

\begin{openproblem}
	17.68. Let $C$ be a cyclic subgroup of a group $G$ and $G = {CFC}$ where $F$ is a finite cyclic subgroup. Is it true that $\left| {G : C}\right|$ is finite? Cf. 12.62. V. D. Mazurov
\end{openproblem}

\plainproblemsection

\begin{openproblem}
	17.69. Let $G$ be a group of prime exponent acting freely on a non-trivial abelian group. Is $G$ cyclic? V. D. Mazurov
\end{openproblem}

\plainproblemsection

\begin{openproblem}
	17.70. Let $\alpha$ be an automorphism of prime order $q$ of an infinite free Burnside group $G = B\left( {q, p}\right)$ of prime exponent $p$ such that $\alpha$ cyclically permutes the free generators of $G$ . Is it true that $\alpha$ fixes some non-trivial element of $G$ ?
	
	Editors’ comment: this is proved for $q = 2$ in (V. S. Atabekyan, H. T. Aslanyan, Proc. Yerevan State Univ. Physic. Math. Sci., 53, no. 3 (2019), 147-149).
	
	V. D. Mazurov
\end{openproblem}

\plainproblemsection

\begin{openproblem}
	17.71. a) Let $\alpha$ be a fixed-point-free automorphism of prime order $p$ of a periodic group $G$ . Is it true that $G$ does not contain a non-trivial $p$ -element? V. D. Mazurov a) No, it may contain a non-trivial $p$ -element (E. B. Durakov, A. I. Sozutov, Algebra Logic, 52, no. 5 (2013), 422-425).
\end{openproblem}

\plainproblemsection

\begin{openproblem}
	17.71. b) Let $\alpha$ be a fixed-point-free automorphism of prime order $p$ of a periodic group $G$ . Suppose that $G$ does not contain a non-trivial $p$ -element. Is $\alpha$ a splitting automorphism? V. D. Mazurov
\end{openproblem}

\plainproblemsection

\begin{openproblem}
	17.72. a) Let ${AB}$ be a Frobenius group with kernel $A$ and complement $B$ . Suppose that ${AB}$ acts on a finite group $G$ so that ${GA}$ is also a Frobenius group with kernel $G$ and complement $A$ . Is the nilpotency class of $G$ bounded in terms of $\left| B\right|$ and the class of ${C}_{G}\left( B\right)$ ? V. D. Mazurov
	
	a) Yes, it is (N. Yu. Makarenko, P. Shumyatsky, Proc. Amer. Math. Soc., 138 (2010), 3425-3436).
\end{openproblem}

\plainproblemsection

\begin{openproblem}
	17.72. b) Let ${AB}$ be a Frobenius group with kernel $A$ and complement $B$ . Suppose that ${AB}$ acts on a finite group $G$ so that ${GA}$ is also a Frobenius group with kernel $G$ and complement $A$ . Is the exponent of $G$ bounded in terms of $\left| B\right|$ and the exponent of ${C}_{G}\left( B\right)$ ? V. D. Mazurov
\end{openproblem}

\plainproblemsection

\begin{openproblem}
	*17.73. a) Let $G$ be a finite simple group of Lie type defined over a field of characteristic $p$ , and $V$ an absolutely irreducible $G$ -module over a field of the same characteristic. Is it true that in the case $G = {U}_{4}\left( q\right)$ the split extension of $V$ by $G$ must contain an element whose order is distinct from the order of any element of $G$ ? $V.D.{Mazurov}$
	
	*a) Yes, it is true (M. A. Grechkoseeva, S. V. Skresanov, Siberian Electron. Math. Rep., 17 (2020), 585-589).
\end{openproblem}

\plainproblemsection

\begin{openproblem}
	17.73. Let $G$ be a finite simple group of Lie type defined over a field of characteristic $p$ , and $V$ an absolutely irreducible $G$ -module over a field of the same characteristic. Is it true that in the cases
	
	b) $G = {S}_{2n}\left( q\right) , n \geq  3$ ;
	
	c) $G = {O}_{{2n} + 1}\left( q\right) , n \geq  3$ ;
	
	d) $G = {O}_{2n}^{ + }\left( q\right) , n \geq  4$ ;
	
	e) $G = {O}_{2n}^{ - }\left( q\right) , n \geq  4$ ;
	
	f) $G = {}^{3}{D}_{4}\left( q\right) , q \neq  2$ ;
	
	g) $G = {E}_{6}\left( q\right)$ ;
	
	h) $G = {}^{2}{E}_{6}\left( q\right)$ ;
	
	i) $G = {E}_{7}\left( q\right)$ ;
	
	j) $G = {G}_{2}\left( {2}^{m}\right)$
	
	the split extension of $V$ by $G$ must contain an element whose order is distinct from the order of any element of $G$ ? V. D. Mazurov
	
	Yes, it is true: b)-i) (M. A. Grechkoseeva, J. Algebra Appl., 14, no. 4 (2015), Article ID 1550056); j) (A. V. Vasil’ev, A. M. Staroletov, Algebra Logic, 52, no. 1 (2013), $1 - {14})$ .
\end{openproblem}

\plainproblemsection

\begin{openproblem}
	17.74. Let $G$ be a finite simple group of Lie type defined over a field of characteristic $p$ whose Lie rank is at least three, and $V$ an absolutely irreducible $G$ -module over a field of characteristic that does not divide $p$ . It is true that the split extension of $V$ by $G$ must contain an element whose order is distinct from the order of any element of $G$ ? The case of $G = {U}_{n}\left( {p}^{m}\right)$ is of special interest.
	
	Yes, it is true: for linear groups (A. V. Zavarnitsine, Siberian Math. J., 49 (2008), 246- 256); for other classical groups (M. A. Grechkoseeva, J. Algebra, 339, no. 1 (2011), 304-319); for exceptional gorups (M. A. Grechkoseeva, J. Algebra Appl., 14, no. 4 (2015), Article ID 1550056).
\end{openproblem}

\plainproblemsection

\begin{openproblem}
	17.75. Can the Monster $M$ act on a nontrivial finite 3-group $V$ so that all elements of $M$ of orders41,47,59,71have no nontrivial fixed points in $V$ ?
\end{openproblem}

\plainproblemsection

\begin{openproblem}
	17.76. Does there exist a finite group $G$ , with $\left| G\right|  > 2$ , such that there is exactly one element in $G$ which is not a commutator? D. MacHale
\end{openproblem}

\plainproblemsection

\begin{openproblem}
	17.77. Let $k$ be a positive integer such that there is an insoluble finite group with exactly $k$ conjugacy classes. Is it true that a finite group of maximal order with exactly $k$ conjugacy classes is insoluble?
	
	No, it is not; see Table 1 in (A. Vera-Lopez, J. Sangroniz, Math. Nachr., 280, no. 5-6 (2007), 676-694).
\end{openproblem}

\plainproblemsection

\begin{openproblem}
	17.78. Does there exist a finitely generated group without free subsemigroups generating a proper variety containing ${\mathfrak{A}}_{p}\mathfrak{A}$ ? O. Macedońska
\end{openproblem}

\plainproblemsection

\begin{openproblem}
	17.79. Does there exist a finitely generated torsion group of unbounded exponent generating a proper variety? O. Macedońska
	
	Yes, moreover, there is a continuum of such groups (V. S. Atabekyan, Infinite simple groups satisfying an identity, Dep. VINITI no. 5381-V86, Moscow, 1986 (Russian)). Another example was suggested by D. Osin in a letter of 31 August 2013: a free group $G$ in the variety $\mathfrak{M}$ defined by the law ${x}^{n}y = y{x}^{n}$ is a central extension of a free Burnside group of exponent $n$ such that the centre of $G$ is a free abelian group of countable rank (I. S. Ashmanov, A. Yu. Ol'shanskiĭ, Izv. Vyssh. Uchebn. Zaved. Mat.,1985, no. 11 (1985), ${48} - {60}$ (Russian)). Let ${x}_{1},{x}_{2},\ldots$ be a basis in $Z\left( G\right)$ . By adding to $G$ the relations ${x}_{i}^{i} = 1$ for all $i$ we obtain a periodic group of unbounded exponent generating a proper variety (contained in $\mathfrak{M}$ ).
\end{openproblem}

\plainproblemsection

\begin{openproblem}
	17.80. Let $\left\lbrack  {u,{}_{n}v}\right\rbrack   \mathrel{\text{:=}} \left\lbrack  {u,\overset{n}{\overbrace{v,\ldots , v}}}\right\rbrack$ . Is the group ${M}_{n} = \left\langle  {x, y \mid  x = \left\lbrack  {x,{}_{n}y}\right\rbrack  , y = \left\lbrack  {y,{}_{n}x}\right\rbrack  }\right\rangle$ infinite for every $n > 2$ ? (See also 11.18.) O. Macedońska
\end{openproblem}

\plainproblemsection

\begin{openproblem}
	17.81. Given normal subgroups ${R}_{1},\ldots ,{R}_{n}$ of a group $G$ , let $\left\lbrack  \left\lbrack  {{R}_{1},\ldots ,{R}_{n}}\right\rbrack  \right\rbrack   \mathrel{\text{:=}}$ $\prod \left\lbrack  {\mathop{\bigcap }\limits_{{i \in  I}}{R}_{i},\mathop{\bigcap }\limits_{{j \in  J}}{R}_{j}}\right\rbrack$ , where the product is over all $I \cup  J = \{ 1,\ldots , n\} , I \cap  J = \varnothing$ .
	
	Let $G$ be a free group, and let ${R}_{i} = {\left\langle  {r}_{i}\right\rangle  }^{G}$ be the normal closures of elements ${r}_{i} \in  G$ . It is known (B. Hartley, Yu. Kuzmin, J. Pure Appl. Algebra, 74 (1991), 247-256) that the quotient $\left( {{R}_{1} \cap  {R}_{2}}\right) /\left\lbrack  {{R}_{1},{R}_{2}}\right\rbrack$ is a free abelian group. On the other hand, the quotient $\left( {{R}_{1} \cap  \cdots  \cap  {R}_{n}}\right) /\left\lbrack  \left\lbrack  {{R}_{1},\ldots ,{R}_{n}}\right\rbrack  \right\rbrack$ has, in general, non-trivial torsion for $n \geq  4$ . Is this quotient always torsion-free for $n = 3$ ? This is known to be true if ${r}_{1},{r}_{2},{r}_{3}$ are not proper powers in $G$ . R. Mikhailov
\end{openproblem}

\plainproblemsection

\begin{openproblem}
	17.82. Is it true that in every finitely presented group the intersection of the derived series has trivial abelianization? R. Mikhailov
	
	No, not always (V. A. Roman'kov, Siberian Math. J., 52, no. 2 (2011), 348-351).
\end{openproblem}

\plainproblemsection

\begin{openproblem}
	17.83. Does there exist a group such that every central extension of it is residually nilpotent, but there exists a central extension of a central extension of it which is not residually nilpotent? R. Mikhailov
\end{openproblem}

\plainproblemsection

\begin{openproblem}
	17.84. An associative algebra $A$ is said to be Calabi-Yau of dimension $d$ (for short, CY $d$ ) if there is a natural isomorphism of $A$ -bimodules ${\operatorname{Ext}}_{A\text{-bimod }}^{d}\left( {A, A \otimes  A}\right)  \cong  A$ and ${\operatorname{Ext}}_{A\text{-bimod }}^{n}\left( {A, A \otimes  A}\right)  = 0$ for $n \neq  d$ . By Kontsevich’s theorem, the complex group algebra $\mathbb{C}G$ of the fundamental group $G$ of a 3-dimensional aspherical manifold is CY3. Is every group with CY3 complex group algebra residually finite? R. Mikhailov
\end{openproblem}

\plainproblemsection

\begin{openproblem}
	17.85. Let $\mathcal{V}$ be a variety of groups such that any free group in $\mathcal{V}$ has torsion-free integral homology groups in all dimensions. Is it true that $\mathcal{V}$ is abelian?
	
	It is known that free metabelian groups may have 2-torsion in the second homology groups, and free nilpotent groups of class 2 may have 3-torsion in the fourth homology groups. R. Mikhailov
\end{openproblem}

\plainproblemsection

\begin{openproblem}
	17.86. (Simplest questions related to the Whitehead asphericity conjecture). Let $\left\langle  {{x}_{1},{x}_{2},{x}_{3} \mid  {r}_{1},{r}_{2},{r}_{3}}\right\rangle$ be a presentation of the trivial group.
	
	a) Prove that the group $\left\langle  {{x}_{1},{x}_{2},{x}_{3} \mid  {r}_{1},{r}_{2}}\right\rangle$ is torsion-free. R. Mikhailov
\end{openproblem}

\plainproblemsection

\begin{openproblem}
	17.86. (Simplest questions related to the Whitehead asphericity conjecture). Let $\left\langle  {{x}_{1},{x}_{2},{x}_{3} \mid  {r}_{1},{r}_{2},{r}_{3}}\right\rangle$ be a presentation of the trivial group.
	
	b) Let $F = F\left( {{x}_{1},{x}_{2},{x}_{3}}\right)$ and ${R}_{i} = {\left\langle  {r}_{i}\right\rangle  }^{F}$ . Is it true that the group $F/\left\lbrack  {{R}_{1},{R}_{2}}\right\rbrack$ is residually soluble? R. Mikhailov
	
	b) No, not always (V. A. Roman'kov, Siberian Math. J., 52, no. 2 (2011), 348-351).
\end{openproblem}

\plainproblemsection

\begin{openproblem}
	17.87. Construct a group of intermediate growth with finitely generated Schur multiplier. R. Mikhailov
\end{openproblem}

\plainproblemsection

\begin{openproblem}
	17.88. Compute ${K}_{0}\left( {{\mathbb{F}}_{2}G}\right)$ and ${K}_{1}\left( {{\mathbb{F}}_{2}G}\right)$ , where $G$ is the first Grigorchuk group, ${\mathbb{F}}_{2}G$ its group algebra over the field of two elements, and ${K}_{0},{K}_{1}$ the zeroth and first $K$ -functors. R. Mikhailov
\end{openproblem}

\plainproblemsection

\begin{openproblem}
	17.89. By Bousfield's theorem, the free pronilpotent completion of a non-cyclic free group has uncountable Schur multiplier. Is it true that the free prosolvable completion (that is, the inverse limit of quotients by the derived subgroups) of a non-cyclic free group has uncountable Schur multiplier? R. Mikhailov
\end{openproblem}

\plainproblemsection

\begin{openproblem}
	17.90. (G. Baumslag). A group is parafree if it is residually nilpotent and has the same lower central quotients as a free group. Is it true that ${H}_{2}\left( G\right)  = 0$ for any finitely generated parafree group $G$ ? R. Mikhailov
\end{openproblem}

\plainproblemsection

\begin{openproblem}
	17.91. Let $d\left( X\right)$ denote the derived length of a group $X$ .
	
	a) Does there exist an absolute constant $k$ such that $d\left( G\right)  - d\left( M\right)  \leq  k$ for every finite soluble group $G$ and any maximal subgroup $M$ ?
	
	b) Find the minimum $k$ with this property. V. S. Monakhov
\end{openproblem}

\plainproblemsection

\begin{openproblem}
	17.92. What are the non-abelian composition factors of a finite non-soluble group all of whose maximal subgroups are Hall subgroups? V. S. Monakhov Answer: ${L}_{2}\left( 7\right) ,{L}_{2}\left( {11}\right) ,{L}_{5}\left( 2\right)$ (N.V.Maslova, Siberian Math. J.,53, no. 5 (2012), 853-861).
\end{openproblem}

\plainproblemsection

\begin{openproblem}
	17.93. (Well-known problem). Let $G$ be a compact topological group which has elements of arbitrarily high orders. Must $G$ contain an element of infinite order?
	
	J. Mycielski 17.94. (Well-known problem). Can the free product $Z * G$ of the infinite cyclic group $Z$ and a nontrivial group $G$ be the normal closure of a single element?
	
	Editors' comments: the negation is also known as Kervaire's conjecture. It was proved for torsion-free groups (A. Klyachko, Commun. Algebra, 21, no. 7 (1993), 2555-2575) and for residually finite groups (M. Gerstenhaber, O. S. Rothaus, Proc. Nat. Acad. Sci. U.S.A., 48 (1962) 1531-1533). J. Mycielski
\end{openproblem}

\plainproblemsection

\begin{openproblem}
	17.95. Let $G$ be a permutation group on the finite set $\Omega$ . A partition $\rho$ of $\Omega$ is said to be $G$ -regular if there exists a subset $S$ of $\Omega$ such that ${S}^{g}$ is a transversal of $\rho$ for all $g \in  G$ . The group $G$ is said to be synchronizing if $\left| \Omega \right|  > 2$ and there are no non-trivial proper $G$ -regular partitions on $\Omega$ .
	
	a) Are the following affine-type primitive groups synchronizing: ${2}^{p} \cdot  \operatorname{PSL}\left( {2,{2p} + 1}\right)$ where both $p$ and ${2p} + 1$ are prime, $p \equiv  3\left( {\;\operatorname{mod}\;4}\right)$ and $p > {23}$ ? ${2}^{101} \cdot  {He}$ ?
	
	b) For which finite simple groups $S$ are the groups $S \times  S$ acting on $S$ by(g, h): $x \mapsto  {g}^{-1}{xh}$ non-synchronizing? Peter M. Neumann
\end{openproblem}

\plainproblemsection

\begin{openproblem}
	17.96. Does there exist a variety of groups which contains only countably many subvarieties but in which there is an infinite properly descending chain of subvarieties?
	
	Peter M. Neumann
\end{openproblem}

\plainproblemsection

\begin{openproblem}
	17.97. Is every variety of groups of exponent 4 finitely based? Peter M. Neuman
\end{openproblem}

\plainproblemsection

\begin{openproblem}
	17.98. A variety of groups is said to be small if it contains only countably many non-isomorphic finitely generated groups.
	
	a) Is it true that if $G$ is a finitely generated group and the variety $\operatorname{Var}\left( G\right)$ it generates is small then $G$ satisfies the maximal condition on normal subgroups?
	
	b) Is it true that a variety is small if and only if all its finitely generated groups have the Hopf property? Peter M. Neumann
\end{openproblem}

\plainproblemsection

\begin{openproblem}
	17.99. Consider the group $B = \left\langle  {a, b \mid  \left( {{ba}{b}^{-1}}\right) a{\left( ba{b}^{-1}\right) }^{-1} = {a}^{2}}\right\rangle$ introduced by Baumslag in 1969. The same relation is satisfied by the functions $f\left( x\right)  = {2x}$ and $g\left( x\right)  = {2}^{x}$ under the operation of composition in the group of germs of monotonically increasing to $\infty$ continuous functions on $\left( {0,\infty }\right)$ , where two functions are identified if they coincide for all sufficiently large arguments. Is the representation $a \rightarrow  f, b \rightarrow  g$ of the group $B$ faithful? A. Yu. Ol'shanskii
\end{openproblem}

\plainproblemsection

\begin{openproblem}
	17.100. Conjecture: A finite group is not simple if it has an irreducible complex character of odd degree vanishing on a class of odd length.
	
	If true, this implies the solvability of groups of odd order, so a proof independent of CFSG is of special interest. V. Pannone
\end{openproblem}

\plainproblemsection

\begin{openproblem}
	17.101. According to P. Hall, a group $G$ is said to be homogeneous if every isomorphism of its finitely generated subgroups is induced by an automorphism of $G$ . It is known (Higman-Neumann-Neumann) that every group is embeddable into a homogeneous one. Is the same true for the category of representations of groups? A representation(V, G)is finitely generated if $G$ is a finitely generated group, and $V$ a finitely generated module over the group algebra of $G$ . B. I. Plotkin
\end{openproblem}

\plainproblemsection

\begin{openproblem}
	17.102. We say that two subsets $A, B$ of an infinite group $G$ are separated if there exists an infinite subset $X$ of $G$ such that $1 \in  X, X = {X}^{-1}$ , and ${XAX} \cap  B = \varnothing$ . Is it true that any two disjoint subsets $A, B$ of an infinite group $G$ satisfying $\left| A\right|  < \left| G\right|$ , $\left| B\right|  < \left| G\right|$ are separated? This is so if $A, B$ are finite, or $G$ is Abelian. I. V. Protasov
\end{openproblem}

\plainproblemsection

\begin{openproblem}
	17.103. Does there exist a continuum of sets $\pi$ of primes for which every finite group possessing a Hall $\pi$ -subgroup is a ${D}_{\pi }$ -group?D. O. Revin
\end{openproblem}

\plainproblemsection

\begin{openproblem}
	17.104. Let $\Gamma$ be a finite non-oriented graph on the set of vertices $\left\{  {{x}_{1},\ldots ,{x}_{n}}\right\}$ and let ${S}_{\Gamma } = \left\langle  {{x}_{1},\ldots ,{x}_{n} \mid  {x}_{i}{x}_{j} = {x}_{j}{x}_{i} \Leftrightarrow  \left( {{x}_{i},{x}_{j}}\right)  \in  \Gamma ;{\mathfrak{A}}^{2}}\right\rangle$ be a presentation of a partially commutative metabelian group ${S}_{\Gamma }$ in the variety of all metabelian groups. Is the universal theory of the group ${S}_{\Gamma }$ decidable? V. N. Remeslennikov, E. I. Timoshenko
\end{openproblem}

\plainproblemsection

\begin{openproblem}
	17.105. An equation over a pro- $p$ -group $G$ is an expression $v\left( x\right)  = 1$ , where $v\left( x\right)$ is an element of the free pro- $p$ -product of $G$ and a free pro- $p$ -group with basis $\left\{  {{x}_{1},\ldots ,{x}_{n}}\right\}$ ; solutions are sought in the affine space ${G}^{n}$ . Is it true that a free pro- $p$ -group is equationally Noetherian, that is, for any $n$ every system of equations in ${x}_{1},\ldots ,{x}_{n}$ over this group is equivalent to some finite subsystem of it? N. S. Romanovskii
\end{openproblem}

\plainproblemsection

\begin{openproblem}
	17.107. Does $G = {\mathrm{{SL}}}_{2}\left( \mathbb{C}\right)$ contain a 2-generated free subgroup that is conjugate in $G$ to a proper subgroup of itself?
	
	A 6-generated free subgroup with this property has been found (D. Calegari, N. M. Dunfield, Proc. Amer. Math. Soc., 134, no. 11 (2006), 3131-3136). M. Sapir
\end{openproblem}

\plainproblemsection

\begin{openproblem}
	*17.108. Is the group $\left\langle  {a, b, t \mid  {a}^{t} = {ab},{b}^{t} = {ba}}\right\rangle$ linear?
	
	If not, this would be an easy example of a non-linear hyperbolic group. M. Sapir
	\par\smallskip\hrule\smallskip
	*Yes, this group is linear. Indeed, as noticed by M. Sapir, this group is the mapping torus of an irreducible atoroidal self-monomorphism of a free group; thus it is virtually special, and hence $\mathbb{Z}$ -linear by Theorem B in (M. F. Hagen, D. T. Wise, Duke Math. J., 165, no. 9 (2016), 1753-1813). (M. F. Hagen, Letter of 6 August 2018.)
	\par\smallskip\hrule\smallskip
\end{openproblem}

\plainproblemsection

\begin{openproblem}
	17.109. A non-trivial group word $w$ is uniformly elliptic in a class $\mathcal{C}$ if there is a function $f : \mathbb{N} \rightarrow  \mathbb{N}$ such that the width of $w$ in every $d$ -generator $\mathcal{C}$ -group $G$ is bounded by $f\left( d\right)$ (i. e. every element of the verbal subgroup $w\left( G\right)$ is equal to a product of $f\left( d\right)$ values of $w$ or their inverses). If $\mathcal{C}$ is a class of finite groups, this is equivalent to saying that in every finitely generated pro- $\mathcal{C}$ group $G$ the verbal subgroup $w\left( G\right)$ is closed. A. Jaikin-Zapirain (Revista Mat. Iberoamericana, 24 (2008), 617-630) proved that $w$ is uniformly elliptic in finite $p$ -groups if and only if $w \notin  {F}^{\prime \prime }{\left( {F}^{\prime }\right) }^{p}$ , where $F$ is the free group on the variables of $w$ . Is it true that $w$ is uniformly elliptic in $\mathcal{C}$ if and only if $w \notin  {F}^{\prime \prime }{\left( {F}^{\prime }\right) }^{p}$ for every prime $p$ in the case where $\mathcal{C}$ is the class of all
	
	a) finite soluble groups?
	
	b) finite groups?
	
	See also Ch. 4 of (D. Segal, Words: notes on verbal width in groups, LMS Lect. Note Series, 361, Cambridge Univ. Press, 2009. D. Segal
\end{openproblem}

\plainproblemsection

\begin{openproblem}
	17.110. a) Is it true that for each word $w$ , there is a function $h : \mathbb{N} \times  \mathbb{N} \rightarrow  \mathbb{N}$ such that the width of $w$ in every finite $p$ -group of Prüfer rank $r$ is bounded by $h\left( {p, r}\right)$ ?
	
	b) If so, can $h\left( {p, r}\right)$ be made independent of $p$ ?
	
	An affirmative answer to a) would imply A. Jaikin-Zapirain's result (ibid.) that every word has finite width in each $p$ -adic analytic pro- $p$ group, since a pro- $p$ group is $p$ -adic analytic if and only if the Prüfer ranks of all its finite quotients are uniformly bounded. D. Segal
\end{openproblem}

\plainproblemsection

\begin{openproblem}
	17.111. Let $G$ be a finite group, and $p$ a prime divisor of $\left| G\right|$ . Suppose that every maximal subgroup of a Sylow $p$ -subgroup of $G$ has a $p$ -soluble supplement in $G$ . Must $G$ be $p$ -soluble? A. N. Skiba
	
	Yes, it must (GuoHua Qian, Science China Mathematics, 56, no. 5 (2013), 1015-1018).
\end{openproblem}

\plainproblemsection

\begin{openproblem}
	17.112. A subgroup $A$ of a group $G$ is said to be $G$ -permutable in $G$ if for every subgroup $B$ of $G$ there exists an element $x \in  G$ such that $A{B}^{x} = {B}^{x}A$ . A subgroup $A$ is said to be hereditarily $G$ -permutable in $G$ if $A$ is $E$ -permutable in every subgroup $E$ of $G$ containing $A$ . Which finite non-abelian simple groups $G$ possess
	
	a) a non-trivial $G$ -permutable subgroup?
	
	b) a non-trivial hereditarily $G$ -permutable subgroup?
	
	A. N. Skiba, V. N. Tyutyanov
\end{openproblem}

\plainproblemsection

\begin{openproblem}
	17.113. Do there exist for $p > 3$ 2-generator finite $p$ -groups with deficiency zero (see 8.12) of arbitrarily high nilpotency class? J. Wiegold
\end{openproblem}

\plainproblemsection

\begin{openproblem}
	17.114. Does every generalized free product (with amalgamation) of two non-trivial groups have maximal subgroups? J. Wiegold
\end{openproblem}

\plainproblemsection

\begin{openproblem}
	17.115. Can a locally free non-abelian group have non-trivial Frattini subgroup? J. Wiegold 17.117. (Well-known problem). If groups $A$ and $B$ have decidable elementary theories ${Th}\left( A\right)$ and ${Th}\left( B\right)$ , must ${Th}\left( {A * B}\right)$ be decidable? O. Kharlampovich
\end{openproblem}

\plainproblemsection

\begin{openproblem}
	17.116. Let $n\left( G\right)$ be the maximum of positive integers $n$ such that the $n$ -th direct power of a finite simple group $G$ is 2-generated. Is it true that $n\left( G\right)  \geq  \sqrt{\left| G\right| }$ ?
	
	A. Erfanian, J. Wiegold
	
	Yes, it is, even with $n\left( G\right)  > 2\sqrt{\left| G\right| }$ (A. Maróti, M. C. Tamburini, Commun. Algebra, 41, no. 1 (2013), 34-49).
\end{openproblem}

\plainproblemsection

\begin{openproblem}
	17.118. Suppose that a finite $p$ -group $G$ has a subgroup of exponent $p$ and of index $p$ . Must $G$ also have a characteristic subgroup of exponent $p$ and of index bounded in terms of $p$ ? E. I. Khukhro
\end{openproblem}

\plainproblemsection

\begin{openproblem}
	17.119. Suppose that a finite soluble group $G$ admits a soluble group of automor-phisms $A$ of coprime order such that ${C}_{G}\left( A\right)$ has rank $r$ . Let $\left| A\right|$ be the product of $l$ not necessarily distinct primes. Is there a linear function $f$ such that $G/{F}_{f\left( l\right) }\left( G\right)$ has $\left( {\left| A\right| , r}\right)$ -bounded rank, where ${F}_{f\left( l\right) }\left( G\right)$ is the $f\left( l\right)$ -th Fitting subgroup?
	
	An exponential function $f$ with this property was found in (E. Khukhro, V. Mazurov, Groups St. Andrews 2005, vol. II, Cambridge Univ. Press, 2007, 564- 585). It is also known that $\left| {G/{F}_{{2l} + 1}\left( G\right) }\right|$ is bounded in terms of $\left| {{C}_{G}\left( A\right) }\right|$ and $\left| A\right|$ (Hartley-Isaacs). E. I. Khukhro
\end{openproblem}

\plainproblemsection

\begin{openproblem}
	17.120. Is a residually finite group all of whose subgroups of infinite index are finite necessarily cyclic-by-finite? N. S. Chernikov
\end{openproblem}

\plainproblemsection

\begin{openproblem}
	17.121. Let $G$ be a group whose set of elements is the real numbers and which is "nicely definable" (see below). Does $G$ being ${\aleph }_{\omega }$ -free imply it being free if "nicely definable" means
	
	a) being ${F}_{\sigma }$ ?
	
	b) being Borel?
	
	c) being analytic?
	
	d) being projective $\mathbb{L}\left\lbrack  \mathbb{R}\right\rbrack$ ?
	
	See https://arxiv.org/pdf/math/0212250.pdf for justification of the restrictions. S. Shelah
\end{openproblem}

\plainproblemsection

\begin{openproblem}
	17.122. The same questions as 17.121 for an abelian group $G$ . S. Shelah
\end{openproblem}

\plainproblemsection

\begin{openproblem}
	17.123. Do there exist finite groups ${G}_{1},{G}_{2}$ such that ${\pi }_{e}\left( {G}_{1}\right)  = {\pi }_{e}\left( {G}_{2}\right) , h\left( {{\pi }_{e}\left( {G}_{1}\right) }\right)  <$ $\infty$ , and each non-abelian composition factors of each of the groups ${G}_{1},{G}_{2}$ is not isomorphic to a section of the other?
	
	For definitions see Archive, 13.63.W. J. Shi
\end{openproblem}

\plainproblemsection

\begin{openproblem}
	17.124. Is the set of finitely presented metabelian groups recursively enumerable?
	
	V. Shpilrain
\end{openproblem}

\plainproblemsection

\begin{openproblem}
	17.125. Does every finite group $G$ contain a pair of conjugate elements $a, b$ such that $\pi \left( G\right)  = \pi \left( {\langle a, b\rangle }\right)$ ? This is true for soluble groups.
	
	Comment of 2013: it was proved in (A. Lucchini, M. Morigi, P. Shumyatsky, Forum Math., $\mathbf{{24}}\left( {2012}\right) ,{875} - {887})$ that every finite group $G$ contains a 2-generator subgroup $H$ such that $\pi \left( G\right)  = \pi \left( H\right)$ . P. Shumyatsky
\end{openproblem}

\plainproblemsection

\begin{openproblem}
	17.126. Suppose that $G$ is a residually finite group satisfying the identity ${\left\lbrack  x, y\right\rbrack  }^{n} = 1$ . Must $\left\lbrack  {G, G}\right\rbrack$ be locally finite?
	
	An equivalent question: let $G$ be a finite soluble group satisfying the identity ${\left\lbrack  x, y\right\rbrack  }^{n} = 1$ ; is the Fitting height of $G$ bounded in terms of $n$ ? Cf. Archive,13.34
	
	P. Shumyatsky
\end{openproblem}

\plainproblemsection

\begin{openproblem}
	17.127. Suppose that a finite soluble group $G$ of derived length $d$ admits an elementary abelian $p$ -group of automorphisms $A$ of order ${p}^{n}$ such that ${C}_{G}\left( A\right)  = 1$ . Must $G$ have a normal series of $n$ -bounded length with nilpotent factors of(p, n, d)-bounded nilpotency class?
	
	This is true for $p = 2$ . An affirmative answer would follow from an affirmative answer to 11.125. P. Shumyatsky
\end{openproblem}

\plainproblemsection

\begin{openproblem}
	17.128. Let $T$ be a finite $p$ -group admitting an elementary abelian group of auto-morphisms $A$ of order ${p}^{2}$ such that in the semidirect product $P = {TA}$ every element of $P \smallsetminus  T$ has order $p$ . Does it follow that $T$ is of exponent $p$ ? E. Jabara
\end{openproblem}

\plainproblemsection

\begin{openproblem}
	*18.1. Given a group $G$ of finite order $n$ , does there necessarily exist a bijection $f$ from $G$ onto a cyclic group of order $n$ such that for each element $x \in  G$ , the order of $x$ divides the order of $f\left( x\right)$ ? An affirmative answer in the case where $G$ is solvable was given by F. Ladisch. I. M. Isaacs
	
	*Yes, it does (M. Amiri, J. Pure Applied Algebra, 228, no. 7 (2024), 107632).
\end{openproblem}

\plainproblemsection

\begin{openproblem}
	18.2. Let the group $G = {AB}$ be the product of two Chernikov subgroups $A$ and $B$ each of which has an abelian subgroup of index at most 2. Is $G$ soluble? B. Amberg
\end{openproblem}

\plainproblemsection

\begin{openproblem}
	18.3. Let the group $G = {AB}$ be the product of two central-by-finite subgroups $A$ and $B$ . By a theorem of N. Chernikov, $G$ is soluble-by-finite. Is $G$ metabelian-by-finite? B. Amberg
\end{openproblem}

\plainproblemsection

\begin{openproblem}
	18.4. (A. Rhemtulla, S. Sidki). a) Is a group of the form $G = {ABA}$ with cyclic subgroups $A$ and $B$ always soluble? (This is known to be true if $G$ is finite.)
	
	b) The same question when in addition $A$ and $B$ are conjugate in $G$ . $B$ . Amberg
\end{openproblem}

\plainproblemsection

\begin{openproblem}
	18.5. Let the (soluble) group $G = {AB}$ with finite torsion-free rank ${r}_{0}\left( G\right)$ be the product of two subgroups $A$ and $B$ . Is the equation ${r}_{0}\left( G\right)  = {r}_{0}\left( A\right)  + {r}_{0}\left( B\right)  - {r}_{0}\left( {A \cap  B}\right)$ valid in this case? It is known that the inequality $\leq$ always holds. B. Amberg
\end{openproblem}

\plainproblemsection

\begin{openproblem}
	18.6. Is every finitely generated one-relator group residually amenable?
	
	G. N. Arzhantseva
\end{openproblem}

\plainproblemsection

\begin{openproblem}
	18.7. Let $n \geq  {665}$ be an odd integer. Is it true that the group of outer automorphisms Out $\left( {B\left( {m, n}\right) }\right)$ of the free Burnside group $B\left( {m, n}\right)$ for $m > 2$ is complete, that is, has trivial center and all its automorphisms are inner? V. S. Atabekyan
\end{openproblem}

\plainproblemsection

\begin{openproblem}
	18.8. Let $\mathcal{X}$ be a class of finite simple groups such that $\pi \left( \mathcal{X}\right)  = \operatorname{char}\left( \mathcal{X}\right)$ . A formation of finite groups $\mathfrak{F}$ is said to be $\mathcal{X}$ -saturated if a finite group $G$ belongs to $\mathfrak{F}$ whenever the factor group $G/\Phi \left( {{O}_{\mathcal{X}}\left( G\right) }\right)$ is in $\mathfrak{F}$ , where ${O}_{\mathcal{X}}\left( G\right)$ is the largest normal subgroup of $G$ whose composition factors are in $\mathcal{X}$ . Is every $\mathcal{X}$ -saturated formation $\mathcal{X}$ -local in the sense of Förster? (See the definition in (P. Förster, Publ. Sec. Mat. Univ. Autònoma Barcelona, 29, no. 2-3 (1985), 39-76).) A. Ballester-Bolinches
\end{openproblem}

\plainproblemsection

\begin{openproblem}
	18.9. Does there exist a subgroup-closed saturated formation $\mathfrak{F}$ of finite groups properly contained in ${\mathfrak{E}}_{\pi }$ , where $\pi  = \operatorname{char}\left( \mathcal{F}\right)$ , satisfying the following property: if $G \in  \mathfrak{F}$ , then there exists a prime $p$ (depending on the group $G$ ) such that the wreath product ${C}_{p}\wr G$ belongs to $\mathfrak{F}$ , where ${C}_{p}$ is the cyclic group of order $p$ ?
	
	A. Ballester-Bolinches
	
	Yes, it exists (S. F. Kamornikov, Trudy. Inst. Mat. (Minsk), 24, no. 1 (2016), 30-33).
\end{openproblem}

\plainproblemsection

\begin{openproblem}
	18.10. A formation $\mathfrak{F}$ of finite groups satisfies the Wielandt property for residuals if whenever $U$ and $V$ are subnormal subgroups of $\langle U, V\rangle$ in a finite group $G$ , then the $\mathfrak{F}$ -residual $\langle U, V{\rangle }^{\mathfrak{F}}$ of $\langle U, V\rangle$ coincides with $\left\langle  {{U}^{\mathfrak{F}},{V}^{\mathfrak{F}}}\right\rangle$ . Does every Fitting formation $\mathfrak{F}$ satisfy the Wielandt property for residuals? A. Ballester-Bolinches
\end{openproblem}

\plainproblemsection

\begin{openproblem}
	18.11. Let $G$ and $H$ be subgroups of the automorphism group $\operatorname{Aut}\left( {F}_{n}\right)$ of a free group ${F}_{n}$ of rank $n \geq  2$ . Is it true that the free product $G * H$ embeds in the automorphism group $\operatorname{Aut}\left( {F}_{m}\right)$ for some $m$ ? V. G. Bardakov
\end{openproblem}

\plainproblemsection

\begin{openproblem}
	18.12. Let $G$ be a finitely generated group of intermediate growth. Is it true that there is a positive integer $m$ such that every element of the derived subgroup ${G}^{\prime }$ is a product of at most $m$ commutators?
	
	This assertion is valid for groups of polynomial growth, since they are almost nilpotent by Gromov's theorem. On the other hand, there are groups of exponential growth (for example, free groups) for which this assertion is not true. ${VG}$ Bardakov
\end{openproblem}

\plainproblemsection

\begin{openproblem}
	18.13. (D. B. A. Epstein). Is it true that the group $H = \left( {{\mathbb{Z}}_{3} \times  \mathbb{Z}}\right)  * \left( {{\mathbb{Z}}_{2} \times  \mathbb{Z}}\right)  =$ $\left\langle  {x, y, z, t \mid  {x}^{3} = {z}^{2} = \left\lbrack  {x, y}\right\rbrack   = \left\lbrack  {z, t}\right\rbrack   = 1}\right\rangle$ cannot be defined by three relators in the generators $x, y, z, t$ ?
	
	It is known that the relation module of the group $H$ has rank 3 (K. W. Gruenberg, P. A. Linnell, J. Group Theory, 11, no. 5 (2008), 587-608). An affirmative answer would give a solution of the relation gap problem. V. G. Bardakov, M. V. Neshchadim
\end{openproblem}

\plainproblemsection

\begin{openproblem}
	*18.14. For an automorphism $\varphi  \in  \operatorname{Aut}\left( G\right)$ of a group $G$ , let ${\left\lbrack  e\right\rbrack  }_{\varphi } = \left\{  {{g}^{-1}{g}^{\varphi } \mid  g \in  G}\right\}$ .
	
	Conjecture: If ${\left\lbrack  e\right\rbrack  }_{\varphi }$ is a subgroup for every $\varphi  \in  \operatorname{Aut}\left( G\right)$ , then the group $G$ is nilpotent. If in addition $G$ is finitely generated, then $G$ is abelian.
	
	V. G. Bardakov, M. V. Neshchadim, T. R. Nasybullov
	
	*The conjecture is disproved (C. Nicotera, Arch. Math., 123, no. 3 (2024), 225- 232).
\end{openproblem}

\plainproblemsection

\begin{openproblem}
	18.15. For an automorphism $\varphi  \in  \operatorname{Aut}\left( G\right)$ of a group $G$ , let ${\left\lbrack  e\right\rbrack  }_{\varphi } = \left\{  {{g}^{-1}{g}^{\varphi } \mid  g \in  G}\right\}$ . Is it true that if a group $G$ has trivial center, then there is an inner automorphism $\varphi$ such that ${\left\lbrack  e\right\rbrack  }_{\varphi }$ is not a subgroup?
	
	V. G. Bardakov, M. V. Neshchadim, T. R. Nasybullov
	
	No, not always; see V. V. Bludov's Example 2 in (D. Gonçalves, T. Nasybullov, Com-mun. Algebra, 47, no. 3 (2019), 930-944).
\end{openproblem}

\plainproblemsection

\begin{openproblem}
	18.16. Is it true that any definable endomorphism of any ordered abelian group is of the form $x \mapsto  {rx}$ , for some rational number $r$ ? O. V. Belegradek
\end{openproblem}

\plainproblemsection

\begin{openproblem}
	18.17. Is there a torsion-free group which is finitely presented in the quasi-variety of torsion-free groups but not finitely presentable in the variety of all groups?
	
	O. V. Belegradek
	
	Yes, there is (A.I. Budkin, Siberian Electron. Math. Rep., 14 (2017), 937-945, http: //semr.math.nsc.ru).
\end{openproblem}

\plainproblemsection

\begin{openproblem}
	18.18. Mal'cev proved that the set of sentences that hold in all finite groups is not computably enumerable, although its complement is. Is it true that both the set of sentences that hold in almost all finite groups and its complement are not computably enumerable? O. V. Belegradek
\end{openproblem}

\plainproblemsection

\begin{openproblem}
	18.19. Is any torsion-free, relatively free group of infinite rank not ${\aleph }_{1}$ -homogeneous? This is true for group varieties in which all free groups are residually finite (O. Belegradek, Arch. Math. Logic, 51 (2012), 781-787). O. V. Belegradek
\end{openproblem}

\plainproblemsection

\begin{openproblem}
	18.20. Characters $\varphi$ and $\psi$ of a finite group $G$ are said to be semiproportional if they are not proportional and there is a normal subset $M$ of $G$ such that ${\left. \varphi \right| }_{M}$ is proportional to ${\left. \psi \right| }_{M}$ and ${\left. \varphi \right| }_{G \smallsetminus  M}$ is proportional to ${\left. \psi \right| }_{G \smallsetminus  M}$ .
	
	Conjecture: If $\varphi$ and $\psi$ are semiproportional irreducible characters of a finite group, then $\varphi \left( 1\right)  = \psi \left( 1\right)$ . V. A. Belonogov
\end{openproblem}

\plainproblemsection

\begin{openproblem}
	18.21. (B.H.Neumann and H.Neumann). Fix an integer $d \geq  2$ . If $\mathfrak{V}$ is a variety of groups such that all $d$ -generated groups in $\mathfrak{V}$ are finite, must $\mathfrak{V}$ be locally finite?
	
	G. M. Bergman
\end{openproblem}

\plainproblemsection

\begin{openproblem}
	18.22. (H. Neumann). Is the Kostrikin variety of all locally finite groups of given prime exponent $p$ determined by finitely many identities?
\end{openproblem}

\plainproblemsection

\begin{openproblem}
	*18.23. The normal covering number of the symmetric group ${S}_{n}$ of degree $n$ is the minimum number $\gamma \left( {S}_{n}\right)$ of proper subgroups ${H}_{1},\ldots ,{H}_{\gamma \left( {S}_{n}\right) }$ of ${S}_{n}$ such that every element of ${S}_{n}$ is conjugate to an element of ${H}_{i}$ , for some $i = 1,\ldots ,\gamma \left( {S}_{n}\right)$ . Write $n = {p}_{1}^{{\alpha }_{1}}\cdots {p}_{r}^{{\alpha }_{r}}$ for primes ${p}_{1} < \cdots  < {p}_{r}$ and positive integers ${\alpha }_{1},\ldots ,{\alpha }_{r}$ .
	
	Conjecture:
	
	$$
	\gamma \left( {S}_{n}\right)  = \left\{  \begin{array}{ll} \frac{n}{2}\left( {1 - \frac{1}{{p}_{1}}}\right) & \text{ if }r = 1\text{ and }{\alpha }_{1} = 1 \\  \frac{n}{2}\left( {1 - \frac{1}{{p}_{1}}}\right)  + 1 & \text{ if }r = 1\text{ and }{\alpha }_{1} \geq  2 \\  \frac{n}{2}\left( {1 - \frac{1}{{p}_{1}}}\right) \left( {1 - \frac{1}{{p}_{2}}}\right)  + 1 & \text{ if }r = 2\text{ and }{\alpha }_{1} + {\alpha }_{2} = 2 \\  \frac{n}{2}\left( {1 - \frac{1}{{p}_{1}}}\right) \left( {1 - \frac{1}{{p}_{2}}}\right)  + 2 & \text{ if }r \geq  2\text{ and }{\alpha }_{1} + \cdots  + {\alpha }_{r} \geq  3 \end{array}\right.
	$$
	
	This is the strongest form of the conjecture. We would be also interested in a proof that this holds for $n$ sufficiently large. The result for $r \leq  2$ , which includes the first three cases above, is proved for $n$ odd (D. Bubboloni, C. E. Praeger, J. Combin. Theory (A),118 (2011),2000-2024). When $r \geq  3$ we know that ${cn} \leq  \gamma \left( {S}_{n}\right)  \leq  \frac{2}{3}n$ for some positive constant $c$ (D. Bubboloni, C. E. Praeger, P. Spiga, J. Algebra,390 (2013) 199-215). We showed that the conjectured value for $\gamma \left( {S}_{n}\right)$ is an upper bound, by constructing a normal covering for ${S}_{n}$ with this number of conjugacy classes of maximal subgroups, and gave further evidence for the truth of the conjecture in other cases (D. Bubboloni, C. E. Praeger, P. Spiga, Int. J. Group Theory, 3, no. 2 (2014), 57-75).
	
	*A. Maróti showed that the conjecture is incorrect for odd $n$ ; his counterexample is presented in $§2$ of (D. Bubboloni, C. E. Praeger, P. Spiga, Monatsh. Math.,191 (2020), 229-247). D. Bubboloni, C. E. Praeger, P. Spiga
\end{openproblem}

\plainproblemsection

\begin{openproblem}
	18.24. For a group $G$ , a function $\phi  : G \rightarrow  \mathbb{R}$ is a quasimorphism if there is a least non-negative number $D\left( \phi \right)$ (called the defect) such that $\left| {\phi \left( {gh}\right)  - \phi \left( g\right)  - \phi \left( h\right) }\right|  \leq  D\left( \phi \right)$ for all $g, h \in  G$ . A quasimorphism is homogeneous if in addition $\phi \left( {g}^{n}\right)  = {n\phi }\left( g\right)$ for all $g \in  G$ . Let $\phi$ be a homogeneous quasimorphism of a free group $F$ . For any quasimorphism $\psi$ of $F$ (not required to be homogeneous) with $\left| {\phi  - \psi }\right|  <$ const, we must have $D\left( \psi \right)  \geq  D\left( \phi \right) /2$ . Is it true that there is some $\psi$ with $D\left( \psi \right)  = D\left( \phi \right) /2$ ?
	
	M. Burger, D. Calegari
\end{openproblem}

\plainproblemsection

\begin{openproblem}
	18.25. Let $\operatorname{form}\left( G\right)$ be the formation generated by a finite group $G$ . Suppose that $G$ has a unique composition series $1 \vartriangleleft  {G}_{1} \vartriangleleft  {G}_{2} \vartriangleleft  G$ and the consecutive factors of this series are ${Z}_{p}, X,{Z}_{q}$ , where $p$ and $q$ are primes and $X$ is a non-abelian simple group. Is it true that $\operatorname{form}\left( G\right)$ has infinitely many subformations if and only if $p = q$ ?
	
	V. P. Burichenko
\end{openproblem}

\plainproblemsection

\begin{openproblem}
	18.26. Suppose that a finite group $G$ has a normal series $1 < {G}_{1} < {G}_{2} < G$ such that the groups ${G}_{1}$ and ${G}_{2}/{G}_{1}$ are elementary abelian $p$ -groups, $G/{G}_{2} \cong  {A}_{5} \times  {A}_{5}$ (where ${A}_{5}$ is the alternating group of degree 5), ${G}_{1}$ and ${G}_{2}/{G}_{1}$ are minimal normal subgroups of $G$ and $G/{G}_{1}$ , respectively. Is it true that form(G)has finitely many subformations? V. P. Burichenko
\end{openproblem}

\plainproblemsection

\begin{openproblem}
	18.27. Does there exist an algorithm that determines whether there are finitely many subformations in $\operatorname{form}\left( G\right)$ for a given finite group $G$ ?
\end{openproblem}

\plainproblemsection

\begin{openproblem}
	18.28. Is it true that any subformation of every one-generator formation $\operatorname{form}\left( G\right)$ is also one-generator? V. P. Burichenko
\end{openproblem}

\plainproblemsection

\begin{openproblem}
	18.29. Let $\mathfrak{F}$ be a Fitting class of finite soluble groups which contains every soluble group $G = {AB}$ , where $A$ and $B$ are abnormal $\mathfrak{F}$ -subgroups of $G$ . Is $\mathfrak{F}$ a formation?
	
	A. F. Vasil’ev
\end{openproblem}

\plainproblemsection

\begin{openproblem}
	18.30. A subgroup $H$ of a group $G$ is called $\mathbb{P}$ -subnormal in $G$ if either $H = G$ , or there is a chain of subgroups ${H}_{0} \subset  {H}_{1} \subset  \cdots  \subset  {H}_{n - 1} \subset  {H}_{n} = G$ such that $\left| {{H}_{i} : {H}_{i - 1}}\right|$ is a prime for all $i = 1,\ldots , n$ . Must a finite group be soluble if every Shmidt subgroup of it is $\mathbb{P}$ -subnormal? A. F. Vasil’ev, T. I. Vasil’eva, V. N. Tyutyanov Yes, it must (V. N. Tyutyanov, Probl. Fiz. Mat. Tekhn., 2015, no. 1 (22), 88-91 (Russian)).
\end{openproblem}

\plainproblemsection

\begin{openproblem}
	*18.31. Let $\pi$ be a set of primes. Is it true that in any ${D}_{\pi }$ -group $G$ (see Archive, 3.62) there are three Hall $\pi$ -subgroups whose intersection coincides with ${O}_{\pi }\left( G\right)$ ?
	
	E. P. Vdovin, D. O. Revin
	
	*No, it is not true, for example, for $G$ being an extension of ${L}_{2}\left( {27}\right)$ by a field automorphism of order 3, in which $H$ is an extension of a Borel subgroup of ${L}_{2}\left( {27}\right)$ by a field automorphism of order 3 (V.I. Zenkov, Abstracts of Int. Conf. "Mal'tsev Meeting 2020", Novosibirsk, 2020, p. 149).
\end{openproblem}

\plainproblemsection

\begin{openproblem}
	18.32. Is every Hall subgroup of a finite group pronormal in its normal closure? E. P. Vdovin, D. O. Revin
	
	No, not always (M. N. Nesterov, Siberian Math. J., 58, no. 1 (2017), 128-133).
\end{openproblem}

\plainproblemsection

\begin{openproblem}
	18.33. A group in which the derived subgroup of every 2-generated subgroup is cyclic is called an Alperin group. Is there a bound for the derived length of finite Alperin groups?
	
	G. Higman proved that finite Alperin groups are soluble, and finite Alperin $p$ - groups have bounded derived length, see Archive 17.46. B. M. Veretennikov Yes, there is: it is at most 6. Indeed, in (P. Longobardi, M. Maj, H. Smith, Rend. Semin. Mat. Univ. Padova, 115 (2006), 29-40) it was proved that finite Alperin groups of odd order are metabelian, and any finite Alperin group is supersoluble, so the elements of odd order form a metabelian normal subgroup, while finite Alperin 2-groups have derived length at most 4 by (B. Wilkens, J. Group Theory, 17, no. 1 (2014), 151-174).
\end{openproblem}

\plainproblemsection

\begin{openproblem}
	18.34. Let $G$ be a topological group and $S = {G}_{0} \leq  \ldots  \leq  {G}_{n} = G$ a subnormal series of closed subgroups. The infinite-length of $S$ is the cardinality of the set of infinite factors ${G}_{i}/{G}_{i - 1}$ . The virtual length of $G$ is the supremum of the infinite lengths taken over all such series of $G$ . It is known that if $G$ is a pronilpotent group with finite virtual length then every closed subnormal subgroup is topologically finitely generated (N. Gavioli, V. Monti, C. M. Scoppola, J. Austral. Math. Soc., 95, no. 3 (2013), 343-355). Let $G$ be a pro- $p$ group (or more generally a pronilpotent group). If every closed subnormal subgroup of $G$ is tolopologically finitely generated, is it true that $G$ has finite virtual length? N. Gavioli, V. Monti, C. M. Scoppola
\end{openproblem}

\plainproblemsection

\begin{openproblem}
	18.35. A family $\mathcal{F}$ of group homomorphisms $A \rightarrow  B$ is separating if for every nontrivial $a \in  A$ there is $f \in  \mathcal{F}$ such that $f\left( a\right)  \neq  1$ , and discriminating if for any finitely many nontrivial elements ${a}_{1},\ldots ,{a}_{n} \in  A$ there is $f \in  \mathcal{F}$ such that $f\left( {a}_{i}\right)  \neq  1$ for all $i = 1,\ldots , n$ . Let $G$ be a relatively free group of rank 2 in the variety of metabelian groups. Let $T$ be a metabelian group in which the centralizer of every nontrivial element is abelian. If $T$ admits a separating family of surjective homomorphisms $T \rightarrow  G$ must it also admit a discriminating family of surjective homomorphisms $T \rightarrow  G$ ?
	
	A positive answer would give a metabelian analogue of a classical theorem of B. Baumslag. A. Gaglione, D. Spellman
\end{openproblem}

\plainproblemsection

\begin{openproblem}
	*18.36. There are groups of cardinality at most ${2}^{{\aleph }_{0}}$ , even nilpotent of class 2, that cannot be embedded in $S \mathrel{\text{:=}} \operatorname{Sym}\left( \mathbb{N}\right)$ (V. A. Churkin, Algebra and model theory (Novosibirsk State Tech. Univ.), 5 (2005), 39-43 (Russian)). One condition which might characterize subgroups of $S$ is the following. Consider the metric $d$ on $S$ where for all $x \neq  y$ we define $d\left( {x, y}\right)  \mathrel{\text{:=}} {2}^{-k}$ when $k$ is the least element of the set $\mathbb{N}$ such that ${k}^{x} \neq  {k}^{y}$ . Then(S, d)is a separable topological group, and so each subgroup of $S$ (not necessarily a closed subgroup of(S, d)) is a separable topological group under the induced metric. Is it true that every group $G$ for which there is a metric ${d}^{\prime }$ such that $\left( {G,{d}^{\prime }}\right)$ is a separable topological group is (abstractly) embeddable in $S$ ?
	
	Note that any such group $G$ can be embedded as a section in $S$ . J. D. Dixon
	\par\smallskip\hrule\smallskip
	*No, there exists an uncountable connected Polish group whose every abstract homomorphism into $S$ is trivial (A. Kaïchouh, F. Le Maître, Bull. London Math. Soc., 47 (2015), 996-1009). (S. Corson, Letter of 5 June 2024).
	\par\smallskip\hrule\smallskip
\end{openproblem}

\plainproblemsection

\begin{openproblem}
	18.37. (Well-known problem). Is every locally graded group of finite rank almost locally soluble? M. Dixon
\end{openproblem}

\plainproblemsection

\begin{openproblem}
	*18.38. Let ${E}_{\pi }$ denote the class of finite groups that contain a Hall $\pi$ -subgroup. Does the inclusion ${E}_{{\pi }_{1}} \cap  {E}_{{\pi }_{2}} \subseteq  {E}_{{\pi }_{1} \cap  {\pi }_{2}}$ hold for arbitrary sets of primes ${\pi }_{1}$ and ${\pi }_{2}$ ? A. V. Zavarnitsine
	
	*Yes, it does (A. Buturlakin, N. Yang, Preprint, 2025, https://arxiv.org/abs/ 2501.05865).
\end{openproblem}

\plainproblemsection

\begin{openproblem}
	18.39. Conjecture: Let $G$ be a hyperbolic group. Then every 2-dimensional rational homology class is virtually represented by a sum of closed surface subgroups, that is, for any $\alpha  \in  {H}_{2}\left( {G;\mathbb{Q}}\right)$ there are finitely many closed oriented surfaces ${S}_{i}$ and injective homomorphisms ${\rho }_{i} : {\pi }_{1}\left( {S}_{i}\right)  \rightarrow  G$ such that $\mathop{\sum }\limits_{i}\left\lbrack  {S}_{i}\right\rbrack   = {n\alpha }$ , where $\left\lbrack  {S}_{i}\right\rbrack$ denotes the image of the fundamental class of ${S}_{i}$ in ${H}_{2}\left( G\right)$ . D. Calegari
\end{openproblem}

\plainproblemsection

\begin{openproblem}
	18.40. The commutator length ${cl}\left( g\right)$ of $g \in  \left\lbrack  {G, G}\right\rbrack$ is the least number of commutators in $G$ whose product is $g$ , and the stable commutator length is $\operatorname{scl}\left( g\right)  \mathrel{\text{:=}}$ $\mathop{\lim }\limits_{{n \rightarrow  \infty }}{cl}\left( {g}^{n}\right) /n$ .
	
	Conjecture: Let $G$ be a hyperbolic group. Then the stable commutator length takes on rational values on $\left\lbrack  {G, G}\right\rbrack$ . D. Calegari
\end{openproblem}

\plainproblemsection

\begin{openproblem}
	18.41. Let $F$ be a free group of rank 2 .
	
	a) It is known that $\operatorname{scl}\left( g\right)  = 0$ only for $g = 1$ , and $\operatorname{scl}\left( g\right)  \geq  1/2$ for all $1 \neq  g \in  \left\lbrack  {F, F}\right\rbrack$ . Is $1/2$ an isolated value?
	
	b) Are there any intervals $J$ in $\mathbb{R}$ such that the set of values of ${scl}$ on $\left\lbrack  {F, F}\right\rbrack$ is dense in $J$ ?
	
	c) Is there some $T \in  \mathbb{R}$ such that every rational number $\geq  T$ is a value of $\operatorname{scl}\left( g\right)$ for $g \in  \left\lbrack  {F, F}\right\rbrack$ ? D. Calegari
\end{openproblem}

\plainproblemsection

\begin{openproblem}
	18.42. For an orientation-preserving homeomorphism $f : {S}^{1} \rightarrow  {S}^{1}$ of a unit circle ${S}^{1}$ , let $\widetilde{f}$ be its lifting to a homeomorphism of $\mathbb{R}$ ; then the rotation number of $f$ is defined to be the limit $\mathop{\lim }\limits_{{n \rightarrow  \infty }}\left( {{\widetilde{f}}^{n}\left( x\right)  - x}\right) /n$ (which is independent of the point $\left. {x \in  {S}^{1}}\right)$ . Let $F$ be a free group of rank 2 with generators $a, b$ . For $w \in  F$ and $r, s \in  \mathbb{R}$ , let $R\left( {w, r, s}\right)$ denote the maximum value of the (real-valued) rotation number of $w$ , under all representations from $F$ to the universal central extension of ${\operatorname{Homeo}}^{ + }\left( {S}^{1}\right)$ for which the rotation number of $a$ is $r$ , and the rotation number of $b$ is $s$ .
	
	a) If $r, s$ are rational, must $R\left( {w, r, s}\right)$ be rational?
	
	b) Let $R\left( {w, r-, s - }\right)$ denote the supremum of the rotation number of $w$ (as above) under all representations for which $a$ and $b$ are conjugate to rotations through $r$ and $s$ , respectively (in the universal central extension of ${\operatorname{Homeo}}^{ + }\left( {S}^{1}\right)$ ). Is $R\left( {w, r-, s - }\right)$ always rational if $r$ and $s$ are rational?
	
	c) Weak Slippery Conjecture: Let $w$ be a word containing only positive powers of $a$ and $b$ . A pair of values(r, s)is slippery if there is a strict inequality $R\left( {w,{r}^{\prime },{s}^{\prime }}\right)  <$ $R\left( {w, r-, s - }\right)$ for all ${r}^{\prime } < r,{s}^{\prime } < s$ . Is it true that then $R\left( {w, r-, s - }\right)  = {h}_{a}\left( w\right) r +$ ${h}_{b}\left( w\right) s$ , where ${h}_{a}\left( w\right)$ counts the number of $a$ ’s in $w$ , and ${h}_{b}\left( w\right)$ counts the number of $b$ ’s in $w$ ?
	
	d) Slippery Conjecture: More precisely, is it always (without assuming(r, s)being slippery) true that if $w$ has the form $w = {a}^{{\alpha }_{1}}{b}^{{\beta }_{1}}\cdots {a}^{{\alpha }_{m}}{b}^{{\beta }_{m}}$ (all positive) and $R\left( {w, r, s}\right)  = p/q$ where $p/q$ is reduced, then $\left| {p/q - {h}_{a}\left( w\right) r - {h}_{b}\left( w\right) s}\right|  \leq  m/q$ ?
\end{openproblem}

\plainproblemsection

\begin{openproblem}
	18.43. a) Do there exist elements $u, v$ of a 2-generator free group $F\left( {a, b}\right)$ such that $u, v$ are not conjugate in $F\left( {a, b}\right)$ but for any matrices $A, B \in  {GL}\left( {3,\mathbb{C}}\right)$ we have $\operatorname{trace}\left( {u\left( {A, B}\right) }\right)  = \operatorname{trace}\left( {v\left( {A, B}\right) }\right) ?$
	
	b) The same question if we replace ${GL}\left( {3,\mathbb{C}}\right)$ by ${SL}\left( {3,\mathbb{C}}\right)$ .I. Kapovich
\end{openproblem}

\plainproblemsection

\begin{openproblem}
	18.44. Let $G$ be a finite non-abelian group and $V$ a finite faithful irreducible $G$ - module. Suppose that $M = \left| {G/{G}^{\prime }}\right|$ is the largest orbit size of $G$ on $V$ , and among orbits of $G$ on $V$ there are exactly two orbits of size $M$ . Does this imply that $G$ is dihedral of order 8, and $\left| V\right|  = 9$ ? T. M. Keller
\end{openproblem}

\plainproblemsection

\begin{openproblem}
	18.45. Let $G$ be a finite $p$ -group of maximal class such that all its irreducible characters are induced from linear characters of normal subgroups. Let $S$ be the set of the derived subgroups of the members of the central series of $G$ . Is there a logarithmic bound for the derived length of $G$ in terms of $\left| S\right|$ ? T. M. Keller
\end{openproblem}

\plainproblemsection

\begin{openproblem}
	18.46. Every finite group $G$ can be embedded in a group $H$ in such a way that every element of $G$ is a square of an element of $H$ . The overgroup $H$ can be chosen such that $\left| H\right|  \leq  2{\left| G\right| }^{2}$ . Is this estimate sharp?
	
	It is known that the best possible estimate cannot be better than $\left| H\right|  \leq  {\left| G\right| }^{2}$ (D. V. Baranov, Ant. A. Klyachko, Siber. Math. J., 53, no. 2 (2012), 201-206).
	
	Ant. A. Klyachko
\end{openproblem}

\plainproblemsection

\begin{openproblem}
	18.47. Is it algorithmically decidable whether a group generated by three given class transpositions (for the definition, see 17.57)
	
	a) has only finite orbits on $\mathbb{Z}$ ?
	
	b) acts transitively on the set of nonnegative integers in its support?
	
	A difficult case is the group $\left\langle  {{\tau }_{1\left( 2\right) ,4\left( 6\right) },{\tau }_{1\left( 3\right) ,2\left( 6\right) },{\tau }_{2\left( 3\right) ,4\left( 6\right) }}\right\rangle$ , which acts transitively on $\mathbb{N} \smallsetminus  0\left( 6\right)$ if and only if Collatz’ ${3n} + 1$ conjecture is true.
\end{openproblem}

\plainproblemsection

\begin{openproblem}
	18.48. Is it true that there are only finitely many integers which occur as orders of products of two class transpositions? (For the definition, see 17.57.) S. Kohl
\end{openproblem}

\plainproblemsection

\begin{openproblem}
	*18.49. Let $n \in  \mathbb{N}$ . Is it true that for any $a, b, c \in  \mathbb{N}$ satisfying $1 < a, b, c \leq  n - 2$ the symmetric group ${\mathrm{S}}_{n}$ has elements of order $a$ and $b$ whose product has order $c$ ?
	
	S. Kohl
	
	*Yes, it is (G. A. Miller, Amer. J. Math., 22, no. 2 (1900), 185-190). Another solution is in (J. König, Eur. J. Comb., 57 (2016), 50-56).
\end{openproblem}

\plainproblemsection

\begin{openproblem}
	18.50. Let $n \in  \mathbb{N}$ . Is it true that for every $k \in  \{ 1,\ldots , n!\}$ there is some group $G$ and pairwise distinct elements ${g}_{1},\ldots ,{g}_{n} \in  G$ such that the set $\left\{  {{g}_{\sigma \left( 1\right) }\cdots {g}_{\sigma \left( n\right) } \mid  \sigma  \in  {\mathrm{S}}_{n}}\right\}$ of all products of the ${g}_{i}$ obtained by permuting the factors has cardinality $k$ ? S. Kohl
\end{openproblem}

\plainproblemsection

\begin{openproblem}
	18.51. Given a prime $p$ and $n \in  \mathbb{N}$ , let ${f}_{p}\left( n\right)$ be the smallest number such that there is a group of order ${p}^{{f}_{p}\left( n\right) }$ into which every group of order ${p}^{n}$ embeds. Is it true that ${f}_{p}\left( n\right)$ grows faster than polynomially but slower than exponentially when $n$ tends to infinity? S. Kohl
\end{openproblem}

\plainproblemsection

\begin{openproblem}
	18.52. Is every finite simple group generated by two elements of prime-power orders $m, n$ ? (Here numbers $m, n$ may depend on the group.) The work of many authors shows that it remains to verify this property for a small number of finite simple groups. J. Krempa
	
	Yes, it is (mod CFSG); moreover, every finite simple group is generated by an involution and an element of prime order (C. S. H. King, J. Algebra, 478 (2017), 153-173).
\end{openproblem}

\plainproblemsection

\begin{openproblem}
	18.53. Is there a non-linear simple locally finite group in which every centralizer of a non-trivial element is almost soluble, that is, has a soluble subgroup of finite index?
	
	M. Kuzucuoglu
\end{openproblem}

\plainproblemsection

\begin{openproblem}
	18.54. Can a group be equal to the union of conjugates of a proper finite nonabelian simple subgroup? G. Cutolo
\end{openproblem}

\plainproblemsection

\begin{openproblem}
	18.55. a) Can a locally finite $p$ -group $G$ of finite exponent be the union of conjugates of an abelian proper subgroup?
	
	b) Can this happen when $G$ is of exponent $p$ ?G. Cutolo
\end{openproblem}

\plainproblemsection

\begin{openproblem}
	18.56. Let $G$ be a finite 2-group, of order greater than 2, such that $\left| {H/{H}_{G}}\right|  \leq  2$ for all $H \leq  G$ , where ${H}_{G}$ denotes the largest normal subgroup of $G$ contained in $H$ . Must $G$ have an abelian subgroup of index 4? G. Cutolo
\end{openproblem}

\plainproblemsection

\begin{openproblem}
	18.57. Let $G$ be a finite 2-group generated by involutions in which $\left\lbrack  {x, u, u}\right\rbrack   = 1$ for every $x \in  G$ and every involution $u \in  G$ . Is the derived length of $G$ bounded?
	
	D. V. Lytkina
	
	No, it is not (A. Abdollahi, J. Group Theory, 18, no. 1 (2015), 111-114).
\end{openproblem}

\plainproblemsection

\begin{openproblem}
	18.58. Let $G$ be a group generated by finite number $n$ of involutions in which ${\left( uv\right) }^{4} = 1$ for all involutions $u, v \in  G$ . Is it true that $G$ is finite? is a 2-group? This is true for $n \leq  3$ .
	
	Editors' comment of 2021: the previously published claim of an affirmative solution of this problem proved to be erroneous. D. V. Lytkina
\end{openproblem}

\plainproblemsection

\begin{openproblem}
	18.59. Does there exist a periodic group $G$ such that $G$ contains an involution, all involutions in $G$ are conjugate, and the centralizer of every involution $i$ is isomorphic to $\langle i\rangle  \times  {L}_{2}\left( P\right)$ , where $P$ is some infinite locally finite field of characteristic 2 ? D. V. Lytkina
\end{openproblem}

\plainproblemsection

\begin{openproblem}
	18.60. Let $V$ be an infinite countable elementary abelian additive 2-group. Does Aut $V$ contain a subgroup $G$ such that
	
	(a) $G$ is transitive on the set of non-zero elements of $V$ , and
	
	(b) if $H$ is the stabilizer in $G$ of a non-zero element $v \in  V$ , then $V = \langle v\rangle  \oplus  {V}_{v}$ , where ${V}_{v}$ is $H$ -invariant, $H$ is isomorphic to the multiplicative group ${P}^{ * }$ of a locally finite field $P$ of characteristic 2, and the action $H$ on ${V}_{v}$ is similar to the action of ${P}^{ * }$ on $P$ by multiplication?
	
	Conjecture: such a group $G$ does not exist. If so, then the group $G$ in the previous problem does not exist too. D. V. Lytkina
\end{openproblem}

\plainproblemsection

\begin{openproblem}
	18.61. Is a 2-group nilpotent if all its finite subgroups are nilpotent of class at most 3? This is true if 3 is replaced by 2 . D. V. Lytkina.
\end{openproblem}

\plainproblemsection

\begin{openproblem}
	18.62. The spectrum of a finite group is the set of orders of its elements. Let $\omega$ be a finite set of positive integers. A group $G$ is said to be $\omega$ -critical if the spectrum of $G$ coincides with $\omega$ , but the spectrum of every proper section of $G$ is not equal to $\omega$ .
	
	a) Does there exist a number $n$ such that for every finite simple group $G$ the number of $\omega \left( G\right)$ -critical groups is less than $n$ ?
	
	b) For every finite simple group $G$ , find all $\omega \left( G\right)$ -critical groups. V. D. Maz
\end{openproblem}

\plainproblemsection

\begin{openproblem}
	18.63. Let $G$ be a periodic group generated by two fixed-point-free automorphisms of order 5 of an abelian group. Is $G$ finite? V. D. Mazurov
\end{openproblem}

\plainproblemsection

\begin{openproblem}
	18.64. (K. Harada). Conjecture: Let $G$ be a finite group, $p$ a prime, and $B$ a $p$ -block of $G$ . If $J$ is a non-empty subset of $\operatorname{Irr}\left( B\right)$ such that $\sum \chi \left( 1\right) \chi \left( g\right)  = 0$ for every $p$ -singular element $g \in  G$ , then $J = \operatorname{Irr}\left( B\right)$ . V. D. Mazurov
\end{openproblem}

\plainproblemsection

\begin{openproblem}
	18.65. (R. Guralnick, G. Malle). Conjecture: Let $p$ be a prime different from 5, and $C$ a class of conjugate $p$ -elements in a finite group $G$ . If $\left\lbrack  {c, d}\right\rbrack$ is a $p$ -element for any $c, d \in  C$ , then $C \subseteq  {O}_{p}\left( G\right)$ . V. D. Mazurov
\end{openproblem}

\plainproblemsection

\begin{openproblem}
	18.66. Suppose that a finite group $G$ admits a Frobenius group of automorphisms ${FH}$ with kernel $F$ and complement $H$ such that ${GF}$ is also a Frobenius group with kernel $G$ and complement $F$ . Is the derived length of $G$ bounded in terms of $\left| H\right|$ and the derived length of ${C}_{G}\left( H\right)$ ? N. Yu. Makarenko, E. I. Khukhro, P. Shumyatsky
\end{openproblem}

\plainproblemsection

\begin{openproblem}
	18.67. Suppose that a finite group $G$ admits a Frobenius group of automorphisms ${FH}$ with kernel $F$ and complement $H$ such that ${C}_{G}\left( F\right)  = 1$ . Is the exponent of $G$ bounded in terms of $\left| F\right|$ and the exponent of ${C}_{G}\left( H\right)$ ?
	
	This was proved when $F$ is cyclic (E. I. Khukhro, N. Y. Makarenko, P. Shumyatsky, Forum Math., $\mathbf{{26}}\left( {2014}\right) ,{73} - {112}$ ) and for ${FH} \cong  {\mathbb{A}}_{4}$ when $\left( {\left| G\right| ,3}\right)  = 1$ (P. Shumyatsky, . Algebra, 331 (2011), 482-489). N. Yu. Makarenko, E. I. Khukhro, P. Shumyatsky
\end{openproblem}

\plainproblemsection

\begin{openproblem}
	18.68. What are the nonabelian composition factors of a finite nonsoluble group all of whose maximal subgroups have complements?
	
	Note that finite simple groups with all maximal subgroups having complements are, up to isomorphism, ${L}_{2}\left( 7\right) ,{L}_{2}\left( {11}\right)$ , and ${L}_{5}\left( 2\right)$ (V.M. Levchuk, A. G. Likharev, Siberian Math. J., 47, no. 4 (2006), 659-668). The same groups exhaust finite simple groups with Hall maximal subgroups and composition factors of groups with Hall maximal subgroups (N. V. Maslova, Siberian Math. J., 53, no. 5 (2012), 853-861). It is also proved that in a finite group with Hall maximal subgroups all maximal subgroups have complements (N. V. Maslova, D. O. Revin, Siberian Adv. Math., 23, no. 3 (2013), 196-209). N. V. Maslova, D. O. Revin
\end{openproblem}

\plainproblemsection

\begin{openproblem}
	18.69. Does there exist a relatively free group $G$ containing a free subsemigroup and having $\left\lbrack  {G, G}\right\rbrack$ finitely generated?
	
	Note that $G$ cannot be locally graded (Publ. Math. Debrecen,81, no. 3-4 (2012), ${415} - {420}$ .) O. Macedońska
\end{openproblem}

\plainproblemsection

\begin{openproblem}
	18.70. Is every finitely generated Coxeter group conjugacy separable?
	
	Some special cases were considered in (P.-E. Caprace, A. Minasyan, Illinois J. Math., 57, no. 2 (2013), 499-523). A. Minasyan
\end{openproblem}

\plainproblemsection

\begin{openproblem}
	*18.71. (N. Aronszajn). Suppose that $W\left( {x, y}\right)  = 1$ in an open subset of $G \times  G$ , where $G$ is a connected topological group. Must $W\left( {x, y}\right)  = 1$ for all $\left( {x, y}\right)  \in  G \times  G$ ?
	
	For locally compact groups the answer is yes. J. Mycielski
	
	*No, not necessarily: for any odd integer $n > {10}^{10}$ there is a connected topological group such that the identity ${x}^{n} = 1$ holds in some neighborhood of unity, but not in the entire group (E. Reznichenko, I. Zyabrev, Preprint, 2024, https://arxiv.org/ abs/2406.05203).
\end{openproblem}

\plainproblemsection

\begin{openproblem}
	18.72. Is it true that an existentially closed subgroup of a nonabelian free group of finite rank is a nonabelian free factor of this group? A. G. Myasnikov, V. A. Roman'kov
\end{openproblem}

\plainproblemsection

\begin{openproblem}
	18.73. a) Does every finitely generated solvable group of derived length $l \geq  2$ embed into a 2-generated solvable group of length $l + 1$ ?
	
	Comment of 2021: It is proved that any countable solvable group of derived length $l$ with torsion-free abelianization embeds in a 2-generated solvable group of derived length $l + 1$ (V. A. Roman’kov, Proc. Amer. Math. Soc., $\mathbf{{149}}\left( {2021}\right) ,{4133} - {4143}$ ).
	
	*b) Or at least, into some $k$ -generated $\left( {l + 1}\right)$ -solvable group, where $k = k\left( l\right)$ ?
	
	A. Yu. Ol'shanskiĭ
	
	*b) It is proved that any finitely generated solvable group of derived length $l$ can be embedded in a 4-generated solvable group of derived length $l + 1$ (V. A. Roman’kov, Proc. Amer. Math. Soc., 149 (2021), 4133-4143).
\end{openproblem}

\plainproblemsection

\begin{openproblem}
	18.74. Let $G$ be a finitely generated elementary amenable group which is not virtually nilpotent. Is there a finitely generated metabelian non-virtually-nilpotent section in
	
	$G$ ? A. Yu. Ol'shanskii
\end{openproblem}

\plainproblemsection

\begin{openproblem}
	18.75. Does every finite solvable group $G$ have the following property: there is a number $d = d\left( G\right)$ such that $G$ is a homomorphic image of every group with $d$ generators and one relation?
	
	This property holds for finite nilpotent groups and does not hold for every nonsolvable finite group; see (S. A. Zaǐtsev, Moscow Univ. Math. Bull., 52, no. 4 (1997), ${42} - {44})$ . A. Yu. Ol'shanskiĭ
	
	Yes, it does (N. Nikolov, D. Segal, Bull. London Math. Soc., 39, no. 2 (2007), 209- 213).
\end{openproblem}

\plainproblemsection

\begin{openproblem}
	18.76. Let $A$ be a division ring, $G$ a subgroup of the multiplicative group of $A$ , and $E$ an extension of the additive group of $A$ by $G$ such that $G$ acts by multiplication in $A$ . Is it true that $E$ splits? This is true if $G$ is finite. E. A. Palyutin.
\end{openproblem}

\plainproblemsection

\begin{openproblem}
	18.77. Let $G$ be a finite $p$ -group and let ${p}^{e}$ be the largest degree of an irreducible complex representation of $G$ . If $p > e$ , is it necessarily true that $\cap  \ker \Theta  = 1$ , where the intersection runs over all irreducible complex representations $\Theta$ of $G$ of degree ${p}^{e}$ ?
	
	D. S. Passman
\end{openproblem}

\plainproblemsection

\begin{openproblem}
	18.78. Let ${K}^{t}G$ be a twisted group algebra of the finite group $G$ over the field $K$ . If ${K}^{t}G$ is a simple $K$ -algebra, is $G$ necessarily solvable? This is known to be true if ${K}^{t}G$ is central simple. D. S. Passman
\end{openproblem}

\plainproblemsection

\begin{openproblem}
	18.79. Let $K\left\lbrack  G\right\rbrack$ be the group algebra of the finitely generated group $G$ over the field $K$ . Is the Jacobson radical $\mathcal{J}K\left\lbrack  G\right\rbrack$ equal to the join of all nilpotent ideals of the ring? This is known to be true if $G$ is solvable or linear. D. S. Passman
\end{openproblem}

\plainproblemsection

\begin{openproblem}
	18.80. (I. Kaplansky). For $G \neq  1$ , show that the augmentation ideal of the group algebra $K\left\lbrack  G\right\rbrack$ is equal to the Jacobson radical of the ring if and only if char $K = p > 0$ and $G$ is a locally finite $p$ -group. D. S. Passman
\end{openproblem}

\plainproblemsection

\begin{openproblem}
	18.81. Let $G$ be a finitely generated $p$ -group that is residually finite. Are all maximal subgroups of $G$ necessarily normal? D. S. Passman
\end{openproblem}

\plainproblemsection

\begin{openproblem}
	18.82. Is there a function $f : \mathbb{N} \rightarrow  \mathbb{N}$ such that for any prime $p$ , if ${p}^{f\left( n\right) }$ divides the order of a finite group $G$ , then ${p}^{n}$ divides the order of Aut $G$ ? $\;R.M.$ Patne Yes, there is. This was established in (J. A. Green, Proc. Roy. Soc. London Ser. A., 237 (1956), 574-581), which improved a previous result with a function depending on $p$ in (W. Ledermann, B. H. Neumann, Proc. Roy. Soc. London. Ser. A, 235 (1956), 235-246).
\end{openproblem}

\plainproblemsection

\begin{openproblem}
	18.83. A generating system $X$ of a group $G$ is fast if there is an integer $n$ such that every element of $G$ can be expressed as a product of at most $n$ elements of $X$ or their inverses. If not, we say that it is slow. For instance, in $\left( {\mathbb{Z}, + }\right)$ , the squares are fast, but the powers of 2 are slow.
	
	Do there exist countable infinite groups without an infinite slow generating set? Uncountable ones do exist. B. Poizat
\end{openproblem}

\plainproblemsection

\begin{openproblem}
	18.84. Let $\pi$ be a set of primes. We say that a finite group is a $B{S}_{\pi }$ -group if every conjugacy class in this group any two elements of which generate a $\pi$ -subgroup itself generates a $\pi$ -subgroup. Is every normal subgroup of a $B{S}_{\pi }$ -group a $B{S}_{\pi }$ -group?
	
	In the case $2 \notin  \pi$ , an affirmative answer follows from (D. O. Revin, Siberian Math. J., 52, no. 2 (2011), 340-347). D. O. Revin
\end{openproblem}

\plainproblemsection

\begin{openproblem}
	18.85. A subset of a group is said to be rational if it can be obtained from finite subsets by finitely many rational operations, that is, taking union, product, and the submonoid generated by a set.
	
	Conjecture: every finitely generated solvable group in which all rational subsets form a Boolean algebra is virtually abelian.
	
	This is known to be true if the group is metabelian, or polycyclic, or of finite rank (G. A. Bazhenova). V. A. Roman'kov
\end{openproblem}

\plainproblemsection

\begin{openproblem}
	18.86. Is the group $G = \langle a, b \mid  \left\lbrack  {\left\lbrack  {a, b}\right\rbrack  , b}\right\rbrack   = 1\rangle$ , which is isomorphic to the group of all unitriangular automorphisms of the free group of rank 3, linear? V. A. Roman'kov Yes, it is linear, since it embeds in the holomorph Hol ${F}_{2}$ of the free group ${F}_{2}$ , which can be seen by adding a new generator $c = \left\lbrack  {a, b}\right\rbrack$ , so that $G = \left\langle  {a, b, c \mid  {a}^{b} = }\right.$ ${ac},{c}^{b} = c\rangle$ , while $\operatorname{Hol}{F}_{2}$ was shown to be linear in Corollary 3 in (V. G. Bardakov, O. V. Bryukhanov, Vestnik Novosibirsk Univ. Ser. Mat. Mekh. Inf., 7, no. 3 (2007), 45-58 (Russian)) (O. V. Bryukhanov, Letter of 27 January 2014). Another proof can be found in (V. A. Roman'kov, J. Siberian Federal Univ. Math. Phys., 6, no. 4 (2013), 516-520).
\end{openproblem}

\plainproblemsection

\begin{openproblem}
	18.87. A system of equations with coefficients in a group $G$ is said to be independent if the matrix composed of the sums of exponents of the unknowns has rank equal to the number of equations.
	
	a) The Kervaire-Laudenbach Conjecture (KLC): every independent system of equations with coefficients in an arbitrary group $G$ has a solution in some over-group $\bar{G}$ . This is true for every locally residually finite group $G$ (M. Gerstenhaber and O. S. Rothaus).
	
	b) KLC - nilpotent version: every independent system of equations with coefficients in an arbitrary nilpotent group $G$ has a solution in some nilpotent overgroup $\bar{G}$ .
	
	c) KLC - solvable version: every independent system of equations with coefficients in an arbitrary solvable group $G$ has a solution in some solvable overgroup $\bar{G}$ .
	
	V. A. Roman'kov
\end{openproblem}

\plainproblemsection

\begin{openproblem}
	18.88. Can a finitely generated infinite group of finite exponent be the quotient of a residually finite group by a locally finite normal subgroup?
	
	If not, then there exists a hyperbolic group that is not residually finite. M. Sapir
\end{openproblem}

\plainproblemsection

\begin{openproblem}
	18.89. Consider the set of balanced presentations $\left\langle  {{x}_{1},\ldots ,{x}_{n} \mid  {r}_{1},\ldots ,{r}_{n}}\right\rangle$ with fixed $n$ generators ${x}_{1},\ldots ,{x}_{n}$ . By definition, the group $A{C}_{n}$ of Andrews-Curtis moves on this set of balanced presentations is generated by the Nielsen transformations together with conjugations of relators. Is $A{C}_{n}$ finitely presented?
	
	Comment of 2023: It has been proved that $A{C}_{2}$ is not finitely presented (V. A. Roman'kov, Preprint, 2023, https://arxiv.org/pdf/2305.11838.pdf).J.Swan, A. Lisitsa
\end{openproblem}

\plainproblemsection

\begin{openproblem}
	18.90. Prove that the group $Y\left( {m, n}\right)  = \left\langle  {{a}_{1},{a}_{2},\ldots ,{a}_{m} \mid  {a}_{k}^{n} = 1\left( {1 \leq  k \leq  }\right. }\right.$ $m),{\left( {a}_{k}^{i}{a}_{l}^{i}\right) }^{2} = e\left( {1 \leq  k < l \leq  m,1 \leq  i \leq  \frac{n}{2}}\right) \rangle$ is finite for every pair(m, n).
	
	Some known cases are (where ${C}_{n}$ denotes a cyclic group of order $n$ ): $Y\left( {m,2}\right)  =$ ${C}_{2}^{m};Y\left( {m;3}\right)  = {A}_{m + 2}$ (presentation by Carmichael); $Y\left( {m;4}\right)$ has order ${2}^{\frac{m\left( {m + 3}\right) }{2}}$ , nilpo-tency class 3 and exponent $4;Y\left( {2, n}\right)$ is the natural extension of the augmentation ideal $\omega \left( {C}_{n}\right)$ of ${GF}\left( 2\right) \left\lbrack  {C}_{n}\right\rbrack$ by ${C}_{n}$ (presentation by Coxeter); $Y\left( {3, n}\right)  \cong  {SL}\left( {2,\omega \left( {C}_{n}\right) }\right)$ .
	
	These groups have connections with classical groups in characteristic 2. When $n$ is odd, the group $Y\left( {m, n}\right)$ has the following presentation, where ${\tau }_{ij}$ are transpositions: $y\left( {m, n}\right)  = \left\langle  {a,{\mathbb{S}}_{m} \mid  {a}^{n} = e,\left\lbrack  {{\tau }_{12}^{{a}^{i}},{\tau }_{12}}\right\rbrack   = e\left( {1 \leq  i \leq  \frac{n}{2}}\right) ,{\tau }_{12}^{1 + a + \cdots  + {a}^{n - 1}} = }\right.$ $e,{\tau }_{i, i + 1}a{\tau }_{i, i + 1} = {a}^{-1}\left( {2 \leq  i \leq  m - 1}\right) \rangle$ , and it is known that, for example, $y\left( {3,5}\right)  \cong$ ${SL}\left( {2,{16}}\right)  \cong  {\Omega }^{ - }\left( {4,4}\right) , y\left( {4,5}\right)  \cong  {Sp}\left( {2,{16}}\right)  \cong  \Omega \left( {5,4}\right) , y\left( {5,5}\right)  \cong  {SU}\left( {4,{16}}\right)  \cong  {\Omega }^{ - }\left( {6,4}\right) ,$ $y\left( {6,5}\right)  \cong  {4}^{6}{\Omega }^{ - }\left( {6,4}\right) , y\left( {3,7}\right)  \cong  {SL}{\left( 2,8\right) }^{2} \cong  {\Omega }^{ + }\left( {4,8}\right) , y\left( {4,7}\right)  \cong  {Sp}\left( {4,8}\right)  \cong  \Omega \left( {5,8}\right)$ . Extensive computations by Felsch, Neubüser and O'Brien confirm this trend. Different types reveal Bott periodicity and a connection with Clifford algebras.S. Sidki
\end{openproblem}

\plainproblemsection

\begin{openproblem}
	18.91. A subgroup $H$ of a group $G$ is said to be propermutable in $G$ if there is a subgroup $B \leq  G$ such that $G = {N}_{G}\left( H\right) B$ and $H$ permutes with every subgroup of $B$ .
	
	a) Is there a finite group $G$ with subgroups $A \leq  B \leq  G$ such that $A$ is proper-mutable in $G$ but $A$ is not propermutable in $B$ ?
	
	b) Let $P$ be a non-abelian Sylow 2-subgroup of a finite group $G$ with $\left| P\right|  = {2}^{n}$ . Suppose that there is an integer $k$ such that $1 < k < n$ and every subgroup of $P$ of order ${2}^{k}$ is propermutable in $G$ , and also, in the case of $k = 1$ , every cyclic subgroup of $P$ of order 4 is propermutable in $G$ . Is it true that then $G$ is 2-nilpotent?
	
	A. N. Skiba
	
	a) Yes, there is (A. A. Pypka, D. Yu. Storozhenko, Dopov. Nac. Akad. Nauk Ukrain., 2017, no. 7, 18-20 (Ukrainian)).
	
	b) Yes, it is (Kh. A. Al-Sharo, Finite groups with given systems of weakly $S$ -proper-mutable subgroups, J. Group Theory, 19, no. 5 (2016), 871-887).
\end{openproblem}

\plainproblemsection

\begin{openproblem}
	18.92. A non-empty set $\theta$ of formations is called a complete lattice of formations if the intersection of any set of formations in $\theta$ belongs to $\theta$ and $\theta$ has the largest element (with respect to inclusion). If $L$ is a complete lattice, then an element $a \in  L$ is said to be compact if $a \leq  \bigvee X$ for any $X \subseteq  L$ implies that $a \leq  \bigvee {X}_{1}$ for some finite ${X}_{1} \subset  X$ . A complete lattice is called algebraic if every element is the join of a (possibly infinite) set of compact elements.
	
	a) Is there a non-algebraic complete lattice of formations of finite groups?
	
	b) Is there a non-modular complete lattice of formations of finite groups?
\end{openproblem}

\plainproblemsection

\begin{openproblem}
	18.93. Let $\mathcal{M}$ be a one-generated saturated formation, that is, the intersection of all saturated formations containing some fixed finite group. Let $\mathcal{F}$ be a subformation of $\mathcal{M}$ such that $\mathcal{F} \neq  \mathcal{F}\mathcal{F}$ .
	
	a) Is it true that then $\mathcal{F}$ can be written in the form $\mathcal{F} = {\mathcal{F}}_{1}\cdots {\mathcal{F}}_{t}$ , where ${\mathcal{F}}_{i}$ is a non-decomposable formation for every $i = 1,\ldots , t$ ?
	
	b) Suppose that $\mathcal{F} = {\mathcal{F}}_{1}\cdots {\mathcal{F}}_{t}$ , where ${\mathcal{F}}_{i}$ is a non-decomposable formation for every $i = 1,\ldots , t$ . Is it true that then all factors ${\mathcal{F}}_{i}$ are uniquely determined?
	
	A. N. Skiba
\end{openproblem}

\plainproblemsection

\begin{openproblem}
	18.94. Let $G$ be a group without involutions, $a$ an element of it that is not a square of any element of $G$ , and $n$ an odd positive integer. Is it true that the quotient $G/\left\langle  {\left( {a}^{n}\right) }^{G}\right\rangle$ does not contain involutions?
	
	No, not always, as shown by an example of V. I. Trofimov; another example: the quotient of $G = \left\langle  {a, b \mid  {a}^{2} = {b}^{2}}\right\rangle$ by $\left\langle  {a}^{G}\right\rangle$ has order 2 .
\end{openproblem}

\plainproblemsection

\begin{openproblem}
	18.95. Suppose that a group $G = {AB}$ is a product of an abelian subgroup $A$ and a locally quaternion group $B$ (that is, $B$ is a union of an increasing chain of finite generalized quaternion groups). Is $G$ soluble?
	
	Yes, it is (B. Amberg, Ya. Sysak, On products of groups with abelian subgroups of small index, J. Group Theory, 20, no. 6 (2017), 1061-1072).
\end{openproblem}

\plainproblemsection

\begin{openproblem}
	18.96. Suppose that a periodic group $G$ contains involutions and the centralizer of each involution is locally finite. Is it true that $G$ has a nontrivial locally finite normal subgroup? N. M. Suchkov
\end{openproblem}

\plainproblemsection

\begin{openproblem}
	18.97. Let $G$ be a periodic Zassenhaus group, that is, a two-transitive permutation group with trivial stabilizer of every three points. Suppose that the stabilizer of a point is a Frobenius group with locally finite kernel $U$ containing an involution. Is it true that $U$ is a 2-group? This was proved for finite groups by Feit. N. M. Suchkov
\end{openproblem}

\plainproblemsection

\begin{openproblem}
	18.98. The work of many authors shows that most finite simple groups are generated by two elements of orders 2 and 3; for example, see the survey http://math.nsc.ru/ conference/groups2013/slides/MaximVsemirnov\_slides.pdf. Which finite simple groups cannot be generated by two elements of orders 2 and 3 ?
	
	In particular, is it true that, among classical simple groups of Lie type, such exceptions, apart from $\operatorname{PSU}\left( {3,{5}^{2}}\right)$ , arise only when the characteristic is 2 or 3 ?
	
	The question of which finite simple groups are(2,3)-generated remains open only for the orthogonal groups (M. A. Pellegrini, M. C. Tamburini Bellani, J. Algebra, 598 (2022), 156-193). M. C. Tamburini
\end{openproblem}

\plainproblemsection

\begin{openproblem}
	18.99. Let $\operatorname{cd}\left( G\right)$ be the set of all complex irreducible character degrees of a finite group $G$ . Huppert’s Conjecture: if $H$ is a nonabelian simple group and $G$ is any finite group such that $\operatorname{cd}\left( G\right)  = \operatorname{cd}\left( \mathrm{H}\right)$ , then $G \cong  H \times  A$ , where $A$ is an abelian group.
	
	H. P. Tong-Viet
\end{openproblem}

\plainproblemsection

\begin{openproblem}
	18.100. For a given class of groups $\mathfrak{X}$ , let ${\left( \mathfrak{X},\infty \right) }^{ * }$ denote the class of groups in which every infinite subset contains two distinct elements $x$ and $y$ that satisfy $\left\langle  {x,{x}^{y}}\right\rangle   \in  \mathfrak{X}$ . Let $G$ be a finitely generated soluble-by-finite group in the class ${\left( \mathfrak{X},\infty \right) }^{ * }$ , let $m$ be a positive integer, and let $\mathfrak{F},{\mathfrak{E}}_{m},\mathfrak{E},\mathfrak{N}$ , and $\mathfrak{P}$ denote the classes of finite groups, groups of exponent dividing $m$ , groups of finite exponent, nilpotent groups, and polycyclic groups, respectively.
	
	a) If $\mathfrak{X} = \mathfrak{E}\mathfrak{N}$ , then is $G$ in $\mathfrak{E}\mathfrak{N}$ ?
	
	b) If $\mathfrak{X} = {\mathfrak{E}}_{m}\mathfrak{N}$ , then is $G$ in ${\mathfrak{E}}_{m}\left( {\mathfrak{F}\mathfrak{N}}\right)$ ?
	
	c) If $\mathfrak{X} = \mathfrak{N}\left( {\mathfrak{P}\mathfrak{F}}\right)$ , then is $G$ in $\mathfrak{N}\left( {\mathfrak{P}\mathfrak{F}}\right)$ ?N. Trabelsi
\end{openproblem}

\plainproblemsection

\begin{openproblem}
	18.102. (E. C. Dade). Let $C$ be a Carter subgroup of a finite solvable group $G$ , and let $\ell \left( C\right)$ be the number of primes dividing $\left| C\right|$ counting multiplicities. It was proved in (E. C. Dade, Illinois J. Math., 13 (1969), 449-514) that there is an exponential function $f$ such that the nilpotent length of $G$ is at most $f\left( {\ell \left( C\right) }\right)$ . Is there a linear (or at least a polynomial) function $f$ with this property? A. Turull
\end{openproblem}

\plainproblemsection

\begin{openproblem}
	18.103. A group is said to be minimax if it has a finite subnormal series each of whose factors satisfies either the minimum or the maximum condition on subgroups. Is it true that in the class of nilpotent minimax groups only finitely generated groups may have faithful irreducible primitive representations over a finitely generated field of characteristic zero? A. V. Tushev
\end{openproblem}

\plainproblemsection

\begin{openproblem}
	18.104. Construct an example of a pro- $p$ group $G$ and a proper abstract normal subgroup $K$ such that $G/K$ is perfect. J. Wilson
\end{openproblem}

\plainproblemsection

\begin{openproblem}
	18.105. Let $R$ be the formal power series algebra in commuting indeterminates ${x}_{1},\ldots {x}_{n},\ldots$ over the field with $p$ elements, with $p$ a prime. (Thus its elements are (in general infinite) linear combinations of monomials ${x}_{1}^{{a}_{1}}\ldots {x}_{n}^{{a}_{n}}$ with $n,{a}_{1},\ldots ,{a}_{n}$ non-negative integers.) Let $I$ be the maximal ideal of $R$ and $J$ the (abstract) ideal generated by all products of two elements of $I$ . Is there an ideal $U$ of $R$ such that $U < I$ and $U + J = I$ ?
	
	If so, then ${\mathrm{{SL}}}_{3}\left( R\right)$ and the kernel of the map to ${\mathrm{{SL}}}_{3}\left( {R/U}\right)$ provide an answer to the previous question. J. Wilson
\end{openproblem}

\plainproblemsection

\begin{openproblem}
	18.106. Let $R$ be a group that acts coprimely on the finite group $G$ . Let $p$ be a prime, let $P$ be the unique maximal $R{C}_{G}\left( R\right)$ -invariant $p$ -subgroup of $G$ and assume that ${C}_{G}\left( {{O}_{p}\left( G\right) }\right)  \leq  {O}_{p}\left( G\right)$ . If $p > 3$ then $P$ contains a nontrivial characteristic subgroup that is normal in $G$ (P. Flavell, J. Algebra, $\mathbf{{257}}\left( {2002}\right) ,{249} - {264}$ ). Does the same result hold for the primes 2 and 3 ? P. Flavell
\end{openproblem}

\plainproblemsection

\begin{openproblem}
	18.107. Is there a finitely generated infinite residually finite $p$ -group such that every subgroup of infinite index is cyclic?
	
	The answer is known to be negative for $p = 2$ . Note that in (M. Ershov, A. Jaikin-Zapirain, J. Reine Angew. Math., 677 (2013), 71-134) it is shown that for every prime $p$ there is a finitely generated infinite residually finite $p$ -group such that every finitely generated subgroup of infinite index is finite. A. Jaikin-Zapirain
\end{openproblem}

\plainproblemsection

\begin{openproblem}
	18.108. A group $G$ is said to have property $\left( \tau \right)$ if its trivial representation is an isolated point in the subspace of irreducible representations with finite images (in the topological unitary dual space of $G$ ). Is it true that there exists a finitely generated group satisfying property $\left( \tau \right)$ that homomorphically maps onto every finite simple group?
	
	The motivation is the theorem in (E. Breulliard, B. Green, M. Kassabov, A. Lubotzky, N. Nikolov, T. Tao) saying that there are $\epsilon  > 0$ and $k$ such that every finite simple group $G$ contains a generating set $S$ with at most $k$ elements such that the Cayley graph of $G$ with respect to $S$ is an $\epsilon$ -expander. In (M. Ershov, A. Jaikin-Zapirain, M. Kassabov, Mem. Amer. Math. Soc., 1186, 2017) it is proved that there exists a group with Kazhdan's property (T) that maps onto every simple group of Lie type of rank $\geq  2$ . (A group $G$ is said to have Kazhdan’s property (T) if its trivial representation is an isolated point in the unitary dual space of $G$ .) A. Jaikin-Zapirain
\end{openproblem}

\plainproblemsection

\begin{openproblem}
	18.109. Is it true that a group satisfying Kazhdan's property (T) cannot homomor-phically map onto infinitely many simple groups of Lie type of rank 1 ?
	
	It is known that a group which maps onto ${PS}{L}_{2}\left( q\right)$ for infinitely many $q$ does not have Kazhdan's property (T). A. Jaikin-Zapirain
\end{openproblem}

\plainproblemsection

\begin{openproblem}
	18.110. The non-p-soluble length of a finite group $G$ is the number of non-p-soluble factors in a shortest normal series each of whose factors either is $p$ -soluble or is a direct product of non-abelian simple groups of order divisible by $p$ . For a given prime $p$ and a given proper group variety $\mathfrak{V}$ , is there a bound for the non- $p$ -soluble length of finite groups whose Sylow $p$ -subgroups belong to $\mathfrak{V}$ ?
	
	Comment of 2021: the existence of such a bound is proved for $p = 2$ (F. Fumagalli, F. Leinen, O. Puglisi, Proc. London Math. Soc. (3), 125, no. 5 (2022), 1066-1082).
\end{openproblem}

\plainproblemsection

\begin{openproblem}
	18.111. Let $G$ be a discrete countable group, given as a central extension $0 \rightarrow  \mathbb{Z} \rightarrow$ $G \rightarrow  Q \rightarrow  0$ . Assume that either $G$ is quasi-isometric to $\mathbb{Z} \times  Q$ , or that $\mathbb{Z} \rightarrow  G$ is a quasi-isometric embedding. Does that imply that $G$ comes from a bounded cocycle on $Q$ ? I. Chatterji, G. Mislin
\end{openproblem}

\plainproblemsection

\begin{openproblem}
	18.112. Is it true that the orders of all elements of a finite group $G$ are powers of primes if, for every divisor $d\left( {d > 1}\right)$ of $\left| G\right|$ and for every subgroup $H$ of $G$ of order coprime to $d$ , the order $\left| H\right|$ divides the number of elements of $G$ of order $d$ ?
	
	The converse is true (W. J. Shi, Math. Forum (Vladikavkaz), 6 (2012), 152-154). W. J. Shi
	
	Yes, it is true (A. A. Buturlakin, R. Shen, W. Shi, Siberian Math. J., 58, no. 3 (2017), 405-407).
\end{openproblem}

\plainproblemsection

\begin{openproblem}
	18.113. Let $\mathfrak{M}$ be a finite set of finite simple nonabelian groups. Is it true that a periodic group saturated with groups from $\mathfrak{M}$ (see 14.101) is isomorphic to one of groups in $\mathfrak{M}$ ? The case where $\mathfrak{M}$ is one-element is of special interest. A. K. Shlëpkin
\end{openproblem}

\plainproblemsection

\begin{openproblem}
	18.114. Does there exist an irreducible ${5}^{\prime }$ -subgroup $G$ of ${GL}\left( V\right)$ for some finite ${\mathbb{F}}_{{5}^{ - }}$ space $V$ such that the number of conjugacy classes of the semidirect product ${VG}$ is equal to $\left| V\right|$ but $G$ is not cyclic?
	
	Such a subgroup exists if 5 is replaced by $p = 2,3$ but does not exist for primes $p > 5$ (J. Group Theory,14 (2011), 175-199). P. Schmid
\end{openproblem}

\plainproblemsection

\begin{openproblem}
	*18.115. Let $G$ be a finite simple group, and let $X, Y$ be isomorphic simple maximal subgroups of $G$ . Are $X$ and $Y$ conjugate in Aut $G$ ?P. Schmid
	
	*No, not always. For example, there are two Aut $G$ -classes of ${M}_{12}$ in ${E}_{6}\left( 5\right)$ (P. B. Kleidman, R. A. Wilson, J. London Math. Soc., 42 (1990), 555-561); other counterexamples in ${E}_{6}\left( q\right)$ for certain $q$ include two classes of ${PS}{L}_{2}\left( {11}\right)$ , and two of ${PS}{L}_{2}\left( {19}\right)$ (D. A. Craven, Invent. Math.,234, no. 2 (2023), 637-719). (D. A. Craven Letter of 8 September 2021).
\end{openproblem}

\plainproblemsection

\begin{openproblem}
	18.116. Given a finite 2-group $G$ of order ${2}^{n}$ , does there exist a finite set $S$ of rational primes such that $\left| S\right|  \leq  n$ and $G$ is a quotient group of the absolute Galois group of the maximal 2-extension of $\mathbb{Q}$ unramified outside $S \cup  \{ \infty \}$ ?
	
	This is true for $p$ -groups, $p$ odd, as noted by Serre. P. Schmid
\end{openproblem}

\plainproblemsection

\begin{openproblem}
	18.117. Is every element of a nonabelian finite simple group a commutator of two elements of coprime orders?
	
	The answer is known to be "yes" for the alternating groups (P. Shumyatsky, Forum Math., 27, no. 1 (2015), 575-583) and for ${PSL}\left( {2, q}\right)$ (M. A. Pellegrini, P. Shumyatsky, Arch. Math., 99 (2012), 501-507). Without the coprimeness condition, this is Ore's conjecture proved in (M. W. Liebeck, E. A. O'Brien, A. Shalev, Pham Huu Tiep, J. Eur. Math. Soc., 12, no. 4 (2010), 939-1008). P. Shumyatsky
\end{openproblem}

\plainproblemsection

\begin{openproblem}
	18.118. Let $G$ be a profinite group such that $G/Z\left( G\right)$ is periodic. Is $\left\lbrack  {G, G}\right\rbrack$ necessarily periodic? P. Shumyatsky
\end{openproblem}

\plainproblemsection

\begin{openproblem}
	18.119. A multilinear commutator word is any commutator of weight $n$ in $n$ distinct variables. Let $w$ be a multilinear commutator word and let $G$ be a finite group. Is it true that every Sylow $p$ -subgroup of the verbal subgroup $w\left( G\right)$ is generated by $w$ -values? P. Shumyatsky
\end{openproblem}

\plainproblemsection

\begin{openproblem}
	18.120. Let $P = {AB}$ be a finite $p$ -group factorized by an abelian subgroup $A$ and a class-two subgroup $B$ . Suppose, if necessary, $A \cap  B = 1$ . Then is it true that $\langle A,\left\lbrack  {B, B}\right\rbrack  \rangle  = A{B}_{0}$ , where ${B}_{0}$ is an abelian subgroup of $B$ ? E. Jabara
\end{openproblem}

\plainproblemsection

\begin{openproblem}
	19.1. An element $g$ of a group $G$ is a non-near generator of $G$ if for every subset $S \subseteq  G$ such that $\left| {G : \langle g, S\rangle }\right|  < \infty$ it follows that $\left| {G : \langle S\rangle }\right|  < \infty$ . The set of all non-near generators of $G$ forms a characteristic subgroup of $G$ called the lower near Frattini subgroup of $G$ , denoted by $\lambda \left( G\right)$ . A subgroup $M \leq  G$ is nearly maximal in $G$ if it is maximal with respect to being of infinite index in $G$ . The intersection of all nearly maximal subgroups forms a characteristic subgroup called the upper near Frattini subgroup of $G$ , denoted by $\mu \left( G\right)$ . In general, $\lambda \left( G\right)  \leq  \mu \left( G\right)$ . If $\lambda \left( G\right)  = \mu \left( G\right)$ , then this subgroup is called the near Frattini subgroup of $G$ , denoted by $\psi \left( G\right)$ .
	
	a) Is it true that $\psi \left( G\right)  = 1$ if $G$ is the knot group of any product of knots?
	
	b) Is it true that $\psi \left( G\right)  = 1$ if $G$ is a cable knot group? M. K. Azarian
\end{openproblem}

\plainproblemsection

\begin{openproblem}
	19.2. For a subgroup $L$ of a group $G$ , let ${L}_{G}$ denote the largest normal subgroup of $G$ contained in $L$ . If $M \vartriangleleft  G$ , then we say that $G$ nearly splits over $M$ if there is a subgroup $N \leq  G$ such that $\left| {G : N}\right|  = \infty ,\left| {G : {MN}}\right|  < \infty$ , and ${\left( M \cap  N\right) }_{G} = 1$ .
	
	Let $G$ be any group, and $H$ a normal subgroup of prime order. Is it true that $\psi \left( G\right)  \cap  H = 1$ if and only if $G$ nearly splits over $H$ ? M. K. Azarian
\end{openproblem}

\plainproblemsection

\begin{openproblem}
	19.3. Let $G = A{ * }_{H}B$ be the generalized free product of groups $A$ and $B$ with amalgamated subgroup $H$ . Is $\psi \left( G\right)  = 1$ in the following cases?
	
	a) $H$ is finite cyclic and ${H}_{G} = 1$ .
	
	b) $H$ is finite cyclic and either $\lambda \left( A\right)  \cap  {H}_{G} = 1$ or $\lambda \left( B\right)  \cap  {H}_{G} = 1$ .
	
	c) ${H}_{G} = 1$ and $H$ satisfies the minimum condition on subgroups.
	
	d) $\lambda \left( G\right)  \cap  H = 1$ .
	
	e) $A$ and $B$ are free groups, $H$ is finitely generated and at least one of $\left| {A : H}\right|$ or $\left| {B : H}\right|$ is infinite.
	
	f) $H$ is infinite cyclic and is a retract of $A$ and $B$ .
	
	g) $G$ nearly splits over $H$ , and $H$ is a normal subgroup of $G$ of prime order.
	
	M. K. Azarian
\end{openproblem}

\plainproblemsection

\begin{openproblem}
	19.4. a) If $G$ is the free product of infinitely many finitely generated free groups with cyclic amalgamation, then is $\psi \left( G\right)  = 1$ ?
	
	b) If $G$ is the free product of infinitely many finitely generated free abelian groups with cyclic amalgamation, then is $\psi \left( G\right)  = 1$ ? M. K. Azarian
\end{openproblem}

\plainproblemsection

\begin{openproblem}
	19.5. If $G$ is the free product of infinitely many finitely generated abelian groups with amalgamated subgroup $H$ , then is $\psi \left( G\right)$ equal to the torsion subgroup of $H$ ?
	
	M. K. Azarian
\end{openproblem}

\plainproblemsection

\begin{openproblem}
	19.6. Let $G = A{ * }_{H}B$ be the generalized free product of groups $A$ and $B$ with amalgamated subgroup $H$ .
	
	a) Conjecture: If both $A$ and $B$ are nilpotent, then $\mu \left( G\right)  \leq  H$ .
	
	b) Conjecture: If $G$ is residually finite, and $H$ satisfies a nontrivial identical relation, then $\lambda \left( G\right)  \leq  H$ . M. K. Azarian
\end{openproblem}

\plainproblemsection

\begin{openproblem}
	19.7. The group of virtual coloured braids $V{P}_{n}, n \geq  2$ , is generated by elements ${\lambda }_{ij}$ , $1 \leq  i \neq  j \leq  n$ , and is defined by the relations ${\lambda }_{ij}{\lambda }_{kl} = {\lambda }_{kl}{\lambda }_{ij},{\lambda }_{ki}{\lambda }_{kj}{\lambda }_{ij} = {\lambda }_{ij}{\lambda }_{kj}{\lambda }_{ki}$ , where different letters denote different indices.
	
	a) Construct a normal form of words in the group $V{P}_{n}$ for $n \geq  4$ .
	
	b) Is the group $V{P}_{n}$ linear for $n \geq  4$ ? It is known that the group $V{P}_{3}$ is linear.
	
	V. G. Bardakov
\end{openproblem}

\plainproblemsection

\begin{openproblem}
	*19.8. A word in an alphabet $A = \left\{  {{a}_{1},{a}_{2},\ldots ,{a}_{n}}\right\}$ is called a palindrome if it reads the same from left to right and from right to left. Let $k$ a non-negative integer. A word in the alphabet $A$ is called an almost $k$ -palindrome if it can be transformed into a palindrome by changing $\leq  k$ letters in it. (So an almost 0 -palindrome is a palindrome.) Let elements of a free group ${F}_{2} = \langle x, y\rangle$ be represented as words in the alphabet $\left\{  {{x}^{\pm 1},{y}^{\pm 1}}\right\}$ . Do there exist positive integers $m$ and $c$ such that every element in ${F}_{2}$ is a product of $\leq  c$ almost $m$ -palindromes?
	
	It is known that for $m = 0$ there is no such a number $c$ . V. G. Bardakov
\end{openproblem}

\plainproblemsection

\begin{openproblem}
	19.9. Let $G$ be a finitely generated linear group.
	
	a) Is it true that the non-abelian tensor square $G \otimes  G$ is a linear group?
	
	b) In particular, is this true for the braid group ${B}_{n}$ for $n > 3$ ?
	
	It is known that ${B}_{3} \otimes  {B}_{3}$ is linear, and there is a countable group $G$ such that $G \otimes  G$ is not linear. See the definition of the non-abelian tensor square in (R. Brown, J.-L. Loday, Topology, 26, no. 3 (1987), 311-335). V. G. Bardakov, M. V. Neshchadim
\end{openproblem}

\plainproblemsection

\begin{openproblem}
	19.10. a) Is it possible to embed a finitely generated non-abelian free pro- $p$ group as an open subgroup of a simple totally disconnected locally compact group?
	
	b) If the answer is yes, can we require that the simple envelope is also compactly generated? Y. Barnea
\end{openproblem}

\plainproblemsection

\begin{openproblem}
	19.11. Does there exist a constant $c$ such that the number of conjugacy classes in a finite group $G$ is always at least $c{\log }_{2}\left| G\right|$ ?
	
	Editors’ comment of 2021: It is proved that every group $G$ contains at least $\varepsilon \log \left| G\right| /{\left( \log \log \left| G\right| \right) }^{8}$ conjugacy classes for some fixed $\varepsilon  > 0$ (L. Pyber, J. London Math. Soc. (2),46, no. 2 (1992), 239-249). It is also proved that for every $\varepsilon  > 0$ there exists $\delta  > 0$ such that every finite group $G$ of order at least 3 has at least $\delta {\log }_{2}\left| G\right| /{\left( {\log }_{2}{\log }_{2}\left| G\right| \right) }^{3 + \varepsilon }$ conjugacy classes (B. Baumeister, A. Maróti, H. P. Tong-Viet, Forum Math., 29, no. 2 (2017), 259-275). E. Bertram 19.12. A finite group $G$ is called conjugacy-expansive if for every normal subset $N$ and conjugacy class $C$ of $G$ the normal set ${NC}$ contains at least as many conjugacy classes of $G$ as $N$ does. Is it true that every finite simple group is conjugacy-expansive?
	
	M. Bezerra, Z. Halasi, A. Maróti, S. Sidki
\end{openproblem}

\plainproblemsection

\begin{openproblem}
	19.13. The well-known Baer-Suzuki theorem states that if every two conjugates of an element $a$ of a finite group $G$ generate a finite $p$ -subgroup, then $a$ is contained in a normal $p$ -subgroup. Does such a theorem hold in the class of periodic groups for $p > 2$ ?
	
	A counterexample is known for $p = 2$ ; see Archive 11.11a).
	\par\smallskip\hrule\smallskip
	*No, there are no such integers (M. Staiger, J. Algebra, 659 (2024), 475-481).
	\par\smallskip\hrule\smallskip
\end{openproblem}

\plainproblemsection

\begin{openproblem}
	19.14. Let $G$ be a finite group such that the quasivariety generated by all its Sylow subgroups contains only finitely many subquasivarieties. Is it true that the quasiva-riety generated by $G$ also contains only finitely many subquasivarieties?
	
	A. I. Budkin
\end{openproblem}

\plainproblemsection

\begin{openproblem}
	19.15. For a class of groups $M$ , let $L\left( M\right)$ denote the class of all groups $G$ in which the normal closure $\left\langle  {a}^{G}\right\rangle$ of any element $a \in  G$ is contained in $M$ . Let ${qA}$ be the quasivariety generated by a finite nilpotent group $A$ . Can the quasivariety $L\left( {qA}\right)$ contain a non-nilpotent group? A. I. Budkin
\end{openproblem}

\plainproblemsection

\begin{openproblem}
	19.16. (C. Drutu and M. Sapir). Is every free-by-cyclic group of the form ${F}_{n} \rtimes  \mathbb{Z}$ linear? J. O. Button
\end{openproblem}

\plainproblemsection

\begin{openproblem}
	19.17. (Well-known problem). Suppose that $G$ is a finitely presented group such that the set of first Betti numbers over all the finite index subgroups of $G$ is unbounded above. (Here, the first Betti number is the maximum $n$ for which there is a surjective homomorphism onto ${\mathbb{Z}}^{n}$ .)
	
	a) Must $G$ have a finite index subgroup with a surjective homomorphism to a non-abelian free group?
	
	b) Can $G$ even be soluble? J. O. Button
\end{openproblem}

\plainproblemsection

\begin{openproblem}
	19.18. Given an infinite set $\Omega$ , define an algebra $A$ (the reduced incidence algebra of finite subsets) as follows. Let ${V}_{n}$ be the set of functions from the set of $n$ -element subsets of $\Omega$ to the rationals $\mathbb{Q}$ . Now let $A = \bigoplus {V}_{n}$ , with multiplication as follows: for $f \in  {V}_{n}, g \in  {V}_{m}$ , and $\left| X\right|  = m + n$ , let $\left( {fg}\right) \left( X\right)  = \sum f\left( Y\right) g\left( {X \smallsetminus  Y}\right)$ , where the sum is over the $n$ -element subsets $Y$ of $X$ . If $G$ is a permutation group on $\Omega$ , let ${A}^{G}$ be the algebra of $G$ -invariants in $A$ .
	
	Suppose that $G$ has no finite orbits on $\Omega$ . It is known that then ${A}^{G}$ is an integral domain (M. Pouzet, Theor. Inf. App. 42, no. 1 (2008), 83-103). Is it true that the quotient of ${A}^{G}$ by the ideal generated by constant functions is also an integral
	
	domain? P. J. Cameron
\end{openproblem}

\plainproblemsection

\begin{openproblem}
	19.19. A finite transitive permutation group $G$ is said to have the road closure property if, given any orbit $O$ of $G$ on 2-sets, and any proper block of imprimitivity for $G$ acting on $O$ , the graph with edge set $O \smallsetminus  B$ is connected. Such a group must be primitive, and basic (not contained in a wreath product with the product action); it cannot have an imprimitive subgroup of index 2. In addition, such a group cannot be one of the permutation groups arising from triality (whose socle is ${D}_{4}\left( q\right)$ and intersects the point stabiliser in the parabolic subgroup corresponding to the three leaves in the Coxeter-Dynkin diagram for ${D}_{4}$ ).
	
	Road closure conjecture: these necessary conditions are sufficient for the road closure property. P. J. Cameron
\end{openproblem}

\plainproblemsection

\begin{openproblem}
	19.20. For a finite group $G$ , let $\operatorname{End}\left( G\right)$ denote the semigroup of endomorphisms of $G$ , and $\operatorname{PIso}\left( G\right)$ the semigroup of partial isomorphisms of $G$ (isomorphisms between subgroups of $G$ ). If $G$ is abelian, then $\left| {\operatorname{End}\left( G\right) }\right|  = \left| {\operatorname{PIso}\left( G\right) }\right|$ . Is the converse true?
\end{openproblem}

\plainproblemsection

\begin{openproblem}
	19.21. Can a non-discrete, compactly generated, topologically simple, locally compact group be amenable and non-compact?
	
	(A positive answer implies a negative answer to 19.107.) P.-E. Caprace
\end{openproblem}

\plainproblemsection

\begin{openproblem}
	19.22. Let $\mathcal{S}$ denote the class of nondiscrete compactly generated, topologically simple totally disconnected locally compact groups. Two topological groups are called locally isomorphic if they contain isomorphic open subgroups. Is the number of local isomorphism classes of groups in $\mathcal{S}$ uncountable? P.-E. Caprace
\end{openproblem}

\plainproblemsection

\begin{openproblem}
	*19.23. For a group $G$ , let ${\operatorname{Tor}}_{1}\left( G\right)$ be the normal closure of all torsion elements of $G$ , and then by induction let ${\operatorname{Tor}}_{i + 1}\left( G\right)$ be the inverse image of ${\operatorname{Tor}}_{1}\left( {G/{\operatorname{Tor}}_{i}\left( G\right) }\right)$ . The torsion length of $G$ is defined to be either the least positive integer $l$ such that $G/{\operatorname{Tor}}_{l}\left( G\right)$ is torsion-free, or $\omega$ if no such integer exists (since $G/\bigcup {\operatorname{Tor}}_{i}\left( G\right)$ is always torsion-free).
	
	Does there exists a finitely generated, or even finitely presented, soluble group with torsion length greater than 2 ? M. Chiodo, R. Vyas
\end{openproblem}

\plainproblemsection

\begin{openproblem}
	*19.24. For a group $G$ , let $\operatorname{Tor}\left( G\right)$ be the normal closure of all torsion elements of $G$ . Does there exist a finitely presented group $G$ such that $G/\operatorname{Tor}\left( G\right)$ is not finitely presented? Such a group must necessarily be non-hyperbolic.
	
	*Yes, such groups exist: one soluble example and another virtually torsion-free are constructed in (I. J. Leary, A. Minasyan, J. Group Theory, 26, no. 4 (2023), 741-750).
\end{openproblem}

\plainproblemsection

\begin{openproblem}
	19.25. Let $G$ and $H$ be finite groups of the same order with $\mathop{\sum }\limits_{{g \in  G}}\varphi \left( \left| g\right| \right)  =$ $\mathop{\sum }\limits_{{h \in  H}}\varphi \left( \left| h\right| \right)$ , where $\varphi$ is the Euler totient function. Suppose that $G$ is simple. Is $H$ necessarily simple? B. Curtin, G. R. Pourgholi
\end{openproblem}

\plainproblemsection

\begin{openproblem}
	19.26. (Y. O. Hamidoune). Suppose that $A$ and $B$ are finite subsets of a group $G$ such that $\left| A\right|  \geq  2$ and $\left| B\right|  \geq  2$ , and let $A \cdot  {}_{2}B$ denote the set of elements of $G$ which can be expressed in the form ${ab}$ for at least two different $\left( {a, b}\right)  \in  A \times  B$ . Is it true that $\left| A\right|  + \left| B\right|  - \left( {1/2}\right) \left| {AB}\right|  - \left( {1/2}\right) \left| {A{ \cdot  }_{2}B}\right|  \leq  \max \left\{  {2,\left| {gH}\right|  : H \leq  G, g \in  G,{gH} \subseteq  A{ \cdot  }_{2}B}\right\}  ?$ This is proved if $G$ is abelian. W. Dicks
\end{openproblem}

\plainproblemsection

\begin{openproblem}
	19.27. (Well-known question). A finitely generated group $G$ that acts on a tree in such a way that all vertex and edge stabilizers are infinite cyclic groups is called a generalized Baumslag-Solitar group. The Bass-Serre theory gives finite presentations of such groups. Is the isomorphism problem soluble for generalized Baumslag-Solitar groups? F. A. Dudkin
\end{openproblem}

\plainproblemsection

\begin{openproblem}
	19.28. Let $G$ be a group, and $\varphi$ an automorphism of $G$ . Elements $x, y \in  G$ are said to be $\varphi$ -conjugate if $x = {z}^{-1}{y\varphi }\left( z\right)$ for some $z \in  G$ . The $\varphi$ -conjugacy is an equivalence relation; the number of $\varphi$ -conjugacy classes is denoted by $R\left( \varphi \right)$ .
	
	Conjecture: If a finitely generated residually finite group $G$ has an automorphism $\varphi$ such that $R\left( \varphi \right)$ is finite, then $G$ has a soluble subgroup of finite index.
	
	The conjecture is proved if $\varphi$ has prime order (E. Jabara, J. Algebra,320, no. 10 (2008), 3671-3679). If $G$ is infinitely generated, then $G$ does not have to be almost solvable (K. Dekimpe, D. Gonçalves, Bull. London Math. Soc., 46, no. 4 (2014), 737- 746). A. L. Fel'shtyn, E. V. Troitsky
\end{openproblem}

\plainproblemsection

\begin{openproblem}
	19.29. Let $\Phi  \in  \operatorname{Out}G = \operatorname{Aut}G/\operatorname{Inn}G$ . Two automorphisms $\varphi ,\psi  \in  \Phi$ are said to be isogradient if $\varphi  = {\widehat{g}}^{-1}\psi \widehat{g}$ for some inner automorphism $\widehat{g}$ . The number of isogradiency classes in $\Phi$ is denoted by $S\left( \Phi \right)$ .
	
	Conjecture: If a finitely generated residually finite group $G$ has an outer automorphism $\Phi$ such that $S\left( \Phi \right)$ is finite, then $G$ has a soluble subgroup of finite index.
	
	A. L. Fel'shtyn, E. V. Troitsky
\end{openproblem}

\plainproblemsection

\begin{openproblem}
	19.30. An element $g$ of a finite group $G$ is said to be vanishing if $\chi \left( g\right)  = 0$ for some irreducible complex character $\chi  \in  \operatorname{Irr}\left( G\right)$ . Must a finite group and a finite simple group be isomorphic if they have equal orders and the same set of orders of vanishing
	
	elements? M. Foroudi Ghasemabadi, A. Iranmanesh,
\end{openproblem}

\plainproblemsection

\begin{openproblem}
	19.31. Let $\omega \left( {G, S}\right)$ be the exponential growth rate of a group $G$ with a finite generating set $S$ (see 14.7). Let ${MCG}\left( {\sum }_{g}\right)$ be the mapping class group of the orientable surface ${\sum }_{g}$ of genus $g$ . Is there a constant $C > 1$ such that $\omega \left( {{MCG}\left( {\sum }_{g}\right) , S}\right)  \geq  C$ for every $g > 0$ and every finite generating set $S$ of ${MCG}\left( {\sum }_{g}\right)$ ?
	
	If the genus $g$ is fixed, then such a constant $C > 1$ does exist (J. Mangahas, Geom. Funct. Anal. 19, no. 5 (2010), 1468-1480). $K$ . Fujiwara
\end{openproblem}

\plainproblemsection

\begin{openproblem}
	19.32. Let $p \geq  {673}$ be a prime and $r \geq  2$ be an integer. Let $B$ be the free group in the Burnside variety of exponent $p$ on ${a}_{1},\ldots ,{a}_{r}$ . Must every non-identical one-variable equation $w\left( {{a}_{1},\ldots ,{a}_{r}, x}\right)  = 1$ over $B$ have at most finitely many solutions in $B$ ? A. Gaglione
\end{openproblem}

\plainproblemsection

\begin{openproblem}
	19.33. Conjecture: Let $G$ be a finite group, $p$ a prime number, and $P$ a Sylow $p$ - subgroup of $G$ . Suppose that an irreducible ordinary character $\chi$ of $G$ has degree divisible by $p$ . If the restriction ${\chi }_{P}$ of $\chi$ to $P$ has a linear constituent, then ${\chi }_{P}$ has at least $p$ different linear constituents.
	
	This conjecture has been verified for symmetric, alternating, $p$ -solvable, and sporadic simple groups. E. Giannelli
\end{openproblem}

\plainproblemsection

\begin{openproblem}
	19.34. Let $G$ be a finitely generated group such that for every element $g \in  G$ the set of commutators $\{ \left\lbrack  {x, g}\right\rbrack   \mid  x \in  G\}$ is a subgroup of $G$ . Is it true that $G$ is residually nilpotent?
	
	The answer is affirmative if $G$ is finite. D. Gonçalves, T. Nasybullov
\end{openproblem}

\plainproblemsection

\begin{openproblem}
	*19.35. Let $G$ be a finite group of order $n$ . Is it true that for every factorization $n = {a}_{1}\cdots {a}_{k}$ there exist subsets ${A}_{1},\ldots ,{A}_{k}$ such that $\left| {A}_{1}\right|  = {a}_{1},\ldots ,\left| {A}_{k}\right|  = {a}_{k}$ and $G = {A}_{1}\cdots {A}_{k}?$ M. H. Hooshmand
	
	*No, it is not true. A counterexample with $k = 3$ is given by the alternating group on 4 letters $G = {A}_{4}$ and $\left( {{a}_{1},{a}_{2},{a}_{3}}\right)  = \left( {2,3,2}\right)$ (G. M. Bergman, Letter of 19 December 2019, https://math.berkeley.edu/~gbergman/papers/gp\_factzn.pdf).
\end{openproblem}

\plainproblemsection

\begin{openproblem}
	*19.36. Let $G$ be a periodic group and let $\mathcal{I} = \left\{  {x \in  G \mid  {x}^{2} = 1 \neq  x}\right\}$ be the set of its involutions. Let $D$ be a non-empty set of odd integers greater than 1 ; then $G$ is called a group with $D$ -involutions if $G = \langle \mathcal{I}\rangle$ and for $x, y \in  \mathcal{I}$ the order of ${xy}$ is in the set $\{ 1,2\}  \cup  D$ and all these values actually occur. It is clear that if $G$ is a group with $D$ -involutions, then $\mathcal{I}$ is a single conjugacy class.
	
	Conjecture: If $G$ is a group with $\{ 3,5\}$ -involutions, then $G \simeq  {A}_{5}$ or $G \simeq$ ${PSU}\left( {3,4}\right)$ . E. Jabara
	
	*The conjecture is proved, even without using the hypothesis that the group $G$ is periodic (E. Bettio, J. Group Theory, 24 (2021), 1055-1067).
\end{openproblem}

\plainproblemsection

\begin{openproblem}
	*19.37. a) Does there exist an absolute constant $k$ such that for any nilpotent injector $H$ of an arbitrary finite group $G$ there are $k$ conjugates of $H$ the intersection of which is equal to the Fitting subgroup $F\left( G\right)$ of $G$ ?
	
	b) Can one choose $k = 3$ as such a constant? This is true for finite soluble groups (D. S. Passman, Trans. Amer. Math. Soc., 123, no. 1 (1966), 99-111; A. Mann, Proc. Amer. Math. Soc., 53, no. 1 (1975), 262-264). S. F. Kamornikov
	
	*Yes, one can choose $k = 3$ as such a constant, mod CFSG (V.I. Zenkov, Siberian Math. J., 62, no. 4 (2021), 621-637).
\end{openproblem}

\plainproblemsection

\begin{openproblem}
	*19.38. Suppose that $H$ is a subgroup of a finite soluble group $G$ that covers all Frattini chief factors of $G$ and avoids all complemented chief factors of $G$ . Is it true that there are elements $x, y \in  G$ such that $H \cap  {H}^{x} \cap  {H}^{y} = {H}_{G}$ , where ${H}_{G}$ is the largest normal subgroup of $G$ contained in $H$ ?
	
	This is true if $H$ is a prefrattini subgroup of $G$ . S. F. Kamornikov
	
	*Yes, it is true (S. F. Kamornikov, O. L. Shemetkova, J. Algebra, 641 (2024), 1-8).
\end{openproblem}

\plainproblemsection

\begin{openproblem}
	19.39. A group is said to be $\mathfrak{X}$ -critical if it does not belong to $\mathfrak{X}$ but all its proper subgroups belong to $\mathfrak{X}$ . Let $\mathfrak{F}$ be a soluble hereditary formation of finite groups for which all $\mathfrak{F}$ -critical groups are either groups of prime order or $\mathfrak{U}$ -critical groups, where $\mathfrak{U}$ is the formation of all supersoluble finite groups. Must $\mathfrak{F}$ be a saturated formation?
	
	S. F. Kamornikov
\end{openproblem}

\plainproblemsection

\begin{openproblem}
	*19.40. Does Thompson’s group $F$ (see 12.20) have the Howson property, that is, is the intersection of any two finitely generated subgroups of $F$ finitely generated? I. Kapovich
	\par\smallskip\hrule\smallskip
	*No, it does not. Otherwise any subgroup of $F$ would have this property. But the wreath product $\mathbb{Z}\wr \mathbb{Z}$ , which does not have the Howson property (A. S. Kirkinskii, Algebra Logic, 20, no. 1 (1981), 24-36), can be embedded in $F$ (S. Cleary, Pacific J. Math., 228, no. 1 (2006), 53-61). (D. Robertson, Letter of 24 April 2018.)
	\par\smallskip\hrule\smallskip
\end{openproblem}

\plainproblemsection

\begin{openproblem}
	19.41. Let $\phi$ be an automorphism of a free group ${F}_{n}$ of finite rank, and let ${F}_{n}{ \rtimes  }_{\phi }\mathbb{Z}$ be the split extension of ${F}_{n}$ by $\mathbb{Z}$ with $\mathbb{Z}$ acting as $\langle \phi \rangle$ . Is the group ${F}_{n}{ \rtimes  }_{\phi }\mathbb{Z}$ conjugacy
	
	separable? I. Kapovich
\end{openproblem}

\plainproblemsection

\begin{openproblem}
	19.42. Suppose that $H$ is a word-hyperbolic subgroup of a word-hyperbolic group $G$ such that the inclusion of $H$ to $G$ extends to a continuous $H$ -equivariant map $j : \partial H \rightarrow  \partial G$ between their hyperbolic boundaries. If such an extension exists, it is unique and $j$ is called the Cannon-Thurston map.
	
	a) Is it true that for every point $p \in  \partial G$ its full preimage ${j}^{-1}\left( p\right)$ is finite?
	
	b) Moreover, is it true that there is a number $N = N\left( {G, H}\right)  < \infty$ such that for every $p \in  \partial G$ the full preimage ${j}^{-1}\left( p\right)$ consists of at most $N$ points?
	
	The answer to both questions is "yes" in all the cases where the Cannon-Thurston map is known to exist and where it has been possible to analyze the multiplicity of this map. This includes the original set-up considered by Cannon and Thurston where $H$ is the surface group, and $G$ is the fundamental group of a closed hyperbolic 3-manifold fibering over the circle with that surface as a fiber. In this case, $j : {\mathbb{S}}^{1} \rightarrow  {\mathbb{S}}^{2}$ is a uniformly finite-to-one continuous surjective 'Peano curve'. I. Kapovich
\end{openproblem}

\plainproblemsection

\begin{openproblem}
	19.43. Suppose that $\varphi$ is an automorphism of a finite soluble group $G$ . Is the Fitting height of $G$ bounded in terms of $\left| \varphi \right|$ and $\left| {{C}_{G}\left( \varphi \right) }\right|$ ? E. I. Khukhro
\end{openproblem}

\plainproblemsection

\begin{openproblem}
	19.44. By definition a profinite group has finite rank at most $r$ if every subgroup of it can be (topologically) generated by $r$ elements. Suppose that for every element $g$ of a profinite group $G$ there is a closed subgroup ${E}_{g}$ of finite rank such that for every $x \in  G$ all sufficiently long Engel commutators $\left\lbrack  {x, g,\ldots , g}\right\rbrack$ belong to ${E}_{g}$ , that is, for every $x \in  G$ there is a positive integer $n\left( {x, g}\right)$ such that $\left\lbrack  {x, g,\ldots , g}\right\rbrack   \in  {E}_{g}$ whenever $g$ is repeated $\geq  n\left( {x, g}\right)$ times. Is it true that $G$ has a normal subgroup $N$ of finite rank with locally nilpotent quotient $G/N$ ? E. I. Khukhro
\end{openproblem}

\plainproblemsection

\begin{openproblem}
	19.45. Let $G$ be a group generated by 3 class transpositions (see the definition in 17.57), and let $m$ be the least common multiple of the moduli of the residue classes interchanged by the generators of $G$ . Assume that $G$ does not setwisely stabilize any union of residue classes modulo $m$ except for $\varnothing$ and $\mathbb{Z}$ , and assume that the integers $0,1,\ldots ,{42}$ all lie in the same orbit under the action of $G$ on $\mathbb{Z}$ . Is the action of $G$ on $\mathbb{N} \cup  \{ 0\}$ necessarily transitive?
	
	The bound 42 cannot be replaced by a smaller number, since the finite group $\left\langle  {{\tau }_{0\left( 2\right) ,1\left( 2\right) },{\tau }_{0\left( 3\right) ,2\left( 3\right) },{\tau }_{0\left( 7\right) ,6\left( 7\right) }}\right\rangle$ acts transitively on the set $\{ 0,\ldots ,{41}\}$ , as well as on the set of residue classes modulo 42 . S. Kohl
\end{openproblem}

\plainproblemsection

\begin{openproblem}
	19.46. Does the group $\mathrm{{CT}}\left( \mathbb{Z}\right)$ have finitely generated infinite periodic subgroups? (See the definition of $\mathrm{{CT}}\left( \mathbb{Z}\right)$ in 17.57).S. Kohl
\end{openproblem}

\plainproblemsection

\begin{openproblem}
	19.47. Let $K$ be a finite extension of degree $n$ of a field $k$ of odd characteristic. The multiplicative group ${K}^{ * }$ embeds into the group of all $k$ -linear automorphisms ${\operatorname{Aut}}_{k}\left( K\right)$ by the rule $t\left( x\right)  = {tx}$ for all $x \in  K$ . For a fixed basis of $K$ over $k$ , the group ${\operatorname{Aut}}_{k}\left( K\right)$ is isomorphic to ${GL}\left( {n, k}\right)$ . The image of ${K}^{ * }$ is called a nonsplit maximal torus corresponding to the extension $K/k$ and is denoted by $T\left( {K/k}\right)$ . A subgroup of ${GL}\left( {n, k}\right)$ is said to be rich in transvections if it contains all elementary transvections.
	
	Let $H$ be a subgroup of ${GL}\left( {n, k}\right)$ which contains $T\left( {K/k}\right)$ and a one-dimensional transformation. Is $H$ rich in transvections?
\end{openproblem}

\plainproblemsection

\begin{openproblem}
	19.48. A system of additive subgroups ${\sigma }_{ij},1 \leq  i, j \leq  n$ , of a field $K$ is called a net (or a carpet) of order $n$ if ${\sigma }_{ir}{\sigma }_{rj} \subset  {\sigma }_{ij}$ for all $i, r, j$ . A net that does not contain the diagonal is called an elementary net. A net $\sigma  = \left( {\sigma }_{ij}\right)$ is said to be irreducible if all ${\sigma }_{ij}$ are nontrivial. An elementary net $\sigma$ is said to be closed if the elementary net subgroup $E\left( \sigma \right)$ does not contain additional elementary transvections.
	
	Let $R$ be a principal ideals domain with $1 \in  R$ , let $k$ be the field of fractions of $R$ , and $K$ an algebraic extension of the field $k$ . Let $\sigma  = \left( {\sigma }_{ij}\right)$ be an irreducible elementary net $\sigma  = \left( {\sigma }_{ij}\right)$ over $K$ such that all ${\sigma }_{ij}$ are $R$ -modules. Is the net $\sigma$ closed?
	
	V. A. Koibaev
\end{openproblem}

\plainproblemsection

\begin{openproblem}
	*19.49. A skew brace is a set $B$ equipped with two operations + and - such that $\left( {B, + }\right)$ is an additively written (but not necessarily abelian) group, $\left( {B, \cdot  }\right)$ is a multiplicatively written group, and $a \cdot  \left( {b + c}\right)  = {ab} - a + {ac}$ for any $a, b, c \in  B$ .
	
	Let $A$ be a skew brace with left-orderable multiplicative group. Is the additive group of $A$ left-orderable? V. Lebed, L. Vendramin
	
	*No, not always (T. Nasybullov, J. Algebra, 540 (2019), 156-167).
\end{openproblem}

\plainproblemsection

\begin{openproblem}
	*19.50. A finite graph is said to be integral if all eigenvalues of its adjacency matrix are integers.
	
	a) Let $G$ be a finite group generated by a normal subset $R$ consisting of involutions. Is it true that the Cayley graph $\operatorname{Cay}\left( {G, R}\right)$ is integral?
	
	b) Let ${A}_{n}$ be the alternating group of degree $n$ , let $S = \{ \left( {123}\right) ,\left( {124}\right) ,\ldots ,\left( {12n}\right) \}$ and $R = S \cup  {S}^{-1}$ . Is it true that the Cayley graph $\operatorname{Cay}\left( {{A}_{n}, R}\right)$ is integral?
	
	D. V. Lytkina
	
	*a) Yes, it is true (D. O. Revin, Letter of 21 April 2018; see also the reference for part (b) below; A. Abdollahi, Letter of 3 May 2018). Both proofs suggested are based on character theory; here is the second one. It suffices to show that the eigenvalues of $\operatorname{Cay}\left( {G, R}\right)$ are rational, since the eigenvalues of a simple graph are algebraic integers. It is known that every eigenvalue of $\operatorname{Cay}\left( {G, R}\right)$ has the form ${\theta }_{\chi } = \frac{1}{\chi \left( 1\right) }\mathop{\sum }\limits_{{r \in  R}}\chi \left( r\right)$ for some complex irreducible character $\chi$ of $G$ (implicit on pages 175-177 in P. Diaconis, M. Shahshahani, Z. Wahrscheinlichkeitstheorie Verw. Gebiete, 57 (1981), 159-179, see also Theorem 9 in M. R. Murty, J. Ramanujan Math. Soc., 18, no. 1 (2003) 1-20). Since the value of any complex character on an involution is an integer, it follows that ${\theta }_{\chi }$ is rational.
	
	*b) Yes, it is true (W. Guo, D. V. Lytkina, V. D. Mazurov, D. O. Revin, Algebra Logic, 58, no. 4 (2019), 297-305).
\end{openproblem}

\plainproblemsection

\begin{openproblem}
	19.51. A finite group is monomial if each irreducible character of it is induced from a linear character of some subgroup. Monomial groups are soluble (Taketa). The group is normally (subnormally) monomial if each irreducible character is induced from a linear character of some normal (respectively, subnormal) subgroup. Metabelian groups are normally monomial, and abelian-by-nilpotent groups are subnormally monomial. In the other direction, it is known that there exist normally monomial groups of arbitrarily large derived length.
	
	a) For a given prime $p$ , do there exist normally monomial finite $p$ -groups of arbitrarily large derived length?
	
	b) Do there exist subnormally monomial groups of arbitrarily large nilpotence length? A. Mann
\end{openproblem}

\plainproblemsection

\begin{openproblem}
	19.52. The Gruenberg-Kegel graph (or the prime graph) ${GK}\left( G\right)$ of a finite group $G$ has vertex set consisting of all prime divisors of the order of $G$ , and different vertices $p$ and $q$ are adjacent in ${GK}\left( G\right)$ if and only if the number ${pq}$ is an element order in $G$ . Is there a finite non-solvable group $G$ such that ${GK}\left( G\right)$ does not contain 3-cocliques and is not isomorphic to the Gruenberg-Kegel graph of any finite solvable group?
	
	It is known that there are no examples of such groups $G$ among almost simple groups. N. V. Maslova
\end{openproblem}

\plainproblemsection

\begin{openproblem}
	19.53. Let $G$ be a group generated by elements $x, y, z$ such that
	
	$$
	{x}^{3} = {y}^{2} = {z}^{2} = {\left( xy\right) }^{3} = {\left( yz\right) }^{3} = 1\;\text{ and }\;{g}^{12} = 1\text{ for all }g \in  G.
	$$
	
	Is it true that $\left| G\right|  \leq  {12}$ ?
	
	This is true if $G$ is finite (D. V. Lytkina, V. D. Mazurov, Siberian Math. J.,56, no. 3 (2015), 471-475). V. D. Mazurov
\end{openproblem}

\plainproblemsection

\begin{openproblem}
	19.54. What are the chief factors of a finite group in which every 2-maximal subgroup is not $m$ -maximal for any $m \geq  3$ ?
	
	A subgroup $H$ of a group $G$ is said to be $m$ -maximal if there is a chain of subgroups $H = {H}_{0} < {H}_{1} < \ldots  < {H}_{m - 1} < {H}_{m} = G$ in which ${H}_{i}$ is maximal in ${H}_{i + 1}$ for every $i$ . Note that for every $m \geq  3$ there is a finite group in which some 2-maximal subgroup is $m$ -maximal. V. S. Monakhov
\end{openproblem}

\plainproblemsection

\begin{openproblem}
	*19.55. Suppose that in a finite group $G$ every maximal subgroup $M$ is supersoluble whenever $\pi \left( M\right)  = \pi \left( G\right)$ , where $\pi \left( G\right)$ is the set of all prime divisors of the order of $G$ .
	
	a) What are the non-abelian composition factors of $G$ ?
	
	b) Determine the exact upper bounds for the nilpotency length, the rank, and the $p$ -length of $G$ if $G$ is soluble. V. S. Monakhov
	
	*a) Every nonabelian finite simple group can occur as a composition factor of $G$ (A. Moretó, Monatsh. Math., 195, no. 3 (2021), 497-500).
	
	*b) There is not any bound for the nilpotency length or the rank, but the $p$ -length is at most 1 for every prime $p$ (A. Moretó, Monatsh. Math.,195, no. 3 (2021), 497-500).
\end{openproblem}

\plainproblemsection

\begin{openproblem}
	19.56. A $\star$ -commutator is a commutator $\left\lbrack  {x, y}\right\rbrack$ of two elements $x, y$ of coprime prime-power orders. Is a finite group $G$ soluble if $\left| {ab}\right|  \geq  \left| a\right| \left| b\right|$ for any $\star$ -commutators $a$ and $b$ of coprime orders? V. S. Monakhov
\end{openproblem}

\plainproblemsection

\begin{openproblem}
	19.57. What are the non-abelian composition factors of a finite group in which every maximal subgroup is simple or $p$ -nilpotent for some fixed odd prime $p \in  \pi \left( G\right)$ ? V. S. Monakhov, I. N. Tyutyanov
\end{openproblem}

\plainproblemsection

\begin{openproblem}
	19.58. What are the non-abelian composition factors of a finite group in which every maximal subgroup is simple or $p$ -decomposable for some fixed odd prime $p \in  \pi \left( G\right)$ ? V. S. Monakhov, I. N. Tyutyanov
\end{openproblem}

\plainproblemsection

\begin{openproblem}
	19.59. Does the group of isometries of ${\mathbb{Q}}^{3}$ have a free non-abelian subgroup such that every non-trivial element acts without nontrivial fixed points in ${\mathbb{Q}}^{3}$ ?
	
	The answer is yes if the field of rational numbers $\mathbb{Q}$ is replaced by the field $\mathbb{R}$ of real numbers (see G. Tomkowicz, S. Wagon, The Banach-Tarski Paradox, Cambridge Univ. Press, 2016.) J. Mycielski
\end{openproblem}

\plainproblemsection

\begin{openproblem}
	19.60. Let $R$ be a principal ideal domain, let $d$ be a positive integer, and let ${X}_{1},\ldots ,{X}_{m}$ be members of $\mathrm{{SL}}\left( {d, R}\right)$ (or of $\mathrm{{GL}}\left( {d, R}\right)$ ). Is it decidable whether or not these $m$ matrices freely generate a free group? If it is, design an effective algorithm for computing the answer. Peter M. Neumann
\end{openproblem}

\plainproblemsection

\begin{openproblem}
	19.61. Let $\mathfrak{A} = \left\{  {{\mathfrak{A}}_{r} \mid  r \in  \Phi }\right\}$ be an elementary carpet of type $\Phi$ over a commutative ring $K$ (see 7.28), and let $\Phi \left( \mathfrak{A}\right)  = \left\langle  {{x}_{r}\left( {\mathfrak{A}}_{r}\right)  \mid  r \in  \Phi }\right\rangle$ be its carpet subgroup. Define the closure of the carpet $\mathfrak{A}$ to be the set of additive subgroups $\overline{\mathfrak{A}} = \left\{  {{\overline{\mathfrak{A}}}_{r} \mid  r \in  \Phi }\right\}$ , where ${\overline{\mathfrak{A}}}_{r} = \left\{  {t \in  K \mid  {x}_{r}\left( t\right)  \in  \Phi \left( \mathfrak{A}\right) }\right\}$ . Is the closure $\overline{\mathfrak{A}}$ of a carpet $\mathfrak{A}$ always a carpet?
	
	An affirmative answer is known if $\Phi  = {A}_{l},{D}_{l},{E}_{l}$ . Ya. N. Nuzhin
\end{openproblem}

\plainproblemsection

\begin{openproblem}
	19.62. Let $\mathfrak{A} = \left\{  {{\mathfrak{A}}_{r} \mid  r \in  \Phi }\right\}$ be an elementary carpet of type $\Phi$ of rank $l \geq  2$ (see 7.28). For $p \in  \Phi$ , define a set of additive subgroups ${\mathfrak{B}}_{p} = \sum {C}_{{ij},{rs}}{\mathfrak{A}}_{r}^{i}{\mathfrak{A}}_{s}^{j}$ , where the sum is taken over all natural numbers $i, j$ and roots $r, s \in  \Phi$ such that ${ir} + {js} = p$ . It is known that the set $\mathfrak{B} = \left\{  {{\mathfrak{B}}_{p} \mid  p \in  \Phi }\right\}$ is a carpet called the derived carpet of $\mathfrak{A}$ . It is also known that for $\Phi  = {A}_{l}$ the set $\mathfrak{B}$ is a closed (admissible) carpet, which means that its carpet subgroup does not contain new root elements. Is every derived carpet of type $\Phi$ over a commutative ring closed (admissible)? Ya. N. Nuzhin
\end{openproblem}

\plainproblemsection

\begin{openproblem}
	19.63. Let $\mathfrak{A} = \left\{  {{\mathfrak{A}}_{r} \mid  r \in  \Phi }\right\}$ be an elementary carpet of type $\Phi$ over a commutative ring $K$ (see 7.28) and let ${\mathfrak{A}}_{r}^{2} = \left\{  {{t}^{2} \mid  t \in  {\mathfrak{A}}_{r}}\right\}$ . Are the inclusions ${\mathfrak{A}}_{r}^{2}{\mathfrak{A}}_{-r} \subseteq  {\mathfrak{A}}_{r}, r \in  \Phi$ , sufficient for the carpet $\mathfrak{A}$ to be closed (admissible)? Ya. N. Nuzhin
\end{openproblem}

\plainproblemsection

\begin{openproblem}
	19.64. Let $G$ be a group, and $\left( {{g}_{1},\ldots ,{g}_{n}}\right)$ a tuple of its elements. The type of this tuple in $G$ , denoted $T{p}^{G}\left( {{g}_{1},\ldots ,{g}_{n}}\right)$ , is the set of all first order formulas in free variables ${x}_{1},\ldots ,{x}_{n}$ in the standard group theory language which are true on $\left( {{g}_{1},\ldots ,{g}_{n}}\right)$ in $G$ . Two groups $G$ and $H$ are called isotypic if for every tuple of elements $\bar{h} = \left( {{h}_{1},\ldots ,{h}_{n}}\right)$ in $H$ there is a tuple $\bar{g} = \left( {{g}_{1},\ldots ,{g}_{n}}\right)$ in $G$ , such that $\operatorname{Tp}\left( \bar{h}\right)  = \operatorname{Tp}\left( \bar{g}\right)$ and vice versa, for every tuple $\bar{g}$ in $G$ there is a tuple $\bar{h}$ in $H$ such that ${Tp}\left( \bar{h}\right)  = {Tp}\left( \bar{g}\right)$ .
	
	Is it true that every two isotypic finitely generated groups are isomorphic?
	
	The answer is positive if one of the finitely generated groups is free (R. Sklinos), abelian (G. Zhitomirski), virtually polycyclic, metabelian, free solvable (A. Myasnikov, N. Romanovskii), co-Hopfian (in particular homogeneous), finitely presented Hopfian or geometrically Noetherian Hopfian (R. Sklinos). In particular, torsion-free hyperbolic groups, braid groups, linear groups, mapping class groups of compact surfaces, limit groups, quasi-cyclic groups, $\operatorname{Out}\left( {F}_{n}\right) , n > 2$ , irreducible arithmetic lattices (which are not virtually free groups) in semi-simple Lie groups, are of such kind (R. Sklinos).
\end{openproblem}

\plainproblemsection

\begin{openproblem}
	19.65. It is well known that varieties of groups form a free semigroup $N$ (A. L. Shmel'kin, DAN SSSR, 149 (1963), 543-545 (Russian)); B. H. Neuman, H. Neumann, P. M. Neumann, Math. Z., 80 (1962), 44-62). Varieties of linear representations over an infinite field of zero characteristic also form a free semigroup $M$ , and $N$ acts freely on $M$ (B. I. Plotkin, Siberian Math. J.,13, no. 5 (1972),713-729) A similar theorem holds for Lie algebras (V. A. Parfenov, Algebra Logika, 6, no. 4 (1967), 61-73 (Russian); L. A. Simonyan, Siberian Math. J., 29, no. 2 (1988), 276-283).
	
	Are there other varieties of algebras $\Theta$ where the same situation takes place?
	
	B. I. Plotkin
\end{openproblem}

\plainproblemsection

\begin{openproblem}
	19.66. We say that a variety $\Theta$ is of Tarski type if any two non-abelian $\Theta$ -free groups of finite rank are elementarily equivalent.
	
	a) Find examples of Tarski type varieties distinct from the variety of all groups.
	
	b) Is it true that the Burnside variety ${B}_{n}$ of all groups of exponent $n$ , where $n$ is big enough, is of Tarski type?
	
	c) Is it true that the $n$ -Engel variety ${E}_{n}$ of all groups satisfying the identity $\left\lbrack  {\left\lbrack  {\left\lbrack  {x, y}\right\rbrack  , y}\right\rbrack  ,\ldots , y}\right\rbrack   \equiv  1$ , where $n$ is big enough, is of Tarski type?
\end{openproblem}

\plainproblemsection

\begin{openproblem}
	*19.67. Let $G \leq  \operatorname{Sym}\left( \Omega \right)$ , where $\Omega$ is finite. The 2-closure ${G}^{\left( 2\right) }$ of the group $G$ is defined to be the largest subgroup of $\operatorname{Sym}\left( \Omega \right)$ containing $G$ which has the same orbits as $G$ in the induced action on $\Omega  \times  \Omega$ . Is it true that if $G$ is solvable, then every composition factor of ${G}^{\left( 2\right) }$ is either a cyclic or an alternating group?
	
	I. Ponomarenko
	
	*No, it is not true (S.V. Skresanov, Algebra Logic, 58, no. 3 (2019), 249-253).
\end{openproblem}

\plainproblemsection

\begin{openproblem}
	19.68. For a finite group $G$ and a permutation group $K$ , let ${b}_{G}\left( K\right)$ denote the number of conjugacy classes of regular subgroups of $K$ isomorphic to $G$ . Does there exist a function $f$ such that ${b}_{G}\left( K\right)  \leq  {n}^{f\left( r\right) }$ for every abelian group $G$ of order $n$ and rank $r$ , and every group $K$ such that ${K}^{\left( 2\right) } = K$ ? (See 19.67 for the definition of ${K}^{\left( 2\right) }$ .)
	
	I. Ponomarenko
\end{openproblem}

\plainproblemsection

\begin{openproblem}
	19.69. (P. Wesolek). The acronym tdlc stands for totally disconnected and locally compact, and tdlcsc for tdlc and second countable. Let Res(G) denote the intersection of all open normal subgroups of a topological group $G$ . The class of elementary tdlcsc groups is defined as the smallest class $\mathcal{E}$ of tdlcsc groups such that
	
	(1) $\mathcal{E}$ contains all second countable profinite groups and countable discrete groups;
	
	(2) $\mathcal{E}$ is closed under taking closed subgroups, Hausdorff quotients, directed unions of open subgroups, and group extensions.
	
	This class admits a well-behaved, ordinal-valued decomposition rank $\xi$ defined recursively as follows: $\xi \left( {\{ 1\} }\right)  = 1$ , and if $G \in  \mathcal{E}$ is a union of an increasing sequence $\left( {O}_{i}\right)$ of compactly generated open subgroups, then $\xi \left( G\right)  = \mathop{\sup }\limits_{i}\left\{  {\xi \left( {\operatorname{Res}\left( {\mathrm{O}}_{\mathrm{i}}\right) }\right) }\right\}   + 1$ .
	
	What is the supremum of the decomposition ranks of elementary tdlcsc groups? In particular, is it countable?C. Reid
\end{openproblem}

\plainproblemsection

\begin{openproblem}
	19.70. (P. Wesolek and P.-E. Caprace). Let $\mathcal{S}$ denote the class of nondiscrete compactly generated, topologically simple tdlc groups. Let $G$ be a non-elementary tdlcsc group. Is there a compactly generated closed subgroup $H \leq  G$ such that $H$ has a continuous quotient in $\mathcal{S}$ ? C. Reid
\end{openproblem}

\plainproblemsection

\begin{openproblem}
	19.71. (G. Willis). Can a group $G$ in $\mathcal{S}$ be such that every element of $G$ normalizes a compact open subgroup of $G$ ? C. Reid
\end{openproblem}

\plainproblemsection

\begin{openproblem}
	19.72. Let $G$ be a group in $\mathcal{S}$ . Can every element of $G$ have trivial contraction group? C. Reid
\end{openproblem}

\plainproblemsection

\begin{openproblem}
	19.73. (P.-E. Caprace and N. Monod). Is there a compactly generated, locally compact group that is topologically simple, but not abstractly simple?
	
	C. Reid, S. M. Smith
\end{openproblem}

\plainproblemsection

\begin{openproblem}
	*19.74. A subgroup $H$ of a group $G$ is called pronormal if $H$ and ${H}^{g}$ are conjugate in $\left\langle  {H,{H}^{g}}\right\rangle$ for every $g \in  G$ . A subgroup $H$ of a group $G$ is called abnormal if $g \in  \left\langle  {H,{H}^{g}}\right\rangle$ for every $g \in  G$ .
	
	Does there exist an infinite group that does not contain nontrivial proper pronor-mal subgroups?
	
	The question is equivalent to the following: does there exist an infinite simple group that does not contain proper abnormal subgroups? D. O. Revin
	
	*Yes, it does (S. Corson, Preprint, 2025, https://arxiv.org/abs/2505.07330).
\end{openproblem}

\plainproblemsection

\begin{openproblem}
	*19.75. Let $P$ be a finite 2-group of exponent ${2}^{e}$ such that the rank of every abelian subgroup is at most $r$ . Is it true that $\left| P\right|  \leq  {2}^{r\left( {e + 1}\right) }$ ? This bound would be sharp (for a direct product of quaternion groups). B. Sambale
	
	*No, it is not true, as follows from the examples constructed in (A. Yu. Ol'shanski', Math. Notes, 23 (1978), 183-185). (A. Mann, Letter of 23 April 2018.)
\end{openproblem}

\plainproblemsection

\begin{openproblem}
	19.76. A semigroup presentation is called tree-like if all relations have the form $a = {bc}$ where $a, b, c$ are letters and no two relations share the left-hand side or the right-hand side. Is it decidable whether the semigroup given by a finite tree-like presentation contains an idempotent?
	
	This is equivalent to the question whether the closure of a finitely generated subgroup of R. Thompson’s group $F$ contains an isomorphic copy of $F$ . M. Sapir
\end{openproblem}

\plainproblemsection

\begin{openproblem}
	19.77. The Stable Small Cancellation Conjecture: If $F$ is a free group of rank $r \geq  2$ , then, for any fixed $n > 0$ , there exists a generic subset $S$ of ${F}^{n}$ such that, for any automorphism $\phi$ of $F$ , the set $\phi \left( S\right)$ satisfies the small cancellation condition ${C}^{\prime }\left( {1/6}\right)$ .
	
	P. E. Schupp 19.78. Does there exist an infinite simple subgroup of ${SL}\left( {2,\mathbb{Q}}\right)$ ?
	
	This is a special case of a question which Serge Cantat asked me about "nonalgebraic" simple subgroups of ${GL}\left( {n,\mathbb{C}}\right)$ . It is also a special case of 15.57. J.-P. Serre
\end{openproblem}

\plainproblemsection

\begin{openproblem}
	19.79. (T. Springer). Let $G$ be a group. Suppose that for every integer $n > 0$ the group $G$ has a unique (up to isomorphism) irreducible complex linear representation of dimension $n$ . (Note that $G = {SL}\left( {2,\mathbb{Q}}\right)$ has these properties, by a theorem of Borel-Tits (Ann. Math., $\mathbf{{97}}\left( {1973}\right) ,{499} - {571}$ ).) Is it true that $G$ has a normal subgroup $N$ such that $G/N$ is isomorphic to ${SL}\left( {2,\mathbb{Q}}\right)$ ? J.-P. Serre
\end{openproblem}

\plainproblemsection

\begin{openproblem}
	*19.80. For a periodic group $G$ , let ${\pi }_{e}\left( G\right)$ denote the set of orders of elements of $G$ . A periodic group $G$ is said to be an $O{C}_{n}$ -group if ${\pi }_{e}\left( G\right)  = \{ 1,2,\ldots , n\}$ . Is it true that every $O{C}_{7}$ -group is isomorphic to the alternating group ${A}_{7}$ ?
	
	For finite groups, the answer is affirmative.W. J. Shi
	
	*Yes, it is true (E. Jabara, A. S. Mamontov, Siberian Math. J., 62, no. 1 (2021), 94-104).
\end{openproblem}

\plainproblemsection

\begin{openproblem}
	*19.81. (Well-known problem). Is the conjugacy problem in the braid group ${B}_{n}$ in the class NP (that is, decidable in nondeterministic polynomial time with respect to the maximum of the lengths $\left| u\right| ,\left| v\right|$ , where $u, v$ are the input braid words)? A stronger question: given two conjugate elements of ${B}_{n}$ represented by braid words of lengths $\leq  m$ , is there a conjugator whose length is bounded by a polynomial function of $m$ ? V. Shpilrain
	
	*Yes, there is such a conjugator, even for any mapping class groups (J. Tao, Geom. Funct. Anal., 23, no. 1 (2013), 415-466).
\end{openproblem}

\plainproblemsection

\begin{openproblem}
	19.82. Let ${h}^{ * }\left( G\right)$ denote the generalized Fitting height of a finite group $G$ defined as the minimum number $k$ such that ${F}_{k}^{ * }\left( G\right)  = G$ , where ${F}_{1}^{ * }\left( G\right)  = {F}^{ * }\left( G\right)$ is the generalized Fitting subgroup of $G$ , and by induction ${F}_{i + 1}^{ * }\left( G\right)$ is the inverse image of ${F}^{ * }\left( {G/{F}_{i}^{ * }\left( G\right) }\right)$ . If $G$ is soluble, then ${h}^{ * }\left( G\right)  = h\left( G\right)$ is the Fitting height of $G$ .
	
	Does every finite group $G$ contain a soluble subgroup $K$ such that ${h}^{ * }\left( G\right)  = h\left( K\right)$ ?
	
	P. Shumyatsky
\end{openproblem}

\plainproblemsection

\begin{openproblem}
	19.83. An element $g$ of a group $G$ is almost Engel if there is a finite set $\mathcal{E}\left( g\right)$ such that for every $x \in  G$ all sufficiently long commutators $\left\lbrack  {x,{}_{n}g}\right\rbrack$ belong to $\mathcal{E}\left( g\right)$ , that is, for every $x \in  G$ there is a positive integer $n\left( {x, g}\right)$ such that $\left\lbrack  {x,{}_{n}g}\right\rbrack   \in  \mathcal{E}\left( g\right)$ whenever $n \geq  n\left( {x, g}\right)$ . By a linear group we understand a subgroup of ${GL}\left( {m, F}\right)$ for some field $F$ and a positive integer $m$ .
	
	Is the set of almost Engel elements in a linear group always a subgroup?
	
	P. Shumyatsky
\end{openproblem}

\plainproblemsection

\begin{openproblem}
	*19.84. Let $\mathbb{P}$ be the set of all primes, and let $\sigma  = \left\{  {{\sigma }_{i} \mid  i \in  I}\right\}$ be some partition of $\mathbb{P}$ into disjoint subsets. A finite group $G$ is said to be $\sigma$ -primary if $G$ is a ${\sigma }_{i}$ -group for some $i;\sigma$ -nilpotent if $G$ is a direct product of $\sigma$ -primary groups; $\sigma$ -soluble if every chief factor of $G$ is $\sigma$ -primary. A subgroup $A$ of a finite group $G$ is said to be $\sigma$ -subnormal in $G$ if there is a chain $A = {A}_{0} \leq  {A}_{1} \leq  \cdots  \leq  {A}_{n} = G$ such that for every $i$ either ${A}_{i - 1} \trianglelefteq  {A}_{i}$ or ${A}_{i}/{\left( {A}_{i - 1}\right) }_{{A}_{i}}$ is $\sigma$ -primary, where ${\left( {A}_{i - 1}\right) }_{{A}_{i}}$ is the largest normal subgroup of ${A}_{i}$ contained in ${A}_{i - 1}$ .
	
	Suppose that a subgroup $A$ of a finite group $G$ is $\sigma$ -subnormal in $\left\langle  {A,{A}^{x}}\right\rangle$ for all $x \in  G$ . Is it true that then $A$ is $\sigma$ -subnormal in $G$ ?
	
	An affirmative answer is known if $\sigma  = \{ \{ 2\} ,\{ 3\} ,\ldots \}$ (Wielandt). $A.N.{Skiba}$
	
	*No, not always. For example, in $G = {S}_{5}$ with partition $\sigma  = \{ 2,3\}  \cup  \{ 5\}$ the subgroup $A = \langle \left( {12}\right) \rangle$ is $\sigma$ -subnormal in $\left\langle  {A,{A}^{x}}\right\rangle$ for every $x \in  G$ , but it is not $\sigma$ - subnormal in $G$ . (V. N. Tyutyanov, Letter of 28 August 2019.)
\end{openproblem}

\plainproblemsection

\begin{openproblem}
	*19.85. Suppose that every Schmidt subgroup of a finite group $G$ is $\sigma$ -subnormal in $G$ (see 19.84). Is it true that then there is a normal $\sigma$ -nilpotent subgroup $N \leq  G$ such that $G/N$ is cyclic?
	
	An affirmative answer is known if $\sigma  = \{ \{ 2\} ,\{ 3\} ,\ldots \}$ .A. N. Skiba
	
	*Yes, it is true (X. Yi, S. F. Kamornikov, J. Algebra, 560 (2020), 181-191).
\end{openproblem}

\plainproblemsection

\begin{openproblem}
	19.86. Suppose that a finite group $G$ has a ${\sigma }_{i}$ -Hall subgroup for every $i \in  I$ . Suppose that a subgroup $A \leq  G$ is such that $A \cap  H$ is a ${\sigma }_{i}$ -Hall subgroup of $A$ for every $i \in  I$ and every ${\sigma }_{i}$ -Hall subgroup $H$ of $G$ (see 19.84). Is it true that then $A$ is $\sigma$ -subnormal in $G$ ?
	
	An affirmative answer is known if $\sigma  = \{ \{ 2\} ,\{ 3\} ,\ldots \}$ and if $G$ is $\sigma$ -soluble.
	
	A. N. Skiba
\end{openproblem}

\plainproblemsection

\begin{openproblem}
	*19.87. Suppose that for every Sylow subgroup $P$ of a finite group $G$ and every maximal subgroup $V$ of $P$ there is a $\sigma$ -soluble subgroup $T$ such that ${VT} = G$ . Is it true that then $G$ is $\sigma$ -soluble? A. N. Skiba
	
	*Yes, it is true (A.-M. Liu, W. Guo, I. N. Safonova, A. N. Skiba, to appear in J. Algebra, 2021); another solution using CFSG is in (S. F. Kamornikov, V. N. Tyutyanov, Trudy Inst. Mat. Mekh. UrO RAN, 27, no. 1 (2021), 98-102 (Russian); X. Yi, S. F. Kamornikov, V. N. Tyutyanov, Probl. Physics, Math. Techn., 46, no. 1 (2021), ${50} - {53})$ .
\end{openproblem}

\plainproblemsection

\begin{openproblem}
	*19.88. Suppose that for every Sylow subgroup $P$ of a finite group $G$ and every maximal subgroup $V$ of $P$ there is a $\sigma$ -nilpotent subgroup $T$ such that ${VT} = G$ . Is it true that then $G$ is $\sigma$ -nilpotent? A. N. Skiba
	\par\smallskip\hrule\smallskip
	*Yes, it is true (A.-M. Liu, W. Guo, I. N. Safonova, A. N. Skiba, to appear in J. Algebra, 2021); another solution using CFSG is in (S. F. Kamornikov, V. N. Tyutyanov, Trudy Inst. Mat. Mekh. UrO RAN, 27, no. 1 (2021), 98-102 (Russian); X. Yi, S. F. Kamornikov, V. N. Tyutyanov, Probl. Physics, Math. Techn., 46, no. 1 (2021), ${50} - {53})$ .
	\par\smallskip\hrule\smallskip
\end{openproblem}

\plainproblemsection

\begin{openproblem}
	19.89. A permutation group is subdegree-finite if every orbit of every point stabiliser is finite. Let $\Omega$ be countably infinite, and let $G$ be a closed and subdegree-finite subgroup of $\operatorname{Sym}\left( \Omega \right)$ regarded as a topological group under the topology of pointwise convergence.
	
	Conjecture: If the minimal degree of $G$ is infinite, then there is some subset of $\Omega$ whose setwise stabiliser in $G$ is trivial.
	
	Note that this conjecture implies Tom Tucker's well-known Infinite Motion Conjecture for graphs. S. M. Smith
\end{openproblem}

\plainproblemsection

\begin{openproblem}
	19.90. A skew brace is a set $B$ equipped with two operations + and - such that $\left( {B, + }\right)$ is an additively written (but not necessarily abelian) group, $\left( {B, \cdot  }\right)$ is a multiplicatively written group, and $a \cdot  \left( {b + c}\right)  = {ab} - a + {ac}$ for any $a, b, c \in  B$ .
	
	*a) Is there a skew brace with soluble additive group but non-soluble multiplicative group?
	
	b) Is there a skew brace with non-soluble additive group but nilpotent multiplicative group?
	
	c) Is there a finite skew brace with soluble additive group but non-soluble multiplicative group?
	
	*d) Is there a finite skew brace with non-soluble additive group but nilpotent multiplicative group? A. Smoktunowicz, L. Vendramin
	\par\smallskip\hrule\smallskip
	*a) Yes, there is (T. Nasybullov, J. Algebra, 540 (2019), 156-167).
	
	*d) No, there is not (C. Tsang, Q. Chao, Int. J. Algebra Comput., 30, no. 2 (2020), 253-265).
	\par\smallskip\hrule\smallskip
\end{openproblem}

\plainproblemsection

\begin{openproblem}
	19.91. Let $G$ be a finite group with an abelian Sylow $p$ -subgroup $A$ . Suppose that $B$ is a strongly closed elementary abelian subgroup of $A$ . Without invoking the Classification Theorem for Finite Simple Groups (CFSG), prove that $G$ has a normal subgroup $N$ such that $B = {\Omega }_{1}\left( {A \cap  N}\right)$ .
	
	For $p = 2$ , this is a corollary of a theorem of Goldschmidt. For $p$ odd, this has been proved by Flores and Foote (Adv. Math., 222 (2009), 453-484), but their proof relies on CFSG. A CFSG-free proof in the special case when $p = 3$ and A has 3- rank 3 would already be quite interesting. This case arises in Aschbacher's treatment of the $e\left( G\right)  = 3$ problem, and a proof would provide an alternative to part of his argument. (For this application, one could assume that all proper simple sections of $G$ are known.) R. Solomon
\end{openproblem}

\plainproblemsection

\begin{openproblem}
	19.92. Let $A$ be the algebra of $3 \times  3$ skew-Hermitian matrices over the real octonions, where multiplication is given by bracket product. Determine $\operatorname{Aut}\left( A\right)$ .
	
	The question might be of interest for octonion algebras over a different base field or ring. The question was inspired by an observation of John Faulkner concerning a possible connection between $A$ or a related algebra and the Dwyer-Wilkerson 2- compact group ${BDI}\left( 4\right)$ . R. Solomon
\end{openproblem}

\plainproblemsection

\begin{openproblem}
	19.93. Conjecture: There exists a function $f : \mathbb{N} \times  \mathbb{N} \rightarrow  \mathbb{N}$ such that, if $G$ is a finite $p$ -group with $d$ generators and $G$ has no epimorphic images isomorphic to the wreath product ${C}_{p}\wr {C}_{p}$ , then each factor of the $p$ -lower central series of $G$ has order bounded above by $f\left( {p, d}\right)$ .
	
	It is interesting to compare this conjecture with the celebrated characterization of Aner Shalev of finitely generated $p$ -adic analytic pro- $p$ -groups.
\end{openproblem}

\plainproblemsection

\begin{openproblem}
	19.94. By definition two elements ${g}_{1}$ and ${g}_{2}$ of a free metabelian group $F$ with basis $\left\{  {{x}_{1},\ldots ,{x}_{n}}\right\}$ have the same distribution if the equations ${g}_{1}\left( {{x}_{1},\ldots ,{x}_{n}}\right)  = g$ and ${g}_{2}\left( {{x}_{1},\ldots ,{x}_{n}}\right)  = g$ have the same number of solutions in any finite metabelian group $G$ for every $g \in  G$ . Is it true that elements ${g}_{1}$ and ${g}_{2}$ have the same distribution if and only if they are conjugate by some automorphism of $F$ ? E. I. Timoshenko
\end{openproblem}

\plainproblemsection

\begin{openproblem}
	19.95. Let ${G}_{\Gamma }$ be a partially commutative soluble group of derived length $n \geq  3$ with defining graph $\Gamma$ (the definition is similar to the case of $n = 2$ , see 17.104). Is it true that the centralizer of any vertex of the graph is generated by its adjacent vertices?
	
	E. I. Timoshenko
\end{openproblem}

\plainproblemsection

\begin{openproblem}
	19.96. Let $M$ be a free metabelian group of finite rank $r \geq  2$ , and let $P$ be the set of its primitive elements. (An element of a relatively free group is said to be primitive if it can be included in a basis of the group.)
	
	a) Is $P$ a first-order definable set?
	
	b) Is the set of bases of the group $M$ a first-order definable set?
	
	E. I. Timoshenko
\end{openproblem}

\plainproblemsection

\begin{openproblem}
	19.97. Let $G$ be an $m$ -generated group the elementary theory of which $\operatorname{Th}\left( G\right)$ coincides with the elementary theory ${Th}\left( F\right)$ of a free soluble group $F$ of derived length $n \geq  3$ of finite rank $r \geq  2$ . Must the groups $G$ and $F$ be isomorphic?
	
	The answer is affirmative if $m \leq  r$ or $n = 2$ E. I. Timoshenko
\end{openproblem}

\plainproblemsection

\begin{openproblem}
	*19.98. A connected graph $\sum$ is a symmetrical extension of a graph $\Gamma$ by a graph $\Delta$ if there exist a vertex-transitive group $G$ of automorphisms of $\sum$ and an imprimitivity system $\sigma$ of $G$ on the set of vertices of $\sum$ such that the quotient graph $\sum /\sigma$ is isomorphic to $\Gamma$ and blocks of $\sigma$ generate in $\sum$ subgraphs isomorphic to $\Delta$ .
	
	(a) Let $\Gamma$ be a locally finite Cayley graph of a finitely presented group, and $\Delta$ a finite graph. Are there only finitely many symmetrical extensions of $\Gamma$ by $\Delta$ ?
	
	(b) Let $\Gamma$ be a locally finite graph which has the property of $k$ -contractibility for some positive integer $k$ (see the definition in (V.I. Trofimov, Proc. Steklov Inst. Math.,279, suppl. 1 (2012),107-112); note that any $\Gamma$ from (a) is such a graph) and let $\Delta$ be a finite graph. Are there only finitely many symmetrical extensions of $\Gamma$ by $\Delta$ ? V. I. Trofimov
	
	*(a) No, not always, as follows from the construction in the proof of Theorem $\mathrm{H}$ in (M. de la Salle, R. Tessera, J. Topology, 12 (2019), 705-743).
	
	*(b) No, not always as follows from the construction in the proof of Theorem H in (M. de la Salle, R. Tessera, J. Topology, 12 (2019), 705-743).
\end{openproblem}

\plainproblemsection

\begin{openproblem}
	19.99. Is it true that for any positive integer $t$ there is a positive integer $n\left( t\right)$ such that any finite group with at least $n\left( t\right)$ conjugate classes of soluble maximal subgroups of Fitting height at most $t$ is itself a soluble group of Fitting height at most $t$ ?
	
	A. F. Vasil’ev, T. I. Vasil’eva
\end{openproblem}

\plainproblemsection

\begin{openproblem}
	*19.100. Suppose that a finite group $G$ admits a factorization $G = {AB} = {BC} = {CA}$ , where $A, B, C$ are abnormal supersoluble subgroups. Is $G$ supersoluble?
	
	A. F. Vasil’ev, T. I. Vasil’eva
	\par\smallskip\hrule\smallskip
	*Yes, it is (L. S. Kazarin, V. N. Tyutyanov, Preprint, 2023).
	\par\smallskip\hrule\smallskip
\end{openproblem}

\plainproblemsection

\begin{openproblem}
	*19.101. The maximum length of a chain of nested centralizers of a group is called its $c$ -dimension. Let $G$ be a locally finite group of finite $c$ -dimension $k$ , and let $S$ be the preimage in $G$ of the socle of $G/R$ , where $R$ is the locally solvable radical of $G$ . Is it true that the factor group $G/S$ contains an abelian subgroup of index bounded by a function of $k$ ? A. V. Vasil’ev
	
	*Yes, it is true (A. A. Buturlakin, J. Algebra Appl., 18, no. 12 (2019), 1950223, 12 pp).
\end{openproblem}

\plainproblemsection

\begin{openproblem}
	*19.102. A subgroup $H$ of a free group $F$ is called inert if $r\left( {H \cap  K}\right)  \leq  r\left( K\right)$ for every $K \leq  F$ ; and compressed if $r\left( H\right)  \leq  r\left( K\right)$ for every $H \leq  K \leq  F$ .
	
	Is it true that compressed subgroups are inert?
	
	*Yes, it is true (A. Jaikin-Zapirain, Preprint, 2024, https://arxiv.org/pdf/2403.09515).
\end{openproblem}

\plainproblemsection

\begin{openproblem}
	*19.103. Is it true that an intersection of compressed subgroups is compressed?
	
	It is known that arbitrary intersections of inert subgroups are inert. E. Ventur:
	
	*Yes, it is true (A. Jaikin-Zapirain, Preprint, 2024, https://arxiv.org/pdf/2403.09515).
\end{openproblem}

\plainproblemsection

\begin{openproblem}
	*19.104. Is there an algorithm which decides whether a given subgroup of $F$ is inert? An algorithm to decide whether a given $H$ is compressed is known.
	
	*Yes, there is (A. Jaikin-Zapirain, Preprint, 2024, https://arxiv.org/pdf/2403.09515).
\end{openproblem}

\plainproblemsection

\begin{openproblem}
	19.105. Is it true that fixed subgroups of endomorphisms of $F$ are inert?
	
	It is easy to see that retracts of free groups are compressed, and that they are inert if and only if fixed subgroups of endomorphisms of free groups are inert. Therefore the question is equivalent to the following: are retracts of free groups inert? E. Ventura
\end{openproblem}

\plainproblemsection

\begin{openproblem}
	19.106. Let $G$ be a uniformly locally finite group (which means that there is a function $f$ on the natural numbers such that any subgroup generated by $n$ elements has size at most $f\left( n\right)$ ). Suppose that for any two definable subgroups $H$ and $K$ , the intersection $H \cap  K$ has finite index either in $H$ or in $K$ . Is $G$ necessarily nilpotent-by-finite? Or even finite-by-abelian-by-finite? $F$ . Wagner
\end{openproblem}

\plainproblemsection

\begin{openproblem}
	19.107. Let $\mathcal{{AE}}$ be the smallest class of locally compact groups such that
	
	(1) $\mathcal{{AE}}$ contains the compact groups and the discrete amenable groups;
	
	(2) $\mathcal{{AE}}$ is closed under taking closed subgroups, Hausdorff quotients, directed unions of open subgroups, and group extensions.
	
	Is every amenable locally compact group an element of $\mathcal{{AE}}$ ?
	
	(A negative answer would follow from a positive answer to 19.21.) P. Wesolek
\end{openproblem}

\plainproblemsection

\begin{openproblem}
	19.108. Let $P$ be a finite $p$ -group, where $p$ is an odd prime. Let $\chi$ be a complex irreducible character of $P$ . If $\chi \left( x\right)  \neq  0$ for some $x \in  P$ , is it true that the order of $x$ must divide $\left| P\right| /\chi {\left( 1\right) }^{2}$ ? T. Wilde
\end{openproblem}

\plainproblemsection

\begin{openproblem}
	*19.109. A subgroup $H$ of a finite group $G$ is called pronormal if for any $g \in  G$ the subgroups $H$ and ${H}^{g}$ are conjugate by an element of $\left\langle  {H,{H}^{g}}\right\rangle$ . A maximal subgroup of a maximal subgroup is called second maximal. Is it true that in a non-abelian finite simple group $G$ all maximal subgroups are Hall subgroups if and only if every second maximal subgroup of $G$ is pronormal in $G$ ? V. I. Zenkov
	
	*No, not always, for example, in $S{L}_{2}\left( {2}^{11}\right)$ every second maximal subgroup is pronormal (V. N. Tyutyanov, Letter of 23 November 2018).
\end{openproblem}

\plainproblemsection

\begin{openproblem}
	19.110. (G. M. Bergman and A. Magidin). Do there exist varieties of groups in which the relatively free group of rank 2 is finite, and the relatively free group of rank 3 is infinite? P. Zusmanovich
\end{openproblem}

\plainproblemsection

\begin{openproblem}
	19.111. Do there exist infinite groups all of whose proper subgroups are cyclic of order $p$ that do not satisfy any nontrivial group identity except ${x}^{p} = 1$ and its consequences?
	
	Note that there exist infinite groups all of whose proper subgroups are infinite cyclic that do not satisfy any nontrivial group identity (due to A. Ol'shanskiǐ; see P. Zusmanovich, J. Algebra, 388 (2013), 268-286, Remark after Theorem 6.1).
\end{openproblem}

\plainproblemsection

\begin{openproblem}
	20.1. For $k \geq  1$ , a group $G$ is said to be totally $k$ -closed if in each of its faithful permutation representations, say on a set $\Omega , G$ is the largest subgroup of $\operatorname{Sym}\left( \Omega \right)$ which leaves invariant each of the $G$ -orbits in the induced action on the set of ordered $k$ -tuples ${\Omega }^{k}$ .
	
	Are there any finite insoluble totally 2-closed groups with nontrivial Fitting subgroup?
	
	The finite totally 2-closed groups which are either soluble or have trivial Fitting subgroup are known.
\end{openproblem}

\plainproblemsection

\begin{openproblem}
	20.2. Are there any nonabelian simple groups of Lie type which are totally 3-closed?
	
	There are exactly six nonabelian simple totally 2-closed groups - all are sporadic groups, the largest being the Monster.
	
	M. Arezoomand, M. A. Iranmanesh, C. E. Praeger, G. Tracey
\end{openproblem}

\plainproblemsection

\begin{openproblem}
	20.3. For each $k$ , each group of order $k$ is totally $k$ -closed, while the alternating group ${A}_{k}$ is not totally $k - 2$ -closed. Is there a fixed integer $k$ such that all finite simple groups which are not alternating groups are totally $k$ -closed?
	
	M. Arezoomand, M. A. Iranmanesh, C. E. Praeger, G. Tracey
\end{openproblem}

\plainproblemsection

\begin{openproblem}
	20.4. A finite group $G$ is said to be cut (or inverse semi-rational) if $\langle x\rangle  = \langle y\rangle$ implies that $x$ is conjugate to $y$ or to ${y}^{-1}$ for all $x, y \in  G$ .
	
	a) Let $\mathbb{Q}\left( G\right)$ denote the field extension of the rationals obtained by adjoining all entries of the ordinary character table of $G$ . Is there $c > 0$ such that $\left| {\mathbb{Q}\left( G\right)  : \mathbb{Q}}\right|  \leq  c$ for all cut groups? This is true if one assumes in addition that $G$ is solvable (J. F. Tent, J. Algebra, 363 (2012), 73-82).
	
	b) Is a Sylow 3-subgroup of a cut group also a cut group?
	
	c) Let ${\mathrm{O}}_{p}\left( G\right)$ denote the largest normal $p$ -subgroup of $G$ . Let $G$ be a solvable cut group. Is it true that for $p \in  \{ 5,7\}$ the exponent of ${\mathrm{O}}_{p}\left( G\right)$ divides $p$ ?
	
	For additional information and some positive results see (Adv. Group Theory Appl., 8B, 2020, 157-160 or https://arxiv.org/abs/2001.02637).A.Bächle
\end{openproblem}

\plainproblemsection

\begin{openproblem}
	20.5. (Well-known problem). Let $G$ be a finite group, and $\mathrm{V}\left( {\mathbb{Z}G}\right)$ the group of normalized units of the integral group ring of $G$ . Do the spectra of $G$ and $\mathrm{V}\left( {\mathbb{Z}G}\right)$ coincide? That is, is it true that, for any integer $n$ , there is an element of order $n$ in $\mathrm{V}\left( {\mathbb{Z}G}\right)$ if and only if there is an element of order $n$ in $G$ ?
	
	The answer is "yes" if $G$ is solvable (M. Hertweck, Comm. Algebra 36 (2008), 3585-3588).
\end{openproblem}

\plainproblemsection

\begin{openproblem}
	20.6. (W. Kimmerle). The prime graph (or Gruenberg-Kegel graph) $\Gamma \left( X\right)$ of a group $X$ has vertices labeled by primes appearing as orders of elements in $X$ ; two distinct primes $p$ and $q$ are adjacent in $\Gamma \left( X\right)$ if and only if $X$ contains an element of order ${pq}$ . Denote by $\mathrm{V}\left( {\mathbb{Z}G}\right)$ the group of normalized units of the integral group ring of a group $G$ . Is it true that for each finite group $G$ the prime graphs of $G$ and $\mathrm{V}\left( {\mathbb{Z}G}\right)$ coincide?
	
	This question has been reduced to almost simple groups (W. Kimmerle, A. Kono-valov, Internat. J. Algebra Comput., 27 (2017), 619-631).
\end{openproblem}

\plainproblemsection

\begin{openproblem}
	20.7. (W. Boone and G. Higman) Does every finitely generated group with solvable word problem embed into a finitely presented simple group? It is known that every such group embeds into a simple subgroup of a finitely presented group (W. Boone, G. Higman, J. Austral. Math. Soc., 18, no. 1 (1974), 41-53). J. Belk
\end{openproblem}

\plainproblemsection

\begin{openproblem}
	20.8. Suppose $K < H < F$ are free groups of finite rank such that $\operatorname{rank}\left( H\right)  <$ $\operatorname{rank}\left( K\right)$ , but all proper subgroups of $H$ which contain $K$ have ranks $\geq  \operatorname{rank}\left( K\right)$ . Then is the inclusion of $H$ in $F$ the only homomorphism $H \rightarrow  F$ fixing all elements
	
	of $K$ ? G. M. Bergman
\end{openproblem}

\plainproblemsection

\begin{openproblem}
	20.9. In the group algebra of a free group over a field, does every element whose support in the group has cardinality more than 1 generate a proper 2-sided ideal?
	
	This question is one of several related questions in (G. M. Bergman, Commun. Algebra, 49, no. 9 (2021), 3760-3776). G. M. Bergman
\end{openproblem}

\plainproblemsection

\begin{openproblem}
	20.10. (a) If $\mathcal{U}$ and ${\mathcal{U}}^{\prime }$ are nonprincipal ultrafilters on $\mathbb{N}$ , can every group which can be written as a homomorphic image of an ultraproduct of groups with respect to $\mathcal{U}$ also be written as a homomorphic image of an ultraproduct of groups with respect to ${\mathcal{U}}^{\prime }$ ?
	
	*(b) If the answer to (a) is negative, is it at least true that for any two nonprincipal ultrafilters $\mathcal{U}$ and ${\mathcal{U}}^{\prime }$ on $\mathbb{N}$ , there exists a nonprincipal ultrafilter ${\mathcal{U}}^{\prime \prime }$ on $\mathbb{N}$ such that every group which can be written as a homomorphic image of an ultraproduct of groups with respect to $\mathcal{U}$ or with respect to ${\mathcal{U}}^{\prime }$ can be written as a homomorphic image of an ultraproduct with respect to ${\mathcal{U}}^{\prime \prime }$ ?
	
	A positive answer to (b) would imply that the class of groups which can be written as homomorphic images of nonprincipal countable ultraproducts of groups is closed under finite direct products. See (G. M. Bergman, Pacific J. Math., 274 (2015), 451- 495).
	\par\smallskip\hrule\smallskip
	*(b) The affirmative answer is consistent with the ZFC axioms of set theory (S. M. Corson, Preprint, 2025, https://arxiv.org/pdf/2503.09228).
	\par\smallskip\hrule\smallskip
\end{openproblem}

\plainproblemsection

\begin{openproblem}
	20.11. Let $F \leq  H$ be free groups such that there exists a free group $G$ of finite rank with $F \leq  G \leq  H$ , and let $r$ be the least of the ranks of such groups $G$ . Which, if any, of the following statements must hold?
	
	(i) There is a largest $G$ of rank $r$ between $F$ and $H$ .
	
	$\left( {\mathrm{i}}^{\prime }\right)$ For any two ${G}_{1},{G}_{2}$ of rank $r$ between $F$ and $H$ , the subgroup $\left\langle  {{G}_{1},{G}_{2}}\right\rangle$ has rank $r$ .
	
	(ii) There is a smallest $G$ of rank $r$ between $F$ and $H$ .
	
	$\left( {\mathrm{{ii}}}^{\prime }\right)$ For any two ${G}_{1},{G}_{2}$ of rank $r$ between $F$ and $H$ , the subgroup ${G}_{1} \cap  {G}_{2}$ has rank $r$ .
	
	If (i) holds for all such $F$ and $H$ , then so does (i’). If (ii) holds for all $F, H$ , then so does (ii'). The converse of the former implication holds because subgroups of a free group of any fixed finite rank satisfy ACC, but I don't see a way to get the converse of the other implication. G. M. Bergman
\end{openproblem}

\plainproblemsection

\begin{openproblem}
	20.12. Let us say that a group $G$ has the unique $n$ -fold product property (u.- $n$ -p.) if for every $n$ -tuple of finite nonempty subsets ${A}_{1},\ldots ,{A}_{n} \subseteq  G$ there exists $g \in  G$ which can be written in one and only one way as $g = {a}_{1}\cdots {a}_{n}$ with ${a}_{i} \in  {A}_{i}$ . It is easy to see that u.- $n$ -p. implies u.-m-p. for $n \geq  m$ (since some of the ${A}_{i}$ can be $\{ 1\}$ ).
	
	Are the conditions u.-n-p. $\left( {n \geq  2}\right)$ all equivalent to u.-2-p., the usual unique product condition? G. M. Bergman
\end{openproblem}

\plainproblemsection

\begin{openproblem}
	*20.13. (a) If an abelian group can be written as a homomorphic image of a non-principal countable ultraproduct of not necessarily abelian groups ${G}_{i}$ , must it be a homomorphic image of a nonprincipal countable ultraproduct of abelian groups? See (G. M. Bergman, Pacific J. Math., 274 (2015), 451-495).
	
	(b) If an abelian group can be written as a homomorphic image of a direct product of an infinite family of not necessarily abelian finite groups, can it be written as a homomorphic image of a direct product of finite abelian groups?
	
	To see that neither question is trivial, choose for each $n > 0$ a finite group ${G}_{n}$ which is perfect but has elements which cannot be written as products of fewer than $n$ commutators. Then both the direct product of the ${G}_{n}$ and any nonprincipal ultra-product of those groups will have elements which are not products of commutators; hence its abelianization $A$ will be nontrivial. There is no evident family of abelian groups from which to obtain $A$ as an image of a nonprincipal ultraproduct, nor a family of finite abelian groups from which to obtain $A$ as an image of a direct product.
	\par\smallskip\hrule\smallskip
	*(a) Yes, it must (S. M. Corson, Preprint, 2025, https://arxiv.org/pdf/2503.09228).
	
	*(b) Yes it can (S.M. Corson, Preprint, 2025, https://arxiv.org/pdf/2503.09228).
	\par\smallskip\hrule\smallskip
\end{openproblem}

\plainproblemsection

\begin{openproblem}
	20.14. (a) Do there exist a variety $\mathfrak{V}$ of groups and a group $G \in  \mathfrak{V}$ such that the coproduct in $\mathfrak{V}$ of two copies of $G$ is embeddable in $G$ , but the coproduct of three such copies is not? See (G. M. Bergman, Indag. Math., 18 (2007), 349-403).
	
	Given an embedding $G{ * }_{\mathfrak{V}}G \rightarrow  G$ , one might expect the induced map
	
	$$
	G{ * }_{\mathfrak{V}}\left( {G{ * }_{\mathfrak{V}}G}\right)  \rightarrow  G{ * }_{\mathfrak{V}}G \rightarrow  G
	$$
	
	to be an embedding. But this is not automatic, because in a general group variety $\mathfrak{V}$ , a map $G{ * }_{\mathfrak{V}}A \rightarrow  G{ * }_{\mathfrak{V}}B$ induced by an embedding $A \rightarrow  B$ is not necessarily again an embedding.
	
	(b) If there exist $\mathfrak{V}$ and $G$ as in (a), does there in fact exist an example with $\mathfrak{V}$ the variety of groups generated by $G$ ? For this and related questions, see (G. M. Bergman, Algebra Number Theory, 3 (2009), 847-879). G. M. Bergman
\end{openproblem}

\plainproblemsection

\begin{openproblem}
	20.15. Let $\kappa$ be an infinite cardinal. If a residually finite group $G$ is embeddable in the full permutation group of a set of cardinality $\kappa$ , must it be embeddable in the direct product of $\kappa$ finite groups? See (G. M. Bergman, Indag. Math.,18 (2007), 349-403).
	
	The converse is true: any such direct product is residually finite and embeddable in the indicated permutation group. Also, the statement asked for becomes true if one replaces "direct product of $\kappa$ finite groups" by "direct product of ${2}^{\kappa }$ finite groups", since $G$ has cardinality at most ${2}^{\kappa }$ , hence homomorphisms to that many finite groups can be chosen which, together, separate each element of $G$ from $e$ . G. M. Bergman
\end{openproblem}

\plainproblemsection

\begin{openproblem}
	20.16. By Hartley's theorem (based on CFSG), if a simple locally finite group does not contain a particular finite group as a section, then it is isomorphic to a group of Lie type over a locally finite field.
	
	a) Is the same true for a definably simple locally finite group? A group is said to be definably simple if it does not contain proper definable normal subgroups.
	
	b) The same question for a definably simple locally finite group with bounded derived lengths of its solvable subgroups. A. V. Borovik
\end{openproblem}

\plainproblemsection

\begin{openproblem}
	20.17. The normal covering number $\gamma \left( G\right)$ of a finite non-cyclic group $G$ is the minimum number of proper subgroups of $G$ such that $G$ is the union of their conjugates.
	
	*(ae) What is the exact value of $\mathop{\liminf }\limits_{{n\text{ even }}}\gamma \left( {S}_{n}\right) /n$ ?
	
	(ao) What is the exact value of $\mathop{\liminf }\limits_{{n\text{ odd }}}\gamma \left( {S}_{n}\right) /n$ ?
	
	Comment of 2025: It belongs to $\left\lbrack  {1/{12},1/6}\right\rbrack$ for odd (S. Eberhard, C. Mellon, Preprint, 2024, https://arxiv.org/abs/2410.06999).
	
	(b) What are the exact values of $\mathop{\liminf }\limits_{{n\text{ even }}}\gamma \left( {A}_{n}\right) /n$ and $\mathop{\liminf }\limits_{{n\text{ odd }}}\gamma \left( {A}_{n}\right) /n$ ?
	
	Comment of 2025: These belong to $\left\lbrack  {1/{18},1/6}\right\rbrack$ for even, and $\left\lbrack  {1/6,4/{15}}\right\rbrack$ for odd (S. Eberhard, C. Mellon, Preprint, 2024, https://arxiv.org/abs/2410.06999).
	
	*(c) What is the exact value of $\mathop{\limsup }\limits_{{n\text{ even }}}\gamma \left( {S}_{n}\right) /n$ ?
	
	It is known that $\mathop{\limsup }\limits_{{n\text{ odd }}}\gamma \left( {S}_{n}\right) /n = 1/2$ .
	
	*(d) What are the exact values of $\mathop{\limsup }\limits_{{n\text{ odd }}}\gamma \left( {A}_{n}\right) /n$ and $\mathop{\limsup }\limits_{{n\text{ even }}}\gamma \left( {A}_{n}\right) /n$ ?
	
	D. Bubboloni, C. E. Praeger, P. Spiga
	
	*(ae) It is $1/6$ (S. Eberhard, C. Mellon, Preprint,2024, https://arxiv.org/abs/ 2410.06999).
	
	*(c) It is $1/4$ (S. Eberhard, C. Mellon, Preprint,2024, https://arxiv.org/abs/ 2410.06999).
	
	*(d) These are $1/3$ for odd, and $1/4$ for even (S. Eberhard, C. Mellon, Preprint, 2024, https://arxiv.org/abs/2410.06999).
\end{openproblem}

\plainproblemsection

\begin{openproblem}
	20.18. Let ${R}_{{p}^{k}}$ be the variety of class 2 nilpotent groups of exponent ${p}^{k}$ , where $p$ is a prime number and $k \geq  2$ . It is true that for every $p$ and $k$ there are infinitely many subquasivarieties of ${R}_{{p}^{k}}$ each of which is generated by a finite group with derived subgroup of exponent ${p}^{k}$ and does not have an independent basis of quasi-identities?
	
	A. I. Budkin
\end{openproblem}

\plainproblemsection

\begin{openproblem}
	20.19. A subgroup $H$ of a group $G$ is called commensurated if for all $g \in  G$ , the index $\left| {H : H \cap  {gH}{g}^{-1}}\right|$ is finite. Can a non-abelian free group (or a non-elementary hyperbolic group) contain two infinite commensurated subgroups $A, B$ with a trivial intersection? The answer is negative if $A$ or $B$ is normal. P.-E. Caprace 20.20. Let $E\left( {l, m, n}\right)$ be a group on generators $a, b$ with the presentation $\langle a, b \mid$ ${a}^{l} = 1,{\left( ab\right) }^{m} = {b}^{n}\rangle$ , where $l, m, n$ satisfy $1/l + 1/m + 1/n < 1$ . Are there any values for $l, m, n$ such that $E\left( {l, m, n}\right)$ is shortlex automatic?
	
	It is known that $E\left( {6,2,6}\right)$ is not shortlex automatic.C. Chalk
\end{openproblem}

\plainproblemsection

\begin{openproblem}
	20.21. (G. Verret). Does there exist a finite group $G$ with two normal subgroups $K$ and $L$ , each with index 12 in $G$ , such that $K$ is isomorphic to $L$ , but $G/K$ is isomorphic to ${C}_{12}$ , while $G/L$ is isomorphic to ${A}_{4}$ ? M. Conder
\end{openproblem}

\plainproblemsection

\begin{openproblem}
	20.22. Let $G$ be a non-free and non-cyclic one-relator group in which every subgroup of infinite index is free. Is $G$ a surface group? B. Fine, G. Rosenberger, L. Wienke
\end{openproblem}

\plainproblemsection

\begin{openproblem}
	*20.23. Is there a constant $\delta$ with $0 < \delta  < 1$ and an integer $N$ such that whenever $A, B$ , and $C$ are conjugacy classes in the alternating group ${Alt}\left( n\right)$ each of size at least ${\left| Alt\left( n\right) \right| }^{\delta }$ with $n \geq  N$ , then ${ABC} = {Alt}\left( n\right)$ ?
	
	*Yes, there are (D. Dona, Preprint, 2025, https://arxiv.org/abs/2505.06012).
\end{openproblem}

\plainproblemsection

\begin{openproblem}
	20.24. Does the group $S{L}_{2}\left( {\mathbb{Z}\left\lbrack  \sqrt{2}\right\rbrack  }\right)$ admit a faithful transitive amenable action?
	
	This is the same as asking if it admits any co-amenable subgroup except the finite index subgroups. Y. Glasner, N. Monod
\end{openproblem}

\plainproblemsection

\begin{openproblem}
	20.25. Let $n = {k}^{2}$ . Let $\sigma  \in  {S}_{n}$ be the product of $k$ disjoint cycles, of lengths $1,3,5,\ldots ,{2k} - 1$ . Find the limiting proportion of odd entries in the $\sigma$ -column of the character table of ${S}_{n}$ .
	
	For some background, see pages 1005-1006 in (D. Gluck, Proc. Amer. Math. Soc., 147 (2019), 1005-1011). D. Gluck
\end{openproblem}

\plainproblemsection

\begin{openproblem}
	20.26. Let $\phi$ be the Euler totient function. Does there exist a constant $a > 0$ such that $\left| {\operatorname{Aut}\left( G\right) }\right|  \geq  \phi {\left( \left| G\right| \right) }^{a}$
	
	a) for every finite group $G$ ?
	
	b) for every finite nilpotent group $G$ ?
	
	If such a constant $a$ exists, then $a \leq  \frac{40}{41}$ (J. González-Sánchez, A. Jaikin-Zapirain, Forum Math. Sigma, 3, Article ID e7, 11 p., electronic only (2015)). Also cf. 12.77.
	
	J. González-Sánchez, A. Jaikin-Zapirain
\end{openproblem}

\plainproblemsection

\begin{openproblem}
	*20.27. Let $G$ be a finite group, $p$ a prime number, and let ${\left| {x}^{G}\right| }_{p}$ denote the maximum power of $p$ that divides the class size of an element $x \in  G$ . Suppose that there exists a $p$ -element $g \in  G$ such that ${\left| {g}^{G}\right| }_{p} = \mathop{\max }\limits_{{x \in  G}}{\left| {x}^{G}\right| }_{p}$ . Is it true that $G$ has a normal $p$ -complement?
	
	A partial answer is in (https://arxiv.org/abs/1812.03641).I.B.Gorshkov
	
	*No, not necessarily, a counterexample is given by SmallGroup(192, 945) (B. Sam-bale, Letter of 16 February 2022).
\end{openproblem}

\plainproblemsection

\begin{openproblem}
	20.28. Let $L$ be a non-abelian finite simple group, and let $H\left( L\right)  = M\left( L\right) .L$ be the universal perfect central extension, where $M\left( L\right)$ is the Schur multiplier of $L$ . Suppose that $G$ is a finite group such that the set of class sizes of $G$ is the same as the set of class sizes of $H\left( L\right)$ . Is it true that $G \simeq  H\left( L\right)  \times  A$ , where $A$ is an abelian group?
	
	This is proved for $L = {Al}{t}_{5}$ in (J. Algebra Appl.,21, no. 11 (2022), Article ID 2250226, 8 p.). I. B. Gorshkov
\end{openproblem}

\plainproblemsection

\begin{openproblem}
	20.29. Let $S$ be a non-abelian finite simple group. Is it true that for any $n \in  \mathbb{N}$ , if the set of class sizes of a centreless finite group $G$ is the same as the set of class sizes of the direct power ${S}^{n}$ , then $G \simeq  {S}^{n}$ ?
	
	This is proved for $S = {Al}{t}_{5}^{2}$ in (J. Algebra Appl.,21, no. 11 (2022), Article ID 2250226, 8 p.). I. B. Gorshkov
\end{openproblem}

\plainproblemsection

\begin{openproblem}
	20.30. (P. M. Neumann and M. R. Vaughan-Lee). Let $G$ be a perfect and centreless finite group, and let $n$ be the maximum size of a conjugate class in $G$ . Is it true that $\left| G\right|  \leq  {n}^{2}?$
\end{openproblem}

\plainproblemsection

\begin{openproblem}
	20.31. Let $\omega \left( G\right)$ denote the set of element orders of a finite group $G$ , and $h\left( G\right)$ the number of pairwise nonisomorphic finite groups $H$ with $\omega \left( H\right)  = \omega \left( G\right)$ . Find $h\left( L\right)$ , where
	
	a) $L$ is the symmetric group of degree 10 ;
	
	b) $L$ is the automorphism group of the simple sporadic Janko group ${J}_{2}$ ;
	
	c) ${PSL}\left( {2, q}\right)  < L \leq  \operatorname{Aut}\left( {{PSL}\left( {2, q}\right) }\right)$ .
	
	See the current status in (M. A. Grechkoseeva, V. D. Mazurov, W. J. Shi, A. V. Vasil’ev, N. Yang, Commun. Math. Stat., 11, no. 2 (2023), 169-194; Subsection 4.2).
	
	M. A. Grechkoseeva, A. V. Vasil'ev
\end{openproblem}

\plainproblemsection

\begin{openproblem}
	20.32. (E. Rapaport Strasser). The Lovász conjecture states that a vertex-transitive connected graph is Hamiltonian. In the special case of Cayley graphs, one can ask the following. Consider a set $A$ which generates a group $G$ and is symmetric $\left( {x \in  A\text{implies}{x}^{-1} \in  A}\right)$ . Is there a list ${a}_{1},{a}_{2},\ldots ,{a}_{n}$ of elements of $A$ such that ${a}_{1},{a}_{1}{a}_{2},{a}_{1}{a}_{2}{a}_{3},\ldots ,{a}_{1}{a}_{2}\cdots {a}_{n}$ is a complete list of the elements of $G$ ? B. Green
\end{openproblem}

\plainproblemsection

\begin{openproblem}
	20.33. (A. Bauer). Does Higman's Embedding Theorem relativize in the following way? Is it the case that for every subset $X \subseteq  \mathbb{N}$ , there is a finitely generated group ${G}_{X}$ that has an $X$ -computable presentation (that is, there is a finite generating set relative to which the set of relations is computably enumerable with an $X$ oracle), and such that any finitely generated group has an $X$ -computable presentation if and only if it can be embedded as a finitely generated subgroup of a quotient of a free product of finitely many copies of ${G}_{X}$ by the normal closure of a finite subset?
	
	Higman’s Embedding Theorem says that for computable $X$ , one may take ${G}_{X} = \mathbb{Z}$ .
	
	J. Grochow
\end{openproblem}

\plainproblemsection

\begin{openproblem}
	20.34. Let $f$ be an inner screen of a saturated formation $\mathfrak{F}$ and suppose that a finite group $A$ acts faithfully and $f$ -hypercentrally on a finite group $G$ . Is it true that $G \rtimes  A \in  \mathfrak{F}$ for all $G \in  \mathfrak{F}$ if and only if $\mathfrak{F}$ is a Fitting class?W. Guo
\end{openproblem}

\plainproblemsection

\begin{openproblem}
	20.35. (Well-known problem). Let $S$ be a connected orientable hyperbolic surface of finite type and complexity at least 2. Does its mapping class group ${MCG}\left( S\right)$ have a non-elementary hyperbolic quotient?
	
	An affirmative answer is known when $S$ is a closed genus-2 surface (https:// arxiv.org/pdf/2005.00567.pdf), and in the same paper it is shown that the answer will be affirmative provided certain hyperbolic groups are residually finite. See also a closely related question 16.108 about braid groups. M. Hagen
\end{openproblem}

\plainproblemsection

\begin{openproblem}
	20.36. The free Burnside group of exponent four on three generators, $B\left( {3,4}\right)$ , has order ${2}^{69}$ as shown by Bayes, Kautsky, and Wamsley (1974). Their proof is based on a theorem of Sanov (1940) which shows that $B\left( {n,4}\right)$ is finite. Sanov’s proof for $B\left( {3,4}\right)$ uses more than ${2}^{32}$ fourth powers, because the subgroup of $B\left( {3,4}\right)$ generated by two of its generators and the square of the third has order ${2}^{32}$ . It is also known that $B\left( {3,4}\right)$ needs at least 105 relations to define it, as shown by Havas and Newman (1980).
	
	a) Can $B\left( {3,4}\right)$ be defined with fewer than a million fourth powers?
	
	b) Can $B\left( {3,4}\right)$ be defined with fewer than a thousand fourth powers?
	
	c) What is the smallest number of fourth powers which define $B\left( {3,4}\right)$ ?
\end{openproblem}

\plainproblemsection

\begin{openproblem}
	20.37. Is it true that for every finite group $G$ and every factorization $\left| G\right|  = {ab}$ there exist subsets $A, B \subseteq  G$ with $\left| A\right|  = a$ and $\left| B\right|  = b$ such that $G = {AB}$ ? Cf. 19.35.
\end{openproblem}

\plainproblemsection

\begin{openproblem}
	20.38. Suppose that $H$ is an almost simple group of Lie type and $G$ is a finite group such that $G$ and $H$ have the same sets of degrees of irreducible complex characters. Must there exist an abelian normal subgroup $A$ of $G$ such that $G/A$ is isomorphic to $H$ ? A. Iranmanesh
\end{openproblem}

\plainproblemsection

\begin{openproblem}
	20.39. Let $F$ be a non-abelian finitely generated free group, $1 \neq  w \in  F$ , and $n \geq  1$ . Is the group $\left\langle  {F, t \mid  {t}^{n} = w}\right\rangle$ linear of degree 2 over a field of characteristic 0 if $w$ is not a proper power in $F$ ?
	
	This question is motivated by the following well-known question: is it true that the free $\mathbb{Q}$ -group ${F}^{\mathbb{Q}}$ is linear over a field? (See 13.39(b).) A. Jaikin-Zapirain
\end{openproblem}

\plainproblemsection

\begin{openproblem}
	20.40. Let $G$ be a $\kappa$ -existentially closed group of cardinality $\lambda  > \kappa$ , where $\kappa$ is a regular cardinal. Is it true that $\left| {\operatorname{Aut}\left( G\right) }\right|  = {2}^{\lambda }$ ?
	
	The answer is known to be affirmative if $\lambda  = \kappa$ (B. Kaya, M. Kuzucuoğlu, P. Longobardi, M. Maj, Preprint, 2024, https://arxiv.org/pdf/2409.00545.
	
	B. Kaya, M. Kuzucuoğlu
\end{openproblem}

\plainproblemsection

\begin{openproblem}
	20.41. (L. Ciobanu, B. Fine, G. Rosenberger). Suppose that $G$ is a non-cyclic residually finite group in which every subgroup of finite index (including the group itself) is defined by a single defining relation, while all infinite index subgroups are free. Is it true that $G$ is either free or isomorphic to the fundamental group of a compact surface?
	
	See Archive 7.36 for a negative solution of a similar question without the assumption on infinite index subgroups. D. Kielak
\end{openproblem}

\plainproblemsection

\begin{openproblem}
	20.42. (Y. Cornullier, A. Mann, A. Thom). Does there exist a finitely generated residually finite group that is not (elementary) amenable and satisfies a nontrivial group law? S. Kionke
\end{openproblem}

\plainproblemsection

\begin{openproblem}
	20.43. Is every family of finite groups satisfying a common group law uniformly amenable? S. Kionke, E. Schesler
\end{openproblem}

\plainproblemsection

\begin{openproblem}
	20.44. The definition of $\mathrm{{CT}}\left( \mathbb{Z}\right)$ is given in 17.57. Is it true that a finitely generated subgroup of $\mathrm{{CT}}\left( \mathbb{Z}\right)$ either has only finitely many orbits on $\mathbb{Z}$ or there is a set of representatives for its orbits on $\mathbb{Z}$ which has positive density? S. Kohl
\end{openproblem}

\plainproblemsection

\begin{openproblem}
	20.45. Let $n$ be a positive integer, and let $G \leq  \mathrm{{GL}}\left( {n,\mathbb{Z}}\right)$ be finitely generated. Given a bound $b \in  \mathbb{N}$ , let ${e}_{b}$ be the number of elements of $G$ all of whose matrix entries have absolute value $\leq  b$ . Does the limit $\mathop{\lim }\limits_{{b \rightarrow  \infty }}\ln {e}_{b}/\ln b$ always exist?S. Kohl
\end{openproblem}

\plainproblemsection

\begin{openproblem}
	20.46. A group action on a compact space is said to be topologically free if the set of points with trivial stabilizer is dense. Let $G$ be a locally compact group, and ${\partial }_{sp}G$ its Furstenberg boundary (the largest minimal and strongly proximal compact $G$ -space). Let ${\Gamma }_{1}$ and ${\Gamma }_{2}$ be two lattices in $G$ , both acting faithfully on ${\partial }_{sp}G$ .
	
	a) Is it possible that the ${\Gamma }_{1}$ -action on ${\partial }_{sp}G$ is topologically free, but the ${\Gamma }_{2}$ -action on ${\partial }_{sp}G$ is not topologically free?
	
	b) If yes, can this also happen if ${\partial }_{sp}G = G/H$ is a homogeneous $G$ -space?
	
	A. Le Boudec
\end{openproblem}

\plainproblemsection

\begin{openproblem}
	20.47. Let $G = {F}_{n}$ be a finitely generated free group. Let $X$ be a compact $G$ -space on which the $G$ -action is faithful, minimal, and strongly proximal. Does it follow that the action is topologically free? A. Le Boudec, N. Matte Bon
\end{openproblem}

\plainproblemsection

\begin{openproblem}
	20.48. Suppose that $P$ is a finite $p$ -group with a non-trivial partition (which is equivalent to having proper Hughes subgroup ${H}_{p}\left( P\right)  \mathrel{\text{:=}} \left\langle  {g \in  P \mid  {g}^{p} \neq  1}\right\rangle   \neq  P)$ . If $P$ admits a fixed-point-free automorphism of prime order, must $P$ be of exponent $p$ ?
	
	M. Lewis
\end{openproblem}

\plainproblemsection

\begin{openproblem}
	20.49. Is it true that any finite group contains a 2-generated subgroup with the same exponent?
	
	An affirmative answer is known for soluble groups. It is also known that any finite group contains a 3-generated subgroup with the same exponent (E. Detomi, A. Lucchini, J. London Math. Soc. (2), 87, no. 3 (2013), 689-706). A. Lucchini
\end{openproblem}

\plainproblemsection

\begin{openproblem}
	20.50. For a positive integer $m$ , let ${G}_{m}$ be the largest group generated by $m$ involutions such that ${\left( xy\right) }^{4} = 1$ for any two involutions $x, y \in  {G}_{m}$ . What is the order of ${G}_{4}^{\prime }$ ?
	
	It is known that the order of ${G}_{3}$ is equal to ${2}^{11}$ . Also cf. 18.58. D. V. Lytkina
\end{openproblem}

\plainproblemsection

\begin{openproblem}
	20.51. Is it true that for every odd integer $t > 3$ there exists a finite non-abelian group $G$ of odd order with exactly $t$ conjugacy classes? D. MacHale
\end{openproblem}

\plainproblemsection

\begin{openproblem}
	20.52. A famous result of Burnside states that if $k\left( G\right)$ is the number of conjugacy classes of a finite group $G$ of odd order, then $\left| G\right|  - k\left( G\right)$ is divisible by 16 . Is it true that for every integer $m > 0$ there exists a finite non-abelian group $G\left( m\right)$ of odd order such that $\left| {G\left( m\right) }\right|  - k\left( {G\left( m\right) }\right)  = {16m}$ ? D. MacHale
\end{openproblem}

\plainproblemsection

\begin{openproblem}
	20.53. A nonisotropic unitary graph $\Gamma$ is distance-regular with intersection array $\{ q\left( {q - 1}\right) ,\left( {q + 1}\right) \left( {q - 2}\right) , q + 1;1,1, q\left( {q - 2}\right) \}$ for some prime power $q$ . The group $G = \operatorname{Aut}\left( \Gamma \right)$ acts transitively on the vertex set and on the edge set of $\Gamma$ . It is known that $\Gamma$ is distance-transitive if $q = 3$ . Does there exist a distance-regular graph with such an intersection array if $q$ is not a prime power? A. A. Makhnev
\end{openproblem}

\plainproblemsection

\begin{openproblem}
	20.54. A distance-regular graph of diameter 3 with the second eigenvalue ${\theta }_{1} = {a}_{3}$ is called a Shilla graph. For a Shilla graph $\Gamma$ the number $a = {a}_{3}$ divides $k$ and we set $b = b\left( \Gamma \right)  = k/a$ . Koolen and Park proved that there are 12 feasible intersection arrays of Shilla graphs with $b = 3$ . At present it is proved that a Shilla graph with $b = 3$ has intersection array $\{ {12},{10},3;1,3,8\}$ (Doro graph), $\{ {12},{10},5;1,1,8\}$ (nonisotropic unitary graph for $q = 4$ ), or $\{ {15},{12},6;1,2,{10}\}$ . The automorphisms of the last graph were found by A. Makhnev and N. Zyulyarkina (Doklady Maths., 84, no. 1 (2011), 510-514). Does the graph with intersection array $\{ {15},{12},6;1,2,{10}\}$ exist?
	
	A. A. Makhnev, D. O. Revin
\end{openproblem}

\plainproblemsection

\begin{openproblem}
	20.55. Are there soluble finite groups $G$ and $H$ , of derived lengths 2 and 4 and having identical character tables?
	
	Pairs of nonisomorphic soluble finite groups with identical character tables and with derived lengths $n$ and $n + 1$ for any $n \geq  2$ were constructed in (S. Mattarei, J. Algebra, 175 (1995) 157-178). S. Mattarei
\end{openproblem}

\plainproblemsection

\begin{openproblem}
	20.56. Is a periodic group $G$ locally finite if it is generated by involutions and the centralizer of every involution in $G$ is locally finite? V. D. Mazurov
\end{openproblem}

\plainproblemsection

\begin{openproblem}
	20.57. Is a periodic group $G$ locally finite if every finite subgroup of $G$ is contained in a subgroup of $G$ isomorphic to a finite alternating group? V. D. Mazurov
\end{openproblem}

\plainproblemsection

\begin{openproblem}
	20.58. Let $\omega \left( G\right)$ denote the set of element orders of a finite group $G$ . A finite group $G$ is said to be recognizable (by spectrum) if every finite group $H$ with $\omega \left( H\right)  = \omega \left( G\right)$ is isomorphic to $G$ .
	
	*a) Is it true that for every $n$ there is a recognizable group that is the $n$ -th direct power of a nonabelian simple group?
	
	b) Is it true that there is a nonabelian simple group $L$ such that for every $n$ there is a recognizable group whose socle is the $k$ -th direct power of $L$ for some $k \geq  n$ ?
	
	V. D. Mazurov, A. V. Vasil’ev
	
	*a) Yes, it is true (N. Yang, I. Gorshkov, A. Staroletov, A. V. Vasil'ev, Annali Matem. Pura Appl., 202 (2023), 2699-2714).
\end{openproblem}

\plainproblemsection

\begin{openproblem}
	20.59. A subgroup $H$ is called a virtual retract of a group $G$ if $H$ is a retract of a finite-index subgroup of $G$ . Is it true that every finitely generated subgroup of a finitely generated virtually free group is a virtual retract?
	
	For motivation and partial results see (A. Minasyan, Int. Math. Res. Notes, 2021, no. 17 (2021), 13434-13477). A. Minasyan
\end{openproblem}

\plainproblemsection

\begin{openproblem}
	20.60. We say that $G$ is a virtually compact special group if $G$ has a finite-index subgroup which is isomorphic to the fundamental group of a compact special complex (in the sense of F. Haglund, D. T. Wise, Geom. Funct. Anal., 17, no. 5 (2008), 1551- 1620). Let $G$ be a virtual retract of a finitely generated right-angled Artin group. Must $G$ be a virtually compact special group?
	
	An affirmative answer would provide an algebraic characterization of the class of virtually compact special groups as the class groups admitting finite index subgroups that are virtual retracts of right-angled Artin groups. A. Minasyan
\end{openproblem}

\plainproblemsection

\begin{openproblem}
	20.61. Suppose that $G$ is a virtually compact special group. Is it true that the centralizer of any element in $G$ is itself virtually compact special? A. Minasyan
\end{openproblem}

\plainproblemsection

\begin{openproblem}
	20.62. (J. Dixmier) A group is said to be unitarisable if its every uniformly bounded representation on a Hilbert space is unitarisable. Is every unitarisable group amenable? N. Monod
\end{openproblem}

\plainproblemsection

\begin{openproblem}
	20.63. a) Prove that every unitarisable group has trivial cost.
	
	This is known to be true for residually finite groups (I. Epstein, N. Monod, Int. Math. Res. Notes, 2009, no. 22 (2009), 4336-4353).
	
	b) At least prove that every unitarisable group has vanishing first ${L}^{2}$ Betti number.
	
	N. Monod
\end{openproblem}

\plainproblemsection

\begin{openproblem}
	20.64. Let $G$ be a non-amenable group.
	
	a) Prove that ${G}^{n}$ is non-unitarisable for some $n$ .
	
	b) At least prove that ${G}^{\infty }$ , the direct sum (restricted product), is non-unitarisable.
	
	N. Monod
\end{openproblem}

\plainproblemsection

\begin{openproblem}
	20.65. Given a complex irreducible character $\chi  \in  \operatorname{Irr}\left( \mathrm{G}\right)$ of a finite group $G$ , let $\mathbb{Q}\left( \chi \right)$ denote the field extension of $\mathbb{Q}$ obtained by adjoining to $\mathbb{Q}$ all the values of $\chi$ . We say that a finite group $G$ is $k$ -rational if $\left| {\mathbb{Q}\left( \chi \right)  : \mathbb{Q}}\right|$ divides $k$ for every $\chi  \in  \operatorname{Irr}\left( \mathrm{G}\right)$ . Does there exist a real-valued function $f$ such that if $p$ is the order of a cyclic composition factor of a $k$ -rational group $G$ , then $p \leq  f\left( k\right)$ ?
	
	If $k = 1$ , then we know that $p \leq  {11}$ by a theorem of J. G. Thompson (J. Algebra, 319 (2008), 558-594). A. Moretó
\end{openproblem}

\plainproblemsection

\begin{openproblem}
	20.66. A Schmidt(p, q)-group is a finite non-nilpotent group all of whose proper subgroups are nilpotent and whose Sylow $p$ -subgroup is normal. The $N$ -critical graph ${\Gamma }_{Nc}\left( G\right)$ of a finite group $G$ is a directed graph on the vertex set of all prime divisors of $\left| G\right|$ in which(p, q)is an edge of ${\Gamma }_{Nc}\left( G\right)$ if and only if $G$ has a Schmidt(p, q)-subgroup. Suppose that a finite group $G$ is such that $G = {AB} = {AC} = {BC}$ , where $A, B, C$ are subgroups of $G$ . Is ${\Gamma }_{Nc}\left( G\right)  = {\Gamma }_{Nc}\left( A\right)  \cup  {\Gamma }_{Nc}\left( B\right)  \cup  {\Gamma }_{Nc}\left( C\right)$ ? This is true if $A, B, C$ are soluble. V. I. Murashka, A. F. Vasil’ev
\end{openproblem}

\plainproblemsection

\begin{openproblem}
	20.67. Suppose that $G$ is a finite group, $p$ is a prime, and $B$ is a Brauer $p$ -block of $G$ with defect group $D$ . Let $\operatorname{cd}\left( B\right)$ be the set of degrees of the irreducible complex characters in $B$ .
	
	a) Is it true that the derived length of $D$ is bounded by $\left| {\operatorname{cd}\left( B\right) }\right|$ ?
	
	b) (A. Jaikin-Zapirain). Is is even true that $\left| {\operatorname{cd}\left( D\right) }\right|  \leq  \left| {\operatorname{cd}\left( B\right) }\right|$ ?G. Navarro
\end{openproblem}

\plainproblemsection

\begin{openproblem}
	20.68. The submonoid membership problem for a group $G$ generated by an alphabet $A$ asks, for given words ${x}_{1},{x}_{2},\ldots ,{x}_{n}$ and a word $w$ over $A$ , whether $w$ belongs to the submonoid of $G$ generated by the ${x}_{i}$ . Does there exist a hyperbolic one-relator group with undecidable submonoid membership problem? C.-F. Nyberg Brodda
\end{openproblem}

\plainproblemsection

\begin{openproblem}
	20.69. Is the submonoid membership problem decidable for every one-relator group with torsion?
\end{openproblem}

\plainproblemsection

\begin{openproblem}
	20.70. As a strengthening of the Burnside restriction, for every pair(k, n)of positive integers, let a group $G$ satisfy condition ${C}_{k, n}$ if every $k$ -generated subgroup of $G$ is finite of order at most $n$ .
	
	a) Does there exist $k \geq  2$ such that for any $n$ all groups with condition ${C}_{k, n}$ are locally finite?
	
	b) In particular, is it true that for any $n$ the condition ${C}_{2, n}$ implies local finiteness?
	
	c) Find possibly more pairs(k, n)for which groups with condition ${C}_{k, n}$ are locally finite. (For example, all groups with condition ${C}_{2,{20}}$ are metabelian, and therefore locally finite.) A. Yu. Ol'shanskii
\end{openproblem}

\plainproblemsection

\begin{openproblem}
	20.71. (J. Lauri). A card of a finite simple undirected graph $G$ of order $n = \left| {V\left( G\right) }\right|$ is an induced subgraph of order $n - 1$ . Let $k$ be2,3,4. For a connected graph $G$ with $k$ isomorphism types of cards, can $\operatorname{Aut}\left( G\right)$ have more than $k$ orbits on the vertex set $V\left( G\right)$ ?
	
	If a graph $G$ has $k$ isomorphism types of cards, then the group $\operatorname{Aut}\left( G\right)$ of auto-morphisms of $G$ has obviously at least $k$ orbits on $V\left( G\right)$ . It is known that if all cards are mutually isomorphic, then $\operatorname{Aut}\left( G\right)$ is transitive on $V\left( G\right)$ . Examples are known of graphs $G$ with 5 isomorphism types of cards for which $\operatorname{Aut}\left( G\right)$ has 6 orbits on $V\left( G\right)$ , and of graphs with 6 isomorphism types for which $\operatorname{Aut}\left( G\right)$ has 7 orbits on $V\left( G\right)$ .
	
	V. Pannone
\end{openproblem}

\plainproblemsection

\begin{openproblem}
	20.72. Is there a variety of groups $\Theta$ such that the group of automorphisms of the category of free finitely generated groups ${\Theta }^{0}$ contains outer automorphisms?
	
	E. Plotkin
\end{openproblem}

\plainproblemsection

\begin{openproblem}
	20.73. Can every countable group $G$ be factorized $G = {AB}$ into infinite subsets $A, B$ such that every element $g \in  G$ has a unique representation $g = {ab}$ for $a \in  A, b \in  B$ ?
	
	This is true if $G$ is topologizable. I. V. Protasov
\end{openproblem}

\plainproblemsection

\begin{openproblem}
	20.74. a) Conjecture: there are only finitely many nonabelian finite simple groups $G$ that have a normal subset $S$ closed under inversion such that $\left| S\right|  > \left| G\right| /{\log }_{2}\left| G\right|$ and ${S}^{2} \neq  G$ . Comment of 2024: This is true for the class of alternating groups (M. Larsen, P. H. Tiep, Preprint, 2023, arXiv:2305.04806) and for any class of groups of Lie type of bounded Lie rank (S. V. Skresanov, Preprint, 2024, https://arxiv.org/abs/ 2406.12506).
	
	*b) The same conjecture for the groups $\operatorname{PSL}\left( {2, p}\right)$ .L. Pyber
	
	*b) The conjecture is proved for the groups $\operatorname{PSL}\left( {2, p}\right)$ (S. V. Skresanov, Preprint, 2024, https://arxiv.org/abs/2406.12506).
\end{openproblem}

\plainproblemsection

\begin{openproblem}
	20.75. Let $G$ be a finite $p$ -group and assume that all abelian normal subgroups of $G$ can be generated by $k$ elements. Is it true that every abelian subgroup of $G$ can be generated by ${2k}$ elements?
	
	The $k$ -th direct power of ${D}_{16}$ shows that this bound would be best possible.
	
	L. Pyber
\end{openproblem}

\plainproblemsection

\begin{openproblem}
	20.76. Let $G$ be a finite $p$ -group and assume that all abelian normal subgroups of $G$ have order at most ${p}^{k}$ . Is it true that every abelian subgroup of $G$ has order at most ${p}^{2k}$ ? L. Pyber
\end{openproblem}

\plainproblemsection

\begin{openproblem}
	20.77. (O.I. Tavgen'). A group is said to be boundedly generated if it is a product of finitely many cyclic groups. Is it true that every boundedly generated residually finite group is linear over a field of characteristic zero?
	
	This is true for residually finite-soluble groups (L. Pyber, D. Segal, J. Reine Angew. Math., 612 (2007), 173-211). L. Pyber
\end{openproblem}

\plainproblemsection

\begin{openproblem}
	20.78. For an irreducible complex character $\chi$ of a finite group $G$ , the codegree of $\chi$ is defined by $\operatorname{cod}\left( \chi \right)  = \left| {G : \ker \chi }\right| /\chi \left( 1\right)$ . Let $\operatorname{Cod}\left( G\right)$ be the set of irreducible character codegrees of $G$ .
	
	Conjecture: If $G$ has an element of order $m$ , then $m$ divides some member of $\operatorname{Cod}\left( G\right)$ .
	
	The conjecture is proved when $m$ is a prime power (G. Qian, Arch. Math.,97 (2011), 99-103); when $m$ is square-free (I. M. Isaacs, Arch. Math., 97 (2011), 499- 501); or when $G$ is solvable (G. Qian, Bull. London Math. Soc., $\mathbf{{53}}\left( {2021}\right) {820} - {824}$ ).
	
	G. Qian
\end{openproblem}

\plainproblemsection

\begin{openproblem}
	20.79. Conjecture: Suppose that $G$ is a non-abelian simple group and $H$ is a finite group such that $\operatorname{Cod}\left( G\right)  = \operatorname{Cod}\left( H\right)$ . Then $G \cong  H$ .G. Qian
\end{openproblem}

\plainproblemsection

\begin{openproblem}
	20.80. Let $G$ be a finite group, and $\chi$ an irreducible complex character of $G$ . We call $\chi$ a $\mathcal{P}$ -character if $\chi$ is a constituent of ${\left( {1}_{H}\right) }^{G}$ for some maximal subgroup $H$ of $G$ ; and $\chi$ is said to be monomial if $\chi$ is induced by a linear character of a subgroup of $G$ .
	
	Conjecture: $G$ is solvable if and only if all its $\mathcal{P}$ -characters are monomial.
	
	The necessity part of the conjecture is true (G. Qian, Y. Yang, Commun. Algebra, 46, no. 1 (2018), 167-175).G. Qian
\end{openproblem}

\plainproblemsection

\begin{openproblem}
	20.81. Does there exist a class $\mathfrak{X}$ of finite groups satisfying the following conditions: (1) $\mathfrak{X}$ is closed with respect to taking subgroups and homomorphic images;
	
	(2) the product of two normal $\mathfrak{X}$ -subgroups of an arbitrary group is always an $\mathfrak{X}$ -group;
	
	(3) for every positive integer $n$ , there exist a finite group and its conjugacy class $D$ such that any $n$ elements of $D$ generate an $\mathfrak{X}$ -subgroup, whereas $\langle D\rangle  \notin  \mathfrak{X}$ .
	
	One can prove that there are no classes $\mathfrak{X}$ satisfying (1)-(3) that are closed with respect to extensions. D. O. Revin
\end{openproblem}

\plainproblemsection

\begin{openproblem}
	20.82. Two groups are said to be isospectral if they have the same set of element orders. Suppose that $G$ is a finite group such that every finite group isospectral to $G$ is isomorphic to $G$ . Is it true that the quotient of $G$ by its socle is solvable?
\end{openproblem}

\plainproblemsection

\begin{openproblem}
	20.83. Let $\mathfrak{X}$ be a class of finite groups that is closed with respect to taking subgroups, homomorphic images, and extensions. A subgroup $H$ of a finite group $G$ is said to be $\mathfrak{X}$ -submaximal if there exists an embedding of $G$ into a group ${G}^{ * }$ such that $G$ is subnormal in ${G}^{ * }$ and $H$ coincides with the intersection of $G$ and an $\mathfrak{X}$ -maximal subgroup of ${G}^{ * }$ . Suppose that all $\mathfrak{X}$ -submaximal subgroups of a characteristic subgroup $N$ of a finite group $G$ are conjugate in $N$ . Does it follow that ${HN}/N$ is an $\mathfrak{X}$ -submaximal subgroup of $G/N$ for every $\mathfrak{X}$ -submaximal subgroup $H$ of $G$ ?
	
	For normal subgroups $N$ , this is not true even if $N$ is an $\mathfrak{X}$ -group or if $N$ does not contain nontrivial $\mathfrak{X}$ -subgroups. A positive answer is known in the case where $N$ coincides with the $\mathfrak{F}$ -radical of $G$ for a Fitting class $\mathfrak{F}$ .
	
	D. O. Revin, A. V. Zavarnitsine
\end{openproblem}

\plainproblemsection

\begin{openproblem}
	20.84. Does there exist a nilpotent group of class 3 with a non-trivial 5-th dimension subgroup? E. Rips
\end{openproblem}

\plainproblemsection

\begin{openproblem}
	20.85. Let $F$ be a free group. An element $\omega  \in  F$ is said to be primitive if there is a minimal generating system of $F$ that contains $\omega$ , almost primitive if it is primitive in each finitely generated proper subgroup of $F$ containing $\omega$ , tame almost primitive if, whenever ${\omega }^{\alpha }$ is contained in a subgroup $H$ of $F$ with $\alpha  \geq  1$ minimal, either ${\omega }^{\alpha }$ is primitive in $H$ or the index of $H$ in $F$ is just $\alpha$ . In (B.Fine, A. Moldenhauer, G. Rosenberger, L. Wienke, Topics in Infinite Group Theory: Nielsen Methods, Covering Spaces, and Hyperbolic Groups, De Gruyter, Berlin, 2021) it is shown that $u = \left\lbrack  {{a}_{1},{b}_{1}}\right\rbrack  \left\lbrack  {{a}_{2},{b}_{2}}\right\rbrack  \cdots \left\lbrack  {{a}_{g},{b}_{g}}\right\rbrack$ is tame almost primitive in the free group on ${a}_{1},{b}_{1},\ldots ,{a}_{g},{b}_{g}$ with $g \geq  1$ , and $v = {c}_{1}^{2}\cdots {c}_{p}^{2}$ is tame almost primitive in the free group on ${c}_{1},\ldots ,{c}_{p}$ with $p \geq  2$ .
	
	Are there tame almost primitive elements in free groups other than $u, v$ , and their product ${uv}$ in the free group on ${a}_{1},{b}_{1},\ldots ,{a}_{g},{b}_{g},{c}_{1},\ldots ,{c}_{p}$ ?
	
	G. Rosenberger, L. Wienke
\end{openproblem}

\plainproblemsection

\begin{openproblem}
	20.86. Suppose $G$ is a finite group and $p$ a prime such that the number ${s}_{p}\left( G\right)  = 1 + {kp}$ of Sylow $p$ -subgroups $P$ of $G$ is greater than 1 . By a theorem of Frobenius the set ${G}_{p}$ of all $p$ -elements in $G$ has cardinality $\left| {G}_{p}\right|  = \left| P\right|  \cdot  {f}_{p}\left( G\right)$ for some positive integer ${f}_{p}\left( G\right)$ , and one easily gets that ${f}_{p}\left( G\right)  = 1 + \ell \left( {p - 1}\right)  \leq  k - \frac{k - 1}{p - 1}$ . Usually, $\ell  < k$ (but $\ell  = k$ if $k < p$ ). Is always ${\ell }^{p} \geq  {k}^{p - 1}$ ?
	
	Recently P. Gheri showed that ${f}_{p}{\left( G\right) }^{p} \geq  {s}_{p}{\left( G\right) }^{p - 1}$ if $G$ is $p$ -solvable (Ann. Mat. Pura Appl. (4), 200 (2021), 1231-1243). P. Schmid
\end{openproblem}

\plainproblemsection

\begin{openproblem}
	20.87. What are the non-abelian composition factors of finite groups in which the order of every element is divisible by at most two primes?
	
	For the case of one prime, see (https://arxiv.org/pdf/2003.09445.pdf).
	
	W. J. Shi
\end{openproblem}

\plainproblemsection

\begin{openproblem}
	20.88. Let $u$ be an element of a free group ${F}_{r}$ . Is it true that there is $v \in  {F}_{r}$ (that depends on $u$ ) that cannot be a subword of any cyclically reduced word $\varphi \left( u\right)$ , where $\varphi$ is an automorphism of ${F}_{r}$ ? V. Shpilrain
\end{openproblem}

\plainproblemsection

\begin{openproblem}
	20.89. An element $g$ of a group $G$ is said to be almost Engel if there is a finite subset $E\left( g\right)$ of $G$ such that for every $x \in  G$ all sufficiently long commutators $\left\lbrack  {\ldots \left\lbrack  {\left\lbrack  {x, g}\right\rbrack  , g}\right\rbrack  ,\ldots , g}\right\rbrack$ belong to $E\left( g\right)$ , that is, there is a positive integer $n\left( {x, g}\right)$ such that $\left\lbrack  {\ldots \left\lbrack  {\left\lbrack  {x, g}\right\rbrack  , g}\right\rbrack  ,\ldots , g}\right\rbrack   \in  E\left( g\right)$ if $g$ is repeated $\geq  n\left( {x, g}\right)$ times. An element $g$ is Engel if we can take $E\left( g\right)  = \{ 1\}$ . The set of Engel elements of a linear group is a subgroup by a well-known result of Gruenberg (J. Algebra, 3 (1966), 291-303). Is the set of almost Engel elements of a linear group a subgroup?
	
	A linear group in which all elements are almost Engel is finite-by-hypercentral (P. Shumyatsky, Monatsh. Math., 186 (2018), 711-719). P. Shumyatsky
\end{openproblem}

\plainproblemsection

\begin{openproblem}
	20.90. A profinite group in which all centralizers of non-trivial elements are pronilpo-tent is called a CN-group. Find an example of a finitely generated infinite profinite CN-group which is not prosoluble.
	
	The structure of profinite CN-groups is described in (Israel J. Math., 235, no. 1 (2020), 325-347). P. Shumyatsky
\end{openproblem}

\plainproblemsection

\begin{openproblem}
	20.91. A group $G$ is said to be stable if for any first-order formula $\phi \left( {\bar{x},\bar{y}}\right)$ (in the first-order language of groups $\left\{  {\cdot ,{}^{-1},1}\right\}$ there exists a natural number $n$ such that whenever there exist sequences of tuples of $G,{\left( {a}_{i}\right) }_{i < m},{\left( {b}_{i}\right) }_{i < m}$ with $G \vDash  \phi \left( {{a}_{i},{b}_{j}}\right)$ if and only if $i < j$ , then $m \leq  n$ . Sela proved that all torsion-free hyperbolic groups are stable.
	
	a) Is there a non-stable hyperbolic group?
	
	b) Is there a non-stable virtually free group?
	
	c) Is $S{L}_{2}\left( \mathbb{Z}\right)$ non-stable? R. Sklinos
\end{openproblem}

\plainproblemsection

\begin{openproblem}
	20.92. This question is about existence of an analogue of the Lazard correspondence for pre-Lie algebras and braces. A pre-Lie algebra $A$ is a vector space with a bilinear operation $\left( {x, y}\right)  \rightarrow  {xy}$ satisfying $\left( {xy}\right) z - x\left( {yz}\right)  = \left( {yx}\right) z - y\left( {xz}\right)$ for every $x, y, z \in  A$ . A pre-Lie algebra $A$ is said to be left nilpotent if, for some $n \in  \mathbb{N}, A \cdot  \left( {A \cdot  \left( {A\cdots A}\right) }\right) ) = 0$ where $A$ appears $n$ times in the product. Recall that a set $A$ with binary operations + and $\circ$ is a left brace if $\left( {A, + }\right)$ is an abelian group, $\left( {A, \circ  }\right)$ is a group, and $a \circ  \left( {b + c}\right)  + a =$ $a \circ  b + a \circ  c$ for every $a, b, c \in  A$ . Let $p$ be a prime, and ${\mathbb{F}}_{p}$ the field of $p$ elements. A left brace $A$ is called an ${\mathbb{F}}_{p}$ -brace if its additive group is an ${\mathbb{F}}_{p}$ -vector space such that $a * \left( {\alpha b}\right)  = \alpha \left( {a * b}\right)$ for all $a, b \in  A,\alpha  \in  {\mathbb{F}}_{p}$ , where $a * b = a \circ  b - a - b$ . The idea of a connection between braces and pre-Lie algebras comes from a paper by W. Rump (2014).
	
	a) Let $A$ be an ${\mathbb{F}}_{p}$ -brace of cardinality ${p}^{k}$ for some $k$ . Is it true that when $p$ is sufficiently large relative to $k$ , the set $A$ with the same additive operation + and with the operation - defined as $a \cdot  b =  - \mathop{\sum }\limits_{{i = 0}}^{{p - 2}}\frac{1}{{2}^{i}}\left( {\left( {{2}^{i}a}\right)  * b}\right)$ is a pre-Lie algebra?
	
	b) Let $k$ be a natural number, and let $p$ be a prime number such that $p > {2}^{k}$ . Is there a bijective correspondence between ${\mathbb{F}}_{p}$ -braces of cardinality ${p}^{k}$ and left nilpotent pre-Lie algebras over ${\mathbb{F}}_{p}$ of cardinality ${p}^{k}$ ?
	
	An affirmative answer to any of the above questions would have consequences for the theory of set-theoretic solutions of the Yang-Baxter equation and for the theory of Hopf-Galois extensions. A. Smoktunowicz
\end{openproblem}

\plainproblemsection

\begin{openproblem}
	20.93. A group $G$ is called a Shunkov group if for any finite subgroup $H \leq  G$ any two conjugate elements of prime order in ${N}_{G}\left( H\right) /H$ generate a finite subgroup. Are the following well-known results of the theory of finite groups true in the class of (periodic) Shunkov groups?
	
	a) The Baer-Suzuki theorem (see 11.11).
	
	b) The Burnside-Brauer-Suzuki theorem on the existence of a normal section of order 2 in a group with a non-trivial Sylow 2-subgroup containing only one involution (see 4.75).
	
	c) Glauberman’s ${Z}^{ * }$ -theorem (see 10.62 and 11.13). A. I. Sozutov
\end{openproblem}

\plainproblemsection

\begin{openproblem}
	20.94. Does there exist an infinite periodic simple group saturated (see the definition in 14.101) with finite Frobenius groups? A. I. Sozutov
\end{openproblem}

\plainproblemsection

\begin{openproblem}
	20.95. Is a periodic group a Frobenius group (see 6.53) if it is saturated with finite Frobenius groups, contains an involution, and does not contain non-cyclic subgroups of order 4 ? A. I. Sozutov
\end{openproblem}

\plainproblemsection

\begin{openproblem}
	*20.96. Is a periodic group a Frobenius group if it has a proper non-trivial normal abelian subgroup that contains the centralizer of each of its non-identity elements? A. I. Sozutov
	
	*Yes, it is (D. V. Lytkina, V. D. Mazurov, Siberian Math. J., 64, no. 6 (2023), 1350- 1353).
\end{openproblem}

\plainproblemsection

\begin{openproblem}
	20.97. Is a 2-group locally finite if the centralizer of every involution is locally finite? N. M. Suchkov
\end{openproblem}

\plainproblemsection

\begin{openproblem}
	*20.98. Let $G$ be a group of permutations of the set of positive integers $\mathbb{N}$ isomorphic to the additive group of rational numbers. Must there be an element $g \in  G$ such that the set $\{ a - {ag} \mid  a \in  \mathbb{N}\}$ is infinite? N. M. Suchkov
	
	*Yes, it must (N. M. Suchkov, A. A. Shlepkin, D. A. Taysnyov, Siberian Math. J., 65 (2024), 1390-1394).
\end{openproblem}

\plainproblemsection

\begin{openproblem}
	20.99. Conjecture: Let ${a}_{1}{G}_{1},\ldots ,{a}_{k}{G}_{k}, k > 1$ , be finitely many pairwise disjoint left cosets in a group $G$ with $\left\lbrack  {G : {G}_{i}}\right\rbrack   < \infty$ for all $i = 1,\ldots , k$ . Then
	
	$$
	\gcd \left( {\left\lbrack  {G : {G}_{i}}\right\rbrack  ,\left\lbrack  {G : {G}_{j}}\right\rbrack  }\right)  \geq  k\;\text{ for some }1 \leq  i < j \leq  k.
	$$
	
	The conjecture is known to hold for $k = 2,3,4$ .Z.-W. Sun
\end{openproblem}

\plainproblemsection

\begin{openproblem}
	20.100. Conjecture: Let $n$ be a positive integer, and let $G$ be a group containing no elements of order among $2,\ldots , n + 1$ . Then, for any $A \subseteq  G$ with $\left| A\right|  = n$ , we may write $A = \left\{  {{a}_{1},\ldots ,{a}_{n}}\right\}$ with ${a}_{1},{a}_{2}^{2},\ldots ,{a}_{n}^{n}$ pairwise distinct.
	
	The conjecture is known to hold for $n \leq  3$ .Z.-W. Sun
\end{openproblem}

\plainproblemsection

\begin{openproblem}
	20.101. Let $G$ be a finitely generated branch group. Are all finite-index maximal subgroups of $G$ necessarily normal?
	
	Cf. 18.81. A. Thillaisundaram
\end{openproblem}

\plainproblemsection

\begin{openproblem}
	20.102. (Y. Barnea, A. Shalev). Let $G$ be a finitely generated pro- $p$ group. Must $G$ be $p$ -adic analytic if it has finite Hausdorff spectrum with respect to
	
	a) the $p$ -power series?
	
	b) the iterated $p$ -power series?
	
	c) the lower $p$ -series?
	
	d) the Frattini series?
	
	e) the dimension subgroup series? A. Thillaisundaram
\end{openproblem}

\plainproblemsection

\begin{openproblem}
	20.103. Let ${G}_{\Gamma }$ be a partially commutative soluble group of derived length $n \geq  3$ with defining graph $\Gamma$ (the definition is similar to the case of $n = 2$ , see 17.104). Is ${G}_{\Gamma }$ a torsion-free group? E. I. Timoshenko
\end{openproblem}

\plainproblemsection

\begin{openproblem}
	20.104. (Well-known questions). Suppose that a group $G$ is finitely generated and decomposable into a direct product $G = {G}_{1} \times  {G}_{2}$ .
	
	a) Is it true that the elementary theory of $G$ is decidable if and only if the elementary theories of the groups ${G}_{1}$ and ${G}_{2}$ are decidable?
	
	b) Is it true that the universal theory of $G$ is decidable if and only if the universal theories of the groups ${G}_{1}$ and ${G}_{2}$ are decidable?
	
	An affirmative answer to question a) for almost soluble groups follows from (G. A. Noskov, Izv. Akad. Nauk SSSR, Ser. Mat., 47, no. 3 (1983), 498-517).
	
	E. I. Timoshenko
\end{openproblem}

\plainproblemsection

\begin{openproblem}
	20.105. Is there a perfect locally nilpotent $p$ -group, for some prime $p$ , whose proper subgroups are hypercentral? N. Trabelsi
\end{openproblem}

\plainproblemsection

\begin{openproblem}
	20.106. Let $G$ be a residually finite 2 -group and let $x \in  G$ be a left 3-Engel element of order 2 . Is $\left\langle  {x}^{G}\right\rangle$ locally nilpotent?
	
	It is known that in any group a 3-Engel element of odd order belongs to the locally nilpotent radical (E. Jabara, G. Traustason, Proc. Amer. Math. Soc., 147, no. 5 (2019), 1921-1927). G. Traustason
\end{openproblem}

\plainproblemsection

\begin{openproblem}
	20.107. Let $G$ be a group of exponent 8 and $x \in  G$ a left 3-Engel element of order 2. Is $\left\langle  {x}^{G}\right\rangle$ locally finite? G. Traustason
\end{openproblem}

\plainproblemsection

\begin{openproblem}
	20.108. The holomorph $\operatorname{Hol}\left( G\right)$ of a group $G$ can be defined as the normalizer of the subgroup of left translations in the group of all permutations of the set $G$ . The multiple holomorph $\operatorname{NHol}\left( G\right)$ of $G$ is the normalizer of the holomorph. Set $T\left( G\right)  =$ $\mathrm{{NHol}}\left( G\right) /\mathrm{{Hol}}\left( G\right)$ .
	
	a) Is there a centerless group $G$ for which $T\left( G\right)$ is not an elementary abelian 2-group?
	
	b) Is there a finite centerless group $G$ for which $T\left( G\right)$ is not an elementary abelian 2-group?
	
	c) (A. Caranti). Is there a finite $p$ -group $G$ for which the order of $T\left( G\right)$ has a prime divisor not dividing $\left( {p - 1}\right) p$ ? C. Tsang
\end{openproblem}

\plainproblemsection

\begin{openproblem}
	20.109. A skew brace is a set $B$ equipped with two operations + and - such that $\left( {B, + }\right)$ is an additively written (but not necessarily abelian) group, $\left( {B, \cdot  }\right)$ is a multiplicatively written group, and $a \cdot  \left( {b + c}\right)  = {ab} - a + {ac}$ for any $a, b, c \in  B$ .
	
	Is there a finite skew brace with perfect additive group and almost simple multiplicative group? C. Tsang
\end{openproblem}

\plainproblemsection

\begin{openproblem}
	20.110. Are there residually finite hereditarily just infinite groups that are
	
	a) amenable but not solvable?
	
	b) amenable but not elementary amenable?
	
	c) of intermediate word growth?
	
	d) of intermediate subgroup growth?
	
	All examples that we know are either linear (hence Tits Alternative applies), or have a quotient with property (T) (hence cannot be amenable). M. Vannacci
\end{openproblem}

\plainproblemsection

\begin{openproblem}
	20.111. Are there residually finite hereditarily just infinite groups admitting a self-similar action on a rooted tree that are not linear? M. Vannacci
\end{openproblem}

\plainproblemsection

\begin{openproblem}
	20.112. Let $\mathfrak{F}$ be a hereditary saturated formation and let ${\mathrm{w}}^{ * }\mathfrak{F}$ denote the class of all finite groups $G$ for which $\pi \left( G\right)  \subseteq  \pi \left( \mathfrak{F}\right)$ and the normalizers of all Sylow subgroups of $G$ are $\mathfrak{F}$ -subnormal in $G$ . It is known that ${\mathrm{w}}^{ * }\mathfrak{F}$ is a formation. Must ${\mathrm{w}}^{ * }\mathfrak{F}$ be a saturated formation? A. F. Vasil’ev, T. I. Vasil’eva
\end{openproblem}

\plainproblemsection

\begin{openproblem}
	20.113. Let $H\left( {q, c}\right)  = \langle a, b, c, d \mid  \left\lbrack  {b, a}\right\rbrack   = \left\lbrack  {d, c}\right\rbrack$ , of exponent $q$ , nilpotent of class $c\rangle$ .
	
	a) Is it true that the Schur multiplier $M\left( {H\left( {8,{12}}\right) }\right)$ has exponent 32?
	
	b) Is it true that the Schur multiplier $M\left( {H\left( {7,{13}}\right) }\right)$ has exponent 49?
	
	The difficulty is that these groups are too big to compute using current versions of the $p$ -Quotient Algorithm, which use 32 bit arithmetic. So to tackle these groups it would help to have a version of the $p$ -Quotient Algorithm using 64 bit arithmetic.
	
	M. R. Vaughan-Lee
\end{openproblem}

\plainproblemsection

\begin{openproblem}
	20.114. Is there a positive integer $c$ such that for any $q$ the exponent of the Schur multiplier of a finite group of exponent $q$ divides ${q}^{c}$ ? M. R. Vaughan-Lee
\end{openproblem}

\plainproblemsection

\begin{openproblem}
	20.115. Let $\chi$ be a complex irreducible character of a finite group $G$ . If $\chi \left( x\right)  \neq  0$ for some $x \in  G$ , must the order $o\left( x\right)$ of $x$ divide $\left| G\right| /\chi \left( 1\right)$ ?
	
	This is known to be true if $G$ is solvable, and it is known that ${\left( o\left( x\right) \chi \left( 1\right) \right) }^{4}$ divides ${\left| G\right| }^{5}$ for arbitrary $G$ . T. Wilde 20.116. (Well-known problem). Does a profinite torsion group have finite exponent? This problem reduces to the case of pro- $p$ groups (cf. W. Herfort, Arch. Math., 33 (1980), 404-410). Also cf. Archive 3.41. J. S. Wilson
\end{openproblem}

\plainproblemsection

\begin{openproblem}
	20.117. Let $G$ be a pro- $p$ group that is a 3-dimensional Poincaré duality group. Is $G$ coherent?
	
	A group is said to be coherent if each of its finitely generated subgroups is finitely presented, and in the question the coherency is used in the pro- $p$ sense. The question is a famous problem for abstract groups, but has a positive answer for 3-manifold groups. P. Zalesski
\end{openproblem}

\plainproblemsection

\begin{openproblem}
	20.118. Let $G$ be a pro- $p$ group that is a 3-dimensional Poincaré duality group. Can it contain a direct product ${F}_{2} \times  {F}_{2}$ of free pro- $p$ groups of rank 2 ?
	
	If it can, then the answer to 20.117 is negative, but in the abstract case it is known that it can not. P. Zalesski
\end{openproblem}

\plainproblemsection

\begin{openproblem}
	20.119. (G. Wilkes). A pro- $p$ group is said to be accessible if there is a number $n = n\left( G\right)$ such that any finite proper reduced graph of pro- $p$ groups with finite edge groups having fundamental group isomorphic to $G$ has at most $n$ edges (J. Algebra, 525 (2019), 1-18).
	
	Is every finitely presented pro- $p$ group accessible?P. Zalesski
\end{openproblem}

\plainproblemsection

\begin{openproblem}
	20.120. Let $G$ be a locally nilpotent group with an automorphism (of infinite order). Can $G$ be simple as a group with automorphism? E. I. Zelmanov
\end{openproblem}

\plainproblemsection

\begin{openproblem}
	20.121. Let $G$ be a finite almost simple group such that $\operatorname{Soc}\left( G\right)$ is a non-soluble group of Lie type over a field of characteristic $p$ . Let $R$ be a Sylow $q$ -subgroup of $G$ (for some $q$ ) and let ${\operatorname{Min}}_{G}\left( R\right)$ be the subgroup of $R$ generated by all minimal by inclusion intersections of the form $R \cap  {R}^{g}$ , where $g \in  G$ . Is it true that for $p > 3$ the subgroup ${\operatorname{Min}}_{G}\left( R\right)$ is non-trivial if and only if the following hold: $G = \operatorname{Aut}\left( {{L}_{2}\left( p\right) }\right)$ , where $p$ is a Mersenne prime, ${\operatorname{Min}}_{G}\left( R\right)  = R$ , and $q = 2$ .
	
	It is known that this is not always the case for $p = 2,3$ .V. I. Zenkov
\end{openproblem}

\plainproblemsection

\begin{openproblem}
	20.122. For nilpotent subgroups $A, B, C$ of a finite group $G$ , let ${\operatorname{Min}}_{G}\left( {A, B, C}\right)$ be the subgroup of $A$ generated by all minimal by inclusion intersections of the form $A \cap  {B}^{x} \cap  {C}^{y}$ , where $x, y \in  G$ , and let $\mathop{\min }\limits_{G}\left( {A, B, C}\right)$ be the subgroup of ${\operatorname{Min}}_{G}\left( {A, B, C}\right)$ generated by all intersections of this kind of minimal order.
	
	a) Is it true that $\mathop{\min }\limits_{G}\left( {A, B, C}\right)  \leq  F\left( G\right)$ ?
	
	b) Is it true that ${\operatorname{Min}}_{G}\left( {A, B, C}\right)  \leq  F\left( G\right)$ ?
	
	c) The same questions for soluble groups.V. I. Zenkov
\end{openproblem}

\plainproblemsection

\begin{openproblem}
	20.123. A finite group is called a ${D}_{\pi }$ -group if any two of its maximal $\pi$ -subgroups are conjugate.
	
	a) Is it true that for any finite ${D}_{\pi }$ -group $G$ and a $\pi$ -Hall subgroup $H$ of $G$ , there are elements $x, y, z \in  G$ such that ${O}_{\pi }\left( G\right)  = H \cap  {H}^{x} \cap  {H}^{y} \cap  {H}^{z}$ ?
	
	b) Suppose that $G$ is a finite ${D}_{\pi }$ -group in which all simple non-abelian composition factors are sporadic or alternating groups, and let $H$ be a Hall $\pi$ -subgroup of $G$ . Is it true that $H \cap  {H}^{x} \cap  {H}^{y} = {O}_{\pi }\left( G\right)$ for some $x, y \in  G$ ?
	
	c) Suppose that $G$ is a finite ${D}_{\pi }$ -group with trivial soluble radical in which all simple non-abelian composition factors are sporadic groups, and let $H$ be a Hall $\pi$ -subgroup of $G$ . Is it true that $H \cap  {H}^{g} = {O}_{\pi }\left( G\right)$ for some $g \in  G$ ? V. I. Zenkov
\end{openproblem}

\plainproblemsection

\begin{openproblem}
	20.124. A Rota-Baxter operator on a group $G$ is a mapping $B : G \rightarrow  G$ such that $B\left( g\right) B\left( h\right)  = B\left( {{gB}\left( g\right) {hB}{\left( g\right) }^{-1}}\right)$ for all $g, h \in  G$ . Let $F$ be a non-abelian free group. Is there a Rota-Baxter operator on $F$ such that its image is equal to the derived
	
	subgroup $\left\lbrack  {F, F}\right\rbrack$ ? V. G. Bardakov
\end{openproblem}

\plainproblemsection

\begin{openproblem}
	20.125. Does there exist a non-abelian group $G$ and a Rota-Baxter operator $B$ : $G \rightarrow  G$ such that $B$ is surjective but not injective?
\end{openproblem}

\plainproblemsection

\begin{openproblem}
	20.126. A brace $\left( {G;+, \circ  }\right)$ is non-empty set $G$ with two binary operations $+ , \circ$ such that $\left( {G, + }\right)$ is an additively written abelian group, $\left( {G, \circ  }\right)$ is a multiplicatively written group, and $a \circ  \left( {b + c}\right)  + a = \left( {a \circ  b}\right)  + \left( {a \circ  c}\right)$ for all $a, b, c \in  G$ . Does there exist a brace with finitely generated group $\left( {G, + }\right)$ such that
	
	a) the group $\left( {G, \circ  }\right)$ is non-solvable?
	
	b) the group $\left( {G, \circ  }\right)$ contains a non-abelian free group?
\end{openproblem}

\subsection*{References}
\begin{enumerate}[label={[\arabic*]},leftmargin=*,itemsep=0.4em]
	\item Mikl\'{o}s Ab\'{e}rt, \emph{Some Questions}, unpublished problem list, 2 November 2010. \url{https://users.renyi.hu/~abert/questions.pdf}
\end{enumerate}
\begin{enumerate}[label={[\arabic*]},leftmargin=*,itemsep=0.4em]
	\item E.\ I.\ Khukhro and V.\ D.\ Mazurov (eds.), \emph{Unsolved Problems in Group Theory: The Kourovka Notebook}, No.\ 20, Sobolev Institute of Mathematics, Novosibirsk, 2022.
\end{enumerate}

\chapter{Geometry}

\clearpage
\section{Algebraic Geometry}

\plainproblemsection

\begin{openproblem}
	Does there exist a complex abelian fourfold that is the Jacobian of a curve
	and whose Mumford--Tate group is isogenous to a \(\mathbb{Q}\)-form of
	\[
	\mathbb{G}_m\times \operatorname{SL}_2^3?
	\]
\end{openproblem}

\plainproblemsection

\begin{openproblem}
	Let \(F\) be a totally real field of degree \(d\), and let
	\((k_1,\ldots,k_d)\) be the weight of a Hilbert modular form for
	\(\operatorname{GL}_2(F)\). Determine the relation between this weight and
	the restrictions of the associated mod-\(\mathfrak p\) Galois representation
	to the decomposition groups at the primes of \(F\) above \(p\), where
	\(\mathfrak p\) is a prime above \(p\) in the coefficient field. The case
	\(F=\mathbb{Q}\) is known.
\end{openproblem}

\plainproblemsection

\begin{openproblem}
	Determine the zeta function of the moduli space of stable curves of genus
	\(g\) over \(\mathbb{Z}\).
\end{openproblem}

\plainproblemsection

\begin{openproblem}
	Let \((X,\lambda)\) be a generic supersingular polarized abelian variety of
	dimension \(g\geq 2\), and let \(k\) be an algebraically closed field
	containing a field of definition of \(X\). Prove or disprove
	\[
	\operatorname{Aut}\bigl((X,\lambda)\otimes k\bigr)=\{\pm1\}.
	\]
	
	In addition, suppose \(g\geq3\), \(p\) is prime, and a supersingular curve of
	genus \(g\) exists in characteristic \(p\). Must there exist a supersingular
	curve \(C/\overline{\mathbb{F}}_p\) of genus \(g\) with
	\(\operatorname{Aut}(C)=\{1\}\)?
\end{openproblem}

\plainproblemsection

\begin{openproblem}
	Let \(\mathrm{Spor}\) be the set of orders of sporadic finite simple groups
	and define
	\[
	\zeta_{\mathrm{Spor}}(s)
	=\sum_{n\in\mathrm{Spor}}n^{-s}.
	\]
	Prove, without using the classification of finite simple groups, that this
	series converges for \(\operatorname{Re}s>0\).
\end{openproblem}

\plainproblemsection

\begin{openproblem}
	Let \(\mathbb{C}\) be an algebraically closed complete normed field of
	characteristic \(p>0\), let
	\[
	D^\bullet=\{z\in\mathbb{C}:0<|z|<1\},
	\]
	and let \(\sigma\colon D^\bullet\to D^\bullet\) be \(z\mapsto z^q\), where
	\(q\) is a power of \(p\). Suppose that \(\mathcal M\) is a rank-\(n\)
	rigid-analytic vector bundle on \(D^\bullet\) equipped with an isomorphism
	\[
	\tau\colon \sigma^*\mathcal M\longrightarrow\mathcal M.
	\]
	Does \(\mathcal M\) extend to a vector bundle \(\overline{\mathcal M}\) on
	\(D=D^\bullet\cup\{0\}\) such that, for some integer \(r\), \(\tau\) extends
	to a homomorphism
	\[
	\sigma^*\overline{\mathcal M}
	\longrightarrow
	\overline{\mathcal M}\bigl(r[0]\bigr)?
	\]
	
	Equivalently in the free case, let \(R\) be the ring of Laurent series over
	\(\mathbb C\) convergent on \(D^\bullet\), and let \(R'\subset R\) consist
	of series with finite principal part. Is it true that for every
	\(A\in\operatorname{GL}_n(R)\) there exists
	\(B\in\operatorname{GL}_n(R)\) such that
	\[
	B^{-1}A\,{}^\sigma B\in\operatorname{GL}_n(R')?
	\]
\end{openproblem}

\plainproblemsection

\begin{openproblem}
	Let \(F\) be a totally real field, let \(p\) be inert in \(F\), and let
	\(B/F\) be a quaternion algebra unramified outside \(p\) and the
	archimedean places and ramified at every archimedean place. Let
	\(\mathcal O_B\) be a maximal order stable under a positive involution
	\(*\), and define
	\[
	G(R)=
	\bigl\{g\in(\mathcal O_B\otimes R)^\times:g^*g\in R^\times\bigr\}.
	\]
	Determine whether the double-coset space
	\[
	G(\mathbb Q)\backslash G(\mathbb A_f)/G(\widehat{\mathbb Z})
	\]
	is related to a space of Hilbert modular forms of specified weights and
	level, and, if so, give a geometric explanation of this relation.
\end{openproblem}

\plainproblemsection

\begin{openproblem}
	A hyperbolic curve of type \((g,r)\) in characteristic \(p>2\) is called
	\emph{hyperbolically ordinary} if it admits a nilpotent ordinary indigenous
	bundle. Is every hyperbolic curve of every type \((g,r)\), with
	\(2g-2+r>0\), hyperbolically ordinary in every characteristic \(p>2\)?
\end{openproblem}

\plainproblemsection

\begin{openproblem}
	Let \(X_0\) be a hyperbolic curve over a finite field \(k\), and let
	\(\mathcal P_0\) be a nilpotent ordinary indigenous bundle on \(X_0\).
	The pair \((X_0,\mathcal P_0)\) has a canonical lifting
	\((X,\mathcal P)\) over the Witt ring \(W(k)\). Determine when, if ever,
	this canonical lifting is defined over a number field.
\end{openproblem}

\plainproblemsection

\begin{openproblem}
	Let \(k\) be algebraically closed of characteristic \(p>0\). Isomorphism
	classes of \(\mathrm{BT}_1\)-group schemes of type \((d,f)\) are naturally
	parametrized by \(W_X\backslash W_G\), where
	\(W_G\simeq\mathfrak S_d\) and \(W_X\) is the parabolic subgroup determined
	by \(f\).
	
	For decompositions
	\[
	(d,f)=(d_1+d_2,f_1+f_2),
	\]
	describe purely in terms of Weyl groups the map
	\[
	(W_{X_1}\backslash W_{G_1})
	\times(W_{X_2}\backslash W_{G_2})
	\longrightarrow W_X\backslash W_G
	\]
	induced by the product of group schemes. Also describe purely in terms of
	Weyl groups the specialization order \(w'\preccurlyeq w\) on
	\(W_X\backslash W_G\).
\end{openproblem}

\plainproblemsection

\begin{openproblem}
	Let \(S\) be a K3 surface over an algebraically closed field of
	characteristic \(p>0\). Are the following conditions equivalent?
	
	\begin{enumerate}[label=(\roman*)]
		\item \(S\) is a Zariski surface;
		\item \(S\) is unirational;
		\item \(S\) is uniruled;
		\item \(S\) is supersingular in the sense of Shioda, so
		\(\rho(S)=22\);
		\item \(S\) is supersingular in the sense of Artin, so its formal Brauer
		group has infinite height;
		\item \(S\) is rationally connected.
	\end{enumerate}
	
	Several implications and special cases are known, but the full equivalence
	in positive characteristic remains open.
\end{openproblem}

\plainproblemsection

\begin{openproblem}
	Is every elliptic surface in characteristic \(p>0\) liftable to
	characteristic zero?
	
	More generally, construct non-liftable Calabi--Yau varieties in positive
	characteristic, where a Calabi--Yau variety is a nonsingular complete
	variety \(X\) with trivial canonical bundle and
	\[
	H^i(X,\mathcal O_X)=0
	\qquad(1\leq i\leq\dim X-1).
	\]
	Examples are known in some dimensions and characteristics, while K3
	surfaces are liftable.
\end{openproblem}

\plainproblemsection

\begin{openproblem}
	Determine an arithmetic counterpart of the theory of the Jones polynomial
	knot invariant.
\end{openproblem}

\plainproblemsection

\begin{openproblem}
	Let \(R\) be a regular local ring of dimension at least \(2\),
	\(S=\operatorname{Spec}R\), and \(U=S\setminus\{0\}\). Suppose that
	\(X_U\to U\) is either a \(p\)-divisible group or an abelian scheme. Does
	there exist a \(p\)-divisible group or abelian variety \(Y\to S\) whose
	restriction \(Y_U\) is isogenous to \(X_U\)?
\end{openproblem}

\plainproblemsection

\begin{openproblem}
	Let \(K\) be the fraction field of a complete discrete valuation ring
	\(\mathcal O_K\), let \(\kappa\) be its residue field, and let \(T/K\) be a
	torus. A finite splitting extension \(L/K\) determines nonnegative rational
	invariants
	\[
	c_1(\rho_T)\leq\cdots\leq c_d(\rho_T)
	\]
	from the cokernel of the differential of the canonical map between the
	N\'eron models of \(T\) and \(T_L\).
	
	\begin{enumerate}[label=(\roman*)]
		\item Obtain estimates for these invariants, relate them to the invariants
		of the dual integral representation, and describe their variation in
		families.
		\item When \(\kappa\) is imperfect, relate
		\(\sum_i c_i(\rho_T)\) to an appropriate conductor.
		\item For an abelian variety \(A/K\), determine whether
		\(c(A_2)=c(A_1)+c(A_3)\) for every exact sequence
		\(0\to A_1\to A_2\to A_3\to0\), outside the cases where this is known,
		and relate the invariants of \(A\) to those of its dual \(A^t\).
	\end{enumerate}
\end{openproblem}

\plainproblemsection

\begin{openproblem}
	Let \(k\) be algebraically closed of characteristic \(p>0\), and let
	\(\mathcal A_g\) be the moduli space of principally polarized abelian
	varieties. For \(x\in\mathcal A_g(k)\), write \(\mathcal H\cdot x\) for its
	Hecke orbit, \(\mathcal H_\ell\cdot x\) for its \(\ell\)-power Hecke orbit,
	and \(Z(x)\) for the central leaf determined by the polarized
	\(p\)-divisible group of \(x\).
	
	Prove the Oort Hecke-orbit conjectures:
	
	\begin{enumerate}[label=(\roman*)]
		\item the Zariski closure of \(\mathcal H\cdot x\) is the locus of points
		whose Newton polygons equal or lie above that of \(x\);
		\item for every prime \(\ell\neq p\), the orbit
		\(\mathcal H_\ell\cdot x\) is Zariski dense in \(Z(x)\).
	\end{enumerate}
	
	Several special cases are known, but the general assertions remain open.
\end{openproblem}

\plainproblemsection

\begin{openproblem}
	Let \(k'/k\) be an analytic extension of complete non-archimedean fields.
	Extension of scalars is defined for quasi-separated rigid spaces over
	\(k\). Can it be extended naturally to all rigid spaces, compatibly with
	fiber products and open immersions?
	
	Equivalently, if \(i\colon U\hookrightarrow X\) is an open immersion of
	separated rigid spaces, must the induced morphism
	\[
	i'\colon U'\longrightarrow X'
	\]
	be an open immersion? This is known when \(i\) is quasi-compact or
	Zariski-open.
\end{openproblem}

\plainproblemsection

\begin{openproblem}
	For every genus \(g\) and every prime \(p\), does there exist a
	supersingular smooth curve of genus \(g\) in characteristic \(p\)? The
	case \(p=2\) is known, so the question concerns the remaining
	characteristics.
\end{openproblem}

\plainproblemsection

\begin{openproblem}
	Fix a prime power \(q\) and a nonsupersingular Newton polygon \(\beta\).
	For an \(\mathbb F_q\)-isogeny class \(I\) of abelian varieties with Newton
	polygon \(\beta\), let \(N_I\) be the number of
	\(\mathbb F_q\)-isomorphism classes of principally polarized abelian
	varieties in \(I\). Prove, or give an intrinsic heuristic explanation for,
	the predicted average
	\[
	N_I
	=
	e^{O(1)}q^{\operatorname{sdim}(\beta)-d(\beta)}
	\]
	as \(I\) ranges over these isogeny classes, with the expected leading
	constant \(2/v_g+o(1)\). Clarify asymptotically the interaction of
	isomorphism classes, isogeny classes, principal polarizations, and
	\(p\)-power and prime-to-\(p\) isogenies.
\end{openproblem}

\plainproblemsection

\begin{openproblem}
	For positive integers \(n\) and \(p\), is
	\[
	\mathcal R(\mathbb S^n,\mathbb S^p)
	\]
	dense in
	\(\mathcal C^\infty(\mathbb S^n,\mathbb S^p)\) in the
	\(\mathcal C^\infty\)-topology?
\end{openproblem}

\plainproblemsection

\begin{openproblem}
	Let \(X\) be a compact connected nonsingular real algebraic variety of odd
	dimension \(n\). Can every smooth map \(X\to\mathbb S^n\) be approximated
	by regular maps in the \(\mathcal C^\infty\)-topology?
\end{openproblem}

\plainproblemsection

\begin{openproblem}
	Characterize the pairs \((X,Y)\) of nonsingular real algebraic varieties
	for which every smooth map \(X\to Y\) can be approximated by regular maps
	in the \(\mathcal C^\infty\)-topology.
\end{openproblem}

\plainproblemsection

\begin{openproblem}
	Let \(X\) be a compact nonsingular real algebraic curve and let \(Y\) be a
	rationally connected nonsingular real algebraic variety. Can every smooth
	map \(X\to Y\) be approximated by regular maps in the
	\(\mathcal C^\infty\)-topology?
\end{openproblem}

\plainproblemsection

\begin{openproblem}
	Let \(X\) be a compact nonsingular real algebraic curve and let
	\[
	Y=\{(x,y,z)\in\mathbb R^3:x^4+y^4+z^4=1\}.
	\]
	Can every smooth map \(X\to Y\) be approximated by regular maps in the
	\(\mathcal C^\infty\)-topology?
\end{openproblem}

\plainproblemsection

\begin{openproblem}
	Let \(X\) be a compact nonsingular real algebraic variety, let \(p\geq1\),
	and let \(f\colon X\to\mathbb S^p\) be smooth. Suppose
	\[
	f^*H^p(\mathbb S^p;\mathbb Z/2)
	\subset H_{\mathrm{alg}}^p(X;\mathbb Z/2).
	\]
	For \(k\geq0\), can \(f\) be approximated by \(k\)-regulous maps in the
	\(\mathcal C^k\)-topology? The answer is affirmative for \(k=0\) in
	several special dimensions.
\end{openproblem}

\plainproblemsection

\begin{openproblem}
	Let \(k\geq0\) and \(p\geq1\). Can every \(k\)-regulous map
	\(X\to\mathbb S^p\) from a nonsingular real algebraic variety be
	approximated, in the \(\mathcal C^k\)-topology, by nice
	\(k\)-regulous maps?
\end{openproblem}

\plainproblemsection

\begin{openproblem}
	Let \(Z\) be a nonsingular Zariski-closed subvariety of a nonsingular real
	algebraic variety \(X\), and let \(f\colon Z\to\mathbb R\) be
	\(k\)-regulous. Does there exist a \(k\)-regulous function
	\(F\colon X\to\mathbb R\) such that \(F|_Z=f\)? The case \(k=0\) is known,
	so the open part concerns \(k>0\).
\end{openproblem}

\plainproblemsection

\begin{openproblem}
	Consider a polynomial linear system on \(\mathbb R^n\),
	\[
	\sum_{j=1}^p f_{ij}y_j=g_i,
	\qquad i=1,\ldots,m,
	\]
	where \(f_{ij}\) and \(g_i\) are real polynomials. Characterize those
	systems for which the existence of a continuous solution
	\((y_1,\ldots,y_p)\) implies the existence of a regulous solution.
\end{openproblem}

\plainproblemsection

\begin{openproblem}
	Let \(X\) be a compact real algebraic variety and let
	\[
	\Gamma_{\mathbb F}(X)
	=
	K_{\mathbb F}^{(\mathrm{rk})}(X)/
	K_{\mathbb F\text{-}\mathrm{str}}(X),
	\qquad
	\mathbb F\in\{\mathbb R,\mathbb C,\mathbb H\},
	\]
	where the numerator is generated by topological
	\(\mathbb F\)-vector bundles whose rank strata are algebraically
	constructible and the denominator is generated by bundles admitting
	stratified-algebraic structures. Is \(\Gamma_{\mathbb F}(X)\) finite? The
	answer is known when \(\dim X\leq8\).
\end{openproblem}

\plainproblemsection

\begin{openproblem}
	Let \(X\) be a nonsingular real algebraic variety and let \(M\subset X\)
	be a compact smooth submanifold. If the unoriented bordism class of the
	inclusion \(M\hookrightarrow X\) is algebraic, must \(M\) admit a weak
	algebraic approximation in \(X\)? An affirmative result is known in the
	stable range
	\[
	2\dim M+1\leq\dim X,
	\]
	so the remaining question lies outside that range.
\end{openproblem}

\plainproblemsection

\begin{openproblem}
	Let \(X\) be one of
	\(\mathbb R^n\), \(\mathbb S^n\), or \(\mathbb P^n(\mathbb R)\). Does every
	compact smooth submanifold of \(X\) admit an algebraic approximation in
	\(X\)?
\end{openproblem}

\plainproblemsection

\begin{openproblem}
	Let \(\mathcal L_n\subset\mathcal M_2\) be the locus of genus-two curves
	admitting a degree-\(n\) elliptic subcover. Determine \(\mathcal L_n\) for
	\(n=5\) and \(n=7\), and determine the relation between the two elliptic
	curves associated with such a cover.
\end{openproblem}

\plainproblemsection

\begin{openproblem}
	Determine all automorphism groups that can occur for algebraic curves of
	genus \(g\leq10\) in every characteristic. Complete lists are known in
	the smallest genera but not throughout this range, even in characteristic
	zero.
\end{openproblem}

\plainproblemsection

\begin{openproblem}
	Compute explicit equations for the Hurwitz curves of genus \(14\), and, if
	possible, for those of genus \(17\).
\end{openproblem}

\plainproblemsection

\begin{openproblem}
	Find invariants that classify the isomorphism classes of hyperelliptic
	curves of genus \(3\). Equivalently, determine the field of invariants of
	binary octavics.
\end{openproblem}

\plainproblemsection

\begin{openproblem}
	Construct a fast algorithm that, given a hyperelliptic curve \(X_g\) of
	genus \(g\), determines \(\operatorname{Aut}(X_g)\). The intended problem is
	to avoid the inefficient general Gr\"obner-basis approach and to obtain an
	effective implementation at least for small genus.
\end{openproblem}

\plainproblemsection

\begin{openproblem}
	For each of the nontrivial solvable monodromy groups listed in the supplied
	source for a genus-two cover of \(\mathbb P^1\), determine the corresponding
	locus in \(\mathcal M_2\), express it in terms of the absolute invariants
	\(i_1,i_2,i_3\), and determine its dimension.
\end{openproblem}

\plainproblemsection

\begin{openproblem}
	Let \(p\in\mathcal H_g\) be a moduli point of a hyperelliptic curve. If
	\(|\operatorname{Aut}(p)|>2\), is the field of moduli of \(p\) always a
	field of definition?
	
	More generally, determine whether the field of moduli of an arbitrary
	hyperelliptic curve is a field of definition and, when it is, construct an
	explicit model over that field. Only special cases of the first assertion
	are known.
\end{openproblem}

\plainproblemsection

\begin{openproblem}
	Find necessary and sufficient conditions, stated in terms of invariants of
	binary forms, for a hyperelliptic curve to have no automorphisms other than
	the hyperelliptic involution.
\end{openproblem}

\plainproblemsection

\begin{openproblem}
	Give an algorithm which, for a moduli point \(p\in\mathcal H_g\) satisfying
	\(|\operatorname{Aut}(p)|=2\), decides whether the field of moduli is a
	field of definition.
\end{openproblem}

\plainproblemsection

\begin{openproblem}
	For a genus-\(g\) curve \(X\) with an action of a finite group \(G\) of
	signature \(\sigma\), describe the locus
	\(\mathcal M_g(G,\sigma)\) in terms of theta constants.
\end{openproblem}

\plainproblemsection

\begin{openproblem}
	For \(g\geq4\), describe the locus
	\(\mathcal M_g(G,\sigma)\) of genus-\(g\) curves with prescribed
	automorphism group \(G\) and signature \(\sigma\) in terms of vanishing
	theta-nulls.
\end{openproblem}

\plainproblemsection

\begin{openproblem}
	Using the known structure of automorphism groups of algebraic curves of
	genus at most \(48\), determine the possible isogeny decompositions of
	their Jacobians.
\end{openproblem}

\plainproblemsection

\begin{openproblem}
	Develop an efficient algorithm for determining the degree of a rational
	map
	\[
	\phi\colon k^n\dashrightarrow k^m
	\]
	by combining symbolic and numerical methods.
\end{openproblem}

\plainproblemsection

\begin{openproblem}
	Given a parametrizable algebraic surface \(X\), design an algorithm that
	computes a parametrization of \(X\).
\end{openproblem}

\plainproblemsection

\begin{openproblem}
	For each of the partial compactifications
	\[
	\mathcal M_{g,n},\qquad
	\mathcal M_{g,n}^{\mathrm{ct}},\qquad
	\mathcal M_{g,n}^{\mathrm{rt}},\qquad
	\mathcal M_{g,n}^{\leq k\mathrm{\;rat}},
	\]
	determine whether its homotopy type, cohomological dimension, and
	quasi-coherent or constructible cohomological dimension attain the bounds
	predicted by the affine-stratification number.
	In particular, prove or disprove the master conjecture that each such
	space can be covered by at most \(a+1\) affine open subsets, or at least
	stratified by affine varieties of codimension at most \(a\). The parts
	already covered by Harer's theorem are excluded.
\end{openproblem}

\plainproblemsection

\begin{openproblem}
	Find an explicit cell decomposition of \(\mathcal M_g\). More specifically,
	does there exist a natural real-analytic function
	\[
	f\colon\mathcal M_g\longrightarrow\mathbb R
	\]
	that is bounded below and proper and whose Levi form has at most \(g-2\)
	nonpositive eigenvalues at every point? Formulate and answer the analogous
	question on \(\mathcal M_{g,1}\), with \(g-1\) in place of \(g-2\).
	Also determine whether \(\det\Delta\) is a Morse function on
	\(\mathcal M_g\).
\end{openproblem}

\plainproblemsection

\begin{openproblem}
	Determine the Lusternik--Schnirelmann category of \(\mathcal M_g\), and in
	particular compute the cup-length of its cohomology ring.
\end{openproblem}

\plainproblemsection

\begin{openproblem}
	Let \(X\) be a compact K\"ahler manifold. Is the Hochschild cohomology of
	the Fukaya category of \(X\) naturally isomorphic to the quantum
	cohomology of \(X\)? Is the cyclic homology of the Fukaya category
	isomorphic to the Gromov--Witten theory of \(X\)?
\end{openproblem}

\plainproblemsection

\begin{openproblem}
	Determine whether there are stable torsion classes in
	\(AH^*(\mathcal M_{g,n})\), the image of the Chow ring in cohomology.
	Can such classes be represented by visible algebraic cycles? Do all stable
	classes, including torsion, live on the orbifold
	\(\mathcal M_g\)?
\end{openproblem}

\plainproblemsection

\begin{openproblem}
	Let
	\[
	\mathcal I_{g,n}
	=
	\ker\bigl(
	\Gamma_{g,n}\longrightarrow
	\operatorname{Sp}_{2g}(\mathbb Z)
	\bigr)
	\]
	be the Torelli group. Do the even Miller--Morita--Mumford classes
	\(\kappa_{2i}\) restrict to zero in \(H^*(\mathcal I_{g,n})\)? Does this
	vanishing hold stably as \(g\to\infty\)?
\end{openproblem}

\plainproblemsection

\begin{openproblem}
	Let \(\Gamma_\infty\) be the stable mapping class group and
	\(\operatorname{Aut}(F_\infty)\) the stable automorphism group of free
	groups. Determine whether the natural square relating
	\[
	B\Gamma_\infty^+,\quad
	B\operatorname{Sp}_\infty(\mathbb Z)^+,\quad
	B\operatorname{Aut}(F_\infty)^+,\quad
	B\operatorname{GL}_\infty(\mathbb Z)^+
	\]
	is homotopy cartesian after rational localization. Also determine whether
	replacing the lower-right term by Waldhausen's \(A(*)\) and the upper-right
	term by a symplectic analogue makes the square homotopy cartesian
	integrally.
\end{openproblem}

\plainproblemsection

\begin{openproblem}
	Let \(\mathcal M_g^{\mathrm{cc}}\) be the quotient of the compact-type
	moduli space obtained by contracting all rational components, embedded in
	the torus bundle \(D_g^{\mathrm{pr}}\) over the period domain. Is the
	orbifold universal cover of \(\mathcal M_g^{\mathrm{cc}}\)
	\(n_g\)-connected for a sequence \(n_g\to\infty\)?
	Does \(\mathcal M_g^{\mathrm{cc}}\) itself represent a natural moduli
	problem?
\end{openproblem}

\plainproblemsection

\begin{openproblem}
	Give a conceptual explanation of the intersection-number part of Faber's
	conjecture for the tautological ring of \(\mathcal M_g\), beyond its known
	derivation by Virasoro methods, and prove the Morita conjecture.
\end{openproblem}

\plainproblemsection

\begin{openproblem}
	For the locus
	\(\mathcal M_{g,n}^{\leq k\mathrm{\;rational\;components}}\), determine
	whether the top tautological group is
	\[
	R^{g-1+k}\cong\mathbb Q,
	\]
	identify its generators, and determine when a one-dimensional socle should
	be expected.
\end{openproblem}

\plainproblemsection

\begin{openproblem}
	Find operations relating the stable classes on \(B\Gamma_\infty\).
	In particular, explain at the level of infinite loop spaces the
	cohomological description
	\[
	H_{\mathrm{spec}}^*(\mathbb{CP}_{-1}^{\infty})
	\cong \mathbb Z[c_1]\,u,
	\]
	and determine the stable maps
	\(\mathbb{CP}_{-1}^{\infty}\to\mathbb{CP}_{-1}^{\infty}\).
\end{openproblem}

\plainproblemsection

\begin{openproblem}
	Compute the intersection cohomology of the moduli spaces
	\(\overline{\mathcal M}_{g,n}(\mathbb P^r,d)\) of stable maps, and determine
	whether they carry a satisfactory tautological ring.
\end{openproblem}

\plainproblemsection

\begin{openproblem}
	The Hodge class \(\lambda_g\) vanishes rationally but not integrally on
	\(\mathcal A_g\), and its order in integral cohomology is known only up to
	a factor of two. Determine its exact order.
\end{openproblem}

\plainproblemsection

\begin{openproblem}
	For Hurwitz spaces of fixed degree \(d\), determine whether their
	cohomology stabilizes as the genus \(g\to\infty\). More generally, identify
	a Harer-type stability principle for suitable Hurwitz schemes and explain
	whether stable point-counting limits over finite fields arise from such a
	principle.
\end{openproblem}

\plainproblemsection

\begin{openproblem}
	For a finite group \(G\), determine the corresponding homological
	stability theory for maps \(C\to BG\). Clarify its relation to the known
	stabilization, for connected \(G\), of
	\[
	E\operatorname{Diff}(F_{g,1})
	\times_{\operatorname{Diff}(F_{g,1})}
	\operatorname{Map}_{\partial}(F_{g,1},BG)
	\]
	under gluing of tori.
\end{openproblem}

\plainproblemsection

\begin{openproblem}
	Fix a curve \(C\) and consider all its unbranched covers as points in
	\(\mathcal M_g\). Do these points become uniformly dense, with respect to
	the Teichm\"uller metric, in universally defined regions
	\(U_g\subset\mathcal M_g\) as \(g\to\infty\)?
\end{openproblem}

\plainproblemsection

\begin{openproblem}
	For \(\mathcal M_{g,n}\), compute and compare:
	
	\begin{enumerate}[label=(\alph*)]
		\item the signature of its cohomology ring;
		\item the \(L\)-genus defined from Thom's rational Pontryagin classes;
		\item the orbifold \(L\)-genus defined from orbifold Chern classes.
	\end{enumerate}
	
	Determine also the corresponding invariants of the tautological ring.
\end{openproblem}

\plainproblemsection

\begin{openproblem}
	Determine the full algebraic structure carried by the homology of the free
	loop space \(LM\), combining the coproduct
	\[
	H_*(LM)\longrightarrow H_*(LM)\otimes H_*(LM)
	\]
	with the loop product
	\[
	H_*(LM)\otimes H_*(LM)\longrightarrow H_*(LM).
	\]
	In particular, reformulate the interaction of the Hopf- and
	Frobenius-algebra structures in the category of spaces over \(M\).
\end{openproblem}

\plainproblemsection

\begin{openproblem}
	For a global field \(F\), realize the arithmetico-geometric components of
	the Langlands spectral decomposition as motives, rather than merely as
	pieces of cohomology. In particular, obtain a genuinely motivic realization
	for the function-field construction by Drinfeld shtukas.
\end{openproblem}

\plainproblemsection

\begin{openproblem}
	Determine the multiplicities in the spectral decomposition predicted by
	the Langlands program, including the function-field case.
\end{openproblem}

\plainproblemsection

\begin{openproblem}
	Prove the Langlands functoriality conjecture in its general form.
\end{openproblem}

\plainproblemsection

\begin{openproblem}
	Prove the Riemann hypothesis and its modern generalizations involving
	automorphic forms. In the function-field setting, clarify the relation
	between Deligne's theorem and the Grothendieck conjecture.
\end{openproblem}

\subsection*{References}

\begin{enumerate}[label={[\arabic*]},leftmargin=*,itemsep=0.4em]
	\item T. Shaska, ``Some open problems in computational algebraic
	geometry,'' \emph{Albanian Journal of Mathematics} \textbf{1} (2006),
	1--17.
	\item J. Giansiracusa and D. Maulik, \emph{A List of Open Problems and
		Questions on the Moduli Space of Curves}, workshop problem list,
	September 20, 2005.
	\item V. Lafforgue, \emph{Some Open Problems in Algebraic Geometry and the
		Langlands Program}, lecture slides, University of California, Berkeley,
	November 5, 2018.
	\item R. Hain and E. Looijenga, ``Mapping class groups and moduli spaces of
	curves,'' in \emph{Algebraic Geometry---Santa Cruz 1995}, Proceedings of
	Symposia in Pure Mathematics, vol.~62, American Mathematical Society,
	1997, 97--142.
	\item S. Morita, \emph{Cohomological Structure of the Mapping Class Group
		and Beyond}, preprint, 2005.
	\item M. Roth and R. Vakil, ``The affine stratification number and the
	moduli space of curves,'' in \emph{Algebraic Structures and Moduli
		Spaces}, CRM Proceedings \& Lecture Notes, vol.~38, American
	Mathematical Society, 2004, 213--227.
	\item B. Edixhoven, ``The weight in Serre's conjectures on modular forms,''
	\emph{Inventiones Mathematicae} \textbf{109} (1992), 563--594.
	\item S. A. DiPippo and E. W. Howe, ``Real polynomials with all roots on
	the unit circle and abelian varieties over finite fields,''
	\emph{Journal of Number Theory} \textbf{73} (1998), 426--450;
	corrigendum, \textbf{80} (2000), 182.
\end{enumerate}

\clearpage
\section{Differential Geometry}

\plainproblemsection

\begin{openproblem}
	Mark Kac famously asked whether one can hear the shape of a drum, meaning
	whether the Laplace spectrum of a manifold determines its isometry class.
	Milnor answered this in the negative by constructing isospectral,
	nonisometric sixteen-dimensional flat tori. This led to the broader
	question of which geometric properties are audible.
	
	\textbf{Question.}
	Can one hear the presence of an orbifold singularity? More precisely, does
	there exist a pair of isospectral orbifolds such that one has singular
	points while the other is a manifold?
	
	Related constructions give orbit spaces with equivalent
	\(G\)-invariant spectra such that one is an orbifold while the other has a
	non-orbifold singularity.
\end{openproblem}

\plainproblemsection

\begin{openproblem}
	Can a complete metric whose Ricci curvature is bounded below be
	\(L^\infty\)-approximated by a metric with two-sided Ricci-curvature
	bounds? More precisely, for every \(\ell,k>0\), do there exist constants
	\(C,L,K>0\) with the following property? If \((M,g)\) is a complete
	Riemannian manifold satisfying
	\[
	\operatorname{inj}(M,g)\geq\ell,
	\qquad
	\operatorname{Ric}(g)\geq k,
	\]
	then \(M\) admits a Riemannian metric \(h\) such that
	\[
	\operatorname{inj}(M,h)\geq L,\qquad
	|\operatorname{Ric}(h)|\leq K,\qquad
	C^{-1}g\leq h\leq Cg.
	\]
	
	The motivation comes from the Kato square-root problem on manifolds. The
	corresponding problem is known for metrics having both upper and lower
	Ricci bounds and a positive injectivity-radius bound, and stability is
	known under the \(L^\infty\)-perturbations appearing above.
\end{openproblem}

\plainproblemsection

\begin{openproblem}
	By the Bieberbach theorem, every closed flat \(n\)-manifold is finitely
	covered by a flat \(n\)-torus. Its deck group is identified with its
	holonomy group \(H\), a finite subgroup of \(O(n)\). The orthogonal
	\(H\)-representation on \(\mathbb R^n\) is reducible:
	\(\mathbb R^n\) decomposes as an orthogonal direct sum of nontrivial
	\(H\)-invariant subspaces. The known proof is algebraic and uses the
	classification of finite simple groups. The result has geometric
	consequences, including the existence of nonhomothetic flat deformations
	of every closed flat manifold.
	
	\textbf{Question.}
	Give a purely geometric proof that the holonomy representation of every
	closed flat manifold is reducible.
\end{openproblem}

\plainproblemsection

\begin{openproblem}
	For a tangent two-plane \(\sigma\) in a four-manifold, its biorthogonal
	curvature is the average of the sectional curvatures of \(\sigma\) and
	\(\sigma^\perp\). Positive biorthogonal curvature is intermediate between
	positive sectional curvature and positive scalar curvature.
	
	The manifold \(\mathbb S^2\times\mathbb T^2\) clearly admits metrics of
	nonnegative biorthogonal curvature, but it is not known whether such
	metrics can be deformed to have strictly positive biorthogonal curvature.
	A solution would be important for the classification of nonsimply
	connected closed four-manifolds with positive biorthogonal curvature. The
	simply connected case has been classified up to homeomorphism.
	
	\textbf{Question.}
	Does \(\mathbb S^2\times\mathbb T^2\) admit a Riemannian metric of positive
	biorthogonal curvature?
\end{openproblem}

\plainproblemsection

\begin{openproblem}
	Let \(\Sigma\) be a smooth complete \((m+n)\)-dimensional Riemannian
	manifold, let \(\Gamma\subset\Sigma\) be a smooth oriented
	\((m-1)\)-dimensional submanifold, and let \(T\) be an area-minimizing
	integer rectifiable \(m\)-current spanning \(\Gamma\). Write
	\(\operatorname{Reg}_b(T)\subset\Gamma\) for the relatively open set of
	boundary regular points and
	\[
	\operatorname{Sing}_b(T)
	=\Gamma\setminus\operatorname{Reg}_b(T)
	\]
	for the boundary singular set.
	
	The problem is relevant in codimension \(n\geq2\), since
	\(\operatorname{Sing}_b(T)=\varnothing\) in codimension one. In higher
	codimension, \(\operatorname{Reg}_b(T)\) is known to be dense in
	\(\Gamma\), and there are two-dimensional currents with crossing boundary
	singularities \(P_k\) which accumulate at a branching singularity \(P\).
	This contrasts with the interior theory, where the singularities of a
	two-dimensional area-minimizing current are isolated.
	
	\begin{enumerate}[label=(\alph*)]
		\item Does \(\operatorname{Sing}_b(T)\) have zero
		\((m-1)\)-dimensional Hausdorff measure?
		\item If so, is its Hausdorff dimension at most \(m-2\)?
		\item Can one construct an example of the preceding type in which the
		accumulating singularities \(P_k\) are themselves of branching type?
		\item White's conjecture: if a two-dimensional area-minimizing current
		spans a real-analytic closed curve, must it have only finitely many
		interior and boundary singularities, and hence finite topological type?
		\item Does there exist a smooth closed curve
		\(\Gamma\subset\mathbb R^{2+n}\) which bounds an area-minimizing
		two-current of infinite topological type?
	\end{enumerate}
\end{openproblem}

\plainproblemsection

\begin{openproblem}
	Flag manifolds are homogeneous manifolds obtained as quotients of complex
	groups by Borel subgroups. A basic example is the three-dimensional
	complex full flag manifold
	\[
	F_{1,2}=\operatorname{SL}(3,\mathbb C)/B,
	\]
	where \(B\) is the subgroup of upper triangular matrices. Each of the real
	forms \(\operatorname{SL}(3,\mathbb R)\) and
	\(\operatorname{SU}(2,1)\) has a unique compact real
	three-dimensional orbit, modeling flat path geometries and spherical CR
	structures, respectively.
	
	\textbf{Question.}
	Which manifolds can be modeled on an orbit of a real form in a space of
	flags?
\end{openproblem}

\plainproblemsection

\begin{openproblem}
	On a real three-manifold, a flag structure can be specified by two complex
	lines in the complexified tangent bundle. A special case arises from a
	totally real immersion into
	\[
	F_{1,2}=\operatorname{SL}(3,\mathbb C)/B.
	\]
	Intersecting the natural foliations of the flag manifold with the immersed
	manifold produces the two complex line fields. This is a special flat-type
	structure, although it is not locally equivalent to either of the
	real-form orbits described above.
	
	\textbf{Question.}
	Determine the homotopy classification of totally real immersions of real
	three-manifolds into \(F_{1,2}\). Gromov's \(h\)-principle applies, so the
	problem is expected to be homotopy-theoretic.
\end{openproblem}

\plainproblemsection

\begin{openproblem}
	In real projective space with its standard metric, the projective
	subspaces obtained by projecting equatorial spheres are the only
	volume-minimizing cycles in their homology classes.
	
	For \(\mu>0\), the Berger metric \(g_\mu\) on
	\(\mathbb S^{2n+1}\) is obtained by multiplying the standard metric in the
	Hopf direction \(H=JN\) by \(\mu\), while leaving the orthogonal
	distribution unchanged:
	\[
	g_\mu(H,H)=\mu,\qquad
	g_\mu|_{H^\perp}=g|_{H^\perp},\qquad
	g_\mu(H,H^\perp)=0.
	\]
	Passing to the antipodal quotient gives the corresponding Berger metric on
	\(\mathbb{RP}^{2n+1}\).
	
	In dimension three, the projected equatorial two-spheres are minimal for
	every \(\mu\), although they are totally geodesic only when \(\mu=1\).
	They are precisely the area-minimizing projective planes.
	
	\textbf{Question.}
	For \(2n+1>3\) and \(0<k<2n+1\), are the projective subspaces obtained by
	projecting \(k\)-dimensional equatorial spheres minimal in
	\((\mathbb{RP}^{2n+1},g_\mu)\)? Are they the only
	\(k\)-dimensional volume-minimizing cycles in their homology classes?
	
	There is also a dimension-three identification
	\[
	(T^r\mathbb S^2,g^S)
	\cong
	(\mathbb{RP}^3,4g_{r^2}),
	\]
	which preserves the natural circle-bundle structures and is special to
	dimension three.
\end{openproblem}

\plainproblemsection

\begin{openproblem}
	Let \(G\) be a compact connected Lie group with a bi-invariant Riemannian
	metric. Then \(G\) is a symmetric space, so its curvature tensor is
	parallel. Along a geodesic with tangent field \(X\), the Jacobi equation
	has constant coefficients in a parallel orthonormal frame, and the Jacobi
	operator is
	\[
	R(X,\mathord\cdot)X=\frac14\operatorname{ad}_X^2.
	\]
	Since \(\operatorname{ad}_X\) is skew-symmetric, the nonzero eigenvalues
	occur in pairs. Hence conjugate points have even order and every geodesic
	segment has even Morse index.
	
	\textbf{Question.}
	If a left-invariant Riemannian metric on a compact connected Lie group has
	the property that every geodesic segment has even Morse index, must the
	metric be bi-invariant?
	
	Equivalently, for every pair of nonconjugate points \(p,q\in G\), the
	energy functional
	\[
	E(\gamma)=\frac12\int|\dot\gamma|^2\,dt
	\]
	on the Hilbert manifold of \(H^1\)-curves joining \(p\) to \(q\) is a
	perfect Morse function. Metrics with this property are called pointwise
	taut.
\end{openproblem}

\plainproblemsection

\begin{openproblem}
	The classification of closed manifolds admitting metrics of positive
	scalar curvature is complete for simply connected manifolds of dimension
	at least five, but remains largely open for manifolds with finite
	fundamental group.
	
	A connected closed \(d\)-manifold \(M\) is called \(p\)-toral if there
	exist classes
	\[
	c_1,\ldots,c_d\in H^1(M;\mathbb F_p)
	\]
	such that
	\[
	c_1\smile\cdots\smile c_d\neq0
	\quad\text{in }H^d(M;\mathbb F_p);
	\]
	otherwise \(M\) is \(p\)-atoral. The torus \(\mathbb T^d\) is
	\(p\)-toral for every prime \(p\), as is every closed manifold oriented
	bordant to \(\mathbb T^d\) over
	\(B(\mathbb Z/p)^d\). Moreover,
	\(\mathbb{RP}^{2m+1}\) is \(2\)-toral. If \(p\) is odd, every
	\(p\)-toral manifold is orientable. Moreover, if
	\[
	\dim M>
	\operatorname{rank}H^1(B\pi_1(M);\mathbb F_p),
	\]
	then \(M\) is \(p\)-atoral, because degree-one classes square to zero for
	odd \(p\).
	
	\textbf{Question.}
	If \(M\) admits a metric of positive scalar curvature, must it be
	\(p\)-atoral for every odd prime \(p\)?
	
	The answer is affirmative in dimensions at most two. In dimension three
	it follows from prime decomposition and geometrization. A positive answer
	is also known in dimensions at least five when \(\pi_1(M)\) is an
	elementary abelian group of odd prime order and \(M\) is \(p\)-atoral.
	
	There are also \(p\)-toral nonspin manifolds of dimension \(d\geq4\)
	with finite fundamental group of odd order. Such examples can be
	constructed using the bordism observation above. Their universal covers
	are closed, simply connected, and nonspin, and hence admit
	positive-scalar-curvature metrics for \(d\geq5\).
\end{openproblem}

\plainproblemsection

\begin{openproblem}
	A map between metric spaces is a coarse embedding if there are control
	functions \(\omega,\Omega\), with \(\omega(D)\to\infty\) as
	\(D\to\infty\), such that pairs of points at distance \(D\) are mapped to
	pairs whose distances lie between \(\omega(D)\) and \(\Omega(D)\).
	Coarse embeddability arose in convex geometry and later became important
	in geometric group theory, notably for fundamental groups and their
	large-scale topology.
	
	Large-scale analogues of topological dimension are nondecreasing under
	coarse embeddings. Other noninteger-valued invariants include variants of
	conformal dimension for boundaries of hyperbolic spaces and sharp metric
	cotype.
	
	\textbf{Question.}
	Find further numerical invariants of metric spaces that are nondecreasing
	under coarse embeddings.
\end{openproblem}

\plainproblemsection

\begin{openproblem}
	Coarse embeddings can sometimes be approximated by harmonic maps.
	
	\textbf{Question.}
	Find applications of the harmonic-map approximation of coarse embeddings.
\end{openproblem}

\plainproblemsection

\begin{openproblem}
	Let \(G\) be a solvable Lie group. Prove that solvsolitons are the only
	local maxima, among left-invariant Riemannian metrics on \(G\), of the
	Ricci pinching functional
	\[
	F(g)
	=
	\frac{\operatorname{Scal}(g)^2}
	{|\operatorname{Ric}(g)|^2}.
	\]
	A solvsoliton is a left-invariant Riemannian metric satisfying
	\[
	\operatorname{Ric}(g)=cI+D,
	\]
	where \(D\) is a derivation of the Lie algebra of \(G\).
\end{openproblem}

\plainproblemsection

\begin{openproblem}
	Bryant observed that certain finite-type classification problems in
	geometry can be formulated as the integration of a Lie algebroid
	\(A\to X\) to a Lie groupoid \(\mathcal G\rightrightarrows X\). Examples
	include metrics of constant sectional curvature, metrics of Hessian type,
	Bochner--K\"ahler metrics, and symplectic connections with special
	holonomy. The integrating groupoid presents the stack of the corresponding
	moduli problem.
	
	For a Poisson manifold \((M,\pi)\), the cotangent bundle \(T^*M\) has a
	natural Lie-algebroid structure. When this algebroid is integrable, its
	source-simply-connected integration is a symplectic groupoid whose arrow
	space carries a multiplicative symplectic form.
	
	In the finite-type examples above, the relevant algebroid arises as the
	cotangent algebroid of a Poisson manifold. Nevertheless, the original
	classification problem may have no evident Poisson structure, and any
	natural symplectic or Poisson structure present does not appear to arise
	directly from the multiplicative form on the integrating groupoid. The
	associated moduli spaces still exhibit a ``symplectic nature.''
	
	\textbf{Question.}
	Explain the existence and role of this symplectic nature of the
	groupoid/algebroid and its relevance to moduli spaces of finite-type
	geometric structures.
\end{openproblem}

\plainproblemsection

\begin{openproblem}
	Do there exist complete embedded minimal surfaces of finite genus and
	Morse index \(4\)? More specifically, in the one-parameter deformation of
	Costa's surface, does the Morse index remain constant?
	
	Complete embedded minimal surfaces of index \(3\) do not exist.
\end{openproblem}

\plainproblemsection

\begin{openproblem}
	\begin{enumerate}[label=(\alph*)]
		\item Construct congruence lattices which are simultaneously
		\(\ell^1\)- and \(\ell^{1*}\)-isospectral in every dimension.
		\item Do there exist families of more than two lens spaces which are
		\(p\)-isospectral for every \(p\)?
		\item Give a procedure which passes from \(p\)-isospectral lens spaces in
		dimension \(m\), for every \(p\), to such spaces in dimension \(m+1\).
		\item Establish connections between these questions and toric geometry.
	\end{enumerate}
\end{openproblem}

\plainproblemsection

\begin{openproblem}
	Let \(M^n\) be an irreducible full homogeneous submanifold of a sphere,
	with \(n\geq2\). If its normal holonomy group does not act transitively,
	must \(M\) be an orbit of an \(s\)-representation?
\end{openproblem}

\plainproblemsection

\begin{openproblem}
	For homogeneous Riemannian manifolds with nontrivial nullity, determine
	whether there are examples which:
	
	\begin{enumerate}[label=(\alph*)]
		\item are not topologically trivial;
		\item admit a presentation \(G/H\) with \(G\) nonsolvable;
		\item are K\"ahler;
		\item occur in every dimension \(d\geq5\).
	\end{enumerate}
\end{openproblem}

\plainproblemsection

\begin{openproblem}
	Let \(X\) be a simply connected noncompact homogeneous three-manifold
	diffeomorphic to \(\mathbb R^3\). CMC spheres exist exactly for mean
	curvatures
	\[
	H\in(H(X),\infty),
	\]
	where \(H(X)\) is the critical mean curvature of \(X\). Whenever they
	exist, CMC spheres are unique up to congruence and therefore have a
	well-defined center of symmetry. For every \(x\in X\), the CMC spheres
	centered at \(x\) form a real-analytic one-parameter family. They are
	known to be Alexandrov embedded, but their embeddedness is not known in
	general.
	
	\textbf{Questions.}
	Do the CMC spheres centered at \(x\) form a foliation of
	\(X\setminus\{x\}\)? Can such a foliation be used to prove that all CMC
	spheres in \(X\) are embedded?
\end{openproblem}

\plainproblemsection

\begin{openproblem}
	Does every complete embedded minimal surface in \(\mathbb R^3\) have to be
	proper? This is the Calabi--Yau problem.
\end{openproblem}

\plainproblemsection

\begin{openproblem}
	Prove the Hoffman--Meeks conjecture: if a complete embedded minimal surface
	has genus \(g\), \(k\) ends, and finite total curvature, then
	\[
	k\leq g+2.
	\]
\end{openproblem}

\plainproblemsection

\begin{openproblem}
	A submetry is a map with equidistant fibers. Every submetry from a round
	sphere is defined by an algebra of polynomials which is both maximal and
	Laplacian, meaning stable under the Laplacian, and conversely every maximal
	Laplacian algebra defines a spherical submetry.
	
	\textbf{Question.}
	Is every Laplacian algebra of polynomials maximal?
	
	For a quotient submetry
	\[
	\mathbb S^{n-1}\longrightarrow\mathbb S^{n-1}/G
	\]
	arising from a representation \(\rho:G\to O(n)\), the Laplacian algebra is
	the algebra of invariant polynomials. An affirmative answer would imply
	that the invariant algebra of
	\[
	\rho\oplus\rho:G\longrightarrow O(2n)
	\]
	is the Laplacian algebra generated by the invariants of \(\rho\) together
	with
	\[
	\delta(x,y)=\sum_i x_i y_i.
	\]
	This is known when the Laplacian algebra is generated by quadratic
	polynomials, using a generalization of Weyl's first fundamental theorem
	for \(O(n)\).
\end{openproblem}

\plainproblemsection

\begin{openproblem}
	Let \(N^n\) be a closed Riemannian manifold, \(n\geq3\), and consider the
	space \(W_{\mathrm{imm}}^{3,2}(\Sigma^2,N^n)\) of Sobolev immersions of a
	closed surface. An admissible family \(\mathcal A\) is invariant under
	homeomorphisms of this immersion space isotopic to the identity. Its
	minimax value, or width, is
	\[
	\beta_{\mathcal A}
	=
	\inf_{A\in\mathcal A}\sup_{\Phi\in A}\operatorname{Area}(\Phi)>0.
	\]
	
	There exists a smooth surface \(S\) with
	\(\operatorname{genus}(S)\leq\operatorname{genus}(\Sigma)\) and a smooth,
	possibly branched, minimal immersion \(\Phi_{\mathcal A}:S\to N\) such
	that
	\[
	\beta_{\mathcal A}
	=
	\operatorname{Area}(\Phi_{\mathcal A}).
	\]
	If \(\mathcal A\) is homotopic of dimension \(d\), meaning it is given by
	continuous maps from \(\mathbb S^d\) into the immersion space, the known
	upper estimate is
	\[
	\operatorname{Index}(\Phi_{\mathcal A})\leq d.
	\]
	
	\textbf{Question.}
	Prove the complementary lower estimate
	\[
	d\leq
	\operatorname{Index}(\Phi_{\mathcal A})
	+\operatorname{Null}(\Phi_{\mathcal A}),
	\]
	where the nullity is the dimension of the space of Jacobi fields along
	\(\Phi_{\mathcal A}\).
\end{openproblem}

\plainproblemsection

\begin{openproblem}
	In the setting of the preceding minimax construction, prove that every
	closed Riemannian manifold \(N^n\), \(n\geq3\), contains infinitely many
	distinct minimal branched two-dimensional immersions.
\end{openproblem}

\plainproblemsection

\begin{openproblem}
	Frankel's theorem states that, in a complete connected Riemannian
	\(n\)-manifold of positive sectional curvature, two closed totally
	geodesic submanifolds of dimensions \(n_1\) and \(n_2\) must intersect
	whenever \(n_1+n_2\geq n\).
	
	\textbf{Question.}
	Does Frankel's intersection theorem remain true for reversible
	(symmetric) Finsler metrics of positive flag curvature?
	
	The theorem holds in dimension two but fails for nonreversible Finsler
	metrics on \(\mathbb S^2\).
\end{openproblem}

\plainproblemsection

\begin{openproblem}
	Sectional curvature connects local Riemannian geometry with the global
	topology of a manifold. The Bott--Grove--Halperin conjecture predicts that
	nonnegative sectional curvature imposes strong restrictions on rational
	homotopy. It would in particular imply the Hopf conjecture on the Euler
	characteristic and the conjectural bound
	\[
	\sum_i b_i(M)\leq2^{\dim M}
	\]
	for the total Betti number of a compact nonnegatively curved manifold.
	
	Partial results are known for manifolds with an isometric action of
	cohomogeneity two and for real-analytic metrics having an entire Grauert
	tube. For a real-analytic Riemannian manifold, the adapted complex
	structure is defined near the zero section in \(TM\) by complexifying the
	geodesic flow; an entire Grauert tube means that this structure extends to
	all of \(TM\). In this special setting, Morse theory on the based loop
	space and estimates for Jacobi fields yield polynomial growth of its
	cohomology.
	
	\textbf{Conjecture.}
	Every compact simply connected Riemannian manifold \(M\) of nonnegative
	sectional curvature is rationally elliptic; equivalently,
	\[
	\dim_{\mathbb Q}\bigl(\pi_*(M)\otimes\mathbb Q\bigr)<\infty.
	\]
\end{openproblem}

\plainproblemsection

\begin{openproblem}
	\textbf{Fundamental groups of complete affine manifolds.}
	Which virtually polycyclic groups occur as fundamental groups of compact
	complete affine manifolds?
\end{openproblem}

\plainproblemsection

\begin{openproblem}
	\textbf{Chern conjecture for affine manifolds.}
	Prove that every compact affine manifold \(M\) has zero Euler
	characteristic:
	\[
	\chi(M)=0.
	\]
\end{openproblem}

\plainproblemsection

\begin{openproblem}
	\textbf{Markus conjecture.}
	Prove that a compact orientable affine manifold is complete if and only if
	it admits a parallel volume form.
\end{openproblem}

\plainproblemsection

\begin{openproblem}
	\textbf{Complete affine structures on products.}
	Prove that, for \(g\geq2\), the manifold
	\[
	\Sigma_g\times\mathbb S^1
	\]
	does not admit a complete affine structure.
\end{openproblem}

\plainproblemsection

\begin{openproblem}
	\textbf{Strominger--Yau--Zaslow conjecture.}
	Let \(X\) and \(\check X\) be a mirror pair of Calabi--Yau \(n\)-folds.
	Prove that there exist special Lagrangian torus fibrations
	\[
	f\colon X\longrightarrow B
	\qquad\text{and}\qquad
	\check f\colon\check X\longrightarrow B
	\]
	over a common base \(B\), with general fiber an \(n\)-torus. Moreover,
	prove that the two fibrations are dual: whenever \(X_b\) and
	\(\check X_b\) are nonsingular torus fibers, there are canonical
	identifications
	\[
	X_b\cong H^1(\check X_b,\mathbb R/\mathbb Z)
	\qquad\text{and}\qquad
	\check X_b\cong H^1(X_b,\mathbb R/\mathbb Z).
	\]
\end{openproblem}

% Supplemental problems selected from S.-T. Yau's 1982 problem section.
% Problems known to be solved or disproved are omitted; if only part of an
% original item remains open, only that part is retained.

\plainproblemsection

\begin{openproblem}
	Does \(\mathbb S^2\times\mathbb S^2\) admit a Riemannian metric of positive
	sectional curvature? More generally, determine which higher-rank compact
	symmetric spaces admit metrics of positive sectional curvature, and classify
	the compact four-manifolds that admit such metrics.
\end{openproblem}

\plainproblemsection

\begin{openproblem}
	Which exotic spheres admit Riemannian metrics of positive sectional
	curvature?
\end{openproblem}

\plainproblemsection

\begin{openproblem}
	Let \(M^n\) be a compact Riemannian manifold of nonnegative sectional
	curvature. Is it true that
	\[
	b_i(M)\leq b_i(\mathbb T^n)=\binom{n}{i}
	\]
	for every \(i\)?
\end{openproblem}

\plainproblemsection

\begin{openproblem}
	Does every compact Riemannian manifold of positive sectional curvature admit
	a smooth effective \(\mathbb S^1\)-action?
\end{openproblem}

\plainproblemsection

\begin{openproblem}
	Let \(E\to M\) be a vector bundle over a compact manifold \(M\) that admits a
	metric of positive sectional curvature. Does the total space of \(E\) admit a
	complete metric of nonnegative sectional curvature?
\end{openproblem}

\plainproblemsection

\begin{openproblem}
	Let \(M\) be a compact Riemannian manifold of positive sectional curvature.
	Is every abelian subgroup of \(\pi_1(M)\) cyclic?
\end{openproblem}

\plainproblemsection

\begin{openproblem}
	\textbf{Hopf conjecture.}
	Prove that every compact even-dimensional Riemannian manifold of positive
	sectional curvature has positive Euler characteristic.
\end{openproblem}

\plainproblemsection

\begin{openproblem}
	Characterize the groups that occur as fundamental groups of compact
	negatively curved manifolds. In particular, classify the negatively curved
	compact Kähler surfaces, and determine whether the outer automorphism group
	of the fundamental group must be finite.
\end{openproblem}

\plainproblemsection

\begin{openproblem}
	\textbf{Hopf--Singer sign conjecture.}
	If \(M^{2n}\) is a compact Riemannian manifold of negative sectional
	curvature, must
	\[
	(-1)^n\chi(M)>0?
	\]
\end{openproblem}

\plainproblemsection

\begin{openproblem}
	Develop a higher-dimensional analogue of the Cohn--Vossen inequality.
	More specifically, if \(M\) is a complete manifold of finite volume and
	bounded curvature, determine conditions under which its Euler characteristic
	equals the integral of the Gauss--Bonnet--Chern form. Also determine an
	appropriate equality statement for complete noncompact manifolds of
	nonnegative curvature.
\end{openproblem}

\plainproblemsection

\begin{openproblem}
	Develop a useful notion of curvature for piecewise-linear manifolds that
	supports both suitable pinching theorems and curvature formulas for
	characteristic classes.
\end{openproblem}

\plainproblemsection

\begin{openproblem}
	Determine the possible Pontryagin and Stiefel--Whitney classes of compact
	negatively curved manifolds. In particular, does every such manifold have a
	finite cover that is spin?
\end{openproblem}

\plainproblemsection

\begin{openproblem}
	Let \((M_1,g_1)\) and \((M_2,g_2)\) be compact Einstein manifolds of negative
	sectional curvature, with
	\[
	\pi_1(M_1)\cong\pi_1(M_2)
	\qquad\text{and}\qquad
	\dim M_1\geq3.
	\]
	Must \((M_1,g_1)\) and \((M_2,g_2)\) be isometric?
\end{openproblem}

\plainproblemsection

\begin{openproblem}
	Suppose a complete noncompact manifold has nonnegative sectional curvature
	and has strictly positive sectional curvature at some point. Although the
	manifold is then diffeomorphic to Euclidean space, can its given metric always
	be deformed through complete nonnegatively curved metrics to a metric of
	positive sectional curvature?
\end{openproblem}

\plainproblemsection

\begin{openproblem}
	Let \(\Omega\) be a closed \(4k\)-form on a compact manifold \(M\)
	representing a Pontryagin class. Find necessary and sufficient conditions for
	the existence of a Riemannian metric on \(M\) whose Chern--Weil curvature form
	representing that Pontryagin class is exactly \(\Omega\). Formulate and solve
	the analogous problem for Chern forms.
\end{openproblem}

\plainproblemsection

\begin{openproblem}
	Let \(T_{ij}\) be a symmetric tensor on a compact manifold. Determine
	necessary and sufficient conditions for the existence of a metric \(g\)
	satisfying
	\[
	\operatorname{Ric}_{ij}(g)-\frac{1}{2}R(g)g_{ij}=T_{ij},
	\]
	allowing \(T\) to depend on \(g\) and its first derivatives. If the manifold
	has boundary, determine the appropriate boundary conditions. The Lorentzian
	four-dimensional case is the prescribed energy--momentum problem for the
	Einstein equations.
\end{openproblem}

\plainproblemsection

\begin{openproblem}
	Characterize the fundamental groups of complete manifolds of positive Ricci
	curvature. The compact case is strongly constrained by the
	Cheeger--Gromoll splitting theorem, whereas the noncompact case permits
	substantially more complicated behavior.
\end{openproblem}

\plainproblemsection

\begin{openproblem}
	Construct an explicit Ricci-flat metric on a K3 surface. More generally,
	determine whether there exists a compact Ricci-flat four-manifold whose
	finite covers are neither tori nor K3 surfaces. In particular, can
	\(\mathbb S^4\) or \(\mathbb S^2\times\mathbb S^2\) admit a Ricci-flat
	metric?
\end{openproblem}

\plainproblemsection

\begin{openproblem}
	Classify compact four-dimensional Einstein manifolds of negative scalar
	curvature. In particular, does \(\mathbb S^4\) admit an Einstein metric with
	negative Einstein constant?
\end{openproblem}

\plainproblemsection

\begin{openproblem}
	Find dimension-dependent constants \(c_n<C_n\) for which a suitable
	two-sided pinching condition
	\[
	c_n g\leq \operatorname{Ric}_g\leq C_n g
	\]
	forces the underlying manifold to admit an Einstein metric.
\end{openproblem}

\plainproblemsection

\begin{openproblem}
	Classify complete three-dimensional Riemannian manifolds of nonnegative
	scalar curvature.
\end{openproblem}

\plainproblemsection

\begin{openproblem}
	Classify compact four-dimensional Einstein manifolds of positive scalar
	curvature.
\end{openproblem}

\plainproblemsection

\begin{openproblem}
	If a compact aspherical manifold admits a metric of nonnegative scalar
	curvature, must that metric be flat? More generally, determine whether a
	compact \(K(\pi,1)\)-manifold admits a metric of nonnegative scalar curvature
	only in the flat case.
\end{openproblem}

\plainproblemsection

\begin{openproblem}
	Classify compact hypersurfaces in \(\mathbb R^{n+1}\) having constant scalar
	curvature. Must every such hypersurface be congruent to a round sphere? The
	answer is known to be affirmative under convexity.
\end{openproblem}

\plainproblemsection

\begin{openproblem}
	Let \(\Lambda\subset\mathbb S^n\) be closed with
	\[
	\dim_{\mathrm H}\Lambda<\frac n2-1.
	\]
	Does \(\mathbb S^n\setminus\Lambda\) admit a complete metric of constant
	positive scalar curvature conformal to the round metric?
\end{openproblem}

\plainproblemsection

\begin{openproblem}
	Let \(X\) be a Fano manifold whose tangent bundle is nef. Must \(X\) be a
	rational homogeneous manifold?
\end{openproblem}

\plainproblemsection

\begin{openproblem}
	Let \(M\) be a complete noncompact Kähler manifold of positive holomorphic
	bisectional curvature. Must \(M\) be biholomorphic to \(\mathbb C^n\)?
	In the nonnegative case, determine when \(M\) is biholomorphic to the total
	space of a holomorphic vector bundle over a compact Hermitian symmetric
	space.
\end{openproblem}

\plainproblemsection

\begin{openproblem}
	Is every complete simply connected Kähler manifold of negative holomorphic
	bisectional curvature Stein? Investigate also whether the compact case
	imposes corresponding negativity properties on the holomorphic tangent
	bundle.
\end{openproblem}

\plainproblemsection

\begin{openproblem}
	Let \(M\) be a complete Kähler manifold of finite volume and bounded
	curvature. Must \(M\) be biholomorphic to a Zariski-open subset of a
	projective variety? If the holomorphic bisectional curvature is negative,
	must the automorphism group of \(M\) be finite?
\end{openproblem}

\plainproblemsection

\begin{openproblem}
	Let \(M\) be a compact Kähler manifold of negative sectional curvature and
	complex dimension greater than one. Is its complex structure unique among
	the complex structures compatible with the underlying smooth manifold?
\end{openproblem}

\plainproblemsection

\begin{openproblem}
	Let \(M\) be a complete simply connected Kähler manifold with sectional
	curvature \(K\leq-1\). Does \(M\) carry sufficiently many bounded
	holomorphic functions to separate points? In particular, can \(M\) be
	realized holomorphically as a bounded domain? This includes the case in
	which \(M\) is the universal cover of a compact Kähler manifold of negative
	sectional curvature.
\end{openproblem}

\plainproblemsection

\begin{openproblem}
	Is every complete Ricci-flat Kähler manifold biholomorphic to a Zariski-open
	subset of a compact Kähler manifold?
\end{openproblem}

\plainproblemsection

\begin{openproblem}
	Classify compact complex surfaces that admit scalar-flat Kähler metrics.
\end{openproblem}

\plainproblemsection

\begin{openproblem}
	Study the canonical Kähler--Einstein metric on Teichmüller space and
	determine its precise relation to the Bergman metric. More generally, if an
	irreducible bounded domain covers a compact Kähler manifold, under what
	conditions do its Kähler--Einstein and Bergman metrics coincide?
\end{openproblem}

\plainproblemsection

\begin{openproblem}
	Find further Chern-number inequalities for compact Kähler--Einstein
	manifolds of negative scalar curvature. In complex dimension four, is the
	top Chern number \(c_4[M]\) necessarily positive?
\end{openproblem}

\plainproblemsection

\begin{openproblem}
	Let \(\{M_t\}\) be a degenerating holomorphic family of canonically polarized
	Kähler manifolds, and let \(g_t\) be the canonical Kähler--Einstein metric on
	\(M_t\). Describe the metric and analytic behavior of \(g_t\) near a
	degenerate fiber.
\end{openproblem}

\plainproblemsection

\begin{openproblem}
	Find topological criteria determining whether a given smooth four-manifold
	admits a complex structure. In particular, characterize compact complex
	surfaces covered by the unit ball, and determine when two complex structures
	on the same smooth four-manifold can be joined by a deformation.
\end{openproblem}

\plainproblemsection

\begin{openproblem}
	In the space of smooth immersions of a fixed compact surface into
	\(\mathbb R^3\), is the set of infinitesimally rigid immersions generic?
	Describe its complement. The same questions may be asked within the class of
	surfaces of revolution.
\end{openproblem}

\plainproblemsection

\begin{openproblem}
	Let \(M\) be a complete Riemannian manifold with bounded Ricci curvature and
	positive injectivity radius. Does \(M\) admit an isometric embedding into
	some Euclidean space whose mean-curvature vector is uniformly bounded?
\end{openproblem}

\plainproblemsection

\begin{openproblem}
	Develop a higher-dimensional form of the Weyl isometric-embedding problem.
	In particular, determine conditions under which a compact \(n\)-manifold of
	positive sectional curvature admits an isometric embedding into
	\(\mathbb R^{n(n+1)/2}\).
\end{openproblem}

\plainproblemsection

\begin{openproblem}
	Call an isometric embedding into Euclidean space elliptic if, for every
	normal direction, the corresponding second fundamental form has at least two
	nonzero eigenvalues of the same sign. Are two elliptic isometric embeddings
	of a fixed compact manifold necessarily congruent? What is the correct
	higher-dimensional analogue of the Cohn--Vossen rigidity theorem?
	
	In dimension two, is a complete immersed surface in \(\mathbb R^3\) rigid
	whenever it has finite area and bounded nonpositive Gaussian curvature?
\end{openproblem}

\plainproblemsection

\begin{openproblem}
	Does there exist a complete hypersurface in \(\mathbb R^n\) whose Ricci
	curvature is everywhere less than \(-1\)? More generally, can hyperbolic
	\(n\)-space be isometrically immersed in \(\mathbb R^{2n-1}\), possibly after
	allowing a precisely specified weak class of singularities?
\end{openproblem}

\plainproblemsection

\begin{openproblem}
	Find nontrivial sufficient conditions under which a complete negatively
	curved surface admits an isometric immersion into \(\mathbb R^3\). Such
	conditions may involve the rate at which the Gaussian curvature tends to
	zero at infinity and should be related to the Dirichlet problem for
	prescribed Gaussian curvature.
\end{openproblem}

\plainproblemsection

\begin{openproblem}
	Classify compact real-analytic Weingarten surfaces, that is, surfaces whose
	Gaussian and mean curvatures satisfy a nonsingular relation
	\(\Phi(K,H)=0\). In particular, classify the higher-genus examples and
	determine whether the genus-one examples must be tubes or surfaces of
	revolution. Also determine intrinsic characterizations of algebraic
	surfaces through relations among their principal curvatures.
\end{openproblem}

\plainproblemsection

\begin{openproblem}
	Let \(h\colon\mathbb R^3\to\mathbb R\) be smooth. Find natural conditions on
	\(h\) that guarantee the existence of a closed immersed surface of a
	prescribed genus whose mean curvature at \(x\) equals \(h(x)\). Formulate the
	corresponding prescribed-curvature problem.
\end{openproblem}

\plainproblemsection

\begin{openproblem}
	Let \(S\) be the boundary of a convex body in \(\mathbb R^3\), and suppose
	the intrinsic radius of \(S\) is at most \(1\). Determine the largest possible
	area of \(S\) and characterize the equality case.
\end{openproblem}

\plainproblemsection

\begin{openproblem}
	\textbf{Milnor's principal-curvature problem.}
	Let \(\Sigma\) be a complete noncompact surface immersed in
	\(\mathbb R^3\), with principal curvatures \(\kappa_1\) and \(\kappa_2\).
	Prove that at least one of the following occurs:
	\[
	\inf_\Sigma|\kappa_1-\kappa_2|=0,\qquad
	K\ \text{changes sign},\qquad
	K=0\ \text{somewhere}.
	\]
\end{openproblem}

\plainproblemsection

\begin{openproblem}
	\textbf{Carathéodory conjecture.}
	Prove that every smooth closed convex surface in \(\mathbb R^3\) has at
	least two umbilic points.
\end{openproblem}

\plainproblemsection

\begin{openproblem}
	\textbf{Smooth planar inverse spectral problem.}
	Let \(\Omega_1,\Omega_2\subset\mathbb R^2\) be bounded domains with smooth
	boundary. If their Dirichlet Laplacians have identical spectra, including
	multiplicities, must the domains be congruent? Known planar counterexamples
	with corners do not settle the smooth-boundary problem.
\end{openproblem}

\plainproblemsection

\begin{openproblem}
	Suppose the Laplace spectra of two compact Riemannian manifolds, or of two
	bounded domains with the same boundary condition, agree except at finitely
	many eigenvalues. Must the full spectra agree? What can be concluded when
	the exceptional set is infinite but has density zero?
\end{openproblem}

\plainproblemsection

\begin{openproblem}
	\textbf{Pólya conjecture.}
	Let \(\Omega\subset\mathbb R^2\) be a bounded domain of area
	\(|\Omega|\), and let \(\lambda_k^D\) be its Dirichlet eigenvalues. Prove
	\[
	\lambda_k^D\geq\frac{4\pi k}{|\Omega|}.
	\]
	Formulate and prove the corresponding sharp reverse inequality for the
	Neumann spectrum.
\end{openproblem}

\plainproblemsection

\begin{openproblem}
	Describe the discrete and continuous spectra of complete finite-volume
	manifolds whose sectional curvature is bounded between two negative
	constants. Determine when the discrete spectrum is nonempty, its asymptotic
	distribution, and its relation to the closed geodesics. Seek an appropriate
	\(L^2\)-index theorem for natural elliptic operators on such manifolds.
\end{openproblem}

\plainproblemsection

\begin{openproblem}
	Let \((M,g)\) be a compact smooth Riemannian surface, and let
	\(\varphi_\lambda\) be an eigenfunction with eigenvalue \(\lambda\).
	Find the sharp upper bound for the length
	\[
	\mathcal H^1\bigl(\{\varphi_\lambda=0\}\bigr)
	\]
	in terms of \(\sqrt{\lambda}\) and the geometry of \(M\). In particular,
	establish the expected linear bound in \(\sqrt{\lambda}\) for smooth metrics.
\end{openproblem}

\plainproblemsection

\begin{openproblem}
	For a closed surface of fixed genus \(g\) and a fixed eigenvalue index \(k\),
	determine the largest possible multiplicity of \(\lambda_k\). Construct
	explicit metrics that attain the maximal multiplicity whenever it is finite
	and attained.
\end{openproblem}

\plainproblemsection

\begin{openproblem}
	For the Hodge Laplacian on differential forms over a closed Riemannian
	manifold, obtain effective estimates for the first positive eigenvalue in
	terms of computable geometric data. Determine which combinations of
	diameter, curvature bounds, injectivity radius, and topology yield sharp
	estimates in each degree.
\end{openproblem}

\plainproblemsection

\begin{openproblem}
	\textbf{Schiffer--Pompeiu problem.}
	Let \(\Omega\subset\mathbb R^2\) be a smooth bounded domain. Suppose there
	is a nonconstant Neumann eigenfunction \(u\) satisfying
	\[
	\Delta u+\lambda u=0\quad\text{in }\Omega,\qquad
	\partial_\nu u=0\quad\text{on }\partial\Omega,
	\]
	and \(u\) is constant on \(\partial\Omega\). Must \(\Omega\) be a disk?
\end{openproblem}

\plainproblemsection

\begin{openproblem}
	For a Riemannian metric on \(\mathbb S^2\), determine geometric hypotheses
	that give a nontrivial lower bound for the spectral gap
	\(\lambda_4-\lambda_1\), and find the sharp bound under those hypotheses.
\end{openproblem}

\plainproblemsection

\begin{openproblem}
	Determine the lower-order terms in the Weyl asymptotic expansion for the
	eigenvalue counting function of a compact Riemannian manifold, with or
	without boundary, under the weakest natural dynamical hypotheses.
\end{openproblem}

\plainproblemsection

\begin{openproblem}
	Characterize the increasing sequences of positive real numbers that occur,
	with multiplicity, as the Laplace spectrum of a compact Riemannian manifold.
	Beyond the Weyl law, determine the necessary relations among the
	eigenvalues that arise from the finite-dimensional nature of the space of
	metrics.
\end{openproblem}

\plainproblemsection

\begin{openproblem}
	Let \(\Omega\) be a geodesically convex domain in a Riemannian manifold and
	let \(\varphi_1\) be its first positive Dirichlet eigenfunction. Find
	geometric hypotheses under which \(\log\varphi_1\) is geodesically concave,
	thereby extending the Euclidean log-concavity theorem.
\end{openproblem}

\plainproblemsection

\begin{openproblem}
	Determine the point spectrum of the Laplacian on complete noncompact
	manifolds under curvature hypotheses. In particular:
	\begin{enumerate}[label=\textup{(\roman*)},leftmargin=2.2em]
		\item can a complete manifold of nonnegative sectional curvature have
		nonempty \(L^2\) point spectrum;
		\item under what additional hypotheses does a complete simply connected
		manifold with \(-C\leq K\leq-1\) have no \(L^2\) eigenvalues;
		\item must a complete finite-volume manifold with
		\(-C\leq K\leq0\) have infinitely many \(L^2\) eigenvalues?
	\end{enumerate}
\end{openproblem}

\plainproblemsection

\begin{openproblem}
	Does every compact Riemannian manifold carry infinitely many geometrically
	distinct closed geodesics?
\end{openproblem}

\plainproblemsection

\begin{openproblem}
	\textbf{Blaschke conjecture for projective-type spaces.}
	Prove that a Blaschke manifold modeled topologically on
	\(\mathbb{CP}^n\), \(\mathbb{HP}^n\), or the Cayley plane is, after a
	constant rescaling, isometric to the corresponding compact rank-one
	symmetric space.
\end{openproblem}

\plainproblemsection

\begin{openproblem}
	\textbf{Lichnerowicz conjecture for compact harmonic manifolds.}
	Is every compact harmonic manifold either flat or locally symmetric of rank
	one? Equivalently, classify the remaining compact harmonic manifolds whose
	universal cover is noncompact.
\end{openproblem}

\plainproblemsection

\begin{openproblem}
	Let \(M\) be a compact Riemannian manifold with finite fundamental group.
	Must \(M\) possess a nonhyperbolic closed geodesic?
\end{openproblem}

\plainproblemsection

\begin{openproblem}
	For a Riemannian manifold diffeomorphic to \(\mathbb S^n\), obtain
	nontrivial lower bounds for the number of geometrically distinct embedded
	closed geodesics. The classical two-dimensional lower bound is three.
\end{openproblem}

\plainproblemsection

\begin{openproblem}
	Find higher-dimensional analogues of the Loewner and Pu systolic
	inequalities. Identify the correct homotopy or cohomology hypotheses, the
	optimal constants, and the equality metrics.
\end{openproblem}

\plainproblemsection

\begin{openproblem}
	Prove that every Riemannian metric on a manifold diffeomorphic to
	\(\mathbb S^3\) admits at least four distinct embedded minimal two-spheres.
\end{openproblem}

\plainproblemsection

\begin{openproblem}
	Does every compact smooth manifold admit a smooth embedding into some round
	sphere \(\mathbb S^N\) whose image is a minimal submanifold?
\end{openproblem}

\plainproblemsection

\begin{openproblem}
	Classify complete embedded minimal surfaces of finite genus in
	\(\mathbb R^3\). The classification should describe the possible end
	structures and the geometric asymptotics of each end.
\end{openproblem}

\plainproblemsection

\begin{openproblem}
	Does every smooth regular Jordan curve in \(\mathbb R^3\) bound only
	finitely many stable minimal surfaces?
\end{openproblem}

\plainproblemsection

\begin{openproblem}
	Does there exist a smooth simple Jordan curve in \(\mathbb R^3\) that bounds
	a nontrivial continuous family of minimal disks?
\end{openproblem}

\plainproblemsection

\begin{openproblem}
	Let \(\gamma\subset\mathbb S^3=\partial\mathbb B^4\) be a smooth knot that
	bounds a smoothly embedded disk in \(\mathbb B^4\). Is \(\gamma\) isotopic
	in \(\mathbb S^3\) to a knot that bounds an embedded minimal disk in
	\(\mathbb B^4\)?
\end{openproblem}

\plainproblemsection

\begin{openproblem}
	Describe the moduli space of embedded minimal surfaces of fixed genus in
	\(\mathbb S^3\). Which conformal structures occur? Determine how the answer
	changes for minimal immersions into \(\mathbb S^N\), \(N>3\).
\end{openproblem}

\plainproblemsection

\begin{openproblem}
	\textbf{Yau's first-eigenvalue conjecture.}
	Let \(M^n\subset\mathbb S^{n+1}\) be a closed embedded minimal hypersurface.
	Is the first positive eigenvalue of its Laplace--Beltrami operator equal to
	\(n\)?
\end{openproblem}

\plainproblemsection

\begin{openproblem}
	Does some round sphere \(\mathbb S^N\) contain a closed minimal surface whose
	Gaussian curvature is everywhere negative? Determine the possible
	dimensions, genera, and codimensions.
\end{openproblem}

\plainproblemsection

\begin{openproblem}
	Let \(u\colon\mathbb R^n\to\mathbb R\) be an entire solution of the minimal
	surface equation. Must \(u\) have polynomial growth? If not, characterize
	the possible growth rates.
\end{openproblem}

\plainproblemsection

\begin{openproblem}
	Classify topologically the seven-dimensional area-minimizing cones in
	\(\mathbb R^8\).
\end{openproblem}

\plainproblemsection

\begin{openproblem}
	\textbf{Chern conjecture.}
	Consider closed minimal hypersurfaces of a round sphere having constant
	scalar curvature. Is the set of scalar-curvature values attained by such
	hypersurfaces discrete?
\end{openproblem}

\plainproblemsection

\begin{openproblem}
	Let \(M\) be a closed hyperbolic three-manifold and let \(\Sigma\) be a
	closed surface for which an immersion \(f\colon\Sigma\to M\) induces an
	injection on fundamental groups. The homotopy class contains a minimal
	immersion. Under what generic hypotheses on \(M\) is that minimal immersion
	unique?
\end{openproblem}

\plainproblemsection

\begin{openproblem}
	For complete minimal hypersurfaces in \(\mathbb R^{n+1}\), find a
	higher-dimensional analogue of the theorem limiting the number of omitted
	values of the Gauss map of a complete minimal surface in \(\mathbb R^3\).
	Determine the optimal exceptional set in the appropriate Grassmannian.
\end{openproblem}

\plainproblemsection

\begin{openproblem}
	Let \(H\) be an area-minimizing hypersurface in a Riemannian manifold \(M\).
	Can the metric of \(M\) be perturbed so that a representative of the same
	homology class remains area minimizing and has no singularities? Formulate
	an analogous regularization problem in the piecewise-linear category.
\end{openproblem}

\plainproblemsection

\begin{openproblem}
	Let \(\Omega^k\subset\mathbb R^n\) be a compact minimal submanifold with
	boundary and codimension greater than two. Determine whether the sharp
	Euclidean isoperimetric inequality
	\[
	\operatorname{Vol}_k(\Omega)^{k-1}
	\leq \frac{1}{k^k\omega_k}\,
	\operatorname{Vol}_{k-1}(\partial\Omega)^k,
	\]
	holds, where \(\omega_k\) is the volume of the unit ball in
	\(\mathbb R^k\). Determine also the equality case, without assuming that
	\(\Omega\) is area minimizing.
\end{openproblem}

\plainproblemsection

\begin{openproblem}
	Let \(f\colon\Sigma\to M^3\) be an area-minimizing immersion in its homotopy
	class, without assuming that \(f\) is homotopic to an embedding. Relate the
	self-intersection complexity of \(f\), and especially the number of triple
	points, to the topology of \(\Sigma\), \(M\), and the induced homomorphism on
	fundamental groups.
\end{openproblem}

\plainproblemsection

\begin{openproblem}
	Which elements of \(\pi_i(\mathbb S^n)\) can be represented by harmonic
	maps \(\mathbb S^i\to\mathbb S^n\)? More generally, determine which homotopy
	classes of maps into a compact manifold with finite fundamental group admit
	harmonic representatives.
\end{openproblem}

\plainproblemsection

\begin{openproblem}
	Does every closed affine manifold have Euler characteristic zero?
\end{openproblem}

\plainproblemsection

\begin{openproblem}
	\textbf{Cosmic censorship.}
	For asymptotically flat vacuum initial data \((M^3,g,h)\) satisfying the
	Einstein constraint equations, describe the singularities in the maximal
	Cauchy development. Under suitable generic hypotheses, are all singularities
	hidden behind event horizons, so that no naked singularity is visible from
	null infinity?
\end{openproblem}

\plainproblemsection

\begin{openproblem}
	Give an intrinsic and geometrically robust definition of total angular
	momentum for an asymptotically flat spacetime, especially at null infinity.
	Determine its relation to total mass and its transformation under the
	asymptotic symmetry group.
\end{openproblem}

\plainproblemsection

\begin{openproblem}
	Let \(P\to M^4\) be a principal \(\operatorname{SU}(2)\)-bundle over a
	compact oriented four-manifold. Find necessary and sufficient conditions,
	in terms of the topology of \(P\), the intersection form of \(M\), and the
	choice of Riemannian metric, for \(P\) to admit an irreducible anti-self-dual
	connection.
\end{openproblem}

\plainproblemsection

\begin{openproblem}
	\textbf{\(11/8\)-conjecture.}
	Let \(M\) be a closed smooth spin four-manifold with intersection form \(Q_M\).
	Prove
	\[
	b_2(M)\geq\frac{11}{8}\,|\sigma(M)|,
	\]
	where \(\sigma(M)\) is the signature. Equivalently, in the simply connected
	indefinite case, determine whether the intersection form is realized by the
	standard combinations arising from K3 surfaces and
	\(\mathbb S^2\times\mathbb S^2\).
\end{openproblem}

\subsection*{References}

\begin{enumerate}[label={[\arabic*]},leftmargin=*,itemsep=0.4em]
	\item S.-T. Yau, ``Problem section,'' in \emph{Seminar on Differential
		Geometry}, Annals of Mathematics Studies, vol.~102, Princeton University
	Press and University of Tokyo Press, 1982, 669--706.
\end{enumerate}

\clearpage
\section{Geometric Group Theory}

\plainproblemsection

\begin{openproblem}
	Suppose that \(G\) admits a finite \(K(G,1)\) and contains no
	Baumslag--Solitar subgroup
	\[
	BS(m,n)=\langle x,y\mid x^{-1}y^m x=y^n\rangle.
	\]
	Must \(G\) be hyperbolic? If \(G\) embeds in a hyperbolic group, must \(G\)
	be hyperbolic? In the special case of a finite nonpositively curved
	two-complex \(Y=K(G,1)\), does the existence of a flat in
	\(\widetilde Y\) force a subgroup \(\mathbb Z^2<G\)?
\end{openproblem}

\plainproblemsection

\begin{openproblem}
	Suppose that \(G\) admits a finite \(K(G,1)\), contains no
	\(\mathbb Z^2\), and has the property that whenever \(x\) has infinite
	order and \(x^m\) is conjugate to \(x^n\), one has \(|m|=|n|\). Must
	\(G\) be hyperbolic?
\end{openproblem}

\plainproblemsection

\begin{openproblem}
	If \(G\) is word-hyperbolic, does the Rips complex \(P_d(G)\) admit a
	\(G\)-equivariant negatively curved metric for all sufficiently large
	\(d\)?
\end{openproblem}

\plainproblemsection

\begin{openproblem}
	Does every one-ended word-hyperbolic group contain the fundamental group
	of a closed hyperbolic surface?
\end{openproblem}

\plainproblemsection

\begin{openproblem}
	For a given \(n\), does there exist an \(n\)-dimensional hyperbolic group
	whose every infinite-index subgroup is free? More generally, for a given
	\(k\leq n-2\), can such a group have no quasiconvex subgroup of
	codimension at most \(k\)?
\end{openproblem}

\plainproblemsection

\begin{openproblem}
	Let \(H\) be a finitely presented subgroup of a word-hyperbolic group
	\(G\), assume \(H\) has finite index in its normalizer, and suppose the
	intersection of any \(n\) distinct conjugates of \(H\) is finite for some
	\(n>0\). Must \(H\) be quasiconvex in \(G\)? The question remains even
	when \(H\) is malnormal.
\end{openproblem}

\plainproblemsection

\begin{openproblem}
	Let \(X\) be a finite two-complex with fundamental group \(G\), and let
	\(X_H\) be the cover corresponding to a finitely presented subgroup
	\(H<G\). If the injectivity radius \(I(x)\) on \(X_H\) tends to infinity as
	\(x\to\infty\), must \(H\) be quasi-isometrically embedded in \(G\)?
\end{openproblem}

\plainproblemsection

\begin{openproblem}
	Let \(G\) be word-hyperbolic and \(H<G\) finitely presented. If every
	\(g\in G\) has a positive power lying in \(H\), must \(H\) have finite
	index in \(G\)?
\end{openproblem}

\plainproblemsection

\begin{openproblem}
	Can a one-ended hyperbolic group which is not virtually a surface group
	have the property that every infinite-index subgroup is free?
\end{openproblem}

\plainproblemsection

\begin{openproblem}
	Can a finite-index subgroup of a one-ended hyperbolic group be isomorphic
	to an infinite-index subgroup of the same group? The torsion-free case is
	excluded by the co-Hopf property recorded in the source.
\end{openproblem}

\plainproblemsection

\begin{openproblem}
	Prove a full combination theorem for relatively hyperbolic groups,
	extending the known partial versions.
\end{openproblem}

\plainproblemsection

\begin{openproblem}
	Is every word-hyperbolic group residually finite?
\end{openproblem}

\plainproblemsection

\begin{openproblem}
	For a group \(G\), let \(\operatorname{rank}(G^n)\) be the least number of
	generators of the direct power \(G^n\). If \(G\) is word-hyperbolic, must
	\[
	\lim_{n\to\infty}\operatorname{rank}(G^n)=\infty?
	\]
	Also find natural classes such as CAT(0) or automatic groups for which this
	conclusion fails.
\end{openproblem}

\plainproblemsection

\begin{openproblem}
	For a word-hyperbolic group \(G\), is there an algorithm to compute
	\[
	\widetilde H^i(\partial G)
	\cong H^{i+1}(G;\mathbb ZG)?
	\]
	In particular, can one decide algorithmically whether
	\(\widetilde H^i(\partial G)\cong\widetilde H^i(\mathbb S^2)\) for every
	\(i\)?
\end{openproblem}

\plainproblemsection

\begin{openproblem}
	Suppose a hyperbolic group \(G\) is a graph of hyperbolic groups and every
	edge-to-vertex inclusion is a quasi-isometric embedding. Describe the
	point-preimages of each Cannon--Thurston map
	\(\partial V\to\partial G\) associated with a vertex group \(V\), and
	prove that the map is finite-to-one.
\end{openproblem}

\plainproblemsection

\begin{openproblem}
	Can the fundamental group of a hyperbolic three-manifold fibering over the
	circle contain a subgroup which is locally free but not free?
\end{openproblem}

\plainproblemsection

\begin{openproblem}
	Establish a Kan--Thurston theorem using only locally CAT\((-1)\) or
	word-hyperbolic groups: for every finite simplicial complex \(X\), construct
	a locally CAT\((-1)\) polyhedral complex \(Y\) and a map \(Y\to X\) that is
	surjective on fundamental groups and induces homology isomorphisms for
	every local coefficient system on \(X\).
\end{openproblem}

\plainproblemsection

\begin{openproblem}
	Give a proof of Johannson's theorem on
	\(\operatorname{Out}(\pi_1M)\) for Haken three-manifolds analogous to the
	Rips--Sela theorem. More generally, if a CAT(0) group \(G\) has infinite
	\(\operatorname{Out}(G)\), must \(G\) admit an infinite-order Dehn twist?
\end{openproblem}

\plainproblemsection

\begin{openproblem}
	If \(G\) admits a finite-dimensional \(K(G,1)\), must it act properly
	discontinuously by isometries on a complete CAT(0) space?
\end{openproblem}

\plainproblemsection

\begin{openproblem}
	Does there exist a group of cohomological dimension \(2\) and geometric
	dimension \(3\)?
\end{openproblem}

\plainproblemsection

\begin{openproblem}
	Is every subcomplex of an aspherical two-complex itself aspherical?
\end{openproblem}

\plainproblemsection

\begin{openproblem}
	Let \(S\) be a closed surface of genus at least \(2\), let
	\(V=S\times S\), and let \(\Delta\subset V\) be the diagonal. If
	\(\widetilde V\) is a nontrivially branched finite cover of \(V\) along
	\(\Delta\), prove that its natural piecewise-hyperbolic CAT(0) metric cannot
	be replaced by a \(C^2\)-smooth Riemannian metric of nonpositive
	sectional curvature.
\end{openproblem}

\plainproblemsection

\begin{openproblem}
	Suppose \(G\) acts geometrically on CAT(0) spaces \(X\) and \(Y\). Does
	there exist a compact metric space \(Z\) with cell-like maps
	\[
	Z\longrightarrow\partial X,
	\qquad
	Z\longrightarrow\partial Y?
	\]
\end{openproblem}

\plainproblemsection

\begin{openproblem}
	Let \(G\) act geometrically on a CAT(0) space, or suppose that \(G\) is
	automatic. For \(a,b\in G\), does there exist \(n>0\) such that
	\(\langle a^n,b^n\rangle\) is either free or abelian?
\end{openproblem}

\plainproblemsection

\begin{openproblem}
	Do CAT(0) groups, automatic groups, or biautomatic groups satisfy the Tits
	alternative?
\end{openproblem}

\plainproblemsection

\begin{openproblem}
	Does every Artin group admit a finite \(K(G,1)\)?
\end{openproblem}

\plainproblemsection

\begin{openproblem}
	Let a finitely generated group \(G\) act properly by isometries on a proper
	CAT(0) space. Can \(G\) be an infinite torsion group? If the action is also
	cocompact, can \(G\) contain an infinite torsion subgroup?
\end{openproblem}

\plainproblemsection

\begin{openproblem}
	Let \(G\) be a Coxeter group acting properly on a CAT(0) space \(X\), and
	let \(H\) be a special subgroup of \(G\). Must \(H\) act cocompactly on some
	closed convex subset of \(X\)?
\end{openproblem}

\plainproblemsection

\begin{openproblem}
	Classify Coxeter groups up to abstract group isomorphism.
\end{openproblem}

\plainproblemsection

\begin{openproblem}
	Classify Artin groups up to abstract group isomorphism.
\end{openproblem}

\plainproblemsection

\begin{openproblem}
	Are all finite-type Artin groups CAT(0)? Are all Artin groups automatic?
\end{openproblem}

\plainproblemsection

\begin{openproblem}
	Let \(G\) act geometrically on a CAT(0) space \(X\) with isolated flats,
	meaning that \(X\) contains no isometrically embedded triplane. For every
	finitely generated subgroup \(H<G\), is the inclusion \(H\hookrightarrow
	G\) a quasi-isometric embedding if and only if \(H\) is quasiconvex relative
	to this action?
\end{openproblem}

\plainproblemsection

\begin{openproblem}
	For torsion-free quasiconvex subgroups \(A,B\) of a quasiconvex Kleinian
	group in \(\mathbb H^n\), prove an inequality bounding the volume of the
	convex core of \(A\cap B\) in terms of the convex-core volumes of \(A\) and
	\(B\), analogous to the Hanna Neumann inequality for free Fuchsian groups.
\end{openproblem}

\plainproblemsection

\begin{openproblem}
	If a finitely generated group \(G\) has homological dimension \(1\), must
	\(G\) be free?
\end{openproblem}

\plainproblemsection

\begin{openproblem}
	Are all limit groups CAT(0)? Characterize which amalgams
	\(F_m*_{\mathbb Z}F_n\) and which one-relator groups are limit groups.
\end{openproblem}

\plainproblemsection

\begin{openproblem}
	For \(n>1\), is the subset of \(F_n^n\) consisting of ordered free bases of
	\(F_n\) first-order definable? It is conjectured that it is not. Also
	determine whether the only definable subgroups of \(F_n\) are cyclic
	subgroups and \(F_n\) itself.
\end{openproblem}

\plainproblemsection

\begin{openproblem}
	For an element of a free group having commutator genus \(g\), let \(f(g)\)
	be the uniform bound on the number of Nielsen-equivalence classes of its
	expressions as a product of \(g\) commutators. Find an explicit upper bound
	for \(f(g)\).
\end{openproblem}

\plainproblemsection

\begin{openproblem}
	Is \(\operatorname{PSL}_2(\mathbb R)\) a maximal closed subgroup of
	\(\operatorname{Homeo}_+(\mathbb S^1)\)? Is there a proper closed subgroup
	of \(\operatorname{Homeo}_+(\mathbb S^1)\) acting transitively on unordered
	\(k\)-tuples for \(k\neq3\)?
\end{openproblem}

\plainproblemsection

\begin{openproblem}
	Let \(G_1,G_2<\operatorname{Isom}(\mathbb H^n)\) be finitely generated
	discrete groups. If their sets of axes coincide, must \(G_1\) and \(G_2\)
	be commensurable?
\end{openproblem}

\plainproblemsection

\begin{openproblem}
	For a closed hyperbolic manifold
	\(M=\mathbb H^n/\Gamma\), \(n\geq3\), determine whether \(\Gamma\) admits a
	faithful nondiscrete representation into
	\(\operatorname{Isom}_+(\mathbb H^n)\), and characterize those \(M\) for
	which such a representation exists.
\end{openproblem}

\plainproblemsection

\begin{openproblem}
	Let \(G<\operatorname{Isom}(\mathbb H^n)\) be a finitely generated
	Kleinian group. Is
	\[
	\delta(G)\leq\operatorname{vcd}(G),
	\]
	with equality exactly when \(G\) preserves a totally geodesic
	\(\mathbb H^k\subset\mathbb H^n\) with compact quotient? Here
	\(\delta(G)\) is the exponent of convergence.
\end{openproblem}

\plainproblemsection

\begin{openproblem}
	Is there an \(\varepsilon>0\) such that every finitely generated Kleinian
	group \(G\) with \(\delta(G)<\varepsilon\) is virtually free?
\end{openproblem}

\plainproblemsection

\begin{openproblem}
	Does there exist a finitely generated discrete subgroup of
	\(\operatorname{SO}(n,1)\) whose action on its limit set is nonergodic, or
	nonrecurrent?
\end{openproblem}

\plainproblemsection

\begin{openproblem}
	If a group \(G\) is virtually of type \(FP\) over \(\mathbb F_p\) and
	\(g\in G\) has order \(p\), is the centralizer \(C_G(g)\) also virtually of
	type \(FP\) over \(\mathbb F_p\)?
\end{openproblem}

\plainproblemsection

\begin{openproblem}
	Does there exist a group of finite virtual cohomological dimension which
	does not act with finite stabilizers on an acyclic complex of dimension
	equal to its virtual cohomological dimension?
\end{openproblem}

\plainproblemsection

\begin{openproblem}
	If \(G\) is of type \(FP\) over \(\mathbb Q\), must there be a uniform
	bound on the orders of its finite subgroups?
\end{openproblem}

\plainproblemsection

\begin{openproblem}
	Starting with a finitely presented group and alternately taking maximal
	splittings over finite and two-ended subgroups, can the process fail to
	terminate? Prove strong accessibility for the remaining general
	hyperbolic and CAT(0) cases not covered by the partial result recorded in
	the source.
\end{openproblem}

\plainproblemsection

\begin{openproblem}
	Does there exist a finitely presented one-ended group \(G\) with
	\[
	G\cong G*\mathbb Z?
	\]
\end{openproblem}

\plainproblemsection

\begin{openproblem}
	Does there exist a finitely presented torsion-free group which does not
	split over a virtually abelian subgroup but has infinitely many splittings
	over \(F_2\)?
\end{openproblem}

\plainproblemsection

\begin{openproblem}
	If a finitely presented group acts nontrivially on an \(\mathbb R\)-tree,
	must it act nontrivially on a simplicial tree?
\end{openproblem}

\plainproblemsection

\begin{openproblem}
	For a finitely presented group \(G\), compare:
	
	\begin{enumerate}[label=(\alph*)]
		\item some finite-index subgroup of \(G\) has a nontrivial action on a
		simplicial tree;
		\item for a finite complex \(X\) with \(\pi_1X=G\), some covering space of
		\(X\) has at least two ends.
	\end{enumerate}
	
	To what extent are these two properties equivalent?
\end{openproblem}

\plainproblemsection

\begin{openproblem}
	For \(n\geq3\), does there exist a group of type \(F_n\) but not
	\(F_{n+1}\) which contains no subgroup \(\mathbb Z^{n-1}\)?
\end{openproblem}

\plainproblemsection

\begin{openproblem}
	Does there exist a torsion-free group \(G\) of infinite cohomological
	dimension and a constant \(n_0\) such that every subgroup \(H<G\) of finite
	cohomological dimension satisfies
	\[
	\operatorname{cd}H\leq n_0?
	\]
\end{openproblem}

\plainproblemsection

\begin{openproblem}
	Can the Pr\"ufer \(p\)-group \(\mathbb Z_{p^\infty}\) be embedded in a group
	of type \(FP_\infty\),
	or in a group of type \(FP_3\)? More generally, can every group of type
	\(F_n\) be embedded in a group of type \(F_{n+1}\) for \(n\geq2\)?
\end{openproblem}

\plainproblemsection

\begin{openproblem}
	Compute the asymptotic dimensions of CAT(0) groups,
	\(\operatorname{Out}(F_n)\), mapping class groups, nonuniform lattices, and
	Thompson's group. Does there exist a group of finite type with infinite
	asymptotic dimension?
\end{openproblem}

\plainproblemsection

\begin{openproblem}
	Let \(G=F/N\) be finitely presented. Using the notation of the supplied
	source, can the number \(d_G(N/[N,N])\) of \(G\)-orbits of two-cells needed
	to kill first homology be strictly smaller than the number \(d_F(N)\)
	needed to kill the fundamental group? For the associated family
	\(G_n=H*_{C^n}H\), with \(C\) finite and nontrivial, does
	\[
	d_{F*F}(N_n)\longrightarrow\infty?
	\]
\end{openproblem}

\plainproblemsection

\begin{openproblem}
	Let \(K\) and \(L\) be simple-homotopy-equivalent finite two-complexes.
	Can one pass from \(K\) to \(L\) by elementary collapses and expansions of
	one- and two-cells together with slides of two-cells? This is open even
	when \(K\) is contractible and \(L\) is a point.
\end{openproblem}

\plainproblemsection

\begin{openproblem}
	Does there exist a finite aspherical two-complex \(X\) such that
	\(\pi_1(X)\) is coherent and \(\chi(X)\geq2\)?
\end{openproblem}

\plainproblemsection

\begin{openproblem}
	Let \(X\) be a nonpositively curved symmetric space. Characterize the
	groups \(\Gamma\) for which the space of conjugacy classes of faithful
	discrete representations
	\[
	\Gamma\longrightarrow\operatorname{Isom}(X)
	\]
	is compact, and those for which it is noncompact.
\end{openproblem}

\plainproblemsection

\begin{openproblem}
	Does there exist a finite three-complex \(X\) which is not homotopy
	equivalent to a two-complex but satisfies
	\[
	H^3(X;\mathcal G)=0
	\]
	for every local coefficient system \(\mathcal G\)?
\end{openproblem}

\plainproblemsection

\begin{openproblem}
	Describe the quasi-isometry group \(\operatorname{QI}(\mathbb R^n)\).
	Classify the quasi-isometries of the three-dimensional group
	\(\mathrm{Sol}\) and of the Gromov--Thurston negatively pinched manifolds.
\end{openproblem}

\plainproblemsection

\begin{openproblem}
	Classify, up to quasi-isometry, the mapping-torus groups of automorphisms
	\(\phi\colon F_n\to F_n\). Solve the analogous problem for automorphisms of
	\(\mathbb Z^n\).
\end{openproblem}

\plainproblemsection

\begin{openproblem}
	For which finitely generated groups \(G\) is \(G\times\mathbb Z\)
	quasi-isometric to \(G\)? In particular, is this true for Thompson's group?
\end{openproblem}

\plainproblemsection

\begin{openproblem}
	On a nonorientable compact hyperbolic surface, can every measured geodesic
	lamination with two-sided leaves be approximated by a simplicial measured
	lamination with two-sided leaves? More generally, describe the mapping-
	class-group action on projective measured laminations outside the locus
	containing one-sided leaves.
\end{openproblem}

\plainproblemsection

\begin{openproblem}
	Does \(\operatorname{Out}(F_n)\) have the congruence subgroup property?
\end{openproblem}

\plainproblemsection

\begin{openproblem}
	Does there exist a faithful representation
	\[
	\pi_1(\Sigma_g)\longrightarrow\operatorname{MCG}(\Sigma_h)
	\]
	whose nonidentity elements are all pseudo-Anosov? Equivalently in bundle
	form, does there exist a surface bundle over a surface whose fundamental
	group is word-hyperbolic, or whose total space is hyperbolic?
\end{openproblem}

\plainproblemsection

\begin{openproblem}
	If the fixed subgroup of an automorphism
	\(\alpha\in\operatorname{Aut}(F_n)\) is cyclic, bound the word length of a
	generator of \(\operatorname{Fix}(\alpha)\) in terms of the complexity of
	\(\alpha\), or of a relative train-track representative.
\end{openproblem}

\plainproblemsection

\begin{openproblem}
	Which automorphisms of \(F_n\) preserve a right-invariant order? Apart from
	the existence of periodic elements, determine all obstructions.
\end{openproblem}

\plainproblemsection

\begin{openproblem}
	Does \(\operatorname{Out}(F_n)\), for \(n>2\), contain a right-orderable
	subgroup of finite index?
\end{openproblem}

\plainproblemsection

\begin{openproblem}
	Is every residually torsion-free nilpotent group Haagerup, equivalently
	\(a\)-\(T\)-menable?
\end{openproblem}

\plainproblemsection

\begin{openproblem}
	Can a mapping class group, or a finite-index subgroup of one, embed in
	\(\operatorname{Out}(F_n)\)? Can it embed in the mapping class group of a
	punctured surface?
\end{openproblem}

\plainproblemsection

\begin{openproblem}
	Describe the dynamics of the natural action of
	\(\operatorname{Out}(F_n)\) on the quotient, by inner automorphisms, of the
	space of triples of distinct ends of \(F_n\).
\end{openproblem}

\plainproblemsection

\begin{openproblem}
	Does there exist a finitely generated free subgroup of a mapping class
	group whose nonidentity elements are pseudo-Anosov but which is not
	Schottky? Does there exist a nonfree subgroup all of whose nonidentity
	elements are pseudo-Anosov? Formulate and answer the corresponding
	questions for \(\operatorname{Out}(F_n)\).
\end{openproblem}

\plainproblemsection

\begin{openproblem}
	Does every automorphism \(h\colon F_n\to F_n\) preserve a finite-index
	subgroup \(K\) such that the induced action on
	\(K/[K,K]\) has a root of unity as an eigenvalue?
\end{openproblem}

\plainproblemsection

\begin{openproblem}
	Does every closed three-manifold fibering over \(\mathbb S^1\) with fiber
	of genus at least \(2\) have a finite cover with first Betti number greater
	than \(1\)?
\end{openproblem}

\plainproblemsection

\begin{openproblem}
	For \(g\geq3\), is the natural map
	\[
	\operatorname{MCG}(\Sigma_g)
	\longrightarrow
	\operatorname{QI}\bigl(\operatorname{MCG}(\Sigma_g)\bigr)
	\]
	an isomorphism? Must every quasi-isometry of the mapping class group send
	maximal flats to maximal flats? Answer the analogous questions for
	\(\operatorname{Out}(F_n)\).
\end{openproblem}

\plainproblemsection

\begin{openproblem}
	For which triples \((g,b,n)\) is the mapping class group
	\(\operatorname{MCG}(\Sigma_{g,b,n})\) linear?
\end{openproblem}

\plainproblemsection

\begin{openproblem}
	Let \(K\) be a finite complex with amenable fundamental group. Is there a
	uniform bound on the Betti numbers of all finite covers of \(K\)?
\end{openproblem}

\plainproblemsection

\begin{openproblem}
	For every finite graph \(G\), is it true that \(G\) is planar, or a double
	cover of \(G\) is planar, or no finite cover of \(G\) is planar?
\end{openproblem}

\plainproblemsection

\begin{openproblem}
	Does there exist a finitely presented slender group which is not
	polycyclic-by-finite?
\end{openproblem}

\plainproblemsection

\begin{openproblem}
	Is \(\operatorname{SL}_2(K)\) finitely generated when
	\[
	K=\mathbb Z[X,X^{-1}]
	\quad\text{or}\quad
	K=F[X,X^{-1},Y,Y^{-1}]
	\]
	for a field \(F\)?
\end{openproblem}

\plainproblemsection

\begin{openproblem}
	Let \(L\) be a flag triangulation of \(\mathbb S^{2k-1}\), let \(f_i\) be
	the number of \(i\)-simplices of \(L\), and set
	\[
	\chi
	=
	1-\frac12f_0+\frac14f_1-\frac18f_2+\cdots.
	\]
	Prove
	\[
	(-1)^k\chi\geq0.
	\]
\end{openproblem}

\plainproblemsection

\begin{openproblem}
	Does there exist a group with a balanced presentation, trivial first
	homology, and unsolvable word problem? Is there a sequence of perfect
	groups with balanced presentations among which the trivial group cannot
	be recognized?
\end{openproblem}

\plainproblemsection

\begin{openproblem}
	Let \(G\) be a one-relator group with cyclically reduced relator \(W\), and
	let \(P\) be the submonoid generated by all prefixes of \(W\). Is the
	membership problem for \(P\) in \(G\) decidable?
\end{openproblem}

\plainproblemsection

\begin{openproblem}
	Construct finitely presented infinite groups, or prove their nonexistence,
	with each of the following properties:
	
	\begin{enumerate}[label=(\alph*)]
		\item nonamenable and containing no nonabelian free subgroup;
		\item satisfying a bounded-exponent law \(x^n=1\);
		\item every proper subgroup cyclic;
		\item every proper subgroup finite cyclic of one fixed prime order;
		\item every element has roots of every degree;
		\item only finitely many conjugacy classes;
		\item intermediate growth.
	\end{enumerate}
	
	Examples are known only for the first type.
\end{openproblem}

\plainproblemsection

\begin{openproblem}
	Do there exist finitely generated, non-virtually-cyclic, amenable groups
	of each of the monster types consisting of bounded torsion, cyclic proper
	subgroups, fixed-prime-order proper subgroups, divisible elements, or
	finitely many conjugacy classes?
\end{openproblem}

\plainproblemsection

\begin{openproblem}
	Assuming the continuum hypothesis, does there exist a finitely presented
	group with continuum many nonhomeomorphic asymptotic cones?
\end{openproblem}

\plainproblemsection

\begin{openproblem}
	Does there exist a CAT(0) group with two nonhomeomorphic asymptotic cones?
\end{openproblem}

\plainproblemsection

\begin{openproblem}
	How many nonhomeomorphic asymptotic cones do Thompson's group \(F\), the
	Grigorchuk groups of intermediate growth, and other self-similar groups
	have?
\end{openproblem}

\plainproblemsection

\begin{openproblem}
	Does there exist a finitely presented group all of whose asymptotic cones
	are locally isometric while not all are globally isometric?
\end{openproblem}

\plainproblemsection

\begin{openproblem}
	Does there exist a non-virtually-cyclic amenable group all of whose
	asymptotic cones are locally isometric to an \(\mathbb R\)-tree?
\end{openproblem}

\plainproblemsection

\begin{openproblem}
	Does there exist a finitely generated group \(G\) with an asymptotic cone
	whose fundamental group is countable, or finitely generated, and
	nonabelian? Can such a fundamental group be finite and nonabelian?
\end{openproblem}

\plainproblemsection

\begin{openproblem}
	Extend the structure theory of subgroups acting on tree-graded asymptotic
	cones to subgroups of further classes of groups with tree-graded
	asymptotic cones. In particular, describe the pieces, their stabilizers,
	and the induced action on the associated transversal tree.
\end{openproblem}

\plainproblemsection

\begin{openproblem}
	Is the growth of every non-elementary lacunary hyperbolic group
	exponential? Is it uniformly exponential?
\end{openproblem}

\plainproblemsection

\begin{openproblem}
	Can a finitely generated non-virtually-cyclic subgroup of exponential
	growth in a lacunary hyperbolic group satisfy a nontrivial law? A negative
	answer is known under an additional injectivity-radius hypothesis.
\end{openproblem}

\plainproblemsection

\begin{openproblem}
	Prove an analogue of the Bestvina--Feighn combination theorem for
	lacunary hyperbolic groups.
\end{openproblem}

\plainproblemsection

\begin{openproblem}
	Is every linear lacunary hyperbolic group hyperbolic?
\end{openproblem}

\plainproblemsection

\begin{openproblem}
	Let \(\phi\colon F_n\to F_n\) be injective and nonsurjective. Can the
	ascending HNN extension \(\operatorname{HNN}_\phi(F_n)\) be both
	hyperbolic and nonlinear? In particular, is
	\[
	\langle x,y,t\mid txt^{-1}=xy,\;tyt^{-1}=yx\rangle
	\]
	linear?
\end{openproblem}

\plainproblemsection

\begin{openproblem}
	Do there exist non-residually-finite hyperbolic groups obtained as multiple
	ascending HNN extensions of free groups? In particular, is the group
	\[
	\left\langle
	x,y,t,u
	\ \middle|\
	txt^{-1}=xy,\quad
	tyt^{-1}=yx,\quad
	uxu^{-1}=[x,y],\quad
	uyu^{-1}=[x^2y,y^2x]
	\right\rangle
	\]
	residually finite?
\end{openproblem}

\plainproblemsection

\begin{openproblem}
	Is every subgroup of Thompson's group \(F\) either elementary amenable or
	large enough to contain a copy of \(F\)?
\end{openproblem}

\plainproblemsection

\begin{openproblem}
	Is Thompson's group \(F\) automatic?
\end{openproblem}

\plainproblemsection

\begin{openproblem}
	Is Thompson's group \(F\) amenable?
\end{openproblem}

\plainproblemsection

\begin{openproblem}
	Does Thompson's group \(F\) admit an equivariant or nonequivariant coarse
	embedding into Hilbert space with compression function \(\rho\) satisfying
	\[
	\lim_{n\to\infty}\frac{\rho(n)}{\sqrt n}=\infty?
	\]
\end{openproblem}

\plainproblemsection

\begin{openproblem}
	Does there exist a finitely generated group whose Hilbert-space
	compression gap has size smaller than \(\log x\), for example of order
	\(\log\log x\)?
\end{openproblem}

\plainproblemsection

\begin{openproblem}
	Does every finitely generated diagram group have the Burillo property;
	that is, are its word metric and diagram metric quasi-isometric?
\end{openproblem}

\plainproblemsection

\begin{openproblem}
	Does Thompson's group \(F\) contain a distorted copy of itself? Does it
	contain a finitely generated subgroup with superpolynomial, or merely
	superquadratic, distortion? Does it contain a finitely generated subgroup
	with nonrecursive distortion?
\end{openproblem}

\plainproblemsection

\begin{openproblem}
	Is there an algorithm which, given a finite graph \(K\), decides whether
	the right-angled Artin group \(A(K)\) contains the fundamental group of a
	closed hyperbolic surface?
\end{openproblem}

\plainproblemsection

\begin{openproblem}
	Is there an algorithm which, given a \(K\)-dissection diagram on a surface
	\(S\), decides whether the induced homomorphism
	\[
	\pi_1(S)\longrightarrow A(K)
	\]
	is injective?
\end{openproblem}

\plainproblemsection

\begin{openproblem}
	Does the group
	\[
	\langle x,y,t\mid txt^{-1}=xy,\;tyt^{-1}=yx\rangle
	\]
	contain the fundamental group of a closed hyperbolic surface?
\end{openproblem}

\plainproblemsection

\begin{openproblem}
	Characterize the injective nonsurjective endomorphisms
	\(\phi\colon F_k\to F_k\) for which the ascending HNN extension
	\(\operatorname{HNN}_\phi(F_k)\) contains a closed hyperbolic surface
	subgroup.
\end{openproblem}

\plainproblemsection

\begin{openproblem}
	For bond percolation on a Cayley graph, let \(p_c\) and \(p_u\) denote the
	critical and uniqueness probabilities. Is
	\[
	p_c\neq p_u
	\]
	for every Cayley graph of a nonamenable group?
\end{openproblem}

\plainproblemsection

\begin{openproblem}
	If a finitely generated group has superlinear growth, must the critical
	percolation probability of every Cayley graph satisfy \(p_c<1\)? Determine
	also whether the properties \(p_c=1\) and \(p_c=p_u\) are invariant under
	change of generators or quasi-isometry.
\end{openproblem}

\plainproblemsection

\begin{openproblem}
	Determine \(p_c\), \(p_u\), and the exponential-decay threshold
	\(p_{\mathrm{exp}}\) for Cayley graphs of
	\(\operatorname{SL}_2(\mathbb Z)\), Thompson's group \(F\), Grigorchuk
	groups, and other automata groups.
\end{openproblem}

\plainproblemsection

\begin{openproblem}
	For a hyperbolic group and supercritical bond percolation \(p>p_c\), does
	an infinite open cluster almost surely separate the Cayley graph?
\end{openproblem}

\plainproblemsection

\begin{openproblem}
	What growth rates can occur for the expected distortion of an open
	percolation cluster in a Cayley graph? Answer the same question
	specifically for hyperbolic groups.
\end{openproblem}

\plainproblemsection

\begin{openproblem}
	Do Cayley graphs with isometric asymptotic cones have the same
	percolation critical exponents? In particular, does every Cayley graph of
	a non-elementary hyperbolic group have the mean-field critical exponents
	\[
	\beta=1,\qquad\gamma=1,\qquad\delta=2?
	\]
\end{openproblem}

\plainproblemsection

\begin{openproblem}
	Characterize the compacta which occur as boundaries of hyperbolic groups.
	In particular, for which \(k\) does a \(k\)-dimensional stable Menger space
	occur as such a boundary?
\end{openproblem}

\plainproblemsection

\begin{openproblem}
	Can Osajda's construction of stable Menger boundaries from thick
	right-angled hyperbolic buildings be extended after removing the
	right-angled hypothesis?
\end{openproblem}

\plainproblemsection

\begin{openproblem}
	Which two-dimensional spaces occur as boundaries of hyperbolic groups?
	Study in particular groups with no virtual splitting and boundaries with
	no local cut points, cut arcs, or separating Cantor sets.
\end{openproblem}

\plainproblemsection

\begin{openproblem}
	Does there exist a torsion-free hyperbolic group \(G\) such that
	\[
	\frac{\operatorname{cd}_{\mathbb Q}G}
	{\operatorname{cd}_{\mathbb Z}G}
	<\frac23?
	\]
\end{openproblem}

\plainproblemsection

\begin{openproblem}
	Describe the ideal boundaries produced by the strict hyperbolization
	construction of Charney and Davis.
\end{openproblem}

\plainproblemsection

\begin{openproblem}
	Does there exist a relatively hyperbolic group whose parabolic subgroups
	are nilpotent of class at least \(3\) and whose Bowditch boundary is a
	sphere?
\end{openproblem}

\plainproblemsection

\begin{openproblem}
	Let a group \(G\) act topologically transitively as a convergence group on
	a compact metrizable space \(Z\). Does there exist a Gromov-hyperbolic
	space \(X\) with \(\partial X\cong Z\) such that the action extends to a
	uniformly quasi-isometric quasi-action on \(X\)?
\end{openproblem}

\plainproblemsection

\begin{openproblem}
	Find topological restrictions on the ideal boundaries of CAT\((-1)\)
	cubical complexes.
\end{openproblem}

\plainproblemsection

\begin{openproblem}
	Do isomorphic Coxeter groups necessarily have homeomorphic boundaries?
\end{openproblem}

\plainproblemsection

\begin{openproblem}
	Does there exist a Coxeter group with an \(n\)-dimensional boundary whose
	rational homological dimension is \(1\)?
\end{openproblem}

\plainproblemsection

\begin{openproblem}
	Give conditions on a Coxeter diagram which ensure that the boundary is
	\(n\)-connected and locally \(n\)-connected.
\end{openproblem}

\plainproblemsection

\begin{openproblem}
	Can exotic homology manifolds occur as ideal boundaries of Coxeter groups?
\end{openproblem}

\plainproblemsection

\begin{openproblem}
	Find further universality phenomena for boundaries of hyperbolic groups
	and hyperbolic buildings.
\end{openproblem}

\plainproblemsection

\begin{openproblem}
	Let \(C(N,\Delta)\) be the hyperbolic Coxeter construction associated with
	a flag triangulation \(\Delta\) of a closed manifold \(N\). Is
	\(\partial C(N,\Delta)\) a topological invariant of \(N\)?
\end{openproblem}

\plainproblemsection

\begin{openproblem}
	Suppose closed three-manifolds \(N_1,N_2\), equipped with flag
	triangulations, satisfy
	\[
	\partial C(N_1,\Delta_1)\cong
	\partial C(N_2,\Delta_2).
	\]
	Must every prime connected-sum summand of either \(N_i\) appear as a
	summand of the other? Determine the higher-dimensional analogue.
\end{openproblem}

\plainproblemsection

\begin{openproblem}
	Prove that a compactum which is a \(Z\)-boundary of a group cannot be a
	Boltyansky compactum, extending the result recorded in the source for
	Markov compacta with isomorphic building blocks.
\end{openproblem}

\plainproblemsection

\begin{openproblem}
	For one-ended CAT(0) groups, does local connectedness of the visual
	boundary imply uniqueness of that boundary?
\end{openproblem}

\plainproblemsection

\begin{openproblem}
	If a CAT(0) group \(G\) does not split over a small subgroup, must its
	visual boundary be unique?
\end{openproblem}

\plainproblemsection

\begin{openproblem}
	If a group \(G\) acts geometrically on CAT(0) cube complexes \(X_1\) and
	\(X_2\), must \(\partial X_1\) and \(\partial X_2\) be homeomorphic?
\end{openproblem}

\plainproblemsection

\begin{openproblem}
	Which topological invariants distinguish the different CAT(0) boundaries
	of a group? Which topological, shape-theoretic, or coarser properties of
	these boundaries are invariant under quasi-isometry?
\end{openproblem}

\plainproblemsection

\begin{openproblem}
	Let \(G\) act geometrically on CAT(0) spaces \(X_1\) and \(X_2\). Does the
	diagonal action on \(X_1\times X_2\) admit a convex core? If so, can a
	common cell-like model for \(\partial X_1\) and \(\partial X_2\) be chosen
	as the boundary of that core?
\end{openproblem}

\plainproblemsection

\begin{openproblem}
	Define a topology on the set of quasi-geodesics in a proper or cocompact
	CAT(0) space which:
	
	\begin{enumerate}[label=(\alph*)]
		\item is an increasing union of compact metrizable spaces;
		\item contains the visual boundary;
		\item is invariant under quasi-isometry;
		\item carries natural measure classes which are quasi-preserved.
	\end{enumerate}
\end{openproblem}

\plainproblemsection

\begin{openproblem}
	Prove Swenson's conjecture that a CAT(0) group has no infinite torsion
	subgroup.
\end{openproblem}

\plainproblemsection

\begin{openproblem}
	Let \(G\) be hyperbolic and let \(\partial_{EZ}G\) be an equivariant
	\(EZ\)-structure boundary. Is \(\partial_{EZ}G\) equivariantly homeomorphic
	to the Gromov boundary \(\partial G\)?
\end{openproblem}

\plainproblemsection

\begin{openproblem}
	Can a group admit two \(Z\)-structure boundaries which are not cell-like
	equivalent? Is the property of splitting over a two-ended subgroup
	determined by a Bestvina boundary?
\end{openproblem}

\plainproblemsection

\begin{openproblem}
	Give a complete characterization of the compact metrizable spaces which
	occur as boundaries of proper cocompact CAT(0) spaces.
\end{openproblem}

\plainproblemsection

\begin{openproblem}
	Extend to arbitrary finitely generated one-ended groups the following
	quasi-isometry rigidity properties known in the finitely presented case:
	quasi-isometric invariance of the JSJ decomposition, the characterization
	of coarse separation by quasilines, and nonseparation by quasi-rays.
\end{openproblem}

\plainproblemsection

\begin{openproblem}
	Are splittings over \(\mathbb Z^n\), and the corresponding JSJ
	decompositions, invariant under quasi-isometry?
\end{openproblem}

\plainproblemsection

\begin{openproblem}
	If the Cayley graph of a finitely generated group is separated by a
	sequence of quasicircles, must the group be virtually a surface group?
\end{openproblem}

\plainproblemsection

\begin{openproblem}
	Let \(G\) be finitely generated with
	\(\operatorname{asdim}G\geq n\). If a subset \(X\) of its Cayley graph has
	\(\operatorname{asdim}X\leq n-2\) and coarsely separates the graph, must
	\(G\) split over a subgroup \(H\) satisfying
	\(\operatorname{asdim}H\leq n-1\)?
\end{openproblem}

\plainproblemsection

\begin{openproblem}
	Can every quasiconformal homeomorphism of \(\mathbb R^4\) be approximated
	by diffeomorphisms? Affirmative results are known in dimensions \(2,3\)
	and at least \(5\), leaving dimension \(4\).
\end{openproblem}

\plainproblemsection

\begin{openproblem}
	Let \(f\colon B^n\to B^n\) be quasiconformal. Can \(f\) be approximated by
	globally quasiconformal diffeomorphisms \(f_j\colon B^n\to B^n\)? Can the
	maps \(f_j\) be chosen uniformly \(K\)-quasiconformal?
\end{openproblem}

\plainproblemsection

\begin{openproblem}
	Find broad classes of metric spaces on which the infinitesimal metric
	condition for quasiconformality implies the corresponding local condition.
\end{openproblem}

\plainproblemsection

\begin{openproblem}
	Which boundaries, besides those of Fuchsian buildings, have the Loewner
	property?
\end{openproblem}

\plainproblemsection

\begin{openproblem}
	Does a nonsmoothable closed simply connected four-manifold admit an
	Ahlfors \(4\)-regular, linearly locally contractible metric?
\end{openproblem}

\plainproblemsection

\begin{openproblem}
	Develop an analytic theory on the ideal boundaries of relatively
	hyperbolic groups analogous to the existing theory for hyperbolic groups.
\end{openproblem}

\plainproblemsection

\begin{openproblem}
	Determine the general hypotheses under which a quasiconformal
	homeomorphism between metric spaces must be quasisymmetric.
\end{openproblem}

\plainproblemsection

\begin{openproblem}
	Which metric fractals are quasisymmetrically co-Hopfian? Answer this
	question in particular for hyperbolic-group boundaries, round Sierpi\'nski
	carpets, and Menger spaces.
\end{openproblem}

\plainproblemsection

\begin{openproblem}
	Prove Heinonen's conjecture that every Loewner boundary of a hyperbolic
	group is quasisymmetrically co-Hopfian.
\end{openproblem}

\plainproblemsection

\begin{openproblem}
	If a hyperbolic group \(G\) has connected boundary with no local cut
	points, does \(\partial G\) carry a natural measure class invariant under
	all quasisymmetric self-homeomorphisms?
\end{openproblem}

\plainproblemsection

\begin{openproblem}
	Develop analytic structures on visual boundaries of CAT(0) spaces despite
	the absence of a canonical metric. Determine whether the natural
	\(\operatorname{Isom}(X)\)-action behaves regularly with respect to the
	pseudometrics obtained by pushing visual spheres into \(X\).
\end{openproblem}

\plainproblemsection

\begin{openproblem}
	Let \(D\) be a connected hyperbolic-group boundary with no local cut
	points, and suppose \(D\) is not Loewner. Does there exist a nontrivial
	quasisymmetrically invariant closed equivalence relation on \(D\) such that
	the quotient is Hausdorff and is a relative boundary for the group?
\end{openproblem}

\plainproblemsection

\begin{openproblem}
	Clarify the relations among the following notions of conformality on
	hyperbolic-group boundaries: metric \(1\)-quasiconformality, preservation
	of curve modulus, infinitesimally linear quasisymmetry, and preservation
	of a measurable conformal structure arising from a Poincar\'e inequality.
\end{openproblem}

\plainproblemsection

\begin{openproblem}
	For homeomorphisms of Hilbert space, does quasiconformality imply
	quasisymmetry and the property that balls are mapped to quasiballs?
\end{openproblem}

\plainproblemsection

\begin{openproblem}
	Construct exotic quasi-isometric embeddings of quaternionic hyperbolic
	\(n\)-space into quaternionic hyperbolic \(m\)-space with \(m\) as small as
	possible, in particular \(m=n+1\). Answer the analogous question for
	complex hyperbolic spaces.
\end{openproblem}

\plainproblemsection

\begin{openproblem}
	Write the solvable models of two higher-rank symmetric spaces as
	\(A_1N_1\) and \(A_2N_2\). Are all quasi-isometric embeddings between the
	two symmetric spaces either isometries or algebraic embeddings of the form
	\(\phi\colon A_1N_1\longrightarrow A_2N_2\)?
\end{openproblem}

\plainproblemsection

\begin{openproblem}
	Does every hyperbolic group admit a uniformly quasiconformal discrete
	action on some sphere \(\mathbb S^n\)?
\end{openproblem}

\plainproblemsection

\begin{openproblem}
	Prove Cannon's conjecture in either of its equivalent forms: if a
	hyperbolic group \(G\) has \(\partial G\cong\mathbb S^2\), then \(G\) acts
	geometrically on \(\mathbb H^3\), and a visual metric on \(\partial G\) is
	quasisymmetric to the standard metric on \(\mathbb S^2\).
\end{openproblem}

\plainproblemsection

\begin{openproblem}
	Prove Wall's conjecture that every Poincar\'e-duality group of dimension
	\(3\) is the fundamental group of a closed three-manifold.
\end{openproblem}

\plainproblemsection

\begin{openproblem}
	If a hyperbolic group has boundary homeomorphic to the Sierpi\'nski
	carpet, must it act geometrically on a convex subset of \(\mathbb H^3\)?
\end{openproblem}

\plainproblemsection

\begin{openproblem}
	Let \(G\) be hyperbolic relative to virtually \(\mathbb Z^2\) subgroups and
	suppose its Bowditch boundary is \(\mathbb S^2\). Must \(G\) be
	commensurable with the fundamental group of a finite-volume hyperbolic
	three-manifold?
\end{openproblem}

\plainproblemsection

\begin{openproblem}
	If a hyperbolic group has Sierpi\'nski-carpet boundary, is every visual
	metric quasisymmetric to a round Sierpi\'nski-carpet metric?
\end{openproblem}

\plainproblemsection

\begin{openproblem}
	Prove Cannon's conjecture under substantial additional hypotheses, for
	example when the group is a three-manifold group, a \(PD_3\) group
	splitting over a surface group, or a group acting on a CAT(0) cube complex.
\end{openproblem}

\plainproblemsection

\begin{openproblem}
	Derive the relative Cannon conjecture from the absolute conjecture by
	hyperbolic Dehn filling. In particular, prove that sufficiently long
	fillings are \(PD_3\) groups and that three-manifold realizations of all
	such fillings imply a three-manifold realization of the original
	relatively hyperbolic group.
\end{openproblem}

\plainproblemsection

\begin{openproblem}
	For a hyperbolic group \(G\), prove the conjectural formula
	\[
	\operatorname{ACD}(\partial G)
	=
	\inf_{G\curvearrowright X}
	\dim_H(\partial X,d_{\mathrm{vis}}),
	\]
	where the infimum runs over geometric actions on hyperbolic metric spaces.
	If the infimum is attained, must it be attained by a visual metric?
\end{openproblem}

\plainproblemsection

\begin{openproblem}
	Compute the Ahlfors-regular conformal dimension of the standard
	Sierpi\'nski carpet.
\end{openproblem}

\plainproblemsection

\begin{openproblem}
	Find hypotheses ensuring that a \(Q\)-Loewner boundary of a hyperbolic
	group satisfies a \(1\)-Poincar\'e inequality.
\end{openproblem}

\plainproblemsection

\begin{openproblem}
	For the convergence action on the quotient sphere obtained from a
	surface-group-invariant filling lamination, is the action topologically
	conjugate to a conformal action? Equivalently, does there exist an
	equivariant continuous Cannon--Thurston map
	\[
	\mathbb S^1\longrightarrow\mathbb S^2?
	\]
\end{openproblem}

\plainproblemsection

\begin{openproblem}
	Determine the Poisson boundary of a free group equipped with an arbitrary
	probability measure.
\end{openproblem}

\plainproblemsection

\begin{openproblem}
	Let \((x_n)\) be a random walk on a proper metric space and let
	\[
	A=\lim_{n\to\infty}\frac{d(x_0,x_n)}{n}.
	\]
	Is there almost surely a horofunction \(h\) such that
	\[
	\lim_{n\to\infty}-\frac{h(x_n)}{n}=A?
	\]
\end{openproblem}

\plainproblemsection

\begin{openproblem}
	The horofunction compactification of a proper metric space carries an
	incidence geometry at infinity built from halfspaces and stars. Under what
	conditions do homomorphisms, or quasi-isometries, induce
	incidence-preserving maps between these geometries?
\end{openproblem}

\plainproblemsection

\begin{openproblem}
	For the compactification of a finitely generated group constructed from
	its first \(\ell^2\)-cohomology, or another function-space cohomology,
	is the associated incidence geometry at infinity always trivial in the
	sense described in the source?
\end{openproblem}

\plainproblemsection

\begin{openproblem}
	Characterize the stationary hitting measure on the space of projective
	measured foliations produced by a random walk on a mapping class group.
	Is it absolutely continuous with respect to the visual Lebesgue measure?
\end{openproblem}

\plainproblemsection

\begin{openproblem}
	Determine the Poisson boundary of Culler--Vogtmann Outer space.
\end{openproblem}

\plainproblemsection

\begin{openproblem}
	Develop a meaningful structure theory and a useful boundary theory for
	lacunary hyperbolic groups.
\end{openproblem}

\plainproblemsection

\begin{openproblem}
	All asymptotic cones of a relatively hyperbolic group have cut points. To
	what extent does the converse hold? Characterize groups whose asymptotic
	cones have cut points in terms of relative hyperbolicity or an appropriate
	weaker structure.
\end{openproblem}

\plainproblemsection

\begin{openproblem}
	Develop analytic tools for asymptotic cones, beyond their topology and
	bilipschitz geometry.
\end{openproblem}

\plainproblemsection

\begin{openproblem}
	Let \(G\) be the fundamental group of a closed hyperbolic \(n\)-manifold
	and consider a central extension
	\[
	1\longrightarrow\mathbb Z_p\longrightarrow\Gamma
	\longrightarrow G\longrightarrow1.
	\]
	Is \(\Gamma\) residually finite? Equivalently, does some finite-index
	\(G'<G\) make the restriction
	\[
	H^2(G;\mathbb Z_p)\longrightarrow H^2(G';\mathbb Z_p)
	\]
	zero? Answer the analogous question for complex- and quaternionic-
	hyperbolic manifold groups.
\end{openproblem}

\plainproblemsection

\begin{openproblem}
	For a closed hyperbolic manifold group \(G\), does there exist a
	finite-index subgroup \(G'<G\) for which
	\[
	H^3(G;\mathbb Z_2)\longrightarrow H^3(G';\mathbb Z_2)
	\]
	is zero?
\end{openproblem}

\plainproblemsection

\begin{openproblem}
	Does every Gromov-hyperbolic Coxeter group admit a discrete embedding in
	\(\operatorname{Isom}(\mathbb H^n)\) for all sufficiently large \(n\),
	without requiring Coxeter generators to act as reflections?
\end{openproblem}

\plainproblemsection

\begin{openproblem}
	Can a convex-cocompact subgroup of \(\operatorname{PU}(2,1)\) have limit
	set homeomorphic to the Sierpi\'nski carpet?
\end{openproblem}

\plainproblemsection

\begin{openproblem}
	Let \(G<\operatorname{Isom}(\mathbb H^n)\) be discrete, torsion-free,
	finitely generated, and contain no abelian subgroup of rank at least \(2\).
	For its conical limit set \(\Lambda_c(G)\), determine whether:
	
	\begin{enumerate}[label=(\alph*)]
		\item \(\operatorname{cd}_{\mathbb Z}G
		\leq\dim_H\Lambda_c(G)+1\);
		\item equality forces the limit set of \(G\) to be a round sphere;
		\item \(\dim_H\Lambda_c(G)<2\) forces geometric finiteness;
		\item \(\dim_H\Lambda_c(G)<1\) forces \(G\) to be a classical
		Schottky-type group.
	\end{enumerate}
\end{openproblem}

\plainproblemsection

\begin{openproblem}
	Let a finitely generated discrete group act by isometries on a
	Gromov-hyperbolic space, with connected limit set. Must that limit set be
	locally connected?
\end{openproblem}

\plainproblemsection

\begin{openproblem}
	Let
	\(\rho_1,\rho_2\colon G\to\operatorname{Isom}(\mathbb H^n)\) be discrete
	faithful representations and suppose their translation-length functions
	satisfy
	\[
	C^{-1}\leq
	\frac{\ell_{\rho_1}(g)}{\ell_{\rho_2}(g)}
	\leq C
	\qquad(g\in G).
	\]
	Is there a \(\rho_2\circ\rho_1^{-1}\)-equivariant quasiconformal map
	between the two limit sets? Can its dilatation be bounded by a function of
	\(C\)?
\end{openproblem}

\plainproblemsection

\begin{openproblem}
	Give a constructive proof of Rips compactness for discrete faithful
	representation spaces. More precisely, for a finitely presented,
	non-virtually-abelian group \(G\) which does not split over a virtually
	abelian subgroup, find an explicit bound, in terms of the presentation and
	\(n\), for
	\[
	\inf_{x\in\mathbb H^n}
	\max_i d\bigl(x,\rho(g_i)x\bigr)
	\]
	over all \([\rho]\) in the discrete faithful representation space
	\(D_n(G)\).
\end{openproblem}

\plainproblemsection

\begin{openproblem}
	Find new algebraic restrictions on finitely generated abstract Kleinian
	groups beyond the Haagerup and K\"ahler-group restrictions recorded in the
	source.
\end{openproblem}

\plainproblemsection

\begin{openproblem}
	Prove that every arithmetic lattice in
	\(\operatorname{Isom}(\mathbb H^n)\), for \(n\geq4\), is noncoherent.
\end{openproblem}

\plainproblemsection

\begin{openproblem}
	Let \(G<\operatorname{Isom}(\mathbb H_{\mathbb H}^n)\) be discrete and
	have Kazhdan's property \(T\). Must \(G\) preserve a totally geodesic
	subspace and act on that subspace as a lattice?
\end{openproblem}

\plainproblemsection

\begin{openproblem}
	Let \(\Delta\) be a developable triangle of groups whose cell groups have
	property \(T\), and suppose every link in the universal cover has first
	positive eigenvalue greater than \(1/2\). Must \(\pi_1(\Delta)\) have
	property \(T\)?
\end{openproblem}

\plainproblemsection

\begin{openproblem}
	Generalize the Bestvina--Feighn combination theorem from graphs of groups
	to complexes of groups.
\end{openproblem}

\plainproblemsection

\begin{openproblem}
	Prove a complex-hyperbolic analogue of Vinberg's finiteness theorem: does
	there exist \(N\) such that, for \(n\geq N\), no lattice in
	\(\operatorname{PU}(n,1)\) is generated by complex reflections?
\end{openproblem}

\plainproblemsection

\begin{openproblem}
	Develop a satisfactory theory of geometric finiteness for group actions on
	CAT(0) spaces, including a boundary-dynamical characterization of
	undistorted subgroups.
\end{openproblem}

\plainproblemsection

\begin{openproblem}
	Extend the relation between Anosov representations and dynamics on limit
	sets from surface groups to representations of general hyperbolic groups.
\end{openproblem}

\plainproblemsection

\begin{openproblem}
	Generalize holomorphic chain patterns on
	\(\partial\mathbb H_{\mathbb C}^n\) so as to prove rigidity for embeddings
	of lattices in \(\operatorname{PU}(n,1)\) into other higher-rank Lie
	groups.
\end{openproblem}

\plainproblemsection

\begin{openproblem}
	Obtain rigidity theorems for embeddings of real-hyperbolic lattices into
	higher-rank semisimple Lie groups using their boundary maps.
\end{openproblem}

\plainproblemsection

\begin{openproblem}
	Let \(X\) be a compact polyhedron and \(G\) a group of simple
	homotopy-equivalences of \(X\). In the cases not excluded by the
	counterexamples recorded in the source, does there exist a compact
	\(X'\simeq X\) on which \(G\) is realized by homeomorphisms?
\end{openproblem}

\plainproblemsection

\begin{openproblem}
	Is contractibility algorithmically decidable for finite two- or
	three-dimensional cell complexes? It is undecidable in dimensions at
	least \(4\).
\end{openproblem}

\plainproblemsection

\begin{openproblem}
	Let \(G\) be word-hyperbolic and assume it does not split over any
	virtually cyclic subgroup. Can an infinite-index subgroup of \(G\) be
	isomorphic to a finite-index subgroup?
\end{openproblem}

\plainproblemsection

\begin{openproblem}
	Formulate and prove a group-free analogue of McMullen's inflexibility
	theorem for closed connected subsets of the boundary
	\(\mathbb S^2=\partial\mathbb H^3\).
\end{openproblem}

\plainproblemsection

\begin{openproblem}
	Let \(\Lambda\subset\mathbb S^2\) be a quasicircle, let \(K\) be its convex
	hull in \(\mathbb H^3\), and let \(Q_\Lambda\) be the
	Thurston--Reimann bilipschitz extension associated with the random
	Beltrami construction in the source. Is its local bilipschitz constant at
	\(p\in\operatorname{int}K\) bounded by
	\[
	L(Q_\Lambda,p)
	\leq
	1+C_1e^{-C_2d(p,\partial_\infty K)}
	\]
	for universal positive constants \(C_1,C_2\)?
\end{openproblem}

\plainproblemsection

\begin{openproblem}
	Are braid groups CAT(0)?
\end{openproblem}

\plainproblemsection

\begin{openproblem}
	Extend Rips' theory of group actions on \(\mathbb R\)-trees to
	higher-dimensional buildings, including products of \(\mathbb R\)-trees.
\end{openproblem}

\plainproblemsection

\begin{openproblem}
	Let \(Y\) be a compact finite-dimensional locally CAT(0) metric space of
	rank \(n\geq2\). Must \(\widetilde Y\) either split nontrivially as a
	metric product or be a Euclidean building? This is known in dimensions
	\(2\) and \(3\).
\end{openproblem}

\plainproblemsection

\begin{openproblem}
	Compute cogrowth for significant subgroup inclusions. In particular:
	
	\begin{enumerate}[label=(\alph*)]
		\item prove exponential cogrowth of
		\(\operatorname{SL}_n(\mathbb Z)\) in
		\(\operatorname{SL}_{n+1}(\mathbb Z)\);
		\item compute cogrowth of special subgroups of Coxeter groups;
		\item if the Schreier graph \(\Gamma_G/H\) is hyperbolic, determine whether
		the cogrowth must be constant, linear, or exponential.
	\end{enumerate}
\end{openproblem}

\plainproblemsection

\begin{openproblem}
	Let \(X\) be a contractible two-complex and let
	\(Y\subset X^{(1)}\) be a connected subgraph with its path metric. Must
	\(Y\) have coarsely trivial homotopy groups \(\pi_m\) for every \(m\geq2\)?
	Does this coarse Whitehead conjecture hold when \(X\) is the Cayley complex
	of a hyperbolic group?
\end{openproblem}

\subsection*{References}

\begin{enumerate}[label={[\arabic*]},leftmargin=*,itemsep=0.4em]
	\item M. Bestvina, \emph{Questions in Geometric Group Theory}, major
	revision August 22, 2000; updated July 2004.
	\item M. Sapir, ``Some group theory problems,'' survey, 2007,
	arXiv:0704.2899.
	\item M. Kapovich, ``Problems on boundaries of groups and Kleinian
	groups,'' problem list, May 25, 2007.
\end{enumerate}

\chapter{Topology}

\clearpage
\section{Knot Theory}

\paragraph{Notation for Laver-table problems.}
For a shelf \((Q,\triangleright)\), write \(H_R^k(Q)\) for its rack
cohomology.  The \(n\)-th Laver table is denoted by \(A_n\).

\plainproblemsection

\begin{openproblem}
	Determine what topological information about arbitrary braids can be
	extracted from colorings by Laver tables weighted by rack \(2\)- or
	\(3\)-cocycles.
\end{openproblem}
% Source: LDT02

\plainproblemsection

\begin{openproblem}
	Formulate a precise framework that unifies the analogies between number
	theory and knot theory, including the analogy between Galois theory and
	Alexander theory.
\end{openproblem}
% Source: LDT02

\paragraph{Notation for Iwasawa-type problems.}
Let \(L\) be a link in a rational homology \(3\)-sphere \(M\), let
\(X=M\setminus\nu L\), and let
\(\sigma\colon\pi_1(X)\twoheadrightarrow\mathbb Z\).
For a prime \(p\), let \(M_{\sigma,p^n}\) be the corresponding
\(p^n\)-fold cyclic branched cover.  When
\(H_1(M_{\sigma,p^n};\mathbb Z)\) is finite for every \(n\), the source
records integers \(\lambda_{L,\sigma},\mu_{L,\sigma},\nu_{L,\sigma}\)
satisfying
\[
v_p\!\left(\#H_1(M_{\sigma,p^n};\mathbb Z)\right)
=\lambda_{L,\sigma}n+\mu_{L,\sigma}p^n+\nu_{L,\sigma}
\]
for all sufficiently large \(n\).

\plainproblemsection

\begin{openproblem}
	Remove the finiteness assumption in the preceding formula by finding a
	\(p\)-adic growth formula for the torsion subgroup of
	\(H_1(M_{\sigma,p^n};\mathbb Z)\).
\end{openproblem}
% Source: LDT02

\plainproblemsection

\begin{openproblem}
	Determine the possible values of
	\(\lambda_{L,\sigma}\), \(\mu_{L,\sigma}\), and \(\nu_{L,\sigma}\).
\end{openproblem}
% Source: LDT02

\plainproblemsection

\begin{openproblem}
	In the analogy between primes in number fields and links, give a precise
	meaning to the ``number of components'' of a prime.
\end{openproblem}
% Source: LDT02

\plainproblemsection

\begin{openproblem}
	Let
	\[
	\theta\colon B_L\hookrightarrow A_L
	\]
	be the natural homomorphism from the link module to the Alexander module,
	and let \(M_L\subset A_L\) be the submodule generated by meridional
	elements.  Determine when
	\[
	Y_L=B_L/\theta^{-1}(M_L)
	\]
	is a pseudonull module over
	\(\Lambda=\mathbb Z[t_1^{\pm1},\ldots,t_r^{\pm1}]\).
\end{openproblem}
% Source: LDT02

\plainproblemsection

\begin{openproblem}
	Develop Iwasawa-type growth formulas for \(p\)-adic Lie towers of branched
	covers over a link, and provide substantial families of examples.
\end{openproblem}
% Source: LDT02

\paragraph{Notation for profinite-knot problems.}
Let \(\mathcal K\) be the monoid of oriented knots and
\(\widehat{\mathcal K}\) the profinite-knot monoid constructed in the
source.  Write
\[
h\colon\mathcal K\longrightarrow\widehat{\mathcal K}
\]
for the natural map and
\(\operatorname{Frac}\widehat{\mathcal K}\) for its fraction group.

\plainproblemsection

\begin{openproblem}
	Is \(h\colon\mathcal K\to\widehat{\mathcal K}\) injective?
\end{openproblem}
% Source: LDT02

\plainproblemsection

\begin{openproblem}
	Establish an analogue of the Alexander--Markov theorem for profinite
	links.
\end{openproblem}
% Source: LDT02

\plainproblemsection

\begin{openproblem}
	Is the action of
	\[
	G_{\mathbb Q}
	=\operatorname{Gal}(\overline{\mathbb Q}/\mathbb Q)
	\]
	on \(\operatorname{Frac}\widehat{\mathcal K}\) faithful?
\end{openproblem}
% Source: LDT02

\plainproblemsection

\begin{openproblem}
	Does there exist a natural homology, cohomology, or fundamental-group
	theory of a pro-variety \(X/\mathbb Q\) whose Galois action realizes the
	action of \(G_{\mathbb Q}\) on
	\(\operatorname{Frac}\widehat{\mathcal K}\)?
\end{openproblem}
% Source: LDT02

\paragraph{Notation for algebraic \(K\)-theory problems.}
For a representation of a knot or link group in
\(\operatorname{SL}_2(F)\), the source associates a value in the
Milnor--Witt group \(K_2^{MW}(F)\), called its \(K_2\)-value.

\plainproblemsection

\begin{openproblem}
	Find applications of the \(K_2\)-value to knot theory or number theory.
\end{openproblem}
% Source: LDT02

\plainproblemsection

\begin{openproblem}
	Give an arithmetic interpretation of the \(K_2\)-value, including for the
	hyperbolic holonomy of the figure-eight knot.
\end{openproblem}
% Source: LDT02

\plainproblemsection

\begin{openproblem}
	Let \(F\) be a global field.  Find a representation
	\[
	\pi_1(S^3\setminus L)\longrightarrow\operatorname{SL}_2(F)
	\]
	whose \(K_2^{MW}(F)\)-value detects nontrivial global information rather
	than information arising only from local fields.
\end{openproblem}
% Source: LDT02

\plainproblemsection

\begin{openproblem}
	For a hyperbolic knot whose holonomy is defined over a quadratic field
	\(F\), determine an arithmetic relation between its \(K_2\)-value and
	geometric invariants such as volume.  Compare the \(K_2\)-value with the
	Birch--Tate conjectural relation
	\[
	\zeta_F(-1)
	\stackrel{?}{=}
	(-1)^{r_1}\frac{\lvert K_2(\mathcal O_F)\rvert}
	{\lvert K_3(\mathcal O_F)\rvert}
	\]
	when \(F\) is totally real.
\end{openproblem}
% Source: LDT02

\plainproblemsection

\begin{openproblem}
	For a knotted surface \(K\subset S^4\) and a representation
	\[
	\pi_1(S^4\setminus K)\longrightarrow\operatorname{SL}_2(F),
	\]
	construct invariants with values in \(K_2(F)\), \(K_3(F)\), or \(K_4(F)\).
\end{openproblem}
% Source: LDT02

\paragraph{Notation for knot-space problems.}
Let \(\mathcal K_{\mathrm{long}}\) be the space of smooth long immersions
\(\mathbb R\to\mathbb R^3\), let \(\Sigma\) be its discriminant, and let
\(H^1_{\mathrm{ft}}\) denote the finite-type part of
\(H^1(\mathcal K_{\mathrm{long}}\setminus\Sigma;\mathbb Q)\).

\plainproblemsection

\begin{openproblem}
	Give an axiomatic description of \(H^1_{\mathrm{ft}}\) analogous to the
	Birman--Lin axiomatization of finite-type knot invariants.
\end{openproblem}
% Source: LDT02

\plainproblemsection

\begin{openproblem}
	Find cycles in \(\mathcal K_{\mathrm{long}}\setminus\Sigma\), other than the rotation,
	dragging, and Hatcher cycles, that are defined independently of the knot.
\end{openproblem}
% Source: LDT02

\plainproblemsection

\begin{openproblem}
	Evaluation on the rotation, dragging, and Hatcher cycles gives maps from
	\(H^1(\mathcal K_{\mathrm{long}}\setminus\Sigma;\mathbb Q)\) to knot invariants.  Does each
	of these maps send a finite-type \(1\)-cocycle to a finite-type knot
	invariant?
\end{openproblem}
% Source: LDT02

\plainproblemsection

\begin{openproblem}
	Does every \(1\)-cocycle admitting the Gauss-diagram description specified
	in the source represent a finite-type cohomology class?
\end{openproblem}
% Source: LDT02

\plainproblemsection

\begin{openproblem}
	Determine the relation between the integral formulas for knot-space
	\(1\)-cocycles and their Gauss-diagram formulas.
\end{openproblem}
% Source: LDT02

\paragraph{Notation for cable-link problems.}
For the \((p,q)\)-cable \(K^{(p,q)}\) of a knot \(K\), let
\(\alpha(K^{(p,q)})\) be its arc index and let
\(\alpha_c(K^{(p,q)})\) be the canonical arc index defined from a grid
diagram in the source.  If \(G\) is a grid diagram of \(K\), write
\(\alpha(G^{(p,q)})\) for the arc number of the induced cable diagram.

\plainproblemsection

\begin{openproblem}
	Is
	\[
	\alpha(K^{(p,q)})=\alpha_c(K^{(p,q)})
	\]
	for every cable link under consideration?
\end{openproblem}
% Source: LDT02

\plainproblemsection

\begin{openproblem}
	If \(G\) and \(G'\) are minimal grid diagrams of the same knot \(K\), must
	\[
	\alpha(G^{(p,q)})=\alpha((G')^{(p,q)})?
	\]
\end{openproblem}
% Source: LDT02

\plainproblemsection

\begin{openproblem}
	If \(G\) is a minimal grid diagram of \(K\), must
	\[
	\alpha_c(K^{(p,q)})=\alpha(G^{(p,q)})?
	\]
\end{openproblem}
% Source: LDT02

\paragraph{Notation for virtual-knot problems.}
For a virtual-knot diagram \(D\), let \(s_n(D)\) be the \(n\)-th state
number defined in the source, and let \(s_n(K)\) be the minimum over all
diagrams representing \(K\).

\plainproblemsection

\begin{openproblem}
	Characterize the virtual-knot diagrams with \(r\) real crossings that
	attain each of the upper bounds
	\[
	s_1(D)=\frac{2^{r+1}-(-1)^{r+1}}{3},\qquad
	s_2(D)=2^{r-1},\qquad
	s_3(D)=3\cdot2^{r-3}.
	\]
\end{openproblem}
% Source: LDT02

\plainproblemsection

\begin{openproblem}
	For \(n\geq2\), find lower bounds for \(s_n(K)\) in terms of algebraic
	invariants such as the Jones or Miyazawa polynomial.
\end{openproblem}
% Source: LDT02

\plainproblemsection

\begin{openproblem}
	Determine \(s_n(K)\), for \(n\geq1\), for the virtual knots in Green's
	table.
\end{openproblem}
% Source: LDT02

\paragraph{Notation for homology-cobordism maps.}
For a prime power \(p^r\), let
\[
\Sigma_{p^r}\colon\mathcal C\longrightarrow\Theta_{\mathbb Q}^{3}
\]
send a knot-concordance class to the rational homology-cobordism class of
its \(p^r\)-fold cyclic branched cover.

\plainproblemsection

\begin{openproblem}
	Is every smoothly slice knot ribbon?
\end{openproblem}
% Source: LDT03

\plainproblemsection

\begin{openproblem}
	Is the smooth knot-concordance group isomorphic to
	\[
	\mathcal C\cong\mathbb Z^{\infty}\oplus
	(\mathbb Z/2\mathbb Z)^{\infty}?
	\]
\end{openproblem}
% Source: LDT03

\plainproblemsection

\begin{openproblem}
	Does Levine's homomorphism from the smooth knot-concordance group to the
	algebraic concordance group split?
\end{openproblem}
% Source: LDT03

\plainproblemsection

\begin{openproblem}
	Determine which other isomorphism types of subgroups or direct summands can
	occur in the kernel and cokernel of
	\(\Sigma_{p^r}\colon\mathcal C\to\Theta_{\mathbb Q}^{3}\).
\end{openproblem}
% Source: LDT03

\plainproblemsection

\begin{openproblem}
	\label{ldt01:prob.VK=1}
	Find a non-trivial knot $K$ with
	$V_K(t)=1$.
\end{openproblem}
% Source: LDT01
\plainproblemsection

\begin{openproblem}
	\label{ldt01:prob.char_VK}
	Characterize those elements of $\Z[t,t^{-1}]$ of
	the form $V_K(t)$.
\end{openproblem}
% Source: LDT01
\plainproblemsection

\begin{openproblem}
	Find a 3-dimensional topological interpretation
	of the Jones polynomial of links.
\end{openproblem}
% Source: LDT01
\plainproblemsection

\begin{openproblem}
	\label{ldt01:prob.roberts1}
	Why is the Jones polynomial
	a polynomial?
\end{openproblem}
% Source: LDT01
\plainproblemsection

\begin{openproblem}
	\label{ldt01:prob.roberts2}
	Is there a relationship between values of
	Jones polynomials
	at roots of unity and branched cyclic coverings of a knot?
\end{openproblem}
% Source: LDT01
\plainproblemsection

\begin{openproblem}
	\label{ldt01:prob.roberts3}
	Is there a relationship between the Jones polynomial
	of a knot and
	the counting of points in varieties defined over
	finite fields?
\end{openproblem}
% Source: LDT01
\plainproblemsection

\begin{openproblem}
	\label{ldt01:prob.roberts4}
	Define the Jones polynomial
	intrinsically using homology of local systems.
\end{openproblem}
% Source: LDT01
\plainproblemsection

\begin{openproblem}
	\label{ldt01:prob.roberts5}
	Study the relation between the Jones polynomial
	and Gromov-Witten theory.
\end{openproblem}
% Source: LDT01
\plainproblemsection

\begin{openproblem}
	\label{ldt01:prob.desc_0_VK}
	Describe the set of zeros of the Jones polynomial
	of all (alternating) knots.
\end{openproblem}
% Source: LDT01
\plainproblemsection

\begin{openproblem}
	\label{ldt01:prob.vol_VK-1}
	Find the relationship between
	the hyperbolic volume
	of knot complements
	and $\log V_K(-1)$
	(resp. $\log V_K(-1)/\log \mbox{\rm deg} V_K(t)$).
\end{openproblem}
% Source: LDT01
\plainproblemsection

\begin{openproblem}
	Understand Khovanov's categorification
	of the Jones polynomial.
\end{openproblem}
% Source: LDT01
\plainproblemsection

\begin{openproblem}
	Categorify other knot polynomials.
\end{openproblem}
% Source: LDT01
\plainproblemsection

\begin{openproblem}
	\label{ldt01:prob.VK_span}
	Does the Jones polynomial
	$V$ admit only finitely many values
	of given span?
	What about the Q polynomial
	or the skein,
	Kauffman polynomials
	(when fixing the span in both variables)?
\end{openproblem}
% Source: LDT01
\plainproblemsection

\begin{openproblem}
	\label{ldt01:prob.QKunit}
	Why are the unit norm complex numbers $\alpha$ for which the value
	$Q_K(\alpha)$
	has maximal norm statistically concentrated around
	$e^{11\pi \sqrt{-1}/25}$\,?
\end{openproblem}
% Source: LDT01
\plainproblemsection

\begin{openproblem}
	\label{ldt01:prob.PWK=FK}
	Let K be a non-trivial knot,
	and let $W_K$ be a Whitehead double of K.
	
	Is then
	$$
	\deg_m P_{W_K}(l,m) = 2\deg_z F_K(a,z) +2 \,?
	$$
\end{openproblem}
% Source: LDT01
\plainproblemsection

\begin{openproblem}
	\label{ldt01:prob.sign_PL_FL}
	Is for any alternating link $L$,
	\[
	\sigma(L)\ge \mbox{min\,deg}_l \big(P_L(l,m) \big)
	\ge \mbox{min\,deg}_a \big( F_L(a^{-1},z) \big)\,?
	\]
\end{openproblem}
% Source: LDT01
\plainproblemsection

\begin{openproblem}
	\label{ldt01:prob.Conway_cr}
	If $\nabla_{k}$ is the coefficient of $z^k$ in
	the Conway polynomial
	and $c(L)$ is the crossing number
	of a link $L$, is then
	\[
	\big|\nabla_{k}(L)\big|\le\frac{c(L)^k}{2^k\,k!}\,?
	\]
\end{openproblem}
% Source: LDT01
\plainproblemsection

\begin{openproblem}
	\label{ldt01:stoi.prm}
	Does $\mbox{min\,deg}_a \big( F_L(a^{-1},z) \big)\le 1-\chi(L)$
	hold
	for any link $L$?
	If $u(K)$ is the unknotting number of a knot $K$,
	does\break $\mbox{min\,deg}_a \big( F_K(a^{-1},z) \big)\le 2u(K)$
	hold
	for any knot $K$?
\end{openproblem}
% Source: LDT01
\plainproblemsection

\begin{openproblem}
	\label{ldt01:conj.vol_conj}
	For any knot $K$,
	\begin{equation}
		\label{ldt01:eq.vol_conk}
		2 \pi \cdot \lim_{N \to \infty} \frac{\log | J_N(K) |}{N} =
		v_3  || S^3 - K ||,
	\end{equation}
	where $|| \cdot ||$ denotes the
	simplicial volume
	and $v_3$ denotes the
	hyperbolic volume
	of the regular ideal tetrahedron.
\end{openproblem}
% Source: LDT01
\plainproblemsection

\begin{openproblem}
	Justify the above arguments rigorously.
\end{openproblem}
% Source: LDT01
\plainproblemsection

\begin{openproblem}
	\label{ldt01:conj.vol_CS}
	{\rm (H. Murakami, J. Murakami, M. Okamoto, T. Takata,
		\newline Y.~Yokota )}
	\
	For a hyperbolic link $L$,
	$$
	2 \pi \sqrt{-1} \cdot \lim_{N \to \infty} \frac{\log J_N(L)}{N}
	= {\rm CS}(S^3-L) + \sqrt{-1} \mbox{\rm vol}(S^3-L)
	$$
	for an appropriate choice of a branch of the logarithm,
	where CS and vol denote
	the Chern-Simons invariant
	and the
	hyperbolic volume respectively.
	Moreover,
	\begin{equation}
		\label{ldt01:eq.limJN1Jn}
		\lim_{N \to \infty}
		\frac{J_{N+1}(L)}{J_N(L)}
		= \exp \Big( \frac{1}{2\pi\sqrt{-1}}
		\big({\rm CS}(S^3-L) + \sqrt{-1} \mbox{\rm vol}(S^3-L) \big) \Big).
	\end{equation}
\end{openproblem}
% Source: LDT01
\plainproblemsection

\begin{openproblem}
	\label{ldt01:prob.CS_torus_knot}
	For a torus knot $K$,
	calculate ${\rm CS}(S^3-K)$ (giving an appropriate definition of it)
	and calculate $\lim \frac{\log J_N(K)}{N}$
	(fixing an appropriate choice of a branch of the logarithm).
\end{openproblem}
% Source: LDT01
\plainproblemsection

\begin{openproblem}
	\label{ldt01:conj.AS1_tf}
	$\cA(S^1;\Z)$ is
	torsion free.
\end{openproblem}
% Source: LDT01
\plainproblemsection

\begin{openproblem}
	\label{ldt01:conj.torsion_ws}
	Let $R$ be a commutative ring with 1, say $\Z/2\Z$.
	Every weight system
	$\cA(S^1;R)^{(d)}/\mbox{FI} \to R$
	is induced by some Vassiliev invariant
	$R \K \to R$.
\end{openproblem}
% Source: LDT01
\plainproblemsection

\begin{openproblem}
	\label{ldt01:conj.V_dist_knots}
	Vassiliev invariants distinguish oriented
	knots.
	(See Conjecture \ref{ldt01:conj.K_dist_knots} for an equivalent statement
	of this conjecture.)
\end{openproblem}
% Source: LDT01
\plainproblemsection

\begin{openproblem}
	\label{ldt01:prob.VdistKfromO}
	Does there exists a non-trivial oriented knot
	which can not be distinguished from the trivial knot
	by Vassiliev invariants?
	(See Problem \ref{ldt01:prob.KdistKfromO} for an equivalent problem.)
\end{openproblem}
% Source: LDT01
\plainproblemsection

\begin{openproblem}
	\label{ldt01:prob.sign_Vinv}
	Is the knot signature
	the limit of a sequence of
	Vassiliev invariants?
\end{openproblem}
% Source: LDT01
\plainproblemsection

\begin{openproblem}
	\label{ldt01:prob.Okuda_fish}
	Describe the set
	\begin{equation}
		\label{ldt01:eq.fish_set}
		\Big\{ \big( \frac{v_2(K)}{n^2}, \frac{v_3(K)}{n^3} \big)
		\in \R \times \R \ \Big| \
		\mbox{$K$ has a knot diagram with $n$ crossings}
		\Big\}.
	\end{equation}
\end{openproblem}
% Source: LDT01
\plainproblemsection

\begin{openproblem}
	\label{ldt01:conj.bound_v3}
	Let $v_3$ be as above.
	If a knot $K$ has a diagram with $n$ crossings,
	then
	$$
	|v_3(K)| \le \Big\lfloor \frac{n(n^2-1)}{24} \Big\rfloor.
	$$
\end{openproblem}
% Source: LDT01
\plainproblemsection

\begin{openproblem}
	\label{ldt01:prob.dimV}
	Determine the dimension
	of the space of primitive Vassiliev invariants of each degree $d$.
	Equivalently,
	determine the dimension of the
	space
	$\cA(S^1;\Q)\conn^{(d)}$ for each $d$.
\end{openproblem}
% Source: LDT01
\plainproblemsection

\begin{openproblem}
	\label{ldt01:prob.top_pres_muinv}
	Milnor's $\overline\mu$-invariants of string links
	can be defined similarly as above (see ).
	Find a topological presentation of a $\overline\mu$-invariant
	of string links
	(not assuming the vanishing of the lower $\overline\mu$-invariants).
	\begin{itemize}
		\item[\rm(1)]
		Show that $\mbox{lk}(L_{12\cdots n-1},L_n)$
		is well-defined in an appropriate sense.
		\item[\rm(2)]
		Identify it with $\overline{\mu}_{12\cdots n-1, n}(L)$.
	\end{itemize}
\end{openproblem}
% Source: LDT01
\plainproblemsection

\begin{openproblem}
	Let $I$ denote an oriented interval.
	\begin{itemize}
		\item[\rm(1)]
		Determine the dimensions of
		$\vcA(S^1;\Q)^{(d)}$ and $\vcA(I;\Q)^{(d)}$ for each $d$.
		\item[\rm(2)]
		Determine the dimensions of
		$\vcA(S^1;\Q)^{(d)}/\ora{FI}$ and $\vcA(I;\Q)^{(d)}/\ora{FI}$ for each $d$.
		\item[\rm(3)]
		Determine the dimensions of
		$\vcA(S^1;\Q)^{(d)}/\mbox{(weak \ora{FI})}$ and \newline
		$\vcA(I;\Q)^{(d)}/\mbox{(weak \ora{FI})}$ for each $d$.
	\end{itemize}
\end{openproblem}
% Source: LDT01
\plainproblemsection

\begin{openproblem}
	{\rm}\quad
	\label{ldt01:conj.Vinv_ckvk}
	Every Vassiliev invariant of classical knots
	can be extended to
	a finite type invariant
	of long virtual knots.
	(See also Problem \ref{ldt01:prob.Kinv_vk}.)
\end{openproblem}
% Source: LDT01
\plainproblemsection

\begin{openproblem}
	Calculate
	$\cF_d(\Z\K,{\mathfrak m})/\cF_{d+1}(\Z\K,{\mathfrak m})$,
	letting $\mathfrak m$ be a local move such as
	\begin{itemize}
		\item[\rm(1)]
		a $\#$ move,
		\item[\rm(2)]
		a pass move,
		\item[\rm(3)]
		a $\Delta$ move,
		\item[\rm(4)]
		an $n$-gon move.
	\end{itemize}
\end{openproblem}
% Source: LDT01
\plainproblemsection

\begin{openproblem}
	\label{ldt01:prob.D_link_htpc}
	Find necessary and sufficient conditions for two
	$\mu$-component links ($\mu>2$) to be $\Delta$ link homotopic.
\end{openproblem}
% Source: LDT01
\plainproblemsection

\begin{openproblem}
	\label{ldt01:prob.grad_loop_fil}
	Let $R$ be a commutative ring with $1$, say, $\Z$ or $\Q$.
	\begin{itemize}
		\item[\rm(1)]
		Describe the spaces
		$\cF_l(R (\MK); \mbox{loop}) / \cF_{l+1}(R (\MK); \mbox{loop})$.
		\item[\rm(2)]
		Describe the spaces
		$\cF_l(R \K; \DD) / \cF_{l+1}(R \K; \DD)$.
		\item[\rm(3)]
		Describe the image of the above map
		$\cF_l( R \K; \DD) \to \cF_l( R (\MK); \mbox{loop} )$.
	\end{itemize}
\end{openproblem}
% Source: LDT01
\plainproblemsection

\begin{openproblem}
	\label{ldt01:conj.p_loop_move}
	Take $(M_1,K_1)$ and $(M_2,K_2)$ of the above sort.
	Then, there exists a $(\Z/p\Z)$-equivariant isomorphism
	$\phi: H_1(\Sigma^p_{(M_1,K_1)};\Z)$ $\to H_1(\Sigma^p_{(M_2,K_2)};\Z)$
	preserving the linking pairing
	if and only if
	$(M_1,K_1)$ is equivalent to $(M_2,K_2)$
	by a finite sequence of mod $p$ loop moves.
\end{openproblem}
% Source: LDT01
\plainproblemsection

\begin{openproblem}
	$\{ K \ \simsubC{d} O \} / \simsubC{d+1}$ is
	torsion free
	for each $d$.
\end{openproblem}
% Source: LDT01
\plainproblemsection

\begin{openproblem}
	\label{ldt01:conj.Cd_string_link}
	Two $m$-strand string links $L$ and $L'$ are
	$C_d$-equivalent
	if and only if $v(L) = v(L')$
	for any $A$-valued
	finite type invariant
	$v$ of degree $< d$
	for any abelian group $A$.
\end{openproblem}
% Source: LDT01
\plainproblemsection

\begin{openproblem}
	\label{ldt01:prob.GH_vknot}
	Establish the Goussarov-Habiro theory
	for virtual knots.
\end{openproblem}
% Source: LDT01
\plainproblemsection

\begin{openproblem}
	\label{ldt01:prob.MK_HL}
	Describe the abelian group \newline
	$\{ (M,K) \simsubHL{d} (S^3,\mbox{unknot}) \} / \! \simsubHL{d+1}$
	for each $d$.
\end{openproblem}
% Source: LDT01
\plainproblemsection

\begin{openproblem}
	\label{ldt01:prob.HN_Vinv}
	Is there a similar statement for
	finite type invariants
	of links?
	Let $I$ be an ideal in the algebra $V$ of finite type invariants of links.
	Let $Z$ be the set of links that are annihilated by all members of $I$,
	and let $J$ be the ideal in $V$ of all invariants that vanish on $Z$.
	Clearly, $J$ always contains the radical of $I$.
	Are they always equal?
\end{openproblem}
% Source: LDT01
\plainproblemsection

\begin{openproblem}
	\label{ldt01:quest.An_Vn}
	Find a minimal finite subset $A_n$ of $V_n$ such that span$(A_n)=V_n$.
\end{openproblem}
% Source: LDT01
\plainproblemsection

\begin{openproblem}
	\label{ldt01:prob.cal_ZK}
	For each oriented knot $K$,
	calculate the Kontsevich invariant $\hZ(K)$ for all degrees.
\end{openproblem}
% Source: LDT01
\plainproblemsection

\begin{openproblem}
	\label{ldt01:conj.K_dist_knots}
	The Kontsevich invariant distinguishes oriented knots.
	(See Conjecture \ref{ldt01:conj.V_dist_knots} for an equivalent statement
	of this conjecture.)
\end{openproblem}
% Source: LDT01
\plainproblemsection

\begin{openproblem}
	\label{ldt01:prob.KdistKfromO}
	Does there exists a non-trivial oriented knot $K$
	such that $\hZ(K)$ $= \hZ(O)$
	for the trivial knot $O$?
	(See Problem \ref{ldt01:prob.VdistKfromO} for an equivalent problem.)
\end{openproblem}
% Source: LDT01
\plainproblemsection

\begin{openproblem}
	\label{ldt01:conj.Kdist_rev}
	For every oriented knot \(K\), is
	\[
	\widehat Z(K)=\widehat Z(-K),
	\]
	where \(-K\) denotes \(K\) with the opposite orientation?
\end{openproblem}
% Source: LDT01
\plainproblemsection

\begin{openproblem}
	Characterize those elements of $\hat\cA(S^1)\conn$ of the form $\log \hZ (K)$,
	or those elements of $\cB\conn$ of the form $\log_\sqcup \hZ(K)$.
\end{openproblem}
% Source: LDT01
\plainproblemsection

\begin{openproblem}
	\label{ldt01:prob.roberts_tpK}
	Give a good topological construction of the Kontsevich integral.
\end{openproblem}
% Source: LDT01
\plainproblemsection

\begin{openproblem}
	\label{ldt01:prob.Kinv_ff}
	Construct the Kontsevich invariant
	(i.e.\ a universal Vassiliev invariant)
	with coefficients in a finite field.
\end{openproblem}
% Source: LDT01
\plainproblemsection

\begin{openproblem}
	\label{ldt01:prob.Kinv_vk}
	Construct the ``Kontsevich invariant''
	(i.e.\ a universal
	finite type invariant)
	of virtual knots in $\vcA(I)$.
	(See also Conjecture \ref{ldt01:conj.Vinv_ckvk}.)
\end{openproblem}
% Source: LDT01
\plainproblemsection

\begin{openproblem}
	\label{ldt01:prob.csi_vknot}
	Construct a series of
	configuration space integrals
	whose value is in $\vcA(I)$
	so that it gives all
	finite type invariants
	of virtual knots.
\end{openproblem}
% Source: LDT01
\plainproblemsection

\begin{openproblem}
	\label{ldt01:prob.khs}
	Find another way to kill the
	hidden strata,
	so that the above three approaches
	can naturally present the mapping degree of the same map.
\end{openproblem}
% Source: LDT01
\plainproblemsection

\begin{openproblem}
	\label{ldt01:prob.present_ass}
	Find a combinatorial direct presentation of
	an associator
	for all degrees,
	in particular, an associator with rational coefficients.
\end{openproblem}
% Source: LDT01
\plainproblemsection

\begin{openproblem}
	\label{ldt01:drinfeldprob}
	Construct a rational Drinfel'd associator in
	the context of rational homotopy theory.
\end{openproblem}
% Source: LDT01
\plainproblemsection

\begin{openproblem}
	\label{ldt01:prob.roberts6}
	What is graph cohomology
	the cohomology of?
\end{openproblem}
% Source: LDT01
\plainproblemsection

\begin{openproblem}
	\label{ldt01:prob.bott}
	Give a geometric construction of these homology classes
	coming from Lie algebras.
\end{openproblem}
% Source: LDT01
\plainproblemsection

\begin{openproblem}
	\label{ldt01:prob.P_K}
	Find a topological construction of
	the 2-loop polynomial
	$P^\theta_K$.
\end{openproblem}
% Source: LDT01
\plainproblemsection

\begin{openproblem}
	Find a topological construction of the polynomial $P'_K$ given above.
\end{openproblem}
% Source: LDT01
\plainproblemsection

\begin{openproblem}
	Find a topological construction of
	the loop-degree $l$ part of
	the rational Z invariant
	$Z^{\rm rat}(K) \in \cA^{\Q[t^{\pm1},1/\Delta_K(t)]}(\emptyset;\Q)$
	of a knot $K$,
	for each $l$.
\end{openproblem}
% Source: LDT01
\plainproblemsection

\begin{openproblem}
	Find a basis of the space
	$\cA^{\Q[t^{\pm1},1/A(t)]}(\emptyset;\Q)^{({\rm loop}\ l)}$,
	for each $l$,
	where $A(t)$ is a polynomial with $A(1)=1$ and $A(t)=A(t^{-1})$.
	In particular, find a basis of the space
	$\cA^{\Q[t^{\pm1}]}(\emptyset;\Q)^{({\rm loop}\ l)}$.
\end{openproblem}
% Source: LDT01
\plainproblemsection

\begin{openproblem}
	\label{ldt01:kohno_prob.1}
	Construct explicitly a universal invariant of
	finite type
	for links in $\Sigma \times [0, 1]$ with values in ${\mathcal A}_{\Sigma}$.
\end{openproblem}
% Source: LDT01
\plainproblemsection

\begin{openproblem}
	\label{ldt01:prob.qPalg}
	Give a deformation quantization
	of the Poisson algebra ${\mathcal A}_{\Sigma}$
	which descends to a deformation quantization of $C({\mathcal M}^G(\Sigma))$.
\end{openproblem}
% Source: LDT01
\plainproblemsection

\begin{openproblem}
	\label{ldt01:prob.qCMG}
	Clarify the relation between
	a deformation quantization
	of $C({\mathcal M}^G(\Sigma))$
	at a special parameter and
	the space of conformal blocks
	in WZW models.
\end{openproblem}
% Source: LDT01
\plainproblemsection

\begin{openproblem}
	\label{ldt01:prob.ker_tau}
	Determine the image and the kernel of the above map $\tau$.
\end{openproblem}
% Source: LDT01
\plainproblemsection

\begin{openproblem}
	\label{ldt01:prob.hol_cb}
	Compute the holonomy of
	the space of conformal blocks
	of the twisted WZW model.
	In particular, determine the action of
	the braid group of $\Sigma$
	on the space of conformal blocks for each $G$ flat connection on $\Sigma$.
\end{openproblem}
% Source: LDT01
\plainproblemsection

\begin{openproblem}
	\label{ldt01:prob.PnS2AnS}
	Let $P_n(\Sigma)$
	denote the pure braid group of $\Sigma$ with $n$ strings.
	Does there exist an injective multiplicative homomorphism
	$$
	\theta: P_n(\Sigma) \rightarrow {\mathcal A}_n(\Sigma)
	$$
	defined over $\Q$?
\end{openproblem}
% Source: LDT01
\plainproblemsection

\begin{openproblem}
	Calculate
	$S_{2,\infty}(M)$ for each oriented 3-manifold $M$.
	Find a convenient methodology to calculate it.
\end{openproblem}
% Source: LDT01
\plainproblemsection

\begin{openproblem}
	\label{ldt01:prob.prz.it_tS2}
	Incompressible tori and 2-spheres in $M$ yield
	torsion
	in $S_{2,\infty}(M)$.
	It is a question of fundamental importance whether other surfaces
	can yield torsion as well.
\end{openproblem}
% Source: LDT01
\plainproblemsection

\begin{openproblem}
	If every closed incompressible surface in $M$ is parallel to $\partial M$,
	then $S_{2,\infty}(M)$
	is torsion free.
\end{openproblem}
% Source: LDT01
\plainproblemsection

\begin{openproblem}
	\label{ldt01:prob.prz.cS2}
	Compute
	$S_{2,\infty}(F_{0,3} \times S^1)$.
\end{openproblem}
% Source: LDT01
\plainproblemsection

\begin{openproblem}
	Let $F$ be a surface and $I$ an interval.
	Describe the algebra
	$S_{2,\infty}(F \times I)$.
\end{openproblem}
% Source: LDT01
\plainproblemsection

\begin{openproblem}
	Calculate the skein homology
	based on the Kauffman bracket skein relation.
\end{openproblem}
% Source: LDT01
\plainproblemsection

\begin{openproblem}
	We define the {\it $sl_3$ skein module}
	$S^{sl_3}(M)$
	of an oriented 3-manifold $M$
	by the defining relations of the $sl_3$ linear skein.
	Calculate $S^{sl_3}(M)$ of each 3-manifold $M$.
\end{openproblem}
% Source: LDT01
\plainproblemsection

\begin{openproblem}
	Calculate
	$S_3(M)$ for each oriented 3-manifold $M$.
	Find a convenient methodology to calculate it.
\end{openproblem}
% Source: LDT01
\plainproblemsection

\begin{openproblem}
	Let $F$ be a surface and $I$ an interval.
	Describe the algebra $S_3(F \times I)$.
\end{openproblem}
% Source: LDT01
\plainproblemsection

\begin{openproblem}
	Calculate
	$S_{3,\infty}(M)$ for each oriented 3-manifold $M$.
	Find a convenient methodology to calculate it.
\end{openproblem}
% Source: LDT01
\plainproblemsection

\begin{openproblem}
	Calculate the higher skein modules based on the Kauffman skein relation
	$W_i^{3,\infty}(M)$ and $\hat W^{3,\infty}(M)$
	(see below for their definitions).
\end{openproblem}
% Source: LDT01
\plainproblemsection

\begin{openproblem}
	Calculate
	$HS^q(M)$ for each 3-manifold $M$.
\end{openproblem}
% Source: LDT01
\plainproblemsection

\begin{openproblem}
	\label{ldt01:p_4.1}
	\mbox{}
	\begin{enumerate}
		\item[\rm(i)]
		Find generators of ${\cal S}_{4,\infty}(S^3,R)$.
		\item[\rm(ii)]
		For which parameters of the $(4,\infty)$ skein and framing relations,
		trivial links are linearly independent in ${\cal S}_{4,\infty}(S^3;R)$?
		\item[\rm(iii)]
		For which parameters of the $(4,\infty)$ skein and framing relations,
		the trivial knot is not representing
		a torsion
		element of ${\cal S}_{4,\infty}(S^3,R)$?
	\end{enumerate}
\end{openproblem}
% Source: LDT01
\plainproblemsection

\begin{openproblem}
	\label{ldt01:prob.skeinmod_7col}
	For which coefficients of the $(4,\infty)$ skein relation
	is the number of Fox $7$-colorings measured by the
	$(4,\infty)$ skein module?
\end{openproblem}
% Source: LDT01
\plainproblemsection

\begin{openproblem}
	Calculate $W_k(M)$ for each 3-manifold $M$.
\end{openproblem}
% Source: LDT01
\plainproblemsection

\begin{openproblem}
	Define a skein module of 3-manifolds,
	and calculate it.
\end{openproblem}
% Source: LDT01
\plainproblemsection

\begin{openproblem}
	\label{ldt01:prob.cl-qdl}
	Classify
	the isomorphism classes of connected quandles of order $n$
	for each positive integer $n$.
\end{openproblem}
% Source: LDT01
\plainproblemsection

\begin{openproblem}
	\label{ldt01:prob.rep_QK}
	Describe (the number of) representations
	of a knot quandle to a fixed connected quandle of finite order,
	say, by using knot invariants known so far,
	or by reducing the problem to the case of smaller target quandles.
\end{openproblem}
% Source: LDT01
\plainproblemsection

\begin{openproblem}
	Let $h_X$ be as above.
	Then, $\log h_X$ is not a
	Vassiliev invariant,
	unless it is constant.
\end{openproblem}
% Source: LDT01
\plainproblemsection

\begin{openproblem}
	\label{ldt01:prob.computeH2QX}
	Compute
	$H_2^Q(X)$ for each connected quandle $X$.
	More generally,
	find a convenient methodology to compute quandle (co)homology groups.
\end{openproblem}
% Source: LDT01
\plainproblemsection

\begin{openproblem}
	\label{ldt01:prob.HiQSnm}
	Compute
	$H_i^Q({\mathfrak S}_n^m)$ of ${\mathfrak S}_n^m$
	which denotes the quandle of the $n$th symmetric group
	with the binary operation given by $x \ast y = y^{-m} x y^m$.
\end{openproblem}
% Source: LDT01
\plainproblemsection

\begin{openproblem}
	\label{ldt01:prob.compute_PaK}
	Compute
	the quandle cocycle invariant $\Phi_\alpha(K)$ of each knot $K$
	for a second cohomology class $\alpha$ of a connected quandle.
\end{openproblem}
% Source: LDT01
\plainproblemsection

\begin{openproblem}
	Find relations between
	quandle cocycle invariants
	and knot invariants known so far,
	such as quantum invariants.
\end{openproblem}
% Source: LDT01
\plainproblemsection

\begin{openproblem}
	Compute
	$H^2_Q(X;A)$ for each $X$-module $A$.
\end{openproblem}
% Source: LDT01
\plainproblemsection

\begin{openproblem}
	Let the notation be as above.
	Then, extending the definition of the quandle cocycle invariant,
	define a knot invariant associated with $\alpha$,
	which is, roughly speaking,
	an invariant obtained by counting
	representations of a knot quandle $Q(K)$ to $X$
	with information whether
	each representation can lift to a representation $Q(K) \to Y$.
\end{openproblem}
% Source: LDT01
\plainproblemsection

\begin{openproblem}
	\label{ldt01:prob.qquandle}
	Define a quantum quandle.
\end{openproblem}
% Source: LDT01
\plainproblemsection

\begin{openproblem}
	\label{ldt01:FRSprob1}Is there a natural quandle space whose cohomology groups are the
	quandle cohomology groups?
\end{openproblem}
% Source: LDT01
\plainproblemsection

\begin{openproblem}
	\label{ldt01:FRSprob2}
	$H_3(R_p) \cong \Z\oplus \Z/p\Z$ for $p$ prime.
\end{openproblem}
% Source: LDT01
\plainproblemsection

\begin{openproblem}
	\label{ldt01:prob.TLfaithful}
	Is the representation of the braid group
	inside the Temperley-Lieb algebra
	faithful?
\end{openproblem}
% Source: LDT01
\plainproblemsection

\begin{openproblem}
	Is
	the Burau representation
	of $B_4$ faithful?
\end{openproblem}
% Source: LDT01
\plainproblemsection

\begin{openproblem}
	\label{ldt01:prob.B6V62}
	Is
	the action of $B_6$ on $V^6_2$ faithful?
\end{openproblem}
% Source: LDT01
\plainproblemsection

\begin{openproblem}
	\label{ldt01:prob.BMWalg}
	Generalise Lawrence's construction to obtain
	the irreducible representations of the BMW algebra.
\end{openproblem}
% Source: LDT01
\plainproblemsection

\begin{openproblem}
	\label{ldt01:prob.repBn_BMW}
	Find
	a larger family of irreducible representations of $B_n$
	which includes those coming from the BMW algebra.
\end{openproblem}
% Source: LDT01
\plainproblemsection

\begin{openproblem}
	Classify
	all irreducible representations of $B_n$.
\end{openproblem}
% Source: LDT01
\plainproblemsection

\begin{openproblem}
	\label{ldt01:prob.repBn_barQ}
	Is
	there a faithful representation of $B_n$
	into a group of matrices over $\bar\Q$?
\end{openproblem}
% Source: LDT01
\plainproblemsection

\begin{openproblem}
	\label{ldt01:benedetti.2}
	For every $W$, for every (pL)-class $\alpha_1$
	as above, $f_1^*$ is an isomorphism. This means, in
	particular, that
	finite type invariants
	of Legendrian knots should be
	definitely not sensitive to geometric (rigid) properties of the
	contact structures like ``tightness''.
\end{openproblem}
% Source: LDT01
\plainproblemsection

\begin{openproblem}
	\label{ldt01:prob.morton}
	From a knot diagram find an explicit such
	homomorphism to some permutation group or
	establish that the knot is trivial.
\end{openproblem}
% Source: LDT01
\plainproblemsection

\begin{openproblem}
	Find a new proof of the existence of a universal Vassiliev invariant of knots,
	presenting them by KTG's and their operations.
\end{openproblem}
% Source: LDT01
\plainproblemsection

\begin{openproblem}
	\label{ldt01:conj.KTG_ae}
	For each compact Lie group $G$, level $k$, and every KTG $K: \Gamma \to \R^3$,
	there exists a collection of measures $\mu_{{}_K}$
	on the space of gauge equivalence classes of $G$-connections on $\Gamma$
	satisfying the following conditions.
	\begin{itemize}
		\item
		It is well-behaved under KTG operations.
		\item
		It is ``localized'' near connections that extend to $S^3-K$.
		\item
		A half-twist framing change acts by $e^{\sqrt{-1} H \hbar /2}$,
		where $H$ is the Schr\"odinger operator on $G$.
		\item
		It recovers quantum invariants by
		$$
		I_R(K) = \int h_R(A) d \mu_{{}_K}(A),
		$$
		where $h_R(A)$ denotes the holonomy of $A$ in $R$.
		Here, $R$ is a set of representations of $G$ associated to edges of $\Gamma$
		and appropriate intertwiners associated to vertices of $\Gamma$.
	\end{itemize}
\end{openproblem}
% Source: LDT01
\plainproblemsection

\begin{openproblem}
	Construct an invariant of KTG's
	from configuration space integrals
	in a natural way.
\end{openproblem}
% Source: LDT01
\plainproblemsection

\begin{openproblem}
	\label{ldt01:prob.shadow_vol}
	Find a condition on shadow diagrams which is satisfied
	by shadow diagrams from alternating knots; and gives a lower bound
	on the hyperbolic volume.
\end{openproblem}
% Source: LDT01
\plainproblemsection

\begin{openproblem}
	\label{ldt01:conj.autoG}
	Any automorphism of $G$ is either the identity or the mirror map,
	that is,
	any automorphism of $G$ is induced by a diffeomorphism of the ambient space.
\end{openproblem}
% Source: LDT01
\plainproblemsection

\begin{openproblem}
	\label{ldt01:prob.roberts11}
	Extend Kuperberg's work on webs.
\end{openproblem}
% Source: LDT01
\plainproblemsection

\begin{openproblem}
	\label{ldt01:prob.roberts12}
	Extend the theory of measured laminations
	to higher rank groups.
\end{openproblem}
% Source: LDT01
\plainproblemsection

\begin{openproblem}
	\label{ldt01:prob.roberts13}
	What is the generating function for $q$-spin net
	evaluations?
\end{openproblem}
% Source: LDT01
\plainproblemsection

\begin{openproblem}
	\label{ldt01:prob.det_sign}
	If $n=4k+1$ with $k>0$, is there a knot with
	determinant
	$n$ and
	signature
	$4$?
\end{openproblem}
% Source: LDT01
\plainproblemsection

\begin{openproblem}
	\label{ldt01:prob.C2_solv}
	Is ${\cal C}_2$ solvable?
	Does ${\cal C}_2$ contain a free group?
\end{openproblem}
% Source: LDT01
\plainproblemsection

\begin{openproblem}
	\label{ldt01:prob.sign_chi}
	Do positive links
	of given signature
	$\sigma$ have bounded (below) maximal
	Euler characteristic $\chi$?
\end{openproblem}
% Source: LDT01
\plainproblemsection

\begin{openproblem}
	\label{ldt01:prob.pri_ach_sli}
	If a prime knot $K$ can be transformed into its mirror image by one
	crossing change, is $K$ achiral
	or (algebraically?) slice?
\end{openproblem}
% Source: LDT01
\plainproblemsection

\begin{openproblem}
	\label{ldt01:prob.sto.sq}
	Let $n$ be an odd natural number, different from 1, 9, and 49, such that
	$n$ is the sum of two squares.
	Is there a prime alternating achiral
	knot of determinant
	$n$?
\end{openproblem}
% Source: LDT01

\plainproblemsection

\begin{openproblem}
	\label{ldt01:prob.roberts_wqg}
	What are quantum groups?
\end{openproblem}
% Source: LDT01

\plainproblemsection

\begin{openproblem}
	(Flapan \& Gordon) Characterize all graphs $G$ that have the following Property: every imbedding of $G$ in ${\mathbb{R}}^{3}$ contains a subgraph which is a knotted, simple closed curve,.
\end{openproblem}
% Source: Kirby

\subsection*{References}
\begin{enumerate}[label={\arabic*.},leftmargin=2.2em]
	\item T. Ohtsuki (ed.), ``Problems on invariants of knots and 3-manifolds,'' \emph{Geometry \& Topology Monographs} \textbf{4} (2002), 377--572; revised 2004.
	\item T. Ohtsuki (ed.), \emph{Problems on Low-Dimensional Topology, 2014}, RIMS, Kyoto University, 2014.
	\item O. \c{S}avk, ``A survey of the homology cobordism group,'' arXiv:2209.10735, version 4, 2024.
	\item R. C. Kirby, ``Problems in low-dimensional topology,'' in W. H. Kazez (ed.), \emph{Geometric Topology}, Part 2, AMS/IP Studies in Advanced Mathematics \textbf{2.2}, American Mathematical Society and International Press, 1997, 35--473.
\end{enumerate}

\clearpage
\section{Surfaces and Mapping Class Groups}

\plainproblemsection

\begin{openproblem}
	(Thurston) Characterize the subgroups of ${\pi }_{0}\left( {\operatorname{Diff}\left( F\right) }\right)$ which act as translations of some complex geodesic in Teichmüller space.
\end{openproblem}
% Source: Kirby
\plainproblemsection

\begin{openproblem}
	(Meeks) For what genus $g$ does every periodic diffeomorphism $h: {F}_{g} \rightarrow  {F}_{g}$ have an invariant circle?
\end{openproblem}
% Source: Kirby
\plainproblemsection

\begin{openproblem}
	(Mess) (A) Find a finite presentation for the Torelli group, ${\kappa }_{g}, g \geq  3$.
	
	(B) Given $g$, what is the largest $k$ for which ${\kappa }_{g}$ admits a classifying space with finite $k$ - skeleton.
\end{openproblem}
% Source: Kirby
\plainproblemsection

\begin{openproblem}
	(Congruence Subgroup Conjecture) Every subgroup of finite index of ${\Gamma }_{S}$ contains a congruence subgroup, where $S$ is a compact, orientable surface, perhaps with boundary.
\end{openproblem}
% Source: Kirby
\plainproblemsection

\begin{openproblem}
	(Ivanov) (A) Is it true that ${H}^{1}\left( G\right)  = 0$ for any subgroup $G$ of finite index in ${\Gamma }_{S}$?
	
	(B) Does ${\Gamma }_{S}$ satisfy Kazhdan’s property $T$?
\end{openproblem}
% Source: Kirby
\plainproblemsection

\begin{openproblem}
	(A) (Penner) Is it possible that all nontrivial (i.e. $\neq  1$ ) elements of a normal subgroup of ${\Gamma }_{S}$ are pseudo-Anosov diffeomorphisms?
	
	(B) (Ivanov) If the subgroup $H$ of ${\pi }_{1}\left( S\right)$ is characteristic, then the kernel of the natural homomorphism
	
	$$
	{\Gamma }_{S} = \operatorname{Out}\left( {{\pi }_{1}\left( S\right) }\right)  \rightarrow  \operatorname{Out}\left( {{\pi }_{1}\left( S\right) /H}\right)
	$$
	
	is a normal subgroup. Call the class of normal subgroups of ${\Gamma }_{S}$ that arise by this construction $C$.
	
	\emph{Question.} Is it true that any normal subgroup of ${\Gamma }_{S}$ is commensurable with a normal subgroup in $C$?
\end{openproblem}
% Source: Kirby
\plainproblemsection

\begin{openproblem}
	(Mess) In the mapping class group ${\Gamma }_{g}, g \geq  3$, of a closed surface of genus g, a Dehn twist $\tau$ on a non-separating curve is a product of 3 commutators (by the lantern relationship ).
	
	(A) Is 3 the minimum possible?
	
	(B) What is the genus of ${\tau }^{n}$ as a function of $n$, where the genus is the number of commutators needed to express ${\tau }^{n}$?
	
	(C) For genus 2, the abelianization of the mapping class group ${\Gamma }_{g}$ is $\mathbb{Z}/{10}\mathbb{Z}$, generated by $\tau$. So ${\tau }^{10k}$ is a product of commutators; how many as a function of $k$?
	
	(D) What is the genus of a Dehn twist about a separating curve on a surface of genus at least 3? This might depend on the genus of the surface bounded by the separating curve (on a genus 2 surface, a Dehn twist on a separating curve has order 5 in the first homology group).
\end{openproblem}
% Source: Kirby
\plainproblemsection

\begin{openproblem}
	(Mess) Suppose there is an inclusion of ${\pi }_{1}\left( {F}_{g}\right), g \geq  2$, into a mapping class group ${\Gamma }_{h}^{1}$ of a one pointed surface of genus $h$. Is it ever the case that this inclusion lifts to ${\Gamma }_{h,1}$, the mapping class group of a surface of genus $h$ with one boundary component (diffeomorphisms modulo isotopies which fix the boundary)?
\end{openproblem}
% Source: Kirby
\plainproblemsection

\begin{openproblem}
	(Ivanov) (A) Conjecture (Mostow-Margulis superrigidity): If $G$ is an irreducible arithmetic group of rank $\geq  2$, then any homomorphism $G \rightarrow  {\Gamma }_{S}$ has finite image.
	
	(B) \emph{Conjecture.} Irreducible lattices in isometry groups of symmetric spaces of $R$ -rank $\geq  2$ never occur as subgroups of mapping class groups.
\end{openproblem}
% Source: Kirby
\plainproblemsection

\begin{openproblem}
	(Ivanov) Let ${d}_{W}$ (   )be the word metric on ${\Gamma }_{S}$ with respect to some given finite set of generators and let $\tau  \in  {\Gamma }_{S}$ be a Dehn twist. What is the growth rate of ${d}_{W}\left( {{t}^{n},1}\right)$?
\end{openproblem}
% Source: Kirby
\plainproblemsection

\begin{openproblem}
	(Mess) Consider surface bundles over surfaces where both fiber $F$ and base $B$ have genus $\geq  2$ and where ${\pi }_{1}\left( B\right)$ injects in the mapping class group of the fiber. Does such a bundle have a multisection (a submanifold which finitely covers the base)?
\end{openproblem}
% Source: Kirby
\plainproblemsection

\begin{openproblem}
	(Mess) (A) Let $E$ be a surface bundle over a surface, with base ${F}_{g}$ a closed, oriented surface of genus $g$ and fiber ${F}_{{g}^{\prime }}$. What is the minimum $g$ for which $\sigma \left( E\right)  \neq  0$?
	
	(B) Let $g\left( n\right)$ be the least genus for which $\sigma \left( E\right)  = {4n}$. The $\mathop{\lim }\limits_{{n \rightarrow  \infty }}\frac{g\left( n\right) }{n}$ exists, is finite, and is non-zero. What is the limit?
\end{openproblem}
% Source: Kirby
\plainproblemsection

\begin{openproblem}
	(Cooper \& Mess) Does every torsion free 3-manifold group act faithfully by homeomorphisms of ${\mathbb{R}}^{2}$?
\end{openproblem}
% Source: Kirby
\plainproblemsection

\begin{openproblem}
	(Lima) Are there commuting homeomorphisms of the 2-ball ${B}^{2}$ without a common fixed point?
\end{openproblem}
% Source: Kirby

\subsection*{References}
\begin{enumerate}[label={\arabic*.},leftmargin=2.2em]
	\item R. C. Kirby, ``Problems in low-dimensional topology,'' in W. H. Kazez (ed.), \emph{Geometric Topology}, Part 2, AMS/IP Studies in Advanced Mathematics \textbf{2.2}, American Mathematical Society and International Press, 1997, 35--473.
	\item R. \.{I}nan\c{c} Baykur, R. C. Kirby, and D. Ruberman (eds.), \emph{\(K3\): A New Problem List in Low-Dimensional Topology}, Mathematical Surveys and Monographs \textbf{295}, American Mathematical Society, 2026.
\end{enumerate}

\clearpage
\section{Three-Manifolds}

\paragraph{Notation for tribranched-surface problems.}
Let \(M\) be a compact \(3\)-manifold.  An essential tribranched surface
\(\Sigma\subset M\) is understood in the sense of the source.  It is
\emph{strongly essential} when, for every component
\(N\) of \(M\setminus\Sigma\), the natural map
\(\pi_1(N)\to\pi_1(M)\) is injective.

\plainproblemsection

\begin{openproblem}
	There are essential surfaces that cannot be obtained by the
	Culler--Shalen construction.  Can such surfaces be obtained by the
	essential-tribranched-surface construction described in the source?
\end{openproblem}
% Source: LDT02

\plainproblemsection

\begin{openproblem}
	Determine what topological information about a non-Haken \(3\)-manifold
	can be extracted from the essential tribranched surfaces it contains.
\end{openproblem}
% Source: LDT02

\plainproblemsection

\begin{openproblem}
	Does successive cutting along essential tribranched surfaces produce a
	structure analogous to a Haken hierarchy?
\end{openproblem}
% Source: LDT02

\plainproblemsection

\begin{openproblem}
	Give a topological or geometric characterization of strongly essential
	tribranched surfaces.
\end{openproblem}
% Source: LDT02

\plainproblemsection

\begin{openproblem}
	For \(n\geq3\), find an effective sufficient condition ensuring that the
	\(\operatorname{SL}_n\)-character variety of a \(3\)-manifold group has
	dimension at least \(1\).
\end{openproblem}
% Source: LDT02

\plainproblemsection

\begin{openproblem}
	Do the Kuperberg--Kontsevich--Thurston invariant and the
	Le--Murakami--Ohtsuki invariant agree for every rational homology
	\(3\)-sphere?
	\[
	Z^{\mathrm{KKT}}(M)\stackrel{?}{=}Z^{\mathrm{LMO}}(M).
	\]
\end{openproblem}
% Source: LDT02

\plainproblemsection

\begin{openproblem}
	In the correction term for the degree-\(n\) KKT invariant, is
	\[
	\delta_n=0
	\]
	for every \(n>1\)?
\end{openproblem}
% Source: LDT02

\plainproblemsection

\begin{openproblem}
	Let \(X\) be a closed oriented \(4\)-manifold with \(\chi(X)=0\), let
	\(\gamma\) be a unit vector field, and let
	\(\beta_1,\beta_2,\beta_3\) be generic sections of the normal bundle
	\(T^vX\) of \(\gamma\).  The signed count
	\[
	\bigl\langle\!\bigl\langle
	\beta_1,\beta_2,\beta_3
	\bigr\rangle\!\bigr\rangle
	\]
	of points where their span has dimension \(1\) equals
	\(3\operatorname{Sign}(X)\) and is divisible by \(6\).  Give a topological
	interpretation of this divisibility.
\end{openproblem}
% Source: LDT02

\plainproblemsection

\begin{openproblem}
	\label{ldt01:prob.QHS_tauG_LMO}
	Find pairs of non-homeomorphic rational homology 3-spheres
	that can be distinguished
	by their quantum $G$ invariants
	$\tau_{r}^{G}$ or their quantum $PG$
	invariants $\tau_{r}^{PG}$ for some level $r$ and some simply connected
	compact simple Lie group $G$
	but not by their LMO invariants.
\end{openproblem}
% Source: LDT01
\plainproblemsection

\begin{openproblem}
	\label{ldt01:prob.QHS_alltauG}
	Do the family of
	quantum $G$ invariants
	$\tau_{r}^{G}$ or the family of quantum $PG$ invariants
	$\tau_{r}^{PG}$, $G$ running through all simply connected compact simple
	Lie groups and $r$ running through all allowed levels,
	separate rational homology 3-spheres?
	How well do these families of invariants separate closed
	oriented 3-manifolds?
\end{openproblem}
% Source: LDT01
\plainproblemsection

\begin{openproblem}
	Find a 3-dimensional topological interpretation of quantum invariants
	of 3-manifolds.
\end{openproblem}
% Source: LDT01
\plainproblemsection

\begin{openproblem}
	{\rm}\quad
	\label{ldt01:conj.GP98}
	For non-vanishing $\tau^G_r(M)$,
	the absolute value
	$| \tau_r^G(M) |$ only depends on
	the fundamental group $\pi_1(M)$.
\end{openproblem}
% Source: LDT01
\plainproblemsection

\begin{openproblem}
	\label{ldt01:conj.AEC}
	The
	asymptotic expansion of $Z_k^G(M)$ of a closed oriented 3-manifold $M$
	is given by
	\begin{align*}
		Z_k^G(M) \, \underset{k\to\infty}{\sim} \,
		& e^{- \pi \sqrt{-1} (\mboxss{\rm dim } G)(1+b^1(M))/4 } \\*
		& \times \int_{[A] \in {\mathcal M}} \!\!\!
		e^{2\pi\sqrt{-1}r \mboxss{CS$(A)$}}
		r^{(h^1_A - h^0_A)/2}
		e^{- 2\pi\sqrt{-1} (I_A/4 + (h^0_A + h^1_A)/8)}
		\tau_M(A)^{1/2} \\*
		& \times \exp \left( \sum_{l=1}^\infty \frac{c^l k^{-l}}{ (2l)! (3l)! } \!\!
		\sum_{e(\Gamma)=-l}
		\!\! \frac{Z_\Gamma(M,A)}{ | \mbox{\rm Aut}(\Gamma) | } \right),
	\end{align*}
	putting $r = k + h^{{}^\vee}$,
	where the right hand side can be given in the mathematical viewpoint in
	certain cases,
	as mentioned above, but which needs further interpretation in general.
\end{openproblem}
% Source: LDT01
\plainproblemsection

\begin{openproblem}
	\label{ldt01:conj.AEC_Jorgen}
	Let $\{ c_0=0, c_1, \cdots, c_m \}$ be
	the set of values of the Chern-Simons functional of flat $G$ connections
	on a closed oriented 3-manifold $M$.
	There exist
	$d_j \in \Q$,
	$\tilde I_j \in \Q/\Z$,
	$v_j \in \R_+$,
	and $a^e_j \in \C$
	for $j=0,1,\cdots, m$ and $e = 1,2,3,\cdots$ such that ($r = k +
	h^{{}^\vee}$)
	$$
	Z_{k}^{G}(M) \ \underset{r\to\infty}{\sim} \
	\sum_{j=0}^m e^{2\pi\sqrt{-1}r c_j}
	r^{d_j} e^{\pi\sqrt{-1} \tilde I_j /4} v_j
	\big(1 + \sum_{e=1}^\infty a^e_j r^{-e} \big),
	$$
	that is, for all $E = 0, 1, 2,\ldots$, there exists a constant $c_E$ such that
	$$
	\Big| Z_{k}^{G}(M) -
	\sum_{j=0}^m e^{2\pi\sqrt{-1}r c_j}
	r^{d_j} e^{\pi\sqrt{-1} \tilde I_j /4} v_j
	\big(1 + \sum_{e=1}^E a^e_j r^{-e}\big) \Big| \le c_E r^{d-E-1}
	$$
	for all $r=2,3,4,\cdots$.
	Here, $d = \mbox{max} \{ d_0, \cdots, d_m \}$.
\end{openproblem}
% Source: LDT01
\plainproblemsection

\begin{openproblem}
	\label{ldt01:prob.pexp_O_H}
	If such an expansion exists, understand how it is related to the expansion of
	Ohtsuki and the expansion of Habiro.
\end{openproblem}
% Source: LDT01
\plainproblemsection

\begin{openproblem}
	\label{ldt01:conj.dj}
	Let ${\mathcal M}_j$ be the union of components of the moduli
	space of flat connections ${\mathcal M}$ which has Chern-Simons
	value $c_j$. Then
	\[d_j = \frac{1}{2}\max_{A\in {\mathcal M}_j} (h^1_A - h^0_A),\]
	where $\max$ here means the maximum value that $(h^1_A - h^0_A)$
	assumes on a Zariski open subset of ${\mathcal M}_j$.
\end{openproblem}
% Source: LDT01
\plainproblemsection

\begin{openproblem}
	\label{ldt01:conj.gr}
	Let $d = \max \{d_0, \ldots, d_n\}.$
	Then
	$|Z_{r}^{G}(M)| = O(r^d).$
\end{openproblem}
% Source: LDT01
\plainproblemsection

\begin{openproblem}
	\label{ldt01:conj.rt}
	There is a construct of the {\em right} measure, say $\tau_M(A)^{1/2}$ for
	$A\in {\mathcal M}_i$, from the
	square root of the Reidemeister torsion generalizing the
	non-degenerate case explained above and such that
	\[e^{\pi\sqrt{-1} \tilde I_j /4} v_j = \int_{A\in {\mathcal M}_i}
	e^{\pi\sqrt{-1} (-2 I_A + h^0_A + h^1_A) /4}\tau_M(A)^{1/2}.\]
\end{openproblem}
% Source: LDT01
\plainproblemsection

\begin{openproblem}
	\label{ldt01:conj.vol_olim}
	For every closed \(3\)-manifold \(M\), is
	\[
	2\pi i\,
	\underset{N\to\infty}{\operatorname{o\!-\!lim}}
	\frac{\log\tau_N^{SU(2)}(M)}{N}
	=
	\operatorname{CS}(M)+i\,\operatorname{vol}(M)?
	\]
	Here \(\operatorname{o\!-\!lim}\) is the optimistic limit used in the
	source; for a nonhyperbolic \(M\), the volume is interpreted as
	\(v_3\lVert M\rVert\), with \(v_3\) the volume of a regular ideal
	tetrahedron and \(\lVert M\rVert\) the simplicial volume.
\end{openproblem}
% Source: LDT01
\plainproblemsection

\begin{openproblem}
	\label{ldt01:prob.olim_Sf}
	Calculate $\mbox{o-$\lim$} \frac{\log \tau^{SU(2)}_N(M)}{N}$
	for Seifert fibered 3-manifolds $M$.
\end{openproblem}
% Source: LDT01
\plainproblemsection

\begin{openproblem}
	\label{ldt01:prob.inv_vol}
	Find a series of invariants of a 3-manifold
	(depending on roots of unity)
	that grows as its hyperbolic volume
	(or its simplicial volume).
\end{openproblem}
% Source: LDT01
\plainproblemsection

\begin{openproblem}
	\label{ldt01:prob.vol_conj_ncL}
	Find a correct generalization of
	the volume conjecture
	to other non-compact Lie groups.
\end{openproblem}
% Source: LDT01
\plainproblemsection

\begin{openproblem}
	\label{ldt01:prob.topCS}
	Define the Chern-Simons invariant
	${\rm CS}(M)$
	as a topological invariant of any closed oriented 3-manifold $M$,
	and of any knot (link) complement $M$ in a closed 3-manifold.
\end{openproblem}
% Source: LDT01
\plainproblemsection

\begin{openproblem}
	\label{ldt01:prob.Cstr_3mfd}
	Give a ``complex structure'' to the set of 3-manifolds.
	More precisely,
	find an embedding (or, an immersion) of the set of 3-manifolds
	to some complex variety such that
	its restriction to the set
	$\{ N_{K; (p,q)} \ | \ p^2+q^2 > \!\! > 0 \}$
	can be extended to a holomorphic map of
	the above mentioned complex parameter
	for any (hyperbolic) knot $K$ in any 3-manifold $N$.
\end{openproblem}
% Source: LDT01
\plainproblemsection

\begin{openproblem}
	\label{ldt01:bene.0}
	Generalize the construction of the QHI for flat principal $G$-bundles,
	for Lie groups $G$ different from $B$.
\end{openproblem}
% Source: LDT01
\plainproblemsection

\begin{openproblem}
	\label{ldt01:bene.1}
	Fix $(W,L)$ and vary $\rho$. Study $K_N$ as a function of the bundle,
	that is as a function defined on the character variety of $W$ with
	respect to $B$: regularity, fibers, and so on.
\end{openproblem}
% Source: LDT01
\plainproblemsection

\begin{openproblem}
	\label{ldt01:prob.bene2}
	Specialize Problem \ref{ldt01:bene.1} to bundles coming from the ordinary
	cohomology as above. For real additive ones, analyze the behaviour of
	the QHI with respect to Thurston's norm. Are they constant on the
	faces of the corresponding unit sphere?
\end{openproblem}
% Source: LDT01
\plainproblemsection

\begin{openproblem}
	\label{ldt01:prob.bene3}
	Understand the `phase factor' (i.e.\ the
	ambiguity due to $N$-th roots of unity) of the state sum
	$H_N({\cal T})$. Possibly derive from it an invariant for $(W,L,\rho)$
	endowed with some extra-structure, thus refining
	$K_N(W,L,\rho)$.
\end{openproblem}
% Source: LDT01
\plainproblemsection

\begin{openproblem}
	\label{ldt01:prob.bene4} Determine a suitable $(2+1)$ `decorated' cobordism theory
	supporting a (non purely topological) QFT containing the already
	defined QHI. Study in particular the behaviour of the QHI with respect
	to connected sums.
\end{openproblem}
% Source: LDT01
\plainproblemsection

\begin{openproblem}
	\label{ldt01:prob.bene5} Develop a 4-dimensional theory of QHI based on
	Turaev's shadow theory.
\end{openproblem}
% Source: LDT01
\plainproblemsection

\begin{openproblem}
	\label{ldt01:prob.bene6} Determine the actual relationship between $K_N(S^3,\cdot)$ and the
	coloured Jones polynomial
	$J_N(\cdot)$ (evaluated at
	$\omega=\exp(2i\pi/N)$ and normalized by $J_N(\rm{unknot})=1$), as
	functions of links.
\end{openproblem}
% Source: LDT01
\plainproblemsection

\begin{openproblem}
	(Real Volume Conjecture for QHI)\quad
	\label{ldt01:benedetti.conj}
	\noindent For any triple $(W,L,\rho)$ one
	has:
	$$ \lim_{N\to \infty} \ (2 \pi/N^2) \log
	(\vert K_N(W,L,\rho) \vert)= {\rm Im} \ R \bigl({\mathfrak c}_{\cal I}(W,L,\rho) \bigr).$$
\end{openproblem}
% Source: LDT01
\plainproblemsection

\begin{openproblem}
	\label{ldt01:prob.cal_tauSO3}
	For each rational homology 3-sphere $M$,
	calculate
	$\tau^{SO(3)}(M)$ and $\tau^{PSU(N)}(M)$ for all degrees.
\end{openproblem}
% Source: LDT01
\plainproblemsection

\begin{openproblem}
	\label{ldt01:prob.roberts14}
	Explain
	the appearance of modular forms in the Witten invariants.
\end{openproblem}
% Source: LDT01
\plainproblemsection

\begin{openproblem}
	\label{ldt01:prob.ch_tauSO3}
	Characterize
	those elements of $\Z[[q-1]]$ of the form $\tau^{SO(3)}(M)$
	of integral homology 3-spheres $M$.
\end{openproblem}
% Source: LDT01
\plainproblemsection

\begin{openproblem}
	\label{ldt01:conj.IslM}
	For a new indeterminate \(t\), set
	\[
	R_1'=\varprojlim_n
	R_1[t]\big/\bigl((t-q)(t-q^2)\cdots(t-q^n)\bigr).
	\]
	Does there exist an invariant \(I^{sl}(M)\in R_1'\) of integral homology
	\(3\)-spheres such that
	\[
	I^{sl}(M)\big|_{t=q^n}=I^{sl_n}(M)
	\]
	for every \(n\geq1\), with \(I^{sl_1}(M)=1\)?
\end{openproblem}
% Source: LDT01
\plainproblemsection

\begin{openproblem}
	\label{ldt01:prob.classify_TQFT}
	Find (and classify) all TQFT's.
\end{openproblem}
% Source: LDT01
\plainproblemsection

\begin{openproblem}
	\label{ldt01:prob.classify_modu_cat}
	Find (and classify) all
	modular categories.
\end{openproblem}
% Source: LDT01
\plainproblemsection

\begin{openproblem}
	\mbox{}\newline
	\noindent
	(1)\quad
	Characterize the power series of the form $P_{(V,Z)}(t)$.\newline
	\noindent
	(2)\quad
	For each power series $P(t)$
	(satisfying the characterization of (1)),
	classify all TQFT's $(V,Z)$ such that $P_{(V,Z)}(t) = P(t)$.
\end{openproblem}
% Source: LDT01
\plainproblemsection

\begin{openproblem}
	Find other spin TQFT's.
\end{openproblem}
% Source: LDT01
\plainproblemsection

\begin{openproblem}
	Formulate and find spin${}^c$ TQFT's.
\end{openproblem}
% Source: LDT01
\plainproblemsection

\begin{openproblem}
	\label{ldt01:prob.HQFT_spin}
	\begin{itemize}
		\item[\rm(1)]
		Extend HQFT's to spin and spin${}^c$ settings.
		\item[\rm(2)]
		Find algebra structures behind spin and spin${}^c$ HQFT's in dimension $1+1$.
	\end{itemize}
\end{openproblem}
% Source: LDT01
\plainproblemsection

\begin{openproblem}
	\label{ldt01:prob.HQFT_spin2}
	Study (spin and spin${}^c$) HQFT's
	with the target space $K(H,2)$ in dimensions $1+1, 2+1$, and $3+1$
	for $H = \Z^N$.
\end{openproblem}
% Source: LDT01
\plainproblemsection

\begin{openproblem}
	Find a geometric construction of
	a TQFT
	using $H^0({\cal M}_\Sigma, {\cal L}^{\otimes k})$.
	Namely, find a geometric way to associate a vector in
	$H^0({\cal M}_\Sigma, {\cal L}^{\otimes k})$
	to a 3-manifold $M$ with $\partial M = \Sigma$.
\end{openproblem}
% Source: LDT01
\plainproblemsection

\begin{openproblem}
	\label{ldt01:prob.EGonHMSLk}
	Study this action of  the finite group ${\mathcal E}(\Sigma)$ on
	$H^0({\cal M}_\Sigma, {\cal L}^{\otimes k})$, and describe the induced
	decompositions of this vector space according to the characters of  ${\mathcal E}(\Sigma)$. Also  relate these decompositions
	to decompositions of $V(\Sigma)$ for the TQFT $(V,Z)$
	derived from
	the quantum group $U_q(sl_2)$ at a $(k+N)$-th root of unity.
\end{openproblem}
% Source: LDT01
\plainproblemsection

\begin{openproblem}
	For a given TQFT $(V,Z)$,
	determine whether the image of $\tilde {\mathfrak M}_g$
	in $\mbox{End}\big( V(\Sigma_g) \big)$ is finite.
\end{openproblem}
% Source: LDT01
\plainproblemsection

\begin{openproblem}
	\label{ldt01:prob.NT_TQFT}
	
	Is there a relation between the Nielsen-Thurs\-ton classification
	of mapping classes of $\Sigma_g$
	and their images on $V(\Sigma_g)$ for TQFT's $(V,Z)$?
\end{openproblem}
% Source: LDT01
\plainproblemsection

\begin{openproblem}
	\label{ldt01:prob.kerler1}
	{\rm [Cyclotomic integer TQFT's]}
	\begin{enumerate}
		\item Find explicit/computable bases for the ${\cal V}_p(\Sigma_g)$ as free
		modules over $\Z[\zeta_p]$.
		\item Show that ${\cal V}_p$ can be extended to all cobordisms as a
		half-projective TQFT with $\mbox{\sf x}=(\zeta_p-1)^{\frac {p-3}2}\in \mbox{\sf R}
		=\Z[\zeta_p]$.
		\item Determine the structure of the ${\cal V}_p^{[j]}(M)$ and in how far
		they have lifts from $\FF_p$ to $\Z$, analogous to the Ohtsuki invariants
		for closed 3-manifolds.
		\item Find a universal TQFT that combines all ${\cal V}_p$, at least
		perturbatively, into one.
	\end{enumerate}
\end{openproblem}
% Source: LDT01
\plainproblemsection

\begin{openproblem}
	\label{ldt01:prob.kerler2}
	{\rm [Homological TQFT's]}
	\begin{enumerate}
		\item Find the irreducible components and ring structure
		(w.r.t $\oplus$ and
		$\otimes$) of $\ThKeQc^*$.
		\item Determine whether all strictly homological TQFT's lie in $\ThKeQc^*$.
		\item Identify the homological TQFT's that arise from the gauge theory of
		higher rank groups (such as $PSU(n)$ in ) with elements in
		$\ThKeQc^*$.
		\item Identify the irreducible factors of the constant orders
		${\cal V}^{[0]}_p$ of the cyclotomic
		integer expansion of the Reshetikhin-Turaev theory
		with elements in  $\ThKeQc^*$.
	\end{enumerate}
\end{openproblem}
% Source: LDT01
\plainproblemsection

\begin{openproblem}
	\label{ldt01:prob.kerler3}
	{\rm [Length = 1  TQFT's]}
	\begin{enumerate}
		\item Describe and construct
		algebraic $L=1$-extensions of $\Gamma_g$-representations
		to TQFT's, preferably as  ``simple''  generalizations
		of the Frohman-Nicas-$U(1)$-theory.
		\item Produce a classification of $L=1$-TQFT's in the sense of an extension
		theory of $\ThKeQc^*$.
		\item Identify the
		$\Gamma_g$-representations on relative $SU(2)$-moduli space
		from  with these TQFT's, and find similar, higher rank theories.
		\item Identify the ${\cal V}^{[0]}_p$ as $L=1$-theories, if possible.
	\end{enumerate}
\end{openproblem}
% Source: LDT01
\plainproblemsection

\begin{openproblem}
	\label{ldt01:prob.kerler4}
	{\rm [$q/l$-solvable and Casson  TQFT's]}
	\begin{enumerate}
		\item
		Lift the 1/1-solvable TQFT's of Casson type over $\FF_p$
		to a universal 1/1-solvable TQFT's of Casson type over $\Z$.
		\item
		Describe the {\em resulting } invariant
		$\Xi_{\Z}$ for 3-manifolds with $b_1(M)\geq 1$.
		\item
		Develop a  perturbation theory for general $q/l$-solvable TQFT's.
		\item
		Relate those with the various, standard resolutions of $\Gamma_g$.
		\item
		Relate them also to the traditional finite type theory for closed 3-manifolds.
		\item
		Describe the Reshetikhin-Turaev theories
		in this pattern.
	\end{enumerate}
\end{openproblem}
% Source: LDT01
\plainproblemsection

\begin{openproblem}
	\label{ldt01:prob.kerler5}
	{\rm [3-dim cobordisms from Hopf algebras]}
	\begin{enumerate}
		\item
		Find further relations on ${\cal A}lg$, besides the ones arising from
		the axiomatics of Hopf algebras,
		that would make ${\cal A}lg\to\ThKeCob^{\bullet}$ an isomorphism.
		\item
		Find  relations on ${\cal A}lg$ such that the maps
		$${\rm Aut}_{{\cal A}lg}(A^{\otimes g})\to \Gamma_g
		\cong {\rm Aut}_{\ThKeCob^{\bullet}}(\Sigma_{g,1}) $$ are isomorphisms.
		\item
		Relate this to obstructions, such as Steinberg and Whitehead
		groups, via stratified function spaces.
		\item
		What are the analogous algebraic structures in higher dimensions.
	\end{enumerate}
\end{openproblem}
% Source: LDT01
\plainproblemsection

\begin{openproblem}
	\label{ldt01:prob.kerler6}
	{\rm [Extended and half-projective  TQFT's]}
	\begin{enumerate}
		\item
		Describe in how far an ETQFT $\ThKeV$ with circle category $\cal C$
		can differ from $\ThKeV_{\cal C}$, thus introducing a equivalence notion that
		would establish a bijective correspondence between the class of
		ETQFT's and the class of modular tensor categories.
		\item
		Find an extended notion of half-projectivity that includes also
		surfaces that are both disconnected  and have boundary.
		\item
		Find constructions and axioms of ETQFT's
		that apply to more relaxed notions of boundedness or modularity.
	\end{enumerate}
\end{openproblem}
% Source: LDT01
\plainproblemsection

\begin{openproblem}
	\label{ldt01:prob.kerler7}
	{\rm [Non-semisimple vs.\ semisimple TQFT's, the double conjecture]}
	\begin{enumerate}
		\item
		Clarify the difference in the content of ${\cal V}_{\overline{\cal C}}$
		and ${\cal V}_{\cal C}$! \
		Are there homological TQFT's $\cal H$  such that ${\cal V}_{\cal C}$ is
		in some essential way equivalent to
		${\cal V}_{\overline{\cal C}}\otimes{\cal H}$?
		\item
		Find a  construction of  ${\cal W}_{\cal B}$  that generalizes
		the Turaev-Viro TQFT's
		to non-semisimple ${\cal B}$'s, similar to the
		way  generalized the Reshetikhin-Turaev construction.
		In the case of quantum groups
		and closed 3-mani\-folds this should reproduce
		a version of the Kuperberg invariant.
		\item
		What is the relation between ${\cal W}_{\cal B}$ and
		${\cal V}_{D({\cal B})}$? Are they in some sense isomorphic TQFT's?
	\end{enumerate}
\end{openproblem}
% Source: LDT01
\plainproblemsection

\begin{openproblem}
	\mbox{}
	\begin{itemize}
		\item[\rm(1)]
		Find (and classify) all semi-simple monoidal categories
		(with finitely many isomorphism classes of simple objects).
		\item[\rm(2)]
		Find (and classify) (finite dimensional)
		fusion rule algebras
		and sets of $6j$-symbols.
		\item[\rm(3)]
		Find (and classify) all
		subfactors (of finite depth).
	\end{itemize}
\end{openproblem}
% Source: LDT01
\plainproblemsection

\begin{openproblem}
	\label{ldt01:kawah_prob.4}
	Suppose we have a three-dimensional TQFT.
	Can we determine whether it arises from a
	fusion rule algebra
	and $6j$-symbols?
	If yes, can we describe all fusion rule algebras with $6j$-symbols
	producing the TQFT?
\end{openproblem}
% Source: LDT01
\plainproblemsection

\begin{openproblem}
	\label{ldt01:kawah_prob.5}
	Suppose we have two fusion rule algebras with
	$6j$-symbols
	and that two TQFT's arising from them are isomorphic.
	What relation do we have for the two sets of $6j$-symbols?
\end{openproblem}
% Source: LDT01
\plainproblemsection

\begin{openproblem}
	\label{ldt01:kawah_prob.6}
	Suppose we have a TQFT arising from a fusion rule algebra with
	$6j$-symbols.
	Using a fusion rule subalgebra and $6j$-symbols restricted on it,
	we can construct another TQFT.
	What relation do we have for these TQFT's?
\end{openproblem}
% Source: LDT01
\plainproblemsection

\begin{openproblem}
	\label{ldt01:kawah_prob.3}
	Suppose we have a semisimple
	ribbon category
	$C$
	with finitely many isomorphism classes of simple objects.
	If the $S$-matrix is invertible, we can construct the
	Reshetikhin-Turaev invariant
	and the state-sum invariant
	from $C$ and the latter is the square of the
	absolute value of the former.  If the $S$-matrix is
	not invertible, do we still have a similar description
	of the state-sum invariant?
\end{openproblem}
% Source: LDT01
\plainproblemsection

\begin{openproblem}
	\label{ldt01:kawah_prob.1}
	Let \(C_1\) be a semisimple ribbon category with finitely many isomorphism
	classes of simple objects and with noninvertible \(S\)-matrix.  Does
	\(C_1\) admit an essentially unique minimal extension to a modular
	category?  If so, how are the TQFTs associated with \(C_1\) and its
	minimal modular extension related?
\end{openproblem}
% Source: LDT01
\plainproblemsection

\begin{openproblem}
	\label{ldt01:kawah_prob.2}
	Let \(C_1\) be a semisimple ribbon category with degenerate \(S\)-matrix,
	and let \(C_2\) be the modular tensor category obtained from the
	alternative modularization construction described in the source.  How
	are the TQFTs associated with \(C_1\) and \(C_2\) related?
\end{openproblem}
% Source: LDT01
\plainproblemsection

\begin{openproblem}
	\label{ldt01:kawah_prob.7}
	Some fusion-rule algebras with \(6j\)-symbols do not appear to arise from
	quantum groups.  Determine the corresponding TQFTs.  If such examples
	occur in a parametrized family, does differentiation with respect to the
	parameter, after a suitable reparametrization, produce an interesting
	invariant, possibly of Vassiliev type?
\end{openproblem}
% Source: LDT01
\plainproblemsection

\begin{openproblem}
	\label{ldt01:prob.sato1}
	Find a subfactor
	which can distinguish lens spaces $L(7,1)$ and $L(7,2)$.
	Moreover, find a subfactor to classify 3-manifolds as well as possible.
\end{openproblem}
% Source: LDT01
\plainproblemsection

\begin{openproblem}
	\label{ldt01:prob.sato2}
	Construct a well-defined
	state sum type invariant
	from a strongly amenable
	subfactor.
\end{openproblem}
% Source: LDT01
\plainproblemsection

\begin{openproblem}
	\label{ldt01:prob.sato3}
	Let us consider the Turaev-Viro-Ocneanu invariant
	from a subfactor
	with a degenerate braiding. Then, find a description of this
	invariant as a Reshetikhin-Turaev invariant.
\end{openproblem}
% Source: LDT01
\plainproblemsection

\begin{openproblem}
	Can the Casson invariant
	of an integral homology 3-sphere $M$
	be characterized by the signature
	of a certain 4-manifold bounded by $M$?
\end{openproblem}
% Source: LDT01
\plainproblemsection

\begin{openproblem}
	\label{ldt01:prob.sf_CWL}
	Relate this surgery formula for
	the Casson-Wal\-ker-Lescop invariant
	with that of Lescop.
\end{openproblem}
% Source: LDT01
\plainproblemsection

\begin{openproblem}
	\label{ldt01:prob.spinCWL}
	Define an invariant $\lambda$ of a pair $(M,\sigma)$
	of a closed 3-manifold $M$ and a spin structure $\sigma$ on $M$
	such that
	$$
	\lambda_{\rm CWL}(M) = \sum_\sigma \lambda(M,\sigma)
	$$
	for any closed 3-manifold $M$,
	where the sum runs over all spin structures $\sigma$ on $M$.
\end{openproblem}
% Source: LDT01
\plainproblemsection

\begin{openproblem}
	\label{ldt01:prob.Gdf_ftinv}
	By presenting 3-manifolds by surgery along\break framed links in $S^3$,
	we can regard an invariant of 3-manifolds as an invariant of framed links.
	Establish a Gauss diagram formula
	for the link invariant derived from
	each finite type invariant of 3-manifolds.
\end{openproblem}
% Source: LDT01
\plainproblemsection

\begin{openproblem}
	$\cF^{\rm as}_d(\Z\M) / \cF^{\rm as}_{d+1}(\Z\M)$
	(resp. $\cF^{\rm b}_d(\Z\M) / \cF^{\rm b}_{d+1}(\Z\M)$)
	is torsion free
	for each $d$.
\end{openproblem}
% Source: LDT01
\plainproblemsection

\begin{openproblem}
	\label{ldt01:conj.Aphi_tf}
	$\cA(\emptyset;\Z)$ is
	torsion free.
\end{openproblem}
% Source: LDT01
\plainproblemsection

\begin{openproblem}
	\label{ldt01:conj.Fdist_hS3}
	Finite type invariants
	distinguish integral homology\break 3-spheres.
	(See Conjecture \ref{ldt01:conj.LMOdist_hS3}.)
\end{openproblem}
% Source: LDT01
\plainproblemsection

\begin{openproblem}
	\label{ldt01:prob.dimF}
	Determine the dimension
	of the space of primitive
	finite type invariants
	of integral homology 3-spheres of each degree $d$.
	Equivalently,
	determine the dimension of the space
	$\cA(\emptyset;\Q)\conn^{(d)}$
	for each $d$.
\end{openproblem}
% Source: LDT01
\plainproblemsection

\begin{openproblem}
	\label{ldt01:prob.Lambda_gen}
	Describe Vogel's algebra $\Lambda$,
	say, by giving complete sets of generators and relations of $\Lambda$.
\end{openproblem}
% Source: LDT01
\plainproblemsection

\begin{openproblem}
	Find a constructive combinatorial presentation of each
	finite type invariant
	of integral homology 3-spheres, and, in particular,
	of the Casson invariant,
	by localizing
	configuration space integrals.
\end{openproblem}
% Source: LDT01
\plainproblemsection

\begin{openproblem}
	\label{ldt01:prob.roberts15}
	What is the space of $3$-manifolds?
\end{openproblem}
% Source: LDT01
\plainproblemsection

\begin{openproblem}
	\label{ldt01:conj.ZHS_Yd}
	$\{ M \ \simsubY{2d} S^3 \} / \simsubY{2d+1}$ is
	torsion free
	for each $d$.
\end{openproblem}
% Source: LDT01
\plainproblemsection

\begin{openproblem}
	\label{ldt01:conj.hc_Yd}
	Let $F$ be an oriented compact surface.
	Two homology cylinders $C$ and $C'$ over $F$ are
	$Y_{d}$-equivalent
	if and only if
	$v(C) = v(C')$ for any $A$-valued
	finite type invariant $v$
	of $\cF_\star^{\rm Y}$-degree $< d$
	for any abelian group $A$.
\end{openproblem}
% Source: LDT01
\plainproblemsection

\begin{openproblem}
	\label{ldt01:prob.q_sigma}
	Classify the monoid (for orthogonal sum)
	of isomorphism classes of quadratic forms $q_\sigma$.
\end{openproblem}
% Source: LDT01
\plainproblemsection

\begin{openproblem}
	\label{ldt01:prob.spin_Yd}
	Describe the quotient set \newline
	$\{ \mbox{spin closed 3-manifolds} \}/ \!\!\underset{{}^{Y_{d}^s}}{\sim}$,
	in particular,
	for $d=2,3$.
\end{openproblem}
% Source: LDT01
\plainproblemsection

\begin{openproblem}
	\label{ldt01:prob.spinc_Yd}
	Describe the quotient set \newline
	$\{\mbox{spin${}^c$ closed 3-manifolds}\}/ \!\!\underset{{}^{Y_{d}^c}}{\sim}$,
	in particular,
	for $d=2,3$.
\end{openproblem}
% Source: LDT01
\plainproblemsection

\begin{openproblem}
	\label{ldt01:prob.calLMO}
	For each rational homology 3-sphere $M$,
	calculate $\LMO(M)$ for all degrees.
\end{openproblem}
% Source: LDT01
\plainproblemsection

\begin{openproblem}
	\label{ldt01:conj.LMOdist_hS3}
	The LMO invariant distinguishes
	integral homology 3-spheres.
	(See Conjecture \ref{ldt01:conj.Fdist_hS3}.)
\end{openproblem}
% Source: LDT01
\plainproblemsection

\begin{openproblem}
	Does there exist an integral/rational homology 3-sphere $M$
	such that $\LMO(M) = \LMO(S^3)$?
\end{openproblem}
% Source: LDT01
\plainproblemsection

\begin{openproblem}
	Characterize those elements of $\hat\cA(\emptyset)\conn$ which are of the form
	$\log \LMO (M)$
	for integral/rational homology 3-spheres.
\end{openproblem}
% Source: LDT01
\plainproblemsection

\begin{openproblem}
	Construct the LMO invariant with coefficients in a finite field.
\end{openproblem}
% Source: LDT01
\plainproblemsection

\begin{openproblem}
	Construct the LMO invariant
	(or the theory of finite type
	invariants)
	in arrow diagrams.
\end{openproblem}
% Source: LDT01
\plainproblemsection

\begin{openproblem}
	\label{ldt01:prob.spinLMO}
	Define the LMO invariant
	$\LMO(M,\sigma)$ of the pair of a closed
	3-manifold $M$ and a spin structure $\sigma$ of $M$ such that
	$\LMO(M) = \sum_\sigma \LMO(M,\sigma)$, where the sum runs over
	all spin structures on $M$. There is also a similar problem for
	spin${}^c$ structures.
\end{openproblem}
% Source: LDT01
\plainproblemsection

\begin{openproblem}
	\label{ldt01:prob.xiLMO}
	For every element $\xi\in H^1(M,\Z)$ construct an extension of
	$\LMO(M,\xi)$ of the LMO invariant
	such that when $\xi=0$ one
	recovers the usual LMO invariant.
\end{openproblem}
% Source: LDT01
\plainproblemsection

\begin{openproblem}
	\label{ldt01:prob.tor_KTinv}
	Do configuration spaces of
	have torsion
	in $\Z$-homology?
	Does such torsion deduce a torsion invariant of homology 3-spheres?
\end{openproblem}
% Source: LDT01
\plainproblemsection

\begin{openproblem}
	\label{ldt01:conj.shadow_num}
	The shadow number of a 3-manifold is quasi-linear in its Gromov norm.
	That is, there exist constants $c_1$ and $c_2$ such that
	$$
	c_1 || M || \le \mbox{(shadow number of $M$)} \le c_2 || M ||
	$$
	for any 3-manifold $M$,
	where $|| M ||$ denotes the Gromov norm
	of $M$.
\end{openproblem}
% Source: LDT01
\plainproblemsection

\begin{openproblem}
	Construct a universal Reshetikhin-Turaev invariant
	and
	a universal Turaev-Viro-Ocneanu invariant
	of closed 3-manifolds,
	in terms of the KTG algebra.
\end{openproblem}
% Source: LDT01
\plainproblemsection

\begin{openproblem}
	\label{ldt01:conj.turaev.Alex}
	A pair (a finitely generated abelian group $H$ of rank 1,
	an element  $\Delta(t) \in \Z [H/\mbox{Tors}\, H]=\Z[t^{\pm 1}]$)
	(where $t$ is a generator of $H/\mbox{Tors}\, H$)
	can be realized as the pair
	$(H_1(M)$, the Alexander polynomial $ \Delta_M$ of $M$)
	for a closed connected oriented 3-manifold $M$ if and only if
	$\Delta(t) = t^k \Delta(t^{-1})$ with even $k\in \Z$ and
	$\Delta(1) =\pm \vert \mbox{Tors}\, H\vert$.
\end{openproblem}
% Source: LDT01
\plainproblemsection

\begin{openproblem}
	(Jaco) Let ${M}^{3}$ be a compact 3-manifold. Are there at most a finite number of homotopy equivalent 3-manifolds? What if we restrict to the case where $M$ is orientable and sufficiently large? (See.)
\end{openproblem}
% Source: Kirby
\plainproblemsection

\begin{openproblem}
	(Stallings) Is an irreducible h-cobordism from ${\mathbb{RP}}^{2}$ to itself a product?
\end{openproblem}
% Source: Kirby
\plainproblemsection

\begin{openproblem}
	(Jaco) A manifold $M$ is said to have a manifold compactification if there exist a compact manifold $N$ and an imbedding $\varphi: M \rightarrow  N$ with $\varphi \left( {\operatorname{int}M}\right)  = \operatorname{int}N$. Let $M\left( G\right)$ be the unique covering space determined by the conjugacy class of a subgroup $G$ of ${\pi }_{1}\left( M\right)$.
	
	\emph{Question.} If $M$ is a compact, ${P}^{2}$ -irreducible,3-manifold and $G$ is a finitely generated subgroup, then does $M\left( G\right)$ admit a manifold compactification?
\end{openproblem}
% Source: Kirby
\plainproblemsection

\begin{openproblem}
	Classify imbeddings of orientable surfaces ${F}_{g}$ in ${S}^{3}$.
\end{openproblem}
% Source: Kirby
\plainproblemsection

\begin{openproblem}
	(Thurston) Is there a reasonable real-valued function $C$ on the set of 3- manifolds which measures the complexity of ${\pi }_{1}\left( {M}^{3}\right)?C$ should have the following properties:
	
	(A) if ${M}_{1}$ is a $k$ -fold cover of ${M}_{2}$, then $C\left( {M}_{1}\right)  = {kC}\left( {M}_{2}\right)$;
	
	(B) $C\left( {{M}_{1}\# {M}_{2}}\right)  = C\left( {M}_{1}\right)  + C\left( {M}_{2}\right)$;
	
	and perhaps
	
	(C) if $f: {M}_{1} \rightarrow  {M}_{2}$ has positive degree, then $C\left( {M}_{1}\right)  \geq  C\left( {M}_{2}\right)$;
	
	(D) $C\left( M\right)  =$ volume(M)if $M$ is either $\widetilde{SL}\left( {2,\mathbb{R}}\right)$ or hyperbolic 3-space modulo a discrete group.
	
	Note that (A) and (B) imply that $C\left( M\right)  = 0$ if $M$ fibers over ${S}^{1}$ with monodromy of finite order, or if $M$ fibers over a torus, or if $M$ is covered by ${S}^{3}$.
\end{openproblem}
% Source: Kirby
\plainproblemsection

\begin{openproblem}
	Which immersed 2-spheres in ${\mathbb{R}}^{3}$ bound immersed 3-balls?
\end{openproblem}
% Source: Kirby
\plainproblemsection

\begin{openproblem}
	Under what conditions does a closed, orientable 3-manifold $M$ smoothly imbed in ${S}^{4}$?
\end{openproblem}
% Source: Kirby
\plainproblemsection

\begin{openproblem}
	Let $X$ be an acyclic 2-complex and ${M}_{0}^{3}$ an abstract regular neighborhood of $X;\partial {M}_{0} = {S}^{2}$, so cap off to get a homology 3-sphere ${M}^{3}$. Find an effective way to compute the Rohlin invariant of ${M}^{3}$ in terms of $X$ and its regular neighborhood.
\end{openproblem}
% Source: Kirby
\plainproblemsection

\begin{openproblem}
	Which closed, orientable 3-manifolds belong to ${\mathcal{L}}_{1}$? to ${\mathcal{L}}_{k}, k > 1$?
\end{openproblem}
% Source: Kirby
\plainproblemsection

\begin{openproblem}
	To what extent does the Seifert matrix on ${F}^{4}$ determine the topology of the singularity? That is, does it determine ${Kup}$ to diffeomorphism? ${Fup}$ to diffeomorphism? ${Up}$ to isotopy?
\end{openproblem}
% Source: Kirby
\plainproblemsection

\begin{openproblem}
	Let
	\[
	G(t,z_0,z_1,z_2)
	\]
	be a family of isolated hypersurface singularities whose Milnor number is
	constant for \(t\) in a connected open set containing \(0\) and \(1\).
	Are the links of \(G_0\) and \(G_1\) pairwise diffeomorphic, and are their
	Milnor fibrations isomorphic?
\end{openproblem}
% Source: Kirby
\plainproblemsection

\begin{openproblem}
	Let \(F^4\) be the Milnor fiber of an isolated complex
	hypersurface singularity of complex dimension \(2\).  Is
	\[
	\sigma(F^4)\leq0,
	\]
	with equality only for a nonsingular point?
\end{openproblem}
% Source: Kirby
\plainproblemsection

\begin{openproblem}
	Let \(G\) be a finitely generated \(3\)-manifold group.
	\begin{enumerate}[label=\textup{(\Alph*)}]
		\item Does \(G\) admit a faithful representation in
		\(\operatorname{GL}_4(\mathbb R)\)?
		\item Is the Frattini subgroup of \(G\) trivial?
	\end{enumerate}
\end{openproblem}
% Source: Kirby
\plainproblemsection

\begin{openproblem}
	Let ${PL}\left( {M}^{3}\right)$ be the group of ${PL}$ homeomorphisms of a compact 3-manifold.
	
	(A) (Tollefson) Does ${PL}\left( M\right)$ have only finitely many conjugacy classes of finite cyclic subgroups of given order? (If yes, for ${S}^{3}$, then the Smith Conjecture holds.) What about finite subgroups?
	
	(B) (Giffen) Suppose $M$ admits no ${S}^{1}$ action. Is it possible for ${PL}\left( M\right)$ to contain an infinite torsion subgroup?
	
	(C) (Thurston) Is there a bound to the order of finite subgroups of ${PL}\left( M\right)$? Note that this is true for surfaces ${F}_{g}$ since a finite subgroup represents faithfully in ${GL}\left( {{2g} - 2,\mathbb{Z}/3\mathbb{Z}}\right)  =$ $\operatorname{Aut}\left( {{H}_{1}\left( {{F}_{g};\mathbb{Z}/3\mathbb{Z}}\right) }\right)$.
\end{openproblem}
% Source: Kirby
\plainproblemsection

\begin{openproblem}
	(Nielsen) Let $h$ be a homeomorphism of a 3-manifold $M$ such that the ${n}^{\text{th }}$ iterate ${h}^{n}$ is homotopic to the identity. When does there exist a map $g$ homotopic to $h$ such that ${g}^{n} =$ identity?
\end{openproblem}
% Source: Kirby
\plainproblemsection

\begin{openproblem}
	(Mess \& Rubinstein (Generalized Smale Conj.)) A strong version of the Geometrization Conjecture would give information about the diffeomorphism group of a 3-orbifold. For comparison, instead of just proving that each compact 2 manifold admits a metric of constant curvature, one can also show that the inclusion of the identity component of the isometry group (of any one of its constant curvature metrics) into the identity component of the diffeomorphism group is a homotopy equivalence.
	
	It is folk knowledge that the conjectures below would follow naturally if there was a natural flow on the space of metrics, given by a parabolic differential equation or otherwise, which deformed the space of all metrics on a given manifold, which admits one of the standard geometries, to the space of metrics locally isometric to one of the eight model geometries. Although there are partial results on conjectures A and B below, a uniform approach is desirable.
	
	(A) Suppose that $M$ is a compact, boundaryless,3-orbifold, modeled on one of Thurston’s eight geometries.
	
	(A1) \emph{Conjecture.} if $M$ is modeled on ${S}^{3},{H}^{3}$, or Solv, then the inclusion of the isometry group $I\left( M\right)$ of $M$ into $\operatorname{Diff}\left( M\right)$ is a homotopy equivalence.
	
	(B) This is a reformulation of the conjecture: let $G$ be a finite group acting isometrically on a geometric manifold $M$. Then the inclusion of the \{identity component of the centralizer of $G$ in the isometry group of $M\}$ into the $\{$ identity component of the centralizer of $G$ in the diffeomorphism group of $M\}$ is a homotopy equivalence.
\end{openproblem}
% Source: Kirby
\plainproblemsection

\begin{openproblem}
	(Kontsevich) Let ${M}^{3}$ be a connected 3-manifold with non-empty boundary, and let $\operatorname{BDiff}\left( {M\text{rel}\partial M}\right)$ be the classifying space for the group of diffeomorphisms of $M$ which are the identity on $\partial M$.
	
	\emph{Conjecture.} BDiff $\left( {M\text{rel}\partial M}\right)$ has the homotopy type of a finite ${CW}$ -complex.
	
	A particularly interesting case occurs when $M$ is equal to the complement of the unlink of $n$ components in ${S}^{3}, n > 1$, that is, the connected sum of $n$ solid tori.
\end{openproblem}
% Source: Kirby
\plainproblemsection

\begin{openproblem}
	Let \(M\) be a compact \(3\)-manifold and write
	\[
	\mathcal H(M)=\pi_0\operatorname{Homeo}(M).
	\]
	Let \(\mathcal D(M)\) be the subgroup generated by Dehn homeomorphisms,
	and let \(\mathcal D_{>0}(M)\) be the subgroup generated by those supported
	on \(D^2\), \(S^2\), or a two-sided \(\mathbb{RP}^2\).  Determine whether
	the following statements hold:
	\begin{enumerate}[label=\textup{(\Alph*)}]
		\item \(\mathcal D(M)\) has finite index in \(\mathcal H(M)\).
		\item \(\mathcal D_{>0}(M)\) has finite index in the kernel of
		\[
		\mathcal H(M)\longrightarrow\operatorname{Out}(\pi_1(M)).
		\]
		\item The image of
		\[
		\mathcal H(M)\longrightarrow
		\operatorname{Out}_{\partial M}(\pi_1(M))
		\]
		has finite index.
		\item The sequence
		\[
		1\longrightarrow\mathcal D_{>0}(M)\longrightarrow\mathcal H(M)
		\longrightarrow\operatorname{Out}_{\partial M}(\pi_1(M))
		\longrightarrow1
		\]
		is almost exact, meaning that each image has finite index in the
		corresponding kernel.
		\item If \(M\) is closed, irreducible, and non-Haken, then
		\(\mathcal H(M)\) is finite.
		\item \(\mathcal H(M)\) is finitely presented.
		\item \(\mathcal H(M)\) is virtually torsion-free, virtually of finite
		cohomological dimension, and virtually geometrically finite.
	\end{enumerate}
\end{openproblem}
% Source: Kirby
\plainproblemsection

\begin{openproblem}
	(Christy) (A) Which hyperbolic 3-manifolds have Anosov flows?
	
	(B) Does every irreducible, atoroidal 3-manifold with an Anosov flow have a hyperbolic metric?
	
	(C) Given an integer $N$, does there exist a hyperbolic 3-manifold with at least $N$ Anosov flows which are topologically inequivalent.?
\end{openproblem}
% Source: Kirby
\plainproblemsection

\begin{openproblem}
	(Christy) The geodesic flow (which is Anosov) on the tangent ${S}^{1}$ bundle of a closed, hyperbolic surface without boundary can be constructed by Dehn surgery on closed orbits of the suspension flow of a mapping torus (a ${T}^{2}$ bundle over ${S}^{1}$ ),.
	
	\emph{Question.} Does every Anosov flow on a closed compact 3-manifold arise in this way? If so, can it even arise from the mapping torus with monodromy $\left( \begin{array}{ll} 2 & 1 \\  1 & 1 \end{array}\right)$?
\end{openproblem}
% Source: Kirby
\plainproblemsection

\begin{openproblem}
	(Thurston) Does every closed hyperbolic 3-manifold admit a 2-dimensional foliation without Reeb components?
\end{openproblem}
% Source: Kirby
\plainproblemsection

\begin{openproblem}
	Let \(M\) be a hyperbolic \(3\)-manifold equipped with a
	Reebless foliation.  Each lifted leaf is a topological plane in
	\(\mathbb H^3\) and is conformally equivalent to the open unit disk
	\(\mathbb D\).  For a conformal parametrization
	\[
	f\colon\mathbb D\longrightarrow\mathbb H^3
	\]
	of such a leaf, does \(f\) extend continuously to a map from
	\(\overline{\mathbb D}\) to the compactification
	\(\mathbb H^3\cup S^2_\infty\)?
\end{openproblem}
% Source: Kirby
\plainproblemsection

\begin{openproblem}
	(Mess) Given a compact, hyperbolic 3-manifold $M$, show that a finite iteration of the following operation results in a link complement in ${S}^{3}$: remove the shortest geodesic and put a complete hyperbolic structure on the resulting manifold.
\end{openproblem}
% Source: Kirby
\plainproblemsection

\begin{openproblem}
	(Cooper) Given $K > 0$, is there a hyperbolic, rational homology 3-sphere $M$ such that the minimum over all points $p \in  M$ of the injectivity radius at $p$ is bigger than $K$?
\end{openproblem}
% Source: Kirby
\plainproblemsection

\begin{openproblem}
	(Neumann) What fields occur as invariant trace fields of hyperbolic 3- manifolds?
\end{openproblem}
% Source: Kirby
\plainproblemsection

\begin{openproblem}
	(Neumann) Let $k$ be a number field with exactly one conjugate pair of complex imbeddings. Then there exist arithmetic 3-manifolds $M$ with invariant trace field $k$.
	
	(A) \emph{Conjecture.} The Chern-Simons invariant ${CS}\left( M\right)$ is rational iff $k = \bar{k}$.
	
	(B) What is the number theoretic significance of ${CS}\left( M\right)$?
\end{openproblem}
% Source: Kirby
\plainproblemsection

\begin{openproblem}
	(J. D. S. Jones) Is the Chern-Simons invariant of a closed hyperbolic 3- manifold $M$ rational?
	
	(B) Are the known values of the Chern-Simons invariant and the volume of closed hyperbolic 3-manifolds consistent with the conjecture that ${K}_{3}\left( \mathbb{C}\right)$ is generated by the products from ${K}_{1}\left( \mathbb{C}\right)$, the torsion summand $\mathbb{Q}/\mathbb{Z}$, and the Borel classes?
	
	The papers  and  contain very interesting related results.
\end{openproblem}
% Source: Kirby
\plainproblemsection

\begin{openproblem}
	(A) Conjecture (Neumann \& Reid): If ${M}^{3} = {H}^{3}/\Gamma$ is a knot complement, then the commensurator of $\Gamma,\operatorname{com}\left( \Gamma \right)$, is equal to the normalizer of $\Gamma,\operatorname{norm}\left( \Gamma \right)$, except in the cases of the figure-8 knot and the two Aitchison-Rubinstein knots.
	
	(B) Is there a knot other than one of the above three whose complement is hyperbolic with cusp parameter in $\mathbb{Q}\left( \sqrt{-1}\right)$ or $\mathbb{Q}\left( \sqrt{-3}\right)$?
	
	(C) Are there knots other than the two dodecahedral knots for which the cusp field (field generated by the cusp parameter) is a proper subfield of the trace field?
\end{openproblem}
% Source: Kirby
\plainproblemsection

\begin{openproblem}
	(Mess) In each of the geometries Nil, Solv, and $\widetilde{SL}\left( {2,\mathbb{R}}\right)$, answer the following questions for a 3-manifold modeled on the geometry:
	
	(A) Are closed geodesics unique in their homotopy classes?
	
	(B) Are shortest geodesics unique in their homotopy classes?
	
	(C) Describe the different knots formed by simple closed geodesics in compact manifolds modeled on these spaces. In particular, under what circumstances does the complement in the universal cover of the preimage of a simple closed geodesic have free fundamental group?
\end{openproblem}
% Source: Kirby
\plainproblemsection

\begin{openproblem}
	(Freedman, Luo \& Hass) \emph{Conjecture.} Let ${M}^{3}$ be irreducible and non-compact, with $\partial M = \varnothing$, and with a homogeneously regular, complete Riemannian metric. Then an essential, proper imbedded plane is proper homotopic to a least area imbedded plane, where least area means that any Jordan curve in the plane bounds a least area disk which also lies in the plane.
\end{openproblem}
% Source: Kirby
\plainproblemsection

\begin{openproblem}
	(Hass) (A) Let $f: F \rightarrow  {M}^{3}$ be a one-sided proper map of a surface into a Riemannian 3-manifold which induces an isomorphism on ${\pi }_{1}$; also assume $f\left( F\right)$ has least area among all such surfaces.
	
	\emph{Conjecture.} $f$ is an imbedding.
	
	(B) Let $f: F \rightarrow  M$ be a one-sided, least area (in its homotopy class), ${\pi }_{1}$ -injective proper map which is homotopic to an imbedding.
	
	\emph{Conjecture.} $f$ is an imbedding.
	
	(C) A simple special case of this is a conjecture of Meeks: Let ${M}^{3}$ be ${S}^{1} \times  {B}^{2}$ with a metric in which the boundary is convex, and let $K$ be a(2,1)-torus knot in ${S}^{1} \times  {S}^{1}$. Let $F$ be a least area Möbius band bounded by $K$.
	
	\emph{Conjecture.} $F$ is imbedded.
\end{openproblem}
% Source: Kirby
\plainproblemsection

\begin{openproblem}
	Let \(M\) be a Riemannian \(3\)-manifold.  Construct a smooth
	flow on embedded surfaces in \(M\) that preserves embeddedness and carries
	each incompressible surface to a minimal surface.
\end{openproblem}
% Source: Kirby
\plainproblemsection

\begin{openproblem}
	(Harnack) What is the maximal number $N$ of components of a non-singular, real algebraic surface in ${\mathbb{RP}}^{3}$ of degree $m$?
\end{openproblem}
% Source: Kirby
\plainproblemsection

\begin{openproblem}
	(G. Martin) Given a Kleinian group $\Gamma$ and a triangle subgroup $\Delta$ with an invariant hemisphere $\Pi$ in ${H}^{3}$, is it true that for all $g \in  \Gamma$ either $g\left( \Pi \right)  = \Pi$ or $g\left( \bar{\Pi }\right)  \cap  \bar{\Pi } = \varnothing$ (i.e. the 2-orbifold $\Pi /\Delta$ is imbedded in ${H}^{3}/\Gamma$ )?
\end{openproblem}
% Source: Kirby
\plainproblemsection

\begin{openproblem}
	(Mess) Is there any example in any dimension of a finitely generated, infinite, torsion group acting effectively on a compact manifold? Or of an action of such a group on a noncompact manifold which fails to be properly discontinuous?
\end{openproblem}
% Source: Kirby
\plainproblemsection

\begin{openproblem}
	(Hass) Determine which finitely generated, 3-manifold groups are LERF. In particular, is the figure-8 knot group LERF? What about hyperbolic 3-manifold groups?
\end{openproblem}
% Source: Kirby
\plainproblemsection

\begin{openproblem}
	(Scott) A Poincaré duality group $G$ of dimension $n$ (PDn-group) is a finitely presented group whose homology and cohomology satisfy Poincaré duality over $\mathbb{Z}\left\lbrack  G\right\rbrack$ with a fundamental class in dimension $n$. The only known PDn-groups are fundamental groups of closed aspherical n-manifolds, and every PD2-group is the fundamental group of a closed aspherical surface ( and ).
	
	(A) Is every PD3-group G the fundamental group of a closed, aspherical 3-manifold?
	
	(B) \emph{Conjecture.} If $G$ has non-trivial center, then $G$ is the fundamental group of an aspherical Seifert fibered space.
	
	(B) is a special case of (A), and is motivated by the theorem that if a closed 3-manifold has fundamental group with center $\mathbb{Z}$, then it is a Seifert fibered space.
	
	(C) \emph{Conjecture.} There does not exist a finitely presented, infinite, torsion group.
\end{openproblem}
% Source: Kirby
\plainproblemsection

\begin{openproblem}
	For a group \(G\), define its transfinite lower central series
	by
	\[
	G_1=G,\qquad G_{\alpha+1}=[G,G_\alpha],
	\qquad
	G_\lambda=\bigcap_{\alpha<\lambda}G_\alpha
	\]
	when \(\lambda\) is a limit ordinal.  The \emph{length} of \(G\) is the
	least ordinal \(\alpha\) for which \(G_\alpha=G_{\alpha+1}\).
	\begin{enumerate}[label=\textup{(\Alph*)}]
		\item Find a closed \(3\)-manifold group \(G\) and a free group \(F\)
		such that
		\[
		F/F_k\cong G/G_k
		\qquad\text{for every }k\geq1,
		\]
		but the length of \(G\) is not \(\omega\).
		\item Which ordinals occur as the length of the fundamental group of a
		compact \(3\)-manifold?  More generally, what restrictions hold for the
		transfinite lower-central quotients \(G/G_\alpha\) of a closed
		\(3\)-manifold group?
		\item Is the length of a \(3\)-manifold group invariant under homology
		bordism?
		\item Among closed \(3\)-manifold groups, is having length greater than
		\(\omega\) generic in any natural sense, or is it rare?
	\end{enumerate}
\end{openproblem}
% Source: Kirby
\plainproblemsection

\begin{openproblem}
	Call a group \(G\) almost parafree if there is a homomorphism
	\(\varphi\colon F\to G\), where \(F\) is a finitely generated free group,
	such that
	\[
	F/F_n\cong G/G_n
	\]
	for every integer \(n>0\), where \(G_n\) denotes the lower central series.
	Let \(\mathcal E\) denote the stronger condition that there is an
	epimorphism \(e\colon G\to F\) inducing an isomorphism
	\[
	e_*\colon H_1(G)\longrightarrow H_1(F).
	\]
	
	If \(M\) is a closed \(3\)-manifold whose fundamental group is almost
	parafree, is \(M\) homology bordant to a \(3\)-manifold \(N\) for which
	\(\pi_1(N)\) satisfies condition \(\mathcal E\)?
\end{openproblem}
% Source: Kirby
\plainproblemsection

\begin{openproblem}
	(Rubinstein) Is there an algorithm to decide whether a 3-manifold $M$ has a 2-sided, immersed, incompressible surface?
\end{openproblem}
% Source: Kirby
\plainproblemsection

\begin{openproblem}
	(A) Formulate an integer valued notion of complexity of a closed 3-manifold, say $C\left( {M}^{3}\right)$. It should have the property that there are only finitely many 3-manifolds $F\left( N\right)$ with complexity less than any given integer $N$. The number of tetrahedra in a minimal combinatorial triangulation of ${M}^{3}$ is an example.
	
	Let $\mathcal{S}$ be a collection of closed 3-manifolds; we say that $\mathcal{S}$ is almost all 3-manifolds if $\mathop{\lim }\limits_{{N \rightarrow  \infty }}\frac{\left\lbrack  \mathcal{S} \cap  N\right\rbrack  }{F\left( N\right) } = 1$ where by definition $\left\lbrack  {\mathcal{S} \cap  N}\right\rbrack$ is the number of 3-manifolds in $\mathcal{S}$ of complexity less than $N$. Define equivalence of two notions of complexity, $C$ and ${C}^{\prime }$, to mean that $\mathcal{S}$ is almost all for $C$ iff $\mathcal{S}$ is almost all for ${C}^{\prime }$.
	
	(B) What other notions of complexity are equivalent to the example in (A)?
	
	(C) Let $\mathcal{H}$ be the set of Haken 3-manifolds. In the universe of irreducible 3-manifolds, is $\mathcal{H}$ almost all 3-manifolds? If not, what is the limit $\mathop{\lim }\limits_{{N \rightarrow  \infty }}\frac{\left\lbrack  H \cap  N\right\rbrack  }{F\left( N\right) }$?
\end{openproblem}
% Source: Kirby
\plainproblemsection

\begin{openproblem}
	(Boileau) \emph{Conjecture.} $A$ closed, orientable 3-manifold ${M}^{3}$ has a unique unstabilized Heegaard splitting if it is modeled on one of the four geometries,
	
	(1) ${S}^{3}$
	
	(2) ${E}^{3}$
	
	(3) ${S}^{2} \times  \mathbb{R}$
	
	(4) Nil.
\end{openproblem}
% Source: Kirby
\plainproblemsection

\begin{openproblem}
	What can be said about Heegaard splittings of closed, orientable 3-manifolds ${M}^{3}$ with one of the other four geometries,
	
	(1) ${H}^{2} \times  \mathbb{R}$
	
	(2) $\widetilde{PSL}\left( {2,\mathbb{R}}\right)$
	
	(3) Solv
	
	(4) ${H}^{3}$.
\end{openproblem}
% Source: Kirby
\plainproblemsection

\begin{openproblem}
	(Moriah) Is there a general theory describing Heegaard splittings of arbitrary 3-manifolds, as opposed to many special cases?
\end{openproblem}
% Source: Kirby
\plainproblemsection

\begin{openproblem}
	(Boileau) Does every closed, oriented, irreducible, 3-manifold have a finite cover which has a unique Heegaard splitting?
\end{openproblem}
% Source: Kirby
\plainproblemsection

\begin{openproblem}
	(Moriah) Any two Heegaard splittings of a given 3-manifold $M$ become equivalent (isotopic) after stabilizing sufficiently many times by connected summing with the genus one Heegaard splitting of ${S}^{3}$. Is there a universal bound $N$ such that two splittings of $M$ become equivalent when the higher genus splitting is stabilized $N$ times and the other is stabilized appropriately?
\end{openproblem}
% Source: Kirby
\plainproblemsection

\begin{openproblem}
	Let \(M\) be a closed, irreducible, atoroidal
	\(3\)-manifold.  Is the number of isotopy classes of genus-\(g\)
	Heegaard splittings of \(M\) bounded by a polynomial in \(g\)?
\end{openproblem}
% Source: Kirby
\plainproblemsection

\begin{openproblem}
	(Moriah) Find more examples of minimal genus Heegaard splittings of 3- manifolds for which the Heegaard genus is greater than the minimal number of generators in a presentation of ${\pi }_{1}\left( {M}^{3}\right)$.
\end{openproblem}
% Source: Kirby
\plainproblemsection

\begin{openproblem}
	(Moriah) Given any closed, orientable 3-manifold with any metric, show that the Heegaard surface of any Heegaard splitting is isotopic to a minimal surface.
\end{openproblem}
% Source: Kirby
\plainproblemsection

\begin{openproblem}
	(Adams) What is the minimal volume for a closed, orientable, hyperbolic 3-manifold $M$ of Heegaard genus $g$?
\end{openproblem}
% Source: Kirby
\plainproblemsection

\begin{openproblem}
	(Hass) Let $f: F \rightarrow  {M}^{3}$ be a 2-sided immersion of a surface such that ${f}_{ * }: {\pi }_{1}\left( F\right)  \rightarrow  {\pi }_{1}\left( M\right)$ is not injective.
	
	Simple loop conjecture for 3-manifolds: There exists an essential, simple loop on $F$ whose image is null homotopic in $M$.
\end{openproblem}
% Source: Kirby
\plainproblemsection

\begin{openproblem}
	(Haken) Let $f: {S}^{2} \rightarrow  {M}^{3}$ be a smooth immersion into an orientable 3- manifold $M$. Assume that the double point set in $f\left( {S}^{2}\right)$ is an immersed circle (its preimage must be two circles) with ${2k}$ triple points (there must be an even number). Let $N$ be a regular neighborhood of $f\left( {S}^{2}\right)$. ${\pi }_{1}\left( N\right)$ is either trivial or is $\mathbb{Z}/3\mathbb{Z}$, and in either case $\partial N$ is a union of 2-spheres.
	
	\emph{Conjecture.} $N$ is a punctured ${S}^{3}$ or $L\left( {3,1}\right)$.
\end{openproblem}
% Source: Kirby
\plainproblemsection

\begin{openproblem}
	(Y. Rong) Let ${M}^{3}$ be closed and orientable.
	
	(A) Are there only finitely many irreducible 3-manifolds $N$ such that there exists a degree one map $M \rightarrow  N$?
	
	(B) Does there exist an integer ${N}_{M}$ such that if
	
	$$
	M\overset{deg1}{ \rightarrow  }{M}_{2}\overset{deg1}{ \rightarrow  }{M}_{3}\overset{deg1}{ \rightarrow  }\cdots \overset{deg1}{ \rightarrow  }{M}_{k}
	$$
	
	is a sequence of degree 1 maps with $k > {N}_{M}$, then the sequence must contain a homotopy equivalence?
\end{openproblem}
% Source: Kirby
\plainproblemsection

\begin{openproblem}
	(Rourke) (A) Given two oriented, framed links, ${L}_{1}$ and ${L}_{2}$, in ${S}^{3}$ which define the same 3-manifold, is it possible to get from ${L}_{1}$ to ${L}_{2}$ by a sequence of coherent $K$ -moves, where a coherent $K$ -move is an ordinary $K$ -move in which all strands passing through the $\pm  1$ unknot are oriented in the same direction?
	
	(B) Given two (oriented), framed, closed braids, ${L}_{1}$ and ${L}_{2}$, in ${S}^{3}$ which define the same 3-manifold, is it possible to get from ${L}_{1}$ to ${L}_{2}$ using the Markov moves and braided $K$ -moves, where a braided $K$ -move is a $K$ -move in which the $\pm  1$ unknot lies in a plane transverse to the braid and encircles the leftmost’s strands in the braid, $0 \leq  s \leq  n$, for an n-strand braid (of course the $\pm  1$ unknot can be oriented and tilted to form part of the braid).
	
	(C) Since any 3-manifold can be obtained by surgery on a pure, framed closed braid, suppose two such give the same 3-manifold; is it possible to move from one braid to the other by a sequence of braided $K$ -moves (which preserve pureness) and conjugacy of braids?
\end{openproblem}
% Source: Kirby
\plainproblemsection

\begin{openproblem}
	(Auckly) Define the surgery number, $S\left( M\right)$, of a 3-manifold to be the minimal number of components in a framed link needed to define $M$. In the same way, define the Dehn surgery number, ${S}_{D}\left( M\right)$.
	
	(A) Is it true that the connected sum of $n$ non-trivial homology 3 -spheres has ${S}_{D} \geq  n$?
	
	(B) Find an irreducible, atoroidal, homology 3-sphere, $M$, with $S\left( M\right)  > 2$ or ${S}_{D}\left( M\right)  > 2$, or even arbitrarily large.
\end{openproblem}
% Source: Kirby
\plainproblemsection

\begin{openproblem}
	(Hoste) Every open 3-manifold (paracompact) can be obtained by surgery on a possibly infinite link imbedded in ${S}^{3} - X$ where $X$ is a closed subset of a fixed tame Cantor set in ${S}^{3}$.
	
	\emph{Question.} What is the calculus relating the different links that give the same 3-manifold?
\end{openproblem}
% Source: Kirby
\plainproblemsection

\begin{openproblem}
	(Auckly) (A) Is every closed 3-manifold $M$ with a representation $\varphi$: ${\pi }_{1}\left( M\right)  \rightarrow  {SU}\left( 2\right)$ given by a framed link $L$ which has a sublink $S$ satisfying the following two conditions:
	
	1. for any meridian $m$ of $S,\varphi \left( m\right)  =  \pm  I \in  {SU}\left( 2\right)$.
	
	2. $L - S$ can be divided into sublinks ${L}_{1} \cup  {L}_{2} \cup  \ldots  \cup  {L}_{k}, k \geq  1$, such that each ${L}_{i}$ lies in a 3-ball ${B}_{i}$ and the $\left\{  {B}_{i}\right\}$ are pairwise disjoint; then we require that, for each $i \in  \{ 1,\ldots, k\}$, all meridians of ${L}_{i}$ are taken by $\varphi$ into some $U\left( 1\right)$ in ${SU}\left( 2\right)$ which may vary with $i$.
	
	(B) Is every closed 3-manifold $M$, with a representation $\varphi: {\pi }_{1}\left( M\right)  \rightarrow  {SU}\left( 2\right)$, bordant via ${W}^{4}$ to a connected sum of lens spaces such that $\varphi$ extends over the bordism to $\Phi: {\pi }_{1}\left( W\right)  \rightarrow  {SU}\left( 2\right)?$
\end{openproblem}
% Source: Kirby
\plainproblemsection

\begin{openproblem}
	\begin{enumerate}[label=\textup{(\Alph*)}]
		\item Does every closed hyperbolic \(3\)-manifold \(M\) admit a
		nontrivial representation
		\[
		\pi_1(M)\longrightarrow \operatorname{SU}(2)?
		\]
		\item Find a closed \(3\)-manifold \(M\) and an isolated representation
		\(\rho\colon\pi_1(M)\to\operatorname{SU}(2)\) such that
		\[
		H^1(M;\operatorname{ad}\rho)\ne0.
		\]
	\end{enumerate}
\end{openproblem}
% Source: Kirby
\plainproblemsection

\begin{openproblem}
	(Stern) Find an example of a Floer sphere, i.e. find an irreducible homology 3-sphere, ${\sum }^{3}$, other than ${S}^{3}$, with vanishing Floer homology groups.
\end{openproblem}
% Source: Kirby
\plainproblemsection

\begin{openproblem}
	(Braam) The Casson invariant of homology three spheres is defined as the geometric oriented intersection of two subvarieties of the space of flat connections on a surface. Can a homology theory be defined such that these subvarieties are cycles in this theory?
\end{openproblem}
% Source: Kirby
\plainproblemsection

\begin{openproblem}
	(Garoufalidis) Does there exist a closed 3-manifold ${M}^{3}$ such that $Z\left( M\right)  =$ $Z\left( {S}^{3}\right)$ (M might be called a Witten sphere or quantum sphere)?
\end{openproblem}
% Source: Kirby
\plainproblemsection

\begin{openproblem}
	(The Closing Lemma) If $V$ is a ${C}^{r}$ vector field on a closed manifold with a non-trivial recurrent orbit (the orbit is contained in its backward or forward closure, ) is there a ${C}^{r}$ -perturbation ${V}_{0}$ of $V$ with a closed orbit?
\end{openproblem}
% Source: Kirby
\plainproblemsection

\begin{openproblem}
	(Freedman) Find a compactly supported vector field on ${\mathbb{R}}^{3}$ which generates a volume preserving flow ${\psi }_{t}: {\mathbb{R}}^{3} \rightarrow  {\mathbb{R}}^{3}$ so that for some closed loop, $\gamma$
	
	$$
	\text{Whitney norm}\left( {{\psi }_{t}\left( \gamma \right) }\right)  \geq  {c}_{0}{e}^{{c}_{1}t}
	$$
	
	for some positive constants ${c}_{0}$ and ${c}_{1}$.
	
	Are such flows generic?
\end{openproblem}
% Source: Kirby
\plainproblemsection

\begin{openproblem}
	(G. Kuperberg) \emph{Conjecture.} Every closed, orientable m-manifold ${M}^{m}$ with zero Euler characteristic has a minimal flow. In particular, ${S}^{3}$.
\end{openproblem}
% Source: Kirby
\plainproblemsection

\begin{openproblem}
	(S. Matsumoto) A minimal set $X$ of a flow $\phi$ is called isolated if $X$ admits a neighborhood $U$ with the following property: if the whole orbit of a point $x$ is contained in $U$, then $x$ is a point of $X$.
	
	\emph{Question.} Does a smooth flow on a 3-dimensional manifold admit an isolated minimal set which is not a closed orbit (homeomorphic to ${S}^{1}$ or a point)?
\end{openproblem}
% Source: Kirby
\plainproblemsection

\begin{openproblem}
	(G. Kuperberg) Is there a ${C}^{2}$, volume preserving flow on ${S}^{3}$ with no closed orbit? Can it be ${C}^{\infty }$ or even real analytic?
\end{openproblem}
% Source: Kirby

\subsection*{References}
\begin{enumerate}[label={\arabic*.},leftmargin=2.2em]
	\item T. Ohtsuki (ed.), ``Problems on invariants of knots and 3-manifolds,'' \emph{Geometry \& Topology Monographs} \textbf{4} (2002), 377--572; revised 2004.
	\item T. Ohtsuki (ed.), \emph{Problems on Low-Dimensional Topology, 2014}, RIMS, Kyoto University, 2014.
	\item R. C. Kirby, ``Problems in low-dimensional topology,'' in W. H. Kazez (ed.), \emph{Geometric Topology}, Part 2, AMS/IP Studies in Advanced Mathematics \textbf{2.2}, American Mathematical Society and International Press, 1997, 35--473.
	\item R. \.{I}nan\c{c} Baykur, R. C. Kirby, and D. Ruberman (eds.), \emph{\(K3\): A New Problem List in Low-Dimensional Topology}, Mathematical Surveys and Monographs \textbf{295}, American Mathematical Society, 2026.
\end{enumerate}

\clearpage
\section{Four-Manifolds}

\paragraph{Notation for homology-cobordism problems.}
Let \(\Theta_{\mathbb Z}^{3}\) and \(\Theta_{\mathbb Q}^{3}\) denote the
smooth oriented integral and rational homology-cobordism groups of
\(3\)-manifolds, respectively.  Let
\(\Theta_{\mathrm{SF}}^{3}\subset\Theta_{\mathbb Z}^{3}\) be the subgroup
generated by Seifert fibered homology spheres, and let
\[
\psi\colon\Theta_{\mathbb Z}^{3}\longrightarrow
\Theta_{\mathbb Q}^{3}
\]
be the natural map.

\plainproblemsection

\begin{openproblem}
	Is
	\[
	\Theta_{\mathbb Z}^{3}\cong\mathbb Z^{\infty}?
	\]
\end{openproblem}
% Source: LDT03

\plainproblemsection

\begin{openproblem}
	Study \(\Theta_{\mathbb Z}^{3}\) as an ordered group using the preorder
	induced by negative-definite cobordisms.  Describe the resulting
	filtrations, subgroups, and quotients.
\end{openproblem}
% Source: LDT03

\plainproblemsection

\begin{openproblem}
	Does the family
	\[
	\bigl\{\Sigma(p,q,pqn+1): n\ \text{odd}\bigr\}
	\]
	generate a subgroup, or a direct summand, isomorphic to
	\(\mathbb Z^{\infty}\) in \(\Theta_{\mathbb Z}^{3}\)?
\end{openproblem}
% Source: LDT03

\plainproblemsection

\begin{openproblem}
	Determine the relations among the homology-cobordism invariants
	\[
	\bar\mu,\qquad w,\qquad \kappa o_k,\qquad d,\qquad R.
	\]
\end{openproblem}
% Source: LDT03

\plainproblemsection

\begin{openproblem}
	Which integral homology \(3\)-spheres bound smooth contractible
	\(4\)-manifolds, and which bound smooth integral homology \(4\)-balls?
\end{openproblem}
% Source: LDT03

\plainproblemsection

\begin{openproblem}
	Does there exist a Seifert fibered integral homology sphere that bounds a
	smooth integral homology \(4\)-ball but no smooth contractible
	\(4\)-manifold?
\end{openproblem}
% Source: LDT03

\plainproblemsection

\begin{openproblem}
	Can a Seifert fibered integral homology sphere with more than three
	singular fibers bound an integral homology \(4\)-ball?
\end{openproblem}
% Source: LDT03

\plainproblemsection

\begin{openproblem}
	Does every Seifert fibered integral homology sphere bound a cork?
\end{openproblem}
% Source: LDT03

\plainproblemsection

\begin{openproblem}
	Does every integral homology \(3\)-sphere bound both a cork and a
	contractible homology plane?
\end{openproblem}
% Source: LDT03

\plainproblemsection

\begin{openproblem}
	Which integral homology \(3\)-spheres admit a smooth embedding in \(S^4\)?
\end{openproblem}
% Source: LDT03

\plainproblemsection

\begin{openproblem}
	Does there exist a Seifert fibered integral homology sphere that bounds an
	integral homology \(4\)-ball but does not smoothly embed in \(S^4\)?
\end{openproblem}
% Source: LDT03

\plainproblemsection

\begin{openproblem}
	Do graph integral homology \(3\)-spheres generate
	\(\Theta_{\mathbb Z}^{3}\)?
\end{openproblem}
% Source: LDT03

\plainproblemsection

\begin{openproblem}
	Does
	\[
	\Theta_{\mathbb Z}^{3}/\Theta_{\mathrm{SF}}^{3}
	\]
	contain a direct summand isomorphic to \(\mathbb Z^{\infty}\)?
\end{openproblem}
% Source: LDT03

\plainproblemsection

\begin{openproblem}
	Do integral Dehn surgeries on knots in \(S^3\) generate
	\(\Theta_{\mathbb Z}^{3}\)?
\end{openproblem}
% Source: LDT03

\plainproblemsection

\begin{openproblem}
	Does \(\Theta_{\mathbb Z}^{3}\) contain an element of finite order
	\(n\geq2\)?  More generally, is
	\(\Theta_{\mathbb Z}^{3}\) torsion-free?
\end{openproblem}
% Source: LDT03

\plainproblemsection

\begin{openproblem}
	Which rational homology \(3\)-spheres bound rational homology
	\(4\)-balls?
\end{openproblem}
% Source: LDT03

\plainproblemsection

\begin{openproblem}
	Does \(\Theta_{\mathbb Q}^{3}\) contain a direct summand isomorphic to
	\(\mathbb Q^n\) for some \(n\geq1\)?
\end{openproblem}
% Source: LDT03

\plainproblemsection

\begin{openproblem}
	Does \(\Theta_{\mathbb Q}^{3}\) contain an element of finite order
	\(n>2\)?
\end{openproblem}
% Source: LDT03

\plainproblemsection

\begin{openproblem}
	Does \(\ker\psi\) contain a subgroup, or a direct summand, isomorphic to
	\(\mathbb Z^{\infty}\)?
\end{openproblem}
% Source: LDT03

\plainproblemsection

\begin{openproblem}
	Does \(\operatorname{coker}\psi\) contain a direct summand isomorphic to
	\[
	\mathbb Z^{\infty}\oplus
	(\mathbb Z/2\mathbb Z)^{\infty}?
	\]
	Which other isomorphism types of subgroups or direct summands occur in
	\(\operatorname{coker}\psi\)?
\end{openproblem}
% Source: LDT03

\plainproblemsection

\begin{openproblem}
	Determine whether the Brieskorn spheres
	\[
	\Sigma(2,3,6n+1)
	\]
	bound rational homology \(4\)-balls for odd \(n\geq5\), and whether they
	bound integral homology \(4\)-balls for even \(n\geq6\).
\end{openproblem}
% Source: LDT03

\plainproblemsection

\begin{openproblem}
	Determine exactly which even indefinite unimodular forms
	\[
	2kE_8\oplus mH
	\]
	occur as the intersection forms of smooth, simply connected, closed
	\(4\)-manifolds, where
	\[
	H=
	\begin{pmatrix}
		0&1\\
		1&0
	\end{pmatrix}.
	\]
\end{openproblem}
% Source: Kirby
\plainproblemsection

\begin{openproblem}
	Find a homology 3-sphere $H$ of Rohlin invariant one such that $H\# H$ bounds a PL acyclic 4-manifold.
\end{openproblem}
% Source: Kirby
\plainproblemsection

\begin{openproblem}
	(Casson) Which rational homology 3-spheres bound rationally acyclic 4-manifolds?
\end{openproblem}
% Source: Kirby
\plainproblemsection

\begin{openproblem}
	(Freedman) (A) Let $f: \left( {M,\partial M}\right)  \rightarrow  X,\partial X$ ) be a degree 1 normal map from a smooth (or TOP) 4-manifold to a Poincaré space. Suppose $f \mid  \partial$ is a $\mathbb{Z}\left\lbrack  {{\pi }_{1}\left( X\right) }\right\rbrack$ - equivalence. A surgery obstruction
	
	$$
	\sigma \left( f\right)  \in  {\Gamma }_{4}\left( {\mathbb{Z}\left\lbrack  {{\pi }_{1}\left( X\right) }\right\rbrack   \rightarrow  \mathbb{Z}\left\lbrack  {{\pi }_{1}\left( X\right) }\right\rbrack  }\right)  \cong  {L}_{4}\left( {{\pi }_{1}\left( X\right) }\right)
	$$
	
	is defined. If $\sigma \left( f\right)  = 0$, is $f$ normally bordant to a homotopy equivalence rel $\partial$?
	
	(B) Can (A) be reduced to an equivalent question about link concordance?
\end{openproblem}
% Source: Kirby
\plainproblemsection

\begin{openproblem}
	(M. Cohen) Does there exist a 4-dimensional h-cobordism with nontrivial Whitehead torsion?
\end{openproblem}
% Source: Kirby
\plainproblemsection

\begin{openproblem}
	(Thurston) If a closed, orientable 4-manifold ${M}^{4}$ is a $K\left( {\pi,1}\right)$, must the Euler characteristic be $\geq  0$?
\end{openproblem}
% Source: Kirby
\plainproblemsection

\begin{openproblem}
	Classify the Cappell--Shaneson homotopy
	\(\mathbb{RP}^{4}\)'s up to diffeomorphism.  In particular, determine
	which of these manifolds are mutually diffeomorphic.
\end{openproblem}
% Source: Kirby
\plainproblemsection

\begin{openproblem}
	(Cappell \& Shaneson) (A) There is a homotopy equivalence $h: {S}^{2} \times$ ${\mathbb{RP}}^{2} \rightarrow  {S}^{2} \times  {\mathbb{RP}}^{2}$ constructed by
	
	$$
	{S}^{2} \times  {\mathbb{RP}}^{2} \rightarrow  \left( {{S}^{2} \times  {\mathbb{RP}}^{2}}\right)  \vee  {S}^{4}\overset{{id} \vee  \alpha }{ \rightarrow  }{S}^{2} \times  {\mathbb{RP}}^{2}
	$$
	
	where $\alpha$ generates ${\pi }_{4}\left( {S}^{2}\right)$. h is not homotopic to a ${PL}$ homeomorphism because it has a nontrivial normal invariant (the induced normal map ${h}^{-1}\left( {\mathbb{RP}}^{2}\right)  \rightarrow  {\mathbb{RP}}^{2}$ has a nontrivial Kervaire invariant). Is $h$ homotopic to a homeomorphism?
	
	(B) Construct a homotopy ${S}^{2} \times  {\mathbb{RP}}^{2}$ by removing a normal ${B}^{3}$ -bundle of ${\mathbb{RP}}^{1}$ and sewing in a suitable $\left( {{T}^{3} - {B}^{3}}\right)$ -bundle over ${S}^{1}$. Is the manifold or the homotopy equivalenc exotic?
\end{openproblem}
% Source: Kirby
\plainproblemsection

\begin{openproblem}
	(Y. Matsumoto) Is Scharlemann’s exotic ${S}^{1} \times  {S}^{3}\# {S}^{2} \times  {S}^{2}$ an exotic manifold or an exotic self-homotopy equivalence of ${S}^{1} \times  {S}^{3}\# {S}^{2} \times  {S}^{2}$?
\end{openproblem}
% Source: Kirby
\plainproblemsection

\begin{openproblem}
	(Thurston) Can a homology 4-sphere ever be a $K\left( {\pi,1}\right)$? Who knows an example of a rational homology 4-sphere which is a $K\left( {\pi,1}\right)$?
\end{openproblem}
% Source: Kirby
\plainproblemsection

\begin{openproblem}
	Does every simply connected, closed 4-manifold have a handlebody decomposition without 1-handles? Without 1- and 3-handles?
\end{openproblem}
% Source: Kirby
\plainproblemsection

\begin{openproblem}
	Does every integral unimodular symmetric bilinear form contain
	a characteristic element \(\alpha\) such that
	\[
	\alpha\cdot x\equiv x\cdot x\pmod 2
	\quad\text{for all }x,
	\qquad
	\alpha\cdot\alpha=\sigma,
	\]
	where \(\sigma\) is the signature of the form?
\end{openproblem}
% Source: Kirby
\plainproblemsection

\begin{openproblem}
	(Gordon) Let ${\sum }^{3}$ be a homology 3 -sphere which bounds an acyclic 4- manifold ${V}^{4}$ such that ${\pi }_{1}\left( {\sum }^{3}\right)  \rightarrow  {\pi }_{1}\left( {V}^{4}\right)$ is surjective. Let $K$ be a knot in ${\sum }^{3}$. Define $K$ to be homotopically ribbon in $V$ if there is a smoothly imbedded ${B}^{2}$ in $V,\partial {B}^{2} = K$, such that ${\pi }_{1}\left( {{\sum }^{3} - K}\right)  \rightarrow  {\pi }_{1}\left( {V - {B}^{2}}\right)$ is surjective.
	
	(A) Does $K$ slice in $V$ imply $K$ homotopically ribbon in $V$?
	
	(B) Does (A) hold at least for contractible $V$?
\end{openproblem}
% Source: Kirby
\plainproblemsection

\begin{openproblem}
	Let \(f\colon S^{2}\hookrightarrow\mathbb{CP}^{2}\) be a smooth
	embedding representing the positive generator of
	\(H_{2}(\mathbb{CP}^{2};\mathbb{Z})\).  Is the pair
	\[
	\bigl(\mathbb{CP}^{2},f(S^{2})\bigr)
	\]
	diffeomorphic to
	\(\bigl(\mathbb{CP}^{2},\mathbb{CP}^{1}\bigr)\)?
	More strongly, is \(f(S^{2})\) smoothly isotopic to
	\(\mathbb{CP}^{1}\)?
\end{openproblem}
% Source: Kirby
\plainproblemsection

\begin{openproblem}
	Let \(K\colon S^{2}\hookrightarrow S^{4}\) be a smooth
	2-knot.  Remove a tubular neighborhood \(S^{2}\times D^{2}\) and reglue it
	by the nontrivial diffeomorphism of \(S^{2}\times S^{1}\) that does not
	extend over \(S^{2}\times D^{2}\).  Is the resulting smooth
	4-manifold diffeomorphic to \(S^{4}\)?
\end{openproblem}
% Source: Kirby
\plainproblemsection

\begin{openproblem}
	(Y. Matsumoto) Does there exist a smooth, compact ${W}^{4}$, homotopy equivalent to ${S}^{2}$, which is spineless, i.e., contains no ${PL}$ imbedded ${S}^{2}$ representing the generator of ${H}_{2}\left( W\right)$?
\end{openproblem}
% Source: Kirby
\plainproblemsection

\begin{openproblem}
	\begin{enumerate}[label=\textup{(\Alph*)}]
		\item Does every simply connected closed \(4\)-manifold admit a basis
		of \(H_2\) represented by PL embedded \(2\)-spheres?
		\item Which even elements of
		\[
		\pi_2(\mathbb{RP}^{2}\times T^{2})\cong\mathbb Z
		\]
		can be represented by PL embedded \(2\)-spheres?
	\end{enumerate}
\end{openproblem}
% Source: Kirby
\plainproblemsection

\begin{openproblem}
	(Schoenflies) \emph{Conjecture.} If ${S}^{3}$ is ${PL}$ imbedded in ${S}^{4}$, then its closed complements are PL 4-balls.
\end{openproblem}
% Source: Kirby
\plainproblemsection

\begin{openproblem}
	What is ${\pi }_{0}\left( {\operatorname{Diff}\left( {S}^{4}\right) }\right)$? Indeed, is $\operatorname{Diff}\left( {S}^{4}\right)  \simeq  O\left( 5\right)$?
\end{openproblem}
% Source: Kirby
\plainproblemsection

\begin{openproblem}
	Let \(S\) be a simply connected complex algebraic surface,
	and let \(\alpha\in H_2(S;\mathbb Z)\) be represented by an algebraic
	curve of genus \(g\).  When
	\[
	\alpha\cdot\alpha<0,
	\]
	must \(g\) be the minimum genus among all smoothly embedded surfaces
	representing \(\alpha\)?
\end{openproblem}
% Source: Kirby
\plainproblemsection

\begin{openproblem}
	A simply connected, complex analytic surface $S$ may contain exceptional curves (complex ${\mathbb{CP}}^{1}$ ’s with self-intersection -1) in which case $S$ is diffeomorphic to ${M}^{4}\# r\left( {-{\mathbb{CP}}^{2}}\right)$.
	
	Can $S$ be decomposed as a connected sum in any other way?
\end{openproblem}
% Source: Kirby
\plainproblemsection

\begin{openproblem}
	(A. Kas) Let ${F}^{2}$ be a closed orientable 2-manifold. Let ${\gamma }_{1},\ldots,{\gamma }_{n}$ be imbedded circles in ${F}^{2}$. let ${\tau }_{i}: {F}^{2} \rightarrow  {F}^{2}\left( {i = 1,\ldots, n}\right)$ be a left-handed Dehn (Lickor-ish) twist about the circle ${\gamma }_{i}$. Assume that ${\tau }_{n} \circ  \ldots  \circ  {\tau }_{2} \circ  {\tau }_{1}: {F}^{2} \rightarrow  {F}^{2}$ is isotopic to the identity.
	
	\emph{Conjecture.} Let $E \subset  {H}_{1}\left( {{F}^{2};\mathbb{Q}}\right)$ be the subspace spanned by the ${\gamma }_{i}\left( {i = 1,\ldots, n}\right)$. Then the alternating intersection form is nondegenerate on $E$.
\end{openproblem}
% Source: Kirby
\plainproblemsection

\begin{openproblem}
	Let \(g_{1},g_{2}\colon F\looparrowright\mathbb{CP}^{2}\) be
	semialgebraic immersions of the same degree, in the sense used in the
	source.  Is there an ambient isotopy of \(\mathbb{CP}^{2}\) carrying
	\(g_{1}(F)\) to \(g_{2}(F)\)?
\end{openproblem}
% Source: Kirby
\plainproblemsection

\begin{openproblem}
	Construct a locally flat proper embedding of the punctured
	Poincar\'e homology sphere into \(\mathbb{R}^{4}\).
\end{openproblem}
% Source: Kirby
\plainproblemsection

\begin{openproblem}
	Fix an atlas on an exotic \(\mathbb R^4\), and identify its
	underlying topological manifold with \(\mathbb R^4\).  Let
	\[
	B_r=\{x\in\mathbb R^4:\lVert x\rVert<r\}
	\]
	with the induced smooth structure.
	\begin{enumerate}[label=\textup{(\Alph*)}]
		\item Determine the largest \(\rho\) such that \(B_r\) is diffeomorphic
		to the standard \(\mathbb R^4\) for \(r<\rho\), and describe what
		happens at the limiting sphere.
		\item Determine, for \(\rho<r<s\), when \(B_r\) is diffeomorphic to the
		original exotic \(\mathbb R^4\) or to \(B_s\).
		\item Does every smoothly embedded \(S^3\) in the exotic
		\(\mathbb R^4\) bound a smooth \(4\)-ball?  Alternatively, can every
		such sphere be engulfed in a standard smoothly embedded
		\(\mathbb R^4\)?
	\end{enumerate}
\end{openproblem}
% Source: Kirby
\plainproblemsection

\begin{openproblem}
	Let \(R\) be an exotic \(\mathbb R^4\).
	\begin{enumerate}[label=\textup{(\Alph*)}]
		\item Construct a useful complete Riemannian metric on \(R\), and
		describe the cut locus of a point for such a metric.
		\item Does there exist an exotic \(\mathbb R^4\) that cannot be split
		by a proper smooth embedding of \(\mathbb R^3\) into two exotic pieces?
	\end{enumerate}
\end{openproblem}
% Source: Kirby
\plainproblemsection

\begin{openproblem}
	Does there exist an exotic smooth structure on \(S^4\) or
	on \(S^3\times S^1\)?
\end{openproblem}
% Source: Kirby
\plainproblemsection

\begin{openproblem}
	(Freedman) Is a positive untwisted double of the Borromean rings topologically slice?
\end{openproblem}
% Source: Kirby
\plainproblemsection

\begin{openproblem}
	(Freedman) Let $X$ be the cone on the unlink of $n$ components in ${S}^{3}$. Suppose $X$ is imbedded properly in ${B}^{4}$ and is locally flat except at the cone point $*$. Suppose the local homotopy at $*$ is free. Does this imply that the imbedding is flat, i.e. has a neighborhood homeomorphic to a neighborhood of the standard imbedding of $X$ in ${B}^{4}$?
\end{openproblem}
% Source: Kirby
\plainproblemsection

\begin{openproblem}
	(Freedman) Find a homotopy theoretic criterion for when ${M}^{3}/Y \subset  {\mathbb{R}}^{4}$ has a one-sided mapping cylinder neighborhood, where $Y$ is an acyclic set in the 3-manifold $M$.
\end{openproblem}
% Source: Kirby
\plainproblemsection

\begin{openproblem}
	Is every closed, simply connected smooth \(4\)-manifold
	homotopy equivalent to a connected sum of complex surfaces, allowing either
	orientation on each summand?
\end{openproblem}
% Source: Kirby
\plainproblemsection

\begin{openproblem}
	Let a closed smooth \(4\)-manifold be written as
	\[
	M=W_1\cup_N W_2,
	\]
	where \(N\) is a closed \(3\)-manifold.  Suppose that properly embedded
	disks \(D_i\subset W_i\) have the same boundary in \(N\).  Perform a
	\(k\)-fold twist along a neighborhood of \(D_1\) and a
	\((-k)\)-fold twist along a neighborhood of \(D_2\), so that the
	boundary gluing is preserved; call the resulting operation a
	\emph{generalized Gluck twist}.  Can any two simple-homotopy-equivalent
	closed smooth \(4\)-manifolds be related by a finite sequence of such
	generalized Gluck twists?
\end{openproblem}
% Source: Kirby
\plainproblemsection

\begin{openproblem}
	(Mandelbaum) What (minimal) knowledge of homotopy groups, intersection pairings, etc. determines the homotopy type of a closed, compact 4-manifold?
\end{openproblem}
% Source: Kirby
\plainproblemsection

\begin{openproblem}
	Describe the Fintushel-Stern involution on ${S}^{4}$ in equations. (See.)
\end{openproblem}
% Source: Kirby
\plainproblemsection

\begin{openproblem}
	(Melvin) Let ${M}^{4}$ be a smooth closed orientable 4-manifold which supports an effective action of a compact connected Lie group $G$. Suppose that ${\pi }_{1}M$ is a free group.
	
	\emph{Question.} Is $M$ diffeomorphic to a connected sum of copies of ${S}^{1} \times  {S}^{3},{S}^{2} \times  {S}^{2}$ and ${S}^{2}\widetilde{ \times  }{S}^{2}$?
\end{openproblem}
% Source: Kirby
\plainproblemsection

\begin{openproblem}
	Classify closed 4-manifolds which fiber
	
	(A) over a circle with fiber an ${S}^{1}$ -manifold,
	
	(B) over a surface.
\end{openproblem}
% Source: Kirby
\plainproblemsection

\begin{openproblem}
	(Melvin) Let $P \subset  {S}^{4}$ be the standardly imbedded ${\mathbb{RP}}^{2}$ (e.g. $P = q\left( {\mathbb{RP}}^{2}\right)$, where $q: {\mathbb{CP}}^{2} \rightarrow  {S}^{4}$ is the quotient map by complex conjugation) and $K \subset  {S}^{4}$ be an odd twist spun knot. Denote by $\left( {{S}^{4}, P\# K}\right)$ the pairwise connected sum $\left( {{S}^{4}, P}\right) \# \left( {{S}^{4}, K}\right)$.
	
	Is $\left( {{S}^{4}, P\# K}\right)$ pairwise diffeomorphic to $\left( {{S}^{4}, P}\right)$?
\end{openproblem}
% Source: Kirby
\plainproblemsection

\begin{openproblem}
	(Hillman) Minimize the Euler characteristic over all closed 4-manifolds $M$ with ${\pi }_{1}\left( {M}^{4}\right)  = G$ given.
\end{openproblem}
% Source: Kirby
\plainproblemsection

\begin{openproblem}
	(Hass) Let ${M}^{4}$ be closed, smooth and satisfy ${\pi }_{2}\left( M\right)  = 0$ but ${\pi }_{3}\left( M\right)  \neq  0$, i.e. $L\left( {p, q}\right)  \times  {S}^{1}$. Is there a smooth imbedded 3-manifold ${L}^{3}$, with finite cover ${S}^{3}$, representing a non-zero element of ${\pi }_{3}\left( M\right)$?
\end{openproblem}
% Source: Kirby
\plainproblemsection

\begin{openproblem}
	(Hughes) (A) Find representatives for each regular homotopy class of immersions of ${S}^{n}$ in ${\mathbb{R}}^{n + k}$.
	
	(B) Find representatives for all bordism classes of immersions of $n$ -manifolds in ${\mathbb{R}}^{n + k}$.
	
	(C) For a surface in ${\mathbb{R}}^{3}$, a neighborhood of a double curve is an immersed ${B}^{1} \vee  {B}^{1}$ -bundle over ${S}^{1}$. In general a $k$ -tuple set will have an immersed ${B}^{n} \vee  { \cdot  }^{k} \cdot   \vee  {B}^{n}$ -bundle. Does the multiple point set with this structure determine the bordism class of the immersion?
	
	(D) Can one find explicit coordinates for Boy’s surface, i.e. find a smooth function from ${S}^{2}$ to ${\mathbb{R}}^{3}$ taking ${S}^{2}$ onto Boy’s surface as a 2-1 cover.
	
	(E) The number of quadruple points of an immersed ${S}^{3}$ in ${\mathbb{R}}^{4}$ is a $\mathbb{Z}/2\mathbb{Z}$ invariant under regular homotopy. Is it a $\mathbb{Z}/{24}\mathbb{Z}\left( { = {\pi }_{3}^{s}}\right)$ or even a $\mathbb{Z}$ invariant?
\end{openproblem}
% Source: Kirby
\plainproblemsection

\begin{openproblem}
	(A) Do the cyclic branched covers of 2-spheres in ${S}^{4}$ imbed in ${S}^{5}$?
	
	(B) Does every mapping torus of a 3-manifold imbed in ${S}^{5}$?
\end{openproblem}
% Source: Kirby
\plainproblemsection

\begin{openproblem}
	(A) Find a smooth, closed, spin, signature zero 4-manifold ${X}^{4}$ which does not imbed punctured in ${S}^{5}$.
	
	(B) Find an ${X}^{4}$ such that $X\# k{S}^{2} \times  {S}^{2}$ imbeds smoothly in ${S}^{5}$ but $X$ does not.
\end{openproblem}
% Source: Kirby
\plainproblemsection

\begin{openproblem}
	(A) What 4-manifolds have a symplectic structure?
	
	(B) Does every contact structure on a 3-manifold ${M}^{3}$ extend to symplectic structure on a bounding 4-manifold?
\end{openproblem}
% Source: Kirby
\plainproblemsection

\begin{openproblem}
	(A) Find a differential geometric invariant which distinguishes the ends of smooth non-compact 4-manifolds. For example, if ${X}^{4}$ is a simply connected topological manifold with a definite intersection form, then ${X}^{4}$ -point is smooth but its end is not standard; can this be detected in a direct differential geometric way?
	
	(B) Find a differential geometric proof of Rohlin's theorem.
\end{openproblem}
% Source: Kirby
\plainproblemsection

\begin{openproblem}
	How do metrics (e.g. Riemannian, Lorentz, constant curvature) behave under standard topological constructions such as connected sum, plumbing, handle addition? Same question for $\eta$ -invariants, moduli spaces, etc.
\end{openproblem}
% Source: Kirby
\plainproblemsection

\begin{openproblem}
	(Hopf) Does there exist a metric of strictly positive sectional curvature on ${S}^{2} \times  {S}^{2}$?
\end{openproblem}
% Source: Kirby
\plainproblemsection

\begin{openproblem}
	(A) There exists an anti-self-dual Einstein metric on the Kummer surface. Describe it explicitly.
	
	(B) If ${M}^{4}$ is compact, closed and has an Einstein metric, then $\chi \left( M\right)  \geq  \frac{3}{2}\left| {\sigma \left( M\right) }\right|$. Are there any other topological restrictions?
	
	(C) Does $\# p{\mathbb{CP}}^{2}\# q\left( {-{\mathbb{CP}}^{2}}\right)$ have an Einstein metric?
\end{openproblem}
% Source: Kirby
\plainproblemsection

\begin{openproblem}
	(Weinberger) (A) Is every 3-dimensional, ANR, homology manifold stably resolvable?
	
	(B) Construct 4-dimensional, nonresolvable, ANR homology manifolds.
	
	(C) Are these $X$ ’s homogeneous, or are all homogeneous, ANR, homology manifolds actually manifolds?
\end{openproblem}
% Source: Kirby
\plainproblemsection

\begin{openproblem}
	Let $f: X \rightarrow  Y$ be a cellular map between polyhedra and suppose $X$ or $Y$ is a 4-manifold. Can $f$ be approximated by a homeomorphism?
\end{openproblem}
% Source: Kirby
\plainproblemsection

\begin{openproblem}
	(Freedman) \emph{Conjecture.} All compact 4-manifolds are Hölder continuous.
\end{openproblem}
% Source: Kirby
\plainproblemsection

\begin{openproblem}
	Does there exist a closed, non-smoothable 4-manifold ${M}^{4}$ having a triangulation (necessarily non-combinatorial)?
\end{openproblem}
% Source: Kirby
\plainproblemsection

\begin{openproblem}
	(Ancel) Definition: $X$ is a pseudo-spine of a compact manifold $M$ if $X$ is a compact subset of int $M$ and $M - X$ is homeomorphic to $\partial M \times  \lbrack 0,1)$. In contrast, a topological spine $X$ is a compact subset of int $M$ such that $M$ is homeomorphic to the mapping cylinder of a map from $\partial M$ to $X$,.
	
	(A) Does the Mazur 4-manifold  have disjoint pseudo-spines? More generally, does every compact contractible n-manifold have disjoint pseudo-spines?
	
	(B) Does every compact, contractible 4-manifold have an arc pseudo-spine?
	
	The Mazur 4-manifold has a handlebody decomposition consisting of a single 0-handle, 1-handle and 2-handle. This suggests splitting question (B) into the following two questions:
	
	(B.1) Does every ${PL}$ compact contractible 4-manifold have a handlebody decomposition with no 3- or 4-handles?.
	
	(B.2) Does every compact contractible 4-manifold that has a handlebody decomposition with no 3- or 4- handles have an arc pseudo-spine?
	
	(C) If ${M}^{4}$ is a compact 4-manifold which is homotopy equivalent to a compact surface ${F}^{2}$, does $M$ have a pseudo-spine homeomorphic to ${F}^{2}$?
	
	(D) If a compact 4-manifold ${M}^{4}$ is homotopically equivalent to $X\left( {{n}_{1}\ldots {n}_{k}}\right),{n}_{i} \neq  0$, then does ${M}^{4}$ have a pseudo-spine that is homeomorphic to $X\left( {{n}_{1}\ldots {n}_{k}}\right)$?
\end{openproblem}
% Source: Kirby
\plainproblemsection

\begin{openproblem}
	Is every closed 4-manifold ${M}^{4}$ the union of a smooth 4-manifold $Y$ and an acyclic topological manifold $Z$, joined along their common homology sphere boundary? Can acyclic be improved to contractible?
\end{openproblem}
% Source: Kirby
\plainproblemsection

\begin{openproblem}
	(A) (Existence) Does every quasiconformal 4-manifold have a smooth structure?
	
	(B) (Uniqueness) Given smooth 4-manifolds $M$ and $N$, a homeomorphism ${h}_{0}: M \rightarrow  N$, closed sets $C \subset  D \subset  M$ and open sets $U \subset  V \subset  M$ such that $C \subset  U$ and $D \subset  V$. Assume ${h}_{0}$ is quasiconformal on $U$. Does there exist an isotopy ${h}_{t}: M \rightarrow  N$ such that ${h}_{1}$ is quasiconformal in a neighborhood of $D$, and ${h}_{t} = {h}_{0}$ near $C$ and on $N - V$?
\end{openproblem}
% Source: Kirby
\plainproblemsection

\begin{openproblem}
	Show that all Casson handles with exactly one kink at each stage are exotic (in the sense that the attaching circle does not bound a smooth 2-ball in the Casson handle).
\end{openproblem}
% Source: Kirby
\plainproblemsection

\begin{openproblem}
	An exotic smooth structure on ${\mathbb{R}}^{4}$ crossed with ${\mathbb{R}}^{1}$ is diffeomorphic to ${\mathbb{R}}^{5}$.
	
	(A) How can we usefully see the exotic ${\mathbb{R}}^{4}$ in ${\mathbb{R}}^{5}$?
	
	The diffeomorphism induces a smooth, non-zero vector field (the image of the tangents to ${\mathbb{R}}^{1}$ ) on ${\mathbb{R}}^{5}$, whose orbit space is diffeomorphic to the original exotic ${\mathbb{R}}^{4}$. So an exotic ${\mathbb{R}}^{4}$ can be described by 5 differentiable functions of 5 variables (the components of the vector field).
	
	(B) (Arnold) Can the vector field be improved beyond being ${C}^{\infty }$? Can its components be analytic? Trigonometric? Polynomial? Can it be written explicitly?
\end{openproblem}
% Source: Kirby
\plainproblemsection

\begin{openproblem}
	Which smooth structures on ${\mathbb{R}}^{4}$ have compatible Stein structures?
\end{openproblem}
% Source: Kirby
\plainproblemsection

\begin{openproblem}
	(Gompf) (A) Is there a closed 4-manifold covered by an exotic ${\mathbb{R}}^{4}$?
	
	(B) Is there a smooth involution on the standard ${\mathbb{R}}^{4}$, topologically equivalent to Id. $\times   - {Id}$. on ${\mathbb{R}}^{2} \times  {\mathbb{R}}^{2}$, whose quotient is an exotic ${\mathbb{R}}^{4}$?
\end{openproblem}
% Source: Kirby
\plainproblemsection

\begin{openproblem}
	(Gompf) Freedman \& Taylor  constructed a universal half space $H$, that is, a smoothing on $\frac{1}{2}{\mathbb{R}}^{4} = \left\{  {\left( {{x}_{1},{x}_{2},{x}_{3},{x}_{4}}\right)  \mid  {x}_{4} \geq  0}\right\}$ such that any other smooth 4-manifold homeomorphic to $\frac{1}{2}{\mathbb{R}}^{4}$ smoothly imbeds in $H$. They prove $H$ is unique up to diffeomorphism, and that int $H = U$ is a universal smoothing of ${\mathbb{R}}^{4}$.
	
	(A) Is $U$ unique (as a smoothing of ${\mathbb{R}}^{4}$ into which all other smoothings of ${\mathbb{R}}^{4}$ smoothly imbed) up to diffeomorphism?
	
	(B) If $R \leq  {R}^{\prime }$, does it follow that all of $R$ must smoothly imbed in ${R}^{\prime }$?
	
	The class of $U$ is the unique maximal element in this partial ordering, and, assuming that (B) is true,(A) is equivalent to the class of $U$ contains only $U$. The unique minimal element is the class of ${\mathbb{R}}^{4}$ and it is one of uncountably many classes each of which has uncountably many elements.
	
	(C) Do all compact equivalence classes have uncountably many elements?
\end{openproblem}
% Source: Kirby
\plainproblemsection

\begin{openproblem}
	(Teichner) Let $f: {S}^{2} \rightarrow  {\mathbb{RP}}^{2} \rightarrow  {\mathbb{RP}}^{2} \times  {B}^{2}$ be the composition of the covering map and the inclusion. Is ${2f}$ homotopic to a locally flat topological imbedding? A smooth imbedding?
\end{openproblem}
% Source: Kirby
\plainproblemsection

\begin{openproblem}
	(Teichner) Does $\left( {*{\mathbb{RP}}^{4}}\right) \# \left( {*{\mathbb{CP}}^{2}}\right)$ have a smooth structure?
\end{openproblem}
% Source: Kirby
\plainproblemsection

\begin{openproblem}
	(Weinberger) If ${X}^{4}$ is a closed 4-manifold and is homotopy equivalent to a hyperbolic 4-manifold, does $X$ have a smooth structure? Is $X$ then a hyperbolic manifold?
\end{openproblem}
% Source: Kirby
\plainproblemsection

\begin{openproblem}
	Let \(M\) and \(N\) be homotopy-equivalent closed
	\(4\)-manifolds with the same Kirby--Siebenmann invariant.
	\begin{enumerate}[label=\textup{(\Alph*)}]
		\item If \(M\) and \(N\) are spin, must they be stably homeomorphic?
		Must they be homeomorphic?
		\item Let \(S\) be a spherical \(3\)-manifold with
		\(\pi_1(S)=\pi\), and form \(M\) and \(N\) from \(S\times S^1\) by
		surgery on \(\{*\}\times S^1\), using the two possible normal framings.
		If
		\[
		\pi=\Delta^*(2,2,2n+1),
		\]
		are \(M\) and \(N\) diffeomorphic?
	\end{enumerate}
\end{openproblem}
% Source: Kirby
\plainproblemsection

\begin{openproblem}
	(Gompf \& Taylor) Are there any exotic smooth structures on closed, ori-entable 4-manifolds that can be distinguished by classical (non gauge theoretic) invariants?
\end{openproblem}
% Source: Kirby
\plainproblemsection

\begin{openproblem}
	Do all closed, smooth 4-manifolds have more than one smooth structure? Easier is whether all algebraic surfaces have more than one smooth structure.
\end{openproblem}
% Source: Kirby
\plainproblemsection

\begin{openproblem}
	Does every non-compact, smooth 4-manifold have an uncountable number of smoothings?
\end{openproblem}
% Source: Kirby
\plainproblemsection

\begin{openproblem}
	(Poénaru) Define a manifold to be geometrically simply connected if it has a handle body decomposition without 1-handles. Call it geometrically simply connected at long distance (GSCLD) if its interior has an exhaustion by compact geometrically simply connected manifolds.
	
	Let $\Delta$ be a homotopy 3-ball.
	
	Conjecture (Poénaru): If $\Delta  \times  I$ is GSCLD, then $\Delta  \times  I$ is geometrically simply connected.
\end{openproblem}
% Source: Kirby
\plainproblemsection

\begin{openproblem}
	Smooth Poincaré \emph{Conjecture.} A smooth homotopy 4-sphere $\sum$ is diffeo- morphic to ${S}^{4}$.
\end{openproblem}
% Source: Kirby
\plainproblemsection

\begin{openproblem}
	concerns one possible source, the Gluck construction, of potential counterexamples to this Conjecture. Another source, potential counterexamples to the Andrews-Curtis Conjecture (i.e. add 1 and 2-handles to ${B}^{4}$ according to any finite presentation of the trivial group and then double), suggests breaking this Problem into two parts:
	
	(A) Can every homotopy 4-sphere be described without using 3-handles?
	
	(B) Is the double of a homotopy 4-sphere without 3-handles diffeomorphic to ${S}^{4}$?
	
	(C) Can $B$ be constructed so that $\partial B = {S}^{3}$? So that $B$ is diffeomorphic to ${B}^{4}$?
\end{openproblem}
% Source: Kirby
\plainproblemsection

\begin{openproblem}
	Form a 5-dimensional $h$ -cobordism, ${W}^{5}$, between two contractible 4-manifolds, which is trivial over the boundary, as follows: choose a link in ${S}^{3} = \partial {B}^{4}$ consisting of two components, ${K}_{0},{K}_{1}$, each being the unknot, which link algebraically once. Add two 2- handles to ${B}^{4}$ along each unknot with framing zero and call the result ${X}^{4}.{X}^{4}$ has two obvious smoothly imbedded 2-spheres, namely the cores of the 2-handles union the obvious disks, ${D}_{0},{D}_{1}$, in ${B}^{4}$ which the unknots bound. Starting with ${X}^{4} \times  I$, add a 3-handle to one 2-sphere in ${X}^{4} \times  0$ and another 3-handle to the other 2-sphere in ${X}^{4} \times  1$. The result is ${W}^{5}$. The portion of $\partial {W}^{5}$ which is $\left( {\partial {X}^{4}}\right)  \times  I$ is the (obviously trivial) side of the $h$ -cobordism $W$; the bottom and top of the $h$ -cobordism consist of the two contractible 4-manifolds, ${\partial }_{0}W$ and ${\partial }_{1}W$, made by adding a 2-handle to ${B}^{4}$ along ${K}_{1}$ (resp. ${K}_{0}$ ) and excising a thickened ${D}_{0}$ (resp. ${D}_{1}$ ) from ${B}^{4}$ (or equivalently putting a ${dot}$ on ${K}_{0}$ (resp. ${K}_{1}$ ) in the language of the framed link calculus ).
	
	(A) Under what conditions on ${K}_{0} \cup  {K}_{1}$ is $W$ a product h-cobordism relative to the given product structure on the boundary?
	
	(B) Same question: under what conditions on ${L}_{0} \cup  {L}_{1}$ does the product structure on the sides of ${W}^{5}$ extend to a product structure on ${W}^{5}$?
	
	(C) Can $U$ always be chosen to be ${B}^{4}$ so that (B) is the general case? In other words, in the construction in the sketch, can 2-handles be slid so that all the 1-handles are geometrically cancelled?
\end{openproblem}
% Source: Kirby
\plainproblemsection

\begin{openproblem}
	(Shkolnikov) Suppose that ${M}^{4}$ is closed, has no 1 or 3-handles and its signature is zero. Does it follow that $M$ is diffeomorphic to $\# k\left( {{S}^{2} \times  {S}^{2}}\right) \# l\left( {{S}^{2}\widetilde{ \times  }{S}^{2}}\right)$?
\end{openproblem}
% Source: Kirby
\plainproblemsection

\begin{openproblem}
	(11/8-Conjecture) Given a smooth, closed, spin 4-manifold, ${X}^{4}$, with ${H}_{1}\left( {{X}^{4};\mathbb{Z}}\right)  = 0$, then ${\beta }_{2}/\left| \sigma \right|  \geq  {11}/8$.
\end{openproblem}
% Source: Kirby
\plainproblemsection

\begin{openproblem}
	(3/2-Conjecture) Given an irreducible, simply connected, closed, smooth, spin 4-manifold ${X}^{4}$, then $\chi \left( X\right) /\left| {\sigma \left( X\right) }\right|  \geq  3/2$.
\end{openproblem}
% Source: Kirby
\plainproblemsection

\begin{openproblem}
	(Mikhalkin) Is it possible to find closed, smooth, (simply connected would be nice) 4-manifolds $X$, not of the form ${M}^{4}\#  \pm  k{\mathbb{CP}}^{2}$, so that the fraction ${b}_{2}\left( X\right) /\left| {\sigma \left( X\right) }\right|$ is arbitrarily close to 1?
\end{openproblem}
% Source: Kirby
\plainproblemsection

\begin{openproblem}
	(Akbulut) Does Rohlin's Theorem still hold for smooth, spin 4-manifolds with simply connected periodic ends? More specifically, suppose ${X}^{4}$ is spin, $\sigma \left( {X}^{4}\right)  \equiv  8{\;\operatorname{mod}\;}$ 16, and $\partial {X}^{4}$ is a homology 3-sphere. Can there be a simply connected, homology bordism $W$ with $\partial W = \partial X \cup   - \partial X$, so that ${X}^{4}$ union $\{$ countably many copies of $W\}$ would contradict an extended Rohlin's Theorem for periodic ends?
\end{openproblem}
% Source: Kirby
\plainproblemsection

\begin{openproblem}
	We will say that an intersection form is of type(1, n)if it is equivalent (over $\mathbb{R})$ to $\langle 1\rangle  \oplus  n\langle  - 1\rangle$.
	
	Are there any smooth, closed, irreducible, simply connected 4-manifolds of type(1, n) other that ${\mathbb{CP}}^{2},{S}^{2} \times  {S}^{2}$, the Barlow surface of type(1,8), the Dolgachev surfaces of type (1,9), or blowups of any of these?
	
	If so, are any of them algebraic surfaces?
\end{openproblem}
% Source: Kirby
\plainproblemsection

\begin{openproblem}
	(A) \emph{Conjecture.} Every irreducible, smooth, closed, simply connected 4- manifold ${X}^{4}$ except ${S}^{4}$ has an almost complex structure.
	
	(B) \emph{Conjecture.} A smooth, closed, simply connected 4-manifold with even ${b}_{2}^{ + }$ and even ${b}_{2}^{ - }$ decomposes as a connected sum (or is ${S}^{4}$ ).
\end{openproblem}
% Source: Kirby
\plainproblemsection

\begin{openproblem}
	(Yau) Classify smooth, oriented,4-manifolds ${X}^{4}$ with boundary which have a projection to ${B}^{2}$ which is a surface bundle over ${B}^{2} - 0$ and whose fiber over 0 consists of smooth imbedded surfaces in ${X}^{4}$ with normal crossings.
\end{openproblem}
% Source: Kirby
\plainproblemsection

\begin{openproblem}
	(Y. Matsumoto) A ${C}^{\infty }$ -Lefschetz fibration of fiber genus $g$ over a surface $F$ is determined up to ${C}^{\infty }$ -isomorphism by the conjugacy class of the monodromy representation
	
	$$
	\rho: {\pi }_{1}\left( {F - \left\{  {{p}_{1},{p}_{2},\cdots,{p}_{s}}\right\}  }\right)  \rightarrow  {\Gamma }_{g},
	$$
	
	where $\left\{  {{p}_{1},\cdots,{p}_{s}}\right\}$ is the set of singular loci, and ${\Gamma }_{g}$ is the mapping class group of genus $g$. In what follows, $F$ will be ${S}^{2}$ and in this case, the monodromy representation is described by the $s$ -tuple of elements of ${\Gamma }_{g},\rho \left( {p}_{1}\right),\rho \left( {p}_{2}\right),\cdots,\rho \left( {p}_{s}\right)$, each of which is a right handed Dehn twist about a simple closed curve on the general fiber. This satisfies $\rho \left( {p}_{1}\right) \rho \left( {p}_{2}\right) \cdots \rho \left( {p}_{s}\right)  = 1$. Conversely if the latter relation is satisfied, then we obtain a ${C}^{\infty }$ -Lefschetz fibration of fiber genus $g$ over ${S}^{2}$. (Here we use the same notation ${p}_{i}$ for the small loop around the point ${p}_{i}$.)
	
	The mapping class group of genus $2,{\Gamma }_{2}$, is generated by the 5 generators ${\zeta }_{1},{\zeta }_{2},\cdots,{\zeta }_{5}$, which are represented by full right-handed Dehn twists about canonical curves, (see ). We have the following relations
	
	$$
	{\left( {\zeta }_{1}{\zeta }_{2}{\zeta }_{3}{\zeta }_{4}{\zeta }_{5}^{2}{\zeta }_{4}{\zeta }_{3}{\zeta }_{2}{\zeta }_{1}\right) }^{4} = 1
	$$
	
	$$
	{\left( {\zeta }_{1}{\zeta }_{2}{\zeta }_{3}{\zeta }_{4}\right) }^{10} = 1
	$$
	
	The first relation is an immediate consequence of the relations in Birman's book, and the second one is algebraically derived from them. These two relations correspond to ${C}^{\infty } -$ Lefschetz fibrations ${M}_{1} \rightarrow  {S}^{2}$ and ${M}_{2} \rightarrow  {S}^{2}$ of fiber genus two, respectively. The total spaces ${M}_{1}$ and ${M}_{2}$ are simply connected 4-manifolds. They have the same signature,-24, and Euler number, 36, and therefore  are homeomorphic to $5{\mathbb{CP}}^{2}\# {29}\left( {-{\mathbb{CP}}^{2}}\right)$. However, they are not ${C}^{\infty }$ -isomorphic as Lefschetz fibrations, because the monodromy representation of ${M}_{1} \rightarrow  {S}^{2}$ is onto while that of ${M}_{2} \rightarrow  {S}^{2}$ is not.
	
	\emph{Question.} Are ${M}_{1},{M}_{2}$, and $5{\mathbb{CP}}^{2}\# {29}\left( {-{\mathbb{CP}}^{2}}\right)$ mutually diffeomorphic? (Probably not.)
\end{openproblem}
% Source: Kirby
\plainproblemsection

\begin{openproblem}
	(Stern) For fixed odd $r > 1$, the following pairs $H\left( r\right)$ and ${H}^{\prime }\left( r\right)$ of complex algebraic surfaces of general type are homotopy equivalent but deformation inequivalent. They cannot be distinguished by their Seiberg-Witten invariants; D. Gomprecht, Fintushel \& Stern and Morgan \& Szabó have shown they also have the same Donaldson polynomials. (A) \emph{Question.} Are $H\left( r\right)$ and ${H}^{\prime }\left( r\right)$ diffeomorphic?
	
	(B) \emph{Question.} Are there restrictions on self-diffeomorphisms $f$ of a minimal Kähler surface $X$ (of nonnegative Kodaira dimension) with canonical class ${K}_{X}$ beyond the conditions ${f}^{ * }{K}_{X} =  \pm  {K}_{X}?$
\end{openproblem}
% Source: Kirby
\plainproblemsection

\begin{openproblem}
	(Yau) \emph{Conjecture.} The minimal, complex surfaces which are not Kähler, are the manifolds of class ${VI}{I}_{0}$ as listed by $M$. Inoue.
\end{openproblem}
% Source: Kirby
\plainproblemsection

\begin{openproblem}
	Suppose ${X}^{4}$ is a simply connected, algebraic surface.
	
	(A) Is ${X}^{4}\# {\mathbb{CP}}^{2}$ diffeomorphic to a connected sum of copies of ${\mathbb{CP}}^{2}$ and $- {\mathbb{CP}}^{2}$?
	
	(B) If ${X}^{4}$ is also spin, is ${X}^{4}\# {S}^{2} \times  {S}^{2}$ diffeomorphic to a connected sum of copies of ${K3}$ and ${S}^{2} \times  {S}^{2}$?
	
	(C) If $X$ and $Y$ are simply connected algebraic surfaces (or symplectic 4-manifolds) other than ${\mathbb{CP}}^{2}$, does $X\#  - Y$ decompose as in (A) for the non-spin case, or (B) for the spin case?
	
	(D) How about $X\# Y$?
\end{openproblem}
% Source: Kirby
\plainproblemsection

\begin{openproblem}
	(Mikhalkin) Let ${X}^{4}$ be a simply connected, closed, symplectic 4-manifold, and let $\tau: X \rightarrow  X$ be an anti-symplectic $\left( {\tau \left( \omega \right)  =  - \omega }\right)$ involution with a smooth, non-empty, imbedded surface as fixed point set.
	
	\emph{Question.} Is $X/\tau$ completely decomposable, i.e. is
	
	$$
	X/\tau  = \# r{\mathbb{CP}}^{2}\# s\left( {-{\mathbb{CP}}^{2}}\right) \;\text{ or }\;\# n\left( {{S}^{2} \times  {S}^{2}}\right)?
	$$
\end{openproblem}
% Source: Kirby
\plainproblemsection

\begin{openproblem}
	(A) Given an arbitrary smooth 4-manifold, ${X}^{4}$, is the self-intersection of a smoothly imbedded ${S}^{2}$ bounded below by a constant which depends only on ${X}^{4}$?
	
	(B) What is the smoothly imbedded 2-sphere with minimal (most negative) self intersection in the $- {E}_{8}$ plumbing?
	
	(C) Same question for the plumbing of only two 2-spheres, each with self intersection -2. Can one improve on -6?
	
	(D) Same question for the K3 surface.
\end{openproblem}
% Source: Kirby
\plainproblemsection

\begin{openproblem}
	Does there exist a simply connected, algebraic surface having elements $\alpha  \in$ ${H}_{2}\left( {X;\mathbb{Z}}\right)$ such that $\alpha  \cdot  \alpha  =  - 2$, but no smoothly imbedded 2-sphere representing any such class?
\end{openproblem}
% Source: Kirby
\plainproblemsection

\begin{openproblem}
	(Catanese) Does there exist an algebraic surface for which in $f\left( {C \cdot  C}\right)  =$ $- \infty$ where the infimum is taken over all irreducible, complex curves $C$?
\end{openproblem}
% Source: Kirby
\plainproblemsection

\begin{openproblem}
	(Ruberman) Does every smooth, closed, simply connected 4-manifold contain a configuration of 2-spheres such that the complement is a $K\left( {\pi,1}\right)$?
\end{openproblem}
% Source: Kirby
\plainproblemsection

\begin{openproblem}
	Let $X$ be a simply connected 4-manifold and $\sum$ an imbedded surface. If $\sum  \cdot  \sum$ is non-zero and is square free (all prime divisors appear once), does it follow that ${\pi }_{1}\left( {X - \sum }\right)$ is trivial?
\end{openproblem}
% Source: Kirby
\plainproblemsection

\begin{openproblem}
	In ${\mathbb{CP}}^{2}$, let ${F}^{2}$ (smoothly imbedded) represent $\alpha  \in  {H}_{2}\left( {{\mathbb{CP}}^{2};\mathbb{Z}}\right)  = \mathbb{Z}$, suppose genus(F)is minimal and suppose that ${\pi }_{1}\left( {{\mathbb{CP}}^{2} - F}\right)  = \mathbb{Z}/\alpha \mathbb{Z}$.
	
	\emph{Question.} Is $F$ smoothly isotopic to an algebraic curve?
\end{openproblem}
% Source: Kirby
\plainproblemsection

\begin{openproblem}
	(Taylor) Let ${M}^{4}$ be an orientable 4-manifold.
	
	(A) Is there an imbedded (locally flat) surface $F$, dual to ${w}_{2}\left( M\right)$, such that ${H}_{1}\left( {F;\mathbb{Z}/2\mathbb{Z}}\right)  \rightarrow$ ${H}_{1}\left( {M;\mathbb{Z}/2\mathbb{Z}}\right)$ is the zero map?
	
	(B) Given $x \in  {H}_{2}\left( {M;\mathbb{Z}/2\mathbb{Z}}\right), x$ dual to ${w}_{2}\left( M\right)$, then for which covers of $M,\pi: \widetilde{M} \rightarrow  M$, does there exist $y \in  {H}_{2}\left( {\widetilde{M};\mathbb{Z}/2\mathbb{Z}}\right)$ such that ${\pi }_{ * }\left( y\right)  = x$? In particular
	
	$\left( {B}^{\prime }\right)$ does $y$ exist for the universal cover of $M$?
	
	$\left( {B}^{\prime \prime }\right)$ does $y$ exist for the cover corresponding to the kernel of ${\pi }_{1}\left( M\right)  \rightarrow  {H}_{1}\left( {M;\mathbb{Z}/2\mathbb{Z}}\right)$?
\end{openproblem}
% Source: Kirby
\plainproblemsection

\begin{openproblem}
	(A) (Montesinos) Is every closed, smooth, orientable 4-manifold ${X}^{4}$ an irregular, simple,4-fold cover of ${S}^{4}$, branched along a closed surface in ${S}^{4}$?
	
	(B) (Akbulut) Is every simply connected 4-manifold ${X}^{4}$ a cyclic branched cover of ${S}^{4}$, or $\# r\left( {\mathbb{CP}}^{2}\right) \# s\left( {-{\mathbb{CP}}^{2}}\right)$, or $\# n\left( {{S}^{2} \times  {S}^{2}}\right)$, branched over a smooth imbedded surface?
	
	(C) (Akbulut) Does every irreducible, simply connected, closed, smooth 4-manifold ${X}^{4}$ have a branched cover which is a symplectic manifold.
\end{openproblem}
% Source: Kirby
\plainproblemsection

\begin{openproblem}
	\emph{Conjecture.} Let ${M}^{3}$ be a homology 3 -sphere which does not bound a smooth, contractible 4-manifold; then for any homology 3-sphere ${N}^{3}, M\# N$ also does not bound a smooth, contractible 4-manifold.
\end{openproblem}
% Source: Kirby
\plainproblemsection

\begin{openproblem}
	(Ruberman) Given a smooth ${X}^{4}$, if ${H}_{2}\left( {X;\mathbb{Z}}\right)$ has an orthogonal splitting then it can be realized by a splitting of ${X}^{4}$ by a homology 3-sphere $\sum$, that is, ${X}^{4} = {M}_{1}\mathop{\bigcup }\limits_{\sum }{M}_{2}$.
	
	Find an $X$ and a splitting of ${H}_{2}\left( X\right)$ where $\sum$ cannot be chosen to be a Seifert fibered 3-manifold.
\end{openproblem}
% Source: Kirby
\plainproblemsection

\begin{openproblem}
	(Stern) Which $\sum \left( {p, q, r}\right)$ imbed in the standard K3 surface?
	
	\emph{Conjecture.} Only finitely many.
\end{openproblem}
% Source: Kirby
\plainproblemsection

\begin{openproblem}
	Find examples of aspherical 4-manifolds with ${\pi }_{1}$ -injective, immersed,3- manifolds with infinite fundamental group. Are these common?
\end{openproblem}
% Source: Kirby
\plainproblemsection

\begin{openproblem}
	(Kotschick) For a finitely presentable group $\Gamma$, define ${q}^{CAT}\left( \Gamma \right)$ to be the minimum Euler characteristic of a closed oriented CAT 4-manifold ${X}^{4}$ with ${\pi }_{1}\left( X\right)  = \Gamma$.
	
	(A) Is ${q}^{TOP}\left( \Gamma \right)  = {q}^{DIFF}\left( \Gamma \right)$?
	
	Let ${p}^{CAT}\left( \Gamma \right)$ be the minimal of $\chi \left( {X}^{4}\right)  - \sigma \left( {X}^{4}\right)$ over CAT 4-manifolds $X$ as above.
	
	(B) Is ${p}^{TOP}\left( \Gamma \right)  = {p}^{DIFF}\left( \Gamma \right)$?
	
	(C) Is there an integer $k = k\left( {{\Gamma }_{1},{\Gamma }_{2}}\right)$ such that for all CAT manifolds ${X}^{4}$ as above with ${\pi }_{1}\left( X\right)  = {\Gamma }_{1} * {\Gamma }_{2}$, the connected sum $X\# k\left( {{S}^{2} \times  {S}^{2}}\right)$ CAT splits as a connected sum ${X}_{1}\# {X}_{2}$ with ${\pi }_{1}\left( {X}_{i}\right)  = {\Gamma }_{i}$?
	
	(D) What can be said about ${q}^{\text{Symplectic }}$ and ${p}^{\text{Sympl }}$?
\end{openproblem}
% Source: Kirby
\plainproblemsection

\begin{openproblem}
	(Weinberger) Let $G$ be a group on Milnor’s list  of possible groups which can act on ${S}^{3}$. Show that unless $G$ is a subgroup of ${SO}\left( 4\right)$, $\mathbb{Z} \times  G$ is not the fundamental group of any smooth 4-manifold with zero Euler characteristic.
\end{openproblem}
% Source: Kirby
\plainproblemsection

\begin{openproblem}
	(Fukaya) Is there a bundle over ${S}^{2}$, with a 4-manifold as fiber, which is topologically trivial but not smoothly trivial?
\end{openproblem}
% Source: Kirby
\plainproblemsection

\begin{openproblem}
	(Fintushel \& Stern) Does there exist a pseudo-free, smooth, ${S}^{1}$ action on ${S}^{5}$ with more than three multiple orbits?
\end{openproblem}
% Source: Kirby
\plainproblemsection

\begin{openproblem}
	\begin{enumerate}[label=\textup{(\Alph*)}]
		\item Does a fake \(\mathbb{CP}^{2}\) admit a topological involution?
		\item Does a \(K3\) surface admit a periodic diffeomorphism of odd order
		that acts trivially on homology?
	\end{enumerate}
\end{openproblem}
% Source: Kirby
\plainproblemsection

\begin{openproblem}
	(A) Which closed, smooth, simply connected 4-manifolds $X$ have big diffeomorphism groups with respect to some $\kappa  \in  {H}_{2}\left( {X;\mathbb{Z}}\right)$?
	
	(B) Here is a possible partial answer to (A). Let ${\kappa }_{1},\ldots,{\kappa }_{l}$ be the basic classes on a smooth, closed ${b}_{1} = 0,{b}_{2}^{ + } \geq  3,4$ -manifold $X$ of simple type, and let $H$ be the subspace of ${H}_{2}\left( {X;\mathbb{Z}}\right)$ spanned by the basic classes.
	
	\emph{Conjecture.} The image of ${\operatorname{Diff}}_{H}\left( X\right)  = \operatorname{Diff}\left( X\right)$ in ${\operatorname{Aut}}_{H}\left( {{H}_{2}\left( {X;\mathbb{Z}}\right) }\right)$ has finite index.
\end{openproblem}
% Source: Kirby
\plainproblemsection

\begin{openproblem}
	For a closed PL manifold \(M\), let \(C^{\mathrm{PL}}(M)\)
	and \(C^{\mathrm{TOP}}(M)\) denote its PL and topological pseudo-isotopy
	spaces, respectively.
	\begin{enumerate}[label=\textup{(\Alph*)}]
		\item For every closed simply connected smooth \(4\)-manifold \(M\), do
		\(C^{\mathrm{PL}}(M)\) and \(C^{\mathrm{TOP}}(M)\) have the same homotopy
		type?
		\item Is \(C^{\mathrm{PL}}(S^3)\) contractible?
		\item Is the inclusion
		\[
		\operatorname{SO}(5)\longrightarrow
		\operatorname{SDiff}(S^4)
		\]
		a homotopy equivalence?
		\item Is
		\[
		C^{\mathrm{TOP}}(S^3)\simeq K(\mathbb Z/2\mathbb Z,2)?
		\]
	\end{enumerate}
\end{openproblem}
% Source: Kirby
\plainproblemsection

\begin{openproblem}
	(Kronheimer) Let ${X}^{4}$ be an oriented, closed 4-manifold with non-zero Donaldson invariants. Define $\parallel \alpha \parallel$ to be
	
	$$
	\parallel \alpha \parallel  = {2g} - 2 - {F}_{g} \cdot  {F}_{g}
	$$
	
	where $\alpha  \in  {H}_{2}\left( {X;\mathbb{Z}}\right)$ and ${F}_{g}$ is an imbedded surface of minimal genus $g$ which represents $\alpha$.
	
	\emph{Question.} Is $\parallel$ I $a$ seminorm on the positive cone in ${H}_{2}$?
\end{openproblem}
% Source: Kirby
\plainproblemsection

\begin{openproblem}
	(A) Is the Donaldson series for a smooth, closed, simply connected 4- manifold ${X}^{4}$ determined by the following data: the intersection form on ${H}_{2}\left( {X;\mathbb{Z}}\right)$, and the minimal genera of smoothly imbedded surfaces representing an appropriately chosen finite set $\mathcal{S}$ of elements in ${H}_{2}\left( {X;\mathbb{Z}}\right)$?
	
	(B) \emph{Conjecture.} Any basic class $K$ satisfies the equality
	
	$$
	K \cdot  K = {2\chi } + {3\sigma }
	$$
	
	(C) If $X$ is a minimal, algebraic surface of general type, does the following equality hold:
	
	$$
	D = {2}^{{c}_{1}^{2} + 3 - \left( {1 - {b}_{1} + {b}_{2}^{ + }}\right) /2}\exp \left( \frac{\mathcal{Q}}{2}\right) \cosh \left( {K}_{X}\right) \text{ for }{b}_{2}^{ + } \equiv  3\;\left( {\;\operatorname{mod}\;4}\right)
	$$
	
	with cosh replaced by sinh if ${b}_{2}^{ + } \equiv  1\;\left( {\;\operatorname{mod}\;4}\right)$, where ${K}_{X}$ is the canonical class?
	
	(D) Given a class $\alpha  \in  {H}_{2}\left( {X;\mathbb{Z}}\right)$, is there a smoothly imbedded surface $\sum$ of genus $g$ representing $\alpha$ which satisfies
	
	$$
	{2g} - 2 = \alpha  \cdot  \alpha  + \mathop{\max }\limits_{s}\left| {{K}_{s} \cdot  \alpha }\right|?
	$$
\end{openproblem}
% Source: Kirby
\plainproblemsection

\begin{openproblem}
	If ${X}^{4}$ is irreducible, simply connected, admits an almost complex structure and satisfies ${b}_{2}^{ + } > 0$, does it have non-zero Donaldson invariants? Seiberg-Witten invariants?
\end{openproblem}
% Source: Kirby
\plainproblemsection

\begin{openproblem}
	(Kronheimer \& Mrowka) Do all simply connected, closed, smooth 4- manifolds with ${b}_{2}^{ + } \geq  3$ have simple type?
\end{openproblem}
% Source: Kirby
\plainproblemsection

\begin{openproblem}
	(Kotschick) \emph{Conjecture.} The parity of Donaldson polynomials is a homotopy invariant; that is, if $f: X \rightarrow  Y$ is a homotopy equivalence, then ${f}^{ * }\left( {q}_{Y}\right)  = {q}_{X}$ in ${\operatorname{Sym}}^{d}\left( {{H}^{2}\left( {X;\mathbb{Z}/2\mathbb{Z}}\right) }\right)$.
\end{openproblem}
% Source: Kirby
\plainproblemsection

\begin{openproblem}
	(Kotschick) (A) \emph{Conjecture.} There is no orientable, connected and simply connected, smooth 4-manifold for which Donaldson's polynomial invariants are defined and non-zero for both choices of orientation.
	
	(B) \emph{Conjecture.} If two complex algebraic surfaces with finite fundamental groups are orientation reversing diffeomorphic (with respect to their complex orientations), then they are homeomorphic to a geometrically ruled surface, and in particular are simply connected.
\end{openproblem}
% Source: Kirby
\plainproblemsection

\begin{openproblem}
	(Ruan) The first betti number ${\beta }_{1}$ of a complex surface is even iff the surface is Kähler. A Kähler form is also a symplectic form. If two Kähler manifolds are deformation equivalent then the associated symplectic manifolds are symplectic deformation equivalent.
	
	\emph{Question.} Is the converse true?
\end{openproblem}
% Source: Kirby
\plainproblemsection

\begin{openproblem}
	(Ruan) Let $X$ and $Y$ be symplectic 4-manifolds with compatible almost complex structures. Assume that ${c}_{1}\left( X\right)  = {c}_{1}\left( Y\right)$.
	
	\emph{Conjecture.} $X$ and $Y$ are diffeomorphic iff $X \times  {S}^{2}$ and $Y \times  {S}^{2}$ (with the product symplectic structure) are symplectic deformation equivalent.
\end{openproblem}
% Source: Kirby
\plainproblemsection

\begin{openproblem}
	(Taubes) A generic, self dual or ASD closed 2-form is a symplectic 2- form on the complement of a 1-manifold (in ${X}^{4}$ ). For such degenerate symplectic forms, formulate a theory of pseudo-holomorphic curves (smooth imbedded surfaces on which the symplectic form restricts to a volume form).
\end{openproblem}
% Source: Kirby
\plainproblemsection

\begin{openproblem}
	(Gompf) (A) \emph{Conjecture.} Any closed, minimal, symplectic 4-manifold is irreducible.
	
	(B) If ${X}^{4}$ is irreducible, simply connected, has ${b}_{2}^{ + } > 0$, and has an almost complex structure, then does it have a symplectic structure?
\end{openproblem}
% Source: Kirby
\plainproblemsection

\begin{openproblem}
	(A) (McDuff) Let ${X}^{4}$ be a closed, minimal, symplectic 4-manifold containing a symplectic surface ${F}^{2}$ satisfying ${c}_{1}\left( {F}^{2}\right)  > 0$. Does it follow that $X$ must be rational or ruled?
	
	(B) (Gompf) If $X$ is closed and symplectic, and ${3\sigma }\left( X\right)  > \chi \left( X\right)  =$ ordinary Euler characteristic, does it follow that $X$ is ruled?
\end{openproblem}
% Source: Kirby
\plainproblemsection

\begin{openproblem}
	Call an open symplectic \(4\)-manifold \((M,\omega)\)
	\emph{convex} if it admits a complete vector field \(X\) such that
	\(\mathcal L_X\omega=-\omega\) and an exhaustion
	\[
	K_1\subset K_2\subset\cdots\subset M
	\]
	by compact sets invariant under the positive flow of \(X\).  Determine:
	\begin{enumerate}[label=\textup{(\Alph*)}]
		\item whether every exotic \(\mathbb R^4\) admits a convex symplectic
		structure;
		\item whether the standard \(\mathbb R^4\) admits inequivalent convex
		symplectic structures;
		\item whether \(S^2\times\mathbb R^2\) admits a convex symplectic
		structure;
		\item whether the space of symplectic structures on \(\mathbb R^4\) that
		are standard at infinity is contractible, or at least connected.
	\end{enumerate}
\end{openproblem}
% Source: Kirby
\plainproblemsection

\begin{openproblem}
	Does there exist a simply connected closed \(4\)-manifold,
	other than a manifold of either form
	\[
	\mathop{\#}_{n}(S^2\times S^2)
	\qquad\text{or}\qquad
	\mathop{\#}_{p}\mathbb{CP}^{2}
	\mathbin{\#}
	\mathop{\#}_{q}(-\mathbb{CP}^{2}),
	\]
	that admits a metric of positive scalar curvature and is either non-spin,
	or spin with zero signature?
\end{openproblem}
% Source: Kirby
\plainproblemsection

\begin{openproblem}
	(Suciu) Consider a finite union $\mathcal{H}$ of complex planes in ${\mathbb{C}}^{3}$ through the origin. We have the Milnor fibration
	
	$$
	{\mathbb{C}}^{3} \smallsetminus  \mathcal{H} \rightarrow  {\mathbb{C}}^{ \star  }
	$$
	
	given by taking the product of the linear forms which define the planes.
	
	Questions: Does the combinatorial data determine the homology of the Milnor fiber? the characteristic polynomial of the monodromy?
\end{openproblem}
% Source: Kirby
\plainproblemsection

\begin{openproblem}
	(Suciu) Let $\mathcal{H}$ be an arrangement of planes in ${\mathbb{C}}^{3}$ through the origin. Does the combinatorial data determine whether the complement of $\mathcal{H}$ is a $K\left( {\pi,1}\right)$?
\end{openproblem}
% Source: Kirby
\plainproblemsection

\begin{openproblem}
	(Suciu) Consider a finite union $\mathcal{L}$ of projective lines in ${\mathbb{CP}}^{2}$, and an integer $n \geq  2$. Let $Y\left( {\mathcal{L}, n}\right)$ be the minimal desingularization of the branched cover of ${\mathbb{CP}}^{2}$ along $\mathcal{L}$, defined by the map ${\pi }_{1}\left( {{\mathbb{CP}}^{2} \smallsetminus  \mathcal{L}}\right)  \rightarrow  {H}_{1}\left( {{\mathbb{CP}}^{2} \smallsetminus  \mathcal{L};\mathbb{Z}/n\mathbb{Z}}\right)$.
	
	Questions: Does the combinatorial data determine the homotopy type of $Y\left( {\mathcal{L}, n}\right)$? the topological type? or even the diffeomorphism type?
\end{openproblem}
% Source: Kirby

\plainproblemsection

\begin{openproblem}
	Does every finite acyclic \(2\)-complex admit a PL embedding
	in \(\mathbb{R}^{4}\)?
\end{openproblem}
% Source: Kirby

\subsection*{References}
\begin{enumerate}[label={\arabic*.},leftmargin=2.2em]
	\item O. \c{S}avk, ``A survey of the homology cobordism group,'' arXiv:2209.10735, version 4, 2024.
	\item R. C. Kirby, ``Problems in low-dimensional topology,'' in W. H. Kazez (ed.), \emph{Geometric Topology}, Part 2, AMS/IP Studies in Advanced Mathematics \textbf{2.2}, American Mathematical Society and International Press, 1997, 35--473.
	\item R. \.{I}nan\c{c} Baykur, R. C. Kirby, and D. Ruberman (eds.), \emph{\(K3\): A New Problem List in Low-Dimensional Topology}, Mathematical Surveys and Monographs \textbf{295}, American Mathematical Society, 2026.
\end{enumerate}

% <OPEN-PROBLEMS-IN-TOPOLOGY-1990>
% Selected material from Parts IV, VI, VII, and VIII.
% Source-local problem numbers are intentionally omitted.

\clearpage
\section{Topological Groups and Homogeneous Spaces}

\paragraph{Topological Groups and Homogeneous Spaces.}

\plainproblemsection

\begin{openproblem}
	Is it consistent with ZFC that there exists p \(\in\) \(\omega\)\(^*\) such that every countably compact group is p-compact?
\end{openproblem}


\plainproblemsection

\begin{openproblem}
	Is it a theorem of ZFC that there exist two countably compact groups whose product is not countably compact?
\end{openproblem}


\plainproblemsection

\begin{openproblem}
	For every cardinal \(\alpha\leq 2^{\mathfrak c}\), does there exist a topological group \(G\) such that \(G^\gamma\) is countably compact for every \(\gamma<\alpha\), while \(G^\alpha\) is not countably compact?
\end{openproblem}


\plainproblemsection

\begin{openproblem}
	For \(p, q\in\omega^*\), are the conditions \(p\leq_{C, S}q\) and \(p\leq_{C, G}q\) equivalent? Equivalently, is \(\leq_{C, S}=\leq_{C, G}\)?
\end{openproblem}


\plainproblemsection

\begin{openproblem}
	For every \(q\in\omega^*\), does every \(q\)-compact space embed as a closed subspace of a \(q\)-compact topological group?
\end{openproblem}


\plainproblemsection

\begin{openproblem}
	Let \(M\) be a model of ZFC in which \(\leq_{C, S}=\leq_{C, G}\). Must \(M\) also satisfy the assertion that every \(q\)-compact space embeds as a closed subspace of a \(q\)-compact topological group?
\end{openproblem}


\plainproblemsection

\begin{openproblem}
	For \(p\in\omega^*\), are the following conditions equivalent: \(p\) is minimal under \(\leq_{C, S}\); \(p\) is minimal under \(\leq_{C, G}\); and there exists \(q\in\omega^*\) such that \(p=_{C, S}q\) and \(q\) is Rudin--Keisler minimal?
\end{openproblem}


\plainproblemsection

\begin{openproblem}
	For \(p, q\in\omega^*\), write \(p\leq_{C, G}q\) when every \(q\)-compact topological group is \(p\)-compact. Does Martin's Axiom imply that \(\omega^*\) is not downward directed under \(\leq_{C, G}\)?
\end{openproblem}


\plainproblemsection

\begin{openproblem}
	Is the equality \(\leq_{G-F, S}=\leq_{G-F, G}\) valid?
\end{openproblem}


\plainproblemsection

\begin{openproblem}
	Let \(p\in\omega^*\). (a) Does \(\omega\cup\{p\}\), with the topology inherited from \(\beta\omega\), embed as a closed subspace of a \(p\)-sequential group? (b) Does every \(p\)-sequential space embed as a closed subspace of a \(p\)-sequential group?
\end{openproblem}


\plainproblemsection

\begin{openproblem}
	Let \(M\) be a model of ZFC in which \(\leq_{G-F, S}=\leq_{G-F, G}\). Must \(M\) also satisfy the assertion that every \(p\)-sequential space embeds as a closed subspace of a \(p\)-sequential group?
\end{openproblem}


\plainproblemsection

\begin{openproblem}
	Let \(X\) be a compact homogeneous space of countable tightness. Must there exist \(p\in\omega^*\) such that \(X\) is \(p\)-sequential?
\end{openproblem}


\plainproblemsection

\begin{openproblem}
	Can every space \(X\) be embedded as a closed subspace of a topological group \(G\) such that \(t(G)=t(X)\)?
\end{openproblem}


\plainproblemsection

\begin{openproblem}
	Does there exist a topological group \(G\) such that \(\omega=t(G)<t(G\times G)\)?
\end{openproblem}


\plainproblemsection

\begin{openproblem}
	Is it consistent with the axioms of ZFC that for p \(\in\) \(\omega\) the class of p-sequential groups is closed under finite products? Countable products?
\end{openproblem}


\plainproblemsection

\begin{openproblem}
	(Arkhangel'skii ) (a) Is there a Lindelöf group G such that G \(\times\) G is not Lindelöf? (b) Can every Lindelöf space be embedded as a closed subspace of a Lindelöf group? What about the Sorgenfrey line?
\end{openproblem}


\plainproblemsection

\begin{openproblem}
	(a) Do there exist a Lindelöf group \(G\) and a space \(X\) such that \(X\times X\) is not Lindelöf and some closed subspace of \(G\) maps continuously onto \(X\)? (b) Is the Sorgenfrey line a continuous image of a Lindelöf group?
\end{openproblem}


\plainproblemsection

\begin{openproblem}
	Is there a pseudocompact, homogeneous space X such that X \(\times\) X is not pseudocompact?
\end{openproblem}


\plainproblemsection

\begin{openproblem}
	Is there a countably compact, homogeneous space X such that X \(\times\) X is (a) not countably compact? (b) not pseudocompact?
\end{openproblem}


\plainproblemsection

\begin{openproblem}
	(van Douwen) For what cardinal numbers \(\alpha\) is there a compact, homogeneous space with cellularity greater than \(\alpha\)? What about the case \(\alpha\) = c?
\end{openproblem}


\plainproblemsection

\begin{openproblem}
	Does every countably compact space embed as a retract into a countably compact homogeneous space?
\end{openproblem}


\plainproblemsection

\begin{openproblem}
	For what non-empty spaces X does there exist a free (CCAG, CCAG)-group over X? For some X? For all X?
\end{openproblem}


\plainproblemsection

\begin{openproblem}
	Does every countably compact space X admit a free (CCAG, CCAG)-group over X in which X is closed?
\end{openproblem}


\plainproblemsection

\begin{openproblem}
	Let \(\alpha\) be an infinite cardinal. Do there exist totally bounded abelian torsion groups \(G_0\) and \(G_1\), neither having a proper dense subgroup, such that \[w(G_0)=|G_0|=\alpha,\qquad w(G_1)=|G_1|=2^{<\alpha}?\]What is the answer when \(\alpha\) is a strong limit cardinal?
\end{openproblem}


\plainproblemsection

\begin{openproblem}
	Does every pseudocompact group G of uncountable weight have a proper dense pseudocompact subgroup? What if G is Abelian? Connected and Abelian?
\end{openproblem}


\plainproblemsection

\begin{openproblem}
	Let \(G\) be a dense pseudocompact subgroup of \(\mathbb T^{\mathfrak c^+}\). Must \(G\) have a proper dense subgroup?
\end{openproblem}


\plainproblemsection

\begin{openproblem}
	Let \(K\) be an infinite compact group of weight \(\alpha\), and let \(P(X)\) denote the \(P\)-space modification of \(X\). Is the identity \[d\!\left(P(\{0,1\}^{\alpha})\right)=d(P(K))=(\log\alpha)^\omega\]a theorem of ZFC?
\end{openproblem}


\plainproblemsection

\begin{openproblem}
	Let \(G\) be a dense pseudocompact subgroup of a compact group \(K\). Must \(K\) contain a dense countably compact subgroup \(C\) such that \(|C|\leq |G|\)? Can \(C\) be chosen to contain \(G\)?
\end{openproblem}


\plainproblemsection

\begin{openproblem}
	Classify the self-dual locally compact Abelian groups.
\end{openproblem}


\plainproblemsection

\begin{openproblem}
	Must a compact group in which each element has finite order have the property that the orders of its elements are bounded?
\end{openproblem}


\plainproblemsection

\begin{openproblem}
	(M. G. Tkachenko (ZFC)) Can the free Abelian group on c-many generators be given a countably compact topological group topology?
\end{openproblem}


\plainproblemsection

\begin{openproblem}
	For which cardinals \(\alpha\) does there exist a topological group \(G(\alpha)\) of weight \(\alpha\) such that every topological group of weight \(\alpha\) is topologically isomorphic to a subgroup of \(G(\alpha)\)?
\end{openproblem}


\plainproblemsection

\begin{openproblem}
	For which cardinals \(\alpha\) does there exist an abelian topological group \(G(\alpha)\) of weight \(\alpha\) such that every abelian topological group of weight \(\alpha\) is topologically isomorphic to a subgroup of \(G(\alpha)\)? Is \(\alpha=\omega\) such a cardinal?
\end{openproblem}


\plainproblemsection

\begin{openproblem}
	Does there exist an abelian group \(G\) admitting no connected topological-group topology such that every unconditionally closed subgroup \(H\leq G\) satisfies \(|G/H|\geq 2^\omega\)?
\end{openproblem}


\plainproblemsection

\begin{openproblem}
	Let \(\alpha\) be an infinite cardinal, and let \(X\) be a compact space with \(w(X)\leq 2^\alpha\) that is a continuous image of a \(\sigma\)-compact topological group. Must \(d(X)\leq\alpha\)? What if \(w(X)\leq\alpha^+\)?
\end{openproblem}


\plainproblemsection

\begin{openproblem}
	For each group G, let '(G) denote the set of (Hausdorff) topological group topologies, partially ordered by inclusion. Is it true for each \(\alpha\) that the partially ordered sets '(F(\(\alpha\))) and '(FA(\(\alpha\))) are non-isomorphic?
\end{openproblem}


\plainproblemsection

\begin{openproblem}
	Does every Abelian group which is not of bounded order admit a minimally almost periodic topological group topology? What about the countable case?
\end{openproblem}


\plainproblemsection

\begin{openproblem}
	For which compact metrizable spaces \(X\) does the homeomorphism group \(\operatorname{Homeo}(X)\) admit a unique complete separable metrizable topological-group topology?
\end{openproblem}


\plainproblemsection

\begin{openproblem}
	Can a counterexample in ZFC be constructed to the assertion that every countably compact topological semigroup with two-sided cancellation is a topological group?
\end{openproblem}


\plainproblemsection

\begin{openproblem}
	Is there a minimal topological field whose topology is not locally bounded? In particular, can such a topology be placed on \(\mathbb Q\)?
\end{openproblem}


\plainproblemsection

\begin{openproblem}
	(Kowalsky ) Is every minimal Hausdorff field topology on a (commutative) field induced by an absolute value or by a non-Archimedean valuation?
\end{openproblem}


\subsection*{References}
\begin{enumerate}[label={\arabic*.},leftmargin=2.2em]
	\item W. W. Comfort, ``Problems on topological groups and other homogeneous spaces,'' in J. van Mill and G. M. Reed (eds.), \emph{Open Problems in Topology}, North-Holland, 1990.
	\item E. Pearl, ``Open problems in topology,'' \emph{Topology and its Applications} \textbf{136} (2004), 37--85, \url{https://doi.org/10.1016/S0166-8641(03)00183-4}.
	\item M. Megrelishvili, ``Every topological group is a group retract of a minimal group,'' \emph{Topology and its Applications} \textbf{155} (2008), 2105--2127, \url{https://doi.org/10.1016/j.topol.2007.04.028}.
\end{enumerate}

\clearpage
\section{Algebraic and Geometric Topology}

\paragraph{Classification of Incomplete Metric Spaces.}

\plainproblemsection

\begin{openproblem}
	Determine the countable ordinals \(\alpha\) for which \(\mathcal R_\alpha\)-absorbing sets exist.
\end{openproblem}


\plainproblemsection

\begin{openproblem}
	Let \(X\) be a separable absolute retract that is a countable union of strong \(Z\)-sets. Under what conditions is \(X\) an \(\mathcal F_0(X)\)-absorbing set?
\end{openproblem}


\plainproblemsection

\begin{openproblem}
	Let \(X\subset\mathbb R^k\) be a finite-dimensional \(\mathcal C\)-absorbing set. Is every \(\mathcal C\)-absorbing set \(Y\) homeomorphic to \(X\)?
\end{openproblem}


\plainproblemsection

\begin{openproblem}
	Let \(X\) be a \(\mathcal C\)-absorbing set in \(\mathbb R^\infty\), and let \(h: A\to B\) be a homeomorphism between \(Z\)-sets in \(X\). Can the homotopy hypothesis in the \(Z\)-set unknotting theorem be weakened so that the resulting extension theorem applies to the classification of finite-dimensional \(\mathcal C\)-absorbing sets?
\end{openproblem}


\plainproblemsection

\begin{openproblem}
	Let \(X\) and \(Y\) be \(\mathcal C\)-absorbing sets in an absolute retract \(M\). Under what conditions does an ambient homeomorphism of \(M\), arbitrarily close to the identity, send \(X\) onto \(Y\)? In particular, what happens for \(M=\mathbb R^\infty\) or \(M=I^\infty\)?
\end{openproblem}


\plainproblemsection

\begin{openproblem}
	When is a separable incomplete \(\lambda\)-convex set \(Y\) an \(\mathcal F_0(Y)\)-absorbing set?
\end{openproblem}


\plainproblemsection

\begin{openproblem}
	Let Y be an absolute retract \(\lambda\)-convex set. Is every Z-set in Y astrong Z-set?
\end{openproblem}


\plainproblemsection

\begin{openproblem}
	Let \(W\in(\mathcal R_\omega)_\sigma\) be an infinite-dimensional convex subset of a complete metric linear space. Is \(W\) homeomorphic to \(\sigma\)?
\end{openproblem}


\plainproblemsection

\begin{openproblem}
	Let \(E\) be an \(\mathcal R\)-universal \(\sigma\)-compact metric linear space that is an absolute retract. Is \(E\) homeomorphic to \(\Sigma\)?
\end{openproblem}


\plainproblemsection

\begin{openproblem}
	Let \(E\) be a \(\sigma\)-compact metric linear space that is an absolute retract. Is \(E\) a \(\mathcal C(E)\)-absorbing set? Is it sufficient to verify the strong \(\mathcal C(E)\)-universality of \(E\)?
\end{openproblem}


\plainproblemsection

\begin{openproblem}
	Let \(W=\operatorname{conv}\{x_i\}_{i=1}^{\infty}\) be the infinite-dimensional convex hull of countably many vectors in a nonlocally convex metric linear space. Is every finite hull \(\operatorname{conv}\{x_i\}_{i=1}^{n}\) a strong \(Z\)-set in \(W\)?
\end{openproblem}


\plainproblemsection

\begin{openproblem}
	Let \(W\) be an \(\mathcal R\)-universal, \(\sigma\)-compact, convex absolute retract in a complete metric linear space, and suppose every compact subset of \(W\) is a \(Z\)-set. Is \(W\) homeomorphic to \(\Sigma\)?
\end{openproblem}


\plainproblemsection

\begin{openproblem}
	Let \(G\in(\mathcal R_\omega)_\sigma\) be an infinite-dimensional contractible metric group. Is \(G\) homeomorphic to \(\sigma\)? Equivalently, is \(G\) strongly \(\mathcal R_\omega\)-universal?
\end{openproblem}


\plainproblemsection

\begin{openproblem}
	Let \(G\) be an infinite-dimensional contractible metric group belonging to \((\mathcal R_\omega)_\sigma\). Must \(G\) be \(\mathcal R_\omega\)-universal? In particular, must \(G\) contain a topological disc?
\end{openproblem}


\plainproblemsection

\begin{openproblem}
	Let \(G\) be the additive group generated by a linearly independent arc in \(\ell^2\), and suppose \(G\) is contractible. Is \(G\) homeomorphic to \(\sigma\)?
\end{openproblem}


\plainproblemsection

\begin{openproblem}
	Let \(H\) be a nonlocally compact separable complete metric group that is an absolute retract. Does \(H\) contain a subgroup that is an \(\mathcal R_\omega\)-absorbing set in \(H\)?
\end{openproblem}


\plainproblemsection

\begin{openproblem}
	Let \(G\) be an \(\mathcal R\)-universal, \(\sigma\)-compact absolute-retract metric group. Is \(G\) homeomorphic to \(\Sigma\)?
\end{openproblem}


\plainproblemsection

\begin{openproblem}
	Let \(G\) be the additive group generated by a linearly independent copy of the Hilbert cube in \(\ell^2\), and suppose \(G\) is contractible. Is \(G\) homeomorphic to \(\Sigma\)?
\end{openproblem}


\plainproblemsection

\begin{openproblem}
	Let \(\operatorname{LIP}_{\partial}(I^n)\) be the group of Lipschitz homeomorphisms of \(I^n\) that fix its boundary. Is \(\operatorname{LIP}_{\partial}(I^n)\) homeomorphic to \(\Sigma\)?
\end{openproblem}


\plainproblemsection

\begin{openproblem}
	Let \(E\) be a \(\sigma\)-compact pre-Hilbert space, and let \(\mathcal C(E)\) be its associated class of compacta. Is \[\{C\times I: C\in\mathcal C(E)\}\subseteq\mathcal C(E)?\]Likewise, is \[\{C\times D: C, D\in\mathcal C(E)\}\subseteq\mathcal C(E)?\]
\end{openproblem}


\plainproblemsection

\begin{openproblem}
	Does every homeomorphism between finite-dimensional compacta of a nonlocally convex \(\sigma\)-compact metric linear space E extend to a homeomorphism of the whole space E?
\end{openproblem}


\plainproblemsection

\begin{openproblem}
	Classify topologically the spaces cF for Borelian filters F.
\end{openproblem}


\plainproblemsection

\begin{openproblem}
	Let \(\mathcal F\) be a filter on \(\mathbb N\), and define \[c_{\mathcal F}=\bigl\{(x_i)\in\mathbb R^{\mathbb N}:\text{for every }\varepsilon>0\text{ there is }S\in\mathcal F\text{ such that }|x_i|<\varepsilon\text{ for all }i\in S\bigr\}.\]Let \(\mathcal F_0(c_{\mathcal F})\) be the class of spaces homeomorphic to closed subspaces of \(c_{\mathcal F}\). Is \(c_{\mathcal F}\) strongly \(\mathcal F_0(c_{\mathcal F})\)-universal? Is \(c_{\mathcal F}\) homeomorphic to \(c_{\mathcal F}^{\infty}\)?
\end{openproblem}


\plainproblemsection

\begin{openproblem}
	Characterize the spaces \(\sigma_n\) and \(\sigma_{nm}\) topologically. In particular, is every \(\mathcal R_n\)-absorbing set homeomorphic to \(\sigma_n\), and is every \(\mathcal R_{nm}\)-absorbing set homeomorphic to \(\sigma_{nm}\)?
\end{openproblem}


\plainproblemsection

\begin{openproblem}
	Let \(X\) be a separable \(n\)-dimensional completely metrizable space in \(AE(n)\), and suppose every map into \(X\) from a separable \(n\)-dimensional completely metrizable space can be approximated arbitrarily closely by \(Z\)-embeddings. Is \(\sigma_n\) representable in \(X\)?
\end{openproblem}


\plainproblemsection

\begin{openproblem}
	Are \(\sigma_n\) and \(\sigma_{nm}\) representable in the \(n\)-dimensional Menger cube \(M_n^k\), where \(k\geq 2n+1\geq m\geq n\)?
\end{openproblem}


\plainproblemsection

\begin{openproblem}
	Let \(M_n^k\) be the Menger cube, with \(k\geq2n+1\). Can \(M_n^k\) be written as \(X\cup Y\), where \(X\cong\sigma_n^k\), \(Y\cong N_n^k\), and both \(X\) and \(Y\) are locally \(n\)-homotopy negligible in \(M_n^k\)?
\end{openproblem}


\plainproblemsection

\begin{openproblem}
	For \(k\geq2n+1\), let \(f_n: M_n^k\to I^\infty\) be the \((n-1)\)-soft map in Dranishnikov's construction. Are \[f_n^{-1}((-1,1)^\infty)\quad\text{and}\quad f_n^{-1}(I^\infty\setminus(-1,1)^\infty)\]homeomorphic to \(N_n^k\) and \(\sigma_n^k\), respectively?
\end{openproblem}


\plainproblemsection

\begin{openproblem}
	Let \(X\) be one of \(\sigma_n^k,\sigma_{nm}^k, N_n^k\), where \(k\geq2n+1\geq m\geq n\), and let \(f: X\to Y\) be a \(UV^{n-1}\)-map onto \(Y\in AE(n)\). Prove that for every open cover \(\mathcal U\) of \(Y\) there is an open cover \(\mathcal B\) of \(Y\) such that every homeomorphism \(h: A\to B\) between \(Z\)-sets in \(X\), with \(f\circ h\) \(\mathcal B\)-close to \(f|_A\), extends to a homeomorphism \(H: X\to X\) satisfying \(H|_A=h\) and with \(f\circ H\) \(\mathcal U\)-close to \(f\).
\end{openproblem}


\paragraph{Finite-Dimensional Manifolds and Generalized Manifolds.}

\plainproblemsection

\begin{openproblem}
	Can a \(p\)-adic group act effectively on a manifold? Equivalently, can a compact manifold \(M\) admit a homeomorphism \(h\) such that every orbit \(\{h^n(x): n\in\mathbb Z\}\) has small diameter and \(\{h^n: n\in\mathbb Z\}\) is relatively compact in the homeomorphism group of \(M\)?
\end{openproblem}


\plainproblemsection

\begin{openproblem}
	Is every PL embedding of \(S^{n-1}\) in \(\mathbb R^n\) PL standard; equivalently, does its image bound a PL \(n\)-ball?
\end{openproblem}


\plainproblemsection

\begin{openproblem}
	Does the topological \(4\)-sphere admit a smooth structure not diffeomorphic to the standard smooth structure on \(S^4\)?
\end{openproblem}


\plainproblemsection

\begin{openproblem}
	(A problem of Hopf) Given a closed, orientable manifold M, is every (absolute) degree one map f: M \(\to\) M a homotopy equivalence?
\end{openproblem}


\plainproblemsection

\begin{openproblem}
	(Whitehead Conjecture; Whitehead ) Every subcomplex of an aspherical 2-complex is itself aspherical.
\end{openproblem}


\plainproblemsection

\begin{openproblem}
	If \(K\) is a subcomplex of a contractible \(2\)-complex, is \(\pi_1(K)\) locally indicable?
\end{openproblem}


\plainproblemsection

\begin{openproblem}
	(Borel Rigidity Conjecture) Every homotopy equivalence N \(\to\) M between closed, aspherical manifolds is homotopic to a homeomorphism.
\end{openproblem}


\plainproblemsection

\begin{openproblem}
	(Zeeman Conjecture; Zeeman ) If X is a contractible finite 2- complex, then X \(\times\) I is collapsible.
\end{openproblem}


\plainproblemsection

\begin{openproblem}
	(Codimension 1 manifold factor problem (generalized Moore problem)) If X \(\times\) Y is a manifold, is X \(\times\) \(\mathbb R^{1}\) a manifold?
\end{openproblem}


\plainproblemsection

\begin{openproblem}
	Is every generalized \(4\)-manifold homotopy equivalent to a topological \(4\)-manifold?
\end{openproblem}


\plainproblemsection

\begin{openproblem}
	(Kervaire Conjecture (also known as the Kervaire-Laudenbach Conjecture)) If A is a group for which the normalizer of some element r in the free product A \(^*\) Z is A \(^*\) Z itself, then A is trivial.
\end{openproblem}


\plainproblemsection

\begin{openproblem}
	Does there exist a homogeneous compact absolute retract of dimension \(n\) for some \(2<n<\infty\)?
\end{openproblem}


\plainproblemsection

\begin{openproblem}
	(Homogeneous ENRs versus generalized manifolds) If X is a homogeneous, locally compact ENR (= finite dimensional ANR, is X a generalized manifold?
\end{openproblem}


\plainproblemsection

\begin{openproblem}
	Does every compact Euclidean neighborhood retract \(X\) contain a point \(x_0\) such that \(H^*(X, X\setminus\{x_0\})\) is finitely generated?
\end{openproblem}


\plainproblemsection

\begin{openproblem}
	Is every homogeneous generalized manifold necessarily a genuine manifold?
\end{openproblem}


\plainproblemsection

\begin{openproblem}
	Do all finite dimensional H-spaces have the homotopy type of a closed manifold?
\end{openproblem}


\plainproblemsection

\begin{openproblem}
	If M is a compact manifold, is the group Homeo(M) of all self-homeomorphisms an ANR?
\end{openproblem}


\plainproblemsection

\begin{openproblem}
	Let \(G\) be a discrete convergence group acting on \(S^2\), and let \(L(G)\) be its limit set. If \(L(G)\) is connected, must it be locally connected?
\end{openproblem}


\plainproblemsection

\begin{openproblem}
	(Freedman and Skora) Must a K(G, 1)-manifold M, where G is finitely generated, have only a finite number of ends? What if M is covered by \(\mathbb R^n\)?
\end{openproblem}


\plainproblemsection

\begin{openproblem}
	Must the Euler characteristic of every closed aspherical \(2n\)-manifold, when nonzero, have sign \((-1)^n\)?
\end{openproblem}


\plainproblemsection

\begin{openproblem}
	Suppose \(W\) is a contractible \(n\)-manifold such that, for every compact \(C\subset W\), there is an essential map \(S^{n-1}\to W\setminus C\). Must \(W\cong\mathbb R^n\)?
\end{openproblem}


\plainproblemsection

\begin{openproblem}
	Let \(X\) be a compact \(n\)-dimensional space with a strongly convex metric without ramifications. Must \(X\) be an \(n\)-cell? If \(X\) is a generalized manifold with boundary, must \(X\setminus\partial X\) be homogeneous?
\end{openproblem}


\plainproblemsection

\begin{openproblem}
	Is there a complex dominated by a 2-complex but not homotopy equivalent to a 2-complex?
\end{openproblem}


\plainproblemsection

\begin{openproblem}
	Is every finitely presented perfect group (perfect = trivial abelianization) the normal closure of a single element?
\end{openproblem}


\plainproblemsection

\begin{openproblem}
	Is every cell-like map between topological \(4\)-manifolds a limit of homeomorphisms?
\end{openproblem}


\plainproblemsection

\begin{openproblem}
	Can a cell-like map defined on \(\mathbb R^n\) have infinite dimensional image if all point-inverses are contractible? absolute retracts? cells? starlike sets? 1-dimensional compacta?
\end{openproblem}


\plainproblemsection

\begin{openproblem}
	Let \(\mathcal G\) be an upper-semicontinuous decomposition of a compact space \(X\) into simple closed curves. Must \(\dim(X/\mathcal G)\leq\dim X\)?
\end{openproblem}


\plainproblemsection

\begin{openproblem}
	Could there be a decomposition G of an n-manifold M into closed connected manifolds (of varying dimensions) with dim(M/G) >n?
\end{openproblem}


\plainproblemsection

\begin{openproblem}
	Can an open map from a compact manifold onto a space, with one-dimensional solenoids as its point-inverses, raise dimension?
\end{openproblem}


\plainproblemsection

\begin{openproblem}
	Does every generalized \(3\)-manifold \(X\) admit a cell-like resolution by a \(3\)-manifold? Does \(X\times\mathbb R\) admit such a resolution?
\end{openproblem}


\plainproblemsection

\begin{openproblem}
	(The Approximation Problem for 3- and 4-manifolds) Which cell-like maps p: M \(\to\) X from a manifold onto a finite dimensional space can be approximated by homeomorphisms? Is it sufficient to know that, given any two disjoint, tame 2-cells B1, and B2 \(\subseteq\) M, there are maps \(\mu\)i: Bi \(\to\) X approximating p|Bi with \(\mu\)1(B1) \(\cap\)\(\mu\)2(B2)= \(\varnothing\)?
\end{openproblem}


\plainproblemsection

\begin{openproblem}
	Is every decomposition of \(\mathbb R^n\) into points and countably many starlike-equivalent sets shrinkable?
\end{openproblem}


\plainproblemsection

\begin{openproblem}
	Let \(f: S^n\to X\) be a map such that every nondegenerate point-inverse \(f^{-1}(f(x))\) is a standardly embedded \(n\)-cell. Can \(f\) be approximated by homeomorphisms? What if there are countably many nondegenerate point-inverses and each is a standardly embedded \((n-2)\)-cell?
\end{openproblem}


\plainproblemsection

\begin{openproblem}
	Let \(\mathcal G\) be an upper-semicontinuous decomposition of \(\mathbb R^n\). Suppose each element has arbitrarily small neighborhoods whose frontiers are \((n-1)\)-spheres disjoint from every nondegenerate element of \(\mathcal G\). Must \(\mathcal G\) be shrinkable? What if the neighborhoods are Euclidean patches?
\end{openproblem}


\plainproblemsection

\begin{openproblem}
	Let \(A\subset\mathbb R^n\) be an \(n\)-dimensional annulus. Does \(A\) admit a parametrization as \(S^{n-1}\times I\) for which the associated decomposition into points and fibre arcs is shrinkable?
\end{openproblem}


\plainproblemsection

\begin{openproblem}
	Is a countable cell-like decomposition \(\mathcal G\) of \(\mathbb R^n\) shrinkable if every nondegenerate \(g\in\mathcal G\) lies in an affine \((n-1)\)-hyperplane?
\end{openproblem}


\plainproblemsection

\begin{openproblem}
	Let \(n\geq4\), and let \(\{C_i\}\) be any sequence of nondegenerate cellular subsets of \(\mathbb R^n\). Does there exist a nonshrinkable cellular decomposition of \(\mathbb R^n\) whose nondegenerate elements form a null sequence \(\{g_i\}\) with \(g_i\cong C_i\) for every \(i\)?
\end{openproblem}


\plainproblemsection

\begin{openproblem}
	Is there a nonshrinkable decomposition of n-space into points and straight line segments? Into convex sets?
\end{openproblem}


\plainproblemsection

\begin{openproblem}
	If \(X\) is a cell-like image of a \(3\)-manifold \(M\), does \(X\) embed in \(M\times\mathbb R\)?
\end{openproblem}


\plainproblemsection

\begin{openproblem}
	If X is a cellular subset of 4-space and G is a cell-like decomposition of X such that dim(X/G) \(\leq\) 1, is the trivial extension of G over 4-space shrinkable? What if X is an arc?
\end{openproblem}


\plainproblemsection

\begin{openproblem}
	Is every simple upper-semicontinuous decomposition of \(\mathbb R^4\) shrinkable?
\end{openproblem}


\plainproblemsection

\begin{openproblem}
	Let \(f: S^4\to S^4\) be one-to-one over the complement of a Cantor set \(K\subset S^4\). Must \(f\) be cell-like? What if \(f\) is one-to-one over the complement of a noncompact zero-dimensional set?
\end{openproblem}


\plainproblemsection

\begin{openproblem}
	Can every cellular map between finite \(4\)-complexes be approximated by homeomorphisms?
\end{openproblem}


\plainproblemsection

\begin{openproblem}
	If \(X\) and \(Y\) are \(n\)-dimensional \(LC^{n-1}\)-compacta that are shape equivalent, must they be cell-like equivalent?
\end{openproblem}


\plainproblemsection

\begin{openproblem}
	(Ferry) If X and Y are shape equivalent LCk compacta, are they UVk-equivalent?
\end{openproblem}


\plainproblemsection

\begin{openproblem}
	Suppose \(X\) and \(Y\) are cell-like equivalent, locally \(1\)-connected compacta. Can their cell-like equivalence be represented by a finite zigzag of cell-like maps in which every intermediate compactum is also locally \(1\)-connected? What is the answer when \(X\) and \(Y\) are homotopy-equivalent but not simple-homotopy-equivalent polyhedra?
\end{openproblem}


\plainproblemsection

\begin{openproblem}
	Let \(X\) be the inverse limit of a sequence of homotopy equivalences \(S^2\leftarrow S^2\leftarrow\cdots\). Is \(X\) cell-like equivalent to \(S^2\)?
\end{openproblem}


\plainproblemsection

\begin{openproblem}
	Let \(K\subset\mathbb R^n\) be a \(k\)-cell. Under what conditions can \(K\) be squeezed to a \((k-1)\)-cell by a map \(f:\mathbb R^n\to\mathbb R^n\) such that \(f|_K\) is conjugate to the projection \(B^k\to B^{k-1}\), while \(f|_{\mathbb R^n\setminus K}\) is a homeomorphism onto \(\mathbb R^n\setminus f(K)\)? What if \(K\) is cellular? What if every Cantor set in \(K\) is tame?
\end{openproblem}


\plainproblemsection

\begin{openproblem}
	Let \(f: M\to X\) be a cell-like map from an \(n\)-manifold onto a finite-dimensional space. Can \(f\) be approximated by a cell-like map \(F: M\to X\) such that every point-inverse \(F^{-1}(x)\) is one-dimensional? In particular, can this be done for \(n\in\{3,4,5\}\)?
\end{openproblem}


\plainproblemsection

\begin{openproblem}
	Does \(\mathbb R^n\) admit a decomposition into \(k\)-cells, with \(k>0\), or into copies of a fixed nontrivial compact absolute retract?
\end{openproblem}


\plainproblemsection

\begin{openproblem}
	Is there a decomposition of Bn into simple closed curves? of a compact contractible space? of a cell-like set?
\end{openproblem}


\plainproblemsection

\begin{openproblem}
	Does there exist a \(1\)- or \(2\)-dimensional compactum with the disjoint arcs property that is not an absolute neighborhood retract?
\end{openproblem}


\plainproblemsection

\begin{openproblem}
	Let \(M\) be an \((n+k)\)-manifold and \(\mathcal G\) an upper-semicontinuous decomposition of \(M\) into closed connected orientable \(n\)-manifolds. Must the decomposition space \(B=M/\mathcal G\) be an ANR? What if all elements of \(\mathcal G\) are pairwise homeomorphic?
\end{openproblem}


\plainproblemsection

\begin{openproblem}
	Let \(M\) be an \((n+k)\)-manifold and \(\mathcal G\) an upper-semicontinuous decomposition of \(M\) into closed connected orientable \(n\)-manifolds. Must the decomposition space \(M/\mathcal G\) be finite-dimensional? What if every element of \(\mathcal G\) is a simple closed curve?
\end{openproblem}


\plainproblemsection

\begin{openproblem}
	For which \(n\) and \(k\) does \(S^{n+k}\) admit an upper-semicontinuous decomposition into \(n\)-spheres, \(n\)-tori, fixed products of spheres, or closed \(n\)-manifolds? Does \(\mathbb R^{n+k}\) ever admit a decomposition into closed \(n\)-manifolds for \(n>0\)?
\end{openproblem}


\plainproblemsection

\begin{openproblem}
	When n and k are both odd, does every closed (n + k)-manifold M admitting a decomposition into closed n-manifolds have Euler characteristic zero?
\end{openproblem}


\plainproblemsection

\begin{openproblem}
	If G is a usc decomposition of an (n + k)-manifold M into n-spheres, where 2 <n +1 <k < 2n +2, is M/G a generalized k-manifold? What if into homology n-spheres?
\end{openproblem}


\plainproblemsection

\begin{openproblem}
	Let \(M^{n+3}\) be decomposed upper-semicontinuously into closed connected orientable \(n\)-manifolds, and write \(B=M/\mathcal G\). Is the set of points at which \(B\) fails to be a generalized \(3\)-manifold locally finite?
\end{openproblem}


\plainproblemsection

\begin{openproblem}
	Let \(M^4\) be decomposed upper-semicontinuously into simple closed curves, with quotient map \(p: M\to B\). If the degeneracy set \(K(B)\) of the local \(1\)-winding functions is empty---equivalently, if the one-dimensional cohomology sheaf of \(p\) is Hausdorff---must \(B\) be a generalized \(3\)-manifold?
\end{openproblem}


\plainproblemsection

\begin{openproblem}
	When n =3 and k =1 does there exist a decomposition G of a connected M containing homotopy inequivalent elements?
\end{openproblem}


\plainproblemsection

\begin{openproblem}
	Does there exist a compact \(5\)-manifold \(W\) with boundary components \(M_0\) and \(M_1\), where \(\pi_1(M_0)\cong1\) and \(\pi_1(M_1)\cong A_5\), such that \(W\) admits a decomposition into closed \(4\)-manifolds containing \(M_0\) and \(M_1\) as decomposition elements?
\end{openproblem}


\plainproblemsection

\begin{openproblem}
	Let \(N\) be a closed \(n\)-manifold. Does some \((n+k)\)-manifold admit an upper-semicontinuous decomposition into copies of \(N\) whose quotient map is not an approximate fibration? Are there examples beyond homology-sphere factors and manifolds that regularly cover themselves cyclically? Is there such an example with \(N\) a surface of negative Euler characteristic?
\end{openproblem}


\plainproblemsection

\begin{openproblem}
	For which closed \(n\)-manifolds \(N\) and integers \(k\) must an upper-semicontinuous decomposition of an \((n+k)\)-manifold into copies of \(N\) have an approximate-fibration quotient map? What if \(\pi_1(N)\) is finite and \(k=2\), if \(N\) is covered by \(S^n\), if \(N\) is hyperbolic, or if all decomposition elements are locally flat?
\end{openproblem}


\plainproblemsection

\begin{openproblem}
	Let \(\Sigma\subset\mathbb R^n\) be a topologically embedded \((n-1)\)-sphere. Suppose that, for every \(\varepsilon>0\), there is a map \(f:\Sigma\to\operatorname{Int}\Sigma\) moving each point by less than \(\varepsilon\), where \(\operatorname{Int}\Sigma\) is the bounded complementary component. Must \(\overline{\operatorname{Int}\Sigma}\) be an \(n\)-cell?
\end{openproblem}


\plainproblemsection

\begin{openproblem}
	An embedded \((n-1)\)-sphere \(\Sigma\subset\mathbb R^n\) is locally spanned in \(\operatorname{Int}\Sigma\) if, for every \(p\in\Sigma\) and \(\varepsilon>0\), there is an \((n-1)\)-cell \(D\subset\Sigma\) with \(p\in\operatorname{Int}D\) such that, for every \(\gamma>0\), there are an \((n-1)\)-cell \(E\subset\operatorname{Int}\Sigma\) of diameter less than \(\varepsilon\) and a homeomorphism \(\partial D\to\partial E\) moving points by less than \(\gamma\). Must a locally spanned \(\Sigma\) bound a topological \(n\)-cell?
\end{openproblem}


\plainproblemsection

\begin{openproblem}
	Let \(\Sigma\subset\mathbb R^n\) be an embedded \((n-1)\)-sphere. Suppose there is \(\delta>0\) such that, for every \(p\in\Sigma\), a round \(n\)-cell \(B_p\) of radius \(\delta\) contains \(p\) and satisfies \(\operatorname{Int}B_p\subset\operatorname{Int}\Sigma\). Must \(\overline{\operatorname{Int}\Sigma}\) be an \(n\)-cell?
\end{openproblem}


\plainproblemsection

\begin{openproblem}
	Let \(C\) be the closure of a complementary component of an embedded \((n-1)\)-sphere in \(\mathbb R^n\). Can every map \(\psi: B^2\to C\) be approximated arbitrarily closely by a map \(\psi': B^2\to C\) such that \(\dim(\psi'(B^2)\cap\operatorname{Fr}C)\leq0\)?
\end{openproblem}


\plainproblemsection

\begin{openproblem}
	Let \(\Sigma\subset\mathbb R^n\) be an embedded \((n-1)\)-sphere such that, for every \(p, q\in\Sigma\), some self-homeomorphism \(h\) of \(\mathbb R^n\) satisfies \(h(\Sigma)=\Sigma\) and \(h(p)=q\). Must \(\Sigma\) be flat? What if \(\Sigma\) is strongly homogeneously embedded, meaning that every self-homeomorphism of \(\Sigma\) extends to one of \(\mathbb R^n\)?
\end{openproblem}


\plainproblemsection

\begin{openproblem}
	Is every strongly homogeneous Cantor set in \(\mathbb R^{3}\) (\(\mathbb R^{4}\)) tame?
\end{openproblem}


\plainproblemsection

\begin{openproblem}
	Is there a wild Cantor set Kin \(\mathbb R^{4}\) defined by contractible objects?
\end{openproblem}


\plainproblemsection

\begin{openproblem}
	If \(X\) is a wild Cantor set in \(S^n\), do arbitrarily small homeomorphisms \(h: S^n\to S^n\) exist for which \(X\cap h(X)=\varnothing\)?
\end{openproblem}


\plainproblemsection

\begin{openproblem}
	Does the Mazur \(4\)-manifold---a compact contractible \(4\)-manifold with nonsimply connected boundary---have two disjoint spines?
\end{openproblem}


\plainproblemsection

\begin{openproblem}
	Are objects X \(\subseteq\) \(\mathbb R^n\) that can be instantaneously pushed off themselves geometrically tame?
\end{openproblem}


\plainproblemsection

\begin{openproblem}
	Let \(P\subset\mathbb R^n\) be compact and let \(U\supset P\) be open. Must \(U\) contain a Cantor set \(C\) such that every loop in \(\mathbb R^n\setminus U\) that is contractible in \(\mathbb R^n\setminus C\) is also contractible in \(\mathbb R^n\setminus P\)?
\end{openproblem}


\plainproblemsection

\begin{openproblem}
	Let \(\lambda: X\hookrightarrow M\) be a closed embedding of a generalized \(3\)-manifold in a \(4\)-manifold. Can \(\lambda\) be approximated by \(1\)-LCC embeddings? What is the answer for generalized manifolds with boundary?
\end{openproblem}


\plainproblemsection

\begin{openproblem}
	Which homology \(n\)-spheres \(K\) bound acyclic \((n+1)\)-manifolds \(N\) for which inclusion induces an isomorphism \(\pi_1(K)\to\pi_1(N)\)?
\end{openproblem}


\plainproblemsection

\begin{openproblem}
	Let \(X\subset\mathbb R^n\) be cell-like. Does there exist an arc \(\alpha\subset\mathbb R^n\) such that \(\mathbb R^n\setminus\alpha\cong\mathbb R^n\setminus X\)?
\end{openproblem}


\plainproblemsection

\begin{openproblem}
	Can there be a codimension 3 cell Din \(\mathbb R^n\) (n>5) such that all 2-cells in Dare wildly embedded in \(\mathbb R^n\) but each arc (each Cantor set) there is tame?
\end{openproblem}


\plainproblemsection

\begin{openproblem}
	Can every \(n\)-dimensional compact absolute retract be embedded in \(\mathbb R^{2n}\)?
\end{openproblem}


\plainproblemsection

\begin{openproblem}
	Can every \(S^n\)-like continuum be embedded in \(\mathbb R^{2n}\)?
\end{openproblem}


\plainproblemsection

\begin{openproblem}
	Does \(S^4\) contain a possibly wild \(2\)-sphere \(\Sigma\) such that \(S^4\setminus\Sigma\) is topologically, but not smoothly, \(S^1\times\mathbb R^3\)?
\end{openproblem}


\paragraph{Shape Theory.}

\plainproblemsection

\begin{openproblem}
	Let \(f: X\to Y\) be a closed map of metrizable spaces with \(\operatorname{cdim}_{\mathbb Z}X\leq n\). Suppose there are \(m, k\in\mathbb N\) such that, for every \(y\in Y\), the total cohomology of \(f^{-1}(y)\) has rank at most \(m\), and its torsion subgroup has cardinality at most \(k\). Must \(Y\) have finite cohomological dimension?
\end{openproblem}


\plainproblemsection

\begin{openproblem}
	Does there exist an infinite-dimensional compact ANR \(X\) and a nontrivial coefficient group \(G\) such that \(\operatorname{cdim}_G X<\infty\)?
\end{openproblem}


\plainproblemsection

\begin{openproblem}
	Let \(PS^n\) be the Postnikov-system space used for the weak homotopy model of \(S^n\). Suppose every map \(f: A\to PS^n\), with \(A\) closed in a compactum \(X\), extends over \(X\). Must \(\operatorname{cdim}_{\mathbb Z}X\leq n\)?
\end{openproblem}


\plainproblemsection

\begin{openproblem}
	Let \(f: M^n\to X\) map a closed \(n\)-manifold onto a finite-dimensional ANR, and suppose all fibres of \(f\) have the same shape. Must the Čech cohomology of the fibres be finitely generated?
\end{openproblem}


\plainproblemsection

\begin{openproblem}
	For each Euclidean space \(\mathbb R^n\), does there exist a finite polyhedron \(P_n\subset\mathbb R^n\) such that a compactum \(X\subset\mathbb R^n\) has trivial shape if and only if every map \(X\to P_n\) is null-homotopic?
\end{openproblem}


\plainproblemsection

\begin{openproblem}
	For each \(n\in\mathbb N\), does there exist a finite polyhedron \(P_n\) such that a compactum \(X\) of shape dimension at most \(n\) has trivial shape if and only if every map \(X\to P_n\) is null-homotopic?
\end{openproblem}


\plainproblemsection

\begin{openproblem}
	If \(X\) is movable and has shape dimension \(\operatorname{Sd}X=2\), must \(\operatorname{Sd}(X\times S^1)=3\)?
\end{openproblem}


\plainproblemsection

\begin{openproblem}
	(J. Krasinkiewicz, D. R. McMillan ) Is a movable continuum pointed movable?
\end{openproblem}


\plainproblemsection

\begin{openproblem}
	If the wedge X \(\vee\) Y of two continua is movable, is X pointed movable?
\end{openproblem}


\plainproblemsection

\begin{openproblem}
	(Dydak and Segal ) Is X \(\in\) AP pointed 1-movable?
\end{openproblem}


\plainproblemsection

\begin{openproblem}
	(Dydak and Segal ) Is X \(\in\) AP regularly movable?
\end{openproblem}


\plainproblemsection

\begin{openproblem}
	Is every movable subcontinuum of \(\mathbb R^3\) regularly movable?
\end{openproblem}


\plainproblemsection

\begin{openproblem}
	If \(X\) is a movable continuum and \(\operatorname{pro}\text{-}\pi_1(X)\) is trivial, must \(X\) be regularly movable?
\end{openproblem}


\plainproblemsection

\begin{openproblem}
	(J. Krasinkiewicz ) Does every non-movable continuum X contain a non-movable curve?
\end{openproblem}


\plainproblemsection

\begin{openproblem}
	(D. R. McMillan ) Suppose A \(\times\) B is embeddable in a 3-manifold and each continuum A and B is non-degenerate. Is A movable?
\end{openproblem}


\plainproblemsection

\begin{openproblem}
	If Sd X =2 and X is an ANR, is there a 2-dimensional polyhedron of the same shape as X?
\end{openproblem}


\plainproblemsection

\begin{openproblem}
	Let \(X\subset\mathbb R^4\) be a compactum shape equivalent to a polyhedron. Must \(X\) be shape equivalent to a finite polyhedron?
\end{openproblem}


\plainproblemsection

\begin{openproblem}
	(K. Borsuk ) If X is movable and Sd X \(\leq\) n, isthere a compactum Y \(\subset\) E2n of the same shape as X?
\end{openproblem}


\plainproblemsection

\begin{openproblem}
	Let \(X\) be an ANR-divisor with trivial \(\operatorname{pro}\text{-}\pi_1(X)\). Is there a polyhedron having the same shape as \(X\)?
\end{openproblem}


\plainproblemsection

\begin{openproblem}
	If X and Y are ANR-divisors is X \(\times\) Y an ANR-divisor?
\end{openproblem}


\plainproblemsection

\begin{openproblem}
	If X is shape dominated by an ANR-divisor is it an ANRdivisor?
\end{openproblem}


\plainproblemsection

\begin{openproblem}
	Let \(f: M^n\to X\) be a proper surjection whose fibres are closed \(k\)-manifolds, and suppose \(\dim X<\infty\). Must \(X\) be an ANR?
\end{openproblem}


\plainproblemsection

\begin{openproblem}
	Let \(f: X\to Y\) be a map of compacta. Suppose that, for every \(n>0\), there is \(f_n: Y\to X\) such that \(\operatorname{dist}(f\circ f_n,\operatorname{id}_Y)<1/n\). If \(X\) is an ANR, must \(Y\) be an ANR?
\end{openproblem}


\plainproblemsection

\begin{openproblem}
	Let X be a homogeneous ANR of finite dimension. Is it a generalized manifold?
\end{openproblem}


\plainproblemsection

\begin{openproblem}
	Is there a homogeneous contractible ANR of finite dimension?
\end{openproblem}


\plainproblemsection

\begin{openproblem}
	Let \(X\) be a finite-dimensional ANR and let \(k\geq1\). Must there exist \(x\in X\) such that \(H_k(X, X\setminus\{x\};\mathbb Z)\) is finitely generated?
\end{openproblem}


\plainproblemsection

\begin{openproblem}
	Let \(X\) be an ANR such that, for some \(n\), \[H_k(X, X\setminus\{x\})=0\quad(k\neq n),\qquad H_n(X, X\setminus\{x\})\cong\mathbb Z\]for every \(x\in X\). Must \(X\) be finite-dimensional?
\end{openproblem}


\plainproblemsection

\begin{openproblem}
	Is every shape equivalence a strong shape equivalence?
\end{openproblem}


\plainproblemsection

\begin{openproblem}
	(J. Dydak and J. Segal ) Suppose A is a closed subset of a compact space X and f and g are maps from X to K coinciding on A, where K is a CW complex. If the inclusion from A to X is a shape equivalence and f|A = g|A, is f homotopic to g rel. A?
\end{openproblem}


\plainproblemsection

\begin{openproblem}
	(J. Dydak and J. Segal ) Suppose A is a closed subset of a compact space X and f: A \(\to\) Y is a shape equivalence. Is the natural projection p: X \(\to\) Y \(\cup\)f X a shape equivalence?
\end{openproblem}


\plainproblemsection

\begin{openproblem}
	Suppose a map f: X \(\to\) Y is a shape equivalence and x is a point in X. Is the map f from (X, x) to (Y, f(x)) a pointed shape equivalence?
\end{openproblem}


\plainproblemsection

\begin{openproblem}
	Let the shape dimension of a compactum X be finite. Is there a shape equivalence f: X \(\to\) Y (f is a map) where the dimension of Y is finite?
\end{openproblem}


\plainproblemsection

\begin{openproblem}
	Suppose f: M \(\to\) X is a shape domination (i. e. fg = idX for some shape morphism g: X \(\to\) M)and Sd X \(\geq\) n. Is f a shape isomorphism?
\end{openproblem}


\plainproblemsection

\begin{openproblem}
	Given a space X, is there a strong shape equivalence f from X to Y such that for any space Z the natural function from [Z, Y ] (= the set of homotopy classes of maps from Z to Y )to SSh(Z, Y ) (= the set of strong shape morphisms from Z to Y ) is a bijection?
\end{openproblem}


\plainproblemsection

\begin{openproblem}
	Suppose i: X \(\to\) Y is the inclusion of continua such that both X and Y/X are of polyhedral shape. Is i a shape equivalence if it induces isomorphisms of all homotopy pro-groups?
\end{openproblem}


\plainproblemsection

\begin{openproblem}
	Let \(f: X\to Y\) be a cell-like map, and let \(\{A_n\}_{n=1}^{\infty}\) be compact subsets of \(Y\). Assume that \(f|_{f^{-1}(A)}: f^{-1}(A)\to A\) is a shape equivalence whenever \(A\) is contained in some \(A_n\), and that \(f^{-1}(y)\) is a singleton for every \(y\in Y\setminus\bigcup_{n=1}^{\infty}A_n\). Must \(f\) be a shape equivalence?
\end{openproblem}


\plainproblemsection

\begin{openproblem}
	Is the cardinality of the set of shape classes of continua which are like a closed manifold M necessarily uncountable?
\end{openproblem}


\plainproblemsection

\begin{openproblem}
	For a closed manifold M must the set [M, M] of homotopy classes of maps of M into itself be infinite?
\end{openproblem}


\plainproblemsection

\begin{openproblem}
	If (X, x0) is a movable P-like compactum with finitely generated first shape group \(\pi_1\) (X, x0), thenis X an FANR?
\end{openproblem}


\plainproblemsection

\begin{openproblem}
	If X is P-like and the first shape group \(\pi_1\) (X, x0) is finitely generated is the nth shape group \(\pi\)n (X, x0) countable for n=1,2,...?
\end{openproblem}


\paragraph{Algebraic Topology and Transformation Groups.}

\plainproblemsection

\begin{openproblem}
	Let \(F\to E\to B\) be a fibration of connected spaces, with \(P\pi_1(E)=1\) and \(F\) of finite type. If the sequence \[H_*(F;\mathbb Z)\longrightarrow E^+\longrightarrow B^+\]is an isomorphism in the sense of the plus-construction problem, must \(\pi_1(B)=1\)?
\end{openproblem}


\plainproblemsection

\begin{openproblem}
	Let \([2]: S^n\to S^n\) have degree \(2\). When is \(\Omega^q[2]:\Omega^qS^n\to\Omega^qS^n\) homotopic to the \(H\)-space squaring map? In particular, what happens for \(q=2\) and \(n=2^k-1\), and for \(q=3, n=5\)?
\end{openproblem}


\plainproblemsection

\begin{openproblem}
	For which \(n, r\) is the fibre \(W_{n, r}\) of the iterated suspension \(S^n\to\Omega^rS^{n+r}\), localized at \(p=2\), an \(H\)-space? At an odd prime, with \(n\) even and \(S^n\) replaced by \(S^{pn-1}\) in the James construction, is the corresponding fibre a double loop space?
\end{openproblem}


\plainproblemsection

\begin{openproblem}
	Let \(p>2\), and suppose \(\Phi:\Omega^2S^{2np+1}\to S^{2np-1}\) restricts to a degree-\(p\) map on the bottom cell and the composite \[\Omega^2S^{2n+1}\xrightarrow{\Omega H_p}\Omega^2S^{2np+1}\xrightarrow{\Phi}S^{2np-1}\] is null-homotopic. Is the composite \[\Omega^3S^{2np+1}\xrightarrow{\Omega\Phi}\Omega S^{2np-1}\xrightarrow{\Sigma^3}\Omega^3S^{2np+1}\] the triple loop of the degree-\(p\) map?
\end{openproblem}


\plainproblemsection

\begin{openproblem}
	Let \(\pi\) be a finite group and \(G\) a compact Lie group. Is \([B\pi, BG]\) finite? If \(\pi_p\) ranges over the Sylow \(p\)-subgroups of \(\pi\), is the restriction map \[[B\pi, BG]\longrightarrow\prod_p[B\pi_p, BG]\]injective?
\end{openproblem}


\plainproblemsection

\begin{openproblem}
	Find conditions on \(X\) and \(Y\) under which \(X_{K(n)}\simeq Y_{K(n)}\) for every \(n\geq1\) implies \(X\simeq Y\), where \(K(n)\) is Morava \(K\)-theory.
\end{openproblem}


\plainproblemsection

\begin{openproblem}
	Is every suspension spectrum harmonic; that is, local with respect to \(\bigvee_{n\geq0}K(n)\)?
\end{openproblem}


\plainproblemsection

\begin{openproblem}
	Describe the equivariant cobordism ring \(MU_G^*(\mathrm{pt})\), and formulate an equivariant analogue of Landweber's exact functor theorem.
\end{openproblem}


\plainproblemsection

\begin{openproblem}
	Construct equivariant versions of \(K(n)^*\) and \(E(n)^*\), and prove a completion theorem for them.
\end{openproblem}


\plainproblemsection

\begin{openproblem}
	Find good infinite-loop-space models representing the complex-oriented cohomology theories \(K(n)^*\) and \(E(n)^*\).
\end{openproblem}


\plainproblemsection

\begin{openproblem}
	In the homotopy-fibre construction in which \(W(n)\) is the difference between \(S^{2n-1}\) and \(\Omega^2S^{2n+1}\), and \(X(n)\) is the difference between \(W(n)\) and \(\Omega^4W(n+1)\), do the successive differences between the \(X(n)\) approximate the subalgebra \(A_2\) of the Steenrod algebra? Do further successive differences approximate \(A_3, A_4,\ldots\)?
\end{openproblem}


\plainproblemsection

\begin{openproblem}
	Let \(C_n\) be the fibre of \(\Omega^{2n+1}S^{2n+1}\to\Omega^\infty P^{2n}\wedge J\). Find a stable space \(B_n\) and a \(v_2\)-equivalence \(C_n\to QB_n\).
\end{openproblem}


\plainproblemsection

\begin{openproblem}
	Prove or disprove the Freyd generating hypothesis in the stable homotopy category.
\end{openproblem}


\plainproblemsection

\begin{openproblem}
	Give an algebraic analysis of the rationalized stable category of \(G\)-spaces.
\end{openproblem}


\plainproblemsection

\begin{openproblem}
	Let \(G\) be a compact Lie group, let \(H\subseteq J\subseteq K\subseteq G\) be closed subgroups, and let \(p\) be prime. For a closed subgroup \(L\subseteq G\), let \(q(L, p)\) be the prime ideal of the Burnside ring \(A(G)\) determined by reduction modulo \(p\) of the \(L\)-mark. If \(q(H, p)=q(K, p)\), must \(q(J, p)=q(H, p)\)?
\end{openproblem}


\plainproblemsection

\begin{openproblem}
	Let \(B_G\Pi\) be the classifying \(G\)-space for principal \((G,\Pi)\)-bundles, and let \(h_G^*\) be a \(G\)-equivariant cohomology theory. Calculate \(h_G^*(B_G\Pi)\) in significant cases.
\end{openproblem}


\plainproblemsection

\begin{openproblem}
	For the canonical \(G\)-map \(\alpha: B_G\Pi\to\operatorname{Map}(EG, B\Pi)\) associated with the Borel construction on \((G,\Pi)\)-bundles, determine its homotopical properties when \(G\) is an arbitrary finite group.
\end{openproblem}


\plainproblemsection

\begin{openproblem}
	Let \(\pi\) be finite and let \(G(X)\) be the monoid of self-homotopy equivalences of a finite complex \(X\). Determine the homotopy-theoretic structure of \(\operatorname{Map}_0(B\pi, BG(X))\).
\end{openproblem}


\plainproblemsection

\begin{openproblem}
	Let \(\pi\) be a finite group, let \(X\) be a finite complex, and let \(G(X)\) be the monoid of self-homotopy equivalences of \(X\). Determine the homotopy-theoretic structure of \(\operatorname{Map}\bigl(B\pi, EG(X)\times_{G(X)}X\bigr)\).
\end{openproblem}


\plainproblemsection

\begin{openproblem}
	Let \(A\) be the ring of integers in a finite extension of \(\mathbb Q_p\), and let \(V_A\) be the coefficient ring of the universal \(p\)-typical formal \(A\)-module. Is there a spectrum \(S_A\) such that \(BP_*(S_A)=V_A\)?
\end{openproblem}


\plainproblemsection

\begin{openproblem}
	For \(n>0\), is the first possible Adams--Novikov differential for \(T(n)\), where \(BP_*T(n)=BP_*[t_1,\ldots, t_n]\), nontrivial?
\end{openproblem}


\plainproblemsection

\begin{openproblem}
	For \(p>2\), find the least \(k\) such that \(\beta_1^k=0\) in the stable homotopy groups of spheres.
\end{openproblem}


\plainproblemsection

\begin{openproblem}
	For a finite group \(G\), let \(F_n\subset A(G)^\wedge\) be the kernel of \[A(G)^\wedge=[BG, S^0]\longrightarrow[BG, L_nS^0],\]where \(L_nS^0\) is the Bousfield localization with respect to \(v_{n-1}BP\). Is a virtual finite \(G\)-set \(X\) in \(F_n\) exactly when \(\#(X^H)=0\) for every subgroup \(H\leq G\) generated by at most \(n\) elements?
\end{openproblem}


\plainproblemsection

\begin{openproblem}
	Compute \(K(n)_*(\Omega^k\Sigma^kX)\), either as a functor of \(K(n)_*X\) or by proving convergence of a suitable Eilenberg--Moore spectral sequence.
\end{openproblem}


\plainproblemsection

\begin{openproblem}
	Relate the \(v_n\)-torsion or \(v_n\)-periodic behaviour of \(\alpha\in[X, S^0]_j\) to that of its root invariant \(R(\alpha)\in[X, S^0]_{j+N}\).
\end{openproblem}


\plainproblemsection

\begin{openproblem}
	Determine the higher associativity of the pullback of a sphere extension of a Lie group along a degree-\(k\) map.
\end{openproblem}


\plainproblemsection

\begin{openproblem}
	Does there exist a space admitting an \(A_n\)-structure but no \(A_n\)-primitive \(A_n\)-structure? In particular, does such a space exist for \(n=2\)?
\end{openproblem}


\plainproblemsection

\begin{openproblem}
	Is there a three-connected homotopy associative H-space?
\end{openproblem}


\plainproblemsection

\begin{openproblem}
	When does the Bar construction functor induce a weak equivalence from the space of homeomorphisms between compact Lie groups to the mapping space between their classifying spaces?
\end{openproblem}


\plainproblemsection

\begin{openproblem}
	Justify the equivariant theory (homotopy theory, simple homotopy theory, algebraic K-theory, etc.) for non-compact Lie groups, or make clear the essential obstruction.
\end{openproblem}


\plainproblemsection

\begin{openproblem}
	Let \(Y\) be an \(H\)-space, and suppose \(f: Y\to Z\) factors as \[Y\xrightarrow{(\bar\Delta\wedge1)\bar\Delta}Y\wedge Y\wedge Y\longrightarrow Z.\]If \(X\) is the fibre of \(f\), does \(X\) split as \(\Omega Y\times\Omega Z\) as a homotopy-commutative \(H\)-space?
\end{openproblem}


\plainproblemsection

\begin{openproblem}
	Let \(X\) be a simply connected finite \(H\)-space, and let \(A=H^*(X;\mathbb Z_2)\) be its cohomology Hopf algebra over the Steenrod algebra \(A(2)\). Are there irreducible Hopf algebras over \(A(2)\) such that \(A\) splits as their tensor product?
\end{openproblem}


\plainproblemsection

\begin{openproblem}
	Can a finite loop space \(\Omega B\) have \[H^*(B;\mathbb Q)=\mathbb Q[x_{n_1},\ldots, x_{n_r}],\qquad (n_1,\ldots, n_r)=(4,4,48,8,812,12,16,16,20,24,24,28)?\]
\end{openproblem}


\plainproblemsection

\begin{openproblem}
	Let \(X\) be a simply connected finite \(H\)-space. Determine the action of \(A(2)\) on \(PH^*(\Omega X;\mathbb Z_2)\).
\end{openproblem}


\plainproblemsection

\begin{openproblem}
	If \(X\) is a finite loop space, must \(H^*(X;\mathbb Z)\) be isomorphic to the integral cohomology of a Lie group?
\end{openproblem}


\plainproblemsection

\begin{openproblem}
	Construct explicit free actions of finite groups on homology \(3\)-spheres of the type arising from the space-form surgery problem. Which groups can occur as the fundamental group of the underlying homology \(3\)-sphere?
\end{openproblem}


\plainproblemsection

\begin{openproblem}
	Let \(\pi\) be finite. Does the reduced surgery obstruction map \[\bar\sigma:\Omega_n(B\pi\times G/\mathrm{TOP})\longrightarrow L_n'(\mathbb Z\pi)/L_n(\mathbb Z)\]factor through \(\bigoplus_{i=1}^{3}H_i(\pi;\mathbb Z/2)\)?
\end{openproblem}


\plainproblemsection

\begin{openproblem}
	Determine the relationship between \(L\)-theory and Hermitian \(K\)-theory in the sense of Quillen, for rings with involution.
\end{openproblem}


\plainproblemsection

\begin{openproblem}
	Let \(L/K\) be a local field extension, let \(G(L/K)\) be its finite Galois group, and let \(W_K\) denote the Deligne--Langlands local-constant homomorphism. Is \(W_K\) trivial on \(I O(G(L/K))^N\cdot R_{\mathrm{Sp}}(G(L/K))\) for some integer \(N\)?
\end{openproblem}


\plainproblemsection

\begin{openproblem}
	The Atiyah--Patodi--Singer eta invariant gives nonsingular pairings that detect the \(K\)-theory and equivariant unitary bordism of spherical space forms. What are the corresponding results for space forms of constant curvature \(0\) or \(-1\), and why is the eta invariant so effective in this setting?
\end{openproblem}


\plainproblemsection

\begin{openproblem}
	For a spherical space-form group \(G\), prove topologically the conjectural splitting \[\widetilde\Omega_*^U(BG)\cong bu_*(BG)\oplus\mathbb Z[X_4, X_6,\ldots].\]Analytic proofs are known for \(G=\mathbb Z_4\) and \(G=\{\pm1,\pm i,\pm j,\pm h\}\).
\end{openproblem}


\plainproblemsection

\begin{openproblem}
	Let \(G\) and \(G'\) be compact simple Lie groups. If the homogeneous spaces \(G/H\) and \(G'/H'\) are homeomorphic, must they be diffeomorphic?
\end{openproblem}


\plainproblemsection

\begin{openproblem}
	Let \(X^n\subset\mathbb{CP}^{n+r}\) be a complete intersection. Is its diffeomorphism type determined by its dimension, total degree, intersection pairing or Arf invariant, and its Stiefel--Whitney and Pontryagin classes?
\end{openproblem}


\plainproblemsection

\begin{openproblem}
	For \(n\notin\{2,4,6\}\), is every compact Kähler manifold homotopy equivalent to \(\mathbb{CP}^n\) biholomorphic to \(\mathbb{CP}^n\)?
\end{openproblem}


\plainproblemsection

\begin{openproblem}
	For PL actions of a finite group, determine the first nontrivial obstructions to replacing an equivariant homotopy equivalence by an isovariant homotopy equivalence when the usual gap hypothesis on fixed-point dimensions is weakened. Can these obstructions be expressed as linking phenomena?
\end{openproblem}


\plainproblemsection

\begin{openproblem}
	In the equivariant-to-isovariant homotopy problem for PL finite group actions, can the exceptional fixed-set dimension pair \((1,4)\) be removed from the gap hypothesis?
\end{openproblem}


\plainproblemsection

\begin{openproblem}
	If \(G=\prod_{i=1}^{k}\mathbb Z_p\) acts freely on \(S^{n_1}\times\cdots\times S^{n_\ell}\), must \(k\leq\ell\)?
\end{openproblem}


\plainproblemsection

\begin{openproblem}
	Does there exist an action on \(D^4\) having exactly one fixed point?
\end{openproblem}


\plainproblemsection

\begin{openproblem}
	Let \(M^4\) be a closed topological \(4\)-manifold with a topological circle action. Must its Kirby--Siebenmann invariant vanish?
\end{openproblem}


\plainproblemsection

\begin{openproblem}
	Is every topological circle action on a \(4\)-manifold concordant to a smooth action?
\end{openproblem}


\plainproblemsection

\begin{openproblem}
	Let \(M^4\) be a closed topological \(4\)-manifold that is homotopy equivalent but not homeomorphic to \(\mathbb{CP}^2\). Does \(M\) admit a nontrivial involution?
\end{openproblem}


\plainproblemsection

\begin{openproblem}
	Classify \(4\)-dimensional \(h\)-cobordisms up to homeomorphism or diffeomorphism. In particular, if \(\pi\) is the fundamental group of a closed \(3\)-manifold \(M^3\), which Whitehead torsions are realized by \(h\)-cobordisms having one end equal to \(M\)?
\end{openproblem}


\plainproblemsection

\begin{openproblem}
	Which elements of the projective class group can be realized as the finiteness obstructions for tame ends of topological 4-manifolds?
\end{openproblem}


\plainproblemsection

\begin{openproblem}
	For every natural number \(n\), does there exist \(J(n)\) such that whenever a finite group \(G\) admits an effective \(n\)-dimensional homotopy representation, \(G\) has an abelian normal subgroup \(A\) with \(|G/A|<J(n)\)?
\end{openproblem}


\plainproblemsection

\begin{openproblem}
	Give a purely algebraic definition of the equivariant finiteness-obstruction homomorphism \(s\) arising from the geometric realization of \(\operatorname{Pic}A(G)\) as a subgroup of the homotopy-representation group. Can \(s\) be described as a generalized Swan homomorphism?
\end{openproblem}


\plainproblemsection

\begin{openproblem}
	Classify the \(3\)-dimensional homotopy representations of finite groups.
\end{openproblem}


\paragraph{Knot Theory.}

\plainproblemsection

\begin{openproblem}
	Is every slice knot a ribbon knot?
\end{openproblem}


\plainproblemsection

\begin{openproblem}
	Does the Jones polynomial detect knottedness?
\end{openproblem}


\plainproblemsection

\begin{openproblem}
	Can the diagrammatic approach to links provide insight into properties of the Potts model?
\end{openproblem}


\plainproblemsection

\begin{openproblem}
	What does the phenomenon of phase transition mean for the calculation of link polynomials of large links?
\end{openproblem}


\plainproblemsection

\begin{openproblem}
	Can topological information be extracted from the dynamical behaviour of a quasi-physical system associated with a knot?
\end{openproblem}


\plainproblemsection

\begin{openproblem}
	Is the Potts partition function, viewed as a function on alternating link diagrams, a topological invariant of those diagrams?
\end{openproblem}


\plainproblemsection

\begin{openproblem}
	If \(K\) is a reduced alternating diagram, is \([K](A, B, d)\) an ambient-isotopy invariant of \(K\)?
\end{openproblem}


\plainproblemsection

\begin{openproblem}
	For every pair of distinct links, does some representation of the abstract tensor formalism distinguish them? Can such a representation always be obtained from a Drinfeld double, hence from a quantum group?
\end{openproblem}


\paragraph{Infinite-Dimensional Topology.}

\plainproblemsection

\begin{openproblem}
	Let \(f: X\to Y\) be a cell-like map between locally compact spaces. If \(X\) is an absolute neighborhood retract, under what additional conditions must \(Y\) also be an absolute neighborhood retract?
\end{openproblem}


\plainproblemsection

\begin{openproblem}
	If f: X \(\to\) Y is approximately right invertible, and X is a compact ANR, must Y be an ANR? an FANR?
\end{openproblem}


\plainproblemsection

\begin{openproblem}
	If f: X \(\to\) Y is a refinable map and X is a compact ANR, must Y be an ANR? an FANR?
\end{openproblem}


\plainproblemsection

\begin{openproblem}
	Let \(f: D^n\to Y\) be a cell-like map whose nondegenerate point-inverses are arcs. Must \(Y\) be an absolute retract?
\end{openproblem}


\plainproblemsection

\begin{openproblem}
	Let f: Q \(\to\) Y be a CE map with Y an AR. Suppose that Y \(\times\) F \(\approx\) Q for some finite-dimensional compactum F. Must Y \(\times\) In \(\approx\) Q for some n? What about n =2 or n =1?
\end{openproblem}


\plainproblemsection

\begin{openproblem}
	Let f: In \(\to\) Y be a CE map. Can dim Y be greater than n when n =5? What about n =4?
\end{openproblem}


\plainproblemsection

\begin{openproblem}
	Let f: In \(\to\) Y be CE map. If all point inverses are cells of dimension \(\leq\) 6 must dim Y be \(\leq\) n? What if all point inverses are cells of dimension \(\leq\) 5? \(\leq\) 1?
\end{openproblem}


\plainproblemsection

\begin{openproblem}
	If f: Q \(\to\) Y is a CE map onto an ANR with the DD2Pand if each point inverse is finite dimensional, must Y \(\approx\) Q?
\end{openproblem}


\plainproblemsection

\begin{openproblem}
	Give an explicit construction of a CE dimension-raising map.
\end{openproblem}


\plainproblemsection

\begin{openproblem}
	Classify the Taylor examples: which compacta occur as cell-like images of the Hilbert cube?
\end{openproblem}


\plainproblemsection

\begin{openproblem}
	Does every infinite-dimensional compact ANR contain a closed \(n\)-dimensional subset for every \(n\geq0\)?
\end{openproblem}


\plainproblemsection

\begin{openproblem}
	Let \(X\) be a compact AR such that \(H_*(U, U\setminus A)=0\) for every finite-dimensional compact \(A\subset X\) and open \(U\subset X\). If \(X\) has the disjoint \(2\)-cells property, must it have the disjoint \(n\)-cells property for every \(n\)?
\end{openproblem}


\plainproblemsection

\begin{openproblem}
	Characterize the metric compacta that occur as cell-like images of finite-dimensional ANRs or manifolds.
\end{openproblem}


\plainproblemsection

\begin{openproblem}
	Give an explicit construction of a Dranishnikov compactum, that is, an infinite-dimensional compactum of finite integral cohomological dimension.
\end{openproblem}


\plainproblemsection

\begin{openproblem}
	Must a product of two A-weakly infinite dimensional compacta be A-weakly infinite-dimensional?
\end{openproblem}


\plainproblemsection

\begin{openproblem}
	Must every AWID compactum have Property C?
\end{openproblem}


\plainproblemsection

\begin{openproblem}
	Is there a compact metric space of Pol type (i. e., neither strongly infinite dimensional nor countably infinite dimensional) containing no strongly infinite dimensional subspace? What about strongly countable dimensional subspaces?
\end{openproblem}


\plainproblemsection

\begin{openproblem}
	Does there exist an absolute weakly infinite-dimensional compactum \(X\) such that \(\dim_{\mathbb Z}X<\infty\)?
\end{openproblem}


\plainproblemsection

\begin{openproblem}
	If dimZ(X \(\times\)X)= 3, can a cell-like mapping f: X \(\to\) Y raise dimension?
\end{openproblem}


\plainproblemsection

\begin{openproblem}
	Does there exist a generalized homology theory \(h_*\) for which \(h_*(\mathbb{CP}^{\infty})=0\)?
\end{openproblem}


\plainproblemsection

\begin{openproblem}
	Is dim(X)= dimS (X) where S denotes stable cohomotopy?
\end{openproblem}


\plainproblemsection

\begin{openproblem}
	Does there exist an infinite-dimensional compact ANR \(X\) and a group \(G\) such that \(\dim_GX<\infty\)?
\end{openproblem}


\plainproblemsection

\begin{openproblem}
	Let \(X\) be a separable metrizable space with finite integral cohomological dimension. Does \(X\) admit a metrizable compactification of finite integral cohomological dimension?
\end{openproblem}


\plainproblemsection

\begin{openproblem}
	Is there for each n a universal metric compactum X for the class of metric compacta of integral cohomological dimension n?
\end{openproblem}


\plainproblemsection

\begin{openproblem}
	Let M \(\to\) f X \(\leftarrow\)g N be a diagram of CE maps, where both M and N are compact Q-manifolds and X is a metric compactum. Is there a homeomorphism h: M \(\to\) N such that g \(\circ\) h is arbitrarily close to f?
\end{openproblem}


\plainproblemsection

\begin{openproblem}
	If X and Y are shape equivalent UV1 compacta, is there a finite diagram where each \(\leftarrow\)\(\to\) is an hereditary shape equivalence either from Xi to Xi+1 or from Xi+1 to Xi?
\end{openproblem}


\plainproblemsection

\begin{openproblem}
	Characterize ANR divisors. Is the property of being an ANR divisor invariant under shape domination?
\end{openproblem}


\plainproblemsection

\begin{openproblem}
	When is the one-point compactification of a locally compact ANR an ANR?
\end{openproblem}


\plainproblemsection

\begin{openproblem}
	Is there a theory like that of Chapman and Siebenmann for completing a non-compact Q-manifold into a compact one by adding a shape Z-set?
\end{openproblem}


\plainproblemsection

\begin{openproblem}
	Are there analogues of Chapman's Complement Theorem for shape \(Z\)-sets?
\end{openproblem}


\plainproblemsection

\begin{openproblem}
	Let \(X\) be an FANR and a \(Z\)-set in the Hilbert cube \(Q\), and let \(h: X\to X\) be homotopic to the identity in every neighborhood of \(X\) in \(Q\). Does there exist a nested sequence \(M_i\supseteq M_{i+1}\) of compact \(Q\)-manifold neighborhoods with \(\bigcap_iM_i=X\), together with an extension \(H: Q\to Q\) satisfying \(H(M_i)=M_i\) for all \(i\)?
\end{openproblem}


\plainproblemsection

\begin{openproblem}
	Let \(X, Y\) be \(Z\)-sets in \(Q\). If \(f: Q\setminus X\to Q\setminus Y\) is proper and a weak proper homotopy equivalence, must it be a proper homotopy equivalence? Equivalently, is every shape equivalence of metric compacta a strong shape equivalence?
\end{openproblem}


\plainproblemsection

\begin{openproblem}
	Let \(X, Y\) be connected \(Z\)-sets in \(Q\), with base rays chosen at the ends of their complements. If a proper base-ray-preserving map \(f: Q\setminus X\to Q\setminus Y\) is invertible in base-ray-preserving weak proper homotopy theory, must \(f\) be a proper homotopy equivalence?
\end{openproblem}


\plainproblemsection

\begin{openproblem}
	Let \(i: X\hookrightarrow Y\) be an embedding of compacta and a shape equivalence. If \(f, g: Y\to P\) map to an ANR and \(f|_X=g|_X\), must \(f\simeq g\) relative to \(X\)?
\end{openproblem}


\plainproblemsection

\begin{openproblem}
	Can one choose a representative in the shape class of every \(UV^1\) compactum so that strong-shape theory and homotopy theory agree on the resulting class?
\end{openproblem}


\plainproblemsection

\begin{openproblem}
	Let \((X, x)\) be a pointed connected compactum for which \(\operatorname{pro}\text{-}\pi_i(X, x)\) is stable for every \(i\). Must \(X\) be shape equivalent to a locally contractible compactum?
\end{openproblem}


\plainproblemsection

\begin{openproblem}
	Let X be compact, locally connected and dominated by a finite complex. Must X be homotopy equivalent with one? What if X is locally 1-connected?
\end{openproblem}


\plainproblemsection

\begin{openproblem}
	Let \(\Omega\) be the set of all countable ordinals. Does there exist a function \(\beta\):\(\Omega\) \(\to\) \(\Omega\) such that if f: X \(\to\) Y is an hereditary shape equivalence between two countable dimensional compacta, then ind(Y ) \(\leq\) \(\beta\)(ind(X))?
\end{openproblem}


\plainproblemsection

\begin{openproblem}
	If X and Y are shape equivalent continua with pro-\(\pi_1\) profinite, are X and Y CE equivalent? Are they UV \(\omega\)-equivalent?
\end{openproblem}


\plainproblemsection

\begin{openproblem}
	Let X be a finitely dominated compactum with Euler characteristic \(\chi\)(X)= 0. Is the Nielsen number of the identity map of X always zero?
\end{openproblem}


\plainproblemsection

\begin{openproblem}
	Let \(f: X\to Y\) be a map between compact metric spaces for which Euler characteristics are defined, and suppose \(\chi(f^{-1}(y))=0\) for every \(y\in Y\). Under what conditions must \(\chi(X)=\chi(Y)\)?
\end{openproblem}


\plainproblemsection

\begin{openproblem}
	Let \(X\) be the identity component of the homeomorphism group of a closed PL manifold of dimension at least \(5\), and let \(X_0\subset X\) be the subspace of PL homeomorphisms. Is \(X\setminus X_0\) hazy in \(X\)?
\end{openproblem}


\plainproblemsection

\begin{openproblem}
	Let X be a non trivial homogeneous contractible compactum. Is X an AR? Is X a Hilbert cube?
\end{openproblem}


\plainproblemsection

\begin{openproblem}
	When are homogeneous spaces ANR's?
\end{openproblem}


\plainproblemsection

\begin{openproblem}
	Is Q the only homogeneous non-degenerate compact AR?
\end{openproblem}


\plainproblemsection

\begin{openproblem}
	Is Q the only homogeneous continuum homeomorphic with its own cone?
\end{openproblem}


\plainproblemsection

\begin{openproblem}
	Let f: A \(\to\)X be a CE map from an ANR A onto X. If X is not strongly infinite-dimensional, must X be an ANR?
\end{openproblem}


\plainproblemsection

\begin{openproblem}
	For \(n\geq3\), is the Banach--Mazur compactum \(Q(n)\) a Hilbert cube?
\end{openproblem}


\plainproblemsection

\begin{openproblem}
	Let M be a compact Q-manifold, and let f: M \(\to\) M be a map such that f2 is homotopic to the identity. When is f homotopic to an involution?
\end{openproblem}


\plainproblemsection

\begin{openproblem}
	Let \(\pi\) be a group such that there exist compact K(\(\pi\), 1) Q-manifolds M and N. Must they be homeomorphic?
\end{openproblem}


\plainproblemsection

\begin{openproblem}
	Let \(\alpha\) be an open cover of N acompact Q-manifold. What conditions imply that any \(\alpha\)-equivalence f: M \(\to\) N is homotopic to a homeomorphism?
\end{openproblem}


\plainproblemsection

\begin{openproblem}
	Let M be a compact Q-manifold and U a finite open cover of M by contractible open sets such that the intersections of subcollections of U are either contractible or void. Is M homeomorphic with N(U) \(\times\) Q?
\end{openproblem}


\plainproblemsection

\begin{openproblem}
	Is there a Q-manifold version of Quinn's End theorems?
\end{openproblem}


\plainproblemsection

\begin{openproblem}
	Let p: M \(\to\) B be a locally trivial bundle, where B is a locally compact polyhedron and the fibers are Q-manifolds. Does there exist a locally compact polyhedron P, a PL map q: P \(\to\)B, and a fiber preserving homeomorphism h: M \(\to\)P \(\times\)Q?
\end{openproblem}


\plainproblemsection

\begin{openproblem}
	Let p: M \(\to\) B be a Hurewicz fibration where M is a compact Q-manifold. Is B an ANR?
\end{openproblem}


\plainproblemsection

\begin{openproblem}
	Let p: M \(\to\) \(S^{2}\) be an approximate fibration, where M is acompact Q-manifold. Must the fiber have the shape of a finite complex?
\end{openproblem}


\plainproblemsection

\begin{openproblem}
	Let f: M \(\to\) B be a map, with M acompact Q-manifold and B a compact polyhedron. What is the obstruction to homotoping f to an approximate fibration?
\end{openproblem}


\plainproblemsection

\begin{openproblem}
	Let \(X\subset Q\) be a connected pointlike compactum that is the closure of its interior, and suppose \(Q\setminus X\) is pointlike. Does \(Q\setminus X\) have a compactification \(Y\) whose remainder is a \(Q\)-manifold and to which every homeomorphism of the pair \((Q, X)\) extends?
\end{openproblem}


\plainproblemsection

\begin{openproblem}
	Suppose that A is a closed finite-dimensional subset of Q that can be instantly isotoped off itself. Must A be a Z-set?
\end{openproblem}


\plainproblemsection

\begin{openproblem}
	Does there exist a Cantor set \(C\subset Q\) and \(\varepsilon>0\) such that \(h(C)\cap C=\varnothing\) for every homeomorphism \(h: Q\to Q\) that is \(\varepsilon\)-close to the identity?
\end{openproblem}


\plainproblemsection

\begin{openproblem}
	Is every homeomorphism of s the composition of two conjugates of homeomorphisms that extend to Q?
\end{openproblem}


\plainproblemsection

\begin{openproblem}
	Is there a Cech-cohomology version of the fibred general position theory of Toruńczyk and West along the lines of Daverman and Walsh that detects the Q-manifold bundles among the ANR-fibrations with compact fibers?
\end{openproblem}


\plainproblemsection

\begin{openproblem}
	Is there a tameness condition at infinity that will extend the ``general position fibrations are bundles'' theorem of Toruńczyk and West to more bases?
\end{openproblem}


\plainproblemsection

\begin{openproblem}
	Let \(p: B\times Q\to B\) be projection and let \(A\subset B\times Q\) be closed. Under what conditions is \[p|_{(B\times Q)\setminus A}:(B\times Q)\setminus A\to B\]a Hurewicz fibration, an ANR-fibration, or a bundle?
\end{openproblem}


\plainproblemsection

\begin{openproblem}
	Let \(B\) be a metrizable space and \(\varepsilon>0\). Does there exist \(\delta>0\) such that, whenever \(M, N\) are compact \(Q\)-manifolds, \(p: N\to B\) is \(UV^1\), and \(f: M\to N\) is a \(p^{-1}(\delta)\)-equivalence, the map \(f\) is \(p^{-1}(\varepsilon)\)-homotopic to a homeomorphism?
\end{openproblem}


\plainproblemsection

\begin{openproblem}
	Can the fibred end theorem for compactifying a bundle of noncompact \(Q\)-manifolds be extended to \(Q\)-manifold fibrations over complete metrizable ANRs, or more generally over arbitrary separable metrizable spaces?
\end{openproblem}


\plainproblemsection

\begin{openproblem}
	Let M be a compact manifold or polyhedron, and let N acompact Q-manifold. Classify the embeddings f: M \(\to\) N such that the pair (N, M) is homotopically stratified in the sense of Quinn. (i) Start with M = \(S^{1}\), N = Q as in Toruńczyk and West, (ii) How does the case of M a Q-manifold differ from the finite-dimensional case?
\end{openproblem}


\plainproblemsection

\begin{openproblem}
	Classify \(Q\)-manifold approximate fibrations over compact \(Q\)-manifolds.
\end{openproblem}


\plainproblemsection

\begin{openproblem}
	Extend the program of Hughes, Taylor, and Williams over bases B that are compact polyhedra, not compact manifolds.
\end{openproblem}


\plainproblemsection

\begin{openproblem}
	Let p: M \(\to\) B be an approximate fibration of a compact Q-manifold over an aspherical manifold. If p is homotopic to a bundle projection, may it be approximated by them?
\end{openproblem}


\plainproblemsection

\begin{openproblem}
	Let M be a compact Q-manifold, and let Qs be a Lipschitz-homogeneous Hilbert cube. Are all Lipschitz Qs-structures on M equivalent (via bi-Lipschitz homeomorphisms)?
\end{openproblem}


\plainproblemsection

\begin{openproblem}
	Give a characterization of the Lipschitz-homogeneous convex compacta in normed linear spaces.
\end{openproblem}


\plainproblemsection

\begin{openproblem}
	For a standard Lipschitz-homogeneous Qs, is there a class of Lipschitz Z-sets (i. e., such that every bi-Lipschitz homeomorphism between them extends to one of Qs) large enough to contain homeomorphs of all finite polyhedra?
\end{openproblem}


\plainproblemsection

\begin{openproblem}
	Give a Q-manifold version of Quinn's Finite Structure Spectrum in Ends II.
\end{openproblem}


\plainproblemsection

\begin{openproblem}
	Can a \(Q\)-manifold version of Quinn's finite-structure spectrum be used to give a common framework for the controlled topology theories of Quinn, Anderson--Munkholm, Chapman, Hughes--Taylor--Williams, and Farrell--Jones?
\end{openproblem}


\plainproblemsection

\begin{openproblem}
	Let M be a compact Q-manifold. Let Cz(M) denote the smallest integer k such that M may be covered with k open sets each homeomorphic with Q \(\times\) [0, 1). Is it always true that Cz (M \(\times\) \(S^n\))= Cz (M)+1, where \(S^n\) is the n-sphere?
\end{openproblem}


\plainproblemsection

\begin{openproblem}
	Let \(Q=\prod_{i>0}[-1,1]_i\), and let \(h: Q\to Q\) be a basepoint-free involution. Must \(h\) be topologically conjugate to the linear antipodal action \(x\mapsto-x\)?
\end{openproblem}


\plainproblemsection

\begin{openproblem}
	Let \(h: Q\to Q\) be a based-free homeomorphism of finite order. Must the corresponding cyclic action be topologically conjugate to a standard linear action?
\end{openproblem}


\plainproblemsection

\begin{openproblem}
	Does there exist a sequence of principal \(\mathbb Z_2\)-bundles of CW-complexes and bundle maps, each restricting to the identity on \(S^\infty\), such that the induced maps of base pairs \[f_i^{-1}:(B_i,\mathbb{RP}^\infty)\to(B_{i-1},\mathbb{RP}^\infty)\]satisfy: each pair is relatively finite and relatively one-connected, each bundle map is null-homotopic, and every finite composite of the \(f_i\) is essential as a map of pairs?
\end{openproblem}


\plainproblemsection

\begin{openproblem}
	If \(\alpha\) is a standard action of a finite cyclic group G on Q, is HG(Q) a Hilbert manifold? Conversely, if HG(Q) is a Hilbert manifold and \(\alpha\) is a based-free action on Q, is \(\alpha\) standard?
\end{openproblem}


\plainproblemsection

\begin{openproblem}
	Under what conditions can a non-free action \(\alpha\) of a compact group on a Q-manifold M be factored as a diagonal action where N is a finite-dimensional manifold, polyhedron, or ANR? (Is there any difference between these questions?)
\end{openproblem}


\plainproblemsection

\begin{openproblem}
	Let \(\alpha\) be a semi-free action of a finite group G on Q with fixed point set F a Hilbert cube Z-set. When is \(\alpha\) equivalent with the product \(\sigma\) \(\times\) idF, where \(\sigma\) is the standard action of G on Q? What if F = In?
\end{openproblem}


\plainproblemsection

\begin{openproblem}
	Is there an effective action by a p-adic group on an n-manifold?
\end{openproblem}


\plainproblemsection

\begin{openproblem}
	Does there exist a smooth minimal diffeomorphism on separable Hilbert space or on each connected separable Hilbert manifold? What about a minimal smooth flow?
\end{openproblem}


\plainproblemsection

\begin{openproblem}
	Does every compact connected Q-manifold with zero Euler characteristic admit a minimal homeomorphism? Does it admit a minimal flow?
\end{openproblem}


\plainproblemsection

\begin{openproblem}
	Let G be the group of PL homeomorphisms of the unit interval with singularities in Z[1/6] and slopes in the multiplicative group generated by 2/3. Is G finitely generated?
\end{openproblem}


\plainproblemsection

\begin{openproblem}
	Given a non-compact QG-manifold M, what is the obstruction to compactifying it to a QG-manifold by the addition of an equivariant Z-set?
\end{openproblem}


\plainproblemsection

\begin{openproblem}
	Let X \(\subset\) M be a compact G-ANR in a QG-manifold. Does X always have a mapping cylinder neighborhood N? If not, what is the obstruction, and what about the possibility of the boundary \(\partial\)(N) of N not being a QG-manifold but some other G-space?
\end{openproblem}


\plainproblemsection

\begin{openproblem}
	If M is a compact QG-manifold, is the space HG(M) of equivariant homeomorphisms of M a G-ANR?
\end{openproblem}


\plainproblemsection

\begin{openproblem}
	For a compact \(Q_G\)-manifold \(M\), determine the effect of stabilization by \(Q_G\) on \(\operatorname{Wh}_G(M)\).
\end{openproblem}


\plainproblemsection

\begin{openproblem}
	Describe the subgroup of \(\operatorname{Wh}_G(K)\) represented by \(G\)-CW pairs \((L, K)\) whose stabilized inclusion is homotopic to an equivariant homeomorphism. In particular, if such an element is represented by a smooth \(G\)-\(h\)-cobordism \((W^{n+1}; M^n, N^n)\), all nonempty components of fixed-point sets have dimension at least \(6\), and none has codimension \(2\) in another, must the cobordism be equivariantly homeomorphic to \((M\times I; M\times\{0\}, M\times\{1\})\)?
\end{openproblem}


\plainproblemsection

\begin{openproblem}
	If (N, M) is a pair of compact QG-manifolds such that M is an equiv-ariant Z-set in N and the inclusion is an equivariant homotopy equivalence, is there a finite-dimensional G-h-cobordism (W, M) that stabilizes to (N, M), where the G-action is locally smoothable?
\end{openproblem}


\plainproblemsection

\begin{openproblem}
	Let G = Zp and let AG(Q) denote the space of G-actions on Q with the uniform convergence topology. Is AG(Q) connected? Path connected? Locally connected? Locally contractible? An ANR? A Hilbert manifold?
\end{openproblem}


\plainproblemsection

\begin{openproblem}
	What about the subspace AG(Q, A) of Zp actions with prescribed fixed point set A?
\end{openproblem}


\plainproblemsection

\begin{openproblem}
	Let \(\mathcal A_A^G(Q)\) be the space of \(\mathbb Z_p\)-actions on \(Q\) whose fixed-point set is a \(Z\)-set homeomorphic to \(A\). Is this space connected, path connected, locally connected, locally contractible, an ANR, or a Hilbert manifold?
\end{openproblem}


\plainproblemsection

\begin{openproblem}
	Let C be a convex Hilbert cube in a locally convex linear space X. Let Gl(C) denote the restrictions to C of those invertible, continuous linear transformations T of X such that T (C)= C. Give necessary and sufficient conditions that two transformations S, T \(\in\) Gl(C) be topologically conjugate on C.
\end{openproblem}


\plainproblemsection

\begin{openproblem}
	Let \(C\) be a convex Hilbert cube in a locally convex linear space, and let \(S, T\in\operatorname{GL}(C)\) be periodic. When are \(S\) and \(T\) topologically conjugate if the fixed-point sets of \(S^k\) and \(T^k\) are infinite-dimensional for every \(k\)?
\end{openproblem}


\plainproblemsection

\begin{openproblem}
	Let \(C\) be a convex Hilbert cube in a locally convex linear space, and let \(S, T\in\operatorname{GL}(C)\). Suppose the only finite orbit of either map is the origin. When are \(S\) and \(T\) topologically conjugate on \(C\)?
\end{openproblem}


\plainproblemsection

\begin{openproblem}
	Let X be a locally compact, convex subset of 2, and let G be a discrete group acting freely and properly discontinuously on X. Under what conditions is the action of G linearizeable, i. e., topologically conjugate with the restriction to some locally compact, convex subset C of 2 of a subgroup of Gl(C)?
\end{openproblem}


\plainproblemsection

\begin{openproblem}
	For a given finite or discrete group G classify the free properly discontinuous actions of G on contractible Q-manifolds.
\end{openproblem}


\plainproblemsection

\begin{openproblem}
	Is every finitely presented group semistable at infinity?
\end{openproblem}


\plainproblemsection

\begin{openproblem}
	Let \(M^n\) be a compact \(n\)-manifold. Is its homeomorphism group \(\mathcal H(M^n)\), with the uniform topology, a Hilbert manifold?
\end{openproblem}


\plainproblemsection

\begin{openproblem}
	Is every open subset of \(H_{\partial}(B^n)\) dominated by a CW-complex?
\end{openproblem}


\plainproblemsection

\begin{openproblem}
	Let \(M^n\) be a compact manifold, let \(\mathcal H(M)\) be its homeomorphism group, and let \(\overline{\mathcal H(M)}\) be its closure in the uniform topology on continuous self-maps of \(M\). Does there exist a continuous self-map of \(\overline{\mathcal H(M)}\), arbitrarily close to the identity, whose image lies in \(\mathcal H(M)\)?
\end{openproblem}


\plainproblemsection

\begin{openproblem}
	Let \(M\) be a compact manifold, and let \(\overline{\mathcal H}(M)\) be the closure of its homeomorphism group in the uniform topology on continuous self-maps. Is \(\overline{\mathcal H}(M)\) an absolute neighborhood retract?
\end{openproblem}


\plainproblemsection

\begin{openproblem}
	Let \(M\) be a compact \(Q\)-manifold, and let \(\overline{\mathcal H}(M)\) be the closure of its homeomorphism group in the uniform topology on continuous self-maps of \(M\). Is \(\overline{\mathcal H}(M)\) a Hilbert manifold?
\end{openproblem}


\plainproblemsection

\begin{openproblem}
	Let \(M^n\) be a compact manifold, let \(\mathcal R(M)\) be its space of retractions, and put \(\mathcal R_0(M)=\mathcal R(M)\setminus\{\operatorname{id}_M\}\). Is \(\mathcal R(M)\) an ANR? If \(\partial M=\varnothing\), is \(\mathcal R_0(M)\) a Hilbert manifold? In that case, is \(\mathcal R(M)\) itself a Hilbert manifold?
\end{openproblem}


\plainproblemsection

\begin{openproblem}
	Let \(M\) be a compact manifold, let \(\mathcal R(M)\) be its space of retractions with the compact-open topology, and let \(\mathcal H(M)\) be its homeomorphism group. Is \(\mathcal R(M)\setminus\{\operatorname{id}_M\}\) locally homeomorphic to \(\mathcal H(M)\)?
\end{openproblem}


\plainproblemsection

\begin{openproblem}
	If Y is not a smooth manifold, can one put a smooth structure on C(X, Y ) in a ``natural'' way?
\end{openproblem}


\plainproblemsection

\begin{openproblem}
	Can \(C(X, Y)\) be endowed naturally with a topology making it a \(Q^\infty\)-manifold?
\end{openproblem}


\plainproblemsection

\begin{openproblem}
	Can \(PL(X, Y)\) be endowed naturally with a topology making it an \(\mathbb R^\infty\)-manifold?
\end{openproblem}


\plainproblemsection

\begin{openproblem}
	Under what conditions is \(\operatorname{LIP}(X, Y)\) a \(\Sigma\)-manifold? Is it sufficient that \(Y\) be a locally compact absolute local neighborhood extensor for metric spaces and locally Lipschitz maps?
\end{openproblem}


\plainproblemsection

\begin{openproblem}
	For compact metric \(X\), metric \(Y\), and \(k>0\), let \(k\text{-}\operatorname{LIP}(X, Y)\) be the space of Lipschitz maps with Lipschitz constant at most \(k\). Under what conditions is \(k\text{-}\operatorname{LIP}(X, Y)\) a Hilbert cube manifold?
\end{openproblem}


\plainproblemsection

\begin{openproblem}
	If X and Y are Euclidean polyhedra, is an (2,\(\Sigma\),\(\sigma\))-manifold triple?
\end{openproblem}


\plainproblemsection

\begin{openproblem}
	If X is a Lipschitz n-manifold, is HLIP(X) a \(\Sigma\)-manifold?
\end{openproblem}


\plainproblemsection

\begin{openproblem}
	If X is a Q-manifold, what metric conditions on X guarantee that HLIP (X) be a \(\Sigma\)-manifold? A cap set for H(X)?
\end{openproblem}


\plainproblemsection

\begin{openproblem}
	For a polyhedron \(X\), is \(\mathcal H_{\mathrm{fd}}(X\times Q)\) an \(\ell^2\times\sigma\)-manifold?
\end{openproblem}


\plainproblemsection

\begin{openproblem}
	Are and (2,\(\Sigma\),\(\sigma\))-manifold triples when X is a PL manifold?
\end{openproblem}


\plainproblemsection

\begin{openproblem}
	Let X be a PL manifold. Is RoPL(X) (respectively, RPL(X \(\times\)Q))a \(\sigma\)-manifold or a fd-cap set for Ro(X) (respectively, R(X \(\times\)Q))?
\end{openproblem}


\plainproblemsection

\begin{openproblem}
	Let X be a Q-manifold. Is there a metric condition on X ensuring that RLIP (X) be a \(\Sigma\)-manifold or cap set for RLIP (X)?
\end{openproblem}


\plainproblemsection

\begin{openproblem}
	Let X be a Lipschitz \(n\)-manifold. Is RoLIP (X) a \(\Sigma\)-manifold or a cap set for Ro(X)?
\end{openproblem}


\plainproblemsection

\begin{openproblem}
	Let \(\lambda\) denote Lebesgue product measure on Q. Let H\(\lambda\)(Q) denote the subgroup of H(Q) of measure preserving homeomorphisms (under the compact-open topology). Is H\(\lambda\)(Q) locally contractible? An AR? A Hilbert space? What about H\(\lambda\)(In)?
\end{openproblem}


\plainproblemsection

\begin{openproblem}
	Let F be the group of PL homeomorphisms of the line \(\mathbb R^{1}\) generated by p(t) and q(t), where p(t)= t if t \(\leq\)0, p(t)= 2t, if 0 \(\leq\)t \(\leq\)1, and p(t)= t+1, if t \(\geq\)1, while q(t) is analogously defined using 1 and 2 as the singularities, again with slope 2 between them. Is F amenable?
\end{openproblem}


\plainproblemsection

\begin{openproblem}
	Is every infinite dimensional F -space X homeomorphic with X \(\times\) s? with X \(\times\) Q? with X \(\times\) R?
\end{openproblem}


\plainproblemsection

\begin{openproblem}
	Are compact subsets negligible in every infinite-dimensional \(F\)-space?
\end{openproblem}


\plainproblemsection

\begin{openproblem}
	If \(K\) is a compact convex subset of an \(F\)-space, must \(K\) be an absolute retract?
\end{openproblem}


\plainproblemsection

\begin{openproblem}
	For every \(\varepsilon>0\), does there exist an open cover \(\mathcal U\) of \(\ell^1\) such that, at each point \(p\), the sum of the diameters of the members of \(\mathcal U\) containing \(p\) is less than \(\varepsilon\)?
\end{openproblem}


\plainproblemsection

\begin{openproblem}
	Let W be a convex subset of a Hilbert space. Under what conditions is W locally homotopy-negligible in its closure?
\end{openproblem}


\plainproblemsection

\begin{openproblem}
	Is every locally contractible closed additive subgroup of a Hilbert space an ANR?
\end{openproblem}


\plainproblemsection

\begin{openproblem}
	Is every equiconnected space a retract of a convex subset of an F -space?
\end{openproblem}


\plainproblemsection

\begin{openproblem}
	Let \(U\) be Urysohn's universal separable metric space, let \(U_0\) be an isometric copy of \(U\) in a Banach space, containing \(0\), and put \(V=\overline{\operatorname{span}}(U_0)\). Is \(V\) universal for separable Banach spaces up to linear isometry? If not, which separable Banach spaces embed linearly and isometrically in \(V\)?
\end{openproblem}


\plainproblemsection

\begin{openproblem}
	Let \(U_0\) be an isometric copy of Urysohn's universal separable metric space in a Banach space, containing zero, and let \(U\) be its closed linear span. Does \(U\) have a Schauder basis?
\end{openproblem}


\plainproblemsection

\begin{openproblem}
	Let \(M\subset N\) be \(s\)-manifolds and let \(R\subset M\). Suppose \(M\) has local codimension \(1\) at every point of \(M\setminus R\). Must it also have local codimension \(1\) at the points of \(R\) when \(R\) is a point, a compactum, or a \(Z\)-set in both \(M\) and \(N\)?
\end{openproblem}


\plainproblemsection

\begin{openproblem}
	Let \(M\subset N\) be \(s\)-manifolds and \(R\subset M\). Suppose \(M\) has local codimension \(k>2\) at every point of \(M\setminus R\). Must it have local codimension \(k\) at the points of \(R\) when \(R\) is (a) a point, (b) compact, or (c) a \(Z\)-set in both \(M\) and \(N\)?
\end{openproblem}


\plainproblemsection

\begin{openproblem}
	Let \(M\) be a separable \(C^\infty\) Hilbert manifold modeled on \(\ell^2\). Can every self-homeomorphism of \(M\) be approximated by diffeomorphisms?
\end{openproblem}


\plainproblemsection

\begin{openproblem}
	Let M and K be s-manifolds with K \(\subset\) M and K a Zset in M. Under what conditions on the pair (M, K) does there exist an embedding h of M in s such that the topological boundary of h(M) is h(K)?
\end{openproblem}


\plainproblemsection

\begin{openproblem}
	Are \((C^p, b^*)\)-manifolds stable? Do they embed as open subsets of the model space? Are they classified by homotopy type?
\end{openproblem}


\plainproblemsection

\begin{openproblem}
	When does a Hilbert manifold support a topological group structure?
\end{openproblem}


\plainproblemsection

\begin{openproblem}
	If M is a Hilbert manifold that admits the structure of an associative H-space, does it admit a topological group structure?
\end{openproblem}


\plainproblemsection

\begin{openproblem}
	Let \(\mathcal G\) be a decomposition of \(\sigma\) such that \(\sigma/\mathcal G\not\cong\sigma\) but \(\sigma/\mathcal G\) is a countable union of finite-dimensional compacta. Is \((\sigma/\mathcal G)\times\mathbb R^2\cong\sigma\)? If so, is \((\sigma/\mathcal G)\times\mathbb R\cong\sigma\)?
\end{openproblem}


\plainproblemsection

\begin{openproblem}
	Let \(\mathcal G\) be a cell-like upper-semicontinuous decomposition of \(\sigma\) such that \(\sigma/\mathcal G\) is not homeomorphic to \(\sigma\) but is a countable union of finite-dimensional compacta. Let \(\mathcal H\) be the cell-like decomposition of \(\sigma\times\mathbb R\) whose nondegenerate elements are \(g\times\{0\}\), \(g\in\mathcal G\). Is \((\sigma\times\mathbb R)/\mathcal H\cong\sigma\)?
\end{openproblem}


\plainproblemsection

\begin{openproblem}
	Find intrinsic conditions on metrics of \(\sigma\)-manifolds and \(\Sigma\)-manifolds so that their metric completions are l2-manifolds.
\end{openproblem}


\plainproblemsection

\begin{openproblem}
	Let \(G\) be a complete metrizable topological group that is an ANR. Must \(G\) be a manifold modeled on a Fréchet space? In particular, if \(G\) is separable and not locally compact, must it be an \(s\)-manifold?
\end{openproblem}


\plainproblemsection

\begin{openproblem}
	Let \(H\) be a nonseparable Hilbert space. If \(X\times s\cong H\), must \(X\cong H\)?
\end{openproblem}


\plainproblemsection

\begin{openproblem}
	Let \(X\) be a topologically complete separable metric ANR and let \(Y\subset X\) be an \(s\)-manifold. Under what conditions must \(X\) itself be an \(s\)-manifold?
\end{openproblem}


\plainproblemsection

\begin{openproblem}
	Let \(G\) be a locally contractible separable metrizable group that is a countable union of compact finite-dimensional subsets and is not locally compact. Must \(G\) be a \(\sigma\)-manifold?
\end{openproblem}


\plainproblemsection

\begin{openproblem}
	Let \(X=\varprojlim\{X_n, f_n\}\), where every \(X_n\) is an absolute retract and every bonding map \(f_n\) is cell-like. Give practical conditions ensuring \(X\cong Q\). In particular, is it sufficient that every point-inverse of every projection \(X\to X_n\) be infinite-dimensional?
\end{openproblem}


\plainproblemsection

\begin{openproblem}
	Mogilski's characterization of \(\sigma\)-manifolds uses the ANR property, a countable union by finite-dimensional compacta, the strong \(Z\)-set property for compact subsets, and a relative approximation-by-embeddings property. Can these conditions be replaced by a single condition analogous to the disjoint-disks property?
\end{openproblem}


\plainproblemsection

\begin{openproblem}
	Do the following conditions characterize the \(n\)-dimensional Nöbeling spaces: \(X\) is completely metrizable, \(X\in\mathcal C^{\, n-1}\), and \(X\) satisfies the \(n\)-discrete \(n\)-cells property?
\end{openproblem}


\plainproblemsection

\begin{openproblem}
	Fix \(n>1\) and \(0\leq k\leq n\). Let \(\mathcal M_n^k\) be the class of tame compacta in \(\mathbb R^n\) of dimension at most \(k\), let \(B_n^k\) be an \(\mathcal M_n^k\)-absorber, and put \(s_n^k=\mathbb R^n\setminus B_n^{\, n-k-1}\). Give intrinsic characterizations of the spaces \(B_n^k\) and \(s_n^k\).
\end{openproblem}


\plainproblemsection

\begin{openproblem}
	Has every ANR homology n-manifold a CE resolution by a topological n-manifold?
\end{openproblem}


\plainproblemsection

\begin{openproblem}
	Let \(M\) be a compact \(n\)-manifold with \(n>2\). Is the group \(\mathcal H_{\partial}(M)\) of boundary-fixing homeomorphisms a Hilbert manifold?
\end{openproblem}


\plainproblemsection

\begin{openproblem}
	Let \((X, d)\) and \((Y,\rho)\) be metric spaces with \(X\) compact. Under what conditions is the space \(1\text{-}\operatorname{LIP}(X, Y)\) of maps satisfying \(\rho(f(x), f(x'))\leq d(x, x')\) a Hilbert cube manifold?
\end{openproblem}


\plainproblemsection

\begin{openproblem}
	Is the space of all \(\mathbb Z_2\)-actions on a compact \(Q\)-manifold a Hilbert manifold? Is it locally connected? What are the answers for compact finite-dimensional manifolds?
\end{openproblem}


\plainproblemsection

\begin{openproblem}
	Let G be a compact Lie group. Let 2G be its hyperspace of non-void closed subsets with the Hausdorff metric induced from a translation invariant metric on G. Let G act on 2G by, say, left translation. What is the structure of the orbit space 2G/G?
\end{openproblem}


\plainproblemsection

\begin{openproblem}
	Let \(\alpha\) be a free action of \(\mathbb Z_p\) on \(S^3\), and let \(E\) be the space of locally flat unknotted simple closed curves in \(S^3\). Must the induced action \(\alpha^*\) on \(E\) have a fixed point? Does \(E\) contain an \(\alpha^*\)-invariant compact ANR, for example a \(Q\)-manifold, that is homotopy equivalent to \(E\)?
\end{openproblem}


\plainproblemsection

\begin{openproblem}
	Let \(M^n\) be a complete Riemannian manifold of nonpositive sectional curvature, let \(P\) be a compact polyhedron, and let \(k>0\). Is the space of \(k\)-Lipschitz maps \(P\to M\) a \(Q\)-manifold? If so, can the dynamics associated with \(M\) be lifted to this space or to related \(Q\)-manifolds?
\end{openproblem}


\plainproblemsection

\begin{openproblem}
	Let \(M\) be a complete Riemannian manifold of nonpositive sectional curvature, and consider \(Q\)-manifold function spaces of Lipschitz maps from compact polyhedra into \(M\). Can these spaces give a simpler proof of the Farrell--Jones vanishing results for \(\operatorname{Wh}(\pi_1M)\), \(\widetilde K_0(\mathbb Z[\pi_1M])\), and \(K_{-i}(\mathbb Z[\pi_1M])\)?
\end{openproblem}

\plainproblemsection

\begin{openproblem}
	Let \(\mathcal P\) and \(\mathcal P'\) be finite presentations
	of the same deficiency for a group \(\pi\), and let \(K_{\mathcal P}\) and
	\(K_{\mathcal P'}\) be their presentation \(2\)-complexes.  If the two
	complexes are simple-homotopy equivalent, must they be related by
	simple-homotopy moves of dimension at most \(3\)?  Equivalently, does the
	generalized Andrews--Curtis implication \((B)\Rightarrow(C)\) in the source
	hold?
\end{openproblem}
% Source: Kirby

\plainproblemsection

\begin{openproblem}
	(Lickorish) Let $K$ be a contractible finite 2-complex.
	
	(A) Zeeman \emph{Conjecture.} $K \times  I$ collapses to a point.
	
	(B) \emph{Conjecture.} $K$ 3-deforms to a point, i.e., there exists a 3-complex $L$ such that $K \swarrow$ $L \smallsetminus  {pt}$.
	
	(C) \emph{Conjecture.} The unique 5-dim regular neighborhood of ${K}^{2}$ in ${\mathbb{R}}^{5}$ is ${B}^{5}$.
	
	$\left( {\mathbf{D}}_{0}\right)$ Andrews-Curtis \emph{Conjecture.} Any presentation of the trivial group can be changed to the trivial presentation by Andrews-Curtis moves.
\end{openproblem}
% Source: Kirby

\plainproblemsection

\begin{openproblem}
	(J. H. C. Whitehead) Is every subcomplex of an aspherical 2-complex aspherical? Assume finiteness if you wish.
\end{openproblem}
% Source: Kirby

\plainproblemsection

\begin{openproblem}
	(A) (Lickorish) \emph{Conjecture.} Any linear subdivision of an $n$ -simplex collapses simplicially.
	
	(B) (Goodrick) \emph{Conjecture.} Any linear subdivision of a star-like $n$ -cell in ${\mathbb{R}}^{n}$ collapses simplicially.
\end{openproblem}
% Source: Kirby

\plainproblemsection

\begin{openproblem}
	(Freedman) Let $G$ be a nontrivial group, $G * \mathbb{Z}$ the free product, and $(G *$ $\mathbb{Z})/r$ the free product with one relation $r$.
	
	Conjecture (Kervaire): $\left( {G * \mathbb{Z}}\right) /r$ is nontrivial.
\end{openproblem}
% Source: Kirby

\plainproblemsection

\begin{openproblem}
	(Habegger) Suppose $M$ and $N$ are ${PL}$ -homeomorphic cubulated n-manifolds. Are they related by the following set of moves: excise $B$ and replace it by ${B}^{\prime }$, where $B$ and ${B}^{\prime }$ are complementary balls (unions of $n$ -cubes which are homeomorphic to ${B}^{n}$ ) in the boundary of the standard $\left( {n + 1}\right)$ -cube?
\end{openproblem}
% Source: Kirby

\plainproblemsection

\begin{openproblem}
	Let
	\[
	\mathcal P=
	\langle a_1,\ldots,a_k\mid R_1,\ldots,R_k\rangle
	\]
	be a presentation whose standard \(2\)-complex is a generalized dunce
	hat, meaning that for every \(1\leq i\leq k\), the word \(R_i\) freely
	reduces to \(a_i\).  Is the Zeeman conjecture true for every such
	generalized dunce hat?
\end{openproblem}
% Source: Kirby


\subsection*{References}
\begin{enumerate}[label={\arabic*.},leftmargin=2.2em]
	\item T. Dobrowolski and J. Mogilski, ``Problems on topological classification of incomplete metric spaces,'' in J. van Mill and G. M. Reed (eds.), \emph{Open Problems in Topology}, North-Holland, 1990.
	\item R. J. Daverman, ``Problems about finite-dimensional manifolds,'' in J. van Mill and G. M. Reed (eds.), \emph{Open Problems in Topology}, North-Holland, 1990.
	\item J. Dydak and J. Segal, ``Shape theory,'' in J. van Mill and G. M. Reed (eds.), \emph{Open Problems in Topology}, North-Holland, 1990.
	\item G. E. Carlsson, ``Algebraic topology,'' in J. van Mill and G. M. Reed (eds.), \emph{Open Problems in Topology}, North-Holland, 1990.
	\item L. H. Kauffman, ``Problems in knot theory,'' in J. van Mill and G. M. Reed (eds.), \emph{Open Problems in Topology}, North-Holland, 1990.
	\item J. E. West, ``Open problems in infinite-dimensional topology,'' in J. van Mill and G. M. Reed (eds.), \emph{Open Problems in Topology}, North-Holland, 1990.
	\item E. Pearl, ``Open problems in topology,'' \emph{Topology and its Applications} \textbf{136} (2004), 37--85, \url{https://doi.org/10.1016/S0166-8641(03)00183-4}.
	\item I. Agol, ``The virtual Haken conjecture,'' \emph{Documenta Mathematica} \textbf{18} (2013), 1045--1087, \url{https://doi.org/10.4171/DM/421}.
	\item J. W. Morgan and G. Tian, \emph{Ricci Flow and the Poincar\'e Conjecture}, Clay Mathematics Monographs, vol. 3, American Mathematical Society and Clay Mathematics Institute, 2007.
	\item J. Dydak and J. J. Walsh, ``Infinite-dimensional compacta having cohomological dimension two: an application of the Sullivan conjecture,'' \emph{Topology} \textbf{32} (1993), 93--104.
	\item J. G. Ratcliffe and S. T. Tschantz, ``Some examples of aspherical 4-manifolds that are homology 4-spheres,'' \emph{Topology} \textbf{44} (2005), 341--350, \url{https://doi.org/10.1016/j.top.2004.10.006}.
	\item R. C. Kirby, ``Problems in low-dimensional topology,'' in W. H. Kazez (ed.), \emph{Geometric Topology}, Part 2, AMS/IP Studies in Advanced Mathematics \textbf{2.2}, American Mathematical Society and International Press, 1997, 35--473.
\end{enumerate}

\clearpage
\section{Topology Arising from Analysis}

\paragraph{Function-Space Topology.}

\plainproblemsection

\begin{openproblem}
	For every infinite space \(X\), is \(C_p(X)\times\mathbb R\) homeomorphic to \(C_p(X)\)? What if \(X\) is compact?
\end{openproblem}


\plainproblemsection

\begin{openproblem}
	Let \(X\) be an infinite metrizable space. Is \(C_p(X)\) homeomorphic to \(C_p(X)\times C_p(X)\)? Does this hold for every compact metrizable \(X\)?
\end{openproblem}


\plainproblemsection

\begin{openproblem}
	For every space \(X\), is \(C_p(X)\times C_p(X)\) a continuous image of \(C_p(X)\)? What if \(X\) is compact?
\end{openproblem}


\plainproblemsection

\begin{openproblem}
	If \(C_p(X)\) is Lindelöf, must \(C_p(X)\times C_p(X)\) be Lindelöf? What if \(X\) is compact?
\end{openproblem}


\plainproblemsection

\begin{openproblem}
	If \(C_p(X)\) is normal, must \(C_p(X)\times C_p(X)\) be normal? What if \(X\) is compact?
\end{openproblem}


\plainproblemsection

\begin{openproblem}
	Let \(X\) be zero-dimensional, let \(D=\{0,1\}\), and suppose \(C_p(X, D)\) is normal. Must it be collectionwise normal? What if \(X\) is compact?
\end{openproblem}


\plainproblemsection

\begin{openproblem}
	Can the Sorgenfrey line be a closed subspace of a Lindelöf space \(C_p(X)\)? Can it be so represented when \(C_p(X)\) is normal?
\end{openproblem}


\plainproblemsection

\begin{openproblem}
	Can the Sorgenfrey line be a closed subspace of a Lindelöf topological vector space?
\end{openproblem}


\plainproblemsection

\begin{openproblem}
	Can every Lindelöf space be a closed subspace of a Lindelöf topological vector space? Can it be a closed subspace of a Lindelöf space \(C_p(X)\)?
\end{openproblem}


\plainproblemsection

\begin{openproblem}
	Can every paracompact space be a closed subspace of a paracompact topological vector space?
\end{openproblem}


\plainproblemsection

\begin{openproblem}
	Is it possible to represent every normal space as a closed subspace of a normal linear topological space?
\end{openproblem}


\plainproblemsection

\begin{openproblem}
	Let X be a collectionwise normal space. Is X homeomorphic to a closed subspace of a collectionwise normal linear topological space (not necessarily locally convex)?
\end{openproblem}


\plainproblemsection

\begin{openproblem}
	Is it true that every space Y of countable tightness is homeomorphic to a subspace (to a closed subspace) of a linear topological space of countable tightness?
\end{openproblem}


\plainproblemsection

\begin{openproblem}
	If \(C_p(X)\) is hereditarily Lindelöf, must \(C_p(X)^n\) be hereditarily Lindelöf for every \(n\in\mathbb N\)?
\end{openproblem}


\plainproblemsection

\begin{openproblem}
	If \(C_p(X)\) has countable spread, must \(C_p(X)^n\) have countable spread for every \(n\in\mathbb N\)?
\end{openproblem}


\plainproblemsection

\begin{openproblem}
	Let X and Y be t-equivalent spaces. Is it true that dim X = dim Y? What if X and Y are compact?
\end{openproblem}


\plainproblemsection

\begin{openproblem}
	Suppose \(C_p(Y)\) is a linear topological factor of \(C_p(X)\), so that \(C_p(X)\cong C_p(Y)\times M\) for some linear topological space \(M\). Must \(\dim Y\leq\dim X\)?
\end{openproblem}


\plainproblemsection

\begin{openproblem}
	Is it true that weakly topologically equivalent metrizable spaces are always l-equivalent? Is this true at least for compact metrizable spaces?
\end{openproblem}


\plainproblemsection

\begin{openproblem}
	Is every infinite compact space \(t\)-equivalent to a compact space containing a nontrivial convergent sequence?
\end{openproblem}


\plainproblemsection

\begin{openproblem}
	Is every nonempty compact space \(t\)-equivalent to a compact space containing a point of countable character?
\end{openproblem}


\plainproblemsection

\begin{openproblem}
	Let \(Y\) be a compact subspace of an infinite-dimensional space \(C_p(X)\). Must the tightness of \(Y\) be at most the Lindelöf degree of \(C_p(X)\)?
\end{openproblem}


\plainproblemsection

\begin{openproblem}
	Let \(X\) be compact, and suppose \(\aleph_1\) is a caliber of \(C_p(X)\). Must \(X\) be metrizable?
\end{openproblem}


\plainproblemsection

\begin{openproblem}
	Suppose \(X\) and \(Y\) are \(l\)-equivalent and \(X\) is a compact Fréchet--Urysohn space. Must \(Y\) also be Fréchet--Urysohn? What if \(X\) and \(Y\) are \(t\)-equivalent and \(Y\) is compact?
\end{openproblem}


\plainproblemsection

\begin{openproblem}
	Let \(\tau>\aleph_0\), put \(D^\tau=\{0,1\}^\tau\), and let \(I^\tau=[0,1]^\tau\). Is \(D^\tau\) \(t\)-embedded in \(I^\tau\)?
\end{openproblem}


\plainproblemsection

\begin{openproblem}
	Is it true that every compact space is l-embedded (is t-embedded) into some topologically homogeneous compact space?
\end{openproblem}


\plainproblemsection

\begin{openproblem}
	Let \(S=\{0\}\cup\{1/n: n\in\mathbb N^+\}\). For every compact metrizable space \(X\), are \(X\) and \(X\times S\) \(l\)-equivalent? Are they \(t\)-equivalent?
\end{openproblem}


\plainproblemsection

\begin{openproblem}
	Let \(X\) be compact. If \(C_p(C_p(X))=C_{p,2}(X)\) is Lindelöf, must \(C_p(X)\) be Lindelöf?
\end{openproblem}


\paragraph{Topology Arising from Analysis.}

\plainproblemsection

\begin{openproblem}
	For \(r\in(0,1)\), let \(\mu(r)\) be the product measure on \(\{0,1\}^{\mathbb N}\) whose one-coordinate weights are \(1-r\) and \(r\). Are \(\mu(r)\) and \(\mu(s)\) homeomorphic if and only if there exist positive integers \(n, m\) and integers \(a_i, b_j\) satisfying \[0\leq a_i\leq {n\choose i},\qquad 0\leq b_j\leq {m\choose j},\]\[r=\sum_{i=0}^{n}a_i s^i(1-s)^{n-i},\qquad s=\sum_{j=0}^{m}b_j r^j(1-r)^{m-j}?\]
\end{openproblem}


\plainproblemsection

\begin{openproblem}
	For \(t\in(0,1)\), let \(\mu(t)\) be the Bernoulli product measure on \(\{0,1\}^{\mathbb N}\) with one-coordinate weights \(1-t\) and \(t\). Let \(r\in(0,1)\) satisfy \(r^3+r^2-1=0\). If \(E(r)\) is the set of parameters \(t\) for which \(\mu(t)\) is homeomorphic to \(\mu(r)\), and \(F(r)\) is the corresponding mutual-continuous-image class, is either set exactly \(\{r,1-r, r^2,1-r^2\}\)?
\end{openproblem}


\plainproblemsection

\begin{openproblem}
	For \(r\in(0,1)\), are there only finitely many \(s\in(0,1)\) such that the Bernoulli product measures \(\mu(r)\) and \(\mu(s)\) on \(\{0,1\}^{\mathbb N}\) are homeomorphic?
\end{openproblem}


\plainproblemsection

\begin{openproblem}
	Does there exist a Borel set \(M\subset\mathbb R^2\) meeting every straight line in exactly two points? Can such an \(M\) be a \(G_\delta\)-set?
\end{openproblem}


\plainproblemsection

\begin{openproblem}
	If a subset \(M\subset\mathbb R^2\) meets every straight line in exactly two points, must \(M\) be zero-dimensional?
\end{openproblem}


\plainproblemsection

\begin{openproblem}
	Let \(B\subset[0,1]^2\) be Borel, and suppose every horizontal and vertical section of \(B\) is a dense \(G_\delta\)-set. Can \(B\) be filled by pairwise disjoint graphs of Borel isomorphisms of \([0,1]\) onto itself?
\end{openproblem}


\plainproblemsection

\begin{openproblem}
	Let \(a\geq2\) be an integer, let \(A:\mathbb R\to\mathbb R\) be continuous and \(1\)-periodic with \(\lVert A\rVert_\infty=1\), and let \(B>1\), \(b=B^{-1}\). For \[f(x)=\sum_{p=0}^{\infty}b^pA(a^p x),\]determine the Hausdorff dimension \(\gamma\) of the graph of \(f\). More precisely, determine, if possible, a slowly varying function \(L\) such that for \(h(t)=t^\gamma L(t)\) the graph has positive finite \(h\)-Hausdorff measure.
\end{openproblem}


\plainproblemsection

\begin{openproblem}
	Let \(a\geq2\), let \(A:\mathbb R\to\mathbb R\) be continuous and \(1\)-periodic with \(\lVert A\rVert_\infty=1\), and let \(0<B<1\). Define \[T(e^{2\pi i x}, y)=\bigl(e^{2\pi i ax}, B(y-A(x))\bigr)\]and \[K=S^1\times\left[-\frac{B}{1-B},\frac{B}{1-B}\right],\qquad M=\bigcap_{n\geq0}T^n(K).\]Is \(M\) a Sierpiński curve? In particular, is this true when \(A\) is the tent map?
\end{openproblem}


\plainproblemsection

\begin{openproblem}
	Let \(a\geq2\), let \(A:\mathbb R\to\mathbb R\) be a nonconstant continuous \(1\)-periodic function, and let \(0<B<1\). Is the unique continuous solution of \[f(x)=A(x)+B^{-1}f(ax)\]necessarily nonperiodic?
\end{openproblem}


\subsection*{References}
\begin{enumerate}[label={\arabic*.},leftmargin=2.2em]
	\item A. V. Arkhangel'skii, ``Problems in \(C_p\)-theory,'' in J. van Mill and G. M. Reed (eds.), \emph{Open Problems in Topology}, North-Holland, 1990.
	\item R. D. Mauldin, ``Topology arising from analysis,'' in J. van Mill and G. M. Reed (eds.), \emph{Open Problems in Topology}, North-Holland, 1990.
	\item E. Pearl, ``Open problems in topology,'' \emph{Topology and its Applications} \textbf{136} (2004), 37--85, \url{https://doi.org/10.1016/S0166-8641(03)00183-4}.
\end{enumerate}

\clearpage
\section{Topological Dynamics}

\paragraph{Continuum Theory and Topological Dynamics.}

\plainproblemsection

\begin{openproblem}
	Let \(f: M\to M\) be a diffeomorphism of a \(2\)-manifold, and let \(p\) be a hyperbolic saddle fixed point. Suppose one branch \(W_u^+(p)\) of the unstable manifold meets the stable manifold \(W^s(p)\) at a point \(q\), and the intersection at \(q\) is not topologically transverse. Must \(\overline{W_u^+(p)}\) be nonchainable?
\end{openproblem}


\plainproblemsection

\begin{openproblem}
	Let \(f: I\to I\) be continuous. Under what conditions does the inverse-limit space \(\varprojlim(I, f)\) admit a planar embedding for which the shift homeomorphism extends to a diffeomorphism of the plane?
\end{openproblem}


\plainproblemsection

\begin{openproblem}
	Let \(M\subset\mathbb R^2\) be a nonseparating continuum. Must every continuous map \(f: M\to M\) have a periodic point?
\end{openproblem}


\plainproblemsection

\begin{openproblem}
	A homeomorphism \(f: X\to X\) of a compact metric space is expansive if there is \(\delta>0\) such that, for all distinct \(x, y\in X\), one has \(d(f^n(x), f^n(y))\geq\delta\) for some \(n\in\mathbb Z\). Characterize the planar continua that admit expansive homeomorphisms.
\end{openproblem}


\plainproblemsection

\begin{openproblem}
	Is the Mandelbrot set locally connected?
\end{openproblem}


\paragraph{One- and Two-Dimensional Dynamics.}

\plainproblemsection

\begin{openproblem}
	Let \(f: D\to D\) be a homeomorphism onto its image, and let \(O\) be a periodic orbit. Characterize, in terms of the induced map \(f_*:\pi_1(D\setminus O)\to\pi_1(D\setminus O)\), those orbits \(O\) whose existence forces periodic orbits of every period.
\end{openproblem}


\plainproblemsection

\begin{openproblem}
	Can the periodic-orbit forcing theory for a finite invariant orbit of a disc homeomorphism be extended to infinite invariant orbits?
\end{openproblem}


\plainproblemsection

\begin{openproblem}
	For a real-analytic map \(f: D\to D\), is there finite orbit data, analogous to the orbits of the finitely many turning points of an analytic interval map, that determines the global dynamics?
\end{openproblem}


\plainproblemsection

\begin{openproblem}
	Which periodic-orbit structures can be realized by \(C^\infty\) diffeomorphisms of the disc? In particular, does there exist a zero-entropy \(C^\infty\) diffeomorphism having an orbit of period \(3^n\) for every \(n\in\mathbb N\), with each orbit of period \(3^{n+1}\) circling around the orbit of period \(3^n\)?
\end{openproblem}


\plainproblemsection

\begin{openproblem}
	Can the nested periodic-orbit structure consisting of an orbit of period \(3^n\) for every \(n\), with each successive orbit circling around the preceding one, be realized by a real-analytic diffeomorphism of the disc?
\end{openproblem}


\plainproblemsection

\begin{openproblem}
	For almost one-dimensional real-analytic disc diffeomorphisms, a local codimension-one stable manifold of the renormalization operator is known. Can this stable manifold be defined globally by a two-dimensional topological analysis of the renormalization operator?
\end{openproblem}


\plainproblemsection

\begin{openproblem}
	In the space of real-analytic diffeomorphisms of the disc, where ``chaotic'' means positive topological entropy, characterize the boundary of the chaotic maps. Does this boundary have the structure of a stratified manifold?
\end{openproblem}


\plainproblemsection

\begin{openproblem}
	For real-analytic disc diffeomorphisms with infinitely many periodic attractors of Newhouse type, determine geometric properties of their basins of attraction.
\end{openproblem}


\plainproblemsection

\begin{openproblem}
	Does there exist an open set of real-analytic disc diffeomorphisms such that every map in the set has infinitely many periodic attractors?
\end{openproblem}


\plainproblemsection

\begin{openproblem}
	Does there exist a one-parameter family of disc homeomorphisms or diffeomorphisms such that every nearby family contains a parameter whose map has infinitely many periodic attractors?
\end{openproblem}


\plainproblemsection

\begin{openproblem}
	Let \(f:\mathbb R^2\to\mathbb R^2\) be a \(C^1\) orientation-preserving diffeomorphism, and let \(O\) be a periodic orbit. A linked fixed point \(P\) exists. Can \(P\) always be chosen so that its linking number with \(O\) is nonzero?
\end{openproblem}


\plainproblemsection

\begin{openproblem}
	Does there exist a real-analytic diffeomorphism \(f: A\to A\) of the closed annulus, without periodic points, such that for some \(x\in\operatorname{int}A\), the \(\omega\)-limit set \(\omega(x)\) meets both boundary components of \(A\)?
\end{openproblem}


\subsection*{References}
\begin{enumerate}[label={\arabic*.},leftmargin=2.2em]
	\item M. Barge and J. Kennedy, ``Continuum theory and topological dynamics,'' in J. van Mill and G. M. Reed (eds.), \emph{Open Problems in Topology}, North-Holland, 1990.
	\item S. van Strien, ``One-dimensional versus two-dimensional dynamics,'' in J. van Mill and G. M. Reed (eds.), \emph{Open Problems in Topology}, North-Holland, 1990.
	\item E. Pearl, ``Open problems in topology,'' \emph{Topology and its Applications} \textbf{136} (2004), 37--85, \url{https://doi.org/10.1016/S0166-8641(03)00183-4}.
	\item M. Benedicks and L. Carleson, ``The dynamics of the H\'enon map,'' \emph{Annals of Mathematics} \textbf{133} (1991), 73--169.
	\item P. Berger, ``Analytic pseudo-rotations,'' accepted for publication in \emph{Annals of Mathematics}, 2025, \url{https://annals.math.princeton.edu/articles/22134}.
\end{enumerate}

\chapter{Dynamical Systems}

\clearpage
\section{Complex Dynamics}

\plainproblemsection

\begin{openproblem}
	Let
	\[
	f\colon \widehat{\mathbb{C}}\longrightarrow\widehat{\mathbb{C}}
	\]
	be a rational map of degree \(d>1\).  The map \(f\) is called
	\emph{hyperbolic} if its critical points tend under iteration to attracting
	periodic cycles.
	
	The density-of-hyperbolicity conjecture asserts that the set of hyperbolic
	rational maps is open and dense in the space \(\operatorname{Rat}_d\) of all
	rational maps of degree \(d\).
\end{openproblem}

\plainproblemsection

\begin{openproblem}
	Every quadratic polynomial is conjugate to a unique member of the family
	\[
	f_c(z)=z^2+c,\qquad c\in\mathbb{C}.
	\]
	The quadratic density-of-hyperbolicity conjecture asserts that the set of
	parameters \(c\) for which \(f_c\) is hyperbolic is an open dense subset of
	\(\mathbb{C}\).
	
	Equivalently, let
	\[
	M=\bigl\{c\in\mathbb{C}:f_c^n(0)\not\longrightarrow\infty\bigr\}
	\]
	be the Mandelbrot set.  The conjecture states that every component of
	\(\operatorname{int}M\) is hyperbolic.
\end{openproblem}

\plainproblemsection

\begin{openproblem}
	The no-invariant-line-fields conjecture asserts that a rational map \(f\)
	carries no invariant line field on its Julia set, except when \(f\) is double
	covered by an integral torus endomorphism.
	
	Here a line field assigns a real line through the origin in the tangent space
	at each point of a positive-measure subset \(E\) of the Julia set, with
	measurably varying slope.  It is invariant if \(f^{-1}(E)=E\) and \(f'\)
	maps the line at \(z\) to the line at \(f(z)\).
\end{openproblem}

\plainproblemsection

\begin{openproblem}
	The quadratic no-invariant-line-fields conjecture asserts that a quadratic
	polynomial carries no invariant line field on its Julia set.
	
	For a polynomial \(f\), its Julia set is the boundary of the set of points
	which tend to infinity under iteration.  A line field on the Julia set is a
	measurable assignment of a real tangent line on a positive-measure subset
	\(E\); it is invariant if \(f^{-1}(E)=E\) and \(f'\) maps the assigned line
	at \(z\) to the assigned line at \(f(z)\).
\end{openproblem}

\subsection*{References}

\begin{enumerate}[label={[\arabic*]},leftmargin=*,itemsep=0.4em]
	\item C. T. McMullen, \emph{Complex Dynamics and Renormalization}.
	\item R. Ma\~n\'e, P. Sad, and D. Sullivan, ``On the dynamics of rational
	maps,'' \emph{Annales scientifiques de l'\'Ecole Normale Sup\'erieure}
	\textbf{16} (1983), 193--217.
	\item C. McMullen and D. Sullivan, ``Quasiconformal homeomorphisms and
	dynamics III: The Teichm\"uller space of a holomorphic dynamical system,''
	preprint.
\end{enumerate}

\clearpage
\section{Ergodic Theory and Geometric Dynamics}

\plainproblemsection

\begin{openproblem}
	(R. Spatzier) Let \(G\) be a connected semisimple Lie group, let
	\(\Gamma\) be an irreducible lattice in \(G\), and let \(A\) be a closed
	subgroup of a split Cartan subgroup of \(G\), with \(\dim A>1\).
	
	Is every \(C^1\)-small perturbation of the action of \(A\) on \(G/\Gamma\)
	\(C^\infty\)-conjugate to the action of \(A\) defined by a continuous
	homomorphism from \(A\) to the centralizer of \(A\) in \(G\)?
\end{openproblem}

\plainproblemsection

\begin{openproblem}
	In the preceding setting, let \(D\) be the split Cartan subgroup and suppose
	that \(A\) is not contained in a wall of a Weyl chamber of \(D\).
	
	Is every \(C^1\)-small perturbation of the action of \(A\) on \(G/\Gamma\)
	\(C^\infty\)-conjugate to the action of \(A\) defined by a continuous
	homomorphism from \(A\) to \(D\)?
\end{openproblem}

\plainproblemsection

\begin{openproblem}
	The folklore classification conjecture for Anosov diffeomorphisms asserts
	that every Anosov diffeomorphism is topologically conjugate to a hyperbolic
	automorphism of an infranilmanifold.
\end{openproblem}

\plainproblemsection

\begin{openproblem}
	(R. Spatzier) Every irreducible Anosov
	\(\mathbb{Z}^k\)- or \(\mathbb{R}^k\)-action, \(k\geq 2\), is conjectured to
	be \(C^\infty\)-conjugate to an algebraic action.
\end{openproblem}

\plainproblemsection

\begin{openproblem}
	Every action of a connected semisimple Lie group of higher rank, or of a
	lattice in such a group, that contains an element acting nontrivially and
	uniformly partially hyperbolically is conjectured to be
	\(C^\infty\)-conjugate to an algebraic action.  Here higher rank means that
	all simple factors have real rank at least \(2\).
\end{openproblem}

\plainproblemsection

\begin{openproblem}
	(D. Fisher) Do measure-rigidity results for actions of higher-rank abelian
	groups give obstructions to smooth or continuous irreducible actions of a
	higher-rank abelian group on a compact manifold?
\end{openproblem}

\plainproblemsection

\begin{openproblem}
	(H. Furstenberg) Let \(m\) and \(n\) be multiplicatively independent
	positive integers.  Are the Lebesgue measure and the measures with finite
	support the only probability Borel measures on \(\mathbb{R}/\mathbb{Z}\)
	that are ergodic under the semigroup generated by multiplication by \(m\)
	and \(n\)?
\end{openproblem}

\plainproblemsection

\begin{openproblem}
	(A. Katok and R. Spatzier) Let \(M\) be a biquotient of a connected Lie
	group \(G\).  Given an algebraic Anosov action of
	\(\mathbb{R}^k\) or \(\mathbb{Z}^k\), \(k\geq 2\), on \(M\), and an ergodic
	probability Borel measure \(\mu\) on \(M\), is one of the following always
	true?
	\begin{enumerate}[label=\textup{(\roman*)}]
		\item The measure \(\mu\) is a finite sum of measures constructed from Haar
		measures supported on closed subgroups of \(G\).
		\item The support of \(\mu\) fibers equivariantly over a manifold \(N\) such
		that the action on \(N\) reduces to a rank-one action.
	\end{enumerate}
\end{openproblem}

\plainproblemsection

\begin{openproblem}
	(E. Lindenstrauss) Let
	\[
	\mathbb{A}_0=\prod_{v\in S}'\mathbb{Q}_v
	\]
	be a restricted adelic product over a subset \(S\) of the places of
	\(\mathbb{Q}\), let \(D\) be the diagonal subgroup of
	\(\operatorname{SL}(2)\), and let \(\Gamma\) be an irreducible lattice in
	\(\operatorname{SL}(2,\mathbb{A}_0)\).
	
	Classify the finite ergodic \(D(\mathbb{A}_0)\)-invariant measures on
	\(\operatorname{SL}(2,\mathbb{A}_0)/\Gamma\).
\end{openproblem}

\plainproblemsection

\begin{openproblem}
	(L. Silberman) Extend the measure-rigidity results for unipotent flows to
	the adelic setting.
\end{openproblem}

\plainproblemsection

\begin{openproblem}
	(E. Lindenstrauss) Let \(H\) be a connected semisimple Lie subgroup of a
	Lie group \(G\), and let \(P\) be a parabolic subgroup of \(H\).  Suppose
	that \(H\) acts on a space \(X\) preserving a geometric structure.
	
	Under what conditions on \(X\) is every finite \(P\)-invariant measure
	\(H\)-invariant?  Equivalently, which \(H\)-actions are stiff?
\end{openproblem}

\plainproblemsection

\begin{openproblem}
	(F. Ledrappier and O. Sarig) Let \(G\) be a noncompact semisimple Lie group
	of rank one, let \(\Gamma\) be a discrete subgroup of \(G\), and let \(U\)
	be a horospherical subgroup of \(G\).
	
	Classify the \(U\)-ergodic Radon measures on \(G/\Gamma\).  In particular,
	are they all either carried by closed \(U\)-orbits or obtained from the
	harmonic-function construction?
\end{openproblem}

\plainproblemsection

\begin{openproblem}
	(F. Ledrappier and O. Sarig) A \(u_t\)-invariant measure \(\mu\) is called
	\emph{squashable} if the centralizer of \(u_t\) contains an invertible
	nonsingular transformation that does not preserve \(\mu\).
	
	Let \(\Gamma\) be a normal conilpotent subgroup of a uniform lattice in
	\(\operatorname{SL}(2,\mathbb{R})\).  Is the Haar measure nonsquashable?
\end{openproblem}

\plainproblemsection

\begin{openproblem}
	(A. Furman) Consider either of the following actions of a group \(\Gamma\):
	\begin{enumerate}[label=\textup{(\arabic*)}]
		\item \(\Gamma\) is a large subgroup, for example a Zariski-dense subgroup,
		of the automorphism group of a finite-volume nilmanifold;
		\item \(\Gamma\) is a large subgroup of a Lie group \(G\) acting by
		translations on \(G/\Lambda\), where \(\Lambda\) is a lattice in \(G\).
	\end{enumerate}
	Is every ergodic \(\Gamma\)-invariant probability measure either supported
	on a finite \(\Gamma\)-orbit or equal to the Haar measure?
\end{openproblem}

\plainproblemsection

\begin{openproblem}
	(F. Ledrappier) Let \(\varphi_t\colon M\to M\) be an Anosov flow on a
	compact Riemannian manifold, with
	\[
	TM=E^0\oplus E^s\oplus E^u.
	\]
	Suppose that \(E^s\) has a continuous invariant splitting
	\(E^s=E^s_+\oplus E^s_-\), where \(E^s_+\) contracts faster than
	\(E^s_-\), and let \(W^s_+\) be the fast stable foliation obtained by
	integrating \(E^s_+\).
	
	Classify the invariant ergodic measures for \(W^s_+\).  The same question
	may also be asked for Anosov diffeomorphisms.
\end{openproblem}

\plainproblemsection

\begin{openproblem}
	(G. Margulis) Let \(G\) be a Lie group, let \(\Gamma\) be a lattice in
	\(G\), and let \(U=\{u(t)\}\subset G\) be a one-parameter
	\(\operatorname{Ad}\)-unipotent subgroup.  Suppose that \(Ux\) is dense in
	\(G/\Gamma\).
	
	Prove equidistribution of \(\{u(t_n)x\}\) when
	\begin{enumerate}[label=\textup{(\arabic*)}]
		\item \(t_n=\lfloor n^\alpha\rfloor\) for \(\alpha>1\);
		\item \(t_n=\lfloor P(n)\rfloor\), where \(P\) is a polynomial;
		\item \(t_n\) is the \(n\)-th prime number.
	\end{enumerate}
\end{openproblem}

\plainproblemsection

\begin{openproblem}
	(G. Margulis) Let \(V\) be a connected \(\operatorname{Ad}\)-unipotent
	subgroup of a Lie group \(G\), let \(\Gamma\) be a lattice in \(G\), and
	suppose that \(Vx\) is dense in \(G/\Gamma\).  For a suitable sequence of
	sets \(A_n\subset V\), prove, with an effective error term, that
	\[
	\lim_{n\to\infty}
	\frac{1}{\operatorname{Vol}(A_n)}
	\int_{A_n} f(vx)\,\mathrm{d}v
	=
	\int_{G/\Gamma} f\,\mathrm{d}\mu
	\]
	for every \(f\in C_c(G/\Gamma)\), where \(\mathrm{d}v\) is Haar measure on
	\(V\) and \(\mu\) is the probability Haar measure on \(G/\Gamma\).
\end{openproblem}

\plainproblemsection

\begin{openproblem}
	(N. Shah) Let \(L\) be a Lie group, \(G\) a closed subgroup of \(L\), and
	\(\Lambda\) a lattice in \(L\).  For a semisimple element \(a\in G\), let
	\[
	U=\{g\in G:a^{-n}ga^n\longrightarrow e\ \text{as }n\to\infty\}
	\]
	be its expanding horospherical subgroup.  Suppose that \(Gx_0\) is dense
	in \(L/\Lambda\).  Given an analytic curve
	\(\gamma\colon[0,1]\to U\), let \(\nu\) be the image of Lebesgue measure
	under \(t\mapsto\gamma(t)x_0\), and let \(\lambda\) be probability Haar
	measure on \(L/\Lambda\).
	
	Under what conditions on \(\gamma\) does
	\[
	(a^n)_*\nu\longrightarrow\lambda
	\qquad\text{as }n\to\infty?
	\]
\end{openproblem}

\plainproblemsection

\begin{openproblem}
	(N. Shah) In the setting above, suppose that
	\(G=\operatorname{SO}(n,1)\subset L\).  If the analytic curve
	\(\gamma\colon[0,1]\to U\) is not contained in a proper affine subspace or
	an \((n-2)\)-dimensional sphere in \(U\), then it is conjectured that
	\[
	(a^n)_*\nu\longrightarrow\lambda
	\qquad\text{as }n\to\infty
	\]
	for every such ambient Lie group \(L\).
\end{openproblem}

\plainproblemsection

\begin{openproblem}
	(N. Shah) Let \(G\) act linearly on a normed real vector space \(V\), let
	\(p\in V\) have finite stabilizer, and put
	\[
	R_T=\{g\in G:\lVert gp\rVert<T\}.
	\]
	Let \(\Gamma\) be a lattice in \(G\), and let \(\mu_T\) be the projection
	to \(G/\Gamma\) of normalized Haar measure on \(R_T\).
	
	Determine the limiting distribution of \(\mu_T\) as \(T\to\infty\).
\end{openproblem}

\plainproblemsection

\begin{openproblem}
	(Z. Rudnick and P. Sarnak) For an interval
	\([a,b]\subset[0,1]\), define
	\[
	R_2([a,b],N,\alpha)
	=
	\frac{1}{N}
	\#\left\{
	\begin{array}{c}
		1\leq i\neq j\leq N:
		\\[2pt]
		\alpha(i^2-j^2)\in N^{-1}[a,b]\pmod 1
	\end{array}
	\right\}.
	\]
	If \(\alpha\in\mathbb{R}\) is badly approximable by rationals, prove that
	\[
	R_2([a,b],N,\alpha)\longrightarrow b-a
	\qquad\text{as }N\to\infty.
	\]
\end{openproblem}

\plainproblemsection

\begin{openproblem}
	(J. Marklof and A. Str\"ombergsson) Let
	\(X=\Gamma\backslash\mathbb{H}^2\) be a finite-volume hyperbolic surface
	whose cusp has been normalized by
	\[
	\Gamma\cap\{z\mapsto z+a:a\in\mathbb{R}\}
	=
	\{z\mapsto z+a:a\in\mathbb{Z}\}.
	\]
	Let \(u_y(t)\), \(0\leq t\leq1\), be the corresponding closed horocycle,
	let \(f\) be continuous on \(T^1(X)\) with the prescribed growth condition
	at the cusps, and let \(\lambda\) be Liouville measure.  If \(\alpha\) is
	badly approximable and \(0<c_1<c_2\), prove that
	\[
	\frac1M\sum_{m=1}^{M}f\bigl(u_y(\alpha m)\bigr)
	\longrightarrow
	\int_{T^1(X)}f\,\mathrm{d}\lambda
	\]
	uniformly as \(M\to\infty\) and
	\(c_1M^{-2}\leq y\leq c_2M^{-2}\).
\end{openproblem}

\plainproblemsection

\begin{openproblem}
	(H. Oh) Let \(X\) be a finite-volume Riemannian locally symmetric space
	of noncompact type.  Determine the asymptotic number of compact flats in
	\(X\) having volume less than \(T\) as \(T\to\infty\), including the
	optimal rate of convergence in the cases where the leading asymptotic is
	known.
\end{openproblem}

\plainproblemsection

\begin{openproblem}
	(H. Oh) Let \(X=\Gamma\backslash G/K\) be a finite-volume locally
	symmetric space of higher rank, let \(A\) be a Cartan subgroup, and call
	the compact \(A\)-orbits in \(\Gamma\backslash G\) flats.  For a flat \(F\),
	let \(\mu_F\) be its Lebesgue measure and \(\overline{\mu}_F\) its
	normalized Lebesgue measure.  Define
	\[
	\nu_T=
	\frac{\displaystyle\sum_{\operatorname{Vol}(F)<T}\mu_F}
	{\displaystyle\sum_{\operatorname{Vol}(F)<T}\operatorname{Vol}(F)},
	\qquad
	\overline{\nu}_T=
	\frac{\displaystyle\sum_{\operatorname{Vol}(F)<T}\overline{\mu}_F}
	{\#\{F:\operatorname{Vol}(F)<T\}}.
	\]
	Do \(\nu_T\) and \(\overline{\nu}_T\) converge to normalized Haar measure
	on \(\Gamma\backslash G\)?
\end{openproblem}

\plainproblemsection

\begin{openproblem}
	(B. Weiss) Let \(G\) be a semisimple real algebraic group, let \(\Gamma\)
	be a noncocompact arithmetic lattice in \(G\), and let \(D\) be a closed
	subgroup of a maximal \(\mathbb{R}\)-split torus \(A\).  A divergent orbit
	\(Dx\) is called \emph{obvious} if \(D\) is a union of open subsemigroups
	\(D_1,\ldots,D_\ell\) such that, for each \(i\), there are a
	\(\mathbb{Q}\)-defined representation
	\(\rho_i\colon G\to\operatorname{GL}(V_i)\) and \(v_i\in V_i\) for which
	\(\rho_i(dx)v_i\to0\) as \(d\in D_i\) tends to infinity.
	
	The unresolved parts of the conjecture are:
	\begin{enumerate}[label=\textup{(\arabic*)}]
		\item if \(\dim D=\operatorname{rank}_{\mathbb{Q}}G\), then every divergent
		\(D\)-orbit is obvious;
		\item if \(\dim D<\operatorname{rank}_{\mathbb{Q}}G\), then non-obvious
		divergent \(D\)-orbits exist.
	\end{enumerate}
\end{openproblem}

\plainproblemsection

\begin{openproblem}
	(B. Weiss) Construct cones \(D\) in a maximal
	\(\mathbb{R}\)-split torus \(A\), with \(\dim D\geq3\), that have no
	non-obvious divergent trajectories.
\end{openproblem}

\plainproblemsection

\begin{openproblem}
	(B. Weiss) Let \(A\) be a group with \(\dim A\geq2\), let \(A_T\) denote
	the ball of radius \(T\) in \(A\), and let \(\lambda\) be Haar measure on
	\(A\).  Suppose that \(x\in G/\Gamma\) and that \(Dx\) is not divergent
	for every one-parameter subgroup \(D\) of \(A\).
	
	Must there exist a compact set \(K\subset G/\Gamma\) such that
	\[
	\limsup_{T\to\infty}
	\frac{\lambda\{a\in A_T:ax\in K\}}{\lambda(A_T)}>0?
	\]
\end{openproblem}

\plainproblemsection

\begin{openproblem}
	(S. Katok) Construct analogues of both the geometric coding and the
	arithmetic coding for the Weyl-chamber flow, namely the action of the
	diagonal group on
	\[
	\operatorname{SL}(n,\mathbb{R})/\operatorname{SL}(n,\mathbb{Z}).
	\]
\end{openproblem}

\plainproblemsection

\begin{openproblem}
	(A. Katok) Construct periodic orbits for triangular billiards.  In
	particular, determine whether every acute triangle has periodic orbits
	other than the standard one and whether every obtuse triangle has at least
	one periodic orbit.
\end{openproblem}

\plainproblemsection

\begin{openproblem}
	(A. Katok) Find a triangle whose angles are Diophantine modulo \(2\pi\)
	and whose billiard flow is ergodic with respect to Liouville measure.
\end{openproblem}

\plainproblemsection

\begin{openproblem}
	(A. Katok) Does there exist a weakly mixing polygonal billiard?  More
	specifically, are the billiards constructed by Vorobets weakly mixing?
\end{openproblem}

\plainproblemsection

\begin{openproblem}
	(A. Katok) Does the set of ergodic billiard tables have positive measure?
\end{openproblem}

\plainproblemsection

\begin{openproblem}
	(A. Katok) Classify the ergodic invariant probability measures and the
	closed invariant subsets for polygonal billiards.
\end{openproblem}

\plainproblemsection

\begin{openproblem}
	(G. Margulis and A. Selberg) Let
	\[
	G=\operatorname{SL}(2,\mathbb{R})^k,\qquad k\geq2,
	\]
	and let \(U^+\) and \(U^-\) be the upper and lower unipotent subgroups of
	\(G\).  Let \(\Gamma^+\subset U^+\) and \(\Gamma^-\subset U^-\) be
	lattices whose projections to every factor of \(G\) are injective.
	
	If the group \(\langle\Gamma^+,\Gamma^-\rangle\) is discrete, must it be
	an arithmetic lattice in \(G\)?
\end{openproblem}

\plainproblemsection

\begin{openproblem}
	(D. Kleinbock and G. Margulis) An \(m\times n\) real matrix \(A\) is
	called very well approximable if, for some \(\varepsilon>0\), there are
	infinitely many \(\mathbf q\in\mathbb{Z}^n\) and
	\(\mathbf p\in\mathbb{Z}^m\) such that
	\[
	\lVert A\mathbf q-\mathbf p\rVert^m
	<
	\lVert\mathbf q\rVert^{-n-\varepsilon}.
	\]
	Find reasonable and checkable conditions on a smooth map
	\[
	f\colon U\subset\mathbb{R}^k\longrightarrow
	M_{m\times n}(\mathbb{R})
	\]
	that generalize the nondegeneracy condition for vector-valued maps and
	imply that almost every point of \(f(U)\) is not very well approximable.
\end{openproblem}

\plainproblemsection

\begin{openproblem}
	(D. Kleinbock) For \(v>0\), call
	\(\mathbf y=(y_1,\ldots,y_n)\in\mathbb{R}^n\)
	\emph{\(v\)-approximable} if there are infinitely many
	\(\mathbf q=(q_1,\ldots,q_n)\in\mathbb{Z}^n\) and \(p\in\mathbb{Z}\)
	such that
	\[
	\bigl|y_1q_1+\cdots+y_nq_n-p\bigr|<\lVert\mathbf q\rVert^{-v}.
	\]
	Use dynamics to estimate the Hausdorff dimension of the
	\(v\)-approximable vectors in a nondegenerate submanifold of
	\(\mathbb{R}^n\).
\end{openproblem}

\plainproblemsection

\begin{openproblem}
	(A. Pollington) Write
	\(\langle x\rangle=\operatorname{dist}(x,\mathbb{Z})\).  Do there exist
	\(\alpha,\beta\in\mathbb{R}\) such that, for every
	\(\varepsilon>0\),
	\[
	\liminf_{q\to\infty}
	q(\log q)^{2-\varepsilon}
	\langle q\alpha\rangle\langle q\beta\rangle>0?
	\]
\end{openproblem}

\plainproblemsection

\begin{openproblem}
	Littlewood's conjecture asserts that, for every
	\(\alpha,\beta\in\mathbb{R}\),
	\[
	\liminf_{q\to\infty}
	q\langle q\alpha\rangle\langle q\beta\rangle=0,
	\qquad
	\langle x\rangle=\operatorname{dist}(x,\mathbb{Z}).
	\]
\end{openproblem}

\plainproblemsection

\begin{openproblem}
	(G. Margulis) Let \(A\) be the group of all diagonal matrices in
	\(\operatorname{SL}(3,\mathbb{R})\).  Is every bounded \(A\)-orbit in
	\[
	\operatorname{SL}(3,\mathbb{R})/\operatorname{SL}(3,\mathbb{Z})
	\]
	closed?
\end{openproblem}

\plainproblemsection

\begin{openproblem}
	(G. Margulis) Let \(A\) be the diagonal subgroup of
	\(\operatorname{SL}(3,\mathbb{R})\).  Does every compact subset
	\[
	K\subset
	\operatorname{SL}(3,\mathbb{R})/\operatorname{SL}(3,\mathbb{Z})
	\]
	contain only finitely many closed \(A\)-orbits?
\end{openproblem}

\plainproblemsection

\begin{openproblem}
	(G. Margulis) For
	\[
	F(\mathbf x)=
	\prod_{i=1}^{3}\left(\sum_{j=1}^{3}a_{ij}x_j\right),
	\qquad a_{ij}\in\mathbb{R},
	\]
	put
	\[
	\Delta(F)=\det(a_{ij}),
	\qquad
	m(F)=
	\inf_{\mathbf x\in\mathbb{Z}^3\setminus\{\mathbf0\}}
	\left|\frac{F(\mathbf x)}{\Delta(F)}\right|.
	\]
	Show that, for every \(\varepsilon>0\), the set
	\[
	[\varepsilon,\infty)\cap\{m(F):F\text{ as above}\}
	\]
	is finite.
\end{openproblem}

\plainproblemsection

\begin{openproblem}
	(W. Schmidt) For \(0\leq s\leq1\), define
	\[
	C_s=
	\left\{
	(\alpha,\beta)\in\mathbb{R}^2:
	\inf_{q\geq1}
	\max\bigl\{
	q^s\langle q\alpha\rangle,\,
	q^{1-s}\langle q\beta\rangle
	\bigr\}>0
	\right\},
	\]
	where \(\langle x\rangle=\operatorname{dist}(x,\mathbb{Z})\).
	Is
	\[
	C_s\cap C_t\neq\varnothing
	\]
	for every \(s,t\in[0,1]\)?
\end{openproblem}

\plainproblemsection

\begin{openproblem}
	Let \(A\) be the diagonal subgroup of
	\(\operatorname{SL}(3,\mathbb{R})\), and let \(A_1,A_2\subset A\) be
	rays.  Does there exist
	\[
	x\in
	\operatorname{SL}(3,\mathbb{R})/\operatorname{SL}(3,\mathbb{Z})
	\]
	such that \(A_1x\) and \(A_2x\) are bounded but \(Ax\) is not bounded?
\end{openproblem}

\plainproblemsection

\begin{openproblem}
	(H. Davenport and D. J. Lewis) Let \(Q\) be a nondegenerate positive
	definite quadratic form in dimension \(d\) that is not a scalar multiple
	of a rational form.  The unresolved cases \(d=3\) and \(d=4\) of the
	conjecture assert that the gaps between consecutive elements of
	\[
	\{Q(\mathbf x):\mathbf x\in\mathbb{Z}^d\}
	\]
	tend to zero as \(Q(\mathbf x)\to\infty\).
\end{openproblem}

\plainproblemsection

\begin{openproblem}
	(G. Margulis) Let \(Q\) be a nondegenerate indefinite quadratic form of
	dimension \(d\geq3\) that is not a scalar multiple of a rational form.
	Give an effective estimate for \(T=T(\varepsilon)\) such that there is
	\(\mathbf x\in\mathbb{Z}^d\) satisfying
	\[
	0<|Q(\mathbf x)|<\varepsilon,
	\qquad
	\lVert\mathbf x\rVert<T.
	\]
\end{openproblem}

\plainproblemsection

\begin{openproblem}
	(E. S. Barnes and H. P. F. Swinnerton-Dyer) Let
	\[
	Q(x_1,x_2)=ax_1^2+bx_1x_2+cx_2^2
	\]
	be a nondegenerate indefinite quadratic form with rational coefficients
	that does not represent zero over \(\mathbb{Q}\).  For
	\(\mathbf x\in\mathbb{R}^2\), define
	\[
	m(Q,\mathbf x)=
	\inf_{\mathbf z\in\mathbb{Z}^2}|Q(\mathbf x+\mathbf z)|,
	\qquad
	m(Q)=\sup_{\mathbf x\in\mathbb{R}^2}m(Q,\mathbf x).
	\]
	If \(m(Q)\) is isolated, also define
	\[
	m_2(Q)=
	\sup\bigl\{
	m(Q,\mathbf x):
	\mathbf x\in\mathbb{R}^2,\ m(Q,\mathbf x)<m(Q)
	\bigr\}.
	\]
	Prove that \(m(Q)\) is rational and isolated and that both \(m(Q)\) and
	\(m_2(Q)\) are attained at points whose coordinates lie in the splitting
	field of \(Q\).
\end{openproblem}

\plainproblemsection

\begin{openproblem}
	(Y. Bugeaud) For \(0<s<1\), let
	\[
	B_s=
	\left\{
	(\alpha,\beta)\in\mathbb{R}^2:
	\inf_{n\geq1}
	n
	\bigl(\min\{\langle n\alpha\rangle,\langle n\beta\rangle\}\bigr)^s
	\bigl(\max\{\langle n\alpha\rangle,\langle n\beta\rangle\}\bigr)^{2-s}
	>0
	\right\}.
	\]
	Compute the Hausdorff dimension of \(B_s\).
\end{openproblem}

\plainproblemsection

\begin{openproblem}
	(A. Katok) Katok constructed a Finsler metric on the two-dimensional
	sphere whose geodesic flow is ergodic and has only two periodic orbits.
	Do the closed geodesics in this example give rise to scars in the
	semiclassical limits of Laplace eigenfunctions?
\end{openproblem}

\plainproblemsection

\begin{openproblem}
	Let \(X\) be a compact Riemannian manifold with negative sectional
	curvature.  Let
	\[
	-\Delta\varphi_i=\lambda_i\varphi_i,\qquad
	\lVert\varphi_i\rVert_{L^2(X)}=1,
	\]
	and define the probability measures
	\[
	\mathrm{d}\mu_i=|\varphi_i|^2\,\mathrm{d}V,
	\]
	where \(\mathrm{d}V\) is normalized Riemannian volume.  The quantum unique
	ergodicity conjecture asserts that
	\[
	\mu_i\longrightarrow\mathrm{d}V
	\qquad\text{as }i\to\infty.
	\]
\end{openproblem}

\plainproblemsection

\begin{openproblem}
	Let \(X=\Gamma\backslash\widetilde X\) be an arithmetic locally symmetric
	space, and let \(\{\varphi_i\}\) be joint eigenfunctions of the invariant
	differential operators and the Hecke operators.  An arithmetic quantum
	limit is a weak-\(*\) limit of the measures
	\[
	|\varphi_i|^2\,\mathrm{d}V.
	\]
	The arithmetic quantum unique ergodicity conjecture asserts that
	Riemannian volume is the only arithmetic quantum limit.
\end{openproblem}

\plainproblemsection

\begin{openproblem}
	(E. Lindenstrauss) Let
	\(X=\Gamma\backslash\mathbb{H}^2\), where \(\Gamma\) is a noncocompact
	arithmetic lattice.  With
	\(\mathrm{d}\mu_i=|\varphi_i|^2\,\mathrm{d}V\), show that, for all
	\(f,g\in C_c(X)\),
	\[
	\frac{\displaystyle\int_X f\,\mathrm{d}\mu_i}
	{\displaystyle\int_X g\,\mathrm{d}\mu_i}
	\longrightarrow
	\frac{\displaystyle\int_X f\,\mathrm{d}V}
	{\displaystyle\int_X g\,\mathrm{d}V}
	\qquad\text{as }i\to\infty.
	\]
\end{openproblem}

\plainproblemsection

\begin{openproblem}
	(J. Marklof) For a generic negatively curved compact Riemannian manifold
	\(X\), let \(\varphi_i\) be \(L^2\)-normalized Laplace eigenfunctions.
	Prove that, for every \(a<b\),
	\[
	\operatorname{Vol}\{x\in X:a\leq\varphi_i(x)\leq b\}
	\longrightarrow
	\frac{1}{\sqrt{2\pi}}
	\int_a^b e^{-t^2/2}\,\mathrm{d}t
	\qquad\text{as }i\to\infty.
	\]
\end{openproblem}

\plainproblemsection

\begin{openproblem}
	(J. Marklof, after M. Feingold and A. Peres) Let \(X\) be a generic
	negatively curved compact Riemannian surface, normalized to have area
	\(4\pi\).  For \(a\in C_c^\infty(S^*X,\mathbb{R})\) with
	\[
	\int_{S^*X}a\,\mathrm{d}\lambda=0,
	\]
	put
	\[
	\xi_i(a)=
	\lambda_i^{1/4}
	\left|
	\left\langle\operatorname{Op}(a)\varphi_i,\varphi_i\right\rangle
	\right|
	\]
	and
	\[
	V(a)=
	\int_{-\infty}^{\infty}
	\int_{S^*X}
	a(xg^t)a(x)\,\mathrm{d}\lambda(x)\,\mathrm{d}t,
	\]
	where \(g^t\) is the geodesic flow and \(\lambda\) is Liouville measure.
	
	Prove that \(\xi_i(a)\) has a Gaussian limiting distribution of variance
	\(V(a)\).  More explicitly, prove
	\[
	\frac1\lambda
	\sum_{\lambda_i\leq\lambda}\xi_i(a)^2
	\longrightarrow V(a)
	\]
	and, for every interval \(I\subset\mathbb{R}\),
	\[
	\frac1\lambda
	\#\left\{
	\lambda_i\leq\lambda:
	V(a)^{-1/2}\xi_i(a)\in I
	\right\}
	\longrightarrow
	\frac1{\sqrt{2\pi}}\int_I e^{-t^2/2}\,\mathrm{d}t.
	\]
\end{openproblem}

\plainproblemsection

\begin{openproblem}
	(A. Katok) Do periodic orbits in a triangular billiard correspond to
	scars?  More precisely, are there quantum limits supported on periodic
	orbits?
\end{openproblem}

\plainproblemsection

\begin{openproblem}
	(J. Marklof) Classify all quantum limits of the eigenfunctions of a
	polygonal billiard.
\end{openproblem}

\plainproblemsection

\begin{openproblem}
	(J. Marklof) Let \(Q\) be a positive definite binary quadratic form.
	Determine the distribution of
	\[
	\left\{
	\bigl(Q(\mathbf x_1)-Q(\mathbf x_2),
	Q(\mathbf x_2)-Q(\mathbf x_3)\bigr):
	\mathbf x_i\in\mathbb{Z}^2,\ 
	\mathbf x_i\neq\mathbf x_j\ (i\neq j)
	\right\}.
	\]
	More precisely, determine the asymptotics, as \(T\to\infty\), of
	\[
	\begin{aligned}
		N_T((a,b),(c,d))
		=\#\bigl\{(\mathbf x_1,\mathbf x_2,\mathbf x_3)\in(\mathbb{Z}^2)^3:\;&
		a<Q(\mathbf x_1)-Q(\mathbf x_2)<b,\\
		&c<Q(\mathbf x_2)-Q(\mathbf x_3)<d,\\
		&\mathbf x_i\neq\mathbf x_j\ (i\neq j),\
		\lVert\mathbf x_i\rVert<T
		\bigr\}.
	\end{aligned}
	\]
\end{openproblem}

\plainproblemsection

\begin{openproblem}
	(Y. Andr\'e and F. Oort) Let \(\operatorname{Sh}_K(G,X)\) be a Shimura
	variety.  The Andr\'e--Oort conjecture asserts that the Zariski closure of
	any set of special points is a finite union of special subvarieties.
\end{openproblem}

\plainproblemsection

\begin{openproblem}
	Let \(\{x_n\}\) be a sequence of special points on a Shimura variety.
	Suppose that, for every proper special subvariety \(Y\), all sufficiently
	large \(n\) satisfy \(x_n\notin Y\).  Prove that the Galois orbits of
	\(x_n\) become equidistributed, as \(n\to\infty\), with respect to
	normalized Haar measure.
\end{openproblem}

\plainproblemsection

\begin{openproblem}
	(L. Silberman) Give an ergodic-theoretic proof of the equidistribution of
	special points on
	\[
	\operatorname{SL}(2,\mathbb{Z})\backslash\mathbb{H}^2.
	\]
\end{openproblem}

\plainproblemsection

\begin{openproblem}
	Let \(\gamma\) and \(\Gamma\) be convex algebraic ovals, with
	\(\gamma\) inside \(\Gamma\).  The Poncelet map sends a point of
	\(\Gamma\) to the second intersection of \(\Gamma\) with the
	consistently chosen tangent line from that point to \(\gamma\).
	
	Do there exist irreducible algebraic curves of degrees \(n\) and \(m\),
	with \(n+m>4\), each containing such an oval, for which the Poncelet map
	is well defined and conjugate to a rotation of the circle?
\end{openproblem}

\plainproblemsection

\begin{openproblem}
	Consider the ovals
	\[
	\gamma=\{x^2+y^2-1=0\},
	\qquad
	\Gamma_\varepsilon=
	\{p_2(x,y)+\varepsilon p_m(x,y)=0\},
	\]
	where \(\Gamma_0\) is an ellipse surrounding \(\gamma\),
	\(p_m\) has degree \(m\geq3\), \(\varepsilon\) is small, and
	\(p_2+\varepsilon p_m\) is irreducible.  Is the Poncelet map associated
	with these two ovals conjugate to a rotation if and only if
	\(\varepsilon=0\)?
\end{openproblem}

\plainproblemsection

\begin{openproblem}
	Let \(\mathcal U\subset\mathbb{R}^2\) be a neighborhood of the origin,
	let \(f\colon\mathcal U\to\mathbb{R}\) be of class \(C^{n+1}\), with
	\(f(0,0)=0\), and consider
	\[
	\dot x=
	2^n\operatorname{Re}\left(
	\frac{\partial^nf}{\partial\overline z^{\,n}}
	\right),
	\qquad
	\dot y=
	2^n\operatorname{Im}\left(
	\frac{\partial^nf}{\partial\overline z^{\,n}}
	\right),
	\qquad
	\frac{\partial}{\partial\overline z}
	=\frac12\left(
	\frac{\partial}{\partial x}
	+\mathrm{i}\frac{\partial}{\partial y}
	\right).
	\]
	Loewner's conjecture asserts that, if \(f\) is analytic at the origin
	and the origin is an isolated equilibrium of this vector field, then
	the index of the vector field at the origin is at most \(n\).
\end{openproblem}

\subsection*{References}

\begin{enumerate}[label={[\arabic*]},leftmargin=*,itemsep=0.4em]
	\item A. Gorodnik, ``Open problems in dynamics and related fields,''
	\emph{Journal of Modern Dynamics} \textbf{1} (2007), no.~1, 1--35.
	\item Y. Andr\'e, \emph{G-functions and Geometry}, Aspects of Mathematics
	E13, Friedr. Vieweg \& Sohn, Braunschweig, 1989.
	\item E. S. Barnes and H. P. F. Swinnerton-Dyer, ``The inhomogeneous
	minima of binary quadratic forms II,'' \emph{Acta Mathematica}
	\textbf{88} (1952), 279--316.
	\item M. Feingold and A. Peres, ``Distribution of matrix elements of
	chaotic systems,'' \emph{Physical Review A} \textbf{34} (1986),
	591--595.
	\item H. Furstenberg, ``Disjointness in ergodic theory, minimal sets, and
	a problem in Diophantine approximation,'' \emph{Mathematical Systems
		Theory} \textbf{1} (1967), 1--49.
	\item A. Katok and R. Spatzier, ``Invariant measures for higher-rank
	hyperbolic abelian actions,'' \emph{Ergodic Theory and Dynamical
		Systems} \textbf{16} (1996), no.~4, 751--778.
	\item D. Y. Kleinbock and G. A. Margulis, ``Flows on homogeneous spaces
	and Diophantine approximation on manifolds,'' \emph{Annals of
		Mathematics} (2) \textbf{148} (1998), no.~1, 339--360.
	\item G. A. Margulis, ``Problems and conjectures in rigidity theory,'' in
	\emph{Mathematics: Frontiers and Perspectives}, American Mathematical
	Society, Providence, RI, 2000, 161--174.
	\item J. Marklof and A. Str\"ombergsson, ``Equidistribution of Kronecker
	sequences along closed horocycles,'' \emph{Geometric and Functional
		Analysis} \textbf{13} (2003), no.~6, 1239--1280.
	\item F. Oort, ``Canonical liftings and dense sets of CM-points,'' in
	\emph{Arithmetic Geometry (Cortona, 1994)}, Symposia Mathematica
	XXXVII, Cambridge University Press, Cambridge, 1997, 229--234.
	\item Z. Rudnick, P. Sarnak, and A. Zaharescu, ``The distribution of
	spacings between the fractional parts of \(n^2\alpha\),''
	\emph{Inventiones Mathematicae} \textbf{145} (2001), 37--57.
	\item W. M. Schmidt, ``Open problems in Diophantine approximation,'' in
	\emph{Diophantine Approximations and Transcendental Numbers
		(Luminy, 1982)}, Progress in Mathematics 31, Birkh\"auser, Boston,
	1983, 271--287.
	\item Ya. Vorobets, ``Ergodicity of billiards in polygons,''
	\emph{Sbornik: Mathematics} \textbf{188} (1997), no.~3, 389--434.
	\item A. Gasull, ``Some open problems in low dimensional dynamical
	systems,'' arXiv:2012.02524, 2020.
\end{enumerate}

\clearpage
\section{Low-Dimensional Differential Systems}

\plainproblemsection

\begin{openproblem}
	Consider the family of planar cubic rigid systems
	\[
	\begin{cases}
		\dot x=-y+x(a+bx+cy+dx^2+exy),\\
		\dot y=\phantom{-}x+y(a+bx+cy+dx^2+exy).
	\end{cases}
	\]
	Is \(2\) the maximum number of limit cycles in this family?
\end{openproblem}

\plainproblemsection

\begin{openproblem}
	Let \(n\neq m\), and let \(P_n\) and \(Q_m\) be homogeneous
	polynomials of odd degrees \(n\) and \(m\), respectively.  Is
	\((n+m)/2\) the maximum number of limit cycles of
	\[
	\dot x=P_n(x,y),\qquad \dot y=Q_m(x,y)?
	\]
\end{openproblem}

\plainproblemsection

\begin{openproblem}
	\begin{enumerate}[label=\textup{(\roman*)}]
		\item For the cubic family
		\[
		\begin{cases}
			\dot x=ax+by,\\
			\dot y=cx^3+dx^2y+exy^2+fy^3,
		\end{cases}
		\]
		is \(2\) the maximum number of limit cycles?
		\item If it is true, give a simple proof of the uniqueness of the limit
		cycle for
		\[
		\begin{cases}
			\dot x=y,\\
			\dot y=-x^3+dx^2y+y^3.
		\end{cases}
		\]
	\end{enumerate}
\end{openproblem}

\plainproblemsection

\begin{openproblem}
	For the Li\'enard system
	\[
	\dot x=y-F(x),\qquad \dot y=-x,
	\]
	let \(\operatorname{Lie}(n)\) be the maximum number of limit cycles
	when \(F\) is a polynomial of degree \(n\).  Find proofs using Dulac
	functions that
	\[
	\operatorname{Lie}(3)=1
	\qquad\text{and}\qquad
	\operatorname{Lie}(4)=1.
	\]
\end{openproblem}

\plainproblemsection

\begin{openproblem}
	For a general \(T\)-periodic Riccati differential equation
	\[
	\frac{\mathrm{d}x}{\mathrm{d}t}
	=A_2(t)x^2+A_1(t)x+A_0(t),
	\]
	give effective criteria for deciding whether it has a continuum of
	periodic solutions or exactly \(2\), \(1\), or \(0\) limit cycles.
\end{openproblem}

\plainproblemsection

\begin{openproblem}
	Consider the family of trigonometric Abel differential equations
	\[
	\frac{\mathrm{d}x}{\mathrm{d}t}
	=(a_0+a_1\sin t+a_2\cos t)x^3
	+(b_0+b_1\sin t+b_2\cos t)x^2.
	\]
	Is \(3\) the maximum number of \(2\pi\)-periodic limit cycles?
\end{openproblem}

\plainproblemsection

\begin{openproblem}
	Given integers \(p>q\geq2\) and \(m,n\in\mathbb{N}\), determine the
	maximum number of \(2\pi\)-periodic limit cycles of
	\[
	\frac{\mathrm{d}x}{\mathrm{d}t}
	=A_m(t)x^p+B_n(t)x^q,
	\]
	where \(A_m\) and \(B_n\) are \(2\pi\)-periodic trigonometric
	polynomials of degrees \(m\) and \(n\), respectively.
\end{openproblem}

\plainproblemsection

\begin{openproblem}
	For \(m\in\mathbb{N}\), let \(\mathcal M_m\) be the family of planar
	polynomial differential systems having \(m\) monomial vector-field
	terms:
	\[
	(\dot x,\dot y)=\sum_{j=1}^m a_jX_j(x,y),
	\]
	where each \(X_j\) is either
	\((x^{n_j}y^{k_j},0)\) or \((0,x^{n_j}y^{k_j})\).
	Let \(\mathcal H^M[m]\) be the maximum number of limit cycles among
	systems in \(\mathcal M_m\).
	
	\begin{enumerate}[label=\textup{(\roman*)}]
		\item Find upper and lower bounds for \(\mathcal H^M[m]\).
		\item Find the smallest \(m\) for which some planar polynomial
		differential system with \(m\) monomials has at least \(m+1\) limit
		cycles.
	\end{enumerate}
\end{openproblem}

\plainproblemsection

\begin{openproblem}
	Let \(f\) be a continuous, nonzero, \(T\)-periodic function, and let
	\(p>0\).  Find necessary and sufficient conditions on \(f\) for the
	singular differential equation
	\[
	x^p(t)x''(t)=f(t)
	\]
	to have a positive \(T\)-periodic solution.
\end{openproblem}

\plainproblemsection

\begin{openproblem}
	Let \(\mathcal C_n\) be the maximum number of centers of a planar
	polynomial differential system of degree \(n\) having only finitely many
	equilibria.  Determine \(\mathcal C_n\) for \(n\geq4\).
\end{openproblem}

\plainproblemsection

\begin{openproblem}
	Consider the \(k\)-th order rational difference equations
	\[
	x_{n+k}=
	\frac{A_0+A_1x_n+A_2x_{n+1}+\cdots+A_kx_{n+k-1}}
	{B_0+B_1x_n+B_2x_{n+1}+\cdots+B_kx_{n+k-1}},
	\]
	where
	\[
	\sum_{i=0}^kA_i>0,\qquad
	\sum_{i=0}^kB_i>0,\qquad
	A_i,B_i\geq0,\qquad
	A_1^2+B_1^2\neq0.
	\]
	The equation is called \(p\)-periodic if every positive initial
	condition generates a sequence satisfying \(x_{n+p}=x_n\), and \(p\)
	is the smallest positive integer with this property.
	
	Apart from the equations obtained from the following five by the
	standard unfolding and rescaling equivalences,
	\[
	\begin{array}{lll}
		x_{n+1}=x_n &(p=1),&
		x_{n+1}=x_n^{-1} \quad(p=2),\\[2pt]
		x_{n+2}=x_{n+1}/x_n &(p=6),&
		x_{n+2}=(1+x_{n+1})/x_n \quad(p=5),\\[2pt]
		x_{n+3}=(1+x_{n+1}+x_{n+2})/x_n &(p=8),&
	\end{array}
	\]
	do any other periodic equations of this form exist?
\end{openproblem}

\plainproblemsection

\begin{openproblem}
	Consider a Hamiltonian system with a center at the origin and Hamiltonian
	\[
	H(x,y)=H_{2n}(x,y)+H_m(x,y),\qquad m>2n,
	\]
	where \(H_{2n}\) and \(H_m\) are homogeneous polynomials of degrees
	\(2n\) and \(m\), respectively.  Does the period annulus of the origin
	have at most one critical period?
\end{openproblem}

\plainproblemsection

\begin{openproblem}
	Let
	\[
	\dot x=P_{2k+1}(x,y),\qquad
	\dot y=Q_{2\ell+1}(x,y),
	\]
	where \(P_{2k+1}\) and \(Q_{2\ell+1}\) are homogeneous polynomials of
	degrees \(2k+1\) and \(2\ell+1\), respectively.
	\begin{enumerate}[label=\textup{(\roman*)}]
		\item Characterize the centers in this family.
		\item Determine the maximum number of oscillations of the period function
		among the centers in this family.
	\end{enumerate}
\end{openproblem}

\plainproblemsection

\begin{openproblem}
	Let \(\mathcal T(n)\) be the maximum possible number of critical periods
	for a planar polynomial differential system of degree \(n\).  Is there a
	constant \(C>0\) such that
	\[
	\mathcal T(n)\geq Cn^2\log n?
	\]
\end{openproblem}

\plainproblemsection

\begin{openproblem}
	Consider the family of reversible quadratic centers
	\[
	\begin{cases}
		\dot x=-y+xy,\\
		\dot y=x+Dx^2+Fy^2.
	\end{cases}
	\]
	Is \(2\) the maximum number of critical periods in this family?
\end{openproblem}

\plainproblemsection

\begin{openproblem}
	For positive integers \(n\) and \(k\), consider
	\[
	\dot z=\mathrm{i}z+(z\overline z)^nz^{k+1}.
	\]
	Is the period function of the period annulus surrounding the center at
	the origin monotonically decreasing?
\end{openproblem}

\plainproblemsection

\begin{openproblem}
	Determine or improve the entries marked by question marks in the
	following comparison between quadratic systems and piecewise linear
	systems with a straight line of separation:
	\begin{center}
		\small
		\begin{tabularx}{\textwidth}{|Y|>{\centering\arraybackslash}X|>{\centering\arraybackslash}X|}
			\hline
			& Quadratic systems & Piecewise linear systems\\
			\hline
			Limit cycles & \(4\)? & \(3\)?\\
			\hline
			Algebraic limit cycles & \(1\)? & \(2\)?\\
			\hline
			Nonhyperbolic algebraic limit cycles & \(0\)? & \(1\)?\\
			\hline
			Coexistence of algebraic and nonalgebraic limit cycles &
			\(\mathrm{No}\)? & \(\mathrm{No}\)\\
			\hline
			Critical periods & \(2\)? &?\\
			\hline
		\end{tabularx}
	\end{center}
	Only crossing limit cycles are counted in the piecewise linear case.
\end{openproblem}

\plainproblemsection

\begin{openproblem}
	Let \(\mathcal H(n)\) be the maximum number of limit cycles of a planar
	polynomial differential system of degree \(n\).  Let \(\mathcal L(n)\)
	be the maximum number of crossing limit cycles of a planar piecewise
	linear system whose two zones are separated by a branch of an algebraic
	curve of degree \(n\).
	
	Improve, if possible, the following lower bounds:
	\begin{center}
		\small
		\begin{tabularx}{0.9\textwidth}{|>{\centering\arraybackslash}X|>{\centering\arraybackslash}X|}
			\hline
			Polynomial case & Piecewise linear case\\
			\hline
			\(\mathcal H(2)\geq4\) & \(\mathcal L(1)\geq3\)\\
			\hline
			\(\mathcal H(3)\geq13\) & \(\mathcal L(2)\geq4\)\\
			\hline
			\(\mathcal H(n)\geq Kn^2\log n\) &
			\(\mathcal L(n)\geq\lfloor n/2\rfloor\)\\
			\hline
		\end{tabularx}
	\end{center}
\end{openproblem}

\plainproblemsection

\begin{openproblem}
	Does there exist a smooth vector field
	\(\dot{\mathbf x}=F(\mathbf x)\) on \(\mathbb{R}^3\) such that
	\[
	F(\mathbf0)=\mathbf0,
	\qquad
	\operatorname{Re}\lambda<0
	\quad\text{for every eigenvalue \(\lambda\) of \(DF(\mathbf x)\)
		and every \(\mathbf x\in\mathbb{R}^3\),}
	\]
	but the vector field nevertheless has a periodic orbit?
\end{openproblem}

\plainproblemsection

\begin{openproblem}
	Let \(F\colon\mathbb{R}^2\to\mathbb{R}^2\) be a smooth map with a
	fixed point.  If
	\[
	\rho\bigl(|DF(x)|\bigr)<1
	\qquad\text{for every }x\in\mathbb{R}^2,
	\]
	where \(|DF(x)|\) is obtained by taking the absolute value of every
	matrix entry and \(\rho\) denotes spectral radius, is the fixed point
	globally asymptotically stable?
\end{openproblem}

\plainproblemsection

\begin{openproblem}
	Let \(A_0,\ldots,A_n\) be independent standard normal random variables
	and consider
	\[
	A_nx^{(n)}(t)+A_{n-1}x^{(n-1)}(t)+\cdots+
	A_1x'(t)+A_0x(t)=0.
	\]
	Let \(p_n\) be the probability that the zero solution is globally
	asymptotically stable.  Find the asymptotic expansion of \(p_n\) as
	\(n\to\infty\).  Is the sequence \(\{p_n\}\) strictly decreasing?
\end{openproblem}

\subsection*{References}

\begin{enumerate}[label={[\arabic*]},leftmargin=*,itemsep=0.4em]
	\item A. Gasull, ``Some open problems in low dimensional dynamical
	systems,'' arXiv:2012.02524, 2020.
\end{enumerate}

\clearpage
\section{Symbolic Dynamics}

In this section, \(\mathbb{Z}_+\) denotes the nonnegative integers.
A square matrix \(A\) over \(\mathbb{Z}_+\) defines an edge shift of
finite type, denoted by \(\sigma_A\).  A sofic shift is a factor of a
shift of finite type.  For matrices over a semiring \(R\), an elementary
strong shift equivalence is a pair of factorizations
\[
A=UV,\qquad B=VU.
\]
Strong shift equivalence over \(R\) is the equivalence relation generated
by these elementary steps.  Shift equivalence over \(R\) means that, for
some \(\ell\geq1\), there are matrices \(U,V\) such that
\[
A^\ell=UV,\qquad B^\ell=VU,\qquad AU=UB,\qquad VA=BV.
\]

\plainproblemsection

\begin{openproblem}
	The little shift equivalence conjecture asserts that, if a nonnegative
	integer matrix \(A\) has a unique simple nonzero eigenvalue \(n\), then
	\(A\) is strong shift equivalent over \(\mathbb{Z}_+\) to the
	one-by-one matrix \([n]\).
\end{openproblem}

\plainproblemsection

\begin{openproblem}
	Classify shifts of finite type up to topological conjugacy.  In
	particular, give a decision procedure for determining whether two
	nonnegative integer matrices \(A\) and \(B\) define topologically
	conjugate shifts of finite type.
\end{openproblem}

\plainproblemsection

\begin{openproblem}
	Let \(\sigma_A\) be a mixing shift of finite type, let \(G_A\) be its
	dimension group, and let
	\[
	\rho_A\colon\operatorname{Aut}(\sigma_A)
	\longrightarrow\operatorname{Aut}(G_A)
	\]
	be the dimension representation induced by the action of shift
	automorphisms on \(G_A\).  Determine the range of \(\rho_A\).
\end{openproblem}

\plainproblemsection

\begin{openproblem}
	The positive rational shift equivalence conjecture asserts that two
	positive square matrices that are shift equivalent over
	\(\mathbb{Q}_+\) are strong shift equivalent over \(\mathbb{Q}_+\).
\end{openproblem}

\plainproblemsection

\begin{openproblem}
	Let \(S\) be a unital subring of \(\mathbb{R}\), and let
	\(\Lambda=(\lambda_1,\ldots,\lambda_k)\) be a tuple of complex numbers.
	Set
	\[
	\operatorname{tr}(\Lambda^n)=\sum_{i=1}^k\lambda_i^n,
	\qquad
	\operatorname{tr}_n(\Lambda)=
	\sum_{j\mid n}\mu(n/j)\operatorname{tr}(\Lambda^j),
	\]
	where \(\mu\) is the M\"obius function.
	
	The spectral conjecture asserts that \(\Lambda\) is the nonzero spectrum
	of a primitive matrix over \(S\) if and only if:
	\begin{enumerate}[label=\textup{(\arabic*)}]
		\item one \(\lambda_i\) is positive real and strictly larger in modulus
		than every other \(\lambda_j\);
		\item \(\prod_{i=1}^k(t-\lambda_i)\) has all coefficients in \(S\);
		\item if \(S\neq\mathbb{Z}\), then
		\(\operatorname{tr}(\Lambda^n)\geq0\) for every \(n\), and
		\(j\mid n\) together with \(\operatorname{tr}(\Lambda^j)>0\) implies
		\(\operatorname{tr}(\Lambda^n)>0\); if \(S=\mathbb{Z}\), then
		\(\operatorname{tr}_n(\Lambda)\geq0\) for every \(n\).
	\end{enumerate}
\end{openproblem}

\plainproblemsection

\begin{openproblem}
	Let \(S\) be a unital subring of \(\mathbb{R}\), and let \(A\) be a
	square matrix over \(S\) whose nonzero spectrum satisfies the Perron,
	Galois, and trace conditions in the spectral conjecture.  The generalized
	spectral conjecture asserts that \(A\) is strong shift equivalent over
	\(S\) to a primitive matrix.
\end{openproblem}

\plainproblemsection

\begin{openproblem}
	Suppose \(A\) and \(B\) are irreducible nonnegative integer matrices
	with the same spectral radius and:
	\begin{enumerate}[label=\textup{(\arabic*)}]
		\item for every \(n\geq1\),
		\(\operatorname{tr}(A^n)>0\) implies
		\(\operatorname{tr}(B^n)>0\);
		\item the dimension module of \(B\) is a quotient of a closed submodule
		of the dimension module of \(A\).
	\end{enumerate}
	Must the shift of finite type \(\sigma_B\) be a factor of
	\(\sigma_A\)?
\end{openproblem}

\plainproblemsection

\begin{openproblem}
	Let \(S\) and \(T\) be sofic shifts with \(h(S)\geq h(T)\).  Give
	necessary and sufficient conditions for the existence of a factor map
	from \(S\) onto \(T\).
\end{openproblem}

\plainproblemsection

\begin{openproblem}
	Let \(P\) be a stochastic matrix presenting a mixing Markov shift.  Form
	\(P(t)\) by replacing each positive entry \(p\) by \(p^t\), and put
	\[
	q(z)=\det(zI-P(t)).
	\]
	If \(b(z)\) is the irreducible factor of \(q(z)\) satisfying \(b(1)=0\),
	then \(1/b(z)\) is the beta function.  Let \(\Gamma\) be the
	multiplicative group generated by the weights of loops in the graph of
	\(P\), and let \(\Delta\) be the group of ratios of weights of two paths
	having the same length and the same initial and terminal vertices.
	
	The good finitary conjecture asserts that two mixing Markov shifts admit
	a magic-word isomorphism if and only if they have the same beta function,
	the same ratio group \(\Delta\), and the same canonical generator of the
	cyclic quotient group \(\Gamma/\Delta\).
\end{openproblem}

\plainproblemsection

\begin{openproblem}
	Characterize the functions that can occur as stochastic zeta functions
	of mixing Markov shifts.
\end{openproblem}

\plainproblemsection

\begin{openproblem}
	Characterize the functions that can occur as beta functions of mixing
	Markov shifts.
\end{openproblem}

\plainproblemsection

\begin{openproblem}
	Let \(\alpha\) be a \(\mathbb{Z}^2\)-action on a compact metric space.
	A line \(L\subset\mathbb{R}^2\) is called expansive if there are
	\(M,\varepsilon>0\) such that, for every distinct \(x,y\), some
	\(\mathbf n\in\mathbb{Z}^2\) with
	\(\operatorname{dist}(\mathbf n,L)<M\) satisfies
	\[
	\operatorname{dist}\bigl(\alpha^{\mathbf n}x,
	\alpha^{\mathbf n}y\bigr)>\varepsilon.
	\]
	Let \(E_1(\alpha)\) be the set of
	\(\mathbf v\in\mathbb{R}^2\) for which the line
	\(\mathbb{R}\mathbf v\) is expansive.
	
	Which sets can occur as \(E_1(\alpha)\) for a
	\(\mathbb{Z}^2\) shift of finite type?  In particular, which can occur
	when some \(\alpha^{\mathbf n}\) is itself a shift of finite type?
\end{openproblem}

\plainproblemsection

\begin{openproblem}
	Let \(\alpha\) be a \(\mathbb{Z}^2\)-action for which some
	\(\alpha^{\mathbf n}\) is a shift of finite type.
	\begin{enumerate}[label=\textup{(\arabic*)}]
		\item Can \(\alpha\) have infinitely many expansive components?
		\item Can an expansive component have a boundary contained in a line of
		irrational slope?
	\end{enumerate}
\end{openproblem}

\plainproblemsection

\begin{openproblem}
	Let \(L\) be a line of irrational slope in \(\mathbb{R}^2\).  Does
	there exist a \(\mathbb{Z}^2\)-action \(\alpha\) on a compact metric
	space such that
	\[
	\mathbb{R}^2\setminus E_1(\alpha)=L?
	\]
\end{openproblem}

\plainproblemsection

\begin{openproblem}
	Let \(L\) be a line of rational slope in \(\mathbb{R}^2\).  Does
	there exist a \(\mathbb{Z}^2\)-action \(\alpha\) on a compact metric
	space \(X\) such that
	\[
	\mathbb{R}^2\setminus E_1(\alpha)=L
	\]
	and \(\alpha^{\mathbf n}\) is not the identity on \(X\) for every
	nonzero \(\mathbf n\in L\cap\mathbb{Z}^2\)?
\end{openproblem}

\plainproblemsection

\begin{openproblem}
	Suppose \(S\) is an expansive automorphism of an irreducible shift of
	finite type.  Must the dynamical system defined by \(S\) itself be a
	shift of finite type?
\end{openproblem}

\plainproblemsection

\begin{openproblem}
	An expansive automorphism of a one-sided full shift is conjectured to be
	topologically conjugate, rather than merely shift equivalent, to a
	two-sided full shift.
\end{openproblem}

\plainproblemsection

\begin{openproblem}
	Let \(S\) and \(T\) be mixing shifts of finite type.  Is it true that,
	for all sufficiently large positive integers \(i\) and \(j\), there are
	realizations of \(S^i\) and \(T^j\) that commute?
\end{openproblem}

\plainproblemsection

\begin{openproblem}
	Let \(X\) be Ledrappier's three-dot system, the closed
	shift-invariant subgroup of \(\{0,1\}^{\mathbb{Z}^2}\) consisting of
	the configurations \(x\) satisfying
	\[
	x(\mathbf n)+x(\mathbf n+(0,1))+x(\mathbf n+(1,0))=0
	\pmod 2
	\]
	for every \(\mathbf n\in\mathbb{Z}^2\).  Is Haar measure the only
	shift-invariant ergodic Borel probability measure on \(X\) with full
	support?
\end{openproblem}

\plainproblemsection

\begin{openproblem}
	For Ledrappier's three-dot system \(X\), determine all
	shift-invariant Borel probability measures and all shift-invariant
	subsystems of \(X\).
\end{openproblem}

\plainproblemsection

\begin{openproblem}
	(B. Weiss) For \(d>1\), must every \(\mathbb{Z}^d\) sofic shift be a
	factor of a \(\mathbb{Z}^d\) shift of finite type having the same
	entropy?
\end{openproblem}

\plainproblemsection

\begin{openproblem}
	(K. Schmidt) Let \(d>1\), and let
	\(\varphi\colon X\to Y\) be a continuous factor map from a
	\(\mathbb{Z}^d\) shift of finite type onto a
	\(\mathbb{Z}^d\) sofic shift.  Must there be a sofic subshift
	\(W\subset X\) such that
	\[
	h(W)=h(Y),
	\qquad
	\varphi(W)=Y?
	\]
\end{openproblem}

\plainproblemsection

\begin{openproblem}
	(A. Maass) Characterize the stable limit sets of one-dimensional
	cellular automata.  Here a cellular automaton is a block code
	\[
	f\colon X_N\to X_N,\qquad
	X_N=\{0,1,\ldots,N-1\}^{\mathbb{Z}},
	\]
	and the limit set is stable if
	\[
	f^k(X_N)=f^{k+1}(X_N)
	\]
	for some \(k>0\).
\end{openproblem}

\plainproblemsection

\begin{openproblem}
	Let \(T\) be a mixing sofic shift with a receptive fixed point.  When
	does there exist a block code \(f\colon T\to T\) and a shift of finite
	type \(T'\supset T\) such that \(f(T')\subset T\)?
\end{openproblem}

\plainproblemsection

\begin{openproblem}
	Let \(f\) be a surjective block code from a mixing sofic shift \(T\)
	onto itself.  When does there exist a shift of finite type
	\(T'\supset T\) such that \(f(T')\subset T\)?
\end{openproblem}

\plainproblemsection

\begin{openproblem}
	A \(\mathbb{Z}^d\) group shift is an expansive action of
	\(\mathbb{Z}^d\) by continuous automorphisms of a zero-dimensional
	compact metrizable group.  A Bernoulli group shift is the
	\(\mathbb{Z}^d\)-shift on \(G^{\mathbb{Z}^d}\) for a finite group
	\(G\).
	
	Does every nonabelian group shift factor algebraically onto a Bernoulli
	group shift?
\end{openproblem}

\plainproblemsection

\begin{openproblem}
	Two group shifts are weakly algebraically equivalent if each is a
	factor of the other under a continuous group homomorphism.  Is every
	nonabelian group shift weakly algebraically equivalent to a Bernoulli
	group shift?
\end{openproblem}

\plainproblemsection

\begin{openproblem}
	Classify group shifts up to topological conjugacy, with particular
	emphasis on abelian group shifts having completely positive entropy.
\end{openproblem}

\plainproblemsection

\begin{openproblem}
	Let \(\mu\) be a translation-invariant Markov random field on a
	\(\mathbb{Z}^d\) shift \(S\).  Suppose that \(\mu\) is the unique
	Markov random field, among both translation-invariant and
	non-translation-invariant fields, with its conditional probabilities.
	Must \((S,\mu)\) be measurably isomorphic to a Bernoulli shift?
\end{openproblem}

\plainproblemsection

\begin{openproblem}
	(J. Steif) Suppose a \(\mathbb{Z}^d\) shift of finite type has a
	unique measure of maximal entropy and that this measure is Bernoulli.
	Must there exist an independent identically distributed process that
	maps finitarily onto it?
\end{openproblem}

\plainproblemsection

\begin{openproblem}
	Determine whether a decision procedure exists for each of the following:
	\begin{enumerate}[label=\textup{(\arabic*)}]
		\item Given square matrices \(A,B\) over \(\mathbb{Z}_+\), decide whether
		\(\sigma_A\) and \(\sigma_B\) are topologically conjugate, equivalently
		whether \(A\) and \(B\) are strong shift equivalent over
		\(\mathbb{Z}_+\).
		\item Decide whether two given two-sided sofic shifts are topologically
		conjugate.
		\item Decide whether two given one-sided sofic shifts are topologically
		conjugate.
		\item Given an expansive automorphism \(U\) of a shift of finite type
		\(S\), compute the expansive component containing \(S\) for the
		\(\mathbb{Z}^2\)-action generated by \(U\) and \(S\).
		\item Given a surjective block code on a full shift, decide whether it is
		positively expansive.
		\item Given an automorphism of a mixing shift of finite type, decide
		whether it is expansive.
		\item Given an expansive automorphism of a mixing shift of finite type,
		decide whether it is itself a shift of finite type.
	\end{enumerate}
\end{openproblem}

\plainproblemsection

\begin{openproblem}
	Let \(T\) be a mixing one-sided shift of finite type, and let \(S\) be
	a one-sided subshift with \(h(S)<h(T)\).  Give useful necessary and
	sufficient conditions for \(S\) to be topologically conjugate to a
	subshift of \(T\).
\end{openproblem}

\plainproblemsection

\begin{openproblem}
	Let \(S\) be a one-sided subshift, and let \(T\) be the full one-sided
	shift on \(N\) symbols.  Suppose that
	\[
	h(S)<h(T)=\log N
	\]
	and no point of \(S\) has more than \(N\) preimages.  Must \(S\) be
	topologically conjugate to a subshift of \(T\)?
\end{openproblem}

\plainproblemsection

\begin{openproblem}
	Classify one-sided sofic shifts up to topological conjugacy.
\end{openproblem}

\plainproblemsection

\begin{openproblem}
	Let
	\[
	p\colon S\to T,
	\qquad
	p'\colon S\to T'
	\]
	be canonical left-resolving irreducible covers of one-sided irreducible
	sofic shifts, with common domain the one-sided irreducible shift of
	finite type \(S\).  Give a decision procedure for determining whether
	there is an automorphism \(\varphi\) of \(S\) that carries the quotient
	relation of \(p\) to that of \(p'\); equivalently,
	\[
	p(x)=p(y)
	\quad\Longleftrightarrow\quad
	p'(\varphi x)=p'(\varphi y).
	\]
\end{openproblem}

\plainproblemsection

\begin{openproblem}
	Let \(\sigma_N\) denote the full two-sided shift on \(N\) symbols.  Are
	the groups
	\[
	\operatorname{Aut}(\sigma_2)
	\qquad\text{and}\qquad
	\operatorname{Aut}(\sigma_3)
	\]
	isomorphic?
\end{openproblem}

\plainproblemsection

\begin{openproblem}
	Let \(S\) be a mixing shift of finite type.  Let
	\(\operatorname{Aut}_0(S)\) be the kernel of its dimension
	representation, and let \(F_0(S)\) be the subgroup of
	\(\operatorname{Aut}_0(S)\) generated by the finite-order elements.
	The virtual finite-order-generation conjecture asserts that
	\[
	\operatorname{Aut}_0(S)/F_0(S)
	\]
	is finite.
\end{openproblem}

\plainproblemsection

\begin{openproblem}
	Two group actions on compact metric spaces are topologically orbit
	equivalent if a homeomorphism maps every orbit of one action onto an
	orbit of the other.  Classify minimal actions of \(\mathbb{Z}^d\) on
	the Cantor set up to topological orbit equivalence.
\end{openproblem}

\plainproblemsection

\begin{openproblem}
	For \(d>2\), is every minimal \(\mathbb{Z}^d\)-action on the Cantor set
	topologically orbit equivalent to a \(\mathbb{Z}\)-action?
\end{openproblem}

\plainproblemsection

\begin{openproblem}
	Is every minimal action of a countable amenable group on the Cantor set
	topologically orbit equivalent to a \(\mathbb{Z}\)-action?
\end{openproblem}

\plainproblemsection

\begin{openproblem}
	Resolve the following problems concerning topological orbit equivalence:
	\begin{enumerate}[label=\textup{(\arabic*)}]
		\item Classify irreducible shifts of finite type up to topological orbit
		equivalence.
		\item Characterize the unital ordered groups that occur as ordered
		cohomology groups of homeomorphisms of subshifts, irreducible shifts
		of finite type, or compact zero-dimensional spaces.
		\item Classify the unital ordered cohomology groups of irreducible shifts
		of finite type and sofic systems.
		\item For a homeomorphism \(T\) of a zero-dimensional compact space
		\(X\), the cohomology group
		\[
		C(X,\mathbb{Z})/(I-T)C(X,\mathbb{Z})
		\]
		has both the standard order and the winding order.  When do these
		orders agree?
		\item If an irreducible shift of finite type is topologically orbit
		equivalent to a subshift \(S\), must \(S\) be a shift of finite type?
	\end{enumerate}
\end{openproblem}

\plainproblemsection

\begin{openproblem}
	For \(1<r<\infty\), can a \(C^r\) diffeomorphism of a compact
	Riemannian manifold have infinite symbolic extension entropy?
\end{openproblem}

\plainproblemsection

\begin{openproblem}
	For each \(1\leq r\leq\infty\), characterize the entropy structures
	that can occur for \(C^r\) diffeomorphisms of compact Riemannian
	manifolds.
\end{openproblem}

\plainproblemsection

\begin{openproblem}
	Let \(A\) be a finite alphabet and let
	\[
	f\colon A^{\mathbb{Z}^d}\to A^{\mathbb{Z}^d}
	\]
	be a surjective \(d\)-dimensional cellular automaton.  A point is
	temporally periodic if it is periodic under \(f\), spatially periodic if
	its shift orbit is finite, and jointly periodic if both conditions hold.
	
	Must the temporally periodic points be dense?  Must the jointly periodic
	points be dense?
\end{openproblem}

\plainproblemsection

\begin{openproblem}
	For every surjective one-dimensional cellular automaton, are the jointly
	periodic points dense?
\end{openproblem}

\plainproblemsection

\begin{openproblem}
	Let \(S_N\) be the full one-dimensional shift on \(N\) symbols, and let
	\(f\colon S_N\to S_N\) be a cellular automaton.  Let \(P_k(f)\) be the
	number of points fixed by the \(k\)-th power of the shift that are also
	periodic under \(f\), and put
	\[
	\nu_k(f,S_N)=P_k(f)^{1/k},
	\qquad
	\nu(f,S_N)=\limsup_{k\to\infty}\nu_k(f,S_N).
	\]
	If \(f\) is surjective, must \(\nu(f,S_N)>1\)?
\end{openproblem}

\plainproblemsection

\begin{openproblem}
	With \(\nu(f,S_N)\) as defined above, is
	\[
	\nu(f,S_N)\geq\sqrt N
	\]
	for every surjective one-dimensional cellular automaton on \(N\)
	symbols?
\end{openproblem}

\plainproblemsection

\begin{openproblem}
	Does there exist \(N>1\) and a surjective one-dimensional cellular
	automaton \(f\) on \(N\) symbols such that
	\[
	\nu(f,S_N)<N?
	\]
\end{openproblem}

\plainproblemsection

\begin{openproblem}
	Let \(f\) be a \(d\)-dimensional cellular automaton with \(d\geq2\).
	Let \(P\) be the set of spatially periodic points and
	\(f_P=f|_P\).  If the automaton has a quiescent state \(q\), let \(F\)
	be the set of configurations equal to \(q\) outside a finite set and
	let \(f_F=f|_F\).
	
	Determine whether each of the following implications holds:
	\begin{enumerate}[label=\textup{(\arabic*)}]
		\item injectivity of \(f_P\) implies surjectivity of \(f_F\);
		\item surjectivity of \(f_F\) implies surjectivity of \(f_P\);
		\item surjectivity of \(f\) implies surjectivity of \(f_P\).
	\end{enumerate}
\end{openproblem}

\plainproblemsection

\begin{openproblem}
	(M. Gromov) A countable discrete group \(G\) is called
	\emph{surjunctive} if, for every finite set \(A\), every continuous
	shift-commuting map
	\[
	A^G\longrightarrow A^G
	\]
	is either surjective or noninjective.  Which groups are surjunctive?
\end{openproblem}

\plainproblemsection

\begin{openproblem}
	Two Borel systems \(S,T\) of equal entropy are called entropy conjugate
	if there is \(0\leq a<h(S)=h(T)\) and a Borel isomorphism inducing a
	bijection between their invariant ergodic measures of entropy at least
	\(a\), with a measurable conjugacy on each corresponding measured
	system.
	
	If \(S\) and \(T\) are positive recurrent, mixing, countable-state
	Markov shifts of equal entropy, must they be entropy conjugate?
\end{openproblem}

\plainproblemsection

\begin{openproblem}
	(K. Schmidt) For \(\beta>1\), define
	\[
	T_\beta(x)=\{\beta x\},
	\qquad 0\leq x<1,
	\]
	where \(\{\cdot\}\) denotes fractional part, and let
	\(\operatorname{Per}(\beta)\) be the set of periodic points of
	\(T_\beta\).  If \(\beta\) is a Salem number, prove that
	\[
	\operatorname{Per}(\beta)=\mathbb{Q}\cap[0,1).
	\]
\end{openproblem}

\plainproblemsection

\begin{openproblem}
	(K. Thomsen) A dynamical system is intrinsically ergodic if it has a
	unique measure of maximal entropy.  Must every subshift factor of a
	\(\beta\)-shift be intrinsically ergodic?
\end{openproblem}

\plainproblemsection

\begin{openproblem}
	(R. Adler) A renewal shift is a subshift whose points are all
	bi-infinite concatenations of words from a fixed finite collection.
	Is every irreducible sofic shift topologically conjugate to a renewal
	shift?
\end{openproblem}

\plainproblemsection

\begin{openproblem}
	(B. Parry) Let \(G\) be a finite abelian group, let \(\sigma_A\) be an
	irreducible shift of finite type on \(X_A\), and let
	\(f\colon X_A\to G\) be continuous.  The associated skew product is
	\[
	(x,g)\longmapsto(\sigma_Ax,g+f(x))
	\qquad\text{on }X_A\times G.
	\]
	If \(B\) is a corresponding matrix over the positive cone
	\(\mathbb{Z}_+G\) of the integral group ring, define
	\[
	\zeta_f(z)
	=
	\exp\left(
	\sum_{n=1}^{\infty}\frac{\operatorname{tr}(B^n)}{n}z^n
	\right)
	=
	\frac{1}{\det(I-zB)}.
	\]
	Can there be an infinite family of pairwise nonisomorphic skew products
	over one irreducible shift of finite type for which all the functions
	\(\zeta_f\) are equal?
\end{openproblem}

\plainproblemsection

\begin{openproblem}
	Classify two-sided sofic shifts up to topological conjugacy.
\end{openproblem}

\plainproblemsection

\begin{openproblem}
	(K. Matsumoto) A canonical matrix system \((M,I)\) of a subshift
	defines a dimension group \(Z_I\) and an endomorphism
	\(\lambda\colon Z_I\to Z_I\).  Put
	\[
	K_0(M,I)=\operatorname{coker}(I-\lambda),
	\qquad
	K_1(M,I)=\ker(I-\lambda).
	\]
	Which pairs of abelian groups occur as
	\(\bigl(K_0(M,I),K_1(M,I)\bigr)\) for canonical matrix systems of
	subshifts?
\end{openproblem}

\plainproblemsection

\begin{openproblem}
	(K. Matsumoto) Classify subshifts up to flow equivalence.
\end{openproblem}

\plainproblemsection

\begin{openproblem}
	Let \(\tau\) be a substitution on a finite alphabet with substitution
	matrix \(A\).  The substitution is irreducible Pisot if the
	characteristic polynomial of \(A\) is irreducible, has one root greater
	than \(1\), and all other roots have modulus strictly less than \(1\).
	The Pisot conjecture asserts that the substitutive shift associated to
	every irreducible Pisot substitution has pure discrete spectrum.
\end{openproblem}

\plainproblemsection

\begin{openproblem}
	(M. Nivat) Let \(A\) be a nonempty finite alphabet and
	\(x\in A^{\mathbb{Z}^2}\).  Let \(N_x(n_1,n_2)\) be the number of
	distinct \(n_1\times n_2\) rectangular patterns occurring in \(x\).
	If there are positive integers \(n_1,n_2\) such that
	\[
	N_x(n_1,n_2)\leq n_1n_2,
	\]
	must \(x\) be periodic?
\end{openproblem}

\subsection*{References}

\begin{enumerate}[label={[\arabic*]},leftmargin=*,itemsep=0.4em]
	\item M. Boyle, ``Open problems in symbolic dynamics,'' in
	\emph{Geometric and Probabilistic Structures in Dynamics},
	Contemporary Mathematics 469, American Mathematical Society, 2008,
	69--118.
	\item J. van den Berg and J. E. Steif, ``On the existence and
	nonexistence of finitary codings for a class of random fields,''
	\emph{Annals of Probability} \textbf{27} (1999), no.~3, 1501--1522.
	\item M. Boyle, ``Symbolic dynamics and matrices,'' in
	\emph{Combinatorial and Graph-Theoretical Problems in Linear Algebra},
	IMA Volumes in Mathematics and its Applications 50, 1993, 1--38.
	\item M. Boyle and D. Handelman, ``The spectra of nonnegative matrices
	via symbolic dynamics,'' \emph{Annals of Mathematics} \textbf{133}
	(1991), 249--316.
	\item M. Boyle and B. Kitchens, ``Periodic points for onto cellular
	automata,'' \emph{Indagationes Mathematicae} (N.S.) \textbf{10}
	(1999), no.~4, 483--493.
	\item M. Boyle and B. Lee, ``Jointly periodic points in cellular
	automata: computer explorations and conjectures,''
	\emph{Experimental Mathematics}, to appear.
	\item M. Boyle and D. Lind, ``Expansive subdynamics,''
	\emph{Transactions of the American Mathematical Society}
	\textbf{349} (1997), no.~1, 55--102.
	\item M. Boyle and A. Maass, ``Expansive invertible one-sided cellular
	automata,'' \emph{Journal of the Mathematical Society of Japan}
	\textbf{52} (2000), no.~4, 725--740; erratum,
	\emph{Journal of the Mathematical Society of Japan} \textbf{56}.
	\item T. Downarowicz and S. Newhouse, ``Symbolic extensions in smooth
	dynamical systems,'' \emph{Inventiones Mathematicae} \textbf{160}
	(2005), no.~3, 453--499.
	\item M. Gromov, ``Endomorphisms of symbolic algebraic varieties,''
	\emph{Journal of the European Mathematical Society} \textbf{1}
	(1999), no.~2, 109--197.
	\item B. Hasselblatt, ``Problems in dynamical systems and related
	topics,'' in \emph{Dynamics, Ergodic Theory, and Geometry}, MSRI
	Publications 54, 2007, 273--325.
	\item J. Kari, ``Theory of cellular automata: a survey,''
	\emph{Theoretical Computer Science} \textbf{334} (2005), no.~1--3,
	3--33.
	\item K. Matsumoto, ``Presentations of subshifts and their topological
	conjugacy invariants,'' \emph{Documenta Mathematica} \textbf{4}
	(1999), 285--340.
	\item M. Nasu, \emph{Textile Systems for Endomorphisms and
		Automorphisms of the Shift}, Memoirs of the American Mathematical
	Society 114, no.~546, 1995.
	\item M. Nivat, invited talk at ICALP, Bologna, 1997.
	\item K. Schmidt, ``On periodic expansions of Pisot numbers and Salem
	numbers,'' \emph{Bulletin of the London Mathematical Society}
	\textbf{12} (1980), no.~4, 269--278.
	\item K. Thomsen, ``On the structure of beta shifts,'' in
	\emph{Algebraic and Topological Dynamics}, Contemporary Mathematics
	385, American Mathematical Society, Providence, RI, 2005, 321--332.
	\item R. F. Williams, ``Classification of subshifts of finite type,''
	\emph{Annals of Mathematics} \textbf{98} (1973), 120--153; erratum,
	\emph{Annals of Mathematics} \textbf{99} (1974), 380--381.
\end{enumerate}



\chapter{Number Theory and Additive Combinatorics}
\clearpage
\paragraph{Notation.}
The source writes
\[
[N]=\{1,\ldots,N\}.
\]
If $A$ and $B$ are subsets of an abelian group, then
\[
A+B=\{a+b:a\in A,\ b\in B\}.
\]
For real-valued quantities $X$ and $Y$, both $X\ll Y$ and $X=O(Y)$ mean
$|X|\leq CY$ for some absolute constant $C$.  These conventions were repeated
in every source fragment and are recorded here once.

\section{Sum-Free Sets and Additive Configurations}

\plainproblemsection

\begin{openproblem}
Let $A$ be a set of $n$ positive integers. Does $A$ contain a sum-free set of size at least $n/3 + \Omega \left( n\right)$ , where $\Omega \left( n\right)  \rightarrow  \infty$ as $n \rightarrow  \infty$ ?
\paragraph{Additional information.} This is a pretty old and increasingly notorious problem, first mentioned over 50 years ago in [99] . The best known bounds are in Bourgain's paper [39], where he shows that there is necessarily a sum-free set of size $\frac{n + 2}{3}$ . In fact, he shows that there is a sum-free set of size at least

\[
\frac{n}{3} + \frac{c}{\log n}{\begin{Vmatrix}{1}_{A}\end{Vmatrix}}_{{\ell }^{1}}
\]

where

\[
{\begin{Vmatrix}{1}_{A}\end{Vmatrix}}_{{\widehat{\ell }}^{1}} \mathrel{\text{:=}} {\int }_{0}^{1}\left| {\mathop{\sum }\limits_{{a \in  A}}{e}^{2\pi ia\theta }}\right| {d\theta }.
\]

Thus a structural description of sets for which ${\begin{Vmatrix}{1}_{A}\end{Vmatrix}}_{{\widehat{\ell }}^{1}} \leq  K\log n$ would be very relevant. In fact, Eberhard, Manners and I (unpublished) worked out a proof that Problem 1 has a positive solution assuming that any set with this property has symmetric difference $o\left( n\right)$ with a $\pm$ union of ${O}_{K}\left( 1\right)$ progressions.

Let $f : \mathbf{R}/\mathbf{Z} \rightarrow  \mathbf{R}$ be the function taking the value -1 on $\left\lbrack  {\frac{1}{3},\frac{2}{3}}\right\rbrack$ and $\frac{1}{2}$ elsewhere. Does there exist some $\theta  \in  \mathbf{R}/\mathbf{Z}$ such that

(1.1)

\[
\mathop{\sum }\limits_{{a \in  A}}f\left( {a\theta }\right)  <  - \Omega \left( n\right) ?
\]

(Here $\Omega \left( n\right)  \rightarrow  \infty$ as $n \rightarrow  \infty$ .) Is it in fact the case that

(1.2)

\[
{\int }_{0}^{1}\left| {\mathop{\sum }\limits_{{a \in  A}}f\left( {a\theta }\right) }\right| {d\theta } > \Omega \left( n\right) ?
\]

Observe that (1.2) implies (1.1), since ${\int }_{0}^{1}\mathop{\sum }\limits_{{a \in  A}}f\left( {a\theta }\right) {d\theta } = 0$ , and that (1.1) implies a positive solution to Problem 1, since the set $\{ a \in  A : f\left( {a\theta }\right)  =  - 1\}$ is sum-free.

Eberhard, Manners and I convinced ourselves, using Bourgain's ideas, that (1.2) is also a consequence of a sufficiently good structural understanding of sets for which ${\begin{Vmatrix}{1}_{A}\end{Vmatrix}}_{{\widehat{\ell }}^{1}} \leq  K\log n$ .

Note that $f$ is piecewise constant, not identically zero and $\int f = 0$ . Perhaps these are the only relevant features. If so, it is perhaps slightly more natural to ask about (1.1) and (1.2) with $f : \mathbf{R}/\mathbf{Z} \rightarrow  \mathbf{R}$ instead being the function taking the value -1 on $\left\lbrack  {\frac{1}{4},\frac{3}{4}}\right\rbrack$ and 1 elsewhere, though this problem is not applicable to Problem 1. It should also be noted that Bourgain [39] did (in a rather deep argument) make progress for the function $f\left( x\right)  = \frac{1}{4} - {1}_{\left\lbrack  \frac{1}{8},\frac{3}{8}\right\rbrack  }$ , but made crucial use of the asymmetry of this function.

Returning to Problem 1 itself, it is known $\left\lbrack  {{85},{88}}\right\rbrack$ that there do exist sets with no sum-free set of size larger than $\left( {\frac{1}{3} + o\left( 1\right) }\right) n$ . However, the $o\left( 1\right)$ term in these results is more-or-less ineffective; it would be interesting to get a reasonable bound.
\end{openproblem}

\plainproblemsection

\begin{openproblem}
Let $A \subset  \mathbf{Z}$ be a set of $n$ integers. Is there a set $S \subset  A$ of size ${\left( \log n\right) }^{100}$ with $S\widehat{ + }S$ disjoint from $A$ ?
\paragraph{Additional information.} Here $S\widehat{ + }S$ denotes the restricted sumset $\left\{  {{s}_{1} + {s}_{2} : {s}_{1},{s}_{2} \in  S,{s}_{1} \neq  {s}_{2}}\right\}$ . Problems of this type are also at least 50 years old, being once again mentioned in [99] (and attributed to joint discussions of Erdős and Moser). It is known from very recent work of Sanders [251] that there is always such an $S$ with $\left| S\right|  \geq  {\left( \log n\right) }^{1 + c}$ . By contrast the best-known upper bound is due to Ruzsa [243], showing that one cannot in general hope to take $\left| S\right|$ bigger than ${e}^{C\sqrt{\log n}}$ .
\end{openproblem}

\plainproblemsection

\begin{openproblem}
If $A$ is a set of $n$ integers, what is the maximum number of affine translates of the set $\{ 0,1,3\}$ that can $A$ contain?
\paragraph{Additional information.} This problem was apparently raised by Ganguly. I heard it from Robin Pemantle. It seems likely that the answer is $\left( {\frac{1}{3} + o\left( 1\right) }\right) {n}^{2}$ , but this is not known. Aaronson's paper [2] should be consulted for more information.
\end{openproblem}

\plainproblemsection

\begin{openproblem}
Suppose that $A \subset  \left\lbrack  {0,1}\right\rbrack$ is open and has measure greater than $\frac{1}{3}$ . Is there a solution to ${xy} = z$ ?
\paragraph{Additional information.} I do not have a reference for this problem; it was suggested during the work that led to [88]. For any $\theta$ , the set ${A}_{\theta } \mathrel{\text{:=}} \left\{  {x : \theta \log x \in  \left( {\frac{1}{3},\frac{2}{3}}\right) \left( {\;\operatorname{mod}\;1}\right) }\right\}$ is product-free, and $\mu \left( {A}_{\theta }\right)  = \frac{t}{{t}^{2} + t + 1}$ where $t = {e}^{1/{3\theta }}$ . This tends to $\frac{1}{3}$ as $\theta  \rightarrow  \infty$ , so $\frac{1}{3}$ cannot be replaced by any smaller number.
\end{openproblem}

\plainproblemsection

\begin{openproblem}
For which values of $k$ is the following true: whenever we partition $\left\lbrack  N\right\rbrack   = {A}_{1} \cup  \cdots  \cup  {A}_{k},\allowbreak\left| {\mathop{\bigcup }\limits_{{i = 1}}^{k}\left( {{A}_{i}\widehat{ + }{A}_{i}}\right) }\right|  \geq  \frac{1}{10}N?$
\paragraph{Additional information.} So far as I know, all that is known is the following: for $k \ll  \log \log N$ , this is true, whilst for $k \gg  N/\log N$ , it need not be. These results are contained in a 1989 paper of Erdős, Sárközy and Sós [107]. The truth of this statement even for constant $k$ answered a question of Roth from at least 30 or so years earlier. Ruzsa [241] constructs partitions ${A}_{1} \cup  \cdots  \cup  {A}_{k}$ of an $N$ -element set, $k = N/m$ , for which the union $\left| {\mathop{\bigcup }\limits_{{i = 1}}^{k}\left( {{A}_{i}\widehat{ + }{A}_{i}}\right) }\right|$ is as small as about ${O}_{m}\left( {N}^{1/2}\right)$ . It ought to be possible to make these into partitions of $\left\lbrack  N\right\rbrack$ by a Freiman isomorphism/random covering argument. Ruzsa does not clarify the $m$ -dependence in his article, though this should not be too hard to do if desired.
\end{openproblem}

\plainproblemsection

\begin{openproblem}
Fix an integer $d$ . What is the largest sum-free subset of ${\left\lbrack  N\right\rbrack  }^{d}$ ?
\paragraph{Additional information.} This problem was raised in an unpublished note of Peter Cameron from 2002 [53], as well as (in the case $d = 2$ ) in [54, Problem 424], where it is attributed to Harut Aydinian. Writing ${f}_{d}\left( N\right)$ for the quantity in question, the most interesting part of the question seems to be determining ${c}_{d} \mathrel{\text{:=}} \mathop{\limsup }\limits_{{N \rightarrow  \infty }}{N}^{-d}{f}_{d}\left( N\right)$ . It is very well-known and easy that ${c}_{1} = \frac{1}{2}$ . Elsholtz and Rackham [96] showed that ${c}_{2} = \frac{3}{5}$ , and gave lower bounds for larger $d$ . In general, it seems to be believed that a ’slice’ $\left\{  {\left( {{x}_{1},\ldots ,{x}_{d}}\right)  : u \leq  {x}_{1} + \cdots  + {x}_{d} < {2u}}\right\}$ is optimal, where $u$ is chosen to maximise the number of points in this set.

Update 2023. Lepsveridze and Sun [177] have determined ${c}_{3},{c}_{4}$ and ${c}_{5}$ , and have confirmed that the slice example mentioned above is asymptotically optimal in these cases. They also raise the (strictly easier) question of determining the measure of the largest measurable sum-free subset of ${\left\lbrack  0,1\right\rbrack  }^{d}$ , and solve this question for $d \leq  5$ .
\end{openproblem}

\plainproblemsection

\begin{openproblem}
Is there $c > 0$ with the following property: whenever $A \subset  \left\lbrack  N\right\rbrack$ is a set of size ${N}^{1 - c}, A - A$ contains a nonzero square? What about $A - A$ containing a prime minus one? Comments. Both problems were originally considered by Sárközy.

It is known [226] that a condition $\left| A\right|  > N{\left( \log N\right) }^{-\omega \left( N\right) },\omega \left( N\right)  \rightarrow  \infty$ , is sufficient for $A - A$ to contain a square.

Update 2020. Bloom and Maynard [31] have shown that we may take $\omega \left( N\right)  =$ $c\log \log \log N$ here, which is currently the best known bound. In the other direction, Ruzsa [237] showed that for $A - A$ to contain a square one cannot have $c > {0.267}$ . Most likely, this is closer to the truth; that is, if $A - A$ contains no square then $\left| A\right|  \leq  {N}^{1 - c}$ for some $c > 0$ . This is a challenging open question, and even the (presumably much easier) analogue for subsets of $\mathbf{Z}/q\mathbf{Z}$ is not known uniformly in $q$ , the problematic cases being when $q$ is a product of primes $\equiv  3\left( {\;\operatorname{mod}\;4}\right)$ . For more on this, see [112].

Turning to shifted primes, Ruzsa and Sanders [244] show that $A - A$ contains $p - 1$ under the weaker condition $\left| A\right|  > N\exp \left( {-c{\left( \log N\right) }^{1/4}}\right)$ , and in June 2019 R. Wang [282] improved the exponent here to $1/3$ . The biggest set known with $A - A$ not containing $p - 1$ has, I believe, size ${N}^{C/\log \log N} = {N}^{o\left( 1\right) }$ ; see [236]. Update 2022. I [142] obtained an upper bound of shape $\left| A\right|  \ll  {N}^{1 - c}$ for this problem, for some $c > 0$ . Thorner and Zaman [275] showed that $c = {10}^{-{18}}$ is permissible, and I [143] showed that one can take $c = \frac{1}{12} - o\left( 1\right)$ assuming GRH.

Returning to the squares, the following presumably easier question is also open: if $A \subset  \left\lbrack  N\right\rbrack$ is a set of size ${N}^{1 - c}$ , does ${100A} - {100A}$ contain a nonzero square?

In the function field setting, I obtained polynomially effective bounds for squares using the polynomial method of Croot-Lev-Pach [138]. However, the method does not work for primes (irreducibles). Thus, I believe the following is open: let $A$ be a set consisting of ${\left( {1.999}\right) }^{n}$ polynomials over ${\mathbf{F}}_{2}$ , all of degree $\leq  n$ . Do there exist ${a}_{1},{a}_{2} \in  A$ such that ${a}_{1} - {a}_{2} - 1$ is irreducible? Update 2022. Whilst this is still open, it ought to follow from the methods of [142].
\end{openproblem}

\plainproblemsection

\begin{openproblem}
Define Ulam’s sequence $1,2,3,4,6,8,{11},{13},{16},{18},{26},{28},{36},\ldots$ as follows: ${u}_{1} = 1,{u}_{2} = 2$ , and ${u}_{n + 1}$ is the smallest number expressible as a sum ${u}_{i} + {u}_{j}$ , $i < j \leq  n$ , in a unique way. Does this sequence have positive density? Can one explain the curious Fourier properties of this sequence?
\paragraph{Additional information.} The first problem was asked by Ulam [276]. The second problem here refers to the observations of Steinerberger [265], where data is presented to suggest that there is a number $\alpha  \approx  {2.5714474995}$ such that the sequence $\alpha {u}_{n}\left( {{\;\operatorname{mod}\;2}\pi }\right)$ has an interesting distribution function. Steinerberger's paper may be consulted for more information.

A (somewhat) related class of sequence may be defined as follows. Start with some ’seed’ set of values ${u}_{1} < \cdots  < {u}_{m}$ . For $n \geq  m$ , define ${u}_{n + 1}$ to be the smallest integer greater than ${u}_{n}$ not of the form ${u}_{i} + {u}_{j}$ for $i, j \leq  n$ . This construction was considered by Dickson. A variant construction considered by Queneau is to define ${u}_{n + 1}$ to be the smallest integer greater than ${u}_{n}$ and not of the form ${u}_{i} + {u}_{j}$ with $i < j \leq  n$ . For both variants, the following question appears to be open: is ${u}_{n + 1} - {u}_{n}$ eventually periodic? For references and computations showing that periodicity (if present) can be very slow to appear, see [111, 152].
\end{openproblem}

\plainproblemsection

\begin{openproblem}
Suppose that $A \subset  \left\lbrack  N\right\rbrack$ has no more than $\varepsilon {N}^{2}$ solutions to $x + y = z$ . It is known [134] that one may remove ${\varepsilon }^{\prime }N$ elements of $A$ to leave a sum-free set, where ${\varepsilon }^{\prime } \rightarrow  0$ as $\varepsilon  \rightarrow  0$ . But is there a bound for ${\varepsilon }^{\prime }$ in terms of $\varepsilon$ which is at least vaguely reasonable?
\paragraph{Additional information.} This concerns the so-called arithmetic removal lemma. Until a few years ago I would have said that one should first consider the model problem over ${\mathbf{F}}_{2}^{n}$ . However, Fox and L. M. Lovász [115] have used the Croot-Lev-Pach method to show that in that case ${\varepsilon }^{\prime }$ can be taken to be ${\varepsilon }^{C}$ , and even obtain the optimal value of $C$ . By contrast it is known that in $\left\lbrack  N\right\rbrack$ there cannot be a polynomial bound. Aaronson [1] has shown, adapting the ideas of Fox and Lovász, that the problem is essentially equivalent to determining the size of the largest additive matching in $\mathbf{Z}/N\mathbf{Z}$ , that is to say collection of triples ${\left( {x}_{i},{y}_{i},{z}_{i}\right) }_{i = 1,\ldots , t}$ with ${x}_{i} + {y}_{j} + {z}_{k} = 0$ iff $i = j = k$ .

It would be remiss not to mention the better-known cousin of this problem, the triangle-removal problem from graph theory. Suppose that a graph on $N$ vertices has $\varepsilon {N}^{3}$ triangles. Then Ruzsa and Szemerédi showed that one may remove ${\varepsilon }^{\prime }{N}^{2}$ edges from the graph so that the remaining graph is triangle-free, where ${\varepsilon }^{\prime } \rightarrow  0$ as $\varepsilon  \rightarrow  0$ . Is there a bound for ${\varepsilon }^{\prime }$ in terms of $\varepsilon$ which is vaguely reasonable? See [114] for the best known bounds, which are of 'tower' type (but with a relatively small tower height).
\end{openproblem}

\plainproblemsection

\begin{openproblem}
Let $A \subset  \mathbf{R}$ be a set of positive measure. Does $A$ contain an affine copy of $\left\{  {1,\frac{1}{2},\frac{1}{4},\ldots }\right\}$ ?
\paragraph{Additional information.} This is a special case of the Erdős similarity problem; see [267] for a survey.
\end{openproblem}

\clearpage
\section{Arithmetic Progressions and Additive Ramsey Theory}

In this section, $r_k(N)$ denotes the largest subset of
$[N]$ containing no nontrivial $k$-term arithmetic progression, and $r_k(G)$
denotes the analogous quantity for a finite abelian group $G$.

\plainproblemsection

\begin{openproblem}
Let $S \subset  \mathbf{N}$ be random. Under what conditions is Roth’s theorem for progressions of length 3 true with common differences in $S$ ?
\paragraph{Additional information.} Of course, one may also ask the question for longer progressions. A set $S$ for which the conclusion of Szemerédi’s theorem for progressions of length $k$ holds, but with the common difference of the progressions restricted to lie in $S$ , is called a set of(k - 1)-recurrence. Thus Problem 10 is asking for conditions under which a random set is a set of 2-recurrence.

The most optimistic conjecture here, suggested in [119] and [118, Problem 31] is that if $n$ is chosen to lie in $S$ with probability $\omega \left( n\right) /n$ , where $\omega \left( n\right)  \rightarrow  \infty$ , then $S$ is almost surely a set of $r$ -recurrence for all $r \geq  1$ .

The best that is known is the result of Briët and Gopi [47], improving on earlier work of Christ and of Frantzikinakis-Lesigne-Wierdl, showing that if $\mathbf{P}\left( {n \in  S}\right)  \gg$ $\omega \left( n\right) \frac{\log n}{{n}^{1/\left\lbrack  {k/2}\right\rbrack  }}$ then Szemerédi’s theorem for progressions of length $k$ holds (almost surely as $n \rightarrow  \infty$ ) with common differences in $S$ .

The case $k = 3$ (Roth’s theorem with differences in $S$ ) is perhaps the most interesting in the first instance, and here the bound of Briët and Gopi is no stronger than the earlier results. A particular model problem would be to show that if $\mathbf{P}\left( {n \in  S}\right)  = {n}^{-{0.51}}$ then Roth’s theorem holds with common difference in $S$ .

Update 2023. Briët and Castro-Silva [46] have advanced the bounds for $k$ odd, showing that if $\mathbf{P}\left( {n \in  S}\right)  \gg  {n}^{-2/k + o\left( 1\right) }$ then Szemerédi’s theorem for progressions of length $k$ holds (almost surely) with common differences in $S$ . In particular, they resolve the model problem mentioned here.

It is a curious fact that, by the Bergelson-Leibman theorem, Roth's theorem is known with common differences restricted to much sparser sets $S$ , for example $S = \left\{  {{1}^{3},{2}^{3},{3}^{3},{4}^{3},\ldots }\right\}$ . However, the proof of that theorem seems to crucially exploit the fact that the derivatives of the sequence are eventually constant.

The conjecture in $\left\lbrack  {{118},{119}}\right\rbrack$ is motivated by the suggestion that the set of correlation functions of the form $\sigma \left( d\right)  \mathrel{\text{:=}} {\mathbf{E}}_{n}{1}_{A}\left( n\right) {1}_{A}\left( {n + d}\right) {1}_{A}\left( {n + {2d}}\right)$ (where here we are taking $k = 3$ for simplicity) should have ’low complexity’, since one expects $\sigma \left( d\right)$ to be approximated in some sense by highly algebraic objects called nilsequences. However, recent examples of Briët-Labib [49] and Briët and me [48] show that the situation in that regard is more complicated than previously thought. Those examples do not, however, seem to impact on Problem 10 itself.

Finally we remark that Altman [10] has also shown that in finite field models the random set $S$ must be significantly larger for Roth’s theorem with common differences in $S$ to hold.
\end{openproblem}

\plainproblemsection

\begin{openproblem}
Is every set \(A\subset\{1,\ldots,N\}\) with
\[
  |A|\geq N^{0.99}
\]
guaranteed to contain a pair \(x,x+d^{2}\) with \(d\neq0\)?
\end{openproblem}

\plainproblemsection

\begin{openproblem}
Let $G$ be an abelian group of size $N$ , and suppose that $A \subset  G$ has density $\alpha$ . Are there at least ${\alpha }^{15}{N}^{10}$ tuples $\left( {{x}_{1},\ldots ,{x}_{5},{y}_{1},\ldots ,{y}_{5}}\right)  \in  {G}^{10}$ such that ${x}_{i} + {y}_{j} \in  A$ whenever $j \in  \{ i, i + 1, i + 2\}$ ?
\paragraph{Additional information.} This is very closely related to the Cayley graph case of simplest unknown instance of Sidorenko's conjecture for graphs, that of the 'Möbius ladder' ${K}_{5,5} \smallsetminus  {C}_{10}$ .
\end{openproblem}

\plainproblemsection

\begin{openproblem}
(4-term APs in uniform sets.). Suppose that $A \subset  \mathbf{Z}/N\mathbf{Z}$ has density $\alpha$ and is Fourier uniform (that is, all Fourier coefficients of ${1}_{A} - \alpha$ are $o\left( N\right)$ ). Does $A$ contain at least $\gg  {\alpha }^{100}{N}^{2}$ 4-term arithmetic progressions?
\paragraph{Additional information.} Gowers [126] has shown that this is not true if 100 is replaced by 4.01. In view of the ’arithmetic regularity lemma’ it is natural to consider sets $A$ coming from a 2-step nilsequence, that is to say $A$ has the form $\left\{  {n : {g}^{n} \in  S}\right\}$ where $g \in  G$ and $S \subset  \Gamma  \smallsetminus  G$ is an open set.

Update 2023. This problem has been considered in an interesting paper of Deng, Tidor and Zhao [82]. First of all, they improve the exponent 4.01 mentioned above to $3 + \frac{1}{2}{\log }_{3}\left( {22}\right)  \approx  {4.406}$ . More interestingly, they make the following conjecture, and show that it implies a negative answer to the main problem.

Conjecture [82, Conjecture 1.6]. For all $N$ , there is a colouring of $\{ 1,\ldots , N\}$ with ${N}^{o\left( 1\right) }$ colours with no symmetrically coloured 4-term progression (that is, progression $x, x + d, x + {2d}, x + {3d}$ in which $x, x + {3d}$ have the same colour, and $x + d, x + {2d}$ have the same colour).
\end{openproblem}

\plainproblemsection

\begin{openproblem}
Define \(W(3,r)\) to be the least integer \(N\) such that every red--blue
colouring of \(\{1,\ldots,N\}\) contains either a red \(3\)-term arithmetic
progression or a blue \(r\)-term arithmetic progression.

Construct explicitly, for some \(r\), a colouring showing that
\[
  W(3,r)\geq 2r^{2}.
\]
\end{openproblem}

\plainproblemsection

\begin{openproblem}
Does there exist a Lipschitz function $f : \mathbf{N} \rightarrow  \mathbf{Z}$ whose graph $\Gamma  = \{ \left( {n, f\left( n\right) }\right)  : n \in  \mathbf{Z}\}  \subset  {\mathbf{Z}}^{2}$ is free of 3-term progressions?
\paragraph{Additional information.} The same question but with the domain restricted to $\left\lbrack  N\right\rbrack$ (for arbitrary large $N$ , and with the Lipschitz constant $\mathop{\max }\limits_{n}\left| {f\left( {n + 1}\right)  - f\left( n\right) }\right|$ not depending on $N$ ) is still interesting. I first heard this question from Jacob Fox in 2005, then again from Natasha Morrison in around 2014. It has recently been considered in print by Brown, Jungić and Poelstra [51].

The answer is yes for 4-term progressions. This is established by Cassaigne, Currie, Schaeffer and Shallit [58], who write in the language of combinatorics on words. They construct an infinite word ${x}_{1}{x}_{2}{x}_{3}\cdots$ in the alphabet $\{ 0,1,3,4\}$ which avoids ’additive cubes’, by which they mean blocks of the form ${b}_{1}{b}_{2}{b}_{3}$ where ${b}_{1},{b}_{2},{b}_{3}$ have the same length and the same sum of elements (for example, ${b}_{1} = {01334}$ , ${b}_{2} = {44030},{b}_{3} = {11144})$ . One may then define $f\left( n\right)  = {x}_{1} + \cdots  + {x}_{n}$ . In their paper, they explicitly mention that Problem 15 (formulated in terms of combinatorics on words) is unsolved.

(The finitary version of) Problem 15 has a somewhat similar flavour to the case $k = 3$ of Problem 14, in the sense that one is asked to construct a set of size $\sim  N$ , free of 3-term progressions, inside an ambient set of size $\sim  {N}^{2}$ , with certain additional properties. Once again it is tempting to start with a Behrend set and modify it, but for similar reasons this seems to be impossible.
\end{openproblem}

\plainproblemsection

\begin{openproblem}
What is the largest subset of $\left\lbrack  N\right\rbrack$ with no solution to $x + {3y} = {2z} + {2w}$ in distinct integers $x, y, z, w$ ?
\paragraph{Additional information.} This question was asked by Ruzsa [239, Section 9]. Writing $f\left( N\right)$ for the number in question, so far as I know the best bounds known are ${N}^{1/2} \ll$ $f\left( N\right)  \ll  N{e}^{-c{\left( \log N\right) }^{1/7}}$ , the lower bound being in Ruzsa’s paper and the upper bound being due to Schoen and Sisask [253] (see in particular Section 9 of their paper).

In a somewhat similar vein, Yufei Zhao (personal communication) asked me whether there is a subset of $\{ 1,\ldots , N\}$ of size ${N}^{1/3 - o\left( 1\right) }$ with no nontrivial solutions to $x + {2y} + {3z} = {x}^{\prime } + 2{y}^{\prime } + 3{z}^{\prime }$ .
\end{openproblem}

\plainproblemsection

\begin{openproblem}
Suppose that $A \subset  {\mathbf{F}}_{3}^{n}$ is a set of density $\alpha$ . Under what conditions on $\alpha$ is $A$ guaranteed to contain a 3-term progression with nonzero common difference in $\{ 0,1{\} }^{n}$ ? For fixed $\alpha$ , must $A$ contain ${ \gg  }_{\alpha }{6}^{n}$ such progressions?
\paragraph{Additional information.} My thanks go to Vsevelod Lev for reminding me of the first problem. It is known, as a consequence of the density Hales-Jewett Theorem, that if $\alpha  \rightarrow  0$ sufficiently slowly as $n \rightarrow  \infty$ then there must be such a progression. However, this may even by true with $\alpha  = {\left( 1 - c\right) }^{n}$ for some $c > 0$ . Regarding the second problem, a recent reference is [154], where it is shown that this is true for sets invariant under the permutation action of ${S}_{n}$ on ${\mathbf{F}}_{3}^{n}$ . Other related problems are also considered in that paper. For both of these questions, one may replace 3 by an arbitrary prime $q$ .

Update 2023. Bhangale, Khot and Minzer [28], with a very difficult argument, showed that a set $A \subset  {\mathbf{F}}_{p}^{n}$ free of progressions with common difference in $\{ 0,1,2{\} }^{n}$ has density ${ \ll  }_{p}{\left( \log \log \log n\right) }^{-{c}_{p}}$ .

Now we turn to a small selection of multidimensional questions.
\end{openproblem}

\plainproblemsection

\begin{openproblem}
Suppose that $G$ is a finite group, and let $A \subset  G \times  G$ be a subset of density $\alpha$ . Is it true that there are ${ \gg  }_{\alpha }{\left| G\right| }^{3}$ triples $x, y, g$ such that $\left( {x, y}\right),\allowbreak\left( {{gx}, y}\right),\allowbreak\left( {x,{gy}}\right)$ all lie in $A$ ?
\paragraph{Additional information.} This problem was raised by Tim Austin: see [13, Question 2]. Austin calls triples of the stated type naïve corners.

Another type of corner is the configuration $\{ \left( {x, y}\right),\allowbreak\left( {{xg}, y}\right),\allowbreak\left( {x,{gy}}\right) \}$ , which Austin calls a BMZ corner (after Bergelson, McCutcheon and Zhang). For these corners, Solymosi [263] has answered the corresponding question.

One may consider the situation in which $G$ is quasirandom in the sense of Gowers [125], for example $G = {\operatorname{PSL}}_{2}\left( {\mathbf{F}}_{p}\right)$ for large $p$ . It is natural to speculate that any subset of $G \times  G$ of size ${\left| G\right| }^{2 - \delta }$ (for some sufficiently small $\delta  > 0$ ) contains a nontrivial corner of both types. Austin [13, Theorem B] proves this for naïve corners, but his bounds for BMZ corners [13, Theorem C] are weaker.
\end{openproblem}

\plainproblemsection

\begin{openproblem}
Find reasonable bounds for instances of the multidimensional Sze-merédi theorem.
\paragraph{Additional information.} The minimum definition of 'reasonable' is probably a density of finite logarithmic type ${\left( {\log }_{m}N\right) }^{-1}$ , where ${\log }_{m}N$ denotes the $m$ -fold iterated logarithm of $N$ , where $N$ is the size of the underlying structure in which one is working (either $G \times  G$ for some abelian group, or $\{ 1,\ldots , n\}  \times  \{ 1,\ldots , n\}$ ). In full generality, obtaining such a bound currently seems hopeless. The most basic open problem at the moment seems to be that we do not currently have such bounds for the largest subset of ${\mathbf{F}}_{p}^{n} \times  {\mathbf{F}}_{p}^{n}$ not containing a ’square’ $\left( {x, y}\right),\allowbreak\left( {x + h, y}\right),\allowbreak\left( {x, y + h}\right),\allowbreak\left( {x + h, y + h}\right)$ . The problem of finding such bounds for 3-dimensional corners $\left( {x, y, z}\right),\allowbreak\left( {x + h, y, z}\right)$ , $\left( {x, y + h, z}\right),\allowbreak\left( {x, y, z + h}\right)$ is also open, and at least as hard.

For two-dimensional corners $\left( {x, y}\right),\allowbreak\left( {x + h, y}\right),\allowbreak\left( {x, y + h}\right)$ we have the bounds of Shkredov [258], who showed that such corners are present in any subset of ${\left\lbrack  N\right\rbrack  }^{2}$ of density ${\left( \log \log N\right) }^{-c}$ , but no lower bound better than ${e}^{-c\sqrt{\log n}}$ is known.

Beyond this, in 2022 Peluse studied ’L-shaped’ configurations $\left( {x, y}\right),\allowbreak\left( {x + h, y}\right)$ , $\left( {x + {2h}, y}\right),\allowbreak\left( {x, y + h}\right)$ , showing in [218] that a density ${ \gg  }_{p}{\left( {\log }_{m}N\right) }^{-1}$ is sufficient to guarantee such configurations in ${\mathbf{F}}_{p}^{n} \times  {\mathbf{F}}_{p}^{n}$ .

Returning to corners, Yufei Zhao has remarked to me that we do not have good bounds for the corresponding problem in which the corners need not be axis parallel. Representing the points of ${\left\lbrack  N\right\rbrack  }^{2}$ by Gaussian integers, this is equivalent to solving $\left( {x - y}\right)  = i\left( {y - z}\right)$ , and so this problem is perhaps closer in spirit to finding 3-term progressions in the integers. It ought to be possible to prove an upper bound of shape ${N}^{2}{\left( \log N\right) }^{-c}$ (in fact, an upper bound of shape ${N}^{2}{\left( \log \log N\right) }^{-c}$ is known [230] for the much harder problem of non-axis parallel squares) but it is unclear whether a lower bound of ${N}^{2 - o\left( 1\right) }$ can be expected. Update 2021: Pilatte [223] has obtained an upper bound of $\ll  {N}^{2}{\left( \log N\right) }^{-1 - c}$ for some $c > 0$ using methods of Bloom and Sisask.

Curiously, the corresponding question with ’skew corners’ $\left( {x, y}\right),\allowbreak\left( {x, y + h}\right),\allowbreak\left( x + h,{y}^{\prime } \right)$ seems completely wide-open, with Pratt [229] remarking that the best-known upper bound comes from Shkredov's work mentioned above, and that the best known lower bound is $\gg  N\left( {\log N}\right) {\left( \log \log N\right) }^{-1/2}$ by work of Petrov [287]. Pratt [229, Conjecture 1.2] asks whether an upper bound ${N}^{1 + o\left( 1\right) }$ might be true, and he links the question to Problem 36. Update 2024. Pohoata and Zakharov [227] improved the lower bound to ${N}^{5/4}$ , and shortly afterwards Beker [26] improved the lower bound to ${N}^{2}{e}^{-c\sqrt{\log N}}$ and the upper bound to $\ll  {N}^{2}{\left( \log N\right) }^{-c}$ .

Finally, Zachary Chase remarked on the following question asked by Jordan Ellenberg: what is the size of the largest subset of ${\left\lbrack  N\right\rbrack  }^{2}$ with no isosceles triangle (triangles of area zero are allowed)? It is between ${N}^{1 - o\left( 1\right) }$ (put a Behrend set on a single vertical line) and ${N}^{2 - o\left( 1\right) }$ , but nothing more seems to be known.

Finally, we turn to some questions about partition regularity.
\end{openproblem}

\plainproblemsection

\begin{openproblem}
Suppose that ${a}_{1},\ldots ,{a}_{k}$ are integers which do not satisfy Rado’s condition: thus if $\mathop{\sum }\limits_{{i \in  I}}{a}_{i} = 0$ then $I = \varnothing$ . It then follows from Rado’s theorem that the equation ${a}_{1}{x}_{1} + \cdots  + {a}_{k}{x}_{k}$ is not partition regular. Write $c\left( {{a}_{1},\ldots ,{a}_{k}}\right)$ for the least number of colours required in order to colour $N$ so that there is no monochromatic solution to ${a}_{1}{x}_{1} + \cdots  + {a}_{k}{x}_{k} = 0$ . Is $c\left( {{a}_{1},\ldots ,{a}_{k}}\right)$ bounded in terms of $k$ only?
\paragraph{Additional information.} This problem, which is known as Rado's boundedness conjecture, dates back to 1933 [232]. The answer was shown to be affirmative for $k = 3$ by Fox and Kleitman [116], who showed that $c\left( {{a}_{1},{a}_{2},{a}_{3}}\right)  \leq  {24}$ (I am not sure this is sharp; it might be interesting to determine the sharp constant). It is open for all $k \geq  4$ . Let me also highlight a question [116, Conjecture 5] of Fox and Kleitman, which they call a ’modular analogue’ of Rado’s Boundedness Conjecture. Let $p$ be a prime, and suppose that ${a}_{1},\ldots ,{a}_{k}$ are integers with $\mathop{\sum }\limits_{{i \in  I}}{a}_{i} \equiv  0\left( {\;\operatorname{mod}\;p}\right)$ only when $I = \varnothing$ . Does there exist an $f\left( k\right)$ -colouring of ${\left( \mathbf{Z}/p\mathbf{Z}\right) }^{ * }$ with no monochromatic solution to ${a}_{1}{x}_{1} + \cdots  + {a}_{k}{x}_{k} = 0$ ? This seems to be open even when $k = 3;\mathrm{I}$ suspect the answer may be negative.

Milićević [94, Conjecture 11.1] conjectures the following 2-adic variant. For any $k \in  \mathbf{N}$ , there exists $K = K\left( k\right)$ such that the following is true. Let $r$ be a positive integer, and let ${a}_{1},\ldots ,{a}_{k} \in  \mathbf{Z}/{2}^{r}\mathbf{Z}$ . Let $d$ be the largest integer such that $\mathop{\sum }\limits_{{i \in  I}}{a}_{i} \equiv$ $0\left( {\;\operatorname{mod}\;{2}^{d}}\right)$ for some non-empty subset $I \subset  \left\lbrack  k\right\rbrack$ . Then there is a $K$ -colouring of $\mathbf{Z}/{2}^{r}\mathbf{Z}$ such that all monochromatic solutions $x = \left( {{x}_{1},\ldots ,{x}_{k}}\right)$ to the equation ${a}_{1}{x}_{1} + \cdots  + {a}_{k}{x}_{k} = 0$ satisfy ${x}_{i} \equiv  0\left( {\;\operatorname{mod}\;{2}^{r - d}}\right)$ for all $i = 1,\ldots , k$ . Milićević remarks that, if true, this would imply the Rado boundedness conjecture by a compactness argument.
\end{openproblem}

\plainproblemsection

\begin{openproblem}
If $\{ 1,\ldots , N\}$ is $r$ -coloured then, for $N \geq  {N}_{0}\left( r\right)$ , there are integers $x, y \geq  3$ such that $x + y,{xy}$ have the same colour. Find reasonable bounds for ${N}_{0}\left( r\right)$ .
\paragraph{Additional information.} The existence of ${N}_{0}\left( r\right)$ was established only recently, in a celebrated paper of Moreira [210] (who in fact established that $x$ can also be the same colour). The main proof in [210] gives no bound, since it uses topological dynamics. However, the 'elementary proof' given in Section 5 of Moreira's paper does in principle give a bound, albeit an extremely weak one, not least because the known bounds for van der Waerden's theorem are so weak.

Let me also mention the famous question of Hindman: if $\mathbf{N}$ is finitely coloured, are there $x, y, x + y,{xy}$ all of the same colour? (In fact, Hindman asked whether for any $k$ there are ${x}_{1},\ldots ,{x}_{k}$ with all sums $\mathop{\sum }\limits_{{i \in  I}}{x}_{i}$ and all products $\mathop{\prod }\limits_{{i \in  I}}{x}_{i}$ the same colour for all nonempty $I \subseteq  \left\lbrack  k\right\rbrack$ , with the preceding question being the case $k = 2$ .)

Update 2022. Bowen and Sabok [41] solved the analogue of Hindman's question (in the case $k = 2$ ) where $\mathbf{N}$ is replaced by $\mathbf{Q}$ , and Bowen [40] has solved the 2-colour case in the integers, showing that if $\mathbf{N}$ is coloured red-blue then there are infinitely many pairs $x, y$ with $x, y, x + y,{xy}$ the same colour. In 2023 Alweiss [11] solved the analogue of the full Hindman question where $\mathbf{N}$ is replaced by $\mathbf{Q}$ , that is to say with $k$ arbitrary.

It seems possible that in any finite colouring of $\mathbf{N}$ , for any $k$ there are ${x}_{1},\ldots ,{x}_{k}$ with all of the elementary symmetric functions of the ${x}_{i}$ being the same colour.
\end{openproblem}

\plainproblemsection

\begin{openproblem}
Is the Pythagorean equation
\[
  x^{2}+y^{2}=z^{2}
\]
partition regular over \(\mathbb N\)? Equivalently, does every finite
colouring of \(\mathbb N\) contain a monochromatic solution
\((x,y,z)\) to this equation?
\end{openproblem}

\plainproblemsection

\begin{openproblem}
What is the smallest subset of $\mathbf{N}$ containing, for each $d = 1,\ldots , N$ , an arithmetic progression of length $k$ with common difference $d$ ?
\paragraph{Additional information.} Writing ${F}_{k}\left( N\right)$ for this quantity, it is conjectured that ${F}_{k}\left( N\right) { \gg  }_{k}$ ${N}^{1 - {c}_{k}}$ , with ${c}_{k} \rightarrow  0$ as $k \rightarrow  \infty$ . One should not expect this problem to be easy, as it was shown by Ruzsa and me [146] to be equivalent to the so-called arithmetic Kakeya conjecture of Katz and Tao (which implies the Kakeya conjecture in Euclidean space).
\end{openproblem}

\plainproblemsection

\begin{openproblem}
Problems about progressions.

(i) Is ${r}_{5}\left( N\right)  \ll  N{\left( \log N\right) }^{-c}$ ?

(ii) Is ${r}_{4}\left( {\mathbf{F}}_{5}^{n}\right)  \ll  {N}^{1 - c}$ ,
where \(N=5^{n}\)?
\paragraph{Additional information.}
Leng, Sah and Sawhney established the weaker bound
\[
  r_k(N)\ll N\exp\!\left(-(\log\log N)^{c_k}\right)
  \qquad (k\geq5).
\]

The finite-field question asks whether methods related to the
Croot--Lev--Pach and Ellenberg--Gijswijt bounds can yield the stated
estimate for four-term progressions.
\end{openproblem}

\clearpage
\section{Additive Bases, Sidon Sets, and Representation Problems}

For a subset $A$ of an abelian group, let $r_A(n)$ be the number of
representations of $n$ as a sum of two elements of $A$, with
$n=x+y=y+x$ counted once.  A Sidon set, or $B_2$-set, satisfies
$r_A(n)\leq1$ for every $n$.

\plainproblemsection

\begin{openproblem}
Write $F\left( N\right)$ for the largest Sidon subset of $\left\lbrack  N\right\rbrack$ . Improve, at least for infinitely many $N$ , the bounds ${N}^{1/2} + O\left( 1\right)  \leq  F\left( N\right)  \leq  {N}^{1/2} + {N}^{1/4} + O\left( 1\right)$ .
\paragraph{Additional information.} Erdős stated this problem many times. The upper bound is due to Lindstrom from 1939. In my first paper [130] I gave an alternative (but equivalent) proof using Fourier analysis, which I was sure could be tweaked to at least give a small improvement, but I never managed this.

Update 2021. Balogh, Füredi and Roy [15] obtained a small improvement, getting an upper bound of $F\left( N\right)  \leq  {N}^{1/2} + {0.998}{N}^{1/4}$ for large $N$ . In 2023, the constant here was further improved to 0.98183 in [57].

It might be remarked that the situation is far less clear in other settings. For example, I am fairly sure that it is not known whether or not there exists a Sidon subset of $\mathbf{Z}/p\mathbf{Z}$ of size $\left( {1 + o\left( 1\right) }\right) \sqrt{p}$ , for all $p$ , or even whether, if $G$ is an abelian group of size $n$ , there always exists a Sidon subset of $G$ of size ${0.01}\sqrt{n}$ .

Another very nice old problem is whether there is a Sidon subset of $\{ 0,1{\} }^{n}$ of size ${N}^{0.51}$ , where $N = {2}^{n}$ . The best-known upper bound, so far as I know, is ${N}^{0.5753}$ in a paper of Cohen, Litsyn and Zémor [68].

The lower bound in Problem 31 comes from taking a Sidon subset of $\mathbf{Z}/q\mathbf{Z}$ of size $\sqrt{q} + O\left( 1\right)$ (for which there are several constructions, for different $q\mathrm{\;s}$ ) and 'unwrapping'. One feels that, after a suitable dilation, it ought to be possible to do this more efficiently. Moreover, this might be quite a general phenomenon. This motivates the following question.
\end{openproblem}

\plainproblemsection

\begin{openproblem}
Let $p$ be a prime and let $A \subset  \mathbf{Z}/p\mathbf{Z}$ be a set of size $\sqrt{p}$ . Is there a dilate of $A$ with a gap of length ${100}\sqrt{p}$ ?
\paragraph{Additional information.} Shakan [257] has used the polynomial method to show that this is true with 100 replaced by 2 , but this appears to be the limit of his method.

One can formulate variants of this problem in which $\sqrt{p}$ is replaced by an arbitrary function $\omega \left( p\right)$ , and one is looking for a dilate with a gap of length ${100p}/\omega \left( p\right)$ . In the regime $\omega \left( p\right)  \sim  {cp}$ , this is Szemerédi’s theorem; in the regime $\omega \left( p\right)  \leq  c\log p$ , this is basically Dirichlet's lower bound for the size of Bohr sets. Even what happens in the regime $\omega \left( p\right)  \sim  {10}\log p$ is unclear.

Tom Sanders pointed out to me that things are much easier in the finite field model ${\mathbf{F}}_{2}^{n}$ . Indeed, if $\left| A\right|  \sim  \sqrt{N}$ , where $N = {2}^{n}$ , then by an averaging argument one may find a coset $H$ of some subspace of size $\sim  \frac{1}{10}n\sqrt{N}$ on which $A$ has at most $\frac{1}{10}n < \dim H$ points; but then these points sit inside a proper coset of $H$ , and so ${A}^{c}$ contains a coset of dimension $\dim H - 1$ , and hence of size $\geq  {100}\sqrt{N}$ provided $n$ is sufficiently large.

As regards the connection to Problem 31, so far as I am aware it is not currently known that there are infinitely many primes $p$ for which $\mathbf{Z}/p\mathbf{Z}$ admits a Sidon set of size $\sqrt{p} + O\left( 1\right)$ , thus a positive solution to this would not immediately imply a better lower bound in Problem 31.
\end{openproblem}

\plainproblemsection

\begin{openproblem}
Are there infinitely many $q$ for which there is a set $A \subset  \mathbf{Z}/q\mathbf{Z}$ , $\left| A\right|  = \left( {\sqrt{2} + o\left( 1\right) }\right) {q}^{1/2}$ , with $A + A = \mathbf{Z}/q\mathbf{Z}$ ?
\paragraph{Additional information.} Potentially yes, but I'm not sure I know how to construct such sets. One could also ask the same question but with $q = p$ restricted to be a prime. Then, the set $S = \left\{  {\left( {x,{x}^{2}}\right)  : x \in  A}\right\}   \cup  \{ \left( {1,0}\right) \}$ determines a line in every direction. It is easy to see that no set of size smaller than $\left( {\sqrt{2} + o\left( 1\right) }\right) {p}^{1/2}$ can have this property. The question of whether this bound is sharp was asked by Granville [77, Problem 5.2] and, more recently, by Caprace and de la Harpe [55, Question 6.1].

Certainly the most famous problem in the general sphere of Sidon sets is the question of Erdős and Turán. If $A \subset  \mathbf{N}$ is a set with ${r}_{A}\left( n\right)  \leq  r$ for some $r$ , is ${r}_{A}\left( n\right)  = 0$ for infinitely many $n$ ? In the contrapositive, if every sufficiently large $n$ is a sum of two elements of $A$ , are there $n$ which can be so written in arbitrarily many ways? In fact, is this true under the weaker assumption that the $n$ th element of $A$ is at most $C{n}^{2}$ ? Rather than including this famous problem as one of this list, I instead state the following vague question which most likely is the key to answering it in the affirmative, but which also seems impossibly hard.
\end{openproblem}

\plainproblemsection

\begin{openproblem}
Suppose that $A \subset  \left\lbrack  N\right\rbrack$ is a set of size $\geq  c\sqrt{N}$ for which ${r}_{A}\left( n\right)  \leq  r$ for all $n$ . What can be said about the structure of $A$ ?
\paragraph{Additional information.} Whilst all known constructions are at least somewhat algebraic, even formulating a conjecture (even in the case $c \approx  1$ and $r = 1$ ) seems hopeless. See the paper [89] of Eberhard and Manners giving an overview and unified view of the known constructions. Their paper is an enjoyable read but does not leave one optimistic about a reasonable answer to Problem 34, even in the case $r = 1$ and $c = 1 - \varepsilon .$

There are some purely real-variable questions concerning autoconvolutions which come from trying to put bounds on various generalisations of Sidon sets. In some of these questions it is quite surprising how wide apart the bounds are.
\end{openproblem}

\plainproblemsection

\begin{openproblem}
In this problem, we consider the class $\mathcal{F}$ of all integrable functions $f : \left\lbrack  {0,1}\right\rbrack   \rightarrow  {\mathbf{R}}_{ \geq  0}$ with $\int f = 1$ . Let $1 < p \leq  \infty$ . Estimate ${c}_{p} \mathrel{\text{:=}} \mathop{\inf }\limits_{{f \in  \mathcal{F}}}\parallel f * f{\parallel }_{p}$ .
\paragraph{Additional information.} The case $p = 2$ received attention in my first paper [130], where (in slightly different language) the bound ${c}_{2} \geq  {0.7559}\ldots$ was obtained. Some remarks were made in that paper about the nature of the optimal functions $f$ , which appear to be extremely regular.

The best-known lower bound for ${c}_{\infty }$ is 0.64, due to due Cloninger and Steiner-berger [65] (note that their $c$ is $2{c}_{\infty }$ ). The best-known upper bound is ${c}_{\infty } \leq$ ${0.75049}\ldots$ and this appears in Matolcsi and Vinuesa [205]. By contrast to the case $p = 2$ , the functions constructed here are highly pathological.

Note that by Young’s inequality we have ${c}_{\infty } \geq  {c}_{2}^{2}$ .

To conclude this section we formulate one of a number of questions of an additive-combinatorial nature which come up in the theory of fast matrix multiplication.
\end{openproblem}

\plainproblemsection

\begin{openproblem}
Do the following exist, for arbitrarily large $n$ ? An abelian group $H$ with $\left| H\right|  = {n}^{2 + o\left( 1\right) }$ , together with subsets ${A}_{1},\ldots ,{A}_{n},{B}_{1},\ldots ,{B}_{n}$ satisfying $\left| {A}_{i}\right| \left| {B}_{i}\right|  \geq  {n}^{2 - o\left( 1\right) }$ and $\left| {{A}_{i} + {B}_{i}}\right|  = \left| {A}_{i}\right| \left| {B}_{i}\right|$ , such that the sets ${A}_{i} + {B}_{i}$ are disjoint from the sets ${A}_{j} + {B}_{k}\left( {j \neq  k}\right)$ ?
\paragraph{Additional information.} This is [71, Problem 4.7]. If true, it would imply that the exponent of matrix multiplication is 2, that is to say two $n \times  n$ matrices can be composed with ${n}^{2 + o\left( 1\right) }$ multiplications. In [71] and other papers on the topic one may find other questions of this broad type, including in nonabelian groups.

Bonus problems on Sidon sets and related matters.

Klurman and Pohoata asked the following. Let $A \subset  \mathbf{Z}$ be a set of size $n$ . Does $A$ either contain an additive or a multiplicative Sidon set of size ${n}^{1/2 + \delta }$ , for some $\delta  > 0$ ? Shkredov [259] show that this is true if the notion of Sidon is replaced by $g$ -Sidon for some $g = O\left( 1\right)$ . He also showed, as did Roche-Newton and Warren [233], and Peluse and I (unpublished), that one cannot take $\delta  > 1/6$ .

Erdős, Sárközy and Sós asked whether there is an infinite Sidon set in $\mathbf{N}$ which is an asymptotic basis of order 3 .

Update 2023. This last question has been resolved in the affirmative by Pilatte $\left\lbrack  {222}\right\rbrack$ .
\end{openproblem}

\plainproblemsection

\begin{openproblem}
Let ${A}_{1},\ldots ,{A}_{100}$ be ’cubes’ in ${\mathbf{F}}_{3}^{n}$ , that is to say images of $\{ 0,1{\} }^{n}$ under a linear automorphism of ${\mathbf{F}}_{3}^{n}$ . Is it true that ${A}_{1} + \cdots  + {A}_{100} = {\mathbf{F}}_{3}^{n}$ ?
\paragraph{Additional information.} I only recently learned about this problem (from Peter Keevash) but it is in fact 25 years old, and appears in a paper of Jaeger, Linial, Payan and Tarsi [165]. Of course they make a more general conjecture (with ${\mathbf{F}}_{3}$ replaced by ${\mathbf{F}}_{p}$ ) and in the case of ${\mathbf{F}}_{3}$ it may well be the case that the result holds with 100 replaced by 3. It was shown by Alon, Linial and Meshulam [8] that if 100 is replaced by $\sim  \log n$ then the result is true, but so far as I am aware nothing better is known.
\end{openproblem}

\plainproblemsection

\begin{openproblem}
What is the size of the smallest set $A \subset  \mathbf{Z}/p\mathbf{Z}$ (with at least two elements) for which no element in the sumset $A + A$ has a unique representation?
\paragraph{Additional information.} Both Andrew Granville (personal communication) and Swastik Kop-party (open problems session, Harvard 2017) asked me this question.

Update 2023. Bedert [24] has made significant progress on this question, showing that the answer lies between $\omega \left( p\right) \log p$ (for some function $\omega \left( p\right)$ tending to $\infty$ with $p)$ and $O\left( {{\log }^{2}p}\right)$ . This improves previous bounds of $c\log p$ and $O\left( \sqrt{p}\right)$ respectively.
\end{openproblem}

\plainproblemsection

\begin{openproblem}
Suppose that $X, Y$ are two finitely-supported independent random variables taking integer values, and such that $X + Y$ is uniformly distributed on its range. Are $X$ and $Y$ themselves uniformly distributed on their ranges?
\paragraph{Additional information.} This attractive question was asked by Emmanuel Amlot on Stackex-change, and then subsequently attracted some attention on Math Overflow [286]. There are certainly many examples of such $X$ and $Y$ , for example $X$ could be uniformly distributed on $\{ n,{2n},{3n},\cdots ,\left( {n - 1}\right) n\}$ and $Y$ on $\{ 0,1,\ldots , n - 1\}$ .
\end{openproblem}

\plainproblemsection

\begin{openproblem}
What is the largest subset $A \subset  {\mathbf{F}}_{7}^{n}$ for which $A - A$ intersects $\{  - 1,0,1{\} }^{n}$ only at0?
\paragraph{Additional information.} This is a notorious open problem, that of the Shannon capacity of the 7-cycle. I am mentioning it here to draw attention to the fact that it is an additive combinatorics problem, not a graph theory problem. The current best bounds seem to be that ${\left( {C}_{1} - o\left( 1\right) \right) }^{n} \leq  \left| A\right|  \leq  {\left( {C}_{2} + o\left( 1\right) \right) }^{n}$ , where the lower bound ${C}_{1} = {\left( {367}\right) }^{1/5} \approx  {3.2578}$ comes from examples (see [228, Section 9.1]) and the upper bound ${C}_{2} = 7\cos \left( {\pi /7}\right) /\left( {1 + \cos \left( {\pi /7}\right) }\right)  \approx  {3.3177}$ is from Lovász’s celebrated paper [196, Corollary 5].
\end{openproblem}

\plainproblemsection

\begin{openproblem}
If $A \subset  \mathbf{Z}/p\mathbf{Z}$ is random, $\left| A\right|  = \sqrt{p}$ , can we almost surely cover $\mathbf{Z}/p\mathbf{Z}$ with ${100}\sqrt{p}$ translates of $A$ ?
\paragraph{Additional information.} This is a problem I posed in Barbados in 2010. I do not know how to answer this even with 100 replaced by 1.01 . Similar questions are interesting with $\sqrt{p}$ replaced by ${p}^{\theta }$ for any $\theta  \leq  \frac{1}{2}$ . A problem in much the same vein, but possibly more natural, is to take a group $G$ and choose $A \subset  G \times  G,\left| A\right|  = \left| G\right|$ at random. For references to related matters see [34].
\end{openproblem}

\plainproblemsection

\begin{openproblem}
Let $r$ be a fixed positive integer, and let $H\left( r\right)$ be the Hamming ball of radius $r$ in ${\mathbf{F}}_{2}^{n}$ . Let $f\left( r\right)$ be the smallest constant such that there exists an infinite sequence of $n$ s together with subspaces ${V}_{n} \leq  {\mathbf{F}}_{2}^{n}$ with ${V}_{n} + H\left( r\right)  = {\mathbf{F}}_{2}^{n}$ and $\left| {V}_{n}\right|  = \left( {f\left( r\right)  + o\left( 1\right) }\right) \frac{{2}^{n}}{\left| H\left( r\right) \right| }$ . Does $f\left( r\right)  \rightarrow  \infty$ ?
\paragraph{Additional information.} This is a very basic problem about linear covering codes, for which [67] (particularly Chapter 12) may be consulted.

The only value known is $f\left( 1\right)  = 1$ . This follows from the existence of the Hamming code: let $n = {2}^{m} - 1$ , and take ${V}_{n}$ to be the kernel of the $m \times  \left( {{2}^{m} - 1}\right)$ matrix in which the columns are the non-zero vectors in ${\mathbf{F}}_{2}^{m}$ . An easy exercise confirms that every element of ${\mathbf{F}}_{2}^{n}$ has a unique representation in ${V}_{n} + H\left( 1\right)$ (that is, the Hamming code is perfect).

Considering a product of $r$ copies of this example with $n = r\left( {{2}^{m} - 1}\right)$ shows that $f\left( r\right)  \leq  {r}^{r}/r! \sim  {e}^{r}$ (and in particular $f\left( r\right)$ is finite). I am not sure whether any substantially better bound is known or not.

The possibility that $f\left( r\right)  = 1$ for all $r$ has not been ruled out, but it is not known whether $f\left( 2\right)  = 1$ (the best-known upper bound is $f\left( 2\right)  \leq  {1.4238}$ ). [79]).

One may ask the same question without the 'linear', that is to say where the ${V}_{n}$ may be arbitrary subsets of ${\mathbf{F}}_{2}^{n}$ . Writing $\widetilde{f}\left( r\right)$ for the corresponding function, we evidently have $\widetilde{f}\left( r\right)  \leq  f\left( r\right)$ . Here is is known [266]) that $\widetilde{f}\left( 2\right)  = 1$ , but again it could be the case that $\widetilde{f}\left( r\right)  \rightarrow  \infty$ .

Finally, one can ask what happens for all $n$ , rather than merely an infinite sequence. For further references on this, see the book [67] cited above.
\end{openproblem}

\plainproblemsection

\begin{openproblem}
Suppose that $A, B \subset  \{ 1,\ldots , N\}$ both have size ${N}^{0.49}$ . Does $A + B$ contain a composite number?
\paragraph{Additional information.} This is one of my favourite open questions. It arises from an old question of Ostmann, called the 'inverse Goldbach problem': do there exist infinite sets $A, B$ such that $A + B$ coincides with the primes from some point on? The answer to this question is surely no, and a positive answer to Problem (58) would imply this by work of Elsholtz [95]. Adam Harper and me [144] showed that a positive solution follows from a reasonable understanding of the inverse large sieve, which is Problem 47 above.

As a general point, one can ask many questions pertaining to sumsets of sets below the square-root threshold, and our understanding is close to negligible for all of them. For example, do there exist $A, B \subset  \mathbf{Z}/p\mathbf{Z}$ with $\left| A\right| ,\left| B\right|  \sim  {p}^{0.49}$ such that $a + b$ is always a quadratic residue?
\end{openproblem}

\plainproblemsection

\begin{openproblem}
Suppose that $A + A$ contains the first $n$ squares. Is $\left| A\right|  \geq  {n}^{1 - o\left( 1\right) }$ ?
\paragraph{Additional information.} Benny Sudakov reminded me of this question, which was considered by Erdős and Newman [104]. For some discussion, see [7, Theorem 1.6]. It seems that the best known results are the relatively simple observations in [104], where it is shown that necessarily $\left| A\right|  \geq  {n}^{2/3 - o\left( 1\right) }$ , whilst in the other direction there do exist such $A$ with $\left| A\right| { \ll  }_{C}n/{\log }^{C}n$ , for any $C$ .
\end{openproblem}

\plainproblemsection

\begin{openproblem}
Let $p$ be a large prime, and let $A$ be the set of all primes less than $p$ . Is every $x \in  \{ 1,\ldots , p - 1\}$ congruent to some product ${a}_{1}{a}_{2}$ ?
\paragraph{Additional information.} This is a problem of Erdős, Odlyzko and Sárközy [105] from 1987. Walker [279] showed that the result if true if one instead considers 6 -fold products ${a}_{1}{a}_{2}\cdots {a}_{6}$ . of at most 6 primes, and Shparlinski [260] has improved the 6 to a 5 . If one wants products of exactly $k$ primes then it seems that the best value of $k$ known is 20, contained in the thesis of Walker.

Update 2023. Matomäki and Teräväinen [206] have made significant progress on this question, reducing $k$ to 3 .
\end{openproblem}

\plainproblemsection

\begin{openproblem}
Let $A$ be the smallest set containing 2 and 3 and such that ${a}_{1}{a}_{2} - 1 \in$ $A$ if ${a}_{1},{a}_{2} \in  A$ . Does $A$ have positive density?
\paragraph{Additional information.} Erdős [100] attributes this to Hofstadter. The answer is probably yes. A proof of this may have to involve some computation (as well as, presumably, theoretical arguments) since the statement fails if $a = 2$ and $b = 3$ are replaced by, say, $a = 9$ and $b = {10}$ , since the size of the ’words’ of a given length (for example, $a\left( {{ab} - 1}\right)  - 1$ has length 3) grows much more quickly than the number of them.
\end{openproblem}

\plainproblemsection

\begin{openproblem}
Fix a base \(g\geq2\), and let \(P_g(X)\) denote the number of positive
integers \(n\leq X\) that are sums of two base-\(g\) palindromes.
Determine the order of growth of \(P_g(X)\). In particular, narrow the
known gap between bounds of the forms
\[
  X\exp\!\left(-c_g\sqrt{\log X}\right)
  \quad\text{and}\quad
  \frac{X}{(\log X)^c}.
\]

For \(g=2,3,4\), is every sufficiently large positive integer a sum of
three base-\(g\) palindromes?
\end{openproblem}

\clearpage
\section{Sumsets, Approximate Groups, and Inverse Additive Theory}

\plainproblemsection

\begin{openproblem}
Suppose that $A \subset  {\mathbf{F}}_{2}^{n}$ be a set of density $\alpha$ . Does ${10A}$ contain a coset of some subspace of dimension at least $n - O\left( {\log \left( {1/\alpha }\right) }\right)$ ?
\paragraph{Additional information.} This is often known as the Polynomial Bogolyubov conjecture (over finite fields). It may well be true with ${3A}$ in place of ${10A}$ , but the latter would be fine for most applications. It implies the polynomial Freiman-Ruzsa conjecture, and seems to be strictly stronger than it; in particular, it is not addressed by the methods of [127]. Currently, the best bounds known are $n - O\left( {{\log }^{4 + o\left( 1\right) }\left( {1/\alpha }\right) }\right)$ by work of Sanders [249]. One may also ask whether ${2A}$ contains ${99}\%$ of some coset of a subspace of codimension $n - O\left( {\log \left( {1/\alpha }\right) }\right)$ , which would trivially imply that ${4A}$ contains such a subspace. In some ways I feel this is the most natural question in this circle of problems. Update 2024. Košciuszko [188, Theorem 9], building on unpublished work of Konyagin, has shown that for every $\eta  > 0$ there is some $m$ such that ${mA} - {mA}$ contains a subspace of dimension at least $n - O\left( {{\log }^{3 + \eta }\left( {1/\alpha }\right) }\right)$ .

One may ask related questions for subsets of $\left\lbrack  N\right\rbrack$ , but these are presumably at least as hard. The basic form of such a question would be whether, if $A \subset$ $\left\lbrack  N\right\rbrack$ has density $\alpha ,{2A} - {2A}$ contains a Bohr set of constant width and dimension $O\left( {\log \left( {1/\alpha }\right) }\right)$ . A consequence of this (equally unsolved) would be that ${2A} - {2A}$ contains an arithmetic progression of size ${N}^{c\log \left( {1/\alpha }\right) }$ .

There is also a natural question for subsets of ${\left( \mathbf{R}/\mathbf{Z}\right) }^{d}$ : if $A \subset  {\left( \mathbf{R}/\mathbf{Z}\right) }^{d}$ is an open set of measure $\alpha$ , does ${10A} - {10A}$ contains a closed subtorus of codimension $O\left( {\log \left( {1/\alpha }\right) }\right)$ ? Qualitatively, results of this type are known [43].
\end{openproblem}

\plainproblemsection

\begin{openproblem}
Suppose that $A \subset  {\mathbf{F}}_{2}^{n}$ is a set of density $\alpha$ . What is the largest size of coset guaranteed to be contained in ${2A}$ ?
\paragraph{Additional information.} Adding just two copies of a set, as opposed to three or more, leads to less structure. It is known that ${2A}$ must contain a coset of dimension ${ \gg  }_{\alpha }n$ , but need not contain one of dimension $n - \sqrt{n}$ . See [133, Chapter 14] for a discussion of both directions. The behaviour as $\alpha  \rightarrow  {\frac{1}{2}}^{ - }$ is also interesting [248]: in that paper (Question 5.1) Sanders asks whether, if $\alpha  \geq  \frac{1}{2} - \frac{K}{\sqrt{n}}$ , then ${2A}$ must contain a coset of codimension ${O}_{K}\left( 1\right)$ .

This problem has an analogue for subsets of $\left\lbrack  N\right\rbrack$ , with ’coset’ replaced by ’arithmetic progression’. In this setting, the lower bound is $\sim  {e}^{c{\left( \log N\right) }^{1/2}}\left\lbrack  {132}\right\rbrack$ and the upper bound $\sim  {e}^{c{\left( \log N\right) }^{2/3}}\left\lbrack  {238}\right\rbrack$ .
\end{openproblem}

\plainproblemsection

\begin{openproblem}
Suppose that $A \subset  {\mathbf{F}}_{2}^{n}$ is a set with an additive complement of size $K$ (that is, for which there is another set $S \subset {\mathbf{F}}_{2}^{n},\allowbreak\left| S\right|  = K$ , with $A + S = {\mathbf{F}}_{2}^{n}$ ).

Does ${2A}$ contain a coset of codimension ${O}_{K}\left( 1\right)$ ? Could it even contain a coset of codimension $O\left( {\log K}\right)$ ?
\paragraph{Additional information.} It is quite mysterious that even the weaker, qualitative, statement is seemingly not known. The stronger statement (with codimension $O\left( {\log K}\right)$ ) implies the polynomial Bogolyubov conjecture stated above (by the Ruzsa covering lemma; apply the statement with ${2A}$ in place of $A$ ).

This problem also has an obvious analogue in $\left\lbrack  N\right\rbrack$ . Even in the case $S =$ $\{ 1,\ldots , K\}$ (in which case $A$ is syndetic with gaps bounded by $K$ ) we do not seem to know much, though it is natural to conjecture that $A - A$ contains a progression of length ${N}^{{c}_{K}}$ . The problem feels closely connected with a notorious question of Katznelson: if $A \subset  \mathbf{Z}$ is syndetic, does $A - A$ contain a Bohr set $\left\{  {n \in  \mathbf{Z} : \parallel {\theta n}{\parallel }_{\mathbf{T}} \leq  \varepsilon }\right\}$ , for some $\varepsilon  > 0$ ? One should, however, note that the infinitary setting for problems about difference sets can be quite different to the finite one. For more on Katznelson's question and its history, see the introduction to [122].

Tom Sanders reminded me that, around 17 years ago, I established [131] the following very specific result related to the above: if $A \subset  \mathbf{Z}/p\mathbf{Z}$ and if, for all $a \in  A$ , at least one of $a + 2, a + 3$ lies in $A$ , then $A - A$ contains a progression of length $\gg  \sqrt{p}$ .
\end{openproblem}

\plainproblemsection

\begin{openproblem}
Suppose that ${\mathbf{F}}_{2}^{n}$ is partitioned in to sets ${A}_{1},\ldots ,{A}_{K}$ . Does $2{A}_{i}$ contain a coset of codimension ${O}_{K}\left( 1\right)$ for some $i$ ?
\paragraph{Additional information.} This would easily imply a positive solution to the previous problem. It does not seem to be known even for $K = 3$ . Note that a Hamming ball $B \mathrel{\text{:=}}$ $\left\{  {x : {x}_{1} + \cdots  + {x}_{n} \leq  \frac{n}{2} - \sqrt{n}}\right\}$ has the property that ${2B}$ does not contain a coset of codimension ${O}_{K}\left( 1\right)$ , and so a subcase of this problem is to show that three Hamming balls ${A}_{1},{A}_{2},{A}_{3}$ (relative to different bases) cannot cover ${\mathbf{F}}_{2}^{n}$ .

The statement above could conceivably hold with $O\left( {\log K}\right)$ in place of ${O}_{K}\left( 1\right)$ , which would again imply the polynomial Bogolyubov and Freiman-Ruzsa conjectures.

As with the preceding problem, there are analogies with Katznelson's question. Indeed, in [122, Question C2] it is shown that an equivalent form of that question is the following: If $\mathbf{N}$ is partitioned into sets ${A}_{1},\ldots ,{A}_{K}$ , does ${A}_{i} - {A}_{i}$ contain a Bohr set?
\end{openproblem}

\plainproblemsection

\begin{openproblem}
Is there an absolute constant $c > 0$ such that, if $A \subset  \mathbf{N}$ is a set of squares of size at least 2, then $\left| {A + A}\right|  \geq  {\left| A\right| }^{1 + c}$ ?
\paragraph{Additional information.} This is quite a well-known problem, and is strictly easier than a very old question of Hardy and Littlewood, which asks whether the squares are a $\Lambda \left( p\right)$ - set for some $2 < p < 4$ . (Nowadays, however, the sumset question probably seems more basic to most readers.) It may even be that $c$ can be taken arbitrarily close to 1, for sufficiently large sets $A$ .

A positive answer to this problem obviously implies that if $P$ is an arithmetic progression of length $n$ then the number of squares in $P$ is at most $O\left( {n}^{\kappa }\right)$ for some $\kappa  < 1$ . Rudin [235] conjectured that in fact this is even true with $\kappa  = \frac{1}{2}$ . Whilst the weaker statement with $\kappa  < 1$ is known (see [33] with $\kappa  = \frac{2}{3} + o\left( 1\right)$ , and,[35] with $\kappa  = \frac{3}{5} + o\left( 1\right)$ ), all known proofs involve serious statements about the number of rational points on various curves. Thus one expects that this problem must necessarily involve such ingredients as well. In particular, Solymosi [262] has observed that an affirmative answer would follow if one could show that the squares do not contain an ’affine cube’ $\left\{  {x + {\omega }_{1}{h}_{1} + \cdots  + {\omega }_{d}{h}_{d} : {\omega }_{1},\ldots ,{\omega }_{d} \in  }\right.$ $\{ 0,1\} \}$ for some $d$ , a problem which feels purely diophantine.

A very comprehensive resource for this circle of problems is the nice paper of Cilleruelo and Granville [64].
\end{openproblem}

\plainproblemsection

\begin{openproblem}
Is every group well-approximated by finite groups?
\paragraph{Additional information.} There are really two questions here, namely is every group sofic? and is every group hyperlinear? Here is a precise statement of the latter question. Suppose that $G$ is a group. Is the following true? For every finite set $A \subset  G$ containing the identity ${e}_{G}$ and for every $\varepsilon  > 0$ , there is some unitary group $\mathrm{U}\left( n\right)$ and a map $\phi  : A \rightarrow  \mathrm{U}\left( n\right)$ such that $\phi \left( {e}_{G}\right)  = {I}_{n}$ , such that $\phi$ is an ’approximate homomorphism’ in the sense that $\begin{Vmatrix}{\phi \left( {{a}_{1}{a}_{2}}\right)  - \phi \left( {a}_{1}\right) \phi \left( {a}_{2}\right) }\end{Vmatrix} < \varepsilon$ for all ${a}_{1},{a}_{2} \in  A$ such that ${a}_{1}{a}_{2} \in  A$ , whilst $\phi$ is ’nontrivial’ in the sense that $\begin{Vmatrix}{\phi \left( a\right)  - {I}_{n}}\end{Vmatrix} \geq  1$ for all $a \in  A, a \neq  {e}_{G}$ . (Here $\parallel  - \parallel$ is the ${\ell }^{2}$ -to- ${\ell }^{2}$ operator norm. Also, the constant 1 is unimportant, since a certain amplification trick shows that the concept is the same if 1 is replaced by any other constant in $\left( {0,\sqrt{2}}\right)$ ; see [221, p. 458].)

Roughly speaking, the notion of sofic replaces the unitary group $\mathrm{U}\left( n\right)$ here with the symmetric group $\operatorname{Sym}\left( n\right)$ , endowed with normalised Hamming distance. It is known that all sofic groups are hyperlinear, but the reverse implication is not known. See [221].

The Higman group $\left\langle  {{a}_{1},{a}_{2},{a}_{3},{a}_{4} \mid  {a}_{i}^{{a}_{i + 1}} = {a}_{i}^{2}, i \in  \mathbf{Z}/4\mathbf{Z}}\right\rangle$ (where here ${x}^{y} \mathrel{\text{:=}} {y}^{-1}{xy}$ ) is not known to be either hyperlinear or sofic. Helfgott and Juschenko [158] showed that if the Higman group is sofic then some very strange functions $f : \mathbf{Z}/p\mathbf{Z} \rightarrow  \mathbf{Z}/p\mathbf{Z}$ must exist. Namely, given $\varepsilon  > 0$ , if $p$ is sufficiently large in terms of $\varepsilon$ then there is $f$ with $f\left( {x + 1}\right)  = {2f}\left( x\right)$ for at least $\left( {1 - \varepsilon }\right) p$ values of $x \in  \mathbf{Z}/p\mathbf{Z}$ , and also $f \circ  f \circ  f \circ  f\left( x\right)  = x$ for all $x$ .
\end{openproblem}

\plainproblemsection

\begin{openproblem}
Suppose that $A$ is a $K$ -approximate group (not necessarily abelian). Is there $S \subset  A,\left| S\right|  \gg  {K}^{-O\left( 1\right) }\left| A\right|$ , with ${S}^{8} \subset  {A}^{4}$ ?
\paragraph{Additional information.} For the definition of $K$ -approximate group, see for example [136]. Such a conclusion is known with $\left| S\right| { \gg  }_{K}\left| A\right|$ (but not with a polynomial bound) by an argument of Sanders. See [44, Problem 6.5].

Shachar Lovett once mentioned to me the following vaguely related question: suppose that $A \subset  {\mathbf{F}}_{2}^{n}$ is a set of density $\frac{1}{3}$ . Is there a set $B$ with ${4B} \mathrel{\text{:=}} B + B +$ $B + B \subset  A + A$ , and with ${4B}$ having density at least $\frac{1}{100}$ in ${\mathbf{F}}_{2}^{n}$ ?
\end{openproblem}

\plainproblemsection

\begin{openproblem}
Given a set $A \subset  \mathbf{Z}$ with $D\left( A\right)  \leq  K$ , find a large structured subset ${A}^{\prime }$ which ’obviously’ has $D\left( {A}^{\prime }\right)  \leq  K + \varepsilon$ (Here, $D\left( A\right)  \mathrel{\text{:=}} \left| {A - A}\right| /\left| A\right|$ ).
\paragraph{Additional information.} To illustrate what is meant here, let me state a result of Eberhard, Manners and me [88, Section 6]: if $\left| {A - A}\right|  \leq  \left( {4 - \varepsilon }\right) \left| A\right|$ , then there is a progression $P$ of length ${ \gg  }_{\varepsilon }\left| A\right|$ on which $A$ has density $> \frac{1}{2}$ . Then ${A}^{\prime } \mathrel{\text{:=}} A \cap  P$ ’obviously’ has $D\left( {A}^{\prime }\right)  \leq  4$ . One can imagine a more general such result. One could also ask for bounds; in the proof of the result just stated for doubling close to 4 , the constant in the ${ \gg  }_{\varepsilon }1$ is essentially ineffective.
\end{openproblem}

\plainproblemsection

\begin{openproblem}
Which finite groups have the smallest biggest product-free sets?
\paragraph{Additional information.} 20 years ago, Kedlaya [175] observed, refining some work of Babai and Sós, that it follows from the classification of finite simple groups that every finite group $G$ of order $n$ has a product-free subset of size $\gg  {n}^{{11}/{14}}$ . It may well be that this exponent is sharp. To show this, it follows by inspection of Kedlaya's paper that one would need to show that the Ree groups ${}^{2}{G}_{2}\left( q\right) , q = {3}^{{2m} + 1}$ , which have order $\sim  {q}^{7}$ , infinitely often have no product-free set of size $\gg  {q}^{{11}/2}$ . Certainly working with these groups in any direct way is likely to be very challenging, although the dimensions of their smallest representations are understood. A good model problem would be to determine the largest product-free subsets of ${\mathrm{{SL}}}_{2}\left( {\mathbf{F}}_{p}\right)$ , since explicit computation in this group and with its representations are quite tractable. For this problem, I believe the best-known upper bound is $O\left( {n}^{8/9}\right)$ , due to Gowers. For more on these problems, see Kedlaya's more recent survey [176].
\end{openproblem}

\plainproblemsection

\begin{openproblem}
Let $K \subset  {\mathbf{R}}^{\mathbf{N}}$ be a balanced compact set (that is, ${\lambda K} \subset  K$ when $\left| \lambda \right|  \leq  1$ ) and suppose that the normalised Gaussian measure ${\gamma }_{\infty }\left( K\right)$ is at least 0.99 . Does ${10K}$ (say) contain a compact convex set $C$ with ${\gamma }_{\infty }\left( C\right)  \geq  {0.01}$ ?
\paragraph{Additional information.} This question is raised by Talagrand [268]. He raised the question again in [269] and offered 1000 dollars for a solution. The answer is no if ${10}\mathrm{K}$ is replaced by ${2K}$ , so there seems to be a strong formal similarity with the questions mentioned above. This is presumably why several people have mentioned it to me over the years.
\end{openproblem}

\plainproblemsection

\begin{openproblem}
Let $p$ be an odd prime and suppose that $f : {\mathbf{F}}_{p}^{n} \times  {\mathbf{F}}_{p}^{n} \rightarrow  \mathbf{C}$ is a function bounded pointwise by 1 . Suppose that ${\mathbf{E}}_{h}{\begin{Vmatrix}{\Delta }_{\left( h, h\right) }f\end{Vmatrix}}_{\square }^{4} \geq  \delta$ . Does $f$ correlate with a function of the form $a\left( x\right) b\left( y\right) c\left( {x + y}\right) {\left( -1\right) }^{q\left( {x, y}\right) }$ ?
\paragraph{Additional information.} We could take $q$ asymmetric by absorbing the symmetric part into the other terms. Here ${\Delta }_{\left( h, h\right) }f\left( {x, y}\right)  = \overline{f\left( {x, y}\right) f\left( {x + h, y + h}\right) }$ , and $\parallel  \cdot  {\parallel }_{\square }$ is the box norm, thus this norm is counting configurations $\left( {x, y}\right),\allowbreak\left( {x + k, y}\right),\allowbreak\left( {x, y + \ell }\right),\allowbreak\left( x + k, y + \ell \right),\allowbreak\left( {x + h, y + h}\right),\allowbreak\left( {x + k + h, y + h}\right),\allowbreak\left( {x + h, y + \ell  + h}\right),\allowbreak\left( {x + k + h, y + \ell  + h}\right)$ . This appears to be the simplest 'Gowers-type norm' for which an inverse theorem is not known, and as such is relevant, though not directly, to Problem 20. For discussion, and a link to group cohomology, see [12, Example 1.6].
\end{openproblem}

\plainproblemsection

\begin{openproblem}
Bounds for the inverse theorem for Gowers norms.
\paragraph{Additional information.} This has been considered one of the biggest open question in the subject and it may be asked for finite fields and over the integers.

Over finite fields, the question is simpler to state and it is natural to make the following ’Polynomial Gowers Inverse Conjecture’. Fix an integer $s \geq  1$ , and let $p > s$ . Suppose that $f : {\mathbf{F}}_{p}^{n} \rightarrow  \mathbf{C}$ is a function with $\left| {f\left( x\right) }\right|  \leq  1$ for all $x$ , and that $\parallel f{\parallel }_{{U}^{s + 1}} \geq  \delta$ (the definition of the Gowers norm may be found in several places, including the works cited below). Then there is a form $\psi$ of degree $s$ such that $\left| {{\mathbf{E}}_{x \in  {\mathbf{F}}_{p}^{n}}f\left( x\right) e\left( {-{2\pi i\psi }\left( x\right) /p}\right) }\right|  \gg  {\delta }^{{O}_{s}\left( 1\right) }$ . Such a statement is not known for any $s \geq  2$ . Indeed, it was shown independently by Tao and me [149] and by Lovett [197] that the case $s = 2$ of this assertion is equivalent to Marton’s Conjecture (Problem 49) for ${\mathbf{F}}_{p}^{n}$ . Update 2023. Polynomially effective bounds for the case $s = 2$ are now known as a consequence of the solution of Marton's conjecture in [127, 128].

For $s = 3$ , quantitative bounds (of double exponential type) were obtained by Gowers and Milićević [129], but for $s \geq  4$ only qualitative bounds are known, from remarkable work of Bergelson, Tao and Ziegler [27].

In the case $s = 3$ , the problem is very closely related to the following attractive question. Suppose that $f : {\mathbf{F}}_{p}^{n} \rightarrow  {\operatorname{Mat}}_{n}\left( {\mathbf{F}}_{p}\right)$ is a map with image in the $n \times  n$ matrices over ${\mathbf{F}}_{p}$ , and that $\operatorname{rank}\left( {f\left( {x + y}\right)  - f\left( x\right)  - f\left( y\right) }\right)  \leq  r$ for all $x, y$ . Is there a linear map $\widetilde{f} : {\mathbf{F}}_{p}^{n} \rightarrow  {\operatorname{Mat}}_{n}\left( {\mathbf{F}}_{p}\right)$ with $f - \widetilde{f}$ taking values in the matrices of rank at most ${10r}$ (say)? See [174, Question 1.5] for further discussion.

The assumption $p > s$ can be removed, but care is required with the statement $\left\lbrack  {{148},{198}}\right\rbrack$ and significant further difficulties arise in the proof [246,273].

Update 2024. Until very recently, the situation for the $\parallel  \cdot  {\parallel }_{{U}^{s + 1}\left\lbrack  N\right\rbrack  }$ -norm over the integers was worse.

In the case $s = 2$ quantitative (but not optimal) bounds were established by Tao and me [147], and Manners [204] established the first quantitative bounds for general $s$ , of double exponential type.

In a significant breakthrough, Leng, Sah and Sawhney [191] established a 'quasi-polynomial’ inverse theorem for the $\parallel  \cdot  {\parallel }_{{U}^{s + 1}\left\lbrack  N\right\rbrack  }$ -norms for all $s \geq  3$ . That is, if $f : \left\lbrack  N\right\rbrack   \rightarrow  \mathbf{C}$ is a 1-bounded function such that $\parallel f{\parallel }_{{U}^{s + 1}\left\lbrack  N\right\rbrack  } \geq  \delta$ , show that there is a polynomial nilsequence $\phi \left( {p\left( n\right) }\right)$ attached to some nilmanifold $G/\Gamma$ of complexity at most $C$ and dimension $\ll  {\log }^{C}\left( {1/\delta }\right)$ , with $\phi$ having Lipschitz constant bounded by 1, and with $\left| {{\mathbf{E}}_{n \leq  N}f\left( n\right) \phi \left( {p\left( n\right) }\right) }\right|  \gg  \exp \left( {-{\log }^{C}\left( {1/\delta }\right) }\right)$ . (For an explanation of the terms here, see for example [190].)

It is possible that a truly polynomial statement may hold, but such a statement has not even been properly formulated at the moment.
\end{openproblem}

\plainproblemsection

\begin{openproblem}
Let $G$ be an abelian group, and consider the space $\Phi \left( G\right)$ of all functions on $G$ which are convex combinations of functions of the form

\[
\phi \left( g\right)  \mathrel{\text{:=}} {\mathbf{E}}_{{x}_{1} + {x}_{2} + {x}_{3} = g}{f}_{1}\left( {{x}_{2},{x}_{3}}\right) {f}_{2}\left( {{x}_{1},{x}_{3}}\right) {f}_{3}\left( {{x}_{1},{x}_{2}}\right)
\]

with ${\begin{Vmatrix}{f}_{i}\end{Vmatrix}}_{\infty } \leq  1$ . Let ${\Phi }^{\prime }\left( G\right)$ be the space of functions defined similarly, but with ${f}_{3}\left( {{x}_{1},{x}_{2}}\right)$ now required to be a function of ${x}_{1} + {x}_{2}$ . Do $\Phi \left( G\right)$ and ${\Phi }^{\prime }\left( G\right)$ coincide?
\paragraph{Additional information.} One would guess that the answer is probably 'no'. The motivation for this problem is that both $\Phi \left( G\right)$ and ${\Phi }^{\prime }\left( G\right)$ are notions of ’quadratically structured function’ that have been considered in the literature; $\Phi \left( G\right)$ is a ’generalised convolution algebra’ as considered by Conlon-Fox-Zhao [73, Section 5], whereas ${\Phi }^{\prime }\left( G\right)$ consists of Tao’s so-called ${\mathrm{{UAP}}}^{2}\left( G\right)$ -functions [270]. (This last equivalence is not obvious; I have some unpublished notes on it.)
\end{openproblem}

\plainproblemsection

\begin{openproblem}
Let \(A\subset\mathrm{SO}(3)\) be open, and let \(\mu\) be normalized Haar
measure. Is it true that
\[
  \mu(A\cdot A)
  \geq
  \min\!\left\{1,\,4\mu(A)\bigl(1-\mu(A)\bigr)\right\}?
\]
\end{openproblem}

\plainproblemsection

\begin{openproblem}
Pick ${x}_{1},\ldots ,{x}_{k} \in  {A}_{n}$ (the alternating group on $n$ letters) at random. Is it true that, almost surely as $n \rightarrow  \infty$ , the random walk on this set of generators and their inverses equidistributes in time $O\left( {n\log n}\right)$ ?
\paragraph{Additional information.} The statement is one equivalent form of what it means to be an expanding set of generators. This may even be true for $k = 2$ , although it is not known if this statement is true for any pair of elements in ${A}_{n}$ .

The statement about equidistribution of the random walk is one of the known equivalent definitions of being an expanding set of generators. Thus an alternative formulation of the question is: do $k$ random elements of ${A}_{n}$ form an expanding set of generators?

At present this seems a pretty hopeless problem. It would imply that (a.s. as $n \rightarrow  \infty$ ) the diameter of ${A}_{n}$ with respect to the generating set ${x}_{1},\ldots ,{x}_{k}$ is $O\left( {n\log n}\right)$ . By contrast, the best known upper bound for this is $O\left( {{n}^{2}{\log }^{C}n}\right)$ , due to Helfgott, Seress and Zuk [159].
\end{openproblem}

\plainproblemsection

\begin{openproblem}
Find bounds in the classification theorem for approximate groups.
\paragraph{Additional information.} The theorem being referred to is [43]. This theorem is ineffective due to the use of an ultraproduct argument.
\end{openproblem}

\plainproblemsection

\begin{openproblem}
Let $A \subset  {\mathbf{F}}_{2}^{n}$ . If $V$ is a subspace of ${\mathbf{F}}_{2}^{n}$ , write $\alpha \left( V\right)$ for the density of $A$ on $V$ . Is there some $V$ of moderately small codimension on which $\alpha$ is stable in the sense that $\left| {\alpha \left( V\right)  - \alpha \left( {V}^{\prime }\right) }\right|  \leq  \varepsilon$ whenever ${V}^{\prime }$ is a codimension 1 subspace of $V$ ?
\paragraph{Additional information.} This is an arithmetic variant of an open question in graph theory raised by Conlon and Fox [72, Lemma 3.6]. In that lemma, they show that every graph $G$ on $n$ vertices contains an $\varepsilon$ -regular subset on ${n}^{\prime }$ vertices, where ${n}^{\prime }/n \geq$ ${2}^{-{\varepsilon }^{-{\left( {10}/\varepsilon \right) }^{4}}}$ , and they go on to say ’while our bound gives a double exponential dependence, we suspect that the truth is more likely to be a single exponential. We leave this as an open problem.'

Manners showed that one can have such a $V$ of codimension exponential in $1/\varepsilon$ . In the other direction, there are examples to show that the codimension must grow like a power of $1/\varepsilon$ .
\end{openproblem}

\plainproblemsection

\begin{openproblem}
Suppose that $A \subset  \mathbf{Z}/p\mathbf{Z}$ is a set of density $\frac{1}{2}$ . Under what conditions on $K$ can $A$ be almost invariant under all maps $\phi \left( x\right)  = {ax} + b,\left| a\right| ,\left| b\right|  \leq  K$ ?
\paragraph{Additional information.} This problem was considered by Eberhard, Mrazović and me (unpublished). By ’almost invariant’ we mean $\left| {A\bigtriangleup \phi \left( A\right) }\right|  = o\left( p\right)$ . One can have $K \rightarrow  \infty$ (but slowly, something like ${\log }^{c}p$ ); this is related to the fact that the affine group over ${\mathbf{F}}_{p}$ is amenable. In the other direction, we showed that $K$ cannot be as large as ${p}^{1/{100}}$ .
\end{openproblem}

\clearpage
\section{Sieve Methods and Diophantine Problems}

\plainproblemsection

\begin{openproblem}
Let $N$ be a large integer. For each prime $p$ with ${N}^{0.51} \leq  p < 2{N}^{0.51}$ , pick some residue $a\left( p\right)  \in  \mathbf{Z}/p\mathbf{Z}$ . Is it true that

\[
\# \{ n \in  \left\lbrack  N\right\rbrack   : n \equiv  a\left( p\right) \left( {\;\operatorname{mod}\;p}\right) \text{ for some }p\}  \gg  {N}^{1 - o\left( 1\right) }?
\]
\paragraph{Additional information.} This problem is formulated in [139, Section 4]. I think of it as a kind of ’Kakeya problem in dimension $2 + \varepsilon$ ’. One can make many other related conjectures. One of my favourites is raised in [9]: if $A \subset  \mathbf{Z}/p\mathbf{Z}$ is a set of size $\lfloor p/2\rfloor$ , does some dilate of $A$ have no gaps of length more than ${p}^{0.49}$ ? To see the relation to problems of the type considered here, note that the complement of such a set would have to have a progression of length ${p}^{0.49}$ , for each common difference $d \in  {\left( \mathbf{Z}/p\mathbf{Z}\right) }^{ * }$ .

A minor variant of Problem 43 asks the same, but now with $p$ allowed to vary in a range such as ${N}^{0.51} \leq  p < {N}^{0.52}$ . Obtaining a lower bound ${N}^{1 - o\left( 1\right) }$ here is certainly no harder than Problem 43, but I do not know how to do it. In fact, for this problem there may even be a lower bound of ${cN}$ .

One can of course ask similar questions with 0.51 replaced by other exponents $\alpha$ . For $\alpha  < \frac{1}{2}$ these questions are straightforward by an inclusion exclusion or Cauchy-Schwarz argument, for much the same reason that it is relatively easy to establish the Kakeya conjecture in the plane. On the other hand I have observed (unpublished) that by arguments of Bourgain, progress for $\alpha  \approx  1$ would imply the Kakeya conjecture.

To conclude these remarks, let me mention a toy problem related to the way in which pairs of residue classes $a\left( {p}_{1}\right) \left( {\;\operatorname{mod}\;{p}_{1}}\right)$ and $a\left( {p}_{2}\right) \left( {\;\operatorname{mod}\;{p}_{2}}\right)$ interact. It came up from a study of Problem 48, but seems to be concerned with the same phenomena that apparently make Problem 43 hard. Suppose that for all primes $p$ with $X \leq$ $p < {2X}$ one has a residue class $a\left( p\right) \left( {\;\operatorname{mod}\;p}\right)$ . For distinct primes ${p}_{1},{p}_{2}$ in this range, let $f\left( {{p}_{1},{p}_{2}}\right)  \in  \mathbf{Z}/{p}_{1}{p}_{2}\mathbf{Z}$ be the unique solution to $f\left( {{p}_{1},{p}_{2}}\right)  \equiv  a\left( {p}_{i}\right) \left( {\;\operatorname{mod}\;{p}_{i}}\right)$ , $i = 1,2$ . What can be said about the function $a$ if $\left| {{\mathbf{E}}_{X \leq  p \leq  {2X}}e\left( \frac{f\left( {{p}_{1},{p}_{2}}\right) }{{p}_{1}{p}_{2}}\right) }\right|  \geq  {0.99}$ ? Perhaps $a\left( \right)$ must be almost constant. More ambitiously, one could replace 0.99 by 0.01 or even smaller quantities.
\end{openproblem}

\plainproblemsection

\begin{openproblem}
Sieve $\left\lbrack  N\right\rbrack$ by removing half the residue classes $\;{\;\operatorname{mod}\;{p}_{i}}$ , for primes $2 \leq  {p}_{1} < {p}_{2} < \cdots  < {p}_{1000} < {N}^{9/{10}}$ . Does the remaining set have size at most $\frac{1}{10}N$ ?
\paragraph{Additional information.} This is raised in [101, Section 6, Problem 3]. Erdős remarks that the answer is affirmative if the primes are all less than ${N}^{1/2}$ , by the large sieve. I must admit that I do not know anything about this problem other than what Erdős wrote nearly 40 years ago; this part of his paper does not appear to have been cited since. The same comment applies to the next problem.
\end{openproblem}

\plainproblemsection

\begin{openproblem}
Can we pick residue classes ${a}_{p}\left( {\;\operatorname{mod}\;p}\right)$ , one for each prime $p \leq  N$ , such that every integer $\leq  N$ lies in at least 10 of them?
\paragraph{Additional information.} This is raised in [101, Section 6, Problem 6]. Erdős remarks that he does not know how to answer it with 10 replaced by 2 .
\end{openproblem}

\plainproblemsection

\begin{openproblem}
What is the largest $y$ for which one may cover the interval $\left\lbrack  y\right\rbrack$ by residue classes ${a}_{p}\left( {\;\operatorname{mod}\;p}\right)$ , one for each prime $p \leq  x$ ?
\paragraph{Additional information.} This is the Jacobsthal problem, and it is somewhat notorious. It is known [113] that $y \gg  x\frac{\log x\log \log \log x}{\log \log x}$ , and any improvement upon this would lead to a better bound for the largest gap between consecutive primes. The best upper bound is $y \ll  {x}^{2}$ , due to Iwaniec [164]. It seems very likely that one must have $y \ll  {x}^{1 + o\left( 1\right) }$ . A proof of this would not give a better upper bound on gaps between primes, merely on the capability of one method for producing them.

Erdős and Ruzsa [106] and Hildebrand [161] mention the following elegant problem of a similar type: can one cover $\left\lbrack  x\right\rbrack$ with residue classes $a\left( p\right) \left( {\;\operatorname{mod}\;p}\right) , p \leq  x$ , at most one for each prime $p$ , and with $\sum \frac{1}{p} \leq  K$ ? Erdős and Ruzsa suggest that in fact the answer is no, and moreover that the uncovered set should have size at least $c\left( K\right) x$ .
\end{openproblem}

\plainproblemsection

\begin{openproblem}
Suppose that a large sieve process leaves a set of quadratic size. Is that set quadratic?
\paragraph{Additional information.} Questions of this type are known as the inverse large sieve problem. There are many questions here, but the following very particular instance is probably the simplest. Suppose that $A \subset  \mathbf{N}$ is a set with the property that $\left| {A\left( {\;\operatorname{mod}\;p}\right) }\right|  \leq  \frac{1}{2}\left( {p + 1}\right)$ for all sufficiently large $p$ . The large sieve then implies that $\left| {A \cap  \left\lbrack  X\right\rbrack  }\right|  \ll  {X}^{1/2}$ , and this is clearly sharp because one could take $A$ to be the set of squares or the image of $\mathbf{Z}$ under some other quadratic map $\phi  : \mathbf{Q} \rightarrow  \mathbf{Q}$ . However, is it true that either $\left| {A \cap  \left\lbrack  X\right\rbrack  }\right|  \ll  \frac{{X}^{1/2}}{{\log }^{100}X}$ , or $A$ is contained in the image of $\mathbf{Z}$ under a quadratic map $\phi  : \mathbf{Q} \rightarrow  \mathbf{Q}$ ? For background on this question see $\left\lbrack  {{144},{160},{280},{281}}\right\rbrack$ .
\end{openproblem}

\plainproblemsection

\begin{openproblem}
Suppose that a small sieve process leaves a set of maximal size. What is the structure of that set?
\paragraph{Additional information.} It is not at all clear to me what a good formulation of this question, the inverse small sieve problem, should be. One could, for example, consider a linear sieve in which precisely one residue class $a\left( p\right) \left( {\;\operatorname{mod}\;p}\right)$ is removed from $\left\lbrack  X\right\rbrack$ for each prime $p \leq  \sqrt{X}$ (or even for somewhat larger primes as well), leaving a set $S$ . The Selberg sieve or the Rosser-Iwaniec sieve yield a bound $\left| S\right|  \leq  (2 +$ $o\left( 1\right) )X/\log X$ . Selberg [255] observed that one cannot hope to beat such a bound without eliminating the possibility of a Siegel zero, a notorious open problem in number theory. (To see roughly why, imagine that all primes were 3(mod 4), which would correspond to a very extreme type of Siegel zero. Then there would be $\sim  {4N}/\log N$ values of $k \leq  N$ for which ${4k} + 3$ is a prime $> {N}^{0.9}$ . Suppose without loss of generality that at least ${2N}/\log N$ of these $k$ s are odd. Then, setting $a\left( 2\right)  = 0$ and $a\left( p\right)  =  - \frac{3}{4}\left( {\;\operatorname{mod}\;p}\right)$ for $p$ odd, we see that $k \neq  a\left( p\right) \left( {\;\operatorname{mod}\;p}\right)$ for all $p \leq  \sqrt{N}$ .)

One possible form of the inverse small sieve problem would then to be ask whether examples of this type (that is, 'coming from a Siegel zero') are essentially the only ones; if that could be shown, then for instance under the GRH one could improve the upper bound on $\left| S\right|$ .

More realistic, but still seemingly very difficult, might be to ask about sieves of dimension $\kappa  < 1$ , where one only takes a fraction $\kappa$ of all primes $p$ . In the case $\kappa  = \frac{1}{2}$ , for example, the optimal bound on $S$ is known and one might ask for a characterisation of the extremal examples. This has some resemblance to the inverse large sieve problem, since the extremal examples are related to sums of squares. See [256, Sections 16, 19] for more on this.
\end{openproblem}

\plainproblemsection

\begin{openproblem}
Is every $n \leq  N$ the sum of two integers, all of whose prime factors are at most ${N}^{\varepsilon }$ ?
\paragraph{Additional information.} This was asked by Erdős, several times I think, but certainly in a collection [102] from 1981 one finds the following problem, which is essentially equivalent: ’An old conjecture of mine states that if $f\left( n\right)$ is the least integer not of the form $a + b$ with $P\left( {a, b}\right)  \leq  n$ then for every $k$ and for $n > {n}_{0}\left( k\right)$ we have $f\left( n\right)  > {n}^{k}$ . This conjecture does not look hard but I could not get anywhere with it'.

The reason that the problem 'does not look hard' is perhaps that the set of integers less than $N$ , all of whose prime factors are at most ${N}^{\varepsilon }$ , has positive density (depending on $\varepsilon$ ). However, the problem seems closely related to other notorious binary problems such as the Goldbach conjecture. I believe the record is still Balog's exponent $4/9\sqrt{e} \approx  {0.2695}$ , obtained 30 years ago in [14].

Trevor Wooley has pointed out the following, which he says is well-known. If $p \equiv  3\left( {\;\operatorname{mod}\;4}\right)$ , and if $p$ is the sum of two ${p}^{\varepsilon }$ -smooth numbers, then there is some quadratic non-residue ${\;\operatorname{mod}\;p}$ of size at most ${p}^{\varepsilon }$ . Thus, to answer the problem for all $\varepsilon  > 0$ is at least as hard as the notorious Vinogradov least quadratic non-residue problem (at least, the version of that problem restricted to primes $p \equiv  3\left( {\;\operatorname{mod}\;4}\right)$ ). The best unconditional exponent for that problem is $1/4\sqrt{e}$ , which is therefore a serious barrier for our problem. That said, the Vinogradov problem is known under GRH, so it may also be interesting to explore our problem with that assumption.
\end{openproblem}

\plainproblemsection

\begin{openproblem}
Do there exist infinitely many primes for which $p - 2$ has an odd number of prime factors?
\paragraph{Additional information.} The same question may be asked with $p - 1$ (and this is probably more natural) but with $p - 2$ the question is a weak form of the twin prime conjecture. The set of integers $S$ with an odd number of prime factors has density $\frac{1}{2}$ , so one is ’only’ asking for infinitely many primes in a set $\left( {S + 1}\right)$ of density $\frac{1}{2}$ .
\end{openproblem}

\plainproblemsection

\begin{openproblem}
Is there always a sum of two squares between $X - \frac{1}{10}{X}^{1/4}$ and $X$ ?
\paragraph{Additional information.} I am including this mainly because it is in Littlewood's list [193] (Problem 2), Montgomery's list [209] and the first problem paper of Erdős [98] (Problem 15), where he already describes it as 'an old problem'. There is a wellknown, almost trivial, argument showing a bound of $O\left( {X}^{1/4}\right)$ on the left: subtract off the greatest square ${u}^{2}$ less than $X$ , then subtract off the greatest square ${v}^{2}$ less than $X - {u}^{2}$ . Sofia Lindqvist and I [145] have a somewhat different argument, giving a sum of two almost equal squares within $O\left( {X}^{1/4}\right)$ of $X$ .
\end{openproblem}

\plainproblemsection

\begin{openproblem}
Bounds for Waring's problem over finite fields.
\paragraph{Additional information.} To my mind the most interesting questions about the actual Waring's problem (over the integers) concern $G\left( k\right)$ , the least $s$ such that all sufficiently large positive integers are the sum of $s$ non-negative $k$ th powers. It has been known since work of Vinogradov in 1935 that $G\left( k\right)  \ll  k\log k$ , and famously Wooley [283] reduced the implied constant to $1 + o\left( 1\right)$ , still the best known. It has long been a vague dream of additive combinatorialists (suggested by Freiman, and articulated more recently by Gowers [124]) that the structure theory of set addition could have something to say about this problem, but nothing definite has ever been done in this direction, and it is not clear how promising the idea is.

The same question may be asked over finite fields, and one might hope that this provides a fertile testing ground for ideas. I believe a reasonable analogue of $G\left( k\right)$ is ${G}_{\text{fin }}\left( k\right)$ , defined as follows. Fix a prime $p \in  \left( {k,{2k}}\right)$ (in small characteristic, exceptional phenomena occur; see [194]). What is the least $s$ such that, for sufficiently large $d$ , every polynomial over ${\mathbf{F}}_{p}$ of degree $d$ may be written as the sum of at most $s$ $k$ th powers of polynomials of degree $\leq  d/k$ ? The bound ${G}_{\text{fin }}\left( k\right)  \leq  \left( {1 + o\left( 1\right) }\right) k\log k$ , matching what is known over the integers, was obtained by Liu and Wooley [194]. Any improvement to this would be of interest.
\end{openproblem}

\plainproblemsection

\begin{openproblem}
Is a random polynomial with coefficients in $\{ 0,1\}$ and nonzero constant term almost surely irreducible?
\paragraph{Additional information.} More precisely, writing ${p}_{n}$ for the probability that $1 + {a}_{1}x + \cdots  +$ ${a}_{n - 1}{x}^{n - 1} + {x}^{n}$ is irreducible, where the ${a}_{i}$ are i.i.d. Bernouilli random variables, does ${p}_{n} \rightarrow  1$ ?

Unconditionally, the best bound known is ${p}_{n} \gg  1/\log n$ , due to Konyagin [186]. Very recently Breuillard and Varjú [45] have established the conjecture conditional upon the Grand Riemann Hypothesis. It should also be noted that Bary-Soroker and Kozma [20] gave an unconditional proof of the corresponding result in which the ${a}_{i}$ are selected uniformly from $\{ 1,\ldots ,{210}\}$ . This depends on 210 being the product of four primes, and so their argument does not apply with 211.

Update 2023. Bary-Soroker, Koukoulopoulos and Kozma [19, Theorem 1] have made further progression on this question, establishing a corresponding result with $\{ 1,2,\ldots ,{210}\}$ replaced by $\{ 0,1,\ldots , M\}$ for any $M \geq  {34}$ (and so, in particular, answering the question above with $M = {211}$ ). Moreover, they prove a positive lower bound $\delta$ on the irreducibility probability of a random polynomial with coefficients in $\{ 0,1,\ldots , M\}$ for all $1 \leq  M \leq  {33}$ . It is also worth remarking on the short paper [18], where it is shown that a random polynomial with coefficients in $\{  \pm  1\}$ has irreducibility probability tending to 1 along a certain infinite sequence of degrees $n$ .
\end{openproblem}

\plainproblemsection

\begin{openproblem}
Let $d \geq  3$ be an odd integer. Give bounds on $\nu \left( d\right)$ such that if $n > \nu \left( d\right)$ the following is true: given any homogeneous polynomial $F\left( \mathbf{x}\right)  =$ $F\left( {{x}_{1},\ldots ,{x}_{n}}\right)  \in  \mathbf{Z}\left\lbrack  {{x}_{1},\ldots ,{x}_{n}}\right\rbrack$ of degree $d$ , there is some $\mathbf{x} \in  {\mathbf{Z}}^{n} \smallsetminus  \{ \mathbf{0}\}$ such that $F\left( \mathbf{x}\right)  = 0.$
\paragraph{Additional information.} That $\nu \left( d\right)$ exists at all is highly nontrivial and is a famous result of Birch [29]. It was 40 years before Wooley [284] gave the first explicit values of $\nu \left( d\right)$ , which have tower-type growth in $d$ . So far as I am aware it is considered possible that $\nu \left( d\right)  = {d}^{2}$ , but even the local variant of this (that is, finding solutions in ${\mathbf{Q}}_{p}$ rather than $\mathbf{Q}$ ) is wide open (see, for instance,[157]).

It is known that $\nu \left( 3\right)  \leq  {13}$ by work of Heath-Brown [156].
\end{openproblem}

\plainproblemsection

\begin{openproblem}
As Problem 98 illustrates, finding a single solution to a polynomial equation $F\left( {{x}_{1},\ldots ,{x}_{n}}\right)  = C$ (say) can be very difficult, and asymptotically enumerating such solutions inside (say) a box ${\left\lbrack  X\right\rbrack  }^{n}$ is still harder. Nonetheless, one may ask about going yet further, and pose the question of estimating the number of solutions with the ${x}_{i}$ constrained to lie in some set $A \subset  \left\lbrack  X\right\rbrack$ . In particular, what conditions on $A$ ensure that the number of such solutions is roughly ${\alpha }^{n}$ times the number of solutions in $\left\lbrack  X\right\rbrack$ , where $\alpha  \mathrel{\text{:=}} \left| A\right| /X$ is the density of $A$ in $\left\lbrack  X\right\rbrack$ , imagining here that $\alpha  \in  \left( {0,1}\right)$ is fixed and $X$ is large?

When $F$ is linear (and $n \geq  3$ ) such a condition is that $A$ has no large Fourier coefficients, that is to say ${X}^{-1}\left| {\mathop{\sum }\limits_{{x \leq  X}}\left( {{1}_{A}\left( x\right)  - \alpha }\right) e\left( {\theta x}\right) }\right|  = o\left( 1\right)$ uniformly in $\theta$ . Some questions are

(i) What can be said when $\deg F = 2$ and $n = 7$ , at least in the ’generic’ case?

(ii) What is the least value of $n$ for which one can say something in the case $\deg F = 3$ , at least in the ’generic’ case?

(iii) What about specific interesting cases such as $F\left( {{x}_{1},{x}_{2},{x}_{3},{x}_{4}}\right)  = {x}_{1}^{2} + {x}_{2}^{2} +$ ${x}_{3}^{2} + {x}_{4}^{2}?$
\paragraph{Additional information.} Generic here might mean one of two different things. What happens for a ’random’ $F$ with coefficients selected from a large box? What happens for $F$ whose coefficients lie outside some positive codimension algebraic set? I mostly have in mind the latter interpretation, which offers the chance of saying something for particular explicit $F$ , but the former case may also be interesting. In both cases one wants a conclusion for all $A$ , that is to say the allowed $F$ cannot depend on $A$ .

Regarding item (i), in the case $n = 8$ it should probably follow using the methods of [141] that in the generic case it is enough for $A$ not to correlate with progressions, i.e. ${X}^{-1}\left| {\mathop{\sum }\limits_{{x \in  X}}\left( {{1}_{A}\left( x\right)  - \alpha }\right) {1}_{P}\left( x\right) }\right|  = o\left( 1\right)$ for all progressions $P \subset  \left\lbrack  X\right\rbrack$ , but it seems hard to reduce the number of variables to 7 . On the other hand, I would guess that in the generic case $n = 5$ (or perhaps even fewer) variables should suffice for such a statement.

I know very little about (ii). Presumably, some value of $n$ could be provided by adapting the methods of Cook and Magyar [76] but this is not in the literature so far as I am aware.

Finally, (iii) seems very difficult (and it would be very interesting, as it could well shed light on the well-known problem of showing that numbers $\equiv  4\left( {\;\operatorname{mod}\;{24}}\right)$ are ${p}_{1}^{2} + {p}_{2}^{2} + {p}_{3}^{2} + {p}_{4}^{2}$ ). One would certainly need to assume that $A$ has no large quadratic Fourier coefficients, ${X}^{-1}\left| {\mathop{\sum }\limits_{{x \leq  X}}\left( {{1}_{A}\left( x\right)  - \alpha }\right) e\left( {\theta {x}^{2}}\right) }\right|  = o\left( 1\right)$ , and for the analogous question with 5 or more variables a condition of this nature is sufficient (this may be extracted from [52]).

When I prepared the first version of these notes, the following problem seemed not especially well-known. Now it seems to be considered a major unsolved problem in group theory.
\end{openproblem}

\clearpage
\section{Geometric and Incidence Combinatorics}

\plainproblemsection

\begin{openproblem}
How many rotated (about the origin) copies of the 'pyjama set' $\left\{  {\left( {x, y}\right)  \in  {\mathbf{R}}^{2} : \operatorname{dist}\left( {x,\mathbf{Z}}\right)  \leq  \varepsilon }\right\}$ are needed to cover ${\mathbf{R}}^{2}$ ?
\paragraph{Additional information.} Manners [202], solving the 'pyjama problem', proved that there is some finite set of rotations with this property. His proof uses topological dynamics and is ineffective. I am not aware of any nontrivial lower bounds: in particular, is ${\varepsilon }^{-C}$ rotations enough?
\end{openproblem}

\plainproblemsection

\begin{openproblem}
Can the Cohn-Elkies scheme be used to prove the optimal bound for circle-packings?
\paragraph{Additional information.} Cohn and Elkies [69] put forward a general scheme for obtaining upper bounds on the density of sphere packings in ${\mathbf{R}}^{d}$ , and conjectured that it is optimal when $d = 1,2,8,{24}$ . They proved this when $d = 1$ , and famously Viazovska [277] established the case $d = 8$ and, in joint work with Cohn, Kumar, Miller and Radchenko [70], this was adapted to $d = {24}$ .

The case $d = 2$ amounts to constructing a radial function $f : {\mathbf{R}}^{2} \rightarrow  \mathbf{R}$ with $f\left( x\right)  \leq  0$ for $\left| x\right|  \geq  2,\widehat{f}\left( t\right)  \geq  0$ for all $t,\widehat{f}\left( 0\right)  > 0$ , and $\frac{f\left( 0\right) }{\widehat{f}\left( 0\right) } = \frac{\sqrt{3}}{6}$ . Such a function must vanish at the hexagonal lattice $\Lambda$ (except at zero), and its Fourier transform must vanish at the dual lattice ${\Lambda }^{ * }$ (again, except at zero).
\end{openproblem}

\plainproblemsection

\begin{openproblem}
Suppose that $A \subset  {\mathbf{F}}_{p}^{2}$ is a set meeting every line in at most 2 points. Is it true that all except $o\left( p\right)$ points of $A$ lie on a cubic curve?
\paragraph{Additional information.} There is an absolutely vast literature on 'arcs', which is the name given to sets with no-three in a line. However, I could not find this question asked, or really even hinted at, in that literature, which is more concerned with exact quantities for small $p$ . I asked it myself in [137]. It is more natural to formulate the question in projective space.

The largest no-3-in-a-line set (arc) is a conic of size $p + O\left( 1\right)$ , and Voloch [278], building on work of Segre, showed that any arc of size $> \frac{44}{45}p$ lies in a conic. Below the threshold $\frac{1}{2}p$ there are genuinely cubic examples. The would-be solver should also note that the result is false in ${\mathbf{F}}_{q}$ with $q = {p}^{2}$ , where there exist ’hermitian quadrics'.

There are many further problems on no-three-in-a-line sets, on which one would expect a positive solution to Problem 68 to shed some light. One of my favourites is the following. Denoting by ${f}_{d}\left( p\right)$ the size of the largest no-three-in-a-line set in ${\mathbf{F}}_{p}^{d}$ , one has ${f}_{2}\left( p\right)  \sim  p$ and ${f}_{3}\left( p\right)  \sim  {p}^{2}$ . Is ${f}_{4}\left( p\right)  = o\left( {p}^{3}\right)$ , or even $O\left( {p}^{2 + o\left( 1\right) }\right)$ ? The point here is that in dimension 4 and above, low-dimensional algebraic examples are forced to contain whole lines. All that is known is a slight improvement ${f}_{4}\left( p\right)  \leq  \left( {1 - c}\right) {p}^{3}$ on the trivial bound, due to Nagy and Szőnyi [212].

Another problem is to determine the largest subset of ${\mathbf{F}}_{p}^{2}$ with no four on a line; at the moment there is a factor of 2 discrepancy between the lower bounds of $\left( {1 + o\left( 1\right) }\right) p$ and the upper bound $\left( {2 + o\left( 1\right) }\right) p$ .

More generally than in the statement of the problem, one may conjecture that if $A \subset  {\mathbf{F}}_{p}^{2}$ meets every line in $O\left( 1\right)$ points then all but $o\left( p\right)$ points lie on a curve of degree $O\left( 1\right)$ . This may help with the following beautiful problem: let $\operatorname{PG}\left( {2, q}\right)$ be the projective plane over ${\mathbf{F}}_{q}$ . Is there a set $B \subset  \mathrm{{PG}}\left( {2, q}\right)$ (a ’blocking set’) meeting every line in at least 1 , but at most 1000 , points? It seems this problem is originally due to Erdős (see [108]). If $q$ is not prime then, in many cases, there can be such sets: see $\left\lbrack  {{121},\mathrm{p}{283}}\right\rbrack$ for several references. However, as $q$ ranges over sufficiently large primes it is speculated that no such set exists.
\end{openproblem}

\plainproblemsection

\begin{openproblem}
Fix a number $k$ . Let $A \subset  {\mathbf{R}}^{2}$ be a set of $n$ points, with no more than $k$ on any line. Suppose that, for at least $\delta {n}^{2}$ pairs of points $\left( {x, y}\right)  \in  A \times  A$ , the line ${xy}$ contains a third point of $A$ . Is there some cubic curve containing at least ${cn}$ points of $A$ , for some $c = c\left( {k,\delta }\right)  > 0$ ?
\paragraph{Additional information.} Very little is known, except when $\delta  = 1 - O\left( {1/n}\right)$ , in which case the statement follows from the main results of [150]. The statement is also known for sets $A$ lying on curves of higher (but essentially constant) degree [92]; this could well be an important ingredient in any proof. Let us also note the following beautiful result of Barak, Dvir, Wigderson and Yehudayoff [17]: if a set $A \subset  {\mathbf{C}}^{d}$ has $\geq  \delta {n}^{2}$ collinear triples and no $k$ on a line, it has a subset of size ${ \gg  }_{\delta , k}n$ of dimension $\ll  {\delta }^{-2}$ . I have often thought this should be relevant to Problem 69 but have not managed to establish a connection.

Elekes and Szabó [92] make the following conjecture, which is much weaker than Problem 69 (but should hold without the assumption of no more than $k$ points on a line): if $A$ has $c{n}^{2}$ collinear triples, do some 10 points of $A$ lie on a cubic curve? (Note that there is a cubic curve through any 9 points.)

I was led to this question by consideration of the following two beautiful problems, both of which ought to follow from a positive solution to it or closely-related questions.
\end{openproblem}

\plainproblemsection

\begin{openproblem}
Fix integers $k,\ell$ . Is it true that, given a set of $n \geq  {n}_{0}\left( {k,\ell }\right)$ points in ${\mathbf{R}}^{2}$ , either some $k$ of them lie on a line, or some $\ell$ of them are ’mutually visible’, that is to say that line segment joining them contains no other point from the set?
\paragraph{Additional information.} This question was first raised in [172]. It is known for $\ell  \leq  5$ , see [4].

The link to Problem 69 is that if $A$ does not have $\ell$ mutually visible points then one may easily see that it has ${ \gg  }_{\ell }{n}^{2}$ collinear triples.
\end{openproblem}

\plainproblemsection

\begin{openproblem}
Suppose that $A \subset  {\mathbf{R}}^{2}$ is a set of size $n$ with $c{n}^{2}$ collinear 4-tuples. Does it contain 5 points on a line?
\paragraph{Additional information.} This question was asked by Erdős on numerous occasions. Much more should be true (e.g., 100 points on a line). The link to Problem 69 is that the assumption of many collinear 4 tuples certainly implies many collinear triples and so, assuming no 5 on a line, a positive proportion of $A$ lies on a cubic curve. By an iterative argument one may put almost all of $A$ on a union of cubic curves. However, one may verify that such a set cannot, after all, have many collinear 4-tuples. See [92].

A construction of Solymosi and Stojaković [264] shows that one may have as ${e}^{-c\sqrt{\log n}}{n}^{2}$ collinear 4-tuples without a collinear 5-tuple. It ought to be the case that one can modify this construction so that no more than $O\left( 1\right)$ of the points lies on a cubic curve, but I was not able to do this. Indeed, trying to do this led to Problems 73 and 74 below.
\end{openproblem}

\plainproblemsection

\begin{openproblem}
What is the largest subset of the grid ${\left\lbrack  N\right\rbrack  }^{2}$ with no three points in a line? In particular, for $N$ sufficiently large is it impossible to have a set of size ${2N}$ with this property?
\paragraph{Additional information.} Specific cases of this problem date back to Dudeney over 100 years ago. Despite initial appearances, it has a very different flavour to, for example, Problem 68, as lines meeting the grid ${\left\lbrack  N\right\rbrack  }^{2}$ often contain very few points. For a bibliography, see [152, Problem F4]. In particular, we note that there do exist such configurations of ${2N}$ points for $N$ up to around 50 . It was shown in [155] that for arbitrary $N$ one can have $\left( {\frac{3}{2} + o\left( 1\right) }\right) N$ such points, and my personal suspicion is that this is optimal. It is possible that any no-three-in-a-line subset of ${\left\lbrack  N\right\rbrack  }^{2}$ is either small, or has a large subset which reduces $\left( {\;\operatorname{mod}\;p}\right)$ to a set of points on a curve in ${\mathbf{F}}_{p}^{2}$ . It might be interesting to formulate a precise conjecture of this type and see whether it can be used to show that the construction of [155] is optimal. I would find this a lot more convincing than the heuristic given by Guy and Kelly [153].

An interesting related question is the following: is there a subset $A \subset  {\mathbf{Z}}^{2}$ with no three points on a line, and with $\left| {A \cap  {\left\lbrack  -N, N\right\rbrack  }^{2}}\right|  \geq  {cN}$ for some absolute constant $c > 0$ and all sufficiently large $N$ ? The answer to this may well be no, as conjectured by Erde [94, Conjecture 5.1]. This seems plausible to me, due to the vague intuition that examples of large no-three-in-a-line subsets of ${\left\lbrack  -N, N\right\rbrack  }^{2}$ may have to come from constructions $\left( {\;\operatorname{mod}\;p}\right)$ , with the value of $p$ being somehow tied to $N$ . On the other hand, Nagy, Nagy and Woodroofe [214] take the contrary view, and provide numerical evidence based on greedy constructions that $c$ may exist and be at least around 0.8 .
\end{openproblem}

\plainproblemsection

\begin{openproblem}
Let $\Gamma$ be a smooth codimension 2 surface in ${\mathbf{R}}^{n}$ . Must $\Gamma$ intersect some 2-dimensional plane in 5 points, if $n$ is sufficiently large?
\paragraph{Additional information.} Obviously there is a more general problem here, where the codimen-sion of $\Gamma$ is $d$ and one wants to know that some $d$ -dimension plane intersects $\Gamma$ in $f\left( {n, d}\right)$ points.

I intend the question to be local, as opposed to being about global topological properties of $\Gamma$ . Thus for the purposes of this problem one can take $\Gamma  = \{ \left( {x, f\left( x\right) }\right)$ : $\left. {x \in  {\left( -1,1\right) }^{n - 2}}\right\}$ for some smooth function $f : {\left( -2,2\right) }^{n - 2} \rightarrow  {\mathbf{R}}^{2}$ with $f\left( 0\right)  = 0$ , say.

More or less all I know is that $f\left( {n,1}\right)  = 2$ (this is obvious, for an example take a convex hypersurface) and that $f\left( {n, d}\right)  \geq  {2}^{d}$ for large enough $n$ (this is less obvious, but it follows from essentially the same argument which shows that a sufficiently large set in an abelian group contains a ’Hilbert cube’ $x + {\varepsilon }_{1}{u}_{1} + \cdots  + {\varepsilon }_{d}{u}_{d}$ ).

In particular, I do not know whether $f\left( {n,2}\right)$ is bounded independently of $n$ . It is at least true that $f\left( {n, d}\right)$ is finite for all $n, d$ . A proof was sketched by René Thom [274] in a rather obscure paper. The details were worked out in a paper of Chaperon and Mayer [59]. Explicit constructions are given in the work of Dvir and Lovett [84].
\end{openproblem}

\plainproblemsection

\begin{openproblem}
What is the largest subset of ${\left\lbrack  N\right\rbrack  }^{d}$ with no 5 points on a 2-plane?
\paragraph{Additional information.} It was consideration of this problem which led me to Problem 73 above. I do not know if there can be such a set of size ${N}^{d - {100}}$ . A more-or-less equivalent formulation is to find the largest subset of ${\left\lbrack  N\right\rbrack  }^{d}$ with no 5 points on a conic (since conics are planar, and though any 5 points in the plane there is a conic).

Questions of this type seem related to the following question asked by Bays and Breuillard [23]: do there exist arbitrarily large finite sets $A \subset  {\mathbf{R}}^{2}$ with $\mid  A +$ $A{\left|  = {\left| A\right| }^{1 + o\left( 1\right) }\right| }^{2}$ , but which are in ’general position’ in the sense that for every $d$ , $\mathop{\max }\limits_{{C : \deg C = d}}\left| {A \cap  C}\right|$ is bounded? It is important to note that this is required for every $d$ . If one just requires $d = 1$ , for example then an example of Pach [215] suffices.
\end{openproblem}

\plainproblemsection

\begin{openproblem}
Let $X \subset  {\mathbf{R}}^{2}$ be a set of $n$ points. Does there exist a line $\ell$ through at least two points of $X$ such that the numbers of points on either side of $\ell$ differ by at most 100 ?
\paragraph{Additional information.} I read this in Alon [181, Problem 364], but Gil Kalai [171] has directed me to a blog post he prepared on this topic in 2010. In particular, that the answer to the question is affirmative both he and Pinchasi [224, Problem 7.1] refer to as the Kupitz-Perles conjecture. Pinchasi [224] has shown that it is true with 100 replaced by $O\left( {\log \log n}\right)$ ; this is the best-known upper bound for the problem. Alon (see [224]) has given an example to show that it is not true with 100 replaced by 2 .

Conlon and Lim [74] give a negative answer to a variant of this problem with 'pseudolines'; their paper may be consulted for further references.
\end{openproblem}

\plainproblemsection

\begin{openproblem}
Let $A$ be a set of $n$ points in the plane. Can one select a set ${A}^{\prime }$ of $n/2$ points, with the property that any axis-parallel rectangle containing 1000 points of $A$ contains at least one point of ${A}^{\prime }$ ?
\paragraph{Additional information.} This problem was told to me by Nets Katz a number of years ago. I believe it is instance of a well-known problem in discrepancy theory, namely whether weak $\varepsilon$ -nets for axis-parallel boxes of size $O\left( {1/\varepsilon }\right)$ exist. Certainly a search for those terms will yield the relevant papers, of which [6] seems to be particularly pertinent; if I understand correctly (though I should caution that this is far from my expertise), they show that the answer is ’yes’ if 1000 is replaced by $C\log \log n$ . Noga Alon remarked to me that if ${A}^{\prime }$ is required to be a subset of $A$ then the answer is no. See Lemma 3.1 of [216].
\end{openproblem}

\plainproblemsection

\begin{openproblem}
Given $n$ points in the unit disc, must there be a triangle of area at most ${n}^{-2 + o\left( 1\right) }$ determined by them?
\paragraph{Additional information.} This is 'Heilbronn's triangle problem' and it is rather notorious. Kom-lós, Pintz and Szemerédi [183] showed that the o(1)-term is necessary, and in the other direction [182] they showed that there must be a triangle of area at most ${n}^{-8/7 + o\left( 1\right) }$ . I am not aware of any progress since then. It is worth remarking that the problem was a favourite of Klaus Roth [234].

Update 2023. Cohen, Pohoata and Zakharov [66] have improved the bound to ${n}^{-8/7 - c}$ for some $c > 0$ (they obtain $\left. {c = \frac{1}{2000}}\right)$ .
\end{openproblem}

\plainproblemsection

\begin{openproblem}
Suppose that $A$ is an open subset of ${\left\lbrack  0,1\right\rbrack  }^{2}$ with measure $\alpha$ . Are there four points in $A$ determining an axis-parallel rectangle with area $\geq  c{\alpha }^{2}$ ?
\paragraph{Additional information.} This is often known as 'Carbery's rectangle problem', though it appears in a joint paper of Carbery, Christ and Wright. See [56, Section 6]. It is quite easy to show using Cauchy-Schwarz that there must be such a rectangle with area $\gg  {\alpha }^{2}{\left( \log 1/\alpha \right) }^{-1}$ . In the other direction, the example of a diagonal stripe of width $\sim  \alpha$ shows that one could not hope for better. A related discrete problem is as follows: does there exist a set $A \subset  \left\lbrack  N\right\rbrack  ,\left| A\right|  \geq  {100}\sqrt{N}$ , such that if an additive quadruple $x, x + {h}_{1}, x + {h}_{2}, x + {h}_{1} + {h}_{2}$ lies in $A$ then $\left| {h}_{1}\right| \left| {h}_{2}\right|  \leq  N$ ?

This problem and related problems are considered by Keleti [179]. One might instead ask for an axis parallel rectangle $R$ , all four vertices in $A$ , and with $\left| {A \cap  R}\right|  \gg$ ${\alpha }^{C}$ . It is known that there is such an $R$ with $C = 4$ , but this need not be so with $C = 3$ , an example Keleti attributes to Reiman.

In dimension 3 and above, these problems are wide open, even without any restriction the area (volume) of cuboids. In particular one may formulate the 'box problem’: what is the largest density of a subset of ${\left\lbrack  n\right\rbrack  }^{3}$ not containing the eight vertices of a cuboid? The answer is between ${n}^{-1/3}$ and ${n}^{-1/4}$ (up to constants): see [173]. Note that this problem is essentially a hypergraph Túran problem, that of determining $\operatorname{ex}\left( {n,{K}_{2,2,2}^{3}}\right)$ in the language used there. As such, it goes back to questions raised by Erdős in 1965. It seems to be believed (cf. [211, Conjecture 1.4]) that the upper bound represents the truth. A recent reference, giving a new construction of the ${n}^{-1/3}$ lower bound of [173], improved lower bounds in higher dimensions and a useful overview of the literature, is [75].
\end{openproblem}

\plainproblemsection

\begin{openproblem}
Consider a set $S \subset  {\left\lbrack  N\right\rbrack  }^{3}$ with the property that any two distinct elements $s,{s}^{\prime }$ of $S$ are comparable, which means that $s - {s}^{\prime }$ has either two (strictly) positive indices or two (strictly) negative ones. Is $\left| S\right|  \leq  {N}^{2 - \delta }$ for some $\delta  > 0$ ?
\paragraph{Additional information.} This is perhaps the most basic of a host of questions considered by Gowers and Long [123], motivated by a question of Po-Shen Loh. It can be formulated as a question about the maximum independent set in a graph on vertex set ${\left\lbrack  n\right\rbrack  }^{3}$ in which two vertices $x, y$ are joined if $x, y$ are not comparable, which means that $\left( {{x}_{1},{x}_{2},{x}_{3}}\right)$ is joined to (for example) all $\left( {{x}_{1},{x}_{2}^{\prime },{x}_{3}^{\prime }}\right)$ with ${x}_{2}^{\prime } < {x}_{2}$ and ${x}_{3}^{\prime } > {x}_{3}$ . One reason the problem is hard (it seems to me) is that it asks to go below the threshold for which a spectral bound ('Delsarte's method') on the independence number might be possible.

Here is a modular version which, while not as natural, has a similar flavour and more symmetry. Let $S \subset  {\mathbf{F}}_{p}^{3}$ be set consisting of all triples(0, y, z)where $y, - z \in  \left\{  {1,\ldots ,\frac{1}{2}\left( {p - 1}\right) }\right\}$ , together with the other five similar classes (e.g.(z,0, y) satisfying the same property). Suppose that $A - A$ is disjoint from $S$ . Is $\left| A\right|  \leq  {p}^{2 - \delta }$ ?
\end{openproblem}

\clearpage
\section{Harmonic Analysis, Random Structures, and Combinatorial Enumeration}

\plainproblemsection

\begin{openproblem}
Let $A$ be a set of size $n$ integers. Is there some $\theta$ such that $\mathop{\sum }\limits_{{a \in  A}}\cos \left( {a\theta }\right)  \leq   - c\sqrt{n}?$
\paragraph{Additional information.} This is Chowla's cosine problem. First posed explicitly in 1965 [62], it has its genesis in a question from 1948 asked by Ankeny and Chowla.

The best known result is due to Ruzsa [242], who showed that $- {e}^{-c\sqrt{\log n}}$ is attainable. Ruzsa remarked to me in 2001 that almost nothing is known for the more general sum of cosines ${\lambda }_{1}\cos {r}_{1}t + \cdots  + {\lambda }_{n}\cos {r}_{n}t$ in which ${\lambda }_{1},\ldots ,{\lambda }_{n} \in  \left\lbrack  {{0.99},{1.01}}\right\rbrack$ (say).

There is a corresponding problem for sines, but it is somewhat more obscure. The question here is whether there is some $\theta$ such that $\left| {\mathop{\sum }\limits_{{a \in  A}}\sin \left( {a\theta }\right) }\right|  \geq  {n}^{1/2 + c}$ . Montgomery [209, Problem 54] attributes this question to Bohr. So far as I am aware, the best example known is of a set showing that we cannot take $c \geq  \frac{1}{6}$ , due to Bourgain in a somewhat difficult to find paper [38]. By contrast, Konyagin [185] has shown that there is some $\theta$ such that $\left| {\mathop{\sum }\limits_{{a \in  A}}\sin \left( {a\theta }\right) }\right|  \gg  {n}^{1/2}{\left( \log n/\log \log n\right) }^{1/2}$ .

See also the remarks to Problem 1, where related problems are discussed with cos replaced by piecewise constant functions $f$ .
\end{openproblem}

\plainproblemsection

\begin{openproblem}
Let $A \subset  \mathbf{Z}$ be a set of size $n$ . For how many $\theta  \in  \mathbf{R}/\mathbf{Z}$ must we have $\mathop{\sum }\limits_{{a \in  A}}\cos \left( {a\theta }\right)  = 0?$
\paragraph{Additional information.} This is [193, Problem 22]. He says ’probably $n - 1$ , or not much less’. This is false (depending on how one defines 'not much less'): there are examples with at most ${n}^{5/6 + o\left( 1\right) }$ zeros, due to Borwein, Erdélyi, Ferguson and Lockhart [36]. This has recently been improved to $O\left( {{n}^{2/3}{\log }^{2/3}n}\right)$ by Juškevičius and Sahasrabudhe [169].

It is known (see $\left\lbrack  {{97},{245}}\right\rbrack$ that the number of zeros does tend to infinity with $n$ , albeit very slowly (the bound ${\left( \log \log \log n\right) }^{1/2 - o\left( 1\right) }$ is obtained in [245]).
\end{openproblem}

\plainproblemsection

\begin{openproblem}
Describe the rough structure of sets $A \subset  \mathbf{Z}$ with $\left| A\right|  = n$ and ${\begin{Vmatrix}{\widehat{1}}_{A}\end{Vmatrix}}_{1} \leq  K\log n$
\paragraph{Additional information.} This is the 'inverse Littlewood problem'; Littlewood conjectured, and McGehee-Pigno-Smith and Konyagin indepently proved, that ${\begin{Vmatrix}{\widehat{1}}_{A}\end{Vmatrix}}_{1} \gg  \log n$ . It is natural to conjecture that $A$ has symmetric difference $o\left( n\right)$ with a $\pm$ combination of ${O}_{K}\left( 1\right)$ characteristic functions of progressions. More precise conclusions should presumably be possible, but one has to be careful: for instance, if $P$ is a 2-dimensional progression with sidelength $\sim  {e}^{\sqrt{\log n}}$ then ${\begin{Vmatrix}{\widehat{1}}_{S}\end{Vmatrix}}_{1} \ll  \log n$ .

A solution to this problem would be extremely relevant to Problem 1, as mentioned there, as well as to the further problems (1.1), (1.2) mentioned in the remarks to that problem. It may also be relevant (and it would be good to have a formal deduction of this type) to the currently unresolved Strong Littlewood Conjecture, which states that ${\begin{Vmatrix}{\widehat{1}}_{A}\end{Vmatrix}}_{1} \geq  \left( {1 + o\left( 1\right) }\right) {\begin{Vmatrix}{\widehat{1}}_{P}\end{Vmatrix}}_{1} = \left( {\frac{4}{{\pi }^{2}} + o\left( 1\right) }\right) \log n$ , where $P$ is a progression of length $n$ . (In fact, I believe that the even stronger statement ${\begin{Vmatrix}{\widehat{1}}_{A}\end{Vmatrix}}_{1} \geq  {\begin{Vmatrix}{\widehat{1}}_{P}\end{Vmatrix}}_{1}$ is thought to hold.)

There are also interesting questions of this type in finite fields. For example, if $A \subset  {\mathbf{F}}_{2}^{n}$ has $\mathop{\sum }\limits_{r}\left| {{\widehat{1}}_{A}\left( r\right) }\right|  \leq  M$ then Sanders [250] showed ${1}_{A}$ is a $\pm$ sum of $\exp \left( {M}^{3 + o\left( 1\right) }\right)$ indicator functions of cosets. However, the true dependence may well be polynomial.

Let us also mention the ’Littlewood-Gowers problem’: if $A \subset  \mathbf{Z}/p\mathbf{Z}$ has density $\frac{1}{2}$ , what is the smallest possible value of $\mathop{\sum }\limits_{r}\left| {{\widehat{1}}_{A}\left( r\right) }\right|$ ? The example of an arithmetic progression shows that it can be $\ll  \log p$ , but the best known lower bound is $\gg  {\left( \log p\right) }^{1/2 - \varepsilon }$ , due to Sanders [247].
\end{openproblem}

\plainproblemsection

\begin{openproblem}
Let
\[
  P_N(z)=\sum_{j=0}^{N-1}\varepsilon_jz^j,
  \qquad \varepsilon_j\in\{-1,1\}.
\]

(i) Do there exist such polynomials satisfying, uniformly for \(|z|=1\),
\[
  |P_N(z)|=(1+o(1))\sqrt N?
\]

(ii) Is there a universal constant \(c>0\) such that every such polynomial
satisfies the Golay merit-factor bound
\[
  \int_0^1
  \left|P_N\!\left(e^{2\pi i t}\right)\right|^4\,dt
  \geq (1+c)N^2?
\]

(iii) For infinitely many \(q\), does there exist a set
\(A\subset\mathbb Z/q\mathbb Z\) with \(|A|\sim q^{1/3}\) whose nontrivial
Fourier coefficients are all \(O(q^{1/6})\)?
\end{openproblem}

\plainproblemsection

\begin{openproblem}
Let $c > 0$ . Let $A$ be a set of $n$ (distinct) integers. Does there exist $\theta$ such that no interval of length $\frac{1}{n}$ in $\mathbf{R}/\mathbf{Z}$ contains more than ${n}^{c}$ of the numbers ${\theta a}\left( {\;\operatorname{mod}\;1}\right) , a \in  A$ ?
\paragraph{Additional information.} This is only known for $c > \frac{1}{3}$ ; see [187]. The problem is raised as [209, Problem 17 (4)], where it is attributed to Komlós and Ruzsa. This problem feels not unrelated to questions like Problem 43, but I do not know a direct connection.
\end{openproblem}

\plainproblemsection

\begin{openproblem}
Let $p\left( k\right)$ be the limit as $n \rightarrow  \infty$ of the probability that a random permutation on $\left\lbrack  n\right\rbrack$ preserves some set of size $k$ . Is $p\left( k\right)$ a decreasing function of $k$ ? Is $p\left( k\right)  = \left( {C + o\left( 1\right) }\right) {k}^{-\alpha }{\left( \log k\right) }^{-3/2}$ for some absolute constant $C$ ?
\paragraph{Additional information.} I first learned of this problem from a paper of Britnell and Wildon [50]. In a joint paper with Eberhard and Ford [87] it was proven that $p\left( k\right)$ is bounded above and below by constant multiples of ${k}^{-\alpha }{\left( \log k\right) }^{-3/2}$ , which of course makes the second question very reasonable. This second question is a model for the following natural question: how many distinct elements are there in the $N \times  N$ multiplication table?

It is quite easy to see, given known results about the Poisson behaviour of permutations, that

\[
p\left( k\right)  = 1 - \mathop{\sum }\limits_{{\left( {{a}_{1},\ldots ,{a}_{k}}\right)  \in  {\Omega }_{k}}}\mathop{\prod }\limits_{{i = 1}}^{k}{e}^{-1/i}\frac{{\left( 1/i\right) }^{{a}_{i}}}{{a}_{i}!},
\]

where ${\Omega }_{k}$ is the (finite) set of all tuples of non-negative integers for which ${a}_{1}^{\prime } +$ $2{a}_{2}^{\prime } + \cdots  + k{a}_{k}^{\prime } \neq  k$ whenever ${a}_{i}^{\prime } \leq  {a}_{i}$ . For example, ${\Omega }_{1} = \{ \left( 0\right) \}$ , so $p\left( 1\right)  = 1 - \frac{1}{e}$ .

Here is a related model problem (the relation becomes clearer upon studying [87]). Pick a random set of integers $A$ by including $i$ in $A$ with probability $\frac{1}{i\log 2}$ . (Thus $A$ is a random lacunary sequence, with growth rate like that of the powers of 2). Let $\sum \left( A\right)$ be the set of all finite sums of distinct elements of $A$ . Is it true that $\frac{1}{n}\left| {\sum \left( A\right)  \cap  \left\lbrack  n\right\rbrack  }\right|  \sim  C{\left( \log n\right) }^{-1/2}$ as $n \rightarrow  \infty$ ?
\end{openproblem}

\plainproblemsection

\begin{openproblem}
I have a string $x \in  \{ 0,1{\} }^{n}$ . Let $\widetilde{x}$ be the random string obtained by deleting bits from $x$ independently at random with probability $\frac{1}{2}$ (thus, for example, if $n = 8$ and $x = {00110110}$ , it might be the case that $\widetilde{x} = {0111}$ , generated by deleting bits2,3,5,8.) An instance of $\widetilde{x}$ is called a ’trace’. How many independent traces ${\widetilde{x}}_{1},\ldots ,{\widetilde{x}}_{m}$ are needed before one can reconstruct $x$ with probability 0.9 ?
\paragraph{Additional information.} This problem is known as the trace reconstruction problem. I heard it from Yuval Peres in 2012, but it goes back to work of Levenshtein from the early 2000s. Independent work of Nazarov-Peres [213] and of De-O'Donnell-Servedio [80] on the topic gives the best known upper bound, showing that $m = {e}^{C{n}^{1/3}}$ suffices. Update 2020. Chase [61] has improved this to ${e}^{{n}^{1/5}{\log }^{C}n}$ .

Simple examples (see $\left\lbrack  {{21},§{4.2}}\right\rbrack$ ) show that $m$ must grow at least linearly in $n$ , and very recently Zachary Chase [60], improving on work of Holden and Lyons [162], improved this to $m \geq  {n}^{3/2}{\left( \log n\right) }^{-{16}}$ .

A related wide-open problem is that of reconstruction a binary string of length $n$ from the multiset of all $\left( \begin{array}{l} n \\  k \end{array}\right)$ substrings of length $k$ , the so-called $k$ -deck. In [189], it is shown that $k \geq  C\sqrt{n}$ suffices (see also [254] for a different proof of a bound weaker by only a logarithm). In [83], it is shown that $k \geq  {e}^{c\sqrt{\log n}}$ may not be enough (I thank Zachary Chase for these references).
\end{openproblem}

\plainproblemsection

\begin{openproblem}
Is every set $\Lambda  \subset  \mathbf{Z}$ either a Sidon set, or a set of analyticity?
\paragraph{Additional information.} This problem, known as the Dichotomy Problem, was raised in the 1960s in the subject of commutative harmonic analysis. It is important to note that the notion of Sidon set here is not the one featured in Section 4, but rather the harmonic analysis notion: $\Lambda$ is Sidon if and only if the Fourier algebra $A\left( \Lambda \right)  =$ $\left\{  {{\left( \widehat{f}\left( \lambda \right) \right) }_{\lambda  \in  \Lambda } : f \in  {L}^{1}\left( \mathbf{T}\right) }\right\}$ coincides with ${c}_{0}\left( \Lambda \right)$ , the algebra of sequences tending to zero. By contrast, $\Lambda$ is a set of analyticity if only analytic functions $F$ act on $A\left( \Lambda \right)$ . Much of the classical literature on these problems is difficult to penetrate for a modern reader (at least, it is for me). For more on the problem, I recommend, at least to the reader with passable French, the relatively recent paper of Kahane and Katznelson [170]. In that paper they show that a random set $\Lambda$ in which $\mathbf{P}\left( {n \in  \Lambda }\right)  = {p}_{n}$ is almost surely Sidon if $n{p}_{n}$ is bounded, and almost surely a set of analyticity if $n{p}_{n} \rightarrow  \infty$ .

A beautiful open question is whether every Sidon subset of $\mathbf{Z}$ (in the commutative harmonic analysis sense) is a finite union of independent sets, that is to say sets $A \subset  \mathbf{Z}$ , all of whose finite subset sums are distinct (for example, the powers of two). In celebrated work from the 1980s, Pisier [225] showed that $S \subset  \mathbf{Z}$ is Sidon if and only if it has the following property $\left( \star \right)$ : there is a $\delta  > 0$ such that, if ${S}^{\prime } \subset  S$ is finite, then ${S}^{\prime }$ contains an independent set $A$ with $\left| A\right|  \geq  \delta \left| {S}^{\prime }\right|$ . The open question is then the completely combinatorial question of whether every set with property $\left( \star \right)$ is a finite union of independent sets. I should say that I am not sufficiently expert on these topics to say with any certainty what the relationship (if any) between this question and the dichotomy problem is.

I included this problem here in memory of Jean Bourgain, who once told me he considered it a beautiful open question, and lamented that it might never be solved since to a large extent the subject had fallen out of fashion. Let me conclude these comments with another question of Bourgain, asked to Péter Varjú in 2013 [288]: Is it possible to find $n$ points in the unit square such that the $1/n$ -neighborhood of any line contains no more than $C$ of them for some absolute constant $C$ ? More information on this question may be found in [81].
\end{openproblem}

\plainproblemsection

\begin{openproblem}
In how many ways (asymptotically) $Q\left( n\right)$ may $n$ non-attacking queens be placed on an $n \times  n$ chessboard?
\paragraph{Additional information.} This is the $n$ queens problem. Though very recreational in its statement, it certainly hides some interesting mathematics. This is most clearly seen for the 'modular' version of the problem, in which the chessboard is toroidal and one defines $T\left( n\right)$ to be the number of ways to place $n$ non-attacking queens on an $n \times  n$ board.

There has been significant recent progress on the problem of estimating $Q\left( n\right)$ and $T\left( n\right)$ . Regarding the former, Simkin [261] showed that $Q\left( n\right)  = {\left( \left( 1 + o\left( 1\right) \right) n{e}^{-\alpha }\right) }^{n}$ , where $\alpha  \approx  {1.94} \times  {10}^{-3}$ is a constant given in terms of an entropy optimisation problem (the explicit solution of which remains an open problem). Using very different methods, in a monumental paper Bowtell and Keevash [42] showed that $T\left( n\right)  = {\left( \left( 1 + o\left( 1\right) \right) n{e}^{-3}\right) }^{n}$ when $n \equiv  1$ or 5 modulo 6 (there are no configurations in the other cases).

The problem of obtaining an asymptotic (rather than an asymptotic for the log) remains open. In the toroidal case, the problem is equivalent to the following. Identify ${\left( \mathbf{Z}/n\mathbf{Z}\right) }^{n}$ with the space of functions $f : \{ 1,\ldots , n\}  \rightarrow  \mathbf{Z}/n\mathbf{Z}$ , and let $S \subset  {\left( \mathbf{Z}/n\mathbf{Z}\right) }^{n}$ be the set of bijections. How many pairs $\left( {{\pi }_{1},{\pi }_{2}}\right)  \in  {\left( \mathbf{Z}/n\mathbf{Z}\right) }^{n}$ are there with ${\pi }_{1},{\pi }_{2},{\pi }_{1} + {\pi }_{2},{\pi }_{1} - {\pi }_{2} \in  S$ ?

This is, of course, very reminiscent of questions about four-term arithmetic progressions, and one suspects that the Gowers ${U}^{3}$ -norm of ${1}_{S}$ or related objects will come into play. For further comments, and for an ingenious solution to the problem in which only ${\pi }_{1},{\pi }_{2}$ and ${\pi }_{1} - {\pi }_{2}$ are required to lie in $S$ , see [90].
\end{openproblem}

\end{document}
